\documentclass[11pt,letterpaper]{report}
\usepackage{scientific_report}
\usepackage{mathtools}
\usepackage{bm}
\usepackage{mathrsfs}
\usepackage{microtype}
\usepackage{longtable}
\usepackage{cleveref}

\setlength{\headheight}{22pt}
\renewcommand{\chaptermark}[1]{\markboth{\thechapter.\ #1}{}}
\makeatletter
\renewcommand{\@pnumwidth}{2em}
\renewcommand{\@tocrmarg}{3em}
\makeatother

\hypersetup{
  pdftitle={Gauge-Covariant Variational Free Energy and Exact Coarse-Graining},
  pdfauthor={Robert C. Dennis},
  pdfsubject={Agents as local sections of belief and model associated bundles at one base point: independent interaction links written in the exponential chart, whose Lie-algebra generators are latent and inferred under a normalized prior, one collective evidence bound, attention and selective communication, fiber information geometry and inference histories, fast beliefs and slow models, exact coarse-graining of the agent network, and formation witnesses},
  pdfkeywords={variational free energy, ELBO, associated bundle, gauge covariance, Lie-algebra links, latent links, loop defect, attention, selective communication, Fisher information, inference histories, two-timescale dynamics, Markov kernel, exact coarse-graining, meta-agent formation}
}

% ---- shared macros: chapters must use these and must not redefine them ----
\newcommand{\KL}{D_{\mathrm{KL}}}
\newcommand{\E}{\mathbb E}
\newcommand{\R}{\mathbb R}
\newcommand{\given}{\mid}
\newcommand{\Tr}{\operatorname{Tr}}
\newcommand{\Sym}{\mathbb S}
\newcommand{\PSD}{\mathbb S_{+}}
\newcommand{\GL}{\mathrm{GL}}
\newcommand{\Fenergy}{\mathcal F}
\newcommand{\Lelbo}{\mathcal L}

% ---- compact, semantic result headings with one chapter-scoped sequence ----
% Every numbered statement has a typed counter and a stable label for cleveref.
\newcounter{statementsequence}[chapter]
\newcounter{definition}[chapter]
\newcounter{lemma}[chapter]
\newcounter{proposition}[chapter]
\newcounter{theorem}[chapter]
\newcounter{corollary}[chapter]
\newcounter{conjecture}[chapter]
\newcounter{openproblem}[chapter]
\newcounter{hypothesis}[chapter]
\newcounter{remark}[chapter]
\renewcommand{\thedefinition}{\thechapter.\arabic{definition}}
\renewcommand{\thelemma}{\thechapter.\arabic{lemma}}
\renewcommand{\theproposition}{\thechapter.\arabic{proposition}}
\renewcommand{\thetheorem}{\thechapter.\arabic{theorem}}
\renewcommand{\thecorollary}{\thechapter.\arabic{corollary}}
\renewcommand{\theconjecture}{\thechapter.\arabic{conjecture}}
\renewcommand{\theopenproblem}{\thechapter.\arabic{openproblem}}
\renewcommand{\thehypothesis}{\thechapter.\arabic{hypothesis}}
\renewcommand{\theremark}{\thechapter.\arabic{remark}}
\crefname{definition}{definition}{definitions}
\crefname{lemma}{lemma}{lemmas}
\crefname{proposition}{proposition}{propositions}
\crefname{theorem}{theorem}{theorems}
\crefname{corollary}{corollary}{corollaries}
\crefname{conjecture}{conjecture}{conjectures}
\crefname{openproblem}{open problem}{open problems}
\crefname{hypothesis}{hypothesis}{hypotheses}
\crefname{remark}{remark}{remarks}
\Crefname{definition}{Definition}{Definitions}
\Crefname{lemma}{Lemma}{Lemmas}
\Crefname{proposition}{Proposition}{Propositions}
\Crefname{theorem}{Theorem}{Theorems}
\Crefname{corollary}{Corollary}{Corollaries}
\Crefname{conjecture}{Conjecture}{Conjectures}
\Crefname{openproblem}{Open Problem}{Open Problems}
\Crefname{hypothesis}{Hypothesis}{Hypotheses}
\Crefname{remark}{Remark}{Remarks}
\newcommand{\syncstatementcounter}[1]{%
  \stepcounter{statementsequence}%
  \setcounter{#1}{\numexpr\value{statementsequence}-1\relax}%
  \refstepcounter{#1}%
}
\newcommand{\resultheading}[4]{%
  \syncstatementcounter{#2}%
  \label{#4}%
  \paragraph{#1~\csname the#2\endcsname\ (#3).}%
}
\newcommand{\definitionheading}[2]{\resultheading{Definition}{definition}{#1}{#2}}
\newcommand{\lemmaheading}[2]{\resultheading{Lemma}{lemma}{#1}{#2}}
\newcommand{\propositionheading}[2]{\resultheading{Proposition}{proposition}{#1}{#2}}
\newcommand{\theoremheading}[2]{\resultheading{Theorem}{theorem}{#1}{#2}}
\newcommand{\corollaryheading}[2]{\resultheading{Corollary}{corollary}{#1}{#2}}
\newcommand{\conjectureheading}[2]{\resultheading{Conjecture}{conjecture}{#1}{#2}}
\newcommand{\openproblemheading}[2]{\resultheading{Open Problem}{openproblem}{#1}{#2}}
\newcommand{\hypothesisheading}[2]{\resultheading{Hypothesis}{hypothesis}{#1}{#2}}
\newcommand{\remarkheading}[2]{\resultheading{Remark}{remark}{#1}{#2}}
% Compile-only aliases: a construction is a definition and a requirement is a
% standing hypothesis. New text uses the two target macros directly.
\newcommand{\constructionheading}[2]{\definitionheading{#1}{#2}}
\newcommand{\requirementheading}[2]{\hypothesisheading{#1}{#2}}

% ---- status tag: one per theorem-type statement or claim-bearing remark ----
\definecolor{statusfg}{RGB}{120,20,20}
\newcommand{\status}[1]{\mbox{{\small\textsc{\textcolor{statusfg}{[#1]}}}}}

\begin{document}
\hypersetup{pageanchor=false}
\makereporttitle
  {Gauge-Covariant Variational Free Energy\\and Exact Coarse-Graining}
  {Agents as Local Sections at One Base Point: Collective Evidence, Attention, Inference Histories, and Formation}
  {Robert C. Dennis}
  {Independent Researcher}
  {September 2026}

\hypersetup{pageanchor=true}
\pagenumbering{roman}
\tableofcontents
\listoftables
\clearpage
\pagenumbering{arabic}

\part{Single-Point Foundations}
\chapter{Program, Scope, and Standing Conventions}
\label{ch:introduction}

This manuscript asks what collective inference is available to agents that are local sections of a
belief and a model associated bundle of one principal bundle, placed at one base point and interacting
through independent links under one normalized joint law, and which parts of that inference survive
exact coarse-graining of the agents into meta-agents. The answer has four main parts. The whole
population has one evidence bound, an identity in the extended reals whose gap vanishes exactly at the
posterior as a measure (\Cref{thm:elbo-extended-gap,cor:elbo-bound-tightness}). Each agent's local
objective is a conditional free energy whose differential is that of the collective free energy
(\Cref{thm:obs-local-global-potential}). Coarse channels of the agent network preserve the evidence and
decompose the free energy exactly (\Cref{thm:rg-exact-coarse-vfe}), with an effective action for the
retained coordinates (\Cref{thm:rg-effective-action}). A finite-depth family forms successive
meta-agents under free-energy descent with learning, in a stated scope
(\Cref{thm:finite-depth-adaptive-formation}). \Cref{sec:intro-reading-map} maps these results to the
chapters, and \Cref{sec:intro-conventions} fixes, once for the whole manuscript, the defaults and the
claims it declines.

\section{Program aim}
\label{sec:intro-program-aim}

The program is a participatory multi-agent theory. Agents are local sections of associated bundles,
collectives form through the agents' coupled inference, and parents act back on their children through
normalized laws rather than through quantities that read the children's beliefs. Three requirements
follow from that aim, and the manuscript is organized around them. Every collective object must be a
complete probability object, constructed from the agents' mechanisms and interaction records or selected
explicitly, never an additional agent and never a sum of local objectives that happens to look like one.
Every statement about a parent must hold for the normalized parent law that coarse-graining actually
produces. And formation of a parent must be a theorem about a declared family, with its scope stated,
rather than a reading of a formula.

The present phase fixes one base point $c_*$ and settles the theory there before any transport across the
base is considered. At $c_*$ the exact normalized parent law exists for every coarse channel of the
agent network (\Cref{thm:rg-pointwise-parent-datum}), and \Cref{thm:finite-depth-adaptive-formation} gives
a formation witness under a supplied hierarchy of interaction strengths. On a homogeneous four-cycle a
random initial source assignment already makes a collective form from every Gaussian initial law, with
source learning idle (\Cref{prop:phil-seeded-contrast-formation}). At each sufficiently small fixed
temperature, slow source learning selects one matching in the assignment law for almost every Gaussian
initial mean, while the frozen-prior pairing difference vanishes at long times
(\Cref{conj:phil-contrast-formation}). How a
parent's own dynamics are determined by the network it summarizes is
\Cref{open:rg-pointwise-parent-dynamics}, the defining open problem of the phase.

\section{Objects and the local-first architecture}
\label{sec:intro-objects}

Agent $i$ is the object $\mathcal A^i=(\mathcal C_i;q_i^b,q_i^m,u_i)$ of \Cref{def:geo-agent}: a belief
section and a model section of the associated bundles $\mathcal E_b$ and $\mathcal E_m$ of one principal
$G$-bundle, expressed in one common frame $u_i$, a gauge choice by \Cref{prop:geo-common-frame-gauge}. The
two channels may carry inequivalent representations, and cross morphisms between them are additional
structure (\Cref{def:geo-cross-morphisms}). Interactions are normalized record kernels $K_a$ on the scopes
of an interaction multigraph, and the links $\Theta_e^x$ on its edge copies are independent data with loop
holonomy and per-edge invariant classes (\Cref{prop:geo-edge-twist-invariance}). The primitive descriptive
data are the agents, their local recognition laws, their evaluated mechanisms, and the record kernels.
Population laws enter only by construction (\Cref{prop:gen-exact-normalization}) or by an explicit
selection (\Cref{hyp:prob-measurable-coupling-selector}). The architecture is
\begin{equation}
\left\{
\mathcal A^i,G_{\theta,i},\mathbb Q_{i,o,X},K_a
\right\}_{i\in V,a\in\mathfrak A}
\longrightarrow
\begin{cases}
\mathbb P_{\theta,V} & \text{normalized construction},\\
\mathbb Q_{V,o,X}\in\operatorname{Cpl}(\mathbb Q_{\bullet,o,X})
& \text{measurable selection},
\end{cases}
\longrightarrow
\Lelbo_V,
\label{eq:intro-local-to-global-architecture}
\end{equation}
so the constructed joint and the selected coupling precede every local and collective objective. On the
evidence domain of \Cref{hyp:elbo-evidence-domain} there is one identity,
\begin{equation}
\log p_{\theta,V}(o\given X)
=\Lelbo_V(\mathbb Q_{V,o,X};X,o)
+\KL\bigl(\mathbb Q_{V,o,X}\big\Vert\boldsymbol\Pi_{\theta,V,o,X}\bigr),
\label{eq:elbo-intro-identity}
\end{equation}
with equality of the bound and the log evidence exactly when the recognition coupling is the posterior as
a measure, which matching coordinate marginals does not achieve (\Cref{cor:elbo-bound-tightness}). The
identity is invariant under passive re-expression of the frames when the kernels are covariant
(\Cref{prop:gen-product-evidence-invariance,cor:elbo-independent-frame-invariance}). Excluding the
posterior from a recognition family makes every pointwise gap positive, but it does not make the optimized
gap positive or attained (\Cref{sec:restrict-principle}).

\section{Results and reading map}
\label{sec:intro-reading-map}

Part I builds the single-point foundations. \Cref{ch:geometry} fixes the geometry at $c_*$: the active gauge
group is one copy of $G$, links carry per-edge twist classes, and three notions of agreement are separated.
\Cref{ch:probability} fixes kernels before densities, proves the chain rule used everywhere, and
distinguishes sample, law-point, and configuration levels. \Cref{ch:generative} builds the normalized joint
on arbitrary record cycles and states the typing prohibition. \Cref{ch:elbo} derives the evidence bound and
the local-collective potential identity, with observations entering as interaction records
(\Cref{thm:obs-agent-interaction-equivalence}). \Cref{ch:selective-communication} places attention inside
the fixed joint as a latent source label (\Cref{prop:obs-attention-elbo}), builds selective communication,
and records its correspondence with transformer attention (\Cref{tab:obs-transformer-correspondence}).
\Cref{ch:relational-inference} develops fiber information geometry, Fisher contraction under Markov
morphisms (\Cref{thm:cg-fisher-contraction}), and inference histories measured by a Fisher duration that is
invariant under reparameterization (\Cref{thm:hist-fisher-clock-invariance}). \Cref{ch:two-timescales} types
fast beliefs and slow models (\Cref{def:obs-fast-slow-blocks}) and proves the reduced slow flow
(\Cref{prop:obs-rate-separated-flow}).

Part II coarse-grains at one base point. \Cref{ch:coarsegraining} separates Markov coarse channels, energy
restriction, and parameter aggregation (\Cref{sec:cg-three-operations}) and proves the exact coarse free
energy and the pointwise parent datum (\Cref{thm:rg-exact-coarse-vfe,thm:rg-pointwise-parent-datum}).
\Cref{ch:network-state} gives the gauge-retained network state, the kernel--holonomy fixed sectors
(\Cref{prop:cg-kernel-holonomy}), and the holonomy-conditioned marginal law
(\Cref{thm:cg-holonomy-kl-marginal}). \Cref{ch:agent-network-rg} is exact coarse-graining of the agent
network: effective action, residual, finite-step multipliers, and the native growth bound
(\Cref{thm:rg-effective-action,thm:rg-bounded-action-calculus,prop:rg-interaction-native-growth-bound}),
with the retained state collected in \Cref{tab:rg-retained-state} and a parent-sensitive mechanism with
path laws and an exact finite-memory cost (\Cref{sec:rg-memory-parent-mechanism}).
\Cref{ch:full-graph-meta-agent-vfe} assembles the multiscale free energy and closes a tower of priors
(\Cref{sec:full-graph-closing-tower}).

Part III turns to formation and interpretation. \Cref{ch:finite-depth-formation} constructs, for every
finite depth, a family in which successive coherence episodes come from linear relaxation of a prepared mean
through a supplied hierarchy of interaction strengths; source jumps and learning are perturbations the proof
controls, not the cause of formation, and a separate witness produces a four-cycle endogenously
(\Cref{prop:rg-formation-endogenous-cycle}). \Cref{ch:philosophy} states what the mathematics does and does
not support, and \Cref{ch:related-work} positions the construction. The appendices hold measure-theoretic
lemmas (\Cref{app:measure-lemmas}), the formation proofs (\Cref{app:finite-depth-formation-proofs}), the
notation (\Cref{app:notation-contract}), the open problems (\Cref{app:claim-ledger}), and the numerical index
(\Cref{app:numerical-provenance}).

Two supplements accompany the manuscript, and no result here depends on them. Supplement G develops the
multivariate-Gaussian realization. Supplement B collects base transport, covers and connection descent,
multi-point designs, infinite networks with their asymptotic apparatus, and the lattice cluster.

A first reading can take \Cref{ch:geometry,ch:probability,ch:generative} for the objects and the joint, then
\Cref{ch:elbo} for the bound and \Cref{ch:selective-communication} for attention, then
\Cref{ch:coarsegraining,ch:agent-network-rg} for exact coarse-graining and
\Cref{ch:finite-depth-formation} for formation; proofs, witnesses, and the ledger can be recovered through
their cross-references on a second reading.

\section{Standing conventions and scope}
\label{sec:intro-conventions}

Everything takes place at one base point. The design is $\{c_*\}$, design and incidence indices are dropped,
and every agent's frame, beliefs, and models are points of the fibers over $c_*$. Base maps and base curves,
neighborhoods, connection descent, transport along the base, section spaces over $\mathcal C$, and continuum
law theory are deferred to the base supplement named in \Cref{sec:intro-reading-map}.

The following defaults hold throughout, and a statement mentions only departures from them. Measurable spaces
are nonempty and standard Borel, and $\mathcal P(\mathsf S)$ denotes the probability laws on a measurable
space $\mathsf S$. Laws and kernels are normalized and measurable. A coarse channel is a Markov
kernel on latent coordinates that does not depend on the recognition law, leaves observations unchanged, and
is vertical over the identity of $c_*$. An admitted observation has finite positive evidence. Structural data
are parameters unless a configuration prior or a structural channel is declared, and the declaration names
its law. The interaction multigraph at $c_*$ is a structural parameter whose edge copies come in reversal
pairs; it contains every copy in the support of a receiving row's source prior and, unless a model adds
further copies, is the reversal closure of the union of these supports. Each support is complete unless
masked, and supports need not be symmetric: a receiver may attend to a source that does not attend to it. Gibbs attention weights are strictly positive on positive-prior edges when their
costs are finite; the extended-cost rows of \Cref{prop:full-graph-gibbs-rows} can assign zero weight to
infinite-cost edges. Dynamic structure is carried by attention weights and link values. Discrete creation and deletion of edges is not part of the model; it is
\Cref{open:obs-edge-creation}. Coarse multigraphs are induced structure. Absolute continuity of a recognition
law with respect to the posterior is not a default: it is clause (H3) of \Cref{hyp:elbo-evidence-domain}, and
its failure is the singular branch of the extended bound.

Three transports must remain distinct. A pointwise comparison $T_{ij}$ relates two frames in the same
principal fiber and involves no curve (\Cref{prop:geo-frame-comparisons}). A link $\Theta_e^x$ is an
independent variable on an edge copy of the interaction multigraph (\Cref{def:geo-graph-links}). Parallel
transport along a base curve is not used. Graph links are not base transport, and graph-link holonomy is
neither the holonomy of a base connection nor a property of the topology of the principal bundle.

The word coarse has two uses that must not be merged. A Markov channel discards sample information while
preserving normalization; a hard restriction substitutes a lower-dimensional configuration into an energy.
They can agree in a special realization, but a shared word transfers no theorem, and parameter aggregation
inside a closed family is a third operation (\Cref{sec:cg-three-operations}). Likewise, an inference history
is an oriented unparameterized orbit in a configuration space, an ordered interaction-record path retained by
exact coarse-graining is a sample coordinate, coarse-graining depth is a scale index, and a Fisher duration is
a statistical length.

\paragraph{Scope.}
The manuscript declines the following claims. None of the inference index, flow parameter, transition index,
ordered path coordinate, coarse-graining depth, or Fisher duration is identified with physical time; such an
identification would need a named target clock, a calibration map with an uncertainty model, causal and
compositional consistency, and discriminating evidence against alternative parameterizations. The contextual
base is not spacetime and carries no Lorentzian structure. Sections, laws, and meta-agents are mathematical
objects: calling a boundary node an agent is a naming decision, no claim of autonomous agency or of what
exists is made, and population laws, population samples, and blocks are not agents. No canonical coupling
selector, partition, or membership selector is asserted. No renormalization semigroup with rescaling, fixed
points, or beta functions is constructed; Supplement B holds that asymptotic apparatus. The identification of
any limiting object with a physical law is \Cref{claim:physical-law-identification}. Declining a claim is not
refuting it, and a refuted statement appears instead as an established negative result. \status{NOT-CLAIMED}

\section{Status of statements}
\label{sec:intro-status}

Every definition, standing hypothesis, lemma, proposition, theorem, corollary, conjecture, open problem, and
claim-bearing remark is a numbered statement. Definitions, standing hypotheses, conjectures, and open problems
carry their status in their heading type. Every lemma, proposition, theorem, corollary, and claim-bearing
remark carries exactly one of three tags, one for all of its conclusions. Text that motivates, summarizes, or
points carries no tag and makes no new claim. The fourth tag of \Cref{tab:status-taxonomy}, for declined
claims, appears once in each scope paragraph and nowhere else. Open problems are collected in \Cref{tab:open-problems}, and
numerical statements with their checks in \Cref{tab:numerical-index}.

\begin{sciencetable}{Status of statements.}
\label{tab:status-taxonomy}
\begin{tabularx}{\textwidth}{>{\raggedright\arraybackslash}p{0.24\textwidth}>{\raggedright\arraybackslash}X}
\toprule
\textbf{Tag or heading} & \textbf{Meaning} \\
\midrule
\status{ESTABLISHED} & Proved in this manuscript. A refutation is an established negative statement. \\
\status{STANDARD} & A known result, cited to a checked primary source with a theorem or section locator. \\
\status{NUMERICAL} & Supported only by a named computation, indexed with its control and seed; computation is not proof. \\
Definition & A declared type, construction, or convention; nothing is proved. \\
Hypothesis & A standing restriction adopted for a chapter or part; it states what it excludes. \\
Conjecture & A precise unproved statement with a named observation that would refute it. \\
Open problem & An unsettled question with what holds, what is unproved, and what would settle it. \\
\status{NOT-CLAIMED} & A claim declined in a scope paragraph; declining is not refuting. \\
\bottomrule
\end{tabularx}
\end{sciencetable}

\section*{Physicist's summary}

The manuscript studies agents that are local sections of belief and model bundles, all at one base point,
coupled through independent links and interaction records under a single normalized joint law. There is one
evidence bound for the whole population, the local objectives are exact conditional pieces of it, and
coarse-graining the agent network into meta-agents is done with Markov channels that preserve the evidence
and split the free energy exactly. Formation of meta-agents is proved for a declared finite-depth family in a
stated regime and conjectured more generally. Physical time, spacetime, agency, and transport across the base
are not claimed; they are either deferred to the base supplement or declined, and the defaults stated above
are not repeated in later chapters.

\chapter{Geometry at One Base Point}
\label{ch:geometry}

This chapter answers what geometric data an agent carries at one base point $c_*$ and which of those
data are invariant. The active gauge group at $c_*$ is one copy of $G$ inside the passive vertex group
$\prod_{i\in V}G$ (\Cref{prop:geo-active-image-point}); every independent link carries a
gauge-invariant conjugacy class, and loop holonomy is the telescoped product of these per-edge twists
(\Cref{prop:geo-edge-twist-invariance}); frame agreement, agreement of an equivariant statistic, and
law coherence up to holonomy stabilizers are distinct (\Cref{prop:geo-three-agreements}).
\Cref{ch:generative} uses passive re-expression for covariance of the generative law,
\Cref{ch:selective-communication} builds attention on link-transported comparisons, and
\Cref{ch:network-state} turns holonomy and the trivializing criterion
(\Cref{prop:geo-trivializing-criterion}) into fixed sectors of coarse channels.

\paragraph{Scope.}
Normalizability of a law on links built from the loop penalties of \Cref{rem:geo-loop-penalties} against
the Haar reference is not asserted for a noncompact or general linear $G$ (the compact unitary case is in
the remark, and a Gaussian reference on the Lie algebra normalizes the generator prior of
\Cref{prop:link-generator-prior} for every $G$), and
whether law coherence bears on the selection of a partition into meta-agents is part of
\Cref{open:cg-natural-selector}. \status{NOT-CLAIMED}

\section{The base point, the principal torsor, and associated law fibers}
\label{sec:geo-base}
\label{sec:geo-bundle}
\label{sec:geo-associated}

\definitionheading{Contextual base and base point}{def:geo-context-base}
Let $\mathcal C$ be a finite-dimensional, second countable, Hausdorff smooth manifold whose points
$c$ are called contexts, and fix one base point $c_*\in\mathcal C$.

\definitionheading{Common principal gauge bundle and agent frames}{def:geo-principal-systems}
Let $\pi$ be a smooth principal right $G$-bundle for a Lie group $G$ and, for each agent $i\in V$,
let $u_i$ be one local frame on an open support $\mathcal C_i\ni c_*$:
\begin{gather}
\pi:\mathscr P_G\longrightarrow\mathcal C,
\label{eq:geo-principal-bundle}\\
u_i:\mathcal C_i\longrightarrow\mathscr P_G,
\qquad
\pi\circ u_i=\operatorname{id}_{\mathcal C_i}.
\label{eq:geo-frame-sections}
\end{gather}
The fiber $(\mathscr P_G)_{c_*}$ is a right $G$-torsor and the values $u_i:=u_i(c_*)$ are points of
it. For any reference point $\sigma\in(\mathscr P_G)_{c_*}$ the frames are
$u_i=\sigma\cdot U_i^{-1}$ for unique $U_i\in G$.

The same frame $u_i$ serves the belief and the model channel; \Cref{prop:geo-common-frame-gauge}
shows that this is a gauge choice. On the belief and model sample fibers $(\mathsf K,\mathscr K)$ and
$(\mathsf M,\mathscr M)$ declare law fibers, measurable representations, which need not be equivalent,
faithful, or of equal dimension, and their pushforward actions on laws,
\begin{gather}
\mathcal B_b\subseteq\mathcal P(\mathsf K),
\qquad
\mathcal B_m\subseteq\mathcal P(\mathsf M),
\label{eq:geo-law-fibers}\\
\rho_b:G\longrightarrow\operatorname{Aut}(\mathsf K),
\qquad
\rho_m:G\longrightarrow\operatorname{Aut}(\mathsf M),
\label{eq:geo-representations}\\
\widehat\rho_b(g)q=(\rho_b(g))_\#q,
\qquad
\widehat\rho_m(g)q=(\rho_m(g))_\#q,
\label{eq:geo-pushforward-actions}
\end{gather}
so that $(\widehat\rho_x(g)q)(A)=q(\rho_x(g)^{-1}A)$.

\hypothesisheading{Selected smooth tier}{hyp:geo-smooth-tier}
Whenever a differential statement is made, $\mathcal B_b$ and $\mathcal B_m$ are finite-dimensional
smooth parametrized-measure models, each closed under its $G$-action, and the actions
$(g,q)\mapsto(\rho_x(g))_\#q$\label{eq:geo-smooth-actions} are smooth. The hypothesis excludes
stratified and infinite-dimensional law fibers; a statistical reading of the tier further requires
differentiability in quadratic mean and the integrability each tensor uses, and a Fisher tensor is a
metric only where it is nondegenerate.

\definitionheading{Belief and model associated bundles}{def:geo-associated-bundles}
On invariant law fibers define the bundles below, with projections $\varpi_x:\mathcal E_x\to\mathcal C$
and the quotient convention of the second line,
\begin{gather}
\mathcal E_b=\mathscr P_G\times_{\widehat\rho_b}\mathcal B_b,
\qquad
\mathcal E_m=\mathscr P_G\times_{\widehat\rho_m}\mathcal B_m,
\label{eq:geo-associated-bundles}\\
(u\cdot g,\beta)\sim(u,\widehat\rho_x(g)\beta).
\label{eq:geo-quotient-convention}
\end{gather}
At $c_*$ the frame $u_i$ identifies $\mathcal B_x$ with the fiber $(\mathcal E_x)_{c_*}$ by
$\beta\mapsto[u_i,\beta]$. This is the standard associated-bundle construction
\citep[Section~I.5]{Kobayashi1963}, smooth under \Cref{hyp:geo-smooth-tier}.

\section{Frames, passive re-expression, and the active gauge group at \texorpdfstring{$c_*$}{c*}}
\label{sec:geo-frame-gauge}

A passive re-expression of agent $i$ is a choice $g_i\in G$ that replaces the frame as in the first
line below. Since $[u_i,\beta]=[u_i'\cdot g_i,\beta]=[u_i',\widehat\rho_x(g_i)\beta]$, associated
coordinates change by $\beta'=\widehat\rho_x(g_i)\beta$, and the represented sample-coordinate maps are
those of the second line:
\begin{gather}
u_i'=u_i\cdot g_i^{-1},
\label{eq:geo-local-reframing}\\
R_i^b=\rho_b(g_i),
\qquad
R_i^m=\rho_m(g_i),
\qquad
k_i'=R_i^bk_i,\quad m_i'=R_i^mm_i.
\label{eq:geo-represented-frame-change}
\end{gather}
The group element, not a represented operator, multiplies the frame, and the tuple $(g_i)_{i\in V}$
ranges over the passive vertex group $\prod_{i\in V}G$.

\subsection{Pointwise frame comparisons and the active gauge group}
\label{sec:geo-regime-one}

\propositionheading{Frame comparisons at one base point}{prop:geo-frame-comparisons}
For $i,j\in V$ there is a unique $T_{ij}\in G$ with the first relation below, for every reference
point $\sigma$. The comparisons satisfy $T_{ii}=e_G$, $T_{ji}=T_{ij}^{-1}$, and $T_{ij}T_{jl}=T_{il}$,
and under passive re-expressions they transform by the second line:
\begin{gather}
u_j=u_i\cdot T_{ij},
\qquad
T_{ij}=U_iU_j^{-1},
\label{eq:geo-frame-comparison}\\
T_{ij}'=g_iT_{ij}g_j^{-1},
\qquad
U_i'=g_iU_i.
\label{eq:geo-regime-one-gauge-law}
\end{gather}
\status{ESTABLISHED}

\noindent\emph{Proof.}
Freeness and transitivity of the torsor give uniqueness of $T_{ij}$, and
$u_j=\sigma\cdot U_j^{-1}=u_i\cdot U_iU_j^{-1}$ gives the formula; the identities follow from
uniqueness. From $u_j'=u_j\cdot g_j^{-1}=u_i'\cdot g_iT_{ij}g_j^{-1}$ and
$u_i'=\sigma\cdot U_i^{-1}g_i^{-1}$ follow both laws in \eqref{eq:geo-regime-one-gauge-law}.
\hfill$\square$

\propositionheading{Active gauge image at one base point}{prop:geo-active-image-point}
Let $F\in\operatorname{Aut}_G((\mathscr P_G)_{c_*})$, the group of right-$G$-equivariant bijections of
the principal fiber over $c_*$, have local functions $F(u_i)=u_i\cdot\varphi_i$. Then
\begin{equation}
\varphi_j=T_{ij}^{-1}\varphi_iT_{ij},
\qquad i,j\in V,
\label{eq:geo-active-image-point}
\end{equation}
and $F\mapsto(\varphi_i)_{i\in V}$ is a group isomorphism onto the subgroup
$\mathcal A_{c_*}\subseteq\prod_{i\in V}G$ cut out by \eqref{eq:geo-active-image-point}. For each
$i_0\in V$ the projection $(\varphi_i)_{i\in V}\mapsto\varphi_{i_0}$ is an isomorphism
$\mathcal A_{c_*}\cong G$. In the frame $u_i$, $F$ changes the coordinates of both associated
fibers by $\beta\mapsto\widehat\rho_x(\varphi_i)\beta$, which is the coordinate change of the
passive re-expression $g_i=\varphi_i$. Hence the active transformations at $c_*$ form one copy
$\mathcal A_{c_*}\cong G$ of the passive vertex group, a proper subgroup of $\prod_{i\in V}G$ when
$G$ is nontrivial and $|V|\geq2$, of dimension $\dim G$ against $|V|\dim G$ for a Lie group. The
local functions at $c_*$ of every gauge automorphism of $\mathscr P_G$ form one element of
$\mathcal A_{c_*}$. \status{ESTABLISHED}

\noindent\emph{Proof.}
Equivariance gives $F(u_j)=F(u_i\cdot T_{ij})=u_i\cdot\varphi_iT_{ij}=u_j\cdot T_{ij}^{-1}\varphi_iT_{ij}$,
and uniqueness in the torsor gives \eqref{eq:geo-active-image-point}. Since
$F(F'(u))=F(u\cdot\varphi')=u\cdot\varphi\varphi'$, the map is a homomorphism. The value
$\varphi_{i_0}$ determines $F$ on the whole torsor, so the map is injective, and for
$\varphi_{i_0}\in G$ the map $u_{i_0}\cdot g\mapsto u_{i_0}\cdot\varphi_{i_0}g$ is an equivariant
bijection, which proves $\mathcal A_{c_*}\cong G$. The coordinate statement is
$F[u_i,\beta]=[u_i\cdot\varphi_i,\beta]=[u_i,\widehat\rho_x(\varphi_i)\beta]$, the passive coordinate
change with $g_i=\varphi_i$. If $G\neq\{e_G\}$ and $i\neq j$, the tuple with $\varphi_i=a\neq e_G$ and
all other entries $e_G$ violates \eqref{eq:geo-active-image-point}, so the subgroup is proper. A gauge
automorphism restricts to an equivariant bijection of each fiber. \hfill$\square$

\subsection{The common frame is a gauge choice}
\label{sec:geo-diagonal}

\propositionheading{Two atlases reduce to one by a passive re-expression}{prop:geo-common-frame-gauge}
Suppose belief and model coordinates were given in two frames $u_i^b,u_i^m$ of the same fiber. Then
there is a unique $h_i\in G$ with the first relation below, and under independent re-expressions
$u_i^{b\prime}=u_i^b\cdot(g_i^b)^{-1}$ and $u_i^{m\prime}=u_i^m\cdot(g_i^m)^{-1}$ it transforms by the
second:
\begin{gather}
u_i^m=u_i^b\cdot h_i,
\label{eq:geo-relative-frame}\\
h_i'=g_i^bh_i(g_i^m)^{-1}.
\label{eq:geo-relative-frame-law}
\end{gather}
In particular the re-expression $g_i^b=e_G$, $g_i^m=h_i$ of the model atlas alone makes the two
frames equal. In two atlases the belief and model comparisons, the local functions
$F(u_i^x)=u_i^x\cdot\varphi_i^x$ of a fiber automorphism $F$, and the representatives of a cross
morphism are related by
\begin{gather}
T_{ij}^m=h_i^{-1}T_{ij}^bh_j,
\label{eq:geo-frame-comparison-relation}\\
\varphi_i^m=h_i^{-1}\varphi_i^bh_i,
\label{eq:geo-diagonal-gauge-functions}\\
\Phi_i^{\mathrm{conv}}=\widehat\rho_m(h_i)\circ\Phi_i^{bm},
\label{eq:geo-map-bundle-frame-conversion}
\end{gather}
where $\Phi_i^{bm}$ is the representative of $\Phi$ with source frame $u_i^b$ and target frame
$u_i^m$, and $\Phi_i^{\mathrm{conv}}$ is its representative with both frames equal to $u_i^b$. All
three relations become identities when $h_i=e_G$. \status{ESTABLISHED}

\noindent\emph{Proof.}
Uniqueness of $h_i$ is the torsor property. Substituting the re-expressions gives
$u_i^{m\prime}=u_i^b\cdot h_i(g_i^m)^{-1}=u_i^{b\prime}\cdot g_i^bh_i(g_i^m)^{-1}$, which is
\eqref{eq:geo-relative-frame-law}. The two expressions $u_j^m=u_j^b\cdot h_j=u_i^b\cdot T_{ij}^bh_j$
and $u_j^m=u_i^m\cdot T_{ij}^m=u_i^b\cdot h_iT_{ij}^m$ give
\eqref{eq:geo-frame-comparison-relation}. Equivariance gives
$F(u_i^m)=F(u_i^b)\cdot h_i=u_i^b\cdot\varphi_i^bh_i=u_i^m\cdot h_i^{-1}\varphi_i^bh_i$. Finally,
$\Phi[u_i^b,\beta]=[u_i^m,\Phi_i^{bm}\beta]=[u_i^b,\widehat\rho_m(h_i)\Phi_i^{bm}\beta]$ by
\eqref{eq:geo-quotient-convention}. \hfill$\square$

No energy, kernel, or evaluator reads $h_i$, so the common frame
\begin{equation}
u_i^b=u_i^m=u_i,
\qquad
h_i=e_G,
\qquad
T_{ij}^b=T_{ij}^m=T_{ij}
\label{eq:geo-diagonal-principal}
\end{equation}
is adopted throughout.

\section{Agents and cross-associated-bundle morphisms}
\label{sec:geo-agents}
\label{sec:geo-cross-morphisms}

\definitionheading{Agent}{def:geo-agent}
Agent $i$ is the section-bearing object $\mathcal A^i$ of the first line below, with the sections of
the second line, and a population sample at $c_*$ is a point $y$ of a product of sample fibers:
\begin{gather}
\mathcal A^i=\bigl(\mathcal C_i;q_i^b,q_i^m,u_i\bigr),
\label{eq:geo-agent-object}\\
q_i^b\in\Gamma(\mathcal C_i,\mathcal E_b|_{\mathcal C_i}),
\qquad
q_i^m\in\Gamma(\mathcal C_i,\mathcal E_m|_{\mathcal C_i}),
\label{eq:geo-agent-sections}\\
y=\bigl((k_1,m_1),\ldots,(k_N,m_N)\bigr)\in\prod_{i\in V}(\mathsf K\times\mathsf M).
\label{eq:geo-population-sample}
\end{gather}
The values of $q_i^b,q_i^m$ at $c_*$ are points of $(\mathcal E_b)_{c_*}$ and $(\mathcal E_m)_{c_*}$, whose
coordinates in the frame $u_i$ are laws in $\mathcal B_b$ and $\mathcal B_m$; the same letters denote
these coordinate laws, which \Cref{def:prob-recognition-marginals} identifies with the local recognition
marginals. Agents need not have equal supports.

\definitionheading{Cross morphisms}{def:geo-cross-morphisms}
Cross morphisms are measurable fiber-preserving maps $\Phi:\mathcal E_b\to\mathcal E_m$ and
$\widetilde\Phi:\mathcal E_m\to\mathcal E_b$, $\varpi_m\circ\Phi=\varpi_b$ and
$\varpi_b\circ\widetilde\Phi=\varpi_m$.\label{eq:geo-cross-channel-map} At $c_*$ they are maps between
the fibers $(\mathcal E_b)_{c_*}$ and $(\mathcal E_m)_{c_*}$. They are not principal-bundle morphisms, gauge automorphisms, same-channel transports, or assumed inverses,
and their existence is additional structure. Under \Cref{hyp:geo-smooth-tier} a differential
statement also assumes them smooth. In the frame $u_i$ their representatives
$\Phi_i:\mathcal B_b\to\mathcal B_m$ and $\widetilde\Phi_i:\mathcal B_m\to\mathcal B_b$ are defined
by $\Phi[u_i,\beta]=[u_i,\Phi_i\beta]$ and $\widetilde\Phi[u_i,\beta]=[u_i,\widetilde\Phi_i\beta]$.

\propositionheading{Representatives and induced cross morphisms}{prop:geo-intertwining-cross-map}
The representatives transform under passive re-expression by the first two lines below, and those of
different agents are related by the third:
\begin{gather}
\Phi_i'=\widehat\rho_m(g_i)\circ\Phi_i\circ\widehat\rho_b(g_i)^{-1},
\label{eq:geo-phi-gauge-law}\\
\widetilde\Phi_i'=\widehat\rho_b(g_i)\circ\widetilde\Phi_i\circ\widehat\rho_m(g_i)^{-1},
\label{eq:geo-tildephi-gauge-law}\\
\Phi_i\circ\widehat\rho_b(T_{ij})=\widehat\rho_m(T_{ij})\circ\Phi_j,
\qquad
\widetilde\Phi_i\circ\widehat\rho_m(T_{ij})=\widehat\rho_b(T_{ij})\circ\widetilde\Phi_j.
\label{eq:geo-cross-map-gluing}
\end{gather}
If maps $\Phi_0:\mathcal B_b\to\mathcal B_m$ and $\widetilde\Phi_0:\mathcal B_m\to\mathcal B_b$
satisfy the first line below for every $g\in G$, then the maps of the second line
\begin{gather}
\Phi_0\circ\widehat\rho_b(g)=\widehat\rho_m(g)\circ\Phi_0,
\qquad
\widetilde\Phi_0\circ\widehat\rho_m(g)=\widehat\rho_b(g)\circ\widetilde\Phi_0,
\label{eq:geo-intertwining}\\
[p,\beta]\longmapsto[p,\Phi_0\beta],
\qquad
[p,\beta]\longmapsto[p,\widetilde\Phi_0\beta]
\label{eq:geo-induced-cross-morphisms}
\end{gather}
are well-defined cross morphisms, whose representatives in every frame are $\Phi_0$ and
$\widetilde\Phi_0$; conversely, $[p,\beta]\mapsto[p,\Phi_0\beta]$ is well defined only if $\Phi_0$
satisfies \eqref{eq:geo-intertwining}. Not every cross morphism is induced in this way, and for a
linear specialization with irreducible, inequivalent representations every equivariant linear map
between the sample fibers vanishes, while cross morphisms still exist. \status{ESTABLISHED}

\noindent\emph{Proof.}
From $[u_i,\beta]=[u_i',\widehat\rho_b(g_i)\beta]$ we get
$\Phi[u_i',\beta']=[u_i,\Phi_i\widehat\rho_b(g_i)^{-1}\beta']=[u_i',\widehat\rho_m(g_i)\Phi_i\widehat\rho_b(g_i)^{-1}\beta']$,
which is \eqref{eq:geo-phi-gauge-law}; interchanging the channels gives
\eqref{eq:geo-tildephi-gauge-law}. Evaluating $\Phi[u_j,\beta]=\Phi[u_i,\widehat\rho_b(T_{ij})\beta]$
in both frames gives \eqref{eq:geo-cross-map-gluing}. The representatives $(p\cdot g,\beta)$ and
$(p,\widehat\rho_b(g)\beta)$ of one class have the same image under
\eqref{eq:geo-induced-cross-morphisms} exactly when $\widehat\rho_m(g)\Phi_0\beta=\Phi_0\widehat\rho_b(g)\beta$,
which proves both directions of the induced statement. For a cross morphism that is not induced, let
$G=SO(2)$ act on the laws on $\R^2$ in both channels by the defining representation, fix
$\sigma\in(\mathscr P_G)_{c_*}$, and put $\Phi[\sigma,\beta]=[\sigma,D_\#\beta]$ with
$D=\operatorname{diag}(2,1)$, which is well defined because every class has a unique representative
with principal component $\sigma$. By \eqref{eq:geo-phi-gauge-law} its representative in the frame
$\sigma\cdot g^{-1}$ is $\widehat\rho_m(g)D_\#\widehat\rho_b(g)^{-1}$, and for the rotation $R_{\pi/2}$
one has $D_\#(R_{\pi/2})_\#\delta_{(1,0)}=\delta_{(0,1)}\neq\delta_{(0,2)}=(R_{\pi/2})_\#D_\#\delta_{(1,0)}$,
so the representative depends on the frame. In the linear specialization, the kernel and image of an
equivariant linear map between irreducible representation spaces are invariant subspaces, so the map is
zero or an isomorphism, and inequivalence excludes an isomorphism. \hfill$\square$

Law-level equivariance is weaker than linear equivariance: an equivariant Markov kernel induces an
equivariant map on laws without supplying a linear map between the sample fibers.

\section{Independent links at \texorpdfstring{$c_*$}{c*}: vertex law, holonomy, edge twists}
\label{sec:geo-regime-two}
\label{sec:geo-trivializing}
\label{sec:phil-agreements}

\definitionheading{Typed graph links}{def:geo-graph-links}
Let $\mathfrak I=(V,E,\bar{\ })$ be a finite interaction multigraph at $c_*$: a finite set $E$ of
oriented edge copies $e:j\to i$ with a fixed-point-free reversal $e\mapsto\bar e$, parallel copies
allowed. For each channel $x\in\{b,m\}$ and each $e:j\to i$ declare a link $\Theta_e^x$ as in the
first line below, with link map $\mathsf t_e^x(u_j\cdot g)=u_i\cdot\Theta_e^xg$ of the fiber; endpoint
notation $\Theta_{ij}^x$ is used only when the copy $j\to i$ is unique. Under passive re-expression the
links obey the vertex law of the second and third lines, with $R_i^x=\rho_x(g_i)$, because
$\mathsf t_e^x(u_j'\cdot g)=u_i\cdot\Theta_e^xg_j^{-1}g=u_i'\cdot g_i\Theta_e^xg_j^{-1}g$. For a walk
$p=(e_0,\ldots,e_{r-1})$ from $j$ to $i$, that is, with $e_a:i_{a+1}\to i_a$, $i_0=i$, and $i_r=j$, the
path product and, when $i=j$, the holonomy are those of the fourth line, and the common-link
specialization of the fifth line declares that the two link maps coincide in the common frame:
\begin{gather}
\Theta_e^x\in G,
\qquad
\Theta_{\bar e}^x=(\Theta_e^x)^{-1},
\label{eq:geo-regime-two-links}\\
\Theta_e^{x\prime}=g_i\Theta_e^xg_j^{-1},
\label{eq:geo-regime-two-gauge-law}\\
\rho_x(\Theta_e^{x\prime})=R_i^x\rho_x(\Theta_e^x)(R_j^x)^{-1},
\label{eq:gen-represented-gauge-links}\\
\Theta_p^x=\Theta_{e_0}^x\Theta_{e_1}^x\cdots\Theta_{e_{r-1}}^x,
\qquad
\mathrm{Hol}^x(p)=\Theta_p^x,
\label{eq:geo-link-holonomy}\\
\Theta_e^b=\Theta_e^m=\Theta_e.
\label{eq:geo-common-link}
\end{gather}

\propositionheading{Per-edge twist invariance at one base point}{prop:geo-edge-twist-invariance}
Fix a reference point $\sigma$ with $u_i=\sigma\cdot U_i^{-1}$, and define the per-edge twist
$\mathring\Theta_e^x$ of $e:j\to i$ by the first line below; a closed walk $p$ based at $i_0$ then has
the holonomy of the second line:
\begin{gather}
\mathsf t_e^x(\sigma)=\sigma\cdot\mathring\Theta_e^x,
\qquad
\mathring\Theta_e^x=U_i^{-1}\Theta_e^xU_j,
\label{eq:geo-edge-twist}\\
\mathrm{Hol}^x(p)=U_{i_0}\mathring\Theta_{e_0}^x\mathring\Theta_{e_1}^x\cdots\mathring\Theta_{e_{r-1}}^xU_{i_0}^{-1},
\qquad
\mathrm{Hol}^x(p)'=g_{i_0}\mathrm{Hol}^x(p)g_{i_0}^{-1}.
\label{eq:geo-edge-twist-holonomy}
\end{gather}
Then $\mathsf t_e^x$ is a right-$G$-equivariant bijection that does not depend on the frames, and
$\mathsf t_{\bar e}^x=(\mathsf t_e^x)^{-1}$. Every twist is unchanged by arbitrary passive
re-expressions; a change of reference point $\sigma'=\sigma\cdot g_0$ conjugates every twist by the
same $g_0$; and an active $F$ with $F(\sigma)=\sigma\cdot\kappa$ sends $\mathsf t_e^x$ to
$F\circ\mathsf t_e^x\circ F^{-1}$ and every twist to $\kappa\mathring\Theta_e^x\kappa^{-1}$. Each edge
copy therefore carries the conjugacy class $[\mathring\Theta_e^x]$, invariant under all passive and
active changes, and two link data sets on the same multigraph are related by one active
automorphism exactly when their twist tuples, computed from one reference point, are related by one
simultaneous conjugation. The conjugacy class of $\mathrm{Hol}^x(p)$ and the condition
$\mathrm{Hol}^x(p)=e_G$ are invariant. \status{ESTABLISHED}

\noindent\emph{Proof.}
Every point of the fiber is $u_j\cdot g$ for a unique $g$, so $\mathsf t_e^x$ is a well-defined
equivariant bijection whose inverse $u_i\cdot g\mapsto u_j\cdot(\Theta_e^x)^{-1}g$ is $\mathsf t_{\bar e}^x$,
and by the computation in \Cref{def:geo-graph-links} re-expressed frames and links define the same
map. Since $\sigma=u_j\cdot U_j$, $\mathsf t_e^x(\sigma)=u_i\cdot\Theta_e^xU_j=\sigma\cdot
U_i^{-1}\Theta_e^xU_j$; the left side depends only on $\mathsf t_e^x$ and $\sigma$, which proves passive
invariance, and $\mathsf t_e^x(\sigma\cdot g_0)=(\sigma\cdot g_0)\cdot g_0^{-1}\mathring\Theta_e^xg_0$.
For the active statement,
$F\mathsf t_e^xF^{-1}(\sigma)=F(\sigma\cdot\mathring\Theta_e^x\kappa^{-1})=\sigma\cdot\kappa\mathring\Theta_e^x\kappa^{-1}$;
an equivariant self-map of the torsor is determined by its value at $\sigma$, and $F\mapsto\kappa$ is a
bijection onto $G$, which gives the criterion for two link data sets. Finally
$\Theta_{e_a}^x=U_{i_a}\mathring\Theta_{e_a}^xU_{i_{a+1}}^{-1}$ telescopes, and $U_{i_0}'=g_{i_0}U_{i_0}$ by
\eqref{eq:geo-regime-one-gauge-law}. \hfill$\square$

\propositionheading{Trivializing criterion}{prop:geo-trivializing-criterion}
\label{prop:gauss-coboundary-trivialization}
Fix a subfamily $I\subseteq V$ and elements $\Xi_e\in G$ on the internal oriented edge copies of $I$.
The following are equivalent:
\begin{equation}
\begin{gathered}
\exists V_i\in G\ (i\in I)
\quad
\Xi_e=V_iV_j^{-1}
\quad\text{for every internal }e:j\to i,\\
\Xi_{e_0}\Xi_{e_1}\cdots\Xi_{e_{r-1}}=e_G
\quad\text{for every closed internal walk }(e_0,\ldots,e_{r-1}).
\end{gathered}
\label{eq:geo-trivializing-criterion}
\end{equation}
When they hold, the potentials $V_i$ are unique up to right multiplication by one element of $G$ on
each connected component of the internal multigraph. The statement and proof use only the group
structure, so they hold for any group in place of $G$, in particular for a subgroup of $\GL(K)$
containing represented links. \status{ESTABLISHED}

\noindent\emph{Proof.}
The first condition telescopes around every closed walk. Conversely, backtracks
$(e,\bar e)$ are closed walks, so $\Xi_{\bar e}=\Xi_e^{-1}$ and reversed walks have inverse products.
Choose a base vertex in each connected component and set $V_i=\Xi_{p}$ for an internal walk $p$ from
that vertex to $i$, with products ordered as in \eqref{eq:geo-link-holonomy}; two choices differ by a
closed walk and agree, and extending a walk to $j$ by $e:j\to i$ gives $V_i=\Xi_eV_j$. If also
$\Xi_e=W_iW_j^{-1}$, then $W_i^{-1}V_i=W_j^{-1}V_j$ along every internal edge, so $W_i^{-1}V_i$ is
constant on each component. \hfill$\square$

\definitionheading{Graph-flat, jointly trivializing, and frame-aligned links}{def:geo-graph-flat-links}
The links of channel $x$ are graph flat on $I\subseteq V$ when \Cref{prop:geo-trivializing-criterion}
holds for $\Xi_e=\Theta_e^x$, and $I$ is jointly trivializing when this holds in both channels. The
links are frame aligned\label{hyp:geo-flat-links} when $\Theta_e^x=T_{ij}=U_iU_j^{-1}$ for every
$e:j\to i$,\label{eq:geo-flat-regime} equivalently when every twist $\mathring\Theta_e^x$ equals $e_G$
and every link map $\mathsf t_e^x$ is the identity of the fiber.

\remarkheading{Frame alignment, graph flatness, and per-edge classes}{rem:geo-flatness-strictness}
Frame alignment implies graph flatness, because $T_{ij}=U_iU_j^{-1}$ is a coboundary, but not
conversely: on two vertices joined by one edge copy $e:2\to1$ and its reversal, every closed walk is
a backtrack, so any $\Theta_e^x\neq T_{12}$ is graph flat with potentials $V_1=\Theta_e^x$,
$V_2=e_G$, while $[\mathring\Theta_e^x]\neq\{e_G\}$. Graph flatness of the links is graph flatness
of the twists, since $\Theta_e^x=V_iV_j^{-1}$ exactly when $\mathring\Theta_e^x=W_iW_j^{-1}$ with
$W_i=U_i^{-1}V_i$. A single link coordinate between distinct vertices carries no invariant: the passive re-expression
$g_i=e_G$, $g_j=\Theta_e^x$ sends $\Theta_e^x$ to $e_G$. On a tree every loop product is trivial,
yet the per-edge classes can all differ from $\{e_G\}$. Graph flatness is a property of the links
alone: the one fiber over $c_*$ is always a trivial torsor, links on a triangle with nontrivial loop
product are not graph flat, and all-identity links are graph flat whatever the topology of
$\mathscr P_G$. \status{ESTABLISHED}

\remarkheading{Invariant loop penalties}{rem:geo-loop-penalties}
Let $p$ be a closed walk based at $i_0$ and $\mathsf U=\rho_x(\mathrm{Hol}^x(p))$. If $\rho_x$ is unitary
of dimension $n$, then
$\frac{1}{2n}\operatorname{Re}\Tr(2I_n-\mathsf U-\mathsf U^*)=\frac{1}{2n}\Tr((I_n-\mathsf U)^*(I_n-\mathsf U))\geq0$,
with equality exactly when $\mathsf U=I_n$, and the value is invariant under passive re-expression,
which conjugates $\mathsf U$ by the unitary $R_{i_0}^x=\rho_x(g_{i_0})$. For $n\geq2$ no norm on
$\mathfrak{gl}(n,\R)$ is invariant under conjugation by all of $\GL(n,\R)$: conjugation by
$\operatorname{diag}(2,1,\ldots,1)$ doubles the matrix unit $E_{12}$. A positive invariant penalty for
general linear links therefore needs further data: if a symmetric positive definite fiber metric
$\mathsf G^{\mathrm{fib}}$ at $i_0$ is declared to transform by
$\mathsf G^{\mathrm{fib}}\mapsto(R_{i_0}^x)^{-\top}\mathsf G^{\mathrm{fib}}(R_{i_0}^x)^{-1}$, then
$\Tr\bigl((\mathsf G^{\mathrm{fib}})^{-1}(I_n-\mathsf U)^\top\mathsf G^{\mathrm{fib}}(I_n-\mathsf U)\bigr)$
is invariant, nonnegative, and zero exactly when $\mathsf U=I_n$. For compact $G$ with $\rho_x$
unitary, the penalty $\frac{1}{2n}\Tr((I_n-\mathsf U)^*(I_n-\mathsf U))$ of the first sentence lies
in $[0,2]$, so for finitely many closed walks $p$ with holonomies
$\mathsf U_p=\rho_x(\mathrm{Hol}^x(p))$ and coefficients $\mathrm c_p\geq0$ the law on links with
density $\exp\bigl(-\sum_p\mathrm c_p\frac{1}{2n}\Tr((I_n-\mathsf U_p)^*(I_n-\mathsf U_p))\bigr)$
against the product, over one orientation of each edge copy, of the Haar probability measure of $G$
has normalizer in $[e^{-2\sum_p\mathrm c_p},1]$. \status{ESTABLISHED}

\noindent\emph{Proof.}
Unitarity gives $(I_n-\mathsf U)^*(I_n-\mathsf U)=2I_n-\mathsf U-\mathsf U^*$, whose trace is real and
equals the squared Frobenius norm of $I_n-\mathsf U$. With $\Delta=I_n-\mathsf U$ and
$\mathsf G=\mathsf G^{\mathrm{fib}}$, passive re-expression sends $\Delta$ to
$R_{i_0}^x\Delta(R_{i_0}^x)^{-1}$, cyclicity of the trace gives invariance, and
$\Tr(\mathsf G^{-1}\Delta^\top\mathsf G\Delta)=\lVert\mathsf G^{1/2}\Delta\mathsf G^{-1/2}\rVert_{\mathrm F}^2$.
A norm $\lVert\cdot\rVert_\star$ with $\lVert E_{12}\rVert_\star>0$ cannot satisfy
$\lVert2E_{12}\rVert_\star=\lVert E_{12}\rVert_\star$. For the compact case,
$\Tr(2I_n-\mathsf U-\mathsf U^*)=2n-2\operatorname{Re}\Tr\mathsf U$ and $\lvert\Tr\mathsf U\rvert\leq n$
for unitary $\mathsf U$, so that penalty lies in $[0,2]$ and the exponent in $[-2\sum_p\mathrm c_p,0]$;
each $\mathsf U_p$ is a continuous function of the links, and integrating against a probability
measure gives the two bounds on the normalizer. \hfill$\square$

In the exponential chart of \Cref{ch:links}, with the fiber metric at each vertex induced by one matrix at
the reference point, one half of the general linear penalty above is the loop defect of
\Cref{def:link-loop-defect} (\Cref{prop:link-loop-defect}); and for $G=\GL(K,\R)$ or $\GL^+(K)$, on a
closed walk that traverses some reversal pair exactly once, every continuous penalty of the holonomy that
is unchanged by passive re-expression takes at every unipotent holonomy its value at trivial holonomy
(\Cref{cor:link-no-invariant-defect}), so a penalty that vanishes only at trivial holonomy needs data
beyond the holonomy, such as the fiber metric.

\definitionheading{Law coherence}{def:geo-law-coherence}
Let $I\subseteq V$ induce a connected internal multigraph. Laws $q_i\in\mathcal B_x$, $i\in I$, are law
coherent on $I$ when
\begin{equation}
q_i=\widehat\rho_x(\Theta_e^x)q_j
\qquad\text{on each internal edge copy }e:j\to i.
\label{eq:phil-law-coherence}
\end{equation}
For a $G$-set $S$ and a map $s:\mathcal B_x\to S$ with $s(\widehat\rho_x(g)q)=g\cdot s(q)$, the laws are
$s$-agreeing when $s(q_i)=\Theta_e^x\cdot s(q_j)$ on each internal edge copy; the mean of laws on a
linear representation space is the basic instance. The links have frame agreement on $I$ when they
are graph flat on $I$.

\propositionheading{Three agreements}{prop:geo-three-agreements}
Let $I$ be as in \Cref{def:geo-law-coherence} and fix $i_0\in I$. (i) The laws are law coherent
exactly when $q_{i_0}$ is fixed by $\widehat\rho_x(\mathrm{Hol}^x(p))$ for every closed internal walk
$p$ based at $i_0$ and $q_i=\widehat\rho_x(\Theta^x_{p_i})q_{i_0}$ for some, hence every, internal walk
$p_i$ from $i_0$ to $i$; for an equivariant statistic the same holds with $s(q_{i_0})$ fixed by the
holonomy, so that frame agreement and mean agreement on edges, the two conditions of the first line
below, are conditions on the links and on the means, and mean agreement holds exactly when the base
mean satisfies the second line and the other means are its transports:
\begin{gather}
\mathrm{Hol}^x(p)=e_G\ \text{for every closed internal walk }p,
\quad
\mu_i=\rho_x(\Theta_e^x)\mu_j\ \text{on each internal }e:j\to i,
\label{eq:phil-two-agreements}\\
\mu_{i_0}\in\bigcap_{p\text{ closed at }i_0}\ker\bigl(\rho_x(\mathrm{Hol}^x(p))-\operatorname{id}\bigr).
\label{eq:phil-agreement-kernel}
\end{gather}
(ii) Law coherence implies $s$-agreement for every equivariant statistic $s$, and for the mean the converse fails. (iii) Frame agreement neither implies
nor is implied by law coherence or by mean agreement. (iv) With several independent cycles, trivial
holonomy or a fixed base mean for one loop is necessary but not sufficient. In particular, for
$G=\GL(2,\R)$ acting on laws on $\R^2$, three agents with the common law $\mathcal N(0,\Sigma)$,
$\Sigma=\begin{psmallmatrix}4&1\\1&1\end{psmallmatrix}$, and links
$\Theta_{12}=\Sigma^{1/2}R(0.7)\Sigma^{-1/2}$ and $\Theta_{23}=\Theta_{31}=I_2$, with $R(\phi)$ the
rotation by $\phi$, are law coherent while the triangle holonomy $\Theta_{12}$ is not the identity, is
not orthogonal, and fixes no nonzero vector. \status{ESTABLISHED}

\noindent\emph{Proof.}
(i) Composing \eqref{eq:phil-law-coherence} along an internal walk $p$ from $j$ to $i$ gives
$q_i=\widehat\rho_x(\Theta^x_p)q_j$, which yields both conditions. Conversely, let the
conditions hold; for $e:j\to i$ and walks $p_i,p_j$ from $i_0$, the walk that follows $p_j$, then $e$,
then $p_i$ backward is closed at $i_0$ with holonomy $(\Theta_{p_i}^x)^{-1}\Theta_e^x\Theta_{p_j}^x$, so
$\widehat\rho_x(\Theta_e^x)q_j=\widehat\rho_x(\Theta^x_{p_i})\widehat\rho_x\bigl((\Theta_{p_i}^x)^{-1}\Theta_e^x\Theta_{p_j}^x\bigr)q_{i_0}=q_i$.
The statistic version is the same argument in $S$, and the vectors left unchanged by the linear map
$\rho_x(\mathrm{Hol}^x(p))$ form $\ker(\rho_x(\mathrm{Hol}^x(p))-\operatorname{id})$. (ii) Apply $s$ to
\eqref{eq:phil-law-coherence}. For the converse take two agents joined by one edge copy with identity
link and the laws $\mathcal N(0,I_2)$ and $\mathcal N(0,2I_2)$: the means agree and the laws do not.
(iii) The same two agents with the laws $\mathcal N(0,I_2)$ and $\mathcal N(v,I_2)$, $v\neq0$, have frame
agreement without mean agreement or law coherence. In the displayed witness,
$\Theta_{12}\Sigma\Theta_{12}^\top=\Sigma^{1/2}R(0.7)R(0.7)^\top\Sigma^{1/2}=\Sigma$, so every transported
law equals $\mathcal N(0,\Sigma)$ and all means vanish, while $\Theta_{12}$ is similar to a rotation by
$0.7$, hence has no eigenvalue $1$, and $\Theta_{12}\Theta_{12}^\top\neq I_2$ because $\Sigma$ is not
a multiple of $I_2$. (iv) Two vertices joined by three parallel copies with $\Theta_e=\Theta_f\neq\Theta_g$
have trivial holonomy on the walk through $e$ and $\bar f$ and nontrivial holonomy on the walk through
$e$ and $\bar g$. \hfill$\square$

\section{Running example: three agents on one loop}
\label{sec:geo-running-example}

Take $G=SO(3)$, $\rho_b$ the defining representation on $\mathsf K=\R^3$, and $\rho_m$ the trivial
representation on $\mathsf M=\R$, so the channels carry inequivalent representations. Let
$V=\{1,2,3\}$, let $\mathfrak I$ be the triangle with edge copies $e_1:2\to1$, $e_2:3\to2$, $e_3:1\to3$
and their reversals, and, with $R_z(\phi)$ and $R_x(\phi)$ the rotations by $\phi$ about the third and
first coordinate axes, declare the belief links of the first line below; model links act trivially.
Choose $\sigma=u_1$ and, for display, $u_2=u_3=u_1$, so $U_i=e_G$ and each twist equals its link. The
per-edge classes are the angles $\pi/2$, $\pi/2$, and $0$, and the loop $p=(e_1,e_2,e_3)$ based at
agent $1$ has the holonomy of the second line,
\begin{gather}
\Theta_{e_1}^b=R_z(\tfrac\pi2),
\qquad
\Theta_{e_2}^b=R_x(\tfrac\pi2),
\qquad
\Theta_{e_3}^b=I_3,
\label{eq:geo-running-links}\\
\mathrm{Hol}^b(p)=R_z(\tfrac\pi2)R_x(\tfrac\pi2),
\qquad
\Tr\mathrm{Hol}^b(p)=0,
\label{eq:geo-running-holonomy}
\end{gather}
a rotation by $2\pi/3$; two quarter turns about one axis would give a rotation by $\pi$ with the same
classes, so the class tuple does not determine the holonomy, while the tuple modulo one simultaneous
conjugation does (\Cref{prop:geo-edge-twist-invariance}). A passive re-expression $g_2$ changes
$\Theta_{e_1}^b$ to $\Theta_{e_1}^bg_2^{-1}$ and $\Theta_{e_2}^b$ to $g_2\Theta_{e_2}^b$ but changes no
twist, recomputed with $U_2'=g_2$, and not the holonomy. The loop is not graph flat, and the isotropic
laws $q_i^b=\mathcal N(0,s^2I_3)$ are law coherent on it, because every rotation fixes them
(\Cref{prop:geo-three-agreements}). One record reads the transported discrepancy of agents $1$ and $2$:
the Markov kernel from $\mathsf K\times\mathsf K$ to $\{0,1\}$ with
\begin{equation}
K_{a_1}\bigl(k_1,k_2,\{1\}\bigr)=\exp\Bigl(-\tfrac12\bigl\lvert k_1-\rho_b(\Theta_{e_1}^b)k_2\bigr\rvert^2\Bigr),
\label{eq:geo-running-record}
\end{equation}
which is gauge inert, since passive re-expressions rotate $k_1-\rho_b(\Theta_{e_1}^b)k_2$ by
$\rho_b(g_1)$.

\section*{Physicist's summary}

At one base point all frames lie in one principal fiber, so they are compared directly by
$T_{ij}=U_iU_j^{-1}$, the only active symmetry is one global rotation of that fiber, and per-agent frame
changes are bookkeeping. Each independent link carries an invariant, its conjugacy class after dressing
by the frames, whereas in lattice gauge theory only closed-loop products are invariant
\citep{WilsonConfinement1974}. Agreement comes in three kinds, flat links, agreement of a statistic such
as the mean, and law coherence, which survives nontrivial holonomy when the base law has a nontrivial
stabilizer.

\chapter{Probability Semantics at \texorpdfstring{$c_*$}{c*}}
\label{ch:probability}

On which measurable objects, and in which order, does every probabilistic statement at $c_*$ rest? Laws
and kernels come first, densities follow as Radon--Nikodym derivatives against references fixed for a
whole family, and logarithms of density ratios enter only under
\Cref{def:prob-log-density-ratio-convention}; this order makes the evidence bound of \Cref{ch:elbo} an
identity in the extended reals. The main results are common domination of a mixed family exactly when its
atoms are countable (\Cref{thm:prob-mixed-family-common-domination}), the relative-entropy chain rule in
$[0,+\infty]$ (\Cref{lem:prob-kl-chain-rule}) that every later decomposition cites, and three levels of
description, at which a law-point energy is not a sample energy and point-mass recognition of law points
has infinite cost (\Cref{prop:prob-law-point-energy,prop:prob-law-point-point-mass}).
\Cref{ch:generative} builds the joint on these spaces, \Cref{ch:elbo} applies the chain rule, and
\Cref{ch:selective-communication,ch:finite-depth-formation} use the levels.

\paragraph{Scope.}
Injectivity of a model evaluator is not assumed, and no regularity of the quotient of presentations with
the same evaluated kernel is supplied; presentations are retained. The coupling selector is applied once,
and its equivariance under frame re-expression (rerunning it on transformed marginals returns the
transformed coupling) is neither assumed nor claimed. Open are a normalized law-point prior
(\Cref{open:prob-law-point-prior}) and likelihoods for record families without a common dominating
measure (\Cref{open:prob-nondominated-likelihood}). \status{NOT-CLAIMED}

\section{Latent and record spaces at \texorpdfstring{$c_*$}{c*}}
\label{sec:prob-design}

\definitionheading{Sample spaces and product sigma-algebras}{def:prob-product-sample-spaces}
For each agent $i\in V$ let $(\mathsf K_i,\mathscr K_i)$ and $(\mathsf M_i,\mathscr M_i)$ be its belief
and model sample fibers at $c_*$, identified with $\mathsf K$ and $\mathsf M$ by the frame $u_i$, and put
$\mathsf Y_i=\mathsf K_i\times\mathsf M_i$. Let $\mathfrak A$ be a finite set of records; each
$a\in\mathfrak A$ has a scope $\partial a\subseteq V$ and a record space $(\mathsf O_a,\mathscr O_a)$. The
latent and observation spaces are
\begin{align}
\mathsf Y&=\prod_{i\in V}\mathsf Y_i,
&\mathscr Y&=\bigotimes_{i\in V}(\mathscr K_i\otimes\mathscr M_i),
\label{eq:prob-latent-space}\\
\mathsf O&=\prod_{a\in\mathfrak A}\mathsf O_a,
&\mathscr O&=\bigotimes_{a\in\mathfrak A}\mathscr O_a,
\label{eq:prob-observation-space}
\end{align}
with points $y=((k_i,m_i))_{i\in V}$ and $o=(o_a)_{a\in\mathfrak A}$, and
$\mathsf Y_{\partial a}=\prod_{i\in\partial a}\mathsf Y_i$. An observation of a single agent is a record
with singleton scope.

A population sample is a point of $\mathsf Y$, and evaluating a sample field at $c_*$ returns
\begin{equation}
\operatorname{Ev}(y)=\bigl((k_i(c_*),m_i(c_*))\bigr)_{i\in V}\in\mathsf Y,
\label{eq:prob-evaluation-map}
\end{equation}
whereas the section values $q_i^b,q_i^m$ at $c_*$ are laws on $\mathsf K_i$ and $\mathsf M_i$, which become
latent coordinates only at the law-point level of \Cref{sec:prob-levels}. Missing or aggregated
observations are represented by an observation kernel into a coarser record space, which changes
$\mathsf O$ and leaves $\mathsf Y$ unchanged.

\section{Structural data and kernel signatures}
\label{sec:prob-structural}

\definitionheading{Structural configuration and kernel signatures}{def:prob-structural-kernel-signatures}
Fix a measurable configuration space $(\mathsf X,\mathscr X)$. Every $X\in\mathsf X$ uses the same agent
set $V$, the same records and scopes, and the same fiber identifications, so $(\mathsf Y,\mathscr Y)$ and
$(\mathsf O,\mathscr O)$ do not vary with $X$. A configuration carries the interaction multigraph
$\mathfrak I$ of \Cref{def:geo-graph-links}, its links $\Theta_e^x$, the frames $u_i$, and the parameters
of the mechanisms. With $(\mathsf I_i,\mathscr I_i)$ an access space for each $i\in V$, the access maps,
local recognition kernels, record kernels, and population laws have the signatures
\begin{align}
\operatorname{Acc}_i&:\mathsf O\times\mathsf X\longrightarrow\mathsf I_i,
\quad\mathbb Q_i:\mathsf I_i\rightsquigarrow\mathsf Y_i,
&\mathbb Q_{i,o,X}&:=\mathbb Q_i(\operatorname{Acc}_i(o,X),\cdot),
\label{eq:prob-local-recognition-kernel}\\
K_a&:(\mathsf X\times\mathsf Y_{\partial a})\rightsquigarrow\mathsf O_a,
&a&\in\mathfrak A,
\label{eq:prob-record-kernel-signature}\\
\mathbb P_{\theta,V}&:(\mathsf X,\mathscr X)\longrightarrow\mathcal P\bigl(\mathsf Y\times\mathsf O,\mathscr Y\otimes\mathscr O\bigr),
&X&\longmapsto\mathbb P_{\theta,V}(\cdot\given X),
\label{eq:prob-generative-signature}\\
\mathbb Q_V&:(\mathsf O\times\mathsf X,\mathscr O\otimes\mathscr X)\longrightarrow\mathcal P\bigl(\mathsf Y,\mathscr Y\bigr),
&(o,X)&\longmapsto\mathbb Q_{V,o,X},
\label{eq:prob-recognition-signature}
\end{align}
where each record coordinate occurs once, $\mathbb P_{\theta,V}$ is constructed in \Cref{ch:generative}, and
$\mathbb Q_{V,o,X}$ is the coupling selected in \Cref{sec:prob-marginals}; neither is primitive data. The
identity access map is the unrestricted interface; a stricter one restricts what the local law may depend
on, not the population law (\Cref{prop:prob-access-restriction}). An agent's generative-model structure is
the pair $(q_i^m,\operatorname{ev}_i)$ of its model law and the evaluator of \Cref{def:prob-model-evaluation}.
The generative kernel reads only $\theta$, $X$, and latent coordinates (\Cref{req:gen-typing-prohibition}),
and a configuration prior, when declared, is a law on $(\mathsf X,\mathscr X)$.

\definitionheading{Model presentations and evaluated kernels}{def:prob-model-evaluation}
For each $i\in V$ and structural value $X$, let $(\mathsf S_i^X,\mathscr S_i^X)$ and
$(\mathsf T_i^X,\mathscr T_i^X)$ contain every source and output coordinate of agent $i$'s generative
factor, and let $(\mathsf M_i,\mathscr M_i)$ be its presentation space, so that a model sample is
$m_i\in\mathsf M_i$ and the model law is $q_i^m\in\mathcal P(\mathsf M_i)$. The evaluator is a map
\begin{equation}
\operatorname{ev}_i:\mathsf M_i\longrightarrow\operatorname{Kern}(\mathsf S_i^X,\mathsf T_i^X),
\qquad
m_i\longmapsto K^X_{i,m_i},
\label{eq:prob-model-evaluation}
\end{equation}
with $(m_i,z)\mapsto K^X_{i,m_i}(z,B)$ measurable for every $B\in\mathscr T_i^X$ and
$K^X_{i,m_i}(z,\mathsf T_i^X)=1$. Each source space retains every observation, structural, or parent
coordinate its factor requires. In the definite case $q_i^m=\delta_{m_i}$. Neither a model law, a
recognition parameter, nor a posterior is an argument of an evaluated kernel.

\section{Reference measures}
\label{sec:prob-reference}

\definitionheading{Reference measures}{def:prob-reference-measures}
Each Euclidean coordinate carries its Borel sigma-algebra and Lebesgue measure, and each countable
coordinate its power set and counting measure. A mixed real coordinate carries one parameter-independent
reference for the whole family of laws placed on it, generative and recognition alike: when the union
$A_*$ of the atom sets of that family is countable, it carries $\nu_{\mathrm{mix}}$ of
\eqref{eq:prob-mixed-family-reference}; when $A_*$ is uncountable no density model is declared on that
coordinate (\Cref{thm:prob-mixed-family-common-domination}). Any other fiber carries an explicitly
declared sigma-finite measure. With $\nu_i^k,\nu_i^m,\nu_a$ the component measures, the product
references are
\begin{equation}
\nu^{\mathsf Y}=\bigotimes_{i\in V}\bigl(\nu_i^k\otimes\nu_i^m\bigr),
\qquad
\nu^{\mathsf O}=\bigotimes_{a\in\mathfrak A}\nu_a.
\label{eq:prob-reference-measures}
\end{equation}

\theoremheading{Common domination for a mixed family}{thm:prob-mixed-family-common-domination}
Let $(\mu_t)_{t\in\mathbb T}$ be any family of Borel probability measures on $\R$ with no
singular-continuous part, written $\mu_t=f_t\operatorname{Leb}+\sum_{a\in A_t}p_t(a)\delta_a$ with
$A_t:=\{a\in\R:\mu_t(\{a\})>0\}$. A single sigma-finite Borel measure dominates every $\mu_t$ if and only
if $A_*:=\bigcup_{t\in\mathbb T}A_t$ is countable. When $A_*=\{a_1,a_2,\ldots\}$, one such
parameter-independent measure is
\begin{equation}
\nu_{\mathrm{mix}}:=\operatorname{Leb}+\sum_{n\geq1}2^{-n}\delta_{a_n}.
\label{eq:prob-mixed-family-reference}
\end{equation}
In particular, the moving-atom family $\mu_t=\tfrac12\mathcal N(0,1)+\tfrac12\delta_t$, $t\in\R$, has no
common sigma-finite dominating measure. \status{ESTABLISHED}

The proof is in \Cref{app:measure-lemmas}.

\remarkheading{Dirac families are not dominated}{sec:exp-dominated-nondominated}
On an uncountable standard Borel space, an uncountable family of Dirac laws admits no common
sigma-finite dominating measure, so density charts cannot be inferred from a general law fiber.
\status{ESTABLISHED}

\noindent\emph{Proof.}
A dominating measure must give positive mass to every atom of the family, and a sigma-finite measure
has at most countably many points of positive mass: on each finite-measure member of a countable cover,
the points of mass at least $1/n$ are finitely many. \hfill$\square$

\section{Kernels first, densities second}
\label{sec:prob-kernels}

\definitionheading{Normalized kernels; densities as derivatives}{def:prob-normalized-kernels}
The objects \eqref{eq:prob-generative-signature} and \eqref{eq:prob-recognition-signature} are
probability kernels. When $\mathbb P_{\theta,V}(\cdot\given X)\ll\nu^{\mathsf Y}\otimes\nu^{\mathsf O}$,
its density, the Radon--Nikodym derivative $p_{\theta,V}(o,y\given X)$, and the evidence, its marginal
density, are determined by
\begin{align}
\mathbb P_{\theta,V}(dy,do\given X)&=p_{\theta,V}(o,y\given X)\nu^{\mathsf Y}(dy)\nu^{\mathsf O}(do),
\label{eq:prob-generative-density}\\
p_{\theta,V}(o\given X)&=\int_{\mathsf Y}p_{\theta,V}(o,y\given X)\nu^{\mathsf Y}(dy).
\label{eq:prob-evidence}
\end{align}
When $\mathbb Q_{V,o,X}\ll\nu^{\mathsf Y}$ its density is written $q_{V,o,X}(y)$. Densities are determined
up to null sets of the reference, and versions that agree outside such a set denote the same density.
Measurability in $(o,X)$ of the evidence and of recognition expectations is
\Cref{prop:prob-kernel-integration-measurability}, and jointly measurable densities are
\Cref{cor:prob-common-reference-joint-density}; an unnormalized energy in a later chapter comes with a finite
positive normalizer.

\hypothesisheading{Regular-conditional convention}{hyp:prob-regular-conditional-convention}
For fixed $(\theta,X)$ the observation marginal and the posterior, a regular conditional probability, are
\begin{align}
\mathbb P_{\theta,V,X}^{\mathsf O}(A)&:=\mathbb P_{\theta,V}(\mathsf Y\times A\given X),
&A&\in\mathscr O,
\label{eq:prob-observation-marginal}\\
\boldsymbol\Pi_{\theta,V,o,X}(B)&=\mathbb P_{\theta,V}(B\given o,X),
&B&\in\mathscr Y,
\label{eq:prob-rcp}
\end{align}
the posterior being unique for $\mathbb P_{\theta,V,X}^{\mathsf O}$-almost every $o$. When the domination of
\eqref{eq:prob-generative-density} holds, \eqref{eq:prob-evidence} is a density of
$\mathbb P_{\theta,V,X}^{\mathsf O}$, and after a change on a $\nu^{\mathsf O}$-null set there is a
measurable set $\mathsf O_{\theta,V,X}^{\mathrm{reg}}$ with
\begin{align}
&\mathbb P_{\theta,V,X}^{\mathsf O}\bigl(\mathsf O_{\theta,V,X}^{\mathrm{reg}}\bigr)=1,
\qquad
0<p_{\theta,V}(o\given X)<\infty
\quad\text{on }\mathsf O_{\theta,V,X}^{\mathrm{reg}},
\label{eq:prob-regular-observations}\\
&\boldsymbol\Pi_{\theta,V,o,X}(B)
=\frac{\int_Bp_{\theta,V}(o,y\given X)\nu^{\mathsf Y}(dy)}{p_{\theta,V}(o\given X)},
\qquad o\in\mathsf O_{\theta,V,X}^{\mathrm{reg}}.
\label{eq:prob-rcp-density}
\end{align}
Every fixed-observation posterior, evidence-bound, or free-energy statement is asserted for
$\mathbb P_{\theta,V,X}^{\mathsf O}$-almost every $o$, represented by
$o\in\mathsf O_{\theta,V,X}^{\mathrm{reg}}$. The hypothesis excludes everywhere-pointwise statements
unless a jointly measurable density, an evidence representative, and a conditional version are declared
as model data; the joint measure alone does not determine a posterior at an exceptional observation
(\Cref{rem:prob-version-witness}).

\openproblemheading{Nondominated likelihood families}{open:prob-nondominated-likelihood}
The density formulas here use, for every record kernel $K_a$, one sigma-finite $\nu_a$ dominating
$\{K_a(X,y_{\partial a},\cdot)\}$ and a jointly measurable density. For a family without one, such as
$K_a(X,y_{\partial a},\cdot)=\delta_{h(y_{\partial a})}$ with $h$ taking uncountably many values
(\Cref{sec:exp-dominated-nondominated}), the joint law and its regular conditional posterior exist, but the
likelihood $dK_a(X,y_{\partial a},\cdot)/d\nu_a(o_a)$ is undefined and the evidence density can fail to
exist. Open is a treatment of the evidence bound, its local disintegrations, and the coarse channels at the
level of joint measures and regular conditional laws, or hypotheses restoring a jointly measurable
likelihood, with a rule for propagating record-null choices of versions. It is settled by such a treatment
with proofs that the exact identities of \Cref{ch:elbo,ch:coarsegraining} hold in it, or by a record family
for which they provably have no measure-level counterpart.

\section{Absolute continuity}
\label{sec:prob-abscont}

\definitionheading{Logarithmic density ratios and relative entropy}{def:prob-log-density-ratio-convention}
Wherever a logarithmic density ratio or its expectation is written, the numerator measure is absolutely
continuous with respect to the denominator measure. Relative entropy is
\begin{equation}
\KL(Q\Vert P)=
\begin{cases}
\displaystyle\int\log\Bigl(\frac{dQ}{dP}(y)\Bigr)Q(dy), & Q\ll P,\\[1.2ex]
+\infty, & \text{otherwise},
\end{cases}
\label{eq:prob-kl-definition}
\end{equation}
the first branch being the divergence of \citet{Kullback1951} and the value $+\infty$ off absolute
continuity the convention of \citet{Csiszar1967}. The negative part of the integrand is $Q$-integrable,
since $\int_{\{dQ/dP<1\}}-\log(dQ/dP)dQ\leq1/e$, so the first branch lies in $[0,+\infty]$.

\propositionheading{Density criterion for absolute continuity}{prop:prob-density-absolute-continuity}
Let $Q$ and $P$ be probability measures with $Q\ll\nu$ and $P\ll\nu$ for a sigma-finite $\nu$, with
densities $q$ and $p$, and put $Z=\{p=0\}$. Then
\begin{align}
&Q\ll P
\iff
Q(Z)=0
\iff
q=0\quad\nu\text{-almost everywhere on }Z,
\label{eq:prob-abscont-criterion}\\
&Q(S)=1,
\qquad
P(S)=0,
\label{eq:prob-direct-null-set}
\end{align}
where the first line holds, and whether or not $Q$ is dominated by $\nu$, a measurable $S$ satisfying the
second gives $Q\not\ll P$ and $\KL(Q\Vert P)=+\infty$. \status{ESTABLISHED}

\noindent\emph{Proof.}
First $P(Z)=\int_Zp(y)\nu(dy)=0$, so $Q\ll P$ gives $Q(Z)=0$. Conversely, if $Q(Z)=0$ and $P(A)=0$, then
$\nu(A\cap\{p>0\})=0$, so $Q(A\cap\{p>0\})=0$, while $Q(A\cap Z)\leq Q(Z)=0$. The second equivalence
follows from $Q(Z)=\int_Zq(y)\nu(dy)$. The last claim is the definition of absolute continuity together with
\eqref{eq:prob-kl-definition}. \hfill$\square$

A law carried by a Lebesgue-null set, such as a proper affine subspace, has no Lebesgue density, and
\eqref{eq:prob-direct-null-set} shows its singularity (\Cref{prop:elbo-subspace-support-singular}).

\section{The relative-entropy chain rule and data processing}
\label{sec:prob-kl-chain-rule}

\lemmaheading{Relative-entropy chain rule, data processing, and bijection invariance}{lem:prob-kl-chain-rule}
Let $P,Q$ be laws on $\mathsf Y_1\times\mathsf Y_2$ with first marginals $P_1,Q_1$ and conditional kernels
$P_{2|1},Q_{2|1}$. Then:
\begin{enumerate}
\item[(a)] the map $y_1\mapsto\KL(Q_{2|1}(y_1,\cdot)\Vert P_{2|1}(y_1,\cdot))$ is measurable and
\begin{equation}
\KL(Q\Vert P)
=\KL(Q_1\Vert P_1)
+\int\KL\bigl(Q_{2|1}(y_1,\cdot)\big\Vert P_{2|1}(y_1,\cdot)\bigr)Q_1(dy_1)
\label{eq:elbo-kl-chain-rule}
\end{equation}
in $[0,+\infty]$, with a right side that does not depend on the chosen conditional versions;
\item[(b)] on a finite product $\mathsf Y=\prod_{b=1}^B\mathsf Y_b$, for any law $Q$ with marginals $Q_b$
and any laws $\rho_b$ on $\mathsf Y_b$,
\begin{equation}
\KL\Bigl(Q\Big\Vert\bigotimes_b\rho_b\Bigr)
=\KL\Bigl(Q\Big\Vert\bigotimes_bQ_b\Bigr)+\sum_b\KL(Q_b\Vert\rho_b)
\label{eq:prob-kl-product-reference}
\end{equation}
in $[0,+\infty]$;
\item[(c)] for a law $\mu$ on $\mathsf Y_1$, the minimum of $\KL(Q\Vert P)$ over all laws $Q$ with
$Q_1=\mu$ is $\KL(\mu\Vert P_1)$, attained by $\mu(dy_1)P_{2|1}(y_1,dy_2)$, uniquely when
$\KL(\mu\Vert P_1)<\infty$;
\item[(d)]\label{thm:cg-kl-dpi-extended} for laws $P,Q$ on $\mathsf Y$ and every Markov kernel
$K:\mathsf Y\rightsquigarrow\mathsf Z$,
\begin{equation}
\KL(QK\Vert PK)\leq\KL(Q\Vert P)
\qquad\text{in }[0,+\infty];
\label{eq:cg-kl-dpi-extended}
\end{equation}
\item[(e)] for a bimeasurable bijection $A:\mathsf Y\to\mathsf Y'$,
\begin{align}
\frac{d(A_\#Q)}{d(A_\#P)}&=\frac{dQ}{dP}\circ A^{-1}
\quad A_\#P\text{-almost surely, if }Q\ll P,
\label{eq:cg-rn-bimeasurable}\\
\KL(A_\#Q\Vert A_\#P)&=\KL(Q\Vert P)
\quad\text{in }[0,+\infty].
\label{eq:cg-kl-bimeasurable-invariance}
\end{align}
\end{enumerate}
\status{ESTABLISHED}

The proof is in \Cref{app:measure-lemmas}. Clauses (a), (b), and (d) need no absolute-continuity or
finiteness hypothesis.

\section{Recognition marginals, couplings, and selection}
\label{sec:prob-marginals}

\definitionheading{Local recognition laws and their coupling class}{def:prob-recognition-marginals}
With $\operatorname{pr}_i:\mathsf Y\to\mathsf Y_i$ the agent-block projection and
$\operatorname{pr}_i^k,\operatorname{pr}_i^m$ the coordinate projections of $\mathsf Y_i$, for each admitted
$(o,X)$ put
\begin{align}
\operatorname{Cpl}(\mathbb Q_{\bullet,o,X})
&:=\Bigl\{\mathbb J\in\mathcal P(\mathsf Y):(\operatorname{pr}_i)_\#\mathbb J=\mathbb Q_{i,o,X}\text{ for every }i\in V\Bigr\},
\label{eq:prob-recognition-coupling-class}\\
q_i^{b;o,X}&:=(\operatorname{pr}_i^k)_\#\mathbb Q_{i,o,X},
\qquad
q_i^{m;o,X}:=(\operatorname{pr}_i^m)_\#\mathbb Q_{i,o,X}.
\label{eq:prob-recognition-marginals}
\end{align}
The belief and model recognition marginals $q_i^{b;o,X},q_i^{m;o,X}$ are the values at $c_*$, in the frame
$u_i$, of agent $i$'s sections $q_i^b,q_i^m$ at the admitted $(o,X)$. They are local data; a member of
\eqref{eq:prob-recognition-coupling-class} is population-level dependence data.

\propositionheading{Existence and convexity of the coupling class}{prop:prob-coupling-existence-convexity}
For every admitted $(o,X)$ the class $\operatorname{Cpl}(\mathbb Q_{\bullet,o,X})$ is nonempty and convex,
and
\begin{equation}
(o,X)\longmapsto\bigotimes_{i\in V}\mathbb Q_{i,o,X}
\label{eq:prob-product-coupling-kernel}
\end{equation}
is a measurable product-coupling kernel. \status{ESTABLISHED}

\noindent\emph{Proof.}
On measurable rectangles the product is a product of measurable functions, and a monotone-class argument
extends measurability; its $i$-th marginal is $\mathbb Q_{i,o,X}$ because the other factors have unit
mass. Marginalization is affine, so convex combinations of couplings are couplings. \hfill$\square$

\definitionheading{Selected population recognition kernel}{hyp:prob-measurable-coupling-selector}
\label{def:elbo-recognition-kernel}
A selection rule $\mathsf S$ is declared, with law spaces carrying their evaluation sigma-algebras, such
that the population recognition kernel
\begin{align}
\mathbb Q_V&:(\mathsf O\times\mathsf X,\mathscr O\otimes\mathscr X)\longrightarrow\mathcal P(\mathsf Y,\mathscr Y),
&(o,X)&\longmapsto\mathbb Q_{V,o,X},
\label{eq:elbo-recognition-kernel}\\
\mathbb Q_{V,o,X}&=\mathsf S_{o,X}\bigl((\mathbb Q_{i,o,X})_{i\in V}\bigr)\in\operatorname{Cpl}(\mathbb Q_{\bullet,o,X}),
\label{eq:prob-selected-recognition-coupling}
\end{align}
has $(o,X)\mapsto\mathbb Q_{V,o,X}(B)$ measurable for every $B\in\mathscr Y$. The law may correlate agents
and channels; the product coupling is one admissible selection, not the default, and copulas, common latent
variables, ordered conditionals, or graphical families supply others.

\hypothesisheading{Constrained coupling selection}{hyp:prob-constrained-coupling-selection}
When the selection must lie in an imposed subfamily
$\operatorname{Cpl}^{\mathrm{res}}(\mathbb Q_{\bullet,o,X})\subseteq\operatorname{Cpl}(\mathbb Q_{\bullet,o,X})$,
every fiber $\operatorname{Cpl}^{\mathrm{res}}(\mathbb Q_{\bullet,o,X})$ is assumed nonempty,
\label{hyp:prob-constrained-coupling-feasibility}
and a selector $\mathsf S$ from these fibers is assumed with
$(o,X)\mapsto\mathsf S_{o,X}((\mathbb Q_{i,o,X})_{i\in V})(B)$ measurable for every $B\in\mathscr Y$. The
hypothesis excludes restricted families whose fibers are empty or admit no measurable selection; the two
conditions are independent, and the product coupling proves neither.

\propositionheading{Coordinate marginals do not determine a sufficiently rich joint}{prop:prob-marginals-do-not-determine-joint}
For a finite family of standard Borel blocks, the class $\operatorname{Cpl}((\mu_i)_{i\in V})$ is a
singleton if and only if at most one marginal $\mu_i$ is non-Dirac. Hence, when two local recognition
laws are non-Dirac, the local marginals do not determine $\mathbb Q_{V,o,X}$. \status{ESTABLISHED}

\noindent\emph{Proof.}
Let $\mu_1,\mu_2$ be non-Dirac. A non-Dirac law charges some measurable set with probability in $(0,1)$,
since otherwise the members or complements of probability one of a countable separating generator would
meet in a set of probability one with at most one point. Choose such sets $A,B$ and, for
$0<|\varepsilon|<1$, put
\begin{equation}
d\mathbb J_\varepsilon(x,x')
=\bigl[1+\varepsilon(\mathbf 1_A(x)-\mu_1(A))(\mathbf 1_B(x')-\mu_2(B))\bigr]d\mu_1(x)d\mu_2(x').
\label{eq:prob-coupling-perturbation}
\end{equation}
The bracket is at least $1-|\varepsilon|>0$ and integrating in either variable removes the perturbation,
so $\mathbb J_\varepsilon$ couples $\mu_1,\mu_2$, while
$\mathbb J_\varepsilon(A\times B)=\mu_1(A)\mu_2(B)+\varepsilon\mu_1(A)(1-\mu_1(A))\mu_2(B)(1-\mu_2(B))$
varies with $\varepsilon$; tensoring with the other marginals gives distinct couplings of the family.
Conversely, if at most one marginal is non-Dirac, every coupling is carried by the Dirac atoms in the
other coordinates, and its law there is the remaining marginal. \hfill$\square$

\propositionheading{Limits of access restriction}{prop:prob-access-restriction}
A restriction of the access maps constrains the local laws and not the population law, and it has an
evidence price that no selector removes.
\begin{enumerate}
\item[(i)] Let $V=\{1,2\}$ with binary blocks, let both local laws be $\operatorname{Ber}(\tfrac12)$ at every
observation, so that every access map is respected, and fix $0<r<1$. The selector that returns the
coupling with $\mathbb Q_{V,o,X}(y_1=y_2)=(1+r)/2$ when a record $o_2$ of agent $2$ equals $1$ and
$(1-r)/2$ otherwise satisfies \eqref{eq:prob-selected-recognition-coupling}, and the population law
depends on $o_2$ while no local law does. A decentralized information claim therefore also requires a
selector indexed by $((\operatorname{Acc}_i(o,X))_{i\in V},X)$ alone.
\item[(ii)] Let $V=\{1,2\}$ with binary latents, the prior with mass $0.45$ on each of $(0,0)$ and $(1,1)$
and $0.05$ on each of $(0,1)$ and $(1,0)$, and singleton records $o_1,o_2$ that report $y_1,y_2$
correctly with probability $0.8$, independently given $y$. Let $\operatorname{Acc}_1(o,X)=o_1$, so that
$\mathbb Q_{1,o,X}=\operatorname{Ber}(t)$ is the same at $o=(1,0)$ and $o=(1,1)$. The posterior marginals
of $y_1$ at these observations, written as the probabilities of $y_1=0$ and $y_1=1$, are
$\Pi_1^{(0)}=(37/89,52/89)$ and $\Pi_1^{(1)}=(13/161,148/161)$, and
every recognition law with this access restriction satisfies
\begin{equation}
\begin{aligned}
\max_{o_2\in\{0,1\}}\KL\bigl(\mathbb Q_{V,(1,o_2),X}\big\Vert\boldsymbol\Pi_{\theta,V,(1,o_2),X}\bigr)
&\geq\max_{o_2\in\{0,1\}}\KL\bigl(\operatorname{Ber}(t)\big\Vert\Pi_1^{(o_2)}\bigr)\\
&\geq\sup_{\xi\in[0,1]}\bigl(-\log Z_\xi\bigr)=0.08849\ldots,
\end{aligned}
\label{eq:prob-access-gap-floor}
\end{equation}
where $Z_\xi=\sum_{y_1}\Pi_1^{(0)}(y_1)^\xi\Pi_1^{(1)}(y_1)^{1-\xi}$. The last bound is sharp: the supremum
is attained at $\xi\approx0.548$, and at $t\approx0.783$ the middle term equals it. Hence the equality
condition of \Cref{cor:elbo-bound-tightness} lies outside the access-restricted class, and the exact
E-coordinate of \Cref{prop:elbo-exact-e-coordinate}, which optimizes over every law with a prescribed
complement marginal, need not be access restricted.
\end{enumerate}
\status{ESTABLISHED}

\noindent\emph{Proof.}
(i) Both couplings have $\operatorname{Ber}(\tfrac12)$ marginals. (ii) Bayes' rule gives the posterior
marginals, and projection onto $y_1$ with \Cref{lem:prob-kl-chain-rule}(d) gives the first inequality. For
$\xi\in[0,1]$ let $\Pi^{[\xi]}=Z_\xi^{-1}(\Pi_1^{(0)})^\xi(\Pi_1^{(1)})^{1-\xi}$. Expanding logarithms gives
$\xi\KL(\operatorname{Ber}(t)\Vert\Pi_1^{(0)})+(1-\xi)\KL(\operatorname{Ber}(t)\Vert\Pi_1^{(1)})
=\KL(\operatorname{Ber}(t)\Vert\Pi^{[\xi]})-\log Z_\xi\geq-\log Z_\xi$, and a maximum dominates every convex
combination. The concave function $\xi\mapsto-\log Z_\xi$ has derivative
$\E_{\Pi^{[\xi]}}\log(\Pi_1^{(1)}/\Pi_1^{(0)})$, which vanishes at its interior maximizer $\xi_\star$, where
the two divergences from $\Pi^{[\xi_\star]}$ are therefore equal; the identity with
$\operatorname{Ber}(t)=\Pi^{[\xi_\star]}$ shows that both equal $-\log Z_{\xi_\star}$, the Chernoff
information \citep[Section~11.9]{CoverThomas2006}. The value and maximizers follow from the closed forms,
and every access-restricted law has a positive gap at one of the two observations. \hfill$\square$

\section{Levels of description}
\label{sec:prob-levels}

\definitionheading{Sample, law-point, and configuration levels}{def:prob-description-levels}
At the sample level the agent block is $\mathsf Y_i$ of \Cref{def:prob-product-sample-spaces}, the
population recognition law $\mathbb Q_{V,o,X}$ is a law of samples, and the belief of agent $i$ is the
recognition marginal $q_i^{b;o,X}$ of \eqref{eq:prob-recognition-marginals}. Generative kernels and the
record kernels \eqref{eq:prob-record-kernel-signature} read sample coordinates. At the law-point level, law
points are latent coordinates, and the law-point blocks are
\begin{equation}
\mathsf Y^{\mathrm{law}}_i=(\mathcal E_b)_{c_*}\times(\mathcal E_m)_{c_*}\ni\bigl(\mathfrak q_i^b,\mathfrak q_i^m\bigr),
\qquad
\mathsf Y^{\mathrm{law}}=\prod_{i\in V}\mathsf Y^{\mathrm{law}}_i,
\label{eq:prob-law-point-space}
\end{equation}
possibly augmented by further latent coordinates; through the frame $u_i$, $\mathfrak q_i^b$ and
$\mathfrak q_i^m$ are laws on $\mathsf K_i$ and $\mathsf M_i$. A generative law at this level is a law on
$\mathsf Y^{\mathrm{law}}\times\mathsf O$ whose kernels may read law points, and a recognition law is a law
over law points. At the configuration level the latent state is
$(w,y)\in\mathsf Y^{\mathrm{cfg}}\times\mathsf Y$, where the configuration space $\mathsf Y^{\mathrm{cfg}}$
carries a measurable presentation map
$\operatorname{pres}:\mathsf Y^{\mathrm{cfg}}\to\mathsf Y^{\mathrm{law}}$ and may carry further latent
coordinates such as random links. A lifted generative law is a law $\mathbb P^{\mathrm{cfg}}_{\theta,V}(dw,dy,do\given X)$, and a lifted
recognition law has the form $\mathbb J(dw,dy)=\mathbb J^{\mathrm{cfg}}(dw)\kappa^{\mathrm{pres}}(w,dy)$ with a
presentation kernel $\kappa^{\mathrm{pres}}:\mathsf Y^{\mathrm{cfg}}\rightsquigarrow\mathsf Y$ whose marginal
at each agent is the law presented by the corresponding component of $\operatorname{pres}(w)$.

The levels define different optimization problems, so an energy enters a free energy only after its level
is named. The source energies of \eqref{eq:obs-attention-likelihood} and every record kernel are
sample-level. The transported divergences \eqref{eq:full-graph-transported-divergences} and the role
divergences \eqref{eq:obs-selective-role-divergence} are measurable $[0,+\infty]$-valued functions of law
points (\Cref{cor:prob-kernel-kl-measurability}), admissible as generative energies when their arguments
are latent law points; read at the sample level they would make a generative factor read recognition
marginals, which \Cref{req:gen-typing-prohibition} excludes. A block admitted by \Cref{hyp:local-interaction-kernels} with
law points is a law-point or configuration block, and the formation model
\eqref{eq:rg-formation-configuration} is configuration-level, with records that read the configuration.

\propositionheading{A law-point energy is not a sample energy}{prop:prob-law-point-energy}
Let $q_i$ be a receiver law and $p_1,p_2$ source laws in $\mathcal B_x$ with
$\KL(q_i\Vert\widehat\rho_x(\Theta_e^x)p_k)<\infty$ for $k=1,2$, whose transported laws have densities
against a common dominating measure that differ on a set of positive $q_i$-measure. Then for $0<s<1$
\[
\KL\bigl(q_i\Vert\widehat\rho_x(\Theta_e^x)(sp_1+(1-s)p_2)\bigr)
<s\KL\bigl(q_i\Vert\widehat\rho_x(\Theta_e^x)p_1\bigr)+(1-s)\KL\bigl(q_i\Vert\widehat\rho_x(\Theta_e^x)p_2\bigr),
\]
and no recognition-independent sample energy $D$ satisfies
$\E_{q_i\otimes q}D=\KL(q_i\Vert\widehat\rho_x(\Theta_e^x)q)$ at $q=p_1$, at $q=p_2$, and at one strict
mixture of them. \status{ESTABLISHED}

\noindent\emph{Proof.}
Pushforward preserves mixtures. Let $\alpha,\beta$ be the $q_i$-densities of the absolutely continuous parts
of the two transported laws. Finite divergences give $q_i\ll$ each transported law, so $\alpha,\beta>0$
$q_i$-almost everywhere and the divergences are $\int-\log\alpha(y)q_i(dy)$ and $\int-\log\beta(y)q_i(dy)$,
well defined because $\int(\log\alpha(y))^+q_i(dy)\leq\int\alpha(y)q_i(dy)\leq1$. The mixture has relative
density $s\alpha+(1-s)\beta$, and $-\log(s\alpha+(1-s)\beta)\leq-s\log\alpha-(1-s)\log\beta$, strictly on the
set $\{\alpha\neq\beta\}$ of positive $q_i$-measure; integrating gives the strict inequality. An energy $D$
contributes $\E_{q_i\otimes q}D$, affine in $q$ along mixtures, so matching the divergence at $p_1$ and $p_2$
forces the affine value at the mixture, which the strict inequality excludes. \hfill$\square$

\propositionheading{Point-mass recognition of law points has infinite cost}{prop:prob-law-point-point-mass}
At the law-point level let $\mathbb P^{\mathrm{law}}_{\theta,V}(\cdot\given X)$ be a generative law with
latent marginal $\mathbb P^{\mathrm{law},\mathsf Y}_{\theta,V}$ and posterior
$\boldsymbol\Pi^{\mathrm{law}}_{\theta,V,o,X}\ll\mathbb P^{\mathrm{law},\mathsf Y}_{\theta,V}$ at an observation
with finite positive evidence; the domination holds when every record kernel has a density against a fixed
reference. If $\bar{\mathfrak q}$ is not an atom of the prior, then
\begin{equation}
\KL\bigl(\delta_{\bar{\mathfrak q}}\big\Vert\boldsymbol\Pi^{\mathrm{law}}_{\theta,V,o,X}\bigr)=+\infty,
\label{eq:prob-law-point-point-mass}
\end{equation}
and the extended evidence bound of $\delta_{\bar{\mathfrak q}}$ is $-\infty$. Consequently, against a
non-atomic law-point prior, a row or an update computed at fixed law points is not the recognition
optimum of a finite evidence bound at the law-point level: with the law points latent it is a generative
posterior quantity conditional on them, and with the law points declared as parameters in $(\theta,X)$ it
is a function of parameters whose change is an M-coordinate subject to
\Cref{def:elbo-generalized-m-acceptance}. \status{ESTABLISHED}

\noindent\emph{Proof.}
Singletons are measurable, and domination gives
$\boldsymbol\Pi^{\mathrm{law}}_{\theta,V,o,X}(\{\bar{\mathfrak q}\})=0$, so \eqref{eq:prob-direct-null-set} holds for
$\delta_{\bar{\mathfrak q}}$ with $S=\{\bar{\mathfrak q}\}$. The extended bound
\eqref{eq:elbo-extended-definition} subtracts this divergence from a finite log evidence, and a recognition
optimum with finite bound has finite divergence. \hfill$\square$

\propositionheading{Configuration-lift identity}{prop:prob-configuration-lift-identity}
Let $\kappa^{\mathrm{cfg}}:\mathsf Y\rightsquigarrow\mathsf Y^{\mathrm{cfg}}$ be a Markov kernel, and let the lifted
generative law be $\mathbb P_{\theta,V}(dy,do\given X)\kappa^{\mathrm{cfg}}(y,dw)$, so that its
$(y,o)$ marginal is the sample-level law and its records read only sample coordinates. Then at every
$o\in\mathsf O^{\mathrm{reg}}_{\theta,V,X}$ the evidence is unchanged, the lifted posterior may be selected
as $\boldsymbol\Pi_{\theta,V,o,X}(dy)\kappa^{\mathrm{cfg}}(y,dw)$, and every lifted recognition law $\mathbb J$
with sample marginal $\bar{\mathbb Q}$ and conditional kernel $\mathbb J(dw\given y)$ satisfies
\begin{equation}
\Lelbo^{\mathrm{ext,cfg}}_V(\mathbb J;X,o)
=\Lelbo^{\mathrm{ext}}_V(\bar{\mathbb Q};X,o)
-\E_{\bar{\mathbb Q}}\Bigl[\KL\bigl(\mathbb J(\cdot\given Y)\big\Vert\kappa^{\mathrm{cfg}}(Y,\cdot)\bigr)\Bigr]
\label{eq:prob-configuration-lift-identity}
\end{equation}
in $[-\infty,\log p_{\theta,V}(o\given X)]$, where $\Lelbo^{\mathrm{ext,cfg}}_V$ is
\eqref{eq:elbo-extended-definition} for the lifted law. The lift never raises the bound of the induced
sample law, and when that bound is finite it leaves it unchanged exactly when
$\mathbb J(\cdot\given y)=\kappa^{\mathrm{cfg}}(y,\cdot)$ for $\bar{\mathbb Q}$-almost every $y$.
\status{ESTABLISHED}

\noindent\emph{Proof.}
The $(y,o)$ marginal of the lifted law is $\mathbb P_{\theta,V}$, so the evidence is unchanged, and
disintegrating along $o$ gives the stated posterior. \Cref{lem:prob-kl-chain-rule}(a), with $y$ as the
first coordinate, gives the lifted divergence as $\KL(\bar{\mathbb Q}\Vert\boldsymbol\Pi_{\theta,V,o,X})$ plus
the conditional term in $[0,+\infty]$, and the conditional term vanishes exactly under the stated
almost-everywhere equality. Substitution in
\eqref{eq:elbo-extended-definition} subtracts two terms in $[0,+\infty]$ from a finite log evidence.
\hfill$\square$

\openproblemheading{A normalized law-point prior}{open:prob-law-point-prior}
No chapter constructs a normalized generative law on $\mathsf Y^{\mathrm{law}}$ for a general class of law
points, which the law-point energies \eqref{eq:full-graph-transported-divergences} and
\eqref{eq:obs-selective-role-divergence} need to be evidence-bound terms rather than a separately declared
potential. The law fibers supply no canonical sigma-finite reference ($\mathcal P(\mathsf K)$ is in general
infinite dimensional), and the Haar measure of a noncompact $G$ is infinite, so a Haar prior on a free orbit
is not normalizable; the formation model supplies one normalized configuration-level prior,
\eqref{eq:rg-formation-physical-prior}, for finite-dimensional Gaussian presentations with a proper,
non-invariant reference. It is settled for a class of law points by a sigma-finite reference on the law-point
or presentation space, a finite normalizer for the exponentiated energies, the transformation law of the prior
under the active gauge action, record kernels that read law points, and a compatibility relation between the
law-point posterior and the sample-level recognition marginals.

\section*{Physicist's summary}

Everything probabilistic at $c_*$ is a law or a kernel, with densities against references fixed per family,
and one chain rule in $[0,+\infty]$ does all later bookkeeping. Local recognition laws do not determine their
coupling, restricted access has an exact evidence price, and attention's transported divergence is a
law-point energy.

\chapter[The Fixed Normalized Generative Law]{The Fixed Normalized\\Generative Law}
\label{ch:generative}

This chapter answers how one normalized joint law of latent samples and records is built at $c_*$ and
what passive re-expression of the frames does to it. The joint is normalized for every finite record
hypergraph, cyclic ones included (\Cref{prop:gen-exact-normalization}), local normalization of
reciprocal latent conditionals supplies no joint (\Cref{prop:gen-gibbs-counterexample}), and under
kernel covariance the evidence and the bound of a transported coupling are invariant
(\Cref{prop:gen-product-evidence-invariance,cor:elbo-independent-frame-invariance}). Its evidence is
bounded in \Cref{ch:elbo}, and attention and coarse channels act on it in
\Cref{ch:selective-communication,ch:coarsegraining}.

\paragraph{Scope.}
The cross morphisms supply no generative factor, and a model with factors in both directions needs its
own normalized joint. No quotient of the parameter space by passive re-expression is formed, and
no parameter count subtracts the dimension of the vertex group, whose action can have nontrivial
stabilizers. \status{NOT-CLAIMED}

\section{Parameters, the directed latent law, and the typing prohibition}
\label{sec:gen-typing-setup}
\label{sec:gen-composition}
\label{sec:gen-typing}

\definitionheading{Parameters, structural data, and generative factors}{def:gen-parameter-split}
Fix the agent set $V=\{1,\ldots,N\}$ and the spaces and references of \Cref{ch:probability}. Structural
data $X\in\mathsf X$ contain a finite directed acyclic graph $\Gamma$ on $V$ with parent sets
$\operatorname{pa}(i)$, the interaction multigraph $\mathfrak I$ with its links, and the frames; the
numerical parameters that a learning coordinate may move are denoted $\theta$. The split between $X$
and $\theta$ is declared, not derived, and every typing statement below applies to both. For each agent
$i$ declare the Markov kernels of the first line below, its model and state factors, which are its model
and belief priors and are generative factors, not components of the agent of \Cref{def:geo-agent}; at a
root, where $\operatorname{pa}(i)=\varnothing$, the first is a proper model law and the second a bridge
from the agent's own model coordinate to its belief coordinate. The state factor is the evaluated
kernel of \Cref{def:prob-model-evaluation} with source
$\mathsf S_i^X=\prod_{j\in\operatorname{pa}(i)}\mathsf K_j$, the empty product being a singleton
$\{\ast\}$, and target $\mathsf T_i^X=\mathsf K_i$ (second line):
\begin{gather}
P_{\theta,i}^{m}(dm_i\given m_{\operatorname{pa}(i)},X),
\qquad
P_{\theta,i}^{k}(dk_i\given k_{\operatorname{pa}(i)},m_i,X),
\label{eq:gen-kernels}\\
\operatorname{ev}_i(m_i)=K^X_{i,m_i},
\qquad
K^X_{i,m_i}(k_{\operatorname{pa}(i)},B):=P_{\theta,i}^{k}(B\given k_{\operatorname{pa}(i)},m_i,X).
\label{eq:gen-evaluated-state-kernel}
\end{gather}
With $\mathsf Y_{\operatorname{pa}(i)}=\prod_{j\in\operatorname{pa}(i)}\mathsf Y_j$, the agent mechanism is the
ordered composite below, which reads the whole parent block $y_{\operatorname{pa}(i)}$ and has unit mass:
\begin{gather}
G_{\theta,i}:(\mathsf X\times\mathsf Y_{\operatorname{pa}(i)})\rightsquigarrow\mathsf Y_i,
\label{eq:gen-agent-design-mechanism}\\
G_{\theta,i}\bigl(d(k_i,m_i)\given y_{\operatorname{pa}(i)},X\bigr)
:=P_{\theta,i}^{m}(dm_i\given m_{\operatorname{pa}(i)},X)K^X_{i,m_i}(k_{\operatorname{pa}(i)},dk_i).
\label{eq:gen-agent-mechanism}
\end{gather}

\Cref{def:obs-fast-slow-blocks} divides these coordinates into fast recognition and slow model
coordinates. At a root the state factor induces, for a model law $s$ on
$\mathsf M_i$, the belief law
\begin{equation}
\widetilde\Phi^{\mathrm{gen}}_i(s)(B)=\int_{\mathsf M_i}P_{\theta,i}^{k}(B\given m,X)s(dm),
\qquad B\in\mathscr K_i,
\label{eq:gen-root-law-bridge}
\end{equation}
which can serve as the representative at $c_*$ of the cross morphism $\widetilde\Phi$ of
\Cref{def:geo-cross-morphisms} when that identification is declared; at a nonroot the evaluated kernel
is indexed by the parent belief tuple and is a kernel-valued bridge, not a map of law fibers.

\definitionheading{Finite directed latent law}{cons:gen-finite-directed-law}
Choose a topological ordering $v_1,\ldots,v_N$ of $\Gamma$. The latent law is the ordered composition of
the first line below, where the arrow denotes ordered application of kernels as in the second line, for
bounded measurable $f$; every conditioning coordinate is introduced before its kernel is applied:
\begin{gather}
\mathbb P_{\theta,V}^{\mathsf Y}(dy\given X)
:=\overrightarrow{\prod_{l=1}^{N}}
G_{\theta,v_l}\bigl(dy_{v_l}\given y_{\operatorname{pa}(v_l)},X\bigr),
\label{eq:gen-composition}\\
\int f(y)\mathbb P_{\theta,V}^{\mathsf Y}(dy\given X)
=\int G_{\theta,v_1}(dy_{v_1}\given X)\cdots
\int G_{\theta,v_N}\bigl(dy_{v_N}\given y_{\operatorname{pa}(v_N)},X\bigr)f(y).
\label{eq:gen-iterated-integral}
\end{gather}
When the factors have densities against the component references, the densities are
\begin{gather}
g_{\theta,i}(k_i,m_i\given y_{\operatorname{pa}(i)},X)
=p_{\theta,i}^{m}(m_i\given m_{\operatorname{pa}(i)},X)p_{\theta,i}^{k}(k_i\given k_{\operatorname{pa}(i)},m_i,X),
\label{eq:gen-density}\\
p_{\theta,V}^{\mathsf Y}(y\given X)=\prod_{i\in V}g_{\theta,i}(k_i,m_i\given y_{\operatorname{pa}(i)},X),
\qquad
\mathbb P_{\theta,V}^{\mathsf Y}(dy\given X)=p_{\theta,V}^{\mathsf Y}(y\given X)\nu^{\mathsf Y}(dy).
\label{eq:gen-design-density}
\end{gather}

\hypothesisheading{Typing prohibition}{req:gen-typing-prohibition}
Each factor in \eqref{eq:gen-kernels} and each record kernel is a function only of $\theta$, $X$, and its
displayed latent coordinates; no generative kernel takes as an argument the recognition law
$\mathbb Q_{V,o,X}$, one of its marginals or parameters, or the posterior $\boldsymbol\Pi_{\theta,V,o,X}$. In
the two-argument form $\Fenergy_V(\mathbb Q;\theta)$ of the free energy \citep{Neal1998}, every evaluation
holds the generative argument $(\theta,X)$ fixed while the recognition argument varies, which excludes the
composite $\mathbb Q\mapsto\Fenergy_V(\mathbb Q;\theta(\mathbb Q))$ inside one evaluation: a generative
factor that reads the recognition law being scored, such as a transition built from a live peer
recognition marginal, is a map from recognition laws to kernels rather than a factor of one joint. A
parameter may depend on the observation or on an earlier recognition law if it is held fixed while the
bound is evaluated, so slow evolution of an agent's model or prior is admissible as recognition dynamics
of the model channel under one fixed joint or as accepted M-coordinates on $\theta$
(\Cref{prop:elbo-exact-m-coordinate,def:elbo-generalized-m-acceptance,prop:elbo-evidence-monotonicity}),
each accepted value defining a new fixed joint to which the evidence identity applies afresh.

Coupling through latents and loops of inference updates under a fixed joint remain admissible
(\Cref{sec:full-graph-closing-tower}), and a family of joints indexed by recognition laws designates no
target (\Cref{prop:gen-no-distinguished-target}).

\section{Normalized record extension and the undirected alternative}
\label{sec:gen-normalization}
\label{sec:gen-undirected}

\propositionheading{Normalization by latent composition and record attachment}{prop:gen-exact-normalization}
The iterated integral \eqref{eq:gen-iterated-integral} defines a probability kernel
$X\mapsto\mathbb P_{\theta,V}^{\mathsf Y}(\cdot\given X)$. If the record coordinates are conditionally
independent given $(X,y)$, the population law is the kernel of the first line below, and the second
line holds for every finite record hypergraph, cyclic or not, and every $(\theta,X)$:
\begin{gather}
\mathbb P_{\theta,V}(dy,do\given X)
=\mathbb P_{\theta,V}^{\mathsf Y}(dy\given X)\prod_{a\in\mathfrak A}K_a(X,y_{\partial a},do_a),
\label{eq:gen-record-extension}\\
\mathbb P_{\theta,V}(\mathsf Y\times\mathsf O\given X)=1.
\label{eq:gen-normalization}
\end{gather}
If $K_a(X,y_{\partial a},do_a)=\exp\bigl(-E_{a,o}(y_{\partial a})\bigr)\nu_a(do_a)$ with a jointly measurable
negative log density $E_{a,o}$ and $\nu^{\mathsf O}=\bigotimes_{a}\nu_a$, then
\begin{equation}
p_{\theta,V}(o,y\given X)
=p_{\theta,V}^{\mathsf Y}(y\given X)\exp\Bigl(-\sum_{a\in\mathfrak A}E_{a,o}(y_{\partial a})\Bigr)
\label{eq:gen-record-density}
\end{equation}
is its density. \status{ESTABLISHED}

\noindent\emph{Proof.}
Integrating a bounded measurable function against a probability kernel preserves boundedness and
measurability in the remaining variables, so induction along the topological order defines a measurable
kernel in $X$, and $f\equiv1$ gives unit mass; finite kernel composition gives measurability of
\eqref{eq:gen-record-extension}. By Tonelli's theorem, integrate the record coordinates first, in any
order: each $K_a$ has unit mass in $o_a$, so the finite product disappears and leaves
$\mathbb P_{\theta,V}^{\mathsf Y}(dy\given X)$ whatever the overlaps of the scopes, since the records are
conditionally independent children of one latent configuration. The same computation with the
Radon--Nikodym derivatives proves \eqref{eq:gen-record-density}. \hfill$\square$

Each record kernel is owned once by the joint, while the coordinate-local free energies of \Cref{ch:elbo}
include every factor incident to a coordinate.

\paragraph{Running example.}
The agents of \Cref{sec:geo-running-example} carry the directed joint with
$\operatorname{pa}(1)=\varnothing$, $\operatorname{pa}(2)=\operatorname{pa}(3)=\{1\}$, independent
$\mathcal N(0,1)$ model coordinates, $k_1\given m_1\sim\mathcal N(0,e^{-m_1}I_3)$, and children
transported from agent $1$,
$k_2\given k_1,m_2\sim\mathcal N\bigl(\rho_b(\Theta^b_{\bar e_1})k_1,e^{-m_2}I_3\bigr)$ and
$k_3\given k_1,m_3\sim\mathcal N\bigl(\rho_b(\Theta^b_{e_3})k_1,e^{-m_3}I_3\bigr)$. Two binary records are
attached: $a_1$ with scope $\{1,2\}$ as in \eqref{eq:geo-running-record}, and $a_2$ with scope $\{2,3\}$
and success probability $\exp(-\tfrac12\lvert k_2-\rho_b(\Theta^b_{e_2})k_3\rvert^2)$. The records close the
latent tree into the triangle, and \Cref{prop:gen-exact-normalization} normalizes the joint.
At the conditional means of the children the residual read by $a_2$ is
$\rho_b(\Theta^b_{\bar e_1})k_1-\rho_b(\Theta^b_{e_2}\Theta^b_{e_3})k_1
=\rho_b(\Theta^b_{e_1})^{-1}(I_3-\rho_b(\mathrm{Hol}^b(p)))k_1$, whose length is $\sqrt3$ times that of
the component of $k_1$ orthogonal to the rotation axis of the holonomy \eqref{eq:geo-running-holonomy},
so the records see the loop holonomy through an invariant.

\definitionheading{Normalized undirected model}{def:gen-undirected-model}
For nonnegative potentials on the agents, on the edges of an undirected graph with edge set
$E^{\mathrm u}$, and on the records, put $\widetilde p_X$ and $Z_X$ as in the first two lines below.
When $0<Z_X<\infty$, a condition that can depend on the parameters, the undirected model is the law with
the density of the third line:
\begin{gather}
\widetilde p_X(o,y)
=\prod_{i\in V}\psi_i(k_i,m_i;X)
\prod_{\{i,j\}\in E^{\mathrm u}}\psi_{ij}(y_i,y_j;X)
\prod_{a\in\mathfrak A}\psi_a(o_a,y_{\partial a};X),
\label{eq:gen-gibbs-potential}\\
Z_X
=\int\widetilde p_X(o,y)\nu^{\mathsf O}(do)\nu^{\mathsf Y}(dy),
\label{eq:gen-gibbs-partition}\\
p_{\theta,V}^{\mathrm{Gibbs}}(o,y\given X)=Z_X^{-1}\widetilde p_X(o,y).
\label{eq:gen-gibbs-density}
\end{gather}

\propositionheading{Local normalization is insufficient}{prop:gen-gibbs-counterexample}
Normalized agent potentials and an everywhere positive, finite, continuous edge potential need not give a
finite $Z_X$. \status{ESTABLISHED}

\noindent\emph{Proof.}
Take two agents with scalar belief coordinates $k_1,k_2\in\R$ under Lebesgue measure, singleton model and
record spaces with unit point mass, and potentials with product exponent
$-\tfrac12k^\top\mathsf W_{\mathrm c}k$,
\begin{gather}
\psi_i(k_i,m_i;X)=(2\pi)^{-1/2}e^{-k_i^2/2},
\qquad
\psi_{12}(y_1,y_2;X)=e^{\mathrm ck_1k_2},
\qquad \mathrm c\geq1,
\label{eq:gen-gibbs-counterexample}\\
\mathsf W_{\mathrm c}=\begin{psmallmatrix}1&-\mathrm c\\-\mathrm c&1\end{psmallmatrix},
\qquad
\operatorname{spec}\mathsf W_{\mathrm c}=\{1-\mathrm c,1+\mathrm c\};
\label{eq:gen-gibbs-counterexample-matrix}
\end{gather}
in orthonormal sum and difference coordinates the integral factorizes, and the factor in the direction of
the eigenvalue $1-\mathrm c\leq0$ diverges. \hfill$\square$

\remarkheading{Three uses of a reciprocal energy}{rem:gen-appended-energy}
Let $U$ be a measurable energy on $\mathsf Y$, bounded below, that couples coordinates reciprocally. If $U$
enters a potential \eqref{eq:gen-gibbs-potential} with $Z_X=+\infty$, no law is defined and there is no
evidence to bound. If it enters with $0<Z_X<\infty$, the undirected model is a legitimate joint to which
\Cref{ch:elbo} applies. If $U$ is appended to the free energy of a joint $\mathbb P_{\theta,V}$ that does not
contain it, the evidence of that joint is unaffected, and for every recognition law $Q$ in the classical
domain of \Cref{hyp:elbo-evidence-domain} with $\E_Q\lvert U\rvert<\infty$,
\begin{equation}
\Fenergy_V[Q;X,o]+\E_QU=\widetilde\Fenergy_V[Q;X,o]-\log Z_U,
\qquad
Z_U=\int e^{-U(y)}\mathbb P_{\theta,V}(dy,do\given X)\in(0,\infty),
\label{eq:gen-appended-energy}
\end{equation}
where $\widetilde\Fenergy_V$ is the free energy of the tilted joint
$Z_U^{-1}e^{-U(y)}\mathbb P_{\theta,V}(dy,do\given X)$. The appended objective is therefore an evidence bound of
that tilted joint only after $\log Z_U$, which depends on $(\theta,X)$, is restored. \status{ESTABLISHED}

\noindent\emph{Proof.}
Since $U$ is bounded below, $0<Z_U\leq\sup e^{-U}<\infty$, and the tilted joint has density
$p_{\theta,V}(o,y\given X)e^{-U(y)}/Z_U$. With $\Fenergy_V[Q;X,o]=\E_Q[\log q_{V,o,X}(Y)-\log p_{\theta,V}(o,Y\given X)]$,
adding $\E_QU$ and splitting $\log Z_U$ off the tilted log density gives \eqref{eq:gen-appended-energy}; all
expectations are finite under the stated integrability. \hfill$\square$

The directed construction is used because its normalization does not depend on the parameters.

\section{Covariance under passive frame re-expression}
\label{sec:gen-gauge}

A passive re-expression $(g_i)_{i\in V}$ acts on samples through $R_i^x=\rho_x(g_i)$ of
\eqref{eq:geo-represented-frame-change}, and on the latent and joint spaces by the first two lines
below, so records are gauge inert. The structural data transform to $X'$ with the frames
$u_i\cdot g_i^{-1}$ and, by the vertex law \eqref{eq:geo-regime-two-gauge-law}, the links of the third
line; a structural constraint tying parallel copies $e,e':j\to i$ to one link is preserved, because both
copies transform by the same pair (fourth line). Active invariance is invariance under the one copy
$\mathcal A_{c_*}\cong G$ of \Cref{prop:geo-active-image-point}.
\begin{gather}
R=\bigoplus_{i\in V}\bigl(R_i^b\oplus R_i^m\bigr):\mathsf Y\longrightarrow\mathsf Y,
\label{eq:gen-rep-actions}\\
T(y,o)=(Ry,o),
\label{eq:gen-contextual-frame-change}\\
\Theta_e^{x\prime}=g_i\Theta_e^xg_j^{-1},
\qquad e:j\to i,\ x\in\{b,m\},
\label{eq:gen-gauge-links}\\
\Theta_e^x=\Theta_{e'}^x
\quad\Longrightarrow\quad
g_i\Theta_e^xg_j^{-1}=g_i\Theta_{e'}^xg_j^{-1}.
\label{eq:gen-shared-link-admissibility}
\end{gather}

\hypothesisheading{Kernel covariance}{hyp:gen-kernel-covariance}
The transformed parameters $(\theta',X')$ are such that every factor obeys the first two lines below,
with the empty parent tuple at roots, and every record kernel obeys the third, where $R_{\partial a}$ is
the restriction of $R$ to the incident blocks; for a record with scope $\{i\}$ the third line reads as
the fourth:
\begin{gather}
P_{\theta',i}^{m}\bigl(\cdot\given R_{\operatorname{pa}(i)}^mm_{\operatorname{pa}(i)},X'\bigr)
=(R_i^m)_\#P_{\theta,i}^{m}\bigl(\cdot\given m_{\operatorname{pa}(i)},X\bigr),
\label{eq:gen-gauge-pushforward-model}\\
P_{\theta',i}^{k}\bigl(\cdot\given R_{\operatorname{pa}(i)}^bk_{\operatorname{pa}(i)},R_i^mm_i,X'\bigr)
=(R_i^b)_\#P_{\theta,i}^{k}\bigl(\cdot\given k_{\operatorname{pa}(i)},m_i,X\bigr),
\label{eq:gen-gauge-pushforward-state}\\
K_a\bigl(X',R_{\partial a}y_{\partial a},\cdot\bigr)=K_a\bigl(X,y_{\partial a},\cdot\bigr),
\qquad a\in\mathfrak A,
\label{eq:gen-gauge-record-kernel}\\
K_a\bigl(X',(R_i^bk_i,R_i^mm_i),\cdot\bigr)=K_a\bigl(X,(k_i,m_i),\cdot\bigr).
\label{eq:gen-gauge-pushforward-obs}
\end{gather}
The hypothesis excludes mechanisms whose coordinate form is compatible with no re-expression of the
parameters.

\propositionheading{Joint measure covariance and evidence invariance}{prop:gen-product-evidence-invariance}
Under \Cref{hyp:gen-kernel-covariance}, with $R$ and $T$ as in
\eqref{eq:gen-rep-actions}--\eqref{eq:gen-contextual-frame-change}, the first line below holds as
probability measures, without a density or reference hypothesis, and the observation marginals agree, so
the second holds:
\begin{gather}
\mathbb P_{\theta',V}(\cdot\given X')=T_\#\mathbb P_{\theta,V}(\cdot\given X),
\label{eq:gen-gauge-joint-pushforward}\\
p_{\theta',V}(o\given X')=p_{\theta,V}(o\given X)
\quad\text{for }\nu^{\mathsf O}\text{-almost every }o.
\label{eq:gen-gauge-evidence-invariance}
\end{gather}
If latent densities exist and $R_\#\nu^{\mathsf Y}\ll\nu^{\mathsf Y}$, then with $j_R$ of the first line
below the second holds $(\nu^{\mathsf Y}\otimes\nu^{\mathsf O})$-almost everywhere, so composing a density
with $R^{-1}$ represents the pushforward for every density exactly when $j_R=1$ almost everywhere. If every
represented block is an invertible linear map between Euclidean spaces of equal dimension, with Lebesgue
reference (equal within each channel, with $\dim\mathsf K\neq\dim\mathsf M$ allowed), then
$j_R=\lvert\det R\rvert^{-1}$ and the third line holds:
\begin{gather}
j_R(y'):=\bigl(d(R_\#\nu^{\mathsf Y})/d\nu^{\mathsf Y}\bigr)(y'),
\label{eq:gen-gauge-reference-rn}\\
p_{\theta',V}(o,y'\given X')=p_{\theta,V}(o,R^{-1}y'\given X)j_R(y'),
\label{eq:gen-gauge-rn-density}\\
p_{\theta',V}(o,y'\given X')=\lvert\det R\rvert^{-1}p_{\theta,V}(o,R^{-1}y'\given X).
\label{eq:gen-gauge-lebesgue-density}
\end{gather}
Finally, the root law bridges \eqref{eq:gen-root-law-bridge} of $\theta$ and $\theta'$ satisfy
$\widetilde\Phi^{\mathrm{gen}\prime}_i=\widehat\rho_b(g_i)\circ\widetilde\Phi^{\mathrm{gen}}_i\circ\widehat\rho_m(g_i)^{-1}$,
the law \eqref{eq:geo-tildephi-gauge-law} of a cross-morphism representative. \status{ESTABLISHED}

\noindent\emph{Proof.}
Apply \eqref{eq:gen-gauge-pushforward-model} and \eqref{eq:gen-gauge-pushforward-state} along the
topological ordering: finite kernel induction gives
$\mathbb P_{\theta',V}^{\mathsf Y}(\cdot\given X')=R_\#\mathbb P_{\theta,V}^{\mathsf Y}(\cdot\given X)$. For bounded
measurable $F$ on $\mathsf Y\times\mathsf O$, write $\mathbb P_{\theta',V}(\cdot\given X')$ by
\eqref{eq:gen-record-extension}, substitute $y'=Ry$ in the latent integral, and replace the record kernels
at $(X',R_{\partial a}y_{\partial a})$ by those at $(X,y_{\partial a})$ by \eqref{eq:gen-gauge-record-kernel};
this gives $\int F(y,o)\mathbb P_{\theta',V}(dy,do\given X')=\int F(T(y,o))\mathbb P_{\theta,V}(dy,do\given X)$,
which is \eqref{eq:gen-gauge-joint-pushforward}, and projecting to $o$ gives
\eqref{eq:gen-gauge-evidence-invariance}.
For $Q=q\nu^{\mathsf Y}$ and measurable $B$,
$R_\#Q(B)=\int_{R^{-1}B}q(y)\nu^{\mathsf Y}(dy)=\int_Bq(R^{-1}y')(R_\#\nu^{\mathsf Y})(dy')$, so
$R_\#Q=(q\circ R^{-1})j_R\nu^{\mathsf Y}$; applied to the joint density this is
\eqref{eq:gen-gauge-rn-density}, and linear change of variables gives $j_R=\lvert\det R\rvert^{-1}$. For the
bridge, \eqref{eq:gen-gauge-pushforward-state} at a root gives
$\widetilde\Phi^{\mathrm{gen}\prime}_i(\widehat\rho_m(g_i)s)=\int P_{\theta',i}^{k}(\cdot\given R_i^mm,X')s(dm)
=\widehat\rho_b(g_i)\widetilde\Phi^{\mathrm{gen}}_i(s)$. \hfill$\square$

\corollaryheading{Independent channel-frame invariance}{cor:elbo-independent-frame-invariance}
Under \Cref{hyp:gen-kernel-covariance}, let $\mathbb Q_{V,o,X}$ be selected before the re-expression and
transported as a whole, $\mathbb Q'_{V,o,X'}=R_\#\mathbb Q_{V,o,X}$. Then for $\mathbb P_{\theta,V,X}^{\mathsf O}$-almost
every $o$ the posterior transforms as $\boldsymbol\Pi_{\theta',V,o,X'}=R_\#\boldsymbol\Pi_{\theta,V,o,X}$, and
\begin{equation}
\KL(\mathbb Q'_{V,o,X'}\Vert\boldsymbol\Pi_{\theta',V,o,X'})
=\KL(\mathbb Q_{V,o,X}\Vert\boldsymbol\Pi_{\theta,V,o,X}),
\quad
\Lelbo'^{\mathrm{ext}}_V(\mathbb Q'_{V,o,X'};X',o)=\Lelbo^{\mathrm{ext}}_V(\mathbb Q_{V,o,X};X,o)
\label{eq:elbo-gauge-invariance}
\end{equation}
in the extended reals, the singular branch included, with $\Lelbo^{\mathrm{ext}}_V$ as in
\eqref{eq:elbo-extended-definition}; on the classical domain of both sides the same holds for $\Lelbo_V$.
\status{ESTABLISHED}

\noindent\emph{Proof.}
By \eqref{eq:gen-gauge-joint-pushforward}, for measurable $B\subseteq\mathsf Y$ and $A\subseteq\mathsf O$,
$\mathbb P_{\theta',V}(B\times A\given X')=\mathbb P_{\theta,V}(R^{-1}B\times A\given X)
=\int_A(R_\#\boldsymbol\Pi_{\theta,V,o,X})(B)\mathbb P^{\mathsf O}_{\theta,V,X}(do)$, and the observation marginals
agree, so $R_\#\boldsymbol\Pi_{\theta,V,o,X}$ is a version of the new posterior. \Cref{lem:prob-kl-chain-rule}(e)
gives equality of the divergences in $[0,+\infty]$. The evidence agrees by
\eqref{eq:gen-gauge-evidence-invariance}, and substitution in \eqref{eq:elbo-extended-definition} subtracts
equal quantities from equal finite numbers; the classical functional equals the extended one on its domain.
\hfill$\square$

The running example satisfies \Cref{hyp:gen-kernel-covariance} with $\theta'=\theta$, since $\rho_m$ is
trivial and its isotropic state covariances, link-transported means, and record lengths are compatible
with \eqref{eq:gen-gauge-links}.

\section*{Physicist's summary}

The generative model is one normalized joint of models, beliefs given parents, and records, which may
close arbitrary cycles; reciprocal latent conditionals define no joint. No factor may read the
recognition law it is scored against, and passive frame changes preserve evidence and bound.

\chapter[One Evidence Bound and Local-Collective Free Energy]{One Evidence Bound and\\Local-Collective Free Energy}
\label{ch:elbo}
\label{ch:local-collective-elbo}

\Cref{ch:generative} fixed one normalized joint law and \cref{ch:probability}
one correlated recognition coupling of the whole population. This chapter asks
what exact relation binds the two, and how the objective of one agent, or of a
block of agents, is related to that of the population. The answer is one
extended evidence bound, defined for every probability law on the latent space
(\cref{def:elbo-extended,thm:elbo-extended-gap}), and the local--global
potential identity (\cref{thm:obs-local-global-potential}): with the outside
marginal held fixed, the change of the conditional local free energy of a block
equals the change of the collective free energy, so the agents' local
objectives form an exact potential game (\cref{cor:obs-potential-game}). Every
local object is a disintegration of the one bound through
\cref{lem:prob-kl-chain-rule}. The bound is the objective of
\cref{ch:selective-communication,ch:relational-inference,ch:two-timescales},
and \cref{ch:coarsegraining,ch:full-graph-meta-agent-vfe} transport it along
evidence-preserving channels.

\paragraph{Scope.}
The chapter does not claim that a maximizer of the bound over a restricted
recognition family exists, that the bound is tight for any particular
restricted family, that the alternating scheme of exact coordinates converges
(\cref{open:elbo-alternating-convergence}), or that a block of agents together
with its local free energy is a meta-agent. No criterion internal to one bound
that selects a recognition structure is known
(\cref{open:restrict-intrinsic-criterion}). \status{NOT-CLAIMED}

\section{Recognition coupling and total correlation}
\label{sec:elbo-recognition}

The population recognition law is the coupling
$\mathbb Q_{V,o,X}\in\operatorname{Cpl}(\mathbb Q_{\bullet,o,X})$ selected by
\cref{hyp:prob-measurable-coupling-selector}, and every restriction used later
is a subfamily $\mathcal Q'$ of \cref{prop:restrict-principle}. Its marginals
do not determine it once two of them are non-Dirac
(\cref{prop:prob-marginals-do-not-determine-joint}), so they cannot replace it
in an exact entropy or evidence calculation; the next two statements give the
error its exact value and sign.

\corollaryheading{Extended total-correlation chain identity}{thm:elbo-total-correlation-chain}
Let $Y=(Y_j)_{j\in\mathcal J}$ take values in a finite product of standard
Borel spaces, let $Q$ be a probability law with marginals $Q_j$, and let
$(\varrho_j)_{j\in\mathcal J}$ be arbitrary probability laws. With the total
correlation
\begin{equation}
\operatorname{TC}(Q):=\KL\Bigl(Q\Bigm\Vert\bigotimes_{j}Q_j\Bigr)\in[0,+\infty],
\label{eq:obs-total-correlation}
\end{equation}
the identity
\begin{equation}
\KL\Bigl(Q\Bigm\Vert\bigotimes_j\varrho_j\Bigr)
=\operatorname{TC}(Q)+\sum_j\KL(Q_j\Vert\varrho_j)
\label{eq:elbo-total-correlation-chain}
\end{equation}
holds in $[0,+\infty]$, including every singular branch. The infinite branch
occurs for a continuous diagonal: if $U\sim\operatorname{Unif}[0,1]$ and
$Q=\operatorname{law}(U,U)$, both marginals are uniform and
$\operatorname{TC}(Q)=+\infty$. \status{ESTABLISHED}

\noindent\emph{Proof.}
The identity is \eqref{eq:prob-kl-product-reference} of
\cref{lem:prob-kl-chain-rule} with its first term named. The diagonal is null
for the uniform law on the square and carries all the mass of $Q$, so $Q$ is
singular with respect to $Q_1\otimes Q_2$. \hfill$\square$

\propositionheading{Finite-density total correlation and the two signs}{prop:elbo-total-correlation-signs}
Partition the latent coordinates into blocks $j$, let $\mathbb Q_{V,o,X}$
have density $q_{V,o,X}$ with respect to $\nu^{\mathsf Y}$ and block marginal
densities $q_j$, and suppose the joint and block differential entropies are
finite. Then the total correlation
$\operatorname{TC}(\mathbb Q_{V,o,X})\geq0$ of the blocks,
\eqref{eq:obs-total-correlation}, satisfies
\begin{equation}
\E_{\mathbb Q_{V,o,X}}\bigl[\log q_{V,o,X}(Y)\bigr]
=\sum_j\E_{\mathbb Q_{V,o,X}}\bigl[\log q_j(Y_j)\bigr]
+\operatorname{TC}(\mathbb Q_{V,o,X}),
\label{eq:elbo-entropy-defect}
\end{equation}
with equality of the two log-density expectations exactly when the blocks are
independent. If in addition
$\E_{\mathbb Q_{V,o,X}}\lvert\log p_{\theta,V}(o,Y\given X)\rvert<\infty$, the
substituted free energy and its pseudo-ELBO,
\begin{equation}
\widetilde\Fenergy
:=\E_{\mathbb Q_{V,o,X}}\Bigl[-\log p_{\theta,V}(o,Y\given X)
+\sum_j\log q_j(Y_j)\Bigr],
\qquad
\widetilde\Lelbo:=-\widetilde\Fenergy,
\label{eq:elbo-substituted-free-energy}
\end{equation}
satisfy
\begin{equation}
\widetilde\Fenergy=\Fenergy_V-\operatorname{TC}(\mathbb Q_{V,o,X}),
\qquad
\widetilde\Lelbo=\Lelbo_V(\mathbb Q_{V,o,X};X,o)+\operatorname{TC}(\mathbb Q_{V,o,X}),
\label{eq:elbo-total-correlation-signs}
\end{equation}
where $\Lelbo_V$ and $\Fenergy_V$ are the classical functionals of
\cref{thm:elbo-exact-identity}. The pseudo-ELBO exceeds the ELBO by the total
correlation and need not bound the log evidence. \status{ESTABLISHED}

\noindent\emph{Proof.}
Take $\varrho_j=\mathbb Q_{V,o,X,j}$ in
\cref{thm:elbo-total-correlation-chain}; the entropy hypotheses make every
term finite and allow the logarithm of $q_{V,o,X}/\prod_jq_j$ to be split,
which is \eqref{eq:elbo-entropy-defect}. Subtracting
\eqref{eq:elbo-substituted-free-energy} from
$\Fenergy_V=\E_{\mathbb Q_{V,o,X}}[-\log p_{\theta,V}+\log q_{V,o,X}]$ leaves
$\operatorname{TC}$, and negation gives the second identity. \hfill$\square$

\section{Evidence domain}
\label{sec:elbo-hypotheses}

\hypothesisheading{Evidence-domain hypotheses}{hyp:elbo-evidence-domain}
\label{hyp:prob-regular-observation}
\label{hyp:prob-classical-split-elbo}
The extended bound of \cref{sec:elbo-identity} uses (H1)--(H2); its classical
split into two expected log densities also uses (H3)--(H4).

\textbf{(H1) Types and version.} The joint
$\mathbb P_{\theta,V}(\cdot\given X)$ of \cref{ch:generative} has the density
$p_{\theta,V}(o,y\given X)$ relative to its reference, and
$\boldsymbol\Pi_{\theta,V,o,X}$ is the version \eqref{eq:prob-rcp-density} on
the regular set $\mathsf O^{\mathrm{reg}}_{\theta,V,X}$ of
\cref{hyp:prob-regular-conditional-convention}, so every posterior and bound
below holds for $\mathbb P^{\mathsf O}_{\theta,V,X}$-almost every $o$. A
comparison of finitely many parameter values uses an observation in the
intersection of their regular sets. No density of the recognition law is
required.

\textbf{(H2) Positive finite evidence.}
$0<p_{\theta,V}(o\given X)<\infty$; this excludes the observations at which
the posterior formula would divide by zero or by infinity.

\textbf{(H3) Absolute continuity.}
$\mathbb Q_{V,o,X}\ll\boldsymbol\Pi_{\theta,V,o,X}$. Since the posterior is
dominated by $\nu^{\mathsf Y}$, the recognition law then has a density
$q_{V,o,X}=d\mathbb Q_{V,o,X}/d\nu^{\mathsf Y}$, with a jointly measurable
version for parameterized kernels under
\cref{thm:prob-kernel-rn-measurable-version,cor:prob-common-reference-joint-density}.
Clause (H3) excludes every recognition law that charges a posterior-null set;
such a law lies on the singular branch of the extended bound, where the bound
equals $-\infty$ (\cref{prop:elbo-subspace-support-singular}).

\textbf{(H4) Log-integrability.} The generative and recognition halves
\begin{align}
&\E_{\mathbb Q_{V,o,X}}\bigl\lvert\log p_{\theta,V}(o,Y\given X)\bigr\rvert<\infty,
\label{eq:elbo-log-integrability}\\
&\E_{\mathbb Q_{V,o,X}}\bigl\lvert\log q_{V,o,X}(Y)\bigr\rvert<\infty,
\label{eq:prob-integrability}
\end{align}
both hold. Clause (H4) excludes recognition laws whose tails make the two
classical expectations individually infinite. It is not implied by
(H1)--(H3), and the extended bound is well defined without it;
\cref{prop:elbo-integrability-witnesses} gives a Cauchy witness whose formal
difference is $-\infty$ and a countable witness whose two expectations are
both $-\infty$.

\propositionheading{Posterior-null support has infinite cost}{prop:elbo-subspace-support-singular}
Let $Q$ be a probability law on the latent space with $Q(S)=1$ and
$\boldsymbol\Pi_{\theta,V,o,X}(S)=0$ for some measurable $S$. Then (H3) fails
for $Q$, $\KL(Q\Vert\boldsymbol\Pi_{\theta,V,o,X})=+\infty$, and
$\Lelbo^{\mathrm{ext}}_V(Q;X,o)=-\infty$. In particular this holds when the
posterior is dominated by Lebesgue measure on $\R^d$ and $Q$ is supported on a
proper affine subspace. A proper closed support alone does not suffice: the
uniform law on $[0,\tfrac12]$ has divergence $\log2$ from the uniform law on
$[0,1]$. \status{ESTABLISHED}

\noindent\emph{Proof.}
The first claims are the null-set clause of
\cref{prop:prob-density-absolute-continuity} and
\eqref{eq:elbo-extended-definition}. A proper affine subspace is Lebesgue
null, and the uniform law on $[0,\tfrac12]$ has density $2$ on its support
relative to the uniform law on $[0,1]$. \hfill$\square$

\section{The extended evidence bound and its record presentation}
\label{sec:elbo-identity}
\label{sec:local-interaction-record-model}
\label{sec:local-collective-vfe}

\hypothesisheading{Agent state and interaction records}{hyp:local-interaction-kernels}
Agent $i$ has a block space $(\mathsf Y_i,\mathscr Y_i)$, and
$(\mathsf Y,\mathscr Y)=\prod_{i\in V}(\mathsf Y_i,\mathscr Y_i)$. The block may
contain auxiliary latent variables and law points of $(\mathcal E_b)_{c_*}$ and
$(\mathcal E_m)_{c_*}$, which are then latent coordinates at the law-point or
configuration level of \cref{def:prob-description-levels}, not the recognition
marginal of agent $i$. The baseline is the latent population kernel
$\mathbb P^{\mathsf Y}_{\theta,V}:(\mathsf X,\mathscr X)\rightsquigarrow(\mathsf Y,\mathscr Y)$
of \cref{cons:gen-finite-directed-law}. Each record $a\in\mathfrak A$ occurs
once and has a finite scope $\partial a\subseteq V$, a record space
$(\mathsf O_a,\mathscr O_a)$, and a record kernel $K_a$ of the type
\eqref{eq:prob-record-kernel-signature}, all of whose values
$K_a(X,y_{\partial a},\cdot)$ are dominated by one $\sigma$-finite measure
$\nu_a$ with a jointly measurable negative log density $E_{a,o}$ in
$(-\infty,+\infty]$,
\begin{equation}
K_a(X,y_{\partial a},do_a)=e^{-E_{a,o}(X,y_{\partial a})}\nu_a(do_a),
\label{eq:obs-interaction-density}
\end{equation}
where $E_{a,o}$ depends on $o$ only through $o_a$. A record shared by agents
$i$ and $j$ is one element of $\mathfrak A$, and writing both $E_{ij}$ and
$E_{ji}$ asserts two directed records. Binary success records
$K_a(y_{\partial a},\{1\})=e^{-E(y_{\partial a})}$ with measurable
$E:\mathsf Y_{\partial a}\to[0,+\infty]$ are admitted. The kernels and
$\mathbb P^{\mathsf Y}_{\theta,V}$ read neither a recognition law, a
variational parameter, nor a posterior (\cref{req:gen-typing-prohibition}),
and under a common vertex gauge the coordinate presentation of each $K_a$ is
equivariant under the simultaneous change of the incident frames, as in
\eqref{eq:gen-gauge-record-kernel}, so the kernel itself is frame independent.
The hypothesis excludes record families without a common dominating measure
(\cref{open:prob-nondominated-likelihood}) and every generative factor that
reads recognition.

Correlated baseline states are permitted; a graph-local Markov blanket needs in
addition the corresponding locality of the regular conditionals of
$\mathbb P^{\mathsf Y}_{\theta,V}$. Frame independence can be realized by giving
record $a$ its own frame and a link $\Theta^x_{a\leftarrow i}$ on each leg, so
that the density reads $\rho_x(\Theta^x_{a\leftarrow i})y^x_i$; comparing the
belief and model channels still requires the cross morphisms of
\cref{sec:geo-cross-morphisms}. From here on one $X\in\mathsf X$ is fixed and
suppressed from $K_a$, $E_{a,o}$ and $\mathbb P^{\mathsf Y}_{\theta,V}$.

\definitionheading{Record likelihood, slice measure, and extended bound}{def:elbo-extended}
\label{thm:obs-collective-vfe}
Assume \cref{hyp:local-interaction-kernels} and (H1)--(H2), and fix
$o\in\mathsf O^{\mathrm{reg}}_{\theta,V,X}$. The record likelihood and the
interaction energy are
\begin{align}
L_o(y)&:=\prod_{a\in\mathfrak A}e^{-E_{a,o}(y_{\partial a})}\in[0,+\infty),
\label{eq:obs-complete-record-likelihood}\\
H_o(y)&:=-\log L_o(y)=\sum_{a\in\mathfrak A}E_{a,o}(y_{\partial a})
\in(-\infty,+\infty];
\label{eq:obs-interaction-energy}
\end{align}
$H_o$ has one term for each record, never equals $-\infty$ because each factor
density is finite, and can be negative. With $\nu^{\mathsf O}=\bigotimes_a\nu_a$
the joint, normalized on every finite record hypergraph, cyclic or not, by
\cref{prop:gen-exact-normalization}, is
\begin{equation}
\mathbb P_{\theta,V}(dy,do\given X)
=\mathbb P^{\mathsf Y}_{\theta,V}(dy)\prod_{a\in\mathfrak A}K_a(y_{\partial a},do_a)
=\mathbb P^{\mathsf Y}_{\theta,V}(dy)L_o(y)\nu^{\mathsf O}(do).
\label{eq:obs-interaction-joint}
\end{equation}
The evidence and the posterior slice are
\begin{align}
p_{\theta,V}(o\given X)&=\int_{\mathsf Y}L_o(y)\mathbb P^{\mathsf Y}_{\theta,V}(dy)\in(0,+\infty),
\label{eq:obs-global-evidence}\\
\boldsymbol\Pi_{\theta,V,o,X}(dy)&=p_{\theta,V}(o\given X)^{-1}L_o(y)\mathbb P^{\mathsf Y}_{\theta,V}(dy),
\label{eq:obs-global-posterior}
\end{align}
and the slice measure is
\begin{equation}
M_{\theta,V,o,X}(S):=\int_Sp_{\theta,V}(o,y\given X)\nu^{\mathsf Y}(dy)
=\int_SL_o(y)\mathbb P^{\mathsf Y}_{\theta,V}(dy)
=p_{\theta,V}(o\given X)\boldsymbol\Pi_{\theta,V,o,X}(S).
\label{eq:elbo-slice-measure}
\end{equation}
For every probability law $Q$ on $(\mathsf Y,\mathscr Y)$ the extended ELBO and
free energy, which are well-defined extended reals because
$\log p_{\theta,V}(o\given X)$ is finite and relative entropy lies in
$[0,+\infty]$, are
\begin{align}
\Lelbo_V^{\mathrm{ext}}(Q;X,o)&:=\log p_{\theta,V}(o\given X)
-\KL\bigl(Q\Vert\boldsymbol\Pi_{\theta,V,o,X}\bigr)
\in\bigl[-\infty,\log p_{\theta,V}(o\given X)\bigr],
\label{eq:elbo-extended-definition}\\
\Fenergy_V^{\mathrm{ext}}[Q;X,o]&:=-\Lelbo_V^{\mathrm{ext}}(Q;X,o)
=-\log p_{\theta,V}(o\given X)+\KL\bigl(Q\Vert\boldsymbol\Pi_{\theta,V,o,X}\bigr),
\label{eq:elbo-extended-free-energy}
\intertext{and with $\theta$, $V$ and $X$ fixed the collective free energy is written}
\Fenergy^{\mathrm{ext}}_o(Q)&:=-\Lelbo^{\mathrm{ext}}_o(Q):=\Fenergy_V^{\mathrm{ext}}[Q;X,o].
\label{eq:obs-collective-vfe}
\end{align}
No domination, density, or mean-field condition on $Q$ enters.

\corollaryheading{Extended gap identity and energy form}{thm:elbo-extended-gap}
For every probability law $Q$ on the latent space,
\begin{equation}
\log p_{\theta,V}(o\given X)-\Lelbo_V^{\mathrm{ext}}(Q;X,o)
=\Fenergy_V^{\mathrm{ext}}[Q;X,o]+\log p_{\theta,V}(o\given X)
=\KL\bigl(Q\Vert\boldsymbol\Pi_{\theta,V,o,X}\bigr),
\label{eq:elbo-extended-gap}
\end{equation}
so $\Lelbo_V^{\mathrm{ext}}\leq\log p_{\theta,V}(o\given X)$ with equality
exactly when $Q=\boldsymbol\Pi_{\theta,V,o,X}$ as measures. If moreover
$Q\ll\mathbb P^{\mathsf Y}_{\theta,V}$,
$\KL(Q\Vert\mathbb P^{\mathsf Y}_{\theta,V})<\infty$, and the negative part of
$H_o$ is $Q$-integrable, then
\begin{equation}
\Fenergy^{\mathrm{ext}}_o(Q)=\KL\bigl(Q\Vert\mathbb P^{\mathsf Y}_{\theta,V}\bigr)+\E_QH_o
\quad\text{in }(-\infty,+\infty].
\label{eq:obs-collective-energy-form}
\end{equation}
No formula $\log p_{\theta,V}=\Lelbo_V^{\mathrm{ext}}+\KL$ is asserted on the
singular branch, where it would read $-\infty+\infty$. \status{ESTABLISHED}

\noindent\emph{Proof.}
The gap identity adds or subtracts the finite number
$\log p_{\theta,V}(o\given X)$ in
\eqref{eq:elbo-extended-definition}--\eqref{eq:elbo-extended-free-energy}, and
the equality case is the zero case of relative entropy. For the energy form
let $r'=dQ/d\mathbb P^{\mathsf Y}_{\theta,V}$. If $Q\ll M_{\theta,V,o,X}$, then
$dQ/dM_{\theta,V,o,X}=r'e^{H_o}$ $Q$-almost surely, and
\cref{prop:elbo-relative-log-representation} gives
$\Fenergy^{\mathrm{ext}}_o(Q)=\int(\log r'+H_o)dQ$; the first term integrates
to the finite $\KL(Q\Vert\mathbb P^{\mathsf Y}_{\theta,V})$ and the second is
quasi-integrable, so the integral splits. If $Q\not\ll M_{\theta,V,o,X}$, some
set null for $M_{\theta,V,o,X}$ has positive $Q$-mass; since
$M_{\theta,V,o,X}$ and $\mathbb P^{\mathsf Y}_{\theta,V}$ agree in null sets on
$\{L_o>0\}$ and $Q\ll\mathbb P^{\mathsf Y}_{\theta,V}$, that mass lies in
$\{H_o=+\infty\}$, and both sides equal $+\infty$. \hfill$\square$

\theoremheading{Classical finite evidence identity}{thm:elbo-exact-identity}
\label{cor:elbo-bound-tightness}
Under \cref{hyp:elbo-evidence-domain}, the extended bound has the separately
integrable representation
\begin{align}
\Lelbo_V(\mathbb Q_{V,o,X};X,o)
&:=\E_{\mathbb Q_{V,o,X}}\bigl[\log p_{\theta,V}(o,Y\given X)-\log q_{V,o,X}(Y)\bigr]
=\Lelbo_V^{\mathrm{ext}}(\mathbb Q_{V,o,X};X,o),
\label{eq:elbo-definition}\\
\Fenergy_V[\mathbb Q_{V,o,X};X,o]&:=-\Lelbo_V(\mathbb Q_{V,o,X};X,o),
\label{eq:elbo-free-energy}
\end{align}
and for $\mathbb P^{\mathsf O}_{\theta,V,X}$-almost every observation
\begin{equation}
\log p_{\theta,V}(o\given X)=\Lelbo_V(\mathbb Q_{V,o,X};X,o)
+\KL\bigl(\mathbb Q_{V,o,X}\Vert\boldsymbol\Pi_{\theta,V,o,X}\bigr).
\label{eq:elbo-identity}
\end{equation}
Consequently $\Lelbo_V\leq\log p_{\theta,V}(o\given X)$, with equality if and
only if $\mathbb Q_{V,o,X}=\boldsymbol\Pi_{\theta,V,o,X}$ as probability
measures on $(\mathsf Y,\mathscr Y)$. Equality of all coordinate marginals is
not sufficient. \status{ESTABLISHED}

\noindent\emph{Proof.}
The posterior and the slice measure have the same null sets, so (H3) gives
$\mathbb Q_{V,o,X}\ll M_{\theta,V,o,X}$, and since the two measures have the
densities $q_{V,o,X}$ and $p_{\theta,V}(o,\cdot\given X)$ relative to
$\nu^{\mathsf Y}$, $d\mathbb Q_{V,o,X}/dM_{\theta,V,o,X}=q_{V,o,X}/p_{\theta,V}(o,\cdot\given X)$
$\mathbb Q_{V,o,X}$-almost surely.
\Cref{prop:elbo-relative-log-representation} gives
$\Lelbo_V^{\mathrm{ext}}=-\int\log(q_{V,o,X}/p_{\theta,V}(o,\cdot\given X))d\mathbb Q_{V,o,X}$,
and (H4) splits the integral, which is \eqref{eq:elbo-definition}. The bound is
then finite, so \eqref{eq:elbo-extended-definition} rearranges to
\eqref{eq:elbo-identity}, and the equality case is that of
\cref{thm:elbo-extended-gap}. Distinct measures with equal coordinate marginals
are supplied by \cref{prop:prob-marginals-do-not-determine-joint}.
\hfill$\square$

\section{Restriction principle}
\label{sec:restrict-principle}

\propositionheading{Restriction principle}{prop:restrict-principle}
Fix $(\theta,X)$ and $o$ as in \cref{def:elbo-extended}, and write
$P^\star=\boldsymbol\Pi_{\theta,V,o,X}$, so that for every law $Q$
\begin{equation}
\Lelbo_V^{\mathrm{ext}}(Q;X,o)=\log p_{\theta,V}(o\given X)-\KL(Q\Vert P^\star).
\label{eq:restrict-exact-identity}
\end{equation}
For arbitrary sets of probability laws $\mathcal Q'\subseteq\mathcal Q$,
\begin{equation}
\sup_{Q\in\mathcal Q'}\Lelbo_V^{\mathrm{ext}}(Q;X,o)
=\log p_{\theta,V}(o\given X)-\inf_{Q\in\mathcal Q'}\KL(Q\Vert P^\star),
\label{eq:restrict-sup-identity}
\end{equation}
the supremum over $\mathcal Q'$ is at most the supremum over $\mathcal Q$, and
when both infima are finite
\begin{equation}
\sup_{\mathcal Q}\Lelbo_V^{\mathrm{ext}}-\sup_{\mathcal Q'}\Lelbo_V^{\mathrm{ext}}
=\inf_{\mathcal Q'}\KL(Q\Vert P^\star)-\inf_{\mathcal Q}\KL(Q\Vert P^\star)\geq0.
\label{eq:restrict-loss}
\end{equation}
No attainment of the suprema or infima is required. \status{ESTABLISHED}

\noindent\emph{Proof.}
Identity \eqref{eq:restrict-exact-identity} is
\eqref{eq:elbo-extended-definition}. Taking suprema turns a finite constant
minus a nonnegative extended divergence into that constant minus its infimum,
including the value $-\infty$. Inclusion gives monotonicity, and subtraction is
legitimate only when both infima are finite. \hfill$\square$

The evidence in \eqref{eq:restrict-exact-identity} is constant because the
generative kernel does not read $Q$ (\cref{req:gen-typing-prohibition}); the
next two propositions show what fails when a model moves with $Q$.

\remarkheading{Orientation, separation, and attainment}{rem:restrict-attainment}
The infimum in \eqref{eq:restrict-sup-identity} is a reverse information
projection, since $Q$ is the first argument of the divergence. Excluding
$P^\star$ from a family does not make the optimum positive or attained: for
$P^\star=\mathcal N(0,1)$ the family $\{\mathcal N(\mu,1):\mu>0\}$ has gaps
$\mu^2/2$ with unattained infimum zero. A positive optimized gap is the
condition $\inf_{Q\in\mathcal Q}\KL(Q\Vert P^\star)>0$. A family constrained to
product laws need not be convex, because mixtures of product laws need not be
product laws. Strict convexity of relative entropy on the ambient law space
therefore does not establish uniqueness on that family; on a convex feasible
family it does give at most one finite minimizer. \status{ESTABLISHED}

\propositionheading{No distinguished target}{prop:gen-no-distinguished-target}
Let $\mathcal Q$ be a nonempty family of recognition laws, and replace one
factor of \eqref{eq:gen-composition} by a nonconstant $Q$-indexed family, giving
a normalized joint $P_{\theta,Q}$ for each $Q\in\mathcal Q$. Assume at the
observation that $0<p_{\theta,Q}(o\given X)<\infty$, that
$Q\ll P_{\theta,Q}(\cdot\given o,X)$, and that the log densities defining the
ELBO are integrable. Then for each fixed $Q$
\begin{equation}
\log p_{\theta,Q}(o\given X)=\Lelbo(Q;X)
+\KL\bigl(Q\Vert P_{\theta,Q}(\cdot\given o,X)\bigr),
\label{eq:gen-moving-target}
\end{equation}
and with $\mathsf e(Q)=\log p_{\theta,Q}(o\given X)$ two members satisfy
\begin{equation}
\Lelbo(Q';X)-\Lelbo(Q;X)
=\mathsf e(Q')-\mathsf e(Q)
-\bigl[\KL\bigl(Q'\Vert P_{\theta,Q'}(\cdot\given o,X)\bigr)
-\KL\bigl(Q\Vert P_{\theta,Q}(\cdot\given o,X)\bigr)\bigr].
\label{eq:gen-elbo-difference}
\end{equation}
On a level set of $\mathsf e$ the change of the bound is the negative change of
the divergence; off it the two changes differ by $\mathsf e(Q')-\mathsf e(Q)$.
The family selects no joint as the model. \status{ESTABLISHED}

\noindent\emph{Proof.}
Apply \cref{thm:elbo-exact-identity} to each fixed pair $(Q,P_{\theta,Q})$ and
subtract. No member is distinguished because no selection rule was declared.
\hfill$\square$

\propositionheading{The stronger consequences fail}{prop:gen-moving-target-witness}
Nonconstant dependence of $P_{\theta,Q}$ on $Q$ does not imply that $\mathsf e$
varies, that no ascent improves a fixed quantity, or that no single member's
evidence bounds the family of bounds. Explicitly, for one binary latent $y$, one
binary observation, and $Q_{\mathrm p}=\operatorname{Ber}(\mathrm p)$ with
$\mathrm p\in(0,1)$, the joint
\begin{equation}
P_{\theta,Q_{\mathrm p}}(o,y\given X)=\tfrac12\operatorname{Ber}\bigl(\tfrac12\mathrm p+\tfrac14\bigr)(y)
\label{eq:gen-moving-target-witness}
\end{equation}
is normalized and injective in $\mathrm p$, has $\mathsf e(Q_{\mathrm p})=\log\tfrac12$
for every $\mathrm p$, and has gaps
\begin{equation}
\tfrac9{10}\log\tfrac97+\tfrac1{10}\log\tfrac13=0.1163218\ldots,
\qquad
\tfrac7{10}\log\tfrac76+\tfrac3{10}\log\tfrac34=0.0216009\ldots,
\qquad
0
\label{eq:gen-moving-target-values}
\end{equation}
at $\mathrm p=\tfrac9{10},\tfrac7{10},\tfrac12$. \status{ESTABLISHED}

\noindent\emph{Proof.}
The observation marginal is $\tfrac12$ for every $\mathrm p$, and the posterior
is $\operatorname{Ber}(\tfrac12\mathrm p+\tfrac14)$, so the gaps are the
Bernoulli relative entropies displayed. The bound increases to the fixed
evidence along the three values, and the member at $\mathrm p=\tfrac12$ has
evidence equal to the supremum of the bounds. \hfill$\square$

\propositionheading{ELBO values do not order unrelated joints}{prop:restrict-unrelated-joint-elbo}
Let $P_\theta$ and $P_{\theta'}$ be normalized joints over the same observation,
possibly with different latent inventories, with recognition families
$\mathcal Q$ and $\mathcal Q'$ on their latent spaces, and let $o$ lie in the
regular set of each. Suppose both infima of the posterior relative entropy
over the respective recognition families are finite. If no evidence-preserving
channel relates the two joints, their optimized bounds satisfy
\begin{equation}
\Lelbo^\star-\Lelbo'^\star
=\bigl[\log p_\theta(o\given X)-\log p_{\theta'}(o\given X)\bigr]
-\Bigl[\inf_{\mathcal Q}\KL-\inf_{\mathcal Q'}\KL\Bigr],
\label{eq:restrict-cross-model}
\end{equation}
which combines an evidence difference, a model-comparison quantity, with a
tightness difference, a recognition-quality quantity. The sign of the left side
orders neither the models nor the recognition families: with binary $O$ and $Y$,
constants $\mathrm c_1\leq0\leq\mathrm c_2$, and the normalized joints
$P(o,1)=e^{\mathrm c_1-\mathrm c_2}$, $P(o,0)=e^{\mathrm c_1}(1-e^{-\mathrm c_2})$,
$P(\bar o,0)=1-e^{\mathrm c_1}$, $P(\bar o,1)=0$ under the singleton family
$\{\delta_1\}$, the pairs $(\mathrm c_1,\mathrm c_2)=(-1,0),(0,2)$ give bounds
$-1>-2$ with the reverse evidence order, and the pairs $(0,1),(-2,0)$ give bounds
$-1>-2$ with the reverse tightness order. \status{ESTABLISHED}

\noindent\emph{Proof.}
Apply \eqref{eq:restrict-sup-identity} to each model and subtract. In the
witness, $\log P(O=o)=\mathrm c_1$ and $P(Y=1\given o)=e^{-\mathrm c_2}$, so the
minimized gap of $\{\delta_1\}$ is $\mathrm c_2$ and the optimized bound is
$\mathrm c_1-\mathrm c_2$. An unused deterministic coordinate appended to either
model, with its point mass appended to the recognition law, makes the latent
inventories literally different and changes no evidence, divergence, or bound.
\hfill$\square$

\openproblemheading{An intrinsic criterion}{open:restrict-intrinsic-criterion}
No criterion internal to one evidence bound at a fixed model is known that
selects a nondegenerate recognition structure. Along nested recognition
families the optimized bound is monotone by inclusion
(\cref{prop:restrict-principle}), so minimizing the gap selects the largest
family; an arbitrary replacement joint supplies no common evidence
(\cref{prop:restrict-unrelated-joint-elbo}); and evidence-preserving channels,
which do order bounds across latent spaces, do so without selecting a
nondegenerate partition (\cref{sec:cg-elbo}). The problem is settled
positively by a functional on admissible channels or partitions, expressible
relative to one fixed joint and not monotone toward an endpoint, whose
optimizer is nondegenerate, and negatively by a proof of endpoint degeneracy
for a precisely stated class of such functionals.

\section{The factorwise decomposition}
\label{sec:elbo-factorization}

\propositionheading{Exact factorwise decomposition}{prop:elbo-factorwise-decomposition}
\label{hyp:elbo-factorwise-integrability}
Assume (H1)--(H3) of \cref{hyp:elbo-evidence-domain} and factorwise
log-integrability: the absolute logarithm of every component density of
\eqref{eq:gen-density} and every record energy $E_{a,o}$ is
$\mathbb Q_{V,o,X}$-integrable. Then (i) the generative half
\eqref{eq:elbo-log-integrability} of (H4) holds, and (H4) holds exactly when the
recognition half \eqref{eq:prob-integrability} does; (ii) under (H4) the free
energy has exactly one decomposition into factor terms,
\begin{align}
\Fenergy_V[\mathbb Q_{V,o,X};X,o]
={}&\sum_{i\in V}\E_{\mathbb Q_{V,o,X}}\bigl[-\log p^m_{\theta,i}(m_i\given m_{\operatorname{pa}(i)},X)\bigr]
+\sum_{i\in V}\E_{\mathbb Q_{V,o,X}}\bigl[-\log p^k_{\theta,i}(k_i\given k_{\operatorname{pa}(i)},m_i,X)\bigr]
\notag\\
&+\sum_{a\in\mathfrak A}\E_{\mathbb Q_{V,o,X}}\bigl[E_{a,o}(Y_{\partial a})\bigr]
+\E_{\mathbb Q_{V,o,X}}\bigl[\log q_{V,o,X}(Y)\bigr],
\label{eq:elbo-factorization}
\end{align}
with $\operatorname{pa}(i)=\varnothing$ at the roots, $2N+|\mathfrak A|$
generative factors, and a final term that is one full-joint recognition log
density, not a sum of marginal entropies; (iii) factorwise integrability does
not imply the recognition half, and the generative half does not imply
factorwise integrability. \status{ESTABLISHED}

\noindent\emph{Proof.}
By \eqref{eq:gen-density} and \cref{hyp:local-interaction-kernels},
\begin{equation}
\log p_{\theta,V}(o,y\given X)
=\sum_{i\in V}\bigl[\log p^m_{\theta,i}(m_i\given m_{\operatorname{pa}(i)},X)
+\log p^k_{\theta,i}(k_i\given k_{\operatorname{pa}(i)},m_i,X)\bigr]
-\sum_{a\in\mathfrak A}E_{a,o}(y_{\partial a}),
\label{eq:elbo-audit}
\end{equation}
a finite sum, so the triangle inequality gives (i). Under (H4), taking
$\mathbb Q_{V,o,X}$-expectations of the negative of \eqref{eq:elbo-audit}
term by term and adding $\E_{\mathbb Q_{V,o,X}}[\log q_{V,o,X}]$ gives
\eqref{eq:elbo-factorization} from \eqref{eq:elbo-definition}. For (iii),
\cref{sec:app-integrability-witnesses} gives a model with bounded factor
logarithms whose recognition density has an infinite mean logarithm, and a
root law $\mathrm c/(m(\log m)^2)$ whose state bridge cancels its logarithm, so
that the complete log joint is integrable while one factor is not.
\hfill$\square$

\section{Observations as interaction records}
\label{sec:local-observation-equivalence}

\definitionheading{Operational environment node}{def:obs-operational-environment-node}
An operational environment node is a standard Borel state space with a
normalized state law and a measurable message kernel.

\theoremheading{Observation--interaction kernel equivalence}{thm:obs-agent-interaction-equivalence}
Every normalized observation kernel $K(do\given y)$ from a measurable space to a
standard Borel space is the message of an operational environment node: there
are $U\sim\operatorname{Unif}[0,1]$ and a measurable map $F$ such that
\begin{equation}
O=F(Y,U),
\qquad
\Pr\bigl(F(y,U)\in A\bigr)=K(A\given y).
\label{eq:obs-randomization}
\end{equation}
Conversely, marginalizing any environment state and message policy produces a
normalized observation kernel. \status{ESTABLISHED}

\noindent\emph{Proof.}
The forward direction is the randomization lemma for kernels
\citep{Kallenberg2021}. By Kuratowski's theorem a standard Borel space admits a
Borel isomorphism $\phi$ onto a Borel subset of $\R$. The function
$G(y,x)=K(\phi^{-1}(-\infty,x]\given y)$ is measurable in $y$ and a
distribution function in $x$, so the quantile
$F_0(y,u)=\inf\{x:G(y,x)\geq u\}$ on $(0,1)$ satisfies
$\{F_0\leq x\}=\{(y,u):u\leq G(y,x)\}$; hence $F_0$ is jointly measurable and
$\Pr(F_0(y,U)\leq x)=G(y,x)$. The value $F_0(y,U)$ lies in the image of $\phi$
almost surely; setting $F=\phi^{-1}\circ F_0$ there and a fixed point elsewhere
gives \eqref{eq:obs-randomization}. Conversely, a state law
$\varrho_{\mathrm{env}}(du)$ and a message policy
$\kappa_{\mathrm{env}}(y,u,A)$ give
$K(A\given y)=\int\kappa_{\mathrm{env}}(y,u,A)\varrho_{\mathrm{env}}(du)$, which
is measurable in $y$, countably additive in $A$, and normalized.
\hfill$\square$

\propositionheading{Conditioning information and retained statistics}{prop:obs-conditioning-information}
Let $(Y,O)$ have joint law $P$ on standard Borel spaces with marginals $P^Y$ and
$P^O$, let $I(Y;O)=\KL(P\Vert P^Y\otimes P^O)$, and for a measurable statistic
$S=T(O)$ let $I(Y;O\given S)=\int\KL(P_{Y,O\given s}\Vert P_{Y\given s}\otimes P_{O\given s})P^S(ds)$.
The information carried by the record is
\begin{equation}
I(Y;O)=\int\KL\bigl(P(dY\given o)\Vert P^Y\bigr)P^O(do)\in[0,+\infty].
\label{eq:obs-conditioning-information}
\end{equation}
Moreover
\begin{equation}
I(Y;O)=I(Y;S)+I(Y;O\given S)
\quad\text{in }[0,+\infty],
\label{eq:obs-conditioning-statistic-chain}
\end{equation}
so replacing $O$ by $S$ leaves the posterior unchanged exactly when
$I(Y;O\given S)=0$, up to the choice of conditional versions. Deleting the record
makes $S$ constant and returns the prior; for
$Y\sim\operatorname{Ber}(1/2)$ and $O=Y$, $I(Y;O)=\log2$ and
$P(dY\given O=o)=\delta_o$. A kernel that does not depend on $y$ carries zero
information. \status{ESTABLISHED}

\noindent\emph{Proof.}
Disintegrating $\KL(P\Vert P^Y\otimes P^O)$ over $O$ with
\cref{lem:prob-kl-chain-rule}, whose $O$-marginals agree, gives
\eqref{eq:obs-conditioning-information}. Since $o\mapsto(T(o),o)$ is a
bimeasurable bijection onto its graph, relative entropy is unchanged by passing
to $(Y,S,O)$ \eqref{eq:cg-kl-bimeasurable-invariance}. Disintegrating
$\KL(P_{Y,S,O}\Vert P^Y\otimes P_{S,O})$ over $(Y,S)$ gives
$I(Y;S)+\E_{P_{Y,S}}\KL(P_{O\given Y,S}\Vert P_{O\given S})$, because under the
reference $O$ is independent of $Y$ given $S$ with law $P_{O\given S}$; for
each $s$, disintegrating
$\KL(P_{Y,O\given s}\Vert P_{Y\given s}\otimes P_{O\given s})$ over $Y$ and
integrating in $s$ identifies the second term with $I(Y;O\given S)$, which
proves \eqref{eq:obs-conditioning-statistic-chain}; the discrete case is
\citet[Sec.~2.5]{CoverThomas2006}. Posterior sufficiency is the vanishing of
the conditional term, the Bernoulli values are direct, and a $y$-independent
kernel makes $P=P^Y\otimes P^O$. \hfill$\square$

\remarkheading{Environment nodes in the variational problem}{rem:obs-environment-fiber}
\Cref{thm:obs-agent-interaction-equivalence} is exact for the joint law of
$(Y,O)$, and it keeps the record variable $O$ and hence the conditioning
$\sigma$-algebra $\sigma(O)$, but it is not exact for a variational problem
over the enlarged state $(Y,U)$. For
$P^O$-almost every $o$ the posterior of $(Y,U)$ is concentrated on the level set
$\{F(y,u)=o\}$. For $Y\sim\mathcal N(0,1)$ and
$F(y,u)=y+F_{\mathcal N}^{-1}(u)$, with $F_{\mathcal N}$ the standard normal
distribution function, which realizes $K(do\given y)=\mathcal N(y,1)$, that
level set is the Lebesgue-null curve $\{u=F_{\mathcal N}(o-y)\}$, so every
recognition law on $(Y,U)$ with a Lebesgue density has relative entropy
$+\infty$ by \cref{prop:elbo-subspace-support-singular}, whereas the
fiber-restricted law $q(dy)\delta_{F_{\mathcal N}(o-y)}(du)$, the image of $q$
under the injective map $y\mapsto(y,F_{\mathcal N}(o-y))$ that also carries the
posterior of $Y$ to that of $(Y,U)$, has relative entropy
$\KL(q\Vert P(dY\given o))$. A finite variational problem therefore integrates
$U$ out or restricts recognition to that fiber. \status{ESTABLISHED}

\section{Local conditional free energy and the potential identity}
\label{sec:local-conditional-vfe}

Fix a nonempty block $B\subseteq V$, one agent or any aggregate, and write
$y=(y_B,y_{-B})$; the probability model does not change with $B$.

\lemmaheading{Block disintegration of the posterior}{lem:obs-block-posterior}
Let $\mathfrak A_B=\{a\in\mathfrak A:\partial a\cap B\ne\varnothing\}$ and define
the incident and outside likelihoods and energies
\begin{align}
L_{B,o}(y_B;y_{-B})&=\prod_{a\in\mathfrak A_B}e^{-E_{a,o}(y_{\partial a})},
&L_{\bar B,o}(y_{-B})&=\prod_{a\notin\mathfrak A_B}e^{-E_{a,o}(y_{\partial a})},
\label{eq:obs-block-incident-likelihood}\\
H_{B,o}(y_B;y_{-B})&=-\log L_{B,o}=\sum_{a\in\mathfrak A_B}E_{a,o}(y_{\partial a}),
&H_{\bar B,o}(y_{-B})&=-\log L_{\bar B,o}(y_{-B}),
\label{eq:obs-block-incident-energy}
\end{align}
so that $L_o=L_{B,o}L_{\bar B,o}$ and $H_o=H_{B,o}+H_{\bar B,o}$ in
$(-\infty,+\infty]$. Fix a measurable version of the baseline disintegration
\begin{equation}
\mathbb P^{\mathsf Y}_{\theta,V}(dy_B,dy_{-B})
=\mathbb P^{\mathsf Y}_{\theta,V,-B}(dy_{-B})\mathbb P^{\mathsf Y}_{\theta,V,B}(dy_B\given y_{-B}),
\label{eq:obs-baseline-block-disintegration}
\end{equation}
and define the local evidence and the likelihood of the outside state,
\begin{align}
Z_B(y_{-B})&=\int L_{B,o}(y_B;y_{-B})\mathbb P^{\mathsf Y}_{\theta,V,B}(dy_B\given y_{-B}),
\label{eq:obs-local-evidence}\\
\widetilde L_{B,o}(y_{-B})&=\int L_o(y_B,y_{-B})\mathbb P^{\mathsf Y}_{\theta,V,B}(dy_B\given y_{-B}).
\label{eq:obs-full-complement-likelihood}
\end{align}
Then (i) the evidence and the outside posterior marginal are
\begin{equation}
p_{\theta,V}(o\given X)=\int\widetilde L_{B,o}d\mathbb P^{\mathsf Y}_{\theta,V,-B},
\qquad
\boldsymbol\Pi_{\theta,V,o,X,-B}(dy_{-B})
=p_{\theta,V}(o\given X)^{-1}\widetilde L_{B,o}(y_{-B})\mathbb P^{\mathsf Y}_{\theta,V,-B}(dy_{-B});
\label{eq:obs-posterior-outside}
\end{equation}
(ii) the set $\Omega_{B,o}$ of outside states at which
$0<\widetilde L_{B,o}<\infty$, $0<L_{\bar B,o}<\infty$, $0<Z_B<\infty$, and
$\widetilde L_{B,o}=L_{\bar B,o}Z_B$ has $\boldsymbol\Pi_{\theta,V,o,X,-B}$-measure
one, though it need not be $\mathbb P^{\mathsf Y}_{\theta,V,-B}$-full; (iii) the
kernel
\begin{equation}
\boldsymbol\Pi_{\theta,V,o,X,B}(dy_B\given y_{-B})
=Z_B(y_{-B})^{-1}L_{B,o}(y_B;y_{-B})\mathbb P^{\mathsf Y}_{\theta,V,B}(dy_B\given y_{-B})
\qquad(y_{-B}\in\Omega_{B,o}),
\label{eq:obs-local-posterior}
\end{equation}
extended by any measurable probability kernel off $\Omega_{B,o}$, satisfies
\begin{equation}
\boldsymbol\Pi_{\theta,V,o,X}(dy_B,dy_{-B})
=\boldsymbol\Pi_{\theta,V,o,X,-B}(dy_{-B})\boldsymbol\Pi_{\theta,V,o,X,B}(dy_B\given y_{-B}).
\label{eq:obs-posterior-block-disintegration}
\end{equation}
The posterior determines this kernel, and every pointwise local quantity built
from it, only $\boldsymbol\Pi_{\theta,V,o,X,-B}$-almost everywhere; two baseline
versions give local posteriors that agree there. \status{ESTABLISHED}

\noindent\emph{Proof.}
Tonelli's theorem applied to
\eqref{eq:obs-global-evidence}--\eqref{eq:obs-global-posterior} with
\eqref{eq:obs-baseline-block-disintegration} gives (i); defining $\widetilde L_{B,o}$ directly avoids the product
$L_{\bar B,o}Z_B$, which can be $0\cdot\infty$ off $\Omega_{B,o}$. For (ii),
finite evidence makes $\widetilde L_{B,o}<\infty$ off a
$\mathbb P^{\mathsf Y}_{\theta,V,-B}$-null set, and by (i) the outside posterior
charges only $\{\widetilde L_{B,o}>0\}$. There $L_{\bar B,o}$ is finite because
each factor density is, positive because $L_{\bar B,o}=0$ would force
$\widetilde L_{B,o}=0$, and independent of $y_B$, so it leaves the conditional
integral and $Z_B=\widetilde L_{B,o}/L_{\bar B,o}\in(0,\infty)$. For (iii), a
bounded measurable $f$ has
$\int fd\boldsymbol\Pi_{\theta,V,o,X}
=p_{\theta,V}(o\given X)^{-1}\int L_{\bar B,o}\int fL_{B,o}d\mathbb P^{\mathsf Y}_{\theta,V,B}d\mathbb P^{\mathsf Y}_{\theta,V,-B}$,
and on $\Omega_{B,o}$ the inner integral is $Z_B$ times the integral of $f$
against \eqref{eq:obs-local-posterior}; (i) and (ii) turn the outer weight into
$\boldsymbol\Pi_{\theta,V,o,X,-B}$. Almost-everywhere uniqueness is uniqueness
of disintegrations; two baseline versions agree
$\mathbb P^{\mathsf Y}_{\theta,V,-B}$-almost everywhere, hence
$\boldsymbol\Pi_{\theta,V,o,X,-B}$-almost everywhere by (i). \hfill$\square$

\definitionheading{Local conditional free energy}{thm:obs-local-multiagent-elbo}
For $y_{-B}\in\Omega_{B,o}$ and every probability law $r_B(\cdot\given y_{-B})$
on $\mathsf Y_B$, the local conditional ELBO and free energy of the block are
\begin{equation}
\Fenergy^{\mathrm{ext}}_{B,o}(r_B;y_{-B}):=-\Lelbo^{\mathrm{ext}}_{B,o}(r_B;y_{-B})
:=-\log Z_B(y_{-B})
+\KL\bigl(r_B(\cdot\given y_{-B})\Vert\boldsymbol\Pi_{\theta,V,o,X,B}(\cdot\given y_{-B})\bigr).
\label{eq:obs-local-vfe}
\end{equation}
For $B=\{i\}$ this is the local free energy of agent $i$. It contains every
record incident to the block and no other.

\corollaryheading{Local Gibbs bound and the exact block E-coordinate}{prop:elbo-exact-e-coordinate}
(i) $\Lelbo^{\mathrm{ext}}_{B,o}(r_B;y_{-B})\leq\log Z_B(y_{-B})$, with
equality exactly at $r_B=\boldsymbol\Pi_{\theta,V,o,X,B}(\cdot\given y_{-B})$.
If $r_B\ll\mathbb P^{\mathsf Y}_{\theta,V,B}(\cdot\given y_{-B})$ with finite
relative entropy and the negative part of $H_{B,o}(\cdot;y_{-B})$ is
$r_B$-integrable,
\begin{equation}
\Fenergy^{\mathrm{ext}}_{B,o}(r_B;y_{-B})
=\KL\bigl(r_B(\cdot\given y_{-B})\Vert\mathbb P^{\mathsf Y}_{\theta,V,B}(\cdot\given y_{-B})\bigr)
+\E_{r_B(\cdot\given y_{-B})}H_{B,o}(Y_B;y_{-B}).
\label{eq:obs-local-energy-form}
\end{equation}

(ii) If $Q(dy_B,dy_{-B})=Q_{-B}(dy_{-B})r_B(dy_B\given y_{-B})$ and
$Q_{-B}\ll\boldsymbol\Pi_{\theta,V,o,X,-B}$, then in $[0,+\infty]$, with the
baseline conditional of \eqref{eq:obs-baseline-block-disintegration} in the
second identity,
\begin{align}
\KL\bigl(Q\Vert\boldsymbol\Pi_{\theta,V,o,X}\bigr)
&=\KL\bigl(Q_{-B}\Vert\boldsymbol\Pi_{\theta,V,o,X,-B}\bigr)
+\E_{Q_{-B}}\KL\bigl(r_B(\cdot\given Y_{-B})\Vert\boldsymbol\Pi_{\theta,V,o,X,B}(\cdot\given Y_{-B})\bigr),
\label{eq:obs-posterior-block-kl-chain}\\
\KL\bigl(Q\Vert\mathbb P^{\mathsf Y}_{\theta,V}\bigr)
&=\KL\bigl(Q_{-B}\Vert\mathbb P^{\mathsf Y}_{\theta,V,-B}\bigr)
+\E_{Q_{-B}}\KL\bigl(r_B(\cdot\given Y_{-B})\Vert\mathbb P^{\mathsf Y}_{\theta,V,B}(\cdot\given Y_{-B})\bigr).
\label{eq:obs-block-kl-chain}
\end{align}

(iii) Let $\mathcal B$ be any block of latent coordinates with complement
$-\mathcal B$, let $P_-$ be the $-\mathcal B$ marginal of
$P=\boldsymbol\Pi_{\theta,V,o,X}$, let $R$ be any law on the $-\mathcal B$
coordinates, and let $\kappa_{\mathcal B}(dY_{\mathcal B}\given Y_{-\mathcal B})$
be a Markov kernel that agrees $P_-$-almost everywhere with a regular
conditional law of $P$ (for an agent block, \eqref{eq:obs-local-posterior}).
Among all laws with $-\mathcal B$ marginal $R$, the divergence from $P$ has
minimum $\KL(R\Vert P_-)\in[0,+\infty]$, attained by
\begin{equation}
\mathbb Q^{\mathrm{new}}_{V,o,X}(dY_{\mathcal B},dY_{-\mathcal B})
=\kappa_{\mathcal B}(dY_{\mathcal B}\given Y_{-\mathcal B})R(dY_{-\mathcal B}),
\label{eq:elbo-e-coordinate}
\end{equation}
which therefore maximizes $\Lelbo_V^{\mathrm{ext}}$ on that class. The minimizer
is unique when $\KL(R\Vert P_-)<\infty$; otherwise every law in the class has
extended bound $-\infty$. If the minimum is finite and
\begin{equation}
\E_{\mathbb Q^{\mathrm{new}}_{V,o,X}}\bigl[\lvert\log p_{\theta,V}(o,Y\given X)\rvert
+\lvert\log q^{\mathrm{new}}_{V,o,X}(Y)\rvert\bigr]<\infty,
\label{eq:elbo-e-coordinate-h4}
\end{equation}
the new law satisfies (H3)--(H4) and maximizes the classical $\Lelbo_V$ among
the laws with marginal $R$ that satisfy \cref{hyp:elbo-evidence-domain}.
Finite complement relative entropy does not imply \eqref{eq:elbo-e-coordinate-h4}.

(iv) If $\Fenergy^{\mathrm{ext}}_o(Q)<\infty$ for $Q=Q_{-B}r_B$, then for
$\mathrm w\in[0,1]$ the damped kernel
$r_B^{\mathrm w}=(1-\mathrm w)r_B+\mathrm w\boldsymbol\Pi_{\theta,V,o,X,B}$
satisfies
\begin{equation}
\Fenergy^{\mathrm{ext}}_o(Q_{-B}r_B^{\mathrm w})
\leq(1-\mathrm w)\Fenergy^{\mathrm{ext}}_o(Q)
+\mathrm w\Fenergy^{\mathrm{ext}}_o(Q_{-B}\boldsymbol\Pi_{\theta,V,o,X,B})
\leq\Fenergy^{\mathrm{ext}}_o(Q),
\label{eq:obs-damped-block-update}
\end{equation}
so replacing $r_B$ by the local posterior with $Q_{-B}$ fixed does not increase
$\Fenergy^{\mathrm{ext}}_o$. \status{ESTABLISHED}

\noindent\emph{Proof.}
For (i), the conditional model with reference
$\mathbb P^{\mathsf Y}_{\theta,V,B}(\cdot\given y_{-B})$ and likelihood
$L_{B,o}(\cdot;y_{-B})$ has evidence $Z_B(y_{-B})\in(0,\infty)$ and posterior
\eqref{eq:obs-local-posterior}, and \cref{thm:elbo-extended-gap} applies to it.
Part (ii) is \cref{lem:prob-kl-chain-rule} with the disintegrations
\eqref{eq:obs-posterior-block-disintegration} and
\eqref{eq:obs-baseline-block-disintegration}. For (iii), the fixed-marginal
clause of \cref{lem:prob-kl-chain-rule}, with the $-\mathcal B$ coordinates
first, gives the minimum $\KL(R\Vert P_-)$, the minimizer
\eqref{eq:elbo-e-coordinate}, and uniqueness when the minimum is finite; then
$R\ll P_-$, the completion of $\kappa_{\mathcal B}$ off a $P_-$-null set is
irrelevant, and $d\mathbb Q^{\mathrm{new}}_{V,o,X}/dP=(dR/dP_-)(Y_{-\mathcal B})$.
When the minimum is infinite, the data-processing clause
\eqref{eq:cg-kl-dpi-extended} applied to the projection makes every divergence
in the class infinite. \Cref{def:elbo-extended} converts divergence into the
extended bound, and \cref{thm:elbo-exact-identity} gives the classical
statement. The countable witness of \cref{sec:app-integrability-witnesses} has
$R=P_-$ and fails \eqref{eq:elbo-e-coordinate-h4}. For (iv), finiteness makes
$Q_{-B}\ll\boldsymbol\Pi_{\theta,V,o,X,-B}$ by data processing, and
\eqref{eq:obs-posterior-block-kl-chain} leaves the outside term unchanged; the
exact replacement sets the conditional term to zero, and convexity of relative
entropy in its first argument, applied for each $y_{-B}$, gives
\eqref{eq:obs-damped-block-update}. \hfill$\square$

The exact full E-step is case (iii) with $\mathcal B$ equal to all
coordinates; a coordinate that optimizes inside a product family does not have
the form \eqref{eq:elbo-e-coordinate} and leaves the gap of
\cref{prop:restrict-principle}. By (iv) a sequential sweep of exact or damped
block updates does not increase the collective free energy; replacing the
conditionals of several correlated blocks at once is not covered, because each
replacement changes the outside marginal of the others.

\theoremheading{Exact local--global potential identity}{thm:obs-local-global-potential}
Let $Q=Q_{-B}r_B$ and $Q'=Q_{-B}r'_B$ have the same outside marginal with
$Q_{-B}\ll\boldsymbol\Pi_{\theta,V,o,X,-B}$, and suppose both have finite
divergence from $\boldsymbol\Pi_{\theta,V,o,X}$. Then
\begin{equation}
\Fenergy^{\mathrm{ext}}_o(Q')-\Fenergy^{\mathrm{ext}}_o(Q)
=\E_{Q_{-B}}\bigl[\Fenergy^{\mathrm{ext}}_{B,o}(r'_B;Y_{-B})
-\Fenergy^{\mathrm{ext}}_{B,o}(r_B;Y_{-B})\bigr].
\label{eq:obs-local-global-potential}
\end{equation}
When in addition every baseline relative entropy and the positive and negative
parts of the energies are integrable,
\begin{equation}
\Fenergy^{\mathrm{ext}}_o(Q)
=\KL\bigl(Q_{-B}\Vert\mathbb P^{\mathsf Y}_{\theta,V,-B}\bigr)
+\E_{Q_{-B}}\bigl[H_{\bar B,o}(Y_{-B})+\Fenergy^{\mathrm{ext}}_{B,o}(r_B;Y_{-B})\bigr].
\label{eq:obs-local-global-decomposition}
\end{equation}
Every exact local update is thus a coordinate update of the collective free
energy. The fixed outside marginal is load bearing. \status{ESTABLISHED}

\noindent\emph{Proof.}
Apply \eqref{eq:obs-posterior-block-kl-chain} to both laws. The finite
divergences make both conditional terms integrable, the outside terms agree,
the constant $-\log p_{\theta,V}(o\given X)$ cancels, and so does
$-\log Z_B(Y_{-B})$ inside the two local functionals. The decomposition follows
from \eqref{eq:obs-block-kl-chain}, \eqref{eq:obs-local-energy-form}, and
\eqref{eq:obs-collective-energy-form}. \hfill$\square$

\corollaryheading{Exact potential game}{cor:obs-potential-game}
Let the blocks be singletons and choose a nonempty strategy set $\mathcal R_i$
for each agent. Suppose every profile $(r_i)\in\prod_{i\in V}\mathcal R_i$
defines a product recognition law $Q=\bigotimes_{i\in V}r_i$ with
$\KL(Q\Vert\boldsymbol\Pi_{\theta,V,o,X})<\infty$ and
$\E_{Q_{-i}}\lvert\log Z_{\{i\}}(Y_{-i})\rvert<\infty$ for every $i$, where
$Q_{-i}=\bigotimes_{j\ne i}r_j$. The outside-averaged local free energy
\begin{equation}
\overline\Fenergy_i(r_i;r_{-i}):=\E_{Q_{-i}}\Fenergy^{\mathrm{ext}}_{\{i\},o}(r_i;Y_{-i})
\label{eq:obs-outside-averaged-local-vfe}
\end{equation}
is finite, and between two such profiles that differ only in $r_i$ it changes by
exactly the change of $\Fenergy^{\mathrm{ext}}_o$. The players $V$, strategy sets
$\mathcal R_i$, and costs $\overline\Fenergy_i$ therefore form an exact potential game
with potential $\Fenergy^{\mathrm{ext}}_o$ in the sense of
\citet[Sec.~2]{MondererShapley1996}, written for costs instead of payoffs.
\status{ESTABLISHED}

\noindent\emph{Proof.}
Finite divergence gives $Q_{-i}\ll\boldsymbol\Pi_{\theta,V,o,X,-i}$ by data
processing, and \eqref{eq:obs-posterior-block-kl-chain} makes the averaged
conditional divergence finite; with the hypothesis on $\log Z_{\{i\}}$ this
makes \eqref{eq:obs-outside-averaged-local-vfe} finite. A unilateral change
leaves $Q_{-i}$ fixed, and \cref{thm:obs-local-global-potential} with
$B=\{i\}$ gives the equality of changes. \hfill$\square$

\Cref{fig:obs-local-global-potential} shows which records enter a block
potential; block-gradient dynamics are in \cref{sec:local-natural-gradient}.

\begin{figure}[htbp]
\centering
\begin{tikzpicture}[>=stealth,scale=0.85,every node/.style={transform shape}]
  \node[draw,circle,fill=lightblue,minimum size=0.85cm] (a1) at (0,1.5) {$1$};
  \node[draw,circle,fill=lightblue,minimum size=0.85cm] (a2) at (2.1,2.3) {$2$};
  \node[draw,circle,fill=lightgreen,minimum size=0.85cm] (a3) at (4.2,1.5) {$3$};
  \node[draw,circle,fill=lightgreen,minimum size=0.85cm] (a4) at (5.2,-0.2) {$4$};
  \node[draw,circle,fill=lightblue,minimum size=0.85cm] (a5) at (1.2,-0.5) {$5$};
  \node[draw,rectangle,fill=lightorange,minimum size=0.52cm] (f12) at (1.05,1.9) {$o_{a}$};
  \node[draw,rectangle,fill=lightorange,minimum size=0.52cm] (f235) at (2.65,0.45) {$o_{a'}$};
  \node[draw,rectangle,fill=lightorange,minimum size=0.52cm] (f34) at (4.7,0.65) {$o_{a''}$};
  \draw (a1)--(f12)--(a2);
  \draw (a2)--(f235)--(a3);
  \draw (a5)--(f235);
  \draw (a3)--(f34)--(a4);
  \draw[rounded corners,very thick,sciencegreen]
    (3.55,2.15) rectangle (5.75,-0.85);
  \node[text=sciencegreen] at (4.65,2.48) {block $B=\{3,4\}$};
  \draw[->,very thick,primaryblue] (6.05,0.65)--(7.35,0.65);
  \node[align=center,draw,rounded corners,fill=lightblue,text width=3.5cm]
    at (9.25,0.65)
    {$\Delta\Fenergy_o^{\rm ext}
      =\E_{Q_{-B}}\Delta\Fenergy_{B,o}^{\rm ext}$\\
     incident records: $o_{a'},o_{a''}$};
\end{tikzpicture}
\caption{A block-local free energy contains every record incident to the block:
$o_{a'}$ crosses the boundary and enters the block conditional, while $o_a$ lies
outside and cancels.}
\label{fig:obs-local-global-potential}
\end{figure}

\section{Additive accounting}
\label{sec:local-additive-accounting}

The local free energies of \cref{thm:obs-local-multiagent-elbo} are coordinate
potentials of the collective free energy, not shares of it; an additive
single-count ledger is a different construction.

\corollaryheading{Correlated ledger and incident overcounting}{cor:obs-correlated-ledger}
Suppose $\mathbb P^{\mathsf Y}_{\theta,V}=\bigotimes_{i\in V}\varrho_i$ and let
$Q$ be any law on the latent space. Then (i) in $[0,+\infty]$,
\begin{equation}
\KL\bigl(Q\Vert\mathbb P^{\mathsf Y}_{\theta,V}\bigr)
=\operatorname{TC}(Q)+\sum_i\KL(Q_i\Vert\varrho_i);
\label{eq:obs-correlated-complexity-ledger}
\end{equation}
(ii) whenever \eqref{eq:obs-collective-energy-form} holds,
\begin{equation}
\Fenergy^{\mathrm{ext}}_o(Q)
=\operatorname{TC}(Q)+\sum_i\KL(Q_i\Vert\varrho_i)
+\E_Q\Bigl[\sum_{a\in\mathfrak A}E_{a,o}(Y_{\partial a})\Bigr],
\label{eq:obs-global-ledger}
\end{equation}
in which each record appears once and $\operatorname{TC}(Q)=0$ exactly when $Q$
is a product law, so mean field is a special case of the ledger and not a
premise of it; (iii) if every record has a nonempty scope, the singleton
incident energies satisfy, pointwise in $(-\infty,+\infty]$,
\begin{equation}
H_{\{i\},o}(y_i;y_{-i})=\sum_{a:i\in\partial a}E_{a,o}(y_{\partial a}),
\qquad
\sum_{i\in V}H_{\{i\},o}(y_i;y_{-i})=\sum_{a\in\mathfrak A}\lvert\partial a\rvert E_{a,o}(y_{\partial a}),
\label{eq:obs-singleton-incident-counting}
\end{equation}
so the singleton local free energies count each record $\lvert\partial a\rvert$
times and are not the single-count ledger; (iv) this is a statement about these
local free energies, not a bar to additive decompositions: iterating
\cref{lem:prob-kl-chain-rule} along any ordering of $V$ writes
$\KL(Q\Vert\boldsymbol\Pi_{\theta,V,o,X})$ as one term per agent, each
conditioned on the agents before it. \status{ESTABLISHED}

\noindent\emph{Proof.}
Part (i) is \cref{thm:elbo-total-correlation-chain} with
$\varrho_j=\varrho_i$, and (ii) substitutes it into
\eqref{eq:obs-collective-energy-form}. For (iii), every $E_{a,o}$ lies in
$(-\infty,+\infty]$ and both index sets are finite, so the rearrangement is
valid pointwise. Part (iv) is repeated application of the chain rule, peeling
one agent at a time. \hfill$\square$

The same chain rule accounts for link values once they are latent; the
interaction multigraph itself stays structural (\cref{sec:intro-conventions}).

\propositionheading{Collective free energy with random link values}{prop:elbo-random-link-vfe}
Let the link family $\Theta=(\Theta^x_e)$ on the edge copies of $\mathfrak I$ be
a latent variable with values in a standard Borel space $\mathscr T$. Let
$\nu^{\mathscr T}$ be a probability law on $\mathscr T$, let
$\mathscr S:\mathscr T\to(-\infty,+\infty]$ be a measurable action with
$0<Z_{\mathscr S}=\int e^{-\mathscr S}d\nu^{\mathscr T}<\infty$, and let the link
prior be $Z_{\mathscr S}^{-1}e^{-\mathscr S}\nu^{\mathscr T}$. For each $\Theta$
let the record model of \cref{hyp:local-interaction-kernels} with links $\Theta$
have evidence $p(o\given\Theta)\in(0,\infty)$, measurable in $\Theta$, posterior
$\boldsymbol\Pi_\Theta$, and free energy
$\Fenergy_{\mathrm{matter}}(Q_\Theta;o,\Theta)=-\log p(o\given\Theta)+\KL(Q_\Theta\Vert\boldsymbol\Pi_\Theta)$,
and assume the joint evidence
$p(o)=\int p(o\given\Theta)Z_{\mathscr S}^{-1}e^{-\mathscr S(\Theta)}\nu^{\mathscr T}(d\Theta)$
is finite. For a recognition law $Q(d\Theta,dy)=R(d\Theta)Q_\Theta(dy)$ with
$\KL(R\Vert\nu^{\mathscr T})<\infty$ and $\E_R\lvert\mathscr S\rvert<\infty$, the
extended free energy of the joint model is
\begin{equation}
\Fenergy=\E_R\Fenergy_{\mathrm{matter}}(Q_\Theta;o,\Theta)+\E_R\mathscr S
+\KL(R\Vert\nu^{\mathscr T})+\log Z_{\mathscr S}
\quad\text{in }(-\infty,+\infty].
\label{eq:rg-geometry-random-link-vfe}
\end{equation}
Omitting the link entropy term or the normalizer changes the free energy; when
the links are structural parameters instead of latent variables, only
$\Fenergy_{\mathrm{matter}}$ appears. \status{ESTABLISHED}

\noindent\emph{Proof.}
The posterior of the joint model has conditional law $\boldsymbol\Pi_\Theta$
given $\Theta$ and link marginal
$\boldsymbol\Pi^{\mathscr T}(d\Theta)
=p(o)^{-1}p(o\given\Theta)Z_{\mathscr S}^{-1}e^{-\mathscr S(\Theta)}\nu^{\mathscr T}(d\Theta)$,
so the chain rule
\eqref{eq:elbo-kl-chain-rule} gives
$\KL(Q\Vert\boldsymbol\Pi)=\KL(R\Vert\boldsymbol\Pi^{\mathscr T})+\E_R\KL(Q_\Theta\Vert\boldsymbol\Pi_\Theta)$
in $[0,+\infty]$, with a measurable integrand. Since
$\E_R\lvert\mathscr S\rvert<\infty$, $R$ charges no set where
$\mathscr S=+\infty$, so $R\ll\boldsymbol\Pi^{\mathscr T}$ and
$\log(dR/d\boldsymbol\Pi^{\mathscr T})=\log(dR/d\nu^{\mathscr T})+\mathscr S+\log Z_{\mathscr S}
-\log p(o\given\Theta)+\log p(o)$. The first two terms are $R$-integrable, and so
is the positive part $f$ of $\log p(o\given\cdot)$: for $f_n=\min(f,n)$, the
nonnegative divergence from the law proportional to $e^{f_n}$ times the prior
gives $\E_Rf_n\leq\KL(R\Vert Z_{\mathscr S}^{-1}e^{-\mathscr S}\nu^{\mathscr T})+\log(1+p(o))$,
and that divergence equals
$\KL(R\Vert\nu^{\mathscr T})+\E_R\mathscr S+\log Z_{\mathscr S}$. Integrating
term by term and adding the conditional term, whose integrand is nonnegative,
turns $\Fenergy=-\log p(o)+\KL(Q\Vert\boldsymbol\Pi)$ into
\eqref{eq:rg-geometry-random-link-vfe}; every summand has an integrable negative
part, so no difference of infinities occurs. \hfill$\square$

\Cref{ch:links} supplies, for every Lie group $G$, data that satisfy the hypotheses of this proposition:
the latent variable is the family of Lie-algebra generators of the links, the base law is Gaussian and
the action a nonnegative combination of loop defects, the joint has unit mass on every record hypergraph
(\Cref{prop:link-latent-normalization}), and for a Gaussian recognition law of the generators the link
sector of
\eqref{eq:rg-geometry-random-link-vfe} is the relative entropy from the generator prior
(\Cref{prop:link-gaussian-recognition-vfe}).

\section{Model coordinates and accepted updates}
\label{sec:elbo-coordinates}

In the two-argument form $\Fenergy_V(\mathbb Q;\theta)$ of
\cref{req:gen-typing-prohibition}, block coordinates move $\mathbb Q$ at fixed
$\theta$ and a model coordinate moves $\theta$ at fixed $\mathbb Q$; each
accepted value defines a new fixed joint.

\propositionheading{Exact M-coordinate}{prop:elbo-exact-m-coordinate}
Hold $\mathbb Q_{V,o,X}$ fixed, and suppose the parameterized family supplies a
common regular-observation set or jointly measurable pointwise versions for the
candidates compared, with
$\E_{\mathbb Q_{V,o,X}}\lvert\log q_{V,o,X}(Y)\rvert<\infty$ and
$\E_{\mathbb Q_{V,o,X}}\lvert\log p_{\theta',V}(o,Y\given X)\rvert<\infty$
for every candidate $\theta'$. If a maximizer exists in the declared parameter
space, then every
\begin{equation}
\theta^{\mathrm{new}}\in\operatorname*{argmax}_{\theta'}
\E_{\mathbb Q_{V,o,X}}\bigl[\log p_{\theta',V}(o,Y\given X)\bigr]
\label{eq:elbo-m-coordinate}
\end{equation}
maximizes $\Lelbo_{V,\theta'}(\mathbb Q_{V,o,X};X,o)$ over the same space. The
finiteness hypotheses cannot be dropped. \status{ESTABLISHED}

\noindent\emph{Proof.}
The recognition term $\E_{\mathbb Q_{V,o,X}}[\log q_{V,o,X}]$ is finite and does
not depend on $\theta'$. For the last sentence, on $\{2,3,\ldots\}$ let
$q_n=\mathrm c/(n(\log n)^2)$ with $\mathrm c$ normalizing,
$h_n=\mathbf 1_{\{n=2\}}-q_2$, and $r_n=1+\mathrm c_2h_n$ with
$0<\mathrm c_2<1/q_2$, so that $r$ is bounded, strictly positive, nonconstant,
and $\sum_nq_nr_n=1$. With a trivial record, let the two candidates have latent
laws $p_0=q$ and $p_1=qr$, and let the recognition law be $q$. The entropy of
$q$ is infinite (\cref{sec:app-integrability-witnesses}), so
$\E_q\log p_0=\E_q\log p_1=-\infty$ and both candidates maximize the expected
complete log-likelihood, whereas their extended bounds are
$\E_q\log(p_0/q)=0$ and $\E_q\log r=-\KL(q\Vert p_1)<0$. \hfill$\square$

\definitionheading{Generalized M-coordinate acceptance}{def:elbo-generalized-m-acceptance}
A generalized M-coordinate need not attain a maximum, but it must pass the
acceptance condition
\begin{equation}
\Lelbo^{\mathrm{ext}}_{V,\theta^{\mathrm{new}}}(\mathbb Q_{V,o,X};X,o)
\geq\Lelbo^{\mathrm{ext}}_{V,\theta^{\mathrm{old}}}(\mathbb Q_{V,o,X};X,o),
\label{eq:elbo-acceptance}
\end{equation}
evaluated at an observation in the regular sets of both candidates, or with
jointly measurable versions declared as part of the parameterized model. No
recognition density and no (H4) is required; on the classical domain the
condition is the same inequality for \eqref{eq:elbo-definition}, and
\eqref{eq:elbo-m-coordinate} is a sufficient condition.

\corollaryheading{Evidence-free acceptance test}{cor:elbo-acceptance-evidence-free}
Suppose $\mathbb Q_{V,o,X}\ll M_{\theta^{\mathrm{old}},V,o,X}$, the two slice
measures \eqref{eq:elbo-slice-measure} are mutually absolutely continuous, and
both extended values are finite. Then
\begin{equation}
\Lelbo^{\mathrm{ext}}_{V,\theta^{\mathrm{new}}}(\mathbb Q_{V,o,X};X,o)
-\Lelbo^{\mathrm{ext}}_{V,\theta^{\mathrm{old}}}(\mathbb Q_{V,o,X};X,o)
=\E_{\mathbb Q_{V,o,X}}\biggl[\log\frac{p_{\theta^{\mathrm{new}},V}(o,Y\given X)}{p_{\theta^{\mathrm{old}},V}(o,Y\given X)}\biggr],
\label{eq:elbo-acceptance-evidence-free}
\end{equation}
so the test \eqref{eq:elbo-acceptance} needs no evidence value.
\status{ESTABLISHED}

\noindent\emph{Proof.}
With $r_x=d\mathbb Q_{V,o,X}/dM_{\theta^x,V,o,X}$ for
$x\in\{\mathrm{old},\mathrm{new}\}$, the Radon--Nikodym chain rule gives
$\log r_{\mathrm{old}}-\log r_{\mathrm{new}}
=\log(dM_{\theta^{\mathrm{new}},V,o,X}/dM_{\theta^{\mathrm{old}},V,o,X})$ almost
surely, and this derivative is the ratio of the two joint densities at $o$. By
\cref{prop:elbo-relative-log-representation} the negative part of each
$\log r_x$ is integrable, so finiteness makes each $\log r_x$ integrable, and
subtracting the two representations \eqref{eq:elbo-relative-log} proves the
identity. \hfill$\square$

\propositionheading{Evidence monotonicity}{prop:elbo-evidence-monotonicity}
Suppose (H1)--(H2) hold for the fixed $\mathbb Q_{V,o,X}$ at both
$\theta^{\mathrm{old}}$ and $\theta^{\mathrm{new}}$. If the E-coordinate is exact
at the old parameter, $\mathbb Q_{V,o,X}=\boldsymbol\Pi_{\theta^{\mathrm{old}},V,o,X}$,
and the M-coordinate satisfies \eqref{eq:elbo-acceptance}, then
\begin{equation}
\log p_{\theta^{\mathrm{new}},V}(o\given X)
\geq\Lelbo^{\mathrm{ext}}_{V,\theta^{\mathrm{new}}}(\mathbb Q_{V,o,X};X,o)
\geq\Lelbo^{\mathrm{ext}}_{V,\theta^{\mathrm{old}}}(\mathbb Q_{V,o,X};X,o)
=\log p_{\theta^{\mathrm{old}},V}(o\given X).
\label{eq:elbo-monotonicity}
\end{equation}
Consequently a schedule that places a full exact E phase before each accepted
M phase has nondecreasing evidence values, which converge in
$(-\infty,+\infty]$ and to a finite limit when bounded above. Block
E-coordinates and accepted M-coordinates make the compared bound values
nondecreasing on a common functional; the transfer to the evidence in
\eqref{eq:elbo-monotonicity} uses the full E phase. \status{ESTABLISHED}

\noindent\emph{Proof.}
The first inequality is \cref{thm:elbo-extended-gap} at $\theta^{\mathrm{new}}$,
the second is \eqref{eq:elbo-acceptance}, and the equality is the zero-gap case
at $\theta^{\mathrm{old}}$; the singular-branch sum $\Lelbo^{\mathrm{ext}}+\KL$
is never written. A nondecreasing sequence of extended reals converges. Each
block coordinate does not decrease the bound by
\cref{prop:elbo-exact-e-coordinate}, and each accepted step does not decrease it
by \eqref{eq:elbo-acceptance}; only the final equality of
\eqref{eq:elbo-monotonicity} uses the exact E phase. Under (H3)--(H4) at both
candidates this is the classical conclusion of \citet{Dempster1977} in the
reading of \citet{Neal1998}. \hfill$\square$

Convergence of the parameter sequence needs regularity conditions of the kind
identified for the classical algorithm by \citet{Wu1983}. A finite step along a
natural-gradient or other optimizer direction is a proposal, not an accepted
coordinate, and it can decrease the bound.

\propositionheading{Finite natural-gradient steps need not be monotone}{prop:elbo-finite-step-nonmonotonicity}
Let $\mathsf G\succ0$, write $\nabla\Lelbo_V$ for the gradient of
$\theta'\mapsto\Lelbo_{V,\theta'}(\mathbb Q_{V,o,X};X,o)$, and consider the step
$\theta^+=\theta'+\mathrm h\mathsf G^{-1}\nabla\Lelbo_V$ of length
$\mathrm h>0$. At a nonstationary point the direction is an ascent direction,
\begin{equation}
\frac{d}{d\mathrm h}\Lelbo_{V,\theta^+}\Big|_{\mathrm h=0}
=\nabla\Lelbo_V^{\top}\mathsf G^{-1}\nabla\Lelbo_V>0.
\label{eq:elbo-ascent-direction}
\end{equation}
If $\Lelbo_{V,\theta''}=\Lelbo^\star-\tfrac12(\theta''-\theta^\star)^\top\mathsf H(\theta''-\theta^\star)$
with $\mathsf H\succ0$ for $\theta''$ in a neighborhood $\mathcal U$, $d_{\max}$
is the largest generalized eigenvalue of $(\mathsf H,\mathsf G)$, and
$\mathbf z$ is a corresponding eigenvector with
$\mathbf z^\top\mathsf G\mathbf z=1$, then at every
$\theta'=\theta^\star+\mathrm c_1\mathbf z\in\mathcal U$ with $\mathrm c_1\ne0$
and every step whose endpoint lies in $\mathcal U$,
\begin{equation}
\Lelbo_{V,\theta^+}-\Lelbo_{V,\theta'}
=\tfrac12d_{\max}\mathrm c_1^2\bigl[1-(1-\mathrm hd_{\max})^2\bigr]<0
\qquad\text{whenever}\qquad
\mathrm h>\frac{2}{d_{\max}}.
\label{eq:elbo-step-size-failure}
\end{equation}
\status{ESTABLISHED}

\noindent\emph{Proof.}
The first display is the chain rule. With
$\Delta\theta=\theta'-\theta^\star=\mathrm c_1\mathbf z$,
$\nabla\Lelbo_V=-\mathsf H\Delta\theta$ and
$\mathsf H\mathbf z=d_{\max}\mathsf G\mathbf z$ give
$\theta^+-\theta^\star=(1-\mathrm hd_{\max})\mathrm c_1\mathbf z$; substituting
in the quadratic gives the second display, whose bracket is negative exactly
when $\mathrm hd_{\max}>2$. \hfill$\square$

Monotonicity at a finite step therefore needs a line search, a trust region,
or another rule that verifies \eqref{eq:elbo-acceptance} on the same
functional; a better preconditioner does not remove the requirement, because
\eqref{eq:elbo-step-size-failure} is stated in the preconditioned spectrum.

\section{What remains open}
\label{sec:elbo-open}

\openproblemheading{Convergence of the alternating scheme}{open:elbo-alternating-convergence}
\Cref{prop:elbo-evidence-monotonicity} gives monotone convergence of the
evidence values under a full exact E phase and nothing about the parameters,
and since \eqref{eq:elbo-acceptance} is a weak inequality the null step is an
accepted M-coordinate, so the general scheme of block E-coordinates
\eqref{eq:elbo-e-coordinate} and accepted M-coordinates need not move its
parameter at all. The problem concerns models of
\cref{hyp:local-interaction-kernels} with finitely many latent states,
parameters in an open cone of a Euclidean space, and joint densities twice
continuously differentiable in the parameters, under the scheme that
alternates a full exact E phase with the exact M-coordinate
\eqref{eq:elbo-m-coordinate} whose maximizer is unique. The regularity
conditions of \citet{Wu1983} concern the limit points of the parameter
sequence and its convergence when the stationary points at its limit value
are isolated; isolation is not among the hypotheses of the class, in which
uniqueness of the maximizer leaves open that the evidence is constant along
a curve of parameters on which the complete-data likelihood is not. It is
settled by conditions checkable on the model, including compactness of the
upper level sets of the evidence in the cone, that are met by some model of
the class with parameter dimension at least two whose stationary set at the
limit value of an orbit is not discrete, under which every orbit of the
scheme converges to one stationary point of the extended bound, or by a
model of that class with compact upper level sets whose orbit has more than
one limit point.

\section*{Physicist's summary}

The population carries one generative law and one recognition law, and the
only evidence bound is $\log p(o)-\KL(Q\Vert\boldsymbol\Pi)$, defined for every
$Q$ and equal to $-\infty$ when $Q$ charges what the posterior cannot see;
everything else is bookkeeping of that relative entropy through the chain rule.
Substituting marginal entropies overstates the bound by exactly the total
correlation. Observations are records attached to latent states, and an
environment can be rewritten as boundary nodes emitting them. Each block of
agents has a local free energy built from the records that touch it; with the
rest held fixed, lowering it lowers the collective free energy by the same
amount, so the population is an exact potential game. Summed, the local free
energies count every shared record once per participant. Parameters move only
between evaluations, by accepted steps, and the evidence rises monotonically
when each step follows an exact posterior refresh; convergence is open.

\chapter{Attention and Selective Communication}
\label{ch:selective-communication}

Which agent does a receiver listen to, and what does it hear? Attention enters
the joint of \cref{ch:generative} as a latent source label: the softmax row is
its generative posterior at a latent state, and the optimal recognition row
held constant in the latent state is a different softmax, over averaged
energies (\cref{prop:obs-attention-elbo}), at the exact cost of
\cref{prop:obs-attention-constant-row-gap}. Communication enters as normalized
encoder, transmission, and decoder kernels, and the exact receiver update of a
comparison record fixes the value map by the query and key maps
(\cref{prop:obs-selective-receiver-update}). The row calculus and the reduced
curvature of \cref{prop:obs-attention-reduced-hessian} feed
\cref{ch:two-timescales}, the event law of \cref{def:obs-attention-event-law}
is the attention state that \cref{ch:agent-network-rg} coarse-grains, and the
comparison records supply the endpoint energies of \cref{ch:network-state}.
Rows, events, links, and divergences are indexed by oriented edge copies
\(e:j\to i\); the pair \(ij\) names a unique copy, a pair quantity is the sum
over parallel copies, and the construction serves both channels, with a
model-channel row written \(\beta^m\) where the channel must be named.
Longer proofs are deferred to \cref{app:deferred-proofs}.

This chapter does not claim that transformer attention as a whole is derived
(\cref{tab:obs-transformer-correspondence}; feed-forward blocks and layer
normalization have no counterpart), that declared per-head rows follow from a
label model, or that an interface map, its rank, or its coding establishes a
channel capacity, a reconstruction of the sender's belief, or shared meaning.
Capacity needs an admissible input family and, for a nonzero continuous
Gaussian channel, an input resource constraint; a coding theorem also needs
repeated uses, errors, and rates \citep{Shannon1948}; a task-dependent
distortion or sufficient statistic is further model data. Open are the
label-level multi-head model (\cref{open:obs-multihead-attention}), discrete
creation and deletion of interaction edges (\cref{open:obs-edge-creation}), and
formation with evolving selective interfaces
(\cref{open:obs-selective-formation}).
\status{NOT-CLAIMED}

\section{Attention as a latent source label}
\label{sec:local-attention-elbo}

\definitionheading{Per-receiver source-label model}{def:obs-attention-label-model}
Let \(\mathcal I_{\rm att}\) be a finite set of receiving rows, each an agent
\(i\in V\) or an agent--head pair \((i,h)\) that inherits its agent's latent
state (written \(i\) when no head is needed). Row \(i\) carries a source label
\(J_i\) valued in the incoming edge copies of \(i\), with a structural prior
row \(\pi_i\) independent of recognition whose support is the source set
\(\mathcal N_i\); source sums and indices range over \(\mathcal N_i\). By
\cref{sec:intro-conventions} the interaction multigraph at \(c_*\) is the union
of these supports, complete unless a mask sets a prior weight to zero, and
whether \(\mathcal N_i\) contains the self edge is part of the model. The
labels augment the latent baseline independently,
\begin{equation}
\mathbb P_{\theta,V}^{Y,\mathrm{aug}}(dy,dj)
=\mathbb P_{\theta,V}^{Y}(dy)
\prod_{i\in\mathcal I_{\rm att}}\pi_i(dj_i).
\label{eq:obs-attention-augmented-baseline}
\end{equation}
Each row has a temperature \(\tau_i>0\) and finite measurable energies
\(D_{ij}:\mathsf Y\to\R\) that contain no recognition law and are
gauge-invariant functions of sample-level latent states
(\cref{def:prob-description-levels}) and of the links, and at the selected
record row \(i\) has the record density
\begin{equation}
L_i(o_i\given y,J_i=j)
=\mathrm c_i(o_i,y)\exp[-D_{ij}(y)/\tau_i],
\label{eq:obs-attention-likelihood}
\end{equation}
with \(\mathrm c_i\) independent of \(j\); the kernel off the selected slice is
part of the model. When \(\mathrm c_i(o_i,\cdot)\) depends only on \(y_i\) and
each \(D_{ij}\) only on \((y_i,y_j)\), the record is one hyperedge factor of
\cref{hyp:local-interaction-kernels} with scope \(\{i\}\cup\mathcal N_i\), and
through \eqref{eq:obs-block-incident-likelihood} it enters the local free
energy of the receiver and of every candidate source. Each row carries exactly
one label, so receiver occupancy is fixed; inferring it needs a null label in
the source alphabet or a different joint label model with its own likelihood,
such as \cref{hyp:rg-single-edge-attention}.

The slice form \eqref{eq:obs-attention-likelihood} is inhabited: a binary
success record with success probability \(\exp[-D_{ij}/\tau_i]\) and every
\(D_{ij}\ge0\) realizes it with \(\mathrm c_i=1\), as in
\eqref{eq:obs-interaction-energy}, and so does the comparison record of
\cref{prop:obs-selective-receiver-update} at every record value.

\propositionheading{Exact posterior and recognition attention}{prop:obs-attention-elbo}
Assume \cref{def:obs-attention-label-model} and label exclusivity at the
complete selected record \(o\): the augmented likelihood factors as
\begin{equation}
L_o^{\rm aug}(y,j)
=L_o^Y(y)\prod_{i\in\mathcal I_{\rm att}}
\mathrm c_i(o_i,y)\exp[-D_{ij_i}(y)/\tau_i],
\label{eq:obs-attention-augmented-likelihood}
\end{equation}
with \(L_o^Y:\mathsf Y\to[0,+\infty)\) finite, measurable, and free of every
label, and no other record or factor of the generative law reads any \(J_i\).
(a) On \(\{L_o^Y\prod_k\mathrm c_k(o_k,\cdot)>0\}\) the generative posterior
row given \(y\) and the record is
\begin{equation}
\beta^P_{ij}(y)
=\pi_{ij}e^{-D_{ij}(y)/\tau_i}\Big/\sum_k\pi_{ik}e^{-D_{ik}(y)/\tau_i},
\label{eq:obs-attention-posterior}
\end{equation}
a version that may also be chosen off that set; it is strictly positive
exactly on \(\mathcal N_i\), and if row \(i\) has nonnegative energies, an
included self edge with \(D_{ii}=0\), and equal prior weights, no source has a
larger posterior weight than the self edge. (b) Restrict the augmented
recognition law to the full factorization
\begin{equation}
\mathbb Q(dy,dj)=\mathbb Q_Y(dy)
\prod_{i\in\mathcal I_{\rm att}}\beta_i^Q(dj_i),
\qquad \beta_i^Q(dj_i)\ \text{independent of }y,
\label{eq:obs-attention-recognition-factorization}
\end{equation}
which makes the labels conditionally independent and every row independent of
the latent state, and assume \(\E_{\mathbb Q_Y}|D_{ij}|<+\infty\) for every
source and \(\E_{\mathbb Q_Y}|\log\mathrm c_i(o_i,Y)|<+\infty\). Then row \(i\)
contributes \eqref{eq:obs-attention-full-contribution} to the collective VFE,
whose first term does not depend on \(\beta_i^Q\), with
\eqref{eq:obs-attention-vfe}:
\begin{align}
&-\E_{\mathbb Q_Y}\log\mathrm c_i(o_i,Y)+\Fenergy_i^{\rm att}(\beta_i^Q),
\label{eq:obs-attention-full-contribution}\\
&\Fenergy_i^{\rm att}(\beta^Q_i)
=\KL(\beta^Q_i\Vert\pi_i)
+\frac{1}{\tau_i}\sum_j\beta^Q_{ij}\E_{\mathbb Q_Y}D_{ij}.
\label{eq:obs-attention-vfe}
\end{align}
Its unique minimizer on the open simplex over \(\mathcal N_i\) is
\begin{equation}
\beta^{Q\star}_{ij}
=\pi_{ij}e^{-\E_{\mathbb Q_Y}D_{ij}/\tau_i}\Big/\sum_k\pi_{ik}e^{-\E_{\mathbb Q_Y}D_{ik}/\tau_i},
\label{eq:obs-attention-recognition-optimum}
\end{equation}
and the row functional
\(\sum_j\beta_{ij}\E_{\mathbb Q_Y}D_{ij}+\tau_i\KL(\beta_i\Vert\pi_i)\) equals
\(\tau_i\Fenergy_i^{\rm att}\) with the same minimizer, while the collective VFE
contains \(\Fenergy_i^{\rm att}\) itself. (c) For an arbitrary recognition law
\(\mathbb Q_Y(dy)\mathbb Q_{J\mid Y}(dj\given y)\) with row marginals
\(\mathbb Q_{J_i\mid Y}\), and with \(-\E_{\mathbb Q_Y}\log L_o^Y(Y)\) counted
in the state sector, the attention-label contribution is
\begin{equation}
\begin{split}
&-\sum_{i\in\mathcal I_{\rm att}}\E_{\mathbb Q_Y}\log\mathrm c_i(o_i,Y)
+\E_{\mathbb Q_Y}\operatorname{TC}(\mathbb Q_{J\mid Y})\\
&\quad+\sum_{i\in\mathcal I_{\rm att}}\E_{\mathbb Q_Y}\left[
\KL(\mathbb Q_{J_i\mid Y}\Vert\pi_i)
+\frac{1}{\tau_i}
\E_{\mathbb Q_{J_i\mid Y}}D_{iJ_i}(Y)
\right],
\end{split}
\label{eq:obs-attention-correlated-ledger}
\end{equation}
where
\(\operatorname{TC}(\mathbb Q_{J\mid Y=y})=
\KL\bigl(\mathbb Q_{J\mid Y=y}\Vert\prod_i\mathbb Q_{J_i\mid Y=y}\bigr)\). Its
minimum over unrestricted conditional laws is attained by the product of the
posterior rows \eqref{eq:obs-attention-posterior}, and row marginals
independent of \(y\) do not by themselves give \eqref{eq:obs-attention-vfe},
since a correlated conditional law keeps the nonnegative total-correlation
term.
\status{ESTABLISHED}

\propositionheading{Cost of the constant-row restriction}{prop:obs-attention-constant-row-gap}
Assume the label exclusivity of \cref{prop:obs-attention-elbo}, fix
\(\mathbb Q_Y\) with \(\E_{\mathbb Q_Y}\lvert D_{ij}\rvert<+\infty\) for every
row \(i\) and \(j\in\mathcal N_i\), write \(d_{ij}=\E_{\mathbb Q_Y}D_{ij}\), and
define
\begin{equation}
R_i^Q:=-\log\sum_j\pi_{ij}e^{-d_{ij}/\tau_i},
\qquad
R_i^P(y):=-\log\sum_j\pi_{ij}e^{-D_{ij}(y)/\tau_i}.
\label{eq:obs-attention-row-potentials}
\end{equation}
Over constant rows \eqref{eq:obs-attention-recognition-factorization} the
minimum of \(\Fenergy_i^{\rm att}\) is \(R_i^Q\), attained at
\eqref{eq:obs-attention-recognition-optimum}; over all conditional label laws
\(\mathbb Q_{J\mid Y}\) the minimum of the label terms of
\eqref{eq:obs-attention-correlated-ledger} other than
\(-\sum_i\E_{\mathbb Q_Y}\log\mathrm c_i(o_i,Y)\) is
\(\sum_i\E_{\mathbb Q_Y}R_i^P(Y)\), attained by the product of the posterior
rows \eqref{eq:obs-attention-posterior}. For each row the difference is exactly
\begin{equation}
R_i^Q-\E_{\mathbb Q_Y}R_i^P(Y)
=\E_{\mathbb Q_Y}\KL\bigl(\beta_i^{Q\star}\Vert\beta_i^P(Y)\bigr)\geq0,
\label{eq:obs-attention-jensen-gap}
\end{equation}
which vanishes exactly when every difference \(D_{ij}(Y)-D_{ik}(Y)\) is
\(\mathbb Q_Y\)-almost surely constant. If further
\(D_{ij}=\widetilde D_{ij}\circ\psi_i\) for one measurable map
\(\psi_i:\mathsf Y\to\mathsf Z_i\) into a finite-dimensional real vector space
with finite \(\mathbb Q_Y\)-mean \(\bar\psi_i\), and every
\(\widetilde D_{ij}-\widetilde D_{ik}\) is affine, the recognition row is the
posterior row at the recognition mean,
\begin{equation}
\beta^{Q\star}_{ij}
=\pi_{ij}e^{-\widetilde D_{ij}(\bar\psi_i)/\tau_i}\Big/\sum_k\pi_{ik}e^{-\widetilde D_{ik}(\bar\psi_i)/\tau_i}.
\label{eq:obs-attention-affine-row}
\end{equation}
In a linear sample fiber, the energies
\(\frac12(y_i-\mathring y_{ij})^\top W_i(y_i-\mathring y_{ij})\) with
\(W_i\succ0\) and structural templates \(\mathring y_{ij}\) have affine
differences. If instead \(D_{ij}(y)=\frac12 r_{ij}^\top W_ir_{ij}\) with
\(r_{ij}=y_i-\rho_b(\Theta^b_{ij})y_j\) for a unique edge copy, and \(Y_i\) and
\(Y_j\) are uncorrelated under \(\mathbb Q_Y\) with finite second moments, then
\(\beta_i^{Q\star}\) is the posterior row at the recognition mean with the
logit of source \(j\) lowered by
\(\Tr(W_i\rho_b(\Theta^b_{ij})\operatorname{Cov}(Y_j)
\rho_b(\Theta^b_{ij})^\top)/(2\tau_i)\).
\status{ESTABLISHED}

The proofs of \cref{prop:obs-attention-elbo,prop:obs-attention-constant-row-gap} are in \cref{sec:proofs-attention}.
Coarse-graining a row, a conditional law of a source given a receiver, needs a
receiver law, which the next two definitions supply.

\definitionheading{Receiver occupancy and typed event law}{def:obs-attention-event-law}
For a role mark \(x\in\{P,\Pi,Q\}\), a receiver occupancy \(\alpha^x_i\), a
probability on the rows, and a row \(\beta^x_i\) on \(\mathcal N_i\) give the
marked edge-event law
\begin{equation}
\eta^x_{ij}=\alpha^x_i\beta^x_{ij},
\qquad
\sum_{i,j}\eta^x_{ij}=1.
\label{eq:rg-attention-event-law}
\end{equation}
The mark names the law: \(P\) is the generative joint conditioned on the latent
state and the record, so \(\eta^P\), \(\alpha^P\), and \(\beta^P\) are
functions of \(y\); \(\Pi\) is the posterior slice, with
\(\eta^\Pi_{ij}=\E_{\boldsymbol\Pi}\eta^P_{ij}(Y)\); and \(Q\) is the selected
recognition coupling. The event law carries no head index. The occupancy is
not supplied by \cref{def:obs-attention-label-model}, in which every row
carries its own label, and it enters in one of two typed ways, each use naming
which: a structural receiver-selection law, fixed with the coarse channel
before observation and independent of recognition, or the receiver marginal of
a single joint edge label as in \cref{hyp:rg-single-edge-attention}, with the
label prior as its prior and its posterior from Bayes' formula. On
\(\{\alpha^x_i>0\}\) the row is recovered by disintegration,
\(\beta^x_{ij}=\eta^x_{ij}/\alpha^x_i\); a conditional row alone is not
recursive state.

\definitionheading{Single-edge joint-label attention model}{hyp:rg-single-edge-attention}
This generative model differs from \cref{def:obs-attention-label-model}: the
baseline is augmented by one joint edge label \((i,j)\), a single
receiver--source pair for the selected record, with recognition-independent
prior
\begin{equation}
\mathbb P_{\theta,V}^{\rm att}(dy,di,dj)=a_i\pi_{ij}\mathbb P_{\theta,V}^Y(dy),
\qquad
\sum_i a_i=1,
\quad
\sum_j\pi_{ij}=1,
\label{eq:rg-attention-fixed-label-joint}
\end{equation}
and the augmented likelihood at the complete selected record \(o\) factors as
\begin{equation}
L_o^{\rm edge}(y,i,j)
=L_o^Y(y)\widetilde{\mathrm c}_i(o,y)\exp[-D_{ij}(y)/\tau_i],
\qquad\tau_i>0,
\label{eq:rg-attention-single-edge-likelihood}
\end{equation}
with \(L_o^Y\) and \(\widetilde{\mathrm c}_i\) finite, nonnegative, and
measurable, \(L_o^Y\) free of \((i,j)\), \(\widetilde{\mathrm c}_i\) free of
\(j\), each \(D_{ij}\) finite and measurable, and no other record or factor of
the generative law reading the label; this exclusivity excludes every model in
which several receivers attend at once. In general \(\widetilde{\mathrm c}_i\)
is not the \(\mathrm c_i\) of \eqref{eq:obs-attention-likelihood}: if the
selected receiver's record has slice density
\(\mathrm c_i(o_i,y)\exp[-D_{ij}(y)/\tau_i]\) while every other receiver's
record keeps a label-free kernel \(g_k(o_k\given y)\), then where every
\(g_k(o_k\given y)>0\) the factorization holds with
\(L_o^Y=\prod_kg_k(o_k\given y)\) and
\(\widetilde{\mathrm c}_i=\mathrm c_i/g_i(o_i\given y)\), a likelihood ratio
against the non-event kernel. The evidence-weighted occupancy of this model is
\cref{prop:rg-attention-single-edge-weights}; the per-receiver model collapses
with its declared selection law instead
(\cref{prop:rg-attention-per-receiver-collapse}).

\corollaryheading{Self edges and blocking}{cor:obs-attention-self-edge}
Let \(N\ge2\) agents carry one row each with uniform receiver occupancy
\(\alpha_i=1/N\), partition them into blocks of equal size \(n\) dividing
\(N\), and push an event law with rows independent of the latent state along
the partition by block sums, \(\eta^{\rm c}_{IJ}=\sum_{i\in I,j\in J}\eta_{ij}\),
with \(\alpha^{\rm c}_I=\sum_J\eta^{\rm c}_{IJ}\) and
\(\beta^{\rm c}_{IJ}=\eta^{\rm c}_{IJ}/\alpha^{\rm c}_I\). If every fine row is
uniform over \(j\ne i\), the block event masses are \(n^2/(N(N-1))\) from \(I\)
to \(J\ne I\) and \(n(n-1)/(N(N-1))\) from \(I\) to itself, the block occupancy
is \(n/N\), and
\begin{equation}
\beta^{\rm c}_{II}=\frac{n-1}{N-1},
\qquad
\beta^{\rm c}_{IJ}=\frac{n}{N-1}\quad(J\ne I).
\label{eq:obs-attention-self-edge-coarse-rows}
\end{equation}
If every fine row is uniform over all \(N\) sources, including the self edge,
every coarse row is uniform with weight \(n/N\). A fine no-self convention is
therefore not preserved: the coarse source set of \(I\) contains \(I\).
\status{ESTABLISHED}

\noindent\emph{Proof.}
Without self edges each fine event has mass \(1/(N(N-1))\), a block pair
\(I\ne J\) contains \(n^2\) of them, \(I\times I\) contains \(n(n-1)\), and the
block occupancy is \(n/N\); with self edges each event has mass \(1/N^2\) and
every block pair contains \(n^2\) events. Dividing gives both cases.
\(\square\)

\corollaryheading{Row flows, dissipation, and parameter paths}{cor:obs-attention-replicator}
Hold \(\mathbb Q_Y\), and hence every \(\E_{\mathbb Q_Y}D_{ij}\), fixed. With
row relaxation rate \(\mathsf r_i\geq0\), the Fisher natural gradient of
\eqref{eq:obs-attention-vfe} on the open simplex over \(\mathcal N_i\) is the
replicator equation \eqref{eq:obs-attention-replicator}, with partial
row-free-energy rate \eqref{eq:obs-attention-dissipation}:
\begin{align}
\dot\beta^Q_{ij}
&=-\mathsf r_i\beta^Q_{ij}
\left(\mathsf c_{ij}-\sum_k\beta^Q_{ik}\mathsf c_{ik}\right),
\qquad
\mathsf c_{ij}=\log\frac{\beta^Q_{ij}}{\pi_{ij}}
+\frac{\E_{\mathbb Q_Y}D_{ij}}{\tau_i},
\label{eq:obs-attention-replicator}\\
\left.\frac{d}{dt}\Fenergy_i^{\rm att}\right|_{\beta_i}
&=-\mathsf r_i\operatorname{Var}_{\beta^Q_i}(\mathsf c_{ij})\leq0.
\label{eq:obs-attention-dissipation}
\end{align}
If instead \(\mathbb Q_{Y,\lambda(t)}\) moves along a \(C^1\) curve \(\lambda\) in
an open coordinate domain, with every
\(d_{ij}(\lambda)=\E_{\mathbb Q_{Y,\lambda}}D_{ij}\) of class \(C^1\), and the
row follows \eqref{eq:obs-attention-replicator} with \(d_{ij}(\lambda(t))\) in
place of \(\E_{\mathbb Q_Y}D_{ij}\), then
\(\Fenergy_i^{\rm att}(\beta_i;\lambda)
=\KL(\beta_i\Vert\pi_i)+\tau_i^{-1}\sum_j\beta_{ij}d_{ij}(\lambda)\) satisfies
\begin{equation}
\frac{d}{dt}\Fenergy_i^{\rm att}(\beta_i^Q(t);\lambda(t))
=-\mathsf r_i\operatorname{Var}_{\beta^Q_i}(\mathsf c_{ij})
+\frac1{\tau_i}\sum_j\beta^Q_{ij}\nabla_\lambda d_{ij}(\lambda)\cdot\dot\lambda,
\label{eq:obs-attention-joint-rate}
\end{equation}
whose second term has no sign in general. Along any differentiable path on
\(\{\alpha_i>0\}\) an event law and its row satisfy
\eqref{eq:rg-attention-beta-event}, and a softmax row
\(\beta_j=\pi_je^{-E_j/\tau}/\sum_k\pi_ke^{-E_k/\tau}\) on a strictly positive
prior support, with \((\pi,E,\tau)\) differentiable and \(\tau>0\), satisfies
\eqref{eq:rg-attention-beta-softmax}:
\begin{align}
\dot\beta_{ij}
&=\frac{\dot\eta_{ij}-\beta_{ij}\dot\alpha_i}{\alpha_i},
\label{eq:rg-attention-beta-event}\\
\dot\beta_j&=\beta_j\left(\chi_j-\sum_k\beta_k\chi_k\right),
\qquad
\chi_j=\partial_t\log\pi_j-\frac{\dot E_j}{\tau}
+\frac{E_j\dot\tau}{\tau^2}.
\label{eq:rg-attention-beta-softmax}
\end{align}
\status{ESTABLISHED}

The proof is in \cref{sec:proofs-attention}.
No scale map determines the path in \eqref{eq:rg-attention-beta-softmax}, and
on \(\{\alpha_i=0\}\) the flow variable is \(\eta\), not the quotient row.

\propositionheading{Attention-reduced potential and its curvature}{prop:obs-attention-reduced-hessian}
Fix a row \(i\), an open set \(U\) of a finite-dimensional real vector space,
and \(E_{ij}\in C^2(U)\) for \(j\in\mathcal N_i\), and define
\begin{equation}
R_i(x)=-\log\sum_j\pi_{ij}e^{-E_{ij}(x)/\tau_i},
\qquad
\beta_{ij}(x)=\pi_{ij}e^{-E_{ij}(x)/\tau_i}\Big/\sum_k\pi_{ik}e^{-E_{ik}(x)/\tau_i}.
\label{eq:obs-attention-reduced-potential}
\end{equation}
Then \(R_i\in C^2(U)\), and
\begin{equation}
\nabla R_i=\frac1{\tau_i}\sum_j\beta_{ij}\nabla E_{ij},
\qquad
\nabla^2R_i=\frac1{\tau_i}\sum_j\beta_{ij}\nabla^2E_{ij}
-\frac1{\tau_i^2}\operatorname{Cov}_{\beta_i}(\nabla E_i),
\label{eq:obs-attention-reduced-hessian}
\end{equation}
where \(\operatorname{Cov}_{\beta_i}(\nabla E_i)
=\sum_j\beta_{ij}(\nabla E_{ij}-\bar g)(\nabla E_{ij}-\bar g)^\top\) and
\(\bar g=\sum_j\beta_{ij}\nabla E_{ij}\). With
\(\mathscr J_i(x,\beta)=\KL(\beta\Vert\pi_i)+\tau_i^{-1}\sum_j\beta_jE_{ij}(x)\),
the gradient is the frozen-weight gradient
\(\partial_x\mathscr J_i(x,\beta_i(x))\), and the Hessian is the Schur
complement, over the tangent space \(\{w:\sum_jw_j=0\}\) of the source simplex,
of the joint Hessian of \(\mathscr J_i\) at \((x,\beta_i(x))\). The potential
arises under both label treatments: under the constant-row restriction, with
\(x=\lambda\) the coordinate of a recognition family \(\mathbb Q_{Y,\lambda}\)
whose averaged energies \(E_{ij}(\lambda)=\E_{\mathbb Q_{Y,\lambda}}D_{ij}\)
are \(C^2\), \(R_i(\lambda)\) is the row-optimized value \(R_i^Q\) of
\cref{prop:obs-attention-constant-row-gap}; under exact label integration, with
\(x=y\) and \(E_{ij}=D_{ij}\) on an open set of latent coordinates where every
\(D_{ij}\) is \(C^2\), summing over \(J_i\) in
\eqref{eq:obs-attention-augmented-likelihood} replaces the label factor by
\(\mathrm c_i(o_i,y)e^{-R_i(y)}\), so \(R_i=R_i^P\) enters the negative log
posterior density of the latent state relative to its baseline. In both cases
the correction along a direction \(v\) is
\(-\tau_i^{-2}\operatorname{Var}_{\beta_i}(\nabla E_{ij}\cdot v)\), so
positivity of the frozen-weight curvature \(\sum_j\beta_{ij}\nabla^2E_{ij}\)
does not imply positivity of \(\nabla^2R_i\), and in the following normalized
witness the correction reverses the sign. Let the latent \(X\in\R\) have
baseline \(\mathcal N(0,\tau/\omega)\) with \(\omega,\tau>0\), let one row have
source alphabet \(\{-1,1\}\) with uniform prior, and let the record \(O\in\R\)
have kernel \(\mathcal N(x-j\delta,\tau)\) given \(J=j\), with \(\delta>0\).
This kernel is \(\mathrm c\exp[-D_j^o(x)/\tau]\) with
\(\mathrm c=(2\pi\tau)^{-1/2}\) and \(D_j^o(x)=(o-x+j\delta)^2/2\), so
\eqref{eq:obs-attention-likelihood} holds at every record value, and the
evidence at \(o=0\),
\(\exp[-\delta^2/(2\tau(1+\omega^{-1}))]/\sqrt{2\pi\tau(1+\omega^{-1})}\), is
finite and positive. Exact label integration gives the negative log posterior
density
\begin{equation}
\frac{(\omega+1)x^2}{2\tau}-\log\cosh\frac{\delta x}{\tau}+\text{const},
\qquad
\text{with second derivative }\frac1\tau\left(\omega+1-\frac{\delta^2}{\tau}\right)
\text{ at }x=0,
\label{eq:obs-attention-reduced-witness}
\end{equation}
whereas the prior curvature plus the frozen-weight curvature is
\((\omega+1)/\tau\). Under the constant-row restriction with Gaussian
recognition \(\mathcal N(\mathsf m,\mathsf v)\), \(\mathsf v>0\), the averaged
energies \(((\mathsf m-j\delta)^2+\mathsf v)/2\) have differences affine in
\(\mathsf m\), and the \(\mathsf m\)-curvature of the row-optimized collective
VFE at \(\mathsf m=0\) is again \((\omega+1-\delta^2/\tau)/\tau\). For
\((\omega,\delta,\tau)=(1,2,1)\) the frozen-weight value is \(2\) and the
reduced value \(-2\), although every source energy and the prior are convex,
so the symmetric point is a strict local maximum of the reduced functional
along the mean.
\status{ESTABLISHED}

The proof is in \cref{sec:proofs-attention}.

\openproblemheading{Discrete creation and deletion of interaction edges}{open:obs-edge-creation}
At \(c_*\) the interaction multigraph is the support of the source prior and
belongs to the structural data. Every row of this chapter is strictly positive
on it (\cref{prop:obs-attention-elbo}), and the row flows change no support at
finite time: along \eqref{eq:obs-attention-replicator},
\(\beta^Q_{ij}(t)=\beta^Q_{ij}(0)\exp[-\int_0^t\mathsf r_i
(\mathsf c_{ij}-\bar{\mathsf c}_i)ds]\) with
\(\bar{\mathsf c}_i=\sum_k\beta^Q_{ik}\mathsf c_{ik}\) whenever the costs stay
finite. Kinetics that change the support lie outside the model of this
manuscript, and the problem fixes one rule so that the question is which
normalized free energy, if any, such kinetics descend. Let the edge set be a
multigraph on \(V\) with edge multiplicity at most a fixed bound, let a
comparison energy \(D_{ij}\) be declared for every candidate copy \(j\to i\),
present or absent, and let the rows on the current support follow
\eqref{eq:obs-attention-replicator}. The threshold rule deletes a present copy
at a rate that is a continuous nonincreasing function of its recognition row
weight \(\beta^Q_{ij}\), positive below a threshold in \((0,1)\) and zero
above it, and creates an absent copy at a rate that is a continuous
nonincreasing function of its averaged energy \(\E_{\mathbb Q_Y}D_{ij}\),
positive below a finite threshold and zero above it; at a deletion the
receiver's row is restricted to the remaining copies and renormalized, and at
a creation the new copy enters the row at the deletion threshold and the row
is renormalized. The rule reads the recognition law and the averaged
energies and no law over multigraphs. The problem is settled, for some pair
of such rate functions, by a normalized joint at \(c_*\) in which the edge
set is a structural coordinate, declared as a configuration prior or
structural channel with a named law and coupled to the link, interface, and
label variables, such that the collective VFE decomposition of
\cref{thm:obs-collective-vfe} holds with the edge set as a coordinate, the
event law of \cref{def:obs-attention-event-law} and its pushforward in
\cref{sec:rg-meta-attention} extend to the random support with the
zero-occupancy convention used there, and the recognition law over supports,
evolving by the master equation of the rule with the rows on each support
following \eqref{eq:obs-attention-replicator}, has nonincreasing collective
VFE from every initial support and recognition law; negatively, by a proof
that for every pair of such rate functions no joint of that class has this
property.

\section{Communication as a normalized channel}
\label{sec:selective-kernels}

\definitionheading{Directed communication interface}{def:obs-selective-interface}
For an oriented edge copy \(e:j\to i\), whose independent link is that of
\cref{def:geo-graph-links}, and a channel \(x\in\{b,m\}\), let
\(\mathsf S_j^x\) be the source sample space. Choose transmitted and received
message spaces \(\mathsf H_e^-\) and \(\mathsf H_e^+\), a record space
\(\mathsf O_e\), receiver side information \(v\) in a measurable space, and
encoding, transmission, and decoding kernels, in the sense of the
communication model of \citet{Shannon1948}, jointly measurable in \(v\), whose
composition is the composed value kernel \(K_e^v\):
\begin{align}
&\mathsf E_e^v:\mathsf S_j^x\rightsquigarrow\mathsf H_e^-,
\qquad \mathsf N_e^v:\mathsf H_e^-\rightsquigarrow\mathsf H_e^+,
\qquad \mathsf D_e^v:\mathsf H_e^+\rightsquigarrow\mathsf O_e,
\label{eq:obs-selective-kernel-types}\\
&K_e^v(s,do)=\int\mathsf E_e^v(s,d\xi^-)
\mathsf N_e^v(\xi^-,d\xi^+)\mathsf D_e^v(\xi^+,do).
\label{eq:obs-selective-kernel-composition}
\end{align}

\corollaryheading{Normalization and the kernel-level gauge law}{cor:obs-selective-kernel-gauge}
(a) \(K^v_e\) is a Markov kernel jointly measurable in \(v\), and a source law
\(q_j^x\) gives the law \(q_j^xK_e^v\) on \(\mathsf O_e\), a candidate law when
\(\mathsf O_e\) is a receiver sample space, whose membership in a restricted
receiver family requires a separate closure condition. (b) Let
\(\mathsf S_j^x\) be the channel sample fiber on which \(G\) acts measurably by
\(\rho_x\) as in \eqref{eq:geo-representations}, \(\mathsf O_e\) the receiver's
channel sample fiber, the side information gauge inert, and \(\phi_e^\pm\)
separately chosen bimeasurable recoordinatizations of the message spaces, and
write \(\mathsf g_i\) for the old-to-new bijection of the channel-\(x\) sample
fiber induced at agent \(i\) by the vertex re-expression of
\eqref{eq:geo-regime-two-gauge-law}. The re-expressed kernels
\begin{equation}
\begin{gathered}
\mathsf E_e^{v\prime}(s',\cdot)=(\phi_e^-)_\#\mathsf E_e^v(\mathsf g_j^{-1}s',\cdot),
\qquad
\mathsf N_e^{v\prime}(\xi',\cdot)=(\phi_e^+)_\#\mathsf N_e^v((\phi_e^-)^{-1}\xi',\cdot),\\
\mathsf D_e^{v\prime}(\xi',\cdot)=(\mathsf g_i)_\#\mathsf D_e^v((\phi_e^+)^{-1}\xi',\cdot)
\end{gathered}
\label{eq:obs-selective-kernel-gauge}
\end{equation}
compose to \(K_e^{v\prime}(s',\cdot)=(\mathsf g_i)_\#K_e^v(\mathsf g_j^{-1}s',\cdot)\),
so \(q_j'K_e^{v\prime}=(\mathsf g_i)_\#(q_jK_e^v)\) for
\(q_j'=(\mathsf g_j)_\#q_j\); for a gauge-inert record space the pushforward by
\(\mathsf g_i\) is omitted. (c) For the represented link
\(\rho_x(\Theta^x_e)\), a bimeasurable bijection,
\(F_e^v(s,\cdot):=K_e^v(\rho_x(\Theta^x_e)^{-1}s,\cdot)\) satisfies
\(K_e^v(s,\cdot)=F_e^v(\rho_x(\Theta^x_e)s,\cdot)\), and the link law gives
\(F_e^{v\prime}(s',\cdot)=(\mathsf g_i)_\#F_e^v(\mathsf g_i^{-1}s',\cdot)\); this
separation uses the invertibility of the link and does not determine the link
from \(K_e^v\).
\status{ESTABLISHED}

The proof is in \cref{sec:proofs-attention}.
A kernel may encode a belief-law or model-law point only when that point is a
latent coordinate of \cref{hyp:local-interaction-kernels} at the law-point
level of \cref{def:prob-description-levels}; by
\cref{req:gen-typing-prohibition} it never reads the live recognition law over
that point, and on the model channel a compatible evaluator is additional data.

\section{Factored links and loops}
\label{sec:selective-matrix-links}

\definitionheading{Factored matrix communication}{def:geo-selective-matrix-link}
On the finite-dimensional vector tier, write \(E_i\) for the sample fiber of
agent \(i\) in channel \(x\), distinct from the law fiber
\((\mathcal E_x)_{c_*}\), and suppress \(x\). For each edge copy \(e:j\to i\)
choose linear maps, possibly depending on the copy and with message spaces of
possibly different dimensions,
\begin{equation}
\mathsf L^{\rm enc}_e:E_j\to H_e^-,\qquad
\mathsf L^{\rm tr}_e:H_e^-\to H_e^+,\qquad
\mathsf L^{\rm dec}_e:H_e^+\to E_i,\qquad
\mathsf L_e=\mathsf L^{\rm dec}_e\mathsf L^{\rm tr}_e\mathsf L^{\rm enc}_e,
\label{eq:geo-selective-factorization}
\end{equation}
separately in the belief and model representations, crossing channels only
through the morphisms of \cref{sec:geo-cross-morphisms}. The Markov kernel of
a linear realization is the composed value kernel \(K^v_e\), not
\(\mathsf L_e\). For represented frames \(U_j:\R^{\dim E}\to E_j\), a value
extraction \(\mathsf M^{\rm val}:\R^{\dim E}\to H\) and an output placement
\(\mathsf M^{\rm out}:H\to\R^{\dim E}\) give the shared-head specialization
\begin{equation}
\mathsf L^{\rm enc}_e=\mathsf M^{\rm val}U_j^{-1},\qquad
\mathsf L^{\rm tr}_e=I_H,\qquad
\mathsf L^{\rm dec}_e=U_i\mathsf M^{\rm out},\qquad
\mathsf L_e=U_i\mathsf M^{\rm out}\mathsf M^{\rm val}U_j^{-1}.
\label{eq:geo-selective-attention-link}
\end{equation}
With \(\mathsf G_j=U_j^{-\top}U_j^{-1}\), sending uses
\(U_j\operatorname{im}(\mathsf M^{\rm val})^\top\) and placement uses
\(U_i\operatorname{im}\mathsf M^{\rm out}\); these and the comparison directions
can overlap and need not partition a frame, and
\(\mathsf M^{\rm out}\mathsf M^{\rm val}\) need be neither orthogonal nor
idempotent.

\propositionheading{Endpoint covariance and rank}{prop:geo-selective-covariance}
(a) Under the old-to-new maps \(\mathsf g_i\) of
\cref{cor:obs-selective-kernel-gauge}, here linear, and independent message
coordinate changes \(\phi_e^-\in\GL(H_e^-)\) and \(\phi_e^+\in\GL(H_e^+)\),
\begin{align}
\mathsf L^{\rm enc\prime}_e&=\phi_e^-\mathsf L^{\rm enc}_e\mathsf g_j^{-1},&
\mathsf L^{\rm tr\prime}_e&=\phi_e^+\mathsf L^{\rm tr}_e(\phi_e^-)^{-1},&
\mathsf L^{\rm dec\prime}_e&=\mathsf g_i\mathsf L^{\rm dec}_e(\phi_e^+)^{-1},\nonumber\\
\mathsf L_e'&=\mathsf g_i\mathsf L_e\mathsf g_j^{-1},
\label{eq:geo-selective-covariance}
\end{align}
and the message-basis freedom does not enlarge the structure group. (b) Direct
composition gives
\begin{equation}
\ker\mathsf L^{\rm enc}_e\subseteq\ker\mathsf L_e,\qquad
\operatorname{im}\mathsf L_e\subseteq\operatorname{im}\mathsf L^{\rm dec}_e,\qquad
\operatorname{rank}\mathsf L_e\leq
\min(\operatorname{rank}\mathsf L^{\rm enc}_e,\operatorname{rank}\mathsf L^{\rm tr}_e,
\operatorname{rank}\mathsf L^{\rm dec}_e),
\label{eq:geo-selective-rank}
\end{equation}
so the encoder accesses the quotient \(E_j/\ker\mathsf L^{\rm enc}_e\), the
potential placement space is \(\operatorname{im}\mathsf L^{\rm dec}_e\), and the
reachable linear signal is the possibly smaller \(\operatorname{im}\mathsf L_e\).
An SPD source metric \(\mathsf G_j\) identifies the quotient with
\((\ker\mathsf L^{\rm enc}_e)^{\perp_{\mathsf G_j}}
=\mathsf G_j^{-1}\operatorname{im}(\mathsf L^{\rm enc}_e)^\top\), and
\(\mathsf G_j'=\mathsf g_j^{-\top}\mathsf G_j\mathsf g_j^{-1}\), so a
distinguished sending subspace requires a metric while the kernel and quotient
do not. Every finite-dimensional \(\mathsf L_e\) factors through its rank. (c) A
weighted sum \(\sum_hw_h\mathsf L_h\) of maps \(E_j\to E_i\), each factoring
through a space \(H_h\), has rank at most \(\min(\dim E_i,\sum_h\dim H_h)\),
attained with rank-one factors: for \(\dim E_i=\dim E_j=2\), identity frames,
\(\mathsf L_h=\mathbf e_h\mathbf e_h^\top\), and positive weights, the sum is
\(\operatorname{diag}(w_1,w_2)\).
\status{ESTABLISHED}

\propositionheading{Connection data and communication loops}{prop:geo-selective-loops}
Keep the represented invertible link \(\Theta_e\) of \(e:j\to i\) on the vector
tier, with \(\Theta_{\bar e}=\Theta_e^{-1}\). (a) For any \(\mathsf L_e\) and a
chosen link, \(F_e:=\mathsf L_e\Theta_e^{-1}\) gives
\begin{equation}
\mathsf L_e=F_e\Theta_e,
\qquad
F_e'=\mathsf g_iF_e\mathsf g_i^{-1},
\label{eq:geo-selective-processing-connection}
\end{equation}
and if \(\mathsf L^{\rm enc}_e=\widetilde{\mathsf L}^{\rm enc}_e\Theta_e\) for
some \(\widetilde{\mathsf L}^{\rm enc}_e:E_i\to H_e^-\), then
\(F_e=\mathsf L^{\rm dec}_e\mathsf L^{\rm tr}_e\widetilde{\mathsf L}^{\rm enc}_e\).
The separation does not recover \(\Theta_e\) from communication, and a
receiver-only processing map \(F_e=F_i\) exists exactly when
\(\mathsf L_e\Theta_e^{-1}\) agrees on every incoming edge copy at \(i\). (b)
Communication on the reverse copy is separately specified and need not be an
inverse, an adjoint, or of the same rank. For a closed walk
\(p=(e_0,\ldots,e_{r-1})\) with \(e_a:i_{a+1}\to i_a\) and \(i_r=i_0\), the
communication round trip
\begin{equation}
\mathsf L_p=\mathsf L_{e_0}\cdots\mathsf L_{e_{r-1}}
\qquad\text{satisfies}\qquad
\mathsf L_p'=\mathsf g_{i_0}\mathsf L_p\mathsf g_{i_0}^{-1},
\label{eq:geo-selective-loop}
\end{equation}
even when the maps are singular, so its rank, characteristic polynomial, and
trace are invariant; a singular round trip records processing and possible
loss of distinguishability and is not the link holonomy
\(\mathrm{Hol}^x(p)\) of \eqref{eq:geo-link-holonomy}. (c) For two sites with
identity frames and identity links, \(e:1\to2\) reading coordinate one and
writing coordinate two, and \(f:2\to1\) the reverse,
\begin{equation}
\mathsf L_e=\begin{psmallmatrix}0&0\\1&0\end{psmallmatrix},\qquad
\mathsf L_f=\begin{psmallmatrix}0&1\\0&0\end{psmallmatrix},\qquad
\mathsf L_f\mathsf L_e=\begin{psmallmatrix}1&0\\0&0\end{psmallmatrix},\qquad
\mathsf L_e\mathsf L_f=\begin{psmallmatrix}0&0\\0&1\end{psmallmatrix}
\label{eq:geo-selective-two-site-example}
\end{equation}
give round trips that preserve one direction and discard the other despite
trivial link holonomy, and \(2\mathsf L_f\) in place of \(\mathsf L_f\) adds
gain at the same rank. The shared-head family
\eqref{eq:geo-selective-attention-link} has length-\(r\) product
\(U_{i_0}(\mathsf M^{\rm out}\mathsf M^{\rm val})^rU_{i_r}^{-1}\), which depends
on a path only through its length and endpoints; independent edge maps remove
that restriction. (d) Link holonomy and communication rank are independent: for
\(G=\GL(2)\) in its defining representation and two parallel copies
\(a,b:1\to2\) with \(\Theta_a=I_2\) and
\(\Theta_b=\begin{psmallmatrix}1&1\\0&1\end{psmallmatrix}\), the loop through
\(b\) and \(\bar a\) has nonidentity holonomy, while arbitrary singular
\(\mathsf L_a,\mathsf L_b\) are admissible through
\(F_a=\mathsf L_a\Theta_a^{-1}\) and \(F_b=\mathsf L_b\Theta_b^{-1}\).
\status{ESTABLISHED}

The proofs of \cref{prop:geo-selective-covariance,prop:geo-selective-loops} are in \cref{sec:proofs-attention}.
General matrix link fields also occur in linear lattice gauge theory
\citep{Wetterich2014LinearLattice}, whose compact-group action and continuum
conclusions do not follow from this endpoint algebra.

\section{Comparison, content, and the transformer correspondence}
\label{sec:selective-comparison}

\definitionheading{Role comparison, quadratic disagreement, and mean operator}{def:obs-selective-comparison}
For an edge copy \(e:j\to i\), comparison uses data separate from content
transmission: a comparison space \(H_e^{\rm cmp}\) and query and key maps
\(S^Q_e:E_i\to H_e^{\rm cmp}\) and \(S^K_e:E_j\to H_e^{\rm cmp}\), or, in the
measurable setting, measurement kernels with a common output space. For law
points \(q_i,q_j\) the role divergence is
\begin{equation}
d_e=\KL\bigl((S^Q_e)_\#q_i\Vert(S^K_e)_\#q_j\bigr),
\label{eq:obs-selective-role-divergence}
\end{equation}
a law-point-level energy (\cref{def:prob-description-levels}) admitted where it
is finite; this cospan compares observables and chooses no map back into
\(E_i\). Reference maps with \(S^Q_e=S^Q_{\rm ref}U_i^{-1}\) and
\(S^K_e=S^K_{\rm ref}U_j^{-1}\) reproduce projected query and key comparison;
with independent links the sender can instead be measured by
\(S^K_e=S^K_{\rm ref}U_i^{-1}\Theta_e\), which recovers the reference formula
when \(\Theta_e=U_iU_j^{-1}\) without imposing it. For an SPD state metric
\(\mathsf G_i\) and an SPD comparison precision \(W_e\), the quadratic
disagreement is
\(V_e(k_i,k_j)=\tfrac12(S^Q_ek_i-S^K_ek_j)^\top W_e(S^Q_ek_i-S^K_ek_j)\).
Declared per-head rows \(\beta_{ij}^{(h)}\), value extractions
\(\mathsf M^{\rm val}_h:\R^{\dim E}\to H_h\), and output placements
\(\mathsf M^{\rm out}_h:H_h\to\R^{\dim E}\) define the mean operator
\begin{equation}
\mu_i^{\rm new}=\mu_i+\sum_{j,h}
\beta_{ij}^{(h)}U_i\mathsf M^{\rm out}_h\mathsf M^{\rm val}_hU_j^{-1}\mu_j,
\label{eq:obs-selective-attention-residual}
\end{equation}
which in reference coordinates is value extraction, weighted aggregation,
output placement, and a residual, the operator form of multi-head attention
\citep{vaswani2017attention}. The operator is declared:
\cref{def:obs-attention-label-model} constructs one label per row, and nothing
here selects the content maps.

\propositionheading{Natural-gradient response to a quadratic disagreement}{prop:obs-selective-response}
Fix the maps, an SPD state metric \(\mathsf G_i\), and an SPD precision \(W_e\).
The negative natural gradient of \(V_e\) in \(k_i\) is
\begin{equation}
-\operatorname{grad}_{\mathsf G_i}V_e
=-\mathsf G_i^{-1}S_e^{Q\top}W_eS^Q_ek_i+\mathsf G_i^{-1}S_e^{Q\top}W_eS^K_ek_j.
\label{eq:obs-selective-response}
\end{equation}
Descent of this energy thus determines one receiving map
\(\mathsf G_i^{-1}S_e^{Q\top}W_eS^K_e\) and one self term, so a declared
insertion such as \eqref{eq:obs-selective-attention-residual} is a
natural-gradient step of a quadratic disagreement only when its maps have this
form; identifying \(\mathsf G_i\) with a Fisher metric requires a regular
statistical model that supplies it (\cref{ch:relational-inference}).
\status{ESTABLISHED}

\noindent\emph{Proof.}
Differentiate \(V_e\) in \(k_i\) and multiply by \(\mathsf G_i^{-1}\).
\(\square\)

\corollaryheading{Softmax over role divergences}{cor:obs-selective-softmax}
(a) When the latent coordinates include the law points compared in
\eqref{eq:obs-selective-role-divergence}, at the law-point level of
\cref{def:prob-description-levels}, a latent configuration \(y\) gives the
source energies \(D_{ij}(y)=d_{ij}\). On a finite positive-prior source set
with finite costs and \(\tau_i>0\), the categorical row functional satisfies
\begin{equation}
\sum_j\beta_{ij}d_{ij}+\tau_i\KL(\beta_i\Vert\pi_i)
=\tau_i\KL\bigl(\beta_i\Vert\beta_i^P(y)\bigr)
-\tau_i\log\sum_j\pi_{ij}e^{-d_{ij}/\tau_i},
\label{eq:obs-selective-row-functional}
\end{equation}
so the posterior row \eqref{eq:obs-attention-posterior} is its unique
minimizer: softmax attention over divergences between current law points is
the posterior row at a latent configuration. The recognition optimum
\eqref{eq:obs-attention-recognition-optimum} reduces to it for
\(\mathbb Q_Y=\delta_y\), and when the latent baseline gives \(\{y\}\) zero mass
that point mass is not absolutely continuous with respect to the posterior
slice, so its relative entropy and VFE are infinite
(\cref{prop:elbo-subspace-support-singular}); the law-point prior that such a
latent requires is \cref{open:prob-law-point-prior}. (b) Let the projected key
covariance be the same SPD matrix \(\Sigma^K\) for every source, and put
\(\mu_i^Q=S^Q_{\rm ref}U_i^{-1}\mu_i\) and \(\mu_j^K=S^K_{\rm ref}U_j^{-1}\mu_j\).
For Gaussian law points
\begin{equation}
d_{ij}=d_i^\circ-(\mu_i^Q)^\top(\Sigma^K)^{-1}\mu_j^K
+\tfrac12(\mu_j^K)^\top(\Sigma^K)^{-1}\mu_j^K,
\label{eq:obs-selective-dot-product}
\end{equation}
where \(d_i^\circ\) collects every term independent of \(j\), including the
query covariance terms, so whitening gives a dot-product logit together with a
key-norm bias and \(\log\pi_{ij}\). The bias disappears when the whitened key
norms are constant, or when a source prior cancels it at these comparison
laws, and such compensation in a population model requires its own source
law. Unequal projected key covariances generally retain additional
source-dependent terms.
\status{ESTABLISHED}

The proof is in \cref{sec:proofs-attention}.

\corollaryheading{Content changes routing}{cor:obs-selective-content-kl}
Let \(p(j,y,v)=\pi_jp_j(y)r_j(v\given y)\) and
\(q(j,y,v)=\beta_j\widetilde q(y)s_j(v\given y)\) be normalized laws on a
source label, a comparison state, and a content variable, with \(q\ll p\) and
finite terms. Then
\begin{equation}
\KL(q\Vert p)=\KL(\beta\Vert\pi)+
\sum_j\beta_j\left[
\KL(\widetilde q\Vert p_j)+
\E_{\widetilde q}\KL(s_j(\cdot\given Y)\Vert r_j(\cdot\given Y))\right],
\label{eq:obs-selective-content-kl}
\end{equation}
so the conditional content term generally changes the optimal routing, and
matching \(s_j=r_j\) removes it.
\status{ESTABLISHED}

\noindent\emph{Proof.}
Apply \cref{lem:prob-kl-chain-rule} to the disintegration over \(j\), then over
\(y\) given \(j\). \(\square\)

\subsection{Exact receiver update and transformer attention}
\label{sec:selective-transformer}

\propositionheading{Exact receiver update in the comparison-record model}{prop:obs-selective-receiver-update}
Fix a receiver \(i\) with a finite source set on which \(\pi_{ij}>0\), each
source \(j\) reaching \(i\) through one edge copy \(e:j\to i\), and suppose
that these copies share the comparison space \(H_i^{\rm cmp}\), the query map
\(S^Q_i:=S^Q_e\), and an SPD precision \(W_i\), while the key maps \(S^K_e\) may
depend on the copy. Let the attention record of
\cref{def:obs-attention-label-model} be the comparison record
\begin{equation}
o_i\given k_i,J_i=j,y_{-i}\sim
\mathcal N\bigl(S^Q_ik_i-S^K_ek_j,\tau_iW_i^{-1}\bigr),
\qquad\tau_i>0,
\label{eq:obs-selective-attention-record}
\end{equation}
observed at \(o_i=0\), where \(y_{-i}\) collects every latent coordinate other
than \(k_i\) and \(J_i\) and so fixes the source states \(k_j\). At \(o_i=0\)
its density is \(\mathrm c_i\exp[-V_e(k_i,k_j)/\tau_i]\), with \(V_e\) of
\cref{def:obs-selective-comparison} and
\(\mathrm c_i=(2\pi\tau_i)^{-\dim H_i^{\rm cmp}/2}(\det W_i)^{1/2}\), so it
realizes \eqref{eq:obs-attention-likelihood}. Suppose that it is the only
factor reading \(J_i\), as in \eqref{eq:obs-attention-augmented-likelihood},
and that, given \(y_{-i}\) and every record other than \(o_i\), the law of
\(k_i\) is \(\mathcal N(\mu_i,\Sigma_i)\) with \(\Sigma_i\succ0\). Write
\(\widetilde W_i=W_i/\tau_i\) and
\begin{equation}
\Sigma_i^+=\bigl(\Sigma_i^{-1}+S_i^{Q\top}\widetilde W_iS^Q_i\bigr)^{-1},
\qquad
\Sigma_i^{\rm pred}=\widetilde W_i^{-1}+S^Q_i\Sigma_iS_i^{Q\top},
\qquad
r_{ij}=S^Q_i\mu_i-S^K_ek_j.
\label{eq:obs-selective-receiver-covariances}
\end{equation}
(a) Given \(y_{-i}\) and all records, \(J_i\) has the posterior weights
\(w_{ij}\), normalized over the source set, \(k_i\) has the conditional
components, and the conditional mean of \(k_i\) is
\begin{gather}
w_{ij}\propto\pi_{ij}\exp\bigl[-\tfrac12r_{ij}^\top(\Sigma_i^{\rm pred})^{-1}r_{ij}\bigr],
\qquad
k_i\given J_i=j\sim\mathcal N\bigl(\Sigma_i^+(\Sigma_i^{-1}\mu_i
+S_i^{Q\top}\widetilde W_iS^K_ek_j),\Sigma_i^+\bigr),
\label{eq:obs-selective-receiver-posterior}\\
\E[k_i\given y_{-i},o]=\Sigma_i^+\Sigma_i^{-1}\mu_i+\sum_jw_{ij}\mathsf L_e^{\rm val}k_j,
\quad
\mathsf L_e^{\rm val}=\Sigma_i^+S_i^{Q\top}\widetilde W_iS^K_e,
\quad
\Sigma_i^+\Sigma_i^{-1}=I-\Sigma_i^+S_i^{Q\top}\widetilde W_iS^Q_i.
\label{eq:obs-selective-receiver-update}
\end{gather}
The value map is thus fixed by the query map, the key map, the record
precision, and the receiver covariance, and the skip connection is replaced by
the precision-weighted term \(\Sigma_i^+\Sigma_i^{-1}\mu_i\); equivalently the
update is
\(\mu_i+\sum_jw_{ij}\Sigma_i^+S_i^{Q\top}\widetilde W_i(S^K_ek_j-S^Q_i\mu_i)\).
(b) The weights use \(\Sigma_i^{\rm pred}\) in place of \(\widetilde W_i^{-1}\)
and need not equal the posterior row \eqref{eq:obs-attention-posterior} at
\(k_i=\mu_i\): for a scalar receiver with
\(S^Q_i=S^K_e=\widetilde W_i=\Sigma_i=1\), \(\mu_i=0\), a uniform prior, and
sources \(k_1=0\) and \(k_2=2\), one has \(w_{i2}=e^{-1}/(1+e^{-1})\) and
\(\beta^P_{i2}=e^{-2}/(1+e^{-2})\) at \(k_i=0\), while
\(\mathsf L_e^{\rm val}=\Sigma_i^+\Sigma_i^{-1}=1/2\). (c) The Woodbury identity
\(\Sigma_i^+S_i^{Q\top}\widetilde W_i=\Sigma_iS_i^{Q\top}(\Sigma_i^{\rm pred})^{-1}\)
holds, and for full-row-rank \(S^Q_i\) and \(\widetilde W_i^{-1}\to0\) the value
map tends to \(R_eS^K_e\) with
\(R_e=\Sigma_iS_i^{Q\top}(S^Q_i\Sigma_iS_i^{Q\top})^{-1}\), the
\(\Sigma_i^{-1}\)-minimal right inverse of \(S^Q_i\). (d) With copy-dependent
query maps, precisions, or comparison spaces, the same computation gives
\(w_{ij}\propto\pi_{ij}\mathcal N(0;r_{ij},\Sigma_{ij}^{\rm pred})\) and
component covariances \(\Sigma_{ij}^+\), so the normalizer of the predictive
law enters the weights.
\status{ESTABLISHED}

The proof is in \cref{sec:proofs-attention}.
\Cref{tab:obs-transformer-correspondence} compares the pointwise attention
network with the attention layer of \citet{vaswani2017attention} row by row;
the metric lift \(R_eS^K_e\) of (c) is developed for Gaussian beliefs in
Supplement G.

\begin{sciencetable}{Correspondence between transformer attention and the pointwise attention network at $c_*$.}
\label{tab:obs-transformer-correspondence}
\begin{tabularx}{\textwidth}{>{\raggedright\arraybackslash}p{0.2\textwidth}>{\raggedright\arraybackslash}X>{\raggedright\arraybackslash}p{0.17\textwidth}}
\toprule
\textbf{Transformer} & \textbf{Pointwise attention network} & \textbf{Standing} \\
\midrule
Token position & Agent $i\in V$ at $c_*$, its embedding read as a coordinate of its latent state. & declared \\
Query and key projections & $S^Q_e=S^Q_{\rm ref}U_i^{-1}$, $S^K_e=S^K_{\rm ref}U_j^{-1}$ and the role divergence \eqref{eq:obs-selective-role-divergence}; with a common projected key covariance, a dot-product logit plus a key-norm bias, \eqref{eq:obs-selective-dot-product}. & derived, \cref{cor:obs-selective-softmax} \\
One key map shared by all edges & Independent links enter $S^K_e=S^K_{\rm ref}U_i^{-1}\Theta_e$ only through the represented edge twist $U_i^{-1}\Theta_eU_j$; the shared map is the identity-twist case $\Theta_e=U_iU_j^{-1}$, and a nonidentity twist, including nontrivial link holonomy, has no counterpart in that architecture. & derived, \cref{prop:geo-edge-twist-invariance} \\
Softmax row & Generative posterior of a categorical source label, \eqref{eq:obs-attention-posterior}; over law-point divergences, the posterior row at a latent configuration, \eqref{eq:obs-selective-row-functional}, not the recognition optimum. & derived, \cref{prop:obs-attention-elbo,cor:obs-selective-softmax} \\
Attention mask & Zero prior weight: the source lies outside the support of $\pi_i$, hence outside the interaction multigraph, where the posterior row vanishes. & derived, \cref{prop:obs-attention-elbo} \\
Temperature, inverse square root of key width & The temperature $\tau_i>0$ of \eqref{eq:obs-attention-likelihood}; for $\Sigma^K=\varsigma^2I$ the dot-product part of the logit $-d_{ij}/\tau_i$ is $(\mu_i^Q)^\top\mu_j^K/(\varsigma^2\tau_i)$, and $\varsigma^2\tau_i=\sqrt{\dim H^{\rm cmp}_i}$ reproduces the scaled dot product, a choice not derived here. & declared \\
Value and output maps, residual connection & Declared in \eqref{eq:obs-selective-attention-residual}; in the comparison-record model the exact update fixes the value map $\mathsf L_e^{\rm val}$ and weights the residual by $\Sigma_i^+\Sigma_i^{-1}$. & derived, \cref{prop:obs-selective-receiver-update} \\
Multiple heads & Declared per-head rows $\beta_{ij}^{(h)}$; no label-level construction, and no head index in the event law \eqref{eq:rg-attention-event-law}. & absent, \cref{open:obs-multihead-attention} \\
Feed-forward block, layer normalization & No counterpart is modeled. & absent \\
\bottomrule
\end{tabularx}
\end{sciencetable}

\propositionheading{A shared receiver state couples head labels}{prop:obs-selective-head-coupling}
Give two heads \(h=1,2\) of one receiver their own exclusive labels
\(J_{i,h}\) and comparison records of the form
\eqref{eq:obs-selective-attention-record}, conditionally independent given the
latent state. At fixed \(k_i\) the label posterior factors by head. For a
scalar receiver with \(k_i\sim\mathcal N(0,1)\) given the rest,
\(S^Q=S^K=\widetilde W=1\) for both heads, a uniform prior, and sources with
states \(k_1=0\) and \(k_2=a\ne0\), the posterior mass of
\((J_{i,1},J_{i,2})=(j,l)\) after \(k_i\) is integrated is proportional to
\begin{equation}
\exp\bigl[-\tfrac13(k_j^2+k_l^2)+\tfrac13k_jk_l\bigr],
\label{eq:obs-selective-head-coupling}
\end{equation}
and its odds ratio between agreeing and disagreeing label pairs is
\(e^{a^2/3}\ne1\), which excludes a product law.
\status{ESTABLISHED}

The proof is in \cref{sec:proofs-attention}.

\remarkheading{The Gaussian criterion for head-label coupling}{rem:obs-head-coupling-criterion}
For two heads with independent positive label priors, let their stacked
predictive record covariance $\Sigma^{\mathrm{rec}}$ be positive definite
and independent of the labels. Write $\Lambda^{\mathrm{rec}}_{12}$ for the
off-diagonal block of its inverse, and $\mathsf a_j,\mathsf b_l$ for the
head-specific key means. The positive integrated posterior table $p_{jl}$
satisfies
\begin{equation}
\log\frac{p_{jl}p_{j'l'}}{p_{jl'}p_{j'l}}
=-(\mathsf a_j-\mathsf a_{j'})^\top
\Lambda^{\mathrm{rec}}_{12}(\mathsf b_l-\mathsf b_{l'}).
\label{eq:obs-head-cross-ratio}
\end{equation}
Thus the labels are independent exactly when every displayed pairing
vanishes. Nonzero covariance between projected queries alone does not imply
label dependence. \status{ESTABLISHED}

\noindent\emph{Proof.}
Expand the Gaussian log likelihood with precision
$(\Sigma^{\mathrm{rec}})^{-1}$. Its mixed term is
$-\mathsf a_j^\top\Lambda^{\mathrm{rec}}_{12}\mathsf b_l$; the prior
and one-head terms cancel from the cross ratio. A positive table factors
exactly when all its cross ratios equal one. For example, a receiver
$k_i\sim\mathcal N(0,I_3)$ with query maps
$S^Q_1k_i=(k_{i,1},k_{i,2})$ and
$S^Q_2k_i=(k_{i,1},k_{i,3})$, independent record noises of covariance $I_2$,
and key means $(0,0),(0,a)$ and $(0,0),(0,b)$ has query cross-covariance
$\operatorname{diag}(1,0)$ but every mixed key pairing zero. For
$a,b\ne0$ both labels remain informative. $\square$

\openproblemheading{Label-level multi-head attention}{open:obs-multihead-attention}
The per-head rows of \eqref{eq:obs-selective-attention-residual} are declared,
and the event law \eqref{eq:rg-attention-event-law} carries no head index. Give
a receiver \(i\) at least two heads \(h\), each an exclusive label \(J_{i,h}\)
with a comparison record of the form \eqref{eq:obs-selective-attention-record}
and its own query map \(S^Q_{i,h}\), key maps, precision, and temperature, the
records conditionally independent given the latent state, in the setting of
\cref{prop:obs-selective-receiver-update}, where \(k_i\) has the conditional
law \(\mathcal N(\mu_i,\Sigma_i)\) given the rest. The heads read a common
receiver coordinate when the projected queries \(S^Q_{i,h}k_i\) and
\(S^Q_{i,h'}k_i\) are correlated under that law for some pair \(h\ne h'\);
require also that the integrated positive label table does not factor by
head. These conditions define the coupled class of the problem. For two
heads with label-independent predictive covariance, the second condition
is the nonvanishing of some pairing in
\cref{eq:obs-head-cross-ratio}; \cref{prop:obs-selective-head-coupling}
supplies one instance. The factorizing subclass is excluded, including
heads that read disjoint channels through untied links under a receiver law
that is a product over channels.
What holds in the class is the per-row label posterior at a fixed receiver
state, \cref{prop:obs-attention-elbo} with rows \((i,h)\); the event law of
\cref{def:obs-attention-event-law} carries no head index, so the pushforward
of \cref{prop:rg-meta-attention-pushforward} gives a parent no row per head.
A parent whose receiver state is one coarse coordinate of its block
reads every row of its children through that coordinate, so the exact
recursion multiplies the number of rows read through one coordinate at each
depth unless a head-merging channel, which the recursion does not declare, is
imposed; once declared, its conditional relative entropy enters
\eqref{eq:rg-vfe-chain-rule}. The problem is settled by an exact recursion
with bounded head count in the class: a family of comparison-record models,
closed under the pushforward of \cref{prop:rg-meta-attention-pushforward}
composed with a declared head-merging channel without increase of the number
of heads per receiving row, in which the parent's receiver state is a coarse
coordinate of its block with the pushed-forward conditional law, its per-head
rows are the disintegrations of the merged event law, and its posterior mean
given its records and the rest is the coarse coordinate of its children's
posterior means; negatively, by a member of the class and a partition into
blocks for which the posterior mean of a block's coarse coordinate, as a
function of the source states, is reproduced by no comparison-record model
at the block with fewer heads than the block has rows.

\section{Compositions as records}
\label{sec:selective-records}

Choosing \(s=s(y_j)\) and \(v=v(X,y_{\partial e})\) measurably in
\eqref{eq:obs-selective-kernel-composition} makes the output one owned record
of \cref{hyp:local-interaction-kernels}, and replacing record kernels by these
compositions leaves the joint \eqref{eq:obs-interaction-joint} normalized on
cycles by \cref{prop:gen-exact-normalization}, with distinct sending and
receiving interfaces and under the typing of \cref{req:gen-typing-prohibition}.
Records that share a realized random message retain it once as a latent or
declare their joint kernel, since conditional independence cannot be assumed
after it is marginalized; independent draws from one latent law point are
independent given that point. Fixed-record density formulas keep the family
domination and regular-slice hypotheses of \cref{hyp:local-interaction-kernels},
under which the collective VFE of \cref{thm:obs-collective-vfe} and the
local--global potential identity of \cref{thm:obs-local-global-potential}
apply. Composed records realize rectangular endpoint maps on cycles (a linear
Gaussian realization is in Supplement G) but settle neither a canonical
compression of such maps nor coarse-graining through moving interfaces
(\cref{prop:geo-selective-descent}).

\definitionheading{Learned and random interfaces}{def:obs-selective-interface-learning}
If interfaces or links vary as generative parameters, their score includes
every varying normalizer, baseline, and record term. If they are random
relation coordinates that evolve during inference, they are latent coordinates
with a normalized prior and declared dynamics in the augmented state, and
natural-gradient evolution needs a regular configuration metric or a justified
quotient; endpoint covariance specifies neither the metric nor the learning
rule. In both cases the dynamics move link values and interface maps on the
edge copies of the support of the source prior; creating or deleting edge
copies is \cref{open:obs-edge-creation}.

\section{Information loss and descent}
\label{sec:selective-information}

\propositionheading{Conditional data processing through an interface}{prop:obs-selective-information-chain}
For a source law and a side-information law, let \(S_j\) be the source sample
in \(\mathsf S_j^x\), \(\Upsilon\) the random receiver side information with
values \(v\), \(\xi^-\) and \(\xi^+\) the random messages, and \(O\) the decoded
record. Construction \eqref{eq:obs-selective-kernel-composition} makes
\(S_j\to\xi^-\to\xi^+\to O\) a Markov chain conditional on \(\Upsilon\), and
\begin{equation}
I(S_j;O\given\Upsilon)\leq I(S_j;\xi^+\given\Upsilon)\leq I(S_j;\xi^-\given\Upsilon)
\qquad\text{in }[0,+\infty].
\label{eq:obs-selective-information-chain}
\end{equation}
\status{ESTABLISHED}

The proof is in \cref{sec:proofs-attention}.

\propositionheading{Communicated agreement leaves conditional information}{prop:obs-selective-private-kl}
Let \(q,p\) be laws on one standard Borel space with \(q\ll p\) and finite
relative entropy, and let \(M\) be one measurable measurement into a standard
Borel space with output laws \(q_M,p_M\). Then
\begin{equation}
\KL(q\Vert p)=\KL(q_M\Vert p_M)+
\int\KL\bigl(q(\cdot\given M=z)\Vert p(\cdot\given M=z)\bigr)q_M(dz).
\label{eq:obs-selective-private-kl}
\end{equation}
For \(M=(1,0)\), the laws \(\mathcal N((0,a),I_2)\) and \(\mathcal N((0,b),I_2)\)
have identical communicated laws and full relative entropy \((a-b)^2/2\).
\status{ESTABLISHED}

\noindent\emph{Proof.}
Append the deterministic measurement to both laws, which leaves the relative
entropy unchanged, and apply the disintegration form of
\cref{lem:prob-kl-chain-rule} over \(M\). The example has equal first
coordinates and Gaussian relative entropy \((a-b)^2/2\). \(\square\)

The identity needs one measurement: distinct roles \(S^Q_e\ne S^K_e\) need not
ask the same question, and separate head marginals need not fix their joint
law (\cref{prop:prob-marginals-do-not-determine-joint}).

\propositionheading{Descent through a selective interface}{prop:geo-selective-descent}
Let \(E_j,E_i,E\) and \(\overline E_j,\overline E_i,\overline E\) be
finite-dimensional real vector spaces, and let \(M_j:E_j\to\overline E_j\),
\(M_i:E_i\to\overline E_i\), and \(M:E\to\overline E\) be surjective linear
measurements. (a) Static intertwiner: for a linear \(\mathsf L_e:E_j\to E_i\), a
linear \(\overline{\mathsf L}_e:\overline E_j\to\overline E_i\) with
\(M_i\mathsf L_e=\overline{\mathsf L}_eM_j\) exists exactly when
\begin{equation}
\mathsf L_e\ker M_j\subseteq\ker M_i,
\label{eq:geo-selective-descent}
\end{equation}
and it is then unique. (b) Autonomous linear flow: for
\(\mathcal A\in\operatorname{End}(E)\), a linear
\(\overline{\mathcal A}\in\operatorname{End}(\overline E)\) with
\(Me^{t\mathcal A}=e^{t\overline{\mathcal A}}M\) for every \(t\in\R\),
equivalently with \(M\mathsf u\) solving
\(\dot{\bar{\mathsf u}}=\overline{\mathcal A}\bar{\mathsf u}\) for every
solution of \(\dot{\mathsf u}=\mathcal A\mathsf u\), exists exactly when
\begin{equation}
M\mathcal A=\overline{\mathcal A}M
\label{eq:geo-selective-flow-descent}
\end{equation}
is solvable, that is, exactly when \(\mathcal A\ker M\subseteq\ker M\), and
\(\overline{\mathcal A}\) is then unique. (c) Moving measurement: on an open
interval \((t_0,t_1)\), let \(t\mapsto M_t:E\to\overline E\) be a \(C^1\)
family of surjections and \(t\mapsto\mathcal A_t\) continuous. A family
\(\overline{\mathcal A}_t\in\operatorname{End}(\overline E)\) with
\(\tfrac{d}{dt}(M_t\mathsf u(t))=\overline{\mathcal A}_tM_t\mathsf u(t)\) for
every solution of \(\dot{\mathsf u}=\mathcal A_t\mathsf u\) on \((t_0,t_1)\)
exists exactly when
\begin{equation}
\dot M_t+M_t\mathcal A_t=\overline{\mathcal A}_tM_t
\qquad(t_0<t<t_1)
\label{eq:geo-selective-moving-descent}
\end{equation}
is solvable, that is, exactly when \(\dot M_t+M_t\mathcal A_t\) vanishes on
\(\ker M_t\) for every \(t\); the family is then unique and continuous, with
\(\overline{\mathcal A}_t=(\dot M_t+M_t\mathcal A_t)M_t^+\) for any continuous
right inverse \(M_t^+\) of \(M_t\). (d) If
\(\mathsf L_e\ker M_j\not\subseteq\ker M_i\), no linear map of the retained
coordinates reproduces \(M_i\mathsf L_e\); a coordinate swap with
\(M_i=M_j=(1,0)\) is an instance, in which a hidden coordinate becomes visible.
If \(\mathcal A\ker M\not\subseteq\ker M\), two solutions with the same
retained initial value have different retained velocities, so the retained
coordinates do not evolve autonomously, and the same holds at each time for a
moving measurement that violates \eqref{eq:geo-selective-moving-descent}. (e) A
low-rank interface does not by itself give an instantaneous Markov blanket:
for a joint law of a private state \(\mathsf z^{\rm in}\), an interface
\(\mathsf z^{\rm bd}\), and an exterior \(\mathsf z^{\rm ex}\), the blanket
condition is
\(P(d\mathsf z^{\rm in},d\mathsf z^{\rm ex}\given\mathsf z^{\rm bd})
=P(d\mathsf z^{\rm in}\given\mathsf z^{\rm bd})P(d\mathsf z^{\rm ex}\given\mathsf z^{\rm bd})\),
and if \(\mathsf z^{\rm in}\) is nondegenerate and independent of
\(\mathsf z^{\rm bd}\) while \(\mathsf z^{\rm ex}=\mathsf z^{\rm in}\), the
conditional law of \((\mathsf z^{\rm in},\mathsf z^{\rm ex})\) given
\(\mathsf z^{\rm bd}\) is concentrated on the diagonal and is not a product.
\status{ESTABLISHED}

The proof is in \cref{sec:proofs-attention}.

\openproblemheading{Formation with evolving selective interfaces}{open:obs-selective-formation}
Take a connected block of agents on the vector tier with Gaussian belief laws
whose means are the fast coordinates, linear factored communication
(\cref{def:geo-selective-matrix-link}) through surjective linear interfaces
\(M_i\) of one fixed rank below the fiber dimension, and comparison records
\eqref{eq:obs-selective-attention-record} whose query and key maps read the
receiver and the transported source through the receiver's interface,
\(S^Q_i=M_i\) and \(S^K_e=M_i\Theta_e\), with the labels of
\cref{def:obs-attention-label-model}. The means follow the gradient flow of
the sample-level energy of the current label assignment with identity
mobility and diffusion covariance \(2\mathrm TI dt\) at a temperature
\(\mathrm T>0\), and the assignment is redrawn at a rate-one clock from the
receiver posterior rows, as in \cref{def:rg-formation-cycle-experiment}; the
interfaces are generative parameters in the sense of
\cref{def:obs-selective-interface-learning} that move by the natural-gradient
flow of the collective VFE under a declared metric on interfaces of that
rank, at a mobility \(\epsilon_{\mathrm M}>0\). This is the class of the
problem. Along such a path the \(\dot M_t\) term of
\eqref{eq:geo-selective-moving-descent} is not controlled by a fixed-interface
construction, because the path reads the state, and neither is the label
switching it induces. A parent for the block is a retained mechanism in the
sense of \cref{sec:rg-memory-parent-mechanism}: a Markov process on a state
space of fixed dimension, below the dimension of the block's means and
assignment and independent of \(\mathrm T\), \(\epsilon_{\mathrm M}\), and
the horizon, whose state determines the communicated coordinates
\(M_{i,t}\mu_i\) of every agent of the block. The admitted state space is a
smooth manifold, the stochastic coefficients between jumps are locally
Lipschitz, initialization and jump maps are continuous in their native inputs,
and the communicated-coordinate readout is \(C^1\) on each assignment
stratum. Any additional conditional-law filter obeys
\cref{def:rg-regular-finite-filter}. These are explicit regularity restrictions
on this problem; dimension alone would allow arbitrary Borel encodings of
the full state or history. What holds is the exact
alternative that retains the private coordinates with their conditional laws
at the dimension of the block (\cref{thm:rg-memory-filter} in finite state),
and the pointwise descent criterion of \cref{prop:geo-selective-descent}(c)
for a supplied \(C^1\) interface path. The problem is settled by a parent of
that kind for which the relative entropy of the fine path law of the
communicated coordinates with respect to the parent's path law, on horizons
of order \(1/\epsilon_{\mathrm M}\), on which the interfaces move by an
amount of order one, tends to zero as \(\epsilon_{\mathrm M}\to0\) uniformly
over \(\mathrm T\) in an interval \((0,\mathrm T_0]\); negatively, by a member
of the class in which the learned interface path has a rest point at which no
parent of that kind is projectable, or in which every parent of that kind has
that relative entropy bounded below on those horizons as
\(\epsilon_{\mathrm M}\to0\) and \(\mathrm T\to0\) along a stated relation.

\paragraph{Physicist's summary.}
Attention is inference over a hidden pointer: a categorical label names whom a
receiver listens to, its prior's support is the interaction graph, and its
record's likelihood falls with a comparison energy. Bayes' formula gives the
softmax row, masks are zero prior weights, and the temperature is a free
parameter of the record, of which the inverse-square-root key-width scaling is
one choice. A recognition row held constant over the latent state is a
different softmax, over averaged energies, at an exact expected
relative-entropy cost that penalizes uncertain sources. Integrating the label
out gives a log-sum-exp potential whose curvature is the frozen-weight
curvature minus an attention covariance, which can turn a convex landscape
concave, the seed of symmetry breaking in the two-timescale and formation
chapters. Communication is a normalized encoder--channel--decoder kernel,
covariant under separate recoordinatizations of each message space; it may be
singular, and a singular round trip is processing, not holonomy. For a Gaussian
comparison record the exact receiver update has the transformer shape, but the
value map is fixed by the query and key maps, the residual is precision
weighted, and the weights use the predictive covariance. Heads sharing one
receiver state can have coupled labels; an exact recursion with bounded
head count for the coupled class remains open.
Coarse-graining carries the receiver--source event law, not the conditional
row, and a coarse source set contains the block itself.

\chapter{Fiber Information Geometry and Inference Histories}
\label{ch:relational-inference}

This chapter asks how the collective free energy of \cref{ch:elbo} selects an
inference history of the agents at $c_*$ and what duration the Fisher geometry of their
law fibers assigns to it.  A chosen metric makes the positive ray of the natural
gradient select an oriented orbit, whose Fisher length is an information clock
invariant under reparameterization (\cref{thm:hist-fisher-clock-invariance}).  A
parameter-independent Markov channel projects scores and loses Fisher information by an
exact conditional-covariance defect (\cref{thm:cg-fisher-contraction}), so a same-path
Markov image of a history, such as the law of a meta-agent, accumulates no more duration
(\cref{thm:hist-record-clock-contraction}), while a coarse orbit recomputed from its own
free energy compares with the fine one only through an oriented semiconjugacy, whose
rate cancels from every duration
(\cref{thm:hist-oriented-semiconjugacy,thm:hist-duration-relation}).  A three-agent
parity witness shows that a collective free energy and its Fisher tensor are not
functions of the agents' marginal laws.  \Cref{ch:two-timescales} uses the mobility and
clock results, \cref{ch:coarsegraining} the Fisher-loss theorem, and
\cref{ch:agent-network-rg,ch:full-graph-meta-agent-vfe} the coarse configuration map and
the record clocks.  Longer proofs are deferred to \cref{app:deferred-proofs}.

\paragraph{Scope.}
Law fibers and kernel fibers are distinct, and no manifold, dominating measure, affine
chart, cone, or Fisher metric is inferred from the words belief or model.  The Fisher
metric is chosen: no invariance axiom that characterizes it
(\cref{prop:ig-fisher-characterizations}) is imposed, and no uniqueness is asserted.  The
parity hyperedge (\cref{def:hist-parity-hyperedge}) is supplied generative data, not
derived from pairwise locality or a canonical agent inventory, and the promoted parity
coordinate (\cref{cor:hist-promoted-parity-rank}) is not thereby an agent or an
intervention variable.  No fine-to-coarse semiconjugacy is asserted for coarse maps that
discard directions on which the fine free energy depends.  Open are correlated lifts
beyond the parity witness (\cref{open:hist-parity-lift-scope}), global Fisher quotients
and canonical block data (\cref{open:hist-retained-fisher-quotient}), the quantitative
meta-agent clock ratio (\cref{open:hist-meta-clock-ratio}), and the slow-manifold limit
of the information-discarding case (\cref{open:obs-slow-manifold}).  \status{NOT-CLAIMED}

\paragraph{Notation.}
Tangent vectors are $Z,Z_1,Z_2,Z_3$, and $\mathbf s_Z$ is the score in the direction $Z$;
$\xi_{\mathcal M}$ is the fundamental vector field of $\xi$ in the Lie algebra
$\mathfrak g$ of $G$ on a $G$-manifold $\mathcal M$.  A parameter-independent Markov
kernel is $K$, with law map $K_\star p=pK$ and joint law
$p\otimes K(dy,d\bar y)=p(dy)K(y,d\bar y)$.  A history is $Q$ with curve parameter $t$,
$\mathsf a$ is a positive rate factor, $\mathcal I_6$ indexes the bits of the parity
witness and $\operatorname{par}$ is its parity character, $\mathcal N(\mu,\varsigma^2)$ is
the normal family, and upright letters $\mathrm c,\mathrm c_1,\mathrm c_2,\mathrm q$ are
constants local to one statement.

\section{Fiber statistical tensors}
\label{sec:pb-statistical-tensors}
\label{sec:ig-typing}
\label{sec:ig-notclaimed}

At $c_*$ a frame $u$ in the $G$-torsor $(\mathscr P_G)_{c_*}$ identifies each law fiber
$\mathcal B_x$ with $(\mathcal E_x)_{c_*}$ through $\iota_u:\beta\mapsto[u,\beta]$, and
$\iota_{u\cdot g}=\iota_u\circ\widehat\rho_x(g)$ by \eqref{eq:geo-quotient-convention},
so a tensor on $\mathcal B_x$ defines one on $(\mathcal E_x)_{c_*}$ exactly when it is
invariant under the represented action.

\hypothesisheading{Regular channel models}{hyp:pb-regular-models}
Each law fiber $\mathcal B_x$, $x\in\{b,m\}$, and each coarse fiber $\bar{\mathcal B}_x$
used below is a finite-dimensional smooth parametrized-measure model with smooth
represented action, as in \cref{hyp:geo-smooth-tier}, differentiable in quadratic mean
with a positive-definite Fisher form; every score direction is integrable to the third
power, the second- and third-order score moments depend smoothly on the point, and the
model is dominated enough to differentiate the divergences of
\cref{sec:pb-divergence-jets} under the integral through third order.  The hypothesis
excludes singular statistical strata, families whose support moves without this
differentiability, and law fibers given only by a chart without a sample-space action.

\definitionheading{Fiber statistical tensors and the model-parameter Fisher tensor}{def:pb-fiber-statistical-tensors}
With $\mathbf s_Z$ the score of $\mathcal B_x$ at $p$ in the direction
$Z\in T_p\mathcal B_x$, the Fisher and Amari--Chentsov tensors of the channel and, for
the generative parameter $\theta$, structural data $X$, tangent vectors $Z_1,Z_2$ of the
parameter space of $\theta$, and a probability law $\mathsf r(dX)$ on structural data,
the model-parameter Fisher tensor of the evidence $p_{\theta,V}(o\given X)$ and its
average are
\begin{align}
g^F_{x,p}(Z_1,Z_2)&=\E_p[\mathbf s_{Z_1}\mathbf s_{Z_2}],
\label{eq:pb-fiber-fisher}\\
T^F_{x,p}(Z_1,Z_2,Z_3)&=\E_p[\mathbf s_{Z_1}\mathbf s_{Z_2}\mathbf s_{Z_3}],
\label{eq:pb-fiber-amari}\\
\mathsf I_X(\theta)(Z_1,Z_2)&=\E_\theta\left[\partial_{Z_1}\log p_{\theta,V}(O\given X)
\partial_{Z_2}\log p_{\theta,V}(O\given X)\mathrel{\big|}X\right],
\label{eq:ig-conditional-model-fisher}\\
\mathsf I_{\mathsf r}(\theta)&=\int\mathsf I_X(\theta)\mathsf r(dX).
\label{eq:ig-averaged-model-fisher}
\end{align}
Each of $\mathsf I_X$ and $\mathsf I_{\mathsf r}$ is called a metric only where it is
nondegenerate.

On the latent sample space $\mathsf Y$ an interaction energy is a function, not a metric;
the law fibers and the configuration manifolds of \cref{sec:hist-vfe-orbit-selection}
carry $g^F$; the parameter space of $\theta$ carries $\mathsf I_X$ or
$\mathsf I_{\mathsf r}$, of rank depending on $X$ or $\mathsf r$, one available metric for
$\theta$ in the fast and slow coordinates of \cref{def:obs-fast-slow-blocks}, whose
model-channel recognition coordinates carry $g^F_m$.  Equal components of a quadratic
energy on $\mathsf Y$ and a Fisher metric, as in the Gaussian mean sector of
Supplement~G, are a coincidence of coordinates, not an identification.

\propositionheading{Descent of the statistical tensors}{prop:pb-statistical-tensor-descent}
Under \cref{hyp:pb-regular-models}, every $\widehat\rho_x(g)$ preserves $g^F_x$ and
$T^F_x$, so $(\iota_u^{-1})^*g^F_x$ and $(\iota_u^{-1})^*T^F_x$ are tensors on
$(\mathcal E_x)_{c_*}$ that do not depend on the frame $u$; they are denoted by the same
symbols.  The sample-space origin of the action cannot be replaced by invariance of the
metric alone: for the normal family, whose Fisher metric is
$(d\mu^2+2d\varsigma^2)/\varsigma^2$, the map induced by $z\mapsto(1+z)/(1-z)$ on
$z=\mu/\sqrt2+\mathrm i\varsigma$ preserves $g^F$, fixes $(\mu,\varsigma)=(0,1)$, and
does not preserve $T^F$.  \status{ESTABLISHED}

The proof is in \cref{sec:proofs-fiber-geometry}.

\propositionheading{Invariance characterizations of the Fisher metric}{prop:ig-fisher-characterizations}
On the open simplices of positive laws on all finite sets, the family of Fisher
metrics is, up to one positive factor, the unique family of continuous Riemannian
metrics, one on each simplex, invariant under the congruent Markov embeddings between
simplices \citep[Theorem~11.1]{Cencov1982}, restated in
\citet[Proposition~3.19(3)]{Ay2015}; the uniqueness is a property of the family,
since the embeddings relate simplices of different dimensions.  On the cone of positive finite measures on a
finite set, with coordinates $\mathsf x_i$ and total mass $|\mathsf x|=\sum_i\mathsf x_i$,
the positive-definite metrics invariant under the same maps are exactly those with
components $g_{\mathsf x}(\partial_i,\partial_j)=\mathrm f_2(|\mathsf x|)
+\delta_{ij}|\mathsf x|\mathrm f_1(|\mathsf x|)/\mathsf x_i$ for smooth functions with
$\mathrm f_1>0$ and $\mathrm f_1+\mathrm f_2>0$ \citep{Campbell1986}; on tangent vectors
of the simplex, where $\sum_iZ_i=0$, the $\mathrm f_2$ term vanishes and the metric is
$\mathrm f_1(1)$ times the Fisher metric.  For parametrized measure models on general
sample spaces, every local continuous quadratic form field invariant under sufficient
statistics is a combination of the Fisher form and the square of the differential of
total mass with coefficients that are functions of the total mass, and on statistical
models, where the total mass is one, the Fisher form is the unique weakly continuous such
field up to a constant \citep[Main Theorem~2.10 and Remark~3.23]{Ay2015}; the finite
case of this characterization is Campbell's \citep[Section~1]{Le2017}.  Uniqueness up
to a constant also holds, on separable metrizable sample spaces with their Borel
sigma-algebras, for strongly continuous forms that are monotone under statistics
\citep[Theorem~1.4]{Le2017}, for diffeomorphism-invariant smooth weak Riemannian metrics
on smooth positive probability densities on a closed manifold of dimension at least two
\citep[Theorem in the introduction]{BauerBruverisMichor2016}, whose invariant bilinear
forms on all smooth positive densities are the combinations of the Fisher--Rao form and
the product of the total masses of the two tangent densities with constant
coefficients \citep[Main Theorem]{BauerBruverisMichor2016}, and, on exponential
families and their derived families, for Riemannian metrics invariant under independent
and identically distributed extensions and canonical sufficient statistics
\citep[Theorem~1]{Dowty2018}.  \status{STANDARD}

\section{Divergence jets and the orbit Fisher form}
\label{sec:pb-divergence-jets}

\definitionheading{Regular invariant contrast and its jets}{hyp:pb-contrast}
A regular invariant contrast on a law fiber $\mathcal B$ is a function
$\mathscr D:\mathcal B\times\mathcal B\to\R$, $C^3$ near the diagonal, with
$\mathscr D\geq0$ and, near the diagonal, $\mathscr D(p,q)=0$ exactly when $p=q$,
invariant under the simultaneous represented action,
$\mathscr D(\widehat\rho(g)p,\widehat\rho(g)q)=\mathscr D(p,q)$, and with
positive-definite metric jet $g_{\mathscr D}$ of \eqref{eq:pb-divergence-metric-jet}.  With $(Z,0)$ and $(0,Z)$ the lifts of
$Z\in T_p\mathcal B$ to the two factors, extended as constant coordinate fields in one
chart, the jets of $\mathscr D$ are
\begin{align}
g_{\mathscr D}(Z_1,Z_2)
&=-[(Z_1,0)(0,Z_2)\mathscr D]_{(p,p)},
\label{eq:pb-divergence-metric-jet}\\
\Gamma_{\mathscr D}(Z_1,Z_2,Z_3)
&=-[(Z_1,0)(Z_2,0)(0,Z_3)\mathscr D]_{(p,p)},
\label{eq:pb-divergence-first-connection-jet}\\
\Gamma_{\mathscr D}^*(Z_1,Z_2,Z_3)
&=-[(0,Z_1)(0,Z_2)(Z_3,0)\mathscr D]_{(p,p)},
\label{eq:pb-divergence-dual-connection-jet}\\
T^F_{\mathscr D}
&=\Gamma_{\mathscr D}-\Gamma_{\mathscr D}^*.
\label{eq:pb-divergence-amari-jet}
\end{align}
The two third-order arrays are connection coefficients in the chart, their difference is
a symmetric cubic tensor. Invariance makes the metric and cubic jets define tensors,
and the connection jets define affine connections, on the associated fiber
independently of the frame.  The contrast $(p,q)\mapsto(p-q)^4$
on a one-dimensional fiber is positive off the diagonal with zero second jet, so it is
excluded.

\propositionheading{KL divergence jets}{prop:pb-kl-divergence-jets}
On a fiber satisfying \cref{hyp:pb-regular-models}, the relative entropy
$\KL(p\Vert q)$ is a regular invariant contrast, its metric jet is $g^F$, and its cubic
jet $T^F_{\KL}$ is the Amari--Chentsov tensor $T^F$.  \status{ESTABLISHED}

\corollaryheading{Orbit Fisher form}{cor:pb-orbit-fisher-form}
Assume \cref{hyp:pb-regular-models} in one channel.  For $q\in\mathcal B$ define
\begin{equation}
\mathfrak k_q(\xi,\eta)=g^F_q\bigl(\xi_{\mathcal B}(q),\eta_{\mathcal B}(q)\bigr),
\qquad
\mathfrak g_q=\{\xi\in\mathfrak g:\xi_{\mathcal B}(q)=0\}.
\label{eq:pb-orbit-fisher-form}
\end{equation}
Then $\mathfrak k_q$ is positive semidefinite, its radical is exactly $\mathfrak g_q$,
and $\mathfrak g_q$ is the Lie algebra of the stabilizer $G_q$ of $q$.  For every
$\xi\in\mathfrak g$, as $t\to0$,
\begin{equation}
\KL\left(q\Vert\widehat\rho(\exp t\xi)q\right)
=\frac{t^2}{2}\mathfrak k_q(\xi,\xi)+O(t^3),
\label{eq:pb-orbit-kl-expansion}
\end{equation}
and for $\xi\in\mathfrak g_q$ the left side vanishes for every $t$.  The same expansion
holds for every regular invariant contrast, with $g^F$ replaced by its metric jet
$g_{\mathscr D}$.  \status{ESTABLISHED}

The proofs of \cref{prop:pb-kl-divergence-jets,cor:pb-orbit-fisher-form} are in
\cref{sec:proofs-fiber-geometry}.

In the running example of \cref{sec:geo-running-example}, a loop holonomy $\exp t\xi$
acting on an agent's law gives the loop a transported divergence
$\tfrac{t^2}{2}\mathfrak k_q(\xi,\xi)+O(t^3)$ that vanishes to all orders exactly when
$\xi\in\mathfrak g_q$, the second-order form of law coherence up to holonomy stabilizers
(\cref{def:geo-law-coherence}).

With the link as a variable, the orbit Fisher form is the Hessian of the transported divergence at an
exact link, and the minimizers of that divergence are unions of double cosets of the stabilizers, hence
not isolated when a stabilizer has positive dimension (\cref{prop:link-law-level-energy}).

\section{Markov coarse morphisms and the vertical Fisher defect}
\label{sec:pb-fisher-defect}
\label{sec:cg-fisher-contraction}

\theoremheading{Score projection and Fisher loss}{thm:cg-fisher-contraction}
Let $\{P_\lambda\}$ be a family of laws on $\mathsf Y$, indexed by an open subset of
$\R^d$, that is differentiable in quadratic mean at $\lambda_0$ with centered
square-integrable score $\mathbf s_{\lambda_0}:\mathsf Y\to\R^d$, and let
$K:\mathsf Y\rightsquigarrow\bar{\mathsf Y}$ be a Markov kernel that does not depend on
$\lambda$ and satisfies $K(y,\bar{\mathsf Y})=1$ for every $y$.  Under the joint law
$P_{\lambda_0}\otimes K$, the coarse family $\{P_\lambda K\}$ is differentiable in
quadratic mean at $\lambda_0$ with score
\begin{equation}
\bar{\mathbf s}_{\lambda_0}(\bar Y)
=\E\left[\mathbf s_{\lambda_0}(Y)\mid\bar Y\right]
\quad P_{\lambda_0}K\text{-almost surely},
\label{eq:cg-score-projection}
\end{equation}
and its Fisher information obeys
\begin{equation}
\mathsf I_{\bar{\mathsf Y}}(\lambda_0)\preceq\mathsf I_{\mathsf Y}(\lambda_0),
\qquad
\mathsf I_{\mathsf Y}(\lambda_0)-\mathsf I_{\bar{\mathsf Y}}(\lambda_0)
=\E\operatorname{Cov}\left(\mathbf s_{\lambda_0}(Y)\mid\bar Y\right)\succeq0.
\label{eq:cg-fisher-loss}
\end{equation}
Equality in a direction $z\in\R^d$ holds exactly when $z^\top\mathbf s_{\lambda_0}(Y)$
is $\bar Y$-measurable.  No common dominating measure is assumed.
\status{ESTABLISHED}

The proof is in \cref{sec:proofs-fiber-geometry}.

\remarkheading{Fisher equality is not recovery}{rem:cg-fisher-equality-not-recovery}
Equality in \eqref{eq:cg-fisher-loss} at one parameter value is local score
sufficiency, not recovery of the experiment.  Let $Y=(Y_1,Y_2)$ have independent
Bernoulli coordinates with success probabilities $\tfrac12+\tfrac\lambda4$ and
$\tfrac12+\tfrac{\lambda^2}4$, and let $K$ discard $Y_2$.  At $\lambda=0$ the score of
$Y_2$ vanishes, so the fine and coarse Fisher informations both equal $\tfrac14$, yet no
parameter-independent kernel $K'$ satisfies $(P_\lambda K)K'=P_\lambda$ for every
$\lambda$ in a neighborhood of zero: applying any fixed $K'$ to the
Bernoulli law $P_\lambda K$ makes every output event probability affine in
$\lambda$, whereas the required output event $Y_2=1$ has probability
$\tfrac12+\tfrac{\lambda^2}4$.  \citet{KaganShepp2005} give an insufficient statistic that
preserves Fisher information at every parameter value, and
\citet[Theorem~7]{Pollard2013} explains such examples by the criterion of
\cref{thm:cg-fisher-contraction}: preservation at $\lambda$ holds exactly when the score
is almost surely a function of the statistic.  \status{ESTABLISHED}

At $c_*$ the theorem applies to maps between fibers, for instance from the joint fiber
of a block of agents to the fiber of their meta-agent.

\definitionheading{Markov coarse morphism and its vertical defects}{hyp:pb-markov-morphism}
Let $\bar{\mathcal B}_x$ be a second law fiber in channel $x$, with represented action
$\widehat{\bar\rho}_x$ of $G$, associated bundle $\bar{\mathcal E}_x$, and frame maps
$\bar\iota_u$.  A Markov coarse morphism in channel $x$ at $c_*$ is a map
$\mathscr K:(\mathcal E_x)_{c_*}\to(\bar{\mathcal E}_x)_{c_*}$ with
$\bar\iota_u^{-1}\circ\mathscr K\circ\iota_u=K_\star|_{\mathcal B_x}$ for one, and hence
every, frame $u$, where $K:\mathsf Y_x\rightsquigarrow\bar{\mathsf Y}_x$ is a Markov
kernel that does not depend on the law parameter, satisfies $K(y,\bar{\mathsf Y}_x)=1$
for every $y$, maps $\mathcal B_x$ smoothly into $\bar{\mathcal B}_x$, and is
equivariant, $K_\star\circ\widehat\rho_x(g)=\widehat{\bar\rho}_x(g)\circ K_\star$ for
every $g$; no common dominating measure is assumed.  The vertical Fisher defect and the
Amari defect of $\mathscr K$ are
\begin{align}
\Delta_F^{\mathscr K}
&=g^F_x-(T\mathscr K)^*\bar g^F_x,
\label{eq:pb-vertical-fisher-defect}\\
\Delta_T^{\mathscr K}
&=T^F_x-(T\mathscr K)^*\bar T^F_x.
\label{eq:pb-amari-defect}
\end{align}
Equivariance makes the frame representative independent of $u$, since
$\bar\iota_{u\cdot g}^{-1}\mathscr K\iota_{u\cdot g}
=\widehat{\bar\rho}_x(g)^{-1}K_\star\widehat\rho_x(g)$, and by
\cref{prop:pb-statistical-tensor-descent} both defects are tensors on
$(\mathcal E_x)_{c_*}$.

Restriction to a submodel, precomposition of an energy with a map, a kernel that depends
on the law parameter, a fitted coarse approximation, and a comparison across a mismatch
of generative and recognition parameters are Markov coarse morphisms, and fall under the
results below, only when one kernel that realizes them is exhibited, along the whole history when histories are
compared.  One kernel serves every law on one fixed-precision fiber of a Gaussian
family; Supplement~G shows that Gaussian Galerkin aggregation has this form on each
fixed-precision fiber and no such form across precisions.

\corollaryheading{Vertical Fisher defect}{thm:pb-pullback-fisher-defect}
For a Markov coarse morphism $\mathscr K$ and every tangent vector $Z$ at $p$, under the
joint law $p\otimes K$,
\begin{align}
&\Delta_F^{\mathscr K}\succeq0,
\label{eq:pb-fisher-defect-positive}\\
&\Delta_F^{\mathscr K}(Z,Z)
=\E\operatorname{Var}\left(\mathbf s_Z(Y)\mid\bar Y\right),
\label{eq:pb-fisher-defect-score-variance}
\end{align}
and $\Delta_F^{\mathscr K}(Z,Z)=0$ exactly when $\mathbf s_Z(Y)$ is $\bar Y$-measurable.
\status{ESTABLISHED}

\noindent\emph{Proof.}
In a frame $\mathscr K$ is $K_\star$.  For a smooth curve $p_t$ in $\mathcal B_x$ with
velocity $Z$ at $p$, \cref{thm:cg-fisher-contraction} gives the image curve
$K_\star p_t$ the score $\E[\mathbf s_Z(Y)\mid\bar Y]$, which by uniqueness of the
quadratic-mean score is the score of $\bar{\mathcal B}_x$ in the direction $TK_\star Z$,
with second moment $\bar g^F_x(TK_\star Z,TK_\star Z)$; the defect identity \eqref{eq:cg-fisher-loss} for this one-parameter family is
\eqref{eq:pb-fisher-defect-score-variance}, with the stated equality clause.  $\square$

\theoremheading{Vertical Fisher defect cocycle}{thm:pb-fisher-defect-cocycle}
Let $\mathscr K_{01}:(\mathcal E_0)_{c_*}\to(\mathcal E_1)_{c_*}$ and
$\mathscr K_{12}:(\mathcal E_1)_{c_*}\to(\mathcal E_2)_{c_*}$ be smooth maps between
fibers carrying Fisher tensors $g^F_0,g^F_1,g^F_2$.  Then, unconditionally,
\begin{equation}
\Delta_F^{\mathscr K_{12}\circ\mathscr K_{01}}
=\Delta_F^{\mathscr K_{01}}
+(T\mathscr K_{01})^*\Delta_F^{\mathscr K_{12}},
\label{eq:pb-fisher-defect-cocycle}
\end{equation}
and the Amari defects obey the same identity.  If both maps are Markov coarse morphisms,
both terms of \eqref{eq:pb-fisher-defect-cocycle} are positive semidefinite.  No nonzero
Amari defect is sign definite, since
$\Delta_T^{\mathscr K}(-Z,-Z,-Z)=-\Delta_T^{\mathscr K}(Z,Z,Z)$; Markov coarse-graining
therefore gives no monotonicity for $T^F$, and $\Delta_T^{\mathscr K}=0$ is an additional
property of $\mathscr K$.  \status{ESTABLISHED}

\noindent\emph{Proof.}
Contravariance of pullback gives
$\Delta_F^{\mathscr K_{12}\circ\mathscr K_{01}}
=g_0^F-(T\mathscr K_{01})^*g_1^F+(T\mathscr K_{01})^*[g_1^F-(T\mathscr K_{12})^*g_2^F]$,
which is \eqref{eq:pb-fisher-defect-cocycle}, and the same computation applies to $T^F$.
Positivity is \cref{thm:pb-pullback-fisher-defect} for each Markov arrow, and oddness is
homogeneity of degree three.  $\square$

\section{Configuration manifolds and the exact free-energy lift}
\label{sec:hist-vfe-orbit-selection}

\definitionheading{Configuration manifold at $c_*$}{def:hist-configuration-manifold}
The configuration manifold $\mathcal Q_i$ of agent $i$ at $c_*$ is the product
$(\mathcal E_b)_{c_*}\times(\mathcal E_m)_{c_*}$ of its belief and model law fibers,
identified with $\mathcal B_b\times\mathcal B_m$ through its frame $u_i$, or an embedded
submanifold of that product (a regular submodel), with points $Q_i=(q_i^b,q_i^m)$.  The
collective configuration manifold is $\mathcal Q_V=\prod_{i\in V}\mathcal Q_i$, extended
by further coordinate manifolds, such as open simplices of attention rows, when those
coordinates move.  A configuration manifold carries a metric $\mathsf G$ that must be
specified.  The channel-weighted metric of agent $i$ is
$\mathsf G_i=w_bg^F_b\oplus w_mg^F_m$ with positive weights, under which a $C^1$ curve
$Q_i(t)$ has squared speed
\begin{equation}
\lVert\dot Q_i\rVert_{\mathsf G_i}^2
=w_bg^F_b(\dot q_i^b,\dot q_i^b)+w_mg^F_m(\dot q_i^m,\dot q_i^m);
\label{eq:hist-pointwise-clock-speed}
\end{equation}
on $\mathcal Q_V$ the metric is either block diagonal, $\bigoplus_i\mathsf G_i$, or the
pullback of a joint Fisher metric along an exact lift (\cref{hyp:hist-exact-vfe-lift}).

The channel-weighted and block-diagonal metrics are frame independent by
\cref{prop:pb-statistical-tensor-descent}. The pullback along an exact lift is frame
independent when the lift is equivariant in the sense of \cref{hyp:hist-exact-vfe-lift},
since then $\widehat\rho(g)^*\iota_i^*\mathsf G_{\mathscr R_B}
=\iota_i^*\widehat\rho_{\mathscr R}(g)^*\mathsf G_{\mathscr R_B}=\iota_i^*\mathsf G_{\mathscr R_B}$,
and is otherwise frame relative: it depends on the choice of right inverse
(\cref{prop:hist-general-score-residual}), and equivariant correlated lifts under a
connected group on continuous fibers are \cref{open:hist-parity-lift-scope}. The
base-direction version of these metrics is in Supplement~B.

\theoremheading{Fisher radical and quotient of a pulled-back tensor}{thm:hist-retained-fisher-quotient}
Let $\operatorname{ret}:\mathcal Q\to\mathcal Q'$ be a smooth map of finite-dimensional
manifolds, $g$ a positive-semidefinite symmetric tensor on $\mathcal Q'$, and
$h=\operatorname{ret}^*g$.  Then
\begin{equation}
\operatorname{rad}h_Q
=(d\operatorname{ret}_Q)^{-1}\left(\operatorname{rad}g_{\operatorname{ret}(Q)}\right),
\label{eq:hist-cds-fisher-radical}
\end{equation}
so $\operatorname{rad}h_Q=\ker d\operatorname{ret}_Q$ exactly when
$\operatorname{im}d\operatorname{ret}_Q\cap\operatorname{rad}g_{\operatorname{ret}(Q)}
=\{0\}$.  Where $\operatorname{rad}h$ has constant rank and $h$ is $C^1$,
$\operatorname{rad}h$ is a $C^1$ subbundle and $h$ induces a positive-definite form on
the quotient bundle $T\mathcal Q/\operatorname{rad}h$.  A quotient manifold additionally
requires the null leaves to form a Hausdorff manifold leaf space; when
$\operatorname{rad}h=\ker d\operatorname{ret}$ and the connected leaves form such a space
$\mathcal Q/{\sim}$ with quotient projection a smooth submersion, $h$ is the pullback of a positive-definite tensor on
$\mathcal Q/{\sim}$.  For a positive $C^3$ family on a finite set, with
$\operatorname{ret}$ the law map into the open simplex and $g$ its Fisher metric, $h$ is
the Fisher tensor of the family, of class $C^2$, positive definite exactly where
$\operatorname{ret}$ is an immersion, and there $\operatorname{ret}$ is locally
injective and $h$ has a $C^1$ Levi-Civita connection.  For softmax logits $z\in\R^k$
with $p=\operatorname{softmax}(z)$, $h(z)=\operatorname{diag}(p)-pp^\top$ has radical
$\R\mathbf 1$ at every $z$, and for an objective that depends on $z$ only through $p$
the solutions $Z$ of $hZ=\nabla\Fenergy$ form one class $Z_0+\R\mathbf 1$, whose
Moore--Penrose representative $h^+\nabla\Fenergy$ is Euclidean-orthogonal to
$\mathbf 1$.  \status{ESTABLISHED}

The proof is in \cref{sec:proofs-fiber-geometry}.

\openproblemheading{Global Fisher quotients and canonical block data}{open:hist-retained-fisher-quotient}
\Cref{thm:hist-retained-fisher-quotient} gives a positive-definite quotient bundle on
every constant-rank stratum and a quotient manifold once a Hausdorff leaf space is
supplied, and \cref{prop:hist-law-fisher-agentization-nogo} shows that the law and its
Fisher tensor alone select no splitting into agent blocks.  Open are a criterion,
checkable on the configuration manifolds of this chapter, under which the null foliation
of a pulled-back Fisher tensor has a Hausdorff manifold leaf space, and block data on
that quotient determined by a stated typing of agents and channels.  Such a criterion
with a proof, for instance exhibiting the leaves as the orbits of a free and proper group
action, together with a construction of natural block data from the typing or a
counterexample showing that the typing does not determine it, settles the problem.

\hypothesisheading{Standing configuration tier}{hyp:hist-standing-configuration-tier}
\label{hyp:hist-regular-metric-domain}%
Wherever this chapter asserts a natural-gradient vector field or a Fisher duration on a
configuration manifold $\mathcal Q$, either $\mathcal Q$ is a finite-dimensional
configuration manifold as in \cref{def:hist-configuration-manifold}, a finite product of
regular statistical manifolds or an embedded submodel of one, with a $C^1$
positive-definite metric $\mathsf G$, or $\mathcal Q$ is the quotient of such a manifold
by the null leaves of a $C^1$ positive-semidefinite Fisher tensor of constant rank with a
Hausdorff manifold leaf space, as in \cref{thm:hist-retained-fisher-quotient}, and the
objective is constant on those leaves.  The objective $\Fenergy$ is $C^2$, and every
minimal regular exponential family with its Fisher metric is an instance.  The hypothesis
excludes spaces of sections over a base, which belong to Supplement~B, semidefinite
metrics without a constructed quotient, and configuration metrics called exact joint-law
Fisher metrics without the lift of \cref{hyp:hist-exact-vfe-lift}.

\definitionheading{Exact multi-agent free-energy lift}{hyp:hist-exact-vfe-lift}
Fix a block $B$ of latent coordinates, an outside recognition marginal
$\mathbb Q_{B^c}$, and a manifold $\mathscr R_B$ of admissible conditional recognition
kernels $r_B(dy_B\given y_{B^c})$.  An exact lift of the configuration manifold
$\mathcal Q_i$ consists of a smooth extraction map
$\operatorname{ext}_i:\mathscr R_B\to\mathcal Q_i$ and a smooth right inverse
$\iota_i:\mathcal Q_i\to\mathscr R_B$, $Q_i\mapsto r_B^{Q_i}$,
$\operatorname{ext}_i\circ\iota_i=\operatorname{id}_{\mathcal Q_i}$, such that, for every
$Q_i$, the support, finiteness, and integrability hypotheses of
\cref{thm:obs-local-multiagent-elbo,thm:obs-local-global-potential} hold for
$\mathbb Q_{B^c}r_B^{Q_i}$, the lifted family is $C^2$, and for every chart
$\operatorname{ch}$ of $\mathcal Q_i$ with coordinates $\lambda$ and every compact set
$\mathsf C$ in its image there is $\mathrm e_{\mathsf C}\in L^1(\mathbb Q_{B^c})$ with,
for jointly measurable versions of the derivatives,
\begin{equation}
\sup_{\lambda\in\mathsf C}
\bigl|\partial_\lambda^\alpha\Fenergy^{\rm ext}_{B,o}
(r_B^{\operatorname{ch}^{-1}(\lambda)};y_{B^c})\bigr|
\leq\mathrm e_{\mathsf C}(y_{B^c})
\quad\text{for }\mathbb Q_{B^c}\text{-almost every }y_{B^c}
\text{ and }|\alpha|\leq2.
\label{eq:hist-chartwise-vfe-envelope}
\end{equation}
The lifted objective is the outside-averaged local free energy
\begin{equation}
\overline{\Fenergy}_{B,o}(Q_i)
=\E_{\mathbb Q_{B^c}}\left[\Fenergy^{\rm ext}_{B,o}\left(r_B^{Q_i};Y_{B^c}\right)\right],
\label{eq:hist-outside-averaged-local-vfe}
\end{equation}
a history of it is an exact free-energy history, and the exact recognition Fisher metric
of $\mathcal Q_i$ is $\mathsf G_i=\iota_i^*\mathsf G_{\mathscr R_B}$ on a nondegenerate
stratum, or its quotient as in \cref{thm:hist-retained-fisher-quotient}. The lift is
equivariant when $\mathscr R_B$ carries a declared represented action
$\widehat\rho_{\mathscr R}$ of the structure group, per agent or diagonal and compatible
with the frames of the outside marginal $\mathbb Q_{B^c}$, that preserves
$\mathsf G_{\mathscr R_B}$ and satisfies
$\operatorname{ext}_i\circ\widehat\rho_{\mathscr R}(g)=\widehat\rho(g)\circ\operatorname{ext}_i$
and $\iota_i\circ\widehat\rho(g)=\widehat\rho_{\mathscr R}(g)\circ\iota_i$; without such
an action, $\mathsf G_i$ is defined relative to the frames in which $\iota_i$ is written.

The lift carries every correlation that the configuration does not determine, and
$\operatorname{ext}_i\circ\iota_i=\operatorname{id}$ makes it a lift rather than a family
of laws indexed by configurations.  With the outside marginal held fixed,
\eqref{eq:hist-outside-averaged-local-vfe} differs from the collective free energy along
$Q_i\mapsto\mathbb Q_{B^c}r_B^{Q_i}$ by a term independent of $Q_i$
(\cref{thm:obs-local-global-potential}), so the two have the same differential and, for
one metric, the same natural-gradient ray.

\propositionheading{Joint-versus-weighted-marginal score residual and lift dependence}{prop:hist-general-score-residual}
Let $P_\lambda$ be a regular dominated family of joint laws of blocks
$(Y_1,\dots,Y_n)$ with square-integrable joint score
$\mathbf s=\nabla_\lambda\log p_\lambda(Y)$, embedded marginal scores
$\mathbf s_i=\E[\mathbf s\mid Y_i]$, and scalar weights $w_i$, and define
\begin{align}
\mathsf G_{\mathrm J}&=\E[\mathbf s\mathbf s^\top],
&\mathsf G_i&=\E[\mathbf s_i\mathbf s_i^\top],
&\mathsf G_w&=\sum_iw_i\mathsf G_i,
\label{eq:hist-score-metrics}\\
\mathbf r&=\mathbf s-\sum_i\mathbf s_i,
&\mathfrak K&=\E[\mathbf r\mathbf r^\top],
&\mathfrak C&=\sum_{i\ne j}\E[\mathbf s_i\mathbf s_j^\top].
\label{eq:hist-score-residual-tensors}
\end{align}
The cross tensor $\mathfrak C$ is symmetric and can be indefinite,
\begin{equation}
\mathsf G_{\mathrm J}-\mathsf G_w
=\mathfrak K-\mathfrak C+\sum_i(1-w_i)\mathsf G_i,
\label{eq:hist-joint-weighted-marginal-residual}
\end{equation}
and on a tangent subspace $\mathcal T$, $\mathsf G_{\mathrm J}=\mathsf G_w$ holds exactly
when the right side of \eqref{eq:hist-joint-weighted-marginal-residual} vanishes on
$\mathcal T\times\mathcal T$.  For two blocks and unit weights, in a direction $Z$ with
scores $\mathbf s_Z=Z^\top\mathbf s$ and $\mathbf s_{i,Z}=Z^\top\mathbf s_i$,
\begin{equation}
\lVert\mathbf s_Z\rVert^2
-\lVert\mathbf s_{1,Z}\rVert^2-\lVert\mathbf s_{2,Z}\rVert^2
=\lVert\mathbf s_Z-\mathbf s_{1,Z}-\mathbf s_{2,Z}\rVert^2
-2\langle\mathbf s_{1,Z},\mathbf s_{2,Z}\rangle,
\label{eq:hist-joint-versus-product}
\end{equation}
whose first term is nonnegative and whose second is signed.  If the two blocks are
independent under every $P_\lambda$, then $\mathsf G_{\mathrm J}=\mathsf G_1+\mathsf G_2$,
and $\mathsf G_{\mathrm J}=w_1\mathsf G_1+w_2\mathsf G_2$ forces $w_1=1$ as soon as some
direction has $\mathsf G_1(Z,Z)>0=\mathsf G_2(Z,Z)$, and symmetrically for $w_2$.  No
Loewner order between $\mathsf G_{\mathrm J}$ and $\mathsf G_1+\mathsf G_2$ holds in
general: for the bivariate normal location family with precision
$\begin{psmallmatrix}1&\mathrm c\\\mathrm c&1\end{psmallmatrix}$ and $0<\mathrm c<1$, the
difference is
$\begin{psmallmatrix}\mathrm c^2&\mathrm c\\\mathrm c&\mathrm c^2\end{psmallmatrix}$, with
value $2\mathrm c(1+\mathrm c)>0$ along $(1,1)$ and $2\mathrm c(\mathrm c-1)<0$ along
$(1,-1)$.  Distinct right inverses of one extraction map can induce distinct pullback
metrics: for $\mathcal N((\lambda_1,\lambda_2),\mathrm I_2)$, with Fisher metric
$d\lambda_1^2+d\lambda_2^2$, and the extraction $(\lambda_1,\lambda_2)\mapsto\lambda_1$,
the lifts $\lambda_1\mapsto(\lambda_1,0)$ and $\lambda_1\mapsto(\lambda_1,1)$ induce
$d\lambda_1^2$, whereas $\lambda_1\mapsto(\lambda_1,\lambda_1)$ induces $2d\lambda_1^2$.
\status{ESTABLISHED}

The proof is in \cref{sec:proofs-fiber-geometry}.

\subsection{Canonicity boundary of the lift}
\label{sec:hist-canonical-dependence-selection}

A lift supplies the dependence that the marginal configurations do not determine.  The
following finite results show that naturality under local kernels forces the product
law and that admitting correlated refinements leaves no natural selector.

\definitionheading{Finite marginal selection problem}{def:hist-finite-marginal-selection}
For a finite nonempty typed list $\mathcal X=(\mathcal X_i)_{i=1}^n$ of finite sets, let
$\mathcal J_{\mathcal X}=\Delta(\prod_i\mathcal X_i)$,
$\mathcal J^1_{\mathcal X}=\prod_i\Delta(\mathcal X_i)$, and
$\operatorname{mar}_{\mathcal X}:\mathcal J_{\mathcal X}\to\mathcal J^1_{\mathcal X}$ the
map that extracts the one-coordinate marginals.  A selector is a map
$\operatorname{sel}_{\mathcal X}:\mathcal J^1_{\mathcal X}\to\mathcal J_{\mathcal X}$ with
$\operatorname{mar}_{\mathcal X}\circ\operatorname{sel}_{\mathcal X}=\operatorname{id}$.
A local morphism is a tensor product $K=\bigotimes_iK_i$ of finite Markov kernels of the
same arity, and a kernel $K:\prod_i\mathcal X_i\rightsquigarrow\prod_j\mathcal Y_j$ that
changes arity is marginal compatible when
$\operatorname{mar}_{\mathcal Y}K_\star=\bar K\operatorname{mar}_{\mathcal X}$ for some
map $\bar K:\mathcal J^1_{\mathcal X}\to\mathcal J^1_{\mathcal Y}$.  A family of
selectors is natural on a class of morphisms when
$K_\star\operatorname{sel}_{\mathcal X}=\operatorname{sel}_{\mathcal Y}\bar K$ for every
morphism in the class; the class is part of the data, and not every finite kernel is
admitted.

\theoremheading{Tensor-local naturality selects the product law}{thm:hist-local-markov-product-uniqueness}
\label{cor:hist-maximal-product-preserving-category}%
If a family of selectors is natural under every local morphism, then
$\operatorname{sel}_{\mathcal X}(\mu_1,\ldots,\mu_n)=\bigotimes_{i=1}^n\mu_i$.  Among
classes of marginal-compatible kernels that contain every local morphism, a natural
selector exists exactly when every admitted kernel sends product laws to the product of
their image marginals, and then it is the product selector; the product-preserving
marginal-compatible kernels therefore form the largest such class.  \status{ESTABLISHED}

\noindent\emph{Proof.}
For each $i$ let the kernel from a one-point set send its point to $\mu_i$.  Their
tensor product sends the unique law on the one-point product to $\bigotimes_i\mu_i$,
while the induced marginal map sends the unique marginal tuple to $(\mu_i)_i$, so
naturality forces the product selector.  With it forced, naturality for an added kernel
says exactly that it sends each product law to the product of its image marginals.
$\square$

\theoremheading{Admitted correlated refinements refute an absolute selector}{thm:hist-correlated-refinement-nogo}
For a source bit $y$ and target bits $(y',y'')$, admit the split kernels
$K^{\mathrm{split}}_{\mathrm c}(y',y''\mid y)=\mathbf 1_{\{y'=y\}}(1+\mathrm c)/2$ if
$y''=y$ and $\mathbf 1_{\{y'=y\}}(1-\mathrm c)/2$ if $y''\ne y$, for
$\mathrm c\in\{1/3,1/2\}$.  Both are marginal compatible, and no single-valued selector
is natural under both; the argument uses no local naturality.  In particular
marginalization is not injective on positive laws, so no selector returns every joint
law from its marginals.  \status{ESTABLISHED}

\noindent\emph{Proof.}
The two kernels push the fair source law to laws with fair marginals and atom multisets
$\{1/3,1/3,1/6,1/6\}$ and $\{3/8,3/8,1/8,1/8\}$.  Naturality would require the value of
the target selector at the fair marginal pair to equal both laws, which differ even after
relabeling, and the two positive laws share their marginals.  $\square$

Selection relative to a positive reference law is a different problem. A
finite-relative-entropy minimizer over a convex family of moment-constrained laws,
when it exists, is unique by strict convexity of relative entropy. Both that minimizer
and, for a deterministic coarse map, the reference conditional completion of a
prescribed coarse marginal (\cref{lem:prob-kl-chain-rule}) depend on the reference
law as additional data.

\section{A finite correlated lift: the parity witness}
\label{sec:hist-finite-parity-lift}

\definitionheading{Three-agent parity lift}{def:hist-three-agent-parity-lift}
Let $\mathcal I_6=\{1,2,3\}\times\{b,m\}$, $y=(y_{i,x})\in\{0,1\}^{\mathcal I_6}$, and
$\lambda=(\lambda_{i,x})\in(0,1)^{\mathcal I_6}$, where typed agent permutations preserve
the channel label, and set
\begin{align}
&p_0(y;\lambda)
=\prod_{(i,x)}\lambda_{i,x}^{y_{i,x}}(1-\lambda_{i,x})^{1-y_{i,x}},
\qquad
\mathrm D(\lambda)=\prod_{(i,x)}\lambda_{i,x}(1-\lambda_{i,x}),
\label{eq:hist-parity-product-and-amplitude}\\
&p_\varkappa(y;\lambda)=p_0(y;\lambda)+\varkappa\operatorname{par}(y)\mathrm D(\lambda),
\qquad
\operatorname{par}(y)=(-1)^{\sum_{(i,x)}y_{i,x}},
\qquad|\varkappa|<1.
\label{eq:hist-parity-lift-law}
\end{align}
The map $\iota_\varkappa:(0,1)^6\to\Delta^\circ_{63}$ sends $\lambda$ to
$p_\varkappa(\cdot;\lambda)$.  The parity amplitude $\varkappa$ is fixed data of the
lift, not an inferred coordinate, and its hyperedge correlation is
$\bar\varkappa=\varkappa/64$.

\theoremheading{Positive smooth lift with paired relabeling symmetry}{thm:hist-three-agent-parity-lift}
For every $|\varkappa|<1$, the law \eqref{eq:hist-parity-lift-law} is strictly positive
and normalized, and every proper marginal is the corresponding marginal of $p_0$.
Consequently $\iota_\varkappa$ is a smooth right inverse of the map that extracts the six
one-bit marginals, and it is an immersion.  For each marginal tuple $\lambda$, the
positive laws with those marginals form a $57$-dimensional set, of which varying
$\varkappa$ traces one line.  Let $H_{\mathrm{pair}}=(\mathbb Z_2)^3\rtimes S_3$ act by
complementing both bits of a selected agent and by typed agent permutations, and act on
$\lambda$ by the same complements and permutations.  Then
\begin{equation}
\iota_\varkappa(g\lambda)=g_\#\iota_\varkappa(\lambda)
\qquad\text{for every }g\in H_{\mathrm{pair}},
\label{eq:hist-parity-paired-equivariance}
\end{equation}
whereas the complement $F_{i,x}$ of one bit satisfies
\begin{equation}
(F_{i,x})_\#p_\varkappa(\cdot;\lambda)=p_{-\varkappa}(\cdot;F_{i,x}\lambda).
\label{eq:hist-parity-odd-pseudoscalar}
\end{equation}
Odd complements are therefore equivariant only if $\varkappa$ changes sign with them, or
if $\varkappa=0$.  \status{ESTABLISHED}

The proof is in \cref{sec:proofs-fiber-geometry}.

\corollaryheading{Agent-block local--global free-energy identity}{cor:hist-parity-agent-block-vfe}
Fix an agent block $B_i=\{(i,b),(i,m)\}$ and the four outside parameters
$\lambda_{B_i^c}$, and let $r_{i,\lambda_{B_i}}=p_\varkappa(y;\lambda)/p_{0,B_i^c}
(y_{B_i^c};\lambda_{B_i^c})$ be the conditional law of $y_{B_i}$ given $y_{B_i^c}$.  For
any interior block
parameters $\lambda_{B_i}$ and $\lambda'_{B_i}$ and any strictly positive posterior
$\boldsymbol\Pi_o$ on the Boolean cube, all relative entropies are finite, and for every
block tangent $Z$
\begin{align}
&\Fenergy_o^{\rm ext}\bigl(p_\varkappa(\cdot;\lambda'_{B_i},\lambda_{B_i^c})\bigr)
-\Fenergy_o^{\rm ext}\bigl(p_\varkappa(\cdot;\lambda_{B_i},\lambda_{B_i^c})\bigr)
\notag\\
&\qquad=\E_{p_{0,B_i^c}}\bigl[
\Fenergy_{B_i,o}^{\rm ext}(r_{i,\lambda'_{B_i}};Y_{B_i^c})
-\Fenergy_{B_i,o}^{\rm ext}(r_{i,\lambda_{B_i}};Y_{B_i^c})\bigr],
\label{eq:hist-parity-agent-local-global-vfe}\\
&d\Fenergy_o^{\rm ext}[Z]
=\E_{p_{0,B_i^c}}\bigl[d\Fenergy_{B_i,o}^{\rm ext}[Z]\bigr],
\label{eq:hist-parity-agent-local-global-differential}
\end{align}
with both differentials evaluated at the same lifted conditional law.  The
outside-averaged block free energy is thus the restriction of one collective free energy
up to a constant.  \status{ESTABLISHED}

\noindent\emph{Proof.}
By \cref{thm:hist-three-agent-parity-lift} the outside marginal is exactly
$p_{0,B_i^c}$, independent of $\lambda_{B_i}$, and strict positivity makes
$r_{i,\lambda_{B_i}}$ a positive conditional law.  The finite state space and the
positive posterior make every relative entropy finite, so
\eqref{eq:hist-parity-agent-local-global-vfe} is \cref{thm:obs-local-global-potential}
with this outside marginal, and differentiating the finite smooth identity gives
\eqref{eq:hist-parity-agent-local-global-differential}.  $\square$

\corollaryheading{Parity-family Fisher residual}{cor:hist-parity-fisher-residual}
For the family of \cref{def:hist-three-agent-parity-lift}, every pair of bits is
independent, so $\mathfrak C=0$, and with unit weights
$\mathsf G_{\mathrm J}-\sum_{(i,x)\in\mathcal I_6}\mathsf G_{i,x}=\mathfrak K\succ0$ when
$\varkappa\ne0$.  At $\lambda=\tfrac12\mathbf 1$, and there with $\varkappa=\tfrac12$,
\begin{align}
\mathsf G_{\mathrm J}=\frac{4}{1-\bar\varkappa^2}\mathrm I_6,
\qquad
\mathsf G_{\mathrm J}-\sum_{(i,x)\in\mathcal I_6}\mathsf G_{i,x}
&=\frac{4\bar\varkappa^2}{1-\bar\varkappa^2}\mathrm I_6,
\label{eq:hist-parity-center-fisher}\\
\mathsf G_{\mathrm J}=\frac{65536}{16383}\mathrm I_6,
\qquad
\mathsf G_{\mathrm J}-\sum_{(i,x)\in\mathcal I_6}\mathsf G_{i,x}
&=\frac{4}{16383}\mathrm I_6,
\label{eq:hist-parity-half-fisher-residual}
\end{align}
and at the center a weighted marginal metric equals the joint one on the full tangent
space exactly when
\begin{equation}
w_{i,x}=\frac{1}{1-\bar\varkappa^2}
\qquad\text{for every }(i,x)\in\mathcal I_6,
\label{eq:hist-parity-center-weight-equality}
\end{equation}
so unit weights are exact only at $\varkappa=0$.  Since the one-bit marginal families do not depend on $\varkappa$, no
rule that reads only them returns every joint Fisher tensor.  \status{ESTABLISHED}

The proof is in \cref{sec:proofs-fiber-geometry}.

\corollaryheading{Promoted parity has one hidden direction}{cor:hist-promoted-parity-rank}
For the seven-parameter family $(\lambda,\varkappa)\mapsto p_\varkappa(\cdot;\lambda)$ on
$(0,1)^6\times(-1,1)$, the derivative of the joint law has rank seven everywhere, while
the one-bit marginal map has rank six and
\begin{equation}
\ker d\left(\operatorname{mar}\circ p\right)=\operatorname{span}\{\partial_\varkappa\}.
\label{eq:hist-cds-promoted-parity-kernel}
\end{equation}
Thus $\varkappa$ is identifiable from the joint law and invisible to the one-bit
marginals.  \status{ESTABLISHED}

\noindent\emph{Proof.}
The one-bit success probabilities are exactly $\lambda$, so a vanishing joint derivative
has zero $\lambda$-component, and the remaining derivative is a multiple of
$\partial_\varkappa p_\varkappa=\operatorname{par}\mathrm D$, which is nowhere zero on the
open domain; hence the joint derivative is injective.  The same identity gives rank six
for the marginal map and \eqref{eq:hist-cds-promoted-parity-kernel}.  $\square$

At the center $\lambda=\tfrac12\mathbf 1$, where $\mathrm D=4^{-6}$, the lift
\eqref{eq:hist-parity-lift-law} is
\begin{equation}
p_\varkappa\bigl(y;\tfrac12\mathbf 1\bigr)
=\tfrac1{64}\left(1+\bar\varkappa\operatorname{par}(y)\right).
\label{eq:hist-parity-center-law}
\end{equation}

\definitionheading{Three-agent parity hyperedge}{def:hist-parity-hyperedge}
At the center \eqref{eq:hist-parity-center-law}, take the uniform latent prior
$\operatorname{Unif}(y)=1/64$ and the binary record kernel
\begin{equation}
K_{\mathrm{hyp}}(1\given y)=\tfrac12\left(1+\bar\varkappa\operatorname{par}(y)\right),
\qquad
K_{\mathrm{hyp}}(0\given y)=\tfrac12\left(1-\bar\varkappa\operatorname{par}(y)\right).
\label{eq:hist-parity-hyperedge-kernel}
\end{equation}

\propositionheading{Exact posterior and tight free energy for the hyperedge}{prop:hist-parity-hyperedge-vfe}
The record $o=1$ has evidence $\tfrac12$, and its posterior is exactly
$p_\varkappa(\cdot;\tfrac12\mathbf 1)$.  For the free energy
$\Fenergy_1(Q)=\KL(Q\Vert\operatorname{Unif})-\E_Q\log K_{\mathrm{hyp}}(1\given Y)$ the
posterior attains the tight value $\Fenergy_1=\log2$, the product law of its one-bit
marginals is $\operatorname{Unif}$, and
\begin{equation}
\Fenergy_1(\operatorname{Unif})-\log2=-\tfrac12\log(1-\bar\varkappa^2).
\label{eq:hist-parity-product-lift-excess}
\end{equation}
Consequently, when $\varkappa\ne0$, no function of the one-bit marginals of the
recognition law, the one-bit marginals of the posterior, and the evidence equals the free
energy on every recognition law.  \status{ESTABLISHED}

The proof is in \cref{sec:proofs-fiber-geometry}.

\openproblemheading{Scope beyond the finite parity witness}{open:hist-parity-lift-scope}
\Cref{thm:hist-three-agent-parity-lift} gives one smooth correlated right inverse of the
one-bit marginal map, equivariant under paired complements and typed permutations, and
\cref{thm:hist-correlated-refinement-nogo} refutes an absolute selector on the admitted
split kernels, so that question is closed.  Open are the existence, for a stated regular
family of agent laws and a stated class of morphisms that contains
non-product-preserving kernels, of a smooth right inverse natural under that class, and
the existence of correlated lifts equivariant under a connected structure group acting on
continuous law fibers.  A construction with proof for a stated family and class, or a
no-go theorem for it, settles each part.

\section{Natural-gradient orbits}
\label{sec:hist-natural-gradient-oriented-orbit}

\definitionheading{Oriented unparameterized history and natural-gradient history}{def:hist-oriented-history}
A representative of a history in a configuration manifold $\mathcal Q$ is a $C^1$ curve
$Q:\mathcal T\to\mathcal Q$ on an open interval with nowhere-vanishing velocity, and two
representatives define the same oriented history when
\begin{equation}
\widetilde Q=Q\circ\phi,
\qquad
\phi:\widetilde{\mathcal T}\longrightarrow\mathcal T
\text{ an orientation-preserving $C^1$ diffeomorphism}.
\label{eq:hist-oriented-reparameterization-equivalence}
\end{equation}
The resulting abstract oriented curve $\mathscr H$ comes with its traversal map
$\jmath:\mathscr H\to\mathcal Q$, which retains order and multiplicity; for an embedded
representative $\mathscr H$ is its oriented image.  Under
\cref{hyp:hist-standing-configuration-tier} the natural gradient of $\Fenergy$ is defined
by the first line below, the descent ray at a noncritical $Q$ is
$\mathcal R^-_{\Fenergy}(Q)=\{-\mathsf a\operatorname{grad}^{\mathsf G}\Fenergy(Q):\mathsf
a>0\}$, and a natural-gradient history is an oriented history whose representatives
satisfy the second:
\begin{align}
\mathsf G\bigl(\operatorname{grad}^{\mathsf G}\Fenergy,Z\bigr)
&=d\Fenergy[Z]
\quad\text{for every }Z\in T\mathcal Q,
\label{eq:hist-fisher-gradient-definition}\\
\dot Q(t)
&=-\mathsf a(t)\operatorname{grad}^{\mathsf G}\Fenergy(Q(t)),
\qquad \mathsf a(t)>0 .
\label{eq:hist-vfe-ray-representative}
\end{align}
When $\Fenergy$ is the lifted objective of \cref{hyp:hist-exact-vfe-lift}, such a history
is an exact free-energy history.

\propositionheading{Scalar mobility preserves the oriented orbit}{prop:hist-scalar-mobility-orbit}
On a noncritical domain, multiplying the natural-gradient field by a smooth positive
function changes only the parameterization of its orbits.  A positive-definite
anisotropic mobility can change the orbit: on the Euclidean plane, for
$\Fenergy=\tfrac12(z_1^2+z_2^2)$ and $\mathsf M=\operatorname{diag}(1,2)$, gradient
descent from $(z_1^0,z_2^0)$ with $z_1^0z_2^0\ne0$ follows the line $z_2=(z_2^0/z_1^0)z_1$, whereas
$\dot z=-\mathsf M\nabla\Fenergy$ follows
\begin{equation}
z_2=z_2^0\left(z_1/z_1^0\right)^2 .
\label{eq:hist-anisotropic-orbit-change}
\end{equation}
A natural-gradient history is not in general a geodesic: on the same Euclidean plane the
natural-gradient flow of $\Fenergy=\tfrac12(z_1^2+2z_2^2)$ is $\dot z_1=-z_1$,
$\dot z_2=-2z_2$, whose orbits are the curves \eqref{eq:hist-anisotropic-orbit-change},
while the geodesics are straight lines.  \status{ESTABLISHED}

\noindent\emph{Proof.}
If $\mathscr V=-\operatorname{grad}^{\mathsf G}\Fenergy$ and $\mathscr V'=\mathsf
a\mathscr V$ with $\mathsf a>0$, the integral curves have the same tangent lines and
orientation, and the change of parameter defined by a positive integral of
$\mathsf a^{-1}$ carries one equation into the other.  In the example, $\dot z_1=-z_1$
and $\dot z_2=-2z_2$ keep $z_2/z_1^2$ constant, which is
\eqref{eq:hist-anisotropic-orbit-change}, while ordinary gradient descent keeps
$z_2/z_1$ constant, and for $z_1^0z_2^0\ne0$ the images differ; the same equations are
the Euclidean natural-gradient flow of $\tfrac12(z_1^2+2z_2^2)$.  $\square$

A rate ratio between coordinate blocks, such as fast beliefs and slow models, is an
anisotropic mobility (the example is the ratio $1:2$), so the selected orbit can depend
on it; the rate-separated flow is in \cref{sec:obs-fast-slow}.

\propositionheading{Noncritical free-energy orbits are embedded}{prop:hist-vfe-orbit-embedded}
Let $\mathcal Q$ be finite-dimensional, $\Fenergy$ of class $C^2$, and $\mathsf M$ a
$C^1$ field of positive-definite symmetric forms on the cotangent spaces, for instance
$\mathsf M=\mathsf a\mathsf G^{-1}$ for a $C^1$ metric and a $C^1$ function
$\mathsf a>0$.  Let $\gamma$ be a maximal integral curve of
$\mathscr V=-\mathsf M(d\Fenergy)$ through a noncritical point.  Then $\gamma$ meets no
critical point, $\Fenergy\circ\gamma$ is strictly decreasing, and $\gamma$ is an injective
immersion that is a homeomorphism onto its image, an embedded $C^1$ curve on which
$\Fenergy$ is a strictly monotone continuous coordinate, independent of parameterization
and gauge invariant whenever $\Fenergy$ is.  In particular the multiplicity clause of
\cref{def:hist-oriented-history} is vacuous for noncritical histories, and for the
natural gradient the Fisher clock of \cref{def:hist-fisher-clock} satisfies
$d\mathfrak t=-d\Fenergy/\lVert\operatorname{grad}^{\mathsf G}\Fenergy\rVert_{\mathsf G}$
along the orbit.  \status{ESTABLISHED}

\noindent\emph{Proof.}
The field $\mathscr V$ is $C^1$, so its integral curves are unique, and its zeros are the
critical points of $\Fenergy$; since the constant curve at a zero is an integral curve, a
curve through a nonzero of $\mathscr V$ never reaches one.  Along $\gamma$,
$\tfrac{d}{dt}\Fenergy(\gamma)=-\mathsf M(d\Fenergy,d\Fenergy)<0$, which makes
$\Fenergy\circ\gamma$ strictly decreasing, $\gamma$ injective, and $\dot\gamma\ne0$.  On
the image the inverse of $\gamma$ is $Q\mapsto(\Fenergy\circ\gamma)^{-1}(\Fenergy(Q))$, a
composite of continuous maps, so $\gamma$ is an embedding.  The clock relation is
\eqref{eq:hist-vfe-dissipation-per-length}.  $\square$

\propositionheading{A semidefinite tensor cannot be silently inverted}{prop:hist-semidefinite-gradient-obstruction}
If the proposed metric on a configuration space is only semidefinite, the
natural-gradient equation can have no solution or several.  \status{ESTABLISHED}

\noindent\emph{Proof.}
On $\R^2$ let $\mathsf G=dz_1^2$.  For $\Fenergy=z_2$ the equation
$\mathsf G(\operatorname{grad}\Fenergy,\cdot)=dz_2$ has no solution because the left side
annihilates $\partial_{z_2}$, and for $\Fenergy=z_1$ every
$\partial_{z_1}+\mathrm c\partial_{z_2}$ solves it.  $\square$

A quotient constructed as in \cref{thm:hist-retained-fisher-quotient}, or a separate
regularization, must precede any use of an inverse or a pseudoinverse.

\remarkheading{Discrete histories}{rem:hist-discrete-histories}
In a chart where the metric has a constant coefficient matrix $\mathsf G$, the explicit
Euler step of \eqref{eq:hist-vfe-ray-representative} with $\mathsf a\equiv1$ and step
$\Delta t>0$ is $Q_{n+1}=Q_n-\Delta t\mathsf G^{-1}\nabla\Fenergy(Q_n)$.  If
$\nabla^2\Fenergy\preceq\mathrm c_{\mathsf G}\mathsf G$ on the segment $[Q_n,Q_{n+1}]$,
with $\mathrm c_{\mathsf G}>0$, then
\begin{equation}
\Fenergy(Q_{n+1})
\leq\Fenergy(Q_n)
-\Delta t\left(1-\frac{\mathrm c_{\mathsf G}\Delta t}{2}\right)
\left\|\operatorname{grad}^{\mathsf G}\Fenergy(Q_n)\right\|^2_{\mathsf G},
\label{eq:hist-discrete-descent-bound}
\end{equation}
so every step with $0<\Delta t<2/\mathrm c_{\mathsf G}$ decreases the objective off the
critical set, and the bound is sharp: for a quadratic objective whose Hessian has largest
generalized eigenvalue $\mathrm c_{\mathsf G}$ relative to $\mathsf G$, an iterate on that
eigendirection is multiplied by $1-\Delta t\mathrm c_{\mathsf G}$, so
$\Delta t>2/\mathrm c_{\mathsf G}$ increases the objective there.  If, in the same chart,
$\mathscr V=-\mathsf G^{-1}\nabla\Fenergy$ is Lipschitz with constant
$\mathrm c_{\mathrm L}>0$ on the closed $\delta$-neighborhood of an exact arc $Q([0,t_1])$
and $\lVert\ddot Q\rVert\leq\mathrm c_2$ on the arc, then
\begin{equation}
\lVert Q_n-Q(n\Delta t)\rVert
\leq\frac{\mathrm c_2\Delta t}{2\mathrm c_{\mathrm L}}\left(e^{\mathrm c_{\mathrm L}n\Delta t}-1\right),
\qquad n\Delta t\leq t_1,
\label{eq:hist-euler-global-error}
\end{equation}
for every $\Delta t$ small enough that the right side stays below $\delta$ at
$n\Delta t=t_1$.  Euler iterates follow the selected orbit to first order in $\Delta t$
on compact parameter intervals without lying on it.  \status{ESTABLISHED}

\noindent\emph{Proof.}
Taylor's formula with integral remainder along the segment gives
$\Fenergy(Q_{n+1})\leq\Fenergy(Q_n)+\nabla\Fenergy^\top\Delta Q+\tfrac12\mathrm
c_{\mathsf G}\Delta Q^\top\mathsf G\Delta Q$ with
$\Delta Q=-\Delta t\mathsf G^{-1}\nabla\Fenergy$, which is
\eqref{eq:hist-discrete-descent-bound}; sharpness is the stated eigendirection
computation.  The local truncation error is at most $\mathrm c_2\Delta t^2/2$, so
$e_n=\lVert Q_n-Q(n\Delta t)\rVert$ obeys
$e_{n+1}\leq(1+\Delta t\mathrm c_{\mathrm L})e_n+\mathrm c_2\Delta t^2/2$ while the iterates
stay in the neighborhood, the discrete Gr\"onwall inequality gives
\eqref{eq:hist-euler-global-error}, and induction on $n$ keeps the iterates in the
neighborhood.  $\square$

\section{Fisher duration}
\label{sec:hist-fisher-duration}
\label{sec:pb-curve-taxonomy}
\label{sec:hist-clock-tiers}

\definitionheading{Fisher length and information clock}{def:hist-fisher-clock}
Let $Q:\mathcal T\to\mathcal Q$ be a $C^1$ curve in a configuration manifold with metric
$\mathsf G$.  At $c_*$ its velocity lies in the fibers, $\operatorname{ver}(\dot Q)=\dot
Q\in T_Q\mathcal Q$, so no connection enters.  Its Fisher speed, its length on
$[t_0,t_1]$, its cumulative clock from a chosen origin $t_0$, and, for a curve $\Gamma$
in one law fiber with its Fisher metric, the vertical Fisher length are
\begin{align}
\dot{\mathfrak t}(t)
&=\sqrt{\mathsf G\bigl(\dot Q(t),\dot Q(t)\bigr)},
\label{eq:hist-fisher-speed}\\
\operatorname{len}_{\mathsf G}(Q;[t_0,t_1])
&=\int_{t_0}^{t_1}\dot{\mathfrak t}(t) dt,
\label{eq:hist-fisher-length}\\
\mathfrak t_{Q,t_0}(t)
&=\int_{t_0}^{t}\dot{\mathfrak t}(t') dt',
\label{eq:hist-cumulative-fisher-clock}\\
\operatorname{len}_{g^F}(\Gamma)
&=\int\sqrt{g^F\bigl(\operatorname{ver}(\dot\Gamma),\operatorname{ver}(\dot\Gamma)\bigr)} dt .
\label{eq:pb-vertical-fisher-length}
\end{align}
The belief clock and the model clock of agent $i$ are the Fisher lengths of
$t\mapsto q_i^b(t)$ and $t\mapsto q_i^m(t)$, combined by the channel-weighted speed
\eqref{eq:hist-pointwise-clock-speed}.  Along a history of $\mathcal Q_V$ with
block-diagonal metric $\bigoplus_i\mathsf G_i$, the clock $\mathfrak t_i$ of agent $i$ is
the $\mathsf G_i$-length of $t\mapsto Q_i(t)$.  The origin is part of the clock, and the
parameter $t$ does not enter the length.

\theoremheading{Reparameterization invariance and unit-speed form}{thm:hist-fisher-clock-invariance}
Let $\widetilde Q=Q\circ\phi$ with $\phi:\widetilde{\mathcal T}\to\mathcal T$ an
orientation-preserving $C^1$ diffeomorphism.  Corresponding subarcs have the same length,
and
\begin{equation}
\mathfrak t_{\widetilde Q,\widetilde t_0}(\widetilde t)
=\mathfrak t_{Q,t_0}(\phi(\widetilde t)),
\qquad
t_0=\phi(\widetilde t_0).
\label{eq:hist-clock-reparameterization-law}
\end{equation}
If $\dot{\mathfrak t}>0$, the clock is strictly increasing and gives the unit-speed
representative of the oriented history; on a subinterval where $\dot{\mathfrak t}$
vanishes, as on a critical or null segment, the clock is constant and is not a
coordinate.  If $Q$ satisfies \eqref{eq:hist-vfe-ray-representative}, then on its
noncritical part
\begin{align}
\frac{d\mathfrak t}{dt}
&=\mathsf a\left\|\operatorname{grad}^{\mathsf G}\Fenergy\right\|_{\mathsf G},
\label{eq:hist-vfe-clock-rate}\\
\frac{dQ}{d\mathfrak t}
&=-\frac{\operatorname{grad}^{\mathsf G}\Fenergy}
{\left\|\operatorname{grad}^{\mathsf G}\Fenergy\right\|_{\mathsf G}},
\label{eq:hist-vfe-unit-tangent}\\
\frac{d\Fenergy}{d\mathfrak t}
&=-\left\|\operatorname{grad}^{\mathsf G}\Fenergy\right\|_{\mathsf G}.
\label{eq:hist-vfe-dissipation-per-length}
\end{align}
If the clock is read with a second metric $\mathsf G'$, for instance the unweighted
Fisher metric along a flow whose blocks carry rate weights, then along any $C^1$ history
with velocity $Z\ne0$
\begin{equation}
\left|\frac{d\Fenergy}{d\mathfrak t'}\right|
=\frac{|d\Fenergy(Z)|}{\lVert Z\rVert_{\mathsf G'}}
\leq\left\|\operatorname{grad}^{\mathsf G'}\Fenergy\right\|_{\mathsf G'},
\label{eq:hist-clock-metric-dissipation-bound}
\end{equation}
with equality exactly when $d\Fenergy=0$ or $Z$ is parallel to
$\operatorname{grad}^{\mathsf G'}\Fenergy$.
Along a history of $\mathcal Q_V$ with block-diagonal metric the squared speeds add,
$\dot{\mathfrak t}^2=\sum_i\dot{\mathfrak t}_i^2$, and over any common parameter interval
$\max_i\mathfrak t_i\leq(\sum_i\mathfrak t_i^2)^{1/2}\leq\mathfrak t\leq\sum_i\mathfrak
t_i$.  \status{ESTABLISHED}

The proof is in \cref{sec:proofs-fiber-geometry}.

\remarkheading{Endpoints do not determine an update duration}{prop:hist-realized-length-versus-distance}
For an update $Q$ from $Q(t_0)=Q'$ to $Q(t_1)=Q''$ in one path component of
$\mathcal Q$, the distance $\operatorname{dist}_{\mathsf G}(Q',Q'')$, the infimum of
$\operatorname{len}_{\mathsf G}(\Gamma;[t_0,t_1])$ over curves $\Gamma$ with these
endpoints, is at most $\operatorname{len}_{\mathsf G}(Q;[t_0,t_1])$, because the realized
update is one such curve, and the inequality is strict for a curved path in the Euclidean
Fisher geometry of the two-parameter normal location family.  Duration is also distinct
from the amount of curve parameter: in the unit-variance normal location family the path
$\mu(t)=e^{-t}$ has length $\int_0^\infty|\dot\mu(t)|dt=\int_0^\infty e^{-t}dt=1$ over an
unbounded parameter interval.  \status{ESTABLISHED}

\propositionheading{Finite duration for analytic objectives}{prop:hist-analytic-finite-duration}
Let $\mathcal Q$ be finite-dimensional with a continuous positive-definite metric
$\mathsf G$, and let $\gamma:[0,\infty)\to\mathcal Q$ satisfy
$\dot\gamma=-\operatorname{grad}^{\mathsf G}\Fenergy$.  If $\gamma(t)$ has an
accumulation point $Q_*$ as $t\to\infty$ and $\Fenergy$ is real analytic in a chart
around $Q_*$, then $\gamma(t)\to Q_*$, $Q_*$ is critical, and
$\operatorname{len}_{\mathsf G}(\gamma;[0,\infty))$ is finite.  This holds for every
forward history with relatively compact image when $\mathcal Q$ has a real-analytic atlas
in whose charts $\Fenergy$ is real analytic.  Analyticity cannot be weakened to
smoothness: for the $C^\infty$ function of \citet[Section~2, Eq.~(8)]{absil2005convergence}
on the Euclidean plane, the gradient history from the initial data given there is
bounded, has the unit circle as its $\omega$-limit set, and has infinite duration.
\status{ESTABLISHED}

\propositionheading{Invariant free-energy histories are horizontal}{prop:hist-invariant-history-horizontal}
Let a Lie group $\mathbb G$ act smoothly on a finite-dimensional configuration manifold
$\mathcal Q$ with a positive-definite metric $\mathsf G$, write
$\xi_{\mathcal Q}(Q)=\tfrac{d}{dt}\big|_{t=0}\exp(t\xi)\cdot Q$ and
$V_Q=\{\xi_{\mathcal Q}(Q):\xi\in\operatorname{Lie}\mathbb G\}$, and let $\Fenergy$ be
$C^1$ and $\mathbb G$-invariant.  Then $\operatorname{grad}^{\mathsf G}\Fenergy(Q)$ is
$\mathsf G$-orthogonal to $V_Q$ at every $Q$, every natural-gradient history has velocity
orthogonal to the orbits, and its gauge-quotient speed
$\min_{\xi}\lVert\dot Q-\xi_{\mathcal Q}(Q)\rVert_{\mathsf G}$ equals its Fisher speed.
If $\mathsf G$ is $C^\infty$ and the action is free, proper, and isometric, the quotient
map $\pi:\mathcal Q\to\mathcal Q/\mathbb G$ is a Riemannian submersion for a unique
metric $\bar{\mathsf G}$ on the quotient, $\Fenergy=\bar\Fenergy\circ\pi$, and
$T\pi\operatorname{grad}^{\mathsf G}\Fenergy=\operatorname{grad}^{\bar{\mathsf
G}}\bar\Fenergy\circ\pi$; each natural-gradient history descends to a natural-gradient
history of $\bar\Fenergy$ with the same Fisher duration.  \status{ESTABLISHED}

The proofs of \cref{prop:hist-analytic-finite-duration,prop:hist-invariant-history-horizontal}
are in \cref{sec:proofs-fiber-geometry}.

\section{Orthogonal clocks}
\label{sec:hist-global-clock}

\theoremheading{Exact criterion for an orthogonal clock}{thm:hist-global-clock-exactness}
Let $\mathcal U\subseteq\mathcal Q$ be a connected open set without critical points, with
$\Fenergy$ of class $C^2$ and $\mathsf G$ of class $C^1$.  Define $\mathsf n$,
$\mathbf e_{\Fenergy}$, and $\bm\alpha$ by the first line below, and call a $C^1$
function $\mathfrak t^\perp:\mathcal U\to\R$ that satisfies the second an orthogonal
clock:
\begin{align}
&\mathsf n=\left\|\operatorname{grad}^{\mathsf G}\Fenergy\right\|_{\mathsf G},
\qquad
\mathbf e_{\Fenergy}=-\frac{\operatorname{grad}^{\mathsf G}\Fenergy}{\mathsf n},
\qquad
\bm\alpha=\mathbf e_{\Fenergy}^\flat=-\frac{d\Fenergy}{\mathsf n},
\label{eq:hist-normalized-vfe-one-form}\\
&d\mathfrak t^\perp=\bm\alpha.
\label{eq:hist-global-clock-potential}
\end{align}
An orthogonal clock exists if and only if
\begin{equation}
d\bm\alpha=0,
\qquad
\oint_\gamma\bm\alpha=0
\quad\text{for every piecewise $C^1$ closed curve }\gamma\subset\mathcal U,
\label{eq:hist-global-clock-period-criterion}
\end{equation}
and on a simply connected $\mathcal U$ closedness suffices.  The local obstruction is
\begin{equation}
d\bm\alpha=\frac{1}{\mathsf n^2} d\mathsf n\wedge d\Fenergy,
\label{eq:hist-normalized-form-curvature}
\end{equation}
so $\bm\alpha$ is closed exactly when $\mathsf n$ is locally a function of $\Fenergy$,
and an orthogonal clock is then constant on every connected piece of every level set of
$\Fenergy$ and locally a strictly decreasing function of $\Fenergy$.  Whenever
$\mathfrak t^\perp$ exists, every unit-speed natural-gradient history in $\mathcal U$
obeys $d\mathfrak t^\perp/d\mathfrak t=1$, and the level sets of $\mathfrak t^\perp$ are
$\mathsf G$-orthogonal to $\mathbf e_{\Fenergy}$.  For $\Fenergy=z_1z_2$ on the open first
quadrant with the Euclidean metric,
\begin{equation}
d\bm\alpha
=\frac{z_1^2-z_2^2}{(z_1^2+z_2^2)^{3/2}} dz_1\wedge dz_2,
\label{eq:hist-nonexact-clock-example}
\end{equation}
which vanishes identically on no open set, so no orthogonal clock exists on any open
subset of the quadrant.  \status{ESTABLISHED}

\propositionheading{Unit-speed clocks from a global cross-section}{prop:hist-cross-section-clock}
Let $\mathcal U$ be as in \cref{thm:hist-global-clock-exactness}, and let
$\mathcal S\subset\mathcal U$ be a $C^1$ hypersurface that every maximal integral curve
of $\mathbf e_{\Fenergy}$ in $\mathcal U$ meets exactly once and transversally.  Then the
signed Fisher arc length along each orbit, measured from its crossing with $\mathcal S$,
is a $C^1$ function $\mathfrak t^{\mathcal S}:\mathcal U\to\R$ with
$d\mathfrak t^{\mathcal S}(\mathbf e_{\Fenergy})=1$, so every unit-speed natural-gradient
history obeys $d\mathfrak t^{\mathcal S}/d\mathfrak t=1$.  A level set of $\Fenergy$ that
every orbit meets is such a cross-section, and the level sets of
$\mathfrak t^{\mathcal S}$ are $\mathsf G$-orthogonal to $\mathbf e_{\Fenergy}$ only if
$\bm\alpha$ is closed.  For $\Fenergy=z_1z_2$ on the open first quadrant with the
Euclidean metric, the hyperbola $\{z_1z_2=1\}$ is a global cross-section, so a unit-speed
clock exists on the whole quadrant although no orthogonal one exists on any open subset.
\status{ESTABLISHED}

The proofs of \cref{thm:hist-global-clock-exactness,prop:hist-cross-section-clock} are in
\cref{sec:proofs-fiber-geometry}.

\section{Record clocks}
\label{sec:hist-record-clocks}

\corollaryheading{Same-path record and meta-agent clock contraction}{thm:hist-record-clock-contraction}
Let $t\mapsto P_t$ be a $C^1$ history in a family of laws on $\mathsf Y$ that is
differentiable in quadratic mean along it, with directional score $\mathbf s_t$, and let
$K:\mathsf Y\rightsquigarrow\bar{\mathsf Y}$ be a Markov kernel that does not depend on
$t$ and satisfies $K(y,\bar{\mathsf Y})=1$ for every $y$.  Then the record history
$\bar P_t=P_tK$, the image of the same fine path, satisfies
\begin{align}
\bar{\mathbf s}_t(\bar Y)
&=\E_t\left[\mathbf s_t(Y)\mid\bar Y\right],
\label{eq:hist-record-score-projection}\\
\dot{\mathfrak t}_{P}^2-\dot{\mathfrak t}_{\bar P}^2
&=\E_t\operatorname{Var}_t\left(\mathbf s_t(Y)\mid\bar Y\right)
\geq0,
\label{eq:hist-record-conditional-variance-defect}
\end{align}
and on every common parameter interval
\begin{equation}
\operatorname{len}(\bar P)\leq\operatorname{len}(P),
\qquad
\mathfrak t_{\bar P}(t_1)-\mathfrak t_{\bar P}(t_0)
\leq
\mathfrak t_{P}(t_1)-\mathfrak t_{P}(t_0),
\label{eq:hist-record-clock-contraction}
\end{equation}
with equality, when the fine length on the interval is finite, exactly when the fine
score is $\bar Y$-measurable for almost every $t$ in the interval. Equality of two
infinite lengths gives no such conclusion. The inequality applies to a meta-agent exactly when its law is the image
of the same fine history under one such kernel.  Parameter independence of the kernel
cannot be dropped: for a dominated channel $K_t$ with density
$\mathsf k_t(\bar y\given y)$, assume the joint family $P_t(dy)K_t(y,d\bar y)$ is
differentiable in quadratic mean with square-integrable score
$\mathbf s_t(y)+\partial_t\log\mathsf k_t(\bar y\given y)$. Its record score is
\begin{equation}
\bar{\mathbf s}_t(\bar Y)
=\E_t\left[
\mathbf s_t(Y)
+\partial_t\log\mathsf k_t(\bar Y\given Y)
\mathrel{\big|}\bar Y
\right],
\label{eq:hist-parameter-dependent-channel-score}
\end{equation}
and a constant input followed by a channel that emits a Bernoulli bit of success
probability $(1+e^{-t})^{-1}$ has zero fine speed and record speed
$(2\cosh(t/2))^{-1}>0$.  \status{ESTABLISHED}

The proof is in \cref{sec:proofs-fiber-geometry}.

\begin{figure}[htbp]
\centering
\begin{tikzpicture}[>=stealth]
  \node[draw,rounded corners,fill=lightblue,minimum width=3.3cm,minimum height=0.8cm]
    (fine) at (0,2.1) {fine history $P_t$};
  \node[draw,rounded corners,fill=lightgreen,minimum width=3.8cm,minimum height=0.8cm]
    (image) at (6.2,2.1) {same-path image $P_tK$};
  \draw[->,very thick,primaryblue] (fine)--(image);
  \node[text=primaryblue,font=\small] at (3.1,2.85)
    {parameter-independent Markov channel: $d\mathfrak t_{PK}\leq d\mathfrak t_P$};
  \node[draw,rounded corners,fill=lightblue,minimum width=3.3cm,minimum height=0.8cm]
    (fflow) at (0,0) {$\operatorname{Fl}^{\mathscr V_\ell}_t(Q_0)$};
  \node[draw,rounded corners,fill=lightgreen,minimum width=3.8cm,minimum height=0.8cm]
    (mflow) at (6.2,0) {$\operatorname{Fl}^{\mathscr V_{\ell+1}}_{t'}(\mathsf F_\ell Q_0)$};
  \draw[->,very thick,dashed,sciencegreen] (fflow)--(mflow);
  \node[text=sciencegreen,font=\small] at (3.1,0.75)
    {independent flows require oriented semiconjugacy};
\end{tikzpicture}
\caption{The upper arrow is the same-path comparison of
\cref{thm:hist-record-clock-contraction}, with its conditional-variance defect.  The
lower arrow compares independently generated flows and exists only under the
semiconjugacy condition of \cref{def:hist-oriented-semiconjugacy}.}
\label{fig:hist-markov-clock-and-semiconjugacy}
\end{figure}

\openproblemheading{Quantitative meta-agent clock ratio}{open:hist-meta-clock-ratio}
\Cref{thm:hist-record-clock-contraction} shows that a meta-agent whose law is a Markov
image of its constituents' joint history accumulates no more Fisher duration than that
history, with the exact defect \eqref{eq:hist-record-conditional-variance-defect}.  Open
is the ratio between the meta-agent's duration and the constituents' collective duration
for an $N$-constituent meta-agent of \cref{ch:agent-network-rg}, as a function of $N$ and
of the correlation between the constituents' velocities, for coherent histories (aligned
velocities) and incoherent ones (independent random velocities).  The expected ratio of
arc lengths is not the square root of the expected ratio of relative entropies, so a
divergence computation does not settle it; a theorem that computes or bounds the ratio
for a stated coarse map and a stated law of constituent velocities, including its
dependence on $N$ in both regimes, does.

\section{The coarse configuration map and oriented semiconjugacy}
\label{sec:hist-configuration-coarse-map}
\label{sec:hist-independent-flows}

\definitionheading{Coarse configuration map}{def:hist-coarse-configuration-map}
A coarse configuration map at depth $\ell$ is a smooth map
\begin{equation}
\mathsf F_\ell:\mathcal Q_\ell\longrightarrow\mathcal Q_{\ell+1}
\label{eq:hist-configuration-coarse-map}
\end{equation}
from the configuration manifold of the fine agents at $c_*$ to that of their meta-agents,
specified separately from the sample kernels that may realize it.  The fine and coarse
natural-gradient fields are
$\mathscr V_\ell=-\operatorname{grad}^{\mathsf G_\ell}\Fenergy_\ell$ and
$\mathscr V_{\ell+1}=-\operatorname{grad}^{\mathsf G_{\ell+1}}\Fenergy_{\ell+1}$, with
local flows $\operatorname{Fl}^{\mathscr V_\ell}_t$ and
$\operatorname{Fl}^{\mathscr V_{\ell+1}}_{t'}$.

\definitionheading{Oriented semiconjugacy}{def:hist-oriented-semiconjugacy}
For an open $\mathcal U\subseteq\mathcal Q_\ell$, the pair $(\mathsf F_\ell,\mathsf
a_\ell)$ with $\mathsf a_\ell:\mathcal U\to(0,\infty)$ continuous is an oriented
semiconjugacy of $\mathscr V_\ell$ onto $\mathscr V_{\ell+1}$ over $\mathcal U$ when
\begin{equation}
T_Q\mathsf F_\ell \mathscr V_\ell(Q)
=\mathsf a_\ell(Q) \mathscr V_{\ell+1}(\mathsf F_\ell Q)
\quad\text{for every }Q\in\mathcal U.
\label{eq:hist-oriented-vector-field-semiconjugacy}
\end{equation}

\lemmaheading{The factor is determined and regular off the coarse critical set}{lem:hist-semiconjugacy-factor}
If $\mathscr V_{\ell+1}(\mathsf F_\ell Q)\ne0$, then $\mathsf a_\ell(Q)$ is unique and
\begin{equation}
\mathsf a_\ell(Q)
=\frac{\mathsf G_{\ell+1}
\left(T_Q\mathsf F_\ell \mathscr V_\ell(Q),\mathscr V_{\ell+1}(\mathsf F_\ell Q)\right)}
{\left\lVert\mathscr V_{\ell+1}(\mathsf F_\ell Q)\right\rVert^2_{\mathsf G_{\ell+1}}},
\label{eq:hist-semiconjugacy-factor-formula}
\end{equation}
so $\mathsf a_\ell$ is $C^k$ wherever $Q\mapsto T_Q\mathsf F_\ell\mathscr V_\ell(Q)$,
$Q\mapsto\mathscr V_{\ell+1}(\mathsf F_\ell Q)$, and the pulled-back coefficients of
$\mathsf G_{\ell+1}$ are $C^k$, for instance when $\mathsf F_\ell\in C^{k+1}$ and
$\mathscr V_\ell,\mathscr V_{\ell+1},\mathsf G_{\ell+1}$ are $C^k$; continuity of
$\mathsf a_\ell$ is therefore automatic where $\mathscr V_{\ell+1}\ne0$ on
$\mathsf F_\ell(\mathcal U)$.  If $\mathscr V_{\ell+1}(\mathsf F_\ell Q)=0$, then
\eqref{eq:hist-oriented-vector-field-semiconjugacy} forces
$T_Q\mathsf F_\ell\mathscr V_\ell(Q)=0$ and fixes no value of $\mathsf a_\ell(Q)$; limits
from the noncritical set may determine it or may obstruct any continuous positive
extension.  \status{ESTABLISHED}

\noindent\emph{Proof.}
Pair \eqref{eq:hist-oriented-vector-field-semiconjugacy} with
$\mathscr V_{\ell+1}(\mathsf F_\ell Q)$ and divide by its positive squared norm.
$\square$

\theoremheading{Oriented semiconjugacy and maximal intervals}{thm:hist-oriented-semiconjugacy}
\label{prop:hist-oriented-semiconjugacy}%
Assume \eqref{eq:hist-oriented-vector-field-semiconjugacy} on an open $\mathcal U$ with
$\mathscr V_{\ell+1}$ locally Lipschitz.  Fix $Q\in\mathcal U$, let $\mathcal T_Q$ be the
connected component of $0$ in the set of $t$ for which
$\operatorname{Fl}^{\mathscr V_\ell}_{t''}(Q)\in\mathcal U$ for every $t''$ between $0$
and $t$, and set $\psi_Q(t)=\int_0^t\mathsf a_\ell(\operatorname{Fl}^{\mathscr
V_\ell}_{t''}Q) dt''$.  Then $\psi_Q$ is a strictly increasing $C^1$ diffeomorphism of
$\mathcal T_Q$ onto an open interval $\mathcal T'_Q\ni0$, and
\begin{equation}
\mathcal T'_Q\subseteq\bar{\mathcal T}^{\max}_{\mathsf F_\ell Q},
\qquad
\mathsf F_\ell\left(\operatorname{Fl}^{\mathscr V_\ell}_t(Q)\right)
=\operatorname{Fl}^{\mathscr V_{\ell+1}}_{\psi_Q(t)}\left(\mathsf F_\ell Q\right)
\quad\text{for every }t\in\mathcal T_Q,
\label{eq:hist-oriented-flow-semiconjugacy}
\end{equation}
where $\bar{\mathcal T}^{\max}_{\mathsf F_\ell Q}=(\bar t^-,\bar t^+)$ is the maximal
interval of the coarse flow through $\mathsf F_\ell Q$.  With $\mathcal T_Q=(t^-,t^+)$
and endpoints in the extended real line,
\begin{equation}
\mathcal T'_Q
=\left(
\lim_{t\downarrow t^-}\psi_Q(t),
\lim_{t\uparrow t^+}\psi_Q(t)
\right),
\label{eq:hist-semiconjugacy-endpoint-image}
\end{equation}
and the whole maximal coarse orbit is traversed if and only if the two limits equal
$\bar t^-$ and $\bar t^+$; when the fine orbit is complete in $\mathcal U$ and the coarse
maximal interval is $\R$, this means
$\int_0^{+\infty}\mathsf a_\ell(\operatorname{Fl}^{\mathscr V_\ell}_{t}Q)dt
=\int_{-\infty}^{0}\mathsf a_\ell(\operatorname{Fl}^{\mathscr V_\ell}_{t}Q)dt=+\infty$.
A positive lower bound on $\mathsf a_\ell$ suffices and is not necessary: on $\R$ with
$\mathscr V_\ell=\mathscr V_{\ell+1}=\partial_z$ and
$\mathsf F_\ell=\operatorname{arsinh}$, $\mathsf a_\ell=(1+z^2)^{-1/2}$ has infimum zero
while $\psi_0=\operatorname{arsinh}$ maps $\R$ onto $\R$; with $\mathsf F_\ell=\arctan$,
$\mathsf a_\ell=(1+z^2)^{-1}$ is positive while $\mathcal T'_0=(-\pi/2,\pi/2)$ although
both flows are complete.  Conversely, differentiating a flow relation of the form
\eqref{eq:hist-oriented-flow-semiconjugacy} at $t=0$ gives
\eqref{eq:hist-oriented-vector-field-semiconjugacy}.  \status{ESTABLISHED}

The proof is in \cref{sec:proofs-fiber-geometry}.

\propositionheading{Noncollapse is a separate hypothesis}{prop:hist-noncollapse}
Condition \eqref{eq:hist-oriented-vector-field-semiconjugacy} alone permits total
collapse: with $\mathsf F_\ell$ constant and $\mathscr V_{\ell+1}\equiv0$ it holds for
every positive $\mathsf a_\ell$, while every fine orbit maps to a point and every coarse
duration is zero.  A nonconstant shared oriented history additionally requires
$T_Q\mathsf F_\ell\mathscr V_\ell(Q)\ne0$ on every nontrivial subarc, equivalently
$\mathscr V_{\ell+1}(\mathsf F_\ell Q)\ne0$ there, and full traversal requires the
endpoint equalities of \eqref{eq:hist-semiconjugacy-endpoint-image}; neither follows from
\eqref{eq:hist-oriented-vector-field-semiconjugacy}.  \status{ESTABLISHED}

\noindent\emph{Proof.}
Both sides of \eqref{eq:hist-oriented-vector-field-semiconjugacy} vanish in the example,
and the arctan example of \cref{thm:hist-oriented-semiconjugacy} satisfies it without
full traversal.  $\square$

\theoremheading{The semiconjugacy rate cancels from geometric length}{thm:hist-duration-relation}
Under \eqref{eq:hist-oriented-vector-field-semiconjugacy}, let $Q(t)$ be a fine orbit.
The image speed is
\begin{equation}
\left\lVert\tfrac{d}{dt}\mathsf F_\ell Q(t)\right\rVert_{\mathsf G_{\ell+1}}
=\mathsf a_\ell\left(Q(t)\right)
\left\lVert\mathscr V_{\ell+1}\left(\mathsf F_\ell Q(t)\right)\right\rVert_{\mathsf G_{\ell+1}},
\label{eq:hist-image-speed}
\end{equation}
and for $t_0<t_1$ in $\mathcal T_Q$ the accumulated coarse duration is
\begin{equation}
\mathfrak t^{(\ell+1)}(t_1)-\mathfrak t^{(\ell+1)}(t_0)
=\int_{\psi_Q(t_0)}^{\psi_Q(t_1)}
\left\lVert\mathscr V_{\ell+1}\left(\operatorname{Fl}^{\mathscr V_{\ell+1}}_{t'}(\mathsf F_\ell Q)\right)
\right\rVert_{\mathsf G_{\ell+1}}dt',
\label{eq:hist-coarse-duration}
\end{equation}
the intrinsic $\mathsf G_{\ell+1}$-length of the coarse orbit arc, which does not depend
on $\mathsf a_\ell$.  \status{ESTABLISHED}

\noindent\emph{Proof.}
Equation \eqref{eq:hist-image-speed} is \eqref{eq:hist-oriented-vector-field-semiconjugacy}
with positive homogeneity of the norm and $\mathsf a_\ell>0$.  Substituting it and
changing variables $t'=\psi_Q(t)$, whose Jacobian is
$\mathsf a_\ell(\operatorname{Fl}^{\mathscr V_\ell}_tQ)$, cancels the factor.  $\square$

\remarkheading{Exact duration comparison}{thm:hist-duration-criterion}
Under \eqref{eq:hist-oriented-vector-field-semiconjugacy}, along one orbit and from a
common origin, the coarse and fine durations agree identically if and only if
$\lVert T\mathsf F_\ell\mathscr V_\ell\rVert_{\mathsf G_{\ell+1}}=\lVert\mathscr
V_\ell\rVert_{\mathsf G_\ell}$ almost everywhere along the orbit, they agree up to a
factor $\mathrm c>0$ if and only if the same holds with that factor, and the coarse
duration has no larger increments on every subinterval if and only if
\begin{equation}
\left\lVert T\mathsf F_\ell\mathscr V_\ell\right\rVert_{\mathsf G_{\ell+1}}
\leq\left\lVert\mathscr V_\ell\right\rVert_{\mathsf G_\ell}
\quad\text{along the orbit},
\label{eq:hist-duration-criterion}
\end{equation}
because both durations are integrals of continuous nonnegative densities from a common
origin, and two such integrals agree for every upper limit exactly when the densities
agree almost everywhere.  The condition is scalar and concerns the single direction
$\mathscr V_\ell$; the Loewner condition $\mathsf F_\ell^*\mathsf G_{\ell+1}\preceq\mathsf
G_\ell$ is sufficient and stronger, and \cref{thm:hist-pointwise-contraction-lift}
supplies it for channelwise Markov maps.  Fiberwise Fisher contraction alone gives
neither, for two independent reasons.  The configuration metrics are separate data: on
$\R$ with $\mathscr V_\ell=\mathscr V_{\ell+1}=\partial_z$ and $\mathsf F_\ell$ the
identity, the semiconjugacy holds with $\mathsf a_\ell\equiv1$, and the metrics $dz^2$
and $4dz^2$ double the duration, whether realized with the identity fiber map and unequal
channel weights $1$ and $4$ (zero Fisher defect) or with fibers $\mathcal N(\mu,1)$ and
$\mathcal N(\mu,\tfrac14)$ and the identity parameter map (Fisher defect $1-4=-3$, so no
Markov morphism).  And contraction compares the image of the fine orbit, not a recomputed
coarse orbit: on $\R^2$ with $\mathsf F_\ell$ the identity and the Euclidean metric at
both depths, $\Fenergy_\ell=\tfrac12(z_1^2+z_2^2)$ and
$\Fenergy_{\ell+1}=\tfrac12(z_1^2+4z_2^2)$ give the fine history $(e^{-t},e^{-t})$ of
duration $\sqrt2$ and the recomputed coarse history $(e^{-t},e^{-4t})$ of duration
$\int_0^1\sqrt{1+16z^6} dz>\sqrt2$, and
\eqref{eq:hist-oriented-vector-field-semiconjugacy} fails because
$(z_1,z_2)=\mathsf a(z_1,4z_2)$ forces $\mathsf a=1$ and $z_2=0$.  \status{ESTABLISHED}

\corollaryheading{When fiberwise contraction does lift}{thm:hist-pointwise-contraction-lift}
Let $\mathcal Q_\ell$ and $\mathcal Q_{\ell+1}$ carry the channel-weighted metrics
$\mathsf G_\ell=\sum_xw_xg^F_x$ and $\mathsf G_{\ell+1}=\sum_xw_x\bar g^F_x$ with the same
positive weights, and let $\mathsf F_\ell$ act channelwise by Markov coarse morphisms,
$(\mathsf F_\ell Q)_x=\mathscr K_x(Q_x)$.  Then for every $Z\in T_Q\mathcal Q_\ell$,
\begin{equation}
\lVert Z\rVert^2_{\mathsf G_\ell}
-\lVert T\mathsf F_\ell Z\rVert^2_{\mathsf G_{\ell+1}}
=\sum_xw_x\Delta_F^{\mathscr K_x}(Z_x,Z_x)
\geq0,
\label{eq:hist-pointwise-contraction-lift}
\end{equation}
so the image of every $C^1$ fine curve has no greater speed and no greater duration on
every subinterval, and \eqref{eq:hist-duration-criterion} holds along every orbit of an
oriented semiconjugacy. On a subinterval of finite fine duration, equality holds exactly
when, for every channel and almost every parameter value, the fine score in the
direction $Z_x$ is measurable with respect to the coarse output.
\status{ESTABLISHED}

\noindent\emph{Proof.}
Differentiate the channelwise action, insert it into $\mathsf G_{\ell+1}$, and subtract
from $\mathsf G_\ell$; the difference is the weighted sum of the vertical defects, which
are positive semidefinite with the stated equality case by
\cref{thm:pb-pullback-fisher-defect}.  Monotonicity of the square root and of the
integral gives the speed and duration statements.  $\square$

\section{Semiconjugacy on symmetry quotients}
\label{sec:hist-natural-gradient-semiconjugacy}

Throughout this section $\mathscr V_\ell$ and $\mathscr V_{\ell+1}$ are natural-gradient
fields of $C^2$ objectives for positive-definite metrics.

\definitionheading{Functional compatibility}{hyp:hist-functional-compatibility}
A coarse configuration map $\mathsf F_\ell$ is functionally compatible on an open set
$\mathcal U\subseteq\mathcal Q_\ell$ when there is a $C^1$ function $\chi_\ell$ with
$\chi_\ell'>0$ such that $\Fenergy_\ell=\chi_\ell\circ\Fenergy_{\ell+1}\circ\mathsf
F_\ell$ on $\mathcal U$.

\propositionheading{Semiconjugacy on symmetry quotients}{prop:hist-natural-gradient-sufficiency}
Equal objectives do not intertwine natural gradients: with $\mathsf F_\ell$ the identity
of $\R^2$, $\Fenergy_\ell=\Fenergy_{\ell+1}=\tfrac12(z_1^2+2z_2^2)$, and metrics
$\operatorname{diag}(1,1)$ and $\operatorname{diag}(1,\mathrm c)$ with $\mathrm c\ne1$,
\eqref{eq:hist-oriented-vector-field-semiconjugacy} fails wherever $z_1z_2\ne0$.
Functional compatibility on $\mathcal U$ forces $d\Fenergy_\ell$ to vanish on
$\ker T\mathsf F_\ell$ at every point of $\mathcal U$, so when $\mathsf F_\ell$ is a
submersion, $\Fenergy_\ell$ is constant on each connected piece of each fiber of
$\mathsf F_\ell$ in $\mathcal U$.  If $\mathsf F_\ell$ is a surjective submersion,
functionally compatible on $\mathcal U$, and horizontally conformal with dilation
$\mathsf d_\ell>0$, meaning that
$\mathsf G_{\ell+1}(T\mathsf F_\ell Z_1,T\mathsf F_\ell Z_2)=\mathsf d_\ell^2\mathsf
G_\ell(Z_1,Z_2)$ for $Z_1,Z_2\in\operatorname{Hor}_Q=(\ker T_Q\mathsf F_\ell)^\perp$, then
\eqref{eq:hist-oriented-vector-field-semiconjugacy} holds on $\mathcal U$ with
\begin{equation}
\mathsf a_\ell(Q)
=\chi_\ell'\left(\Fenergy_{\ell+1}(\mathsf F_\ell Q)\right)
\mathsf d_\ell(Q)^2>0 .
\label{eq:hist-conformal-semiconjugacy-factor}
\end{equation}
In particular the quotient map of a free, proper, isometric action that preserves
$\Fenergy_\ell$, with the quotient metric and descended objective of
\cref{prop:hist-invariant-history-horizontal}, is an oriented semiconjugacy with
$\mathsf a_\ell\equiv1$; and at $c_*$ the identity map from the mean coordinate of the
fiber $\{\mathcal N(\mu,1)\}$ to that of $\{\mathcal N(\mu,2)\}$, realized by the kernel
$y\mapsto\mathcal N(y,1)$, is horizontally conformal with $\mathsf d_\ell^2=\tfrac12$ and
gives $\mathsf a_\ell=\tfrac12$ for $\Fenergy_\ell=\Fenergy_{\ell+1}=\mu^2/2$, an
injective case that does not exercise the fiber condition.  The route does not reach
maps that discard directions on which the fine objective depends: if
$\Fenergy_{\ell+1}(Q')=\inf\{\Fenergy_\ell(Q):\mathsf F_\ell Q=Q'\}$, $\mathsf F_\ell$ is
a submersion with positive-dimensional fibers, and each fiber contains at most one
minimizer of $\Fenergy_\ell$, then compatibility with $\chi_\ell=\operatorname{id}$ holds
on no open set; for $\mathsf F_\ell(f,s)=s$ and
$\Fenergy_\ell(f,s)=\tfrac12(f-s)^2+\tfrac12s^2$ it holds for no $\chi_\ell$ on any open
set.  \status{ESTABLISHED}

The proof is in \cref{sec:proofs-fiber-geometry}.

For maps that discard directions on which the fine objective depends, the exact route
runs through the reduced slow flow, which on the set of conditional minimizers is
conjugate to the coarse natural-gradient flow (\cref{prop:obs-rate-separated-flow},
\eqref{eq:obs-reduced-semiconjugacy}); whether the fine or rate-separated flow follows
it is \cref{open:obs-slow-manifold}.

\section{Coarse-graining depth, orbit parameter, and duration}
\label{sec:hist-rg-depth-not-time}

A tower of coarse configurations $Q^{(\ell+1)}=\mathsf F_\ell(Q^{(\ell)})$,
$\ell=0,1,2,\ldots$, is indexed by the coarse-graining depth $\ell$, which is not a
parameter on a history.  A history can be present at every depth, as $Q^{(\ell)}(t)$
with accumulated duration $\mathfrak t^{(\ell)}(t)$; when
$Q^{(\ell+1)}(t)=\mathsf F_\ell(Q^{(\ell)}(t))$ and $\mathsf F_\ell$ acts by Markov coarse
morphisms as in \cref{thm:hist-pointwise-contraction-lift}, the record-clock comparison
applies between adjacent depths, while for a general oriented semiconjugacy the exact
comparison is \cref{thm:hist-duration-criterion}, and independently recomputed
histories also need the noncollapse conditions of \cref{prop:hist-noncollapse}.  The
Markov property of a configuration arrow can fail even for natural meta-agent maps:
no parameter-independent Markov kernel sends every product of two normal laws to
their joint mean-and-variance barycenter, as Supplement~G proves.

\propositionheading{Typed non-identification of depth, orbit parameter, and duration}{prop:hist-coordinate-independence}
Coarse-graining depth $\ell$, an orbit parameter $t$, and Fisher duration
$\mathfrak t^{(\ell)}$ are distinct: $\ell$ is a discrete index, $t$ a chosen oriented
parameter on a selected orbit, and $\mathfrak t^{(\ell)}$ the arc length for a chosen
metric from a chosen origin.  Once the metric, origin, orientation, and parameterized path
are chosen, $\mathfrak t^{(\ell)}$ is a function of $t$, but it is not determined by the
unparameterized orbit or by $\ell$: changing the metric rescales it and changing the
origin shifts it.  Strict monotonicity of $\mathfrak t^{(\ell)}$ does not make it a
regular coordinate: for $\mathsf G=dz^2$ on $\R$ and the path $z(t)=t^2$ on $[0,1]$,
$\mathfrak t(t)=t^2$, whose inverse is not differentiable at the origin, and a zero-speed
subinterval destroys strict monotonicity.  \status{ESTABLISHED}

\noindent\emph{Proof.}
Only the metric, origin, orientation, and parameterized path produce a duration
function.  Metric dependence is homogeneity of the norm in \eqref{eq:hist-fisher-speed},
and the example is $\int_0^t2t''dt''=t^2$ with the square root not differentiable at
zero.  $\square$

\section*{Physicist's summary}

The Fisher metric on an agent's belief and model law fibers is a gauge-invariant ruler and
the second-order term of the relative entropy, so a loop holonomy near the identity
appears as a quadratic form on the Lie algebra that vanishes on the isotropy algebra of
the law.  A Markov channel loses exactly the conditional variance of the score, so a
meta-agent that is a channel image of its constituents' history has an information clock
no faster than theirs, by a ratio that is open.  A chosen metric turns the free energy
into an orbit, Fisher arc length removes its parameterization, and a belief-model rate
ratio, unlike a scalar mobility, can bend the orbit; an orthogonal clock needs the
gradient norm to be locally a function of the free energy, a unit-speed clock only a
global cross-section.  By the parity witness, joint laws, Fisher metrics, and free
energies are not functions of marginals.  A recomputed coarse flow is the image of the
fine flow only under an oriented semiconjugacy, which holds with rate one on symmetry
quotients; information-discarding maps are compared through the reduced slow flow of the
next chapter.

\chapter{Two-Timescale Dynamics}
\label{ch:two-timescales}

This chapter asks what the collective free energy of \cref{ch:elbo} supplies for
separating fast inference of beliefs from slower evolution of models, priors, and
parents: descent for every positive-definite block mobility, local in a product family
(\cref{prop:obs-local-block-descent}); a reduced slow flow conjugate to its
tangent lift on the graph of the conditional optimum
(\cref{lem:obs-envelope-slow-flow,prop:obs-rate-separated-flow}), the exact route of
\cref{ch:relational-inference} for information-discarding maps; and a Markov parent
history under instantaneous within-parent refresh
(\cref{thm:rg-memory-prepared-parent}), whose failure under finite refresh
(\cref{rem:rg-memory-partial-refresh}) is measured in \cref{ch:agent-network-rg}. The
slow block (\cref{def:obs-fast-slow-blocks}) moves by accepted M-coordinates
(\cref{def:elbo-generalized-m-acceptance}) or by the rate-separated flow; the formation
family of \cref{ch:finite-depth-formation} realizes one scalar slow parameter
\eqref{eq:rg-formation-own-law-dynamics}, not slow per-agent model laws, while on
the four-cycle of that chapter, whose link translations this chapter writes
$\check\omega$, slow model means are a supplied rate, and the model
channel lags the belief channel only when the links are slow as well, at a rate of
order at most the model rate
(\cref{prop:obs-slow-model-link-dichotomy}, the held-sheet instance of
\cref{open:obs-slow-manifold}). The finite
model of the last section uses fine states $\xi\in\Xi$, parent values $z\in\mathsf Z$,
a parent map $\operatorname{par}$, a prior $p^0$, record laws $\psi_\xi$, a refresh
fraction $\omega$, and sans-serif kernels acting on row vectors from the right.

\paragraph{Scope.}
The collective free energy does not choose the fast--slow split, the rate ratio, or
the block mobilities, and a rate ratio is an anisotropic mobility that can change the
selected orbit (\cref{prop:hist-scalar-mobility-orbit}). The chapter does not claim
that the rate-separated flow tracks the reduced flow for a recognition family of the
program beyond the held matching sheets of
\cref{prop:obs-slow-model-link-dichotomy}, that the sampled dynamics with source
jumps track it, that a sampled episode with source jumps at the temperature of
\cref{def:rg-formation-cycle-experiment} keeps a matched pair belief-coherent and
model-incoherent on an interval of order $1/\varepsilon_{\mathrm m}$ with probability
bounded below, that the discrete alternating scheme converges, that the parent sweep
converges at a rate uniform in the number of fine states or in the smallest posterior
weight, or that a static parent datum has closed dynamics
(\cref{open:obs-slow-manifold,open:elbo-alternating-convergence,open:rg-pointwise-parent-dynamics}).
\status{NOT-CLAIMED}

\section{Block descent at fixed parameter}
\label{sec:local-natural-gradient}

\propositionheading{Block descent of the collective free energy}{prop:obs-local-block-descent}
Let $\lambda\mapsto Q_\lambda$ be a family of recognition laws on an open domain
$\Lambda=\prod_{i\in V}\Lambda_i$ with block coordinates $\lambda=(\lambda_i)_i$,
such that $F(\lambda):=\Fenergy^{\mathrm{ext}}_o(Q_\lambda)$ is real and $C^1$ on
$\Lambda$, and let $\mathsf G_i(\lambda)$ be continuous fields of symmetric
positive-definite block metrics, with positive block rates absorbed into the
metrics and blocks of zero rate held fixed. Let $\lambda(t)$ solve
\begin{equation}
\dot\lambda_i=-\mathsf G_i^{-1}\nabla_{\lambda_i}F
\label{eq:obs-local-natural-gradient}
\end{equation}
for the moving blocks $i\in V'\subseteq V$ and remain in $\Lambda$.
\begin{enumerate}
\item[(i)] Along the solution
\begin{equation}
\frac{d}{dt}F=-\sum_{i\in V'}(\nabla_{\lambda_i}F)^\top\mathsf G_i^{-1}(\nabla_{\lambda_i}F)\leq0.
\label{eq:obs-global-dissipation}
\end{equation}
\item[(ii)] If the family is a regular product family
$Q_\lambda=\bigotimes_iQ_{i,\lambda_i}$ with
$\E_{Q_{-i}}\lvert\log Z_{\{i\}}(Y_{-i})\rvert<\infty$ for every $i$ and $\lambda$,
then $\nabla_{\lambda_i}F=\nabla_{\lambda_i}\overline\Fenergy_i$ for the
outside-averaged local free energy \eqref{eq:obs-outside-averaged-local-vfe}, so
an agent that descends its own local free energy moves along a block of a
descent flow of the collective one. The Fisher metric of such a family is block
diagonal, so with $\mathsf G_i$ equal to the block Fisher metrics the flow is the
Fisher natural gradient of $F$.
\item[(iii)] Outside product families moving $\lambda_i$ can change $Q_{-i}$,
and the identification in (ii) needs the conditional lift of
\cref{hyp:hist-exact-vfe-lift}, whose outside-averaged objective
\eqref{eq:hist-outside-averaged-local-vfe} has the differential of the collective
free energy restricted to fixed outside marginal. The Fisher metric then need not
be block diagonal: for the bivariate normal location family with unit variances
and correlation $\mathrm c_2$, $0<\lvert\mathrm c_2\rvert<1$, its off-diagonal entry
is $-\mathrm c_2/(1-\mathrm c_2^2)$.
\end{enumerate}
\status{ESTABLISHED}

\noindent\emph{Proof.}
Part (i) is the chain rule for $F$ along the solution. For (ii), a change of
$\lambda_i$ leaves $Q_{-i}$ fixed, and \cref{thm:obs-local-global-potential} with
$B=\{i\}$, whose finiteness hypotheses hold because $F$ is real, shows that $F$
and $\overline\Fenergy_i$ differ by a quantity independent of $\lambda_i$. Scores
of distinct factors of a product law are independent with mean zero, so their
cross covariances vanish. For (iii), the Fisher metric of the location family is
the inverse covariance matrix. $\square$

\remarkheading{Simultaneous block updates}{rem:obs-simultaneous-block-updates}
The flow above, \eqref{eq:obs-damped-block-update}, and
\cref{cor:obs-potential-game} move one block at a time or move continuously; a
discrete schedule that replaces several blocks at once is covered only in the
following cases. (a) Let $B_1,\dots,B_k$ be disjoint blocks with union $U$ such
that no record scope meets two of them and the baseline makes
$Y_{B_1},\dots,Y_{B_k}$ conditionally independent given $Y_{-U}$, as a product
baseline does; for singleton blocks under a product baseline this is an
independent set of the record hypergraph. For a law
$Q=Q_{-U}\bigotimes_lr_{B_l}(\cdot\given y_{-U})$ with
$\Fenergy^{\mathrm{ext}}_o(Q)<\infty$, replacing every $r_{B_l}$ at once by a damped
kernel of \cref{prop:elbo-exact-e-coordinate}(iv) gives the law reached by the
sequential sweep in any order, so it does not increase $\Fenergy^{\mathrm{ext}}_o$;
and under the hypotheses of \cref{cor:obs-potential-game}, changing the strategies
of the agents of such a set at once changes $\Fenergy^{\mathrm{ext}}_o$ by the sum
of their unilateral changes. A chromatic schedule that updates the classes of a
partition of the agents into such sets in turn therefore descends whenever each
unilateral change does. (b) In \cref{cor:obs-potential-game} let the strategy set of
agent $i$ be a regular minimal exponential family
$q_{i,\eta}(dy_i)=\exp\{\langle\eta,\mathsf T_i(y_i)\rangle-\mathsf A_i(\eta)\}\nu_i(dy_i)$,
$\eta\in\mathsf N_i$, with mean parameter $\mu=\nabla\mathsf A_i(\eta)$, and let it
be conditionally conjugate: with the other strategies held,
$\overline\Fenergy_i(q_{i,\eta};r_{-i})=\KL(q_{i,\eta}\Vert q_{i,\eta^*})+\mathrm c_i$
for some $\eta^*\in\mathsf N_i$ and $\mathrm c_i$ independent of $\eta$, so that
$q_{i,\eta^*}$ is the best response. Damping in mean coordinates, to the member with
mean $(1-\mathrm w)\mu+\mathrm w\mu^*$ when that mean belongs to the family, which it
always does for Gaussian families, satisfies the two inequalities of
\eqref{eq:obs-damped-block-update} with the best response in place of the local
posterior. Damping in natural coordinates, to
$\eta_{\mathrm w}=(1-\mathrm w)\eta+\mathrm w\eta^*$, does not increase
$\Fenergy^{\mathrm{ext}}_o$ either, but it need not satisfy the first inequality:
for one agent with a Poisson latent, posterior $\operatorname{Poi}(1)$, natural
parameter the logarithm of the rate, $\eta=-4$, $\eta^*=0$ and $\mathrm w=1/2$, the
relative entropy to the posterior is $1-3e^{-2}=0.594\ldots$, above
$(1-5e^{-4})/2=0.454\ldots$. (c) In the conjugate case the natural gradient of
$\Fenergy^{\mathrm{ext}}_o$ in $\eta_i$ for the Fisher block metric is
$\eta_i-\eta_i^*$, so simultaneous natural-coordinate damping of all agents with
weight $\mathrm w$ is the explicit Euler step of size $\mathrm w$ of
\eqref{eq:obs-local-natural-gradient}, and $\mathrm w=1$ is the simultaneous best
response. For blocks that share records it can increase the free energy. Take three
agents with scalar latents, baseline $\mathcal N(0,5/2)$ for each, and one record
$o_a\sim\mathcal N(y_1+y_2+y_3,5/3)$, so the posterior is Gaussian with precision
$\Lambda=\tfrac25I+\tfrac35\mathbf 1\mathbf 1^\top$ and some mean $\bar y^\star$, and
take product Gaussian strategies. At unit variances and means $\bar y^\star+\mathbf 1$,
\begin{equation}
\KL(Q\Vert\boldsymbol\Pi_{\theta,V,o,X})=\frac{33}{10}+\frac12\log\frac{125}{44},
\label{eq:obs-simultaneous-witness}
\end{equation}
and the simultaneous best responses keep unit variances and move the means to
$\bar y^\star-\tfrac65\mathbf 1$, where the relative entropy is
$\tfrac{594}{125}+\tfrac12\log\tfrac{125}{44}$; each repetition multiplies the
first term by $36/25$, whereas the sequential sweep never increases the free energy.
The update of weight $\mathrm w$ increases it along $\mathbf 1$ exactly when
$\mathrm w>10/11=2/\mathrm c_{\mathsf G}$, where $\mathrm c_{\mathsf G}=11/5$ is the
largest eigenvalue of the Hessian $\Lambda$ relative to the Fisher metric $I$ of the
means at unit variance, which is the threshold of
\eqref{eq:hist-discrete-descent-bound}. A simultaneous step of blocks that share
records is therefore covered only by a step bound of that kind or by a step
accepted on the same functional, as after
\cref{prop:elbo-finite-step-nonmonotonicity}. \status{ESTABLISHED}

\noindent\emph{Proof.}
For (a), the incident likelihood $L_{B_l,o}$ of \eqref{eq:obs-block-incident-likelihood}
reads no coordinate of another $B_m$, and conditional independence gives the baseline
conditional of $Y_{B_l}$ given $Y_{-B_l}$ a version that reads only $y_{-U}$; by
\eqref{eq:obs-local-posterior} and \eqref{eq:obs-local-vfe} the local evidence,
posterior, and free energy of $B_l$ then read only $y_{-U}$. Under $Q$ the conditional
of $Y_{B_1}$ given $Y_{-B_1}$ is $r_{B_1}(\cdot\given y_{-U})$, so
\cref{prop:elbo-exact-e-coordinate}(iv) applies with $B=B_1$, and the damped kernel
again reads only $y_{-U}$, so the new law has the same form and finite free energy;
induction over the blocks, in any order, ends at
$Q_{-U}\bigotimes_lr^{\mathrm w_l}_{B_l}$. For product laws, the integrand of
\eqref{eq:obs-outside-averaged-local-vfe} for agent $i$ does not read $y_j$ for
another agent $j$ of the set, so $\overline\Fenergy_i$ does not depend on $r_j$;
changing the strategies one agent at a time, \cref{cor:obs-potential-game} gives
each change of $\Fenergy^{\mathrm{ext}}_o$ as the change of $\overline\Fenergy_i$ at
the original strategies of the set, and the changes add. For (b), $\mathsf A_i$ is
finite near $\eta$ and $\eta^*$, so $\mathsf T_i$ is integrable under both members,
with means $\mu=\nabla\mathsf A_i(\eta)$ and $\mu^*$, and the Fisher metric of the
natural chart is $\nabla^2\mathsf A_i$
\citep[Theorems~7.1, 7.2 and~8.1]{BarndorffNielsen1978}. By
\cref{cor:obs-potential-game}, $\Fenergy^{\mathrm{ext}}_o$ changes by the change of
$\KL(\cdot\Vert q_{i,\eta^*})$. The mixture
$\mathrm m=(1-\mathrm w)q_{i,\eta}+\mathrm wq_{i,\eta^*}$ has mean
$\mu_{\mathrm w}=(1-\mathrm w)\mu+\mathrm w\mu^*$; if $q_{i,\eta'_{\mathrm w}}$ has
that mean, $\log(q_{i,\eta'_{\mathrm w}}/q_{i,\eta^*})$ is affine in $\mathsf T_i$, so
its expectations under $\mathrm m$ and under $q_{i,\eta'_{\mathrm w}}$ agree, and
\[
(1-\mathrm w)\KL(q_{i,\eta}\Vert q_{i,\eta^*})\geq\KL(\mathrm m\Vert q_{i,\eta^*})
=\KL(\mathrm m\Vert q_{i,\eta'_{\mathrm w}})+\KL(q_{i,\eta'_{\mathrm w}}\Vert q_{i,\eta^*})
\geq\KL(q_{i,\eta'_{\mathrm w}}\Vert q_{i,\eta^*}),
\]
the first inequality by convexity of relative entropy in its first argument. This is
the first inequality of \eqref{eq:obs-damped-block-update} after adding the
constant, and the second follows because the best response minimizes. For a
Gaussian family the means $(\mu,\mathrm M_2)$ of $(Y,YY^\top)$ belong to the family exactly when
$\mathrm M_2-\mu\mu^\top\succ0$, that is, when the matrix with blocks $1$, $\mu^\top$, $\mu$, $\mathrm M_2$ is
positive definite, a convex condition. With $\Delta\eta=\eta-\eta^*$ and
$\mathrm f(t)=\KL(q_{i,\eta^*+t\Delta\eta}\Vert q_{i,\eta^*})
=\mathsf A_i(\eta^*)-\mathsf A_i(\eta^*+t\Delta\eta)+t\langle\nabla\mathsf A_i(\eta^*+t\Delta\eta),\Delta\eta\rangle$,
one has $\mathrm f'(t)=t\Delta\eta^\top\nabla^2\mathsf A_i(\eta^*+t\Delta\eta)\Delta\eta\geq0$ on the segment,
which lies in the convex domain $\mathsf N_i$, and $\eta_{\mathrm w}$ is the point
$t=1-\mathrm w$. In the Poisson family $\mathsf A(\eta)=e^\eta$ and
$\mathrm f$ at $\eta^*=0$ is $1-(1-\eta)e^{\eta}$ evaluated at $\eta=-2$ and
$\eta=-4$. For (c), differentiating the Bregman form gives
$\nabla_{\eta}\KL(q_{i,\eta}\Vert q_{i,\eta^*})=\nabla^2\mathsf A_i(\eta)(\eta-\eta^*)$,
which by \cref{prop:obs-local-block-descent}(ii) is $\nabla_{\eta_i}\Fenergy^{\mathrm{ext}}_o$;
multiplying by the inverse Fisher metric gives $\eta_i-\eta_i^*$, and the Euler step
of size $\mathrm w$ is $\eta_{\mathrm w}$. In the witness the record adds precision
$\tfrac35\mathbf 1\mathbf 1^\top$ to the baseline precision $\tfrac25I$, so
$\Lambda$ has unit diagonal and eigenvalues $\tfrac25,\tfrac25,\tfrac{11}5$ with
$\det\Lambda=44/125$. For a product law with variances $\mathrm v_i$ and mean error $\boldsymbol\delta$,
$\KL(Q\Vert\boldsymbol\Pi_{\theta,V,o,X})
=\frac12[\sum_i\mathrm v_i+\boldsymbol\delta^\top\Lambda\boldsymbol\delta-3-\log\det\Lambda-\sum_i\log\mathrm v_i]$, whose
best response for agent $i$ is $\mathrm v_i=1$ and $\delta_i=-\sum_{j\ne i}\Lambda_{ij}\delta_j$; at
$\boldsymbol\delta=\mathbf 1$ this is $\delta_i=-\tfrac65$, and
$\mathbf 1^\top\Lambda\mathbf 1=\tfrac{33}5$ gives both values. The weighted update
is $\boldsymbol\delta\mapsto\boldsymbol\delta-\mathrm w\Lambda\boldsymbol\delta$, which multiplies the component along
$\mathbf 1$ by $1-\tfrac{11}5\mathrm w$. Each sequential best response is a
unilateral improvement of \cref{cor:obs-potential-game}. $\square$

\section{Fast recognition and slow model coordinates}
\label{sec:obs-fast-slow}

\definitionheading{Fast and slow coordinate blocks}{def:obs-fast-slow-blocks}
Let $\mathcal U_{\rm f}\subseteq\R^{n_{\rm f}}$ and
$\mathcal U_{\rm s}\subseteq\R^{n_{\rm s}}$ be open, and write
$s=(s_{\rm m},\theta)\in\mathcal U_{\rm s}$. A fast--slow coordinatization is a
map $(f,s)\mapsto(Q_{f,s_{\rm m}},\theta)$ into population recognition laws and
generative parameters, together with the real-valued objective
\begin{equation}
F(f,s):=\Fenergy^{\mathrm{ext}}_o(Q_{f,s_{\rm m}};\theta)
=-\log p_{\theta,V}(o\given X)
+\KL\bigl(Q_{f,s_{\rm m}}\Vert\boldsymbol\Pi_{\theta,V,o,X}\bigr),
\label{eq:obs-fast-slow-objective}
\end{equation}
in which the generative argument of \eqref{eq:obs-collective-vfe} is explicit. The
fast coordinate $f$ carries the belief-channel recognition coordinates and, when
present, the source-label rows of \cref{sec:local-attention-elbo}; the slow
coordinate carries the model-channel recognition coordinates $s_{\rm m}$, which move
a recognition law, and the generative parameter $\theta$, which moves the fixed joint
between evaluations and never inside one, so a composite $\theta(Q)$ inside one
evaluation of \eqref{eq:obs-fast-slow-objective} stays excluded by
\cref{req:gen-typing-prohibition}. Because \eqref{eq:obs-fast-slow-objective}
evaluates one observation $o$ at every $\theta$ in a continuum, the coordinatization
also declares, as model data in the sense of
\cref{hyp:prob-regular-conditional-convention}, a version
$(\theta,o,y)\mapsto p_{\theta,V}(o,y\given X)$ jointly measurable in its arguments,
with evidence $0<p_{\theta,V}(o\given X)<\infty$ at the observed $o$ for every
$\theta$ occurring in $\mathcal U_{\rm s}$ and posterior slices computed from that
version by \eqref{eq:prob-rcp-density}; real values and differentiability of $F$
in $\theta$ are statements about this version. The assignment of coordinates to
the two blocks is a declaration.

\lemmaheading{Envelope identity and reduced slow descent}{lem:obs-envelope-slow-flow}
Let $F\in C^1(\mathcal U_{\rm f}\times\mathcal U_{\rm s})$, let
$\mathcal U_0\subseteq\mathcal U_{\rm s}$ be open, and let
$f^\star:\mathcal U_0\to\mathcal U_{\rm f}$ be a $C^1$ map with
$F(f^\star(s),s)=\min_{f\in\mathcal U_{\rm f}}F(f,s)$ for every $s\in\mathcal U_0$.
The reduced functional
\begin{equation}
F_{\rm red}(s):=\min_{f\in\mathcal U_{\rm f}}F(f,s)=F(f^\star(s),s)
\label{eq:obs-reduced-functional}
\end{equation}
satisfies the envelope identity
\begin{equation}
dF_{\rm red}(s)=\partial_sF(f^\star(s),s),\qquad s\in\mathcal U_0.
\label{eq:obs-envelope-identity}
\end{equation}
For every continuous field $\mathsf G_{\rm s}$ of symmetric positive-definite
metrics on $\mathcal U_{\rm s}$, the reduced slow flow
\begin{equation}
\dot s=-\mathsf G_{\rm s}(s)^{-1}\partial_sF(f^\star(s),s)
\label{eq:obs-reduced-slow-flow}
\end{equation}
is the $\mathsf G_{\rm s}$ natural-gradient flow of $F_{\rm red}$, and along the
lifted curve $(f^\star(s(t)),s(t))$ the free energy decreases at the rate
$dF_{\rm red}\cdot\mathsf G_{\rm s}^{-1}dF_{\rm red}$. A $C^1$ branch exists near
$s_0\in\mathcal U_0$ whenever $F\in C^2$, $\partial^2_{ff}F(f^\star(s_0),s_0)$ is
positive definite, and $f^\star$ is continuous at $s_0$. If $\mathcal U_{\rm f}$ is
replaced by a compact set on which $F$ and $\partial_sF$ are continuous and the
conditional minimizer is unique, \eqref{eq:obs-envelope-identity} holds without
differentiability of $f^\star$. \status{ESTABLISHED}

\noindent\emph{Proof.}
At the minimizer $\partial_fF(f^\star(s),s)=0$, so the chain rule gives
$dF_{\rm red}=\partial_fF\cdot Df^\star+\partial_sF=\partial_sF$, and substitution in
\eqref{eq:obs-reduced-slow-flow} gives the flow and its rate. For the branch, the
implicit function theorem applied to $\partial_fF=0$ at $(f^\star(s_0),s_0)$ gives a
$C^1$ map whose graph is the whole critical set near that point, and continuity places
$f^\star(s)$ on it for $s$ near $s_0$. In the compact case, uniqueness and compactness
make $f^\star$ continuous, since a limit point $f'$ of $f^\star(s_n)$ with $s_n\to s$
satisfies $F(f',s)\leq F(f^\star(s),s)$. With $s'=s+\Delta s$, the difference
$F_{\rm red}(s')-F_{\rm red}(s)$ lies between
$F(f^\star(s'),s')-F(f^\star(s'),s)$ and $F(f^\star(s),s')-F(f^\star(s),s)$, and by the
mean value theorem and uniform continuity of $\partial_sF$ on a compact neighborhood
both equal $\partial_sF(f^\star(s),s)\Delta s+o(\lvert\Delta s\rvert)$; this is the
theorem of \citet[Ch.~III]{Danskin1967} in the unique-minimizer case. $\square$

\propositionheading{Rate-separated natural gradient and its reduced flow}{prop:obs-rate-separated-flow}
Let $F\in C^1(\mathcal U_{\rm f}\times\mathcal U_{\rm s})$, let
$\mathsf G_{\rm f}(f,s)$ and $\mathsf G_{\rm s}(s)$ be continuous symmetric
positive-definite block metrics with $\mathsf G_{\rm s}$ independent of $f$, and
for $\varepsilon>0$ set
$\mathsf G_\varepsilon=\varepsilon\mathsf G_{\rm f}\oplus\mathsf G_{\rm s}$. The
$\mathsf G_\varepsilon$ natural gradient of $F$ is the rate-separated flow
\begin{equation}
\varepsilon\dot f=-\mathsf G_{\rm f}^{-1}\partial_fF,
\qquad
\dot s=-\mathsf G_{\rm s}^{-1}\partial_sF,
\label{eq:obs-rate-separated-flow}
\end{equation}
which decreases $F$ at the rate
$\varepsilon^{-1}\partial_fF\cdot\mathsf G_{\rm f}^{-1}\partial_fF
+\partial_sF\cdot\mathsf G_{\rm s}^{-1}\partial_sF$ for every $\varepsilon>0$; its
orbit can depend on $\varepsilon$. Under the hypotheses of
\cref{lem:obs-envelope-slow-flow}, let
$\mathcal S^\star=\{(f^\star(s),s):s\in\mathcal U_0\}$. At every point of
$\mathcal S^\star$ and for every $\varepsilon>0$, the fast component of
\eqref{eq:obs-rate-separated-flow} vanishes and the slow component equals
$-\mathsf G_{\rm s}^{-1}dF_{\rm red}$. The reduced field
$\mathscr V^\star=(Df^\star v,v)$ with $v=-\mathsf G_{\rm s}^{-1}dF_{\rm red}$ is
tangent to $\mathcal S^\star$, and
\begin{equation}
d\operatorname{pr}_{\rm s}\mathscr V^\star
=-\operatorname{grad}^{\mathsf G_{\rm s}}F_{\rm red}\circ\operatorname{pr}_{\rm s}
\quad\text{on }\mathcal S^\star.
\label{eq:obs-reduced-semiconjugacy}
\end{equation}
Thus $(\operatorname{pr}_{\rm s}|_{\mathcal S^\star},1)$ is an oriented
semiconjugacy in the sense of \cref{def:hist-oriented-semiconjugacy}, with the
manifold $\mathcal S^\star$ in place of an open fine domain, and a conjugacy,
because $\operatorname{pr}_{\rm s}|_{\mathcal S^\star}$ is a diffeomorphism onto
$\mathcal U_0$. The speeds satisfy
$\lVert\mathscr V^\star\rVert^2_{\mathsf G_\varepsilon}
=\varepsilon\lVert Df^\star v\rVert^2_{\mathsf G_{\rm f}}+\lVert v\rVert^2_{\mathsf G_{\rm s}}$,
so the $\mathsf G_\varepsilon$-length of a reduced orbit is at least the
$\mathsf G_{\rm s}$-length of its image. These lengths decrease to the latter as
$\varepsilon\downarrow0$ on every compact time interval of the orbit, and on
any interval where the lifted length is finite at some positive
$\varepsilon_0$. The conjugacy and the length limit are properties of the lift
$\mathscr V^\star$ of the reduced field; they say nothing about orbits of
\eqref{eq:obs-rate-separated-flow} off $\mathcal S^\star$, whose relation to the
reduced orbit is \cref{open:obs-slow-manifold}. \status{ESTABLISHED}

\noindent\emph{Proof.}
Because $\mathsf G_\varepsilon^{-1}=\varepsilon^{-1}\mathsf G_{\rm f}^{-1}\oplus\mathsf G_{\rm s}^{-1}$,
the flow is \eqref{eq:obs-rate-separated-flow} with the stated dissipation, and the
orbit statement is \cref{prop:hist-scalar-mobility-orbit}, since
$\mathsf G_\varepsilon$ differs from $\mathsf G_1$ by an anisotropic rescaling. On
$\mathcal S^\star$, $\partial_fF=0$ and $\partial_sF=dF_{\rm red}$ by
\eqref{eq:obs-envelope-identity}, and $\mathsf G_{\rm s}$ does not depend on $f$. The
field $\mathscr V^\star$ is the differential of $s\mapsto(f^\star(s),s)$ applied to
$v$, hence tangent to $\mathcal S^\star$ with slow component $v$. The length
identity is the block form of $\mathsf G_\varepsilon$. For
$0<\varepsilon\le\varepsilon_0$ the lifted speed is bounded by its
$\varepsilon_0$ value, which is integrable under the stated hypothesis and
on compact time intervals by continuity. Dominated convergence gives the
length limit. $\square$

\remarkheading{The slow graph off the optimum and under correlated metrics}{rem:obs-rate-separation-witnesses}
Off $\mathcal S^\star$ the slow component $-\mathsf G_{\rm s}^{-1}\partial_sF(f,s)$
need not be a positive multiple of
$-\operatorname{grad}^{\mathsf G_{\rm s}}F_{\rm red}(s)$: for
$F=[(f-s_1)^2+s_1^2+s_2^2]/2$ with $\mathsf G_{\rm s}=I$,
$\partial_sF=(2s_1-f,s_2)$ and $\operatorname{grad}F_{\rm red}=(s_1,s_2)$ are not
parallel when $f\ne s_1$ and $s_2\ne0$, so the projection semiconjugates the
finite-$\varepsilon$ flow to the reduced flow only where that proportionality
holds. At points of $\mathcal S^\star$ the finite-$\varepsilon$ velocity $(0,v)$ is
tangent to $\mathcal S^\star$ only where $Df^\star v=0$, so $\mathcal S^\star$ is not
invariant in general. A correlated belief--model metric requires the rate
declaration to name its slow manifold. For $F(f,s)=[(f-s)^2+s^2]/2$,
whose conditional minimizer is $f=s$, and the metric $\mathsf G$ with unit diagonal
and off-diagonal entry $\mathrm c_2\in(0,1)$, the row-rated field $\varepsilon\dot f=-(\mathsf G^{-1}dF)_f$,
$\dot s=-(\mathsf G^{-1}dF)_s$ has, in the fast time $t/\varepsilon$, the critical set
\begin{equation}
\bigl\{(\mathsf G^{-1}dF)_f=0\bigr\}=\Bigl\{f=\frac{1+2\mathrm c_2}{1+\mathrm c_2}s\Bigr\},
\label{eq:obs-correlated-slow-manifold}
\end{equation}
which attracts at rate $1/(1-\mathrm c_2)$ and differs from the conditional
minimizer. This field need not descend $F$: at $\mathrm c_2=1/2$,
every $\varepsilon\in(0,7-4\sqrt3)$, and $(f,s)=(1+2\delta,1+\delta)$ with
$\delta=(1+\varepsilon)/4$, it gives
$dF/dt=(1+\varepsilon)^2/(12\varepsilon)-4/3>0$, so for those rates it
is the natural gradient of no metric. The congruence rescaling
$\mathsf G_\varepsilon=\Delta_\varepsilon\mathsf G\Delta_\varepsilon$ with
$\Delta_\varepsilon=\operatorname{diag}(\sqrt\varepsilon,1)$ gives a natural
gradient that descends $F$ for every $\varepsilon$, and the critical set of its
$\varepsilon=0$ fast subsystem is the conditional-minimizer set
$\{\partial_fF=0\}$. \status{ESTABLISHED}

\noindent\emph{Proof.}
The first two statements are direct differentiation. For the correlated metric,
$\mathsf G^{-1}$ has diagonal entries $1/(1-\mathrm c_2^2)$ and off-diagonal entries
$-\mathrm c_2/(1-\mathrm c_2^2)$, and $dF=(f-s,2s-f)$, so $(\mathsf G^{-1}dF)_f$ vanishes on the displayed line, and
its derivative in $f$ is $(1+\mathrm c_2)/(1-\mathrm c_2^2)=1/(1-\mathrm c_2)$.
Substituting the stated point into
$dF/dt=-\varepsilon^{-1}\partial_fF(\mathsf G^{-1}dF)_f-\partial_sF(\mathsf G^{-1}dF)_s$
at $\mathrm c_2=1/2$ gives the displayed value, which is positive exactly when
$\varepsilon^2-14\varepsilon+1>0$. Under congruence, $\mathsf G_\varepsilon$ is
positive definite for every $\varepsilon>0$, and in the fast time the fast
component of its natural gradient is
$-(1-\mathrm c_2^2)^{-1}(\partial_fF-\mathrm c_2\sqrt\varepsilon\partial_sF)$, whose
zero set at $\varepsilon=0$ is $\{\partial_fF=0\}$. $\square$

\propositionheading{Slow model means on the four-cycle and the model-link dichotomy}{prop:obs-slow-model-link-dichotomy}
In \cref{def:rg-formation-cycle-experiment} let the belief means $\bar k$ move at
rate one, the model means $\bar m$ at rate $\varepsilon_{\mathrm m}\in[0,1]$, and the
link translations $\check\omega$ at rate $\varepsilon_{\check\omega}\in[0,1]$: between clock times
\begin{equation}
d\mathbf x=-\bigl(\partial_{\bar k}U_S,\varepsilon_{\mathrm m}\partial_{\bar m}U_S,
\varepsilon_{\check\omega}\partial_{\check\omega} U_S\bigr)dt
+\sqrt{2\mathrm T}\bigl(dW^{\bar k},\varepsilon_{\mathrm m}^{1/2}dW^{\bar m},
\varepsilon_{\check\omega}^{1/2}dW^{\check\omega}\bigr)
\label{eq:obs-cycle-slow-model-kinetics}
\end{equation}
with one standard Brownian motion $W$ in $\R^{12}$, and the clock redraws the
assignment as before, leaving $\mathbf x$ unchanged. The rates are supplied, since
the Fisher metric of $\mathcal N(\bar m_i,\sigma^2_{\mathrm{ev}})$ in $\bar m_i$ is
$1/\sigma^2_{\mathrm{ev}}$ and a Fisher natural gradient would move the model means at
rate $1/2$. For $i\in\{0,2\}$ write $\bar k^\pm_i=(\bar k_{i+1}\pm\bar k_i)/\sqrt2$
and $\bar m^\pm_i=(\bar m_{i+1}\pm\bar m_i)/\sqrt2$, so that
$r^b_i=\sqrt2\bar k^-_i-\check\omega_i$ and $r^m_i=\sqrt2\bar m^-_i-\check\omega_i$.
\begin{enumerate}
\item[(i)] \emph{Held matching.} On the held assignment $S^{(0)}$ a Gaussian law
of $\mathbf x$ stays Gaussian and its mean follows
\eqref{eq:obs-cycle-slow-model-kinetics} without the noise; for
$\varepsilon_{\mathrm m},\varepsilon_{\check\omega}>0$, in the slow time
$\tau=\varepsilon_{\mathrm m}t$, this mean flow is \eqref{eq:obs-rate-separated-flow}
for $F=U_{S^{(0)}}$, which is the collective free energy of
\cref{prop:rg-formation-endogenous-cycle} at the Gaussian law up to terms
independent of the mean, with $\varepsilon=\varepsilon_{\mathrm m}$, $f=\bar k$,
$s=(\bar m,\check\omega)$, $\mathsf G_{\rm f}=I_4$ and
$\mathsf G_{\rm s}=I_4\oplus(\varepsilon_{\mathrm m}/\varepsilon_{\check\omega})I_4$. In the
orthonormal coordinates $(\bar k^-_i,\bar m^-_i,\check\omega_i)$ and
$(\bar k^+_i,\bar m^+_i)$ for $i\in\{0,2\}$, $\check\omega_1$, and $\check\omega_3$, the Hessian
of $U_{S^{(0)}}$ and the mobility are block diagonal, with the blocks
\begin{equation}
\begin{pmatrix}4+3\mathrm T&-2\mathrm T&-2\sqrt2\\-2\mathrm T&4+3\mathrm T&-2\sqrt2\\
-2\sqrt2&-2\sqrt2&4+\mathrm T\end{pmatrix},\qquad
\mathrm T\begin{pmatrix}3&-2\\-2&3\end{pmatrix},\qquad\mathrm T,
\label{eq:obs-cycle-held-blocks}
\end{equation}
so the mean flow is a direct sum of six flows and $r^b_i,r^m_i$ read only the block
of $e_i$. With $\varepsilon_{\check\omega}=\varepsilon_{\mathrm m}$ the split by mobility,
$f=\bar k$ and $s=(\bar m,\check\omega)$, is the only split for which the metric has the
form $\mathsf G_\varepsilon$ with blocks independent of $\varepsilon_{\mathrm m}$;
restricted to a selected-edge block it satisfies the three hypotheses of
\cref{open:obs-slow-manifold} with $\mathrm c_1=4$ for every $\mathrm T>0$, on every
closed ball whose interior contains the reduced orbit, while on the full sheet its
fast-block curvature along the pair means $\bar k^+_i$ is $3\mathrm T$.
\item[(ii)] \emph{Dichotomy at zero temperature.} Hold $S^{(0)}$, drop the
temperature terms, and start the mean flow
$\dot{\mathbf x}=-(\partial_{\bar k}V,\varepsilon_{\mathrm m}\partial_{\bar m}V,\varepsilon_{\check\omega}\partial_{\check\omega} V)$,
$V=V_{S^{(0)}}$, from $\mathbf x^\circ$, where every residual equals $-1$. Unless
$(\varepsilon_{\mathrm m},\varepsilon_{\check\omega})\in\{(0,0),(1,0)\}$, the flow matrix
of a selected-edge block has the simple eigenvalue $0$ on $(1,1,\sqrt2)$, which both
residual rows annihilate, and the two rates
\begin{equation}
\varsigma_\pm=2(1+\varepsilon_{\mathrm m}+\varepsilon_{\check\omega})
\pm2\sqrt{(1-\varepsilon_{\mathrm m})^2+\varepsilon_{\check\omega}^2},
\label{eq:obs-cycle-block-rates}
\end{equation}
and each internal residual is the combination of $e^{-\varsigma_-t}$ and
$e^{-\varsigma_+t}$ with value $-1$ and derivative
$4(\varepsilon_{\mathrm m}+\varepsilon_{\check\omega})$ (model) or $4(1+\varepsilon_{\check\omega})$
(belief) at $t=0$. With a fast link, $\varepsilon_{\check\omega}=1$, one has
$\varsigma_-\ge4-2\sqrt2$ for every $\varepsilon_{\mathrm m}\in[0,1]$ and
$|r^m_i(t)|\le e^{-(4-2\sqrt2)t}$: the link removes the model disagreement on the
belief clock even when the model means are frozen. With fast model means,
$\varepsilon_{\mathrm m}=1$ and $\varepsilon_{\check\omega}\in(0,1]$,
$r^m_i(t)=-e^{-(4+4\varepsilon_{\check\omega})t}$. With slow model means and links,
$\varepsilon_{\check\omega}=\varepsilon_{\mathrm m}\le1/2$, one has
$6\varepsilon_{\mathrm m}-2\varepsilon_{\mathrm m}^2\le\varsigma_-\le6\varepsilon_{\mathrm m}$,
$\varsigma_+\ge4$, and
\begin{equation}
|r^b_i(t)|\le(1+\varepsilon_{\mathrm m})e^{-4t}+\varepsilon_{\mathrm m},\qquad
\bigl(1-\tfrac32\varepsilon_{\mathrm m}\bigr)e^{-6\varepsilon_{\mathrm m}t}-\tfrac32\varepsilon_{\mathrm m}
\le|r^m_i(t)|\le e^{-(6\varepsilon_{\mathrm m}-2\varepsilon_{\mathrm m}^2)t}+\tfrac32\varepsilon_{\mathrm m},
\label{eq:obs-cycle-slow-link-bounds}
\end{equation}
so for $\varepsilon_{\mathrm m}\le1/25$ each matched pair has
$|r^b_i|\le2\varepsilon_{\mathrm m}$ and $|r^m_i|\ge1/2$ at every time from
$\tfrac14\log(4/\varepsilon_{\mathrm m})$ to $1/(12\varepsilon_{\mathrm m})$.
\item[(iii)] \emph{Persistence along every history.} Let $\mathrm T>0$, let
$\mathbf x_t$ be the sampled process from any initial configuration and
assignment, and let $\bar\varepsilon=\max(\varepsilon_{\mathrm m},\varepsilon_{\check\omega})>0$.
For an edge $e_i$ with $|r^m_i(\mathbf x_0)|\ge1$, numbers $\mathrm R\ge1$ and
$\mathrm R'>0$, and $0<t\le[4\bar\varepsilon(12\mathrm R+9\mathrm T\mathrm R')]^{-1}$,
the event that $|r^m_i(\mathbf x_u)|<1/2$ for some $u\le t$ while every residual
stays bounded by $\mathrm R$ and $|\mathbf x_u|\le\mathrm R'$ on $[0,t]$ has
probability at most $2\exp[-1/(192\mathrm T\bar\varepsilon t)]$, whatever the
assignment path, and at most $2e^{-1/(4\mathrm T)}$ at the largest such $t$.
\end{enumerate}
Among the three arms of (ii), a belief-agreeing, model-disagreeing interval of
order $1/\varepsilon_{\mathrm m}$ exists exactly in the arm
$\varepsilon_{\check\omega}=\varepsilon_{\mathrm m}\le1/2$: there it is the interval
of (ii), whose model disagreement (iii) protects along every assignment path up to a
time of order $1/\varepsilon_{\mathrm m}$ while residuals and configuration stay
bounded, and a fast link or fast model means remove the disagreement at a rate at
least $4-2\sqrt2$ whatever the other rate. Small rates alone do not suffice: for
$\varepsilon_{\mathrm m},\varepsilon_{\check\omega}\in[0,1/2]$,
\begin{equation}
4\varepsilon_{\mathrm m}+2\varepsilon_{\check\omega}-2\varepsilon_{\check\omega}^2
\le\varsigma_-\le4\varepsilon_{\mathrm m}+2\varepsilon_{\check\omega},
\qquad
|r^m_i(t)|\le\bigl(1+4(\varepsilon_{\mathrm m}+\varepsilon_{\check\omega})\bigr)
e^{-(4\varepsilon_{\mathrm m}+\varepsilon_{\check\omega})t},
\label{eq:obs-cycle-general-slow-rate}
\end{equation}
so $|r^m_i|\ge1/2$ at a time $\mathrm c_6/\varepsilon_{\mathrm m}$ with
$\mathrm c_6>0$ requires
$\varepsilon_{\check\omega}\le\varepsilon_{\mathrm m}\log(10)/\mathrm c_6$, and an
interval of order $1/\varepsilon_{\mathrm m}$ needs a link rate of order at most
$\varepsilon_{\mathrm m}$. \status{ESTABLISHED}

\noindent\emph{Proof.}
For (i), both endpoints of $e_0$ and of $e_2$ attend them under $S^{(0)}$ and no
receiver attends $e_1$ or $e_3$, so $V_{S^{(0)}}=2\mathcal D_{e_0}+2\mathcal D_{e_2}$
with $2\mathcal D_{e_i}=(\sqrt2\bar k^-_i-\check\omega_i)^2+(\sqrt2\bar m^-_i-\check\omega_i)^2$,
the reference term is invariant under the orthonormal change of coordinates, and
the evaluator term is
$\mathrm T\sum_{i\in\{0,2\}}[(\bar k^-_i-\bar m^-_i)^2+(\bar k^+_i-\bar m^+_i)^2]$;
collecting second derivatives gives \eqref{eq:obs-cycle-held-blocks}, and the
mobility stays diagonal because the new coordinates mix one channel at a time. On a
held assignment \eqref{eq:obs-cycle-slow-model-kinetics} is linear with constant
coefficients, so a Gaussian law stays Gaussian and its mean solves the noiseless
equation. Since $U_S(\mathbf x)=\frac12\mathbf x^\top\nabla^2U_S\mathbf x$, the
identity
$\Fenergy_{\mathrm T}(q)=\mathrm T\KL(q\Vert q^\star)+\mathrm B_{\mathrm T}-\mathrm T\log Z_{\mathrm{opt}}$
of \cref{prop:rg-formation-endogenous-cycle} at $q=\mathcal N(\mu,\Sigma)\otimes\delta_S$
gives $\Fenergy_{\mathrm T}=U_S(\mu)+\frac12\Tr(\nabla^2U_S\Sigma)-\frac{\mathrm T}2\log\det\Sigma$
up to a constant. Dividing the mean equation by $\varepsilon_{\mathrm m}$ gives the
slow-time form with the metric
$\varepsilon_{\mathrm m}I_4\oplus I_4\oplus(\varepsilon_{\mathrm m}/\varepsilon_{\check\omega})I_4$,
whose blocks are independent of $\varepsilon_{\mathrm m}$ exactly when $f$ consists
of the unit-rate coordinates. On a selected-edge block with
$\varepsilon_{\check\omega}=\varepsilon_{\mathrm m}$, $F$ is the quadratic form of the first
block of \eqref{eq:obs-cycle-held-blocks}, positive definite for $\mathrm T>0$ since
its leading minors are $4+3\mathrm T$, $(4+\mathrm T)(4+5\mathrm T)$ and
$\mathrm T(8+\mathrm T)(4+5\mathrm T)$; $\partial^2_{ff}F=4+3\mathrm T\ge4\mathsf G_{\rm f}$;
the reduced matrix is a Schur complement, hence positive definite, so $|s|$ is
nonincreasing along the reduced flow and every closed ball whose interior contains
the initial slow state serves as $\mathcal K$; and the frozen fast flow is linear
with coefficient $4+3\mathrm T$, so every fast state lies in its basin. The
pair-mean block gives the curvature $3\mathrm T$ along $\bar k^+_i$.

For (ii), the block flow matrix is $\operatorname{diag}(1,\varepsilon_{\mathrm m},\varepsilon_{\check\omega})$
times the first block of \eqref{eq:obs-cycle-held-blocks} at $\mathrm T=0$, whose
kernel is spanned by $(1,1,\sqrt2)$, on which $r^b_i=r^m_i=0$. Expanding the
determinant, its characteristic polynomial is
$\varsigma[\varsigma^2-4(1+\varepsilon_{\mathrm m}+\varepsilon_{\check\omega})\varsigma+8(2\varepsilon_{\mathrm m}+\varepsilon_{\check\omega}+\varepsilon_{\mathrm m}\varepsilon_{\check\omega})]$,
whose quadratic factor has discriminant
$16[(1-\varepsilon_{\mathrm m})^2+\varepsilon_{\check\omega}^2]$ and constant term vanishing
only at $(0,0)$, so its roots are \eqref{eq:obs-cycle-block-rates} and the
eigenvalues are distinct outside the two excluded pairs. The matrix is then
diagonalizable and the residual rows annihilate the kernel, so each internal residual
is $\mathrm c_2e^{-\varsigma_-t}+\mathrm c_3e^{-\varsigma_+t}$ with
$\mathrm c_2+\mathrm c_3=-1$ and $\varsigma_-\mathrm c_2+\varsigma_+\mathrm c_3$ equal
to minus the derivative at $t=0$, which is
$a\cdot\operatorname{diag}(1,\varepsilon_{\mathrm m},\varepsilon_{\check\omega})(-2\sqrt2,-2\sqrt2,4)^\top$
for the residual row $a$; hence $\mathrm c_3=(\varsigma_--\mathrm c_5)/(\varsigma_+-\varsigma_-)$
with $\mathrm c_5=4(\varepsilon_{\mathrm m}+\varepsilon_{\check\omega})$ for $r^m_i$ and
$\mathrm c_5=4(1+\varepsilon_{\check\omega})$ for $r^b_i$. For $\varepsilon_{\check\omega}=1$ put
$\mathrm c_4=\sqrt{(1-\varepsilon_{\mathrm m})^2+1}$, so
$\varsigma_\pm=4+2\varepsilon_{\mathrm m}\pm2\mathrm c_4$; $\varsigma_-$ increases in
$\varepsilon_{\mathrm m}$ from $4-2\sqrt2$, and
$\varsigma_-\le4+4\varepsilon_{\mathrm m}\le\varsigma_+$ because
$\mathrm c_4\ge\varepsilon_{\mathrm m}$, so both coefficients of $r^m_i$ are
nonpositive with sum $-1$. For $\varepsilon_{\mathrm m}=1$,
$\varsigma_-=4=\mathrm c_5-4\varepsilon_{\check\omega}$ gives $\mathrm c_3=-1$. For
$\varepsilon_{\check\omega}=\varepsilon_{\mathrm m}\le1/2$ put
$\mathrm c_4=\sqrt{(1-\varepsilon_{\mathrm m})^2+\varepsilon_{\mathrm m}^2}$, so that
$1-\varepsilon_{\mathrm m}\le\mathrm c_4\le1-\varepsilon_{\mathrm m}+\varepsilon_{\mathrm m}^2$
by squaring, $\varsigma_\pm=2+4\varepsilon_{\mathrm m}\pm2\mathrm c_4$ and
$\varsigma_+-\varsigma_-=4\mathrm c_4\ge4(1-\varepsilon_{\mathrm m})$, which gives the
stated brackets; the coefficient of $e^{-\varsigma_+t}$ in $r^m_i$ is
$(\varsigma_--8\varepsilon_{\mathrm m})/(4\mathrm c_4)$, of absolute value at most
$\varepsilon_{\mathrm m}(1+\varepsilon_{\mathrm m})/(2(1-\varepsilon_{\mathrm m}))\le\frac32\varepsilon_{\mathrm m}$,
and the coefficient of $e^{-\varsigma_-t}$ in $r^b_i$ is
$(\varsigma_--4\varepsilon_{\mathrm m})/(4\mathrm c_4)\in[0,\varepsilon_{\mathrm m}]$.
Since each pair of coefficients sums to $-1$, \eqref{eq:obs-cycle-slow-link-bounds}
follows, and the last claim is \eqref{eq:obs-cycle-slow-link-bounds} at the two
times, using $(1+\varepsilon_{\mathrm m})/4\le1$ and
$e^{-1/2}(1-\frac3{50})-\frac3{50}>\frac12$.

For (iii), write $a$ for the row of $r^m_i$, with entries $-1$ and $1$ at
$\bar m_i$ and $\bar m_{i+1}$ and $-1$ at $\check\omega_i$. Since the clock leaves
$\mathbf x$ unchanged, \eqref{eq:obs-cycle-slow-model-kinetics} gives, along every
assignment path,
$r^m_i(\mathbf x_t)-r^m_i(\mathbf x_0)=-\int_0^ta\cdot(\partial_{\bar k}U_{S_u},\varepsilon_{\mathrm m}\partial_{\bar m}U_{S_u},\varepsilon_{\check\omega}\partial_{\check\omega} U_{S_u})(\mathbf x_u)du+\sqrt{2\mathrm T(2\varepsilon_{\mathrm m}+\varepsilon_{\check\omega})}B_t$
with $B$ a standard Brownian motion. With $\mathrm n_j(S)\in\{0,1,2\}$ the number of
receivers attending $e_j$ under $S$,
$\partial_{\bar m_j}U_S=\mathrm n_{j-1}(S)r^m_{j-1}-\mathrm n_j(S)r^m_j+3\mathrm T\bar m_j-2\mathrm T\bar k_j$
and $\partial_{\check\omega_j}U_S=-\mathrm n_j(S)(r^b_j+r^m_j)+\mathrm T\check\omega_j$, so on the
stated event the integrand is bounded by
$\varepsilon_{\mathrm m}(8\mathrm R+8\mathrm T\mathrm R')+\varepsilon_{\check\omega}(4\mathrm R+\mathrm T\mathrm R')\le\bar\varepsilon(12\mathrm R+9\mathrm T\mathrm R')$,
using $3|\bar m_j|+2|\bar k_j|\le\sqrt{13}|\mathbf x|$. For $t$ as stated the drift
moves the residual by at most $1/4$, so a crossing below $1/2$ requires
$\sup_{u\le t}|B_u|\ge[4\sqrt{2\mathrm T(2\varepsilon_{\mathrm m}+\varepsilon_{\check\omega})}]^{-1}\ge[4\sqrt{6\mathrm T\bar\varepsilon}]^{-1}$.
Doob's maximal inequality applied to the martingale $\exp(\theta B_u-\theta^2u/2)$
with $\theta=\mathrm y/t$ gives $\Pr(\sup_{u\le t}B_u\ge\mathrm y)\le e^{-\mathrm y^2/(2t)}$,
and with $\mathrm y^2=1/(96\mathrm T\bar\varepsilon)$ the two-sided bound is
$2e^{-1/(192\mathrm T\bar\varepsilon t)}$; at the largest $t$ the exponent is
$-(12\mathrm R+9\mathrm T\mathrm R')/(48\mathrm T)\le-1/(4\mathrm T)$. The first
closing sentence collects (ii) and (iii), the fast-means case having rate
$4+4\varepsilon_{\check\omega}\ge4$. For \eqref{eq:obs-cycle-general-slow-rate} put
$\mathrm c_4=\sqrt{(1-\varepsilon_{\mathrm m})^2+\varepsilon_{\check\omega}^2}$; since
$2(1-\varepsilon_{\mathrm m})\ge1$, squaring gives
$1-\varepsilon_{\mathrm m}\le\mathrm c_4\le1-\varepsilon_{\mathrm m}+\varepsilon_{\check\omega}^2$,
and $\varsigma_-=2(1+\varepsilon_{\mathrm m}+\varepsilon_{\check\omega})-2\mathrm c_4$
gives the bracket, whose lower end is at least
$4\varepsilon_{\mathrm m}+\varepsilon_{\check\omega}$ because
$\varepsilon_{\check\omega}\le1/2$. At $(0,0)$ the residual is constant. Otherwise
the coefficient of $e^{-\varsigma_+t}$ in $r^m_i$ is
$\mathrm c_3=(\varsigma_--4\varepsilon_{\mathrm m}-4\varepsilon_{\check\omega})/(4\mathrm c_4)$,
and $0\le\varsigma_-\le4\varepsilon_{\mathrm m}+2\varepsilon_{\check\omega}$ with
$\mathrm c_4\ge1/2$ give
$|\mathrm c_3|\le2(\varepsilon_{\mathrm m}+\varepsilon_{\check\omega})$; the other
coefficient is $-1-\mathrm c_3$, so
$|r^m_i(t)|\le(1+2|\mathrm c_3|)e^{-\varsigma_-t}$, which is the second bound. At
$t=\mathrm c_6/\varepsilon_{\mathrm m}$ that bound is at most
$5e^{-\mathrm c_6\varepsilon_{\check\omega}/\varepsilon_{\mathrm m}}$, which is below
$1/2$ when $\varepsilon_{\check\omega}>\varepsilon_{\mathrm m}\log(10)/\mathrm c_6$.
$\square$

In the latent reading of \cref{ch:links}, the link translations are generators of the translation
instance with anchor zero and no loop term, and the supplied rate $\varepsilon_{\check\omega}$ places
their recognition means in the slow block on the held matching sheet; the natural-gradient recognition
flow of the generators assigns belief means, model means, and link means one common rate on an invariant
slice, so the separation used here comes from the supplied mobilities (\cref{rem:link-rates}).

\openproblemheading{Slow-manifold reduction of the collective free energy}{open:obs-slow-manifold}
Consider the rate-separated flow \eqref{eq:obs-rate-separated-flow} under three
hypotheses. First, $F\in C^2$ and the block metric
$\mathsf G_{\rm f}(f,s)\oplus\mathsf G_{\rm s}(s)$ has $C^1$ symmetric
positive-definite blocks with $\mathsf G_{\rm s}$ independent of $f$. Second, there
are a compact set $\mathcal K\subset\mathcal U_0$, a $C^1$ branch $f^\star$ of
conditional minimizers, and $\mathrm c_1>0$ with
$\partial^2_{ff}F(f^\star(s),s)\succeq\mathrm c_1\mathsf G_{\rm f}(f^\star(s),s)$ for
all $s\in\mathcal K$. Third, the reduced flow \eqref{eq:obs-reduced-slow-flow} from
$s_0$ stays in the interior of $\mathcal K$ on $[0,t_1]$, and the initial fast state
lies in the basin of $f^\star(s_0)$ for the frozen fast flow
$\dot f=-\mathsf G_{\rm f}^{-1}\partial_fF(f,s_0)$. Under the second hypothesis the
linearized frozen fast flow has spectrum in $(-\infty,-\mathrm c_1]$, so $f^\star(s)$
is isolated and uniformly asymptotically stable on $\mathcal K$, and Tikhonov's
theorem gives $s_\varepsilon(t)\to s(t)$ for $0\leq t\leq t_1$ and
$f_\varepsilon(t)\to f^\star(s(t))$ for $0<t\leq t_1$ as $\varepsilon\to0$
\citep{Tikhonov1952}; for smooth data Fenichel's theory gives a locally invariant
slow manifold $O(\varepsilon)$-close to a compact normally hyperbolic piece of
$\mathcal S^\star$ \citep{fenichel1979geometric}. Their application at $c_*$ is open
beyond the held matching sheets of \cref{prop:obs-slow-model-link-dichotomy}(i),
whose selected-edge blocks meet the three hypotheses with a curvature constant
uniform in the temperature, so that Tikhonov's theorem applies to them, while the
split by mobility on the full sheet has a constant of order $\mathrm T$:
\cref{prop:obs-attention-reduced-hessian} shows that the second hypothesis can fail
after the source rows are optimized even when every source energy and every
frozen-weight curvature is positive, the sampled dynamics with source jumps of
\cref{ch:finite-depth-formation} lie outside this deterministic statement, and no
explicit bound on the slow orbit or its Fisher duration in terms of $\varepsilon$ is
available. It is settled in the class of the sampled dynamics of
\cref{def:rg-formation-cycle-experiment} under the mobility of
\cref{prop:obs-slow-model-link-dichotomy} with $\varepsilon_{\check\omega}=\varepsilon_{\mathrm m}$
and no parameter learning, where the split by mobility, $f=\bar k$ and
$s=(\bar m,\check\omega)$, and \eqref{eq:obs-reduced-slow-flow} for $F=U_S$ under the
current assignment $S$ define the reduced slow flow of an assignment path: positively,
by a bound on the distance between the slow coordinates of the sampled process and
that reduced flow over a bounded slow-time interval, explicit in
$\varepsilon_{\mathrm m}$, $\mathrm T$, and the number of assignment changes, tending
to zero with $\varepsilon_{\mathrm m}$ and $\mathrm T$, and holding, for every
positive integer $\mathrm n$, with probability tending to one on the event that the
assignment path has at most $\mathrm n$ changes per unit slow time; negatively, by
a positive integer $\mathrm n$, a sequence of mobility ratios and temperatures
tending to zero, and initial states from which, with probability bounded below, the
assignment path has at most $\mathrm n$ changes per unit slow time and that
distance stays bounded below over a bounded slow-time interval. The
learning case, in which the generative parameter is the slow coordinate and the
reduced object is the recognition law at the frozen parameter, lies outside this
class.

\section{Parent-sensitive inference and the prepared parent}
\label{sec:rg-memory-parent-mechanism}

\definitionheading{Parent-sensitive sweep}{def:rg-memory-parent-sweep}
Let $\Xi$ be a nonempty finite set of fine states with a parent map
$\operatorname{par}:\Xi\to\mathsf Z$ that retains the complete decorated parent,
keeping only parent values with nonempty fibers. Take a prior $p^0$ on $\Xi$
with $p^0(\xi)>0$, normalized conditional laws $\psi_\xi$ of an unobserved record
$r$ in a finite set $\mathsf O'$, and record likelihoods $L_o(\xi)>0$ at the
observation, none of which reads recognition. The generative law of $(\xi,r,o)$
has density $p^0(\xi)\psi_\xi(r)L_o(\xi)$, and its evidence and posterior are
\begin{equation}
\mathcal Z_o=\sum_\xi p^0(\xi)L_o(\xi),
\qquad
p(\xi)=\frac{p^0(\xi)L_o(\xi)}{\mathcal Z_o}.
\label{eq:rg-memory-target}
\end{equation}
On an undirected graph of allowed moves take symmetric proposal weights
$\mathsf a_{\xi\xi'}=\mathsf a_{\xi'\xi}\geq0$ with off-diagonal row sums at most
$1/2$. The Metropolis kernel, the parent projection, the posterior lift, the
refresh, and the sweep are
\begin{equation}
\begin{aligned}
\mathsf M_{\xi\xi'}&=\mathsf a_{\xi\xi'}\min\bigl(1,p(\xi')/p(\xi)\bigr)\quad(\xi\ne\xi'),
&\mathsf M_{\xi\xi}&=1-\sum_{\xi'\ne\xi}\mathsf M_{\xi\xi'},\\
\mathsf C_{\xi z}&=\mathbf 1_{\operatorname{par}(\xi)=z},
&\mathsf J_{z\xi}&=\frac{p(\xi)\mathbf 1_{\operatorname{par}(\xi)=z}}{p_Z(z)},
\qquad p_Z=p\mathsf C,\\
\mathsf D&=\mathsf C\mathsf J,
&\mathsf T&=\mathsf M\mathsf D.
\end{aligned}
\label{eq:rg-memory-sweep}
\end{equation}
The sweep depends on the parent through the conditional laws $p(\cdot\given z)$. The
marked version of a kernel $\mathsf K$ on
$\Xi$ is $\widetilde{\mathsf K}((\xi,r),(\xi',r'))=\mathsf K_{\xi\xi'}\psi_{\xi'}(r')$.
The proposal and the update order are declared choices.

\hypothesisheading{Instantaneous within-parent refresh}{hyp:rg-memory-instantaneous-refresh}
In this section the refresh is the exact kernel $\mathsf D=\mathsf C\mathsf J$:
each application resamples the fine state from its exact conditional law given
the current complete parent. This is the adiabatic elimination of the
within-parent fine state, whose timescale is infinitely shorter than the parent's.
It excludes finite
within-fiber mixing, for example the partial refresh
$\mathsf D_\omega=(1-\omega)I+\omega\mathsf C\mathsf J$ with $0\leq\omega<1$, which
is equally normalized, $p$-reversible, and parent preserving.

\theoremheading{Stationarity, exact descent, and finite-component convergence}{thm:rg-memory-inference-descent}
\begin{enumerate}
\item[(i)] For a recognition law $q(\xi)\psi'_\xi(r)$ the free energy is
\begin{equation}
\Fenergy(q\otimes\psi';o)=\KL(q\Vert p)+\E_q\KL(\psi'_\xi\Vert\psi_\xi)
-\log\mathcal Z_o,
\label{eq:rg-memory-fixed-vfe}
\end{equation}
so matching the record conditionals, $\psi'=\psi$, is a specialization of the
joint objective and not the removal of its dependence terms.
\item[(ii)] The kernels $\mathsf M$, $\mathsf D$, $\mathsf T$ are normalized and
preserve $p$, and $\mathsf M$ and $\mathsf D$ are $p$-reversible; their marked
versions are normalized and preserve $p(\xi)\psi_\xi(r)$.
\item[(iii)] $q\mathsf D$ minimizes $\KL(\cdot\Vert p)$ among laws with parent
marginal $q\mathsf C$, with $\KL(q\mathsf D\Vert p)=\KL(q\mathsf C\Vert p_Z)$.
\item[(iv)] Call a component a connected component of the union of the positive
off-diagonal edges of $\mathsf M$ and $\mathsf D$. For every initial $q$, $\mathsf T$
preserves component masses, and for every $p$-stationary kernel $\mathsf K$ on
$\Xi$, with $q'=q\mathsf K$ and reversal
$\mathsf K^*_{\xi'\xi}=p(\xi)\mathsf K_{\xi\xi'}/p(\xi')$,
\begin{equation}
\KL(q\Vert p)-\KL(q\mathsf K\Vert p)
=\sum_{\xi':q'(\xi')>0}q'(\xi')
\KL\bigl(q(\cdot\given\xi')\Vert\mathsf K^*_{\xi'\cdot}\bigr),
\label{eq:rg-memory-reverse-conditional-loss}
\end{equation}
where $q(\xi\given\xi')=q(\xi)\mathsf K_{\xi\xi'}/q'(\xi')$. Hence the marked
sweep maps $q\otimes\psi'$ to $(q\mathsf T)\otimes\psi$ and lowers
\eqref{eq:rg-memory-fixed-vfe} by
$\KL(q\Vert p)-\KL(q\mathsf T\Vert p)+\E_q\KL(\psi'_\xi\Vert\psi_\xi)\geq0$,
which vanishes exactly when $q(\xi)=q(\Xi')p(\xi\given\Xi')$ on each component
$\Xi'$ and $\psi'_\xi=\psi_\xi$ for every $\xi$ with $q(\xi)>0$. For matched
record conditionals, $\psi'=\psi$, this is the statement that $q\mapsto q\mathsf T$
decreases $\KL(q\Vert p)$, strictly unless $q$ has that form.
\item[(v)] On each component, $q\mathsf T^n$ converges geometrically to that
conditional target.
\end{enumerate}
\status{ESTABLISHED}

\noindent\emph{Proof.}
For (i), the posterior of $(\xi,r)$ given $o$ is $p(\xi)\psi_\xi(r)$ because $r$
is unobserved and conditionally independent of $o$, and the evidence is
$\mathcal Z_o$; \cref{def:elbo-extended} and the chain rule
\cref{lem:prob-kl-chain-rule} give the display. For (ii),
$p(\xi)\mathsf M_{\xi\xi'}=\mathsf a_{\xi\xi'}\min(p(\xi),p(\xi'))=p(\xi')\mathsf M_{\xi'\xi}$;
$\mathsf J\mathsf C=I$ gives $p\mathsf D=p$, and
$p(\xi)\mathsf D_{\xi\xi'}=p(\xi)p(\xi')/p_Z(z)$ on a common parent fiber $z$.
Stationarity survives composition; summing the terminal record of a marked kernel
proves normalization and summing the initial record proves invariance. Part (iii) is
\cref{prop:elbo-exact-e-coordinate} for the block of fine coordinates given the
parent: $q\mathsf D$ has parent marginal $q\mathsf C\mathsf J\mathsf C=q\mathsf C$ and
conditionals $p(\cdot\given z)$, and the chain rule over the parent gives the value.
For (iv), the input--output laws $q(\xi)\mathsf K_{\xi\xi'}$ and
$p(\xi)\mathsf K_{\xi\xi'}$ have divergence $\KL(q\Vert p)$, and disintegrating them
over the output with \cref{lem:prob-kl-chain-rule} gives
\eqref{eq:rg-memory-reverse-conditional-loss}. With positive holding, equality forces
$q/p$ to be constant across every positive edge; equality for $\mathsf M$ makes $q$
stationary under $\mathsf M$, and equality for the subsequent $\mathsf D$ makes the
ratio constant on parent fibers, which is constancy on components. The marked sweep
sends $q(\xi)\psi'_\xi(r)$ to
$\sum_{\xi,r}q(\xi)\psi'_\xi(r)\mathsf T_{\xi\xi'}\psi_{\xi'}(r')=(q\mathsf T)(\xi')\psi_{\xi'}(r')$,
so by (i) the free energy before minus after is the displayed sum, whose two
terms are nonnegative, the first by \eqref{eq:rg-memory-reverse-conditional-loss}
and the second as an average of relative entropies, and which vanishes exactly
when both terms do. For (v), positive
holding makes every $\mathsf M$ or $\mathsf D$ edge a positive $\mathsf T$ edge, so on
a component $\Xi'$ of $n$ states $\mathsf T^k>0$ for $k=\max(1,n-1)$. With target
$p_{\Xi'}=p(\cdot\given\Xi')$ and
$\mathrm c_1=\min_{\xi\xi'}(\mathsf T^k)_{\xi\xi'}/p_{\Xi'}(\xi')\in(0,1]$, one has
$\mathsf T^k=\mathrm c_1\mathbf 1p_{\Xi'}+(1-\mathrm c_1)\mathsf K'$ on $\Xi'$ with a
$p_{\Xi'}$-stationary stochastic $\mathsf K'$ (or $\mathsf T^k=\mathbf 1p_{\Xi'}$ when
$\mathrm c_1=1$), and convexity with data processing give
$\KL(q\mathsf T^k\Vert p_{\Xi'})\leq(1-\mathrm c_1)\KL(q\Vert p_{\Xi'})$ for laws $q$ on
$\Xi'$. $\square$

\theoremheading{The prepared parent is Markov}{thm:rg-memory-prepared-parent}
Under \cref{hyp:rg-memory-instantaneous-refresh}, the reduced sweep
$\overline{\mathsf T}:=\mathsf J\mathsf M\mathsf C$ is reversible under $p_Z$ and
\begin{equation}
\mathsf J\mathsf T=\overline{\mathsf T}\mathsf J,
\qquad
(\nu\mathsf J)\mathsf T^n=\nu\overline{\mathsf T}^n\mathsf J.
\label{eq:rg-memory-parent-semiconjugacy}
\end{equation}
For a start $q_0=\nu\mathsf J$ the parent history $Z_t=\operatorname{par}(\xi_t)$
of the chain with kernel $\mathsf T$ is Markov with transition $\overline{\mathsf T}$,
and for an arbitrary start it is so after the first completed refresh.
\status{ESTABLISHED}

\noindent\emph{Proof.}
With $\mathsf T=\mathsf M\mathsf C\mathsf J$ and $\mathsf J\mathsf C=I$,
$\mathsf J\mathsf T=\mathsf J\mathsf M\mathsf C\mathsf J=\overline{\mathsf T}\mathsf J$,
and induction gives the second identity. Moreover $p_Z(z)\overline{\mathsf T}_{zz'}$
is the sum of $p(\xi)\mathsf M_{\xi\xi'}$ over $\operatorname{par}(\xi)=z$ and
$\operatorname{par}(\xi')=z'$, which is symmetric by reversibility of $\mathsf M$. The
final refresh makes the fine law given the whole parent history $\mathsf J_{z\cdot}$,
with $z$ the last parent, so the next parent row is
$\mathsf J_{z\cdot}\mathsf M\mathsf C=\overline{\mathsf T}_{z\cdot}$, and induction over
histories gives the Markov property. $\square$

The refresh holds the fast fine state at its conditional optimum given the slow parent
(\cref{thm:rg-memory-inference-descent}(iii)), so $\mathsf J\mathsf T=\overline{\mathsf T}\mathsf J$
is the discrete instantaneous-equilibration case of \eqref{eq:obs-reduced-semiconjugacy}.

\remarkheading{Finite refresh leaves parent memory}{rem:rg-memory-partial-refresh}
\Cref{hyp:rg-memory-instantaneous-refresh} cannot be dropped. Take
$\Xi=\{1,2,3\}$ with $\operatorname{par}(1)=\operatorname{par}(2)=z_1$ and
$\operatorname{par}(3)=z_2$, uniform $p$, the single proposal
$\mathsf a_{23}=1/2$, the partial sweep
$\mathsf T_\omega=\mathsf M\mathsf D_\omega$, and a start $\nu\mathsf J$ with
$\nu(z_1),\nu(z_2)>0$. Then
\begin{equation}
\mathbb P(Z_2=z_2\given Z_1=z_1,Z_0=z_1)=1/6+\omega/12,
\qquad
\mathbb P(Z_2=z_2\given Z_1=z_1,Z_0=z_2)=1/2-\omega/4.
\label{eq:rg-memory-partial-refresh-witness}
\end{equation}
The dependence on $Z_0$ is $(1-\omega)/3$, equal to $7/30$ at $\omega=3/10$, and it
vanishes only at $\omega=1$, where both values equal
$\overline{\mathsf T}_{z_1z_2}=1/4$. For $\omega<1$ the parent history is not
Markov. It remains a deterministic function of the fine Markov chain with kernel
$\mathsf T_\omega$, so the hidden-history law of \cref{sec:rg-memory-hidden-history}
retains it exactly, and the memory cost of a first-order parent closure is the
defect $D_N^{(1)}$ of \cref{thm:rg-memory-projection}, whose value
$I(Z_2;Z_0\given Z_1)$ at $N=3$ is positive. \status{ESTABLISHED}

\noindent\emph{Proof.}
Only state $2$ moves under $\mathsf M$, to $3$ with probability $1/2$, and
$\mathsf D_\omega$ keeps a state of $\{1,2\}$ with probability $1-\omega/2$. Given
$Z_1=z_1$, the state before the second sweep is $2$ with probability
$1-\omega/2$ if $Z_0=z_2$, since it moved from $3$, and with probability
$1/3+\omega/6$ if $Z_0=z_1$, since after $\mathsf M$ it is $2$ with probability
$1/3$ and the refresh redraws it uniformly with probability $\omega$; multiplying
by $1/2$ gives
\eqref{eq:rg-memory-partial-refresh-witness}. Both conditioning events have positive
probability, so $I(Z_2;Z_0\given Z_1)>0$. $\square$

\section*{Physicist's summary}

Fast and slow are declared, not derived. The collective free energy is then a
potential for every positive-definite mobility. The reduced gradient flow is
conjugate to its tangent lift on the graph of the conditional optimum; the
finite-rate flow shares its slow component there but can leave that graph.
Instantaneous within-parent refresh gives the Markov parent proved here.
On the four-cycle, slow model means lag the beliefs only if the links
are slow too; a fast link carries the model agreement on the belief clock.

\chapter[Lie-Algebra Links as Latent Fields]{Lie-Algebra Links as Latent Fields}
\label{ch:links}

This chapter asks what becomes of the independent links of \Cref{def:geo-graph-links} when they are
inferred jointly with the agents' states instead of being held as structural data or moved as
parameters. Written in the exponential chart of a reference point, every twist is the exponential of a
generator in the Lie algebra, on which passive re-expressions do not act
(\Cref{def:link-exponential-chart,prop:link-generator-gauge}). A Gaussian base law on the generators
multiplied by the exponential of minus a nonnegative loop action is a normalized prior for every
structure group whose Lie algebra carries an inner product, compact or not
(\Cref{prop:link-generator-prior}); on the general linear group the loop action measures the reference
holonomy with a fiber metric, because no continuous class function separates unipotent holonomy from
the identity (\Cref{cor:link-no-invariant-defect,prop:link-loop-defect}). With the generators latent,
the collective bound of \Cref{prop:elbo-random-link-vfe} specializes to a joint that has unit mass on
every record hypergraph and to a free energy whose link sector is the relative entropy from the
generator prior (\Cref{prop:link-latent-normalization,prop:link-gaussian-recognition-vfe}). Its
stationarity condition is a link field equation in which the stress of the prior balances the expected
matter current transported by the adjoint derivative of the chart (\Cref{prop:link-field-equation});
for a comparison record the current has the outer-product form of a delta rule
(\Cref{cor:link-delta-rule}), and without the prior the link problem degenerates
(\Cref{cor:link-prior-necessary}). Held links are the zero-width limit of the latent model, and
statements proved at every held link transfer conditionally on the generators
(\Cref{thm:link-held-limit,prop:link-conditional-transfer}). Matter attached by record extension
induces no link action, while matter normalized jointly with the generators induces minus the logarithm
of its normalizer (\Cref{prop:link-induced-action}). Recognition of the generators is a first-order
natural-gradient history whose explicit Euler step is the Gaussian Bayesian learning rule
(\Cref{prop:link-recognition-flow}); across learning episodes a Gibbs prior on generator histories gives
the most probable history an inertia whose mass is the precision of the generator increment
(\Cref{prop:link-history-prior,prop:link-map-inertia}); and near zero generators the stationarity
condition is a linear equation whose positive operator collects the anchor, the loops, the episodes,
and the matter (\Cref{prop:link-linearized-equation}). These results are used for the coarse links and
the holonomy-conditioned parents of \Cref{ch:network-state,ch:full-graph-meta-agent-vfe}, whose link
statements gain a version with latent generators, and the link translations of
\Cref{ch:finite-depth-formation} are their translation-group instance
(\Cref{prop:link-translation-instance}); Supplement G realizes them for Gaussian matter. Longer proofs
are deferred to \cref{app:deferred-proofs}.

The symbols of \cref{tab:link-notation} are declared for this chapter and its proof section, and those
that recur in \Cref{ch:network-state,ch:full-graph-meta-agent-vfe} are also listed in \cref{tab:notation};
$\mathbf B$, $\mathbf S$, $\mathbf T$, and $\mathbf H$ denote generic square matrices, $\mathbf b$ and
$\mathbf c$ generic vectors, and $\mathbf h$ a generic real function, each bound in the statement or proof
where it appears, and a ring accent marks an object of the generator sector written in the coordinates of
the reference point.

\begin{sciencetable}{Symbols of the latent-link construction.}
\label{tab:link-notation}
\begin{tabularx}{\textwidth}{>{\raggedright\arraybackslash}p{0.33\textwidth}>{\raggedright\arraybackslash}X}
\toprule
\textbf{Symbol} & \textbf{Meaning} \\
\midrule
\(\mathfrak g\), \(\exp\), \(\operatorname{Ad}_g\), \(\mathrm d\rho_x\) & Lie algebra of \(G\), its exponential map, the adjoint action, the differential of \(\rho_x\) at \(e_G\) \\
\(E_+\), \(\mathscr T=\mathfrak g^{E_+}\) & one orientation of each reversal pair; link space of one channel, \(\mathscr T\times\mathscr T\) for two independent channels \\
\(\mathring\omega^x_e\), \(\exp_\sigma\), \(\mathring\Theta^x_p\) & generator of the twist of \(e\); exponential chart at \(\sigma\); reference path product, the reference holonomy of a closed walk \\
\(\mathsf G^{\mathrm{fib}}\), \(\langle\cdot,\cdot\rangle_{\mathsf G}\), \(\lVert\cdot\rVert_{\mathsf G}\) & fiber metric at the reference point; trace form on \(\mathfrak{gl}(\mathbb V)\) and the inner product on \(\mathfrak g\) \\
\(\mathfrak Q^x_p\), \(\mathbb L\), \(\mathrm c_p\), \(\mathbf C\) & loop defect of a closed walk; family of closed walks and their coefficients; cycle map \\
\(\mathbb L_+\), \(\mathbf C_+\), \(z_p\), \(r_{\mathbb L}\), \(\kappa\) & loops with positive coefficient and their cycle map; net traversal vector of a loop; rank of these vectors; number of components of the multigraph \\
\(d\), \(Z\), \(\mathscr T_{\mathrm{ex}}\), \(\mathscr T_{\mathrm{har}}\) & vertex-difference map and cycle space of \(\R^{E_+}\); exact and harmonic generators \\
\(\nu^{\mathscr T}\), \(\mathscr S\), \(Z_{\mathscr S}\), \(\mathbb P^{\mathscr T}_\theta\) & base law, link action, normalizer, generator prior \\
\(\mathring\alpha^x_e\), \(\mathring\Sigma^x_e\); \(\mathring\mu\), \(\mathring\Lambda\) & anchor and prior covariance; mean and precision of a Gaussian recognition law \(R\) of the generators \\
\(\mathbf L(\mathbf B,\mathbf T)\) & Fr\'echet derivative of the matrix exponential \\
\(\mathbf H_o\), \(\mathbf U\), \(\mathbf V\), \(\mathbf J^x_e\) & matter energy, expected action plus matter energy, most-probable-link objective, link current \\
\(\mathbf M\), \(\mathbf M^\star\), \(\mathbf l\) & transport matrix of a comparison record, its rank-one minimizer, expected record energy \\
\(\mathbf F_Q\), \(\mathbf F\) & held-link free energy as a function of the generators, and its infimum over a class \\
\(\mathscr S^{\mathrm{ind}}\), \(\mathfrak Z\) & induced link action, joint normalizer of generators and undirected matter \\
\(\mathfrak i_\lambda\), \(\mathring\Xi\), \(\mathring\vartheta\), \(\mathring\varpi\), \(\mathring\Gamma\) & generator integrand, space of admissible precisions, natural and expectation parameters, log-partition function \\
\(\mathbf m\), \(\mathfrak U\), \(\mathbb P^{\mathrm{hist}}\), \(Z_{\mathrm{hist}}\), \(\mathring Z^{\leftarrow}_n\) & increment precision, kinetic term, generator history prior, its normalizer, backward functions \\
\(\mathscr S^{\mathrm{hist}}\), \(\mathfrak V_n\) & history action and episode potentials \\
\(\mathbf K\), \(\mathbf N_0\), \(\mathring\Phi\), \(\mathbf j\) & linearized link operator, current correction of the Hessian, matter operator, source \\
\bottomrule
\end{tabularx}
\end{sciencetable}

\paragraph{Scope.}
The chapter does not claim that latent links improve the evidence, the prediction of any record, or any
learning outcome relative to held links; that the generator prior is invariant under reference changes
that do not act by multiples of isometries of the fiber metric, or that the loop defect is independent
of the base vertex of a closed walk; that every local minimizer of the most-probable-link objective is
global; that the recognition flow of the generators converges, or that stationary generator histories
are unique outside the ball of \cref{prop:link-linearized-equation}; that the joint M-coordinate of the
anchors and the prior covariances exists when the action has loop terms; that frames are dynamical, so
that time-like links would carry variables; that the linearized link operator is the operator of a
physical field; or that the outer-product form of the link current identifies the model with a
fast-weight or delta-rule architecture. \status{NOT-CLAIMED} The interaction multigraph stays structural, and edge creation and
deletion is \cref{open:obs-edge-creation}; a finite-dimensional family of link laws that survives the
coarse step is \cref{open:link-coarse-closure}; and a slow-manifold reduction with the generator
recognition means in the slow block is part of \cref{open:obs-slow-manifold}.

\section{The exponential chart of the links}
\label{sec:link-chart}

This section writes the independent links of \Cref{def:geo-graph-links} in the exponential chart of
the reference point $\sigma$ of \Cref{prop:geo-edge-twist-invariance}. The twist of an edge copy is
unchanged by every passive re-expression, so a generator of the twist carries no vertex gauge: the
frames $U_i$ carry the passive action, while a change of reference point or an active automorphism
moves all generators by one adjoint map (\Cref{prop:link-generator-gauge}). The chart covers exactly
the link families whose twists are exponentials. For $G=\GL(K,\R)$ this excludes some twists of
positive determinant, while on a closed walk that traverses two reversal pairs once each the
reference holonomies of chart links still reach every element of $\GL^+(K)$
(\Cref{prop:link-chart-image}).

\lemmaheading{Exponential maps and homomorphisms}{lem:link-exp-naturality}
Let $\boldsymbol\varphi:G\to\mathbb Y$ be a smooth homomorphism of Lie groups with differential
$\mathrm d\boldsymbol\varphi:\mathfrak g\to\operatorname{Lie}(\mathbb Y)$ at $e_G$. Then
$\boldsymbol\varphi(\exp\mathring\omega)=\exp(\mathrm d\boldsymbol\varphi\mathring\omega)$ for every
$\mathring\omega\in\mathfrak g$. In particular, let $\operatorname{Ad}_g$ be the differential at $e_G$
of the conjugation $g'\mapsto gg'g^{-1}$, and let $\rho_x:G\to\GL(\mathbb V)$ be a smooth homomorphism
into the general linear group of a finite-dimensional real vector space $\mathbb V$, with differential
$\mathrm d\rho_x$ at $e_G$. Then for all $g\in G$ and $\mathring\omega\in\mathfrak g$, with the matrix
exponential as the exponential map of $\GL(\mathbb V)$,
\begin{equation}
g\exp(\mathring\omega)g^{-1}=\exp(\operatorname{Ad}_g\mathring\omega),
\qquad
\rho_x(\exp\mathring\omega)=\exp(\mathrm d\rho_x\mathring\omega),
\qquad
\mathrm d\rho_x(\operatorname{Ad}_g\mathring\omega)=\rho_x(g)\mathrm d\rho_x(\mathring\omega)\rho_x(g)^{-1}.
\label{eq:link-exp-naturality}
\end{equation}
\status{ESTABLISHED}

\noindent\emph{Proof.}
The exponential map sends $\mathring\omega$ to the value at time one of the integral curve through
the identity of the left-invariant vector field whose value at the identity is $\mathring\omega$, and
$t\mapsto\exp(t\mathring\omega)$ is a smooth homomorphism from $(\R,+)$. Hence
$\boldsymbol\gamma(t)=\boldsymbol\varphi(\exp(t\mathring\omega))$ is a smooth homomorphism from $(\R,+)$ into $\mathbb Y$.
Differentiating $\boldsymbol\gamma(t+t')=\boldsymbol\gamma(t)\boldsymbol\gamma(t')$ in $t'$ at $t'=0$ shows that $\boldsymbol\gamma'(t)$ is the image
of $\boldsymbol\gamma'(0)$ under the differential at the identity of the left multiplication by $\boldsymbol\gamma(t)$, so
$\boldsymbol\gamma$ is the integral curve through the identity of the left-invariant vector field with value
$\boldsymbol\gamma'(0)=\mathrm d\boldsymbol\varphi\mathring\omega$, and uniqueness of integral curves gives
$\boldsymbol\gamma(t)=\exp(t\mathrm d\boldsymbol\varphi\mathring\omega)$; the first claim is the case $t=1$. Conjugation by
$g$ is a smooth homomorphism with differential $\operatorname{Ad}_g$, which gives the first identity of
\cref{eq:link-exp-naturality}. On $\GL(\mathbb V)$ the left-invariant vector field with value
$\mathbf H$ at the identity is $y\mapsto y\mathbf H$, and the matrix exponential satisfies
$\frac{\mathrm d}{\mathrm dt}e^{t\mathbf H}=e^{t\mathbf H}\mathbf H$ with $e^{0}=I_{\mathbb V}$, so it is the
exponential map of $\GL(\mathbb V)$; this gives the second identity. Applying $\rho_x$ to the first
identity for $t\mathring\omega$ and using the second gives
$\rho_x(g)\exp(t\mathrm d\rho_x\mathring\omega)\rho_x(g)^{-1}=\exp(t\mathrm d\rho_x(\operatorname{Ad}_g\mathring\omega))$,
whose derivative at $t=0$ is the third identity. \hfill$\square$

\definitionheading{Exponential chart of the links}{def:link-exponential-chart}
Let $\mathfrak g$ be the Lie algebra of $G$ and $\exp:\mathfrak g\to G$ its exponential map, let
$E_+\subseteq E$ contain exactly one edge copy of each reversal pair of $\mathfrak I$, and let the
reference point $\sigma$ and the frames $u_i=\sigma\cdot U_i^{-1}$ be those of
\Cref{prop:geo-edge-twist-invariance}. For a channel $x$, a generator family is a point
$\mathring\omega^x=(\mathring\omega^x_e)_{e\in E_+}$ of the link space $\mathscr T=\mathfrak g^{E_+}$,
extended to the reversed copies by the first line below. Its twists and links are those of the second
line, and for a walk $p=(e_0,\ldots,e_{r-1})$ with $e_a:i_{a+1}\to i_a$, as in
\Cref{def:geo-graph-links}, its reference path product is that of the third line, called the reference
holonomy of $p$ when $p$ is closed:
\begin{gather}
\mathring\omega^x_{\bar e}=-\mathring\omega^x_e,
\label{eq:link-generator-reversal}\\
\mathring\Theta^x_e=\exp(\mathring\omega^x_e),
\qquad
\Theta^x_e=U_i\exp(\mathring\omega^x_e)U_j^{-1}\quad(e:j\to i),
\label{eq:link-chart}\\
\mathring\Theta^x_p=\mathring\Theta^x_{e_0}\mathring\Theta^x_{e_1}\cdots\mathring\Theta^x_{e_{r-1}}
=\exp(\mathring\omega^x_{e_0})\exp(\mathring\omega^x_{e_1})\cdots\exp(\mathring\omega^x_{e_{r-1}}).
\label{eq:link-reference-holonomy}
\end{gather}
The links satisfy \cref{eq:geo-regime-two-links} because $\exp(-\mathring\omega)=\exp(\mathring\omega)^{-1}$,
and $\mathring\Theta^x_e$ is their twist in the sense of \cref{eq:geo-edge-twist}. The exponential chart
at $\sigma$ is the continuous map $\exp_\sigma:\mathscr T\to G^{E_+}$,
$\mathring\omega^x\mapsto(U_i\exp(\mathring\omega^x_e)U_j^{-1})_{e\in E_+}$. The link space is a real
vector space of dimension $\lvert E_+\rvert\dim G$, standard Borel with its Borel sets, and the chart
covers exactly the link families whose twists lie in $\exp(\mathfrak g)$ (\Cref{prop:link-chart-image}).
When $\rho_x$ is a smooth homomorphism $G\to\GL(\mathbb V)$ on a real vector sample fiber
$\mathbb V=\R^{\dim\mathbb V}$, the represented generators are
$\mathrm d\rho_x(\mathring\omega^x_e)\in\mathfrak{gl}(\mathbb V)$, and
$\rho_x(\mathring\Theta^x_e)=\exp(\mathrm d\rho_x\mathring\omega^x_e)$ by \Cref{lem:link-exp-naturality}.
The two channels carry independent generator families, with link space $\mathscr T\times\mathscr T$;
under the common-link specialization \cref{eq:geo-common-link} one generator family
$\mathring\omega=\mathring\omega^b=\mathring\omega^m$ serves both channels and the index $x$ is dropped.

\propositionheading{Gauge action on generators}{prop:link-generator-gauge}
In \Cref{def:link-exponential-chart}, let the generator family $\mathring\omega^x$ have the links
$\Theta^x=\exp_\sigma(\mathring\omega^x)$. (a) After a passive re-expression $(g_i)_{i\in V}$ the links
$g_i\Theta^x_eg_j^{-1}$ are the links of the same $\mathring\omega^x$ relative to $\sigma$ and the
re-expressed frames $g_iU_i$, so every generator is unchanged. (b) Relative to the reference point
$\sigma'=\sigma\cdot g_0$ the frames are $U_ig_0$, and the links $\Theta^x$ are those of the generator
family $\operatorname{Ad}_{g_0^{-1}}\mathring\omega^x=(\operatorname{Ad}_{g_0^{-1}}\mathring\omega^x_e)_{e\in E_+}$,
so that $\exp_{\sigma'}\circ\operatorname{Ad}_{g_0^{-1}}=\exp_\sigma$; every reference path product
becomes $g_0^{-1}\mathring\Theta^x_pg_0$. (c) An active $F$ with $F(\sigma)=\sigma\cdot\kappa$, as in
\Cref{prop:geo-edge-twist-invariance}, carries $\Theta^x$ to the links of
$\operatorname{Ad}_\kappa\mathring\omega^x$ relative to $\sigma$ and the same frames, and every
reference path product to $\kappa\mathring\Theta^x_p\kappa^{-1}$. (d) The holonomy of a closed walk
$p$ based at $i_0$ and, for a smooth homomorphism $\rho_x:G\to\GL(\mathbb V)$, its represented
reference holonomy are
\begin{equation}
\mathrm{Hol}^x(p)=U_{i_0}\mathring\Theta^x_pU_{i_0}^{-1},
\qquad
\rho_x(\mathring\Theta^x_p)=\exp(\mathrm d\rho_x\mathring\omega^x_{e_0})\exp(\mathrm d\rho_x\mathring\omega^x_{e_1})\cdots\exp(\mathrm d\rho_x\mathring\omega^x_{e_{r-1}}).
\label{eq:link-holonomy-chart}
\end{equation}
Each of (a) to (c) moves all generators by one linear map of $\mathfrak g$, namely the identity,
$\operatorname{Ad}_{g_0^{-1}}$, or $\operatorname{Ad}_\kappa$, and the passive vertex group acts on the
frames alone. \status{ESTABLISHED}

\noindent\emph{Proof.}
(a) The vertex law \cref{eq:geo-regime-two-gauge-law} and $U_i'=g_iU_i$ from
\cref{eq:geo-regime-one-gauge-law} give
\begin{equation*}
g_i\Theta^x_eg_j^{-1}=(g_iU_i)\exp(\mathring\omega^x_e)(g_jU_j)^{-1}.
\end{equation*}
(b) From
$\sigma=\sigma'\cdot g_0^{-1}$ we get $u_i=\sigma'\cdot(U_ig_0)^{-1}$, and
\Cref{lem:link-exp-naturality} gives
$(U_ig_0)\exp(\operatorname{Ad}_{g_0^{-1}}\mathring\omega^x_e)(U_jg_0)^{-1}=U_ig_0g_0^{-1}\exp(\mathring\omega^x_e)g_0g_0^{-1}U_j^{-1}=\Theta^x_e$;
the same identity applied to each factor of \cref{eq:link-reference-holonomy} gives the path products.
(c) By \Cref{prop:geo-edge-twist-invariance} the image twists are
$\kappa\exp(\mathring\omega^x_e)\kappa^{-1}=\exp(\operatorname{Ad}_\kappa\mathring\omega^x_e)$, and a twist
determines its link given the frames through $\Theta^x_e=U_i\mathring\Theta^x_eU_j^{-1}$. (d) The first
identity is \cref{eq:geo-edge-twist-holonomy}, and the second follows from the homomorphism property of
$\rho_x$ and \Cref{lem:link-exp-naturality}. \hfill$\square$

\propositionheading{Real logarithms}{prop:link-real-logarithm}
A real square matrix is the exponential of a real square matrix exactly when it is nonsingular and
each Jordan block belonging to a negative eigenvalue occurs an even number of times
\citep[Theorem~1]{culver1966existence}. \status{STANDARD}

\propositionheading{Image of the exponential chart}{prop:link-chart-image}
In \Cref{def:link-exponential-chart} the following hold. (a) The image of $\exp_\sigma$ is the set of
link families of channel $x$ whose twists all lie in $\exp(\mathfrak g)$; it does not depend on
$\sigma$, and passive re-expressions and active automorphisms map it onto itself. Link families with a
twist outside $\exp(\mathfrak g)$ have no generator family. (b) For $G=\GL(K,\R)$ or $G=\GL^+(K)$ the
set $\exp(\mathfrak{gl}(K,\R))$ is the set of matrices described in \Cref{prop:link-real-logarithm}
and lies in $\GL^+(K)$; for $K\geq2$ it is a proper subset of $\GL^+(K)$, since
$\operatorname{diag}(-1,-2,1,\ldots,1)$ is not the exponential of a real matrix, so the chart misses
every link family with such a twist, and for $G=\GL(K,\R)$ it also misses every link family with a
twist of negative determinant. (c) Every element of $\GL^+(K)$ equals $\exp(\mathbf S)\exp(\mathbf T)$
with $\mathbf S$ symmetric and $\mathbf T$ skew-symmetric, and every element of $SO(K)$ equals
$\exp(\mathbf T)$ with $\mathbf T$ skew-symmetric; hence $\GL^+(K)$ is connected and generated by
$\exp(\mathfrak{gl}(K,\R))$, and for $G=SO(K)$ the chart covers every link family. (d) For every
connected Lie group $G$, $\exp(\mathfrak g)$ generates $G$. (e) Let $G$ be $\GL(K,\R)$ or $\GL^+(K)$,
and let the closed walk $p$ traverse two distinct reversal pairs exactly once each. Then every element
of $\GL^+(K)$ is the reference holonomy $\mathring\Theta_p$ of some generator family, so for $K\geq2$
the reference holonomies of chart links leave $\exp(\mathfrak{gl}(K,\R))$. On a triangle walk
$(e_0,e_1,e_2)$ with every $e_a\in E_+$ and $K=2$, the generators $\operatorname{diag}(0,\log2)$,
$\pi(E_{21}-E_{12})$, and $0$ give $\mathring\Theta_p=\operatorname{diag}(-1,-2)$. \status{ESTABLISHED}

The proof is in \cref{sec:proofs-links}.

\section{Curvature on a noncompact structure group}
\label{sec:link-curvature}

A loop penalty should vanish exactly on the closed walks with trivial represented holonomy. For compact
$G$ and unitary $\rho_x$ the class function of \Cref{rem:geo-loop-penalties} does so; on the general
linear group no continuous class function can (\Cref{cor:link-no-invariant-defect}), because
conjugation contracts every unipotent element to the identity
(\Cref{prop:link-class-function-blindness}). The loop defect defined below measures the reference
holonomy with a fiber metric placed at the reference point. It separates every holonomy from the
identity; with the metric held as a matrix at the reference point, it is invariant under the reference
changes that act by multiples of isometries of that metric, and for the general linear group under no
others (\Cref{prop:link-loop-defect}).

\propositionheading{Continuous class functions are blind to unipotent holonomy}{prop:link-class-function-blindness}
Let $K\geq2$, let $\mathbb H$ be $\GL(K,\R)$ or $\GL^+(K)$, and let $\mathbf h:\mathbb H\to\R$ be
continuous and constant on the conjugacy classes of $\mathbb H$. Then $\mathbf h(I_K+\mathbf T)=\mathbf h(I_K)$
for every nilpotent $\mathbf T$, that is, at every unipotent matrix $I_K+\mathbf T$, which lies in
$\GL^+(K)$ because $\det(I_K+\mathbf T)=1$. \status{ESTABLISHED}

\noindent\emph{Proof.}
Since $\mathbf T^K=0$, the kernels
$\ker\mathbf T\subseteq\ker\mathbf T^2\subseteq\cdots\subseteq\ker\mathbf T^K=\R^K$ increase, and
$\mathbf T$ maps $\ker\mathbf T^{k+1}$ into $\ker\mathbf T^k$. Listing a basis of $\ker\mathbf T$,
extending it to a basis of $\ker\mathbf T^2$, and so on, gives the columns of an invertible $\mathbf B$
with $\mathbf B^{-1}\mathbf T\mathbf B=\mathbf T_0$ strictly upper triangular. Replacing the first column
of $\mathbf B$ by its negative conjugates $\mathbf T_0$ by a diagonal matrix, which keeps it strictly
upper triangular, so $\det\mathbf B>0$ may be assumed. For $t>0$ let
$\mathbf S_t=\operatorname{diag}(1,t,\ldots,t^{K-1})$; then $\mathbf B\mathbf S_t\in\GL^+(K)\subseteq\mathbb H$
and $(\mathbf S_t^{-1}\mathbf T_0\mathbf S_t)_{kl}=t^{l-k}(\mathbf T_0)_{kl}$, which tends to zero as
$t\to0$ because $(\mathbf T_0)_{kl}=0$ unless $l>k$. Hence
$\mathbf h(I_K+\mathbf T)=\mathbf h\bigl((\mathbf B\mathbf S_t)^{-1}(I_K+\mathbf T)\mathbf B\mathbf S_t\bigr)
=\mathbf h(I_K+\mathbf S_t^{-1}\mathbf T_0\mathbf S_t)$ for every $t>0$, and continuity gives
$\mathbf h(I_K+\mathbf T)=\mathbf h(I_K)$. \hfill$\square$

\corollaryheading{No invariant loop defect on the general linear group}{cor:link-no-invariant-defect}
Let $K\geq2$ and $\mathbb H$ be as in \Cref{prop:link-class-function-blindness}. No continuous function
on $\mathbb H$ that is constant on its conjugacy classes vanishes exactly at $I_K$: every such function
that vanishes at $I_K$ vanishes at the shear $I_K+E_{12}\neq I_K$, and so does every function of the
eigenvalues, counted with multiplicity, that vanishes at $I_K$. Consequently, for $G=\GL(K,\R)$ or
$G=\GL^+(K)$ and a closed walk $p$ that traverses some reversal pair exactly once, a loop penalty that
is a continuous function of $\mathrm{Hol}^x(p)$ on $G$ and is unchanged by every passive re-expression
of every link family takes at every unipotent holonomy the value it takes at the trivial holonomy, so it
vanishes at every unipotent holonomy whenever it vanishes at the trivial one. In particular no norm on
$\mathfrak{gl}(K,\R)$ is invariant under conjugation by $\mathbb H$, which for $\mathbb H=\GL(K,\R)$ is
the observation of \Cref{rem:geo-loop-penalties}; the corollary is stronger, since it excludes every
continuous conjugation-invariant penalty and not only those of the form $g'\mapsto\lVert g'-I_K\rVert$.
\status{ESTABLISHED}

\noindent\emph{Proof.}
The first claim is \Cref{prop:link-class-function-blindness} at $\mathbf T=E_{12}$, which is
nilpotent because $E_{12}^2=0$; the shear is upper triangular with unit diagonal, so its eigenvalues are
those of $I_K$. If $p$ traverses the pair of $e$ exactly once, the link family with $\Theta^x_e$ equal
to a given $g'\in G$ or to $g'^{-1}$, according to the direction of the traversal, and with all other
links equal to $e_G$ has $\mathrm{Hol}^x(p)=g'$, and the passive re-expression with $g_{i_0}=g$ replaces
it by $gg'g^{-1}$ by \cref{eq:geo-edge-twist-holonomy}; the penalty is therefore a continuous class
function on $G$. For a norm $\lVert\cdot\rVert_\star$ invariant under conjugation by $\mathbb H$, the
function $g'\mapsto\lVert g'-I_K\rVert_\star$ is continuous, constant on conjugacy classes because
$gg'g^{-1}-I_K=g(g'-I_K)g^{-1}$, and vanishes exactly at $I_K$. \hfill$\square$

\definitionheading{Loop defect and cycle map}{def:link-loop-defect}
In \Cref{def:link-exponential-chart}, let $\rho_x:G\to\GL(\mathbb V)$ be a smooth homomorphism on a real
vector sample fiber $\mathbb V=\R^{\dim\mathbb V}$, and let $\mathsf G^{\mathrm{fib}}$ be a symmetric
positive definite matrix at the reference point: the matrix, in the coordinates of $\sigma$, of an inner
product on the fiber at $c_*$ of the vector bundle associated with $\rho_x$. It plays the role of the
fiber metric of \Cref{rem:geo-loop-penalties}, placed at $\sigma$ instead of at a vertex. On
$\mathfrak{gl}(\mathbb V)$ put
\begin{equation}
\langle\mathbf H,\mathbf H'\rangle_{\mathsf G}=\Tr\bigl((\mathsf G^{\mathrm{fib}})^{-1}\mathbf H^\top\mathsf G^{\mathrm{fib}}\mathbf H'\bigr),
\qquad
\lVert\mathbf H\rVert_{\mathsf G}=\langle\mathbf H,\mathbf H\rangle_{\mathsf G}^{1/2},
\label{eq:link-trace-form}
\end{equation}
and, when $\mathrm d\rho_x$ is injective, write
$\langle\mathring\omega,\mathring\omega'\rangle_{\mathsf G}=\langle\mathrm d\rho_x\mathring\omega,\mathrm d\rho_x\mathring\omega'\rangle_{\mathsf G}$
for the pulled-back inner product on $\mathfrak g$. The loop defect of a closed walk $p$ based at $i_0$
is the first quantity below, with $\mathring\Theta^x_p$ the reference holonomy
\cref{eq:link-reference-holonomy}. For a finite family $\mathbb L$ of closed walks, the cycle map
$\mathbf C:\mathscr T\to\mathfrak g^{\mathbb L}$ is the second, with $p=(e_0,\ldots,e_{r-1})$ and the
reversal convention \cref{eq:link-generator-reversal}:
\begin{gather}
\mathfrak Q^x_p(\mathring\omega^x)=\tfrac12\bigl\lVert I_{\mathbb V}-\rho_x(\mathring\Theta^x_p)\bigr\rVert_{\mathsf G}^2,
\label{eq:link-loop-defect}\\
(\mathbf C\mathring\omega^x)_p=\mathring\omega^x_{e_0}+\mathring\omega^x_{e_1}+\cdots+\mathring\omega^x_{e_{r-1}}.
\label{eq:link-cycle-map}
\end{gather}
Each reversal pair thus enters $(\mathbf C\mathring\omega^x)_p$ with its generator on $E_+$ times the
number of traversals of $p$ along the orientation in $E_+$ minus the number against it, and $\mathbf C$
is linear.

\lemmaheading{The trace form determines the metric up to scale}{lem:link-trace-form-scale}
For symmetric positive definite matrices $\mathsf G'$ and $\mathsf G''$ on $\mathbb V$, write
$\lVert\cdot\rVert_{\mathsf G'}$ and $\lVert\cdot\rVert_{\mathsf G''}$ for \cref{eq:link-trace-form} with
$\mathsf G'$ and $\mathsf G''$ in place of $\mathsf G^{\mathrm{fib}}$. Then
$\lVert\mathbf H\rVert_{\mathsf G'}=\lVert\mathbf H\rVert_{\mathsf G''}$ for every rank-one $\mathbf H$ exactly
when $\mathsf G''$ is a positive multiple of $\mathsf G'$, and in that case the two norms agree on all of
$\mathfrak{gl}(\mathbb V)$. \status{ESTABLISHED}

\noindent\emph{Proof.}
For $\mathbf H=\mathbf b\mathbf c^\top$, cyclicity of the trace gives
$\lVert\mathbf H\rVert^2_{\mathsf G'}=(\mathbf b^\top\mathsf G'\mathbf b)(\mathbf c^\top\mathsf G'^{-1}\mathbf c)$,
and likewise for $\mathsf G''$. If $\mathsf G''=\mathrm c_{\mathrm s}\mathsf G'$ with $\mathrm c_{\mathrm s}>0$,
the factors $\mathrm c_{\mathrm s}^{-1}$ and $\mathrm c_{\mathrm s}$ in \cref{eq:link-trace-form} cancel for
every $\mathbf H$. Conversely, fixing $\mathbf c\neq0$ gives
$\mathbf b^\top\mathsf G''\mathbf b=\mathrm c_{\mathrm s}\mathbf b^\top\mathsf G'\mathbf b$ for every
$\mathbf b$, with
$\mathrm c_{\mathrm s}=(\mathbf c^\top\mathsf G'^{-1}\mathbf c)/(\mathbf c^\top\mathsf G''^{-1}\mathbf c)>0$,
and symmetric matrices with proportional quadratic forms are proportional. \hfill$\square$

\propositionheading{Properties of the loop defect}{prop:link-loop-defect}
In \Cref{def:link-loop-defect}, let $p=(e_0,\ldots,e_{r-1})$ be a closed walk based at $i_0$.

(a) The form \cref{eq:link-trace-form} is an inner product on $\mathfrak{gl}(\mathbb V)$ with
$\lVert\mathbf H\mathbf H'\rVert_{\mathsf G}\leq\lVert\mathbf H\rVert_{\mathsf G}\lVert\mathbf H'\rVert_{\mathsf G}$
and $\lVert I_{\mathbb V}\rVert_{\mathsf G}=(\dim\mathbb V)^{1/2}$. The defect $\mathfrak Q^x_p$ is a
continuous function on $\mathscr T$ that depends on $\mathring\omega^x$ only through the links; it is
nonnegative and vanishes exactly when $\rho_x(\mathring\Theta^x_p)=I_{\mathbb V}$, equivalently when
$\rho_x(\mathrm{Hol}^x(p))=I_{\mathbb V}$.

(b) Passive re-expressions leave $\mathfrak Q^x_p$ unchanged. Under the reference change
$\sigma'=\sigma\cdot g_0$, with $\mathbf R=\rho_x(g_0)$ and the generators
$\operatorname{Ad}_{g_0^{-1}}\mathring\omega^x$ of \Cref{prop:link-generator-gauge}, the defect computed
at $\sigma'$ with the matrix $\mathbf R^\top\mathsf G^{\mathrm{fib}}\mathbf R$, which represents the same
inner product in the coordinates of $\sigma'$, equals $\mathfrak Q^x_p$ for every link family, while the
defect computed at $\sigma'$ with the matrix $\mathsf G^{\mathrm{fib}}$ equals the defect at $\sigma$
computed with $\mathbf R^{-\top}\mathsf G^{\mathrm{fib}}\mathbf R^{-1}$. The latter equals
$\mathfrak Q^x_p$ for every generator family whenever $\mathbf R^\top\mathsf G^{\mathrm{fib}}\mathbf R$ is
a positive multiple of $\mathsf G^{\mathrm{fib}}$, that is, whenever $\mathbf R$ is a nonzero multiple of a
$\mathsf G^{\mathrm{fib}}$-isometry, and, when
$\mathrm d\rho_x(\mathfrak g)=\mathfrak{gl}(\mathbb V)$ and $p$ traverses some reversal pair exactly
once, only then. Likewise, the defect of the rotated walk $(e_1,\ldots,e_{r-1},e_0)$ based at $i_1$ is
that of $p$ computed with
$\rho_x(\mathring\Theta^x_{e_0})^{-\top}\mathsf G^{\mathrm{fib}}\rho_x(\mathring\Theta^x_{e_0})^{-1}$.

(c) For each agent $i$ let
$\mathsf G^{\mathrm{fib}}_i=\rho_x(U_i)^{-\top}\mathsf G^{\mathrm{fib}}\rho_x(U_i)^{-1}$, the matrix of
the same inner product in the coordinates of $u_i$. Under passive re-expression
$\mathsf G^{\mathrm{fib}}_i$ becomes $(R^x_i)^{-\top}\mathsf G^{\mathrm{fib}}_i(R^x_i)^{-1}$, the law of
\Cref{rem:geo-loop-penalties}; the matrices of different agents satisfy
$\mathsf G^{\mathrm{fib}}_i=\rho_x(T_{ij})^{-\top}\mathsf G^{\mathrm{fib}}_j\rho_x(T_{ij})^{-1}$; and, with
$\mathsf U=\rho_x(\mathrm{Hol}^x(p))$,
\begin{equation}
\mathfrak Q^x_p=\tfrac12\Tr\bigl((\mathsf G^{\mathrm{fib}}_{i_0})^{-1}(I_{\mathbb V}-\mathsf U)^\top\mathsf G^{\mathrm{fib}}_{i_0}(I_{\mathbb V}-\mathsf U)\bigr),
\label{eq:link-defect-local}
\end{equation}
which is one half of the invariant penalty of \Cref{rem:geo-loop-penalties} at $i_0$. Conversely, if
matrices $\widetilde{\mathsf G}_{i_0}$ assigned to the frames of agent $i_0$ transform by the law of that
remark, then $\rho_x(U_{i_0})^\top\widetilde{\mathsf G}_{i_0}\rho_x(U_{i_0})$ is unchanged by passive
re-expression, and $\widetilde{\mathsf G}_{i_0}$ is its frame image.

(d) If $\rho_x(\mathring\Theta^x_p)$ is a $\mathsf G^{\mathrm{fib}}$-isometry, then
$\mathfrak Q^x_p=\dim\mathbb V-\Tr\rho_x(\mathring\Theta^x_p)=\dim\mathbb V-\Tr\mathsf U\in[0,2\dim\mathbb V]$,
which is $\dim\mathbb V$ times the plaquette term $1-\Tr\mathsf U/\dim\mathbb V$ of Wilson's lattice
action \citep[Eq.~(3.29)]{WilsonConfinement1974}, in the normalization of
\citet[Ch.~7, Eq.~(7.6)]{Creutz1983}.

(e) Let $\mathring\Omega_a=\mathrm d\rho_x(\mathring\omega^x_{e_a})$ and
$\mathring\varsigma_p=\sum_{a=0}^{r-1}\lVert\mathring\Omega_a\rVert_{\mathsf G}$. Then
\begin{equation}
\Bigl\lvert\mathfrak Q^x_p-\tfrac12\bigl\lVert\mathrm d\rho_x(\mathbf C\mathring\omega^x)_p\bigr\rVert^2_{\mathsf G}\Bigr\rvert
\leq\mathring\varsigma_p\bigl(\exp(\mathring\varsigma_p)-1-\mathring\varsigma_p\bigr)
+\tfrac12\bigl(\exp(\mathring\varsigma_p)-1-\mathring\varsigma_p\bigr)^2,
\label{eq:link-defect-expansion}
\end{equation}
and the right side is at most
$(\tfrac12\exp(1)+\tfrac18\exp(2))\mathring\varsigma_p^3<\tfrac{23}{10}\mathring\varsigma_p^3$ when
$\mathring\varsigma_p\leq1$. When $\mathrm d\rho_x$ is injective and $\mathscr T$ carries the norm
$\lVert\mathring\omega\rVert^2=\sum_{e\in E_+}\lVert\mathring\omega_e\rVert^2_{\mathsf G}$, one has
$\mathring\varsigma_p\leq r\lVert\mathring\omega^x\rVert$, so the bound is
$\tfrac{23}{10}r^3\lVert\mathring\omega^x\rVert^3$ on the ball $r\lVert\mathring\omega^x\rVert\leq1$.
\status{ESTABLISHED}

The proof is in \cref{sec:proofs-links}.

The matrix at the reference point is thus the one datum beyond the group that the defect uses. Carried
along with the reference point, it leaves the defect invariant under every reference change; held as a
matrix, it leaves the defect invariant under the reference changes acting by nonzero multiples of its
isometries, and, when $\mathrm d\rho_x(\mathfrak g)=\mathfrak{gl}(\mathbb V)$ and the walk traverses some
reversal pair exactly once, under no other (\Cref{prop:link-loop-defect}(b)). The symmetry of the measurement is reduced in this way from all
reference changes to those acting by multiples of isometries of $\mathsf G^{\mathrm{fib}}$; on
holonomies that are isometries the defect is Wilson's trace density (part (d)), and one reference matrix
supplies the metric of \Cref{rem:geo-loop-penalties} at every vertex and for every closed walk (part
(c)). A cyclic rotation of a closed walk replaces the matrix by its image under the twist moved across
(part (b)), so the base vertex does not affect the defect when that twist acts by a multiple of an
isometry. The quadratic part of the defect is carried by the cycle map (part (e)), which reappears in
\Cref{prop:link-linearized-equation}.

\section{The generator prior}
\label{sec:link-prior}

The generator chart supplies a link prior for every Lie group: a Gaussian base law on the generators,
which exists whether or not $G$ is compact, multiplied by the exponential of minus a nonnegative loop
action (\Cref{prop:link-generator-prior}). It replaces the Haar reference, whose normalizability the scope
of \Cref{ch:geometry} declines for noncompact and general linear $G$, by a reference on the Lie algebra,
and the abelian links of \Cref{ch:finite-depth-formation} are its translation-group instance
(\Cref{prop:link-translation-instance}).

\definitionheading{Generator prior}{def:link-generator-prior}
In \Cref{def:link-exponential-chart,def:link-loop-defect}, equip $\mathfrak g$ with an inner product,
written $\langle\cdot,\cdot\rangle_{\mathsf G}$ (the pulled-back inner product of
\Cref{def:link-loop-defect} when $\mathrm d\rho_x$ is injective), and $\mathscr T$ with the product inner
product and the Lebesgue measure $\mathrm d\mathring\omega$ normalized by it, under which the cube
spanned by an orthonormal basis has volume one. For each $e\in E_+$ let the anchor
$\mathring\alpha^x_e\in\mathfrak g$ and the covariance $\mathring\Sigma^x_e$, a positive definite
self-adjoint operator on $(\mathfrak g,\langle\cdot,\cdot\rangle_{\mathsf G})$, be generative parameters,
that is, coordinates of $\theta$, and let $\mathbb L$ be a finite family of closed walks with coefficients
$\mathrm c_p\geq0$. The base law, link action, normalizer, and generator prior are
\begin{gather}
\nu^{\mathscr T}=\bigotimes_{e\in E_+}\mathcal N(\mathring\alpha^x_e,\mathring\Sigma^x_e),
\qquad
\mathscr S(\mathring\omega^x)=\sum_{p\in\mathbb L}\mathrm c_p\mathfrak Q^x_p(\mathring\omega^x),
\label{eq:link-base-action}\\
Z_{\mathscr S}=\int_{\mathscr T}e^{-\mathscr S}d\nu^{\mathscr T},
\qquad
Z_{\mathscr S}^{-1}e^{-\mathscr S}\nu^{\mathscr T},
\label{eq:link-generator-prior}
\end{gather}
where $\mathcal N(\mathring\alpha,\mathring\Sigma)$ has the density
$(2\pi)^{-\dim\mathfrak g/2}(\det\mathring\Sigma)^{-1/2}\exp\bigl(-\tfrac12\langle\mathring\omega-\mathring\alpha,\mathring\Sigma^{-1}(\mathring\omega-\mathring\alpha)\rangle_{\mathsf G}\bigr)$
with respect to the normalized Lebesgue measure of $\mathfrak g$. With independent links in the two
channels the link space is $\mathscr T\times\mathscr T$, the base law is the product of the two base
laws, and the action is the sum of the two channel actions; under the common-link specialization one
base law on $\mathscr T$ serves both channels and the action is the sum over channels.

\propositionheading{Normalization and covariance of the generator prior}{prop:link-generator-prior}
Let $G$ be a Lie group, compact or not, and let the data of \Cref{def:link-generator-prior} be given.

(a) The action $\mathscr S$ is finite, continuous, and nonnegative on $\mathscr T$ and has finite
$\nu^{\mathscr T}$-mean, and
\begin{equation}
0<\exp\bigl(-\E_{\nu^{\mathscr T}}\mathscr S\bigr)\leq Z_{\mathscr S}\leq1;
\label{eq:link-normalizer-bounds}
\end{equation}
when $\rho_x(G)$ consists of $\mathsf G^{\mathrm{fib}}$-isometries, also
$Z_{\mathscr S}\geq\exp(-2\dim\mathbb V\sum_p\mathrm c_p)$. The generator prior is a probability law on
the standard Borel space $\mathscr T$ with a strictly positive continuous density with respect to
$\mathrm d\mathring\omega$. Hence $\nu^{\mathscr T}$, $\mathscr S$, and $Z_{\mathscr S}$ satisfy the
hypotheses that \Cref{prop:elbo-random-link-vfe} places on its base law, action, and normalizer, with the
generator family as the latent link variable and the links $\exp_\sigma(\mathring\omega^x)$ as a
continuous function of it. With $\widetilde{\mathscr S}$ the action written as a function of the link
family through $\mathring\Theta^x_p=U_{i_0}^{-1}\mathrm{Hol}^x(p)U_{i_0}$, so that
$\mathscr S=\widetilde{\mathscr S}\circ\exp_\sigma$, the image of the prior under $\exp_\sigma$ is
$Z_{\mathscr S}^{-1}e^{-\widetilde{\mathscr S}}(\exp_\sigma)_\#\nu^{\mathscr T}$, with the same normalizer.

(b) Passive re-expressions leave the generator prior unchanged and carry its image on link families by
the vertex law \cref{eq:geo-regime-two-gauge-law}.

(c) Let $\sigma'=\sigma\cdot g_0$ and $\mathbf R=\rho_x(g_0)$, and let $\operatorname{Ad}_{g_0^{-1}}$ act
on every component of $\mathscr T$. The image of the generator prior under $\operatorname{Ad}_{g_0^{-1}}$
is the generator prior relative to $\sigma'$ with the anchors
$\operatorname{Ad}_{g_0^{-1}}\mathring\alpha^x_e$, the covariances
$\operatorname{Ad}_{g_0^{-1}}\mathring\Sigma^x_e\operatorname{Ad}_{g_0}$, the fiber matrix
$\mathbf R^\top\mathsf G^{\mathrm{fib}}\mathbf R$, the inner product
$\langle\operatorname{Ad}_{g_0}\cdot,\operatorname{Ad}_{g_0}\cdot\rangle_{\mathsf G}$ on $\mathfrak g$, which is the
pullback for $\mathbf R^\top\mathsf G^{\mathrm{fib}}\mathbf R$ when $\mathrm d\rho_x$ is injective, and the
same walks and coefficients; both priors have the same image on link families.

(d) If $\operatorname{Ad}_{g_0}$ is an isometry of $\langle\cdot,\cdot\rangle_{\mathsf G}$, fixes every
anchor, and commutes with every covariance, and $\mathbf R$ is a nonzero multiple of a
$\mathsf G^{\mathrm{fib}}$-isometry, then the image of the generator prior under
$\operatorname{Ad}_{g_0^{-1}}$ is the generator prior itself, so the reference change by $g_0$ leaves the
anchors, covariances, and coefficients unchanged and multiplies the fiber matrix by a positive scalar,
which changes neither the inner product nor the defect. For the pulled-back inner product the isometry
condition follows from the condition on $\mathbf R$, and when $\mathrm d\rho_x(\mathfrak g)=\mathfrak{gl}(\mathbb V)$
it is equivalent to it. \status{ESTABLISHED}

The proof is in \cref{sec:proofs-links}.

The Haar-referenced density of \Cref{rem:geo-loop-penalties} needs the finite invariant measure of a
compact group, whereas the base law of \Cref{def:link-generator-prior} lives on the Lie algebra, so
\cref{eq:link-normalizer-bounds} holds for every Lie group; the price is a set of anchors and covariances
that a reference change transforms (\Cref{prop:link-generator-prior}).

\propositionheading{Translation links}{prop:link-translation-instance}
Let $G=(\R^K,+)$ act on the sample fiber $\R^K$ by translation, and write the group additively.

(a) The affine representation and its differential,
\begin{equation}
\rho(\mathring\omega)=\begin{pmatrix}I_K&\mathring\omega\\0&1\end{pmatrix},
\qquad
\mathrm d\rho(\mathring\omega)=\begin{pmatrix}0&\mathring\omega\\0&0\end{pmatrix},
\label{eq:link-affine-representation}
\end{equation}
form a faithful smooth linear representation on $\R^{K+1}$ with $\mathfrak g=\R^K$ and $\exp$ the
identity, and $\rho(\mathring\omega)(k,1)=(k+\mathring\omega,1)$, so the translations of the sample fiber
act on the affine slice $\R^K\times\{1\}$. A link $\Theta_e\in\R^K$ on $e:j\to i$, acting by
$k\mapsto k+\Theta_e$, has the generator $\mathring\omega_e=\Theta_e-U_i+U_j$, and the reference holonomy
of a closed walk is the loop sum $\mathring\Theta_p=(\mathbf C\mathring\omega)_p=\sum_a\Theta_{e_a}$, in
which the frames cancel.

(b) For every symmetric positive definite $\mathsf G^{\mathrm{fib}}$ on $\R^{K+1}$, the pulled-back norm
and the loop defect are
\begin{equation}
\lVert\mathring\omega\rVert^2_{\mathsf G}
=\bigl((\mathsf G^{\mathrm{fib}})^{-1}\bigr)_{K+1,K+1}(\mathring\omega,0)^\top\mathsf G^{\mathrm{fib}}(\mathring\omega,0),
\qquad
\mathfrak Q_p=\tfrac12\bigl\lVert(\mathbf C\mathring\omega)_p\bigr\rVert^2_{\mathsf G},
\label{eq:link-translation-defect}
\end{equation}
so the left side of \cref{eq:link-defect-expansion} vanishes; for $\mathsf G^{\mathrm{fib}}=I_{K+1}$,
$\mathfrak Q_p=\tfrac12\lvert\sum_a\Theta_{e_a}\rvert^2$.

(c) In the four-cycle experiment of \Cref{def:rg-formation-cycle-experiment} and the gated translation
family of \Cref{def:rg-gated-translation-family}, where $K=1$ and the link translations $\omega_e$ are
the links $\Theta_e$, the configuration prior is $\mathcal N(0,I)$, so the link translations are
independent of the means and have the law $\prod_e\mathcal N(0,1)$. In generator coordinates this is the
base law of \Cref{def:link-generator-prior} for $\mathsf G^{\mathrm{fib}}=I_2$ with unit covariances and
the anchors $\mathring\alpha_e=U_j-U_i$, which vanish when every frame equals $\sigma$ and always have
zero loop sums, and the action is $\mathscr S=0$. The loop $p=(e_3,e_2,e_1,e_0)$ of the four-cycle,
based at agent $0$, has $(\mathbf C\mathring\omega)_p=\sum_i\omega_i$ and
$\mathfrak Q_p=\tfrac12(\sum_i\omega_i)^2$; in these models it enters only through the edge costs, whose
residuals satisfy $\sum_ir^b_i=\sum_ir^m_i=-\sum_i\omega_i$.

(d) In the triangle-chain experiment of \Cref{def:rg-formation-experiment}, with the symbols of
\cref{tab:rg-formation-notation} and $\mathsf G^{\mathrm{fib}}=I_2$, every cyclically oriented triangle
$q$ has $\mathfrak h_q=(\mathbf C\mathring\omega)_q$ and loop energy
$\mathrm J_1h_q^2=\tfrac23\mathrm J_1\mathfrak Q_q$, and for every configuration
\begin{equation}
\frac{\mathbf x^\top\Omega_0(\vartheta)\mathbf x}{2\mathrm T}
=\sum_q\frac{\vartheta\mathrm J_1}{3\mathrm T}\mathfrak Q_q+\frac{\mathbf x^\top\Omega^\perp_0(\vartheta)\mathbf x}{2\mathrm T},
\qquad
\Omega^\perp_0(\vartheta)=\vartheta B^\top\operatorname{diag}(\mathrm J_\ell I)B+P_{\mathcal K}-\mathrm TR_{\mathrm{ev}},
\label{eq:link-triangle-chain-prior}
\end{equation}
where the kernel of $\Omega^\perp_0(\vartheta)$ is the loop space $\mathcal H$ and $\Omega^\perp_0(\vartheta)$ is
positive definite on $\mathcal U\oplus\mathcal K$. Under the configuration prior the loop sums
$\mathfrak h_q$ are therefore independent with law $\mathcal N(0,3\mathrm T/(\vartheta\mathrm J_1))$ and
independent of the projection of $\mathbf x$ onto $\mathcal U\oplus\mathcal K$, and the quadratic energies
$\mathrm E_\vartheta$ and $\mathrm E^{\mathrm{phys}}_\vartheta$ of
\cref{eq:rg-formation-partition-integrals} are $\tfrac13\vartheta\mathrm J_1\sum_q\mathfrak Q_q$ plus a
function of that projection. This prior differs from \Cref{def:link-generator-prior} in that its second
factor is a Gaussian energy of means and links jointly that is constant along $\mathcal H$, so it is not
integrable alone, and the loop defect with coefficient $\vartheta\mathrm J_1/(3\mathrm T)$ supplies all
of the prior's precision along loop coordinates. \status{ESTABLISHED}

The proof is in \cref{sec:proofs-links}.

\remarkheading{The running example in the chart}{rem:link-running-example}
In the running example of \Cref{sec:geo-running-example} take $E_+=\{e_1,e_2,e_3\}$, $\sigma=u_1$, and
$U_i=e_G$, and write $R_x(\phi)=\exp(\phi(E_{32}-E_{23}))$ and $R_z(\phi)=\exp(\phi(E_{21}-E_{12}))$ for
the rotations about the first and third axes. The chart covers every link family of $SO(3)$
(\Cref{prop:link-chart-image}), and the generators of rotation angle at most $\pi$ are
$\mathring\omega^b_{e_1}=\tfrac\pi2(E_{21}-E_{12})$, $\mathring\omega^b_{e_2}=\tfrac\pi2(E_{32}-E_{23})$,
and $\mathring\omega^b_{e_3}=0$. The reference holonomy of $p=(e_1,e_2,e_3)$ is
$\mathring\Theta^b_p=R_z(\tfrac\pi2)R_x(\tfrac\pi2)$, the cyclic permutation matrix that sends each
standard basis vector to the next one, and it equals $\mathrm{Hol}^b(p)$, a rotation by $2\pi/3$ about the
axis spanned by $(1,1,1)$. With $\mathsf G^{\mathrm{fib}}=I_3$ every rotation is an isometry, a rotation
by $\phi$ has defect $4\sin^2(\phi/2)\in[0,4]$, and $\mathfrak Q^b_p=3-\Tr\mathring\Theta^b_p=3$, while the
trivial model representation gives $\mathfrak Q^m_p=0$. The cycle map gives
$(\mathbf C\mathring\omega^b)_p=\tfrac\pi2(E_{21}-E_{12}+E_{32}-E_{23})$, whose quadratic term
$\tfrac12\lVert(\mathbf C\mathring\omega^b)_p\rVert^2_{\mathsf G}=\pi^2/2$ exceeds the defect; here
$\mathring\varsigma_p=\pi\sqrt2>1$, outside the ball of \Cref{prop:link-loop-defect}(e). A passive
re-expression $g_2$ changes $\Theta^b_{e_1}$ and $\Theta^b_{e_2}$ but no generator.
\status{ESTABLISHED}

\noindent\emph{Proof.}
The series gives $R_z(\tfrac\pi2)$ and $R_x(\tfrac\pi2)$ as the quarter turns of \cref{eq:geo-running-links},
whose product is the displayed permutation matrix; it fixes $(1,1,1)$ and has trace $0=1+2\cos\phi$ with
$\phi=2\pi/3$. A rotation by $\phi$ has eigenvalues $1$ and $e^{\pm i\phi}$, so
\Cref{prop:link-loop-defect}(d) gives the defect $3-(1+2\cos\phi)=4\sin^2(\phi/2)$, and
$\mathfrak Q^b_p=3-0$. The four nonzero entries of $(\mathbf C\mathring\omega^b)_p$ have modulus $\pi/2$,
so its squared Frobenius norm is $\pi^2$, and each nonzero generator has norm $\pi/\sqrt2$. Generators
are unchanged by passive re-expression by \Cref{prop:link-generator-gauge}(a). \hfill$\square$

\section{Latent generators in the collective evidence bound}
\label{sec:link-latent}

What joint law and what collective free energy result when the link values of a channel leave the
structural data and the generators of the exponential chart become latent coordinates? The joint has
unit mass on every finite record hypergraph, and its free energy splits into a link sector and the
conditional free energies of the held-link model (\cref{prop:link-latent-normalization}); for a Gaussian
recognition law of the generators every term is finite and the link sector is the relative entropy from
the generator prior (\cref{prop:link-gaussian-recognition-vfe}). Throughout, the generator prior is that
of \cref{def:link-generator-prior} in the following form: the base law $\nu^{\mathscr T}$ is the product,
over the copies $e\in E_+$, of the Gaussian laws
$\mathcal N(\mathring\alpha^x_e,\mathring\Sigma^x_e)$ on $(\mathfrak g,\langle\cdot,\cdot\rangle_{\mathsf G})$,
written $\mathcal N(\mathring\alpha,\mathring\Sigma)$ with a block-diagonal covariance $\mathring\Sigma$; the
action $\mathscr S$ is a finite nonnegative combination of the loop defects $\mathfrak Q^x_p$ of
\cref{def:link-loop-defect}; and $0<Z_{\mathscr S}\leq1$ by \cref{prop:link-generator-prior}. The
definitions and statements of this section and of \cref{sec:link-field,sec:link-held,sec:link-induced} are
written for the link space $\mathscr T$ of one latent channel $x$, which covers the common-link
specialization. When both channels carry independent latent generator families, the link space is
$\mathscr T\times\mathscr T$ with the product of the two base laws and the sum of the two actions
(\cref{def:link-exponential-chart,def:link-generator-prior}); every definition and statement then holds with
these data in place of $\mathscr T$, $\nu^{\mathscr T}$, and $\mathscr S$, with the summands indexed by the
pairs $(x,e)$, by the same proof, since the proofs use only that the link space is a finite-dimensional
inner product space carrying a Gaussian base law that factors over orthogonal summands and an action that
is a finite nonnegative combination of loop defects. The space
$\mathscr T$ carries the orthogonal sum of the inner products of its summands, its norm
$\lVert\cdot\rVert$, and the Lebesgue measure normalized by it; gradients and Hessians on $\mathscr T$
refer to this inner product, $\lVert\cdot\rVert$ is the operator norm on matrices, and upright
$\mathrm c_1,\mathrm c_2,\ldots$ are constants local to one statement.

\definitionheading{Joint law with latent generators}{def:link-latent-model}
Let the structural data $X$ contain the interaction multigraph $\mathfrak I$, an orientation set $E_+$ with
one copy of each reversal pair, the reference point $\sigma$ with the frames $U_i$, $u_i=\sigma\cdot
U_i^{-1}$, and the acyclic graph $\Gamma$ of \cref{def:gen-parameter-split}, but no link values of the
channels $x$ whose links are latent. For those channels the generator family
$\mathring\omega=(\mathring\omega^x_e)$ is a latent coordinate with values in $\mathscr T$, or in
$\mathscr T\times\mathscr T$ when both channels are latent, and it determines the links of the exponential
chart of
\cref{def:link-exponential-chart},
\begin{equation}
\Theta^x_e(\mathring\omega)=U_i\exp(\mathring\omega^x_e)U_j^{-1},
\qquad
\Theta^x_{\bar e}(\mathring\omega)=\Theta^x_e(\mathring\omega)^{-1},
\qquad e:j\to i\text{ in }E_+,
\label{eq:link-generator-links}
\end{equation}
a continuous map $\mathring\omega\mapsto\Theta(\mathring\omega)$ into the link families. Write $(X,\Theta)$ for
$X$ with a link family $\Theta$ adjoined, and let the factors \eqref{eq:gen-kernels} and the record kernels
of \cref{hyp:local-interaction-kernels} be measurable in the link values, so that for each $\Theta$ the
held-link model of \cref{ch:generative,ch:elbo} at $(X,\Theta)$ has the joint
$\mathbb P_{\theta,V}(dy,do\given X,\Theta)=\mathbb P^{\mathsf Y}_{\theta,V}(dy\given X,\Theta)\prod_{a\in\mathfrak A}K_a(X,\Theta,y_{\partial a},do_a)$,
the evidence $p_{\theta,V}(o\given X,\Theta)$, the posterior $\boldsymbol\Pi_{\theta,V,o,X,\Theta}$, and, for
laws $Q$ on $\mathsf Y$, the free energy
$\Fenergy_V[Q;X,\Theta,o]=-\log p_{\theta,V}(o\given X,\Theta)+\KL(Q\Vert\boldsymbol\Pi_{\theta,V,o,X,\Theta})$ of
\cref{def:elbo-extended}. With the generator prior
$\mathbb P^{\mathscr T}_\theta(d\mathring\omega\given X)=Z_{\mathscr S}^{-1}e^{-\mathscr S(\mathring\omega)}\nu^{\mathscr T}(d\mathring\omega)$,
the joint law with latent generators is the law on $\mathscr T\times\mathsf Y\times\mathsf O$
\begin{equation}
\mathbb P_{\theta,V}(d\mathring\omega,dy,do\given X)
=\mathbb P^{\mathscr T}_\theta(d\mathring\omega\given X)
\mathbb P^{\mathsf Y}_{\theta,V}\bigl(dy\given X,\Theta(\mathring\omega)\bigr)
\prod_{a\in\mathfrak A}K_a\bigl(X,\Theta(\mathring\omega),y_{\partial a},do_a\bigr),
\label{eq:link-latent-joint}
\end{equation}
whose evidence $p_{\theta,V}(o\given X)$, posterior $\boldsymbol\Pi_{\theta,V,o,X}(d\mathring\omega,dy)$, and
extended free energy $\Fenergy_V[Q;X,o]$ for laws $Q$ on $\mathscr T\times\mathsf Y$ are those of
\cref{def:elbo-extended} with the latent space $\mathscr T\times\mathsf Y$ and the baseline
$\mathbb P^{\mathscr T}_\theta(d\mathring\omega\given X)\mathbb P^{\mathsf Y}_{\theta,V}(dy\given X,\Theta(\mathring\omega))$.
With $\theta$, $X$ and $o$ held, write
$p(o\given\mathring\omega)=p_{\theta,V}(o\given X,\Theta(\mathring\omega))$,
$\boldsymbol\Pi_{\mathring\omega}=\boldsymbol\Pi_{\theta,V,o,X,\Theta(\mathring\omega)}$,
$\mathbb P^{\mathsf Y}(\cdot\given\mathring\omega)=\mathbb P^{\mathsf Y}_{\theta,V}(\cdot\given X,\Theta(\mathring\omega))$,
and, for records with densities \eqref{eq:obs-interaction-density},
$L_o(\mathring\omega,y)=\prod_{a}e^{-E_{a,o}(X,\Theta(\mathring\omega),y_{\partial a})}$ and
$H_o(\mathring\omega,y)=\sum_aE_{a,o}(X,\Theta(\mathring\omega),y_{\partial a})$. The anchors
$\mathring\alpha^x_e$, the prior covariances $\mathring\Sigma^x_e$, and the loop coefficients are coordinates
of $\theta$, and the walks of $\mathscr S$ belong to $X$. The prior reads $\theta$ and $X$; each agent factor
reads $\theta$, $X$, its parents and $\mathring\omega$; each record reads $\theta$, $X$, $y_{\partial a}$ and
$\mathring\omega$. No factor reads a recognition law, a variational parameter, or the posterior, so
\cref{req:gen-typing-prohibition} holds: moving an anchor or a prior covariance is an M-coordinate
(\cref{def:elbo-generalized-m-acceptance}), while the generators are moved by recognition. Passive
re-expression changes the frames and no generator (\cref{prop:geo-edge-twist-invariance}); a change of
reference point acts on the generator prior as in \cref{prop:link-generator-gauge}.

\propositionheading{Normalization, evidence, and the conditional decomposition}{prop:link-latent-normalization}
Assume \cref{def:link-latent-model}. (a) For every finite record hypergraph, cyclic or not, and every
$(\theta,X)$, \eqref{eq:link-latent-joint} is a probability law on $\mathscr T\times\mathsf Y\times\mathsf O$,
and the generator family is a root block of the directed latent law. (b) If the records have densities
\eqref{eq:obs-interaction-density}, then
$\mathring\omega\mapsto p(o\given\mathring\omega)=\int L_o(\mathring\omega,y)\mathbb P^{\mathsf Y}(dy\given\mathring\omega)$
is measurable and
\begin{equation}
p_{\theta,V}(o\given X)=\int_{\mathscr T}p(o\given\mathring\omega)\mathbb P^{\mathscr T}_\theta(d\mathring\omega\given X);
\label{eq:link-latent-evidence}
\end{equation}
the evidence is positive exactly when $p(o\given\cdot)$ is positive on a set of positive Lebesgue measure,
and if every record density is bounded, $e^{-E_{a,o}}\leq\mathrm c_a$, then
$p(o\given\mathring\omega)\leq\mathrm c_{\mathfrak A}:=\prod_a\mathrm c_a$ for every $\mathring\omega$ and the evidence is at
most $\mathrm c_{\mathfrak A}$; binary success records have $\mathrm c_a=1$. (c) Suppose further that the record
densities are bounded and $p(o\given\mathring\omega)>0$ for every $\mathring\omega$. Then the posterior
disintegrates as
$\boldsymbol\Pi_{\theta,V,o,X}(d\mathring\omega,dy)=\boldsymbol\Pi^{\mathscr T}(d\mathring\omega)\boldsymbol\Pi_{\mathring\omega}(dy)$
with link marginal
$\boldsymbol\Pi^{\mathscr T}(d\mathring\omega)=p_{\theta,V}(o\given X)^{-1}p(o\given\mathring\omega)\mathbb P^{\mathscr T}_\theta(d\mathring\omega\given X)$,
and every law $Q$ on $\mathscr T\times\mathsf Y$, disintegrated as $Q(d\mathring\omega,dy)=R(d\mathring\omega)Q_{\mathring\omega}(dy)$,
satisfies
\begin{equation}
\Fenergy_V[Q;X,o]
=\KL\bigl(R\Vert\mathbb P^{\mathscr T}_\theta(\cdot\given X)\bigr)
+\int_{\mathscr T}\Fenergy_V\bigl[Q_{\mathring\omega};X,\Theta(\mathring\omega),o\bigr]R(d\mathring\omega)
\qquad\text{in }(-\infty,+\infty],
\label{eq:link-decomposition}
\end{equation}
with a measurable integrand bounded below by $-\log\mathrm c_{\mathfrak A}$. When $\KL(R\Vert\nu^{\mathscr T})<\infty$
and $\E_R\lvert\mathscr S\rvert<\infty$, the first term equals
$\KL(R\Vert\nu^{\mathscr T})+\E_R\mathscr S+\log Z_{\mathscr S}$, and \eqref{eq:link-decomposition} is
\eqref{eq:rg-geometry-random-link-vfe} for the latent variable $\mathring\omega$. \status{ESTABLISHED}

The proof is in \cref{sec:proofs-links}.

\lemmaheading{Exponential moments and growth of loop actions}{lem:link-growth}
(a) For a Gaussian law $\mathcal N(\mu,\Sigma)$ on $\mathscr T$ and every $\mathrm c\geq0$,
$\E_{\mathcal N(\mu,\Sigma)}e^{\mathrm c\lVert\mathring\omega\rVert}\leq e^{\mathrm c\lVert\mu\rVert}\bigl(2e^{\mathrm c^2\lVert\Sigma\rVert/2}\bigr)^{\dim\mathscr T}$,
and the same bound holds on each summand of $\mathscr T$ with its dimension.
(b) For a representation $\rho_x$ on a real vector space of finite dimension, the map
$\mathring\xi\mapsto\rho_x(\exp\mathring\xi)=\exp(\mathrm d\rho_x\mathring\xi)$ on
$(\mathfrak g,\langle\cdot,\cdot\rangle_{\mathsf G})$ is smooth, and its value and its first and second
derivatives at $\mathring\xi$ have norms at most $e^{\lVert\mathrm d\rho_x\rVert\lVert\mathring\xi\rVert}$,
$\lVert\mathrm d\rho_x\rVert e^{\lVert\mathrm d\rho_x\rVert\lVert\mathring\xi\rVert}$ and
$\lVert\mathrm d\rho_x\rVert^2e^{\lVert\mathrm d\rho_x\rVert\lVert\mathring\xi\rVert}$. (c) Every real polynomial
$\mathbf h$ in the matrix entries of the represented twists
$\rho_x(\mathring\Theta^x_e)^{\pm1}=\exp(\pm\mathrm d\rho_x\mathring\omega^x_e)$, $e\in E_+$, is smooth on
$\mathscr T$ and satisfies, for constants depending on $\mathbf h$,
\begin{equation}
\lvert\mathbf h(\mathring\omega)\rvert+\lVert\nabla\mathbf h(\mathring\omega)\rVert
+\lVert\nabla^2\mathbf h(\mathring\omega)\rVert\leq\mathrm c_1e^{\mathrm c_2\lVert\mathring\omega\rVert};
\label{eq:link-growth}
\end{equation}
each loop defect of \cref{def:link-loop-defect}, and hence the action $\mathscr S$, is such a polynomial.
\status{ESTABLISHED}

\noindent\emph{Proof.}
(a) Write $\mathring\omega=\mu+\Sigma^{1/2}\mathbf b$ with $\mathbf b$ standard Gaussian in an orthonormal
basis, so that $\lVert\mathring\omega\rVert\leq\lVert\mu\rVert+\lVert\Sigma\rVert^{1/2}\sum_l\lvert\mathbf b_l\rvert$;
the coordinates are independent and, for $\mathrm c'\geq0$,
$\E e^{\mathrm c'\lvert\mathbf b_l\rvert}\leq\E e^{\mathrm c'\mathbf b_l}+\E e^{-\mathrm c'\mathbf b_l}=2e^{\mathrm c'^2/2}$.
(b) The identity $\rho_x(\exp\mathring\xi)=\exp(\mathrm d\rho_x\mathring\xi)$ is
\cref{lem:link-exp-naturality}. The series $\exp\mathbf B=\sum_l\mathbf B^l/l!$ converges
with all its derivatives. The first derivative of $\mathbf B\mapsto\mathbf B^l$ in the direction $\mathbf T$ is a
sum of $l$ products of $l-1$ factors $\mathbf B$ and one factor $\mathbf T$, and its second derivative in the
directions $\mathbf T$ and $\mathbf T'$ is a sum of $l(l-1)$ products of $l-2$ factors $\mathbf B$ and one factor
each of $\mathbf T$ and $\mathbf T'$. Summing the series, the value and the first and second derivatives of the
exponential have norms at most $e^{\lVert\mathbf B\rVert}$, $e^{\lVert\mathbf B\rVert}\lVert\mathbf T\rVert$ and
$e^{\lVert\mathbf B\rVert}\lVert\mathbf T\rVert\lVert\mathbf T'\rVert$; composing with the linear map $\mathrm d\rho_x$
gives (b). (c) By (b) every entry of
$\exp(\pm\mathrm d\rho_x\mathring\omega^x_e)$ and its first two derivatives are bounded by
$\mathrm c e^{\lVert\mathrm d\rho_x\rVert\lVert\mathring\omega\rVert}$, and the Leibniz rule bounds a monomial of
degree $\mathrm c_0$ and its derivatives by $\mathrm c'e^{\mathrm c_0\lVert\mathrm d\rho_x\rVert\lVert\mathring\omega\rVert}$. The
reference-coordinate holonomy of a closed walk is the product of the matrices
$\rho_x(\mathring\Theta^x_e)^{\pm1}$ along the walk (\cref{prop:geo-edge-twist-invariance}), and a loop
defect is a quadratic polynomial in its entries (\cref{def:link-loop-defect}). \hfill$\square$

\propositionheading{Gaussian recognition of the generators}{prop:link-gaussian-recognition-vfe}
Let $R=\mathcal N(\mathring\mu,\mathring\Lambda^{-1})$ be a Gaussian law on $\mathscr T$ with mean
$\mathring\mu$ and positive-definite self-adjoint precision $\mathring\Lambda$; in the mean-field case
$R=\bigotimes\mathcal N(\mathring\mu^x_e,(\mathring\Lambda^x_e)^{-1})$ over the summands of $\mathscr T$.
(a) The relative entropy from the base law is finite and equals
\begin{equation}
\KL(R\Vert\nu^{\mathscr T})
=\tfrac12\Bigl[\Tr\bigl(\mathring\Sigma^{-1}\mathring\Lambda^{-1}\bigr)
+\bigl\langle\mathring\mu-\mathring\alpha,\mathring\Sigma^{-1}(\mathring\mu-\mathring\alpha)\bigr\rangle
-\dim\mathscr T+\log\det\mathring\Sigma+\log\det\mathring\Lambda\Bigr];
\label{eq:link-gaussian-kl}
\end{equation}
in the mean-field case it is the sum of the same expressions over the summands, and in general it equals
$\operatorname{TC}(R)+\sum\KL(R^x_e\Vert\mathcal N(\mathring\alpha^x_e,\mathring\Sigma^x_e))$, with $R^x_e$ the
summand marginals and
$\operatorname{TC}(R)=\frac12[\sum\log\det(\mathring\Lambda^{-1})^x_e+\log\det\mathring\Lambda]$, where
$(\mathring\Lambda^{-1})^x_e$ are the diagonal blocks of $\mathring\Lambda^{-1}$. (b)
$\E_Re^{\mathrm c\lVert\mathring\omega\rVert}<\infty$ for every $\mathrm c\geq0$ and $\E_R\mathscr S<\infty$.
Consequently, if the record densities are bounded and $p(o\given\mathring\omega)>0$ for every
$\mathring\omega$, then for every Markov kernel $\mathring\omega\mapsto Q_{\mathring\omega}$ the free energy of
$Q=R(d\mathring\omega)Q_{\mathring\omega}(dy)$ is \eqref{eq:rg-geometry-random-link-vfe} for the latent variable
$\mathring\omega$, with the finite link sector
\begin{equation}
\KL(R\Vert\nu^{\mathscr T})+\E_R\mathscr S+\log Z_{\mathscr S}=\KL\bigl(R\Vert\mathbb P^{\mathscr T}_\theta(\cdot\given X)\bigr).
\label{eq:link-sector}
\end{equation}
(c) If moreover $Q_{\mathring\omega}\ll\mathbb P^{\mathsf Y}(\cdot\given\mathring\omega)$ for $R$-almost every
$\mathring\omega$, $\int\KL(Q_{\mathring\omega}\Vert\mathbb P^{\mathsf Y}(\cdot\given\mathring\omega))R(d\mathring\omega)<\infty$,
and $\E_Q\lvert H_o\rvert<\infty$, then the free energy is finite and equals
\begin{equation}
\Fenergy_V[Q;X,o]=\KL\bigl(R\Vert\mathbb P^{\mathscr T}_\theta(\cdot\given X)\bigr)
+\int\KL\bigl(Q_{\mathring\omega}\Vert\mathbb P^{\mathsf Y}(\cdot\given\mathring\omega)\bigr)R(d\mathring\omega)
+\E_QH_o.
\label{eq:link-gaussian-energy-form}
\end{equation}
The condition on $H_o$ holds for a product law $Q=R\otimes Q^{\mathsf Y}$ when $\mathsf Y$ is a real vector
space of finite dimension, $Q^{\mathsf Y}$ has finite second moments, and every record satisfies
$\lvert E_{a,o}(X,\Theta(\mathring\omega),y_{\partial a})\rvert\leq\mathrm c_1(1+\lVert y\rVert^2)e^{\mathrm c_2\lVert\mathring\omega\rVert}$.
\status{ESTABLISHED}

\noindent\emph{Proof.}
(a) Both laws have densities with respect to Lebesgue measure on $\mathscr T$, and for a self-adjoint
operator $\mathbf B$, $\E_R\langle\mathring\omega-\mathring\alpha,\mathbf B(\mathring\omega-\mathring\alpha)\rangle
=\Tr(\mathbf B\mathring\Lambda^{-1})+\langle\mathring\mu-\mathring\alpha,\mathbf B(\mathring\mu-\mathring\alpha)\rangle$;
taking the $R$-expectation of the difference of the two log densities gives \eqref{eq:link-gaussian-kl}. In
the mean-field case all operators are block diagonal, so traces and log-determinants are sums over blocks.
The general identity is \cref{thm:elbo-total-correlation-chain} with the product law $\nu^{\mathscr T}$,
and the total correlation of a Gaussian law is the sum of the summand entropies minus the joint entropy,
in which the factors $2\pi e$ cancel because the summand dimensions add up to $\dim\mathscr T$. (b) The
first claim is \cref{lem:link-growth}(a), and \cref{lem:link-growth}(c) gives
$\E_R\mathscr S\leq\mathrm c_1\E_Re^{\mathrm c_2\lVert\mathring\omega\rVert}<\infty$. With (a), the hypotheses of
the last sentence of \cref{prop:link-latent-normalization}(c) hold, and \eqref{eq:link-sector} is
$\log(dR/d\mathbb P^{\mathscr T}_\theta)=\log(dR/d\nu^{\mathscr T})+\mathscr S+\log Z_{\mathscr S}$ integrated
against $R$. (c) The baseline of the latent space is
$\mathbb P^{\mathscr T}_\theta(d\mathring\omega\given X)\mathbb P^{\mathsf Y}(dy\given\mathring\omega)$, and
\cref{lem:prob-kl-chain-rule}(a) writes its divergence from $Q$ as the first two terms of
\eqref{eq:link-gaussian-energy-form}, which are finite, so $Q$ is absolutely continuous with respect to
it; the negative part of $H_o$ is $Q$-integrable, and \eqref{eq:obs-collective-energy-form} of
\cref{thm:elbo-extended-gap}, applied on the latent space $\mathscr T\times\mathsf Y$, gives
\eqref{eq:link-gaussian-energy-form}. For a product law,
$\E_Q\lvert E_{a,o}\rvert\leq\mathrm c_1\E_Re^{\mathrm c_2\lVert\mathring\omega\rVert}\E_{Q^{\mathsf Y}}(1+\lVert y\rVert^2)<\infty$
by (b) and Tonelli's theorem. \hfill$\square$

The relative entropy \eqref{eq:link-sector} is the whole link sector: a loop energy placed in the prior
enters with its normalizer, as \cref{rem:gen-appended-energy} requires of any reciprocal energy that is to
be part of an evidence bound, and in the mean-field case the sector is a sum over edge copies of Gaussian
divergences from the anchors plus the recognition expectation of the loop action.

\section{The link field equation}
\label{sec:link-field}

Which condition does the recognition law of the generators satisfy where the collective free energy is
stationary? For a Gaussian recognition law and a given matter recognition law, the stress of the anchor
plus the expected gradient of the loop action equals the expected matter current, filtered through the
adjoint derivative of the exponential chart (\cref{prop:link-field-equation}); for a comparison record the
current is a precision-weighted prediction error times the source mean, minus the transported source
covariance (\cref{cor:link-delta-rule}). Without the generator prior the stationary problem degenerates,
at the sample level and at the law level (\cref{prop:link-law-level-energy,cor:link-prior-necessary}). The
recognition flow that relaxes to these conditions is \cref{prop:link-recognition-flow}. Matrices are real,
$\langle\mathbf S,\mathbf T\rangle=\Tr(\mathbf S^\top\mathbf T)$ is the trace pairing of square matrices of
one size, gradients of functions of matrices refer to it, and for a representation $\rho_x$ on a sample
fiber written in reference coordinates, $(\mathrm d\rho_x)^\top$ is the adjoint of $\mathrm d\rho_x$ for
$\langle\cdot,\cdot\rangle_{\mathsf G}$ on $\mathfrak g$ and the trace pairing.

\lemmaheading{Derivative of the exponential chart}{lem:link-exp-derivative}
For square matrices $\mathbf B,\mathbf T$ of one size put
\begin{equation}
\mathbf L(\mathbf B,\mathbf T)=\int_0^1e^{t\mathbf B}\mathbf Te^{(1-t)\mathbf B}dt.
\label{eq:link-exp-derivative}
\end{equation}
(a) For all $\mathbf B,\mathbf T$, in the operator norm,
$\lVert e^{\mathbf B+\mathbf T}-e^{\mathbf B}-\mathbf L(\mathbf B,\mathbf T)\rVert\leq\frac12\lVert\mathbf T\rVert^2e^{\lVert\mathbf B\rVert+\lVert\mathbf T\rVert}$,
so the matrix exponential is Fr\'echet differentiable at $\mathbf B$ with derivative
$\mathbf T\mapsto\mathbf L(\mathbf B,\mathbf T)$, and
$\lVert\mathbf L(\mathbf B,\mathbf T)\rVert\leq e^{\lVert\mathbf B\rVert}\lVert\mathbf T\rVert$. (b) For every
$\mathbf S$,
\begin{equation}
\langle\mathbf S,\mathbf L(\mathbf B,\mathbf T)\rangle=\langle\mathbf L(\mathbf B^\top,\mathbf S),\mathbf T\rangle.
\label{eq:link-exp-adjoint}
\end{equation}
(c) If a function $\mathbf h$ is differentiable at $e^{\mathbf B}$ with gradient $\nabla\mathbf h(e^{\mathbf B})$,
then $\mathbf B\mapsto\mathbf h(e^{\mathbf B})$ has gradient $\mathbf L(\mathbf B^\top,\nabla\mathbf h(e^{\mathbf B}))$
at $\mathbf B$, and $\mathring\xi\mapsto\mathbf h(\rho_x(\exp\mathring\xi))$ has gradient
$(\mathrm d\rho_x)^\top\mathbf L\bigl((\mathrm d\rho_x\mathring\xi)^\top,\nabla\mathbf h(\rho_x(\exp\mathring\xi))\bigr)$ at
$\mathring\xi$ for $\langle\cdot,\cdot\rangle_{\mathsf G}$. \status{ESTABLISHED}

\noindent\emph{Proof.}
For square matrices $\mathbf B$ and $\mathbf B'$, Duhamel's formula
$e^{\mathbf B'}-e^{\mathbf B}=\int_0^1e^{t\mathbf B'}(\mathbf B'-\mathbf B)e^{(1-t)\mathbf B}dt$ holds because the
integrand is the $t$-derivative of $e^{t\mathbf B'}e^{(1-t)\mathbf B}$, and with
$\lVert e^{t\mathbf B}\rVert\leq e^{t\lVert\mathbf B\rVert}$ it gives
$\lVert e^{\mathbf B'}-e^{\mathbf B}\rVert\leq\lVert\mathbf B'-\mathbf B\rVert e^{\max(\lVert\mathbf B'\rVert,\lVert\mathbf B\rVert)}$.
With $\mathbf B'=\mathbf B+\mathbf T$,
$e^{\mathbf B+\mathbf T}-e^{\mathbf B}-\mathbf L(\mathbf B,\mathbf T)=\int_0^1(e^{t(\mathbf B+\mathbf T)}-e^{t\mathbf B})\mathbf Te^{(1-t)\mathbf B}dt$,
and the same bound applied to $t(\mathbf B+\mathbf T)$ and $t\mathbf B$ bounds the integrand by
$t\lVert\mathbf T\rVert^2e^{t(\lVert\mathbf B\rVert+\lVert\mathbf T\rVert)+(1-t)\lVert\mathbf B\rVert}\leq t\lVert\mathbf T\rVert^2e^{\lVert\mathbf B\rVert+\lVert\mathbf T\rVert}$,
whose integral gives (a); the bound on $\mathbf L$ follows from
$\lVert e^{t\mathbf B}\rVert\lVert e^{(1-t)\mathbf B}\rVert\leq e^{\lVert\mathbf B\rVert}$.
For (b), cyclicity of the trace and $(e^{t\mathbf B})^\top=e^{t\mathbf B^\top}$ give
$\Tr(\mathbf S^\top e^{t\mathbf B}\mathbf Te^{(1-t)\mathbf B})=\Tr\bigl((e^{t\mathbf B^\top}\mathbf Se^{(1-t)\mathbf B^\top})^\top\mathbf T\bigr)$,
and integration over $t$ gives \eqref{eq:link-exp-adjoint}. Part (c) is the chain rule with (a) and (b);
for the representation, $\rho_x(\exp\mathring\xi)=\exp(\mathrm d\rho_x\mathring\xi)$ (\cref{lem:link-exp-naturality}),
its derivative in the direction $\mathring\eta$ is $\mathbf L(\mathrm d\rho_x\mathring\xi,\mathrm d\rho_x\mathring\eta)$, and
\eqref{eq:link-exp-adjoint}
moves $\mathbf L$ onto the gradient before the adjoint $(\mathrm d\rho_x)^\top$ is applied. \hfill$\square$

The integral form of the Fr\'echet derivative of the matrix exponential is classical \citep[Ch.~10]{Higham2008}.

\lemmaheading{Gaussian expectations and their derivatives}{lem:link-gaussian-gradients}
Let $\mathbf h$ be a twice continuously differentiable function on $\mathscr T$, or on one of its summands,
with
$\lvert\mathbf h(\mathring\omega)\rvert+\lVert\nabla\mathbf h(\mathring\omega)\rVert+\lVert\nabla^2\mathbf h(\mathring\omega)\rVert\leq\mathrm c_1e^{\mathrm c_2\lVert\mathring\omega\rVert}$.
Then $(\mu,\Sigma)\mapsto\E_{\mathcal N(\mu,\Sigma)}\mathbf h$ is continuously differentiable on the product
of the space and the open cone of positive-definite self-adjoint operators, and for every self-adjoint
$\mathring\Delta$
\begin{equation}
\nabla_\mu\E_{\mathcal N(\mu,\Sigma)}\mathbf h=\E_{\mathcal N(\mu,\Sigma)}\nabla\mathbf h,
\qquad
\frac{d}{dt}\E_{\mathcal N(\mu,\Sigma+t\mathring\Delta)}\mathbf h\Big|_{t=0}=\tfrac12\Tr\bigl(\mathring\Delta\E_{\mathcal N(\mu,\Sigma)}\nabla^2\mathbf h\bigr);
\label{eq:link-bonnet-price}
\end{equation}
that is, for the trace pairing of self-adjoint operators the gradient in $\Sigma$ is one half of
$\E_{\mathcal N(\mu,\Sigma)}\nabla^2\mathbf h$. \status{ESTABLISHED}

The proof is in \cref{sec:proofs-links}.

The two identities are the Gaussian derivative identities used for variational Gaussian approximations
in \citet[Eqs.~(18)--(19)]{OpperArchambeau2009}.

\propositionheading{Link field equation}{prop:link-field-equation}
Assume \cref{def:link-latent-model} with bounded record densities and $p(o\given\mathring\omega)>0$ for every
$\mathring\omega$, a $\sigma$-finite reference $\nu^{\mathsf Y}$ with respect to which every
$\mathbb P^{\mathsf Y}(\cdot\given\mathring\omega)$ has a density, and a matter recognition law $Q^{\mathsf Y}$
with density $q^{\mathsf Y}$ and $\E_{Q^{\mathsf Y}}\lvert\log q^{\mathsf Y}\rvert<\infty$. Let the matter energy
$\mathbf H_o(\mathring\omega,y)=-\log p_{\theta,V}(o,y\given X,\Theta(\mathring\omega))$, the negative logarithm
of the density \eqref{eq:gen-record-density} of the held-link joint at $\Theta(\mathring\omega)$, be twice
continuously differentiable in $\mathring\omega$ for every $y$ with
\begin{equation}
\E_{Q^{\mathsf Y}}\sup_{\lVert\mathring\omega\rVert\leq\mathring r}\Bigl(\lvert\mathbf H_o\rvert+\lVert\nabla_{\mathring\omega}\mathbf H_o\rVert
+\lVert\nabla^2_{\mathring\omega}\mathbf H_o\rVert\Bigr)(\mathring\omega,Y)\leq\mathrm c_1e^{\mathrm c_2\mathring r}
\qquad\text{for every }\mathring r\geq0.
\label{eq:link-matter-domination}
\end{equation}
Restrict recognition to the laws $R\otimes Q^{\mathsf Y}$ with $R=\mathcal N(\mathring\mu,\mathring\Lambda^{-1})$
Gaussian on $\mathscr T$, general or mean-field over the summands, and put
$\mathbf U(\mathring\omega)=\mathscr S(\mathring\omega)+\E_{Q^{\mathsf Y}}\mathbf H_o(\mathring\omega,Y)$. Then:

(a) $\mathbf U$ is twice continuously differentiable, with
$\lvert\mathbf U\rvert+\lVert\nabla\mathbf U\rVert+\lVert\nabla^2\mathbf U\rVert\leq\mathrm c_3e^{\mathrm c_4\lVert\mathring\omega\rVert}$,
and the free energy
\begin{equation}
\Fenergy_V[R\otimes Q^{\mathsf Y};X,o]=\KL(R\Vert\nu^{\mathscr T})+\E_R\mathbf U+\log Z_{\mathscr S}
+\E_{Q^{\mathsf Y}}\log q^{\mathsf Y}
\label{eq:link-field-free-energy}
\end{equation}
is finite and continuously differentiable in $(\mathring\mu,\mathring\Lambda)$.

(b) A pair $(\mathring\mu,\mathring\Lambda)$ is stationary for \eqref{eq:link-field-free-energy} exactly when
\begin{gather}
\mathring\Sigma^{-1}(\mathring\mu-\mathring\alpha)+\E_R\nabla\mathscr S
=-\E_{R\otimes Q^{\mathsf Y}}\nabla_{\mathring\omega}\mathbf H_o,
\label{eq:link-field-mean}\\
\mathring\Lambda=\mathring\Sigma^{-1}+\E_R\nabla^2\mathscr S+\E_{R\otimes Q^{\mathsf Y}}\nabla^2_{\mathring\omega}\mathbf H_o,
\label{eq:link-field-precision}
\end{gather}
and in the mean-field case exactly when the summand components of \eqref{eq:link-field-mean} and the
diagonal blocks of \eqref{eq:link-field-precision} hold. When the agent latent law does not read the links,
$\nabla_{\mathring\omega}\mathbf H_o=\sum_a\nabla_{\mathring\omega}E_{a,o}(X,\Theta(\mathring\omega),y_{\partial a})$, the sum over the
records whose kernels read a link, so the right side of \eqref{eq:link-field-mean} is minus the expected
gradient of the incident record energies.

(c) If $\mathbf H_o$ reads $\mathring\omega^x_e$ only through the represented twist
$\rho_x(\mathring\Theta^x_e)=\exp(\mathrm d\rho_x\mathring\omega^x_e)$, as a differentiable function of this matrix on
an open set of matrices that contains the represented twists, and its gradient
$\partial\mathbf H_o/\partial\rho_x(\mathring\Theta^x_e)$ is $Q^{\mathsf Y}$-integrable at every $\mathring\omega$, then
with the link current
\begin{equation}
\mathbf J^x_e(\mathring\omega)=-\E_{Q^{\mathsf Y}}\frac{\partial\mathbf H_o}{\partial\rho_x(\mathring\Theta^x_e)}(\mathring\omega,Y)
\label{eq:link-current}
\end{equation}
the $(x,e)$ component of \eqref{eq:link-field-mean} reads
\begin{equation}
(\mathring\Sigma^x_e)^{-1}(\mathring\mu^x_e-\mathring\alpha^x_e)+\E_R\nabla_{(x,e)}\mathscr S
=\E_R\Bigl[(\mathrm d\rho_x)^\top\mathbf L\bigl((\mathrm d\rho_x\mathring\omega^x_e)^\top,\mathbf J^x_e(\mathring\omega)\bigr)\Bigr]:
\label{eq:link-field-current}
\end{equation}
the expected stress of the generator prior equals the expected matter current transported by the adjoint
derivative of the exponential chart.

(d) Put
$\mathbf V(\mathring\omega)=\frac12\langle\mathring\omega-\mathring\alpha,\mathring\Sigma^{-1}(\mathring\omega-\mathring\alpha)\rangle+\mathbf U(\mathring\omega)$.
Then \eqref{eq:link-field-mean} is $\E_R\nabla\mathbf V=0$, and
$\E_{\mathcal N(\mathring\mu,\mathring\Lambda^{-1})}\nabla\mathbf V\to\nabla\mathbf V(\mathring\mu)$ as
$\mathring\Lambda^{-1}\to0$, uniformly for $\mathring\mu$ in compact sets. Every limit point $\mathring\mu_*$, as
$\mathring\Lambda^{-1}\to0$, of solutions of \eqref{eq:link-field-mean} satisfies $\nabla\mathbf V(\mathring\mu_*)=0$;
and if $\nabla\mathbf V(\mathring\mu_*)=0$ with $\nabla^2\mathbf V(\mathring\mu_*)$ invertible, then for every
positive-definite $\mathring\Lambda_0$ there are $t_0>0$ and a continuously differentiable curve
$t\mapsto\mathring\mu(t)$ on $(-t_0,t_0)$ with $\mathring\mu(0)=\mathring\mu_*$ such that
$(\mathring\mu(t),t^{-2}\mathring\Lambda_0)$ solves \eqref{eq:link-field-mean} for $0<\lvert t\rvert<t_0$. When
$\mathbf U$ is bounded below, the law
$R^\star(d\mathring\omega)\propto e^{-\mathbf U(\mathring\omega)}\nu^{\mathscr T}(d\mathring\omega)$ is the
unique minimizer of $\Fenergy_V[R\otimes Q^{\mathsf Y};X,o]$ over all laws $R$ on $\mathscr T$, and $-\mathbf V$ is
its log density up to a constant, so the limit condition $\nabla\mathbf V=0$ is the stationarity condition
of the most probable link given the matter recognition.

(e) No point mass is admissible: $\KL(\delta_{\mathring\mu}\Vert\nu^{\mathscr T})=+\infty$, so the free
energy of $\delta_{\mathring\mu}\otimes Q^{\mathsf Y}$ is $+\infty$, and \eqref{eq:link-field-free-energy} tends
to $+\infty$ as $\mathring\Lambda^{-1}\to0$ at a given $\mathring\mu$. \status{ESTABLISHED}

The proof is in \cref{sec:proofs-links}.

At a point law the matter current is evaluated at one link value and the condition $\nabla\mathbf V=0$ is the
most-probable-link equation; part (e) types it as a generative quantity at given matter recognition, in
the sense of \cref{prop:prob-law-point-point-mass}, and not as a recognition optimum.

\corollaryheading{The current of a comparison record}{cor:link-delta-rule}
In the setting of \cref{prop:link-field-equation}(c), let $x=b$, let $e:j\to i$ lie in $E_+$, and write samples
and links in the reference gauge, in which every frame is the identity, reached by the passive
re-expression by the inverse frames, which changes no twist and hence no generator
(\cref{prop:geo-edge-twist-invariance}). Suppose that the only term
of $\mathbf H_o$ that reads $\mathring\omega^b_e$ is either (a) a binary success record $a$ with scope $\{i,j\}$,
observed as a success, with
$E_{a,o}=\frac12(k_i-\rho_b(\mathring\Theta^b_e)k_j)^\top W_a(k_i-\rho_b(\mathring\Theta^b_e)k_j)$ and
$W_a\succ0$, or (b) the comparison record \eqref{eq:obs-selective-attention-record} of row $i$ at $o_i=0$ in
the per-receiver source-label model of \cref{def:obs-attention-label-model}, with query map $S^Q_i$, key map
$S^K_e=S^K_{\rm ref}\rho_b(\mathring\Theta^b_e)$ as in \cref{def:obs-selective-comparison}, tempered weight
$\widetilde W_i=W_i/\tau_i$, the source labels counted among the matter coordinates, and a matter recognition
law of the factorization \eqref{eq:obs-attention-recognition-factorization} whose row $\beta^Q_i$ does not
depend on $(\mathring\omega,y)$. Let the belief marginals of $Q^{\mathsf Y}$ have finite second moments, with
means $\mu_i,\mu_j$, covariance $\Sigma_j$ of $k_j$, and cross covariance
$\Sigma_{ij}=\E_{Q^{\mathsf Y}}[(k_i-\mu_i)(k_j-\mu_j)^\top]$, which vanishes for a product law. Then the
current \eqref{eq:link-current} is, in the two cases,
\begin{gather}
\mathbf J^b_e(\mathring\omega)
=W_a\bigl[(\mu_i-\rho_b(\mathring\Theta^b_e)\mu_j)\mu_j^\top+\Sigma_{ij}-\rho_b(\mathring\Theta^b_e)\Sigma_j\bigr],
\label{eq:link-delta-current}\\
\mathbf J^b_e(\mathring\omega)
=\beta^Q_{ij}(S^K_{\rm ref})^\top\widetilde W_i\bigl[(S^Q_i\mu_i-S^K_{\rm ref}\rho_b(\mathring\Theta^b_e)\mu_j)\mu_j^\top
+S^Q_i\Sigma_{ij}-S^K_{\rm ref}\rho_b(\mathring\Theta^b_e)\Sigma_j\bigr],
\label{eq:link-delta-current-attention}
\end{gather}
and \eqref{eq:link-field-current} holds for $e$ with this current. For a product law in case (a) the
current is the precision-weighted prediction error of the receiver times the source mean, minus the
precision-weighted transported source covariance. \status{ESTABLISHED}

\noindent\emph{Proof.}
For a matrix $\mathbf M$ in place of $\rho_b(\mathring\Theta^b_e)$, the gradient of
$\frac12(k_i-\mathbf Mk_j)^\top W_a(k_i-\mathbf Mk_j)$ in $\mathbf M$ for the trace pairing is
$-W_a(k_i-\mathbf Mk_j)k_j^\top$, which is $Q^{\mathsf Y}$-integrable by the second moments, and
$\E_{Q^{\mathsf Y}}[(k_i-\mathbf Mk_j)k_j^\top]=\mu_i\mu_j^\top+\Sigma_{ij}-\mathbf M(\Sigma_j+\mu_j\mu_j^\top)$, which
gives \eqref{eq:link-delta-current}. In case (b), \cref{prop:obs-attention-elbo}(b) with the latent state
$(\mathring\omega,y)$ shows that row $i$ contributes $\tau_i^{-1}\sum_j\beta^Q_{ij}\E D_{ij}$ together with terms
that do not read the links, because $\mathrm c_i$ of \cref{prop:obs-selective-receiver-update} is constant; the
source $j$ reached through $e$ has $D_{ij}=V_e$ of \cref{def:obs-selective-comparison}, whose gradient in
$\mathbf M$ is $-(S^K_{\rm ref})^\top W_i(S^Q_ik_i-S^K_{\rm ref}\mathbf Mk_j)k_j^\top$, and the same expectation gives
\eqref{eq:link-delta-current-attention}. The last sentence is
\cref{prop:link-field-equation}(c). \hfill$\square$

By \cref{cor:link-delta-rule}, the source of the link field equation \eqref{eq:link-field-current} has the
outer-product form of the delta rule that \citet[Sec.~4.2, Eq.~(24)]{Schlag2021} write for fast-weight
programmers, a line of work they trace in their Sec.~2 to \citet{Schmidhuber1992}: a weighted
prediction error of the receiver times the source mean. Three features have no counterpart there and are
consequences of the joint: the transported source covariance, the transport of the current by the
adjoint derivative of the exponential chart, and the average over the recognition law of the generators,
balanced by the stress of the generator prior. The correspondence is one of form between update rules; no
learning architecture is identified with the model.

\propositionheading{Law-level link energies and their stabilizers}{prop:link-law-level-energy}
Let $q_i,q_j$ be laws in the law fiber $\mathcal B_x$ of a channel, let
$D_e(\Theta)=\KL(q_i\Vert\widehat\rho_x(\Theta)q_j)$ for $\Theta\in G$ be the transported divergence
\eqref{eq:full-graph-transported-divergences} of a copy $e:j\to i$ as a function of its link value, and let
$G_q=\{g\in G:\widehat\rho_x(g)q=q\}$ be the stabilizer of $q$, a closed subgroup when
$g\mapsto\widehat\rho_x(g)q$ is continuous for the weak topology on laws, with Lie algebra $\mathfrak g_q$; under
\cref{hyp:pb-regular-models} it is the isotropy algebra of \eqref{eq:pb-orbit-fisher-form}.

(a) $D_e(h\Theta h')=D_e(\Theta)$ for all $\Theta\in G$, $h\in G_{q_i}$, $h'\in G_{q_j}$. Hence the set of
minimizers of $D_e$, and the set of its stationary points when $D_e$ is finite and continuously
differentiable on $G$, are unions of double cosets $G_{q_i}\Theta G_{q_j}$, and a minimizer or stationary
point is not isolated when $\dim\mathfrak g_{q_i}>0$ or $\dim\mathfrak g_{q_j}>0$. At every generator at which
the exponential map of $\mathfrak g$ is a local diffeomorphism, a minimizer or stationary point of the chart
energy $\mathring\omega^x_e\mapsto D_e(U_i\exp(\mathring\omega^x_e)U_j^{-1})$ is likewise not isolated when
$\dim\mathfrak g_{q_i}>0$ or $\dim\mathfrak g_{q_j}>0$.

(b) Let $G=\GL^+(K)$ act on $\R^K$ by its defining representation, and let
$q=\mathcal N(\mu,\Sigma)$ with $\Sigma\succ0$. Then
$G_q=\Sigma^{1/2}\{O\in\mathrm{SO}(K):O\Sigma^{-1/2}\mu=\Sigma^{-1/2}\mu\}\Sigma^{-1/2}$, of dimension
$(K-1)(K-2)/2$ if $\mu\neq0$ and $K(K-1)/2$ if $\mu=0$. The Mahalanobis norm
$\lvert\mu\rvert_\Sigma=(\mu^\top\Sigma^{-1}\mu)^{1/2}$ satisfies
$\lvert g\mu\rvert_{g\Sigma g^\top}=\lvert\mu\rvert_\Sigma$ for every $g$, and for $K\geq2$ a link $\Theta_0$ with
$\widehat\rho_x(\Theta_0)q_j=q_i$ exists exactly when $\lvert\mu_i\rvert_{\Sigma_i}=\lvert\mu_j\rvert_{\Sigma_j}$; the zero
set of $D_e$ is then the double coset $G_{q_i}\Theta_0G_{q_j}=\Theta_0G_{q_j}$. For $K\geq3$ every minimizer and
stationary point of $D_e$ is non-isolated, and for $K=2$ this holds when $\mu_i=0$ or $\mu_j=0$.

(c) If $\widehat\rho_x(\Theta_0)q_j=q_i$ and \cref{hyp:pb-regular-models} holds in the channel, then $D_e$ is
$C^3$ near $\Theta_0$ with zero value and differential at $\Theta_0$, and its Hessian at $\Theta_0$ on the
tangent vector of $t\mapsto\exp(t\mathring\xi)\Theta_0$ is the orbit Fisher form $\mathfrak k_{q_i}(\mathring\xi,\mathring\xi)$ of
\cref{cor:pb-orbit-fisher-form}; it is positive semidefinite with kernel the tangent space of the zero set.
For Gaussian laws in (b),
\begin{equation}
\mathfrak k_q(\mathring\xi,\mathring\xi)=\bigl\lvert\Sigma^{-1/2}\mathring\xi\mu\bigr\rvert^2
+\tfrac12\bigl\lVert\Sigma^{-1/2}\mathring\xi\Sigma^{1/2}+(\Sigma^{-1/2}\mathring\xi\Sigma^{1/2})^\top\bigr\rVert_{\mathrm F}^2,
\qquad\mathring\xi\in\mathfrak{gl}(K,\R).
\label{eq:link-gaussian-orbit-form}
\end{equation}
\status{ESTABLISHED}

The proof is in \cref{sec:proofs-links}.

\corollaryheading{Without the generator prior the link problem degenerates}{cor:link-prior-necessary}
(a) Let $G=\GL^+(K)$, $K\geq2$, act on the belief fiber by its defining representation, and in the reference
gauge let the record of \cref{cor:link-delta-rule}(a) be the only factor reading the link of $e:j\to i$.
For a matter recognition law with finite second moments, $\Sigma_j\succ0$, and $\Sigma_{ij}=0$, let the link be
moved as a parameter without a generator prior, with expected record energy
$\mathbf l(\mathbf M)=\E_{Q^{\mathsf Y}}\frac12(k_i-\mathbf Mk_j)^\top W_a(k_i-\mathbf Mk_j)$ for real
$K\times K$ matrices $\mathbf M$. Then:
(i) $\mathbf l$ is a strictly convex quadratic whose unique stationary point and minimizer is
\begin{equation}
\mathbf M^\star=\mu_i\mu_j^\top(\Sigma_j+\mu_j\mu_j^\top)^{-1}
=\frac{\mu_i\mu_j^\top\Sigma_j^{-1}}{1+\mu_j^\top\Sigma_j^{-1}\mu_j},
\label{eq:link-rank-one}
\end{equation}
of rank at most one and therefore singular; with a cross covariance it is instead the second-moment
regression $\E[k_ik_j^\top]\E[k_jk_j^\top]^{-1}$, which can be invertible. (ii) $\mathbf l$ has no stationary
point on $\GL(K)$; its infimum over $\GL^+(K)$ equals $\mathbf l(\mathbf M^\star)$, is not attained, and every
minimizing sequence in $\GL^+(K)$ converges to $\mathbf M^\star$. (iii) A generator $\mathring\omega^b_e$ is
stationary for $\mathring\omega^b_e\mapsto\mathbf l(\exp\mathring\omega^b_e)$ only where
$\mathbf L(\mathring\omega^b_e,\cdot)$ is singular, and such generators exist: for $K=2$,
$\mu_i=\mu_j=\begin{psmallmatrix}1\\0\end{psmallmatrix}$, $\Sigma_j=W_a=I_2$, the generator
$\mathring\omega^b_e=-(\log3)I_2+2\pi\begin{psmallmatrix}0&-1\\1&0\end{psmallmatrix}$ is stationary, with link
$\exp\mathring\omega^b_e=\frac13I_2\neq\mathbf M^\star$. (iv) If $\mu_j^\top\Sigma_j^{-1}\mu_i\geq0$, then
$\mathbf M^\star$ is a limit of chart links $\exp\mathring\omega^b_e$, so the infimum over chart links equals
$\mathbf l(\mathbf M^\star)$, is not attained, and is approached only along links converging to
$\mathbf M^\star$; if $\mu_j^\top\Sigma_j^{-1}\mu_i<0$, no sequence of chart links converges to $\mathbf M^\star$, the
infimum over chart links exceeds $\mathbf l(\mathbf M^\star)$, and for $K=2$ it is not attained.

(b) If the free energy depends on a link value only through a law-point energy $D_e$ of
\cref{prop:link-law-level-energy}, finite and continuously differentiable on $G$, and the link is moved as a
parameter without a generator prior, its stationary link values form a union of double cosets, none
isolated when $\dim\mathfrak g_{q_i}+\dim\mathfrak g_{q_j}>0$; for nondegenerate Gaussian laws under $\GL^+(K)$,
where $D_e$ is finite and smooth, this holds for every pair of laws when $K\geq3$.

(c) With the generator prior: (i) under the hypotheses of \cref{prop:link-field-equation} with $\mathbf U$
bounded below, $\mathbf V$ attains its minimum on $\mathscr T$, and
$(\mathring\mu,\mathring\Lambda)\mapsto\Fenergy_V[\mathcal N(\mathring\mu,\mathring\Lambda^{-1})\otimes Q^{\mathsf Y};X,o]$ has
compact sublevel sets in the product of $\mathscr T$ and the open cone of positive-definite precisions, or of
block-diagonal ones in the mean-field case, and attains its minimum there; (ii) if $\mathbb P^{\mathsf Y}$ is a
nondegenerate Gaussian law on a real vector space $\mathsf Y$ of finite dimension that does not read the links,
and every record energy is nonnegative, lower semicontinuous in $(\mathring\omega,y)$, and bounded by
$\mathrm c_1(1+\lVert y\rVert^2)e^{\mathrm c_2\lVert\mathring\omega\rVert}$, then the free energy restricted to the product
Gaussian laws $\mathcal N(\mathring\mu,\mathring\Lambda^{-1})\otimes\mathcal N(\mu^{\mathsf Y},\Sigma^{\mathsf Y})$ has compact sublevel
sets in $(\mathring\mu,\mathring\Lambda,\mu^{\mathsf Y},\Sigma^{\mathsf Y})$ and attains its minimum.
\status{ESTABLISHED}

The proof is in \cref{sec:proofs-links}.

At the sample level the degeneracy of (a) is tied to a product matter recognition, since a correlated one
can make the stationary matrix invertible, while the coercivity of (c) comes from the anchor term of the
prior at every width $\mathring\Sigma$. Part (a)(iii) shows that the chart itself creates stationary generators
at its folds: a gradient flow in generator coordinates started at such a generator does not move, although
its link value is not stationary for the link energy.

\section{Held links as a limit, and conditional statements}
\label{sec:link-held}

How does the theory with latent generators relate to the theory in which links are structural data or
parameters? As the prior covariance shrinks at a held anchor, the infimum of the latent-link free energy
converges to the infimum of the held-link free energy at the anchor links, and near-minimizing link
recognition laws concentrate at the anchor (\cref{thm:link-held-limit}); at every width, the conditional
laws given the generators are held-link laws, so statements proved for every link value transfer
conditionally, while statements about mixtures over the generators need not
(\cref{prop:link-conditional-transfer}).

\theoremheading{Held links are the zero-width limit}{thm:link-held-limit}
Hold $\theta$ except for the prior covariances, which are $\mathring\Sigma=\mathring\sigma^2\mathring\Sigma_0$ for
$\mathring\sigma>0$ and a block-diagonal $\mathring\Sigma_0\succ0$, and hold $X$, $o$ and the anchor $\mathring\alpha$.
Write $\nu^{\mathscr T}_{\mathring\sigma}$, $\mathbb P^{\mathscr T}_{\mathring\sigma}$, $Z_{\mathscr S,\mathring\sigma}$ and
$\Fenergy_{\mathring\sigma}$ for the base law, the generator prior, its normalizer, and the free energy of the
latent-link model at width $\mathring\sigma$. Let $\mathcal Q'$ be a nonempty class of laws on $\mathsf Y$, put
$\mathbf F_Q(\mathring\omega)=\Fenergy_V[Q;X,\Theta(\mathring\omega),o]$ for $Q\in\mathcal Q'$ and
$\mathbf F(\mathring\omega)=\inf_{Q\in\mathcal Q'}\mathbf F_Q(\mathring\omega)$, and assume:
(i) the record densities are bounded, $L_o\leq\mathrm c_{\mathfrak A}$, and $p(o\given\mathring\omega)>0$ for every
$\mathring\omega$;
(ii) $\mathbf F$ is finite and lower semicontinuous at $\mathring\alpha$;
(iii) a subclass $\mathcal Q'_0\subseteq\mathcal Q'$ with $\inf_{Q\in\mathcal Q'_0}\mathbf F_Q(\mathring\alpha)=\mathbf F(\mathring\alpha)$
consists of laws for which $\mathbf F_Q$ is continuous at $\mathring\alpha$ and bounded by
$\mathrm c_Qe^{\mathrm c_Q\lVert\mathring\omega\rVert}$.
Then:

(a) As $\mathring\sigma\to0$, the infimum of $\Fenergy_{\mathring\sigma}$ over the product laws $R\otimes Q$ with
$Q\in\mathcal Q'$ and $R$ in any class of laws on $\mathscr T$ that contains $\nu^{\mathscr T}_{\mathring\sigma}$, and its
infimum over all laws $R(d\mathring\omega)Q_{\mathring\omega}(dy)$ whose kernel takes values in $\mathcal Q'$, both
converge to
\begin{equation}
\mathbf F(\mathring\alpha)=\inf_{Q\in\mathcal Q'}\Fenergy_V\bigl[Q;X,\Theta(\mathring\alpha),o\bigr],
\qquad
\Theta^x_e(\mathring\alpha)=U_i\exp(\mathring\alpha^x_e)U_j^{-1},
\label{eq:link-held-limit}
\end{equation}
the infimum over $\mathcal Q'$ of the held-link free energy at the anchor links.

(b) Under (i) alone, every family of laws $R_{\mathring\sigma}(d\mathring\omega)Q_{\mathring\sigma,\mathring\omega}(dy)$ with
$\sup_{\mathring\sigma}\Fenergy_{\mathring\sigma}<\infty$ satisfies
$R_{\mathring\sigma}\{\lVert\mathring\omega-\mathring\alpha\rVert\geq\mathring\epsilon\}\to0$ for every $\mathring\epsilon>0$; under (i)
and (ii),
every such family with kernels in $\mathcal Q'$ and $\Fenergy_{\mathring\sigma}\to\mathbf F(\mathring\alpha)$ satisfies in
addition $\KL(R_{\mathring\sigma}\Vert\mathbb P^{\mathscr T}_{\mathring\sigma})\to0$.

(c) Under (i) alone, if $p(o\given\cdot)$ is continuous at $\mathring\alpha$, then the evidence of the
latent-link model converges to $p_{\theta,V}(o\given X,\Theta(\mathring\alpha))$, and the link marginal
$\boldsymbol\Pi^{\mathscr T}_{\mathring\sigma}$ of its posterior satisfies
$\boldsymbol\Pi^{\mathscr T}_{\mathring\sigma}\{\lVert\mathring\omega-\mathring\alpha\rVert\geq\mathring\epsilon\}\to0$ for every
$\mathring\epsilon>0$.

(d) Under (i), hypothesis (ii) holds when $\mathcal Q'=\{Q_\lambda:\lambda\in\mathbb U\}$ for a subset $\mathbb U$ of a
Euclidean space, $(\mathring\omega,\lambda)\mapsto\mathbf F_{Q_\lambda}(\mathring\omega)$ is lower semicontinuous,
$\mathbf F(\mathring\alpha)<\infty$, and for some $\mathring\epsilon_0>0$ and some $\mathrm c>\mathbf F(\mathring\alpha)$ the set
of $\lambda$ with $\mathbf F_{Q_\lambda}(\mathring\omega)\leq\mathrm c$ for some
$\lVert\mathring\omega-\mathring\alpha\rVert<\mathring\epsilon_0$ has compact closure in $\mathbb U$. \status{ESTABLISHED}

The proof is in \cref{sec:proofs-links}.

The present held-link theory is therefore the zero-width limit of the latent one in value: the prior and
the recognition law narrow together, so the link sector costs nothing in the limit, in contrast with a
recognition law that narrows under a prior of constant width, whose cost diverges
(\cref{prop:link-field-equation}(e)).

\propositionheading{Conditional statements transfer; link-marginal statements need not}{prop:link-conditional-transfer}
Assume the hypotheses of \cref{prop:link-latent-normalization}(c).
(a) For every $\mathring\omega$ the conditional law of the posterior given $\mathring\omega$ is the held-link
posterior $\boldsymbol\Pi_{\theta,V,o,X,\Theta(\mathring\omega)}$, and for every law
$Q=R(d\mathring\omega)Q_{\mathring\omega}(dy)$ the conditional free energy
$-\log p(o\given\mathring\omega)+\KL(Q_{\mathring\omega}\Vert\boldsymbol\Pi_{\mathring\omega})$ is the held-link free energy
$\Fenergy_V[Q_{\mathring\omega};X,\Theta(\mathring\omega),o]$. If $Q'=R(d\mathring\omega)Q'_{\mathring\omega}(dy)$ has the same link
marginal and both free energies are finite, then
\begin{equation}
\Fenergy_V[Q';X,o]-\Fenergy_V[Q;X,o]
=\int\Bigl(\Fenergy_V\bigl[Q'_{\mathring\omega};X,\Theta(\mathring\omega),o\bigr]-\Fenergy_V\bigl[Q_{\mathring\omega};X,\Theta(\mathring\omega),o\bigr]\Bigr)R(d\mathring\omega).
\label{eq:link-transfer-potential}
\end{equation}
(b) Let $\mathbb O$ be a measurable set of link families, and consider a property of a pair consisting of
held-link structural data and a law on $\mathsf Y$, expressed through the held-link joint, evidence,
posterior and free energy. If the property holds for $((X,\Theta),Q'')$ for every $\Theta\in\mathbb O$ and
every $Q''$ in a class $\mathcal Q''$, then it holds for $((X,\Theta(\mathring\omega)),Q_{\mathring\omega})$ at every
$\mathring\omega$ with $\Theta(\mathring\omega)\in\mathbb O$ and $Q_{\mathring\omega}\in\mathcal Q''$, hence for $R$-almost
every $\mathring\omega\in\Theta^{-1}(\mathbb O)$ when the kernel lies in $\mathcal Q''$ $R$-almost surely. For instance, for
$G=\GL^+(K)$ in its defining representation and a closed walk $p$ in which the copies of some reversal pair
occur with nonzero net count, the generators at which $\lvert\det\rho_x(\mathrm{Hol}^x(p))\rvert=1$ form a
hyperplane of $\mathscr T$, so for every $R$ absolutely continuous with respect to Lebesgue measure the
conclusion of \cref{cor:cg-holonomy-noncompact-gaussian} for a block containing $p$ holds conditionally on
$\mathring\omega$ for $R$-almost every $\mathring\omega$.
(c) A property of laws on $\mathsf Y$ transfers from the conditional laws to the link-marginal laws, namely
the matter marginal $\int Q_{\mathring\omega}R(d\mathring\omega)$ of a recognition law, the matter marginal
$\int\boldsymbol\Pi_{\mathring\omega}\boldsymbol\Pi^{\mathscr T}(d\mathring\omega)$ of the posterior, and their images under a coarse
channel that does not read the links, when the set of laws with the property is preserved by mixtures, and
need not otherwise. For $G=\GL^+(1)$ acting on $\R$ by multiplication, two agents with
$k_j\sim\mathcal N(0,1)$ and $k_i\given k_j\sim\mathcal N(\rho_b(\Theta^b_e)k_j,1)$ for the copy $e:j\to i$, trivial
frames, $\mathscr S=0$, $\nu^{\mathscr T}=\mathcal N(0,1)$, and no record, every held-link posterior is Gaussian,
while the $k_i$-marginal of the latent-link posterior has
$\E k_i^4=3(1+2e^2+e^8)>3(1+e^2)^2=3(\E k_i^2)^2$, so it is not Gaussian. \status{ESTABLISHED}

\noindent\emph{Proof.}
(a) The first claim is the disintegration of \cref{prop:link-latent-normalization}(c), and the second is
\cref{def:elbo-extended} at $(X,\Theta(\mathring\omega))$. Finite free energies make both integrals in
\eqref{eq:link-decomposition} finite, because the integrands are bounded below and the divergence term is
common, so subtraction gives \eqref{eq:link-transfer-potential}. (b) By (a) the conditional objects at
$\mathring\omega$ are the held-link objects at $(X,\Theta(\mathring\omega))$, and kernels are determined up to
$R$-null sets. For the example, $\det e^{\mathbf B}=e^{\operatorname{tr}\mathbf B}$ (triangularize $\mathbf B$ over
$\mathbb C$), so by \eqref{eq:geo-edge-twist-holonomy}
$\det\rho_x(\mathrm{Hol}^x(p))$ is the exponential of the linear functional
$\mathring\omega\mapsto\sum_{e\in E_+}\operatorname{tr}\mathring\omega^x_e$ weighted by the net count of each $e$ in $p$, which is
nonzero by assumption, and the set where the determinant has modulus one is the kernel of this functional,
which is Lebesgue null. Off it the holonomy has $\lvert\det\rvert\neq1$, and
\cref{cor:cg-holonomy-noncompact-gaussian} applies at every such link value. (c) The marginals are mixtures
over $\mathring\omega$ of the conditional laws, and a link-independent kernel commutes with mixing, which gives
the positive clause. In the example there are no records, so the posterior is the prior; given the
generator, $(k_i,k_j)$ is Gaussian with $\E[k_i^2\given\mathring\omega]=1+e^{2\mathring\omega}$, and a centered Gaussian
has $\E k^4=3(\E k^2)^2$, so $\E k_i^4=3\E(1+e^{2\mathring\omega})^2=3(1+2e^2+e^8)$ and $\E k_i^2=1+e^2$, by
$\E e^{2\mathring\omega}=e^2$ and $\E e^{4\mathring\omega}=e^8$; strict inequality is strict Jensen, since
$e^{2\mathring\omega}$ is not almost surely constant. \hfill$\square$

A coarse channel that reads the links, such as the tree transports of
\cref{def:full-graph-coarse-links}, is a kernel on $(\mathring\omega,y)$ and is covered by (b) conditionally on
$\mathring\omega$; its link-marginal image is again a mixture.

\section{Directed and undirected matter}
\label{sec:link-induced}

Does matter generate a link action of its own? In the record extension it does not: the generator marginal
is the generator prior, and matter reaches the links only through the evidence. When generators and matter
are normalized jointly as an undirected model, the generator marginal acquires minus the logarithm of the
matter normalizer (\cref{prop:link-induced-action}).

\propositionheading{Induced link action}{prop:link-induced-action}
(a) For the joint \eqref{eq:link-latent-joint},
$\int_{\mathsf Y\times\mathsf O}\mathbb P_{\theta,V}(d\mathring\omega,dy,do\given X)=\mathbb P^{\mathscr T}_\theta(d\mathring\omega\given X)$,
and with record densities the link marginal of the posterior is proportional to
$p(o\given\mathring\omega)\mathbb P^{\mathscr T}_\theta(d\mathring\omega\given X)$.
(b) Let the potentials $\psi_i(\cdot;X,\Theta)$, $\psi_{ij}(\cdot;X,\Theta)$, $\psi_a(\cdot;X,\Theta)$ of
\cref{def:gen-undirected-model} depend jointly measurably on the link family, write
$\widetilde p_{X,\Theta}$ and $Z_{X,\Theta}$ for \eqref{eq:gen-gibbs-potential} and \eqref{eq:gen-gibbs-partition}
at $(X,\Theta)$, and normalize generators and matter jointly: with
\begin{equation}
\mathfrak Z=\int_{\mathscr T}e^{-\mathscr S(\mathring\omega)}Z_{X,\Theta(\mathring\omega)}\nu^{\mathscr T}(d\mathring\omega),
\label{eq:link-joint-undirected-normalizer}
\end{equation}
assume $0<\mathfrak Z<\infty$, and take the law with density
$\mathfrak Z^{-1}e^{-\mathscr S(\mathring\omega)}\widetilde p_{X,\Theta(\mathring\omega)}(o,y)$ with respect to
$\nu^{\mathscr T}\otimes\nu^{\mathsf Y}\otimes\nu^{\mathsf O}$. It is a probability law,
$Z_{X,\Theta(\mathring\omega)}<\infty$ for $\nu^{\mathscr T}$-almost every $\mathring\omega$, its generator marginal is
\begin{equation}
\mathfrak Z^{-1}e^{-\mathscr S(\mathring\omega)-\mathscr S^{\mathrm{ind}}(\mathring\omega)}\nu^{\mathscr T}(d\mathring\omega),
\qquad
\mathscr S^{\mathrm{ind}}(\mathring\omega)=-\log Z_{X,\Theta(\mathring\omega)},
\label{eq:link-induced-marginal}
\end{equation}
with $\mathscr S^{\mathrm{ind}}(\mathring\omega)\in(-\infty,+\infty]$ for $\nu^{\mathscr T}$-almost every
$\mathring\omega$, and on $\{0<Z_{X,\Theta(\mathring\omega)}<\infty\}$ its conditional law given $\mathring\omega$ is the undirected model
\eqref{eq:gen-gibbs-density} at $(X,\Theta(\mathring\omega))$. A record extension with densities is the case
$Z_{X,\Theta}=1$ for every $\Theta$, in which $\mathscr S^{\mathrm{ind}}=0$. The condition $\mathfrak Z<\infty$ is not
implied by $Z_{\mathscr S}<\infty$: for $G=\GL^+(1)$, trivial frames, $\mathscr S=0$, two scalar agents with
standard Gaussian self potentials, and the edge potential $\exp(\rho_b(\Theta_e)k_1k_2)$,
$Z_{X,\Theta}=(1-\rho_b(\Theta_e)^2)^{-1/2}$ when $\rho_b(\Theta_e)<1$ and $Z_{X,\Theta}=+\infty$ when
$\rho_b(\Theta_e)=e^{\mathring\omega_e}\geq1$, a set of positive $\nu^{\mathscr T}$-measure, so $\mathfrak Z=+\infty$.
\status{ESTABLISHED}

\noindent\emph{Proof.}
(a) By Tonelli's theorem, integrate the records, each a kernel of unit mass, and then $y$ against the
probability kernel $\mathbb P^{\mathsf Y}(\cdot\given\mathring\omega)$; the posterior marginal is that of
\cref{prop:link-latent-normalization}(c). (b) Tonelli's theorem gives
$\int e^{-\mathscr S}\widetilde p_{X,\Theta(\mathring\omega)}d(\nu^{\mathsf Y}\otimes\nu^{\mathsf O})=e^{-\mathscr S(\mathring\omega)}Z_{X,\Theta(\mathring\omega)}$,
whose $\nu^{\mathscr T}$-integral is $\mathfrak Z$, so the density integrates to one; finiteness of $\mathfrak Z$ with
$e^{-\mathscr S}>0$ forces $Z_{X,\Theta(\mathring\omega)}<\infty$ almost everywhere, the generator marginal is
$\mathfrak Z^{-1}e^{-\mathscr S}Z_{X,\Theta(\mathring\omega)}\nu^{\mathscr T}$, which is \eqref{eq:link-induced-marginal} with the
value $+\infty$ of $\mathscr S^{\mathrm{ind}}$ where $Z_{X,\Theta(\mathring\omega)}=0$, and dividing the joint density by
the marginal density gives \eqref{eq:gen-gibbs-density}. The densities of a record extension are
conditional densities of unit mass in $(y,o)$ for every link family (\cref{prop:gen-exact-normalization}),
so as potentials they have $Z_{X,\Theta}=1$. In the example,
$\iint(2\pi)^{-1}e^{-(k_1^2+k_2^2)/2+\rho_b(\Theta_e)k_1k_2}dk_1dk_2=\det\begin{psmallmatrix}1&-\rho_b(\Theta_e)\\-\rho_b(\Theta_e)&1\end{psmallmatrix}^{-1/2}$
for $\rho_b(\Theta_e)<1$, the integral diverges for $\rho_b(\Theta_e)\geq1$ by
\cref{prop:gen-gibbs-counterexample}, and the Gaussian law $\nu^{\mathscr T}$ on $\mathfrak g=\R$ charges
$\{\mathring\omega_e\geq0\}$. \hfill$\square$

For Gaussian self potentials and pairwise potentials of the connection-Laplacian form
\eqref{eq:cg-connection-laplacian-energy}, the matter normalizer is a Gaussian determinant, so the induced
action is one half the log-determinant of the connection Laplacian plus the self term, and its expansion in
closed walks produces loop products of the links; Supplement G develops both.

\section{Inference dynamics of the links}
\label{sec:link-inference-dynamics}

The generators are latent coordinates, so they move by the dynamics of every recognition coordinate
of the program, descent of the collective free energy. This section writes that descent for a
Gaussian recognition law of the generators jointly with the agents' recognition coordinates, shows
that it is a natural-gradient history whose explicit Euler step is the Gaussian Bayesian learning
rule (\cref{prop:link-recognition-flow}), gives the exact M-coordinate of the anchor
(\cref{prop:link-anchor-m-coordinate}), and reads the supplied link rate of
\cref{prop:obs-slow-model-link-dichotomy} in this setting (\cref{rem:link-rates}). In this section
and the next two, the links of one channel are latent, or both channels share them under the
common-link specialization \eqref{eq:geo-common-link}, and the channel index is dropped. The link
space $\mathscr T=\mathfrak g^{E_+}$ carries the product inner product
$\langle\mathring\omega,\mathring\omega'\rangle=\sum_{e\in E_+}\langle\mathring\omega_e,\mathring\omega'_e\rangle_{\mathsf G}$
of \cref{def:link-generator-prior}, its norm $\lVert\cdot\rVert$, and the Lebesgue measure of that definition;
self-adjointness, traces, and determinants of operators on $\mathscr T$ refer to this inner product.
Where a representation $\rho$ enters, $\mathrm d\rho$ is injective and $\mathfrak g$ carries the
pulled-back inner product of \cref{def:link-loop-defect}, so that
$\lVert\mathring\omega_e\rVert_{\mathsf G}=\lVert\mathrm d\rho\mathring\omega_e\rVert_{\mathsf G}$. For a
matrix $\mathbf H$ on the represented fiber,
$\mathbf H^*=(\mathsf G^{\mathrm{fib}})^{-1}\mathbf H^\top\mathsf G^{\mathrm{fib}}$ is its adjoint for the
fiber metric, so that $\langle\mathbf H,\mathbf H'\rangle_{\mathsf G}=\Tr(\mathbf H^*\mathbf H')$ in
\eqref{eq:link-trace-form} and left or right multiplication by $\mathbf H$ has as adjoint the left or
right multiplication by $\mathbf H^*$; $\operatorname{pr}_{\mathfrak g}$ is the
$\langle\cdot,\cdot\rangle_{\mathsf G}$-orthogonal projection of matrices onto
$\mathrm d\rho(\mathfrak g)$, read in $\mathfrak g$.

\propositionheading{Natural-gradient recognition flow of the generators}{prop:link-recognition-flow}
In the latent-link model of \cref{def:link-latent-model}, with bounded record densities and
$p(o\given\mathring\omega)>0$ for every $\mathring\omega$, let the recognition law be
$R\otimes Q_\lambda$. Here $R=\mathcal N(\mathring\mu,\mathring\Lambda^{-1})$ is a Gaussian law on
$\mathscr T$ with mean $\mathring\mu\in\mathscr T$ and precision $\mathring\Lambda$ in the open cone
$\mathring\Xi_{\succ0}$ of positive-definite elements of the space $\mathring\Xi$ of self-adjoint
operators that preserve each summand of a given orthogonal decomposition of $\mathscr T$: one summand
admits every positive-definite precision, the summands $\mathfrak g$ of the edge copies give the
products $\bigotimes_{e\in E_+}\mathcal N(\mathring\mu_e,\mathring\Lambda_e^{-1})$, and the lines of an
orthonormal basis give the diagonal precisions. The law $Q_\lambda$ is the agents' recognition law,
with coordinates $\lambda$ in a collective configuration manifold $\mathcal Q_V$
(\cref{def:hist-configuration-manifold}) that satisfies \cref{hyp:hist-standing-configuration-tier}
with metric $\mathsf G_V$. By \eqref{eq:link-decomposition} and
\cref{prop:link-gaussian-recognition-vfe}(b) with $Q_{\mathring\omega}=Q_\lambda$ for every
$\mathring\omega$, the collective free energy at $R\otimes Q_\lambda$ is
\begin{equation}
F(\mathring\mu,\mathring\Lambda,\lambda)
=\E_R\bigl[\mathfrak i_\lambda\bigr]+\tfrac12\log\det\mathring\Lambda+\log Z_{\mathscr S}
-\tfrac12\dim\mathscr T\log(2\pi e),
\label{eq:link-gaussian-free-energy}
\end{equation}
with the generator integrand
\begin{equation}
\mathfrak i_\lambda(\mathring\omega)
=\Fenergy_V\bigl[Q_\lambda;X,\Theta(\mathring\omega),o\bigr]+\mathscr S(\mathring\omega)
+\tfrac12\bigl\langle\mathring\omega-\mathring\alpha,\mathring\Sigma^{-1}(\mathring\omega-\mathring\alpha)\bigr\rangle
+\tfrac12\log\det\bigl(2\pi\mathring\Sigma\bigr).
\label{eq:link-generator-integrand}
\end{equation}
Assume that $F$ is $C^2$ on $\mathscr T\times\mathring\Xi_{\succ0}\times\mathcal Q_V$ and that each
$\mathfrak i_\lambda$ is $C^2$ with
$\lvert\mathfrak i_\lambda\rvert+\lVert\nabla\mathfrak i_\lambda\rVert+\lVert\nabla^2\mathfrak i_\lambda\rVert\leq\mathrm c_1e^{\mathrm c_2\lVert\mathring\omega\rVert}$
for constants $\mathrm c_1,\mathrm c_2$. Write $\mathring\mu\otimes\mathring\mu$ for the operator
$\mathring\omega\mapsto\langle\mathring\mu,\mathring\omega\rangle\mathring\mu$ and
$\operatorname{pr}_{\mathring\Xi}$ for the orthogonal projection onto $\mathring\Xi$ in the trace inner
product.
\begin{enumerate}
\item[(a)] The laws $\mathcal N(\mathring\mu,\mathring\Lambda^{-1})$ with
$\mathring\Lambda\in\mathring\Xi_{\succ0}$ form a minimal regular exponential family with natural
parameter $\mathring\vartheta=(\mathring\Lambda\mathring\mu,-\tfrac12\mathring\Lambda)$, expectation
parameter
$\mathring\varpi=(\mathring\mu,\mathring\Lambda^{-1}+\operatorname{pr}_{\mathring\Xi}(\mathring\mu\otimes\mathring\mu))$,
and log-partition function $\mathring\Gamma$, hence a regular statistical manifold. Its Fisher metric
is smooth and positive definite, and in the coordinates $(\mathring\mu,\mathring\Lambda)$ it is
\begin{equation}
g^F\bigl((\delta\mathring\mu,\delta\mathring\Lambda),(\delta\mathring\mu,\delta\mathring\Lambda)\bigr)
=\langle\delta\mathring\mu,\mathring\Lambda\delta\mathring\mu\rangle
+\tfrac12\Tr\bigl(\mathring\Lambda^{-1}\delta\mathring\Lambda\mathring\Lambda^{-1}\delta\mathring\Lambda\bigr).
\label{eq:link-gaussian-fisher}
\end{equation}
The product $\mathcal Q$ of this family with $\mathcal Q_V$, with the metric
$\mathsf G=g^F\oplus\mathsf G_V$, satisfies \cref{hyp:hist-standing-configuration-tier}.
\item[(b)] The natural-gradient field $-\operatorname{grad}^{\mathsf G}F$ of
\eqref{eq:hist-fisher-gradient-definition} on $\mathcal Q$ is $C^1$, so its maximal integral curves
exist and are unique. Along each of them $\tfrac{d}{dt}F=-\lVert\operatorname{grad}^{\mathsf G}F\rVert_{\mathsf G}^2\leq0$,
and each one through a noncritical point has nowhere-vanishing velocity and satisfies
\eqref{eq:hist-vfe-ray-representative} with $\mathsf a\equiv1$, so it represents a natural-gradient
history (\cref{def:hist-oriented-history}).
\item[(c)] The flow is of first order. Its velocity is a function of the configuration alone and
vanishes exactly at the critical points of $F$. It obeys the force balance
$\mathsf G(\dot Q,\cdot)=-dF$, in which $\mathsf G$ is the friction tensor, and
$\tfrac{d}{dt}F=-\mathsf G(\dot Q,\dot Q)$. The free energy strictly decreases along every
nonconstant solution, so no nonconstant solution is periodic. At a critical point $Q_\star$ with
invertible Hessian the linearized field is $-\mathsf G(Q_\star)^{-1}\nabla^2F(Q_\star)$, which is
diagonalizable with real eigenvalues, so no linearized mode oscillates.
\item[(d)] In the chart $(\mathring\vartheta,\lambda)$, the generator block of the explicit Euler step
of the flow with step $\Delta t>0$ is
\begin{equation}
\mathring\Lambda'=(1-\Delta t)\mathring\Lambda+\Delta t\operatorname{pr}_{\mathring\Xi}\E_R\bigl[\nabla^2\mathfrak i_\lambda\bigr],
\qquad
\mathring\mu'=\mathring\mu-\Delta t\mathring\Lambda'^{-1}\E_R\bigl[\nabla\mathfrak i_\lambda\bigr],
\label{eq:link-blr}
\end{equation}
whenever $\mathring\Lambda'$ is positive definite, for instance when $0<\Delta t<1$ and
$\operatorname{pr}_{\mathring\Xi}\E_R[\nabla^2\mathfrak i_\lambda]$ is positive semidefinite. For one summand
this is the Gaussian form of the Bayesian learning rule \citep[Eq.~(12)]{KhanRue2023} applied to
$\mathfrak i_\lambda$, and for the lines of an orthonormal basis it is its diagonal-precision form, which
\citet[Sec.~4.2, Eq.~(37)]{KhanRue2023} state with the expected gradient and Hessian replaced by their
values at the mean.
\item[(e)] Let $R'=\mathcal N(\mathring\mu',\mathring\Lambda'^{-1})$ with expectation parameter
$\mathring\varpi'$, and let $\widetilde F(\mathring\varpi)$ denote $F$ at fixed $\lambda$ in the
expectation chart. If $\mathring\Lambda'$ is positive definite and
$\nabla^2\widetilde F\preceq\mathrm c_{\mathsf G}\nabla^2\mathring\Gamma^\star$ on the segment from
$\mathring\varpi$ to $\mathring\varpi'$ for some $\mathrm c_{\mathsf G}>0$, where
$\mathring\Gamma^\star$ is the convex conjugate of $\mathring\Gamma$ and $\nabla^2\mathring\Gamma^\star$ is
the Fisher metric in the expectation chart, then
\begin{equation}
F(R')\leq F(R)-\Delta t^{-1}\KL(R\Vert R')-\bigl(\Delta t^{-1}-\mathrm c_{\mathsf G}\bigr)\KL(R'\Vert R),
\label{eq:link-blr-descent}
\end{equation}
so a step with $0<\Delta t\leq1/\mathrm c_{\mathsf G}$ decreases $F$ unless $R'=R$. For one summand
and a quadratic integrand
$\mathfrak i_\lambda(\mathring\omega)=\tfrac12\langle\mathring\omega,\mathbf Q\mathring\omega\rangle-\langle\mathbf b,\mathring\omega\rangle+\mathrm c_3$
with $\mathbf Q$ positive definite, the hypothesis holds with $\mathrm c_{\mathsf G}=1$ everywhere, and
the step $\Delta t=1$ gives $R'=\mathcal N(\mathbf Q^{-1}\mathbf b,\mathbf Q^{-1})$, the minimizer of $F$
over Gaussian $R$.
\end{enumerate}
\status{ESTABLISHED}

The proof is in \cref{sec:proofs-links}.

For links, the latent option of \cref{def:obs-selective-interface-learning} asks of a random relation
coordinate a prior that is a probability law, dynamics in the augmented state, and a regular configuration
metric or a justified quotient, and the present construction discharges these requirements. The generator
prior of \cref{def:link-generator-prior}, in the joint of \cref{def:link-latent-model}, supplies the first for
every Lie group $G$, compact or not (\cref{prop:link-generator-prior,prop:link-latent-normalization}), with
directed matter, and with undirected matter normalized jointly whenever the normalizer
\eqref{eq:link-joint-undirected-normalizer} is finite (\cref{prop:link-induced-action}); the dynamics and the
metric are those of \cref{prop:link-recognition-flow}.

\propositionheading{Exact M-coordinate of the anchor}{prop:link-anchor-m-coordinate}
Let the generator prior of \cref{def:link-generator-prior} have base law
$\nu^{\mathscr T}_{\mathring\alpha}=\bigotimes_{e\in E_+}\mathcal N(\mathring\alpha_e,\mathring\Sigma_e)$
and an action $\mathscr S\geq0$ that does not depend on $\mathring\alpha$, so that the prior is
$\mathbb P^{\mathscr T}_{\mathring\alpha}=Z_{\mathscr S}(\mathring\alpha)^{-1}e^{-\mathscr S}\nu^{\mathscr T}_{\mathring\alpha}$
with $Z_{\mathscr S}(\mathring\alpha)=\int e^{-\mathscr S}d\nu^{\mathscr T}_{\mathring\alpha}$, and let
no other factor of the joint depend on $\mathring\alpha$. Hold the recognition coupling
$\mathbb Q_{V,o,X}$ fixed, let its generator marginal $R$ have finite second moments, and assume the
finiteness hypotheses of \cref{prop:elbo-exact-m-coordinate} for every anchor.
\begin{enumerate}
\item[(a)] Up to terms that do not depend on $\mathring\alpha$, the objective of
\eqref{eq:elbo-m-coordinate} is
\begin{equation}
\Lelbo^{\mathrm{anc}}(\mathring\alpha)=-\tfrac12\E_R\bigl\langle\mathring\omega-\mathring\alpha,\mathring\Sigma^{-1}(\mathring\omega-\mathring\alpha)\bigr\rangle-\log Z_{\mathscr S}(\mathring\alpha),
\label{eq:link-anchor-objective}
\end{equation}
which contains the normalizer of the prior, and
\begin{equation}
\nabla_{\mathring\alpha_e}\Lelbo^{\mathrm{anc}}=\mathring\Sigma_e^{-1}\bigl(\E_R[\mathring\omega_e]-\E_{\mathbb P^{\mathscr T}_{\mathring\alpha}}[\mathring\omega_e]\bigr).
\label{eq:link-anchor-gradient}
\end{equation}
\item[(b)] The function $\Lelbo^{\mathrm{anc}}$ is strictly concave and has exactly one maximizer
$\mathring\alpha^{\mathrm{new}}$, the unique anchor at which the prior mean equals the recognition
mean, $\E_{\mathbb P^{\mathscr T}_{\mathring\alpha^{\mathrm{new}}}}[\mathring\omega]=\E_R[\mathring\omega]$. By
\cref{prop:elbo-exact-m-coordinate} it is the exact M-coordinate on the anchor at the given
covariances.
\item[(c)] If $\mathscr S=0$, as when every loop coefficient $\mathrm c_p$ vanishes, then
$\mathring\alpha^{\mathrm{new}}=\E_R[\mathring\omega]$, which is $\mathring\mu$ for a Gaussian
generator recognition law. If
$\mathscr S(\mathring\omega)=\tfrac12\sum_p\mathrm c_p\lVert\mathbf C_p\mathring\omega\rVert_{\mathsf G}^2$,
where $\mathbf C_p\mathring\omega=(\mathbf C\mathring\omega)_p$ is the loop sum of the cycle map
\eqref{eq:link-cycle-map} and $\mathbf C_p^\top$ is the adjoint of $\mathbf C_p:\mathscr T\to\mathfrak g$,
as for the loop defects of the translation instance in its affine representation
(\cref{prop:link-translation-instance}), then
\begin{equation}
\mathring\alpha^{\mathrm{new}}=\Bigl(I+\mathring\Sigma\sum_p\mathrm c_p\mathbf C_p^\top\mathbf C_p\Bigr)\E_R[\mathring\omega],
\label{eq:link-anchor-quadratic}
\end{equation}
which differs from $\E_R[\mathring\omega]$ exactly when $\mathbf C_p\E_R[\mathring\omega]\neq0$ for
some loop $p$ with $\mathrm c_p>0$.
\item[(d)] If the covariances $\mathring\Sigma_e$ range as well over the positive-definite
self-adjoint operators on $\mathfrak g$, the objective is strictly concave in the coordinates
$(\mathring\Sigma_e^{-1}\mathring\alpha_e,-\tfrac12\mathring\Sigma_e^{-1})_{e\in E_+}$, it has at most
one maximizer, and a maximizer matches the mean and the covariance of every edge generator under
the prior and under $R$. If $\mathscr S=0$ and every $\operatorname{Cov}_R(\mathring\omega_e)$ is
positive definite, the maximizer exists and is
$(\E_R[\mathring\omega_e],\operatorname{Cov}_R(\mathring\omega_e))_{e\in E_+}$. With a loop term it
need not exist: in the translation instance with $\mathsf G^{\mathrm{fib}}=I_2$, on one agent $i$
with one reversal pair of edge copies $e,\bar e:i\to i$, the loop $p=(e)$ with
$\mathrm c_p=\mathrm c>0$, and $R=\mathcal N(\mathring\mu,\mathrm v)$ with $\mathrm v\geq1/\mathrm c$,
the objective has no maximizer over $(\mathring\alpha_e,\mathring\Sigma_e)\in\R\times(0,\infty)$, and
its supremum is approached as $\mathring\Sigma_e\to\infty$.
\end{enumerate}
\status{ESTABLISHED}

The proof is in \cref{sec:proofs-links}.

\remarkheading{The supplied link rate in the latent reading}{rem:link-rates}
Read \cref{def:rg-formation-cycle-experiment} with every frame at the reference point, so that each
twist equals its link (\cref{prop:geo-edge-twist-invariance}).
(i) Its link translations, written $\check\omega_i$ in \cref{prop:obs-slow-model-link-dichotomy}, are
generators of the translation instance (\cref{prop:link-translation-instance}): $(\R,+)$ acting on the
scalar fibers by translation, in its affine representation with $\mathsf G^{\mathrm{fib}}=I_2$, for
which $\mathfrak g=\R$, $\langle\mathrm x,\mathrm y\rangle_{\mathsf G}=\mathrm x\mathrm y$, and every loop
defect is the quadratic $\tfrac12(\sum_a\check\omega_{e_a})^2$. The link block $\mathcal N(0,I_4)$ of the
configuration prior is the generator prior of \cref{def:link-generator-prior} with anchor $0$,
$\mathring\Sigma_e=1$, and no loop term; the loop energy of the experiment enters through its
records. For a law of the configuration, the generator
recognition law $R$ is its $\check\omega$-marginal, and its mean is the generator recognition mean
$\mathring\mu$.
(ii) Moving the link translations at the supplied rate $\varepsilon_{\check\omega}$ is therefore, on
the held matching sheet and for $\varepsilon_{\mathrm m},\varepsilon_{\check\omega}>0$, the
rate-separated flow \eqref{eq:obs-rate-separated-flow} with the generator recognition means in the
slow block, $s=(\bar m,\mathring\mu)$ and
$\mathsf G_{\mathrm s}=I_4\oplus(\varepsilon_{\mathrm m}/\varepsilon_{\check\omega})I_4$, by
\cref{prop:obs-slow-model-link-dichotomy}(i).
(iii) The anchor $\mathring\alpha$ is a coordinate of the generative parameter $\theta$, so it
belongs to the slow block $s=(s_{\mathrm m},\theta)$ of \cref{def:obs-fast-slow-blocks}: it moves the
joint between evaluations, by accepted M-coordinates (\cref{def:elbo-generalized-m-acceptance}), of
which \cref{prop:link-anchor-m-coordinate} gives the exact one, or by the rate-separated flow, and
never inside an evaluation (\cref{req:gen-typing-prohibition}). In the experiment no generative
parameter is learned and the anchor stays at $0$.
(iv) The recognition flow of \cref{prop:link-recognition-flow} does not supply the separation of
(ii). In the notation of \cref{def:rg-formation-cycle-experiment}, take the held sheet $S^{(0)}$, the
Gaussian configuration laws
$\mathcal N(\mu,\Sigma)\otimes\delta_{S^{(0)}}$ on $\R^{12}$ with their Fisher metric, of the form
\eqref{eq:link-gaussian-fisher}, and the free energy $\Fenergy_{\mathrm T}/\mathrm T$ of
\cref{prop:rg-formation-endogenous-cycle}, at the temperature $\mathrm T$ of that proposition. Then
$\nabla^2U_{S^{(0)}}\succeq\mathrm TI$, the natural-gradient flow leaves the slice
$\Sigma=\mathrm T(\nabla^2U_{S^{(0)}})^{-1}$ invariant, and on it $\mu(t)=e^{-t}\mu(0)$: belief means,
model means, and link means relax at one common rate. From
$\mu(0)=\mathbf x^\circ$, every residual is $r^b_i(t)=r^m_i(t)=-e^{-t}$, so at no time is
$\lvert r^b_i\rvert\leq2\varepsilon_{\mathrm m}<\tfrac12\leq\lvert r^m_i\rvert$, and the
belief-agreeing, model-disagreeing interval of \cref{prop:obs-slow-model-link-dichotomy}(ii) does not
occur. \status{ESTABLISHED}

The proof is in \cref{sec:proofs-links}.

On the slice of \cref{rem:link-rates}(iv) the Fisher geometry of the recognition family assigns every
coordinate the same rate, so there the separation in \cref{prop:obs-slow-model-link-dichotomy} comes
from its supplied mobilities alone, and a slow link block is a declaration of the kind that
\cref{def:obs-fast-slow-blocks} allows.

\section{Link histories across learning episodes}
\label{sec:link-histories}

Across learning episodes the generators form a history, and the latent route of
\cref{def:full-graph-episode-joint} gives that history one joint law. This section puts a Gibbs prior
on generator histories whose kinetic term is the loop defect of a time-like loop
(\cref{def:link-history-prior}), proves that it is a latent route with the generators in the parent
readout (\cref{prop:link-history-prior}), and shows that its increment precision is the inertial mass
of the most probable history (\cref{prop:link-map-inertia}). Throughout,
$\mathfrak Q(\mathbf H)=\tfrac12\lVert\mathbf H-I\rVert_{\mathsf G}^2$ for a matrix $\mathbf H$ on the
represented fiber, so that the loop defect \eqref{eq:link-loop-defect} of a closed walk $p$ is
$\mathfrak Q_p=\mathfrak Q(\rho(\mathring\Theta_p))$.

\definitionheading{Generator history prior}{def:link-history-prior}
Index learning episodes by $n=0,\ldots,n_1$ as in \cref{def:full-graph-episode-joint}, let
$\mathring\omega^{(n)}=(\mathring\omega^{(n)}_e)_{e\in E_+}\in\mathscr T$ be the generators of episode
$n$, with $\mathring\omega^{(n)}_{\bar e}=-\mathring\omega^{(n)}_e$, and write
$\mathring\omega^{(0:n_1)}$ for a history. Fix the generator prior of \cref{def:link-generator-prior},
with base law $\nu^{\mathscr T}=\bigotimes_{e\in E_+}\mathcal N(\mathring\alpha_e,\mathring\Sigma_e)$ and
action $\mathscr S=\sum_p\mathrm c_p\mathfrak Q_p$ (the statements below use only that $\mathscr S$ is
finite, measurable, and nonnegative), and a mass $\mathbf m>0$. Frames are static at $c_*$: the copies
of agent $i$ at episodes $n$ and $n+1$ are joined by the identity link. For an edge copy $e:j\to i$,
the time-like loop of $e$ between episodes $n$ and $n+1$ is the closed walk that runs from $j$ at
episode $n+1$ to $j$ at episode $n$, along $e$ at episode $n$, from $i$ at episode $n$ to $i$ at
episode $n+1$, and along $\bar e$ at episode $n+1$. Its reference holonomy
\eqref{eq:link-reference-holonomy}, with the identity twist for each time-like step, is
$\mathring\Theta^{(n+1\leftarrow n)}_e=\exp(-\mathring\omega^{(n+1)}_e)\exp(\mathring\omega^{(n)}_e)$, with
the represented holonomy
\begin{equation}
\rho\bigl(\mathring\Theta^{(n+1\leftarrow n)}_e\bigr)=e^{-\mathrm d\rho\mathring\omega^{(n+1)}_e}e^{\mathrm d\rho\mathring\omega^{(n)}_e}.
\label{eq:link-timelike-holonomy}
\end{equation}
The kinetic term of $e$ between episodes $n$ and $n+1$ is either the time-like loop defect
$\mathfrak U^{(n)}_e=\mathfrak Q(\rho(\mathring\Theta^{(n+1\leftarrow n)}_e))$ or the Gaussian increment
$\mathfrak U^{(n)}_e=\tfrac12\lVert\mathring\omega^{(n+1)}_e-\mathring\omega^{(n)}_e\rVert_{\mathsf G}^2$,
and $\mathfrak U(\mathring\omega^{(n+1)},\mathring\omega^{(n)})=\sum_{e\in E_+}\mathfrak U^{(n)}_e$. The
generator history prior is the law on $\mathscr T^{n_1+1}$
\begin{equation}
\mathbb P^{\mathrm{hist}}\bigl(d\mathring\omega^{(0:n_1)}\bigr)
=Z_{\mathrm{hist}}^{-1}\exp\Bigl[-\sum_{n=0}^{n_1}\mathscr S\bigl(\mathring\omega^{(n)}\bigr)
-\mathbf m\sum_{n=0}^{n_1-1}\mathfrak U\bigl(\mathring\omega^{(n+1)},\mathring\omega^{(n)}\bigr)\Bigr]
\prod_{n=0}^{n_1}\nu^{\mathscr T}\bigl(d\mathring\omega^{(n)}\bigr),
\label{eq:link-history-prior}
\end{equation}
with $Z_{\mathrm{hist}}$ the integral of the exponential against $\bigotimes_n\nu^{\mathscr T}$. The
anchor, the prior covariances, the loop coefficients, and the mass are coordinates of the generative
parameter $\theta$.

\propositionheading{The generator history prior}{prop:link-history-prior}
In \cref{def:link-history-prior}, with either kinetic term:
\begin{enumerate}
\item[(a)] $0<Z_{\mathrm{hist}}\leq Z_{\mathscr S}^{n_1+1}\leq1$, so
$\mathbb P^{\mathrm{hist}}$ is a probability law, absolutely continuous with respect to
$\bigotimes_n\nu^{\mathscr T}$ with density at most $Z_{\mathrm{hist}}^{-1}$; every function
integrable under $\bigotimes_n\nu^{\mathscr T}$, in particular every function bounded by
$\mathrm c_1e^{\mathrm c_2\lVert\mathring\omega^{(0:n_1)}\rVert}$, is integrable under it.
\item[(b)] With the backward functions $\mathring Z^{\leftarrow}_{n_1}\equiv1$ and
\begin{equation}
\mathring Z^{\leftarrow}_n(\mathring\omega)=\int\exp\bigl[-\mathscr S(\mathring\omega')-\mathbf m\mathfrak U(\mathring\omega',\mathring\omega)\bigr]\mathring Z^{\leftarrow}_{n+1}(\mathring\omega')\nu^{\mathscr T}(d\mathring\omega'),
\qquad n<n_1,
\label{eq:link-backward-functions}
\end{equation}
which are measurable with values in $(0,1]$, $\mathbb P^{\mathrm{hist}}$ is the law of the Markov chain
with root and transition kernels
\begin{equation}
\begin{gathered}
\mathbb P^{(0)}_{\mathrm{hist}}(d\mathring\omega)=Z_{\mathrm{hist}}^{-1}e^{-\mathscr S(\mathring\omega)}\mathring Z^{\leftarrow}_0(\mathring\omega)\nu^{\mathscr T}(d\mathring\omega),\\
\mathbb P^{(n+1\leftarrow n)}_{\mathrm{hist}}(d\mathring\omega'\given\mathring\omega)
=\frac{e^{-\mathscr S(\mathring\omega')-\mathbf m\mathfrak U(\mathring\omega',\mathring\omega)}\mathring Z^{\leftarrow}_{n+1}(\mathring\omega')}{\mathring Z^{\leftarrow}_n(\mathring\omega)}\nu^{\mathscr T}(d\mathring\omega').
\end{gathered}
\label{eq:link-history-kernels}
\end{equation}
These kernels are the conditional laws of the Gibbs law \eqref{eq:link-history-prior} and depend on
$n$ and on the horizon $n_1$ through $\mathring Z^{\leftarrow}_{n+1}$. The chain whose kernels normalize each
increment factor separately, started from the law proportional to $e^{-\mathscr S}\nu^{\mathscr T}$,
has the law \eqref{eq:link-history-prior} exactly when $\mathring Z^{\leftarrow}_{n_1-1}$ is constant
$\nu^{\mathscr T}$-almost everywhere, and this fails for one scalar generator with
$\nu^{\mathscr T}=\mathcal N(0,1)$, $\mathscr S=0$, the Gaussian increment, and $n_1=1$. With tower states
$\mathbf Y^{(n)}=(\mathring\omega^{(n)},y^{(n)})$ at depth zero, parent readout
$\chi^{\mathrm{ep}}(\mathbf Y^{(n)})=\mathring\omega^{(n)}$, the agents' latent law of each episode given
by a Markov kernel $\mathbb P^{\mathrm{ag}}_n(dy^{(n)}\given\mathring\omega^{(n)},X)$, and records of
episode $n$ that read $\mathbf Y^{(n)}$ only, the joint is a latent route of
\cref{def:full-graph-episode-joint}, with root
$\mathbb P^{(0)}_{\mathrm{hist}}(d\mathring\omega^{(0)})\mathbb P^{\mathrm{ag}}_0(dy^{(0)}\given\mathring\omega^{(0)},X)$ and transitions
$\mathbb P^{(n\leftarrow n-1)}_{\mathrm{hist}}(d\mathring\omega^{(n)}\given\mathring\omega^{(n-1)})\mathbb P^{\mathrm{ag}}_n(dy^{(n)}\given\mathring\omega^{(n)},X)$,
each a composition of kernels. Hence
\cref{prop:full-graph-episode-joint} applies: the joint has unit mass and is acyclic with forward
cross-episode edges, the root is required, the predictive prior of the generators at episode $n+1$ is
the filtered law of $\mathring\omega^{(n)}$ pushed through $\mathbb P^{(n+1\leftarrow n)}_{\mathrm{hist}}$, and, on the
finite-density domain with observation densities, the free energy decomposes along $n$ as in
\eqref{eq:full-graph-episode-decomposition}.
\item[(c)] The time-like links carry no variable. The twist of the identity link of agent $i$ between
two episodes is $e_G$, so $\mathbb P^{\mathrm{hist}}$ is a law on the spatial generator histories
alone, and each kinetic term is a function of the generators of one edge at two consecutive episodes.
In lattice gauge theory this configuration is the temporal gauge, in which a time-like plaquette
couples two space-like links at consecutive times \citep[Ch.~15, Eq.~(15.9)]{Creutz1983}.
\end{enumerate}
\status{ESTABLISHED}

The proof is in \cref{sec:proofs-links}.

\propositionheading{Inertia of the most probable history}{prop:link-map-inertia}
Let $\mathfrak V_n:\mathscr T\to\R$, $n=0,\ldots,n_1$, be differentiable, let $\mathbf m>0$ and
$\mathfrak U$ be either kinetic term of \cref{def:link-history-prior}, and let
\begin{equation}
\mathscr S^{\mathrm{hist}}\bigl(\mathring\omega^{(0:n_1)}\bigr)=\sum_{n=0}^{n_1}\mathfrak V_n\bigl(\mathring\omega^{(n)}\bigr)
+\mathbf m\sum_{n=0}^{n_1-1}\mathfrak U\bigl(\mathring\omega^{(n+1)},\mathring\omega^{(n)}\bigr).
\label{eq:link-history-action}
\end{equation}
For the history prior, $\mathfrak V_n=\mathscr S+\tfrac12\langle\cdot-\mathring\alpha,\mathring\Sigma^{-1}(\cdot-\mathring\alpha)\rangle$
makes $\mathscr S^{\mathrm{hist}}$ the negative log density of \eqref{eq:link-history-prior} relative to Lebesgue measure up
to a constant; adding to $\mathfrak V_n$ a matter term of episode $n$, such as the negative log
evidence of its records given the generators or their expected energy under a matter recognition law,
gives the corresponding posterior or conditional objective. A most probable history is a minimizer of
$\mathscr S^{\mathrm{hist}}$, with free ends or with ends held at given values.
\begin{enumerate}
\item[(a)] For the Gaussian increment, every critical point satisfies
\begin{equation}
\mathbf m\bigl(2\mathring\omega^{(n)}-\mathring\omega^{(n+1)}-\mathring\omega^{(n-1)}\bigr)+\nabla\mathfrak V_n\bigl(\mathring\omega^{(n)}\bigr)=0,
\qquad 0<n<n_1,
\label{eq:link-map-inertia}
\end{equation}
and, with free ends,
$\mathbf m(\mathring\omega^{(0)}-\mathring\omega^{(1)})+\nabla\mathfrak V_0(\mathring\omega^{(0)})=0$ and
$\mathbf m(\mathring\omega^{(n_1)}-\mathring\omega^{(n_1-1)})+\nabla\mathfrak V_{n_1}(\mathring\omega^{(n_1)})=0$.
With unit step in $n$, \eqref{eq:link-map-inertia} is the central-difference form of
$\mathbf m\ddot{\mathring\omega}=\nabla\mathfrak V$, Newton's equation with mass $\mathbf m$ in the
potential $-\mathfrak V$, the Euclidean sign. If each $\mathfrak V_n$ is continuous and
$\sum_n\mathfrak V_n(\mathring\omega^{(n)})\to\infty$ as
$\lVert\mathring\omega^{(0:n_1)}\rVert\to\infty$, as when each $\mathfrak V_n$ is at least the anchor
quadratic minus a constant, a most probable history exists.
\item[(b)] For the time-like loop defect, write
$\mathring\Theta^{\leftarrow}_e=\exp(-\mathring\omega'_e)\exp(\mathring\omega_e)$ for a later family
$\mathring\omega'$ and an earlier family $\mathring\omega$, so that $\rho(\mathring\Theta^{\leftarrow}_e)$ is
\eqref{eq:link-timelike-holonomy} with $\mathring\omega'$ and $\mathring\omega$ in place of
$\mathring\omega^{(n+1)}$ and $\mathring\omega^{(n)}$, and let $\mathbf L$ be the Fr\'echet derivative
\eqref{eq:link-exp-derivative}. The gradients of $\mathfrak U(\mathring\omega',\mathring\omega)$ in its
earlier and later arguments are
\begin{equation}
\begin{aligned}
\bigl(\partial_{\mathring\omega}\mathfrak U\bigr)_e&=\operatorname{pr}_{\mathfrak g}\mathbf L\bigl((\mathrm d\rho\mathring\omega_e)^*,(e^{-\mathrm d\rho\mathring\omega'_e})^*\bigl(\rho(\mathring\Theta^{\leftarrow}_e)-I\bigr)\bigr),\\
\bigl(\partial_{\mathring\omega'}\mathfrak U\bigr)_e&=-\operatorname{pr}_{\mathfrak g}\mathbf L\bigl(-(\mathrm d\rho\mathring\omega'_e)^*,\bigl(\rho(\mathring\Theta^{\leftarrow}_e)-I\bigr)(e^{\mathrm d\rho\mathring\omega_e})^*\bigr),
\end{aligned}
\label{eq:link-timelike-gradients}
\end{equation}
and every critical point satisfies, for $0<n<n_1$,
\begin{equation}
\mathbf m\bigl[\partial_{\mathring\omega}\mathfrak U\bigl(\mathring\omega^{(n+1)},\mathring\omega^{(n)}\bigr)+\partial_{\mathring\omega'}\mathfrak U\bigl(\mathring\omega^{(n)},\mathring\omega^{(n-1)}\bigr)\bigr]+\nabla\mathfrak V_n\bigl(\mathring\omega^{(n)}\bigr)=0.
\label{eq:link-map-inertia-defect}
\end{equation}
For $\mathrm r>0$ let $\mathrm c_{\mathrm r}=\sqrt2\sup\lVert D^3\mathfrak U(\mathring\omega',\mathring\omega)\rVert$ over
$\lVert\mathring\omega\rVert,\lVert\mathring\omega'\rVert\leq\mathrm r$. On that set
$\lVert\partial_{\mathring\omega}\mathfrak U-(\mathring\omega-\mathring\omega')\rVert$ and
$\lVert\partial_{\mathring\omega'}\mathfrak U-(\mathring\omega'-\mathring\omega)\rVert$ are at most
$\mathrm c_{\mathrm r}\mathrm r\lVert\mathring\omega'-\mathring\omega\rVert$, so at a history in the ball of radius $\mathrm r$
the left side of \eqref{eq:link-map-inertia-defect} differs from that of \eqref{eq:link-map-inertia}
by at most
$\mathbf m\mathrm c_{\mathrm r}\mathrm r(\lVert\mathring\omega^{(n+1)}-\mathring\omega^{(n)}\rVert+\lVert\mathring\omega^{(n)}-\mathring\omega^{(n-1)}\rVert)$.
At a given $\mathring\omega$ and small increments $\delta\mathring\omega$,
$\mathfrak U(\mathring\omega+\delta\mathring\omega,\mathring\omega)=\tfrac12\sum_e\lVert\mathbf L(-\mathrm d\rho\mathring\omega_e,\mathrm d\rho\delta\mathring\omega_e)e^{\mathrm d\rho\mathring\omega_e}\rVert_{\mathsf G}^2+O(\lVert\delta\mathring\omega\rVert^3)$,
so the time-like defect gives the increment the mass tensor
$\delta\mathring\omega\mapsto\mathbf m\sum_e\lVert\mathbf L(-\mathrm d\rho\mathring\omega_e,\mathrm d\rho\delta\mathring\omega_e)e^{\mathrm d\rho\mathring\omega_e}\rVert_{\mathsf G}^2$,
which depends on the generators and reduces to $\mathbf m\lVert\delta\mathring\omega\rVert^2$ at
$\mathring\omega=0$.
\item[(c)] Let $\mathscr T$ be one-dimensional, as for one reversal pair of edge copies of the
translation instance, identified with $\R$ isometrically, take the Gaussian increment, hold the ends
at $\mathring\omega^{(0)}=\mathrm a<\mathrm b=\mathring\omega^{(n_1)}$, and let $\mathfrak V_n=\mathfrak V$
for $0<n<n_1$ be continuous with $\mathfrak V\geq\mathfrak V(\mathrm a)$ on $(-\infty,\mathrm a]$ and
$\mathfrak V\geq\mathfrak V(\mathrm b)$ on $[\mathrm b,\infty)$, as for a double well with minima at
$\mathrm a$ and $\mathrm b$. Then a most probable history exists, and every most probable history is
nondecreasing in $n$; if moreover $\mathfrak V$ is differentiable with
$\mathfrak V'(\mathrm a)=\mathfrak V'(\mathrm b)=0$, every most probable history is strictly
increasing. The most probable path between the wells is a monotone kink and does not oscillate.
\end{enumerate}
\status{ESTABLISHED}

The proof is in \cref{sec:proofs-links}.

The inertia of \cref{prop:link-map-inertia} belongs to the smoothed history, a joint optimization over
$n$, while the recognition flow of \cref{prop:link-recognition-flow} is of first order in its
parameter. Recognition dynamics with velocity variables arise when the generative model itself
declares the motion of its states, as in the generalized coordinates of motion of
\citet[p.~852, Eq.~(13)]{Friston2008DEM}, whose ensemble dynamics couple each state to its motion;
the history prior declares no such variables. The Gaussian increment term of
\eqref{eq:link-history-action} is, up to a constant, the negative log transition density of a random
walk with increment covariance $\mathbf m^{-1}I$, the discrete-time form of the path weight of a
driftless diffusion.

\section{The linearized link equation}
\label{sec:link-linear}

Near zero generators the stationarity condition of the generator history becomes a linear equation
whose operator collects every term of the construction: the anchor, the loops, the episodes, and the
matter. Its positivity makes the linearized history unique, and the neglected terms are of second
order in the sources (\cref{prop:link-linearized-equation}).

\propositionheading{The linearized link equation}{prop:link-linearized-equation}
In the latent-link model of \cref{def:link-latent-model}, let the generators carry the history prior
of \cref{def:link-history-prior} with zero anchor, prior covariance $\mathring\Sigma$, action
$\mathscr S=\sum_p\mathrm c_p\mathfrak Q_p$, mass $\mathbf m$, and either kinetic term, and write samples
in the reference gauge of \cref{cor:link-delta-rule}. At each episode $n$ let $\mathfrak A_n$ be a
finite set of binary success records (\cref{hyp:local-interaction-kernels}), observed as successes;
record $a$ has a scope $\{i,j\}$ joined by an edge copy $e:j\to i$ of $E_+$ or of its reversal, with
$\mathring\omega_{\bar e}=-\mathring\omega_e$, and the energy of \cref{cor:link-delta-rule}(a),
\begin{equation}
E_{a,o}(k;\mathring\omega)=\tfrac12\bigl(k_i-e^{\mathrm d\rho\mathring\omega_e}k_j\bigr)^\top W_a\bigl(k_i-e^{\mathrm d\rho\mathring\omega_e}k_j\bigr),
\qquad W_a\succ0,
\label{eq:link-comparison-record}
\end{equation}
where $i$, $j$, and $e$ are those of $a$. Record $a$ enters the free energy of its episode with a
weight $\beta_a\in[0,1]$ that does not depend on the generators: $\beta_a=1$ for a record that is
always read, and a weight below one arises when a record enters through a recognized source label,
as in \cref{cor:link-delta-rule}(b). Hold a product Gaussian matter recognition under which $k_i$ at
episode $n$ has mean $\mu^{(n)}_i$ and covariance $\Sigma^{(n)}_i$, and let $\mathscr S^{\mathrm{hist}}$ be
the history action \eqref{eq:link-history-action} with
$\mathfrak V_n=\tfrac12\langle\cdot,\mathring\Sigma^{-1}\cdot\rangle+\mathscr S+\sum_{a\in\mathfrak A_n}\beta_a\E[E_{a,o}(k;\cdot)]$,
that is,
\begin{equation}
\begin{aligned}
\mathscr S^{\mathrm{hist}}(\mathring\omega)={}&\sum_{n=0}^{n_1}\Bigl[\tfrac12\bigl\langle\mathring\omega^{(n)},\mathring\Sigma^{-1}\mathring\omega^{(n)}\bigr\rangle+\sum_p\mathrm c_p\mathfrak Q_p\bigl(\mathring\omega^{(n)}\bigr)+\sum_{a\in\mathfrak A_n}\beta_a\E\bigl[E_{a,o}\bigl(k;\mathring\omega^{(n)}\bigr)\bigr]\Bigr]\\
&+\mathbf m\sum_{n=0}^{n_1-1}\mathfrak U\bigl(\mathring\omega^{(n+1)},\mathring\omega^{(n)}\bigr),
\end{aligned}
\label{eq:link-history-objective}
\end{equation}
which, when the agents' latent law does not read the links, is the generator part of the expected
negative log joint density. Its minimizers are the most
probable generator histories given the matter recognition, in the sense of
\cref{prop:link-map-inertia}. For one episode with every $\beta_a=1$, and an agent latent law that
does not read the links, $\mathscr S^{\mathrm{hist}}$ is the function $\mathbf V$ of
\cref{prop:link-field-equation}(d) up to a constant, so $\nabla\mathscr S^{\mathrm{hist}}=0$ is the history
form of the most-probable-link condition $\nabla\mathbf V=0$. On $\mathscr T^{n_1+1}$, with
$\langle\mathring\omega,\mathring\omega'\rangle=\sum_n\langle\mathring\omega^{(n)},\mathring\omega'^{(n)}\rangle$,
define self-adjoint operators $\mathbf K$ and $\mathbf N_0$ and a source $\mathbf j$ by
\begin{equation}
\begin{aligned}
\langle\mathring\omega,\mathbf K\mathring\omega\rangle
={}&\sum_n\Bigl[\bigl\langle\mathring\omega^{(n)},\mathring\Sigma^{-1}\mathring\omega^{(n)}\bigr\rangle+\sum_p\mathrm c_p\bigl\lVert\mathbf C_p\mathring\omega^{(n)}\bigr\rVert_{\mathsf G}^2
+\sum_{a\in\mathfrak A_n}\beta_a\Tr\bigl(W_a\mathrm d\rho\mathring\omega^{(n)}_e\E\bigl[k_jk_j^\top\bigr](\mathrm d\rho\mathring\omega^{(n)}_e)^\top\bigr)\Bigr]\\
&+\mathbf m\sum_{n<n_1}\bigl\lVert\mathring\omega^{(n+1)}-\mathring\omega^{(n)}\bigr\rVert^2,\\
\langle\mathring\omega,\mathbf N_0\mathring\omega\rangle
={}&-\sum_n\sum_{a\in\mathfrak A_n}\beta_a\Tr\bigl(\mathbf J_a^\top(\mathrm d\rho\mathring\omega^{(n)}_e)^2\bigr),
\qquad
\langle\mathbf j,\mathring\omega\rangle=\sum_n\sum_{a\in\mathfrak A_n}\beta_a\Tr\bigl(\mathbf J_a^\top\mathrm d\rho\mathring\omega^{(n)}_e\bigr),
\end{aligned}
\label{eq:link-linear-operator}
\end{equation}
where $e$ and $j$ are those of $a$, $\E[k_jk_j^\top]=\Sigma^{(n)}_j+\mu^{(n)}_j\mu^{(n)\top}_j$,
$\mathbf C_p\mathring\omega=(\mathbf C\mathring\omega)_p$ is the loop sum of the cycle map
\eqref{eq:link-cycle-map}, and, for $a\in\mathfrak A_n$,
\begin{equation}
\mathbf J_a=W_a\bigl[(\mu^{(n)}_i-\mu^{(n)}_j)\mu^{(n)\top}_j-\Sigma^{(n)}_j\bigr]
\label{eq:link-zero-current}
\end{equation}
is the current \eqref{eq:link-delta-current} of record $a$ at zero generators for a product law. Thus
$\mathbf K=\mathring\Sigma^{-1}+\sum_p\mathrm c_p\mathbf C_p^\top\mathbf C_p+\mathbf m(-\Delta_n)+\mathring\Phi$,
where $\mathring\Sigma^{-1}$ and the loop term act episode by episode,
$(-\Delta_n\mathring\omega)^{(n)}=2\mathring\omega^{(n)}-\mathring\omega^{(n+1)}-\mathring\omega^{(n-1)}$ in the
interior, $\mathring\omega^{(0)}-\mathring\omega^{(1)}$ and $\mathring\omega^{(n_1)}-\mathring\omega^{(n_1-1)}$ at
the ends, and $\mathring\Phi$ is the matter operator of the last term of the bracket. Let
$\mathrm c_{\mathring\Sigma}$ be the smallest eigenvalue of $\mathring\Sigma^{-1}$, and let
$\operatorname{cond}(\mathsf G^{\mathrm{fib}})$ be the ratio of the extreme eigenvalues of
$\mathsf G^{\mathrm{fib}}$.
\begin{enumerate}
\item[(a)] $\nabla\mathscr S^{\mathrm{hist}}(0)=-\mathbf j$ and $\nabla^2\mathscr S^{\mathrm{hist}}(0)=\mathbf K+\mathbf N_0$, so to first order in
the generators and the currents the stationarity condition is
\begin{equation}
\mathbf K\mathring\omega=\mathbf j.
\label{eq:link-linearized-equation}
\end{equation}
\item[(b)] $\mathbf K\succeq\mathrm c_{\mathring\Sigma}I$. Hence \eqref{eq:link-linearized-equation}
has the unique solution $\mathring\omega_{\mathrm{lin}}=\mathbf K^{-1}\mathbf j$, with
$\lVert\mathring\omega_{\mathrm{lin}}\rVert\leq\lVert\mathbf j\rVert/\mathrm c_{\mathring\Sigma}$, and
$\lVert\mathbf K^{-1}\mathbf j-\mathbf K^{-1}\mathbf j'\rVert\leq\lVert\mathbf j-\mathbf j'\rVert/\mathrm c_{\mathring\Sigma}$.
\item[(c)] Let $\mathrm c_0=\lVert\mathbf N_0\rVert$ and, for $\mathrm r>0$,
$\mathrm c_{\mathrm r}=\sup_{\lVert\mathring\omega\rVert\leq\mathrm r}\lVert D^3\mathscr S^{\mathrm{hist}}(\mathring\omega)\rVert$. Then
$\mathrm c_0\leq\operatorname{cond}(\mathsf G^{\mathrm{fib}})\max_{n,e\in E_+}\mathrm j^{(n)}_e$ and
$\lVert\mathbf j\rVert\leq\operatorname{cond}(\mathsf G^{\mathrm{fib}})^{1/2}(\sum_{n,e\in E_+}(\mathrm j^{(n)}_e)^2)^{1/2}$,
where $\mathrm j^{(n)}_e$ is the sum of $\beta_a\lVert\mathbf J_a\rVert_{\mathrm F}$ over the records
$a\in\mathfrak A_n$ on $e$ or $\bar e$. If $\mathrm c_0+\mathrm c_{\mathrm r}\mathrm r\leq\mathrm c_{\mathring\Sigma}/2$ and
$\lVert\mathbf j\rVert\leq\mathrm c_{\mathring\Sigma}\mathrm r/4$, then $\mathscr S^{\mathrm{hist}}$ has exactly one critical point
$\mathring\omega_\star$ in the ball of radius $\mathrm r$ about $0$; it is the unique minimizer of $\mathscr S^{\mathrm{hist}}$ on
that ball, $\lVert\mathring\omega_\star\rVert\leq2\lVert\mathbf j\rVert/\mathrm c_{\mathring\Sigma}$, and
\begin{equation}
\lVert\mathring\omega_\star-\mathring\omega_{\mathrm{lin}}\rVert
\leq\frac{2\lVert\mathbf j\rVert}{\mathrm c_{\mathring\Sigma}^2}\Bigl(\mathrm c_0+\frac{\mathrm c_{\mathrm r}\lVert\mathbf j\rVert}{\mathrm c_{\mathring\Sigma}}\Bigr),
\label{eq:link-linearization-error}
\end{equation}
which, at fixed $\mathrm c_{\mathrm r}$, is of second order in the currents.
\end{enumerate}
\status{ESTABLISHED}

The proof is in \cref{sec:proofs-links}.

In the language of lattice gauge theory, \cref{prop:link-linearized-equation} reads as follows, and
the reading adds nothing to it. The spatial term $\sum_p\mathrm c_p\mathbf C_p^\top\mathbf C_p$ is a
discrete curl-curl operator: $\mathbf C_p$ sums the generators around the loop $p$, so the term
penalizes the circulation of the generator field on the cycle space of the multigraph, the magnetic
part of a lattice action. The episode term $\mathbf m(-\Delta_n)$ has the quadratic form
$\mathbf m\sum_n\lVert\mathring\omega^{(n+1)}-\mathring\omega^{(n)}\rVert^2$ of a forward difference in the
episode index; by \cref{prop:link-history-prior}(c) and \cref{prop:link-map-inertia}(b) it is the
second-order part of the time-like loop defect with identity time-like links, the temporal gauge, in
which the time-like plaquettes give the electric term of the Hamiltonian of
\citet{KogutSusskind1975}, as derived in \citet[Ch.~15, Eq.~(15.41)]{Creutz1983}. The anchor
precision $\mathring\Sigma^{-1}$ is a bare mass and the matter operator $\mathring\Phi$ a state-dependent
screening mass, so at linear order $\mathbf K$ has the form of a massive Euclidean vector-field
operator of Proca type, a mass term plus curl-curl plus a squared episode difference, acting on the
spatial components alone, since the history carries no time-like variable. Its curl-curl and episode
terms act as scalar graph operators on each component of $\mathfrak g$ in an orthonormal basis, and
the bracket of $\mathfrak g$ enters only through the terms of order three and higher, which
\cref{prop:link-linearized-equation}(c) bounds. The episode index is an
inference index in the sense of \cref{sec:intro-conventions}, and the signature is Euclidean
because $\mathbf K$ is positive (\cref{prop:link-linearized-equation}(b)).

The loop term sees only part of the link space, and the declared family $\mathbb L$ fixes which part.
The next statement makes this exact for the quadratic loop action of
\cref{prop:link-anchor-m-coordinate}(c) and for the loop term of $\mathbf K$.

\corollaryheading{Cycle-space splitting of the loop sector}{cor:link-cycle-splitting}
Let $\mathfrak I=(V,E,\bar{\ })$ be the interaction multigraph of \cref{def:geo-graph-links}, with
$\kappa$ connected components, write $\mathscr T=\R^{E_+}\otimes\mathfrak g$, let
$\mathbb L_+=\{p\in\mathbb L:\mathrm c_p>0\}$, and let $\mathbf C_+:\mathscr T\to\mathfrak g^{\mathbb L_+}$ be the
cycle map \eqref{eq:link-cycle-map} of $\mathbb L_+$. Let $d:\R^V\to\R^{E_+}$ be $(da)_e=a_i-a_j$ for
$e:j\to i$, let $Z=(\operatorname{im}d)^\perp$ be the cycle space, and for $p\in\mathbb L_+$ let
$z_p\in\R^{E_+}$ count the traversals of each $e\in E_+$ along its orientation minus those against it,
so that $(\mathbf C\mathring\omega)_p=\sum_{e\in E_+}z_{p,e}\mathring\omega_e$. Put
$r_{\mathbb L}=\dim\operatorname{span}\{z_p:p\in\mathbb L_+\}\leq\lvert\mathbb L_+\rvert$, and call
\begin{equation}
\mathscr T_{\mathrm{ex}}=\operatorname{im}d\otimes\mathfrak g,
\qquad
\mathscr T_{\mathrm{har}}=\bigl(Z\cap\{z_p:p\in\mathbb L_+\}^\perp\bigr)\otimes\mathfrak g
\label{eq:link-exact-harmonic}
\end{equation}
the exact and the harmonic generators; $\mathring\omega$ is exact when $\mathring\omega_e=a_i-a_j$ for some
$a\in\mathfrak g^V$.
\begin{enumerate}
\item[(a)] $\ker\mathbf C_+=\mathscr T_{\mathrm{ex}}\oplus\mathscr T_{\mathrm{har}}$, a sum orthogonal for the
inner product of $\mathscr T$, with
\begin{equation}
\dim\mathscr T_{\mathrm{ex}}=\dim\mathfrak g\,(\lvert V\rvert-\kappa),
\qquad
\dim\mathscr T_{\mathrm{har}}=\dim\mathfrak g\,(\lvert E_+\rvert-\lvert V\rvert+\kappa-r_{\mathbb L})
\geq\dim\mathfrak g\,(\lvert E_+\rvert-\lvert V\rvert+\kappa-\lvert\mathbb L_+\rvert).
\label{eq:link-cycle-splitting-dims}
\end{equation}
In the translation instance (\cref{prop:link-translation-instance}) an exact generator shift by
$a_i-a_j$ is the link change $\Theta_e\mapsto\Theta_e+a_i-a_j$ with the frames held, and every loop
defect $\mathfrak Q_p$ with $p\in\mathbb L_+$ is constant along $\ker\mathbf C_+$.
\item[(b)] Under the quadratic action of \cref{prop:link-anchor-m-coordinate}(c), with
$\mathring\Sigma=\bigoplus_{e\in E_+}\mathring\Sigma_e$,
\begin{equation}
\mathring\alpha^{\mathrm{new}}-\E_R[\mathring\omega]\in\mathring\Sigma\operatorname{im}\mathbf C_+^\top,
\qquad
\mathscr T=\mathscr T_{\mathrm{ex}}\oplus\mathscr T_{\mathrm{har}}\oplus\mathring\Sigma\operatorname{im}\mathbf C_+^\top,
\label{eq:link-anchor-coexact}
\end{equation}
and the last summand is orthogonal to the first two for
$\langle\cdot,\mathring\Sigma^{-1}\cdot\rangle$. Hence $\mathring\alpha^{\mathrm{new}}$ and $\E_R[\mathring\omega]$ have
the same components in $\mathscr T_{\mathrm{ex}}$ and in $\mathscr T_{\mathrm{har}}$, and the anchor
update corrects only the component in $\mathring\Sigma\operatorname{im}\mathbf C_+^\top$.
\item[(c)] For the operator $\mathbf K$ of \cref{prop:link-linearized-equation}, if every
$\mathring\omega^{(n)}$ lies in $\ker\mathbf C_+$, the loop term vanishes and
$\langle\mathring\omega,\mathbf K\mathring\omega\rangle=\sum_n\langle\mathring\omega^{(n)},\mathring\Sigma^{-1}\mathring\omega^{(n)}\rangle
+\mathbf m\sum_{n<n_1}\lVert\mathring\omega^{(n+1)}-\mathring\omega^{(n)}\rVert^2+\langle\mathring\omega,\mathring\Phi\mathring\omega\rangle$.
\end{enumerate}
\status{ESTABLISHED}

\noindent\emph{Proof.}
(a) The kernel of $d$ is the functions constant on components, so
$\dim\operatorname{im}d=\lvert V\rvert-\kappa$ and $\dim Z=\lvert E_+\rvert-\lvert V\rvert+\kappa$. Along a
closed walk every arrival at a vertex is followed by a departure, so the net traversals at each vertex
cancel, $d^\top z_p=0$, and $z_p\in Z$. For the matrix $\mathsf Z_+$ with rows $z_p$, $p\in\mathbb L_+$,
this gives $\ker\mathsf Z_+=(\operatorname{span}\{z_p\})^\perp=\operatorname{im}d\oplus\bigl(Z\cap\{z_p\}^\perp\bigr)$,
orthogonally, and the second summand has dimension $\dim Z-r_{\mathbb L}$ because
$\operatorname{span}\{z_p\}\subset Z$. By \eqref{eq:link-cycle-map} and the reversal convention
\eqref{eq:link-generator-reversal}, $\mathbf C_+=\mathsf Z_+\otimes\mathrm{id}_{\mathfrak g}$, and the inner product
of $\mathscr T$ is the Euclidean one on $\R^{E_+}$ times $\langle\cdot,\cdot\rangle_{\mathsf G}$, so tensoring with
$\mathfrak g$ carries the decomposition, its orthogonality, and the dimensions over. In the translation
instance $\mathring\omega_e=\Theta_e-U_i+U_j$ by \cref{prop:link-translation-instance}(a), so with the frames
held an exact shift of $\mathring\omega$ is the stated shift of $\Theta$, and by
\eqref{eq:link-translation-defect} $\mathfrak Q_p=\tfrac12\lVert(\mathbf C\mathring\omega)_p\rVert_{\mathsf G}^2$
depends on $\mathring\omega$ only through $\mathbf C_+\mathring\omega$ when $p\in\mathbb L_+$.
(b) Terms with $\mathrm c_p=0$ vanish in \eqref{eq:link-anchor-quadratic}, so
$\mathring\alpha^{\mathrm{new}}-\E_R[\mathring\omega]=\mathring\Sigma\sum_{p\in\mathbb L_+}\mathrm c_p\mathbf C_p^\top\mathbf C_p\E_R[\mathring\omega]\in\mathring\Sigma\operatorname{im}\mathbf C_+^\top$.
Each $\mathring\Sigma_e$ is self-adjoint and positive, so for $h\in\ker\mathbf C_+$ and
$v\in\operatorname{im}\mathbf C_+^\top$, $\langle\mathring\Sigma v,\mathring\Sigma^{-1}h\rangle=\langle v,h\rangle=0$.
The subspace $\mathring\Sigma\operatorname{im}\mathbf C_+^\top$ therefore meets $\ker\mathbf C_+$ only in zero, and its
dimension $\operatorname{rank}\mathbf C_+=\dim\mathscr T-\dim\ker\mathbf C_+$ makes the sum all of $\mathscr T$; with
(a) this gives \eqref{eq:link-anchor-coexact}, and the components of the difference in the first two
summands vanish. (c) $\mathbf C_p\mathring\omega^{(n)}=0$ for $p\in\mathbb L_+$, and the other loop terms carry
$\mathrm c_p=0$, so only the remaining terms of \eqref{eq:link-linear-operator} survive. $\square$

Attaching a two-cell along each walk of $\mathbb L_+$ makes the declared loop family a two-complex on
the agents, and \cref{cor:link-cycle-splitting} is the Hodge decomposition of its one-cochains with
values in $\mathfrak g$ \citep[Theorem~3.1]{HansenGhrist2019}, applied to the constant sheaf: exact
generators are its gradient part, $\mathring\Sigma\operatorname{im}\mathbf C_+^\top$ its circulating part in the
metric of the prior, and $\mathscr T_{\mathrm{har}}$ the circulation around cycles of $\mathfrak I$ that bound
no declared walk. By \eqref{eq:link-cycle-splitting-dims}, a family of short loops on a multigraph whose
cycle space is mostly spanned by longer cycles leaves most of that space to the anchor, the episodes,
and the matter, so the choice of $\mathbb L$ is a structural declaration on the agent set of the same
kind as the support of the source prior. Away from zero generators the loop defects of a nonabelian
group are not functions of $\mathbf C\mathring\omega$, and the splitting describes only their quadratic part
(\cref{prop:link-loop-defect}(e)).

\section*{Physicist's summary}

The links between agents are treated here as a field that is inferred, not as a background that is held.
Each link is the exponential of a Lie-algebra generator measured in the frame of a reference point;
per-agent frame changes do not touch the generators, and a change of reference point rotates all of them
by one adjoint map. A Gaussian law on the generators, reweighted by a Wilson-type loop penalty, is a
normalized prior for any structure group, so no finite Haar measure is needed; on the general linear
group the loop penalty is measured with a fiber metric, because conjugation-invariant penalties cannot
see unipotent holonomy. The links enter the free energy through one relative entropy from this prior.
Stationarity is a field equation: the restoring force of the prior balances a matter current, which for
a comparison record is a precision-weighted prediction error times the source state, the form of a delta
rule, transported by the derivative of the exponential map. Shrinking the prior to its anchor recovers
the free energy and the evidence of the held links at the anchor. Integrating out matter that is normalized jointly with the links generates an induced
link action, as in induced lattice gauge theory, while matter attached by record extension generates
none. Inference of the generators is first-order gradient flow in the Fisher geometry of their Gaussian
law, with the Bayesian learning rule as its Euler step; inertia appears only in the most probable
history across learning episodes, where the precision of the increment acts as a mass. Linearized about
zero generators, the stationarity operator is a bare mass from the anchor, a curl-curl term from the
loops, a squared episode difference, and a screening mass from the matter, which is the form of a
massive Euclidean vector operator in the episode index. On the cycle space the loop term acts only on the circulating part of the
generators: exact generators and circulation around cycles that bound no declared loop are fixed by the
anchor, the episodes, and the matter alone, and the anchor update corrects only the circulating part.


\part{Exact Coarse-Graining at One Base Point}
\chapter{Coarse Channels}
\label{ch:coarsegraining}

Part II asks what a coarse description of the agent network at $c_*$ keeps of the joint of
\Cref{ch:generative}, of its posterior, and of a recognition law. For one coarse channel the answer
is \Cref{thm:rg-exact-coarse-vfe}: a channel that pushes all three laws preserves the evidence and
the posterior version, and splits the free energy into the coarse free energy plus the expected
relative entropy of the discarded conditional laws. \Cref{thm:rg-pointwise-parent-datum}
specializes it to a candidate parent of a block of agents, and the exact results of
\Cref{ch:network-state,ch:agent-network-rg,ch:full-graph-meta-agent-vfe,ch:finite-depth-formation}
rest on both. The independent-partition selector and its perturbation boundaries are
\cref{prop:cg-finest-independent-partition,prop:cg-independent-selector-instability,prop:cg-partition-information-stability}; their remaining stable selection question is
\cref{open:cg-natural-selector}. The chapter also separates Markov pushforward, precomposition, and Galerkin
restriction (\Cref{sec:cg-three-operations}) and closes with finite labeled families, in which a
retained label carries information of order $N^2$ at a tie among $N$ agents
(\Cref{sec:cg-retained-label-asymptotics}). Fisher contraction under a channel is
\Cref{thm:cg-fisher-contraction} and is not restated. The longer proofs of the chapter are
deferred to \cref{app:deferred-proofs}.

\paragraph{Scope.} The common-channel identity concerns one channel applied to all three laws with
the observation unchanged. It does not compare different generative and recognition channels, a
coarse joint fitted rather than pushed forward, or a channel that also coarsens the observation, even
an invertible or evidence-sufficient one, because each changes the comparison target or the evidence
event; nor is a Fisher tensor of a labeled family used as the metric of a dynamics.
\status{NOT-CLAIMED} A stable intrinsic selector for interacting laws is
\Cref{open:cg-natural-selector}, and equivariant Bayes versions beyond the finite and compact cases
are \Cref{open:cg-equivariant-versions}.

\section{Markov coarse maps}
\label{sec:cg-markov-category}

\definitionheading{Coarse channel}{def:cg-coarse-channel}
\label{eq:cg-markov-pushforward}
A coarse channel between standard Borel spaces $\mathsf X$ and $\mathsf Y$ is a Markov kernel
$K:\mathsf X\rightsquigarrow\mathsf Y$ in the sense of \Cref{sec:intro-conventions}. It sends a law
$P$ on $\mathsf X$ to $(PK)(B)=\int_{\mathsf X}K(x,B)P(dx)$ and a statistical experiment
$(\mathsf X,\{P_\theta\}_{\theta\in\Theta})$ to $(\mathsf Y,\{P_\theta K\}_{\theta\in\Theta})$. A
measurable statistic $c:\mathsf X\to\mathsf Y$ is the deterministic kernel
$K_c(x,B)=\mathbf 1_B(c(x))$. A family $K_\theta$ that varies with the parameter is a
parameter-dependent model transformation, not one channel between experiments.

\propositionheading{Markov kernels form a category}{prop:cg-markov-category}
\label{eq:cg-kernel-composition}
If $K:\mathsf X\rightsquigarrow\mathsf Y$ and $L:\mathsf Y\rightsquigarrow\mathsf Z$ are Markov
kernels, then $(KL)(x,D)=\int_{\mathsf Y}L(y,D)K(x,dy)$ is a Markov kernel, composition is
associative with the deterministic identity as unit, and
\begin{equation}
P(KL)=(PK)L
\label{eq:cg-law-functor}
\end{equation}
for every law $P$ on $\mathsf X$. \status{ESTABLISHED}

\noindent\emph{Proof.}
Kernel integration gives measurability in $x$, and monotone convergence gives countable additivity
in $D$. By Tonelli's theorem $[P(KL)](D)=\int_{\mathsf X}\int_{\mathsf Y}L(y,D)K(x,dy)P(dx)$, which
equals $\int_{\mathsf Y}L(y,D)(PK)(dy)$; this is \eqref{eq:cg-law-functor}, and applied to point
masses and a third kernel it proves associativity. These are the morphisms and the composition of
the Kleisli category of the Giry probability monad \citep{Giry1982}. $\square$

\section{KL contraction, equality, and recovery}
\label{sec:cg-kl-recovery}

Every Markov kernel contracts relative entropy in $[0,+\infty]$ (\Cref{thm:cg-kl-dpi-extended});
the equality case needs finiteness.

\theoremheading{Data processing with its exact equality condition}{thm:cg-dpi-equality}
\label{eq:cg-kl-dpi}
\label{eq:cg-kl-equality}
Let $P,Q$ be laws on $\mathsf X$ with $P\ll Q$ and $\KL(P\Vert Q)<\infty$, let
$K:\mathsf X\rightsquigarrow\mathsf Y$ be a Markov kernel, and put $\mathbb J(dx,dy)=Q(dx)K(x,dy)$,
$r=dP/dQ$, and $\bar r(y)=\E_{\mathbb J}[r(X)\given Y=y]$. Then $d(PK)/d(QK)=\bar r$ and
$\KL(PK\Vert QK)\leq\KL(P\Vert Q)$, with equality exactly when $r(X)=\bar r(Y)$
$\mathbb J$-almost surely. \status{ESTABLISHED}

\noindent\emph{Proof.}
For bounded measurable $f$ on $\mathsf Y$,
$\int f d(PK)=\E_{\mathbb J}[f(Y)r(X)]=\E_{\mathbb J}[f(Y)\bar r(Y)]$, which is the
Radon-Nikodym formula. Conditional Jensen for the strictly convex function $t\mapsto t\log t$
gives $\E_{\mathbb J}[\bar r(Y)\log\bar r(Y)]\leq\E_{\mathbb J}[r(X)\log r(X)]$, and finite divergence
and strict convexity make equality equivalent to $r(X)=\bar r(Y)$ almost surely. This is the
information-loss theorem of \citet{Kullback1951,Csiszar1967} in kernel form. $\square$

\corollaryheading{Pairwise Bayes recovery}{cor:cg-pairwise-bayes-recovery}
In the setting of \Cref{thm:cg-dpi-equality}, let $R_Q(y,dx)$ be a regular conditional law of $X$
given $Y=y$ under $\mathbb J$. Then $QKR_Q=Q$, and at equality in the data-processing inequality
also $PKR_Q=P$. Conversely, any one Markov kernel $R$ that recovers both $P$ and $Q$ forces that
equality. For a dominated experiment $\{P_\theta\}$ with one common dominating law $P_0$,
simultaneous finite equalities $\KL(P_\theta K\Vert P_0K)=\KL(P_\theta\Vert P_0)$ for every $\theta$
make the single kernel $R_{P_0}$ recover every $P_\theta$. \status{ESTABLISHED}

\noindent\emph{Proof.}
Disintegration gives $QKR_Q=Q$. At equality,
$(PKR_Q)(A)=\int\bar r(y)R_Q(y,A)(QK)(dy)=\iint\mathbf 1_A(x)r(x)\mathbb J(dx,dy)=P(A)$. The converse
follows by data processing through $K$ and then through $R$, and the family statement is the pairwise
statement with $Q=P_0$ for each $\theta$. This is the commutative case of the channel-sufficiency
theorem of \citet{Petz1988}, whose reference-dependent reverse map is here the Bayes kernel $R_Q$.
$\square$

\propositionheading{Common recovery by likelihood descent}{prop:cg-descent-common-recovery}
Let $K:\mathsf X\rightsquigarrow\mathsf Y$ be a coarse channel, let $\{P_\theta\}_{\theta\in\Theta}$ be
laws on $\mathsf X$ dominated by a law $P_0$, with likelihood ratios $r_\theta=dP_\theta/dP_0$, and let
$R_{P_0}$ be a regular conditional law of $X$ given $Y$ under $\mathbb J_{P_0}(dx,dy)=P_0(dx)K(x,dy)$;
no relative entropy is assumed finite. The density $r=dP/dP_0$ of a law $P\ll P_0$ descends when it
agrees $\mathbb J_{P_0}$-almost everywhere with a measurable function of the coarse coordinate, which is
then $\bar r=d(PK)/d(P_0K)$:
\begin{equation}
r(x)=\bar r(y)\qquad\text{for }\mathbb J_{P_0}\text{-almost every }(x,y).
\label{eq:cg-likelihood-descent}
\end{equation}
(i) One Markov kernel $R:\mathsf Y\rightsquigarrow\mathsf X$ recovers $P_0$ and every $P_\theta$, that
is, $P_0KR=P_0$ and $P_\theta KR=P_\theta$, exactly when every $r_\theta$ descends, and $R_{P_0}$ is then
such a kernel. The laws $P\ll P_0$ that $R_{P_0}$ recovers are those whose density descends; when
every $r_\theta$ descends, they include every law whose density lies in the $L^1(P_0)$-closed linear
span of the constant $1$ and the $r_\theta$. For the records of \Cref{def:rg-measure-pair}, with $P_0=\varrho$, the family
$\{\bar{\mathsf m}_o\}$ indexed by records $o$, and $K=C$, the likelihood ratios are $L_o/M_o$ and
their pushed densities are $L^c_o/M_o$, so one kernel recovers the reference and every
$\bar{\mathsf m}_o$ exactly when each record likelihood $L_o$ agrees with its effective likelihood
$L^c_o$ almost everywhere under the reference joint. (ii) For the deterministic kernel $K_c$ of a
statistic $c:\mathsf X\to\mathsf Y$, descent reads $r=\bar r\circ c$ $P_0$-almost everywhere, and a
Markov kernel $R$ supported on fibers, $R(y,c^{-1}(\{y\}))=1$ for $c_\#P_0$-almost every $y$, recovers
$P_0$ exactly when it is a version of $R_{P_0}$. A kernel supported on fibers therefore recovers $P_0$
and every $P_\theta$ exactly when it is a version of $R_{P_0}$ and every $r_\theta$ descends, and it is
then unique up to a $c_\#P_0$-null set. (iii) Let $\mathsf X=\mathsf X_1\times\mathsf X_2$ with
coordinates $X_1,X_2$, let $c:\mathsf X_1\to\mathsf Y$ be a statistic, and let the channel be the
deterministic kernel of $c\times\operatorname{id}_{\mathsf X_2}:\mathsf X\to\mathsf Y\times\mathsf X_2$,
which coarsens $X_1$ and retains $X_2$. A kernel $R((y,x_2),\cdot)=R'(y,\cdot)\otimes\delta_{x_2}$, which
does not read the retained coordinate, with $R':\mathsf Y\rightsquigarrow\mathsf X_1$ satisfying
$R'(y,c^{-1}(\{y\}))=1$ for $(c(X_1))_\#P_0$-almost every $y$, recovers $P_0$ and every $P_\theta$ exactly when $R'$ is a regular conditional law of
$X_1$ given $c(X_1)$ under $P_0$, $X_1$ and $X_2$ are conditionally independent given $c(X_1)$ under
$P_0$, and every $r_\theta$ descends. Descent does not imply the conditional independence: for
$\mathsf X_1=\mathsf X_2=\{0,1\}$, $P_0$ uniform on $\{(0,0),(1,1)\}$, the family $\{P_0\}$, and
$\mathsf Y$ a point, $R_{P_0}$ recovers $P_0$, while every kernel that does not read $X_2$ sends $P_0$ to
a law giving mass $1/2$ to the $P_0$-null set $\{x_1\neq x_2\}$. \status{ESTABLISHED}

The proof is in \cref{sec:proofs-coarse-channels}.

Part (i) replaces the simultaneous finite equalities in the family clause of
\Cref{cor:cg-pairwise-bayes-recovery} by descent, a condition on likelihood ratios rather than on
relative entropies; under finite divergence the two agree by \Cref{thm:cg-dpi-equality}, while
\Cref{cor:cg-dpi-infinite-equality-warning} shows that the relative-entropy form fails without
finiteness. For a channel that coarsens a block and leaves the other agents unchanged, part (iii)
separates reconstruction of the block from its coarse coordinate alone from reconstruction that also
reads the other agents.

For a deterministic channel, the equivalence in (ii) between descent of every likelihood ratio and
recovery of the family by $R_{P_0}$ is the factorization theorem of
\citet[Sec.~5, Theorem~1]{HalmosSavage1949} applied to the set $\{P_0\}\cup\{P_\theta\}$, which is
equivalent to $P_0$. For finitely many $\theta$, part (i) characterizes when the coarse experiment is
sufficient for the fine one in the sense of \citet[Sec.~2]{Blackwell1953}, and the conditions of part
(iii) on $P_0$ and the $r_\theta$ are equivalent to adequacy \citep[Sec.~2, Definition~1]{Skibinsky1967}
of the $\sigma$-algebra of $c(X_1)$ for that of $X_1$ with respect to that of $X_2$ and the family
$\{P_0\}\cup\{P_\theta\}$.

\corollaryheading{Infinite equality carries no recovery conclusion}{cor:cg-dpi-infinite-equality-warning}
The finiteness hypothesis of \Cref{thm:cg-dpi-equality} cannot be removed from its equality
characterization: $\KL(PK\Vert QK)=\KL(P\Vert Q)=+\infty$ does not imply that one reverse kernel
recovers $P$ and $Q$. \status{ESTABLISHED}

\noindent\emph{Proof.}
Let $\mathsf X=\{a,b,c\}$ and $\mathsf Y=\{u,v\}$, let the deterministic channel send $a$ to $u$
and $b,c$ to $v$, and set $P=\tfrac12\delta_a+\tfrac12\delta_b$ and
$Q=\tfrac12\delta_b+\tfrac12\delta_c$. Both divergences are $+\infty$, because $P(a)>0=Q(a)$ and
$(PK)(u)>0=(QK)(u)$. A reverse kernel $R$ recovering $Q$ has $R(v,\cdot)=Q$, since $QK=\delta_v$;
then $(PK)R=\tfrac12R(u,\cdot)+\tfrac12Q$ assigns at least $1/4$ to $c$, while $P(c)=0$. $\square$

\section{The common-channel VFE identity and the pointwise parent datum}
\label{sec:cg-elbo}
\label{sec:rg-law-level-vfe}
\label{eq:rg-coarse-channel}

The measure-level case uses the joint and posterior versions of
\cref{prop:prob-measure-evidence}; a record density is needed only for the
scalar free energies below.

Fix the joint $\mathbb P(do,dy)$ on $\mathsf O\times\mathsf Y$ of \Cref{ch:generative} with
$(\theta,V,X)$ suppressed, an admitted observation $o$ with evidence density representative $p(o)$,
the selected posterior-version kernel $o\mapsto\boldsymbol\Pi_o$, a recognition law $\mathbb Q_o$ on
$\mathsf Y$, and a coarse channel $C:\mathsf Y\rightsquigarrow\mathsf Z$ on the latent coordinates,
which reads neither $\mathbb Q_o$ nor the observation. The common pushforward of the three laws is
\begin{equation}
\mathbb P^c(do,dz)=\int_{\mathsf Y}C(y,dz)\mathbb P(do,dy),
\qquad
\boldsymbol\Pi^c_o=\boldsymbol\Pi_oC,
\qquad
\mathbb Q^c_o=\mathbb Q_oC,
\label{eq:rg-common-pushforward}
\end{equation}
and free energies are the extended values $\Fenergy_{\mathbb P}(\mathbb Q_o)=-\log p(o)
+\KL(\mathbb Q_o\Vert\boldsymbol\Pi_o)$ of \Cref{def:elbo-extended}, defined for every law on the
latent space.

\theoremheading{Exact coarse VFE and discarded conditional information}{thm:rg-exact-coarse-vfe}
\label{eq:rg-bridge-lifts}
In the setting of \eqref{eq:rg-common-pushforward}, the following hold. (i) The law $\mathbb P^c$
is normalized and has the observation marginal of $\mathbb P$, so the same evidence representative
$p(o)$ serves both scales. The kernel $o\mapsto\boldsymbol\Pi_oC$ is a selected posterior version of
$\mathbb P^c$, including at the exceptional observations of the fine version. If
$\mathbb Q_o\ll\boldsymbol\Pi_o$, then $\mathbb Q^c_o\ll\boldsymbol\Pi^c_o$. (ii) With the bridge
lifts $\widehat{\mathbb Q}_o(dy,dz)=\mathbb Q_o(dy)C(y,dz)$ and
$\widehat{\boldsymbol\Pi}_o(dy,dz)=\boldsymbol\Pi_o(dy)C(y,dz)$ and their disintegrations over $z$,
the free energies satisfy, in $(-\infty,+\infty]$ and with no absolute-continuity hypothesis,
\begin{equation}
\Fenergy_{\mathbb P}(\mathbb Q_o)
=\Fenergy_{\mathbb P^c}(\mathbb Q^c_o)+\Delta(o),
\qquad
\Delta(o)=\int_{\mathsf Z}
\KL\bigl(\widehat{\mathbb Q}_o(dy\given z)\Vert\widehat{\boldsymbol\Pi}_o(dy\given z)\bigr)
\mathbb Q^c_o(dz)\in[0,+\infty],
\label{eq:rg-vfe-chain-rule}
\end{equation}
so the fine free energy is at least the coarse one. (iii) $\Delta(o)=0$ exactly when the two
discarded conditional laws in \eqref{eq:rg-vfe-chain-rule} agree $\mathbb Q^c_o$-almost surely.
(iv) If $\KL(\mathbb Q_o\Vert\boldsymbol\Pi_o)<\infty$, then
$\Fenergy_{\mathbb P}(\mathbb Q_o)-\Fenergy_{\mathbb P^c}(\mathbb Q^c_o)=\Delta(o)$ is an
ordinary difference of real numbers, and $\Delta(o)=0$ exactly when one Markov kernel
$R:\mathsf Z\rightsquigarrow\mathsf Y$ recovers both laws, $\mathbb Q^c_oR=\mathbb Q_o$ and
$\boldsymbol\Pi^c_oR=\boldsymbol\Pi_o$. This recovery is pairwise; one kernel recovering every
member of a family of laws requires the hypothesis for each member simultaneously.
\status{ESTABLISHED}

The proof is in \cref{sec:proofs-coarse-channels}.

Evidence preservation in part (i) is the exactness condition of Kadanoff's variational
renormalization group on block spins \citep{kadanoff1966scaling}: an exponentiated coupling of fine
and coarse variables whose sum over the coarse variables is one, that is, a Markov kernel from fine
to coarse variables, leaves the free energy unchanged, as stated in
\citet[Eqs.~(8)--(9)]{MehtaSchwab2014}. Part (ii) adds the chain rule of relative entropy for the
recognition law and the posterior.

\corollaryheading{Coarse evidence without record likelihood densities}{cor:cg-measure-evidence}
Let $\mathbb P(do,dy)$ be any joint law, with record marginal
$\mathbb P^{\mathsf O}$, selected posterior $\boldsymbol\Pi_o$, and
recognition kernel $\mathbb Q_o$. For a common coarse channel $C$,
$\boldsymbol\Pi_oC$ is a posterior version of $\mathbb P^c$ and
\begin{equation}
\KL(\mathbb Q_o\Vert\boldsymbol\Pi_o)
=\KL(\mathbb Q_oC\Vert\boldsymbol\Pi_oC)+\Delta(o)
\label{eq:cg-measure-posterior-chain}
\end{equation}
in $[0,+\infty]$, with the backward-conditional term $\Delta(o)$ of
\eqref{eq:rg-vfe-chain-rule}. For $\zeta^{\mathsf O}\ll\mathbb P^{\mathsf O}$,
the joint recognition laws of \cref{prop:prob-measure-evidence} satisfy
\begin{equation}
\KL(\mathbb Q^\zeta\Vert\mathbb P)
=\KL(\mathbb Q^{\zeta,c}\Vert\mathbb P^c)
+\int\Delta(o)\zeta^{\mathsf O}(do).
\label{eq:cg-measure-joint-chain}
\end{equation}
If only the record marginal has a density with finite positive value,
the scalar free-energy identity of \cref{thm:rg-exact-coarse-vfe} follows
from the same density at both scales. These identities compose along
successive coarse channels. They are invariant under changes of the
fine posterior on record-marginal null sets. When
$\mathbb Q_oC\not\ll\boldsymbol\Pi_oC$, the coarse divergence is infinite
and a backward conditional term on coarse null states need not be
version-independent; it is not then an intrinsic finite information loss.
\status{ESTABLISHED}

\noindent\emph{Proof.}
Extending each latent law by the same $C$ leaves its relative entropy
unchanged, since the forward conditional channel agrees. Disintegrate
the extended laws over the coarse state and apply
\cref{lem:prob-kl-chain-rule}(a). Integrate with respect to
$\zeta^{\mathsf O}$ and apply
\eqref{eq:prob-measure-evidence-chain} on both sides. The record marginal
is unchanged, so marginal evidence terms agree and posterior null sets
remain null under $\zeta^{\mathsf O}$. Kernel composition gives the tower
statement. The backward kernel is determined only almost everywhere under
its own coarse marginal, which gives the last qualification.
\hfill$\square$

\corollaryheading{Evidence-preserving coarse channel}{thm:cg-evidence-preserving-channel}
\label{eq:cg-fixed-joint-pushforward}
In the setting of \Cref{thm:rg-exact-coarse-vfe}, $\mathbb P^c(do\times\mathsf Z)=\mathbb P(do\times\mathsf Y)$,
$p^c(o)=p(o)$, and $\boldsymbol\Pi^c_o=\boldsymbol\Pi_oC$ for every declared $o$, and the extended ELBO
of \Cref{def:elbo-extended} satisfies, in $[-\infty,\log p(o)]$,
\begin{equation}
\Lelbo^{\mathrm{ext}}_{\mathbb P^c}(\mathbb Q_oC)
=\log p(o)-\KL(\mathbb Q_oC\Vert\boldsymbol\Pi_oC)
\geq\log p(o)-\KL(\mathbb Q_o\Vert\boldsymbol\Pi_o)
=\Lelbo^{\mathrm{ext}}_{\mathbb P}(\mathbb Q_o).
\label{eq:cg-elbo-monotone}
\end{equation}
When the fine ELBO is finite, the difference of the two ELBOs is $\Delta(o)$: the rise is the
discarded conditional information at unchanged evidence, neither model improvement nor evidence
ascent. For two separately declared models with evidences $p_1(o),p_2(o)$ and finite gaps, the
difference of their ELBOs is $\log p_1(o)-\log p_2(o)$ minus the difference of their gaps, so it
compares variational tightness only when the evidences agree, as they do under one
evidence-preserving common pushforward. \status{ESTABLISHED}

\noindent\emph{Proof.}
The first three identities are part (i) of the theorem. The inequality is data processing in
$[0,+\infty]$ (\Cref{thm:cg-kl-dpi-extended}), and the difference is part (ii) written for
$\Lelbo^{\mathrm{ext}}=-\Fenergy$. The statement on separately declared models is the gap identity
of \Cref{thm:elbo-extended-gap} for each model. $\square$

\remarkheading{Support boundaries of the discarded term}{rem:cg-support-boundaries}
If $\mathbb Q_o\not\ll\boldsymbol\Pi_o$ while $\KL(\mathbb Q^c_o\Vert\boldsymbol\Pi^c_o)<\infty$,
then the fine free energy is $+\infty$, the coarse one is finite, and $\Delta(o)=+\infty$. If
$\mathbb Q^c_o\not\ll\boldsymbol\Pi^c_o$, both free energies are $+\infty$; the conditional law
$\widehat{\boldsymbol\Pi}_o(dy\given z)$ is then undetermined on a $\mathbb Q^c_o$-positive set of
targets, so $\Delta(o)$ depends there on an arbitrary version and the difference of the two
infinite free energies is not a numerical loss. \status{ESTABLISHED}

\noindent\emph{Proof.}
In the first case \eqref{eq:rg-vfe-chain-rule} reads $+\infty=\text{finite}+\Delta(o)$. In the
second, a $\boldsymbol\Pi^c_o$-null set $B$ with $\mathbb Q^c_o(B)>0$ gives
$\widehat{\mathbb Q}_o(\mathsf Y\times B)>0=\widehat{\boldsymbol\Pi}_o(\mathsf Y\times B)$, so both
divergences are infinite, and a disintegration of $\widehat{\boldsymbol\Pi}_o$ over $z$ is
determined only $\boldsymbol\Pi^c_o$-almost surely. $\square$

\propositionheading{Stochastic loss and deterministic minimization}{prop:cg-stochastic-infimum}
\label{eq:rg-thermo-finite-unwarranted-infimum}
\label{eq:rg-thermo-finite-stochastic-infimum-witness}
Let $\mathsf Y$ and $\mathsf Z$ be finite and write
$\Fenergy^c(q')=-\log p(o)+\KL(q'\Vert\boldsymbol\Pi_oC)$ for a law $q'$ on $\mathsf Z$. (i) For a
deterministic statistic $c:\mathsf Y\to\mathsf Z$ and $q'\ll c_\#\boldsymbol\Pi_o$, the fine law
$\mathbb Q_*(y)=q'(c(y))\boldsymbol\Pi_o(y\given c(y))$, defined to be zero where $q'(c(y))=0$
without evaluating a posterior fiber at a null target, has image $q'$ and satisfies
$\Fenergy_{\mathbb P}(\mathbb Q_*)=\Fenergy^c(q')=\min\{\Fenergy_{\mathbb P}(\mathbb Q):
\mathbb QK_c=q'\}$. (ii) For a stochastic channel the statement
$\Fenergy^c(q')=\inf_{\mathbb Q:\mathbb QC=q'}\Fenergy_{\mathbb P}(\mathbb Q)$ can fail: with a
one-point observation alphabet, $\boldsymbol\Pi_o=(1/2,1/2)$,
$C=\begin{psmallmatrix}3/4&1/4\\1/4&3/4\end{psmallmatrix}$, $\mathbb Q=(1,0)$, and
$q'=\mathbb QC=(3/4,1/4)$, the law $\mathbb Q$ is the only fine law with image $q'$, its free
energy is $\log2$, the coarse free energy is $\tfrac34\log\tfrac32+\tfrac14\log\tfrac12$, and the gap
is $2\log2-\tfrac34\log3>0$. \status{ESTABLISHED}

\noindent\emph{Proof.}
(i) Summing $\mathbb Q_*$ over a fiber $c^{-1}(z)$ with $q'(z)>0$ gives $q'(z)$, because the
posterior fiber $\boldsymbol\Pi_o(\cdot\given z)$ is a law on $c^{-1}(z)$; so $\mathbb Q_*$ is a
probability law with image $q'$. By \eqref{eq:rg-vfe-chain-rule}, every $\mathbb Q$ with image $q'$
has $\Fenergy_{\mathbb P}(\mathbb Q)=\Fenergy^c(q')+\Delta\geq\Fenergy^c(q')$, and $\Delta=0$ for
$\mathbb Q_*$ because its occupied reverse fibers are the posterior fibers. (ii) The matrix $C$ has
determinant $1/2$, so $\mathbb Q$ is the unique probability vector with image $q'$. The evidence is
one, $\boldsymbol\Pi_oC=(1/2,1/2)$, and direct evaluation gives the two free energies, whose
difference $2\log2-\tfrac34\log3$ is positive because $2^8>3^3$. $\square$

\corollaryheading{Full pointwise probabilistic datum for a candidate parent}{thm:rg-pointwise-parent-datum}
\label{eq:rg-pointwise-parent-overlap}
\label{eq:rg-pointwise-parent-space}
\label{eq:rg-pointwise-parent-channel}
\label{eq:rg-pointwise-parent-kl-chain}
\label{eq:rg-pointwise-parent-defect}
\label{eq:rg-pointwise-parent-evaluator}
\label{eq:rg-pointwise-parent-evaluator-compatibility}
\label{eq:rg-pointwise-parent-recognition-marginals}
\label{eq:rg-pointwise-parent-prior-marginals}
\label{eq:rg-pointwise-parent-posterior-marginals}
Fix a nonempty finite block $I\subseteq V$ of agents with $c_*\in\bigcap_{i\in I}\mathcal C_i$, a
parent label $A$, and structural data $X$ with $X_A=\chi_A(X)$ kept outside the random channel. Let
the parent state space
$\mathsf Z_A=\mathsf B_A\times\mathsf M_A\times\boldsymbol\Xi_A\times\mathsf H_A$ have nonempty
standard Borel factors for parent belief samples, parent model presentations, the inputs read by
the parent evaluator, and retained marks such as the rooted-holonomy and dressed-boundary data of
\Cref{sec:rg-gauge-cross-scale}. Supply the joint $\mathbb P_I(do,dy\given X)$ of the block with
its selected posterior-version kernel $o\mapsto\boldsymbol\Pi_{I,o,X}$, a recognition law
$\mathbb Q_{I,o,X}\ll\boldsymbol\Pi_{I,o,X}$ at an admitted $o$, and one parent coarse kernel
$C_A:\mathsf Y_I\rightsquigarrow\mathsf Z_A$, a declared input like $I$, not a canonical selection.
Define
\begin{align}
\mathbb P_A(do,dz\given X)&=\int_{\mathsf Y_I}C_A(y,dz)\mathbb P_I(do,dy\given X),
\label{eq:rg-pointwise-parent-generative}\\
\boldsymbol\Pi_{A,o,X}=\boldsymbol\Pi_{I,o,X}C_A,&\qquad
\mathbb Q_{A,o,X}=\mathbb Q_{I,o,X}C_A.
\label{eq:rg-pointwise-parent-triple}
\end{align}
Then the following hold. (i) All three parent laws are normalized, the parent observation marginal
and evidence equal the fine ones, $o\mapsto\boldsymbol\Pi_{A,o,X}$ is the selected parent posterior
version, $\mathbb Q_{A,o,X}\ll\boldsymbol\Pi_{A,o,X}$, and in $[0,+\infty]$
$\KL(\mathbb Q_{I,o,X}\Vert\boldsymbol\Pi_{I,o,X})
=\KL(\mathbb Q_{A,o,X}\Vert\boldsymbol\Pi_{A,o,X})+\Delta_A(o,X)$, where $\Delta_A(o,X)\geq0$ is
the discarded term $\Delta(o)$ of \eqref{eq:rg-vfe-chain-rule} for
$\mathbb P=\mathbb P_I$ and $C=C_A$. Adding $-\log p_X(o)$ gives $\Fenergy_I=\Fenergy_A+\Delta_A$ in
the extended reals, where a finite free energy may be negative. The zero-defect and recovery
statements of \Cref{thm:rg-exact-coarse-vfe} hold with $C=C_A$; a bare equality
$+\infty=+\infty$ supplies neither. (ii) The parent model coordinate is generative only through a
jointly measurable normalized evaluator
$\operatorname{ev}_A:\mathsf M_A\to\operatorname{Kern}(\boldsymbol\Xi_A,\mathsf B_A\times\mathsf
O\times\mathsf H_A)$, $m_A\mapsto K^{X_A}_{A,m_A}$. A selected standard-Borel disintegration of
$\mathbb P_A$ over $(m_A,\xi_A)$ supplies one such evaluator; an evaluator declared independently of
$\mathbb P_A$ is compatible with it only under the separate hypothesis that
$\mathbb P_A(db_A,do,dh_A\given\xi_A,m_A,X)=K^{X_A}_{A,m_A}(\xi_A;db_A,do,dh_A)$ almost surely for
the parent $(m_A,\xi_A)$ marginal. Disintegration neither validates an arbitrary declared
evaluator nor selects its null extension or makes evaluation injective. (iii) With
$\mathbb P^Z_A(D\given X)=\mathbb P_A(\mathsf O\times D\given X)$, the parent recognition, prior, and
posterior marginals in channel $x\in\{b,m\}$ are the coordinate pushforwards
$q_A^x=(\operatorname{pr}_x)_\#\mathbb Q_{A,o,X}$,
$p_A^x=(\operatorname{pr}_x)_\#\mathbb P^Z_A(\cdot\given X)$, and
$\boldsymbol\Pi^x_{A,o,X}=(\operatorname{pr}_x)_\#\boldsymbol\Pi_{A,o,X}$. If both laws of a
marginal pair are non-Dirac, or if $\boldsymbol\Xi_A\times\mathsf H_A$ has at least two points, the
pair does not determine the full law on $\mathsf Z_A$. If $\boldsymbol\Xi_A\times\mathsf H_A$ is a
point and one marginal is a point mass, the coupling class of the pair is a single law.
\status{ESTABLISHED}

The proof is in \cref{sec:proofs-coarse-channels}.

\section{Pushforward, precomposition, and Galerkin restriction}
\label{sec:cg-three-operations}

Three operations share the name coarse-graining. A Markov pushforward
(\Cref{def:cg-coarse-channel}) moves laws and obeys data processing; precomposition identifies
variables inside an energy and defines a new model; Galerkin restriction compresses an operator and
needs further data before it propagates signals. A proof about one of them does not transfer to
another.

\definitionheading{Energy precomposition}{def:cg-energy-precomposition}
For a measurable injection $\iota:\mathsf Z\to\mathsf X$ and an energy $E$ on $\mathsf X$, the
precomposed model is
\begin{equation}
\bar E(z)=E(\iota z),
\qquad
\bar P(dz)=\frac{e^{-\bar E(z)}}{\bar Z}\bar\nu(dz),
\label{eq:cg-energy-precomposition}
\end{equation}
defined once a reference $\bar\nu$ on $\mathsf Z$ and a finite normalizer $\bar Z$ are declared;
the reference, its Radon-Nikodym or Jacobian rule, and the normalizer are part of the construction.

\propositionheading{Hard identification is conditioning}{prop:cg-hard-identification}
(i) On a finite or countable $\mathsf X$, with counting references on $\mathsf Z$ and $\mathsf X$
and $E=-\log p$ for the mass function of a law $P$, the precomposed law is
$\bar P(z)=p(\iota z)/P(\iota(\mathsf Z))$: the law of $X$ conditioned on $X\in\iota(\mathsf Z)$,
transported by $\iota^{-1}$. (ii) Conditioning can increase relative entropy. On
$\mathsf X=\{1,2,3\}$ with $\iota(\mathsf Z)=\{1,2\}$, the laws $P=(1,1,98)/100$ and
$Q=(1,2,97)/100$ have $\KL(P\Vert Q)\approx0.0031$, while their conditioned laws $(1/2,1/2)$ and
$(1/3,2/3)$ have relative entropy $\approx0.0589$. Hence no Markov kernel sends $P$ and $Q$ to their
precomposed laws, and data processing, evidence preservation, and \eqref{eq:cg-elbo-monotone} do not
transfer to precomposition. (iii) When $\iota(\mathsf Z)$ is null for the fine reference, as for
block-constant configurations of a continuous fiber, the fine law defines no conditional law on it,
and different references give different precomposed laws: for $\bar E(z)=z^2$ on $\mathsf Z=\R$,
Lebesgue reference gives $\mathcal N(0,1/2)$ and the reference $e^{-z^2}dz$ gives
$\mathcal N(0,1/4)$. Exact marginalization is a Markov pushforward, and restricting a recognition
family keeps the joint and optimizes the same ELBO over a subset, which is neither pushforward nor
precomposition. \status{ESTABLISHED}

\noindent\emph{Proof.}
(i) Substitution in \eqref{eq:cg-energy-precomposition} gives $\bar P(z)\propto p(\iota z)$, with
normalizer $P(\iota(\mathsf Z))$. (ii) The values are direct evaluation; every Markov kernel
contracts relative entropy (\Cref{thm:cg-kl-dpi-extended}), and the conditioned pair has larger
relative entropy. (iii) Conditioning on a null event is not defined by the law, and the densities
$e^{-z^2}$ and $e^{-2z^2}$ are proportional to the stated Gaussians. The last sentence restates the
definitions. $\square$

\subsection{Galerkin operators and message passing}
\label{sec:cg-galerkin}

\definitionheading{Galerkin data}{def:cg-galerkin-data}
\label{eq:cg-galerkin-data}
Let $S:V_c\to V_f$ be an injective linear map of finite-dimensional real spaces, let $M\succ0$ be a
fine reference form, and let $L$ be a fine bilinear form. Put $M_c=S^\top MS$,
$S^+=M_c^{-1}S^\top M$, $H=M^{-1}L$, and $H_c=S^+HS=M_c^{-1}S^\top LS$. The form $S^\top LS$
is the Galerkin coarse form of aggregation-based algebraic multigrid
\citep[Eqs.~(2.1) and (2.2)]{Notay2010Aggregation}, \citep[Sec.~2]{LivneBrandt2012LAMG}, here for
an arbitrary injective prolongation as in smoothed aggregation \citep{VanekMandelBrezina1996}; the
map $H_c$ is the coarse propagator.

\propositionheading{Exact message-passing residual}{prop:cg-message-passing-residual}
\label{eq:cg-message-residual}
The data of \Cref{def:cg-galerkin-data} satisfy $S^+S=I$ and $S^+H-H_cS^+=S^+H(I-SS^+)$. The
idempotent $SS^+$ is the $M$-orthogonal projector onto $\operatorname{range}S$, with kernel
$\ker S^+=(\operatorname{range}S)^{\perp_M}$. Hence exact coarse message passing, $S^+H=H_cS^+$ on
all fine signals, holds exactly when
\begin{equation}
H\bigl((\operatorname{range}S)^{\perp_M}\bigr)\subseteq(\operatorname{range}S)^{\perp_M}.
\label{eq:cg-galerkin-exactness}
\end{equation}
The operator $H$ is self-adjoint in the $M$-inner product exactly when $L$ is symmetric, and then
\eqref{eq:cg-galerkin-exactness} is equivalent to $H$-invariance of $\operatorname{range}S$.
Without symmetry the two invariance conditions are independent: with $M=I$, $S=e_1$, and
$L=\begin{psmallmatrix}2&1&0\\0&3&0\\0&0&4\end{psmallmatrix}$, the subspace $\operatorname{range}S$
is invariant but the residual $S^+H(I-SS^+)=(0,1,0)$ is nonzero, while $L^\top$ leaves the
complement invariant and gives zero residual although $\operatorname{range}S$ is not invariant.
\status{ESTABLISHED}

\noindent\emph{Proof.}
$S^+S=M_c^{-1}S^\top MS=I$, and substituting $H_c=S^+HS$ gives the residual identity. The map
$SS^+$ is idempotent with range $\operatorname{range}S$ and kernel $\ker S^+$; since $M_c$ is
invertible, $S^+y=0$ exactly when $S^\top My=0$, that is, when $y$ is $M$-orthogonal to
$\operatorname{range}S$. The residual vanishes exactly when $S^+Hy=0$ for every
$y\in\operatorname{range}(I-SS^+)=\ker S^+$, which is \eqref{eq:cg-galerkin-exactness}. Next
$\langle Hx,y\rangle_M=x^\top L^\top y$ and $\langle x,Hy\rangle_M=x^\top Ly$, and for an
$M$-self-adjoint operator a subspace is invariant exactly when its $M$-orthogonal complement is.
The two witnesses are direct computations. $\square$

\section{Gauge and permutation equivariance}
\label{sec:cg-equivariance}

\propositionheading{Equivariant Markov channels}{prop:cg-equivariant-channels}
\label{eq:cg-law-equivariance}
\label{eq:cg-kernel-equivariance}
Let a group $\mathbb G$ act bimeasurably on $\mathsf X$ and $\mathsf Y$, and let
$K:\mathsf X\rightsquigarrow\mathsf Y$ be equivariant, $K(gx,B)=K(x,g^{-1}B)$. Then
$(g_\#P)K=g_\#(PK)$ for every law $P$ on $\mathsf X$, and composites of equivariant kernels are
equivariant. \status{ESTABLISHED}

\noindent\emph{Proof.}
For measurable $B$ the law $(g_\#P)K$ gives $B$ the mass $\int K(gx,B)P(dx)$, which by
equivariance equals $\int K(x,g^{-1}B)P(dx)=[g_\#(PK)](B)$. Substitution in the composition of
\Cref{prop:cg-markov-category} proves closure under composition. $\square$

At $c_*$ a principal gauge transformation acts on belief and model coordinates through $\rho_b(g)$
and $\rho_m(g)$ of one $g\in G$ (\Cref{prop:geo-active-image-point}), and each passive
re-expression $g_i$ acts on both channels; a coarse channel on joint belief and model coordinates is
gauge covariant when it is equivariant for these actions, and for agent relabelings the acting group
is a permutation group.

\propositionheading{Equivariant transport of Bayes recovery kernels}{prop:cg-equivariant-bayes-versions}
\label{eq:cg-bayes-version-transport}
\label{eq:cg-finite-equivariant-bayes}
\label{eq:cg-haar-averaged-bayes}
Let $\mathbb G$ act on $\mathsf X$ and $\mathsf Y$ by bimeasurable bijections, let $K$ be
equivariant as in \Cref{prop:cg-equivariant-channels}, let $Q$ be a law on $\mathsf X$, let $R_Q$ be a regular
conditional law of $X$ given $Y$ under $\mathbb J_Q(dx,dy)=Q(dx)K(x,dy)$, and put
$R_Q^g(y,A)=R_Q(g^{-1}y,g^{-1}A)$ for $g\in\mathbb G$. (i) $R_Q^g$ is a regular conditional law of
$X$ given $Y$ under $\mathbb J_{g_\#Q}$ for every $g$. (ii) If $PKR_Q=P$, then
$(g_\#P)KR_Q^g=g_\#P$, and $\KL((g_\#P)K\Vert(g_\#Q)K)=\KL(PK\Vert QK)$ and
$\KL(g_\#P\Vert g_\#Q)=\KL(P\Vert Q)$. (iii) If $\mathsf X$ and $\mathsf Y$ are finite and
$\mathbb G$ acts by permutations, the rule $R_Q(y,\{x\})=Q(\{x\})K(x,\{y\})/(QK)(\{y\})$ where
$(QK)(\{y\})>0$ and $R_Q(y,\{x\})=1/\lvert\mathsf X\rvert$ where $(QK)(\{y\})=0$ selects a version
for every $Q$, depends Borel measurably on $(Q,y)$, and satisfies
$R_{g_\#Q}(gy,\{gx\})=R_Q(y,\{x\})$ for all $g,Q,y,x$. (iv) If $\mathbb G$ is a topological group
acting jointly measurably and the stabilizer $\mathbb G_Q=\{h:h_\#Q=Q\}$ is a compact metrizable
group with Haar probability $dh$, then $\bar R_Q(y,A)=\int_{\mathbb G_Q}R_Q^h(y,A)dh$ is a version
of $R_Q$ with $\bar R_Q(hy,hA)=\bar R_Q(y,A)$ for $h\in\mathbb G_Q$, and $R_{g_\#Q}:=\bar R_Q^g$ is
well defined on the orbit of $Q$ and satisfies $R_{k_\#Q'}(ky,kA)=R_{Q'}(y,A)$ for every $Q'$ in
that orbit and every $k\in\mathbb G$. \status{ESTABLISHED}

The proof is in \cref{sec:proofs-coarse-channels}.

\corollaryheading{Continuous posterior versions are equivariant}{cor:cg-continuous-version-equivariant}
In the setting of \Cref{prop:cg-equivariant-bayes-versions}, let $\mathsf X$ and $\mathsf Y$ be Polish,
let every $h\in\mathbb G_Q$ act on both by homeomorphisms, and let $S$ be the support of $QK$. If some
version of $R_Q$ is weakly continuous on $S$, then $R(y,\cdot)=R_Q(y,\cdot)$ for $y\in S$ and
$R(y,\cdot)=Q$ for $y\notin S$ is a version of $R_Q$ with $R(hy,hA)=R(y,A)$ for every
$h\in\mathbb G_Q$, every $y\in\mathsf Y$, and every measurable $A$, whether or not the stabilizer is
compact. \status{ESTABLISHED}

\noindent\emph{Proof.}
By \Cref{prop:cg-equivariant-channels}, $h_\#(QK)=(h_\#Q)K=QK$ for $h\in\mathbb G_Q$, so $S$ is a
closed $\mathbb G_Q$-invariant set, and $R$ is a Markov kernel that agrees with $R_Q$ on the full-measure
set $S$. For $h\in\mathbb G_Q$ the kernel $R_Q^h$ is a version of $R_Q$ by part (i) of the proposition,
and $y\mapsto R_Q^h(y,\cdot)=h_\#R_Q(h^{-1}y,\cdot)$ is weakly continuous on $S$, because $h^{-1}$ is a
homeomorphism of $\mathsf Y$ preserving $S$ and pushforward by the homeomorphism $h$ of $\mathsf X$ is
weakly continuous. Two versions agree $QK$-almost everywhere, and the set where two weakly continuous
maps from $S$ into the Hausdorff space of laws agree is closed in $S$; its complement in $S$ is an open
$QK$-null subset of the support, hence empty. So $R_Q^h=R_Q$ on $S$, which is the equivariance identity
at $y\in S$; for $y\notin S$ also $h^{-1}y\notin S$, and $R(h^{-1}y,h^{-1}A)=Q(h^{-1}A)=Q(A)=R(y,A)$.
$\square$

\openproblemheading{Equivariant Bayes versions beyond the finite and compact cases}{open:cg-equivariant-versions}
\Cref{prop:cg-equivariant-bayes-versions} settles equivariant Bayes versions for finite spaces and
along orbits with compact stabilizer, and \Cref{cor:cg-continuous-version-equivariant} settles them for
every stabilizer when the posterior has a weakly continuous version on the support of $QK$. Two cases
remain open in the class of a locally compact second countable group acting continuously on Polish
spaces, an equivariant channel $K$, and laws $Q$ whose stabilizer is noncompact, so that no Haar
probability averages versions, and whose posterior has no weakly continuous version on the support of
$QK$: an everywhere-defined version equivariant under the stabilizer; and a version that is
simultaneously equivariant and jointly measurable across a measurable family of such laws meeting
several orbits, which also needs a measurable choice of orbit representatives. The problem is settled
by an equivariant measurable selection theorem giving, for every member of that class, an
everywhere-defined version equivariant under the stabilizer, jointly measurable across every measurable
family of such laws meeting several orbits, or by a pair $(Q,K)$ of the class with no everywhere-defined
version equivariant under the stabilizer.

\section{Independence partitions and information levels}
\label{sec:cg-intrinsic-selection}

\definitionheading{Independent partition and conditional context}{def:cg-independence-partition}
For a law $P$ on $\prod_{i\in V}\mathsf Y_i$, let $\operatorname{Part}(V)$ be the finite
set of unlabeled partitions of $V$. A partition $\mathfrak p$ is independent when
\begin{equation}
\mathsf T_P(\mathfrak p):=\KL\Bigl(P\Bigm\Vert\bigotimes_{I\in\mathfrak p}P_I\Bigr)=0.
\label{eq:cg-independent-partition}
\end{equation}
If $P$ is instead a law of $(Y_V,\Upsilon)$ with an already specified retained context
$\Upsilon$, conditional independence means
\begin{equation}
\mathsf T_{P\mid\Upsilon}(\mathfrak p):=
\int\KL\Bigl(P(dy_V\given v)\Bigm\Vert
\bigotimes_{I\in\mathfrak p}P(dy_I\given v)\Bigr)P_\Upsilon(dv)=0.
\label{eq:cg-conditionally-independent-partition}
\end{equation}
Both quantities are defined in $[0,+\infty]$ and require no densities.

\propositionheading{Finest independent partition and a law-based objective}{prop:cg-finest-independent-partition}
Every law $P$ of \cref{def:cg-independence-partition} has a unique finest independent
partition $\mathfrak p_*(P)$. It uniquely minimizes
\begin{equation}
\mathscr J^{\mathrm{ind}}_P(\mathfrak p)
:=\mathbf1_{\{\mathsf T_P(\mathfrak p)>0\}}-
\frac{|\mathfrak p|}{N+1},
\label{eq:cg-independent-selector-objective}
\end{equation}
which, whenever the finest partition is intermediate, is not monotone toward either
partition endpoint. The selector is natural under agent relabeling and invariant under
separate bimeasurable coordinate changes of the agents, including their gauge
re-expressions. With $\mathsf T_{P\mid\Upsilon}$ in place of $\mathsf T_P$, the same
claims hold for the unique finest partition that is conditionally independent given
$\Upsilon$. Selecting that context is additional data: taking $\Upsilon=Y_V$ makes
every partition conditionally independent. The objective is a preference for maximal
factorization, not a consequence of maximizing the evidence bound.
\status{ESTABLISHED}

The proof is in \cref{sec:proofs-coarse-channels}. The unconditional target is the
independence-clustering target of \citet[Def.~1]{Ryabko2017IndependenceClustering}; the
statement includes the general-law and conditional-context versions.

\propositionheading{Nondegeneracy and instability under weak coupling}{prop:cg-independent-selector-instability}
For every even $N\geq4$, on $\{-1,1\}^N$ the positive law
\begin{equation}
P_0(y)\propto\exp\left(\mathrm j\sum_{b=1}^{N/2}y_{2b-1}y_{2b}\right),
\text{ with }\mathrm j>0,
\label{eq:cg-independent-pair-law}
\end{equation}
has $\mathfrak p_*(P_0)=\{\{2b-1,2b\}:1\leq b\leq N/2\}$. Thus
\eqref{eq:cg-independent-selector-objective} selects an intermediate partition in a
non-Gaussian family at arbitrarily large finite size. Replacing the exponent by
\begin{equation}
\mathrm j\sum_{b=1}^{N/2}y_{2b-1}y_{2b}
+\mathrm r\sum_{b=1}^{N/2-1}y_{2b}y_{2b+1},
\text{ with }\mathrm r>0,
\label{eq:cg-independent-pair-perturbation}
\end{equation}
gives $\mathfrak p_*(P_{\mathrm r})=\{V\}$ for every $\mathrm r>0$, although
$P_{\mathrm r}\to P_0$ in total variation as $\mathrm r\downarrow0$. Consequently
the exact-independence selector need not be stable under arbitrarily weak interactions.
\status{ESTABLISHED}

The proof is in \cref{sec:proofs-coarse-channels}.

\propositionheading{Exchangeability obstructs universal hard selection}{prop:cg-exchangeable-hard-selector}
Let $N\geq3$, let all agent spaces coincide, and let $P$ be invariant under every
agent permutation. Every deterministic, relabeling-natural hard partition selector
that reads only $P$ returns either the singleton partition or $\{V\}$. This excludes
universal nondegeneracy in that class of selectors. It does not exclude a randomized
membership law, additional structural inputs, or nondegenerate selection on a restricted
class of nonsymmetric laws. \status{ESTABLISHED}

The proof is in \cref{sec:proofs-coarse-channels}.

\definitionheading{Partition information and its threshold family}{def:cg-partition-information}
For a finite-alphabet law $P$ and $I\subseteq V$ with $|I|\geq2$, put
\begin{align}
\mathsf I_I(P)&:=\min_{\substack{\mathfrak p\in\operatorname{Part}(I)\\|\mathfrak p|\geq2}}
\frac{\sum_{J\in\mathfrak p}\mathrm H(P_J)-\mathrm H(P_I)}{|\mathfrak p|-1},\notag\\
\mathscr A_P(\mathrm a)&:=\{I\subseteq V:|I|\geq2,\mathsf I_I(P)>\mathrm a\},
\label{eq:cg-partition-information}
\end{align}
where $\mathrm H$ is Shannon entropy with natural logarithms. The maximal members of
$\mathscr A_P(\mathrm a)$ are the clusters at level $\mathrm a$; retaining the entire
family over all levels leaves the choice of one level unspecified.

\propositionheading{Information-clustering hierarchy}{prop:cg-information-hierarchy}
The maximal threshold clusters of \cref{def:cg-partition-information}, over all levels,
form a hierarchical family: any two clusters are disjoint or one contains the other.
This is the information hierarchy of
\citet[Eqs.~(3.3a)--(3.3b), Thms.~3.1--3.3]{Chan2016InfoClustering}.
\status{STANDARD}

\propositionheading{Information levels and finite-alphabet stability}{prop:cg-partition-information-stability}
The quantity $\mathsf I_I(P)$ is nonnegative and vanishes exactly when $P_I$ has a
nontrivial product partition. Let $P,Q$ be laws on the same finite product alphabet with
$\operatorname{TV}(P,Q)\leq\delta\leq1/2$, and let
$D_I=\prod_{i\in I}|\mathsf Y_i|$. With
$h_2(\delta)=-\delta\log\delta-(1-\delta)\log(1-\delta)$, including its continuous
endpoint values,
\begin{equation}
|\mathsf I_I(P)-\mathsf I_I(Q)|\leq
3h_2(\delta)+2\delta\log D_I.
\label{eq:cg-partition-information-stability}
\end{equation}
In particular, with $\eta=3h_2(\delta)+2\delta\log D_V$,
\begin{equation}
\mathscr A_P(\mathrm a+\eta)\subseteq\mathscr A_Q(\mathrm a)
\subseteq\mathscr A_P(\mathrm a-\eta).
\label{eq:cg-partition-threshold-stability}
\end{equation}
If every subset score has distance greater than $\eta$ from a specified threshold,
its qualified subset family and maximal clusters are unchanged. The bound charges
alphabet and block size; it is not a continuity assertion for unrestricted continuous
laws. Separate bimeasurable re-expressions preserve the relative-entropy form of the
partition numerator, so its information levels are gauge invariant.
\status{ESTABLISHED}

The proof is in \cref{sec:proofs-coarse-channels}.

\openproblemheading{A stable natural selector of partitions and membership laws}{open:cg-natural-selector}
\label{eq:cg-partition-naturality}
A data-dependent hard partition selector $S(X)$ must be natural under agent relabeling: for every
fine permutation $P$ there must exist a coarse relabeling $Q=Q(P,X)$ with
$S(P\cdot X)=PS(X)Q^\top$, that is, the unlabeled partition is equivariant. The quantifier on $Q$
is existential because coarse labels are arbitrary; a universal one would force $S(X)=S(X)Q^\top$
for every $Q$ at $P=\mathrm{id}$, which no partition into two or more blocks satisfies.
Gauge-carrying inputs must enter through invariant assignment scores, while represented messages
transform equivariantly. These conditions type an admissible selector of a partition, a membership
law, or a structural channel. The literal existence of a law-based functional with an
intermediate optimizer is supplied by
\cref{prop:cg-finest-independent-partition,prop:cg-independent-selector-instability};
\cref{prop:cg-exchangeable-hard-selector} rules out universal nondegeneracy for a
deterministic hard selector reading a permutation-invariant law alone. Neither result
supplies a stable soft choice for interacting laws. An unpenalized sum of nonnegative
within-block divergences is minimized by the singleton partition
(\Cref{rem:cg-structural-vs-marginal}), the Gaussian recognition costs of the Gaussian supplement
select trivial endpoints or need an external coefficient, and one common recovery kernel for a
whole statistical experiment, \eqref{eq:cg-holonomy-common-recovery}, is available only under
simultaneous family hypotheses. Remaining is a relabeling-natural, gauge-compatible
functional on a stated class of interacting laws, not monotone under refinement,
whose nondegenerate optimizer is stable under a stated probability-law perturbation
bound and positive score margin, without imposing the number of blocks, an external
relevance variable, or a fitted penalty coefficient. For the information hierarchy of
\cref{def:cg-partition-information}, the choice of one level or a membership law must
also follow from that functional. A construction with these properties, or a theorem
excluding them on the same stated class, settles this remaining question. A statistical
partition alone supplies no restoring dynamics, predictive parent closure, or causal
second formation of \cref{open:rg-pointwise-parent-dynamics,open:full-graph-hierarchy}.

\section{Scale diagrams}
\label{sec:rg-scale-category}

\definitionheading{Scale diagram}{def:rg-scale-diagram}
\label{eq:rg-scale-diagram}
Let $\mathsf S$ be a category whose objects are resolutions $\ell$, and choose one target category
$\mathscr K$: sets and maps for a deterministic state, standard Borel spaces and Markov kernels for
a stochastic state, or a declared category of Banach spaces and bounded linear maps, or a declared
product of such categories, for an operator or mixed state. A scale diagram is a functor
$\mathfrak X:\mathsf S\to\mathscr K$, with objects $\mathfrak X_\ell$ and arrows
$C_{\ell k}\in\operatorname{Hom}_{\mathscr K}(\mathfrak X_\ell,\mathfrak X_k)$ satisfying
$C_{kr}\circ C_{\ell k}=C_{\ell r}$, where $\circ$ is composition in $\mathscr K$, componentwise in
a product. For a stochastic component, kernels act on the right,
\begin{equation}
K_{\ell r}:=K_{\ell k}K_{kr},
\qquad
(\mu K_{\ell k})K_{kr}=\mu K_{\ell r},
\label{eq:rg-right-acting-kernel-composition}
\end{equation}
so that $K_{\ell k}K_{kr}$ represents the abstract composite $K_{kr}\circ K_{\ell k}$. One symbol is
never at once a map, a kernel, and an operator. An expression on a reference space needs declared
isomorphisms, and transformations of reference measures stay explicit.

\propositionheading{When functoriality has content}{prop:rg-scale-functoriality}
(i) On a chain of resolutions $0<1<\cdots<L$, any family of adjacent arrows $C_{\ell,\ell+1}$
extends uniquely to a scale diagram, by $C_{\ell k}=C_{k-1,k}\circ\cdots\circ C_{\ell,\ell+1}$ and
$C_{\ell\ell}=\operatorname{id}$; the composition law imposes no condition. (ii) On a poset of
resolutions that is not a chain, functoriality requires any two routes between the same
resolutions to agree. For a family of kernels generated by one blocking rule indexed by block
factors, functoriality is the semigroup law $K_{b_1b_2}=K_{b_1}K_{b_2}$. (iii) Deterministic arrows
that keep the complete block configuration, or an additive block statistic such as block sums,
satisfy both conditions. Majority with ties resolved by the first block member does not: on six
agents it sends $(0,0,1,0,1,0)$ to $1$ by pairs and then triples, and to $0$ by triples and then
pairs. (iv) Mutually inverse Markov kernels between finite spaces are permutation matrices, so no
stochastic isomorphism exists between finite levels with different numbers of configurations.
\status{ESTABLISHED}

The proof is in \cref{sec:proofs-coarse-channels}.

\section{Labeled families and retained label information}
\label{sec:rg-thermo-finite-experiment}

A finite mixture that keeps its label in the sample space is the elementary model of a random
partition or random link pattern among agents at $c_*$; its Fisher tensor splits into a label part
and a weighted branch part.

\definitionheading{Finite labeled family}{def:rg-thermo-label-family}
\label{def:rg-thermo-finite-score}
\label{hyp:rg-thermo-finite-support}
\label{eq:rg-thermo-label-family}
\label{eq:rg-thermo-finite-fisher-definition}
Let $\Theta\subseteq\R^d$ be open, $1\leq d<\infty$, and let $\mathfrak G$ be a finite nonempty
set of labels $\gamma$, for instance partitions of a block of agents or link patterns. Each label
has a finite nonempty carrier $\mathsf X_\gamma$, and the sample space is the disjoint union of the
$\{\gamma\}\times\mathsf X_\gamma$, so that labels stay distinct when branch laws coincide. A finite
labeled family on $\Theta$ consists of effective log masses $\lambda_\gamma\in C^3(\Theta)$ and
branch laws $p_{\gamma,\theta}$ on $\mathsf X_\gamma$, with
\begin{equation}
P_\theta(\gamma,x)=W_\gamma(\theta)p_{\gamma,\theta}(x),
\qquad
W_\gamma=\frac{e^{\lambda_\gamma}}{\sum_{\delta\in\mathfrak G}e^{\lambda_\delta}},
\label{eq:rg-thermo-finite-log-masses}
\end{equation}
on a fixed support stratum: every branch mass is either identically zero on $\Theta$, a structural
zero that carries no score, or positive and $C^3$ in $\theta$. A common additive function of
$\theta$ in all $\lambda_\gamma$ is a change of representative. When the family comes from a
normalization model, the effective log masses absorb its label weights and branch partitions
(\Cref{prop:rg-geometry-weights}). For a fixed-support finite family $\theta\mapsto R_\theta$ the
score is $\mathbf s_R=\nabla_\theta\log R_\theta$ on positive atoms and the Fisher tensor is
$g^F(R)=\sum_{z:R_\theta(z)>0}R_\theta(z)\mathbf s_R(z)\mathbf s_R(z)^\top$; write
$\xi_\gamma=\nabla\lambda_\gamma$ and $\bar\xi=\sum_\gamma W_\gamma\xi_\gamma$. An observation model
is a record kernel $K^{\mathrm{obs}}(o\given\gamma,x)$ on a finite alphabet $\mathsf O$, of the same
fixed-support $C^3$ type and independent of recognition; it defines the complete observation law
$P^{\mathrm{obs}}_\theta(\gamma,x,o)=P_\theta(\gamma,x)K^{\mathrm{obs}}(o\given\gamma,x)$ and the
evidence $p_\theta(o)=\sum_{\gamma,x}P^{\mathrm{obs}}_\theta(\gamma,x,o)$. The complete latent,
complete observation, and observed-data experiments are $\theta\mapsto P_\theta$,
$\theta\mapsto P^{\mathrm{obs}}_\theta$, and $\theta\mapsto p_\theta$. An extensive labeled family
is a sequence $P_N=(W_{\gamma,N}p_{\gamma,N})$ of finite labeled families indexed by an extensive
size $N\geq1$, such as the number of agents, with carriers that may depend on $N$. On a compact
$\mathsf K\subset\Theta$ it has extensive label growth with bulk log-mass densities
$\mathsf s_\gamma$ if, on an open neighborhood of $\mathsf K$,
\begin{equation}
\|\lambda_{\gamma,N}\|_{C^3(\mathsf K)}\leq\mathrm C_{\mathsf K}N,
\qquad
\lambda_{\gamma,N}/N\longrightarrow\mathsf s_\gamma\quad\text{in }C^3(\mathsf K),
\label{eq:rg-thermo-label-effective-growth}
\end{equation}
where $\|f\|_{C^k(\mathsf K)}$ is the largest supremum over $\mathsf K$ of the ordinary coordinate
derivatives of $f$ through order $k$, including order zero, with fixed norms on vector and matrix
entries.

\propositionheading{Complete labeled Fisher decomposition}{prop:rg-thermo-finite-labeled-fisher}
\label{eq:rg-thermo-finite-branch-score}
For a finite labeled family, the branch score $\mathbf s_{p_\gamma}=\nabla\log p_\gamma$ on the
positive atoms of $\mathsf X_\gamma$ satisfies $\sum_{x}p_\gamma(x)\mathbf s_{p_\gamma}(x)=0$, and
the latent score, the Fisher tensor, and, for every multi-index $\alpha$ with
$\lvert\alpha\rvert\leq2$, the derivatives of its weighted branch term satisfy
\begin{align}
\mathbf s_P(\gamma,x)&=\xi_\gamma-\bar\xi+\mathbf s_{p_\gamma}(x),
\label{eq:rg-thermo-finite-complete-score}\\
g^F(P)&=\operatorname{Cov}_W(\xi_\gamma)+\sum_\gamma W_\gamma g^F(p_\gamma),
\label{eq:rg-thermo-finite-labeled-fisher}\\
\partial^\alpha\Bigl(\sum_\gamma W_\gamma g^F(p_\gamma)\Bigr)
&=\sum_\gamma\sum_{\beta\leq\alpha}\binom{\alpha}{\beta}
(\partial^\beta W_\gamma)(\partial^{\alpha-\beta}g^F(p_\gamma)).
\label{eq:rg-thermo-finite-weighted-fisher-derivatives}
\end{align}
The first term of $g^F(P)$ is the Fisher tensor of the label marginal $W$. A common additive
function in every $\lambda_\gamma$ leaves these identities and their derivatives unchanged, and
$g^F(P)\succeq0$ with no rank assumption. \status{ESTABLISHED}

The proof is in \cref{sec:proofs-coarse-channels}.

\propositionheading{Selected posterior and its complete score}{prop:rg-thermo-finite-posterior-score}
\label{def:rg-thermo-finite-posterior}
\label{eq:rg-thermo-finite-evidence}
\label{eq:rg-thermo-finite-posterior}
\label{eq:rg-thermo-finite-posterior-labels}
For a finite labeled family with an observation model and a record $o$, let
$e_\gamma(o)=\sum_xp_\gamma(x)K^{\mathrm{obs}}(o\given\gamma,x)$ be the branch evidence, so that
$p_\theta(o)=\sum_\gamma W_\gamma e_\gamma(o)$, and where $p_\theta(o)>0$ define the selected
posterior $\boldsymbol\Pi_o(\gamma,x)=W_\gamma p_\gamma(x)K^{\mathrm{obs}}(o\given\gamma,x)/p_\theta(o)$
and its label weights $W^o_\gamma=W_\gamma e_\gamma(o)/p_\theta(o)$. A label with $e_\gamma(o)=0$
stays on the carrier with zero posterior mass and has no conditional
branch; zero evidence leaves the posterior undefined. On occupied labels $W^o$ is the softmax of
$\lambda^o_\gamma=\lambda_\gamma+\log e_\gamma(o)$. On the positive posterior support the complete
score and, within an occupied branch
$p^o_\gamma(x)=p_\gamma(x)K^{\mathrm{obs}}(o\given\gamma,x)/e_\gamma(o)$, the branch score are
\begin{align}
\mathbf s_{\boldsymbol\Pi_o}(\gamma,x)
&=\mathbf s_P(\gamma,x)+\nabla\log K^{\mathrm{obs}}(o\given\gamma,x)-\nabla\log p_\theta(o),
\label{eq:rg-thermo-finite-posterior-score}\\
\nabla\log p^o_\gamma(x)&=\mathbf s_{p_\gamma}(x)+\nabla\log K^{\mathrm{obs}}(o\given\gamma,x)
-\nabla\log e_\gamma(o),
\label{eq:rg-thermo-finite-conditional-posterior-score}
\end{align}
and the decomposition \eqref{eq:rg-thermo-finite-labeled-fisher} holds for the posterior family
$\theta\mapsto\boldsymbol\Pi_{o,\theta}$, with the weights $W^o$ and the branches $p^o_\gamma$.
\status{ESTABLISHED}

\noindent\emph{Proof.}
Summing the posterior over all labeled atoms gives one, and the label weights
are proportional to $e^{\lambda_\gamma}e_\gamma(o)$. Differentiating the logarithm of the positive
posterior mass gives \eqref{eq:rg-thermo-finite-posterior-score}, and the same computation within
one branch gives \eqref{eq:rg-thermo-finite-conditional-posterior-score}. The fixed zero patterns of
branch laws and record kernel keep the occupied labels and atoms fixed on the stratum, so no
division by zero occurs, and the proof of \Cref{prop:rg-thermo-finite-labeled-fisher} applies to the
posterior family. $\square$

\propositionheading{Observed-data information and missing information}{prop:rg-thermo-finite-missing-information}
For a finite labeled family with an observation model, let $\mathsf O^+=\{o:p_\theta(o)>0\}$, fixed
on the stratum, let $\mathbf s^c=\nabla\log P^{\mathrm{obs}}$ be the complete score on the positive
support of the complete observation law, and let
$g^F(p)=\sum_{o\in\mathsf O^+}p_\theta(o)\nabla\log p_\theta(o)\nabla\log p_\theta(o)^\top$ be the
observed-data Fisher tensor. Then
\begin{align}
g^F(P^{\mathrm{obs}})
&=g^F(P)+\sum_\gamma W_\gamma\sum_xp_\gamma(x)g^F\bigl(K^{\mathrm{obs}}(\cdot\given\gamma,x)\bigr),
\label{eq:rg-thermo-finite-complete-observation-fisher}\\
\nabla\log p_\theta(o)
&=\E_{\boldsymbol\Pi_o}\bigl[\mathbf s^c(\gamma,X,o)\bigr],
\qquad
g^F(p)=g^F(P^{\mathrm{obs}})-\sum_{o\in\mathsf O^+}p_\theta(o)
\operatorname{Cov}_{\boldsymbol\Pi_o}(\mathbf s^c).
\label{eq:rg-thermo-finite-missing-information}
\end{align}
In particular $g^F(p)\preceq g^F(P^{\mathrm{obs}})$, and positive definiteness of $g^F(P)$ or of
$g^F(P^{\mathrm{obs}})$ does not imply that of $g^F(p)$: for a one-point alphabet $p_\theta(o)=1$
identically and $g^F(p)=0$ for all branch data. \status{ESTABLISHED}

The proof is in \cref{sec:proofs-coarse-channels}.

Fisher contraction requires a channel that does not depend on the parameter; the constant-input
Bernoulli witness of \Cref{thm:hist-record-clock-contraction} shows that a parameter-dependent
channel can create information that its input does not carry.

\corollaryheading{Stochastic score projection and retained-label loss}{prop:rg-thermo-finite-fisher-loss}
\label{def:rg-thermo-finite-channel}
\label{eq:rg-thermo-finite-channel}
\label{eq:rg-thermo-finite-score-projection}
\label{eq:rg-thermo-finite-fisher-loss}
A label-preserving channel is given by row-stochastic matrices $C_\gamma(x,y)$ from
$\mathsf X_\gamma$ to finite nonempty targets $\mathsf Y_\gamma$, independent of $\theta$ and of
recognition, and lifted as zero between different labels, so that
$P^C(\gamma,y)=\sum_xP(\gamma,x)C_\gamma(x,y)$. The source law may be the latent family or a
positive selected posterior; on the complete observation law the channel acts on $x$ and keeps $o$.
For such a channel, on positive targets, the score projection
$\mathbf s_{P^C}(\gamma,y)=\E_{P\otimes C}[\mathbf s_P(\gamma,X)\given\gamma,y]$ and the Fisher loss
$g^F(P)-g^F(P^C)=\E_{P\otimes C}[(\mathbf s_P-\mathbf s_{P^C})(\mathbf s_P-\mathbf s_{P^C})^\top]\succeq0$
hold, and for $P=(W_\gamma p_\gamma)$
\begin{equation}
g^F(P)-g^F(P^C)=\sum_{\gamma:W_\gamma>0}W_\gamma\bigl[g^F(p_\gamma)-g^F(p_\gamma C_\gamma)\bigr],
\qquad
0\preceq g^F(P)-g^F(P^C)\preceq\sum_\gamma W_\gamma g^F(p_\gamma),
\label{eq:rg-thermo-finite-within-loss}
\end{equation}
while the label marginal and all its derivatives are unchanged. \status{ESTABLISHED}

\noindent\emph{Proof.}
The score projection and the Fisher loss are \Cref{thm:cg-fisher-contraction} for the finite
fixed-support family, or directly: differentiating $P^C$ at a positive target and dividing by its
mass gives the conditional expectation, and the residual $\mathbf s_P-\mathbf s_{P^C}$ has zero
conditional mean given the target, so the cross terms vanish. Row sums give
$\sum_yP^C(\gamma,y)=W_\gamma$ as functions of $\theta$. Subtracting the two decompositions
\eqref{eq:rg-thermo-finite-labeled-fisher} gives the first part of
\eqref{eq:rg-thermo-finite-within-loss}, and $g^F(p_\gamma C_\gamma)\succeq0$ gives the upper bound.
A zero-mass target receives no incidence from a positive atom and is a structural zero. $\square$

\propositionheading{Finite-word score and VFE telescopes}{prop:rg-thermo-account-word}
\label{eq:rg-thermo-account-score-telescope}
\label{eq:rg-thermo-account-vfe-telescope}
\label{eq:rg-thermo-account-vfe-area}
Let $C_1,\ldots,C_m$ be parameter-independent Markov kernels between finite spaces, let $P_0$ be a
fixed-support finite family, and put $P_j=P_0C_1\cdots C_j$ on its fixed positive support, with
score $\mathbf s_j$ and Fisher tensor $g_j=g^F(P_j)$. Then $g_0-g_m=\sum_{j=1}^mD_j$ with
$D_j=g_{j-1}-g_j\succeq0$, and on the path law $P_0(x_0)C_1(x_0,x_1)\cdots C_m(x_{m-1},x_m)$ the
increments $\varepsilon^{\mathrm{sc}}_j=\mathbf s_{j-1}(X_{j-1})-\mathbf s_j(X_j)$ are orthogonal
with $\E[\varepsilon^{\mathrm{sc}}_j(\varepsilon^{\mathrm{sc}}_j)^\top]=D_j$, so
$\E[(\mathbf s_0(X_0)-\mathbf s_m(X_m))(\mathbf s_0(X_0)-\mathbf s_m(X_m))^\top]=\sum_jD_j$. For a
record with positive evidence, a selected posterior $\boldsymbol\Pi_0$, and a recognition law
$\mathbb Q_0\ll\boldsymbol\Pi_0$, push both laws and the joint by the same word and set
$\Fenergy_j=-\log p(o)+\KL(\mathbb Q_j\Vert\boldsymbol\Pi_j)$. Then
$\Fenergy_0-\Fenergy_m=\sum_{j=1}^mK^{\mathrm{KL}}_j$, where $K^{\mathrm{KL}}_j\geq0$ is the
discarded conditional relative entropy of step $j$. If the fine and
coarse families have extensive sizes $N_0$ and $N_m$, the per-agent bookkeeping reads
$\tfrac{N_0}{N_m}\cdot\tfrac{\Fenergy_0}{N_0}=\tfrac{\Fenergy_m}{N_m}+\tfrac{1}{N_m}\sum_{j=1}^mK^{\mathrm{KL}}_j$.
\status{ESTABLISHED}

The proof is in \cref{sec:proofs-coarse-channels}.

\section{Retained-label Fisher asymptotics}
\label{sec:cg-retained-label-asymptotics}

When log masses grow like the number of agents $N$, the softmax calculus of
\Cref{prop:rg-thermo-finite-labeled-fisher} has three regimes with different hypotheses and
coordinates: a strict winner, a tie with converging weights, and a window of width $1/N$ at a tie.
Fix an extensive labeled family with weights $W_{\gamma,N}$, log masses
$\lambda_{\gamma,N}$, gradients $\xi_{\gamma,N}=\nabla\lambda_{\gamma,N}$, Fisher density
$g_N=N^{-1}g^F(P_N)$, and branch densities $\mathcal I_{\gamma,N}=N^{-1}g^F(p_{\gamma,N})$ with limits
$\mathcal I_\gamma$ when they converge; $\mathrm C_{\mathsf K}$ depends on $\mathsf K$ and the family
and may change from line to line, and limits are taken at a fixed label set.

\theoremheading{Strict-phase convergence with both metric derivatives}{thm:rg-thermo-label-strict}
\label{eq:rg-thermo-label-strict-limit}
\label{eq:rg-thermo-label-native-floor-input}
\label{eq:rg-thermo-label-weight-first}
\label{eq:rg-thermo-label-mean-derivative}
\label{eq:rg-thermo-label-weight-second}
\label{eq:rg-thermo-label-losing-derivatives}
\label{eq:rg-thermo-label-strict-centered-covariance}
\label{eq:rg-thermo-label-strict-branch-error}
Assume extensive label growth \eqref{eq:rg-thermo-label-effective-growth} on a neighborhood of a
compact $\mathsf K$, and $\mathcal I_{\gamma,N}\to\mathcal I_\gamma$ in $C^2(\mathsf K)$ for every
label. Suppose a label $\gamma_*$ satisfies
$\mathsf s_{\gamma_*}-\mathsf s_\gamma\geq\delta_{\mathsf K}>0$ on $\mathsf K$ for every
$\gamma\neq\gamma_*$. Then $\|g_N-\mathcal I_{\gamma_*}\|_{C^2(\mathsf K)}\to0$. More precisely, for
$j=0,1,2$ and all large $N$,
\begin{equation}
\Bigl\|D^j\Bigl[N^{-1}\operatorname{Cov}_{W_N}(\xi_{\gamma,N})\Bigr]\Bigr\|_{\mathsf K}
\leq\mathrm C_{\mathsf K}N^{j+1}e^{-\delta_{\mathsf K}N/2},
\label{eq:rg-thermo-label-strict-covariance-bound}
\end{equation}
and the weighted branch error $\sum_\gamma W_{\gamma,N}\mathcal I_{\gamma,N}-\mathcal I_{\gamma_*,N}$
has $C^2(\mathsf K)$ norm at most $\mathrm C_{\mathsf K}N^2e^{-\delta_{\mathsf K}N/2}$. Only three
derivatives of the effective log masses are used. If moreover
$\mathcal I_{\gamma_*}\succeq\mathrm c_{\min}\mathrm{Id}_d$ on $\mathsf K$ for a constant
$\mathrm c_{\min}>0$, then $g_N(\theta)\succeq\tfrac12\mathrm c_{\min}\mathrm{Id}_d$ for all
$\theta\in\mathsf K$ and all large $N$. \status{ESTABLISHED}

The proof is in \cref{sec:proofs-coarse-channels}.

\propositionheading{Squared-size Fisher tensor at a finite-weight tie}{prop:rg-thermo-label-tie-scale}
\label{eq:rg-thermo-label-tie-direction}
\label{eq:rg-thermo-label-tie-limit}
Fix $\theta_0$ and assume \eqref{eq:rg-thermo-label-effective-growth} on a neighborhood of
$\theta_0$. Let $S$ be the nonempty set of labels maximizing $\mathsf s_\gamma(\theta_0)$, suppose
the weights satisfy $W_{\gamma,N}(\theta_0)\to\bar{\mathsf p}_\gamma$, and assume
$g^F(p_{\gamma,N})(\theta_0)=O(N)$ for every label. Write
$\mathsf v_\gamma=\nabla\mathsf s_\gamma(\theta_0)$ and
$\bar{\mathsf v}=\sum_\gamma\bar{\mathsf p}_\gamma\mathsf v_\gamma$. Then $\bar{\mathsf p}$ is
supported on $S$ and $N^{-2}g^F(P_N)(\theta_0)\to\operatorname{Cov}_{\bar{\mathsf p}}(\mathsf v_\gamma)
=\sum_\gamma\bar{\mathsf p}_\gamma(\mathsf v_\gamma-\bar{\mathsf v})(\mathsf v_\gamma-\bar{\mathsf v})^\top$.
Its rank is the dimension of the span of the differences $\mathsf v_\gamma-\mathsf v_\eta$ over
pairs with positive limiting weights, at most
$\min(d,\lvert\operatorname{supp}\bar{\mathsf p}\rvert-1)$. A direction with positive variance has
Fisher information asymptotic to a positive multiple of $N^2$; zero variance gives only $o(N^2)$ in
that direction. \status{ESTABLISHED}

The proof is in \cref{sec:proofs-coarse-channels}.

\theoremheading{Center ratios, smooth weights, and the full Fisher limit}{thm:rg-thermo-label-crossover}
\label{eq:rg-thermo-label-window}
Assume \eqref{eq:rg-thermo-label-effective-growth} and $\mathcal I_{\gamma,N}\to\mathcal I_\gamma$
in $C^2$ on a neighborhood of $\theta_0$. Let $S$ be the maximizing set of
$\mathsf s_\gamma(\theta_0)$, choose $\gamma_0\in S$, and assume the center-ratio condition
\begin{equation}
\exp\bigl[\lambda_{\gamma,N}(\theta_0)-\lambda_{\gamma_0,N}(\theta_0)\bigr]\longrightarrow
\mathrm c_\gamma\in(0,\infty)\quad(\gamma\in S),
\qquad\mathrm c_{\gamma_0}=1.
\label{eq:rg-thermo-label-center-ratios}
\end{equation}
Set $\mathsf v_\gamma=\nabla\mathsf s_\gamma(\theta_0)$ and define the window experiment on the
original labeled carrier by $\theta_N(\mathsf w)=\theta_0+\mathsf w/N$ and
$\mathcal E_N(\mathsf w)=P_{N,\theta_N(\mathsf w)}$. On every compact set of $\mathsf w$ its label
law converges in $C^3$ to
\begin{equation}
\bar{\mathsf p}_\gamma(\mathsf w)=
\frac{\mathrm c_\gamma e^{\mathsf w\cdot\mathsf v_\gamma}}{\sum_{\eta\in S}\mathrm c_\eta e^{\mathsf w\cdot\mathsf v_\eta}}
\quad(\gamma\in S),
\qquad
\bar{\mathsf p}_\gamma(\mathsf w)=0\quad(\gamma\notin S),
\label{eq:rg-thermo-label-limiting-weights}
\end{equation}
and its Fisher tensor in the coordinate $\mathsf w$ converges in $C^2$:
\begin{equation}
g^F_{\mathsf w}(\mathcal E_N)(\mathsf w)=N^{-2}g^F(P_N)(\theta_N(\mathsf w))
\longrightarrow
\operatorname{Cov}_{\bar{\mathsf p}(\mathsf w)}(\mathsf v_\gamma)
=D^2_{\mathsf w}\log\sum_{\gamma\in S}\mathrm c_\gamma e^{\mathsf w\cdot\mathsf v_\gamma}.
\label{eq:rg-thermo-label-crossover-fisher}
\end{equation}
The center-ratio condition is additional data at finite $N$; equality of the bulk densities
supplies no such condition. \status{ESTABLISHED}

The proof is in \cref{sec:proofs-coarse-channels}.

\propositionheading{Ties that diverge or do not converge}{prop:rg-thermo-label-tie-counterexamples}
\label{eq:rg-thermo-label-linear-tie-example}
\label{eq:rg-thermo-label-oscillatory-data}
\label{eq:rg-thermo-label-oscillatory-failure}
Extensive label growth, even with identical parameter-independent branches, implies neither a
finite extensive Fisher density at a tie, nor convergence of the label law, nor $C^2$ convergence
of a uniformly convergent Fisher density. For two labels $0,1$ with one common branch law
independent of $\theta$ and log-odds $\omega_N=\lambda_{1,N}-\lambda_{0,N}$,
\begin{equation}
W_{1,N}=\frac{e^{\omega_N}}{1+e^{\omega_N}},
\qquad
g^F(P_N)=\operatorname{bv}(\omega_N)(\omega'_N)^2,
\qquad
\operatorname{bv}(x)=\frac{e^x}{(1+e^x)^2}=\frac1{4\cosh^2(x/2)},
\label{eq:rg-thermo-label-binary-formula}
\end{equation}
and each of the following families has extensive label growth on every compact interval. (a) For
$\omega_N=N\theta$ the tie at $\theta=0$ has $g^F(P_N)(\theta)=N^2/[4\cosh^2(N\theta/2)]$ and
$g^F(P_N)(0)/N=N/4$, so the extensive density diverges, and $N^{-2}g^F(P_N)$ tends to $1/4$ at zero
and to zero elsewhere. (b) For $\omega_N$ equal to $0$ at even $N$ and $\log2$ at odd $N$, the
Fisher tensor vanishes identically while $W_{1,N}$ alternates between $1/2$ and $2/3$. (c) For
$\omega_N=N^{\mathrm c}\theta$ with $0<\mathrm c<1$, both bulk densities vanish and
\begin{equation}
\frac{g^F(P_N)(0)}N=\frac14N^{2\mathrm c-1},
\label{eq:rg-thermo-label-subextensive-example}
\end{equation}
which is below, at, or above the extensive scale according as $\mathrm c<1/2$, $\mathrm c=1/2$, or
$\mathrm c>1/2$; for $\omega_N=\mathrm c_N+\sqrt N\theta$ with $\mathrm c_N$ alternating between $0$
and $\log2$, the density at zero alternates between $1/4$ and $2/9$. (d) For
$\omega_N=N\theta-\mathrm c\log N$ with $\mathrm c>0$, the tied label is rare,
$W_{1,N}(0)=(1+N^{\mathrm c})^{-1}\to0$, yet
\begin{equation}
\frac{g^F(P_N)(0)}N=\frac{N^{1+\mathrm c}}{(1+N^{\mathrm c})^2}\sim N^{1-\mathrm c},
\label{eq:rg-thermo-label-rare-example}
\end{equation}
so the limiting $N^2$ coefficient is zero while the density diverges for $\mathrm c<1$, tends to one
for $\mathrm c=1$, and vanishes for $\mathrm c>1$. (e) For
$\omega_N(\theta)=N^{-1/6}\sin(N^{1/3}\theta)$ both weights converge uniformly to $1/2$ and
$F_N=N^{-1}g^F(P_N)$ converges uniformly to zero, while
\begin{equation}
F''_N=\frac1N\Bigl[\operatorname{bv}''(\omega_N)(\omega'_N)^4+5\operatorname{bv}'(\omega_N)(\omega'_N)^2\omega''_N
+2\operatorname{bv}(\omega_N)(\omega''_N)^2+2\operatorname{bv}(\omega_N)\omega'_N\omega'''_N\Bigr]
\label{eq:rg-thermo-label-oscillatory-derivative}
\end{equation}
satisfies $F''_N(0)=-\tfrac12-\tfrac18N^{-1/3}\to-\tfrac12$. \status{ESTABLISHED}

The proof is in \cref{sec:proofs-coarse-channels}.

\remarkheading{The retained label is observed information}{rem:cg-label-observation}
The covariance term of \eqref{eq:rg-thermo-finite-labeled-fisher} is the Fisher tensor of the label
marginal, so the $N^2$ scales of \Cref{prop:rg-thermo-label-tie-scale,thm:rg-thermo-label-crossover}
are information about the label. When all carriers are one common set, the experiment that observes
only the configuration, through $(\gamma,x)\mapsto x$, loses exactly the expected conditional
covariance of $\mathbf s_P$ given $X$, which vanishes when $X$ determines the label; for a record
the loss is given by \Cref{prop:rg-thermo-finite-missing-information}. In the binary families of
\Cref{prop:rg-thermo-label-tie-counterexamples} the configuration alone has a law independent of
$\theta$ and zero Fisher information, while for $\omega_N=N\theta$ the labeled window experiment has
Fisher tensor $\operatorname{bv}(\mathsf w)$ in $\mathsf w$ for every $N$. \status{ESTABLISHED}

\noindent\emph{Proof.}
The label marginal has score $\xi_\gamma-\bar\xi$. The map $(\gamma,x)\mapsto x$ is a
parameter-independent statistic, so \Cref{thm:cg-fisher-contraction} gives the loss as the expected
conditional covariance, which vanishes when the score is a function of $X$, in particular when $X$
determines $(\gamma,x)$. With one common $\theta$-independent branch law $p$, the configuration has
law $\sum_\gamma W_\gamma p=p$. For $\omega_N=N\theta$ at $\theta=\mathsf w/N$,
\eqref{eq:rg-thermo-label-binary-formula} gives $N^{-2}\operatorname{bv}(\mathsf w)N^2$. $\square$

\paragraph{Physicist's summary.}
A coarse channel is a noisy detector on latent coordinates, under which relative entropy and Fisher
information can only drop. Applied to the generative joint, the posterior, and the recognition law at
once, it leaves the evidence untouched, and the free energy drops by exactly what it discarded about
the disagreement of recognition and posterior inside each coarse cell, which is bookkeeping, not
learning. Exact independence supplies an intrinsic partition for general laws, but arbitrarily weak
coupling can destroy that partition; the information-level stability bound requires a finite alphabet
and a score margin. These are statistical selection statements, not parent dynamics or formation
(\cref{prop:cg-independent-selector-instability,prop:cg-partition-information-stability,open:cg-natural-selector}). A retained discrete label, such as which partition is present, carries its own Fisher
information: with $N$ agents a strict winner hides it, a tie exposes information of order $N^2$, and
near a tie the natural coordinate is $N(\theta-\theta_0)$, as in the finite-volume rounding of a
first-order transition.

\chapter[Gauge-Retained Network State and Fixed Sectors]{Gauge-Retained Network State\\and Fixed Sectors}
\label{ch:network-state}

Grouping the agents at $c_*$ into blocks does not average away their gauge data: links are independent
group elements and closed walks carry holonomy. This chapter asks what a coarse description of a block
must retain of that data and which part of the block's state survives exact aggregation. A rooted spanning
tree retains tree transports, the based holonomy, and block-dressed boundary edges as one root-framed datum
(\Cref{prop:rg-rooted-datum}). The fixed sector of a block is the kernel of its connection Laplacian, with
positive-definite weights the holonomy-fixed root values (\Cref{prop:cg-kernel-holonomy}), and transported
marginal laws admit a parent law exactly when they occupy one holonomy-stabilized mode
(\Cref{thm:cg-holonomy-kl-marginal}). Cut links compress exactly into rectangular endpoint maps
(\Cref{thm:cg-rectangular-endpoint-closure}), and into one group-valued link with new self terms exactly
when an explicit matrix is invertible (\Cref{lem:cg-cut-completed-square}), with one balanced choice of
link and weight (\Cref{prop:cg-balanced-cut-presentation}) that no presentation at the level of laws
reproduces once the cut has excess (\Cref{prop:cg-no-balanced-law-presentation}) and that admits no continuous
extension to all singular cuts (\Cref{open:cg-cut-compression}); the coarse frame is invisible to
the constituents and cannot be computed from their frames by a symmetric equivariant rule
(\Cref{prop:cg-frame-sector}); and a record at $c_*$ distinguishes loop holonomies at the level of laws
(\Cref{thm:cg-loop-operational-holonomy}). These results supply the retained gauge data of
\Cref{ch:agent-network-rg,ch:full-graph-meta-agent-vfe}, the formation witnesses of
\Cref{ch:finite-depth-formation}, and the testability discussion of \Cref{ch:philosophy}; Supplement G
realizes them for Gaussian families.

\paragraph{Scope.} The fixed-sector rank is an admissibility criterion: it assigns no partition cost,
keeps singletons admissible, and does not require constituent beliefs to coincide. The marginal-law score
takes its block and weights as inputs, selects no partition, supplies no Markov kernel, and does not
recover the dependence structure of a joint law. A cross map parallel for the links neither identifies the
belief and model bundles nor orders their two scores, and the operational record, which sees a loop only
through its represented holonomy at the probe, recovers no link values. \status{NOT-CLAIMED} The obstruction
to a continuous extension of balanced cut links is \Cref{open:cg-cut-compression}.

\section{Levelwise fiber state and scale intertwiners at \texorpdfstring{$c_*$}{c*}}
\label{sec:rg-common-principal-state}

\definitionheading{Levelwise fiber state}{def:rg-geometric-state}
At depth $\ell$ fix a Lie group $G_\ell$ and a right $G_\ell$-torsor $\mathscr P_\ell$, the fiber over $c_*$
of the principal bundle of that depth, containing one frame of each agent or meta-agent of depth $\ell$,
used for both channels. Law fibers $\mathcal B_{\ell,x}$, $x\in\{b,m\}$, carry left actions
$\widehat\rho_{\ell,x}$ of $G_\ell$, which may be inequivalent, with associated fibers
$\mathcal E_{\ell,x}=\mathscr P_\ell\times_{\widehat\rho_{\ell,x}}\mathcal B_{\ell,x}$. The depth-$\ell$
interaction multigraph $\mathfrak I_\ell$ carries links $\Theta^x_e\in G_\ell$ with
$\Theta^x_{\bar e}=(\Theta^x_e)^{-1}$ (\Cref{def:geo-graph-links}), every block $I$ of depth $\ell$ carries
its based holonomy groups $\mathrm{Hol}^x_I$ (\Cref{def:rg-rooted-tree-data}), and declared cross maps
$\Phi_\ell:\mathcal E_{\ell,b}\to\mathcal E_{\ell,m}$ and $\widetilde\Phi_\ell:\mathcal E_{\ell,m}\to\mathcal E_{\ell,b}$
need not be invertible or mutually inverse. The fiber state of depth $\ell$ is
\begin{equation}
\mathfrak X_\ell=\Bigl(G_\ell,\mathscr P_\ell;\ \mathcal E_{\ell,b},\mathcal E_{\ell,m};\
\mathfrak I_\ell,(\Theta^b_e,\Theta^m_e)_e,(\mathrm{Hol}^b_I,\mathrm{Hol}^m_I)_I;\
\Phi_\ell,\widetilde\Phi_\ell;\ \ldots\Bigr),
\label{eq:rg-full-state}
\end{equation}
where the ellipsis stands for law, kernel, or operator components declared by a realization.

\propositionheading{Scale maps and intertwiners at \texorpdfstring{$c_*$}{c*}}{prop:rg-scale-intertwiner}
(i) A map $\Phi_\ell:\mathcal E_{\ell,b}\to\mathcal E_{\ell,m}$, expressed in a frame $u$ by the fiber map
$f_u$ with $\Phi_\ell[u,w]=[u,f_u(w)]$, has one expression $f^{bm}_\ell$ in every frame exactly when
$f^{bm}_\ell\circ\widehat\rho_{\ell,b}(g)=\widehat\rho_{\ell,m}(g)\circ f^{bm}_\ell$ for every $g\in G_\ell$,
which is also the condition that $\Phi_\ell$ commutes with every automorphism of $\mathscr P_\ell$ and
constrains the covariance of a cross map, not its existence; the reverse direction is analogous. For a
homomorphism $\kappa_\ell:G_\ell\to G_{\ell+1}$, points
$p_0\in\mathscr P_\ell$ and $p_1\in\mathscr P_{\ell+1}$, and fiber maps
$q_{\ell,x}:\mathcal B_{\ell,x}\to\mathcal B_{\ell+1,x}$, $x\in\{b,m\}$, consider
\begin{align}
\mathsf P_\ell(p\cdot g)&=\mathsf P_\ell(p)\cdot\kappa_\ell(g),
\label{eq:rg-principal-scale-map}\\
\widehat q_{\ell,x}[p,w]&=[\mathsf P_\ell(p),q_{\ell,x}(w)],
\label{eq:rg-associated-scale-map}\\
q_{\ell,x}\circ\widehat\rho_{\ell,x}(g)&=\widehat\rho_{\ell+1,x}(\kappa_\ell(g))\circ q_{\ell,x}
\qquad\text{for every }g\in G_\ell.
\label{eq:rg-scale-intertwiner}
\end{align}
(ii) Exactly one map $\mathsf P_\ell:\mathscr P_\ell\to\mathscr P_{\ell+1}$ satisfies $\mathsf P_\ell(p_0)=p_1$
and \eqref{eq:rg-principal-scale-map}. (iii) The formula \eqref{eq:rg-associated-scale-map} defines a map
$\widehat q_{\ell,x}:\mathcal E_{\ell,x}\to\mathcal E_{\ell+1,x}$ exactly when \eqref{eq:rg-scale-intertwiner}
holds. (iv) For composable steps the homomorphisms, torsor maps, fiber maps, and induced maps compose, the
rightmost factor acting first, and on law fibers given by pushforward the composite is the pushforward of
the composed kernel, by \eqref{eq:cg-law-functor}. \status{ESTABLISHED}

\noindent\emph{Proof.}
(i) Each point of $\mathcal E_{\ell,b}$ has one representative with first entry $u$, and
$[u\cdot g,w]=[u,\widehat\rho_{\ell,b}(g)w]$ gives
$f_{u\cdot g}=\widehat\rho_{\ell,m}(g)^{-1}\circ f_u\circ\widehat\rho_{\ell,b}(g)$, as in
\eqref{eq:geo-phi-gauge-law}. An automorphism of the torsor, $u\mapsto u\cdot\kappa$, acts by
$[u,w]\mapsto[u,\widehat\rho(\kappa)w]$, so commuting with all of them is the same condition. (ii) The
action is free and transitive, so $\mathsf P_\ell(p_0\cdot g):=p_1\cdot\kappa_\ell(g)$ is well defined,
satisfies \eqref{eq:rg-principal-scale-map}, and is forced by it. (iii) The representatives $(p,w)$ and
$(p\cdot g,\widehat\rho_{\ell,x}(g)^{-1}w)$ of one point have images $[\mathsf P_\ell(p),q_{\ell,x}(w)]$ and
$[\mathsf P_\ell(p),\widehat\rho_{\ell+1,x}(\kappa_\ell(g))q_{\ell,x}(\widehat\rho_{\ell,x}(g)^{-1}w)]$, which
agree for every $(p,w,g)$ exactly under \eqref{eq:rg-scale-intertwiner}. (iv) The composites satisfy the
defining identities by substitution. $\square$

On a chain of depths these composites form a scale diagram (\Cref{def:rg-scale-diagram,prop:rg-scale-functoriality});
at the information-geometric tier an induced map also needs the pushed laws to stay in the coarse model and
$q_{\ell,x}$ to be smooth. For an edge copy $e:j\to i$ of $\mathfrak I_\ell$, the cross-scale morphism
defects and the link-intertwining defects are the ordered pairs
\begin{equation}
\begin{aligned}
\mathbf D^\Phi_\ell&=\bigl(\widehat q_{\ell,m}\circ\Phi_\ell,\ \Phi_{\ell+1}\circ\widehat q_{\ell,b}\bigr),
&
\mathbf D^{\widetilde\Phi}_\ell&=\bigl(\widehat q_{\ell,b}\circ\widetilde\Phi_\ell,\ \widetilde\Phi_{\ell+1}\circ\widehat q_{\ell,m}\bigr),\\
\mathbf D^\Phi_e&=\bigl(\widehat{\mathsf t}^m_e\circ\Phi_\ell,\ \Phi_\ell\circ\widehat{\mathsf t}^b_e\bigr),
&
\mathbf D^{\widetilde\Phi}_e&=\bigl(\widehat{\mathsf t}^b_e\circ\widetilde\Phi_\ell,\ \widetilde\Phi_\ell\circ\widehat{\mathsf t}^m_e\bigr),
\end{aligned}
\label{eq:rg-cross-morphism-defects}
\end{equation}
where $\widehat{\mathsf t}^x_e[p,w]=[\mathsf t^x_e(p),w]$ is induced by the link map $\mathsf t^x_e$ of
\Cref{prop:geo-edge-twist-invariance}. A defect vanishes when its two components are equal; subtracting
them needs a vector or affine realization.

\propositionheading{Covariance of the defect diagrams}{prop:rg-defect-covariance}
The maps $\widehat{\mathsf t}^x_e$ are well defined and act from the frame of $j$ to the frame of $i$ as
$\widehat\rho_x(\Theta^x_e)$. Under passive re-expressions at depths $\ell$ and $\ell+1$, both components
of every pair in \eqref{eq:rg-cross-morphism-defects} acquire the same coordinate changes at source and
target, so the equality locus of each pair is gauge invariant; in a fiberwise-linear realization the
difference of the components transforms by the same left and right laws. \status{ESTABLISHED}

\noindent\emph{Proof.}
Right equivariance of $\mathsf t^x_e$ gives
$[\mathsf t^x_e(p\cdot g),\widehat\rho_x(g)^{-1}w]=[\mathsf t^x_e(p),w]$, and
$\mathsf t^x_e(u_j\cdot g)=u_i\cdot\Theta^x_eg$ gives the coordinate form. Both components of each pair are
frame-free maps between the same two associated fibers, so their frame expressions are conjugated by the
same coordinate changes, and equality of two maps does not depend on the frames. $\square$

\section{Rooted trees, based holonomy, and dressed edges}
\label{sec:rg-gauge-cross-scale}

Fix a depth and suppress its index. Let $\mathcal P$ be a partition of the agents into blocks and
$\mathfrak I_I$ the internal multigraph of a block $I$, with its loops and parallel copies. A passive
re-expression $g=(g_i)$ acts by $\Theta^x_e\mapsto g_i\Theta^x_eg_j^{-1}$ for $e:j\to i$
\eqref{eq:geo-regime-two-gauge-law} and by $z^x_i\mapsto\rho_x(g_i)z^x_i$ on sample coordinates
\eqref{eq:geo-represented-frame-change}.

\definitionheading{Rooted spanning-tree data}{def:rg-rooted-tree-data}
The component-indexed meta-agent set of $\mathcal P$ is
\begin{equation}
\widehat{\mathcal P}=\{(I,\mathsf c):I\in\mathcal P,\ \mathsf c\in\pi_0(\mathfrak I_I)\}.
\label{eq:rg-component-meta-index}
\end{equation}
From here on $I,J$ denote elements of $\widehat{\mathcal P}$, each connected with one root $r_I$; the state
of a disconnected block is the product over its components, and collapsing that product to one root, row,
or feature is an additional coarse channel. Choose a rooted spanning tree in each $I$, nested across
depths. With the path products $\Theta^x_p$ and holonomies $\mathrm{Hol}^x(p)$ of
\eqref{eq:geo-link-holonomy}, the tree transport of $i\in I$ along the tree path $p_{I,i}$ from $i$ to
$r_I$, the holonomy of the fundamental cycle of a non-tree copy $e:j\to i$ of $\mathfrak I_I$, the based
holonomy group of $I$ at a vertex $r$, and the block-dressed edge of a boundary copy $e:j\to i$ from $J$ to
$I$ are
\begin{align}
&\Theta^x_{I\leftarrow i}=\Theta^x_{p_{I,i}},
\qquad
\mathrm{Hol}^x(p_e)=\Theta^x_{I\leftarrow i}\Theta^x_e(\Theta^x_{I\leftarrow j})^{-1},
\label{eq:full-graph-tree-transport-law}\\
&\mathrm{Hol}^x_I(r)=\{\Theta^x_p:p\text{ a closed walk of }\mathfrak I_I\text{ at }r\}\subseteq G,
\qquad\mathrm{Hol}^x_I=\mathrm{Hol}^x_I(r_I),
\label{eq:full-graph-based-holonomy}\\
&\widehat\Theta^x_e=\Theta^x_{I\leftarrow i}\Theta^x_e(\Theta^x_{J\leftarrow j})^{-1}.
\label{eq:full-graph-dressed-edge}
\end{align}
For an action $\rho_x$ of $G$ on the channel-$x$ sample fiber, the root-framed child tuple
$\bar z^x_I=(\rho_x(\Theta^x_{I\leftarrow i})z^x_i)_{i\in I}$ collects the member coordinates in the root
frame, and a law $Q$ on the root sample fiber of $I$ is holonomy stabilized when
\begin{equation}
\widehat\rho_x(\mathsf k)Q=Q\qquad\text{for every }\mathsf k\in\mathrm{Hol}^x_I,
\label{eq:full-graph-holonomy-stabilization}
\end{equation}
which is law coherence up to holonomy stabilizers on the block (\Cref{def:geo-law-coherence}).

\propositionheading{Gauge laws of the rooted datum}{prop:rg-rooted-datum}
The map $p\mapsto\Theta^x_p$ on closed walks at $r_I$ is multiplicative under concatenation and invariant
under insertion or removal of backtracks, so it is a homomorphism on the fundamental group of
$\mathfrak I_I$ at $r_I$, with image $\mathrm{Hol}^x_I$ generated by the fundamental-cycle holonomies
$\mathrm{Hol}^x(p_e)$. Under a passive re-expression $g$ the tree transports, this homomorphism, and, for a
boundary copy $e:j\to i$ from $J$ to $I$, the link and the dressed edge transform as
\begin{align}
&\Theta^x_{I\leftarrow i}\longmapsto g_{r_I}\Theta^x_{I\leftarrow i}g_i^{-1},
\label{eq:rg-tree-transport-law}\\
&\pi_1(\mathfrak I_I,r_I)\longrightarrow G,\quad p\longmapsto\Theta^x_p,
\qquad
\Theta^x_p\longmapsto g_{r_I}\Theta^x_pg_{r_I}^{-1},
\label{eq:rg-full-holonomy-representation}\\
&\Theta^x_e\longmapsto g_i\Theta^x_eg_j^{-1},
\qquad
\widehat\Theta^x_e\longmapsto g_{r_I}\widehat\Theta^x_eg_{r_J}^{-1},
\label{eq:rg-dressed-boundary-edge}
\end{align}
and $\bar z^x_I\mapsto\rho_x(g_{r_I})\bar z^x_I$, one root element acting on every entry. The datum
consisting of $\bar z^x_I$, the homomorphism \eqref{eq:rg-full-holonomy-representation}, and the dressed
edges is therefore root framed and equivariant under one root re-expression per block. The homomorphism
does not depend on the trees, and a change of tree multiplies each tree transport on the left by a based
holonomy. The tree links can be set to the identity by a passive re-expression unique up to one element
per block, after which each non-tree internal link equals the conjugated fundamental-cycle holonomy it
closes. \status{ESTABLISHED}

\noindent\emph{Proof.}
Substitution cancels interior factors of path products and gives the law of $\bar z^x_I$, and a backtrack
contributes $\Theta^x_e\Theta^x_{\bar e}=e$. Every internal copy satisfies
$\Theta^x_e=(\Theta^x_{I\leftarrow i})^{-1}\mathrm{Hol}^x(p_e)\Theta^x_{I\leftarrow j}$, with $\mathrm{Hol}^x(p_e)=e$
on tree edges, so the product along a closed walk at $r_I$ telescopes to a product of fundamental-cycle
holonomies and their inverses; walks and holonomy do not involve the trees, and for a second tree
$\widetilde\Theta^x_{I\leftarrow i}(\Theta^x_{I\leftarrow i})^{-1}$ is the product along a closed walk at
$r_I$. Finally $g_i=g_{r_I}\Theta^x_{I\leftarrow i}$ makes every tree transport, hence every tree link, the
identity, is forced up to the choice of $g_{r_I}$, and turns a non-tree link into
$g_{r_I}\mathrm{Hol}^x(p_e)g_{r_I}^{-1}$. $\square$

The last statement is the maximal-tree gauge fixing of lattice gauge theory \citep{Creutz1977}, in
the spanning-tree form of \citet[Sec.~IV A, Eq.~(4.1)]{DelcampDittrich2017PhaseDiagram}. The
measure tier retains the raw root-framed datum, since quotienting each holonomy separately would lose its
orientation relative to the root features and boundary legs, and an orbit space would need a quotient
theorem not supplied here.

\propositionheading{Two full-law treatments of holonomy for a candidate parent}{prop:rg-parent-holonomy-alternatives}
In the setting of \Cref{thm:rg-pointwise-parent-datum} two alternatives are available, which need not
exclude each other at different retained-coordinate or quotient levels. (i) Groupoid-covariant parent:
supply bimeasurable actions of a typed groupoid on the fine random space, on $\mathsf Z_A$, on the
observation space, and on the structural data; let the fine generative and recognition laws be covariant
and the selected posterior versions compatible with the actions; and require
$C_A(g\cdot y,g\cdot D)=C_A(y,D)$ for every admitted arrow $g$ and measurable $D\subseteq\mathsf Z_A$,
together with covariance of the evaluator for the actions on $(m_A,\xi_A)$ and on $(b_A,o,h_A)$. Then the
parent generative, posterior, and recognition laws are covariant. Invariance at one slice $(o,X)$ holds for
the arrows that fix $(o,X)$ and the selected posterior versions; invariance of $q^b_A$ and $q^m_A$
separately is only a projection consequence and does not replace these hypotheses. (ii) Retention: the
component roots, the based holonomy homomorphisms, and the dressed boundary edges of
\Cref{prop:rg-rooted-datum} are placed in $\mathsf H_A$, and their joint law is pushed forward with no
invariance claim. Neither alternative selects $I$, $A$, a partition, or a membership kernel.
\status{ESTABLISHED}

\noindent\emph{Proof.}
(i) Substituting the invariance of $C_A$ and the covariant fine laws into
\eqref{eq:rg-pointwise-parent-generative} and \eqref{eq:rg-pointwise-parent-triple} gives covariance of the
three parent laws, and evaluator covariance carries it to the evaluator tier; at a fixed slice only arrows
fixing $(o,X)$ and the versions map the slice to itself. (ii) defines the retained coordinates, and the
pushforward is that of \Cref{thm:rg-pointwise-parent-datum}. $\square$

\section{Linear feature tier and tree dependence}

\propositionheading{Linear cross-scale maps and their tree dependence}{prop:rg-linear-cross-scale}
Let $\rho^{\mathbb W}_x:G\to\GL(\mathbb W_x)$ be a linear representation on a feature fiber, choose weights $w_{Ii}\geq0$ with
$\sum_{i\in I}w_{Ii}=1$, and set
\begin{equation}
(\mathsf C_xz)_I=\sum_{i\in I}w_{Ii}\rho^{\mathbb W}_x(\Theta^x_{I\leftarrow i})z_i,
\qquad
(\mathsf S_x\hat z)_i=\rho^{\mathbb W}_x(\Theta^x_{I\leftarrow i})^{-1}\hat z_I
\label{eq:rg-linear-cross-scale-maps}
\end{equation}
for a fine feature $z$ and a coarse feature $\hat z=(\hat z_I)_I$, with
$\rho^{\mathbb W}_{x,f}(g)=\bigoplus_i\rho^{\mathbb W}_x(g_i)$ and $\rho^{\mathbb W}_{x,c}(g)=\bigoplus_I\rho^{\mathbb W}_x(g_{r_I})$. (i)
$\mathsf C_x\mathsf S_x=I$, the averaged feature $(\mathsf C_xz)_I$ is a weighted mean of the entries of the
root-framed tuple, which it compresses, and
$\mathsf C'_x=\rho^{\mathbb W}_{x,c}(g)\mathsf C_x\rho^{\mathbb W}_{x,f}(g)^{-1}$,
$\mathsf S'_x=\rho^{\mathbb W}_{x,f}(g)\mathsf S_x\rho^{\mathbb W}_{x,c}(g)^{-1}$. (ii) At nested depths $0\to1\to2$, with
$\Theta^{x,12}_{A\leftarrow I}$ the tree transport of the coarse vertex $I$ to the root of its depth-two
block $A$, the conditions
$\Theta^{x,02}_{A\leftarrow i}=\Theta^{x,12}_{A\leftarrow I}\Theta^{x,01}_{I\leftarrow i}$ and
$w^{02}_{Ai}=w^{12}_{AI}w^{01}_{Ii}$ give $\mathsf C^{02}_x=\mathsf C^{12}_x\mathsf C^{01}_x$ and
$\mathsf S^{02}_x=\mathsf S^{01}_x\mathsf S^{12}_x$; nested forests alone give neither equality. (iii) For a
fine message operator $\mathcal A_x$ with $\mathcal A'_x=\rho^{\mathbb W}_{x,f}(g)\mathcal A_x\rho^{\mathbb W}_{x,f}(g)^{-1}$, the
fine-to-meta, meta-to-fine, and meta-to-meta operators
$\mathcal A^x_{I\leftarrow j}=(\mathsf C_x\mathcal A_x)_{Ij}$,
$\mathcal A^x_{j\leftarrow I}=(\mathcal A_x\mathsf S_x)_{jI}$, and
$\mathcal A^x_{I\leftarrow J}=(\mathsf C_x\mathcal A_x\mathsf S_x)_{IJ}$ transform between the corresponding
endpoint frames, and for a linear cross map $\Phi_f:\mathbb W_{b,f}\to \mathbb W_{m,f}$ with
$\Phi'_f=\rho^{\mathbb W}_{m,f}(g)\Phi_f\rho^{\mathbb W}_{b,f}(g)^{-1}$ the compressed map $\Phi^c=\mathsf C_m\Phi_f\mathsf S_b$
satisfies $\Phi^{c\prime}=\rho^{\mathbb W}_{m,c}(g)\Phi^c\rho^{\mathbb W}_{b,c}(g)^{-1}$. (iv) A second family of trees gives
transports $\widetilde\Theta^x_{I\leftarrow i}=\mathrm{Hol}^x(\gamma_i)\Theta^x_{I\leftarrow i}$, with
$\gamma_i$ the closed walk from $r_I$ to $i$ along the first tree and back along the second, and
\begin{equation}
\begin{gathered}
(\widetilde{\mathsf C}_xz)_I-(\mathsf C_xz)_I
=\sum_{i\in I}w_{Ii}\bigl(\rho^{\mathbb W}_x(\mathrm{Hol}^x(\gamma_i))-I\bigr)\rho^{\mathbb W}_x(\Theta^x_{I\leftarrow i})z_i,\\
(\widetilde{\mathsf S}_x\hat z)_i=\rho^{\mathbb W}_x(\Theta^x_{I\leftarrow i})^{-1}\rho^{\mathbb W}_x(\mathrm{Hol}^x(\gamma_i))^{-1}\hat z_I.
\end{gathered}
\label{eq:rg-linear-cross-scale-tree-change}
\end{equation}
The difference is not a root re-expression, so $\mathsf C_x$, $\mathsf S_x$, and the message operators
depend on the trees through gauge-invariant quantities: for two vertices joined by two edge copies with
$\mathrm{SO}(2)$ links the identity and the rotation by $\vartheta$, the defining representation, equal
weights, equal features $e_1$, and root at the first vertex, the two trees give averaged features of norms
$1$ and $\lvert\cos(\vartheta/2)\rvert$, although these norms are invariant under every root
re-expression. Tree independence holds for inputs whose root-framed features lie in the fixed space
$\operatorname{Fix}(\rho^{\mathbb W}_x\circ\mathrm{Hol}^x_I)=\{v\in \mathbb W_x:\rho^{\mathbb W}_x(\mathsf k)v=v\text{ for every }\mathsf k\in\mathrm{Hol}^x_I\}$,
for $\mathsf S_x$ on coarse features in that subspace, and for all inputs when
$\rho^{\mathbb W}_x\circ\mathrm{Hol}^x_I$ is trivial. \status{ESTABLISHED}

\noindent\emph{Proof.}
(i) $(\mathsf C_x\mathsf S_x\hat z)_I=\sum_iw_{Ii}\hat z_I=\hat z_I$, and substituting
\eqref{eq:rg-tree-transport-law} gives the covariance laws. (ii) Substitute the two conditions into the
composites. (iii) follows from (i). (iv) The closed walk $\gamma_i$ has path product
$\widetilde\Theta^x_{I\leftarrow i}(\Theta^x_{I\leftarrow i})^{-1}$, which gives
\eqref{eq:rg-linear-cross-scale-tree-change}. In the example the two trees give $e_1$ and
$\tfrac12(e_1+R_\vartheta e_1)$, of norm $\lvert\cos(\vartheta/2)\rvert$. If every
$\rho^{\mathbb W}_x(\Theta^x_{I\leftarrow i})z_i$ is fixed by the represented holonomy, the difference vanishes, and
likewise for $\mathsf S_x$. $\square$

This fixed space is the feature-level counterpart of \eqref{eq:full-graph-holonomy-stabilization}; a tree
choice is a gauge fixing for the retained datum but part of the model for the averaged feature
$\mathsf C_xz$.

\propositionheading{Composite trees and one rooted forest compose the tree transports}{prop:rg-forest-tree-composition}
At nested depths $0\to1\to2$ as in \Cref{prop:rg-linear-cross-scale}(ii), fix a channel $x$ with links
$\Theta^x_e\in G$, let the blocks $I$ of depth one partition $V$, let the depth-one multigraph $\mathfrak I_1$
of \Cref{def:full-graph-coarse-links} carry the dressed marks $\Theta^{x,1}_e=\widehat\Theta^x_e$ as links,
and let the blocks $A$ of depth two partition the vertices of $\mathfrak I_1$, every block connected; $A$ also
denotes the set of agents of its children $I\subseteq A$, with internal multigraph $\mathfrak I_A$ in
$\mathfrak I$. The transports $\Theta^{x,01}_{I\leftarrow i}$, $\Theta^{x,12}_{A\leftarrow I}$, and
$\Theta^{x,02}_{A\leftarrow i}$ are taken along rooted spanning trees of $\mathfrak I_I$, of the internal
multigraph of $A$ in $\mathfrak I_1$, and of $\mathfrak I_A$. (i) For any such trees of the children, with
roots $r_I$, and of $A$ in $\mathfrak I_1$, with root child $I_A$, the copies of these trees form a spanning
tree of $\mathfrak I_A$, the composite tree of $A$. With the composite tree and the root $r_{I_A}$, every agent
$i\in I\subseteq A$ and every copy $e:j\to i$ with $j\in J\subseteq A'$ satisfy
\begin{equation}
\begin{gathered}
\Theta^{x,02}_{A\leftarrow i}=\Theta^{x,12}_{A\leftarrow I}\Theta^{x,01}_{I\leftarrow i},\\
\Theta^{x,12}_{A\leftarrow I}\widehat\Theta^x_e\bigl(\Theta^{x,12}_{A'\leftarrow J}\bigr)^{-1}
=\Theta^{x,02}_{A\leftarrow i}\Theta^x_e\bigl(\Theta^{x,02}_{A'\leftarrow j}\bigr)^{-1},
\end{gathered}
\label{eq:rg-forest-tree-composition}
\end{equation}
so the depth-two link obtained by lifting twice with \eqref{eq:full-graph-coarse-links} is the one-step lift
through the composite trees: for $A'\neq A$ a dressed edge \eqref{eq:full-graph-dressed-edge}, and for $A'=A$
the fundamental-cycle holonomy $\mathrm{Hol}^x(p_e)$ of \eqref{eq:full-graph-tree-transport-law} for the
composite tree, which is $e_G$ on its copies. The transport condition of
\Cref{prop:rg-linear-cross-scale}(ii) therefore holds, and with the weights
$w^{02}_{Ai}=w^{12}_{AI}w^{01}_{Ii}$, which sum to one over $A$, both equalities of that part hold. (ii)
Choose, before the partitions and independently of them, a rooted spanning forest of $\mathfrak I$, that is,
a rooted spanning tree of each connected component of $\mathfrak I$, and let every block be a connected
subtree of it, meaning that the forest copies joining agents of the block connect the block. Each block then
has exactly one agent nearest to the root of its component, measured by the number of copies on the forest
path; call it $r_I$ for a child and $r_A$ for a block of depth two. Give each child the forest copies joining
its agents, rooted at $r_I$, and give $A$ in $\mathfrak I_1$ the forest copies joining distinct children of
$A$, rooted at the child $I_A$ that contains $r_A$. These are rooted spanning trees, $r_{I_A}=r_A$, and the
composite tree of $A$ consists of the forest copies joining agents of $A$, so (i) applies to the forest tree
of $A$ rooted at $r_A$, and
\begin{equation}
\Theta^{x,12}_{A\leftarrow I}=\Theta^{x,02}_{A\leftarrow r_I},
\qquad
\Theta^{x,02}_{A\leftarrow i}=\Theta^{x,02}_{A\leftarrow r_I}\Theta^{x,01}_{I\leftarrow i}.
\label{eq:rg-forest-root-factorization}
\end{equation}
For a longer chain of nested partitions whose blocks are all connected subtrees of the forest, the trees
induced at every depth $\ell$, with the agents as the blocks of depth zero, formed by the forest copies
inside a block of depth $\ell+1$ that join distinct blocks of depth $\ell$ and rooted at the block of depth
$\ell$ that contains its nearest agent, satisfy the hypotheses of (i) at the depths $\ell\to\ell+1\to\ell+2$
for every $\ell$. (iii) If the links of channel $x$ are graph flat on every block of depth two
(\Cref{def:geo-graph-flat-links}), in particular if they are frame aligned,
\eqref{eq:rg-forest-tree-composition} holds with any rooted spanning tree of $\mathfrak I_A$ with root
$r_{I_A}$ in place of the composite tree. The transports in (i) and (ii) are path products of the
group-valued links, whose reversed copies carry the inverse links; a matrix communication
map $\mathsf L_e$ of \Cref{def:geo-selective-matrix-link}, whose reversed copy is specified separately, is
not a transport, and on the vector tier, where $\rho_x$ acts linearly on the sample fibers, its dressings by
the represented transports of (i) satisfy
\begin{equation}
\begin{aligned}
&\rho_x(\Theta^{x,12}_{A\leftarrow I})\bigl[\rho_x(\Theta^{x,01}_{I\leftarrow i})\mathsf L_e
\rho_x(\Theta^{x,01}_{J\leftarrow j})^{-1}\bigr]\rho_x(\Theta^{x,12}_{A'\leftarrow J})^{-1}\\
&\qquad=\rho_x(\Theta^{x,02}_{A\leftarrow i})\mathsf L_e\rho_x(\Theta^{x,02}_{A'\leftarrow j})^{-1},
\end{aligned}
\label{eq:rg-forest-matrix-dressing}
\end{equation}
which inverts only represented transports, so $\mathsf L_e$ may be singular. \status{ESTABLISHED}

\noindent\emph{Proof.}
Deleting a backtrack leaves a path product unchanged, since $\Theta^x_e\Theta^x_{\bar e}=e_G$, and a walk in
a tree reduces in this way to the tree path between its endpoints, because a closed walk of positive length
in a tree backtracks at a vertex farthest from its base point. The argument for (i) uses only that the links
are elements of $G$ whose reversed copies carry the inverses, which holds at every depth by
\Cref{prop:full-graph-coarse-links}. (i) Counting reversal pairs, the trees of the children contain
$\lvert A\rvert$ less the number of children, each pair joining agents of one child, and the tree of $A$ in
$\mathfrak I_1$ contains the number of children less one, each pair joining distinct children; together
they connect the agents of $A$ with $\lvert A\rvert-1$ pairs, so they form a spanning tree. Let the tree path
$p$ from $I$ to $I_A$ in $\mathfrak I_1$ traverse the copies $e_a:j_a\to i_a$ from $I_{a-1}$ to $I_a$ for
$a=1,\ldots,\lvert p\rvert$, oriented along the path, with $I_0=I$, $I_{\lvert p\rvert}=I_A$, and $i_0=i$.
By \eqref{eq:full-graph-coarse-links},
\[
\Theta^{x,12}_{A\leftarrow I}\Theta^{x,01}_{I\leftarrow i}
=\Theta^{x,01}_{I_{\lvert p\rvert}\leftarrow i_{\lvert p\rvert}}\Theta^x_{e_{\lvert p\rvert}}
\bigl[(\Theta^{x,01}_{I_{\lvert p\rvert-1}\leftarrow j_{\lvert p\rvert}})^{-1}
\Theta^{x,01}_{I_{\lvert p\rvert-1}\leftarrow i_{\lvert p\rvert-1}}\bigr]
\Theta^x_{e_{\lvert p\rvert-1}}\cdots\Theta^x_{e_1}
\bigl[(\Theta^{x,01}_{I_0\leftarrow j_1})^{-1}\Theta^{x,01}_{I_0\leftarrow i_0}\bigr].
\]
Each bracket is the path product of a walk in the tree of $I_a$ from $i_a$ through $r_{I_a}$ to $j_{a+1}$, and
the leftmost factor is that of the tree path from $i_{\lvert p\rvert}$ to $r_{I_A}$, so the right side is the
path product of a walk in the composite tree from $i$ to $r_{I_A}$, which is $\Theta^{x,02}_{A\leftarrow i}$.
Substituting this at both endpoints of
$\widehat\Theta^x_e=\Theta^{x,01}_{I\leftarrow i}\Theta^x_e(\Theta^{x,01}_{J\leftarrow j})^{-1}$ gives the
second equality, whose right side is \eqref{eq:full-graph-dressed-edge} for $A'\neq A$ and
\eqref{eq:full-graph-tree-transport-law} for $A'=A$, equal to $e_G$ on tree copies as in
\Cref{prop:full-graph-coarse-links}. The weights are nonnegative with
$\sum_{I\subseteq A}w^{12}_{AI}\sum_{i\in I}w^{01}_{Ii}=1$, so \Cref{prop:rg-linear-cross-scale}(ii)
applies. (ii) If two distinct agents of a block were nearest to the root, the forest path between them would
lie in the block, since the forest copies inside the block connect it and paths in a tree are unique, and it
would pass through their last common ancestor, which is nearer to the root. The forest copies inside a block
connect it and contain no cycle, so they form a spanning tree of its internal multigraph. Inside $A$ lie
$\lvert A\rvert-1$ forest pairs, of which $\lvert A\rvert$ less the number of children lie inside children,
so the number of children less one join distinct children; these pairs connect the children,
because the forest copies inside $A$ connect $A$, so they form a spanning tree of $A$ in $\mathfrak I_1$, and
with the trees of the children they make up the forest copies inside $A$. The agent $r_A$ lies in
$I_A\subseteq A$ and is nearest in $A$, hence in $I_A$, so $r_{I_A}=r_A$. Then (i) applies; putting $i=r_I$,
where $\Theta^{x,01}_{I\leftarrow r_I}=e_G$, gives the first equality of
\eqref{eq:rg-forest-root-factorization}, and substituting it into (i) gives the second. For a longer chain the
same count shows that each induced tree is a spanning tree. For a block of depth $\ell+2$, its induced tree at
depth $\ell$ is the union of the induced trees at depth $\ell$ of its children and of the forest copies
joining distinct children, which form its induced tree at depth $\ell+1$, and the block of depth $\ell$ that
contains its nearest agent is the root at depth $\ell$ of the child containing that agent, whose nearest agent
it is. (iii) Graph-flat links have path product $e_G$ on every closed walk in the block
(\Cref{prop:geo-trivializing-criterion}), and another tree with the same root multiplies each transport on
the left by a based holonomy (\Cref{prop:rg-rooted-datum}), which leaves $\Theta^{x,02}_{A\leftarrow i}$
unchanged; frame-aligned links are graph flat (\Cref{rem:geo-flatness-strictness}). The matrix identity
follows from the first equality of \eqref{eq:rg-forest-tree-composition} at $i$ and at $j$ and the
multiplicativity of $\rho_x$. $\square$

Part (i) supplies what nesting alone leaves open in \Cref{prop:rg-linear-cross-scale}(ii): a depth-zero
tree of $A$ that contains the depth-one tree and is rooted at the root of the root child; another tree with
that root multiplies the direct transports $\Theta^{x,02}_{A\leftarrow i}$ by based holonomies
(\Cref{prop:rg-rooted-datum}). In (i) the trees may be chosen after each partition is realized, and the
composite tree then depends on the intermediate partition. In (ii) the forest is part of the blocking scheme
and is chosen before the partitions, so the direct transports of $A$ depend only on $A$ and the forest; a
random membership mark $R_\ell$ (\Cref{def:full-graph-membership-kernel}) enters (ii) only through realized
partitions whose blocks are connected subtrees of that forest, a realized block that is not such a subtree
is outside (ii), and overlapping membership kernels are outside the proposition, as in
\Cref{def:full-graph-coarse-links}.

Setting the tree links to the identity is the spanning-tree gauge fixing of lattice gauge theory in
the form of \citet[Sec.~IV A, Eq.~(4.1)]{DelcampDittrich2017PhaseDiagram}, whose change of spanning
tree at a fixed root is their Sec.~IV C, Eq.~(4.4); the coarse link as an ordered product of fine links
along a composite path is the macroscopic holonomy of
\citet[Sec.~4, Eq.~(4.1)]{Bahr2012OperatorSpinFoams}. Neither source composes trees across nested
partitions.

\section{Holonomy-conditioned marginal-law modes}
\label{sec:cg-holonomy-law-modes}

\definitionheading{Holonomy-conditioned marginal-law score}{def:cg-holonomy-law-score}
Fix a channel $x$, a connected block $I$ with root $r$, weights $w_i>0$ with $\sum_{i\in I}w_i=1$, and
marginal laws $q^x_i$ of the agents of $I$ on the sample fiber $\mathsf Z^x$, on which $G$ acts by
bimeasurable bijections $\rho_x(g)$ with pushforward $\widehat\rho_x(g)$. Let parent-law families
$\mathscr M^x_i\subseteq\operatorname{Prob}(\mathsf Z^x)$ be equivariant:
$\widehat\rho_x(\Theta^x_p)\mathscr M^x_j=\mathscr M^x_i$ for every walk $p$ of $\mathfrak I_I$ from $j$ to
$i$, and $\mathscr M^{x\prime}_i=\widehat\rho_x(g_i)\mathscr M^x_i$ under a passive re-expression. For walks
$p_i$ from $i$ to $r$ put $\bar q^x_i=\widehat\rho_x(\Theta^x_{p_i})q^x_i$, and with the based holonomy
group $\mathrm{Hol}^x_I(r)$ of \eqref{eq:full-graph-based-holonomy} define the fixed parent family and the
score
\begin{align}
\mathscr Q^x_{I,\mathrm{fix}}(r)&=\bigl\{Q\in\mathscr M^x_r:\widehat\rho_x(\mathsf k)Q=Q
\text{ for every }\mathsf k\in\mathrm{Hol}^x_I(r)\bigr\},
\label{eq:cg-holonomy-fixed-parent}\\
\mathfrak D^x_I&=\inf_{Q\in\mathscr Q^x_{I,\mathrm{fix}}(r)}\sum_{i\in I}w_i\KL(\bar q^x_i\Vert Q),
\qquad\inf\varnothing:=+\infty.
\label{eq:cg-holonomy-kl-score}
\end{align}
A minimizer, when one exists, is a holonomy-conditioned marginal-law mode.

\theoremheading{Invariant holonomy-conditioned marginal-law mode}{thm:cg-holonomy-kl-marginal}
The score $\mathfrak D^x_I$ does not depend on the walks, the root, or the passive frames. If the infimum
is attained, then $\mathfrak D^x_I=0$ exactly when some $Q_\star\in\mathscr Q^x_{I,\mathrm{fix}}(r)$
satisfies $\bar q^x_i=Q_\star$ for every $i\in I$, and the laws $\widehat\rho_x(\Theta^x_{p_i})^{-1}Q_\star$
then form a holonomy-stabilized parallel marginal-law section. If $\mathscr Q^x_{I,\mathrm{fix}}(r)=\varnothing$,
then $\mathfrak D^x_I=+\infty$. If $\mathfrak D^x_I=0$ without attainment, the conclusion is only forward-KL
approximability by fixed parent laws. \status{ESTABLISHED}

\noindent\emph{Proof.}
Relative entropy is invariant under bimeasurable bijections \eqref{eq:cg-kl-bimeasurable-invariance}.
Another walk $p'_i$ from $i$ to $r$ replaces $\bar q^x_i$ by $\widehat\rho_x(\mathsf k)\bar q^x_i$ with
$\mathsf k=\Theta^x_{p'_i}(\Theta^x_{p_i})^{-1}\in\mathrm{Hol}^x_I(r)$, and for every fixed parent $Q$,
$\KL(\widehat\rho_x(\mathsf k)\bar q^x_i\Vert Q)=\KL(\widehat\rho_x(\mathsf k)\bar q^x_i\Vert\widehat\rho_x(\mathsf k)Q)=\KL(\bar q^x_i\Vert Q)$,
so every objective value is path independent. A walk $\delta$ from $r$ to $r'$ conjugates
$\mathrm{Hol}^x_I(r)$ onto $\mathrm{Hol}^x_I(r')$, and pushforward by $\Theta^x_\delta$ maps the two fixed
families bijectively, by equivariance, and the transported laws of $r$ to those of $r'$, preserving every
term. A passive re-expression sends $\Theta^x_p$ to $g_i\Theta^x_pg_j^{-1}$ for $p$ from $j$ to $i$ and
$q^x_i$ to $\widehat\rho_x(g_i)q^x_i$, so it pushes every transported law by $\widehat\rho_x(g_r)$,
conjugates the holonomy group by $g_r$, and carries the fixed family bijectively, and invariance of
relative entropy gives gauge invariance. Relative entropy vanishes exactly between equal laws and every
$w_i$ is positive, which gives the zero case; since $Q_\star$ is holonomy fixed, the reconstructed laws do
not depend on the walks and are law coherent by \Cref{prop:geo-three-agreements}. The empty case is the
convention, and without attainment the infimum supplies $Q_n$ with score below $1/n$ and no minimizer.
$\square$

\propositionheading{Compact holonomy: closed form and attainment}{prop:cg-holonomy-compact-closed-form}
In the setting of \Cref{thm:cg-holonomy-kl-marginal}, let the represented holonomy group
$\mathfrak H_r=\rho_x(\mathrm{Hol}^x_I(r))$ be a compact metrizable group acting jointly measurably on
$\mathsf Z^x$, with Haar probability $d\mathsf k$, and put
\begin{align}
\bar P&=\sum_{i\in I}w_i\bar q^x_i,
\qquad
\bar M(B)=\int_{\mathfrak H_r}\bar P(\mathsf k^{-1}B)d\mathsf k.
\label{eq:cg-haar-symmetrized-mixture}
\intertext{Then $\bar M$ is $\mathfrak H_r$-invariant, and every $\mathfrak H_r$-invariant law $Q$ satisfies,
in $[0,+\infty]$,}
\sum_{i\in I}w_i\KL(\bar q^x_i\Vert Q)&=\sum_{i\in I}w_i\KL(\bar q^x_i\Vert\bar M)+\KL(\bar M\Vert Q).
\label{eq:cg-holonomy-compensation}
\end{align}
Consequently $\mathfrak D^x_I=\sum_iw_i\KL(\bar q^x_i\Vert\bar M)
+\inf_{Q\in\mathscr Q^x_{I,\mathrm{fix}}(r)}\KL(\bar M\Vert Q)$. If $\bar M\in\mathscr M^x_r$, for instance
when $\mathscr M^x_r=\operatorname{Prob}(\mathsf Z^x)$, the infimum is attained at $\bar M$, the score is the
holonomy-symmetrized information radius $\sum_iw_i\KL(\bar q^x_i\Vert\bar M)$, and $\bar M$ is the unique
minimizer when that value is finite. If instead $\mathsf Z^x=\R^K$, the holonomy acts linearly,
$\mathscr M^x_r$ is the nondegenerate Gaussian family, and $\bar M$ has finite second moments, mean
$\bar m$, nonsingular covariance $\bar S$, and $\KL(\bar M\Vert\mathcal N(\bar m,\bar S))<\infty$, then
$\mathcal N(\bar m,\bar S)$ is holonomy invariant and attains the infimum, uniquely when
$\mathfrak D^x_I<\infty$. For a finite holonomy group $d\mathsf k$ is uniform; a holonomy group that is only
relatively compact in a group acting jointly continuously on a Polish fiber has the same fixed laws as its
closure, to which the statement applies; and when every $\bar q^x_i$ is a nondegenerate Gaussian and the
holonomy acts linearly with compact closure, the Gaussian hypotheses hold and the score is finite.
\status{ESTABLISHED}

\noindent\emph{Proof.}
Left invariance of $d\mathsf k$ gives invariance of $\bar M$. Since $d\bar q^x_i/d\bar P\leq w_i^{-1}$,
splitting $\log(d\bar q^x_i/dQ)$ through $\bar P$ gives, for every $Q$, the compensation identity
$\sum_iw_i\KL(\bar q^x_i\Vert Q)=\sum_iw_i\KL(\bar q^x_i\Vert\bar P)+\KL(\bar P\Vert Q)$ in $[0,+\infty]$: the
negative part of the second term is integrable because $t(\log t)^-\leq e^{-1}$, and both sides are
infinite unless $\bar P\ll Q$. Applied with $Q$ and with $\bar M$, it reduces
\eqref{eq:cg-holonomy-compensation} to $\KL(\bar P\Vert Q)=\KL(\bar P\Vert\bar M)+\KL(\bar M\Vert Q)$ for
invariant $Q$, where both sides are infinite unless $\bar P\ll Q$. If $\bar P\ll Q$ with density $f$,
invariance of $Q$ gives $\bar M$ the invariant density $\bar f(z)=\int f(\mathsf k^{-1}z)d\mathsf k$; the
invariant set $\{\bar f=0\}$ is null for $\bar M$ and hence for $\bar P$, so $\bar P\ll\bar M\ll Q$. Averaging
the invariant function $\log\bar f$ over the group by Tonelli's theorem gives
$\int\log\bar fd\bar P=\int\log\bar fd\bar M=\KL(\bar M\Vert Q)$ with finite negative parts, and integrating
$\log f=\log(f/\bar f)+\log\bar f$ against $\bar P$ proves the identity. Attainment and uniqueness follow
because relative entropy vanishes only between equal laws. In the Gaussian case, invariance of $\bar M$ and
linearity give $\mathsf k\bar m=\bar m$ and $\mathsf k\bar S\mathsf k^\top=\bar S$, and
$\log[d\mathcal N(\bar m,\bar S)/d\mathcal N(\mu,\Sigma)]$ is a quadratic polynomial whose $\bar M$- and
$\mathcal N(\bar m,\bar S)$-expectations agree, which gives
$\KL(\bar M\Vert\mathcal N(\mu,\Sigma))=\KL(\bar M\Vert\mathcal N(\bar m,\bar S))+\KL(\mathcal N(\bar m,\bar S)\Vert\mathcal N(\mu,\Sigma))$.
For a relatively compact group, $\mathsf k\mapsto\int\phi(\mathsf kz)Q(dz)$ is continuous for bounded
continuous $\phi$ by dominated convergence, so fixed laws of the group and of its closure agree. For
Gaussian constituents with compact linear holonomy, $\bar M$ is a mixture of Gaussians with covariances in
a compact set of positive-definite matrices, so its density is bounded and its second moments are finite,
and convexity of relative entropy in its second argument bounds each $\KL(\bar q^x_i\Vert\bar M)$ by an
average of Gaussian relative entropies continuous on the compact group. $\square$

\corollaryheading{Noncompact Gaussian holonomy leaves no admissible parent}{cor:cg-holonomy-noncompact-gaussian}
Let $\mathsf Z^x=\R^K$, let the holonomy act by invertible linear maps, and let $\mathscr M^x_r$ be the
nondegenerate Gaussian family. Then $\mathscr Q^x_{I,\mathrm{fix}}(r)\neq\varnothing$ if and only if the
represented holonomy group is conjugate into $\mathrm O(K)$. If one loop holonomy has an eigenvalue off the
unit circle, for instance $\lvert\det\rvert\neq1$, or is not diagonalizable, then $\mathfrak D^x_I=+\infty$
for every configuration of marginals. For a block containing a loop with independent $\GL(K)$ links, the
link configurations with $\lvert\det\rvert=1$ on that loop form a Lebesgue-null set, so the score is
$+\infty$ off a null set of configurations; it distinguishes configurations only after the links are
restricted to a group conjugate into $\mathrm O(K)$ or a parent family with a nonempty fixed set is
declared. \status{ESTABLISHED}

\noindent\emph{Proof.}
$\mathcal N(\mu,\Sigma)$ with $\Sigma\succ0$ is fixed by $\mathsf k$ exactly when $\mathsf k\mu=\mu$ and
$\mathsf k\Sigma\mathsf k^\top=\Sigma$. If one such $\Sigma$ is fixed by the whole group, then
$\Sigma^{-1/2}\mathsf k\Sigma^{1/2}$ is orthogonal for every element; conversely, if
$S^{-1}\mathsf kS\in\mathrm O(K)$ for every element, then $\mathcal N(0,SS^\top)$ is fixed. An element
conjugate into $\mathrm O(K)$ is diagonalizable with eigenvalues of modulus one, and
$\det(\mathsf k\Sigma\mathsf k^\top)=(\det\mathsf k)^2\det\Sigma$; an empty fixed family gives $+\infty$ by
convention. The holonomy of a simple loop is a product of link matrices in which one chosen link $\Theta$
enters once, as $A_1\Theta^{\pm1}A_2$ with $A_1,A_2$ fixed by the other links; this is a diffeomorphism of
$\GL(K)$, so the preimage of the null set $\{\lvert\det\rvert=1\}$ is null, and Fubini's theorem over the
remaining links completes the argument. $\square$

With latent links (\Cref{ch:links}), the generators carry a law, and the score of
\Cref{def:cg-holonomy-law-score} becomes a function of them.

\corollaryheading{Latent noncompact links empty the Gaussian parent family almost surely}{cor:link-holonomy-parent-almost-sure}
Let \(G=\GL^+(K)\) with \(K\geq1\) act on the sample fiber \(\mathsf Z^x=\R^K\) of a channel \(x\) with latent
links by its defining representation, let the parent family be the nondegenerate Gaussian family as in
\Cref{cor:cg-holonomy-noncompact-gaussian}, and let \(I\) be a connected block, as in
\Cref{def:cg-holonomy-law-score}, whose internal multigraph
\(\mathfrak I_I\) is not a forest, that is, has a loop copy, two parallel copies, or a cycle through three or
more agents. Then \(\mathfrak I_I\) has a closed walk \(p=(e_0,\ldots,e_{r-1})\) that traverses some reversal
pair a different number of times in its two orientations, and for every such walk the following hold.
(i) For every choice of frames,
\(\det\mathrm{Hol}^x(p)=\exp\bigl(\sum_{a=0}^{r-1}\Tr\mathring\omega^x_{e_a}\bigr)\), with
\(\mathring\omega^x_{\bar e}=-\mathring\omega^x_e\).
(ii) The generator families with \(\sum_a\Tr\mathring\omega^x_{e_a}=0\) form a hyperplane of \(\mathscr T\), which
is Lebesgue-null, and for every generator family off this hyperplane the fixed parent family
\eqref{eq:cg-holonomy-fixed-parent} is empty and \(\mathfrak D^x_I=+\infty\) for every configuration of
marginals.
(iii) Under every law of the generators that is absolutely continuous with respect to Lebesgue measure on
\(\mathscr T\), in particular the generator prior of \Cref{def:link-generator-prior} whenever it has a
density, as it does when its base law is Gaussian with positive-definite covariances, and every recognition
law absolutely continuous with respect to such a prior, \(\mathfrak D^x_I=+\infty\) almost surely, and its
expectation is \(+\infty\).
(iv) A law of the generators under which \(\mathfrak D^x_I\) is finite with positive probability charges the
hyperplane of (ii); the point mass at \(\mathring\omega=0\) is such a law for nondegenerate Gaussian
marginals.
(v) For \(G=\mathrm{SL}(2,\R)\) acting on \(\R^2\) by its defining representation and a block whose internal
multigraph has cycle rank one, the generator families at which the represented holonomy group is conjugate
into \(\mathrm{SO}(2)\) contain a nonempty open subset of \(\mathscr T\); on it the fixed parent family is
nonempty, and every law of the generators with an everywhere positive density gives it positive
probability. \status{ESTABLISHED}

\noindent\emph{Proof.}
If \(\mathfrak I_I\) is a forest, every copy is a bridge, and a closed walk crosses a bridge as often in one
orientation as in the other, because it returns to the side it started from. A loop copy traversed once,
the walk \((e,\bar e')\) through two parallel copies \(e,e'\), and a cycle traversed once each traverse some
reversal pair once in one orientation and never in the other.
(i) By \eqref{eq:geo-edge-twist-holonomy},
\(\mathrm{Hol}^x(p)=U_{i_0}e^{\mathring\omega^x_{e_0}}\cdots e^{\mathring\omega^x_{e_{r-1}}}U_{i_0}^{-1}\), and
\(\det e^{\mathbf B}=e^{\Tr\mathbf B}\) by Jacobi's formula.
(ii) Collecting terms, \(\sum_a\Tr\mathring\omega^x_{e_a}=\sum_{e\in E_+}\mathrm n_e\Tr\mathring\omega^x_e\), where
\(\mathrm n_e\) is the number of traversals of \(e\) minus that of \(\bar e\). Some \(\mathrm n_e\) is nonzero,
and the family with \(\mathring\omega^x_e=I_K\) and every other generator zero gives the value
\(\mathrm n_eK\neq0\), so the functional is a nonzero linear form and its zero set is a hyperplane. Off the
hyperplane \(\lvert\det\mathrm{Hol}^x(p)\rvert\neq1\) by (i). For the root \(r\) of \(I\), the closed walk at
\(r\) that follows a walk from \(r\) to \(i_0\), then \(p\), then the reverse walk has a holonomy conjugate to
\(\mathrm{Hol}^x(p)\), with the same determinant, so a loop holonomy of the block has an eigenvalue off the
unit circle, and \Cref{cor:cg-holonomy-noncompact-gaussian} gives an empty fixed family and
\(\mathfrak D^x_I=+\infty\) for every configuration of marginals.
(iii) Such a law gives the null hyperplane probability zero, and a law absolutely continuous with respect
to one that is absolutely continuous with respect to Lebesgue measure is itself absolutely continuous with
respect to Lebesgue measure. The generator prior \(Z_{\mathscr S}^{-1}e^{-\mathscr S}\nu^{\mathscr T}\) of
\Cref{def:link-generator-prior} is absolutely continuous with respect to its base law, which has a Lebesgue
density when it is Gaussian with positive-definite covariances.
(iv) The first statement is (ii). At \(\mathring\omega=0\) every twist is the identity, so every holonomy is
the identity by \eqref{eq:geo-edge-twist-holonomy}; the trivial group is compact, and
\Cref{prop:cg-holonomy-compact-closed-form} gives a finite score for nondegenerate Gaussian marginals.
(v) With cycle rank one, \(\mathrm{Hol}^x_I\) is generated by one fundamental-cycle holonomy \(\mathsf k\)
(\Cref{prop:rg-rooted-datum}), which by \eqref{eq:geo-edge-twist-holonomy} is conjugate to an ordered product
of exponentials of the generators along the cycle, so \(\Tr\mathsf k\) is a continuous function of
\(\mathring\omega\). At the generator family with \(\frac\pi2\mathfrak J\) on the non-tree copy,
\(\mathfrak J=\begin{psmallmatrix}0&-1\\1&0\end{psmallmatrix}\), and zero elsewhere, \(\mathsf k\) is conjugate to
the rotation by \(\frac\pi2\), so the set where \(\lvert\Tr\mathsf k\rvert<2\) is open and nonempty. There
\(\mathsf k\in\mathrm{SL}(2,\R)\) has the eigenvalues \(e^{\pm\mathrm i\phi}\) with \(\sin\phi\neq0\); the real and
imaginary parts of an eigenvector for \(e^{\mathrm i\phi}\) form a basis of \(\R^2\) in which \(\mathsf k\) acts as
a rotation, so the group generated by \(\mathsf k\) is conjugate into \(\mathrm{SO}(2)\), and
\Cref{cor:cg-holonomy-noncompact-gaussian} gives a nonempty fixed family. A positive density has positive
integral over a nonempty open set. $\square$

The null-set sentence of \Cref{cor:cg-holonomy-noncompact-gaussian} measures link configurations in
\(\GL(K)\). With latent generators the exceptional set is the explicit hyperplane of (ii), linear in the
generators because \(\log\det e^{\mathbf B}=\Tr\mathbf B\), and it depends neither on the frames, nor on the
trace-free components of the generators, nor on whether they commute. The score is then not a statistic of
the recognition law: it is \(+\infty\) under every law with a density, and only a law that charges the
hyperplane, such as a point mass at a held link family, can make it finite. By (v), on
\(\mathrm{SL}(2,\R)\), which has no determinant direction, a set of generator families of positive
probability keeps a nonempty fixed family.

\remarkheading{Holonomy score on a three-agent loop}{rem:cg-holonomy-score-loop}
On the triangle of the running example (\Cref{sec:geo-running-example}), now with belief sample fiber
$\R^2$ and links in $\mathrm{SO}(2)$ acting by rotation, root the block at the first agent, take the tree
edges joining it to the other two, and let the loop holonomy be a rotation by $\phi$ with $\sin\phi\neq0$.
For transported marginals $\bar q^b_i=\mathcal N(m_i,\sigma^2I_2)$ with weights $w_i$ and the nondegenerate
Gaussian parent family, the fixed family is $\{\mathcal N(0,sI_2):s>0\}$ and
$\mathfrak D^b_I=\log(1+\tfrac{1}{2\sigma^2}\sum_iw_i\lvert m_i\rvert^2)$, attained at
$s=\sigma^2+\tfrac12\sum_iw_i\lvert m_i\rvert^2$, the moment-matched Gaussian of the Haar-symmetrized
mixture; rotations preserve norms, so the score does not depend on the walks. With frame-aligned links on
the same triangle the holonomy is trivial, every nondegenerate Gaussian is fixed, and the score is
$\tfrac12\log\det(I_2+\operatorname{Cov}_w(m)/\sigma^2)$, which never exceeds the value above. Replacing one
link by a matrix of determinant $2$ empties the fixed family and makes the score $+\infty$
(\Cref{cor:cg-holonomy-noncompact-gaussian}). \status{ESTABLISHED}

\noindent\emph{Proof.}
A fixed Gaussian has mean fixed by $R_\phi$, hence zero because $R_\phi-I$ is invertible, and a symmetric
covariance commuting with $R_\phi$; the linear conditions on its off-diagonal entry and the difference of
its diagonal entries have determinant $2\sin^2\phi\neq0$, so the covariance is scalar. In dimension two,
$\KL(\mathcal N(m,\sigma^2I_2)\Vert\mathcal N(0,sI_2))=\sigma^2/s+\lvert m\rvert^2/(2s)-1+\log(s/\sigma^2)$,
whose weighted sum is minimized at the stated $s$ with the stated value. In the flat case moment matching
gives $\mathcal N(\bar m,\sigma^2I_2+\operatorname{Cov}_w(m))$ and the stated value. With $c_1,c_2$ the
eigenvalues of $\operatorname{Cov}_w(m)/\sigma^2$ and
$\sum_iw_i\lvert m_i\rvert^2=\lvert\bar m\rvert^2+\operatorname{tr}\operatorname{Cov}_w(m)$, the inequality of
arithmetic and geometric means gives $\tfrac12\log((1+c_1)(1+c_2))\leq\log(1+\tfrac12(c_1+c_2))$, which is at
most the holonomy value. A link of determinant $2$ gives the loop holonomy a determinant of modulus $2$.
$\square$

\definitionheading{Exact common recovery}{def:cg-common-recovery}
A statistical experiment $\{P_\theta\}_{\theta\in\Theta}$ admits exact common recovery through a
parameter-independent Markov kernel $K_\downarrow$ when a parameter-independent Markov kernel $K_\uparrow$
satisfies
\begin{equation}
P_\theta K_\downarrow K_\uparrow=P_\theta\qquad\text{for every }\theta\in\Theta.
\label{eq:cg-holonomy-common-recovery}
\end{equation}

\remarkheading{A zero score is not common recovery}{rem:cg-marginal-not-joint}
The score \eqref{eq:cg-holonomy-kl-score} depends on the laws of the block only through their transported
marginals. When both marginals of a pair of agents are non-Dirac, there are joint laws $P_1\neq P_2$ with
the same marginals (\Cref{prop:prob-marginals-do-not-determine-joint}), and every kernel $K_\downarrow$ for
which $PK_\downarrow$ depends on $P$ only through its marginals sends them to one law, so the experiment
$\{P_1,P_2\}$ admits no common recovery through it. \status{ESTABLISHED}

\noindent\emph{Proof.}
If $P_1K_\downarrow=P_2K_\downarrow$, a kernel $K_\uparrow$ recovering both would give
$P_1=P_1K_\downarrow K_\uparrow=P_2K_\downarrow K_\uparrow=P_2$. $\square$

\corollaryheading{Parallel cross morphisms preserve fixed marginal-law sections}{cor:cg-holonomy-cross-morphism}
Let $\Phi:\mathcal E_b\to\mathcal E_m$ be a cross map, expressed in the frame of agent $i$ by
$\Phi_i:\mathcal B_b\to\mathcal B_m$, whose link-intertwining defects \eqref{eq:rg-cross-morphism-defects}
vanish on the internal copies of the block, so that
$\widehat\rho_m(\Theta^m_p)\circ\Phi_j=\Phi_i\circ\widehat\rho_b(\Theta^b_p)$ for every walk $p$ from $j$ to
$i$, and assume $\Phi_r(\mathscr M^b_r)\subseteq\mathscr M^m_r$. Then
$\Phi_r(\mathscr Q^b_{I,\mathrm{fix}}(r))\subseteq\mathscr Q^m_{I,\mathrm{fix}}(r)$, and applying $\Phi_i$
agentwise to a holonomy-stabilized parallel belief-law section gives such a model-law section. If
$\widetilde\Phi:\mathcal E_m\to\mathcal E_b$ likewise has vanishing link-intertwining defects and
$\widetilde\Phi_r(\mathscr M^m_r)\subseteq\mathscr M^b_r$, then
$\widetilde\Phi_r(\mathscr Q^m_{I,\mathrm{fix}}(r))\subseteq\mathscr Q^b_{I,\mathrm{fix}}(r)$, with the
analogous statement for sections. \status{ESTABLISHED}

\noindent\emph{Proof.}
Vanishing defects on the edges give the intertwining for walks by composition, using the coordinate form of
\Cref{prop:rg-defect-covariance}. If $Q\in\mathscr Q^b_{I,\mathrm{fix}}(r)$ and $p$ is a closed walk at
$r$, then $\widehat\rho_m(\Theta^m_p)\Phi_r(Q)=\Phi_r(\widehat\rho_b(\Theta^b_p)Q)=\Phi_r(Q)$, and family
compatibility gives the first inclusion. If the belief marginals transport to $Q$, the same square shows
that their images transport to $\Phi_r(Q)$. Interchanging the channels and replacing $\Phi$ by
$\widetilde\Phi$ gives the rest. $\square$

\section{Structural fixed sectors}
\label{sec:cg-partial}

\definitionheading{Connection Laplacian and fixed fiber}{def:cg-connection-laplacian}
Fix a channel and a block $I$ with internal multigraph $\mathfrak I_I$, let the channel's sample fiber be
a real vector space $\mathbb V$ of dimension $K$ with a linear representation $\rho$, and write $\Theta_e$ for the
represented link $\rho(\Theta^x_e)\in\GL(\mathbb V)$. For each unoriented internal copy choose an orientation
$e:j\to i$ and a positive-semidefinite form $W_e$ on $\mathbb V$. The connection Laplacian of the block is the
positive-semidefinite form $L_I$ on $\mathbb V^I$ with
\begin{equation}
z^\top L_Iz=\sum_{e:j\to i}(z_i-\Theta_ez_j)^\top W_e(z_i-\Theta_ez_j),
\label{eq:cg-connection-laplacian-energy}
\end{equation}
in which loops and parallel copies each contribute one term. Its fixed fiber and rank are
$\mathbb V_{I,\mathrm{fix}}=\ker L_I$ and $f_I=\dim\mathbb V_{I,\mathrm{fix}}$. For a root $r$ the represented based holonomy group
$\mathfrak H_r=\rho(\mathrm{Hol}^x_I(r))$ has the fixed space
$\operatorname{Fix}(\mathfrak H_r)=\{v\in \mathbb V:\mathsf kv=v\text{ for every }\mathsf k\in\mathfrak H_r\}$.

\propositionheading{Kernel-holonomy fixed-space isomorphism}{prop:cg-kernel-holonomy}
(i) The energy \eqref{eq:cg-connection-laplacian-energy} does not depend on the chosen orientations when
the reversed copy carries $W_{\bar e}=\Theta_e^\top W_e\Theta_e$. For positive-semidefinite weights,
$z\in\ker L_I$ exactly when
\begin{equation}
W_e^{1/2}(z_i-\Theta_ez_j)=0\qquad\text{for every internal copy }e:j\to i,
\label{eq:cg-psd-section-constraint}
\end{equation}
and $\ker L_I$ is the space $H^0$ of global sections of the cellular sheaf with vertex stalks $\mathbb V$ and
restriction maps $W_e^{1/2}$ at $i$ and $W_e^{1/2}\Theta_e$ at $j$, of which $L_I$ is the degree-zero
Laplacian \citep[Sec.~3.2 and Thm.~3.1]{HansenGhrist2019}; for orthogonal transports and scalar weights
$L_I$ is the vector-diffusion operator of \citet[Sec.~3]{SingerWu2012}. A passive re-expression, with the
weights re-expressed as $W_e\mapsto\rho(g_i)^{-\top}W_e\rho(g_i)^{-1}$, carries $\ker L_I$ to
$\rho_f(g)\ker L_I$, with $\rho_f(g)=\bigoplus_{i\in I}\rho(g_i)$. (ii) If $W_e\succ0$ on every internal
copy, then on every connected component $I_\alpha$ with root $r_\alpha$ root evaluation
$z\mapsto z_{r_\alpha}$ is a linear isomorphism
$\ker L_{I_\alpha}\to\operatorname{Fix}(\mathfrak H_{r_\alpha})$, so
$\ker L_I\cong\bigoplus_\alpha\operatorname{Fix}(\mathfrak H_{r_\alpha})$ and
$f_I=\sum_\alpha\dim\operatorname{Fix}(\mathfrak H_{r_\alpha})$. (iii) With positive-definite weights a
connected block has $f_I=K$ exactly when every represented holonomy element is the identity, equivalently
when the represented links are graph flat on $I$ (\Cref{prop:geo-trivializing-criterion}); this is the
full fixed sector, $0<f_I<K$ is a partial fixed sector, and $f_I=0$ leaves no retained degree. A
nonfaithful representation can make the represented holonomy trivial when the principal holonomy is not.
(iv) With positive-semidefinite weights holonomy may survive in weight-null directions: on a triangle with
$W_{12}=\operatorname{diag}(1,0)$, $\Theta_{12}=\operatorname{diag}(1,2)$, $W_{23}=\operatorname{diag}(0,1)$,
$\Theta_{23}=\operatorname{diag}(2,1)$, and $W_{31}=\Theta_{31}=I$, the loop product is $2I$ and fixes no
nonzero vector, while $\ker L_I$ is the two-dimensional space of constant configurations.
\status{ESTABLISHED}

\noindent\emph{Proof.}
(i) Substituting $\Theta_{\bar e}=\Theta_e^{-1}$ and $W_{\bar e}$ reproduces the term of $e$. Since
$L_I\succeq0$, $z\in\ker L_I$ exactly when $z^\top L_Iz=\sum_e\lvert W_e^{1/2}(z_i-\Theta_ez_j)\rvert^2=0$. The
diagonal blocks of $L_I$ are sums of $W_e$ and $\Theta_e^\top W_e\Theta_e$ and its off-diagonal blocks are
$-W_e\Theta_e$, the blocks of the sheaf Laplacian with the stated restriction maps, whose kernel is $H^0$ by
the discrete Hodge theorem. The energy is unchanged under $z_i\mapsto\rho(g_i)z_i$,
$\Theta_e\mapsto\rho(g_i)\Theta_e\rho(g_j)^{-1}$, and the weight law, so
$L'_I=\rho_f(g)^{-\top}L_I\rho_f(g)^{-1}$. (ii) With positive-definite weights, $z\in\ker L_{I_\alpha}$ exactly
when $z_i=\Theta_ez_j$ on every copy. Such a section is determined by its root value, and transporting that
value around a closed walk returns it, so root evaluation is injective with image in
$\operatorname{Fix}(\mathfrak H_{r_\alpha})$. Conversely, for $u\in\operatorname{Fix}(\mathfrak H_{r_\alpha})$
put $z_v=\rho(\Theta^x_{I_\alpha\leftarrow v})^{-1}u$. Tree constraints hold by construction, and for a
non-tree copy $e:j\to i$ the represented fundamental-cycle holonomy
$\rho(\mathrm{Hol}^x(p_e))=\rho(\Theta^x_{I_\alpha\leftarrow i})\Theta_e\rho(\Theta^x_{I_\alpha\leftarrow j})^{-1}$
of \eqref{eq:full-graph-tree-transport-law} fixes $u$, which gives $\Theta_ez_j=z_i$. The Laplacian is the
direct sum over components. (iii) $\operatorname{Fix}(\mathfrak H_r)=\mathbb V$ exactly when every element is the
identity, the fundamental-cycle holonomies generate the group (\Cref{prop:rg-rooted-datum}), and elements
of the kernel of $\rho$ are represented by the identity. (iv) The loop product is
$\operatorname{diag}(1,2)\operatorname{diag}(2,1)=2I$. The copy $31$ forces $z_3=z_1$, the copy $12$ equal
first coordinates of $z_1$ and $z_2$, and the copy $23$ equal second coordinates of $z_2$ and $z_3$, so the
kernel consists of the constant configurations. $\square$

\propositionheading{Fixed-section compression}{prop:cg-fixed-section-compression}
Let $\Lambda$ be a positive-semidefinite form on the fine configuration space, for instance the connection
Laplacian of the network plus positive-semidefinite self terms. For a partition $\mathcal P$ into blocks
with positive-semidefinite internal weights, let $\iota_I$ be a full-column matrix whose range is
$\ker L_I$, and $\iota_{\mathcal P}=\operatorname{blkdiag}_{I\in\mathcal P}\iota_I$. Then
\begin{equation}
L_I\iota_I=0,
\qquad
\Lambda_c=\iota_{\mathcal P}^\top\Lambda\iota_{\mathcal P}\succeq0,
\qquad
\Lambda_c\succ0\Longleftrightarrow\operatorname{range}(\iota_{\mathcal P})\cap\ker\Lambda=\{0\}.
\label{eq:cg-partial-properness}
\end{equation}
Internal edges contribute nothing to $\Lambda_c$, and cut edges become the rectangular endpoint data of
\Cref{thm:cg-rectangular-endpoint-closure}. With positive-definite weights, the tree construction in the
proof of \Cref{prop:cg-kernel-holonomy} extends a basis of root values to one such matrix.
\status{ESTABLISHED}

\noindent\emph{Proof.}
The first identity holds by construction and the second by congruence; $y^\top\Lambda_cy=0$ exactly when
$\iota_{\mathcal P}y\in\ker\Lambda$, and $\iota_{\mathcal P}$ is injective. An internal term vanishes on
$\operatorname{range}\iota_I=\ker L_I$, and a cut term restricted to the ranges is the energy of
\Cref{thm:cg-rectangular-endpoint-closure} with the blocks of $\iota$ as injections. $\square$

\propositionheading{Nested fixed sections compose exactly}{prop:cg-nested-sections-compose}
Let a larger block $A$ be a union of first-level blocks with positive-semidefinite internal weights, and
let $\iota_{A\leftarrow\mathcal P}=\operatorname{blkdiag}_{I\subset A}\iota_I$ and
$L^{(1)}_A=\iota_{A\leftarrow\mathcal P}^\top L_A\iota_{A\leftarrow\mathcal P}$. Then
$\ker L_A=\iota_{A\leftarrow\mathcal P}\ker L^{(1)}_A$, and for every injection
$\iota^{(2)}_A:\mathbb V_{A,\mathrm{fix}}\hookrightarrow\bigoplus_{I\subset A}\mathbb V_{I,\mathrm{fix}}$ with range $\ker L^{(1)}_A$ the fine
fixed-section injection is $\iota^{\mathrm{fine}}_A=\iota_{A\leftarrow\mathcal P}\iota^{(2)}_A$.
\status{ESTABLISHED}

\noindent\emph{Proof.}
A zero mode of $L_A$ has zero energy in every positive-semidefinite edge summand, so its restriction to each
first-level block lies in $\ker L_I$ and it equals $\iota_{A\leftarrow\mathcal P}y$ for some $y$. Since
$L_A\succeq0$, $y\in\ker L^{(1)}_A$ exactly when $y^\top L^{(1)}_Ay=0$, that is, exactly when
$\iota_{A\leftarrow\mathcal P}y\in\ker L_A$, which proves the first identity; substituting the range of
$\iota^{(2)}_A$ gives the second, and associativity of matrix composition iterates the statement.
$\square$

\remarkheading{Structural fixed sector versus marginal law occupation}{rem:cg-structural-vs-marginal}
Relative entropy between the marginals of two agents needs one of them transported into the frame of the
other and depends on the walk: a second walk changes the transported law by a loop-holonomy action, and
the two values agree for every pair of laws only under stabilization; with
$\mathsf k=\operatorname{diag}(1,-1,-1)\in\mathrm{SO}(3)$ and $q_i=q_j=\mathcal N(e_2,I_3)$, the identity walk
gives zero and a walk differing by $\mathsf k$ gives $2$. An unpenalized sum of nonnegative within-block
pairwise divergences vanishes at, and is minimized by, the singleton partition; it contains only marginal
information and does not test common recovery (\Cref{rem:cg-marginal-not-joint}). The structural fixed
sector and the marginal-law score are independent: on a two-vertex tree with identity transport and
$W=I_K$ the represented holonomy is trivial, yet $\KL(\mathcal N(-ae_1,I_K)\Vert\mathcal N(ae_1,I_K))=2a^2$ is
arbitrarily large; on a triangle with $K=3$, unit weights, identity transports on two edges, and closing
transport $\mathsf k$, the represented holonomy is nontrivial and
$\operatorname{Fix}(\mathsf k)=\operatorname{span}(e_1)$, so $f_I=1<K$, while $\mathsf k$ preserves
$\mathcal N(0,\sigma^2I_3)$ and assigning that law to every vertex makes every transported marginal
divergence zero. Transported divergence thus measures the occupation of the available law modes, not the
dimension of the structural fixed sector. \status{ESTABLISHED}

\noindent\emph{Proof.}
For equal covariances $\KL(\mathcal N(m_1,I)\Vert\mathcal N(m_2,I))=\tfrac12\lvert m_1-m_2\rvert^2$, which
gives $\tfrac12\lvert e_2-\mathsf ke_2\rvert^2=2$ and $2a^2$. A sum of nonnegative terms over pairs inside
blocks is zero for singleton blocks. The fixed space of $\mathsf k$ is $\operatorname{span}(e_1)$, and
$\mathsf k\sigma^2I_3\mathsf k^\top=\sigma^2I_3$. $\square$

\section{Cut links and rectangular sheaf closure}
\label{sec:cg-sheaf-closure}

Let $I$ and $J$ be blocks reduced to root coordinates $z_I,z_J\in \mathbb V$, with represented dressed boundary
edges $\widehat\Theta_e\in\GL(\mathbb V)$ (\Cref{def:rg-rooted-tree-data}) and positive-definite weights $W_e$ on
the cut copies $e$ from $J$ to $I$. The cut energy
$\mathcal E_{\mathrm{cut}}(z_I,z_J)=\sum_e(z_I-\widehat\Theta_ez_J)^\top W_e(z_I-\widehat\Theta_ez_J)$
equals $z_I^\top\Lambda_{II}z_I-2z_I^\top\Lambda_{IJ}z_J+z_J^\top\Lambda_{JJ}z_J$ with
$\Lambda_{II}=\sum_eW_e$, $\Lambda_{IJ}=\sum_eW_e\widehat\Theta_e$, and
$\Lambda_{JJ}=\sum_e\widehat\Theta_e^\top W_e\widehat\Theta_e$. Coarse families admit newly generated
positive-semidefinite self terms. The cut energy carries no factor one half, so $W_e$ and the blocks
$\Lambda_{II}$, $\Lambda_{IJ}$, $\Lambda_{JJ}$ are the coefficients of the quadratic form
$\mathcal E_{\mathrm{cut}}$, and a Gaussian law proportional to $\exp(-\mathcal E_{\mathrm{cut}})$ has
precision blocks $2\Lambda_{II}$, $-2\Lambda_{IJ}$, and $2\Lambda_{JJ}$. Supplement G writes every
quadratic energy with the factor one half, so its weights and precision blocks are twice the $W_e$
and $\Lambda$ blocks here. The excess, the admissible interval, and the balanced weight below are
homogeneous of degree one in the weights, and the links they determine of degree zero, so no
conclusion of this section depends on the convention.

\definitionheading{Anchored and no-new-anchor presentations of a cut}{def:cg-no-new-anchor-pair}
An anchored one-link presentation of the cut is a tuple $(\Theta,W,\Lambda^{\mathrm s}_I,\Lambda^{\mathrm s}_J)$
with $W\succ0$, $\Theta\in\GL(\mathbb V)$, and $\Lambda^{\mathrm s}_I,\Lambda^{\mathrm s}_J\succeq0$ such that, for all
$z_I,z_J$,
\begin{equation}
\mathcal E_{\mathrm{cut}}(z_I,z_J)=(z_I-\Theta z_J)^\top W(z_I-\Theta z_J)
+z_I^\top\Lambda^{\mathrm s}_Iz_I+z_J^\top\Lambda^{\mathrm s}_Jz_J.
\label{eq:cg-cut-anchored-link}
\end{equation}
The no-new-anchor pair subfamily consists of the presentations with
$\Lambda^{\mathrm s}_I=\Lambda^{\mathrm s}_J=0$: one link and one weight, with no new self terms.

\lemmaheading{Completed square for a cut}{lem:cg-cut-completed-square}
One has $\Lambda_{JJ}\succ0$ and, with $\bar\Theta=\Lambda_{II}^{-1}\Lambda_{IJ}$,
\begin{align}
\mathcal E_{\mathrm{cut}}(z_I,z_J)&=(z_I-\bar\Theta z_J)^\top\Lambda_{II}(z_I-\bar\Theta z_J)+z_J^\top\Delta z_J,
\label{eq:cg-cut-completed-square}\\
\Delta&=\Lambda_{JJ}-\Lambda_{IJ}^\top\Lambda_{II}^{-1}\Lambda_{IJ}
=\sum_e(\widehat\Theta_e-\bar\Theta)^\top W_e(\widehat\Theta_e-\bar\Theta)\succeq0.
\label{eq:cg-cut-excess}
\end{align}
(i) An anchored presentation exists if and only if $\det\Lambda_{IJ}\neq0$; with links restricted to
$\GL^+(\mathbb V)$ the condition is $\det\Lambda_{IJ}>0$. (ii) When $\det\Lambda_{IJ}\neq0$, the anchored presentations are exactly
$\Theta=W^{-1}\Lambda_{IJ}$, $\Lambda^{\mathrm s}_I=\Lambda_{II}-W$, and
$\Lambda^{\mathrm s}_J=\Lambda_{JJ}-\Lambda_{IJ}^\top W^{-1}\Lambda_{IJ}$, with $W$ in the Loewner interval
\begin{equation}
\mathfrak W_{\mathrm{cut}}=\{W:\Lambda_{IJ}\Lambda_{JJ}^{-1}\Lambda_{IJ}^\top\preceq W\preceq\Lambda_{II}\}
=\{W\succ0:W\preceq\Lambda_{II},\ \Lambda_{IJ}^\top W^{-1}\Lambda_{IJ}\preceq\Lambda_{JJ}\}.
\label{eq:cg-cut-admissible-weights}
\end{equation}
The choice $W=\Lambda_{II}$ gives $\Lambda^{\mathrm s}_I=0$ and $\Lambda^{\mathrm s}_J=\Delta$, and
$W=\Lambda_{IJ}\Lambda_{JJ}^{-1}\Lambda_{IJ}^\top$ gives $\Lambda^{\mathrm s}_J=0$, the reversed completion with
link $\Lambda_{JJ}^{-1}\Lambda_{IJ}^\top$ from $I$ to $J$ and weight $\Lambda_{JJ}$. (iii) When
$\det\Lambda_{IJ}\neq0$, the anchored presentation is unique if and only if $\Delta=0$, and it is then the
bare pair $(\bar\Theta,\Lambda_{II},0,0)$; otherwise there are infinitely many. The condition $\Delta=0$
itself forces $\det\Lambda_{IJ}\neq0$. (iv) The no-new-anchor pair subfamily contains a presentation of the
cut if and only if $\Delta=0$, that is, if and only if all cut twists $\widehat\Theta_e$ coincide. (v) The
weighted mean $\bar\Theta$ has the gauge law of a dressed edge,
$\bar\Theta\mapsto\rho(g_{r_I})\bar\Theta\rho(g_{r_J})^{-1}$, but it can be singular: equal weights with
twists $I$ and $-I$ give $\bar\Theta=0$. \status{ESTABLISHED}

\noindent\emph{Proof.}
Each $\widehat\Theta_e$ is invertible and $W_e\succ0$, so $\Lambda_{JJ}\succ0$. Completing the square in
$z_I$ gives \eqref{eq:cg-cut-completed-square}, and expanding the sum in \eqref{eq:cg-cut-excess} gives
$\Lambda_{JJ}-\Lambda_{IJ}^\top\Lambda_{II}^{-1}\Lambda_{IJ}$. Matching the three blocks of
\eqref{eq:cg-cut-anchored-link} gives $W\Theta=\Lambda_{IJ}$, $W+\Lambda^{\mathrm s}_I=\Lambda_{II}$, and
$\Theta^\top W\Theta+\Lambda^{\mathrm s}_J=\Lambda_{JJ}$, hence the formulas of (ii); both self terms are
positive semidefinite exactly on the second description in \eqref{eq:cg-cut-admissible-weights}. The link
is invertible exactly when $\Lambda_{IJ}$ is, and $\det\Theta=\det\Lambda_{IJ}/\det W$ has its sign. If
$\det\Lambda_{IJ}\neq0$, then $\Lambda_{IJ}^\top W^{-1}\Lambda_{IJ}\preceq\Lambda_{JJ}$ is equivalent to
$W^{-1}\preceq\Lambda_{IJ}^{-\top}\Lambda_{JJ}\Lambda_{IJ}^{-1}$ and, since inversion reverses the Loewner
order on positive-definite matrices, to $W\succeq\Lambda_{IJ}\Lambda_{JJ}^{-1}\Lambda_{IJ}^\top$; this is the
first description, which contains $W=\Lambda_{II}$ because $\Delta\succeq0$, and proves (i). The reversed
completion is the same computation with $z_I$ and $z_J$ exchanged. Expanding the right side shows
$\Lambda_{II}-\Lambda_{IJ}\Lambda_{JJ}^{-1}\Lambda_{IJ}^\top=\Lambda_{II}\Lambda_{IJ}^{-\top}\Delta\Lambda_{JJ}^{-1}\Lambda_{IJ}^\top$,
so the interval is a single point exactly when $\Delta=0$; otherwise every point of the segment between
its endpoints is admissible and distinct weights give distinct presentations. If $\Delta=0$ then
$\Lambda_{IJ}^\top\Lambda_{II}^{-1}\Lambda_{IJ}=\Lambda_{JJ}\succ0$, so $\Lambda_{IJ}$ is invertible. For (iv),
a presentation without self terms needs $W=\Lambda_{II}$, $W\Theta=\Lambda_{IJ}$, and
$\Theta^\top W\Theta=\Lambda_{JJ}$, that is, $\Delta=0$, and the sum in \eqref{eq:cg-cut-excess} with
$W_e\succ0$ vanishes exactly when every $\widehat\Theta_e=\bar\Theta$. For (v), invariance of the energy
under $z_I\mapsto\rho(g_{r_I})z_I$ and $z_J\mapsto\rho(g_{r_J})z_J$ transforms $W_e$ at the receiving end
and $\widehat\Theta_e$ as in \eqref{eq:rg-dressed-boundary-edge}, and substitution gives the law of
$\bar\Theta$; the example is direct. $\square$

\propositionheading{Balanced presentation of a cut}{prop:cg-balanced-cut-presentation}
In the setting of \Cref{lem:cg-cut-completed-square} with $\det\Lambda_{IJ}\neq0$, write
$W_{\mathrm{lo}}=\Lambda_{IJ}\Lambda_{JJ}^{-1}\Lambda_{IJ}^\top$ and $W_{\mathrm{hi}}=\Lambda_{II}$ for the endpoints
of $\mathfrak W_{\mathrm{cut}}$, and call an anchored presentation balanced when the self term at $J$,
carried to $I$ along the presentation's own link, equals the self term at $I$,
\begin{equation}
\Theta^{-\top}\Lambda^{\mathrm s}_J\Theta^{-1}=\Lambda^{\mathrm s}_I,
\qquad\text{equivalently}\qquad
\Theta^\top\Lambda_{II}\Theta=\Lambda_{JJ},
\label{eq:cg-cut-balance}
\end{equation}
so that the balanced link carries the form $\Lambda_{JJ}$ to the form $\Lambda_{II}$. (i) Exactly one
anchored presentation is balanced. Its weight is the matrix geometric mean of the endpoints,
\begin{equation}
W_{\mathrm{bal}}=W_{\mathrm{lo}}^{1/2}\bigl(W_{\mathrm{lo}}^{-1/2}W_{\mathrm{hi}}W_{\mathrm{lo}}^{-1/2}\bigr)^{1/2}W_{\mathrm{lo}}^{1/2},
\label{eq:cg-cut-balanced-weight}
\end{equation}
the unique positive-definite solution of $WW_{\mathrm{lo}}^{-1}W=W_{\mathrm{hi}}$, and its link is
$\Theta_{\mathrm{bal}}=W_{\mathrm{bal}}^{-1}\Lambda_{IJ}$. (ii) Under root re-expressions the weight transforms
as a weight at $I$, $W_{\mathrm{bal}}\mapsto\rho(g_{r_I})^{-\top}W_{\mathrm{bal}}\rho(g_{r_I})^{-1}$, the link as a
dressed edge, $\Theta_{\mathrm{bal}}\mapsto\rho(g_{r_I})\Theta_{\mathrm{bal}}\rho(g_{r_J})^{-1}$, and each self
term as a form at its own block. The presentation is end symmetric: read from $J$, with link
$\Theta_{\mathrm{bal}}^{-1}$ and weight $\Theta_{\mathrm{bal}}^\top W_{\mathrm{bal}}\Theta_{\mathrm{bal}}$, it is the
balanced presentation of the reversed cut, so that weight is the geometric mean of the endpoints
$\Lambda_{IJ}^\top\Lambda_{II}^{-1}\Lambda_{IJ}$ and $\Lambda_{JJ}$ of $J$. (iii) For scalar weights
$W_e=w_eI$ and orthogonal twists, an anchored presentation is balanced exactly when its link is
orthogonal; then $W_{\mathrm{bal}}=(\Lambda_{IJ}\Lambda_{IJ}^\top)^{1/2}$ and $\Theta_{\mathrm{bal}}$ is the
orthogonal polar factor of $\bar\Theta$. \status{ESTABLISHED}

\noindent\emph{Proof.}
By part (ii) of the lemma an anchored presentation has $\Theta=W^{-1}\Lambda_{IJ}$,
$\Lambda^{\mathrm s}_I=\Lambda_{II}-W$, and $\Lambda^{\mathrm s}_J=\Lambda_{JJ}-\Lambda_{IJ}^\top W^{-1}\Lambda_{IJ}$,
so $\Theta^{-\top}\Lambda^{\mathrm s}_J\Theta^{-1}=W\Lambda_{IJ}^{-\top}\Lambda_{JJ}\Lambda_{IJ}^{-1}W-W
=WW_{\mathrm{lo}}^{-1}W-W$, and the first form of \eqref{eq:cg-cut-balance} reads
$WW_{\mathrm{lo}}^{-1}W=W_{\mathrm{hi}}$. The second form is
$\Lambda^{\mathrm s}_J=\Theta^\top\Lambda^{\mathrm s}_I\Theta$ with the common term $\Theta^\top W\Theta$
removed from both sides, and after the substitution $\Theta=W^{-1}\Lambda_{IJ}$ it reads
$W^{-1}W_{\mathrm{hi}}W^{-1}=W_{\mathrm{lo}}^{-1}$, the same equation. (i) The substitution
$W=W_{\mathrm{lo}}^{1/2}YW_{\mathrm{lo}}^{1/2}$ is a bijection of positive-definite matrices and turns the
equation into $Y^2=M$ with $M=W_{\mathrm{lo}}^{-1/2}W_{\mathrm{hi}}W_{\mathrm{lo}}^{-1/2}\succ0$. A
positive-definite $Y$ with $Y^2=M$ commutes with $M$, so it is diagonal in an eigenbasis of $M$ with
positive eigenvalues whose squares are those of $M$; hence $Y=M^{1/2}$, which is
\eqref{eq:cg-cut-balanced-weight}. Since $\Delta\succeq0$ gives $W_{\mathrm{lo}}\preceq W_{\mathrm{hi}}$, one has
$M\succeq I$, hence $I\preceq M^{1/2}\preceq M$, and congruence by $W_{\mathrm{lo}}^{1/2}$ gives
$W_{\mathrm{lo}}\preceq W_{\mathrm{bal}}\preceq W_{\mathrm{hi}}$; the solution lies in $\mathfrak W_{\mathrm{cut}}$ and
is the weight of exactly one anchored presentation. (ii) As in the proof of part (v) of the lemma, a
root re-expression acts by $W_e\mapsto g_I^{-\top}W_eg_I^{-1}$ and
$\widehat\Theta_e\mapsto g_I\widehat\Theta_eg_J^{-1}$ with $g_I=\rho(g_{r_I})$ and $g_J=\rho(g_{r_J})$, so
$\Lambda_{II}$, $\Lambda_{IJ}$, and $\Lambda_{JJ}$ transform by $g_I^{-\top}(\cdot)g_I^{-1}$,
$g_I^{-\top}(\cdot)g_J^{-1}$, and $g_J^{-\top}(\cdot)g_J^{-1}$, and both endpoints by
$g_I^{-\top}(\cdot)g_I^{-1}$. If $W$ solves $WW_{\mathrm{lo}}^{-1}W=W_{\mathrm{hi}}$, then $g_I^{-\top}Wg_I^{-1}$
solves the transformed equation, so uniqueness gives the law of $W_{\mathrm{bal}}$, and substitution gives
the laws of the link and of the self terms. For end symmetry,
$(z_I-\Theta z_J)^\top W(z_I-\Theta z_J)=(z_J-\Theta^{-1}z_I)^\top\Theta^\top W\Theta(z_J-\Theta^{-1}z_I)$, so an
anchored presentation of the cut, read from $J$, is the anchored presentation
$(\Theta^{-1},\Theta^\top W\Theta,\Lambda^{\mathrm s}_J,\Lambda^{\mathrm s}_I)$ of the reversed cut, whose blocks
are $\Lambda_{JJ}$, $\Lambda_{IJ}^\top$, and $\Lambda_{II}$. The second form of \eqref{eq:cg-cut-balance} is
equivalent to $\Theta^{-\top}\Lambda_{JJ}\Theta^{-1}=\Lambda_{II}$, the balance condition of the reversed cut,
so the balanced presentation read from $J$ is balanced for the reversed cut, and (i) applied to that cut
identifies its weight. (iii) Here $\Lambda_{II}=\Lambda_{JJ}=(\sum_ew_e)I$, so the second form of
\eqref{eq:cg-cut-balance} says that $\Theta$ is orthogonal. For orthogonal $\Theta$ the weight
$W=\Lambda_{IJ}\Theta^\top$ is positive definite exactly when $\Lambda_{IJ}=W\Theta$ is a polar decomposition
of $\Lambda_{IJ}$, and then $W^2=\Lambda_{IJ}\Lambda_{IJ}^\top$ forces $W=(\Lambda_{IJ}\Lambda_{IJ}^\top)^{1/2}$ and
$\Theta=W^{-1}\Lambda_{IJ}$, the polar factor of $\Lambda_{IJ}$ and of $\bar\Theta=\Lambda_{IJ}/\sum_ew_e$.
$\square$

\propositionheading{No balanced presentation of a cut at the level of laws}{prop:cg-no-balanced-law-presentation}
Let $\operatorname{Gau}(\mathbb V)$ be the fiber of nondegenerate Gaussian laws on $\mathbb V$, and let $k$ copies $e$ from $J$ to
$I$ carry dressed edges $\widehat\Theta_e\in\GL(\mathbb V)$ and weights $w_e>0$ with $\sum_ew_e\widehat\Theta_e^{-1}$
invertible. For $q_J=\mathcal N(m_J,\Sigma_J)\in\operatorname{Gau}(\mathbb V)$ the transported laws
$(\widehat\Theta_e)_\#q_J=\mathcal N(\widehat\Theta_em_J,\Sigma_e)$, $\Sigma_e=\widehat\Theta_e\Sigma_J\widehat\Theta_e^\top$,
form the cut of \Cref{lem:cg-cut-completed-square} with weights $W_e=\tfrac12w_e\Sigma_e^{-1}$ and twists
$\widehat\Theta_e$; its blocks $\Lambda_{II}(\Sigma_J)$, $\Lambda_{IJ}(\Sigma_J)$, $\Lambda_{JJ}(\Sigma_J)$, its excess
$\Delta(\Sigma_J)$ \eqref{eq:cg-cut-excess}, and its balanced weight and link $W_{\mathrm{bal}}(\Sigma_J)$ and
$\Theta_{\mathrm{bal}}(\Sigma_J)$ \eqref{eq:cg-cut-balanced-weight} are functions of $\Sigma_J$, with
$\Lambda_{IJ}(\Sigma_J)$ invertible for every $\Sigma_J$. The cut has excess, $\Delta(\Sigma_J)\neq0$ for one and
hence for every $\Sigma_J$, exactly when the dressed edges are not all equal. Let $\mathscr C\subseteq\operatorname{Gau}(\mathbb V)$
be nonempty and closed under rescaling of the covariance, so that $\mathcal N(m,\Sigma)\in\mathscr C$ implies
$\mathcal N(m,t\Sigma)\in\mathscr C$ for every $t>0$. A presentation of the copies at the level of laws on
$\mathscr C$ is a weight $W'>0$, an image map $\mathfrak T:\mathscr C\to\operatorname{Gau}(\mathbb V)$, and self terms $\mathfrak S_I:\operatorname{Gau}(\mathbb V)\to\R$
and $\mathfrak S_J:\mathscr C\to\R$ with
\begin{equation}
\sum_ew_e\KL\bigl(q_I\Vert(\widehat\Theta_e)_\#q_J\bigr)=W'\KL\bigl(q_I\Vert \mathfrak T[q_J]\bigr)+\mathfrak S_I(q_I)+\mathfrak S_J(q_J)
\quad\text{for all }q_I\in\operatorname{Gau}(\mathbb V)\text{ and }q_J\in\mathscr C,
\label{eq:cg-cut-law-presentation}
\end{equation}
and it is balanced when its Gaussian mean part is the link term of the balanced presentation of the cut of
$\Sigma_J$, that is, writing $\mathfrak T[q_J]=\mathcal N(\mu_{\mathfrak T},\Sigma_{\mathfrak T})$,
\begin{equation}
\tfrac12W'\Sigma_{\mathfrak T}^{-1}=W_{\mathrm{bal}}(\Sigma_J)\quad\text{and}\quad\mu_{\mathfrak T}=\Theta_{\mathrm{bal}}(\Sigma_J)m_J
\qquad\text{for every }q_J=\mathcal N(m_J,\Sigma_J)\in\mathscr C.
\label{eq:cg-cut-law-balance}
\end{equation}
If the cut has excess, no presentation at the level of laws on $\mathscr C$ is balanced. Two copies with dressed
edges $I$ and $2I$ and weights $w_1,w_2$ have $\Delta(\Sigma_J)=\frac{w_1w_2}{2(4w_1+w_2)}\Sigma_J^{-1}$, so they
admit no balanced presentation on any such $\mathscr C$. \status{ESTABLISHED}

\noindent\emph{Proof.}
Since $\Sigma_e^{-1}=\widehat\Theta_e^{-\top}\Sigma_J^{-1}\widehat\Theta_e^{-1}$, the blocks are
$\Lambda_{II}(\Sigma_J)=\tfrac12\sum_ew_e\widehat\Theta_e^{-\top}\Sigma_J^{-1}\widehat\Theta_e^{-1}$,
$\Lambda_{IJ}(\Sigma_J)=\tfrac12\bigl(\sum_ew_e\widehat\Theta_e^{-\top}\bigr)\Sigma_J^{-1}$, and
$\Lambda_{JJ}(\Sigma_J)=\tfrac12\bigl(\sum_ew_e\bigr)\Sigma_J^{-1}$, so $\Lambda_{IJ}(\Sigma_J)$ is invertible
exactly when $\sum_ew_e\widehat\Theta_e^{-1}$ is, and each block is homogeneous of degree $-1$ in $\Sigma_J$.
Hence so are the endpoints $W_{\mathrm{lo}}=\Lambda_{IJ}\Lambda_{JJ}^{-1}\Lambda_{IJ}^\top$ and
$W_{\mathrm{hi}}=\Lambda_{II}$, the excess $\Delta$, and, since \eqref{eq:cg-cut-balanced-weight} is jointly
homogeneous of degree one in its two arguments, the balanced weight,
$W_{\mathrm{bal}}(t\Sigma_J)=t^{-1}W_{\mathrm{bal}}(\Sigma_J)$. By part (iv) of the lemma, $\Delta(\Sigma_J)=0$
exactly when all $\widehat\Theta_e$ coincide, a condition that does not involve $\Sigma_J$.

Fix a presentation, a law $q_J=\mathcal N(m_J,\Sigma_J)\in\mathscr C$, and $q_I=\mathcal N(m_I,\Sigma_I)\in\operatorname{Gau}(\mathbb V)$;
write $\mathfrak T[q_J]=\mathcal N(\mu_{\mathfrak T},\Sigma_{\mathfrak T})$, $B=\tfrac12W'\Sigma_{\mathfrak T}^{-1}$, and abbreviate
$\Lambda_{II}=\Lambda_{II}(\Sigma_J)$, $\Lambda_{IJ}=\Lambda_{IJ}(\Sigma_J)$. For Gaussian laws
$\KL(\mathcal N(m_1,\Sigma_1)\Vert\mathcal N(m_2,\Sigma_2))=\tfrac12(m_1-m_2)^\top\Sigma_2^{-1}(m_1-m_2)
+\tfrac12[\Tr(\Sigma_2^{-1}\Sigma_1)-\dim\mathbb V+\log\det\Sigma_2-\log\det\Sigma_1]$, so both sides of
\eqref{eq:cg-cut-law-presentation} are quadratic polynomials in $m_I$, and their difference is
\begin{align*}
0={}&m_I^\top(\Lambda_{II}-B)m_I-2m_I^\top(\Lambda_{IJ}m_J-B\mu_{\mathfrak T})+\Tr[(\Lambda_{II}-B)\Sigma_I]\\
&+\tfrac12\bigl(W'-\textstyle\sum_ew_e\bigr)\log\det\Sigma_I-\mathfrak S_I(q_I)+\mathrm c(q_J),
\end{align*}
where $\mathrm c(q_J)$ collects the terms that read $q_J$ alone. Subtracting this identity for a second law
$q_J'\in\mathscr C$ removes $\mathfrak S_I(q_I)$ and the $\log\det\Sigma_I$ term and leaves a quadratic polynomial in $m_I$
that vanishes for every $m_I\in\mathbb V$ and every $\Sigma_I\succ0$, so its quadratic and linear coefficients
vanish: $M=\Lambda_{II}(\Sigma_J)-B$ and $v=\Lambda_{IJ}(\Sigma_J)m_J-B\mu_{\mathfrak T}$ take the same value at every
$q_J\in\mathscr C$. If the presentation is balanced, \eqref{eq:cg-cut-law-balance} gives
$B=W_{\mathrm{bal}}(\Sigma_J)$ and
$B\mu_{\mathfrak T}=W_{\mathrm{bal}}(\Sigma_J)\Theta_{\mathrm{bal}}(\Sigma_J)m_J=\Lambda_{IJ}(\Sigma_J)m_J$, so $v=0$ and
$M=\Lambda_{II}(\Sigma_J)-W_{\mathrm{bal}}(\Sigma_J)$ for every $q_J\in\mathscr C$. Comparing $q_J$ with its
rescaling $\mathcal N(m_J,t\Sigma_J)\in\mathscr C$, the homogeneity gives $M=t^{-1}M$ for every $t>0$, hence
$M=0$ and $W_{\mathrm{bal}}(\Sigma_J)=W_{\mathrm{hi}}(\Sigma_J)$. Substituting into $WW_{\mathrm{lo}}^{-1}W=W_{\mathrm{hi}}$
gives $W_{\mathrm{hi}}W_{\mathrm{lo}}^{-1}W_{\mathrm{hi}}=W_{\mathrm{hi}}$, so $W_{\mathrm{lo}}=W_{\mathrm{hi}}$, and the
identity $\Lambda_{II}-\Lambda_{IJ}\Lambda_{JJ}^{-1}\Lambda_{IJ}^\top=\Lambda_{II}\Lambda_{IJ}^{-\top}\Delta\Lambda_{JJ}^{-1}\Lambda_{IJ}^\top$
of the proof of the lemma, whose outer factors are invertible, gives $\Delta(\Sigma_J)=0$: the cut has no
excess. For the two copies $I$ and $2I$, $\Sigma_1=\Sigma_J$ and $\Sigma_2=4\Sigma_J$, so
$\Lambda_{II}=\tfrac18(4w_1+w_2)\Sigma_J^{-1}$, $\Lambda_{IJ}=\tfrac14(2w_1+w_2)\Sigma_J^{-1}$, and
$\Lambda_{JJ}=\tfrac12(w_1+w_2)\Sigma_J^{-1}$, whence
$\Delta=\bigl[\tfrac12(w_1+w_2)-\tfrac{(2w_1+w_2)^2}{2(4w_1+w_2)}\bigr]\Sigma_J^{-1}=\frac{w_1w_2}{2(4w_1+w_2)}\Sigma_J^{-1}$.
$\square$

The balance condition is imposed only on the mean part of the link term, so the conclusion does not depend on
whether the image map is gauge covariant or the self terms are end symmetric: a self term that reads one law
cannot carry $\Lambda_{II}-W_{\mathrm{bal}}$, which is homogeneous of degree $-1$ in the covariance of the other
law, and the excess of the cut is exactly what keeps that difference from vanishing.

The exact closed category is larger than one invertible link; its endpoint maps and stalks use the notation
of the comparison maps and spaces of \Cref{ch:selective-communication}, which are an instance.

\theoremheading{Rectangular endpoint-map closure}{thm:cg-rectangular-endpoint-closure}
Choose a minimal factor $W_e=(S^Q_e)^\top S^Q_e$ with $S^Q_e:\mathbb V\to H_e$ of full row rank, and put
$S^K_e=S^Q_e\Theta_e$. Then
$(z_i-\Theta_ez_j)^\top W_e(z_i-\Theta_ez_j)=\lVert S^Q_ez_i-S^K_ez_j\rVert^2_{H_e}$. For linear maps
$\iota_{iI}:\mathbb V_I\to \mathbb V$ from block coordinates, the coarse endpoint maps $S^{Q,c}_{I,e}=S^Q_e\iota_{iI}$ and
$S^{K,c}_{J,e}=S^K_e\iota_{jJ}$ reproduce the Galerkin compression of the cut energy exactly, parallel cut
copies compress losslessly into the direct sum of their stalks $H_e$, and nested blocking is associative.
Parallel copies between the same two coarse vertices also admit a minimal exact compression: with
coboundary rows $\delta_e=(S^{Q,c}_{I,e},-S^{K,c}_{J,e})$, one edge with stalk $H$ and coboundary $\delta$
reproduces their summed energy exactly when $\delta^\top\delta=\sum_e\delta_e^\top\delta_e$; the minimal stalk
dimension is $\operatorname{rank}\sum_e\delta_e^\top\delta_e\leq\dim \mathbb V_I+\dim \mathbb V_J$, attained by any
full-row-rank factor, and two full-row-rank factors differ by an orthogonal change of edge basis.
\status{ESTABLISHED}

\noindent\emph{Proof.}
Expanding the square gives the pair form. Precomposition gives the coarse maps, whose cross block
$-(S^{Q,c}_{I,e})^\top S^{K,c}_{J,e}=-\iota_{iI}^\top W_e\Theta_e\iota_{jJ}$ and diagonal blocks agree with the
compressed energy. Stacking in $\bigoplus_eH_e$ sums the energies, and a second injection $\iota'$ gives
$(S\iota)\iota'=S(\iota\iota')$. Two quadratic forms agree exactly when their symmetric matrices agree, which
gives the compression criterion and the minimal stalk dimension; for full-row-rank factors
$\delta_1,\delta_2$ of one matrix, $O=\delta_2\delta_1^+$ satisfies $O^\top O=I$ and $\delta_2=O\delta_1$.
$\square$

This is the degree-zero cellular-sheaf Laplacian category \citep{HansenGhrist2019}, in which vertex spaces
may have unequal dimensions and automorphism groups $\GL(\mathbb V_I)$ and edge basis changes act on both endpoint
maps and the edge metric. Twice the quadratic disagreement $V_e$ of \Cref{def:obs-selective-comparison} is
the edge term of such an energy once its comparison precision is absorbed into the endpoint maps, as in
the discourse sheaves of \citet{HansenGhrist2021Discourse}; the closure selects no decoder and no law, and
the normalized realization for fixed maps is \eqref{eq:obs-selective-attention-record}.

\propositionheading{No continuous extension of balanced cut links}{open:cg-cut-compression}
There is no rule assigning a matrix-valued link and a positive-semidefinite weight and self terms to
every cut that is continuous in the constituent weights and twists and agrees with
\Cref{prop:cg-balanced-cut-presentation} whenever $\Lambda_{IJ}$ is invertible. This already fails for two
parallel copies with identity weights and twists in $\mathrm{SO}(2)\subset\GL^+(2)$. Consequently adding
gauge covariance, invariance under reversal, or a rank requirement on singular links cannot produce such
a rule on all singular and rectangular cuts. This obstruction concerns extension of the balanced
presentation; the exact rectangular endpoint-map closure of \Cref{thm:cg-rectangular-endpoint-closure}
remains available. \status{ESTABLISHED}

\noindent\emph{Proof.}
For $-1<t<1$, take $W_1=W_2=I_2$ and
\begin{equation}
\widehat\Theta_1(t)=\begin{pmatrix}t&-\sqrt{1-t^2}\\\sqrt{1-t^2}&t\end{pmatrix},
\qquad \widehat\Theta_2(t)=\widehat\Theta_1(t)^\top.
\label{eq:cg-cut-continuous-extension-witness}
\end{equation}
The twists lie in $\mathrm{SO}(2)$ and depend continuously on $t$, while
$\Lambda_{II}=\Lambda_{JJ}=2I_2$ and $\Lambda_{IJ}=2tI_2$. For $t\neq0$, the interval endpoints are
$W_{\mathrm{lo}}=2t^2I_2$ and $W_{\mathrm{hi}}=2I_2$, so the balanced weight and link are
$W_{\mathrm{bal}}=2|t|I_2$ and $\Theta_{\mathrm{bal}}=\operatorname{sign}(t)I_2$.
The same limiting cut at $t=0$ is approached from both sides, but the links tend to $I_2$ and $-I_2$.
A matrix-valued continuous extension would have to equal both limits. $\square$

\section{The frame sector}
\label{sec:cg-frame}

\propositionheading{The frame sector}{prop:cg-frame-sector}
Fix a block $I$ whose constituent frames are $u_i=\sigma\cdot U_i^{-1}$ relative to a reference point
$\sigma$ (\Cref{prop:geo-frame-comparisons}), and a coarse datum $(u_I,q_I)$ with $q_I\in\mathcal B_x$ and
$u_I=\sigma\cdot U_I^{-1}$ a point of the same torsor, so that $u_I=u_i\cdot T_{iI}$ with
$T_{iI}=U_iU_I^{-1}$. (i) The constituent expressions
$\widehat\rho_x(T_{iI})q_I=\widehat\rho_x(U_i)\bar q_I$ of the coarse law, with
$\bar q_I=\widehat\rho_x(U_I)^{-1}q_I$ the coordinate in $\sigma$ of the point $[u_I,q_I]$ of the associated
fiber, depend on $(u_I,q_I)$ only through that point; for the Gaussian law fiber of the defining
representation of $\GL^+(K)$ the constituent expression of $q_I=\mathcal N(\mu_I,\Sigma_I)$ is
$\mathcal N(U_i\bar m,U_i\bar SU_i^\top)$ with $\bar m=U_I^{-1}\mu_I$ and $\bar S=U_I^{-1}\Sigma_IU_I^{-\top}$.
The data with the same constituent expressions form exactly one orbit
$\{(u_I\cdot g,\widehat\rho_x(g)^{-1}q_I):g\in G\}$, on which $G$ acts freely, so the coarse frame occupies a
copy of $G$, of dimension $K^2$ for $G=\GL^+(K)$, that the constituents do not see. If only the represented
frame $\rho_x(U_I)$ is recorded, a nonfaithful representation determines the principal coarse frame only
modulo its kernel. (ii) If $G$ contains an element $v\neq e$ with $v^n=e$, no map $F:G^n\to G$ is both left
equivariant, $F(gg_1,\ldots,gg_n)=gF(g_1,\ldots,g_n)$, and invariant under permutations of its arguments. For
$G=\GL^+(K)$ with $K\geq2$ and every $n\geq2$ such an element exists, a rotation of order $n$ in one
coordinate plane. The torsion hypothesis cannot be dropped: on the additive group $\R^d$ the arithmetic
mean is left equivariant and symmetric. \status{ESTABLISHED}

\noindent\emph{Proof.}
(i) By the quotient convention \eqref{eq:geo-quotient-convention},
$[u_I,q_I]=[u_i\cdot T_{iI},q_I]=[u_i,\widehat\rho_x(T_{iI})q_I]$, so the constituent expression is the
coordinate of $[u_I,q_I]$ in the frame $u_i$ and determines that point. Likewise
$[u_I,q_I]=[\sigma,\widehat\rho_x(U_I)^{-1}q_I]=[u_i\cdot U_i,\bar q_I]=[u_i,\widehat\rho_x(U_i)\bar q_I]$, and
the Gaussian pushforward by the linear map $U_i$ gives the Gaussian form. Two data define the same
point exactly when they differ by $(u_I,q_I)\mapsto(u_I\cdot g,\widehat\rho_x(g)^{-1}q_I)$, and
$u_I\cdot g=u_I$ forces $g=e$. Elements of the kernel of $\rho_x$ give equal represented frames. (ii) Left
multiplication by $v$ permutes the tuple $(e,v,\ldots,v^{n-1})$ cyclically, because $v\cdot v^{n-1}=e$.
Symmetry gives $F$ the same value at the permuted tuple, while equivariance multiplies it by $v$; hence
$vF=F$ and $v=e$, a contradiction. A rotation by $2\pi/n$ in one coordinate plane lies in $\GL^+(K)$ and has
order $n$. The mean on $\R^d$ satisfies $\tfrac1n\sum_i(g+g_i)=g+\tfrac1n\sum_ig_i$. $\square$

A rule $u_I=\sigma\cdot F(a_1,\ldots,a_n)$ with $u_i=\sigma\cdot a_i$ is independent of $\sigma$ exactly when
$F$ is left equivariant, because $\sigma'=\sigma\cdot g_0$ replaces every $a_i$ by $g_0^{-1}a_i$; (ii)
therefore excludes every symmetric rule of this kind for a meta-agent's frame.

\section{An operational record at \texorpdfstring{$c_*$}{c*} sensitive to loop holonomy}

\theoremheading{An operational record sensitive to network-loop holonomy}{thm:cg-loop-operational-holonomy}
Let $p$ be a closed walk of the interaction multigraph based at a probe agent $i_0$ whose belief and model
sample fibers are real vector spaces $\mathbb V_b,\mathbb V_m$ with linear representations $\rho_b,\rho_m$, and write
$z_{i_0}(y)\in \mathbb V_{\mathrm{pr}}=\mathbb V_b\oplus \mathbb V_m$ for the probe coordinates of a latent configuration $y$ and
$\rho_{\mathrm{pr}}=\rho_b\oplus\rho_m$. Let $\Omega$ be a finite set of labels $\omega$ with prior masses
$\psi_\omega>0$, each prescribing link values on the edges of $p$, with represented loop holonomy
$\mathrm{Hol}^{\mathrm{pr}}_\omega=\rho_b(\mathrm{Hol}^b_\omega(p))\oplus\rho_m(\mathrm{Hol}^m_\omega(p))$. Keep
the joint $\mathbb P(dy,do)$ independent of $\omega$, with selected posterior version
$o\mapsto\boldsymbol\Pi_o$, the slice $\boldsymbol\Pi_{\theta,V,o,X}$ at fixed parameters and structural
data; fix $\Sigma_{\mathrm r}\succ0$ on $\mathbb V_{\mathrm{pr}}$, and declare the record joint
\[
P^{\mathrm{rec}}(\omega,dy,do,dr)=\psi_\omega\mathbb P(dy,do)
\mathcal N\bigl(r;(\mathrm{Hol}^{\mathrm{pr}}_\omega-I)z_{i_0}(y),\Sigma_{\mathrm r}\bigr)dr.
\]
(i) $P^{\mathrm{rec}}$ is a normalized joint. At an admitted $o$ and given $\omega$, the record is
distributed as $(\mathrm{Hol}^{\mathrm{pr}}_\omega-I)Z+\epsilon$, with $Z\sim(z_{i_0})_\#\boldsymbol\Pi_o$ and
$\epsilon\sim\mathcal N(0,\Sigma_{\mathrm r})$ independent of $Z$. Two labels $\omega,\omega'$ have the same
record law exactly when $(\mathrm{Hol}^{\mathrm{pr}}_\omega-I)Z$ and $(\mathrm{Hol}^{\mathrm{pr}}_{\omega'}-I)Z$
have the same law. (ii) If $\mathrm{Hol}^{\mathrm{pr}}_{\omega_0}=I$ and $\mathrm{Hol}^{\mathrm{pr}}_{\omega_1}-I$
is invertible, as for $\mathrm{SO}(2)$ links in both channels whose loop rotation angles are not multiples
of $2\pi$, the record laws of $\omega_0$ and $\omega_1$ differ exactly when $(z_{i_0})_\#\boldsymbol\Pi_o$ is
not the point mass at $0$. Under the same condition the statistic $\mathsf q=r^\top\Sigma_{\mathrm r}^{-1}r$
has different laws: its mean is $\dim \mathbb V_{\mathrm{pr}}$ under $\omega_0$ and
$\dim \mathbb V_{\mathrm{pr}}+\E\bigl[Z^\top(\mathrm{Hol}^{\mathrm{pr}}_{\omega_1}-I)^\top\Sigma_{\mathrm r}^{-1}
(\mathrm{Hol}^{\mathrm{pr}}_{\omega_1}-I)Z\bigr]\in(\dim \mathbb V_{\mathrm{pr}},+\infty]$ under $\omega_1$. (iii)
Under a passive re-expression $g$, with the record expressed in the probe frame,
$r\mapsto\rho_{\mathrm{pr}}(g_{i_0})r$ and
$\Sigma_{\mathrm r}\mapsto\rho_{\mathrm{pr}}(g_{i_0})\Sigma_{\mathrm r}\rho_{\mathrm{pr}}(g_{i_0})^\top$, the
record joint is the image of the original under a bimeasurable bijection; the posterior over $\omega$ given
$(o,r)$ and the law of $\mathsf q$ are unchanged. \status{ESTABLISHED}

\noindent\emph{Proof.}
(i) The Gaussian factor integrates to one in $r$, and $\psi$ and $\mathbb P$ are normalized. The label is
independent of $(y,o)$, so given $\omega$ and $o$ the configuration has law $\boldsymbol\Pi_o$ and the
record kernel is $\mathcal N((\mathrm{Hol}^{\mathrm{pr}}_\omega-I)z_{i_0}(y),\Sigma_{\mathrm r})$. The
characteristic function of the record is
$\E[e^{is^\top(\mathrm{Hol}^{\mathrm{pr}}_\omega-I)Z}]e^{-s^\top\Sigma_{\mathrm r}s/2}$, whose Gaussian factor
never vanishes, so two record laws agree exactly when the laws of the two images of $Z$ agree. (ii) Under
$\omega_0$ the image is $0$. Under $\omega_1$ invertibility gives
$(\mathrm{Hol}^{\mathrm{pr}}_{\omega_1}-I)Z=0$ almost surely exactly when $Z=0$ almost surely; a rotation by
$\alpha$ in the plane has $\det(R_\alpha-I)=2-2\cos\alpha$. Given $Z$, the conditional mean of $\mathsf q$ is
$\lvert\Sigma_{\mathrm r}^{-1/2}(\mathrm{Hol}^{\mathrm{pr}}_\omega-I)Z\rvert^2+\dim \mathbb V_{\mathrm{pr}}$, and
Tonelli's theorem gives the two means; the second exceeds the first exactly when the image of $Z$ is
nonzero with positive probability. (iii) The loop product telescopes under
$\Theta^x_e\mapsto g_i\Theta^x_eg_j^{-1}$, so with $U=\rho_{\mathrm{pr}}(g_{i_0})$ one has
$\mathrm{Hol}^{\mathrm{pr}}_\omega\mapsto U\mathrm{Hol}^{\mathrm{pr}}_\omega U^{-1}$ and
$z_{i_0}\mapsto Uz_{i_0}$, hence the new mean is $U(\mathrm{Hol}^{\mathrm{pr}}_\omega-I)z_{i_0}$ and the new
record kernel is the image of the old one under $r\mapsto Ur$. The joint is thus the image under the
bijection that changes the latent coordinates and multiplies $r$ by $U$, which leaves the posterior over
$\omega$ unchanged, and $r^\top U^\top(U\Sigma_{\mathrm r}U^\top)^{-1}Ur=r^\top\Sigma_{\mathrm r}^{-1}r$.
$\square$

\paragraph{Physicist's summary.}
A block of agents at one base point is a small lattice gauge system whose invariants are loop products and
edge twists modulo one global conjugation. A coarse description keeps a root frame, the tree transports to
it, the holonomy of every independent loop, and the tree-dressed boundary links; the tree is a gauge
fixing, as in lattice gauge theory, but averaged member features depend on it whenever the holonomy acts
on them. The configurations that survive exact aggregation form the kernel of the connection Laplacian,
for positive-definite couplings the holonomy-fixed vectors, and the members' marginal laws have one parent
exactly when their transported laws coincide in a holonomy-invariant law, so a planar rotation loop forces
a zero parent mean and a volume-stretching loop allows no Gaussian parent. Links between blocks compress
exactly into rectangular endpoint maps, and into one link with new self-energies when an explicit matrix is
invertible, uniquely once the self-energies are required to balance across the link; at the level of the
laws themselves no such balanced form exists once the copies disagree. The parent frame is
pure gauge for the constituents, and a noisy record of
how the loop rotates the probe's state distinguishes holonomies unless that state is pinned at zero;
Supplement B computes that holonomy from a connection along a base loop.

\chapter{Exact Agent-Network Coarse-Graining}
\label{ch:agent-network-rg}

This chapter asks what one exact coarse-graining step does to the interacting
agent network of \Cref{ch:local-collective-elbo} at $c_*$, and which state must
be kept so that the step closes. The effective action of a coarse agent network
is a conditional partition function (\Cref{thm:rg-effective-action}), real
analytic in bounded perturbations (\Cref{thm:rg-bounded-action-calculus}); in
Hoeffding interaction coordinates the step is an exact nonlinear map between
spaces of different dimension whose finite-step multipliers are bounded in the bounded tier by
$3^{\lvert V_n\rvert}-1$ at every depth, and by that number times the product of the
one-step contraction coefficients of the tilted reverse kernels, which decays with depth for
stochastic arrows (\Cref{prop:rg-interaction-native-growth-bound}); and exact closure keeps induced
hyperedges, bridge kernels, joint attention-event laws, gauge-retained marks, and
path laws with their memory (\Cref{tab:rg-retained-state}). The multiscale free
energy of \Cref{ch:full-graph-meta-agent-vfe} is built from these objects, and the
formation witnesses of \Cref{ch:finite-depth-formation} are exact marked parents
in the sense fixed here.

At depth $\ell$ a coarse channel $C_\ell:\mathsf Y_\ell\rightsquigarrow\mathsf
Y_{\ell+1}$ (\Cref{def:cg-coarse-channel}) maps the product state space of the
finite agent set $V_\ell$ to the next depth; for one arrow the spaces are
$\mathsf Y$ and $\mathsf Z$ with samples $y$ and $z$. Local symbols are
$\varrho_\ell,\mathsf m_\ell,M_\ell,\bar{\mathsf m}_\ell,H_\ell$ for the measure pair;
$C^\dagger,\mathsf U,\mathcal A$ for the reverse kernel, its conditional
expectation, and the action increment; $\mathcal G_\ell,\mathsf g,\mathsf
E_{\ell,B},\mathsf P_{\ell,B},\operatorname{asm}_\ell,\operatorname{hd}_\ell$ for
interaction coordinates; $\mathcal T_\ell,D_\ell,D_{n\leftarrow\ell}$ for the exact
step and its multipliers; $\chi$ and $\mathsf T$ for a deterministic coarse map and
a transition kernel; and a superscript $c$ for coarse objects. The auxiliary
refresh kernels and entropy estimates of
\cref{sec:rg-conditional-refresh-accuracy} concern inference accuracy
relative to a specified joint target. Longer proofs are
deferred to \cref{app:deferred-proofs}.

Scope. The retained-state table of \Cref{sec:rg-closure-theorem} catalogs exact
operations under their admission hypotheses and is not a completeness theorem;
the coarse channel and the child blocks are inputs, not canonical selections; and
the auxiliary path experiment of \Cref{sec:rg-memory-hidden-history} has its own
evidence and free energy, not identified with those of the original sensor model.
\status{NOT-CLAIMED} Open items are the invariant retained subspaces of the
interaction step (\Cref{open:rg-invariant-retained-subspace}), the dynamic closure
of the pointwise parent (\Cref{open:rg-pointwise-parent-dynamics}), and jointly
measurable equivariant conditional versions beyond the compact and finite cases
(\Cref{open:cg-equivariant-versions}).

\section{Exact effective action}
\label{sec:rg-effective-action}
\label{sec:rg-measure-pair-arrows}

\definitionheading{Scale-indexed measure pair}{def:rg-measure-pair}
For each depth $\ell$ let $\varrho_\ell$ be a probability measure on
$\mathsf Y_\ell$, let $\mathsf m_\ell\ll\varrho_\ell$ be a finite positive measure,
and let a depth arrow be a parameter-independent Markov kernel
$C_\ell:\mathsf Y_\ell\rightsquigarrow\mathsf Y_{\ell+1}$ acting on both components:
\begin{align}
\mathsf m_\ell(dy)&=\exp[-H_\ell(y)]\varrho_\ell(dy),
\qquad
0<M_\ell:=\mathsf m_\ell(\mathsf Y_\ell)<\infty,
\qquad
\bar{\mathsf m}_\ell=\mathsf m_\ell/M_\ell,
\label{eq:rg-measure-pair-normalization}\\
\varrho_{\ell+1}&=\varrho_\ell C_\ell,
\qquad
\mathsf m_{\ell+1}=\mathsf m_\ell C_\ell,
\qquad
H_{\ell+1}=-\log\frac{d\mathsf m_{\ell+1}}{d\varrho_{\ell+1}}
\quad\varrho_{\ell+1}\text{-almost everywhere},
\label{eq:rg-measure-pair-action}
\end{align}
with $\exp(-\infty)=0$, so that each action is an extended real function defined
almost everywhere. At a record $o$ and for one arrow
$C:\mathsf Y\rightsquigarrow\mathsf Z$ with its index suppressed, the fine pair is
$(\varrho,\mathsf m_o)$ with $\mathsf m_o(dy)=L_o(y)\varrho(dy)$,
$L_o=\exp(-H_o)$, $0<M_o:=\mathsf m_o(\mathsf Y)<\infty$, and
$\bar{\mathsf m}_o=M_o^{-1}\mathsf m_o$, and the pushed pair, effective likelihood,
and effective action are
\begin{equation}
\varrho^c=\varrho C,
\quad
\mathsf m_o^c=\mathsf m_oC,
\quad
\bar{\mathsf m}_o^c=\bar{\mathsf m}_oC=M_o^{-1}\mathsf m_o^c,
\quad
L_o^c=\frac{d\mathsf m_o^c}{d\varrho^c},
\quad
H_o^c=-\log L_o^c.
\label{eq:rg-reference-evidence-pair}
\end{equation}
When $\varrho$ is the latent prior and $L_o$ the record likelihood,
$\bar{\mathsf m}_o$ is the posterior slice and $M_o$ the evidence.

\propositionheading{Exact composition and normalized law}{prop:rg-measure-pair-composition}
Composable arrows act on the pair componentwise:
$(\varrho_\ell,\mathsf m_\ell)(C_\ell C_{\ell+1})
=((\varrho_\ell C_\ell)C_{\ell+1},(\mathsf m_\ell C_\ell)C_{\ell+1})$. Moreover
$\mathsf m_{\ell+1}\ll\varrho_{\ell+1}$, so that
\eqref{eq:rg-measure-pair-action} is defined, $\mathsf m_{\ell+1}(\mathsf
Y_{\ell+1})=M_\ell$, and $\bar{\mathsf m}_{\ell+1}=\bar{\mathsf m}_\ell C_\ell$; the
same statements hold for the single-arrow pair
\eqref{eq:rg-reference-evidence-pair}. \status{ESTABLISHED}

\noindent\emph{Proof.}
Composition on each component is \Cref{prop:cg-markov-category}. If
$\varrho_{\ell+1}(A)=0$, then $C_\ell(y,A)=0$ for $\varrho_\ell$-almost every $y$,
hence for $\mathsf m_\ell$-almost every $y$, and $\mathsf m_{\ell+1}(A)=0$. A
Markov kernel preserves total mass, and dividing by that mass gives the normalized
law. $\square$

\theoremheading{Conditional-partition formula and composition}{thm:rg-effective-action}
Let $C^\dagger_\varrho(z,dy)$ be a reverse conditional of $Y$ given $Z=z$ under
the joint law $\varrho(dy)C(y,dz)$. Then for $\varrho^c$-almost every $z$
\begin{equation}
\exp[-H_o^c(z)]
=\int\exp[-H_o(y)]C^\dagger_\varrho(z,dy),
\qquad
M_o=\int\exp[-H_o^c(z)]\varrho^c(dz).
\label{eq:rg-conditional-partition}
\end{equation}
For a second channel $C'$ on $\mathsf Z$, the direct effective likelihood of $CC'$
and the staged one agree $\varrho CC'$-almost everywhere when their reverse
conditionals are chosen from one joint tower. If $0\le L_o\le1$, as for a binary
success record, then $0\le L_o^c\le1$, so one effective binary record realizes the
whole coarse action; for a general record density $L_o^c$ need not be bounded by
one. \status{ESTABLISHED}

\noindent\emph{Proof.}
For bounded measurable $f$, $\int f(z)\mathsf m_o^c(dz)=\E_\varrho[f(Z)L_o(Y)]
=\int f(z)\E_\varrho[L_o(Y)\given Z=z]\varrho^c(dz)$ under the joint law, which
identifies the Radon--Nikodym derivative, as in the density formula of
\Cref{thm:cg-dpi-equality} with $L_o=d\mathsf m_o/d\varrho$; $f=1$ gives the
evidence. The tower property proves composition, and a conditional expectation of
a function with values in $[0,1]$ takes values in $[0,1]$. $\square$

\begin{figure}[htbp]
\centering
\begin{tikzpicture}[>=stealth,every node/.style={font=\small}]
  \node[draw,rounded corners,fill=lightblue,minimum height=0.7cm]
    (fine) at (0,0) {fine law $(\varrho,\mathsf m_o,\eta)$};
  \node[draw,rounded corners,fill=lightgreen,minimum height=0.7cm]
    (coarse) at (4.9,0) {coarse law $(\varrho^c,\mathsf m_o^c,\eta^c)$};
  \node[draw,rounded corners,fill=lightorange,minimum height=0.7cm]
    (action) at (12.1,0) {$H_o^c=-\log(d\mathsf m_o^c/d\varrho^c)$};
  \draw[->,very thick,primaryblue] (fine)--node[above]{$C$}(coarse);
  \draw[->,very thick,sciencegreen] (coarse)--node[above]{Radon--Nikodym}(action);
\end{tikzpicture}
\caption{One exact step acts on the reference probability, the evidence measure,
and the joint attention-event law; the effective action is the Radon--Nikodym
derivative of the pushed evidence measure with respect to the pushed reference.}
\label{fig:rg-exact-measure-flow}
\end{figure}

\section{Bounded action calculus and contraction}
\label{sec:rg-local-action-calculus}
\label{sec:rg-action-contraction}

\theoremheading{Exact bounded action calculus}{thm:rg-bounded-action-calculus}
\label{prop:rg-action-bounded-recentering}%
Fix a jointly measurable reverse kernel $C^\dagger(z,dy)$ of the normalized joint
law and its reverse conditional expectation,
\begin{equation}
\bar{\mathsf m}_o(dy)C(y,dz)=\bar{\mathsf m}_o^c(dz)C^\dagger(z,dy),
\qquad
(\mathsf Uh)(z)=\int h(y)C^\dagger(z,dy),
\label{eq:rg-normalized-reverse-kernel}
\end{equation}
whose values on $\bar{\mathsf m}_o^c$-null sets are not used. A bounded $\psi$
perturbs the evidence measure to $e^{-\psi}\mathsf m_o$; on
$\mathfrak B=L^\infty(\bar{\mathsf m}_o;\R)$ the exact coarse action increment
$\mathcal A:\mathfrak B\to L^\infty(\bar{\mathsf m}_o^c;\R)$ is
\begin{equation}
\mathcal A(\psi)(z)
=-\log\frac{d\bigl((e^{-\psi}\mathsf m_o)C\bigr)}{d(\mathsf m_oC)}(z)
=-\log\int e^{-\psi(y)}C^\dagger(z,dy),
\label{eq:rg-local-action-map}
\end{equation}
the difference of the coarse actions of $\mathsf m_o$ and $e^{-\psi}\mathsf m_o$
where both are finite; elsewhere the calculus supplies no action difference and the
Radon--Nikodym ratio is the definition.

(a) $e^{-\lVert\psi\rVert_\infty}\le\mathsf U(e^{-\psi})\le e^{\lVert\psi\rVert_\infty}$,
so $\lVert\mathcal A(\psi)\rVert_\infty\le\lVert\psi\rVert_\infty$, and $\mathcal A$
is real analytic on the whole of $\mathfrak B$.

(b) For $\psi\in\mathfrak B$ the tilted reverse kernel and its operator
\begin{equation}
C^\dagger_\psi(z,dy)=
\frac{e^{-\psi(y)}C^\dagger(z,dy)}{\int e^{-\psi(u)}C^\dagger(z,du)},
\qquad
\mathsf U^\psi h=\frac{\mathsf U(e^{-\psi}h)}{\mathsf U(e^{-\psi})},
\label{eq:rg-tilted-reverse-kernel}
\end{equation}
are positive, unital, and bounded from $L^\infty(\bar{\mathsf m}_o)$ to
$L^\infty(\bar{\mathsf m}_o^c)$, and for bounded $h$ and $k$
\begin{align}
\mathcal A(\psi+h)-\mathcal A(\psi)&=-\log\mathsf U^\psi(e^{-h}),
\label{eq:rg-action-recentering-identity}\\
D\mathcal A(\psi)[h](z)&=\int h(y)C^\dagger_\psi(z,dy)=(\mathsf U^\psi h)(z),
\label{eq:rg-action-first-frechet-general}\\
D^2\mathcal A(\psi)[h,k](z)&=
-\operatorname{Cov}_{C^\dagger_\psi(z,\cdot)}(h,k),
\label{eq:rg-action-second-frechet-general}
\end{align}
which at the origin read $D\mathcal A(0)[h]=\mathsf Uh$ and
$D^2\mathcal A(0)[h,k]=-\operatorname{Cov}_{C^\dagger(z,\cdot)}(h,k)$.

(c) The map is order preserving and additively homogeneous:
$\psi\le\psi'+c$ implies $\mathcal A(\psi)\le\mathcal A(\psi')+c$, and
$\mathcal A(\psi+c)=\mathcal A(\psi)+c$ for constants $c$. Consequently
\begin{equation}
\lVert\mathcal A(\psi)-\mathcal A(\psi')\rVert_\infty
\le\lVert\psi-\psi'\rVert_\infty .
\label{eq:rg-nonlinear-action-sup-contraction}
\end{equation}

(d) Both $\mathcal A$ and $\mathsf U$ descend to the projective action space
$\overline{\mathfrak B}=\mathfrak B/\R1$, which forgets the perturbed mass
$(e^{-\psi}\mathsf m_o)(\mathsf Y)$.

(e) Normalizing the perturbed fine law instead gives
\begin{equation}
\widehat{\mathsf m}^\psi
=\frac{e^{-\psi}\bar{\mathsf m}_o}{\bar{\mathsf m}_o(e^{-\psi})},
\qquad
\widetilde{\mathcal A}(\psi)
:=-\log\frac{d(\widehat{\mathsf m}^\psi C)}{d\bar{\mathsf m}_o^c}
=\mathcal A(\psi)+\log\bar{\mathsf m}_o(e^{-\psi}),
\label{eq:rg-normalized-perturbation-action}
\end{equation}
which differs from $\mathcal A$ by a constant, with derivatives
$D\widetilde{\mathcal A}(0)[h]=\mathsf Uh-\bar{\mathsf m}_o(h)$ and
$D^2\widetilde{\mathcal A}(0)[h,k]=-\operatorname{Cov}_{C^\dagger(z,\cdot)}(h,k)
+\operatorname{Cov}_{\bar{\mathsf m}_o}(h,k)$. \status{ESTABLISHED}

\theoremheading{Conditional-expectation contraction and exact defect}{thm:rg-action-lp-contraction}
For every $1\le p\le\infty$ the operator $\mathsf U$ of
\eqref{eq:rg-normalized-reverse-kernel} extends uniquely to a mean-preserving
contraction $\mathsf U:L^p(\bar{\mathsf m}_o;\R)\to L^p(\bar{\mathsf m}_o^c;\R)$,
and at $p=2$ the loss is a conditional variance:
\begin{align}
\lVert\mathsf Uh\rVert_{L^p(\bar{\mathsf m}_o^c)}&\le\lVert h\rVert_{L^p(\bar{\mathsf m}_o)},
\label{eq:rg-action-lp-contraction}\\
\lVert h\rVert_{L^2(\bar{\mathsf m}_o)}^2-\lVert\mathsf Uh\rVert_{L^2(\bar{\mathsf m}_o^c)}^2
&=\int\operatorname{Var}_{C^\dagger(z,\cdot)}(h)\bar{\mathsf m}_o^c(dz)\ge0,
\label{eq:rg-action-l2-conditional-defect}
\end{align}
with equality exactly when $h(y)=f(z)$ for $\bar{\mathsf m}_o(dy)C(y,dz)$-almost
every $(y,z)$ and some measurable $f$ on $\mathsf Z$; for a deterministic coarse
map $z=\chi(y)$ this means $h=f\circ\chi$ $\bar{\mathsf m}_o$-almost everywhere.
With the oscillation norm on $L^\infty(\bar{\mathsf m}_o)/\R1$ and the contraction
coefficient of \citet{Dobrushin1956} of the reverse kernel,
\begin{align}
\lVert[h]\rVert_{\mathrm{osc}}
&=\inf_{c\in\R}\lVert h-c\rVert_\infty
=\tfrac12\bigl(\operatorname*{ess~sup}h-\operatorname*{ess~inf}h\bigr),
\label{eq:rg-action-oscillation-norm}\\
\delta&=
\operatorname*{ess~sup}_{(z,z')\sim(\bar{\mathsf m}_o^c)^{\otimes2}}
\sup_{A}\bigl\lvert C^\dagger(z,A)-C^\dagger(z',A)\bigr\rvert,
\label{eq:rg-action-dobrushin-coefficient}
\end{align}
the inner supremum running over Borel sets, one has
$\lVert[\mathsf Uh]\rVert_{\mathrm{osc}}\le\delta\lVert[h]\rVert_{\mathrm{osc}}$, and
along a finite composition of arrows the oscillation norm of the composed reverse
operators is at most the product of the one-step coefficients.
\status{ESTABLISHED}

The proofs of \Cref{thm:rg-bounded-action-calculus,thm:rg-action-lp-contraction}
are in \cref{sec:proofs-network-coarsegraining}.

\corollaryheading{Divergence bounds for a bounded action perturbation}{cor:rg-perturbation-divergence-bounds}
In the setting of \Cref{thm:rg-bounded-action-calculus}, which is the bounded tier of
\Cref{rem:rg-bounded-tier-scope}, let $\psi\in\mathfrak B$ perturb the evidence
measure to $e^{-\psi}\mathsf m_o$, whose action on the same
reference is $H_o+\psi$. Then $\widehat{\mathsf m}^\psi$ of
\eqref{eq:rg-normalized-perturbation-action} is equivalent to $\bar{\mathsf m}_o$, and
\begin{equation}
\KL\bigl(\bar{\mathsf m}_o\Vert\widehat{\mathsf m}^\psi\bigr),\quad
\KL\bigl(\widehat{\mathsf m}^\psi\Vert\bar{\mathsf m}_o\bigr)
\le\min\bigl\{2\lVert[\psi]\rVert_{\mathrm{osc}},\lVert[\psi]\rVert_{\mathrm{osc}}^2/2\bigr\},
\label{eq:rg-perturbation-two-sided-kl}
\end{equation}
where both sides read $\psi$ only through its class in the projective action space
$\overline{\mathfrak B}$ of part (d) of \Cref{thm:rg-bounded-action-calculus}. The
offsets read the essential bounds of $\psi$ instead: the log evidence ratio
$\log\bar{\mathsf m}_o(e^{-\psi})=\log\bigl((e^{-\psi}\mathsf m_o)(\mathsf Y)/M_o\bigr)$,
the constant of part (e), lies in
$[-\operatorname*{ess~sup}\psi,-\operatorname*{ess~inf}\psi]$, and for every law
$Q\ll\bar{\mathsf m}_o$ the free energies
$\Fenergy_{\mathsf m_o}(Q)=-\log M_o+\KL(Q\Vert\bar{\mathsf m}_o)$ and
$\Fenergy_{e^{-\psi}\mathsf m_o}(Q)$, of the form of \Cref{def:elbo-extended},
satisfy $\Fenergy_{e^{-\psi}\mathsf m_o}(Q)=\Fenergy_{\mathsf m_o}(Q)+\E_Q\psi$ in
$(-\infty,+\infty]$ with
$\E_Q\psi\in[\operatorname*{ess~inf}\psi,\operatorname*{ess~sup}\psi]$, while both are
$+\infty$ for every other law $Q$. Along arrows $C_0,\ldots,C_{L-1}$ with pairs
$(\varrho_\ell,\mathsf m_\ell)$ as in \Cref{def:rg-measure-pair}, let
$\widetilde{\mathsf m}_0=e^{-\psi_0}\mathsf m_0$ and
$\widetilde{\mathsf m}_{\ell+1}=e^{-\psi_{\ell+1}}(\widetilde{\mathsf m}_\ell C_\ell)$ with
$\psi_\ell\in L^\infty(\bar{\mathsf m}_\ell)$, let
$\widetilde H_\ell=-\log(d\widetilde{\mathsf m}_\ell/d\varrho_\ell)$, and let
$\mathcal A_\ell$ be the action increment \eqref{eq:rg-local-action-map} of $C_\ell$
at $\mathsf m_\ell$. Then each $\widetilde H_\ell-H_\ell$ is essentially bounded, and
\begin{equation}
\widetilde H_{\ell+1}-H_{\ell+1}=\mathcal A_\ell(\widetilde H_\ell-H_\ell)+\psi_{\ell+1},
\qquad
\sum_{j=0}^{L}\operatorname*{ess~inf}\psi_j
\le\widetilde H_L-H_L
\le\sum_{j=0}^{L}\operatorname*{ess~sup}\psi_j,
\label{eq:rg-perturbation-tower}
\end{equation}
the recurrence holding $\mathsf m_{\ell+1}$-almost everywhere and the bounds
$\mathsf m_L$-almost everywhere. Hence the preceding statements hold with
$\mathsf m_L$ in place of $\mathsf m_o$ and $\psi=\widetilde H_L-H_L$, with
$\lVert[\widetilde H_L-H_L]\rVert_{\mathrm{osc}}\le\sum_{j=0}^{L}\lVert[\psi_j]\rVert_{\mathrm{osc}}$
and both offsets in the summed interval. \status{ESTABLISHED}

The proof is in \cref{sec:proofs-network-coarsegraining}.

The divergence bound sees the width of a perturbation and not its offset, which the
evidence and the free energy carry. Outside $\mathfrak B$, for instance for the
polynomial couplings on $\R^K$ described in \Cref{rem:rg-bounded-tier-scope}, the
corollary gives no bound.

\corollaryheading{Contracted widths and transmitted offsets along a perturbed tower}{cor:rg-contracted-perturbation-tower}
In the setting of the tower of \Cref{cor:rg-perturbation-divergence-bounds}, write
$\mathfrak r_\ell=\widetilde H_\ell-H_\ell$ for the propagated residual, let
$C^\dagger_\ell$ be a reverse kernel of the normalized joint law
$\bar{\mathsf m}_\ell(dy)C_\ell(y,dz)$ as in \eqref{eq:rg-normalized-reverse-kernel},
so that $\mathcal A_\ell(\psi)(z)=-\log\int e^{-\psi(y)}C^\dagger_\ell(z,dy)$, and
measure widths by the full essential range under $\bar{\mathsf m}_\ell$,
\begin{equation}
\omega_\ell=2\lVert[\mathfrak r_\ell]\rVert_{\mathrm{osc}}
=\operatorname*{ess~sup}\mathfrak r_\ell-\operatorname*{ess~inf}\mathfrak r_\ell,
\qquad
\omega^{\mathrm{st}}_\ell=2\lVert[\psi_\ell]\rVert_{\mathrm{osc}},
\label{eq:rg-tower-widths}
\end{equation}
so that $\omega_0=\omega^{\mathrm{st}}_0$. Let $\delta_\ell$ be the coefficient
\eqref{eq:rg-action-dobrushin-coefficient} of $C^\dagger_\ell$, and let its
projective diameter in the sense of Birkhoff, $\Delta_\ell\in[0,\infty]$, be the
infimum of the $t\ge0$ for which
\begin{equation}
C^\dagger_\ell(z,A)C^\dagger_\ell(z',B)\le e^{t}C^\dagger_\ell(z',A)C^\dagger_\ell(z,B)
\label{eq:rg-birkhoff-diameter}
\end{equation}
holds, for $(\bar{\mathsf m}_{\ell+1})^{\otimes2}$-almost every $(z,z')$, simultaneously
for all Borel $A,B\subseteq\mathsf Y_\ell$, with
$\Delta_\ell=\infty$ when no $t$ qualifies and $\tanh(\infty)=1$.

(a) For $0\le\ell<L$,
\begin{equation}
\omega_{\ell+1}\le\min\Bigl\{\log\bigl(1+\delta_\ell(e^{\omega_\ell}-1)\bigr),
\tanh(\Delta_\ell/4)\omega_\ell\Bigr\}+\omega^{\mathrm{st}}_{\ell+1}.
\label{eq:rg-contracted-width-recursion}
\end{equation}
When $\omega_\ell>0$, the first term of the minimum is smaller than $\omega_\ell$
exactly when $\delta_\ell<1$, and the second exactly when $\Delta_\ell<\infty$.

(b) The first term is the exact supremum at given coefficient and width: for every
$\delta\in[0,1]$ and $x\ge0$ there are a normalized law on a two-point space, a
channel into a two-point space whose reverse kernel has coefficient $\delta$, and a
$\psi$ with $\operatorname*{ess~sup}\psi-\operatorname*{ess~inf}\psi=x$, for which
$\operatorname*{ess~sup}\mathcal A(\psi)-\operatorname*{ess~inf}\mathcal
A(\psi)=\log(1+\delta(e^x-1))$, and for $\delta>0$ this reverse kernel has infinite
projective diameter. Hence no bound on the width of $\mathcal A_\ell(\psi)$ by a
function of $\delta_\ell$ and the width of $\psi$ alone is smaller than the first
term, and $\delta_\ell<1$ does not imply $\Delta_\ell<\infty$.

(c) If $\delta_\ell\le\delta<1$ for $0\le\ell<L$ and
$\omega^{\mathrm{st}}_\ell\le s$ for $1\le\ell\le L$ with $s<\log(1/\delta)$, read as
$+\infty$ when $\delta=0$, then $\omega_\ell\le\max\{\omega_0,x_{\delta,s}\}$ for
$0\le\ell\le L$, where $x_{\delta,s}$ is the unique $x\ge0$ with
$x-\log(1+\delta(e^x-1))=s$. If instead $\Delta_\ell\le\Delta<\infty$ for
$0\le\ell<L$ and $\omega^{\mathrm{st}}_\ell\le s$ for $1\le\ell\le L$, with no
condition on $s\ge0$, then
$\omega_\ell\le\max\{\omega_0,s/(1-\tanh(\Delta/4))\}$ for $0\le\ell\le L$. Neither bound depends on $L$, and by
\Cref{cor:rg-perturbation-divergence-bounds} both divergences between
$\bar{\mathsf m}_L$ and the normalization of $\widetilde{\mathsf m}_L$ are at most
$\min\{\omega_L,\omega_L^2/8\}$.

(d) If $C_\ell(y,A)=1_A(\chi_\ell(y))$ for a Borel map
$\chi_\ell:\mathsf Y_\ell\to\mathsf Y_{\ell+1}$ and $\bar{\mathsf m}_{\ell+1}$ is not
a point mass, then $\delta_\ell=1$ and $\Delta_\ell=\infty$, and
\eqref{eq:rg-contracted-width-recursion} reduces to the additive recursion
$\omega_{\ell+1}\le\omega_\ell+\omega^{\mathrm{st}}_{\ell+1}$ behind
\eqref{eq:rg-perturbation-tower}. If every arrow has this form, $F$ is a bounded
Borel function on $\mathsf Y_L$, $c_0,\ldots,c_L\ge0$, and
$\psi_j=c_jF\circ\chi_{L-1}\circ\cdots\circ\chi_j$ (with $\psi_L=c_LF$), then
$\omega_\ell=\sum_{j=0}^{\ell}\omega^{\mathrm{st}}_j$ for $0\le\ell\le L$, so the
additive bound is attained at every depth.

(e) For every $\psi\in L^\infty(\bar{\mathsf m}_\ell)$,
$\log\bar{\mathsf m}_{\ell+1}(e^{-\mathcal A_\ell(\psi)})=\log\bar{\mathsf m}_\ell(e^{-\psi})$.
With $\widehat{\mathsf n}_j=\widetilde{\mathsf m}_{j-1}C_{j-1}/\widetilde{\mathsf
m}_{j-1}(\mathsf Y_{j-1})$ for $1\le j\le L$, the log evidence ratio decomposes as
\begin{equation}
\log\frac{\widetilde{\mathsf m}_L(\mathsf Y_L)}{M_L}
=\log\bar{\mathsf m}_L(e^{-\mathfrak r_L})
=\log\bar{\mathsf m}_0(e^{-\psi_0})+\sum_{j=1}^{L}\log\widehat{\mathsf n}_j(e^{-\psi_j}),
\label{eq:rg-tower-offset-decomposition}
\end{equation}
each summand lying in
$[-\operatorname*{ess~sup}\psi_j,-\operatorname*{ess~inf}\psi_j]$. Replacing one
$\psi_j$ by $\psi_j+c$ for a constant $c$ leaves every $\omega_\ell$ and the
normalization of every $\widetilde{\mathsf m}_\ell$ unchanged and lowers the log evidence
ratio by exactly $c$, whatever the arrows, including arrows with $\delta_\ell=0$.
\status{ESTABLISHED}

The proof is in \cref{sec:proofs-network-coarsegraining}.

The factor $\tanh(\Delta_\ell/4)$ in (a) has the form of the contraction ratio $\tanh(\Delta/4)$, in
Hilbert's projective metric, of a positive linear map whose image cone has projective diameter
$\Delta$ \citep[Sec.~4, Lemma~1, Eq.~(9a)]{Birkhoff1957}, stated for nonnegative kernels acting on
nonnegative measures by \citet[Lemma~3.8, Eq.~(11)]{LeGlandOudjane2004}. The proof does not invoke
that lemma: it bounds the width directly for rows defined almost everywhere, with $\Delta_\ell$
defined on rows by the cross-ratio \eqref{eq:rg-birkhoff-diameter}.

Along a stochastic tower the widths, which control the divergences, contract through
\eqref{eq:rg-contracted-width-recursion}, while the offsets, which the evidence and
the free energy carry, pass every arrow with coefficient one.

\section{Finite-network interaction coordinates}
\label{sec:rg-finite-network-interactions}

\definitionheading{Product-referenced network scale}{def:rg-network-scale}
At depth $\ell$ agent $i\in V_\ell$ has state space $\mathsf Y_{\ell i}$ and a
probability law $\nu_{\ell i}$, and the normalized law $\bar{\mathsf m}_\ell$ of
\Cref{def:rg-measure-pair} on $\mathsf Y_\ell=\prod_{i\in V_\ell}\mathsf Y_{\ell i}$
is mutually absolutely continuous, $\bar{\mathsf m}_\ell\sim\nu_\ell$, with the
product reference $\nu_\ell=\bigotimes_{i\in V_\ell}\nu_{\ell i}$. Equivalence
identifies $L^\infty(\bar{\mathsf m}_\ell)$ with $L^\infty(\nu_\ell)$
isometrically, with no bound on the Radon--Nikodym ratio, and the product
structure of $\nu_\ell$ makes the coordinate expectations below well defined on
equivalence classes by Fubini's theorem. For an arrow $C_\ell$, equivalence of
$\bar{\mathsf m}_{\ell+1}=\bar{\mathsf m}_\ell C_\ell$ to a product reference
$\nu_{\ell+1}$ is a separate premise
(\Cref{prop:rg-product-equivalence-not-preserved}). If $V_\ell=\varnothing$, then
$\mathsf Y_\ell$ is a point, $\mathcal G_\ell=\{0\}$, and every map below is zero;
estimates involving $3^{\lvert V_\ell\rvert}-1$ assume $V_\ell\ne\varnothing$. The vertex sets
are fixed and every coordinate is an agent coordinate; a state with a random vertex set or with
edge variables, such as the tower state of \Cref{ch:full-graph-meta-agent-vfe}, is covered only
on its agent coordinates after conditioning on the mark and the graph datum, as in
\Cref{open:full-graph-hierarchy}.

\remarkheading{The bounded tier}{rem:rg-bounded-tier-scope}
Every interaction component below is essentially bounded. The bounded tier contains
every interaction on finite agent spaces and every continuous interaction on
compact ones, hence the natural couplings of categorical agents and of compact
group-valued states. On $\mathsf Y_{\ell i}=\R^K$ it excludes the polynomial
couplings: for a matrix $\mathsf N\ne0$ and every $M>0$ the set where
$\lvert y_i^\top\mathsf Ny_j\rvert>M$ is a nonempty open subset of
$\R^K\times\R^K$, so it has positive mass under every reference with full support,
and $y_i^\top\mathsf Ny_j\notin L^\infty(\nu_\ell)$. Such couplings need a separate
tier, exact Gaussian marginalization or a weighted action space, which this
chapter does not construct. \status{ESTABLISHED}

\propositionheading{Product equivalence is an admitted, not an automatic, scale premise}{prop:rg-product-equivalence-not-preserved}
Let $\mathsf Y_0=\{0,1\}$ with its uniform law $\bar{\mathsf m}_0=\nu_0$, and let
$C$ be the deterministic diagonal-cloning channel $y\mapsto(y,y)$ from $\{0,1\}$ to
$\{0,1\}^2$. Then $\bar{\mathsf m}_1=\bar{\mathsf m}_0C$ is supported on
$\{(0,0),(1,1)\}$, no product probability $\nu_{11}\otimes\nu_{12}$ on
$\{0,1\}^2$ is equivalent to it, and no target Hoeffding decomposition relative to
a product reference is admitted at this target. \status{ESTABLISHED}

\noindent\emph{Proof.}
Equivalence would make both diagonal atoms positive under the product, so each
marginal would charge both $0$ and $1$ and both off-diagonal atoms would have
positive product mass, while $\bar{\mathsf m}_1$ assigns them mass zero.
$\square$

\definitionheading{Hoeffding coordinates}{def:rg-hoeffding-coordinates}
Under \Cref{def:rg-network-scale}, for $B\subseteq V_\ell$ let
$(\E_{\ell,B}f)(y_B)=\int f(y_B,y'_{B^c})\nu_{\ell,B^c}(dy'_{B^c})$ integrate
out the complementary coordinates against the product reference, with
$\E_{\ell,\varnothing}f=\nu_\ell(f)$. The Hoeffding projectors are the
M\"obius inversions
\begin{equation}
\mathsf P_{\ell,B}
=\sum_{B'\subseteq B}(-1)^{\lvert B\rvert-\lvert B'\rvert}\E_{\ell,B'}
=\prod_{i\in B}\bigl(I-\E_{\ell,V_\ell\setminus\{i\}}\bigr)\E_{\ell,B},
\qquad B\subseteq V_\ell,
\label{eq:rg-hoeffding-mobius-projectors}
\end{equation}
where the alternating sum is the definition and the product is the identity of
\Cref{thm:rg-hoeffding-action-isomorphism}. The component spaces
$\mathcal H_{\ell,B}=\mathsf P_{\ell,B}L^\infty(\nu_\ell)$ form the interaction
space $\mathcal G_\ell=\bigoplus^{\ell^1}_{\varnothing\ne B\subseteq
V_\ell}\mathcal H_{\ell,B}$ with norm $\lVert\mathsf g\rVert_{\mathcal
G_\ell}=\sum_{\varnothing\ne B\subseteq V_\ell}\lVert\mathsf g_B\rVert_\infty$, so
every nonempty hyperedge has its component and no bound on $\lvert V_\ell\rvert$
is imposed. With the action quotient
$\overline{\mathfrak B}_\ell=L^\infty(\nu_\ell)/\R1$, normed by
\eqref{eq:rg-action-oscillation-norm}, the assembly and extraction maps are
\begin{equation}
\psi_{\mathsf g}=\sum_{\varnothing\ne B\subseteq V_\ell}\mathsf g_B,
\qquad
\operatorname{asm}_\ell\mathsf g=[\psi_{\mathsf g}],
\qquad
\operatorname{hd}_\ell[f]=(\mathsf P_{\ell,B}f)_{\varnothing\ne B\subseteq V_\ell}.
\label{eq:rg-interaction-assembly-extraction}
\end{equation}

\theoremheading{Hoeffding assembly is an exact finite-network action isomorphism}{thm:rg-hoeffding-action-isomorphism}
For every finite $V_\ell$ the coordinate expectations commute and satisfy
$\E_{\ell,B}\E_{\ell,B'}=\E_{\ell,B\cap B'}$, the product formula
in \eqref{eq:rg-hoeffding-mobius-projectors} holds, and
$\mathsf P_{\ell,B}\mathsf P_{\ell,B'}=\mathbf 1_{B=B'}\mathsf P_{\ell,B}$ with
$\sum_{B\subseteq V_\ell}\mathsf P_{\ell,B}=I$. Each component $\mathsf g_B$
depends only on $y_B$ and has zero $\nu_\ell$-mean in every coordinate of $B$. The
maps of \eqref{eq:rg-interaction-assembly-extraction} are mutually inverse bounded
linear isomorphisms,
\begin{equation}
\operatorname{hd}_\ell\operatorname{asm}_\ell=I_{\mathcal G_\ell},
\qquad
\operatorname{asm}_\ell\operatorname{hd}_\ell=I_{\overline{\mathfrak B}_\ell},
\qquad
\lVert\operatorname{asm}_\ell\rVert\le1,
\qquad
\lVert\operatorname{hd}_\ell\rVert\le3^{\lvert V_\ell\rvert}-1 .
\label{eq:rg-hoeffding-inverse-and-bounds}
\end{equation}
The last bound is not dimension free, and its exponential dependence is sharp: for independent $\{-1,1\}$-valued coordinates with
$\Pr(Y_i=1)=p$ with $1/2\leq p<1$ and $f=\prod_iY_i$ on $n=\lvert V_\ell\rvert$ agents, the sum of
the component norms is $(4p-1)^n-(2p-1)^n$, which tends to $3^n-1$ as $p\uparrow1$.
\status{ESTABLISHED}

The proof is in \cref{sec:proofs-network-coarsegraining}.

\corollaryheading{Gauge covariance of interaction components}{cor:rg-hoeffding-gauge-covariance}
Let a passive change $g=(g_i)$ of \Cref{ch:geometry}, one element per agent acting
on both the belief and the model components of that agent's state, be realized at
depth $\ell$ by coordinate Borel isomorphisms and a Borel permutation $\varsigma$ of
the agents, whose product map $\kappa_\ell:\mathsf Y_\ell\to\widehat{\mathsf
Y}_\ell$ sends $\nu_\ell$ to a product reference $\widehat\nu_\ell$. With
$\mathsf K_\ell f=f\circ\kappa_\ell^{-1}$, its action $\overline{\mathsf K}_\ell$
on the quotient by constants, and its action $\mathsf K^{\mathcal G}_\ell$ on
components,
\begin{equation}
\mathsf K_\ell\E^{\nu}_{\ell,B}=\E^{\widehat\nu}_{\ell,\varsigma B}\mathsf K_\ell,
\qquad
\mathsf K_\ell\mathsf P^{\nu}_{\ell,B}=\mathsf P^{\widehat\nu}_{\ell,\varsigma B}\mathsf K_\ell,
\qquad
\overline{\mathsf K}_\ell\operatorname{asm}_\ell
=\widehat{\operatorname{asm}}_\ell\mathsf K^{\mathcal G}_\ell,
\qquad
\mathsf K^{\mathcal G}_\ell\operatorname{hd}_\ell
=\widehat{\operatorname{hd}}_\ell\overline{\mathsf K}_\ell.
\label{eq:rg-hoeffding-gauge-covariance}
\end{equation}
Preservation of the product reference alone is insufficient: the
product-Haar-preserving shear $(y_1,y_2)\mapsto(y_1,y_1+y_2)$ of the torus
$\mathbb T^2$ sends a function of $y_2$ to one depending on both coordinates and
moves a singleton component into a pair component. If only almost-everywhere
reverse-kernel equivariance is supplied, the covariance holds for $L^\infty$
equivalence classes and not for a selected pointwise version.
\status{ESTABLISHED}

\noindent\emph{Proof.}
Fubini's theorem proves the first identity, and the others follow by the finite
M\"obius sums. The shear maps $f(y_2)$ to $f(y_2-y_1)$, which is not a function of
one coordinate when $f$ is nonconstant. $\square$

\section{The exact interaction step, its residual, and finite-step multipliers}
\label{sec:rg-exact-nonlinear-interaction-rg}

\definitionheading{Exact interaction step}{def:rg-interaction-step}
Assume \Cref{def:rg-network-scale} at depths $\ell$ and $\ell+1$, let
$\mathcal A_\ell$ be the action increment of \Cref{thm:rg-bounded-action-calculus}
for the arrow $C_\ell$ at the base law $\bar{\mathsf m}_\ell$, and let
$\overline{\mathcal A}_\ell:\overline{\mathfrak B}_\ell\to\overline{\mathfrak
B}_{\ell+1}$ be its quotient map, defined on the whole bounded space by additive
homogeneity. The exact interaction step is
\begin{equation}
\mathcal T_\ell=\operatorname{hd}_{\ell+1}\overline{\mathcal A}_\ell\operatorname{asm}_\ell:
\mathcal G_\ell\longrightarrow\mathcal G_{\ell+1},
\qquad
\mathsf g^{\mathrm{ex}}_{\ell+1}=\mathcal T_\ell(\mathsf g_\ell).
\label{eq:rg-exact-nonlinear-interaction-map}
\end{equation}

\lemmaheading{Mass, origin, and derivative of the exact step}{lem:rg-interaction-step}
The sum $\psi_{\mathsf g}$ is the representative of
$\operatorname{asm}_\ell\mathsf g$ with zero $\nu_\ell$-mean, and the perturbed
evidence mass is carried separately and preserved by the arrow,
\begin{equation}
M_\ell(\mathsf g):=M_\ell\bar{\mathsf m}_\ell\bigl(e^{-\psi_{\mathsf g}}\bigr)
=\bigl((e^{-\psi_{\mathsf g}}\mathsf m_\ell)C_\ell\bigr)(\mathsf Y_{\ell+1})
=:M_{\ell+1}(\mathsf g).
\label{eq:rg-interaction-evidence-mass}
\end{equation}
The step fixes the origin, $\mathcal T_\ell(0)=0$, for every model and every
channel. The tilted operator $\mathsf U_\ell^{\psi_{\mathsf g}}$ of
\eqref{eq:rg-tilted-reverse-kernel} is unchanged when a constant is added to
$\psi_{\mathsf g}$, descends to the action quotient, and gives the derivative
\begin{equation}
D\mathcal T_\ell(\mathsf g)
=\operatorname{hd}_{\ell+1}\overline{\mathsf U_\ell^{\psi_{\mathsf g}}}\operatorname{asm}_\ell:
\mathcal G_\ell\longrightarrow\mathcal G_{\ell+1};
\label{eq:rg-nonlinear-interaction-derivative}
\end{equation}
in general only at $\mathsf g=0$ may the tilted operator be replaced by $\mathsf U_\ell$, and
also wherever $\psi_{\mathsf g}$ is constant $C^\dagger(z,\cdot)$-almost everywhere for almost
every $z$, as for $\psi_{\mathsf g}=f\circ\chi$ under a deterministic channel $z=\chi(y)$.
\status{ESTABLISHED}

\noindent\emph{Proof.}
A Markov kernel preserves total mass, which gives
\eqref{eq:rg-interaction-evidence-mass}, and
$\mathcal A_\ell(0)=-\log\mathsf U_\ell(1)=0$. A constant added to $\psi_{\mathsf g}$
multiplies numerator and denominator of $\mathsf U_\ell^{\psi_{\mathsf g}}h$ by one
positive number, and the chain rule with part (b) of
\Cref{thm:rg-bounded-action-calculus} gives
\eqref{eq:rg-nonlinear-interaction-derivative}. If $\psi_{\mathsf g}$ equals a number $c(z)$
$C^\dagger(z,\cdot)$-almost everywhere, the factor $e^{-c(z)}$ cancels in
\eqref{eq:rg-tilted-reverse-kernel} and $\mathsf U_\ell^{\psi_{\mathsf g}}h(z)=\mathsf U_\ell h(z)$;
for a deterministic channel $z=\chi(y)$ and $\psi_{\mathsf g}=f\circ\chi$ this holds for almost
every $z$, since $C^\dagger(z,\cdot)$ is then carried by $\chi^{-1}(\{z\})$. $\square$

\remarkheading{Two agents and one exact step}{rem:rg-two-agent-example}
Let $V_0=\{1,2\}$ with $\mathsf Y_{0i}=\{-1,1\}$ and uniform product base
$\bar{\mathsf m}_0=\nu_0$, and perturb by
$\psi(y)=-\xi_1y_1-\xi_2y_2-\xi_{12}y_1y_2$, whose Hoeffding components are
$-\xi_1y_1$, $-\xi_2y_2$, and $-\xi_{12}y_1y_2$, so $\dim\mathcal G_0=3$. The
channel $C_0(y,dz)=\delta_{y_1}(dz)$ keeps the first agent, so $V_1=\{I\}$,
$\mathsf Y_1=\{-1,1\}$, the pushed base is uniform, and $\dim\mathcal G_1=1$. The
reverse kernel sets $y_1=z$ and leaves $y_2$ uniform, so
$\mathcal A_0(\psi)(z)=-\xi_1z-\log\cosh(\xi_2+\xi_{12}z)$, and the exact step is
$\mathcal T_0(\mathsf g)_{\{I\}}(z)=-\xi_1'z$ with
\[
\xi_1'=\xi_1+\tfrac12\log\frac{\cosh(\xi_2+\xi_{12})}{\cosh(\xi_2-\xi_{12})}
=\xi_1+\operatorname{artanh}\bigl(\tanh\xi_2\tanh\xi_{12}\bigr).
\]
The pair coupling enters the coarse field only through the polarization of the
eliminated agent, $\partial\xi_1'/\partial\xi_{12}=\tanh\xi_2$ at $\xi_{12}=0$. At
the origin $D\mathcal T_0(0)$ keeps the field of agent $1$ and sends the field of
agent $2$ and the pair direction to zero, since
$\mathsf U_0(y_2)=\mathsf U_0(y_1y_2)=0$; at a point with $\xi_2\ne0$ the pair
direction has the nonzero image $-\tanh(\xi_2)z$. The drop from dimension $3$ to
$1$ is the obstruction of \Cref{prop:rg-interaction-finite-network-obstruction},
and the multipliers obey the bound $3^1-1=2$ of
\Cref{prop:rg-interaction-native-growth-bound}. \status{ESTABLISHED}

\noindent\emph{Proof.}
Each monomial has zero mean in each of its coordinates, which gives the components
by \eqref{eq:rg-hoeffding-mobius-projectors}. Averaging $e^{-\psi}$ over $y_2$ at
$y_1=z$ gives $e^{\xi_1z}\cosh(\xi_2+\xi_{12}z)$, and subtracting the mean of
$\mathcal A_0(\psi)$ over $z=\pm1$ gives the coefficient $-\xi_1'$; the artanh form
follows from the addition formula for $\cosh$. Differentiating at $\xi_{12}=0$
gives $\tanh\xi_2$, and \eqref{eq:rg-nonlinear-interaction-derivative} at
$\psi=-\xi_2y_2$ gives $\mathsf U_0^{\psi}(y_1y_2)(z)=z\tanh\xi_2$. $\square$

\propositionheading{A product base and a block-product channel preserve the intra-block sector}{prop:rg-interaction-intra-block-invariance}
Suppose that $\bar{\mathsf m}_\ell=\nu_\ell$ is a product law, that blocks
$I\in V_{\ell+1}$ partition $V_\ell$, each identified with its coarse agent, and
that the channel factors as
\begin{equation}
C_\ell(y,dz)=\prod_{I\in V_{\ell+1}}C_I(y_I,dz_I)
\label{eq:rg-block-product-channel}
\end{equation}
with Markov kernels $C_I$. Then $\bar{\mathsf m}_{\ell+1}=\bar{\mathsf m}_\ell
C_\ell$ is a product law, so the convention $\bar{\mathsf
m}_{\ell+1}=\nu_{\ell+1}$ is admissible at the next depth. Let
$\mathcal G_\ell^{\mathrm{intra}}=\{\mathsf g\in\mathcal G_\ell:\mathsf g_B=0\text{
unless }B\subseteq I\text{ for some }I\}$. For every $\mathsf g\in\mathcal
G_\ell^{\mathrm{intra}}$ and every product reference $\nu_{\ell+1}\sim\bar{\mathsf
m}_{\ell+1}$, the components of $\mathcal T_\ell(\mathsf g)$ vanish on every
$B\subseteq V_{\ell+1}$ with $\lvert B\rvert\ge2$; in particular the singleton
sector $\{\mathsf g:\mathsf g_B=0\text{ for }\lvert B\rvert\ge2\}$, which lies in
$\mathcal G_\ell^{\mathrm{intra}}$, is mapped into the singleton sector of the next
depth. \status{ESTABLISHED}

The proof is in \cref{sec:proofs-network-coarsegraining}.

\definitionheading{Scale-indexed evolution and its derivative cocycle}{def:rg-derivative-cocycle}
Let $(\mathscr X_\ell,\lVert\cdot\rVert_\ell)_{\ell\ge0}$ be Banach spaces and let
$\mathcal T_\ell:\mathscr O_\ell\to\mathscr X_{\ell+1}$ be continuously Fr\'echet
differentiable on an open set $\mathscr O_\ell\subseteq\mathscr X_\ell$. The ordered
evolution $\mathcal T_{n\leftarrow\ell}=\mathcal T_{n-1}\circ\cdots\circ\mathcal
T_\ell$, with $\mathcal T_{\ell\leftarrow\ell}=\operatorname{id}_{\mathscr X_\ell}$,
is defined at the points whose successive images stay in the successive domains.
Along an exact orbit $\xi_{\ell+1}=\mathcal T_\ell(\xi_\ell)$ in those domains, the
finite-step multipliers are
\begin{equation}
D_\ell=D\mathcal T_\ell(\xi_\ell)\in\mathcal L(\mathscr X_\ell,\mathscr X_{\ell+1}),
\qquad
D_{n\leftarrow\ell}=D_{n-1}\cdots D_\ell,
\qquad
D_{\ell\leftarrow\ell}=I_{\mathscr X_\ell},
\label{eq:rg-derivative-cocycle}
\end{equation}
with the rightmost factor acting first.

\propositionheading{Ordered composition of the evolution and its derivative}{prop:rg-cocycle-composition}
For $\ell\le n\le r$ and an exact orbit as in \Cref{def:rg-derivative-cocycle},
\begin{equation}
\mathcal T_{r\leftarrow n}\circ\mathcal T_{n\leftarrow\ell}=\mathcal T_{r\leftarrow\ell},
\qquad
D_{n\leftarrow\ell}=D\mathcal T_{n\leftarrow\ell}(\xi_\ell),
\qquad
D_{r\leftarrow n}D_{n\leftarrow\ell}=D_{r\leftarrow\ell}.
\label{eq:rg-cocycle-composition}
\end{equation}
\status{ESTABLISHED}

\noindent\emph{Proof.}
The first and third identities are associativity. The second follows by induction
on $n-\ell$ from the Fr\'echet chain rule applied to $\mathcal
T_{n+1\leftarrow\ell}=\mathcal T_n\circ\mathcal T_{n\leftarrow\ell}$ at $\xi_\ell$.
$\square$

\definitionheading{Finite-step multipliers of the interaction step}{def:rg-interaction-modes}
Instantiate \Cref{def:rg-derivative-cocycle} with $\mathscr X_\ell=\mathcal G_\ell$,
$\mathcal T_\ell$ the exact interaction step, and an exact interaction orbit
$\mathsf g_{\ell+1}=\mathcal T_\ell(\mathsf g_\ell)$. By
\Cref{lem:rg-interaction-step} the one-step multipliers are
\begin{equation}
D_\ell=D\mathcal T_\ell(\mathsf g_\ell)
=\operatorname{hd}_{\ell+1}\overline{\mathsf U_\ell^{\psi_{\mathsf g_\ell}}}\operatorname{asm}_\ell:
\mathcal G_\ell\longrightarrow\mathcal G_{\ell+1},
\label{eq:rg-interaction-mode-generator}
\end{equation}
and $D_{n\leftarrow\ell}$ are their ordered products over the finite depth.

\propositionheading{Native growth bound for the exact interaction cocycle}{prop:rg-interaction-native-growth-bound}
Equip every action quotient $\overline{\mathfrak B}_k$ with the oscillation norm
\eqref{eq:rg-action-oscillation-norm}. Along an exact interaction orbit with
multipliers \eqref{eq:rg-interaction-mode-generator} and for $n>\ell$, the ordered
product telescopes, and its norm is bounded by the target size:
\begin{align}
D_{n\leftarrow\ell}
&=\operatorname{hd}_n\overline{\mathsf U_{n-1}^{\psi_{\mathsf g_{n-1}}}}\cdots
\overline{\mathsf U_\ell^{\psi_{\mathsf g_\ell}}}\operatorname{asm}_\ell,
\label{eq:rg-interaction-cocycle-telescoping}\\
\lVert D_{n\leftarrow\ell}\rVert_{\mathcal G_\ell\to\mathcal G_n}
&\le3^{\lvert V_n\rvert}-1 .
\label{eq:rg-interaction-native-growth-bound}
\end{align}
The bound does not depend on the number of steps, and it is zero when
$V_n=\varnothing$. With $\delta_k$ the coefficient \eqref{eq:rg-action-dobrushin-coefficient} of
the tilted reverse kernel $C^\dagger_{k,\psi_{\mathsf g_k}}$ of
\eqref{eq:rg-tilted-reverse-kernel} for the arrow $C_k$ at $\bar{\mathsf m}_k$, and $\delta^\circ_k$
that of the untilted reverse kernel $C^\dagger_k$,
\begin{equation}
\lVert D_{n\leftarrow\ell}\rVert_{\mathcal G_\ell\to\mathcal G_n}
\le\bigl(3^{\lvert V_n\rvert}-1\bigr)\prod_{k=\ell}^{n-1}\delta_k,
\qquad
\delta_k\le1-e^{-2\lVert[\psi_{\mathsf g_k}]\rVert_{\mathrm{osc}}}\bigl(1-\delta^\circ_k\bigr),
\label{eq:rg-interaction-dobrushin-growth-bound}
\end{equation}
so the multipliers decay geometrically with depth along an orbit of bounded width whose arrows
have $\delta^\circ_k\le\delta^\circ<1$. \status{ESTABLISHED}

\noindent\emph{Proof.}
By \eqref{eq:rg-hoeffding-inverse-and-bounds},
$\operatorname{asm}_{k+1}\operatorname{hd}_{k+1}=I$ on $\overline{\mathfrak
B}_{k+1}$, so every interior pair cancels in $D_{n-1}\cdots D_\ell$, which gives
\eqref{eq:rg-interaction-cocycle-telescoping}. Each tilted operator is positive and
unital by \Cref{thm:rg-bounded-action-calculus}, so $\mathsf U_k^\psi(f-c)=\mathsf
U_k^\psi f-c$ and $\lVert\mathsf U_k^\psi(f-c)\rVert_\infty\le\lVert
f-c\rVert_\infty$ for every constant $c$, and the infimum over $c$ shows that its
quotient has norm at most one. Multiplying by the bounds on
$\operatorname{asm}_\ell$ and $\operatorname{hd}_n$ in
\eqref{eq:rg-hoeffding-inverse-and-bounds} gives
\eqref{eq:rg-interaction-native-growth-bound}. For the refinement, the tilted reverse kernel
$C^\dagger_{k,\psi}$ is a reverse kernel of the normalized joint law
$\widehat{\mathsf m}^\psi_k(dy)C_k(y,dz)$ with
$\widehat{\mathsf m}^\psi_k=e^{-\psi}\bar{\mathsf m}_k/\bar{\mathsf m}_k(e^{-\psi})$, because that
law equals $\bar{\mathsf m}_{k+1}(dz)C^\dagger_k(z,dy)e^{-\psi(y)}/\bar{\mathsf m}_k(e^{-\psi})$.
For bounded $\psi$ the laws $\widehat{\mathsf m}^\psi_k$ and $\widehat{\mathsf m}^\psi_kC_k$ are
equivalent to $\bar{\mathsf m}_k$ and $\bar{\mathsf m}_{k+1}$, so
\Cref{thm:rg-action-lp-contraction} applied to this joint law gives
$\lVert[\mathsf U_k^\psi h]\rVert_{\mathrm{osc}}\le\delta_k\lVert[h]\rVert_{\mathrm{osc}}$, and the
product of these factors replaces the factor one above. For the second inequality, the
coefficient of a kernel is one minus the essential infimum, over pairs of rows, of the mass of
the minimum of the two rows. With $w$ the essential range of $\psi$, the tilt
$e^{-\psi(y)}/\int e^{-\psi(u)}C^\dagger_k(z,du)$ is at least $e^{-w}$, so every tilted row
dominates $e^{-w}$ times its untilted row, the minimum of two tilted rows has at least $e^{-w}$
times the mass of the minimum of the untilted ones, and $w=2\lVert[\psi]\rVert_{\mathrm{osc}}$ by
\eqref{eq:rg-action-oscillation-norm}. $\square$

\propositionheading{Finite networks admit no typed interaction comparison across a reducing step}{prop:rg-interaction-finite-network-obstruction}
Let every agent space be finite. Then $\mathcal G_\ell\cong
L^\infty(\nu_\ell)/\R1$ has dimension
$\lvert\operatorname{supp}\nu_\ell\rvert-1
=\prod_{i\in V_\ell}\lvert\operatorname{supp}\nu_{\ell i}\rvert-1$. If a step reduces
this number, no space admits linear isomorphisms onto both $\mathcal G_\ell$ and
$\mathcal G_{\ell+1}$, so no comparison identification of the interaction spaces
exists across that step; six binary agents blocked into three binary blocks give
$\dim\mathcal G_\ell=63$ and $\dim\mathcal G_{\ell+1}=7$. If the blocking ratios are
the vertex ratios $b_k=\lvert V_k\rvert/\lvert V_{k+1}\rvert\ge b>1$, then
$\lvert V_n\rvert\le\lvert V_\ell\rvert b^{-(n-\ell)}$, and a nonempty $V_n$ forces
$n-\ell\le\log_b\lvert V_\ell\rvert$, so a finite network admits only finitely many
such steps. \status{ESTABLISHED}

\noindent\emph{Proof.}
By \Cref{thm:rg-hoeffding-action-isomorphism}, $\mathcal G_\ell$ is linearly
isomorphic to $L^\infty(\nu_\ell)/\R1$. On a finite space $L^\infty(\nu_\ell)$ is the
space of functions on the support of $\nu_\ell$, the product of the coordinate
supports, and the quotient lowers the dimension by one; linear isomorphisms
preserve dimension. The vertex ratios telescope to
$\lvert V_\ell\rvert/\lvert V_n\rvert=\prod_{k=\ell}^{n-1}b_k\ge b^{n-\ell}$, and
$\lvert V_n\rvert\ge1$ bounds $n-\ell$. $\square$

On finite networks the depth objects are therefore the exact step, its residual
(\Cref{prop:rg-interaction-residual}), and the multipliers, none of which needs a
comparison isomorphism; blocking sequences with independently declared ratios, and
infinite networks, are treated in Supplement B.

\definitionheading{Retained projection and closure residual}{def:rg-retained-projection}
Let $\mathsf P^{\mathrm{ret}}_\ell:\mathcal G_\ell\to\mathcal G_\ell$ be a bounded
idempotent that intertwines the gauge action of
\Cref{cor:rg-hoeffding-gauge-covariance}. The retained output, the coordinate
residual, and the action residual are
\begin{align}
\mathsf g^{\mathrm{ret}}_{\ell+1}&=\mathsf P^{\mathrm{ret}}_{\ell+1}\mathcal T_\ell(\mathsf g_\ell),
\qquad
r^{\mathcal G}_{\ell+1}(\mathsf g_\ell)
=\bigl(I-\mathsf P^{\mathrm{ret}}_{\ell+1}\bigr)\mathcal T_\ell(\mathsf g_\ell),
\label{eq:rg-interaction-coordinate-residual}\\
\overline r^{\mathcal A}_{\ell+1}(\mathsf g_\ell)
&=\overline{\mathcal A}_\ell\operatorname{asm}_\ell\mathsf g_\ell
-\operatorname{asm}_{\ell+1}\mathsf P^{\mathrm{ret}}_{\ell+1}\operatorname{hd}_{\ell+1}
\overline{\mathcal A}_\ell\operatorname{asm}_\ell\mathsf g_\ell .
\label{eq:rg-interaction-action-residual}
\end{align}

\propositionheading{Residual identity and exact-image invariance}{prop:rg-interaction-residual}
The residuals satisfy $\overline r^{\mathcal A}_{\ell+1}=\operatorname{asm}_{\ell+1}r^{\mathcal G}_{\ell+1}$,
so one vanishes exactly when the other does, and
\begin{equation}
\mathcal T_\ell\bigl(\operatorname{Ran}\mathsf P^{\mathrm{ret}}_\ell\bigr)
\subseteq\operatorname{Ran}\mathsf P^{\mathrm{ret}}_{\ell+1}
\quad\Longleftrightarrow\quad
r^{\mathcal G}_{\ell+1}(\mathsf g)=0
\quad\text{for every }\mathsf g\in\operatorname{Ran}\mathsf P^{\mathrm{ret}}_\ell .
\label{eq:rg-interaction-exact-image-invariance}
\end{equation}
A retained update is exact precisely under this exact-image invariance; otherwise it
is a projection scheme with the displayed residual. For $V_{\ell+1}\ne\varnothing$,
\begin{equation}
\operatorname{hd}_{\ell+1}\overline r^{\mathcal A}_{\ell+1}=r^{\mathcal G}_{\ell+1},
\qquad
\frac{\lVert r^{\mathcal G}_{\ell+1}\rVert_{\mathcal G_{\ell+1}}}{3^{\lvert V_{\ell+1}\rvert}-1}
\le\lVert\overline r^{\mathcal A}_{\ell+1}\rVert_{\overline{\mathfrak B}_{\ell+1}}
\le\lVert r^{\mathcal G}_{\ell+1}\rVert_{\mathcal G_{\ell+1}}.
\label{eq:rg-interaction-residual-norm-control}
\end{equation}
\status{ESTABLISHED}

\noindent\emph{Proof.}
Since $\operatorname{asm}_{\ell+1}\operatorname{hd}_{\ell+1}=I$, the first term of
\eqref{eq:rg-interaction-action-residual} is $\operatorname{asm}_{\ell+1}\mathcal
T_\ell(\mathsf g_\ell)$, so the action residual is $\operatorname{asm}_{\ell+1}$
applied to $(I-\mathsf P^{\mathrm{ret}}_{\ell+1})\mathcal T_\ell(\mathsf g_\ell)$,
and $\operatorname{asm}_{\ell+1}$ is injective. An element lies in the range of an
idempotent exactly when the complementary idempotent annihilates it, which is
\eqref{eq:rg-interaction-exact-image-invariance}. Applying
$\operatorname{hd}_{\ell+1}$ gives the coordinate identity, and the norm bounds are
\eqref{eq:rg-hoeffding-inverse-and-bounds}. $\square$

\propositionheading{Radon--Nikodym covariance of the nonlinear interaction step}{prop:rg-interaction-rn-gauge-covariance}
At adjacent depths let $\kappa_\ell$ and $\kappa_{\ell+1}$ be realizations as in
\Cref{cor:rg-hoeffding-gauge-covariance}, transform the measure pair and the
normalized law by pushforward, $\widehat\varrho_\ell=(\kappa_\ell)_\#\varrho_\ell$,
$\widehat{\mathsf m}_\ell=(\kappa_\ell)_\#\mathsf m_\ell$, and
$\widehat{\bar{\mathsf m}}_\ell=(\kappa_\ell)_\#\bar{\mathsf m}_\ell$, and suppose
that the transformed channel satisfies $\widehat C_\ell(\kappa_\ell
y,\kappa_{\ell+1}A)=C_\ell(y,A)$ for every Borel set $A$. Then, as
Radon--Nikodym equivalence classes,
\begin{align}
\widehat{\mathcal A}_\ell(\mathsf K_\ell\psi)&=\mathsf K_{\ell+1}\mathcal A_\ell(\psi),
\label{eq:rg-action-rn-gauge-covariance}\\
\widehat{\mathsf U}_\ell^{\mathsf K_\ell\psi}\mathsf K_\ell h&=\mathsf K_{\ell+1}\mathsf U_\ell^\psi h,
\label{eq:rg-tilted-derivative-gauge-covariance}\\
\widehat{\mathcal T}_\ell\mathsf K^{\mathcal G}_\ell&=\mathsf K^{\mathcal G}_{\ell+1}\mathcal T_\ell .
\label{eq:rg-nonlinear-interaction-gauge-covariance}
\end{align}
If the retained projections also intertwine,
$\widehat{\mathsf P}^{\mathrm{ret}}_k\mathsf K^{\mathcal G}_k=\mathsf K^{\mathcal
G}_k\mathsf P^{\mathrm{ret}}_k$ for $k=\ell,\ell+1$, then both residuals are
covariant,
\begin{align}
\widehat r^{\mathcal G}_{\ell+1}(\mathsf K^{\mathcal G}_\ell\mathsf g)
&=\mathsf K^{\mathcal G}_{\ell+1}r^{\mathcal G}_{\ell+1}(\mathsf g),
\label{eq:rg-coordinate-residual-gauge-covariance}\\
\widehat{\overline r}{}^{\mathcal A}_{\ell+1}(\mathsf K^{\mathcal G}_\ell\mathsf g)
&=\overline{\mathsf K}_{\ell+1}\overline r^{\mathcal A}_{\ell+1}(\mathsf g).
\label{eq:rg-action-residual-gauge-covariance}
\end{align}
\status{ESTABLISHED}

The proof is in \cref{sec:proofs-network-coarsegraining}.

\remarkheading{Grouping, symmetry, and pullback pairs}{rem:rg-trivial-invariant-pairs}
Under \Cref{def:rg-network-scale} at depths $\ell$ and $\ell+1$, three families of
retained projections satisfy \eqref{eq:rg-interaction-exact-image-invariance}.
(a) Let blocks $I\in V_{\ell+1}$ partition $V_\ell$, let the blocks be grouped into
at least two groups, let the channel be block product
\eqref{eq:rg-block-product-channel}, and let $\bar{\mathsf m}_\ell$ be a product
over the groups. The grouping sector at either depth consists of the elements whose
components vanish on every set that meets two groups; the two grouping sectors
form an invariant pair, and when a group contains two blocks the depth-$\ell$
sector contains components coupling agents of different blocks. (b) Let a finite
group act at both depths by product maps as in
\Cref{cor:rg-hoeffding-gauge-covariance} that preserve the references, the measure
pair, and the channel. The fixed sectors of its two actions on components form an
invariant pair onto which the averaging operators are bounded idempotents; the base
law may couple agents of different blocks, as for a law on $\{-1,1\}^{V_\ell}$
whose density is even under $y\mapsto-y$ with a channel that keeps one agent of
each block, where the fixed sector at depth $\ell$ is spanned by the components of
even order. (c) Let the channel be a deterministic Borel surjection $\chi$ and put
$\chi^*[h]=[h\circ\chi]$ on $\overline{\mathfrak B}_{\ell+1}$. Then
$\overline{\mathcal A}_\ell\chi^*=I$ and $\overline{\mathsf U_\ell}\chi^*=I$ for
every base law, so for every retained projection $\mathsf P^{\mathrm{ret}}_{\ell+1}$
the operator
\begin{equation}
\mathsf P^{\mathrm{ret}}_\ell
=\operatorname{hd}_\ell\chi^*\operatorname{asm}_{\ell+1}\mathsf P^{\mathrm{ret}}_{\ell+1}
\operatorname{hd}_{\ell+1}\overline{\mathsf U_\ell}\operatorname{asm}_\ell
\label{eq:rg-pullback-retained-projection}
\end{equation}
is a bounded idempotent whose range, the pullback sector
$\operatorname{hd}_\ell\chi^*\operatorname{asm}_{\ell+1}\operatorname{Ran}\mathsf P^{\mathrm{ret}}_{\ell+1}$,
is mapped by $\mathcal T_\ell$ onto $\operatorname{Ran}\mathsf P^{\mathrm{ret}}_{\ell+1}$.
Transported with the gauge together with the grouping, the group action, and the
channel, each projection intertwines the gauge action. \status{ESTABLISHED}

\noindent\emph{Proof.}
(a) The joint law $\bar{\mathsf m}_\ell(dy)C_\ell(y,dz)$ is a product over the
groups, so a product of group reverse kernels is a version of the reverse kernel,
and for $\mathsf g$ in the depth-$\ell$ sector $\psi_{\mathsf g}$ is a sum over
groups of functions of the agents of one group, so that $\mathcal A_\ell(\psi_{\mathsf
g})$ is by Fubini's theorem a sum over groups of functions of the coarse agents of
one group. If $B$ meets two groups, choose $j\in B$ outside the group of one summand;
$\E_{\ell+1,B}$ of that summand does not depend on $z_j$, so the factor
$I-\E_{\ell+1,V_{\ell+1}\setminus\{j\}}$ of \eqref{eq:rg-hoeffding-mobius-projectors}
annihilates it, as in the proof of \Cref{prop:rg-interaction-intra-block-invariance}.
(b) \Cref{prop:rg-interaction-rn-gauge-covariance} with every hatted object equal to
the original gives $\mathcal T_\ell\mathsf K^{\mathcal G}_\ell=\mathsf K^{\mathcal
G}_{\ell+1}\mathcal T_\ell$ for each group element, so a fixed $\mathsf g$ has a
fixed image, and the average of a representation of a finite group by bounded
operators is a bounded idempotent onto its fixed space. In the example the
components of a set $B$ of binary agents are the multiples of $\prod_{i\in B}y_i$,
on which the flip acts by $(-1)^{\lvert B\rvert}$, and the channel commutes with the
flips at the two depths. (c) The reverse kernel of a deterministic channel is
carried by the fibers of $\chi$, so $\int e^{-h(\chi(y))}C^\dagger(z,dy)=e^{-h(z)}$
and $\mathsf U_\ell(h\circ\chi)=h$ for $\bar{\mathsf m}_{\ell+1}$-almost every $z$.
Idempotence follows from $\overline{\mathsf U_\ell}\chi^*=I$ and
$\operatorname{asm}_{\ell+1}\operatorname{hd}_{\ell+1}=I$, the same identities show
that $\mathsf P^{\mathrm{ret}}_\ell$ fixes the pullback sector, and
$\mathcal T_\ell\operatorname{hd}_\ell\chi^*\operatorname{asm}_{\ell+1}
=\operatorname{hd}_{\ell+1}\overline{\mathcal A}_\ell\chi^*\operatorname{asm}_{\ell+1}=I$
on $\mathcal G_{\ell+1}$. A gauge change permutes the sets that index components,
conjugates the group action, and turns $\chi$ into
$\kappa_{\ell+1}\chi\kappa_\ell^{-1}$, so the transported projections intertwine by
\eqref{eq:rg-hoeffding-gauge-covariance} and, for (c),
\eqref{eq:rg-tilted-derivative-gauge-covariance} at $\psi=0$. $\square$

\openproblemheading{Invariant retained subspaces}{open:rg-invariant-retained-subspace}
Under a product base law and a block-product channel,
\Cref{prop:rg-interaction-intra-block-invariance} gives
\eqref{eq:rg-interaction-exact-image-invariance} for the retained projections onto
the intra-block sector and the singleton sector of the next depth, so the residual
vanishes on them, and \Cref{rem:rg-trivial-invariant-pairs} adds the grouping
sectors of a base law that is a product over groups of blocks, the fixed sectors of
a symmetry of the measure pair and the channel, and the pullback sectors of a
deterministic channel under every base law. Unproved is an invariant pair beyond
these. The problem is posed for finite agent spaces with fixed product references
of full support, at least two blocks, and deterministic block-product channels each
of whose block maps depends on every coordinate of its block, in each of two classes
of base laws of full support, the product laws and all laws. It is settled by a
proper invariant pair, with fixed retained projections, that is invariant under the
step of every base law in a nonempty open subset of the class and whose depth-$\ell$
range is not contained in the sum of the pullback sector of $\mathcal G_{\ell+1}$
and the grouping sector of any partition of the blocks into at least two groups; or
by a proof that, for base laws outside a closed nowhere dense subset of the class,
every proper invariant pair has its depth-$\ell$ range in the sum of the pullback
sector of $\mathcal G_{\ell+1}$ and the grouping sector of some such partition.

\section{Exact closure generically requires hyperedges}
\label{sec:rg-hypergraph-closure}

\propositionheading{Elimination cliques and hyperedge generation}{prop:rg-hyperedge-generation}
(a) Suppose the complete nonnegative density of the law, baseline included, is
factorized against a product or block-product reference. Marginalizing an internal
variable multiplies every incident factor and integrates their product, giving one
factor on the union of the remaining incident scopes. Iterating represents the
normalized law by generally unnormalized factors on the elimination cliques, whose
product is normalized by the inherited total normalizer; the joint law does not
depend on the elimination order, although the displayed factorization does. A
correlated baseline without such a factorization is kept as one global factor and
pushed at the level of laws.

(b) For a finite coarse agent set $V^c$, a finite-valued coarse action $H_o^c$, and
an anchor $z^\circ$, the anchored components satisfy, for $B\subseteq V^c$,
\begin{equation}
H^{c,\circ}_{o,B}(z_B)
=\sum_{B'\subseteq B}(-1)^{\lvert B\rvert-\lvert B'\rvert}
H_o^c(z_{B'},z^\circ_{V^c\setminus B'}),
\qquad
H_o^c(z)=\sum_{B\subseteq V^c}H^{c,\circ}_{o,B}(z_B).
\label{eq:rg-mobius-potentials}
\end{equation}
The expansion exists for every finite-valued function on a finite product, so
closure of the full hypergraph family is automatic. The anchored components are indexed by the same subsets as the
Hoeffding components \eqref{eq:rg-hoeffding-mobius-projectors} and also depend only
on $z_B$, but their values differ, and only the Hoeffding components enter the
exact step. An extended-valued coarse action is kept as one top-order factor, which
closes without forming undefined alternating sums.

(c) The anchor is a covariant background: if the action is gauge invariant, each
component satisfies $H^{c,\circ}_{o,B}(gz_B;gz^\circ)=H^{c,\circ}_{o,B}(z_B;z^\circ)$,
while an anchor held fixed in one coordinate chart gives covariance of the sum
only.

(d) Pairwise closure fails in general. Eliminating the center $y_0$ of a star of
$\{-1,1\}$-valued agents with field $\xi_0$ and couplings $\xi_1,\xi_2,\xi_3$ gives
the factor $\sum_{y_0=\pm1}\exp[y_0(\xi_0+\sum_{r=1}^3\xi_ry_r)]
=2\cosh(\xi_0+\sum_{r=1}^3\xi_ry_r)$, whose negative logarithm has the coefficient
$2\operatorname{sech}^2(\xi_0)\tanh(\xi_0)\xi_1\xi_2\xi_3+O(\xi^5)$ on $y_1y_2y_3$,
nonzero for all sufficiently small nonzero couplings when $\xi_0\ne0$. At
$\xi_0=0$ this coefficient vanishes identically, because
$\log2\cosh(\sum_r\xi_ry_r)$ is even under $y\mapsto-y$, so hyperedge generation is
generic, not universal. \status{ESTABLISHED}

The proof is in \cref{sec:proofs-network-coarsegraining}.

\section{Fine--meta bridge kernels}
\label{sec:rg-cross-scale-kernels}

\definitionheading{Posterior bridge and bridge kernels}{def:rg-posterior-bridge}
At a record $o$ with posterior slice $\boldsymbol\Pi_o$ on the fine latent space and
a coarse channel $C$, the posterior bridge is the joint law
$\mathbb B_o(dy,dz)=\boldsymbol\Pi_o(dy)C(y,dz)$, with coarse marginal
$\boldsymbol\Pi^c_o=\boldsymbol\Pi_oC$ and full reverse kernel
$\mathbb B_o(dy\given z)$. For a fine agent $j$ and a meta-agent $I$, the
$(Y_j,Z_I)$ marginal of $\mathbb B_o$ disintegrated in both directions,
\begin{equation}
\boldsymbol\Pi_{o,j}(dy_j)K_{I\leftarrow j}(y_j,dz_I)
=\boldsymbol\Pi^c_{o,I}(dz_I)K_{j\leftarrow I}(z_I,dy_j),
\label{eq:rg-agent-meta-bayes-adjoint}
\end{equation}
defines the agent-to-meta and meta-to-agent bridge kernels. For a deterministic
statistic $z=\chi(y)$, the lift and the posterior conditional expectation are
\begin{equation}
(\mathsf Jf)(y)=f(\chi(y)),
\qquad
(\mathsf Cg)(z)=\E_{\boldsymbol\Pi_o}[g(Y)\given\chi(Y)=z],
\label{eq:rg-conditional-operators}
\end{equation}
and for a randomized channel $(\mathsf Jf)(y)=\int f(z)C(y,dz)$ and
$(\mathsf Cg)(z)=\int g(y)\mathbb B_o(dy\given z)$.

\propositionheading{Bridge-kernel identities}{prop:rg-bridge-kernels}
For bounded observables $f$ and $g$ the bridge kernels are adjoint,
$\int f(y_j)(K_{I\leftarrow j}g)(y_j)\boldsymbol\Pi_{o,j}(dy_j)
=\int(K_{j\leftarrow I}f)(z_I)g(z_I)\boldsymbol\Pi^c_{o,I}(dz_I)$. The kernel
$K_{j\leftarrow I}$ does not reconstruct the fine posterior, whose full reverse
kernel generally reads all meta variables: a product baseline and a block-product
channel make the baseline reverse kernel factor over blocks, but a likelihood that
couples blocks can destroy this factorization in the posterior, and
posterior-local refinement requires $Y_I\perp Z_{-I}\given Z_I$ under
$\mathbb B_o$. The operators of \eqref{eq:rg-conditional-operators} are adjoint in
$L^2$; for a deterministic statistic $\mathsf C\mathsf J=I$ in
$L^p(\boldsymbol\Pi^c_o)$, and the operators compose by the tower property along a
nested filtration, while for a randomized channel $\mathsf C\mathsf J=I$ fails in
general. \status{ESTABLISHED}

The proof is in \cref{sec:proofs-network-coarsegraining}.

\section{Exact attention between meta-agents}
\label{sec:rg-meta-attention}

A row-stochastic attention matrix is a conditional law of a source given a receiver
and cannot be coarse-grained without the receiver law. The recursive attention state
is therefore the joint edge-event law $\eta$ of \eqref{eq:rg-attention-event-law}
over the oriented edge copies $e:j\to i$, with receiver occupancy $\alpha$ supplied
by a structural receiver-selection law or by the single-edge model of
\Cref{hyp:rg-single-edge-attention}.

\propositionheading{Exact meta-attention pushforward}{prop:rg-meta-attention-pushforward}
For a posterior bridge $\mathbb B_o$ and a partition of the agents into blocks, put
\begin{equation}
\eta^c_{IJ}(z)
=\E_{\mathbb B_o}\Bigl[\sum_{\substack{e:j\to i\\ i\in I,\ j\in J}}\eta_e(Y)\Bigm|Z=z\Bigr],
\qquad
\alpha^c_I(z)=\sum_J\eta^c_{IJ}(z),
\qquad
\beta^c_{IJ}(z)=\frac{\eta^c_{IJ}(z)}{\alpha^c_I(z)},
\label{eq:rg-meta-attention}
\end{equation}
the last ratio on $\{\alpha^c_I>0\}$. The coarse event law is normalized, pushing
then disintegrating is associative under nested partitions, and the next depth
pushes $\eta^c$, not $\beta^c$. If the channel and the selected posterior versions
are gauge equivariant, $\eta^c$, $\alpha^c$, and $\beta^c$ are invariant as values at
corresponding coarse states and covariant as functions of $z$. Events with both
endpoints in $I$ become meta-self events in $\eta^c_{II}$, so the coarse source set
of $I$ contains $I$ whenever $\eta^c_{II}>0$, even when the fine rows exclude
self-edges; for rows that do not depend on the latent state, $\eta^c$ is the block
sum, and \cref{cor:obs-attention-self-edge} gives the uniform case.
\status{ESTABLISHED}

\noindent\emph{Proof.}
Conditional expectation preserves $\sum_e\eta_e=1$, and the tower property gives
associativity. Equivariance transports $\mathbb B_o$ and leaves the scalar event
probabilities unchanged. When the rows do not depend on the latent state the
conditional expectation is the block sum. $\square$

\propositionheading{Evidence-weighted meta-attention in the single-edge model}{prop:rg-attention-single-edge-weights}
Assume \Cref{hyp:rg-single-edge-attention}, fix the latent state $y$, let
$\mathcal Z_i(y)=\sum_j\pi_{ij}e^{-D_{ij}(y)/\tau_i}$ be the normalizer of the row
$\beta^P_{ij}(y)$ of \eqref{eq:obs-attention-posterior}, and put
$\overline{\mathcal Z}_I(y)=\sum_{k\in I}a_{k\mid I}\widetilde{\mathrm c}_k(o,y)\mathcal Z_k(y)$
for a receiver block $I$ with $a_I=\sum_{k\in I}a_k>0$ and $a_{i\mid I}=a_i/a_I$.
On $\{\overline{\mathcal Z}_I(y)>0\}$ the posterior receiver law conditional on the
receiver lying in $I$ and the posterior block row are
\begin{equation}
\alpha_i'(y)=\frac{a_{i\mid I}\widetilde{\mathrm c}_i(o,y)\mathcal Z_i(y)}{\overline{\mathcal Z}_I(y)},
\qquad
\beta^c_{IJ}(y)=\sum_{i\in I}\alpha_i'(y)\sum_{j\in J}\beta^P_{ij}(y).
\label{eq:rg-attention-evidence-weights}
\end{equation}
At that $y$ they coincide with the block sums and disintegration of
\eqref{eq:rg-meta-attention} applied to the posterior edge-label law
$\eta_{ij}(y)=\alpha^P_i(y)\beta^P_{ij}(y)$, whose receiver occupancy is
$\alpha^P_i(y)\propto a_i\widetilde{\mathrm c}_i(o,y)\mathcal Z_i(y)$. The formula holds
posterior-almost everywhere; on a zero-evidence block the row is an arbitrary
conditional version and is removed from the effective support. The equality
$\alpha_i'=a_{i\mid I}$ holds exactly when $\widetilde{\mathrm c}_i(o,y)\mathcal Z_i(y)$ is
constant on the conditional support in $I$, which guarantees equality with the
pre-observation coarse mixture, although the block rows can also agree accidentally
after source coarsening. \status{ESTABLISHED}

\noindent\emph{Proof.}
Bayes' formula gives $P(i,j\given y,o)\propto a_i\pi_{ij}\widetilde{\mathrm
c}_i(o,y)\exp[-D_{ij}(y)/\tau_i]$, because the label-free likelihood factor cancels.
Summing over $j$ gives $P(i\given y,o)\propto a_i\widetilde{\mathrm c}_i(o,y)\mathcal
Z_i(y)$, and conditioning on $i\in I$ gives $\alpha_i'$. Given $i$ the source
posterior is $\beta^P_{ij}(y)$, and the block sums of $P(i,j\given y,o)$ over
$I\times J$ and over $I\times V$ have the stated ratio. $\square$

\propositionheading{Prior-weighted collapse in the per-receiver model}{prop:rg-attention-per-receiver-collapse}
In the per-receiver label model of \Cref{sec:local-attention-elbo}, under its
exclusivity hypothesis \eqref{eq:obs-attention-augmented-likelihood}, let $I$ be a
set of receivers and let a structural coarse channel draw a receiver $i^\star\in I$
with the law $a_{\cdot\mid I}$, independently of the latent state, the labels, and
the record, and output the source label of $i^\star$. Then the posterior law of
$i^\star$ given $(y,o)$ is $a_{\cdot\mid I}$, and the posterior block row is
\begin{equation}
\Pr(\text{source of }i^\star\in J\given y,o)
=\sum_{i\in I}a_{i\mid I}\sum_{j\in J}\beta^P_{ij}(y),
\label{eq:rg-attention-prior-weighted-row}
\end{equation}
with $\beta^P_{ij}$ of \eqref{eq:obs-attention-posterior}; its receiver occupancy is
the selection law, and \eqref{eq:rg-attention-evidence-weights} does not apply to
this model. The two block rows agree when $\widetilde{\mathrm c}_i\mathcal Z_i$ is
constant on the conditional support in $I$ and otherwise can differ: for two
receivers with $a_{\cdot\mid I}=(1/2,1/2)$, block-source masses $\sum_{j\in
J}\beta^P_{ij}$ equal to $1$ and $0$, and $\widetilde{\mathrm c}_i\mathcal Z_i$ equal
to $3$ and $1$, the evidence-weighted mass of $J$ is $3/4$ and the prior-weighted
mass is $1/2$. If all receiver labels are kept, the exact meta-label is their
vector; the collapse to one categorical row is a further coarse channel, and its
conditional relative entropy enters \eqref{eq:rg-vfe-chain-rule}.
\status{ESTABLISHED}

\noindent\emph{Proof.}
The selector enters neither the baseline nor the likelihood, so its posterior equals
its prior and it is conditionally independent of the labels given $(y,o)$. By
\Cref{prop:obs-attention-elbo} the label posterior factors by receiver with row $i$
equal to $\beta^P_i(y)$, which gives \eqref{eq:rg-attention-prior-weighted-row}. In
the example the evidence weights are $3/4$ and $1/4$. $\square$

\remarkheading{Coarse prior and log-sum-exp energy of one row}{rem:rg-attention-log-sum-exp}
For one categorical row with prior $\pi_j$, row energies $D_j$, temperature $\tau$,
and a source partition $j\mapsto J$, the block posterior is the Gibbs row of the
coarse prior and coarse energy
\begin{equation}
\pi^c_J=\sum_{j\in J}\pi_j,
\qquad
D^c_J=-\tau\log\Bigl[\frac1{\pi^c_J}\sum_{j\in J}\pi_j\exp(-D_j/\tau)\Bigr],
\label{eq:rg-attention-log-sum-exp}
\end{equation}
for every block with $\pi^c_J>0$; a block of zero prior has zero posterior mass and
leaves the effective support. The rule associates for nested partitions of one row.
Merging independently normalized receiver rows needs a receiver law: the
evidence-weighted occupancy of \Cref{prop:rg-attention-single-edge-weights} in the
single-edge model and the selection law of
\Cref{prop:rg-attention-per-receiver-collapse} in the per-receiver model.
\status{ESTABLISHED}

\noindent\emph{Proof.}
By construction $\pi^c_J\exp(-D^c_J/\tau)=\sum_{j\in J}\pi_j\exp(-D_j/\tau)$, so the
block sums of the Gibbs row equal the Gibbs row of $(\pi^c,D^c)$ with the same
normalizer. Applying the rule to a coarser partition of the blocks gives the same
sums, which is associativity. $\square$

\propositionheading{Transported meta-messages as marked moments}{prop:rg-attention-hom-moment}
Let each edge copy $e:j\to i$ with $i\in I$ and $j\in J$ carry its dressed link
$\widehat\Theta^x_e$ of \eqref{eq:rg-dressed-boundary-edge}, let
$\rho^{\mathbb W}_x:G\to\GL(\mathbb W_x)$ be the representation of the linear feature tier of
\eqref{eq:rg-linear-cross-scale-maps}, and assume the conditional
Bochner-integrability condition
$\E_{\mathbb B_o}[\sum_e\eta_e(Y)\lVert\rho^{\mathbb W}_x(\widehat\Theta^x_e)\rVert\given Z]<\infty$
almost everywhere, the sum running over these edge copies. Define
\begin{equation}
\begin{split}
\mathcal M^x_{IJ}(z)
&=\E_{\mathbb B_o}\Bigl[\sum_{\substack{e:j\to i\\ i\in I,\ j\in J}}\eta_e(Y)
\rho^{\mathbb W}_x(\widehat\Theta^x_e)\Bigm|Z=z\Bigr],\\
\overline{\mathcal M}^x_{IJ}&=\frac{\mathcal M^x_{IJ}}{\eta^c_{IJ}}
\quad\text{on }\{\eta^c_{IJ}>0\},
\qquad
\mathcal M^{x,\mathrm{att}}_{IJ}=\frac{\mathcal M^x_{IJ}}{\alpha^c_I}
\quad\text{on }\{\alpha^c_I>0\}.
\end{split}
\label{eq:rg-attention-hom-operator}
\end{equation}
Then $\mathcal M^{x,\mathrm{att}}_{IJ}=\beta^c_{IJ}\overline{\mathcal M}^x_{IJ}$ on
$\{\eta^c_{IJ}>0\}$, and the attention-weighted block vanishes on
$\{\alpha^c_I>0,\eta^c_{IJ}=0\}$, where $\overline{\mathcal M}^x_{IJ}$ is a
null-event conditional left undefined. Under a passive root-frame change with no
relabeling of agents, corresponding conditional versions, and the transformed coarse
state $z'$,
$\mathcal M^{x\prime}_{IJ}(z')=\rho^{\mathbb W}_x(g_{r_I})\mathcal M^x_{IJ}(z)\rho^{\mathbb W}_x(g_{r_J})^{-1}$
for $\boldsymbol\Pi^c_o$-almost every $z$, so these operators map the $J$ root frame
to the $I$ root frame. The conditional mean need not lie in $\rho^{\mathbb W}_x(G)$: the average
of the rotations by $\pi/2$ and $-\pi/2$ is zero. The moment is exact only as the
first operator moment on a common source input: for equally likely marks and
features $(I,v)$ and $(-I,-v)$ the separate means vanish while the mean product is
$v$, so exact closure keeps the marked operator--feature kernel, with
\eqref{eq:rg-attention-hom-operator} as a derived moment. Since neither the
definition nor the covariance uses invertibility, both hold for any measurable edge
mark with values in $\operatorname{End}(\mathbb W_x)$ that obeys the two-endpoint law, under
the corresponding integrability condition. In particular, when $\mathbb W_x$ is the
vector sample fiber of \Cref{sec:selective-matrix-links} and $\rho^{\mathbb W}_x$ is its
representation $\rho_x$,
the selective links of \eqref{eq:geo-selective-factorization} obey that law by
\eqref{eq:geo-selective-covariance}, and their dressed event-weighted first moment is
a covariant operator between root frames that need not be invertible.
\status{ESTABLISHED}

\noindent\emph{Proof.}
The first two statements are the definitions of $\beta^c$ and of the ratios. The
dressed link transforms as $\widehat\Theta^{x\prime}_e=g_{r_I}\widehat\Theta^x_eg_{r_J}^{-1}$
by \eqref{eq:rg-dressed-boundary-edge}, the event probabilities are invariant, and
conditional expectation commutes with the fixed linear maps $\rho^{\mathbb W}_x(g_{r_I})$ and
$\rho^{\mathbb W}_x(g_{r_J})^{-1}$. For a general mark, substituting its two-endpoint law and the tree
transport law \eqref{eq:rg-tree-transport-law} into the dressed mark gives the same
root-frame law. The two witnesses are direct computations. $\square$

\propositionheading{Meta-attention and meta-agent kernels split one conditional divergence}{prop:rg-meta-attention-kl-split}
At a fixed outside context $z_{-I}$ let
$\beta^{\mathbb Q}_I(J\given z_{-I})=\mathbb Q(J\given z_{-I})$ and
$K^{\mathbb Q}_{I\leftarrow J}(dz_I\given z_{-I})=\mathbb Q(dz_I\given J,z_{-I})$, and
define the generative counterparts under $\mathbb P(\cdot\given z_{-I},o)$. If the
recognition row and kernels are absolutely continuous with respect to their
generative counterparts, then, in $[0,+\infty]$,
\begin{equation}
\begin{aligned}
&\KL\bigl(\mathbb Q(J,Z_I\given z_{-I})\Vert\mathbb P(J,Z_I\given z_{-I},o)\bigr)
=\KL\bigl(\beta^{\mathbb Q}_I(\cdot\given z_{-I})\Vert\beta^{\mathbb P}_I(\cdot\given z_{-I},o)\bigr)\\
&\qquad+\sum_J\beta^{\mathbb Q}_I(J\given z_{-I})
\KL\bigl(K^{\mathbb Q}_{I\leftarrow J}(\cdot\given z_{-I})\Vert K^{\mathbb P}_{I\leftarrow J}(\cdot\given z_{-I},o)\bigr).
\end{aligned}
\label{eq:rg-meta-attention-kl-split}
\end{equation}
\status{ESTABLISHED}

\noindent\emph{Proof.}
This is the chain rule \Cref{lem:prob-kl-chain-rule} for the categorical coordinate
$J$ followed by $Z_I$. $\square$

\section{Path laws and memory}
\label{sec:rg-path-space}

On a finite ordered path interval the pointwise coarse channel applied to the joint
path law is again a coarse channel, so \Cref{thm:rg-exact-coarse-vfe} applies on
path space; the order parameter is an inference index, not the Fisher duration of
\Cref{ch:relational-inference}. The coarse path law need not be first-order Markov
(\Cref{rem:rg-lumpability-witnesses}); this section says when it is and gives the
exact non-Markov dynamics and the cost of finite memory.

\theoremheading{Standard-Borel strong lumpability}{thm:rg-strong-lumpability}
Let $\chi:\mathsf Y\to\mathsf Z$ be Borel and surjective, and let $\mathsf T$ be a
Markov kernel on $\mathsf Y$. Consider
\begin{equation}
\mathsf T(y,\chi^{-1}A)=\mathsf T(y',\chi^{-1}A)
\quad\text{whenever }\chi(y)=\chi(y'),
\qquad\text{for every Borel }A\subseteq\mathsf Z.
\label{eq:rg-strong-lumpability}
\end{equation}
If a Markov kernel $\mathsf T^\chi$ on $\mathsf Z$ satisfies
\begin{equation}
\chi_\#(\mu\mathsf T)=(\chi_\#\mu)\mathsf T^\chi
\qquad\text{for every probability }\mu\text{ on }\mathsf Y,
\label{eq:rg-lumped-kernel}
\end{equation}
then \eqref{eq:rg-strong-lumpability} holds and $\mathsf T^\chi$ is unique.
Conversely, if \eqref{eq:rg-strong-lumpability} holds and $\chi$ has a Borel right
inverse $\iota$, then $\mathsf T^\chi(z,A)=\mathsf T(\iota(z),\chi^{-1}A)$ is a
Markov kernel satisfying \eqref{eq:rg-lumped-kernel}. A Borel right inverse exists
for a countable coarse space and for a product projection
$\mathsf Y=\mathsf Z\times\mathsf Y'$ with $\mathsf Y'$ nonempty. For finite spaces
the biconditional is the lumpability criterion of \citet[Thm.~6.3.2]{KemenySnell1976}.
\status{ESTABLISHED}

The proof is in \cref{sec:proofs-network-coarsegraining}.

\remarkheading{Two memory witnesses}{rem:rg-lumpability-witnesses}
Strong lumpability is the every-initial-law condition and is not necessary for a
single initial law. On $\mathsf Y=\{1,2,3\}$ with $\chi(1)=\chi(2)=0$ and $\chi(3)=1$,
let $\mathsf T$ send $1$ to $3$, send $2$ to $1$ or $2$ with equal probability, and
send $3$ to $1$ or $3$ with equal probability. Then $\mathsf T(1,\chi^{-1}\{0\})=0$
while $\mathsf T(2,\chi^{-1}\{0\})=1$, so \eqref{eq:rg-strong-lumpability} fails, but
started at $\delta_3$ the chain stays in $\{1,3\}$, on which $\chi$ is injective, and
the coarse process is Markov. A stationary coarse process can also have memory. On
the deterministic cycle $1\to2\to3\to1$ with uniform stationary law and labels
$\chi=(0,0,1)$, the observed histories $(0,0)$ and $(1,0)$ both have positive
probability and the same last symbol, but predict the next symbols $1$ and $0$
respectively, so first-order Markov closure fails under this stationary preparation.
\status{ESTABLISHED}

\noindent\emph{Proof.}
In the first witness the reachable set from $3$ is $\{1,3\}$, and the coarse
transitions are $0\mapsto1$ surely and $1\mapsto0$ or $1$ with equal probability. In
the second, the history $(0,0)$ is the state pair $(1,2)$, followed by $3$, and the
history $(1,0)$ is $(3,1)$, followed by $2$. $\square$

\definitionheading{Regular finite-dimensional filter class}{def:rg-regular-finite-filter}
For the dynamic-closure problem below, an admitted filter is a deterministic functional of the
retained path taking values in a finite-dimensional smooth manifold, with dimension bounded
independently of $\epsilon$ and the horizon. Between assignment jumps it has a time-homogeneous
finite-dimensional stochastic realization with locally Lipschitz coefficients; its initialization
and jump-update maps are continuous in their declared native inputs. A decoded conditional law of
the hidden coordinates has a density with respect to one declared sigma-finite measure on the
hidden-coordinate space, and the decoder is continuously Fr\'echet differentiable from the filter
manifold into that $L^1$ space. Discrete assignment values are treated as separate strata.
This is an additional regularity restriction for this problem: dimension alone does not rule out
Borel encodings of an entire history or of a probability-valued filter into one real coordinate.
The restriction excludes such arbitrary encodings and singular decoded families without a common
dominating measure, while allowing exact finite-dimensional filters that satisfy these conditions.
Changing the regularity class changes the closure question.

\openproblemheading{Dynamic closure of the pointwise parent}{open:rg-pointwise-parent-dynamics}
The static construction of \Cref{thm:rg-pointwise-parent-datum} supplies no fine
evolution and no parent evolution. In a differentiable special case with fine vector
field $\mathscr V_t$, moving deterministic coarse map $\mathsf F_t$, and candidate
parent field $\overline{\mathscr V}_t$, the exact semiconjugacy defect is
\begin{equation}
\delta_t=\partial_t\mathsf F_t+D\mathsf F_t\mathscr V_t-\overline{\mathscr V}_t\circ\mathsf F_t .
\label{eq:rg-pointwise-parent-semiconjugacy-defect}
\end{equation}
For a Markov evolution the corresponding exact condition is semigroup intertwining; a
generator identity is an alternative criterion only on an invariant domain or core,
with generation and uniqueness hypotheses sufficient to lift it to the semigroups.
Neither condition is proved for the static datum together with any autonomy,
moving-membership, path-law, stochastic, delayed, memory, or nonequilibrium
hypothesis. Without a class the question is decided in both directions by
statements already present: strongly lumpable kernels lump
(\Cref{thm:rg-strong-lumpability}), the prepared parent intertwines under exact
refresh (\Cref{thm:rg-memory-prepared-parent}), a parent carrying the terminal
posterior of \Cref{thm:rg-memory-filter} reproduces every path law, and the
stationary cycle of \Cref{rem:rg-lumpability-witnesses} and the finite refresh of
\Cref{rem:rg-memory-partial-refresh} admit no first-order parent on their coarse
labels. The problem is therefore posed for switching diffusions at $c_*$ of the
type of \Cref{def:rg-formation-cycle-experiment}, indexed by the regime parameter
$\epsilon$: finitely many source assignments, a positive-definite quadratic sheet
energy for each assignment with Hessian bounded uniformly in $\epsilon$, diffusion
of the configuration with identity mobility on the sheet of the current assignment,
and one rate-one clock at which the assignment is redrawn from its conditional law
under the restricted optimizer. The channels keep the assignment and an orthogonal
projection of the configuration, fixed or depending on the assignment, as in part
(b) of \Cref{prop:rg-formation-endogenous-cycle}, and on some assignment whose
weight under the restricted optimizer is bounded below in $\epsilon$ they hide
coordinates on which the receiver rows of that assignment depend. The parents are
time-homogeneous Markov evolutions of the retained assignment and coordinates, or of
these together with a filter in the explicitly restricted class of
\Cref{def:rg-regular-finite-filter}, each activated from the pushforward of the recognition law
at the activation time. Two levels are asked: the retained and assignment path law on
windows of fixed length, uniformly over windows placed up to the switching time,
the time after activation by which the expected number of assignment changes
reaches one; and the complete retained and assignment history over horizons on
which the expected number of assignment changes diverges as $\epsilon\to0$. At
either level the problem is settled by a proof that every member of the class
admits such a parent whose retained and assignment path law is within a total
variation tending to zero with $\epsilon$ of the native one, or by a member for
which every such parent keeps that distance bounded below.

\theoremheading{Exact discrete projection-memory recurrence}{thm:rg-projection-memory}
Let $\mathsf X$ and $\mathsf W$ be vector spaces, let $\mathsf T:\mathsf X\to\mathsf X$,
$\mathsf C:\mathsf X\to\mathsf W$, and $\mathsf J:\mathsf W\to\mathsf X$ be linear
with $\mathsf C\mathsf J=I_{\mathsf W}$, and set $\mathsf P^{\mathrm{res}}=\mathsf
J\mathsf C$ and $\mathsf Q=I-\mathsf P^{\mathrm{res}}$, so that $\mathsf
P^{\mathrm{res}}$ is idempotent. For the fine recursion $\xi_{n+1}=\mathsf T\xi_n$
write $w_n=\mathsf C\xi_n$, so that $\mathsf P^{\mathrm{res}}\xi_n=\mathsf Jw_n$.
Then for every $n\ge0$ the unresolved part and the resolved coordinate obey
\begin{align}
\mathsf Q\xi_n
&=(\mathsf Q\mathsf T\mathsf Q)^n\mathsf Q\xi_0
+\sum_{k=0}^{n-1}(\mathsf Q\mathsf T\mathsf Q)^{n-1-k}\mathsf Q\mathsf T\mathsf Jw_k,
\label{eq:rg-unresolved-solution}\\
w_{n+1}
&=\mathsf C\mathsf T\mathsf Jw_n
+\sum_{k=0}^{n-1}\mathsf C\mathsf T\mathsf Q(\mathsf Q\mathsf T\mathsf Q)^{n-1-k}\mathsf Q\mathsf T\mathsf Jw_k
+\mathsf C\mathsf T\mathsf Q(\mathsf Q\mathsf T\mathsf Q)^n\mathsf Q\xi_0 ,
\label{eq:rg-memory-recurrence}
\end{align}
an exact closed recurrence for $w_n$. The operators
$\mathsf C\mathsf T\mathsf Q(\mathsf Q\mathsf T\mathsf Q)^m\mathsf Q\mathsf T\mathsf
J$, $m\ge0$, are the memory kernels and the last term of
\eqref{eq:rg-memory-recurrence} is the unresolved initial-state term; no closure or
Markov assumption is used. \status{ESTABLISHED}

\noindent\emph{Proof.}
Applying $\mathsf P^{\mathrm{res}}$ and $\mathsf Q$ to $\xi_{n+1}=\mathsf T\xi_n$ with
$\xi_n=\mathsf Jw_n+\mathsf Q\xi_n$ and $\mathsf Q^2=\mathsf Q$ gives
$w_{n+1}=\mathsf C\mathsf T\mathsf Jw_n+\mathsf C\mathsf T\mathsf Q\mathsf Q\xi_n$ and
$\mathsf Q\xi_{n+1}=\mathsf Q\mathsf T\mathsf Q\mathsf Q\xi_n+\mathsf Q\mathsf T\mathsf Jw_n$.
The second relation gives \eqref{eq:rg-unresolved-solution} by induction on $n$, and
substituting it into the first gives \eqref{eq:rg-memory-recurrence}. $\square$

For the slow block of a fast--slow coordinatization
(\Cref{def:obs-fast-slow-blocks}) of a linear evolution,
\eqref{eq:rg-memory-recurrence} is the exact discrete slow dynamics; the
projection-operator method goes back to \citet{Nakajima1958} and
\citet{Zwanzig1960}, in continuous time on ensemble densities.

\corollaryheading{Autonomy of the resolved recursion}{cor:rg-resolved-autonomy}
Let $\mathscr S_0\subseteq\mathsf X$ be a class of initial states. The resolved
recursion is autonomous on $\mathscr S_0$, meaning $w_{n+1}=\mathsf C\mathsf T\mathsf
Jw_n$ for every $n\ge0$ and every $\xi_0\in\mathscr S_0$, exactly when the total
correction in \eqref{eq:rg-memory-recurrence} vanishes for every such $n$ and
$\xi_0$. It suffices that all memory kernels vanish and that
$\mathsf C\mathsf T\mathsf Q(\mathsf Q\mathsf T\mathsf Q)^n\mathsf Q\xi_0=0$ for
$n\ge0$ throughout $\mathscr S_0$. On the resolved class $\mathsf Q\xi_0=0$ the
second family is vacuous, and there $\mathsf Q\mathsf T\mathsf J=0$ is sufficient but
not necessary: on $\mathsf X=\R^2$ with $\mathsf W=\R$, $\mathsf C(u,v)=u$,
$\mathsf J(w)=(w,0)$, and $\mathsf T(u,v)=(u,u)$, one has
$\mathsf Q\mathsf T\mathsf Jw=(0,w)\ne0$ for $w\ne0$ while $\mathsf C\mathsf T\mathsf
Q=0$, so every memory kernel vanishes and $w_{n+1}=w_n$ is exactly autonomous.
\status{ESTABLISHED}

\noindent\emph{Proof.}
The biconditional is \eqref{eq:rg-memory-recurrence} read as an identity, and the two
families are its two correction terms. In the witness
$\mathsf Q(u,v)=(0,v)$ and $\mathsf T(0,v)=(0,0)$ give $\mathsf C\mathsf T\mathsf Q=0$,
while $\mathsf T(w,0)=(w,w)$ gives $\mathsf Q\mathsf T\mathsf Jw=(0,w)$. $\square$

\subsection{The hidden history and its filter}
\label{sec:rg-memory-hidden-history}

\definitionheading{Hidden path experiment}{def:rg-memory-hidden-path}
Let $Y_t$ take values in a finite set $\mathsf Y$ with initial law $p_0$ and
transition matrix $\mathsf T$, and let $Z_t=\chi(Y_t)$ take values in a finite set
$\mathsf Z$, with $\mathsf P_z$ the diagonal projector onto $\chi^{-1}(z)$. The
joint path law over $N$ steps and its observed marginal are
\begin{equation}
p_0(y_0)\prod_{t=0}^{N-2}\mathsf T_{y_ty_{t+1}}\prod_{t=0}^{N-1}\mathbf 1_{\chi(y_t)=z_t},
\qquad
\mathbb P^{\mathrm{path}}_N(z_{0:N-1})
=p_0\mathsf P_{z_0}\mathsf T\mathsf P_{z_1}\cdots\mathsf T\mathsf P_{z_{N-1}}\mathbf 1 .
\label{eq:rg-memory-hidden-path}
\end{equation}
This auxiliary model, the law-level counterpart of \Cref{thm:rg-projection-memory},
is fixed before a recognition law on its paths is varied.

\theoremheading{Exact filtering without a projected Markov assumption}{thm:rg-memory-filter}
For an observed word $z_{0:n}$ of positive probability put
\begin{equation}
\mathsf f_{z_{0:n}}=p_0\mathsf P_{z_0}\mathsf T\mathsf P_{z_1}\cdots\mathsf T\mathsf P_{z_n},
\qquad
\widehat p_{z_{0:n}}=\frac{\mathsf f_{z_{0:n}}}{\mathsf f_{z_{0:n}}\mathbf 1},
\qquad
\mathbb P(Z_{n+1}=z\given z_{0:n})=\widehat p_{z_{0:n}}\mathsf T\mathsf P_z\mathbf 1 .
\label{eq:rg-memory-filter}
\end{equation}
These are the exact terminal posterior and next-symbol prediction. A word of zero
probability has zero messages for all its extensions and undefined conditional
readouts. The terminal posterior determines the law of every future suffix, but its
attainable values need not form a finite set. \status{ESTABLISHED}

\noindent\emph{Proof.}
Summing the joint path law with the final hidden state fixed gives
$\mathsf f_{z_{0:n}}$. Replacing $\mathsf P_z$ by $I$ and using $\mathsf T\mathbf
1=\mathbf 1$ proves normalization and prefix consistency, division by the positive
word mass gives the posterior, and a further matrix product gives any suffix.
Nonnegative zero messages stay zero under multiplication. A posterior is a
probability vector on $\mathsf Y$, and sufficiency does not assert that its
attainable subset is finite. $\square$

\subsection{Finite-memory projection}
\label{sec:rg-memory-projection}
\label{sec:rg-memory-stationary}

\definitionheading{Finite-memory projection of a path law}{def:rg-memory-projection}
Let $\mathbb P^{\mathrm{path}}_N$ be any law on $\mathsf Z^N$, stationary or not. For an
integer $k\ge0$ and a position $t$ let
$\mathbf z_t=z_{\max(0,t-k):t-1}$ be the retained context. On contexts of positive
mass define the rows and the projected law
\begin{equation}
\mathsf T^{(k)}_t(z\given\mathbf z)=\mathbb P^{\mathrm{path}}_N(Z_t=z\given\mathbf z_t=\mathbf z),
\qquad
\widehat{\mathbb P}^{(k)}_N(z_{0:N-1})=\prod_{t=0}^{N-1}\mathsf T^{(k)}_t(z_t\given\mathbf z_t),
\label{eq:rg-memory-projection-product}
\end{equation}
recursively, so that a prefix of zero mass has zero extensions and no row is chosen
on a null context.

\theoremheading{Forward-KL optimality and its exact information defect}{thm:rg-memory-projection}
The projected law \eqref{eq:rg-memory-projection-product} is normalized, has the same
context-plus-symbol marginals as $\mathbb P^{\mathrm{path}}_N$, and is positive on
every word of positive $\mathbb P^{\mathrm{path}}_N$ probability. Among all laws with
time-dependent $k$-context rows it uniquely minimizes the forward path divergence, and
the minimum is
\begin{equation}
D^{(k)}_N:=\KL\bigl(\mathbb P^{\mathrm{path}}_N\Vert\widehat{\mathbb P}^{(k)}_N\bigr)
=\sum_{t=0}^{N-1}I\bigl(Z_t;Z_{0:t-k-1}\given Z_{\max(0,t-k):t-1}\bigr),
\label{eq:rg-memory-information-defect}
\end{equation}
an empty discarded history contributing zero. For a competing law
$\widetilde{\mathbb P}$ with rows $\widetilde{\mathsf T}_t$,
\begin{equation}
\KL\bigl(\mathbb P^{\mathrm{path}}_N\Vert\widetilde{\mathbb P}\bigr)-D^{(k)}_N
=\sum_{t,\mathbf z:\ \mathbb P^{\mathrm{path}}_N(\mathbf z_t=\mathbf z)>0}
\mathbb P^{\mathrm{path}}_N(\mathbf z_t=\mathbf z)
\KL\bigl(\mathsf T^{(k)}_t(\cdot\given\mathbf z)\Vert\widetilde{\mathsf T}_t(\cdot\given\mathbf z)\bigr).
\label{eq:rg-memory-excess-risk}
\end{equation}
The defect decreases in $k$ and vanishes for $k\ge N-1$; the projection is
idempotent and commutes with restriction to a shorter prefix; and
$2\operatorname{TV}(\mathbb P^{\mathrm{path}}_N,\widehat{\mathbb P}^{(k)}_N)^2\le D^{(k)}_N$.
The reverse divergence is not controlled: the stationary cycle of
\Cref{rem:rg-lumpability-witnesses} has $\mathbb P^{\mathrm{path}}_3(000)=0$ but
$\widehat{\mathbb P}^{(1)}_3(000)=1/6$. For a path recognition law $\mathbb
Q^{\mathrm{path}}$ supported on the support of $\mathbb P^{\mathrm{path}}_N$,
\begin{equation}
\KL\bigl(\mathbb Q^{\mathrm{path}}\Vert\widehat{\mathbb P}^{(k)}_N\bigr)
-\KL\bigl(\mathbb Q^{\mathrm{path}}\Vert\mathbb P^{\mathrm{path}}_N\bigr)
=\E_{\mathbb Q^{\mathrm{path}}}\log\frac{\mathbb P^{\mathrm{path}}_N}{\widehat{\mathbb P}^{(k)}_N},
\label{eq:rg-memory-recognition-vfe-difference}
\end{equation}
whose worst case over such laws is set by the extrema of the log ratio, not by its
$\mathbb P^{\mathrm{path}}_N$-average. \status{ESTABLISHED}

\theoremheading{Memory identity on standard Borel records}{thm:rg-memory-borel-projection}
Let $\mathbb P^{\mathrm{path}}_N$ be any law on $\mathsf Z^N$ for a record space
$\mathsf Z$ that need not be finite, and for each position $t$ let $\mathbf f_t$ be a
Borel map from $\mathsf Z^t$, a point when $t=0$, to a statistic space. Let
$\mathsf T^{\mathbf f}_t$ be a regular conditional law of $Z_t$ given
$\mathbf f_t(Z_{0:t-1})$ under $\mathbb P^{\mathrm{path}}_N$, let
$\widetilde{\mathsf T}_t$ be any Markov kernel from the statistic space to
$\mathsf Z$, and let $\widetilde{\mathbb P}$ be the law on $\mathsf Z^N$ under which
$Z_t$ given $Z_{0:t-1}=z_{0:t-1}$ has law
$\widetilde{\mathsf T}_t(\mathbf f_t(z_{0:t-1}),\cdot)$. A conditional mutual
information is, as in \Cref{prop:obs-conditioning-information}, the average over the
conditioning variable of the divergence between the conditional joint law and the
product of the conditional marginals. Then, in $[0,+\infty]$, with no
absolute-continuity or finiteness hypothesis, and whatever versions are chosen,
\begin{equation}
\begin{aligned}
\KL\bigl(\mathbb P^{\mathrm{path}}_N\Vert\widetilde{\mathbb P}\bigr)
&=D^{\mathbf f}_N+\sum_{t=0}^{N-1}\E_{\mathbb P^{\mathrm{path}}_N}
\KL\bigl(\mathsf T^{\mathbf f}_t(\mathbf f_t(Z_{0:t-1}),\cdot)\big\Vert
\widetilde{\mathsf T}_t(\mathbf f_t(Z_{0:t-1}),\cdot)\bigr),\\
D^{\mathbf f}_N&=\sum_{t=0}^{N-1}I\bigl(Z_t;Z_{0:t-1}\given\mathbf f_t(Z_{0:t-1})\bigr).
\end{aligned}
\label{eq:rg-memory-borel-identity}
\end{equation}
Hence $D^{\mathbf f}_N$ is the minimum of the forward path divergence over all such
laws $\widetilde{\mathbb P}$, attained when $\widetilde{\mathsf T}_t=\mathsf T^{\mathbf f}_t$
for every $t$; if $D^{\mathbf f}_N<\infty$, subtracting it leaves the sum of expected
row divergences as the excess risk; and the $t$-th term of $D^{\mathbf f}_N$ vanishes
exactly when the law of $Z_t$ given $Z_{0:t-1}$ has a version that depends on the
history only through $\mathbf f_t$. For $\mathbf f_t(z_{0:t-1})=\mathbf z_t$ of
\Cref{def:rg-memory-projection},
$I(Z_t;Z_{0:t-1}\given\mathbf z_t)=I(Z_t;Z_{0:t-k-1}\given\mathbf z_t)$, zero when the
discarded history is empty, and on a finite $\mathsf Z$ the identity is
\eqref{eq:rg-memory-information-defect} with \eqref{eq:rg-memory-excess-risk}. If at
$\mathbb P^{\mathrm{path}}_N$-almost every history the competitor row at
$\mathbf f_t(z_{0:t-1})$ is $e^{-\psi}$ times the source row there for a bounded
$\psi$, that row divergence is at most
$\min\{2\lVert[\psi]\rVert_{\mathrm{osc}},\lVert[\psi]\rVert_{\mathrm{osc}}^2/2\}$, with
the oscillation norm taken under the source row, as in
\Cref{cor:rg-perturbation-divergence-bounds}. \status{ESTABLISHED}

The defect is what a predictor that reads the history only through $\mathbf f_t$
loses by forgetting, and the row sum is what it loses by predicting badly from what it
keeps. Neither the finite hidden chain of \Cref{def:rg-memory-hidden-path} nor a
finite record set enters, so the identity covers record laws with continuous values.
For the window statistic $\mathbf f_t(z_{0:t-1})=\mathbf z_t$ and a competitor that dominates
$\mathbb P^{\mathrm{path}}_N$, \eqref{eq:rg-memory-borel-identity} is
\citet[Sec.~7.3, Theorem~7.3.1, Eqs.~(7.13) and (7.16)]{Gray1990EIT} on standard alphabets; the
theorem removes the domination hypothesis and allows any Borel statistic of the history.

\theoremheading{Stationary projection and the finite hidden information budget}{thm:rg-memory-duration}
Let the hidden chain of \Cref{def:rg-memory-hidden-path} start from a stationary law
$p_0=\bar p$, $\bar p\mathsf T=\bar p$, restricted to its positive support, which is
closed under $\mathsf T$; arbitrary initial preparations are excluded. Write
$\mathrm h_j=H(Z_j\given Z_{0:j-1})$ in natural logarithms, with
$\mathrm h_0=H(Z_0)$, and $\mathrm c_Y=H(Y_0)$. Then the stationary projection is
normalized, stationary, and prefix consistent, and agrees with the source on every
window of length at most $k+1$. The sequence $\mathrm h_j$ decreases to a limit
$\mathrm h_\infty$, and with $\mathrm h^{\mathrm{ex}}_k=\mathrm h_k-\mathrm h_\infty$,
\begin{equation}
\begin{aligned}
D^{(k)}_N&=\sum_{t=k+1}^{N-1}(\mathrm h_k-\mathrm h_t),
&\lim_{N\to\infty}\frac{D^{(k)}_N}N&=\mathrm h^{\mathrm{ex}}_k,\\
\sum_{j=0}^\infty\mathrm h^{\mathrm{ex}}_j&\le\mathrm c_Y,
&\mathrm h^{\mathrm{ex}}_k&\le\frac{\mathrm c_Y}{k+1},\\
0\le\mathrm h^{\mathrm{ex}}_k-\frac{D^{(k)}_N}N&\le\frac{\mathrm c_Y}N,
&D^{(k)}_N&\le(N-k-1)_+\mathrm h^{\mathrm{ex}}_k,
\end{aligned}
\label{eq:rg-memory-duration-identities}
\end{equation}
an empty sum being zero. The series $\sum_j\mathrm h^{\mathrm{ex}}_j$ is the excess
entropy of the observed process in the sense of \citet{CrutchfieldFeldman2003}, so the
second line bounds it by the hidden-state entropy. No irreducibility or forgetting
assumption is used. \status{ESTABLISHED}

The proofs of
\Cref{thm:rg-memory-projection,thm:rg-memory-borel-projection,thm:rg-memory-duration}
are in \cref{sec:proofs-network-coarsegraining}.

The memory cost of a prepared parent whose within-parent refresh is only partial
(\Cref{sec:rg-memory-parent-mechanism}) is the defect $D^{(1)}_N$ of
\Cref{thm:rg-memory-projection} for the parent process.

\section{Conditional refresh and joint-law accuracy}
\label{sec:rg-conditional-refresh-accuracy}

\propositionheading{Conditional entropy dissipation}{prop:rg-conditional-refresh-entropy}
Let $P^\star$ be a probability law on a finite product and let
$\mathfrak B$ be a finite collection of update blocks, possibly overlapping,
with probabilities $p_B$. For each block take a conditional version
$P^\star_B(\cdot\given y_{-B})$ and a measurable
$\alpha^{\mathrm{ref}}_B(y_{-B})\in[0,1]$. Define an auxiliary inference
kernel that keeps the outside and, inside $B$, uses
\begin{equation}
\mathsf K_B^{\mathrm{ref}}(y,dz)
=\delta_{y_{-B}}(dz_{-B})
\left[\alpha^{\mathrm{ref}}_B(y_{-B})P^\star_B(dz_B\given y_{-B})
+(1-\alpha^{\mathrm{ref}}_B(y_{-B}))\delta_{y_B}(dz_B)\right],
\label{eq:rg-conditional-refresh-kernel}
\end{equation}
and set $\mathsf K^{\mathrm{ref}}=\sum_Bp_B\mathsf K_B^{\mathrm{ref}}$.
Both kernels preserve $P^\star$. For a law $\mu$ with
$e=\KL(\mu\Vert P^\star)<\infty$, put
$d_B(y_{-B})=\KL(\mu_B(\cdot\given y_{-B})\Vert
P^\star_B(\cdot\given y_{-B}))$. Then
\begin{equation}
\KL(\mu\mathsf K^{\mathrm{ref}}\Vert P^\star)
\leq e-\sum_Bp_B\int\alpha^{\mathrm{ref}}_Bd_B d\mu_{-B}.
\label{eq:rg-conditional-refresh-dissipation}
\end{equation}
In particular, an estimate
\begin{equation}
\sum_Bp_B\int\alpha^{\mathrm{ref}}_Bd_B d\mu_{-B}
\geq\kappa e-\mathrm b,\qquad 0<\kappa\leq1,\quad \mathrm b\geq0,
\label{eq:rg-weighted-entropy-hypothesis}
\end{equation}
on a class containing every evolving law gives
$e_{n+1}\leq(1-\kappa)e_n+\mathrm b$. For uniform block probabilities,
a lower refresh $\alpha^{\mathrm{ref}}_B\geq\mathrm a>0$, and
\begin{equation}
\KL(\mu\Vert P^\star)\leq
\mathrm C\sum_{B\in\mathfrak B}\int d_B d\mu_{-B},
\text{ with }1\leq\mathrm C<\infty,
\label{eq:rg-block-entropy-factorization}
\end{equation}
one may take $\kappa=\mathrm a/(|\mathfrak B|\mathrm C)$ and $\mathrm b=0$.
The target law and its conditional laws may be non-Gaussian or singular.
\status{ESTABLISHED}

\noindent\emph{Proof.}
The refresh probability depends only on the held outside; each conditional
mixture therefore preserves its conditional target. Under $\mu$, the
outside marginal also remains fixed. Conditional convexity bounds the
new conditional divergence by
$(1-\alpha^{\mathrm{ref}}_B)d_B$. The chain rule
\cref{lem:prob-kl-chain-rule}(a) cancels the unchanged outside divergence,
giving the single-block inequality. Convexity across the random block
gives \eqref{eq:rg-conditional-refresh-dissipation}. The two stated
dissipation assumptions give the recurrence and its constant.
\hfill$\square$

The factorization in \eqref{eq:rg-block-entropy-factorization} is an
additional dependence estimate. Approximate tensorization is formulated
by \citet[Eq.~(1.4), Prop.~1.1]{CaputoMenzTetali2015}; their sufficient
conditions use bounded conditional ratios. The convexity hypotheses of
\citet[Thms.~3.1--3.2]{AscolaniLavenantZanella2024} supply another
sufficient setting. Neither is a consequence of having accurate
individual conditional updates.

\propositionheading{Target changes and numerical transfer}{prop:rg-refresh-target-transfer}
For a prescribed sequence of equivalent targets $P^\star_n$, suppose
step $n$ satisfies \eqref{eq:rg-weighted-entropy-hypothesis} with
$\kappa_n,\mathrm b_n$, and
$\log(dP^\star_n/dP^\star_{n+1})\leq\mathrm s_n<\infty$ almost surely.
Starting with finite $e_0=\KL(\mu_0\Vert P^\star_0)$, the ideal sequence
satisfies
\begin{equation}
e_{n+1}\leq(1-\kappa_n)e_n+\mathrm b_n+\mathrm s_n,
\qquad e_n=\KL(\mu_n\Vert P^\star_n).
\label{eq:rg-refresh-moving-target}
\end{equation}
If numerical transition kernels differ from these ideal kernels uniformly
in TV by at most $\varepsilon_n$, with the same initial law, their
complete path TV through step $n$ is at most
$\sum_{j<n}\varepsilon_j$, and their endpoint law satisfies
\begin{equation}
\operatorname{TV}(\widehat\mu_n,P^\star_n)
\leq\sum_{j<n}\varepsilon_j+\sqrt{e_n/2}.
\label{eq:rg-refresh-numerical-transfer}
\end{equation}
A common record channel contracts the ideal endpoint KL and the total TV
bound. The ideal KL direction is $\mu_n$ relative to $P^\star_n$.
\status{ESTABLISHED}

\noindent\emph{Proof.}
The change of reference adds
$\E_{\mu_{n+1}}\log(dP^\star_n/dP^\star_{n+1})\leq\mathrm s_n$
to the post-update divergence. Finite initial divergence, the dissipation
bound and the bounded log ratio justify this step successively.
Telescoping the numerical and ideal path kernels charges at most one
$\varepsilon_j$ per transition. Pinsker's inequality and the triangle
inequality give the endpoint bound. Data processing applies to the same
record channel on both laws. \hfill$\square$

\remarkheading{Accuracy assumptions and scientific dynamics}{rem:rg-refresh-accuracy-scope}
Exact block conditionals do not imply
\eqref{eq:rg-block-entropy-factorization}. For
$P^\star=(\delta_{00}+\delta_{11})/2$, singleton refreshes leave
$\mu=\delta_{00}$ fixed: both conditional divergences are zero, while
$\KL(\mu\Vert P^\star)=\log2$. A point-mass initializer against a
nonatomic target has infinite KL and is outside the finite-entropy
recurrence. Numerical TV accuracy alone does not give numerical KL
accuracy. If a parameter is learned from the same evolving state, its
realized path cannot be treated as the prescribed target sequence of
\cref{prop:rg-refresh-target-transfer} without a corresponding joint or
conditional-law argument. The auxiliary kernel in
\cref{prop:rg-conditional-refresh-entropy} does not identify the
scientific dynamics or establish a parent in the class of
\cref{open:rg-pointwise-parent-dynamics}.
\status{ESTABLISHED}

\noindent\emph{Proof.}
The two-state support gives the first conditional and joint divergences
directly; singularity gives the initializer claim. For example, adding
any positive point mass to a nonatomic target can make TV arbitrarily
small while making forward KL infinite. Conditioning on an endogenously
chosen parameter path changes the conditional transition laws; the
prescribed-kernel proof does not assert otherwise. A posterior integration
kernel is separate data from a scientific evolution.
\hfill$\square$

\section{The exact retained state}
\label{sec:rg-closure-theorem}

Exact closure keeps every object of \Cref{tab:rg-retained-state}, listed with its
exact operation, source statement, and admission hypotheses, uniformly in the
number and arrangement of the agents.

\begin{sciencetable}[!t]{The retained state of exact agent-network coarse-graining at $c_*$.}
\label{tab:rg-retained-state}
\label{thm:rg-complete-effective-theory}
\label{cor:rg-complete-analytic-tier}
\begin{tabularx}{\textwidth}{>{\raggedright\arraybackslash}p{0.21\textwidth}>{\raggedright\arraybackslash}X>{\raggedright\arraybackslash}p{0.27\textwidth}}
\toprule
\textbf{Retained object} & \textbf{Exact operation and standing} & \textbf{Admission hypotheses} \\
\midrule
Normalized generative law, posterior versions, recognition law & One channel pushes all three; the variational free energy decomposes by the chain rule. Derived: \Cref{thm:rg-exact-coarse-vfe}, \Cref{prop:rg-measure-pair-composition}. & Recognition-independent channel; finite positive evidence; recognition law absolutely continuous with respect to the posterior. \\
Reference probability, evidence measure, effective action & Radon--Nikodym action by the conditional partition function, composing along towers. Derived: \Cref{thm:rg-effective-action}. & Probability reference pushed with the evidence measure; reverse conditionals from one joint tower. \\
Bounded action calculus & Analytic action increment, tilted derivatives, contraction and oscillation bounds. Derived: \Cref{thm:rg-bounded-action-calculus}, \Cref{thm:rg-action-lp-contraction}. & Bounded perturbations; jointly measurable reverse kernel. \\
Interaction coordinates and the exact step & Hoeffding isomorphism and exact nonlinear step with its closure residual. Derived: \Cref{thm:rg-hoeffding-action-isomorphism}, \Cref{lem:rg-interaction-step}, \Cref{prop:rg-interaction-residual}. Retained projection: declared. & Product-equivalent normalized laws at both depths, on fixed vertex sets (\Cref{def:rg-network-scale}). \\
Finite-step multipliers & Ordered products of tilted reverse operators, bounded by $3^{\lvert V_n\rvert}-1$ times the product of their contraction coefficients. Derived: \Cref{def:rg-interaction-modes}, \Cref{prop:rg-interaction-native-growth-bound}. & An exact orbit in the bounded tier. \\
Induced hyperedges & Elimination cliques and M\"obius components; generically nonpairwise. Derived: \Cref{prop:rg-hyperedge-generation}. & Finite-valued coarse action, or one top-order factor. \\
Bridge kernels & Pair-marginal disintegrations of the posterior bridge. Derived: \Cref{prop:rg-bridge-kernels}. Gauge-covariant versions: declared beyond the compact and finite cases (\Cref{open:cg-equivariant-versions}). & Regular disintegration on standard Borel spaces. \\
Roots, based holonomy, dressed boundary marks & Raw root-framed marks and their joint law, pushed without a holonomy-blind claim. Derived: \Cref{sec:rg-gauge-cross-scale}. Trees: declared. & Rooted spanning forests of the blocks. \\
Joint attention-event law & Event pushforward and disintegration on positive occupancy. Derived: \Cref{prop:rg-meta-attention-pushforward}, \Cref{prop:rg-attention-single-edge-weights}, \Cref{prop:rg-attention-per-receiver-collapse}. & Receiver law of the label model. \\
Marked operator--feature kernels & Event-weighted first operator moments. Derived: \Cref{prop:rg-attention-hom-moment}. & Conditional Bochner integrability. \\
Path laws and memory & Lumped kernel when strongly lumpable; otherwise memory kernels, filters, and finite-memory projections with their defects. Derived: \Cref{thm:rg-strong-lumpability}, \Cref{thm:rg-projection-memory}, \Cref{thm:rg-memory-projection}, \Cref{thm:rg-memory-duration}. & The full path law, or the memory kernels and initial-state term. \\
\bottomrule
\end{tabularx}
\end{sciencetable}

The blocks and the channel enter the table as inputs, and the smallest case in which
three two-dimensional blocks are asked for shows that the law cannot supply them.

\propositionheading{Law and Fisher geometry do not canonically agentize blocks}{prop:hist-law-fisher-agentization-nogo}
At the uniform law on seven outcomes, no rule that depends only on the unlabeled law
and its Fisher metric and is natural under every relabeling of the outcomes can
select an unordered decomposition of the six-dimensional tangent space into three
two-dimensional blocks. Declared typing or other symmetry-breaking data can supply
blocks; the law and its Fisher tensor alone do not. \status{ESTABLISHED}

The proof is in \cref{sec:proofs-network-coarsegraining}.

\section*{Summary for physicists}

One exact coarse-graining step pushes normalized laws through a Markov channel; the
effective action is the Radon--Nikodym derivative of the pushed evidence measure
with respect to a pushed reference probability, the conditional partition function
of the fine action over each coarse state, analytic in bounded changes of the fine
action and never expanding the supremum norm. In Hoeffding coordinates, one
component per set of coupled agents, the step is an exact nonlinear map whose
finite-step multipliers telescope to one extraction, a product of tilted
conditional expectations, and one assembly, so in the bounded tier they are bounded by
$3^{\lvert V_n\rvert}-1$ at every depth and decay with depth when the arrows are stochastic;
the polynomial couplings of Gaussian agents lie outside that tier (\Cref{rem:rg-bounded-tier-scope}). Exactness costs state: eliminated agents
generate hyperedges unless a symmetry forbids them, attention is carried as the
joint event law, transported messages are moments of marked kernels, and the coarse
path law generally has memory, with an exact memory-kernel recurrence and, for
memory length $k$, a finite-memory cost of at most the hidden-state entropy divided
by $k+1$ per step.
Conditional refresh can control a joint law only with an additional
dependence and initialization estimate
(\cref{prop:rg-conditional-refresh-entropy,prop:rg-refresh-target-transfer});
these inference estimates do not establish scientific parent dynamics.

\chapter{The Multiscale Graph VFE and Closing the Tower}
\label{ch:full-graph-meta-agent-vfe}

This chapter asks what probabilistic structure places adaptive event laws and links
on the interaction multigraph, random partitions, and the influence of parents
on their children inside one variational free energy, and how a tower of such depths
closes. The answer is the free energy of one normalized generative law and one
recognition law over the whole tower (\Cref{def:full-graph-tower-vfe}). For an
adaptive tower both laws factor top-down, each depth conditioned on the whole
next-depth state and the structural mark (\Cref{lem:full-graph-adaptive-reversal}),
and the free energy splits into observation, top-prior, partition, graph, and
parent--child terms (\Cref{thm:full-graph-vfe-decomposition}). A finite acyclic tower
needs a declared latent root (\Cref{prop:obs-declared-root-unavoidable});
participation across depths enters through the inference loop, parameter learning
(\Cref{prop:full-graph-aggregate-prior}), and delayed feedback
(\Cref{prop:full-graph-episode-joint,prop:full-graph-episode-lifted-joint}). The formation witnesses of
\Cref{ch:finite-depth-formation} are adaptive towers, and \Cref{ch:philosophy} reads
the three loci of participation.

The depth $\ell$ runs from $0$ to a finite $L$. At depth $\ell$, $Z_\ell$ are the
coordinates, $\mathfrak N_\ell$ the graph datum, $Y_\ell=(Z_\ell,\mathfrak N_\ell)$ the
state on a standard Borel space $\mathsf Y_\ell$ that here includes the graph datum,
$R_\ell\in\mathsf R_\ell$ the structural mark with realized value $\gamma$, and
$\mathfrak h_\ell$ the retained root, holonomy, and boundary marks. The whole tower is
$\mathbf Y$, row normalizers are $\mathcal Z$, and $\mathbb M$ is an input law:
generative, posterior, or recognition. In the episode-indexed models of
\Cref{sec:obs-participatory}, $\widetilde o_n$ is a latent copy of the record $o_n$,
$\epsilon_{\mathrm e}$ is the constant step of the aggregate parameter update, and
$\mathrm T$ is the temperature of the descent kinetics, the symbol of
\Cref{ch:finite-depth-formation}.

Scope. In an adaptive tower the downward kernel is a reversal kernel, and the coupling
it carries is not read as causal downward influence; declining that reading does not
refute a causal reading of a separately declared top-down hierarchy. The composite
potential of \Cref{def:full-graph-composite-potential} is not claimed to be an
evidence bound, and delayed feedback is not claimed to support a sustained
nonequilibrium reading. \status{NOT-CLAIMED} Open items are the persistent hierarchy
(\Cref{open:full-graph-hierarchy}), the accounting across episodes
(\Cref{open:full-graph-episode-accounting}), the dynamic closure of the pointwise
parent (\Cref{open:rg-pointwise-parent-dynamics}), and a natural partition selector
(\Cref{open:cg-natural-selector}).

\section{The full answer before decomposition}
\label{sec:full-graph-vfe-global}

\definitionheading{Tower variational free energy}{def:full-graph-tower-vfe}
Write $X_*=X(c_*)$, fix an admitted observation $o\in\mathsf
O^{\mathrm{reg}}_{\theta,X_*}$ and a finite depth $L$, and let $\mathbf Y$ contain
every random object of the multiscale model: belief and model coordinates at depths
$0,\ldots,L$, links, event laws, retained marks, structural marks, and auxiliary
latents. Declare one normalized generative kernel $\mathbb
P_\theta(do,d\mathbf Y\given X_*)$ and one recognition kernel $\mathbb
Q_\lambda(d\mathbf Y\given o,X_*)$. With posterior slice
$\boldsymbol\Pi_{\theta,o,X_*}$ and evidence density $p_\theta(o\given
X_*)\in(0,\infty)$, the tower variational free energy is
\begin{equation}
\boxed{
\Fenergy^{\mathrm{ext}}_{\mathrm{tower}}[\mathbb Q_\lambda;X_*,o]
=-\log p_\theta(o\given X_*)
+\KL\bigl(\mathbb Q_\lambda(\cdot\given o,X_*)\Vert\boldsymbol\Pi_{\theta,o,X_*}\bigr).}
\label{eq:full-graph-vfe-global}
\end{equation}
By the evidence identity of \Cref{ch:elbo}, on the finite-density domain it equals
\begin{equation}
\Fenergy_{\mathrm{tower}}
=\E_{\mathbb Q_\lambda}\Bigl[\log\frac{q_\lambda(\mathbf Y\given o,X_*)}{p_\theta(o,\mathbf Y\given X_*)}\Bigr].
\label{eq:full-graph-vfe-density}
\end{equation}
Every later term is an accounting decomposition of this object or is named a
composite potential.

\section{A normalized hierarchical graph model}
\label{sec:full-graph-vfe-joint}

\definitionheading{Closed scale datum}{def:full-graph-scale-datum}
At depth $\ell$ let $V_\ell$ be the vertex set and let $Z_\ell=(Z_{i,\ell})_{i\in
V_\ell}$ contain the full probabilistic coordinates of that depth, belief and model
presentations rather than only their marginals. The graph datum lives on a finite
interaction multigraph $\mathfrak I_\ell=(V_\ell,E_\ell,\bar{\ })$
(\Cref{def:geo-graph-links}), at depth $0$ the multigraph of \cref{sec:intro-conventions} and at deeper depths
induced by the link lift of \Cref{def:full-graph-coarse-links}. Put
\begin{equation}
\eta_\ell=(\eta^{b,\ell},\eta^{m,\ell}),
\qquad
\Theta^{x,\ell}=(\Theta^{x,\ell}_e)_{e\in E_\ell},
\qquad
\mathfrak N_\ell=(\eta_\ell,\Theta^{b,\ell},\Theta^{m,\ell},\mathfrak h_\ell),
\label{eq:full-graph-closed-scale-datum}
\end{equation}
where $\eta^{x,\ell}=(\eta^{x,\ell}_e)_{e\in E_\ell}$ is a normalized event law on the
copies, $\Theta^{x,\ell}_e\in G$ is the channel-$x$ link of the copy $e$ with
$\Theta^{x,\ell}_{\bar e}=(\Theta^{x,\ell}_e)^{-1}$, and $\mathfrak h_\ell$ contains
the root, based-holonomy, and dressed boundary marks of
\Cref{sec:rg-gauge-cross-scale}. Each copy carries one link, one event mass, and one
divergence per channel; endpoint notation abbreviates a copy only when it is unique,
and $\eta^{x,\ell}_{ij}$ is the total mass of the copies $j\to i$. Occupancies and
rows are derived from $\eta_\ell$ as in \eqref{eq:rg-attention-event-law} and do not
replace it in the recursive state. The observation and $X_*$ stay outside every
coarse channel.

\definitionheading{Common structural channel}{def:full-graph-structural-channel}
A structural step is one recognition-independent Markov channel
\begin{equation}
C^{\mathrm{struct}}_{\ell+1\leftarrow\ell}:
\mathsf Y_\ell\rightsquigarrow\mathsf R_\ell\times\mathsf Y_{\ell+1},
\label{eq:full-graph-common-structural-channel}
\end{equation}
whose output includes $(R_\ell,\eta_{\ell+1},\mathfrak h_{\ell+1})$, with augmented
transition state
\begin{equation}
\widehat{\mathsf Y}_{\ell+1\leftarrow\ell}
:=(\mathsf Y_\ell,R_\ell,\mathsf Y_{\ell+1}),
\label{eq:full-graph-augmented-transition-state}
\end{equation}
which does not duplicate $R_\ell$. The same channel is applied to the fine generative
marginal, the selected posterior, and the recognition law,
\begin{equation}
\widehat{\mathbb M}^{\mathrm{struct}}_{\ell+1\leftarrow\ell}(dy_\ell,d\gamma,dy_{\ell+1})
:=\mathbb M_\ell(dy_\ell)C^{\mathrm{struct}}_{\ell+1\leftarrow\ell}(y_\ell;d\gamma,dy_{\ell+1}),
\qquad
\mathbb M\in\{\mathbb P,\boldsymbol\Pi,\mathbb Q\},
\label{eq:full-graph-common-structural-pushforwards}
\end{equation}
so the three laws of $R_\ell$ may differ while the channel does not. If the channel is
realized by an auxiliary seed, the same measurable map and seed law define all three
pushforwards; a blocking rule that depends on the input law is not admitted.

\hypothesisheading{Depth-Markov hierarchy}{hyp:full-graph-depth-markov}
The generative law factors top-down as
\begin{align}
\mathbb P_\theta(do,d\mathbf Y\given X_*)
={}&\mathbb P^{\mathrm{obs}}_\theta(do\given Z_0,X_*)\mathbb P_L(dY_L\given X_*)
\nonumber\\
&\times\prod_{\ell=0}^{L-1}
\mathbb P^\ell_R(dR_\ell\given Y_{\ell+1},X_*)
\mathbb P^\ell_{\mathfrak N}(d\mathfrak N_\ell\given Y_{\ell+1},R_\ell,X_*)
K^\ell_\downarrow(dZ_\ell\given Y_{\ell+1},R_\ell,\mathfrak N_\ell,X_*),
\label{eq:full-graph-generative-factorization}
\end{align}
with the kernels composed from depth $L$ downward, and recognition factors along the
same order,
\begin{align}
\mathbb Q_\lambda(d\mathbf Y\given o,X_*)
={}&\mathbb Q_L(dY_L\given o,X_*)\prod_{\ell=0}^{L-1}
\mathbb Q^\ell_R(dR_\ell\given Y_{\ell+1},o,X_*)
\nonumber\\
&\times\mathbb Q^\ell_{\mathfrak N}(d\mathfrak N_\ell\given Y_{\ell+1},R_\ell,o,X_*)
\mathbb Q^\ell_Z(dZ_\ell\given Y_{\ell+1},R_\ell,\mathfrak N_\ell,o,X_*).
\label{eq:full-graph-recognition-factorization}
\end{align}
Each depth-$\ell$ factor reads the complete next-depth state
$Y_{\ell+1}=(Z_{\ell+1},\mathfrak N_{\ell+1})$ and no deeper variable; this excludes
direct dependence on depths $\ell+2,\ldots,L$ given $(Y_{\ell+1},R_\ell)$, leaves
arbitrary dependence within a depth and between adjacent depths, and is not a
mean-field assumption. In a declared top-down hierarchy
$K^\ell_\downarrow$ is where downward influence enters, and the bottom-up coarse map
belongs to recognition or to an induced pushforward, never to a second generative
arrow at the same depth.

\lemmaheading{Depth-Markov reversal of an adaptive tower}{lem:full-graph-adaptive-reversal}
Let $\mathbb M_0$ be a probability law on $\mathsf Y_0$, and generate the adaptive
tower bottom-up through the common structural channels,
\begin{equation}
\mathbb M_0(dy_0)\prod_{\ell=0}^{L-1}
C^{\mathrm{struct}}_{\ell+1\leftarrow\ell}(y_\ell;d\gamma_\ell,dy_{\ell+1}).
\label{eq:full-graph-adaptive-tower-law}
\end{equation}
For every $\ell$, the conditional law of $R_\ell$ given all variables
$(Y_{\ell+1},R_{\ell+1},\ldots,Y_L)$ of greater depth has a version that depends only
on $Y_{\ell+1}$, and the conditional law of $Y_\ell$ given $R_\ell$ and those variables
has a version that depends only on $(R_\ell,Y_{\ell+1})$. Hence
\eqref{eq:full-graph-adaptive-tower-law} factors as in
\eqref{eq:full-graph-generative-factorization} without its observation factor, with
$\mathbb P^\ell_R$, $\mathbb P^\ell_{\mathfrak N}$, and $K^\ell_\downarrow$ selected
versions of these conditional laws. The same holds for the recognition input $\mathbb
Q_0(dy_0\given o,X_*)$, for the selected posterior slice, and for the generative input
$\mathbb P_0(do,dy_0\given X_*)$ whenever its observation kernel reads only $Z_0$, so
\Cref{hyp:full-graph-depth-markov} holds automatically for an adaptive tower. These
factors are reversal, or bridge, kernels in the sense of
\eqref{eq:rg-agent-meta-bayes-adjoint}, not independently chosen structural
channels, and when the parent coordinates are a function of the child coordinates and
the structural mark, $K^\ell_\downarrow$ is supported on the child configurations
compatible with the realized parent. \status{ESTABLISHED}

\noindent\emph{Proof.}
For each $\ell$, the variables $(R_{\ell+1},Y_{\ell+2},\ldots,Y_L)$ are produced from
$Y_{\ell+1}$ by kernels that read neither the observation nor any of
$(Y_0,R_0,\ldots,Y_\ell,R_\ell)$, so they are conditionally independent of the latter
given $Y_{\ell+1}$, which gives both version statements. Regular conditional laws
exist on standard Borel spaces \citep{Kallenberg2021}, and successive disintegration
from depth $L$ downward gives the factorization. A conditional law of $Z_\ell$ given a
value of a function of $(Z_\ell,R_\ell)$ charges only the preimage of that value.
$\square$

\propositionheading{The next-depth graph datum cannot be dropped}{prop:full-graph-dropped-datum}
If the depth-$\ell$ factors of an adaptive tower read only $(Z_{\ell+1},R_\ell)$, the
product $\mathbb P^{\mathrm{red}}$ of the corresponding selected conditional laws of
the tower law $\mathbb P^{\mathrm{tow}}$ of \eqref{eq:full-graph-adaptive-tower-law}
differs from it by
\begin{equation}
\KL\bigl(\mathbb P^{\mathrm{tow}}\Vert\mathbb P^{\mathrm{red}}\bigr)
=\sum_{\ell=0}^{L-1}
I_{\mathbb P^{\mathrm{tow}}}\bigl(Y_\ell;\mathfrak N_{\ell+1}\given Z_{\ell+1},R_\ell\bigr).
\label{eq:full-graph-dropped-datum-cost}
\end{equation}
The cost vanishes exactly when every $\mathfrak N_{\ell+1}$ is conditionally
independent of $Y_\ell$ given $(Z_{\ell+1},R_\ell)$. For a deterministic event
pushforward it is positive as soon as two event laws in the support, at the same
$(Z_{\ell+1},R_\ell)$, are pushed to different coarse event laws: for four fine
vertices with the partition $\{0,1\}\mapsto I$, $\{2,3\}\mapsto J$, constant
coordinates, trivial links and marks, and $\eta_0$ equal with probability $1/2$ each
to the unit mass on a copy $2\to0$ or on a copy $1\to0$, the pushed $\eta_1$ is the
unit mass on $(I,J)$ or on $(I,I)$, and the divergence is $\log2$.
\status{ESTABLISHED}

\noindent\emph{Proof.}
By \Cref{lem:full-graph-adaptive-reversal} the reduced product replaces each
conditional law of $Y_\ell$ given $(R_\ell,Y_{\ell+1})$ by its conditional law given
$(R_\ell,Z_{\ell+1})$, and the chain rule \Cref{lem:prob-kl-chain-rule} turns each
replacement into the displayed conditional mutual information. In the witness the
joint law of $(\mathfrak N_0,\mathfrak N_1)$ has two atoms of mass $1/2$ while the
reduced product is the product of their marginals, which puts mass $1/2$ on pairs
that violate the pushforward. $\square$

\theoremheading{Exact hierarchical decomposition}{thm:full-graph-vfe-decomposition}
Assume \Cref{hyp:full-graph-depth-markov}, which holds for an adaptive tower by
\Cref{lem:full-graph-adaptive-reversal}, together with common domination and the
integrability needed for the finite log-density form. Then
\begin{align}
\Fenergy_{\mathrm{tower}}
={}&-\E_{\mathbb Q_\lambda}\log p^{\mathrm{obs}}_\theta(o\given Z_0,X_*)
+\KL(\mathbb Q_L\Vert\mathbb P_L)
\nonumber\\
&+\sum_{\ell=0}^{L-1}\E_{\mathbb Q_\lambda}\bigl[
\KL(\mathbb Q^\ell_R\Vert\mathbb P^\ell_R)
+\KL(\mathbb Q^\ell_{\mathfrak N}\Vert\mathbb P^\ell_{\mathfrak N})
+\KL(\mathbb Q^\ell_Z\Vert K^\ell_\downarrow)\bigr],
\label{eq:full-graph-vfe-factorized}
\end{align}
where each divergence compares the recognition and generative factors of
\eqref{eq:full-graph-recognition-factorization} and
\eqref{eq:full-graph-generative-factorization} at their common conditioning
variables, the recognition factor also reading $o$, and is averaged over those
variables under $\mathbb Q_\lambda$. These are the observation, top-prior,
partition, graph, and parent--child contributions. \status{ESTABLISHED}

\noindent\emph{Proof.}
Substitute the two factorizations into the Radon--Nikodym ratio of
\eqref{eq:full-graph-vfe-density} and apply the chain rule for conditional relative
entropy, \Cref{lem:prob-kl-chain-rule} and \citet[Sec.~2.5]{CoverThomas2006}, in its
measure-level form, which subtracts no infinities. $\square$

\section{The scale-local two-channel graph energy}
\label{sec:full-graph-vfe-scale-local}

\definitionheading{Transported divergences and row free energies}{def:full-graph-transported-divergences}
For every copy $e:j\to i$ in $E_\ell$ and $x\in\{b,m\}$ put
\begin{equation}
D^{x,\ell}_e:=\KL\bigl(q^x_{i,\ell}\Vert(\Theta^{x,\ell}_e)_\#q^x_{j,\ell}\bigr),
\label{eq:full-graph-transported-divergences}
\end{equation}
where $(\Theta)_\#$ is the pushforward action $\widehat\rho_x(\Theta)$ on channel-$x$
laws; parallel copies carry separate divergences, and these are law-point energies in
the sense of \Cref{sec:prob-levels}. A selectively communicating variant uses the
role divergences \eqref{eq:obs-selective-role-divergence} or a message channel
instead, which changes the energy, and its recursive state keeps the interface and
noise data and any private conditional law or memory that later kernels read
(\Cref{sec:selective-records}, \Cref{prop:geo-selective-descent}). With $E_\ell(i)$
the copies received by $i$, the occupancies and rows of an event law are
\begin{equation}
\alpha^{x,\ell}_i:=\sum_{e\in E_\ell(i)}\eta^{x,\ell}_e,
\qquad
\beta^{x,\ell}_e:=\frac{\eta^{x,\ell}_e}{\alpha^{x,\ell}_i}
\quad(e\in E_\ell(i),\ \alpha^{x,\ell}_i>0),
\label{eq:full-graph-edge-event-disintegration}
\end{equation}
with pair rows $\beta^{x,\ell}_{ij}=\sum_{e:j\to i}\beta^{x,\ell}_e$. The rows are the
conditional source laws of whichever event law is displayed, generative, posterior,
or recognition; a row of zero occupancy is not determined by the joint law. With
normalized copy priors $\pi^{x,\ell}_i$ on $E_\ell(i)$ at every vertex $i$ that carries a row (a
vertex with an empty source set carries no row and has zero occupancy, although $E_\ell(i)$ may
contain the reversals of copies it sends) and temperatures
$\tau^x_{i,\ell}>0$, the row free energies in energy units are
\begin{equation}
\Fenergy^{\mathrm{att},x}_{i,\ell}
=\sum_{e\in E_\ell(i)}\beta^{x,\ell}_eD^{x,\ell}_e
+\tau^x_{i,\ell}\KL\bigl(\beta^{x,\ell}_i\Vert\pi^{x,\ell}_i\bigr),
\label{eq:full-graph-row-free-energies}
\end{equation}
and their dimensionless form is $\Fenergy^{\mathrm{att},x}_{i,\ell}/\tau^x_{i,\ell}$.

\propositionheading{Gibbs rows}{prop:full-graph-gibbs-rows}
With $\mathcal Z^{x,\ell}_i=\sum_{e\in E_\ell(i)}\pi^{x,\ell}_e
e^{-D^{x,\ell}_e/\tau^x_{i,\ell}}$, assume $\mathcal Z^{x,\ell}_i>0$, equivalently that at least
one edge of positive prior mass has finite energy. On the closed simplex of rows supported on
$\{e:\pi^{x,\ell}_e>0\}$, with $0\cdot(+\infty)=0$, the extended-real identity is
$\Fenergy^{\mathrm{att},x}_{i,\ell}=\tau^x_{i,\ell}\KL(\beta^{x,\ell}_i\Vert\beta^{x,\ell,*}_i)
-\tau^x_{i,\ell}\log\mathcal Z^{x,\ell}_i$ with
\begin{equation}
\beta^{x,\ell,*}_e=\frac{\pi^{x,\ell}_ee^{-D^{x,\ell}_e/\tau^x_{i,\ell}}}{\mathcal Z^{x,\ell}_i},
\label{eq:full-graph-gibbs-rows}
\end{equation}
so the minimizer is unique and the minimum is $-\tau^x_{i,\ell}\log\mathcal
Z^{x,\ell}_i$, with zero optimal mass on every infinite-energy edge. If $\mathcal Z^{x,\ell}_i=0$,
every prior-supported row has infinite free energy, and the Gibbs formula defines no probability row
or finite minimizer. When all prior-supported energies are finite, the optimum has positive mass on
every prior-supported edge. The optimal rows move with the transported divergences, and omitting
the row relative entropy changes the objective. The dimensionless form is the row
free energy \eqref{eq:obs-attention-vfe} with prior $\pi^{x,\ell}_i$ and the fixed
energies $D^{x,\ell}_e$ in place of the expected divergences, with the same minimizer
\eqref{eq:obs-attention-recognition-optimum} and minimum $-\log\mathcal
Z^{x,\ell}_i$. The energy-unit form with $\tau^x_{i,\ell}\ne1$ is not an
independently weighted sector of one evidence bound (\Cref{prop:obs-attention-elbo}),
and the composite potential below inherits this scaling. \status{ESTABLISHED}

\noindent\emph{Proof.}
Expanding $\KL(\beta\Vert\beta^*)$ with \eqref{eq:full-graph-gibbs-rows} gives
$\sum_e\beta_e\log(\beta_e/\pi_e)+\sum_e\beta_eD_e/\tau+\log\mathcal Z$, which is the
stated identity after multiplication by $\tau$. Relative entropy is nonnegative and
vanishes only at $\beta=\beta^*$. A row charging an infinite-energy edge has both sides equal to
$+\infty$; otherwise the expansion uses only finite-energy edges. If $\mathcal Z=0$, every edge of
positive prior mass has infinite energy, giving the stated exceptional case. $\square$

\definitionheading{Composite multiscale potential}{def:full-graph-composite-potential}
If a scale-local prior is a normalized Gibbs law containing the self terms and
$\sum_{e\in E_\ell}(\eta^{b,\ell}_eD^{b,\ell}_e+\eta^{m,\ell}_eD^{m,\ell}_e)$, those
terms and the global normalizer belong to the tower variational free energy. If they
are appended to a fixed population free energy $\Fenergy_{\mathrm{pop}}$ of
\Cref{ch:local-collective-elbo} instead, the resulting objective is the composite
potential
\begin{equation}
\mathcal J_{\mathrm{comp}}
=\Fenergy_{\mathrm{pop}}
+\sum_\ell\sum_i\bigl(\alpha^{b,\ell}_i\Fenergy^{\mathrm{att},b}_{i,\ell}
+\alpha^{m,\ell}_i\Fenergy^{\mathrm{att},m}_{i,\ell}\bigr)
+\sum_\ell\Fenergy^\ell_{\mathrm{part}}
+\sum_\ell\Fenergy^\ell_{\mathrm{cross}}
+\sum_\ell\mathcal J^\ell_{\mathrm{link}},
\label{eq:full-graph-composite-potential}
\end{equation}
where the sum over $i$ runs over the vertices that carry a row, a term of zero occupancy is
zero, the partition and cross-scale terms are those of \eqref{eq:full-graph-partition-free-energy}
and \eqref{eq:full-graph-cross-scale-kl}, and $\mathcal J^\ell_{\mathrm{link}}$ is a declared link
regularizer. It is a descent potential and carries no evidence
interpretation of its own.

\section{Partition variables and event pushforwards}
\label{sec:full-graph-vfe-partitions}

Attention rows do not create a parent random variable. A hierarchy meant to emerge
under the same variational principle puts the structural mark $R_\ell$ inside both
the generative and the recognition law.

\definitionheading{Partition factor}{def:full-graph-partition-factor}
A finite-temperature partition term has the form
\begin{equation}
\Fenergy^\ell_{\mathrm{part}}
=\E_{\mathbb Q^\ell_R}\bigl[H^\ell_{\mathrm{block}}(R_\ell;Z_\ell,\mathfrak N_\ell)\bigr]
+\tau^\ell_R\KL\bigl(\mathbb Q^\ell_R\Vert\mathbb P^{\ell,\circ}_R\bigr),
\label{eq:full-graph-partition-free-energy}
\end{equation}
and enters the variational free energy only when the normalized Gibbs factor that it
represents is part of the generative law. The block action
$H^\ell_{\mathrm{block}}$ may trade internal transported mismatch against boundary
retention, evaluator closure, and residual cost. No canonical partition selector
follows from the global free energy alone (\Cref{open:cg-natural-selector}).

\subsection{A random partition catalog}
\label{sec:rg-geometry-normalization}

The structural mark may be random with a finite catalog of partitions, or membership
configurations, of $V_\ell$; each label fixes the branch carrier of the next depth,
$X$ stays outside the channel, no edge is created, and the two normalizations of the
catalog are distinct generative models.

\definitionheading{Finite partition catalog}{def:rg-geometry-ensemble}
Let $\mathfrak G$ be a finite nonempty catalog of values of $R_\ell$. For each
$\gamma\in\mathfrak G$ declare a finite nonempty carrier $\mathsf Y_\gamma$, a
strictly positive probability reference $\nu_\gamma$, a positive raw factor
$\widetilde p_\gamma$, a positive weight $\mathsf a_\gamma$, and a normalized,
possibly vanishing, likelihood $L_\gamma(o\given y)$ on a common finite observation
alphabet, all fixed generative data independent of recognition. On the disjoint union
of the carriers put
\begin{equation}
\mathcal Z_\gamma=\sum_y\nu_\gamma(y)\widetilde p_\gamma(y),
\qquad
p_\gamma(y)=\frac{\nu_\gamma(y)\widetilde p_\gamma(y)}{\mathcal Z_\gamma},
\qquad
\bar{\mathsf a}_\gamma=\frac{\mathsf a_\gamma}{\sum_\delta\mathsf a_\delta}.
\label{eq:rg-geometry-branch}
\end{equation}
The hierarchically normalized and the globally normalized joints are
\begin{align}
\mathbb P_{\mathrm h}(\gamma,y,o)
&=\bar{\mathsf a}_\gamma p_\gamma(y)L_\gamma(o\given y),
\label{eq:rg-geometry-hierarchy}\\
\mathbb P_{\mathrm f}(\gamma,y,o)
&=\frac{\mathsf a_\gamma\nu_\gamma(y)\widetilde p_\gamma(y)L_\gamma(o\given y)}{\mathcal Z_{\mathrm{all}}},
\qquad
\mathcal Z_{\mathrm{all}}=\sum_\delta\mathsf a_\delta\mathcal Z_\delta .
\label{eq:rg-geometry-factor}
\end{align}

\propositionheading{Branch normalization changes label weights}{prop:rg-geometry-weights}
Both joints are normalized and have the same conditional branch laws, while the label
weight of the globally normalized joint is $\mathsf a_\gamma\mathcal
Z_\gamma/\mathcal Z_{\mathrm{all}}$. The label weights of the two joints are the
softmax weights of \Cref{def:rg-thermo-finite-score} with effective log masses
$\log\mathsf a_\gamma$ and $\log\mathsf a_\gamma+\log\mathcal Z_\gamma$
respectively. For the same positive weights the joints coincide exactly when all
$\mathcal Z_\gamma$ coincide. Replacing $\widetilde p_\gamma$ by $\mathrm
c_\gamma\widetilde p_\gamma$ with $\mathrm c_\gamma>0$ leaves the hierarchical joint
unchanged and reweights the global one by $\mathrm c_\gamma$, and a common factor
cancels in both. For a strictly positive probability reference $\nu_{\mathfrak G}$ on
the catalog write $\mathsf a_\gamma=\nu_{\mathfrak G}(\gamma)e^{-H^{\mathrm{lab}}_\gamma}$ and
$\widetilde p_\gamma=e^{-H_\gamma}$; every positive family of weights has this form.
Fix positive evidence $M_o=\sum_{\gamma,y}\mathbb P(\gamma,y,o)$ and a recognition law
$\mathbb Q(\gamma,y)=\mathbb Q^{\mathrm{lab}}(\gamma)\mathbb Q_\gamma(y)$, with
$\mathbb Q_\gamma$ defined where $\mathbb Q^{\mathrm{lab}}(\gamma)>0$. Then, with
$\Fenergy_o=\KL(\mathbb Q\Vert\boldsymbol\Pi_o)-\log M_o$,
\begin{align}
\Fenergy_{\mathrm f,o}(\mathbb Q)
&=\KL(\mathbb Q^{\mathrm{lab}}\Vert\nu_{\mathfrak G})+\E_{\mathbb Q^{\mathrm{lab}}}H^{\mathrm{lab}}_\gamma
+\log\mathcal Z_{\mathrm{all}}
\nonumber\\
&\quad+\sum_\gamma\mathbb Q^{\mathrm{lab}}(\gamma)
\bigl[\KL(\mathbb Q_\gamma\Vert\nu_\gamma)+\E_{\mathbb Q_\gamma}H_\gamma
-\E_{\mathbb Q_\gamma}\log L_\gamma(o\given Y)\bigr],
\label{eq:rg-geometry-vfe-f}\\
\Fenergy_{\mathrm h,o}(\mathbb Q)
&=\Fenergy_{\mathrm f,o}(\mathbb Q)-\log\mathcal Z_{\mathrm{all}}
+\log\sum_\delta\mathsf a_\delta
+\sum_\gamma\mathbb Q^{\mathrm{lab}}(\gamma)\log\mathcal Z_\gamma,
\label{eq:rg-geometry-vfe-h}
\end{align}
with the convention $0\log0=0$; a recognition-positive atom with zero likelihood
gives $+\infty$. Each expansion contains one full recognition entropy, decomposed by
its conditional chain, and no further correlation-entropy term.
\status{ESTABLISHED}

\noindent\emph{Proof.}
Sum either joint over $o$, then $y$, then $\gamma$. Equality of the joints implies
equality of the label marginals, hence $\mathcal Z_\gamma=\mathcal
Z_{\mathrm{all}}/\sum_\delta\mathsf a_\delta$ for every $\gamma$; the converse is
substitution. A branch rescaling changes $\mathcal Z_\gamma$ to $\mathrm
c_\gamma\mathcal Z_\gamma$ and cancels in $p_\gamma$. The softmax statement is the
label marginal of each joint written as the exponential of its log mass divided by the
sum of these exponentials. The expansions follow by writing $\log\mathbb P$ with the
displayed factors and taking the expectation under $\mathbb Q$; the two differ only in
the normalizers, $\log\sum_\delta\mathsf a_\delta+\log\mathcal Z_\gamma$ against
$\log\mathcal Z_{\mathrm{all}}$. $\square$

Reference, raw factor, label weight, evidence, and posterior are distinct objects,
and the score of either joint in a parameter is the complete labeled score of
\Cref{prop:rg-thermo-finite-labeled-fisher}, which contains every varying generative
factor, the $\mathcal Z_\gamma$ of the global model included.

\subsection{Membership kernels and event pushforwards}

\definitionheading{Membership kernels}{def:full-graph-membership-kernel}
Overlapping assignment is a membership Markov kernel
\begin{equation}
\mathfrak m_\ell(I\given i,R_\ell)\ge0,
\qquad
\sum_{I\in V_{\ell+1}}\mathfrak m_\ell(I\given i,R_\ell)=1,
\label{eq:full-graph-membership-kernel}
\end{equation}
of which a hard partition is the deterministic case $\mathfrak m_\ell(I\given
i,R_\ell)=\mathbf 1\{R_\ell(i)=I\}$. Correlated endpoints use one joint endpoint
kernel
\begin{equation}
\mathfrak M_\ell(I,J\given i,j,R_\ell)\ge0,
\qquad
\sum_{I,J}\mathfrak M_\ell(I,J\given i,j,R_\ell)=1,
\label{eq:full-graph-joint-endpoint-kernel}
\end{equation}
whose product form requires a declared conditional independence of the two draws.

\propositionheading{Event pushforwards}{prop:full-graph-event-pushforward}
Fix $y\in\mathsf Y_\ell$ and a realized mark $\gamma\in\mathsf R_\ell$. If every
vertex receives one membership draw and the draws of distinct vertices are
conditionally independent at $(y,\gamma)$, the conditional coarse event law is
\begin{equation}
\eta^{x,\ell+1}_{IJ}(y,\gamma)
=\sum_{\substack{i,j\in V_\ell\\ i\ne j}}\eta^{x,\ell}_{ij}(y)
\mathfrak m_\ell(I\given i,\gamma)\mathfrak m_\ell(J\given j,\gamma)
+\delta_{IJ}\sum_{i\in V_\ell}\eta^{x,\ell}_{ii}(y)\mathfrak m_\ell(I\given i,\gamma),
\label{eq:full-graph-edge-pushforward}
\end{equation}
where the loop term uses the diagonal endpoint kernel, since both endpoints of a loop
follow one draw. With correlated endpoints,
\begin{equation}
\eta^{x,\ell+1}_{IJ}(y,\gamma)
=\sum_{i,j\in V_\ell}\eta^{x,\ell}_{ij}(y)\mathfrak M_\ell(I,J\given i,j,\gamma).
\label{eq:full-graph-correlated-edge-pushforward}
\end{equation}
With the assignment marginal $C^R_\ell(y,d\gamma)=C^{\mathrm{struct}}_{\ell+1\leftarrow\ell}(y;d\gamma\times\mathsf
Y_{\ell+1})$ of the structural channel and an input law $\mathbb M$, the unconditional
event mass averages the whole conditional expression,
\begin{equation}
\overline\eta^{x,\ell+1}_{IJ}[\mathbb M]
=\sum_{i,j\in V_\ell}\int\eta^{x,\ell}_{ij}(y)\mathfrak M_\ell(I,J\given i,j,\gamma)
C^R_\ell(y,d\gamma)\mathbb M(dy).
\label{eq:full-graph-unconditional-edge-pushforward}
\end{equation}
All three coarse event laws sum to one. Conditional endpoint independence permits the
product in \eqref{eq:full-graph-edge-pushforward} but not the factoring of the
expectation of a random event mass and a random assignment, so
\eqref{eq:full-graph-unconditional-edge-pushforward} applies to all three input laws
with no independence assumed between event and membership marks. The random
$\eta_{\ell+1}(y,\gamma)$ is part of $\mathfrak N_{\ell+1}$ with rows recovered by
\eqref{eq:full-graph-edge-event-disintegration}, and its unconditional mean does not
replace it. Loop copies arise at every depth after the first link lift
\eqref{eq:full-graph-coarse-links}; they are induced by blocking and are not a new
support declaration, as in \Cref{prop:rg-meta-attention-pushforward}. For a hard
partition the loop and product
terms agree, and under the posterior bridge the pushforward is
\eqref{eq:rg-meta-attention}. Literal unit-weight membership in several parents has
total mass greater than one; it is finite-measure replication, not a probability
pushforward. \status{ESTABLISHED}

\noindent\emph{Proof.}
Summing the first term of \eqref{eq:full-graph-edge-pushforward} over $I,J$ gives
$\sum_{i\ne j}\eta_{ij}$ and the loop term gives $\sum_i\eta_{ii}$, which add to one;
the correlated and unconditional cases use normalization of $\mathfrak M_\ell$,
$\eta_\ell$, $C^R_\ell$, and $\mathbb M$. The expectation of a product of dependent
random quantities is not the product of expectations. A vertex with unit weight in $k$
parents contributes its mass $k$ times. $\square$

\definitionheading{Chapman--Kolmogorov composition}{def:full-graph-composition}
Nested assignments compose by
\begin{equation}
\mathfrak m_{\ell+2\leftarrow\ell}(A\given i,R_\ell,R_{\ell+1})
=\sum_I\mathfrak m_{\ell+1}(A\given I,R_{\ell+1})\mathfrak m_\ell(I\given i,R_\ell),
\label{eq:full-graph-membership-composition}
\end{equation}
\begin{equation}
\mathfrak M_{\ell+2\leftarrow\ell}(A,A'\given i,j,R_\ell,R_{\ell+1})
=\sum_{I,J}\mathfrak M_{\ell+1}(A,A'\given I,J,R_{\ell+1})\mathfrak M_\ell(I,J\given i,j,R_\ell),
\label{eq:full-graph-joint-membership-composition}
\end{equation}
and the augmented structural channels by
\begin{equation}
C^{\mathrm{struct}}_{\ell+2\leftarrow\ell}
=C^{\mathrm{struct}}_{\ell+2\leftarrow\ell+1}\circ C^{\mathrm{struct}}_{\ell+1\leftarrow\ell},
\label{eq:full-graph-structural-channel-composition}
\end{equation}
where composition integrates every intermediate assignment, event law, and mark, and
the root, holonomy, and boundary marks stay in the target of each composite. An
independently declared multi-depth channel is admitted only when it coincides with the
composite, and the composite is the exact multi-depth step only when the augmented
state carries every mark that the next structural kernel reads.

\section{Coarse links and what the parent retains}
\label{sec:full-graph-vfe-nonflat}

The links are not assumed flat: $\mathfrak h_\ell$ keeps the root-framed datum of
\Cref{prop:rg-rooted-datum} in every connected block with its joint law, and on a
non-flat graph a reduced mean of the dressed marks must carry a declared residual,
whose penalization to zero is a modeling preference, not a consequence of gauge
covariance.

\definitionheading{Coarse links under a hard partition}{def:full-graph-coarse-links}
Let the realized mark $R_\ell=\gamma$ be a hard partition whose blocks are replaced by
their component-indexed meta-agents \eqref{eq:rg-component-meta-index}, so that each
coarse vertex $I$ has one root $r_I$ and one rooted spanning tree. This excludes
overlapping membership kernels, for which a copy would need a root transport into
every block that contains one of its endpoints with positive membership; the event
pushforwards of \Cref{prop:full-graph-event-pushforward} hold for general membership
kernels. The depth-$(\ell+1)$ multigraph $\mathfrak I_{\ell+1}=(V_{\ell+1},E_{\ell+1},\bar{\ })$
has the coarse vertices, the same copies $E_{\ell+1}:=E_\ell$ with the same reversal,
and endpoints $e:J\to I$ whenever $e:j\to i$ with $j\in J$ and $i\in I$; its event law
is $\eta^{x,\ell+1}_e:=\eta^{x,\ell}_e$, and its links are the dressed marks of all
copies,
\begin{equation}
\Theta^{x,\ell+1}_e
:=\widehat\Theta^{x,\ell}_e
=\Theta^{x,\ell}_{I\leftarrow i}\Theta^{x,\ell}_e\bigl(\Theta^{x,\ell}_{J\leftarrow j}\bigr)^{-1},
\qquad e:j\to i,\ j\in J,\ i\in I,
\label{eq:full-graph-coarse-links}
\end{equation}
internal as well as boundary, so that internal copies become loops at their block. The
parent laws $q^x_{I,\ell+1}$ are expressed in the frame of the root $r_I$. The coarse
multigraph is induced structure; no edge is created.

\propositionheading{Properties of the coarse links}{prop:full-graph-coarse-links}
The coarse links are well typed and keep the whole link content of depth $\ell$. An
internal tree copy has the identity link $1_G$. An internal non-tree copy carries the
based holonomy of its fundamental loop, and for every closed walk $p$ in $I$ based at
$r_I$ that traverses the copies $e_0,\ldots,e_{n-1}$ in this order,
$\mathrm{Hol}^{x,\ell}_I(p)=\Theta^{x,\ell+1}_{e_{n-1}}\cdots\Theta^{x,\ell+1}_{e_0}$.
Inverse reversal persists, and the dressed law is the vertex law
\eqref{eq:geo-regime-two-gauge-law} with vertex coordinates $g_{r_I}$, so the
divergences, rows, and row free energies of
\Cref{def:full-graph-transported-divergences} are defined and gauge invariant at
depth $\ell+1$. On a loop the divergence
$\KL(q^x_{I,\ell+1}\Vert(\Theta^{x,\ell+1}_e)_\#q^x_{I,\ell+1})$ vanishes on tree copies
and, on a non-tree copy, exactly when the parent law is fixed by that holonomy, as in
\eqref{eq:full-graph-holonomy-stabilization}. The pair marginal
$\sum_{e:J\to I}\eta^{x,\ell+1}_e$ is \eqref{eq:full-graph-edge-pushforward} for the
hard partition, and no averaged group element replaces the parallel copies. The links
depend on the trees, not only on the roots: changing the root of a fixed tree is a
vertex gauge transformation at depth $\ell+1$, while replacing the tree changes
$\Theta^{x,\ell}_{I\leftarrow i}$ to $\mathrm{Hol}^{x,\ell}_I(p_i)\Theta^{x,\ell}_{I\leftarrow i}$
for a based loop $p_i$ that depends on $i$, and a copy that leaves the tree changes its
loop link from $1_G$ to the holonomy $\mathrm{Hol}^{x,\ell}_I(p)$ of its new fundamental
loop $p$ and its loop divergence from $0$ to
$\KL(q^x_{I,\ell+1}\Vert\mathrm{Hol}^{x,\ell}_I(p)_\#q^x_{I,\ell+1})$. Whenever such a
holonomy is not the identity, the trees are part of the coarse channel rather than a
gauge choice. When $q^x_{I,\ell+1}$ and $q^x_{J,\ell+1}$ lie in the stabilized sector,
the divergences $D^{x,\ell+1}_e$ do not depend on the trees. This link-level tree
dependence complements the feature-level one of
\eqref{eq:rg-linear-cross-scale-tree-change}. \status{ESTABLISHED}

\noindent\emph{Proof.}
The dressed link of an internal copy is $1_G$ on a tree copy, since there
$\Theta_e=\Theta_{I\leftarrow i}^{-1}\Theta_{I\leftarrow j}$, and is the
fundamental-cycle holonomy \eqref{eq:full-graph-tree-transport-law} otherwise.
Inserting $(\Theta_{I\leftarrow k})^{-1}\Theta_{I\leftarrow k}$ at each intermediate
vertex $k$ of a closed walk turns the ordered product of the fine links into that of
the coarse links. Reversal and the vertex law follow from
\eqref{eq:rg-tree-transport-law}, and relative entropy is invariant under a common
bijective pushforward, which gives gauge invariance; a loop divergence vanishes
exactly when $q_I=\mathrm{Hol}_I(p)_\#q_I$. A root change multiplies every transport
of the block by one element, and a tree change multiplies each tree transport on the
left by a based holonomy (\Cref{prop:rg-rooted-datum}), so the link of $e:j\to i$ becomes
$\mathrm{Hol}_I(p_i)\Theta^{x,\ell+1}_e\mathrm{Hol}_J(p_j)^{-1}$, and when
$\mathrm{Hol}_I(p_i)$ fixes $q_I$ and $\mathrm{Hol}_J(p_j)$ fixes $q_J$, invariance under
the pushforward by $\mathrm{Hol}_I(p_i)^{-1}$ gives
$\KL(q_I\Vert(\mathrm{Hol}_I(p_i)\Theta\mathrm{Hol}_J(p_j)^{-1})_\#q_J)=\KL(q_I\Vert\Theta_\#q_J)$. $\square$

\section{How a parent influences its children}
\label{sec:full-graph-vfe-downward}

\propositionheading{Downward coupling in a declared hierarchy}{prop:full-graph-downward-coupling}
Under \Cref{hyp:full-graph-depth-markov} for a declared top-down hierarchy, a
factorization such as
\begin{equation}
K^\ell_\downarrow(dZ_\ell\given Y_{\ell+1},R_\ell,\mathfrak N_\ell,X_*)
=\prod_{I\in V_{\ell+1}}K^\ell_{\downarrow,I}\bigl(dZ_{I,\ell}\given Z_{I,\ell+1},\mathfrak N_{\ell+1},R_\ell,\mathfrak N_\ell,X_*\bigr)
\label{eq:full-graph-downward-kernel}
\end{equation}
makes the parent belief and model presentations arguments of the children's
conditional generative prior, and the corresponding term of
\eqref{eq:full-graph-vfe-factorized} is
\begin{equation}
\Fenergy^\ell_{\mathrm{cross}}
=\E_{\mathbb Q_\lambda}
\KL\bigl(\mathbb Q^\ell_Z(\cdot\given Y_{\ell+1},R_\ell,\mathfrak N_\ell,o,X_*)\Vert
K^\ell_\downarrow(\cdot\given Y_{\ell+1},R_\ell,\mathfrak N_\ell,X_*)\bigr).
\label{eq:full-graph-cross-scale-kl}
\end{equation}
The generative denominator in this cross-scale divergence reads the parent coordinates through
$K^\ell_\downarrow$. The full tower objective can also couple parent coordinates to child structural
marks through $\mathbb P^\ell_R$ and $\mathbb P^\ell_N$ in \Cref{hyp:full-graph-depth-markov}.
A downward kernel that does not depend on the parent coordinates
contributes no parent--child cross-derivative through this generative factor, and
normalization of the downward kernel alone implies no coupling. When
$K^\ell_\downarrow$ depends on the parent coordinates and the recognition
parameterization transmits that dependence, the mixed derivative in a parent and a
child parameter can be nonzero, so that optimization is coupled in both directions
although the generative arrow is top-down; a universal claim of such coupling for a
family requires a proved nonzero cross-derivative in that family. For a binary parent
$b$ and child $y$ with product recognition factors $q_b(1)=u$ and $q_y(1)=v$, the
mixed derivative of $\E_{q_bq_y}[-\log K_\downarrow(y\given b)]$ in $(u,v)$ is
$-\log[K_\downarrow(1\given1)K_\downarrow(0\given0)/(K_\downarrow(0\given1)K_\downarrow(1\given0))]$,
which vanishes exactly when the odds ratio of the kernel is one.
\status{ESTABLISHED}

\noindent\emph{Proof.}
The generative denominator in \eqref{eq:full-graph-cross-scale-kl} is $K^\ell_\downarrow$.
If that kernel does not read the parent coordinates, its negative log density has no parent
argument; the expectation and recognition parameterization may still depend on the parent. For the binary case the expectation is bilinear in $(u,v)$, and its $uv$
coefficient is the displayed alternating sum of $-\log K_\downarrow(y\given b)$ over
the four configurations. $\square$

This is the static sense in which a meta-agent influences its sub-agents without a
generative factor reading a posterior; in an adaptive tower the same kernel is the
reversal kernel of \Cref{lem:full-graph-adaptive-reversal}.

\section{Closing the tower}
\label{sec:full-graph-closing-tower}

A tower in which each depth's prior is supplied by the depth above must terminate, and
within one finite acyclic joint its top cannot be constituted by the agents it
organizes (\Cref{prop:obs-declared-root-unavoidable}). A same-time reciprocal
structure among latents is admissible only as a globally normalized joint
(\Cref{sec:gen-undirected}, \Cref{prop:gen-gibbs-counterexample}): a flat unanchored
energy is improper, an anchored or non-flat energy with its global normalizer is a
valid undirected model, and an energy appended without its normalizer is not an
evidence bound. A same-time deterministic parent $Z_{\ell+1}=\chi_\ell(Z_\ell)$ and an
independently sampled parent that generates $Z_\ell$ are therefore not two free
descriptions of one joint: one uses the hierarchy above, a deterministic pushforward,
or a normalized undirected compatibility factor, and reciprocal same-time influence
requires the last with its partition function. Supplement G works out the anchored
Gaussian star and the convergence of its coordinate ascent.

\subsection{A declared root}
\label{sec:obs-survives}

A truncated tower with a declared top prior $\mathbb P_L$ satisfies the hypotheses of
\Cref{ch:elbo} without amendment; its top is given, not constituted.

\propositionheading{A declared root is unavoidable}{prop:obs-declared-root-unavoidable}
Let a finite directed acyclic generative model be typed as in \Cref{ch:generative}, so
that every observed record is a child of at least one latent vertex. Then at least one
latent has no parents, and its factor in the joint depends only on the parameters and
the structural data. No finite acyclic construction of this type has every prior
constituted by other latents. \status{ESTABLISHED}

\noindent\emph{Proof.}
A finite directed acyclic graph admits a topological order, whose first element has
no incoming edge. That element is not a record, because every record has a latent
parent, so it is a latent, and its factor is conditioned on no other vertex. The typing
hypothesis is needed: in the graph $o\to y$ with observed root $o$, the only latent has
a parent. $\square$

Anchored, globally normalized cyclic models remain available, but their anchors and
partition functions are declared; parameter learning fits only their parameters.

\subsection{Three loci of participation}
\label{sec:obs-participatory}

What a cycle in the model was reaching for is partly available without one, at three
loci, each a sequence of evaluations of fixed normalized joints.

\remarkheading{The inference loop is block coordinate ascent}{rem:full-graph-inference-loop}
In a mean-field recognition family, the coordinate update of a factor reads the current
factors of every variable that shares a factor of the joint with it
\citep[Sec.~10.1.1]{Bishop2006}. In every connected hierarchical model in which the
children's conditionals depend on their parents, parent and child updates therefore
read each other; at a stationary point of the schedule the factor that acts as a child's
prior in its own update is $\exp\E[\log K_\downarrow(z_{\mathrm{child}}\given
\cdot)]$ under the parent's current factor, which is itself formed from the factors of
all its children. The effective prior is thus formed from the population's current
beliefs through a latent inside the joint. This is a schedule of updates against one
fixed, acyclic, normalized joint, not participation in time, and no kernel reads a
posterior, so \Cref{req:gen-typing-prohibition} holds. \status{STANDARD}

\propositionheading{The aggregate recognition law is the optimal top prior}{prop:full-graph-aggregate-prior}
Let episodes $n=1,\ldots,N_{\mathrm e}$ be conditionally independent given $\theta$,
each with its own top latent in a common space carrying the prior density $p_\theta$
relative to a fixed reference measure, let $\mathbb Q_n$ be the recognition law of the
top latent of episode $n$, and let every other factor be free of $\theta$. The
$\theta$-dependent part of the summed bound is
\begin{equation}
\sum_{n=1}^{N_{\mathrm e}}\E_{\mathbb Q_n}[\log p_\theta]
=N_{\mathrm e}\E_{\bar{\mathbb Q}}[\log p_\theta],
\qquad
\bar{\mathbb Q}=\frac1{N_{\mathrm e}}\sum_{n=1}^{N_{\mathrm e}}\mathbb Q_n .
\label{eq:obs-aggregate-prior}
\end{equation}
If $\bar{\mathbb Q}$ has a density $\bar q$ with finite entropy relative to the same
reference, then $\E_{\bar{\mathbb Q}}[\log p]=\E_{\bar{\mathbb Q}}[\log\bar
q]-\KL(\bar{\mathbb Q}\Vert p)$ for every prior density $p$. Over an unrestricted
family of prior densities the bound is maximized exactly at the aggregate recognition
law $p=\bar q$, and within a restricted family its maximizers are the minimizers of
$\KL(\bar{\mathbb Q}\Vert p)$ over the family, when they exist. Each evaluation uses one
fixed normalized joint, and updating $\theta$ at fixed recognition is a coordinate
step of the two-argument objective \citep{Neal1998}, so
\Cref{req:gen-typing-prohibition} holds. \status{ESTABLISHED}

\noindent\emph{Proof.}
The first identity is linearity of expectation in the recognition law. The second is
$\log p=\log\bar q-\log(\bar q/p)$ integrated against $\bar{\mathbb Q}$, and Gibbs'
inequality gives the maximizers. $\square$

This is the variational counterpart of empirical Bayes \citep{Robbins1956}, whose
optimum is the aggregate-posterior prior of \citet{HoffmanJohnson2016} and
\citet{TomczakWelling2018}; only the parameter of the declared prior family is fitted,
by a slow-block update in the typing of \Cref{sec:obs-fast-slow}.

\definitionheading{Episode-indexed joint}{def:full-graph-episode-joint}
Episodes are indexed by $n=0,1,\ldots$ with records $o_n$. In the parameter route each
episode has its own normalized joint $\mathbb P_{\theta_n,V}$, and the parameter is
updated by
\begin{equation}
\theta_{n+1}=\mathcal U_n\bigl(\theta_n,\mathbb Q^{(n)}_{V,o_n,X}\bigr),
\label{eq:obs-episode-update}
\end{equation}
where $\mathbb Q^{(n)}_{V,o_n,X}$ is the recognition coupling selected at episode $n$
and $\mathcal U_n$ is a declared map; within each episode the joint is fixed, and its
generative kernels read $\theta_n$ and never a recognition law of the same episode. In
the latent route one joint spans episodes $0,\ldots,n_1$: episode $n$ has a tower state
$\mathbf Y^{(n)}$ with parent coordinates $Z^{(n)}_{\mathrm{par}}=\chi^{\mathrm{ep}}(\mathbf Y^{(n)})$
for a measurable read-out $\chi^{\mathrm{ep}}$, and
\begin{equation}
\begin{aligned}
\mathbb P^{\mathrm{ep}}\bigl(d\mathbf Y^{(0:n_1)},do_{0:n_1}\given X\bigr)
={}&\mathbb P^{(0)}(d\mathbf Y^{(0)}\given X)
\prod_{n=1}^{n_1}\mathbb P^{(n\leftarrow n-1)}\bigl(d\mathbf Y^{(n)}\given Z^{(n-1)}_{\mathrm{par}},X\bigr)\\
&\times\prod_{n=0}^{n_1}\mathbb P^{\mathrm{obs}}\bigl(do_n\given\mathbf Y^{(n)},X\bigr),
\end{aligned}
\label{eq:full-graph-episode-joint}
\end{equation}
with a declared root in $\mathbb P^{(0)}$; the parent latent of episode $n-1$
parameterizes the conditional law of the tower of episode $n$. The observation
kernel is any Markov kernel and need not have a density: the tower state may carry
a latent copy $\widetilde o_n$ of the record, observed through
$\mathbb P^{\mathrm{obs}}(do_n\given\mathbf Y^{(n)},X)=\delta_{\widetilde o_n}(do_n)$,
in which case the free energy of the joint is its extended form
\eqref{eq:full-graph-vfe-global} and the decomposition
\eqref{eq:full-graph-episode-decomposition} applies as displayed only for discrete
records.

\propositionheading{Delayed feedback in the episode-indexed joint}{prop:full-graph-episode-joint}
In the latent route, \eqref{eq:full-graph-episode-joint} is a normalized joint whose
kernels read structural data and earlier latents only, so
\Cref{req:gen-typing-prohibition} holds; its graph is acyclic with every
cross-episode edge pointing forward in $n$, so
\Cref{prop:obs-declared-root-unavoidable} still requires a declared root in episode
$0$. The prior of the tower at episode $n+1$ given the records so far is the filtered
parent law pushed through the transition,
\begin{equation}
\mathbb P^{\mathrm{ep}}\bigl(\mathbf Y^{(n+1)}\in\cdot\given o_{0:n}\bigr)
=\int\mathbb P^{(n+1\leftarrow n)}(\cdot\given z,X)
\mathbb P^{\mathrm{ep}}\bigl(Z^{(n)}_{\mathrm{par}}\in dz\given o_{0:n}\bigr),
\label{eq:full-graph-episode-predictive}
\end{equation}
so the children's prior at episode $n+1$ is set by the parent inferred from the
population at episode $n$. For every recognition law $\mathbb Q$ on $\mathbf
Y^{(0:n_1)}$, on the finite-density domain with observation densities
$p^{\mathrm{obs}}$,
\begin{equation}
\begin{aligned}
\Fenergy^{\mathrm{ep}}[\mathbb Q]
={}&\sum_{n=0}^{n_1}\E_{\mathbb Q}\bigl[-\log p^{\mathrm{obs}}(o_n\given\mathbf Y^{(n)},X)\bigr]
+\KL\bigl(\mathbb Q^{(0)}\Vert\mathbb P^{(0)}\bigr)\\
&+\sum_{n=1}^{n_1}\E_{\mathbb Q}\KL\bigl(\mathbb Q(\cdot\given\mathbf Y^{(0:n-1)})\Vert
\mathbb P^{(n\leftarrow n-1)}(\cdot\given Z^{(n-1)}_{\mathrm{par}},X)\bigr),
\end{aligned}
\label{eq:full-graph-episode-decomposition}
\end{equation}
where $\mathbb Q^{(0)}$ is the law of $\mathbf Y^{(0)}$ under $\mathbb Q$ and
$\mathbb Q(\cdot\given\mathbf Y^{(0:n-1)})$ the conditional law of $\mathbf Y^{(n)}$.
In the parameter route, if $\mathcal U_n$ maximizes the bound of episode $n$ over the
parameter at fixed $\mathbb Q^{(n)}_{V,o_n,X}$, the bound of episode $n$ at
$\theta_{n+1}$ is at least its value at $\theta_n$. \status{ESTABLISHED}

\noindent\emph{Proof.}
Normalization is composition of Markov kernels, as in
\Cref{prop:gen-exact-normalization}, and acyclicity follows because every kernel reads
only earlier episodes. Given $Z^{(n)}_{\mathrm{par}}$, the tower of episode $n+1$ is
independent of $(\mathbf Y^{(0:n)},o_{0:n})$, because its kernel reads only
$Z^{(n)}_{\mathrm{par}}$ and the records up to $n$ are generated before it; conditioning
on $o_{0:n}$ and integrating out $Z^{(n)}_{\mathrm{par}}$ gives
\eqref{eq:full-graph-episode-predictive}. The decomposition is the chain rule
\Cref{lem:prob-kl-chain-rule} applied along $n$ to the Radon--Nikodym ratio of
$\mathbb Q$ against the joint. The last statement is the definition of a maximizer.
$\square$

\propositionheading{One joint for a declared parameter route}{prop:full-graph-episode-lifted-joint}
Let the episode joints of the parameter route form a Markov kernel
$\theta\mapsto\mathbb P_{\theta,V}(dy,do\given X)$ on the parameter space, with
evidence densities $p_{\theta,V}(o\given X)$ relative to one reference $\nu$ on
$\mathsf O$ and posterior slices $\boldsymbol\Pi_{\theta,V,o,X}$ jointly measurable in
$(\theta,o)$; let the coupling $\mathbb Q^{\theta}_{V,o,X}$ selected by
\Cref{hyp:prob-measurable-coupling-selector} for the joint at $\theta$, so that
$\mathbb Q^{(n)}_{V,o_n,X}=\mathbb Q^{\theta_n}_{V,o_n,X}$, be jointly measurable in
$(\theta,o)$; and let each update $\mathcal U_n$ of
\eqref{eq:obs-episode-update} be measurable, so that
$\mathcal U^{\mathrm{sel}}_n(\theta,o):=\mathcal U_n(\theta,\mathbb Q^{\theta}_{V,o,X})$
is measurable and the route reads $\theta_{n+1}=\mathcal U^{\mathrm{sel}}_n(\theta_n,o_n)$
from a declared $\theta_0$. Write $\theta_n(o_{0:n-1})$ for the parameter the route
reaches after the records $o_{0:n-1}$, and let each realized record be admitted at
that parameter. The lifted joint on the tower states
$\mathbf Y^{(n)}=(\theta_n,y_n,\widetilde o_n)$, $n=0,\ldots,n_1$, and the records,
\begin{equation}
\mathbb P^{\mathrm{lift}}\bigl(d\mathbf Y^{(0:n_1)},do_{0:n_1}\given X\bigr)
=\delta_{\theta_0}(d\theta_0)
\prod_{n=0}^{n_1}\mathbb P_{\theta_n,V}(dy_n,d\widetilde o_n\given X)\delta_{\widetilde o_n}(do_n)
\prod_{n=1}^{n_1}\delta_{\mathcal U^{\mathrm{sel}}_{n-1}(\theta_{n-1},\widetilde o_{n-1})}(d\theta_n),
\label{eq:full-graph-lifted-joint}
\end{equation}
has four properties. (i) It is a latent route of \Cref{def:full-graph-episode-joint}
with parent coordinates $Z^{(n)}_{\mathrm{par}}=(\theta_n,\widetilde o_n)$, transition
$\mathbb P^{(n\leftarrow n-1)}(d\mathbf Y^{(n)}\given z,X)
=\delta_{\mathcal U^{\mathrm{sel}}_{n-1}(z)}(d\theta_n)\mathbb P_{\theta_n,V}(dy_n,d\widetilde o_n\given X)$,
observation kernel $\delta_{\widetilde o_n}(do_n)$, and root
$\delta_{\theta_0}\otimes\mathbb P_{\theta_0,V}$, so \Cref{prop:full-graph-episode-joint}
applies to it. (ii) Given $o_{0:n-1}$, the conditional law of $(y_n,\widetilde o_n)$ is
the episode joint $\mathbb P_{\theta_n(o_{0:n-1}),V}$ of the parameter route, and the
evidence density of the lifted joint relative to $\nu^{\otimes(n_1+1)}$ is the product
of the episode evidences,
\begin{equation}
p^{\mathrm{lift}}(o_{0:n_1}\given X)=\prod_{n=0}^{n_1}p_{\theta_n(o_{0:n-1}),V}(o_n\given X).
\label{eq:full-graph-lifted-evidence}
\end{equation}
(iii) Its posterior slice fixes $\widetilde o_n=o_n$ and $\theta_n=\theta_n(o_{0:n-1})$
and is the product $\bigotimes_n\boldsymbol\Pi_{\theta_n(o_{0:n-1}),V,o_n,X}$ on
$y_{0:n_1}$, so the posterior law of $y_n$ reads $o_{0:n}$ only. (iv) For every
recognition law $\mathbb Q$ on the tower states, the free energy
\eqref{eq:full-graph-vfe-global} of the lifted joint is
\begin{equation}
\Fenergy^{\mathrm{lift}}[\mathbb Q]
=\sum_{n=0}^{n_1}\Bigl[-\log p_{\theta_n,V}(o_n\given X)
+\KL\bigl(\mathbb Q^{y}_n\Vert\boldsymbol\Pi_{\theta_n,V,o_n,X}\bigr)\Bigr]
+\operatorname{TC}(\mathbb Q^{y}),
\qquad\theta_n=\theta_n(o_{0:n-1}),
\label{eq:full-graph-lifted-free-energy}
\end{equation}
when $\mathbb Q$ fixes $\widetilde o_n=o_n$ and $\theta_n=\theta_n(o_{0:n-1})$ almost
surely, and $+\infty$ otherwise, where $\mathbb Q^{y}$ is the law of $y_{0:n_1}$ under
$\mathbb Q$, $\mathbb Q^{y}_n$ its marginal on $y_n$, each bracket is the extended free
energy \eqref{eq:elbo-extended-free-energy} of episode $n$ at the parameter the route
reaches, and $\operatorname{TC}$ is the total correlation
\eqref{eq:obs-total-correlation} across episodes. At product laws, and in particular at
the selected couplings $\bigotimes_n\mathbb Q^{\theta_n}_{V,o_n,X}$, the free energy
is the sum of the episode bounds of the parameter route. \status{ESTABLISHED}

\noindent\emph{Proof.}
Each factor of \eqref{eq:full-graph-lifted-joint} is a Markov kernel in the
coordinates to its left: $\theta\mapsto\mathbb P_{\theta,V}$ by hypothesis, and the
Dirac factors as kernels of measurable maps, $\mathcal U^{\mathrm{sel}}_{n-1}$ being
the composition of $\mathcal U_{n-1}$ with the jointly measurable selection. Finite
composition normalizes the product (\Cref{prop:gen-exact-normalization}), and reading
the factors in the order of \eqref{eq:full-graph-episode-joint} gives (i); every kernel
reads structural data and latent coordinates only, because the selection is evaluated
at the latent copy $\widetilde o_{n-1}$ and not at a recognition law of the joint being
scored. For (ii), integrate \eqref{eq:full-graph-lifted-joint} against a bounded
measurable $f(o_{0:n-1},y_n,\widetilde o_n)$: the factors of episodes $n+1,\ldots,n_1$
have unit mass, the Dirac factors identify $o_k=\widetilde o_k$ for $k<n$ and, by
induction along the recursion, $\theta_n=\theta_n(o_{0:n-1})$, and Tonelli's theorem
leaves the integral of
$\int f(o_{0:n-1},y,\widetilde o)\mathbb P_{\theta_n(o_{0:n-1}),V}(dy,d\widetilde o\given X)$
against the law of $o_{0:n-1}$, which is the stated conditional law. Its record
marginal is $p_{\theta_n(o_{0:n-1}),V}(o_n\given X)\nu(do_n)$, and multiplying the
conditional densities along $n$ gives \eqref{eq:full-graph-lifted-evidence}; the
product is measurable because the evidence density is jointly measurable and
$o_{0:n-1}\mapsto\theta_n(o_{0:n-1})$ is a composition of measurable maps, and it is
finite and positive because each factor is. For (iii), substitute the disintegration
$\mathbb P_{\theta,V}(dy,d\widetilde o\given X)
=\boldsymbol\Pi_{\theta,V,\widetilde o,X}(dy)p_{\theta,V}(\widetilde o\given X)\nu(d\widetilde o)$
into each factor and integrate out the Dirac coordinates: for bounded measurable
$f$ of the tower states and $g$ of the records,
\begin{equation}
\begin{aligned}
\E_{\mathbb P^{\mathrm{lift}}}[fg]
={}&\int g(o)\Bigl[\int f\bigl((\theta_n(o_{0:n-1}),y_n,o_n)_{n\le n_1}\bigr)
\bigotimes_{n=0}^{n_1}\boldsymbol\Pi_{\theta_n(o_{0:n-1}),V,o_n,X}(dy_n)\Bigr]\\
&\times p^{\mathrm{lift}}(o\given X)\nu^{\otimes(n_1+1)}(do),
\end{aligned}
\label{eq:full-graph-lifted-posterior}
\end{equation}
and the bracket, measurable in $o$ by the joint measurability of the slices, is a
version of the conditional expectation of $f$ given the records. For (iv), the
posterior is carried by the set on which $\widetilde o_n=o_n$ and
$\theta_n=\theta_n(o_{0:n-1})$ for every $n$; a $\mathbb Q$ that charges its
complement is not absolutely continuous with respect to the posterior, and its
relative entropy is $+\infty$. Otherwise disintegrate both laws over $y_{0:n_1}$ with
\Cref{lem:prob-kl-chain-rule}(a): the conditional laws of the remaining coordinates
given $y_{0:n_1}$ are the same Dirac mass under both, so
$\KL(\mathbb Q\Vert\boldsymbol\Pi^{\mathrm{lift}})
=\KL(\mathbb Q^{y}\Vert\bigotimes_n\boldsymbol\Pi_{\theta_n,V,o_n,X})$, and the
product-reference identity \eqref{eq:elbo-total-correlation-chain} splits the right
side into $\operatorname{TC}(\mathbb Q^{y})+\sum_n\KL(\mathbb Q^{y}_n\Vert\boldsymbol\Pi_{\theta_n,V,o_n,X})$.
Adding $-\log p^{\mathrm{lift}}(o\given X)$ from \eqref{eq:full-graph-lifted-evidence}
gives \eqref{eq:full-graph-lifted-free-energy}, and the total correlation of a product
law is zero. $\square$

The route's bookkeeping with the parameter read from the law being scored is the
composite that \Cref{req:gen-typing-prohibition} excludes; the lifted joint evaluates
the selection at the latent copy of the record instead, and its evidence is the
sequential empirical-Bayes evidence of the route.

\openproblemheading{Accounting across episodes}{open:full-graph-episode-accounting}
In the parameter route the bounds of different episodes refer to different joints and
records. Within an episode a maximizing update is monotone
(\Cref{prop:full-graph-episode-joint}), and for an update that reads a coupling
selected by a declared rule the lifted joint of
\Cref{prop:full-graph-episode-lifted-joint} is one joint whose evidence is the product
of the episode evidences and whose free energy is the sum of the episode bounds plus
the cross-episode total correlation; an update that reads the recognition law being
scored is excluded by \Cref{req:gen-typing-prohibition}. The lifted joint informs the
parent of episode $n+1$ by the record of episode $n$ through its latent copy. What is
unproved is whether the label coordinates alone can carry the feedback. The problem is
posed for label-sufficient latent routes of \Cref{def:full-graph-episode-joint}: the
episode model, with a finite label coordinate $y^{\mathrm{lab}}$ whose prior is the
parameter, its private latents, and its record kernel, is that of the parameter route
with the posterior slice as its recognition coupling and the aggregate update of
constant step,
$\theta_{n+1}=(1-\epsilon_{\mathrm e})\theta_n+\epsilon_{\mathrm e}\boldsymbol\Pi^{\mathrm{lab}}_{\theta_n,V,o_n,X}$
with $\boldsymbol\Pi^{\mathrm{lab}}$ the label marginal of the posterior slice; the
tower state of episode $n$ consists of the episode model's latents and an auxiliary
preference coordinate; each record is emitted from its own episode's latents by the
record kernel; and the read-out $\chi^{\mathrm{ep}}$ reads the label and the
preference coordinate only. The records either have one value whose density
$p_V(o\given y^{\mathrm{lab}},X)$ given the label does not depend on the label, while
the expected squared total-variation displacement of the label posterior from the
prior, under the evidence at that prior, is bounded below uniformly over priors whose
smallest weight is at least half that of $\theta_0$; or their likelihood vectors
$(p_V(o\given y^{\mathrm{lab}},X))_{y^{\mathrm{lab}}}$ span a cone on which no nonzero
polynomial vanishes. For such routes, with the per-episode discrepancy the total
variation between the route's filtered label law
$\mathbb P^{\mathrm{ep}}(y^{\mathrm{lab}}_{n+1}\in\cdot\given o_{0:n})$ and
$(1-\epsilon_{\mathrm e})\mathbb P^{\mathrm{ep}}(y^{\mathrm{lab}}_{n}\in\cdot\given o_{0:n-1})
+\epsilon_{\mathrm e}\mathbb P^{\mathrm{ep}}(y^{\mathrm{lab}}_{n}\in\cdot\given o_{0:n})$,
the aggregate update applied to the route's own filtered law and label posterior, the
problem is settled by a route, defined for every step, whose discrepancy is
$o(\epsilon_{\mathrm e})$ uniformly over record histories and over episodes $n$ up to
any fixed multiple of $1/\epsilon_{\mathrm e}$, or by a lower bound of order
$\epsilon_{\mathrm e}$ on the supremum of the discrepancy over such histories and
episodes for every such route. Whether a declared update intertwines the fine and
parent evolutions is part of \Cref{open:rg-pointwise-parent-dynamics}; $\theta_n$
belongs to the slow block of \Cref{def:obs-fast-slow-blocks}.

\subsection{The remaining hierarchy question}
\label{sec:full-graph-vfe-target}

The prepared parent of \Cref{sec:rg-memory-parent-mechanism}, the memory cost of
partial refresh in \Cref{sec:rg-memory-projection}, and the successive formation of
\Cref{thm:finite-depth-adaptive-formation} are constructions at stated hypotheses, not
consequences of arbitrary descent.

\openproblemheading{Persistent multiscale hierarchy}{open:full-graph-hierarchy}
Starting from $c_*$ with a finite multigraph datum $\mathfrak N_0$, construct a
normalized partition law and parent--child kernels such that descent of
\eqref{eq:full-graph-vfe-global} produces persistent nested blocks, the pushforward
rules \eqref{eq:full-graph-edge-pushforward} to
\eqref{eq:full-graph-structural-channel-composition} and the link lift
\eqref{eq:full-graph-coarse-links} compose across depths, the parent family has a
controlled closure residual in the sense of \Cref{prop:rg-interaction-residual}, and
non-flat holonomy is kept or stabilized without imposing flatness. What holds: every
candidate hierarchy has the exact accounting of
\Cref{thm:full-graph-vfe-decomposition}; at one depth the fast--slow typing of
\Cref{def:obs-fast-slow-blocks} and its conditional slow-manifold limit,
\Cref{open:obs-slow-manifold}, are available; \Cref{thm:finite-depth-adaptive-formation}
forms nested blocks under a supplied energy hierarchy whose strengths are tied to the
temperature; and \Cref{prop:rg-formation-endogenous-cycle} forms one level. What is
unproved is that descent produces a second level that depends on the first. The
problem is posed in the class of models at $c_*$ whose structural marks are drawn from
a finite partition catalog with a partition factor inside the generative and the
recognition law (\Cref{def:rg-geometry-ensemble,def:full-graph-partition-factor}),
whose recognition coordinates descend \eqref{eq:full-graph-vfe-global} by a reversible
diffusion at temperature $\mathrm T$ with rate-one redraws of the discrete marks from
their conditional laws, as in \Cref{def:rg-formation-cycle-experiment}, and whose
energies, precisions, prior ratios, hyperparameters, and mobilities do not depend on
$\mathrm T$, as $\mathrm T\to0$. It is settled by three results together, or by a
proof that the first of them cannot hold in this class. The first is a lineage of at
least two levels, formed with probability bounded below from every initial law with a
density and a finite second moment, whose second level meets two conditions. It
depends on the first level: for a declared handoff time, a declared coalescence event
whose block is a union of at least two first-level blocks and whose windows are
measured from the handoff, one declared statistic of the handoff state (the
probability of the event, of its occurrence by declared times, or of a declared
block), and a handoff class fixed before the comparison that contains every
first-level handoff and, with probability bounded below, the handoffs of histories
without first-level formation, the event has probability bounded below after
first-level formation, does not have probability tending to one uniformly over the
handoff class at the declared windows, and the statistic exceeds its value after a
counterfactual handoff in the class by a margin bounded below. It is a coherence
episode rather than restored equilibrium or coincidental coherence: a restoring
interaction, evaluated on the law of the coordinates given the realized marks or on
their marginal over the marks and not on a conditional law held fixed by hand, holds
the second-level block's coherence at a declared resolution, with the restoring rate
times the squared resolution divided by $\mathrm T$ tending to infinity, on horizons
short compared with the stationary time of entrance into the event from a set of
stationary mass bounded below, or with a probability exceeding its stationary value by
a margin bounded below. The second result is a nondegenerate intrinsic partition
selector (\Cref{open:cg-natural-selector}), and the third is a separated slow manifold
on which depth-$(\ell+1)$ variables feed back consistently into depth-$\ell$ kernels.
The negative resolution is a proof that in this class no second level produced by
descent, apart from a supplied-scale sequence whose event occurs with probability
tending to one uniformly over the handoff class at the declared windows, both meets
the restoring and horizon conditions and exceeds its counterfactual statistic by
such a margin.

\section*{Summary for physicists}

The multiscale free energy is that of one generative and one recognition law over the
whole tower, not a sum of energies chosen scale by scale. Bayes reversal of a
bottom-up structural channel conditions every depth on the whole next-depth state,
graph datum included. Partitions enter as a normalized factor, normalized branch by
branch or globally, and coarse links are the dressed links of all copies, so internal
loops keep their holonomy. A parent can couple to child coordinates through a downward kernel that
reads it and to child structural marks through the partition and graph factors. A finite acyclic tower has a declared root, and the agents
participate in their own top through coordinate ascent, a learned aggregate prior, and
delayed feedback. A sequence of episodes whose parameter update reads a coupling
selected by a declared rule is itself one joint, with the product of the episode
evidences as its evidence (\Cref{prop:full-graph-episode-lifted-joint}).


\part{Formation and Interpretation}
\chapter{Formation Witnesses}
\label{ch:finite-depth-formation}

Exact parents are normalized pushforwards of the fine law
(\Cref{thm:rg-exact-coarse-vfe}) but in general not autonomous reduced
dynamics (\Cref{thm:rg-projection-memory}). This chapter asks whether, in one
normalized experiment at \(c_*\), coherent components can form in succession
while every parent is activated from the complete preceding recognition law,
retains every loop coordinate and source assignment, and learns its own copy
of a model parameter from its own law. Two witnesses answer yes under
explicit restrictions. \Cref{thm:finite-depth-adaptive-formation} treats a
triangle chain with a supplied hierarchy of interaction energies: at each
finite depth \(L\) the coherent components pass through
\(3\cdot2^{L-1}\to2^{L-1}\to\cdots\to1\) with probability tending to one,
while the discarded conditional losses, parameter covectors, assignment
histories, and continuous paths obey separate vanishing bounds.
\Cref{prop:rg-formation-endogenous-cycle} treats a homogeneous four-cycle
without supplied hierarchy, in which the source assignment decides which
pairs form, with probability bounded below independently of the small
parameter. \Cref{prop:rg-gated-depth-one} bounds what that mechanism can do
on any finite graph: in the gated family of
\Cref{def:rg-gated-translation-family}, the stationary coherent-block
partition is, up to the error of \cref{eq:rg-gated-block-law}, the
partition into the components of a minimum edge cover, and each merge
beyond such a cover costs a factor of the temperature. The proposition and
the conjecture of
the interpretation
chapter (\Cref{prop:phil-seeded-contrast-formation,conj:phil-contrast-formation})
are posed on the four-cycle family of the second witness. Both witnesses are
approximation statements under a stated preparation and regime, weaker than
the exact intertwining that an autonomous parent dynamics would require
(\Cref{open:rg-pointwise-parent-dynamics}). Both realize the general
construction for the translation group \((\R,+)\) on scalar sample fibers:
agents are local sections at \(c_*\) (\Cref{def:geo-agent}), one common frame
per agent lets one independent link translate both channels
(\Cref{def:geo-graph-links}), and the free energy is that of
\Cref{thm:obs-collective-vfe}.

Letters of the notation table that would collide with chapter-local objects
are replaced by the symbols of \cref{tab:rg-formation-notation}; upright
\(\mathrm J,\mathrm Q,\mathrm A,\mathrm T,\mathrm c_L,\mathrm C_L\) are local
scalars, the configuration is bold, and \(\ell'<\ell\) whenever two levels
occur.

\begin{sciencetable}{Chapter-local symbols of the formation families, shared with the proof appendix.}
\label{tab:rg-formation-notation}
\begin{tabularx}{\textwidth}{>{\raggedright\arraybackslash}p{0.28\textwidth}>{\raggedright\arraybackslash}X}
\toprule
\textbf{Symbol} & \textbf{Meaning} \\
\midrule
\(L\), \(N_\triangle=2^{L-1}\), \(N=3N_\triangle\), \(\ell,\ell'\) & depth, number of triangles, number of agents, level indices \\
\(\mathbf x=(\bar k,\bar m,\omega)\), \(\mathbf X_t\) & configuration of law presentations (belief means, model means, link translations) and its process \\
\(\mathfrak d_i=\bar k_i-\bar m_i\), \(\mathfrak D\), \(r_e\) & belief--model discrepancy, its row matrix, edge residual \\
\(\mathfrak h_q\), \(h_q\), \(O_\triangle\) & oriented triangle holonomy, normalized loop coordinate, contrast rows \\
\(\mathcal U,\mathcal H,\mathcal K\); \(W_\ell,\mathcal U_\ell,M_\ell,P_\ell\) & normal, loop, and gauge spaces; retained spaces and projections \\
\(H\), \(H_\ell\), \(\mathsf Q_\ell\) & stiffness, its Schur complement on \(M_\ell\), orthonormal flag columns \\
\(\mathrm J_\ell\), \(\mathrm Q_\ell\), \(\mathrm A\), \(\mathrm T\), \(t_*\) & energy strengths, their exponents, prepared amplitude, temperature, horizon \\
\(\sigma^2_{\mathrm{ev}}\), \(R_{\mathrm{ev}}\), \(\nu^*_i\) & evaluator variance, precision of the evaluator restriction cost, optimal evaluator coupling \\
\(\vartheta\), \(\Omega_0(\vartheta)\), \(\mathsf M\), \(\mathsf m=\mathrm T\mathsf M\) & learned scalar, prior precision, complete Fisher metric and its scaled form \\
\(\kappa\), \(z_\star\), \(o_{\mathcal K}\), \(o_{\mathrm{cal}}\), \(\mathsf o\) & root observation precision, root record, root and calibration observations, observed records \\
\(S\in\mathcal J\), \(\mathcal N_i^\triangle\), \(i_0\), \(\pi_{ij}(\vartheta)\), \(\pi_S(\vartheta)\) & complete source assignment, internal source palette, distinguished receiver, source priors \\
\(\chi\), \(\phi_S\), \(\delta\), \(B_0\), \(B_*\) & source function, source potential, record coefficient, and its two bounds \\
\(\beta_S\), \(\beta^\ell_{SS'}\), \(\bar\beta_{jk}\) & complete posterior source row, inherited proposal row, averaged fast row \\
\(q^\star_\vartheta\), \(q^{\star,\ell}_\vartheta\), \(\boldsymbol\Pi_{\vartheta,\mathsf o}\), \(q^\ell_t\) & restricted optimizer and its level-\(\ell\) push, physical posterior, recognition density \\
\(C_{\ell\leftarrow\ell'}\), \(\mathcal D^{\ell,\mu}_e\), \(\Delta_{\ell':\ell}\) & marked channel, inherited edge statistic, discarded conditional loss \\
\(\Fenergy_{\mathrm T}\), \(\mathrm C_{\mathrm T}\), \(\mathrm E_\vartheta\), \(U_{\vartheta,S}\), \(V_{\ell,S}\), \(\psi_{\ell,S}\), \(g_\ell\) & scaled restricted free energy, its offset, quadratic and sheet energies, parent potential and tilt, parameter covector \\
\(\mathscr L^\ell_\vartheta\) & inference generator at level \(\ell\) \\
\(t_\ell\), \(I_\ell\), \(t^{\mathrm{act}}_\ell\), \(k_L\), \(\eta_L\), \(\varsigma_\epsilon\), \(r_\epsilon\) & window start, window, activation time, entrance constant, tolerance, \(1+\log(1/\epsilon)\), parameter remainder \\
\(\mathrm c_L\), \(\mathrm C_L\), \(\mathsf j_m\) & positive constants depending on \(L\) and the fixed data; proposal marks \\
\(\mathsf B_{sr}\), \(\mathfrak j_\ell\), \(\mathfrak Y_\ell\), \(\mathfrak y_\ell\), \(\mathfrak x_\ell\), \(\varrho_s\), \(\mathrm Q_{\mathrm T}\), \(\mathrm Q_\Delta\) & flag blocks of the residual rows, least eigenvalue of a diagonal Gram block, forcing map and core coefficient of an interface, current-layer transient, strength ratio, exponents of the thermal ratio and of a loss target \\
\((V,E)\), \(\mathfrak n\), \(F\), \(\operatorname{cc}(F)\), \(\operatorname{att}(S)\) & graph of the gated family, its matching number, an edge set, the number of components of \((V,F)\), the edge set attended under \(S\) \\
\(\mathbf a_e\), \(\mathbf a_F\), \(W_F\), \(\mathrm s_F\) & residual pair map of an edge, its stack over \(F\), the kernel of the stack, the least positive singular value of the stack \\
\(H_S\), \(\mathrm k_S\), \(\mathrm h_S\), \(\Omega\), \(Z_S\), \(\mathrm z_S\), \(V_\pi\), \(V_{\min}\) & sheet stiffness, its rank and least positive eigenvalue, reference precision, sheet weight and its constant, source-marginal and gated energies \\
\(\operatorname{Coh}_{\mathrm a}(\mathbf x)\), \(\mathfrak P_{\mathrm a}(\mathbf x)\), \(\mathrm a\) & coherent edges and coherent blocks of a configuration; resolution \\
\bottomrule
\end{tabularx}
\end{sciencetable}

\paragraph{Scope.}
For the finite-depth family this chapter declines the following claims: that
attention or learning drives formation, since coherence is linear relaxation
of the prepared mean, the source record saturates at the coherence scale
(\Cref{rem:rg-formation-saturation}), and learning is calibration dominated up
to \(\epsilon^{2\mathrm Q_1+20}\); that per-agent models are slow, since means
and links have unit mobility and \(\mathfrak d_i\) is a level-one row; that
loop holonomy is of order one or nonabelian, since the group is \((\R,+)\) and
the retained loops are thermally negligible; that there is an inter-agent
model-channel record, since only the belief channel is observed; that parents
are coarse agents with their own belief and model fibers, since a parent
removes only normal coordinates; that source learning creates the supplied
contrasts; that the doubling recursion \(\mathrm Q_\ell=2\mathrm Q_{\ell+1}+20\)
is required, as distinct from the doubly exponential separation of inner base
\(4/3\) that \Cref{prop:rg-formation-gaussian-core} shows the Gaussian core to
require; that the continuous path laws converge in total variation or are exactly lumpable; and
that the construction saves computation or uses depth-uniform resources. For
the four-cycle it declines generic formation, formation of a matching from
every initial assignment, a unique block chosen by the invariant full law,
all-slow parents, exact closure under arbitrary source switching, permanent
membership, repeated coalescence, and a physical posterior conditioned on the
episode. The corresponding constructions are
\Cref{open:rg-formation-regime-sharpness,open:rg-formation-adaptive-family,open:rg-formation-general}.
\status{NOT-CLAIMED}

\section{The pointwise carrier and its restoring flag}
\label{sec:rg-formation-carrier}

\definitionheading{Triangle-chain carrier}{def:rg-formation-carrier}
Fix an integer \(L\ge2\), set \(N_\triangle=2^{L-1}\), and arrange
\(N_\triangle\) disjoint triangles in a chain, with one bridge between each
consecutive pair. Each triangle is cyclically oriented and every bridge has
one fixed orientation, so the graph has \(N=3N_\triangle\) vertices (the
agents) and \(4N_\triangle-1\) oriented edge copies. Bridge number \(j\),
between triangles \(j\) and \(j+1\), has level \(\ell(j)=2+v_2(j)\), where
\(v_2(j)\) is the exponent of two dividing \(j\); triangle edges have level
one. Energy strengths \(\mathrm J_1\ge\cdots\ge\mathrm J_L>0\) are attached to
the levels. The group \((\R,+)\) acts by translation on two scalar sample
fibers at \(c_*\); fix \(0<\sigma^2_{\mathrm{ev}}<1\). A configuration and its
law presentations are
\begin{equation}
\begin{aligned}
\mathbf x&=(\bar k,\bar m,\omega)\in\R^{2N+4N_\triangle-1}=\R^{10N_\triangle-1},&
\Theta_e(k)&=k+\omega_e,\\
q_i^b(\cdot\mid\mathbf x)&=\mathcal N(\bar k_i,1),&
q_i^m(\cdot\mid\mathbf x)&=\mathcal N(\bar m_i,\sigma^2_{\mathrm{ev}}),\\
\mathfrak d_i(\mathbf x)&=\bar k_i-\bar m_i,&
r_e(\mathbf x)&=\bar k_i-\bar k_j-\omega_e\quad(e:j\to i),
\end{aligned}
\label{eq:rg-formation-configuration}
\end{equation}
with \(\omega_{\bar e}=-\omega_e\). A passive frame change \(g_i\) sends
\(\bar k_i,\bar m_i,\omega_e\) to \(\bar k_i+g_i\), \(\bar m_i+g_i\),
\(\omega_e+g_i-g_j\), and re-expresses reference centers and observation
records in the same way. For triangle \(q\) put
\(\mathfrak h_q=\sum_{e\in q}\omega_e\) and \(h_q=\mathfrak h_q/\sqrt3\), and
let
\(O_\triangle=\begin{pmatrix}1/\sqrt2&-1/\sqrt2&0\\1/\sqrt6&1/\sqrt6&-2/\sqrt6\end{pmatrix}\).
Let \(B_1\) stack the \(N\) rows \(\mathfrak d_i\) and the \(2N_\triangle\)
internal contrast rows \(O_\triangle r_q\), where \(r_q\) is the vector of the
three residuals of triangle \(q\); for \(\ell\ge2\) let \(B_\ell\) stack the
bridge residuals of level \(\ell\), and let \(B\) stack all of them. With the
oriented incidence operator \((\partial t)_e=t_i-t_j\),
\begin{equation}
\begin{aligned}
\mathcal K&=\{(t,t,\partial t):t\in\R^N\},&
\mathcal H&=\operatorname{span}\{\text{normalized loop covectors}\},\\
\mathcal U&=\operatorname{range}B^\top,&
H&=B^\top\operatorname{diag}(\mathrm J_\ell I)B+\mathrm J_1P_{\mathcal H}.
\end{aligned}
\label{eq:rg-formation-orthogonal-split}
\end{equation}
The retained spaces and their Euclidean orthogonal projections are
\begin{equation}
W_\ell=\ker(B_1,\ldots,B_\ell),\qquad
\mathcal U_\ell=W_\ell\cap\mathcal U,\qquad
M_\ell=\mathcal U_\ell\oplus\mathcal H,\qquad
P_\ell:\R^{10N_\triangle-1}\to W_\ell,\qquad P_0=I.
\label{eq:rg-formation-projections}
\end{equation}

\lemmaheading{Carrier algebra}{lem:rg-formation-carrier}
In \Cref{def:rg-formation-carrier}, \(\mathfrak d_i\) and \(r_e\) are
invariant under passive frame changes, and every triangle satisfies
\begin{equation}
|r_q|^2=|O_\triangle r_q|^2+h_q^2,
\label{eq:rg-formation-loop-decomposition}
\end{equation}
so the loop energy belongs to the original sum of edge mismatches. The rows of
\(B\) are independent, \(\operatorname{rank}B=6N_\triangle-1\), the normalized
loop covectors are orthonormal and orthogonal to every row of \(B\), and
\(\R^{10N_\triangle-1}=\mathcal U\oplus\mathcal H\oplus\mathcal K\) is an
orthogonal direct sum. The stiffness has kernel \(\mathcal K\) and
\(\mathbf x^\top H\mathbf x=\mathrm J_1\sum_i\mathfrak d_i(\mathbf x)^2
+\sum_e\mathrm J_{\ell(e)}r_e(\mathbf x)^2\). Finally
\(P_\ell P_{\ell'}=P_\ell\) for \(\ell'<\ell\), every \(W_\ell\) contains
\(\mathcal H\oplus\mathcal K\), and
\begin{equation}
\dim W_\ell=4N_\triangle-1+\frac{N_\triangle}{2^{\ell-1}}\qquad(1\le\ell\le L).
\label{eq:rg-formation-dimensions}
\end{equation}
\status{ESTABLISHED}

\noindent\emph{Proof.}
A passive change adds \(g_i-g_i\) to \(\mathfrak d_i\) and
\(g_i-g_j-(g_i-g_j)\) to \(r_e\). Around an oriented triangle the vertex means
telescope, so \(\sum_{e\in q}r_e=-\mathfrak h_q\); the rows of \(O_\triangle\)
are orthonormal and orthogonal to \((1,1,1)/\sqrt3\), which gives
\cref{eq:rg-formation-loop-decomposition}. In a linear relation among the rows
of \(B\), the link columns occur only in contrast and bridge rows, with the
independent patterns \(-O_\triangle\) on the triangle edges and \(-1\) on each
bridge, so every contrast and bridge coefficient vanishes; the model-mean
columns then force the discrepancy coefficients to vanish. A normalized loop
covector has link entries \(1/\sqrt3\) on one triangle and nothing else, so it
is orthogonal to the discrepancy rows, to the contrast rows because
\(O_\triangle(1,1,1)^\top=0\), and to the bridges. Every row of \(B\) and every
loop covector annihilates \((t,t,\partial t)\), and
\((6N_\triangle-1)+N_\triangle+N=10N_\triangle-1\). The quadratic identity is
the sum of \(\mathrm J_\ell|B_\ell\mathbf x|^2\) and
\(\mathrm J_1|P_{\mathcal H}\mathbf x|^2\) combined with
\cref{eq:rg-formation-loop-decomposition}; it vanishes exactly when
\(B\mathbf x=0\) and \(P_{\mathcal H}\mathbf x=0\), that is on \(\mathcal K\).
The spaces \(W_\ell\) decrease, which gives \(P_\ell P_{\ell'}=P_\ell\), and
\(\mathcal H\oplus\mathcal K\) lies in each of them. There are
\(N_\triangle/2^{\ell''-1}\) bridges of level \(\ell''\ge2\), so
\(\operatorname{rank}(B_1,\ldots,B_\ell)=3N_\triangle+2N_\triangle
+N_\triangle(1-2^{1-\ell})\), which gives
\cref{eq:rg-formation-dimensions}. \(\square\)

The restoring estimates of the proof appendix rest on the bounds
\cref{eq:rg-formation-weighted-flag} of \(H\) by weighted energies
\(\sum_\ell\mathrm J_\ell|w_\ell|^2\) along orthonormal flag columns
\(\mathsf Q_\ell\) spanning \(\mathcal U_{\ell-1}\cap\mathcal U_\ell^\perp\).

\section{One normalized experiment with a nonlinear source record}
\label{sec:rg-formation-normalized-model}

\definitionheading{Normalized formation experiment}{def:rg-formation-experiment}
At each vertex the physical model sample is \(m_i\sim\mathcal N(\bar m_i,
\sigma^2_{\mathrm{ev}})\) and the evaluator sends \(m_i\) to
\(\mathcal N(m_i,1-\sigma^2_{\mathrm{ev}})\) on the belief sample \(k_i\).
Among couplings with the prescribed marginals
\(\mathcal N(\bar m_i,\sigma^2_{\mathrm{ev}})\) and \(\mathcal N(\bar k_i,1)\),
the recognition coupling at a fixed configuration is
\begin{equation}
m_i=\bar m_i+\xi_i,\qquad
k_i=\bar k_i+\xi_i+\xi'_i,\qquad
\xi_i\sim\mathcal N(0,\sigma^2_{\mathrm{ev}}),\quad
\xi'_i\sim\mathcal N(0,1-\sigma^2_{\mathrm{ev}}),
\label{eq:rg-formation-evaluator-lift}
\end{equation}
with independent noises; write \(\nu_i^*\) for its law and
\(R_{\mathrm{ev}}=\mathfrak D^\top\mathfrak D/(1-\sigma^2_{\mathrm{ev}})\),
where \(\mathfrak D\mathbf x=\bar k-\bar m\). Let \(\vartheta\in[1,21/10]\) be
the learned scalar and \(0<\mathrm T\le\mathrm J_1(1-\sigma^2_{\mathrm{ev}})/2\).
The configuration prior and its Lebesgue normalizer are
\begin{equation}
\begin{aligned}
\Omega_0(\vartheta)&=\vartheta H+P_{\mathcal K}-\mathrm TR_{\mathrm{ev}},\\
\mathbb P^0_\vartheta(d\mathbf x)
&=Z_0(\vartheta)^{-1}
  \exp\left[-\frac{\mathbf x^\top\Omega_0(\vartheta)\mathbf x}{2\mathrm T}\right]d\mathbf x,\\
Z_0(\vartheta)&=(2\pi\mathrm T)^{(10N_\triangle-1)/2}\det\Omega_0(\vartheta)^{-1/2}.
\end{aligned}
\label{eq:rg-formation-physical-prior}
\end{equation}
Fix \(0<\kappa\le1\) and a nonzero \(z_\star\in\mathcal K\). The observation
kernels are
\begin{equation}
o_{\mathcal K}\mid\mathbf x\sim
  \mathcal N_{\mathcal K}\bigl(P_{\mathcal K}\mathbf x,(\mathrm T/\kappa)I_{\mathcal K}\bigr),
\qquad
o_{\mathrm{cal}}\mid\vartheta\sim\mathcal N(\vartheta,\mathrm T),
\label{eq:rg-formation-observation-kernels}
\end{equation}
with the records \(o_{\mathcal K}=z_\star\) and \(o_{\mathrm{cal}}=2\); the
first is a density on the \(N\)-dimensional carrier \(\mathcal K\). Each
receiver \(i\) has the two candidate sources \(\mathcal N_i^\triangle\) inside
its own triangle, and the complete assignment is
\(S=(S_i)_i\in\mathcal J=\prod_i\mathcal N_i^\triangle\), with
\(|\mathcal J|=2^N\); every bridge is an added reversal pair of copies that keeps
its link and interaction energy but lies outside the source palettes. A distinguished receiver \(i_0\) has sources
labeled \(+\) and \(-\), and the independent source priors are
\begin{equation}
\begin{aligned}
\pi_{i_0,+}(\vartheta)
  &=\operatorname{sig}\bigl(4(\vartheta-3/2)\bigr),&
\pi_{i_0,-}(\vartheta)&=1-\pi_{i_0,+}(\vartheta),\\
\pi_{ij}(\vartheta)&=\tfrac12\quad(i\ne i_0),&
\pi_S(\vartheta)&=\prod_i\pi_{iS_i}(\vartheta),
\end{aligned}
\label{eq:rg-formation-source-priors}
\end{equation}
where \(\operatorname{sig}(u)=(1+e^{-u})^{-1}\). Set
\(B_0=10^{-4}\), \(\delta=B_0/N\), and \(\chi(r)=1-e^{-r^2/2}\). At each
receiver a binary record \(E_i\), conditionally independent given
\((\mathbf x,S)\), has
\begin{equation}
\begin{aligned}
\Pr(E_i=1\mid\mathbf x,S_i=j)
  &=q_0\exp\left[-\delta\chi\left(\sqrt{\mathrm J_1/\mathrm T} r_{ij}(\mathbf x)\right)\right],
       \qquad q_0=\tfrac12,\\
\phi_S(\mathbf x)&=\delta\sum_i
       \chi\left(\sqrt{\mathrm J_1/\mathrm T} r_{iS_i}(\mathbf x)\right),
       \qquad 0\le\phi_S\le B_0,
\end{aligned}
\label{eq:rg-formation-source-record}
\end{equation}
where \(r_{ij}\) is the residual of the edge from source \(j\) to receiver
\(i\), and the observed record is \(E_i=1\) for every \(i\). Since
\(\mathrm J_1r_{ij}^2/(2\mathrm T)=(\mathrm J_1/\mathrm T)
\KL(\mathcal N(\bar k_i,1)\Vert\mathcal N(\bar k_j+\omega_{ij},1))\), the
record is a bounded observation of the standardized edge mismatch. The
complete joint is
\begin{equation}
\begin{aligned}
\mathbb P_\vartheta
={}&\mathbb P^0_\vartheta(d\mathbf x)
 L_{\mathcal K}(do_{\mathcal K}\mid\mathbf x)
 L_{\mathrm{cal}}(do_{\mathrm{cal}}\mid\vartheta)\\
&{}\times\prod_i\left[
 \pi_i(dS_i\mid\vartheta)L_i(dE_i\mid\mathbf x,S_i)
 \mathcal N(dm_i;\bar m_i,\sigma^2_{\mathrm{ev}})
 \mathcal N(dk_i;m_i,1-\sigma^2_{\mathrm{ev}})\right].
\end{aligned}
\label{eq:rg-formation-full-joint}
\end{equation}

\propositionheading{Normalization and the evaluator optimum}{prop:rg-formation-normalized-experiment}
In \Cref{def:rg-formation-experiment}, for every admissible coupling
\(\nu_i\) of the two prescribed marginals,
\begin{equation}
\KL\bigl(\nu_i\Vert
  \mathcal N(\bar m_i,\sigma^2_{\mathrm{ev}})\operatorname{ev}_i\bigr)
=\KL(\nu_i\Vert\nu_i^*)+\frac{\mathfrak d_i^2}{2(1-\sigma^2_{\mathrm{ev}})},
\label{eq:rg-formation-evaluator-optimum}
\end{equation}
so \(\nu_i^*\) is the optimal coupling and the total restriction cost is
\(\mathbf x^\top R_{\mathrm{ev}}\mathbf x/2\). Moreover
\(\mathrm TR_{\mathrm{ev}}\preceq H/2\), hence
\(\Omega_0(\vartheta)\succeq(\vartheta/2)H+P_{\mathcal K}\succ0\) and the prior
is a normalized correlated Gaussian, and the joint
\cref{eq:rg-formation-full-joint} has total mass one, contains no kernel that
takes a recognition law as input, and gives the selected observation a
strictly positive density. \status{ESTABLISHED}

The term \(-\mathrm TR_{\mathrm{ev}}\) of the prior cancels the restriction
cost \(\tfrac{\mathrm T}2\mathbf x^\top R_{\mathrm{ev}}\mathbf x\) that
\cref{eq:rg-formation-evaluator-optimum} adds, so the restricted optimizer
below has a temperature-independent quadratic energy.

\section{Physical posterior, restricted optimizer, and recognition}
\label{sec:rg-formation-laws}

\definitionheading{Three laws of the formation experiment}{def:rg-formation-laws}
Write \(\mathsf o=(z_\star,2,\mathbf 1)\) for the observed records. The
recognition family consists of the laws with a smooth positive density
\(q(\mathbf x,S)\) relative to Lebesgue measure times counting measure and
conditional evaluator coupling \(\prod_i\nu_i^*\); assignments and
configurations may be correlated, and the conditional assignment law need not
be the source posterior. With
\begin{equation}
\begin{aligned}
\widehat K_\vartheta&=\vartheta H+(1+\kappa)P_{\mathcal K},&
\mathrm E_\vartheta(\mathbf x)&=\tfrac12\mathbf x^\top\widehat K_\vartheta\mathbf x
                  -\kappa z_\star^\top\mathbf x,\\
\mathrm E_\vartheta^{\mathrm{phys}}(\mathbf x)
   &=\mathrm E_\vartheta(\mathbf x)-\tfrac{\mathrm T}2\mathbf x^\top R_{\mathrm{ev}}\mathbf x,&
U_{\vartheta,S}(\mathbf x)&=\mathrm E_\vartheta(\mathbf x)+\mathrm T\phi_S(\mathbf x),\\
Z_{\mathrm{opt}}(\vartheta)
   &=\sum_S\pi_S(\vartheta)
      \int e^{-\mathrm E_\vartheta(\mathbf x)/\mathrm T-\phi_S(\mathbf x)}d\mathbf x,&
Z_{\mathrm{phys}}(\vartheta)
   &=\sum_S\pi_S(\vartheta)
      \int e^{-\mathrm E_\vartheta^{\mathrm{phys}}(\mathbf x)/\mathrm T-\phi_S(\mathbf x)}d\mathbf x,
\end{aligned}
\label{eq:rg-formation-partition-integrals}
\end{equation}
the restricted optimizer is
\(q^\star_\vartheta(\mathbf x,S)=\pi_S(\vartheta)e^{-U_{\vartheta,S}(\mathbf x)/\mathrm T}
/Z_{\mathrm{opt}}(\vartheta)\), and the physical posterior
\(\boldsymbol\Pi_{\vartheta,\mathsf o}\) on \((\mathbf x,S)\) is the same
expression with \(\mathrm E_\vartheta^{\mathrm{phys}}\) and
\(Z_{\mathrm{phys}}\). The scaled restricted free energy
\(\Fenergy_{\mathrm T}(q,\vartheta)\) is \(\mathrm T\) times the negative ELBO
of the evaluator restriction, evaluated at the records \(\mathsf o\).

\propositionheading{Laws, restricted free energy, and the physical gap}{prop:rg-formation-laws}
The integrals \cref{eq:rg-formation-partition-integrals} are finite and
positive, both laws have full support, and their common conditional
assignment row at fixed \(\mathbf x\) is the product of the receiver posterior
rows,
\begin{equation}
\beta_S(\mathbf x,\vartheta)=
\frac{\pi_S(\vartheta)e^{-\phi_S(\mathbf x)}}
     {\sum_{S'}\pi_{S'}(\vartheta)e^{-\phi_{S'}(\mathbf x)}}.
\label{eq:rg-formation-full-source-row}
\end{equation}
With the scalars \(\mathrm B_{\mathrm T}\), \(\mathrm C_{\mathrm T}\), and
\(p_\vartheta(\mathsf o)\) of \cref{eq:rg-formation-evidence-scalar},
\begin{equation}
\begin{aligned}
\Fenergy_{\mathrm T}(q,\vartheta)
 &=\mathrm T\sum_S\int q(\mathbf x,S)\log q(\mathbf x,S) d\mathbf x
  +\E_q\bigl[\mathrm E_\vartheta(\mathbf x)+\mathrm T\phi_S(\mathbf x)
       -\mathrm T\log\pi_S(\vartheta)\bigr]
       +\mathrm B_{\mathrm T}(\vartheta)\\
 &=\mathrm T\KL(q\Vert q^\star_\vartheta)+\mathrm C_{\mathrm T}(\vartheta),
\end{aligned}
\label{eq:rg-formation-restricted-vfe}
\end{equation}
and the excess over the unconstrained physical optimum is
\(\mathrm C_{\mathrm T}+\mathrm T\log p_\vartheta(\mathsf o)
=\mathrm T\log(Z_{\mathrm{phys}}/Z_{\mathrm{opt}})\ge0\). Whenever
\(N\mathrm T\le(1-\sigma^2_{\mathrm{ev}})\mathrm J_1\),
\begin{equation}
0\le\mathrm T\log\frac{Z_{\mathrm{phys}}(\vartheta)}{Z_{\mathrm{opt}}(\vartheta)}
\le\frac{2e^{B_0}N\mathrm T^2}{(1-\sigma^2_{\mathrm{ev}})\mathrm J_1},
\qquad
\KL(q^\star_\vartheta\Vert\boldsymbol\Pi_{\vartheta,\mathsf o})
\le\frac{2e^{B_0}N\mathrm T}{(1-\sigma^2_{\mathrm{ev}})\mathrm J_1}.
\label{eq:rg-formation-physical-optimizer-gap}
\end{equation}
\status{ESTABLISHED}

\section{Exact marked parents and their inherited inference}
\label{sec:rg-formation-parents}

\definitionheading{Marked parents and inherited edge statistics}{def:rg-formation-parents}
The parent carrier at level \(\ell\) is \(W_\ell\times\mathcal J\). For
\(\ell'<\ell\) the deterministic Markov kernel
\begin{equation}
C_{\ell\leftarrow\ell'}((y,S),\cdot)=\delta_{(P_\ell y,S)},
\qquad C_\ell=C_{\ell\leftarrow0},
\label{eq:rg-formation-marked-channel}
\end{equation}
extended by the identity on observations, pushes the generative joint, the
physical posterior, the restricted optimizer, and the recognition law, giving
\(q^{\star,\ell}_\vartheta=(C_\ell)_\#q^\star_\vartheta\),
\(\boldsymbol\Pi^\ell_{\vartheta,\mathsf o}\), and \((C_\ell)_\#q\), and
erases no assignment. The evaluator samples \((m_i,k_i)_i\), the records
\((E_i)_i\), and the discarded normal coordinates leave the parent carrier
through the disintegration of \cref{eq:rg-formation-full-joint} given
\((P_\ell\mathbf x,S)\), so their correlations are retained and no mark is
replaced by its value at \(P_\ell\mathbf x\). For an original edge
\(e\) and each \(\mu_\vartheta\in
\{\boldsymbol\Pi_{\vartheta,\mathsf o},q^\star_\vartheta\}\), the inherited
edge statistic is
\begin{equation}
\mathcal D_e^{\ell,\mu}(y,S)
=\E_{\mu_\vartheta}
 \left[\tfrac12r_e(\mathbf X)^2\mid P_\ell\mathbf X=y,S\right],
\label{eq:rg-formation-inherited-edge-statistic}
\end{equation}
the discrepancy observables for \(\mathfrak d_i\) are defined in the same way,
and at \(\ell=0\) both reduce to the configuration residuals. For a density
\(q^\ell\) on the marked parent carrier, the parent free energy is
\begin{equation}
\Fenergy^\ell_{\mathrm T}(q^\ell,\vartheta)
  =\mathrm T\KL(q^\ell\Vert q^{\star,\ell}_\vartheta)+\mathrm C_{\mathrm T}(\vartheta).
\label{eq:rg-formation-parent-vfe}
\end{equation}

\propositionheading{Exact parent potential and discarded conditional cost}{prop:rg-formation-parent-potential}
The kernels \cref{eq:rg-formation-marked-channel} compose,
\(C_{\ell''\leftarrow\ell}\circ C_{\ell\leftarrow\ell'}=C_{\ell''\leftarrow\ell'}\),
and positive smooth versions of every pushed law and conditional exist. The
root factor of \(q^{\star,\ell}_\vartheta\) is independent of the assignment
and of \(\vartheta\). On \(M_\ell\) let \(H_\ell\) be the Schur complement of
\(H|_{\mathcal K^\perp}\) and \(\gamma^{\mathrm{hid}}_{\ell,\vartheta}(dz\mid y)\)
the cumulative hidden Gaussian conditional of the quadratic optimizer, whose
mean is independent of \(\vartheta\) and whose covariance is
\(\mathrm T/\vartheta\) times a fixed inverse precision. Integration at fixed
\(S\) gives the sheet potential
\begin{equation}
\begin{aligned}
V_{\ell,S}(y,\vartheta)
 &=\tfrac12\vartheta y^\top H_\ell y+\mathrm C_\ell^{\mathrm G}(\vartheta)
       +\mathrm T\psi_{\ell,S}(y,\vartheta),\\
\psi_{\ell,S}(y,\vartheta)
 &=-\log\E_{\gamma^{\mathrm{hid}}_{\ell,\vartheta}(\cdot\mid y)}
                  e^{-\phi_S},&
0&\le\psi_{\ell,S}\le B_0,\\
\frac{d}{d\vartheta}\mathrm C_\ell^{\mathrm G}(\vartheta)
 &=\frac{\mathrm T \mathrm d_\ell}{2\vartheta},
\end{aligned}
\label{eq:rg-formation-parent-potential}
\end{equation}
where \(\mathrm d_\ell=\dim W_\ell^\perp\) is the cumulative eliminated
dimension, \(\mathrm C_0^{\mathrm G}=0\), \(\psi_{0,S}=\phi_S\), and the
unchanged root potential is appended on \(W_\ell\). At one common parameter
and one common time,
\begin{equation}
\begin{aligned}
\Delta_{\ell':\ell}(q^{\ell'},\vartheta)
 &=\E_{(C_{\ell\leftarrow\ell'})_\#q^{\ell'}}
   \KL\left(q^{\ell'}(dz\mid y,S)
        \middle\Vert q^{\star,\ell'}_\vartheta(dz\mid y,S)\right),\\
\Fenergy_{\mathrm T}^{\ell'}(q^{\ell'},\vartheta)
 &-\Fenergy_{\mathrm T}^\ell\bigl((C_{\ell\leftarrow\ell'})_\#q^{\ell'},\vartheta\bigr)
   =\mathrm T\Delta_{\ell':\ell}(q^{\ell'},\vartheta),
\end{aligned}
\label{eq:rg-formation-conditional-cost}
\end{equation}
where \(z\) consists only of the discarded normal coordinates, loops and
\(S\) staying in the retained state. \status{ESTABLISHED}

\noindent\emph{Proof.}
Composition follows from \(P_{\ell''}P_\ell=P_{\ell''}\), and finite Gaussian
domination supplies the positive smooth versions. The potential is the
Gaussian integral computed in \Cref{lem:formation-parent-score}. The cost
identity is \Cref{thm:rg-exact-coarse-vfe} for the channel
\(C_{\ell\leftarrow\ell'}\), since the scalar \(\mathrm C_{\mathrm T}\) is
common to both sides and the channel leaves the observations unchanged.
\(\square\)

\definitionheading{Declared inference kinetics}{def:rg-formation-kinetics}
Every microscopic coordinate has unit continuous mobility, and a rate-one
Poisson proposal clock redraws the whole assignment from
\(\beta(\mathbf x,\vartheta)\) without moving any continuous coordinate, a
proposal being allowed to return the current assignment. The microscopic
generator on global Euclidean space times the finite set \(\mathcal J\) is
\begin{equation}
\mathscr L_\vartheta^0\mathsf f(\mathbf x,S)
=\mathrm T\Delta_{\mathbf x}\mathsf f-\nabla_{\mathbf x}U_{\vartheta,S}\cdot\nabla_{\mathbf x}\mathsf f
 +\sum_{S'}\beta_{S'}(\mathbf x,\vartheta)[\mathsf f(\mathbf x,S')-\mathsf f(\mathbf x,S)],
\label{eq:rg-formation-microscopic-generator}
\end{equation}
with diffusion covariance \(2\mathrm TI dt\). The inherited proposal row and
parent generator are
\begin{equation}
\begin{aligned}
\beta^\ell_{SS'}(y,\vartheta)
 &=\E_{q^\star_\vartheta}
     [\beta_{S'}(\mathbf X,\vartheta)\mid P_\ell\mathbf X=y,S],\\
\mathscr L_\vartheta^\ell\mathsf f(y,S)
 &=\mathrm T\Delta_y\mathsf f-\nabla_yV_{\ell,S}\cdot\nabla_y\mathsf f
   +\sum_{S'}\beta^\ell_{SS'}(y,\vartheta)[\mathsf f(y,S')-\mathsf f(y,S)],
\end{aligned}
\label{eq:rg-formation-parent-generator}
\end{equation}
with the unchanged root drift included when \(y\) denotes all of \(W_\ell\).
Self-proposals are kept, so every row sums to one and every exit rate is at
most one; these rows declare the parent kinetics rather than project the fine
process in time.

\propositionheading{Reversible inherited rows}{prop:rg-formation-rate-composition}
The rows \(\beta^\ell\) satisfy detailed balance with \(q^{\star,\ell}_\vartheta\),
and for \(\ell'<\ell\)
\begin{equation}
\beta^{\ell'}_{SS'}(y,\vartheta)\ge e^{-B_0}\pi_{S'}(\vartheta)>0,
\qquad
\E_{q^{\star,\ell'}_\vartheta}
 [\beta^{\ell'}_{SS'}(Y,\vartheta)\mid P_\ell Y=z,S]
 =\beta^\ell_{SS'}(z,\vartheta).
\label{eq:rg-formation-rate-composition}
\end{equation}
\status{ESTABLISHED}

\noindent\emph{Proof.}
The product \(q^\star_\vartheta(\mathbf x,S)\beta_{S'}(\mathbf x,\vartheta)\)
equals the continuous marginal of \(q^\star_\vartheta\) times
\(\beta_S\beta_{S'}\) and is symmetric in \(S,S'\); conditional integration
over a retained fiber gives detailed balance. The lower bound follows from
\cref{eq:rg-formation-full-source-row} and \(0\le\phi_S\le B_0\), and the
second identity is the tower property of conditional expectation. \(\square\)

\section{Complete model scores and dissipation}
\label{sec:rg-formation-learning}

\definitionheading{Own-law parameter dynamics}{def:rg-formation-learning}
Let \(\mathsf M(\vartheta)\) be the Fisher information of the complete
normalized experiment \cref{eq:rg-formation-full-joint} in \(\vartheta\),
before conditioning on the records, and \(\mathsf m(\vartheta)=
\mathrm T\mathsf M(\vartheta)\). With the model metric
\(\mathsf m(\vartheta)/\epsilon\), \(\epsilon>0\), each level \(\ell\) evolves
its own recognition density and its own parameter by
\begin{equation}
\partial_tq_t^\ell=(\mathscr L_{\vartheta_\ell(t)}^\ell)^*q_t^\ell,
\qquad
\dot\vartheta_\ell(t)
  =-\epsilon \mathsf m(\vartheta_\ell(t))^{-1}
        g_\ell(q_t^\ell,\vartheta_\ell(t)),
\qquad
g_\ell=\left.\partial_\vartheta\Fenergy_{\mathrm T}^\ell\right|_{q^\ell}.
\label{eq:rg-formation-own-law-dynamics}
\end{equation}
The parameter path is deterministic, since the feedback is an expectation
under the law of its level. Parents inherit this metric, with no new rate
multiplier at activation.

\propositionheading{Complete metric, covector, and dissipation}{prop:rg-formation-score-dissipation}
The complete Fisher metric is
\begin{equation}
\mathsf M(\vartheta)
=\frac1{\mathrm T}+\frac12\Tr\bigl[(\Omega_0(\vartheta)^{-1}H)^2\bigr]
 +16\pi_{i_0,+}(\vartheta)\bigl[1-\pi_{i_0,+}(\vartheta)\bigr],
\label{eq:rg-formation-model-fisher}
\end{equation}
and the full parameter covector of level \(\ell\) is
\begin{equation}
\begin{aligned}
g_\ell(q^\ell,\vartheta)
={}&\vartheta-2+\tfrac12\E_{q^\ell}y^\top H_\ell y
       +\frac{\mathrm T \mathrm d_\ell}{2\vartheta}
       +\mathrm T\E_{q^\ell}\partial_\vartheta\psi_{\ell,S}(y,\vartheta)\\
 &-\frac{\mathrm T}2\Tr\bigl(\Omega_0(\vartheta)^{-1}H\bigr)
       -\mathrm T\E_{q^\ell}\partial_\vartheta\log\pi_S(\vartheta),
\end{aligned}
\label{eq:rg-formation-complete-score}
\end{equation}
a partial derivative at fixed recognition law, which at \(\ell=0\), where
\(\mathrm d_0=0\) and \(\partial_\vartheta\phi_S=0\), is the microscopic score.
Along \cref{eq:rg-formation-own-law-dynamics}, with
\(\mathfrak r_S=q^\ell_S/q^{\star,\ell}_{\vartheta_\ell}(\cdot,S)\),
\begin{equation}
\begin{aligned}
\frac{d}{dt}\Fenergy_{\mathrm T}^\ell
={}&-\mathrm T^2\sum_S\int q_S^\ell
       \left|\nabla\log\mathfrak r_S\right|^2dy
 -\frac{\mathrm T}2\int\sum_{S,S'}q^{\star,\ell}_{\vartheta_\ell}(y,S)\beta^\ell_{SS'}(y,\vartheta_\ell)
       (\mathfrak r_S-\mathfrak r_{S'})(\log\mathfrak r_S-\log\mathfrak r_{S'}) dy\\
 &-\epsilon \mathsf m(\vartheta_\ell)^{-1}g_\ell(q^\ell,\vartheta_\ell)^2
 \le0.
\end{aligned}
\label{eq:rg-formation-dissipation}
\end{equation}
At the microscopic level the jump term equals
\(-\mathrm T\int q_{\mathbf X}[\KL(q(\cdot\mid\mathbf x)\Vert\beta)
+\KL(\beta\Vert q(\cdot\mid\mathbf x))]d\mathbf x\). \status{ESTABLISHED}

The proofs are in \cref{sec:rg-formation-exact-proof}; the identities hold for
the solutions admitted by \Cref{lem:formation-prepared-continuation}, whose
cutoff passage is \Cref{lem:formation-activation-entropy}.

\section{Preparation and successive activation}
\label{sec:rg-formation-preparation}

\definitionheading{Regime, preparation, and activation schedule}{def:rg-formation-regime}
For \(0<\epsilon<1\) the regime is
\begin{equation}
\mathrm Q_L=0,\qquad \mathrm Q_\ell=2\mathrm Q_{\ell+1}+20,\qquad
\mathrm J_\ell=\epsilon^{-\mathrm Q_\ell},\qquad
\mathrm A^2=\epsilon^{3\mathrm Q_1+20},\qquad
\mathrm T=\mathrm A^2\epsilon^6,\qquad t_*=\frac2\epsilon,
\label{eq:rg-formation-weakest-one-regime}
\end{equation}
with unit mobility and proposal rate and the record coefficient \(\delta\) of
\cref{eq:rg-formation-source-record}; the weakest energy is exactly
\(\mathrm J_L=1\). Put \(\varsigma_\epsilon=1+\log(1/\epsilon)\) and
\begin{equation}
r_\epsilon=\mathrm A^2\mathrm J_1+\mathrm T\varsigma_\epsilon^2
=\epsilon^{2\mathrm Q_1+20}+\epsilon^{3\mathrm Q_1+26}\varsigma_\epsilon^2.
\label{eq:rg-formation-parameter-remainder}
\end{equation}
Let \(\mathbf v_\star\in\mathcal U\) be the unit vector of
\Cref{lem:formation-preparation-margin}, with margin \(\mathrm m_L>0\), and
\(\mathbf v_0\) any vector of the neighborhood of \(\mathbf v_\star\) on which
the margin is at least \(\mathrm m_L/2\). The initial parameter is
\(\vartheta(0)=1\). The root is independent of the non-root coordinates and
the assignments, with law \(\mathcal N_{\mathcal K}(\kappa z_\star/(1+\kappa),
\mathrm T(1+\kappa)^{-1}I_{\mathcal K})\). On every assignment sheet the
reference non-root law has common mean \(\mathrm A\mathbf v_0\), zero loop
mean, and covariance \(\mathrm T(H|_{\mathcal K^\perp})^{-1}\); the positive
assignment masses are near the product prior at \(\vartheta=1\), and the
distinguished source mass \(\lambda_{i_0,+}(0)=\Pr\{S_{i_0}(0)=+\}\) satisfies
\(|\lambda_{i_0,+}(0)-\operatorname{sig}(-2)|<10^{-2}\). The preparation
neighborhood consists of the smooth within-sheet transports of this reference
with bounded log-density ratio, retaining the root factor, whose coupling
satisfies
\begin{equation}
\bigl|\widehat K_1^{1/2}(\mathbf X_0-\mathbf X_0^{\mathrm G})\bigr|
 <\epsilon\sqrt{\mathrm T}
\quad\text{almost surely};
\label{eq:rg-formation-preparation-transport}
\end{equation}
small compactly supported diffeomorphisms in standardized non-root coordinates
supply such transports, the neighborhood may depend on \(\epsilon\) and on
\(L\), and later laws need not stay Gaussian or factor across sheets. Given an
entrance constant \(k_L\), the windows and activation times are
\begin{equation}
t_\ell=\frac{k_L\log(1/\epsilon)}{\mathrm J_\ell},
\qquad I_\ell=[t_\ell,2t_\ell],\qquad
t^{\mathrm{act}}_0=0,\quad t^{\mathrm{act}}_\ell=2t_\ell.
\label{eq:rg-formation-activation-times}
\end{equation}
At \(t^{\mathrm{act}}_\ell\) the complete level-\((\ell-1)\) law is pushed
through \(C_{\ell\leftarrow\ell-1}\), its parameter is copied, and the new law
then evolves by \cref{eq:rg-formation-own-law-dynamics}; an independently
continued source process keeps its own law and parameter.

\remarkheading{The regime in closed form}{rem:rg-formation-regime}
Solving the recursion gives \(\mathrm Q_\ell=20(2^{L-\ell}-1)\), so
\begin{equation}
\mathrm Q_1=20(2^{L-1}-1),\qquad
\frac{\mathrm J_1}{\mathrm J_L}=\epsilon^{-\mathrm Q_1},\qquad
\mathrm A^2=\epsilon^{3\mathrm Q_1+20},\qquad
\mathrm T=\epsilon^{3\mathrm Q_1+26}=\epsilon^{60\cdot2^{L-1}-34}.
\label{eq:rg-formation-regime-closed-form}
\end{equation}
The regime is a low-temperature limit with doubly exponential energy
separation in the depth: \(\mathrm T=\epsilon^{86},\epsilon^{206},
\epsilon^{446}\) for \(L=2,3,4\), every ratio \(\mathrm J_\ell/\mathrm T\)
diverges, \(\mathrm J_{\ell+1}/\mathrm J_\ell=\epsilon^{\mathrm Q_{\ell+1}+20}\)
and \(\mathrm J_{\ell+1}^2/\mathrm J_\ell=\epsilon^{20}\) for \(\ell<L\), and
the relative thermal scale is \(\sqrt{\mathrm T}/\mathrm A=\epsilon^3\). There
is no common acceleration of every energy, no vanishing source rate, and no
decrease of \(\delta\). \status{ESTABLISHED}

The recursion comes from the early estimate of
\Cref{lem:formation-conditional-entropy}, which bounds the retained Fisher
information by the full budget \cref{eq:rg-formation-full-fisher-budget} and
hence the discarded loss on \(I_\ell\) by the term
\((\mathrm A^2/\mathrm T)\mathrm J_{\ell+1}^2/\mathrm J_\ell\) of
\cref{eq:rg-formation-early-entropy-bound}. The linear-Gaussian core shows
that this term is not sharp.

\propositionheading{Gaussian core of the discarded conditional loss}{prop:rg-formation-gaussian-core}
Take \(\delta=0\), \(\vartheta\equiv1\), and an active non-root normal law
\(\mathcal N(\bar{\mathbf x}(t),\mathrm TH_{\ell'}^{-1})\) at level
\(\ell'<\ell\), whose covariance is stationary for the Ornstein--Uhlenbeck
flow, so \(\bar{\mathbf x}(t)=e^{-H_{\ell'}t}\bar{\mathbf x}(0)\). Split the
coordinates into the discarded \(z\) and retained \(y\), with blocks
\(H_{zz}\), \(H_{zy}\), \(H_{yy}\) of \(H_{\ell'}\) and Schur complement
\(H_\ell=H_{yy}-H_{zy}^\top H_{zz}^{-1}H_{zy}\).

(a) \emph{Loss formula and upper bound.} Then
\begin{equation}
\begin{aligned}
\Delta_{\ell':\ell}(t)&=\frac1{2\mathrm T}\mathfrak v^\top H_{zz}^{-1}\mathfrak v,
&
\mathfrak v&=(H_{\ell'}\bar{\mathbf x})_z,\\
\dot{\mathfrak v}&=-H_{zz}\mathfrak v-H_{zy}(H_{\ell'}\bar{\mathbf x})_y,
&
(H_{\ell'}\bar{\mathbf x})_y&=H_\ell\bar{\mathbf x}_y+H_{zy}^\top H_{zz}^{-1}\mathfrak v,
\end{aligned}
\label{eq:rg-formation-gaussian-core-loss}
\end{equation}
and, for \(|\bar{\mathbf x}|\le\mathrm C_L\mathrm A\),
\(\Delta_{\ell':\ell}\le\mathrm C_L(\mathrm A^2/\mathrm T)
[\mathrm J_{\ell'+1}e^{-\mathrm c_L\mathrm J_\ell(t-t^{\mathrm{act}}_{\ell'})}
+\mathrm J_{\ell+1}^4/\mathrm J_\ell^3]\).

(b) \emph{Two-sided asymptotic across an early interface.} Let the strengths
\(\mathrm J_1\ge\cdots\ge\mathrm J_L=1\) depend on \(\epsilon\) with ratios
\(\varrho_s=\mathrm J_{s+1}/\mathrm J_s\to0\) for \(s<L\), put \(\varrho_0=\varrho_L=0\), and
let \(1\le\ell\le L-1\) and \(\ell'=\ell-1\). Write
\(\mathsf B_{sr}=B_s\mathsf Q_r\) for the blocks of the matrix \(\mathsf B\)
of \cref{eq:rg-formation-weighted-flag}, which vanish for \(r>s\) and are
square and invertible for \(r=s\) (\Cref{lem:formation-parent-score}), let
\(\mathfrak j_\ell=\lambda_{\min}(\mathsf B_{\ell\ell}^\top\mathsf B_{\ell\ell})\),
let \(\mathfrak Y_\ell=\mathsf B_{\ell+1,\ell}^\top\mathsf B_{\ell+1,\ell+1}\mathsf B_{\ell+1,\ell+1}^\top\)
be the forcing map of the interface, and for
\(\mathbf v\in\R^{10N_\triangle-1}\) let
\begin{equation}
\mathfrak y_\ell(\mathbf v)=\tfrac12\bigl|(\mathsf B_{\ell\ell}^\top\mathsf B_{\ell\ell})^{-3/2}
\mathfrak Y_\ell B_{\ell+1}P_\ell\mathbf v\bigr|^2,
\qquad
\mathfrak x_\ell(t)=\varrho_\ell^{-2}\exp\bigl[-\tfrac12\mathfrak j_\ell\mathrm J_\ell(t-t^{\mathrm{act}}_{\ell-1})\bigr].
\label{eq:rg-formation-core-coefficient}
\end{equation}
Then \(\mathsf B_{\ell+1,\ell}\) has full row rank for every choice of the
bridge endpoints, so \(\mathfrak Y_\ell\) is injective. There are
\(\mathrm c_0,\mathrm C>0\) depending only on \(L\) and the carrier such
that, for the fine process started from the reference non-root law of
\Cref{def:rg-formation-regime} with mean \(\mathrm A\mathbf v_0\) and
activated at the times \cref{eq:rg-formation-activation-times}, whenever
\(\varrho_{\ell-1}+\varrho_\ell+\varrho_{\ell+1}\le\mathrm c_0\), the active level \(\ell-1\)
satisfies, for \(t^{\mathrm{act}}_{\ell-1}\le t\le t^{\mathrm{act}}_\ell\),
\begin{equation}
\begin{aligned}
\Bigl|\frac{\mathrm T\mathrm J_\ell^3}{\mathrm A^2\mathrm J_{\ell+1}^4}
\Delta_{\ell-1:\ell}(t)-\mathfrak y_\ell\bigl(\bar{\mathbf x}(t)/\mathrm A\bigr)\Bigr|
&\le\mathrm C\bigl[\varrho_{\ell-1}+\varrho_\ell+\varrho_{\ell+1}+\mathfrak x_\ell(t)+\mathfrak x_\ell(t)^2\bigr],\\
|B_{\ell+1}P_\ell\bar{\mathbf x}(t)-\mathrm AB_{\ell+1}P_\ell\mathbf v_0|
&\le\mathrm C\mathrm J_{\ell+1}\mathrm At.
\end{aligned}
\label{eq:rg-formation-core-two-sided}
\end{equation}
If moreover \(k_L\log(1/\epsilon)\varrho_\ell\to0\) and
\(\tfrac12\mathfrak j_\ell k_L\log(1/\epsilon)(1-2\varrho_{\ell-1})-2\log(1/\varrho_\ell)\to\infty\),
then, uniformly on \(I_\ell\),
\begin{equation}
\Delta_{\ell-1:\ell}(t)
=\frac{\mathrm A^2}{\mathrm T}\frac{\mathrm J_{\ell+1}^4}{\mathrm J_\ell^3}
\bigl[\mathfrak y_\ell(\mathbf v_0)+o(1)\bigr],
\qquad
\mathfrak y_\ell(\mathbf v_0)
\ge\frac{\sigma_{\min}(\mathfrak Y_\ell)^2N_\triangle\mathrm m_L^2}{2^{\ell+3}\|\mathsf B_{\ell\ell}\|^6}>0,
\label{eq:rg-formation-core-window}
\end{equation}
where \(\sigma_{\min}\) denotes the least singular value.

(c) \emph{Separation criterion.} With power laws
\(\mathrm J_s=\epsilon^{-\mathrm Q_s}\), \(\mathrm Q_1>\cdots>\mathrm Q_L=0\),
with \(\mathrm A^2/\mathrm T=\epsilon^{-\mathrm Q_{\mathrm T}}\) for some
\(\mathrm Q_{\mathrm T}>0\), which is \(\mathrm Q_{\mathrm T}=6\) in \cref{eq:rg-formation-weakest-one-regime}, and with an entrance
constant \(k_L>4(\mathrm Q_\ell-\mathrm Q_{\ell+1})/\mathfrak j_\ell\) for
every \(\ell<L\), both \(\inf_{I_\ell}\Delta_{\ell-1:\ell}\) and
\(\sup_{I_\ell}\Delta_{\ell-1:\ell}\) are of exact order
\(\epsilon^{3\mathrm Q_\ell-4\mathrm Q_{\ell+1}-\mathrm Q_{\mathrm T}}\). The core loss is
\(O(\epsilon^{\mathrm Q_\Delta})\) on every early window, \(\mathrm Q_\Delta\ge0\), exactly when
\(3\mathrm Q_\ell-4\mathrm Q_{\ell+1}\ge\mathrm Q_{\mathrm T}+\mathrm Q_\Delta\) for every \(\ell\le L-1\), which
forces \(\mathrm Q_1\ge(\mathrm Q_{\mathrm T}+\mathrm Q_\Delta)[(4/3)^{L-1}-1]\), with
equality along the recursion
\(3\mathrm Q_\ell=4\mathrm Q_{\ell+1}+\mathrm Q_{\mathrm T}+\mathrm Q_\Delta\); it vanishes on every
early window exactly when \(3\mathrm Q_\ell-4\mathrm Q_{\ell+1}>\mathrm Q_{\mathrm T}\) for every
\(\ell\le L-1\). \status{ESTABLISHED}

The proof of (b) and (c) is in \cref{sec:rg-formation-core-proof}.

\noindent\emph{Proof of (a).}
Both conditional laws of \(z\) given \(y\) have precision
\(H_{zz}/\mathrm T\), and their means differ by
\(H_{zz}^{-1}(H_{\ell'}\bar{\mathbf x})_z\) independently of \(y\), which gives
the first formula. Differentiating
\(\mathfrak v=(H_{\ell'}\bar{\mathbf x})_z\) along
\(\dot{\bar{\mathbf x}}=-H_{\ell'}\bar{\mathbf x}\) and eliminating
\(\bar{\mathbf x}_z\) give the other two. Put \(\mathfrak u=H_{zz}^{-1/2}\mathfrak v\). Its equation is
\[
\dot{\mathfrak u}
=-\bigl(H_{zz}+H_{zz}^{-1/2}H_{zy}H_{zy}^\top H_{zz}^{-1/2}\bigr)\mathfrak u
-H_{zz}^{-1/2}H_{zy}H_\ell\bar{\mathbf x}_y.
\]
The matrix multiplying \(\mathfrak u\) is symmetric and bounded below by
\(\mathrm c_L\mathrm J_\ell I\). The retained loop block is orthogonal to
\(H_{zy}\) and contributes no forcing; on the future normal coordinates,
\cref{eq:rg-formation-current-future-blocks} bounds the forcing by
\(\mathrm C_L\mathrm A\mathrm J_{\ell+1}^2/\sqrt{\mathrm J_\ell}\).
At \(t_0=t^{\mathrm{act}}_{\ell'}\), the weighted flag gives
\(|\mathfrak u(t_0)|\le\mathrm C_L\mathrm A\sqrt{\mathrm J_{\ell'+1}}\).
Variation of constants therefore gives
\[
|\mathfrak u(t)|\le\mathrm C_L\mathrm A\left[
\sqrt{\mathrm J_{\ell'+1}}e^{-\mathrm c_L\mathrm J_\ell(t-t_0)}
+\frac{\mathrm J_{\ell+1}^2}{\mathrm J_\ell^{3/2}}\right].
\]
Squaring and using \(\Delta_{\ell':\ell}=|\mathfrak u|^2/(2\mathrm T)\)
proves the bound, with the decay constant adjusted. The forcing vanishes for
\(\ell=L\). \(\square\)

\openproblemheading{Sharpness of the doubly exponential regime}{open:rg-formation-regime-sharpness}
Decide whether the recursion \(\mathrm Q_\ell=2\mathrm Q_{\ell+1}+20\) is
required for \Cref{thm:finite-depth-adaptive-formation}, and whether its
constants and preparation neighborhoods can be made uniform in the depth.
What holds is the theorem at every finite \(L\) with constants depending on
\(L\), and, by \Cref{prop:rg-formation-gaussian-core}, that in the Gaussian
core the discarded loss on every early window is of exact order
\((\mathrm A^2/\mathrm T)\mathrm J_{\ell+1}^4/\mathrm J_\ell^3\), smaller than
the term of the general estimate by \((\mathrm J_{\ell+1}/\mathrm J_\ell)^2\),
so that the early interfaces need exactly
\(3\mathrm Q_\ell\ge4\mathrm Q_{\ell+1}+20\) for the core loss at the scale
\(\epsilon^{14}\), a separation still doubly exponential in \(L\) but of inner
base \(4/3\), with \(\mathrm Q_\ell\ge14\) for the term
\(\mathrm C_{\mathcal J}/\mathrm J_\ell\) of the jump remainder
\cref{eq:rg-formation-conditional-jump-remainder}. The doubling enters the
present proof through the full Fisher budget
\cref{eq:rg-formation-full-fisher-budget} in the early estimate of
\Cref{lem:formation-conditional-entropy} and through the history estimate of
\Cref{lem:formation-history-coupling}, which uses it at the top interface in
\cref{eq:rg-formation-averaged-environment-error}. The depth enters through
\(\mathrm Q_1\) in the regime and every entrance condition, through
\(\log|\mathcal J|=N\log2\) in the jump remainder, and through the Frobenius
constants of \cref{eq:rg-formation-weighted-flag}, whose squares grow linearly
in \(N_\triangle\). The problem is posed in the family of the theorem under
(H1), (H3), and (H4), with power-law strengths
\(\mathrm J_\ell=\epsilon^{-\mathrm Q_\ell}\), \(\mathrm Q_L=0\), satisfying
\(3\mathrm Q_\ell\ge4\mathrm Q_{\ell+1}+20\) and \(\mathrm Q_\ell\ge14\) for
\(\ell<L\), with the amplitude, temperature, horizon, and record coefficient
of \cref{eq:rg-formation-weakest-one-regime} formed from \(\mathrm Q_1\), and
with an entrance constant above the thresholds of the theorem and of
\Cref{prop:rg-formation-gaussian-core}. It is settled by conclusions (a) to
(e) for every schedule of this class, with constants and preparation
neighborhoods independent of \(L\), through an early estimate of the retained
Fisher information by the future-layer energy, a history-environment bound
without the top-interface recursion, and depth-uniform flag and jump
constants; or by an admitted preparation and a schedule of the class under
which one of the conclusions fails, or a lower bound showing that one of the
conclusions cannot hold with constants independent of \(L\).

\section{The finite-depth formation theorem}
\label{sec:rg-formation-theorem}
\label{sec:rg-formation-common-regime-proof}

Formation in this family is hierarchical relaxation of the prepared mean. The
non-root configuration starts near \(\mathrm A\mathbf v_0\), where every
residual has order \(\mathrm A\) and every bridge residual
\(r_e(P_{\ell'}\mathrm A\mathbf v_0)\) with \(\ell'<\ell(e)\) stays at least
\(\mathrm A\mathrm m_L/2\) away from zero. The weighted flag makes residuals
of level at most \(\ell\) relax at rate of order \(\mathrm J_\ell\), while
bridges of higher level feel only couplings of order \(\mathrm J_{\ell+1}\).
On \(I_\ell\) the first have decayed by \(\epsilon^{\mathrm c_Lk_L}\) relative
to \(\mathrm A\), whereas the second have moved by at most
\(\mathrm C_L\mathrm A\mathrm J_{\ell+1}t_\ell\to0\) and keep their margin, so
each window joins pairs of coherent components along the bridges of level
\(\ell\) and halves their number. The source record reads residuals at the
thermal resolution \(\sqrt{\mathrm T/\mathrm J_1}\), far below
the coherence scale \(\mathrm A\), so it is saturated there. Source jumps, the
record, and learning are therefore perturbations, and the content of the
theorem is that exact marked parents, activation from the complete law, and
own-law learning can be carried through this relaxation with separately
vanishing errors.

\theoremheading{Finite-depth formation under a supplied energy hierarchy}{thm:finite-depth-adaptive-formation}
Assume
(H1) the experiment of \Cref{def:rg-formation-carrier,def:rg-formation-experiment,def:rg-formation-laws,def:rg-formation-parents}
with the kinetics of \Cref{def:rg-formation-kinetics,def:rg-formation-learning},
at a fixed depth \(L\ge2\);
(H2) the regime \cref{eq:rg-formation-weakest-one-regime};
(H3) an initial law in the preparation neighborhood
\cref{eq:rg-formation-preparation-transport};
(H4) parents activated at the times \cref{eq:rg-formation-activation-times}.
There exist constants \(k_{\mathrm{geom},L},\mathrm c_{\mathrm{early},L},
\mathrm c_{\mathrm{joint},L},\mathrm c_{\mathrm{fast},L},\eta_L^{\max}>0\),
depending only on \(L\) and the fixed data and constructed respectively in
\Cref{lem:formation-coherence} (together with \(\eta_L^{\max}\), a fixed
fraction of the margin \(\mathrm m_L\)),
\Cref{lem:formation-conditional-entropy},
\Cref{lem:formation-terminal-joint-entropy}, and
\Cref{lem:formation-fast-mixing}, with the following property. For every
entrance constant \(k_L\) with
\begin{equation}
\begin{aligned}
k_L&>k_{\mathrm{geom},L},&
\tfrac12\mathrm c_{\mathrm{early},L}k_L&>\mathrm Q_1+20,\\
\mathrm c_{\mathrm{joint},L}k_L&>\mathrm Q_1+8,&
\mathrm c_{\mathrm{fast},L}k_L&>\mathrm Q_1/2+12,
\end{aligned}
\label{eq:rg-formation-common-entrance}
\end{equation}
and every tolerance \(0<\eta_L\le\eta_L^{\max}\), there exist
\(\epsilon_L,\mathrm c_L,\mathrm C_L>0\), depending on \(L\), the fixed data,
\(k_L\), and \(\eta_L\), such that for every \(0<\epsilon<\epsilon_L\)
conclusions (a) to (e) hold for the fine process and for the lineage of
activated parents. The lineage is said to admit all \(L\) successive
activations when the fine process and every activated parent have a unique
nonexplosive solution of \cref{eq:rg-formation-own-law-dynamics} through
\(t_*\), with positive smooth densities and parameter in \([1,21/10]\).

(a) \emph{Well-posedness and parameter.} The lineage admits all \(L\)
successive activations, and every process satisfies, in absolute time,
\begin{equation}
1\le\vartheta_\ell(t)\le21/10,\qquad
|\dot\vartheta_\ell(t)|\le\mathrm C_L\epsilon,\qquad
|\vartheta_\ell(t)-(2-e^{-\epsilon t})|\le\mathrm C_Lr_\epsilon,
\label{eq:rg-formation-parameter-bound}
\end{equation}
with \(r_\epsilon=O(\epsilon^{2\mathrm Q_1+20})\) of
\cref{eq:rg-formation-parameter-remainder}; the comparison curve solves the
calibration part of \cref{eq:rg-formation-own-law-dynamics}.

(b) \emph{Coherence.} For a level-\(\ell\) description at time \(t\), an
original edge \(e\) is coherent when
\(\sqrt{2\mathcal D_e^{\ell,\mu}}\le\eta_L\mathrm A\), with
\cref{eq:rg-formation-inherited-edge-statistic} evaluated at the current
retained state and parameter of that description, and a vertex is internally
coherent when the corresponding observable for \(\mathfrak d_i\) is at most
\(\eta_L\mathrm A\); the coherent components are the connected components of
the graph on the \(N\) vertices whose edges are the coherent edges. For each
of \(\mu_\vartheta=\boldsymbol\Pi_{\vartheta,\mathsf o}\) and
\(\mu_\vartheta=q^\star_\vartheta\), except on an event of probability at most
\begin{equation}
\mathrm C_L(1+t_*)\exp(-\mathrm c_L/\epsilon^6),
\label{eq:rg-formation-coherence-probability}
\end{equation}
the microscopic process and its sequential parent descriptions have the
successive coherent components
\begin{equation}
3\cdot2^{L-1}\longrightarrow2^{L-1}
\longrightarrow2^{L-2}\longrightarrow\cdots\longrightarrow1.
\label{eq:rg-formation-agent-counts}
\end{equation}
On \(I_\ell\) every edge of level at most \(\ell\) is coherent, every vertex is
internally coherent, and every bridge of level greater than \(\ell\) has
residual observable at least a fixed positive multiple of \(\mathrm A\), so the
count on \(I_\ell\) is \(2^{L-\ell}\); the single final component persists
through \(t_*\). Loop coordinates remain random and retained, with no
identity-holonomy constraint.

(c) \emph{Conditional loss and score.} On each window \(I_\ell\), at the
current law and parameter,
\begin{equation}
\Delta_{\ell-1:\ell}
 \le
 \begin{cases}
\mathrm C_L\epsilon^{14},&\ell<L,\\
\mathrm C_L\epsilon^2,&\ell=L,
\end{cases}
\qquad
\bigl|g_{\ell-1}(q,\vartheta)
 -g_\ell\bigl((C_{\ell\leftarrow\ell-1})_\#q,\vartheta\bigr)\bigr|
 \le
 \begin{cases}
\mathrm C_L\mathrm T\epsilon^7,&\ell<L,\\
\mathrm C_L\mathrm T\epsilon,&\ell=L,
\end{cases},
\label{eq:rg-formation-conditional-score-bounds}
\end{equation}
and \(\mathrm T\Delta_{\ell-1:\ell}\) is the discarded restricted free energy.

(d) \emph{Source and kernel movement.} Let \(\lambda^\ell_{i_0,+}(t)\) be the
distinguished receiver's source mass in a fine, parent, or sequential process.
From each activation time at which that process is continued,
\begin{equation}
\lambda^\ell_{i_0,+}(t_*)-
  \lambda^\ell_{i_0,+}(t^{\mathrm{act}}_{\ell'})>\tfrac12,
\label{eq:rg-formation-source-movement}
\end{equation}
so the binary recognition source row changes by total variation greater than
one half. Between \(\vartheta=1\) and the final parameter, the first hidden
physical conditional and the first hidden optimizing conditional each change
by relative entropy greater than \(1/200\), at every fixed retained input and
assignment. Both changes are driven by the learned parameter through
\cref{eq:rg-formation-source-priors} and \(\vartheta H\), the record moving
the distinguished log odds by at most \(\delta\), and both take place after
formation, since \(\vartheta\le1+\mathrm C_L\epsilon\log(1/\epsilon)\) until
the last window closes.

(e) \emph{Histories and paths.} For any source process and its inherited
parent begun from a common marked activation sample at
\(t^{\mathrm{act}}_\ell\), there is a deterministic path
\(\bar{\mathbf x}_\ell:[t^{\mathrm{act}}_\ell,t_*]\to W_\ell\) with
\begin{equation}
\Pr\Bigl\{\sup_{t^{\mathrm{act}}_\ell\le t\le t_*}|P_\ell\mathbf X_t-\bar{\mathbf x}_\ell(t)|
  >\mathrm C_L\mathrm A\epsilon^2\Bigr\}
+\Pr\Bigl\{\sup_{t^{\mathrm{act}}_\ell\le t\le t_*}|\overline{\mathbf X}_t-\bar{\mathbf x}_\ell(t)|
  >\mathrm C_L\mathrm A\epsilon^2\Bigr\}
\le\mathrm C_L(1+t_*)e^{-\mathrm c_L/\epsilon^2}.
\label{eq:rg-formation-common-path-concentration}
\end{equation}
There is a coupling preserving the activation sample under which the first
bound below holds, and the second holds under every coupling of the two
processes:
\begin{equation}
\begin{aligned}
\Pr\bigl\{S_t\ne\overline S_t
       \text{ for some }t\in[t^{\mathrm{act}}_\ell,t_*]\bigr\}
 &\le\mathrm C_L\epsilon^6,\\
\Pr\Bigl\{\sup_{t^{\mathrm{act}}_\ell\le t\le t_*}
       |P_\ell\mathbf X_t-\overline{\mathbf X}_t|>\mathrm C_L\mathrm A\epsilon^2\Bigr\}
 &\le\mathrm C_L(1+t_*)e^{-\mathrm c_L/\epsilon^2}.
\end{aligned}
\label{eq:rg-formation-two-path-bounds}
\end{equation}
The resolution
\(\mathrm C_L\mathrm A\epsilon^2=\mathrm C_L\sqrt{\mathrm T}/\epsilon\) is
coarser than the thermal scale, and a finite number of successive
replacements changes only the constants. \status{ESTABLISHED}

\noindent\emph{Proof.}
The exact statements are
\Cref{prop:rg-formation-normalized-experiment,prop:rg-formation-laws,prop:rg-formation-parent-potential,prop:rg-formation-rate-composition},
and \Cref{tab:rg-formation-lemma-dependencies} lists the lemmas of
\cref{app:finite-depth-formation-proofs} with their inputs and exponents.
Conclusion (a) is
\Cref{lem:formation-prepared-continuation,lem:formation-parameter-comparison},
and (b) is \Cref{lem:formation-coherence}, which uses neither conditional
equilibration nor history averaging. For (c),
\Cref{lem:formation-conditional-entropy} gives the early bound from
\(\mathrm J_{\ell+1}^2/\mathrm J_\ell=\epsilon^{20}\) and the second entrance
condition, and \Cref{lem:formation-terminal-joint-entropy} the terminal bound
by joint diffusion-plus-jump relaxation; conditional relative entropy is
bounded by the full one through the chain rule, and the complete hidden score
bounds the covector difference by
\(\mathrm C_L\mathrm T(\sqrt\Delta+\Delta)\). Conclusion (d) is
\Cref{lem:formation-source-envelope}. For (e),
\Cref{lem:formation-history-coupling} couples the assignment histories and
\Cref{lem:formation-continuous-path} gives
\cref{eq:rg-formation-common-path-concentration}; each of its events concerns
the law of one process, so the union bound gives the second line of
\cref{eq:rg-formation-two-path-bounds} under every coupling. The four
entrance conditions are equivalent to the single threshold
\begin{equation}
k_L>\max\left\{k_{\mathrm{geom},L},\
\frac{2(\mathrm Q_1+20)}{\mathrm c_{\mathrm{early},L}},\
\frac{\mathrm Q_1+8}{\mathrm c_{\mathrm{joint},L}},\
\frac{\mathrm Q_1/2+12}{\mathrm c_{\mathrm{fast},L}}\right\},
\label{eq:rg-formation-joined-entrance-proof}
\end{equation}
whose terms supply, in order, the geometric windows, the early entrance,
terminal joint relaxation, and the source-history estimate; every lemma uses
the same preparation, regime, and entrance choice. The final window has order
\(k_L\log(1/\epsilon)\) and lies before \(t_*-10=2/\epsilon-10\) once
\(\epsilon\) is small; this, the ordering of the windows, and the smallness
conditions inside the lemmas fix \(\epsilon_L\) after \(k_L\) and \(\eta_L\).
\(\square\)

\remarkheading{Saturated record, calibrated learning, and thermally negligible loops}{rem:rg-formation-saturation}
Uniformly in the configuration, the assignment, and the parameter,
\begin{equation}
\begin{aligned}
|\log\beta_S(\mathbf x,\vartheta)-\log\pi_S(\vartheta)|&\le B_0=10^{-4},\\
|H^{-1/2}\mathrm T\nabla_{\mathbf x}\phi_S(\mathbf x)|&\le B_*\sqrt{\mathrm T},
\qquad
B_*=\delta\sqrt{2N}=10^{-4}\sqrt{2/N}\le5.8\times10^{-5},
\end{aligned}
\label{eq:rg-formation-source-saturation}
\end{equation}
with \(H^{-1/2}\) acting on \(\mathcal K^\perp\), by
\cref{eq:rg-formation-full-source-row} with \(0\le\phi_S\le B_0\) and by
\cref{eq:rg-formation-original-relative-bounds}. At a residual of order
\(\mathrm A\) the argument of \(\chi\) has order
\(\epsilon^{-(3+\mathrm Q_1/2)}\), so the record equals its saturated value up
to an exponentially small error. The proof of \Cref{lem:formation-coherence}
uses the record and the assignments only through \(0\le\phi_S\le B_0\) and the
pathwise force bound \cref{eq:rg-formation-path-force-bound}, valid for every
adapted assignment path, and the parameter only through
\(1\le\vartheta\le21/10\), so with \(\delta=0\) and \(\vartheta\equiv1\) it
gives the same components. By
\Cref{lem:formation-prepared-continuation,lem:formation-parameter-comparison}
every law-dependent term of the covector is at most \(\mathrm C_Lr_\epsilon\),
which gives the parameter bounds of \Cref{thm:finite-depth-adaptive-formation}.
The loop energy \(\mathrm J_1|h|^2\) belongs to the strongest term, so
\(\mathrm J_1\E|h|^2\le\mathrm C_L\mathrm T\) by
\cref{eq:rg-formation-actual-moments}, and the loops have root-mean-square size
at most \(\mathrm C_L\sqrt{\mathrm T/\mathrm J_1}
=\mathrm C_L\mathrm A\epsilon^{3+\mathrm Q_1/2}\). \status{ESTABLISHED}

\section{A two-triangle instance and the remaining scope}
\label{sec:rg-formation-example-scope}

\remarkheading{The two-triangle instance}{rem:rg-formation-two-triangle}
For \(L=2\), take triangles \((1,2,3)\) and \((4,5,6)\) with the single bridge
\((3,4)\). The retained carriers are
\begin{equation}
\R^{19}\times\mathcal J
\longrightarrow\R^9\times\mathcal J
\longrightarrow\R^8\times\mathcal J,
\qquad |\mathcal J|=64,
\label{eq:rg-formation-two-triangle-carrier}
\end{equation}
the coherent component counts are \(6\to2\to1\), both loops survive both
maps, and the regime becomes
\begin{equation}
\mathrm J_1=\epsilon^{-20},\qquad \mathrm J_2=1,\qquad
\mathrm A=\epsilon^{40},\qquad \mathrm T=\epsilon^{86}.
\label{eq:rg-formation-two-triangle-regime}
\end{equation}
The first window has order \(\epsilon^{20}\log(1/\epsilon)\), the second
\(\log(1/\epsilon)\), and the source and model changes occur over the longer
horizon \(2/\epsilon\). The source-dependent term is not quadratic: along a
direction with selected internal residual slope \(\lambda\ne0\),
\(\chi(\lambda t)=\lambda^2t^2/2-\lambda^4t^4/8+O(t^6)\), and every assignment
selects an internal edge with nonzero first-normal projection, so the hidden
log density has a nonzero fourth derivative in some first-normal direction;
exact marginalization keeps the full conditional integral and the negative
covariance term of its Hessian, and neither first hidden conditional is a
Gaussian. \status{ESTABLISHED}

\openproblemheading{A family in which learning or holonomy selects the components}{open:rg-formation-adaptive-family}
Construct a normalized family at \(c_*\) in which parameter learning or loop
holonomy changes, at the level of the law, which components form, with exact
marked parents and activation from the complete law. What holds is
\Cref{thm:finite-depth-adaptive-formation}, whose components are fixed by the
supplied contrasts and internal source support, and
\Cref{prop:rg-formation-endogenous-cycle}, whose pairs are seeded by the
random initial assignment without learning, so that in it the source record
selects only along sampled paths; \Cref{prop:phil-seeded-contrast-formation}
shows that its seeded bound is unchanged by learning and that from an initial
law invariant under the cycle automorphisms learning leaves the rows uniform,
so a law-level selection needs an initial law without that invariance. Such a
family needs at least one of a source record standardized at the coherence
scale, for example through \(\chi(r/(\eta_L\mathrm A))\), model means and
model links both slow relative to the belief mobility, since a fast link or
fast model means remove the model disagreement on the belief clock whatever
the other rate (\Cref{prop:obs-slow-model-link-dichotomy}), loop coordinates
of order \(\mathrm A\) or nonabelian links, or interacting model-channel
records, together with new fast-marginal and history estimates; sources on
different bridge levels and source learning that creates the contrasts belong
to the same problem. It is posed on the four-cycle of
\Cref{prop:phil-seeded-contrast-formation}, with its temperature, learning
mobility, and horizon, for Gaussian initial laws
\(\mathcal N(\mathsf m_0,I_{12})\) with \(\mathsf m_0\) outside a
Lebesgue-null set, and for parents of the type of part (b) of
\Cref{prop:rg-formation-endogenous-cycle} activated from the recognition law
of the episode prefix; the law of the formed block means the conditional law,
given contrast formation, of the first block for which it occurs, the
conditioning event having probability bounded below in every arm by part (b)
of \Cref{prop:phil-seeded-contrast-formation}, so that a difference in
whether formation occurs is not selection. It is settled by such parents with
separately vanishing conditional and dynamical errors together with a
total-variation difference, bounded below as \(\epsilon\to0\), between the
conditional laws of the formed block in the learned and in the frozen-row
arm, or between those of two initial means that differ only in their loop
coordinate \(\sum_i\omega_i\), the amplified datum being the initial law or
its loop coordinate and the other arm its null; or by a proof that, for
almost every \(\mathsf m_0\), each of these two pairs of arms gives
conditional laws of the formed block whose difference vanishes as
\(\epsilon\to0\).

A separate slower-mobility regime is resolved by
\cref{prop:phil-early-block-null}: keeping the same contrast event and
horizon but taking \(\epsilon_{\mathrm L}=o(\mathrm T^{1/4})\), the
learned and frozen conditional first-block laws have vanishing total
variation for every fixed Gaussian mean. The original coupled mobility
\(\epsilon_{\mathrm L}=\mathrm T^{1/8}\) and the loop-coordinate
comparison remain outside that estimate. The late coherent-partition
contrast in \cref{prop:phil-selected-pair-persistence} concerns restored
equilibrium coherence and does not close this first-block problem.
For the late coherence event on the joint slow schedule,
\cref{prop:phil-actual-conditioned-parent,prop:phil-actual-parent-dynamics}
establish actual conditioned activation of the existing parent with
separately vanishing conditional-information and dynamical errors.
The activation event is a completed coherence prefix after selection.
It is not the first-contrast event in the original regime posed above.
For the seeded sub-event of that original contrast, at every mobility
\(\epsilon_{\mathrm L}\in[0,1]\) including \(\mathrm T^{1/8}\),
\cref{prop:phil-early-seeded-parent,prop:phil-early-parent-dynamics} give
separately vanishing conditional and dynamical errors when the parent is
activated after hidden thermalization, at a deterministic activation time
\(\theta\) with \(\theta-\log(1/\epsilon)\to\infty\) and
\(\epsilon^4\theta\to0\), on following windows short compared with
\(\mathrm T^{-1}\), while
\cref{prop:phil-early-first-completion} shows that activation at the first
qualifying completion has diverging conditional error, also under
conditioning on the whole contrast event. The learned and frozen first-block
comparison at \(\epsilon_{\mathrm L}=\mathrm T^{1/8}\), the loop-coordinate
comparison, and conditioning on contrast histories outside the seeds remain
open.

\section{An endogenous four-cycle witness}
\label{sec:rg-formation-endogenous-cycle}

\definitionheading{Homogeneous four-cycle experiment}{def:rg-formation-cycle-experiment}
Index four agents by \(i\in\mathbb Z/4\) and orient edge \(e_i\) from \(i\) to
\(i+1\). A configuration is \(\mathbf x=(\bar k,\bar m,\omega)\in\R^{12}\), with
\(q_i^b=\mathcal N(\bar k_i,1)\), \(q_i^m=\mathcal N(\bar m_i,\sigma^2_{\mathrm{ev}})\),
\(\sigma^2_{\mathrm{ev}}=1/2\), four independent link translations
\(\omega_i\), residuals \(r^b_i=\bar k_{i+1}-\bar k_i-\omega_i\) and
\(r^m_i=\bar m_{i+1}-\bar m_i-\omega_i\), and the edge cost
\begin{equation}
\mathcal D_{e_i}(\mathbf x)=\tfrac12\bigl[(r^b_i)^2+(r^m_i)^2\bigr]
=\KL\bigl(q^b_{i+1}\Vert(\Theta_{e_i})_\#q^b_i\bigr)
+\sigma^2_{\mathrm{ev}}\KL\bigl(q^m_{i+1}\Vert(\Theta_{e_i})_\#q^m_i\bigr),
\label{eq:rg-formation-cycle-cost}
\end{equation}
which is the same for both orientations. Every receiver has both neighbors as
sources, \(S\in\mathcal J=\prod_i\{i-1,i+1\}\) with \(|\mathcal J|=16\) and
uniform prior \(\pi_S=1/16\), and \(V_S=\sum_i\mathcal D_{e(i,S_i)}\), where
\(e(i,j)\) joins \(i\) and \(j\). The configuration prior is
\(\mathcal N(0,I_{12})\); given \((\mathbf x,S)\), independent binary records
have \(\Pr(E_i=1\mid\mathbf x,S)=\tfrac12\exp[-\mathcal D_{e(i,S_i)}(\mathbf x)/\mathrm T]\),
and \(E=\mathbf 1\) is observed; the physical evaluator and recognition
coupling are those of \Cref{def:rg-formation-experiment}. With
\(\mathfrak D\mathbf x=\bar k-\bar m\) and
\(R_{\mathrm{ev}}=2\mathfrak D^\top\mathfrak D\), the restricted optimizer is
\(q^\star(\mathbf x,S)\propto\pi_Se^{-U_S(\mathbf x)/\mathrm T}\) with
\(U_S=V_S+\tfrac{\mathrm T}2|\mathbf x|^2+\tfrac{\mathrm T}2\mathbf x^\top R_{\mathrm{ev}}\mathbf x\),
and the physical posterior omits the last term. The kinetics are identity
mobility with diffusion covariance \(2\mathrm TI dt\) and one rate-one clock
at which the complete assignment is redrawn from the product of the receiver
posterior rows \(\beta_S(\mathbf x)\), self-proposals allowed; no generative
parameter is learned. The two perfect-matching assignments are
\(S^{(0)}=(1,0,3,2)\), which selects \(e_0,e_2\) from both ends, and
\(S^{(1)}=(3,2,1,0)\). For \(0<\epsilon\le1/100\) put
\begin{equation}
\mathrm T=\epsilon^8,\quad t_\epsilon=2\log(1/\epsilon),\quad
t^\circ_*=t_\epsilon+1,\quad \varrho_\circ=\tfrac1{100},\quad
\sigma_\circ=\varrho_\circ/\sqrt{48},\quad \mathbf x^\circ=(0,0,\mathbf 1_4).
\label{eq:rg-formation-cycle-regime}
\end{equation}
The preparation class consists of the laws
\(\frac1{16}\mathfrak p_\circ(\mathbf x)(1+\mathfrak z(\mathbf x)) d\mathbf x\)
on \((\mathbf x,S)\), where \(\mathfrak p_\circ\) is the equal mixture of
\(\mathcal N(\mathbf x^\circ,\sigma_\circ^2I)\) and
\(\mathcal N(-\mathbf x^\circ,\sigma_\circ^2I)\) and \(\mathfrak z\) is smooth,
invariant under the cycle automorphisms (with link signs transformed under
reflections), of \(\mathfrak p_\circ\)-mean zero, and with
\(\|\mathfrak z\|_\infty<1/2\).

The preparation names no block: its assignment is uniform and independent of
the configuration, reflections exchange its two centers, which carry opposite
loop orientations, and all initial edge mismatch second moments agree. Under
the source gating an edge that neither endpoint selects carries no
interaction energy, although every edge stays in the palettes of both its endpoints and every
attention weight stays positive.

\propositionheading{Endogenous pairing on a homogeneous four-cycle}{prop:rg-formation-endogenous-cycle}
In \Cref{def:rg-formation-cycle-experiment}, let \(0<\epsilon\le1/100\) and let
the initial law belong to the preparation class. The joint is normalized, the
scaled negative ELBO is
\(\Fenergy_{\mathrm T}(q)=\mathrm T\KL(q\Vert q^\star)+\mathrm B_{\mathrm T}
-\mathrm T\log Z_{\mathrm{opt}}\) with
\(\mathrm B_{\mathrm T}=\mathrm T[6\log(2\pi)+4\log2]\), and the recognition law
of the kinetics has nonincreasing \(\Fenergy_{\mathrm T}\) and is invariant
under the cycle automorphisms at every time. Put
\begin{equation}
p_\epsilon=73728 \epsilon^4t^\circ_*+4t^\circ_*e^{-1/\mathrm T}<\tfrac12,
\qquad
\eta_\epsilon=\tfrac92\varrho_\circ^2\epsilon^8+\frac{3}{4(1-e^{-8})}\epsilon^{32},
\qquad
\bar\eta_\epsilon=\frac{\eta_\epsilon+\mathcal S_2(p_\epsilon)}{1-p_\epsilon},
\label{eq:rg-formation-cycle-constants}
\end{equation}
with \(\mathcal S_2(p)=-p\log p-(1-p)\log(1-p)\). Then:

(a) \emph{Episode.} With probability at least
\(p_0=3e^{-1/4}/512\), the complete assignment equals one perfect matching
throughout \([0,t^\circ_*]\); at time zero the internal belief transported
relative entropy of each matched pair exceeds \(1/4\), so that the contrast is
entered during the episode and is not present at its start; on
\([t_\epsilon,t^\circ_*]\) both internal belief and model transported relative
entropies of each matched pair are at most \(\epsilon^4\); and on
\([1/4,t^\circ_*]\) every exterior belief relative entropy exceeds \(1/2\), the
model ones being larger. This is a coherence episode held by the two attended
edge costs: on the held matching sheet the normal restoring rates of the
kinetics are \(8+\mathrm T\) and \(4+5\mathrm T\), those of the physical law
\(8+\mathrm T\) and \(4+\mathrm T\), so at the residual resolution
\(\epsilon^2\) the least restoring rate times the squared resolution, divided
by the temperature, is \(4\epsilon^{-4}\), and the sheet is selected by the
posterior proposals on the event. The kinetics leave \(q^\star\) invariant,
and started from \(q^\star\) they realize the same event with probability at
most \(4e^{-1/(4\mathrm T)}\), so the episode exceeds its stationary
probability by at least \(p_0-4e^{-1/(4\mathrm T)}>p_0/2\).

(b) \emph{Parent channel.} For each matching \(M\), let \(Z_M\) be the span of
the three orthonormal directions
\((n^+_{e}-n^+_{f})/\sqrt2\), \(n^-_{e}\), \(n^-_{f}\), where \(e,f\) are its
two edges, \(n^\pm_e\) the unit normal covectors of
\cref{eq:rg-formation-cycle-normals}, and \(W_M=Z_M^\perp\). The channel
\(C(\mathbf x,S)=(S,P_{W_S}\mathbf x)\), with \(W_S=W_M\) on the two matchings
and \(W_S=\R^{12}\) otherwise, extended by the identity on observations and
evaluator samples, keeps every assignment, the loop
\(\sum_i\omega_i\), and a fast common-normal coordinate, and pushes the joint,
the physical posterior, the restricted optimizer, and the recognition law to
normalized parent laws.

(c) \emph{Discarded conditional free energy.} For every
\(t\in[t_\epsilon,t^\circ_*]\), the recognition law \(q_t\) of the episode
has mean discarded optimizing conditional relative entropy at most
\(\bar\eta_\epsilon\) and mean discarded physical conditional relative entropy,
with the evaluator samples retained, at most \(\bar\eta_\epsilon+\mathrm T\).

(d) \emph{Parent dynamics.} The inherited parent of \(C\), activated at
\(t_\epsilon\) from the complete recognition law of the episode prefix, has
retained path law on \([t_\epsilon,t^\circ_*]\), including the complete
assignment path, within total variation
\(2p_\epsilon/(1-p_\epsilon)+4\sqrt{8/3} e^{-1/\mathrm T}\) of the retained
native path conditioned on the same prefix.

The estimates hold under the stated source gating, random matching seed,
symmetric circulated preparation, temperature \(\epsilon^8\), Euclidean
kinetic metric, duration-one episode, and fixed generative parameters.
\status{ESTABLISHED}

The proof is in \cref{sec:rg-formation-cycle-proof}. On the center path
\cref{eq:rg-formation-cycle-center}, which holds \(S^{(0)}\) from
\(\mathbf x^\circ\), where every residual equals \(-1\), the selected
residuals decay as \(-e^{-(8+\mathrm T)t}\) while the exterior ones grow
toward \(-3/2\), so a source-cost gap greater than one develops by time
\(1/4\), and the posterior proposals then keep the matching except with
probability at most \(4t^\circ_*e^{-1/\mathrm T}\). Since the full
recognition law is invariant, the pairing is an event of sampled paths seeded
by the initial assignment, not a block chosen by the law.

The witness raises the question whether its mechanism, an edge that carries
energy only while an endpoint attends to it, produces a stationary law that
favors coherence beyond one level, or a second coalescence of the pairs it
forms. The next statement answers the stationary part negatively in the
class to which the four-cycle belongs, the family of
\Cref{def:rg-gated-translation-family} on an arbitrary finite graph; the
second coalescence is posed for models outside that family in
\Cref{open:rg-formation-general}. Under exclusive gating the source label
enters only the weight of a sheet and never its stiffness, so every sheet
restores its attended residual directions at rates of order one and leaves
every other direction, including every direction that translates one
coherent block relative to another, to the thermal reference at rates of
order \(\mathrm T\). This has a spectral and a stationary consequence.

\definitionheading{Gated translation family on a graph}{def:rg-gated-translation-family}
Let \((V,E)\) be a finite simple graph without isolated vertices, with
\(N=|V|\) agents and matching number \(\mathfrak n\), the largest number of
pairwise disjoint edges. Orient each edge \(e\in E\) once as \(e:j\to i\); the
interaction multigraph at \(c_*\) consists of the copies \(e\) and \(\bar e\)
of the edges of \((V,E)\), and the source set \(\mathcal N_i\) of receiver \(i\)
is its set of neighbors in \((V,E)\), with \(e(i,j)\) the edge joining \(i\) and
\(j\). The agents, links, laws, evaluator, and recognition coupling are those
of \Cref{def:rg-formation-cycle-experiment} on this graph: a configuration is
\(\mathbf x=(\bar k,\bar m,\omega)\in\R^{2N+|E|}\) with one link translation
\(\omega_e\) per edge, \(q_i^b=\mathcal N(\bar k_i,1)\),
\(q_i^m=\mathcal N(\bar m_i,1/2)\), and edge \(e:j\to i\) carries the residual
pair and the cost
\begin{equation}
\mathbf a_e\mathbf x=(r^b_e,r^m_e)
=(\bar k_i-\bar k_j-\omega_e,\ \bar m_i-\bar m_j-\omega_e),
\qquad
\mathcal D_e(\mathbf x)=\tfrac12|\mathbf a_e\mathbf x|^2,
\label{eq:rg-gated-residual-pair}
\end{equation}
so \(\mathbf a_e:\R^{2N+|E|}\to\R^2\) is linear and invariant under passive
frame changes as in \Cref{lem:rg-formation-carrier}, and \(\mathcal D_e\) is
the transported relative entropy \cref{eq:rg-formation-cycle-cost}. Each
receiver carries one source label of \Cref{def:obs-attention-label-model}
with source set \(\mathcal N_i\), so a complete assignment is
\(S\in\mathcal J=\prod_i\mathcal N_i\); the labels are exclusive, the prior is
the product \(\pi_S=\prod_i\pi_{iS_i}\) with every \(\pi_{ij}>0\) independent
of \(\mathrm T\), and given \((\mathbf x,S)\) independent binary records have
\(\Pr(E_i=1\mid\mathbf x,S)=\tfrac12\exp[-\mathcal D_{e(i,S_i)}(\mathbf x)/\mathrm T]\),
with \(E=\mathbf 1\) observed. The attended edge set, sheet energy, and sheet
stiffness of \(S\) are
\begin{equation}
\operatorname{att}(S)=\{e(i,S_i):i\in V\},\qquad
V_S=\sum_i\mathcal D_{e(i,S_i)}=\tfrac12\mathbf x^\top H_S\mathbf x,\qquad
H_S=\sum_i\mathbf a_{e(i,S_i)}^\top\mathbf a_{e(i,S_i)}.
\label{eq:rg-gated-sheet-stiffness}
\end{equation}
The configuration prior is \(\mathcal N(0,I_{2N+|E|})\), and with
\(R_{\mathrm{ev}}=2\mathfrak D^\top\mathfrak D\) the restricted optimizer and
the physical posterior are the laws
\begin{equation}
q_\Omega(\mathbf x,S)=\frac{\pi_S}{\sum_{S'}\pi_{S'}Z_{S'}}
\exp\Bigl[-\frac{V_S(\mathbf x)}{\mathrm T}-\tfrac12\mathbf x^\top\Omega\mathbf x\Bigr],
\qquad
Z_S=\int\exp\Bigl[-\frac{V_S(\mathbf x)}{\mathrm T}-\tfrac12\mathbf x^\top\Omega\mathbf x\Bigr]d\mathbf x,
\label{eq:rg-gated-sheet-law}
\end{equation}
with the reference precision \(\Omega=I+R_{\mathrm{ev}}\) for the optimizer
and \(\Omega=I\) for the physical posterior; \(Z_S\) is the sheet weight. For
either \(\Omega\) the kinetics have identity mobility, diffusion covariance
\(2\mathrm TI dt\), drift \(-\nabla_{\mathbf x}U_S\) with
\(U_S=V_S+\tfrac{\mathrm T}2\mathbf x^\top\Omega\mathbf x\), and one rate-one
clock at which the complete assignment is redrawn from the product of the
receiver posterior rows
\(\beta_i(j\mid\mathbf x)=\pi_{ij}e^{-\mathcal D_{e(i,j)}(\mathbf x)/\mathrm T}
\big/\allowbreak\sum_{j'\in\mathcal N_i}\pi_{ij'}e^{-\mathcal D_{e(i,j')}(\mathbf x)/\mathrm T}\),
self-proposals allowed; no generative parameter is learned. For an edge set
\(F\subseteq E\), \(\mathbf a_F\) stacks the maps \(\mathbf a_e\), \(e\in F\),
\(W_F=\ker\mathbf a_F\), \(\mathrm s_F\) is the least positive singular value
of \(\mathbf a_F\), and \(\operatorname{cc}(F)\) is the number of connected
components of the graph \((V,F)\), isolated vertices included. At resolution
\(\mathrm a>0\) the coherent edges of a configuration are
\(\operatorname{Coh}_{\mathrm a}(\mathbf x)=\{e\in E:|\mathbf a_e\mathbf x|\le\mathrm a\}\),
and \(\mathfrak P_{\mathrm a}(\mathbf x)\) is the partition of \(V\) into the
vertex sets of the components of \((V,\operatorname{Coh}_{\mathrm a}(\mathbf x))\),
the coherent blocks.

The four-cycle of \Cref{def:rg-formation-cycle-experiment} is the member with
\((V,E)\) the four-cycle and the uniform prior, and its kinetics are those with
\(\Omega=I+R_{\mathrm{ev}}\). A minimum edge cover of \((V,E)\) is an edge set
covering every vertex with the fewest edges; the proposition identifies these
covers with the spanning forests of stars with \(\mathfrak n\) components.

\propositionheading{Depth one under exclusive source gating}{prop:rg-gated-depth-one}
In \Cref{def:rg-gated-translation-family}, for either reference precision
\(\Omega\), let \(\mathrm k_S=\operatorname{rank}H_S\), let \(\mathrm h_S\) be
the least positive eigenvalue of \(H_S\), \(\mathrm h_*=\min_S\mathrm h_S\),
and \(\mathrm z_S=[\operatorname{pdet}H_S\cdot\det(\Omega|_{W_{\operatorname{att}(S)}})]^{-1/2}\),
where \(\operatorname{pdet}\) is the product of the positive eigenvalues and
\(\Omega|_W=\mathsf Q_W^\top\Omega\mathsf Q_W\) for a matrix \(\mathsf Q_W\)
whose columns form an orthonormal basis of \(W\); let \(\mathcal M\) be the set
of assignments of minimum rank \(\mathrm k_S\).

(a) \emph{Two restoring bands.} For every \(S\) and \(\mathrm T>0\), the
eigenvalues \(\lambda_1\le\cdots\le\lambda_{2N+|E|}\) of \(H_S+\mathrm T\Omega\)
satisfy \(\lambda_k\in[\mathrm T\|\Omega^{-1}\|^{-1},\mathrm T\|\Omega\|]\) for
\(k\le2N+|E|-\mathrm k_S\) and
\(\lambda_k\in[\mathrm h_S,\|H_S\|+\mathrm T\|\Omega\|]\) for
\(k>2N+|E|-\mathrm k_S\). The range of \(H_S\) is
\(W_{\operatorname{att}(S)}^\perp\), the span of the attended residual
covectors, and the prior enters \(q_\Omega\) only through the sheet weights
\(\pi_SZ_S\), never through \(H_S\).

(b) \emph{Gating wells.} With
\[
V_\pi(\mathbf x)=-\mathrm T\sum_i\log\sum_{j\in\mathcal N_i}\pi_{ij}e^{-\mathcal D_{e(i,j)}(\mathbf x)/\mathrm T},
\qquad
V_{\min}(\mathbf x)=\sum_i\min_{j\in\mathcal N_i}\mathcal D_{e(i,j)}(\mathbf x),
\]
one has
\begin{equation}
0\le V_\pi(\mathbf x)-V_{\min}(\mathbf x)
\le\mathrm T\sum_i\log\frac1{\min_{j\in\mathcal N_i}\pi_{ij}}
\qquad\text{for every }\mathbf x,
\label{eq:rg-gated-well-depth}
\end{equation}
the configuration marginal of \(q_\Omega\) has density proportional to
\(\exp[-V_\pi(\mathbf x)/\mathrm T-\tfrac12\mathbf x^\top\Omega\mathbf x]\),
and if every edge of an edge set \(F\) covering \(V\) has zero residual at
\(\mathbf x\), then \(V_{\min}\) vanishes at \(\mathbf x\) and at every
configuration obtained from it by translating the belief and model means of
the agents of one component of \((V,F)\) by a common amount with the links
unchanged.

(c) \emph{Rank, minimum edge covers, and sheet weights.} For every
\(F\subseteq E\), \(\operatorname{rank}\mathbf a_F=|F|+N-\operatorname{cc}(F)\).
Hence \(\mathrm k_S=|\operatorname{att}(S)|+N-\operatorname{cc}(\operatorname{att}(S))\ge2(N-\mathfrak n)\),
with equality exactly when \(\operatorname{att}(S)\) is a minimum edge cover
of \((V,E)\); the minimum edge covers are the spanning forests of stars with
\(\mathfrak n\) components, every one of them is attended by some
assignment, and an edge outside \(\operatorname{att}(S)\) joining two of its
components raises the rank by two. The sheet weight satisfies
\begin{equation}
\Bigl(1+\frac{\mathrm T\|\Omega\|}{\mathrm h_S}\Bigr)^{-\mathrm k_S/2}
\le\frac{Z_S}{(2\pi)^{(2N+|E|)/2}\mathrm z_S\mathrm T^{\mathrm k_S/2}}\le1,
\label{eq:rg-gated-sheet-weight}
\end{equation}
and consequently, for \(0<\mathrm T\le1\) and every integer \(m\ge1\),
\begin{equation}
q_\Omega\{S:\operatorname{cc}(\operatorname{att}(S))\le\mathfrak n-m\}\le\mathrm K_m\mathrm T^m,
\qquad
\mathrm K_m=\Bigl(1+\frac{\|\Omega\|}{\mathrm h_*}\Bigr)^{N-\mathfrak n}
\frac{\sum_{S:\operatorname{cc}(\operatorname{att}(S))\le\mathfrak n-m}\pi_S\mathrm z_S}
{\sum_{S\in\mathcal M}\pi_S\mathrm z_S}.
\label{eq:rg-gated-merge-cost}
\end{equation}

(d) \emph{Stationary block law.} Let \(\mathrm s_*\) be the least, over all
\(S\) and all edges \(e\notin\operatorname{att}(S)\) joining two components of
\(\operatorname{att}(S)\), of the least positive singular value of
\(\mathbf a_eP_{W_{\operatorname{att}(S)}}\), where \(P_W\) is the orthogonal
projection onto \(W\). When no such pair \((S,e)\) exists, put
\(\mathrm s_*=+\infty\) and \(1/\mathrm s_*^2=0\). For \(\mathrm a>0\) and
\(0<\mathrm T\le1\), the
probability under \(q_\Omega\) that \(\mathfrak P_{\mathrm a}(\mathbf x)\) is
not the partition of \(V\) into the components of a minimum edge cover is at
most
\begin{equation}
\mathrm K_1\mathrm T+\frac{|E|\|\Omega\|}{2\mathrm s_*^2}\mathrm a^2+4|E|e^{-\mathrm a^2/(4\mathrm T)}.
\label{eq:rg-gated-block-law}
\end{equation}
\status{ESTABLISHED}

The proof is in \cref{sec:rg-formation-depth-one-proof}. In the four-cycle,
\(\mathfrak n=2\) and the minimum edge covers are the two perfect matchings,
so (c) gives them all but \(O(\mathrm T)\) of the stationary mass of either
law and (d) makes coherence of an unattended bridge at resolution
\(\mathrm a\) cost \(O(\mathrm a^2)\) per bridge. A second coalescence that
is held by an attended bridge at stationarity with probability bounded below
as \(\mathrm T\to0\) therefore needs a model outside the family.

\openproblemheading{Repeated endogenous formation}{open:rg-formation-general}
Construct, at \(c_*\), a second endogenous coalescence with inherited source
interactions and retained marks, in a family where a matching or partition
forms from initially nonmatching source assignments, generative parameters are
learned, and the episode length is not fixed in advance. What holds is
\Cref{prop:rg-formation-endogenous-cycle}, one duration-one episode seeded by
the initial assignment, \Cref{thm:finite-depth-adaptive-formation}, whose
hierarchy is supplied, and \Cref{prop:rg-gated-depth-one}, by which in the
gated translation family of \Cref{def:rg-gated-translation-family} the
stationary coherent-block partition is, up to the error of
\cref{eq:rg-gated-block-law}, the partition into the components of a minimum
edge cover, and every further merge at stationarity costs a factor
\(\mathrm T\). The question includes arbitrary families and initial
laws, mixed-level source support, permanent membership under later forcing,
and control as the depth tends to infinity; the intrinsic selector and the
hierarchy question are
\Cref{open:cg-natural-selector,open:full-graph-hierarchy}. It is posed for
models at \(c_*\) with source-gated attention on a finite connected graph
whose energies, precisions, prior ratios, hyperparameters, and mobilities do
not depend on the temperature \(\mathrm T\), as \(\mathrm T\to0\), outside the
family of \Cref{def:rg-gated-translation-family}. A second coalescence in
such a model is declared by a handoff time, deterministic or a stopping time;
a first-stage event of the history up to the handoff, with its block
partition; a second-stage event, with windows measured from the handoff,
whose block is a union of at least two first-stage blocks; one statistic of
the handoff state, namely the probability of the second-stage event, the
probability that it has occurred by declared times, or the probability of a
declared block; and a handoff class, fixed before the comparison, that
contains every first-stage handoff and, with probability bounded below, the
handoffs of histories without first-stage formation. The problem is settled
by such a coalescence, with learned generative parameters and from initially
nonmatching assignments, whose second-stage event does not occur with
probability tending to one uniformly over the handoff class, whose statistic
after first-stage formation exceeds its value after a counterfactual handoff
in the class by a margin bounded below, whose block is held at a declared
resolution by a restoring interaction whose rate times the squared
resolution, divided by \(\mathrm T\), tends to infinity on the sheets the
kinetics select, on horizons short compared with the stationary recurrence
time of the event measured by its rate of entrance from a set of stationary
mass bounded below, and with parents whose conditional and dynamical errors
vanish separately; or by a proof that in this class every coalescence of
first-stage blocks on such horizons fails the margin or the restoring
clause or has a second-stage event occurring with probability tending to one
uniformly over the handoff class.

\section*{Physicist's summary}

In the triangle chain a supplied hierarchy of stiffnesses
\(\mathrm J_1\gg\cdots\gg\mathrm J_L=1\) makes the prepared displacement relax
like overdamped springs of very different stiffness: the stiffest bonds
equilibrate first and glue their endpoints, and each softer level glues the
clusters already formed (\Cref{thm:finite-depth-adaptive-formation}). Every
parent keeps every loop and source assignment, starts from the complete
current law, and learns its own parameter; the price of each reduction
vanishes in an extreme regime (temperature \(\epsilon^{86}\) at depth two) in
which nothing adaptive decides the outcome. In the homogeneous four-cycle, without
a supplied hierarchy, an edge carries energy only when an endpoint attends to
it; attended edges contract, the prepared circulation stretches the unattended
ones, and the posterior keeps attending to the contracted edges, so a pair
forms with probability bounded away from zero
(\Cref{prop:rg-formation-endogenous-cycle}). The same gating
is why the story stops there in that family
(\Cref{prop:rg-gated-depth-one}): an unattended edge between two formed
blocks feels only the thermal reference, so at stationarity the blocks are
the components of a minimum edge cover and each further merge costs a factor
\(\mathrm T\).

\chapter{Interpretation and Its Limits}
\label{ch:philosophy}

What do the preceding chapters commit to, and which readings of them could an
observation or a proof count against? The chapter's readings are definitions,
hypotheses, and open problems; it states a proposition, a theorem, and a
corollary, and summarizes seven propositions proved in the appendices.
\Cref{prop:phil-seeded-contrast-formation} shows that on the four-agent cycle
of \cref{sec:rg-formation-endogenous-cycle} a collective forms from every
Gaussian initial law with probability bounded below uniformly in the regime
parameter, through a lower bound that is carried by the random initial
assignment seeding the pair, is the same whether or not the source prior
learns, and is a positivity statement of size at most \(3\times10^{-30}\); its proof is in
\cref{sec:rg-formation-cycle-proof}. The main statement,
\cref{conj:phil-contrast-formation}, is an instance of the parameter-learning
locus of \cref{sec:phil-participatory}: at a sufficiently small fixed
temperature, the continuous-time learning rule for the mixture proportions of
the source prior selects one pairing at the level of the law from almost every
Gaussian initial law, while a frozen prior leaves the long-time pairing law
symmetric. The learned prior amplifies the asymmetry of the initial law, as
any hierarchical model with the same learned prior would. It feeds
\cref{open:phil-testability}.

Beyond \cref{sec:intro-conventions}, the chapter declines the claims that a
structural-realist, Neo-Kantian, or noumenal reading follows from the invariance
results or is selected by the mathematics; that the lineage citations below
prove the probability construction or establish its interpretation; that
\cref{def:phil-idle-wheel} follows from constructive empiricism; that
participation entails consciousness or that variational free energy is a
physical action principle; and that a constructed record sensitive to link
holonomy is a canonical observable or an empirical realization. Its open items
are \cref{claim:physical-law-identification,open:phil-testability}.
\status{NOT-CLAIMED}

\section{Scaffolding, idle posits, and structural readings}
\label{sec:phil-base}
\label{sec:phil-noumenon}
\label{sec:phil-realism}

The verdict of the preceding chapters is that the architecture is local-first,
not background-free. Agent \(i\) carries the law points
\begin{equation}
q_i^b\in(\mathcal E_b)_{c_*},
\qquad
q_i^m\in(\mathcal E_m)_{c_*},
\label{eq:phil-fiber-measures}
\end{equation}
which the cross maps \(\Phi\) and \(\widetilde\Phi\) do not identify, and the
base, principal bundle, and measurable spaces are coordination scaffolds that
make agent-relative interfaces comparable and that local inference laws do not
construct. The verdict would change if edges became random structural
coordinates updated by inference (\cref{open:obs-edge-creation}) or links were
recovered from population records, with single-point part
\cref{open:phil-testability}. A Neo-Kantian reading takes local frames, model
evaluations, and access maps as constitutive but revisable conditions of a
record's inferential meaning, in the lineage of Kant, Reichenbach, Cassirer, and
Friedman \citep{Kant1781,Reichenbach1965,Cassirer1923,Friedman2001}; a noumenal
reading takes the base and bundles as structure reached only through
agent-side representations. As scaffolding is no inferential target, the
base's lack of a law decides neither reading.

\definitionheading{Idle-wheel criterion}{def:phil-idle-wheel}
A posit is idle when changing it alters the law of no record and of no
statistic of a record in the model. The idle-wheel criterion removes idle
posits by parsimony. It is a declared rule of reading.

The criterion is close in spirit to constructive empiricism, on which accepting
a theory involves belief only in its empirical adequacy
\citep[11--12]{vanFraassen1980}, but it is declared, not derived. A closed walk
\(p\) and an edge copy \(e\) determine the classes
\begin{equation}
[\mathrm{Hol}^b(p)]_{\rm conj},
\qquad
[\mathrm{Hol}^m(p)]_{\rm conj},
\qquad
[\mathring\Theta^x_e]_{\rm conj}
\label{eq:phil-holonomy}
\end{equation}
of the link holonomies \eqref{eq:geo-link-holonomy} and the per-edge twists,
invariant under passive and active changes
(\cref{prop:geo-edge-twist-invariance}). On a graph-flat family the loop classes
are trivial (\cref{prop:geo-trivializing-criterion}) and give a noumenal reading
no work; a nontrivial one is a referent only within the interaction complex,
made operational at \(c_*\) by \cref{thm:cg-loop-operational-holonomy}.

The interface theory of perception permits an arbitrary measurable relation
between representation and world \citep{Hoffman2015,Hoffman2019}. Here the
principal fiber at \(c_*\) is a \(G\)-torsor, a frame is the choice of a
reference point, and what survives a change of frame is structural data,
\begin{equation}
\mathcal I_{\mathrm{struct}}
=\left\{
\text{ELBO and relative-entropy values},\
[\mathrm{Hol}^b(p)]_{\rm conj},\
[\mathrm{Hol}^m(p)]_{\rm conj},\
\bigl[(\mathring\Theta^b_e,\mathring\Theta^m_e)_{e}\bigr]
\right\},
\label{eq:phil-invariants}
\end{equation}
the last entry being the per-edge twists of both channels modulo one
simultaneous conjugation, invariant by
\cref{prop:gen-product-evidence-invariance},
\eqref{eq:cg-kl-bimeasurable-invariance}, and
\cref{prop:geo-edge-twist-invariance}. Of the three agreements of
\cref{prop:geo-three-agreements}, frame agreement is independent of the other
two, law coherence implies mean agreement but not conversely, and none selects a
partition (\cref{open:cg-natural-selector}). These invariants are compatible
with epistemic structural realism \citep{Worrall1989}, with moderate structural
realism, on which objects share the ontological footing of relations and are
characterized only by them \citep{esfeld2008moderate}, and with ontic structural
realism, on which structure is all that exists \citep{Ladyman2007,Ladyman2014}.
No argument from theory change is given, Worrall's synchronic application of
epistemic structural realism to an accepted theory is not invoked as support,
and nothing here selects among these positions or the Neo-Kantian
reading. Support would need, respectively, the empirical success of the
construction for a target system that the epistemic position presupposes,
which is part of \cref{open:phil-testability}, a typed interpretation of the formalism's objects and
relations in the moderate ontology, and an argument from formal
unidentifiability to ontological elimination with discriminating evidence.

\Cref{ch:links} changes the status of the links, not of the scaffolding. Inferred as Lie-algebra
generators under a declared prior, the links take the formal role of a latent field: stationarity
balances the stress of the prior against a matter current that, for a comparison record, has the
outer-product form of a delta rule (\cref{prop:link-field-equation,cor:link-delta-rule}); near zero
generators the stationarity operator has the form of a massive Euclidean vector operator in the episode
index, with the matter as a screening mass (\cref{prop:link-linearized-equation}); and the held-link
model, at held links whose twists lie in the image of the exponential map, is the zero-width limit of
the latent one (\cref{thm:link-held-limit}). These are correspondences
of form, which that chapter declines to read as a physical field or as a learning architecture. The base,
the principal bundle, the frames, and the reference point remain structural data that inference does not
construct, so the verdict above is unchanged, and links inferred within the model are not links
recovered from the records of a target system, which remains part of \cref{open:phil-testability}.

\section{Participation, formation, and testability}
\label{sec:phil-participatory}
\label{sec:phil-testable}

\hypothesisheading{Operational observational closure}{hyp:phil-observational-closure}
Every datum directly available to an interior agent is a clamped state or a
message of another interior agent or of a boundary environment node, an
operational environment node in the sense of
\cref{def:obs-operational-environment-node}. The hypothesis excludes data that
reach a record kernel other than through such a node. It does not remove the
conditioning record, does not require a boundary node to be autonomous, and does
not exclude an unobserved environment that affects all agents without becoming
the argument of any agent's belief.

By \cref{thm:obs-agent-interaction-equivalence} the hypothesis is free for
kernels but removes no conditioning: a statistic of the record is
posterior-sufficient exactly when the conditional mutual information in
\eqref{eq:obs-conditioning-statistic-chain} vanishes, and a latent-independent
kernel is a zero-information boundary message
\eqref{eq:obs-conditioning-information}, not an agent.

\hypothesisheading{Every boundary node is called an agent}{hyp:phil-boundary-agents}
Every boundary environment node of \cref{hyp:phil-observational-closure} is
called an agent. The hypothesis is a naming decision. It adds to no node a
persistent state, an action, a Markov blanket, or a local free energy, and it
excludes reading the constructed joint \(\mathbb P_{\theta,V}\) or the selected
coupling \(\mathbb Q_{V,o,X}\) as an agent: these are complete joint laws at the
population level. A block conditional is a meta-agent only after its own coarse
state, observation interface, evaluated generative mechanism, recognition law,
and update rule have been constructed, as in
\cref{ch:agent-network-rg,ch:full-graph-meta-agent-vfe}. A strong agent-only
ontology would additionally need a criterion of agency for every boundary node,
and the formalism supplies no unique one.

Participation is reciprocal inferential dependence through records inside one
joint, with no live posterior in a same-step generative factor
(\cref{req:gen-typing-prohibition}). \Cref{sec:full-graph-closing-tower} locates
it in three sequences of fixed joints: the same-time inference loop, block
coordinate ascent as in any connected hierarchical model under mean-field
updates, which resembles Wheeler's participatory proposal \citep{Wheeler1990}
without being derived from it or evidence for it; parameter learning, whose
optimal prior over an unrestricted family is the aggregate recognition law
\eqref{eq:obs-aggregate-prior}; and delayed feedback through
\eqref{eq:obs-episode-update}, which for an update reading a declared
recognition kernel is the free energy of one joint
(\cref{prop:full-graph-episode-lifted-joint}) and whose dynamic closure is
part of \cref{open:rg-pointwise-parent-dynamics}. None of these separates the program
from a hierarchical model with learned priors. Formation with the
counterfactual-handoff dependence of \cref{open:rg-formation-general} could,
provided it is shown to fail for rivals without inference, such as the
bounded-confidence and self-attention particle dynamics recalled in
\cref{sec:rg-formation-endogenous-cycle}, which form clusters as well.
\Cref{thm:finite-depth-adaptive-formation} forms successive collectives under a
supplied energy hierarchy and source support, without showing that source
learning creates the contrasts, and \cref{prop:rg-formation-endogenous-cycle}
forms one collective on a homogeneous four-cycle with probability bounded below
uniformly in its regime parameter, from a symmetric preparation whose random
initial source assignment seeds the pair, and without parameter learning. The
same seed gives a lower bound on the probability that a collective forms from
every Gaussian initial law, and the bound does not depend on whether the source
prior learns.

\propositionheading{Seeded contrast formation with a learning-independent lower bound}{prop:phil-seeded-contrast-formation}
Take the homogeneous four-cycle experiment of
\cref{def:rg-formation-cycle-experiment}, with its regime parameter
\(0<\epsilon\le1/100\), the constants of \cref{eq:rg-formation-cycle-regime},
and its records, evaluator, reference law, and kinetics, and modify it in two
ways. First, the uniform prior on assignments becomes a product
\(\pi_S=\prod_i\pi_{i,S_i}\) of receiver rows, each a generative parameter that
starts uniform and learns from the recognition law \(q_t\) of the kinetics by
\begin{equation}
\dot\pi_{ij}=\epsilon\bigl(q_t(S_i=j)-\pi_{ij}\bigr),
\label{eq:phil-source-learning}
\end{equation}
the natural-gradient flow of the free energy in the prior row under the Fisher
metric of the simplex at mobility \(\epsilon\); its rest point is the recognition
label marginal, as for the aggregate prior \eqref{eq:obs-aggregate-prior}, and the
assignment is redrawn from posterior rows computed with the current prior.
Second, the initial law is \(\mathcal N(\mathsf m_0,I_{12})\) on configurations,
with an arbitrary mean \(\mathsf m_0\in\R^{12}\), times the uniform law on
assignments. Call contrast formation the event that some connected block \(I\)
with \(2\le|I|\le|V|-1\) has, at time zero, an internal edge whose belief
transported relative entropy exceeds \(1/4\), and, on some time interval of
length one inside \([0,2/\epsilon]\), all its internal belief and model
transported relative entropies at most \(\epsilon^4\) while every belief
transported relative entropy on an edge joining it to its complement exceeds
\(1/2\), the bounds of part (a) of \cref{prop:rg-formation-endogenous-cycle};
the first clause makes the event an entrance into the contrast and excludes a
contrast present at time zero that persists. Then:

(a) \emph{Rows.} The kinetics and \cref{eq:phil-source-learning} have exactly
one joint solution on \([0,\infty)\), with continuous deterministic rows, and
\(|\pi_{ij}(t)-1/2|\le\epsilon t/2\) for all \(t\ge0\).

(b) \emph{Formation.} For every \(\mathsf m_0\in\R^{12}\) and every
\(\epsilon\in(0,1/100]\),
\[
\begin{aligned}
\Pr(\text{contrast formation})
&\ge\frac{e^{-1/4}}{16}\sum_{\pm}\mathcal N(\mathsf m_0,I_{12})
\bigl\{|\mathbf x\mp\mathbf x^\circ|<\varrho_\circ\bigr\}\\
&\ge\frac{e^{-1/4}\varrho_\circ^{12}}{737280}
\exp\Bigl[-\tfrac12\bigl(|\mathbf x^\circ-\mathsf m_0|+\varrho_\circ\bigr)^2\Bigr]>0,
\end{aligned}
\]
and the same bound holds when the rows are frozen at \(1/2\). The bound is
carried by the sub-event in which \(I\) is a matched pair of the initial
assignment, held throughout \([0,t^\circ_*]\), and the contrast is that of
part (a) of \cref{prop:rg-formation-endogenous-cycle}. This sub-event is a
coherence episode with the restoring interaction and rates of that part, and
for every fixed row configuration \(\pi\) the kinetics at rows \(\pi\) leave
\(q^\star_\pi(\mathbf x,S)\propto\pi_Se^{-U_S(\mathbf x)/\mathrm T}\) invariant
and, started from it, realize the sub-event with probability at most
\(4e^{-1/(4\mathrm T)}\); so for every \(\mathsf m_0\), in either arm, the
sub-event exceeds every such stationary value by a margin bounded below as
\(\epsilon\to0\), the first display not depending on \(\epsilon\).

(c) \emph{Invariant initial laws.} The initial law is invariant under the cycle
automorphisms, with link signs transformed under reflections, exactly when
\(\mathsf m_0=(\mathrm c_1\mathbf 1_4,\mathrm c_2\mathbf 1_4,0)\) for some
\(\mathrm c_1,\mathrm c_2\in\R\), and then \(\pi_{ij}(t)=1/2\) for every
receiver, source, and \(t\ge0\).
\status{ESTABLISHED}

The proof is in \cref{sec:rg-formation-cycle-proof}. The lower bound in part (b)
is supplied by the preparation's matching seed, an initial assignment drawn
uniformly and independently of a configuration law of full support, and its
proof uses learning only through the bound of part (a) on the rows over the
seed interval \([0,t^\circ_*]\). The proposition does not compare the learned
and frozen probabilities of contrast formation on \([0,2/\epsilon]\), and the
comparison of the formed blocks in the two arms at this mobility is part of
\cref{open:rg-formation-adaptive-family}. The bound is also small: since the
mass of a ball of fixed radius under \(\mathcal N(\mathsf m_0,I_{12})\) is
largest when the ball is centered at \(\mathsf m_0\), the first display is at
most \((e^{-1/4}/8)\mathcal N(0,I_{12})\{|\mathbf x|<\varrho_\circ\}<3\times10^{-30}\) for
every \(\mathsf m_0\), so part (b) is a positivity statement and gives no
estimate of the typical probability of contrast formation. A statement in
which learning creates a contrast must compare learned and frozen priors and concern what the law selects,
not the probability of a sample-path episode.

\theoremheading{Pairing selection by slow source learning}{conj:phil-contrast-formation}
Take the homogeneous four-cycle experiment of
\cref{def:rg-formation-cycle-experiment}, with its records, evaluator, reference
law, and kinetics, at a fixed temperature \(\mathrm T>0\) in place of
\(\epsilon^8\); let the product receiver rows of
\cref{prop:phil-seeded-contrast-formation} start uniform and learn by
\cref{eq:phil-source-learning} with a mobility \(\epsilon_{\mathrm L}>0\) in
place of \(\epsilon\); and let the initial law be
\(\mathcal N(\mathsf m_0,I_{12})\) on configurations times the uniform law on
assignments. Write \(q^{\epsilon_{\mathrm L}}_t\) for the recognition law of the
kinetics and \(q^0_t\) for that of the frozen-prior arm, whose rows stay at
\(1/2\). First, there is \(\mathrm T_0>0\) such
that for every \(\mathrm T\in(0,\mathrm T_0]\) there is
\(\mathrm C_{\mathrm T}<\infty\) for which, for Lebesgue-almost every
\(\mathsf m_0\in\R^{12}\), one of the two perfect matchings
\(M\in\{S^{(0)},S^{(1)}\}\) satisfies
\[
\lim_{\epsilon_{\mathrm L}\to0}\ \sup_{t_{\mathrm L}\le t\le2t_{\mathrm L}}
q^{\epsilon_{\mathrm L}}_t(S\ne M)=0,
\qquad
t_{\mathrm L}=\mathrm C_{\mathrm T}\epsilon_{\mathrm L}^{-1}\log(1/\epsilon_{\mathrm L}).
\]
Second, for every \(\mathrm T>0\) and every \(\mathsf m_0\),
\(q^0_t(S=S^{(0)})-q^0_t(S=S^{(1)})\to0\) as \(t\to\infty\).
\status{ESTABLISHED}

The proof is in \cref{sec:phil-pairing-selection-proof}.
\corollaryheading{Joint selection along a temperature schedule}{cor:phil-joint-selection}
In the model of \cref{conj:phil-contrast-formation}, fix any sequence
\(\mathrm T_n\downarrow0\) within its temperature interval and any
\(\eta_n\in(0,1)\) tending to zero. There are thresholds \(d_n\downarrow0\), independent
of the initial mean, such that for Lebesgue-almost every fixed
\(\mathsf m_0\), all sufficiently large \(n\), and every
\(0<\epsilon_{\mathrm L}\le d_n\), a matching \(M_n(\mathsf m_0)\) satisfies
\[
 \sup_{t_{\mathrm L}\le t\le2t_{\mathrm L}}
 q_t^{\epsilon_{\mathrm L},\mathrm T_n}(S\ne M_n)\le\eta_n,
 \qquad t_{\mathrm L}=2\epsilon_{\mathrm L}^{-1}\log(1/\epsilon_{\mathrm L}).
\]
The same thresholds can make the learned law's relative entropy to its
instantaneous invariant law, and the frozen law's relative entropy to the
uniform-row invariant law, at most \(\eta_n\) throughout this window.
Consequently the frozen matching probabilities tend to \(1/2\) each.
For any prescribed positive caps \(h_n\), in particular
\(h_n=\mathrm T_n^a\) with any \(a>0\), the thresholds can also satisfy
\(d_n\le h_n\). They are existence thresholds, rather than a
proved power-law selection condition; the selected matching may depend on
\(n\).
\status{ESTABLISHED}

The proof and an explicit temperature-dependent tracking bound are in
\cref{sec:phil-joint-selection-proof}. Setting
\(\mathrm T(\epsilon_{\mathrm L})=\mathrm T_n\) for
\(d_{n+1}<\epsilon_{\mathrm L}\le d_n\) gives a joint limit with a schedule
common to almost every initial mean.
\Cref{prop:phil-selected-pair-persistence} identifies what persists on this
late window. At residual resolution \(\mathrm a=\mathrm T^{1/4}\), the
selected pairing remains the coherent partition for a following window of
length \(\mathrm T^{-1/4}\), with probability tending to one when the schedule tolerances satisfy
\(\eta_n\le\mathrm T_n\) and the thresholds satisfy \(d_n\le\mathrm T_n^2\). The frozen arm approaches the equal
mixture of the two persistent pairings. The statement controls the source
redraws of the kinetics; its comparison with a held-sheet Gaussian process
is an estimate on path laws. This is restored equilibrium coherence.

The first contrast block is a different statistic.
\Cref{prop:phil-early-block-null} keeps the contrast event and the early
horizon of \cref{prop:phil-seeded-contrast-formation}, writing
\(\xi=\mathrm T^{1/8}\) and separating the learning mobility
\(\epsilon_{\mathrm L}\) from \(\xi\). If
\(\epsilon_{\mathrm L}=o(\mathrm T^{1/4})\), the learned and frozen laws of
the first block, each conditioned on contrast formation, have total
variation tending to zero for every fixed Gaussian initial mean. Thus the
slow schedules can select different late pairing laws while leaving the
early first-block comparison unchanged. The original regime
\(\epsilon_{\mathrm L}=\xi\) is outside this estimate.

The retained channel in part (b) of
\cref{prop:rg-formation-endogenous-cycle} also applies after this late
coherence. Complete a duration-one coherence prefix at
\(\theta=t_{\mathrm L}+1\), and initialize the parent from the full
recognition law conditioned on that prefix, pushed through the channel.
\Cref{prop:phil-actual-conditioned-parent,prop:phil-actual-parent-dynamics}
give, along \cref{cor:phil-joint-selection} with tolerances
\(\eta,\delta\le\mathrm T\) and mobility cap
\(\epsilon_{\mathrm L}\le\mathrm T^2\), separately vanishing discarded
conditional relative entropy and retained-path total variation over a
following interval \(\mathrm T^{-1/4}\). The parent inherits the rows at activation and then
learns from its own actual law. It retains the assignment, the loop, and
the common fast normal coordinate; no independent Gaussian preparation
is substituted for its activation law. These are averaged guarantees
for this declared past event. They do not assert closure conditional on
every realized history or activate a parent at the first contrast.

The same channel can be activated early, on the seeded sub-event of the
original contrast, at every mobility \(\epsilon_{\mathrm L}\in[0,1]\) and so
at \(\epsilon_{\mathrm L}=\xi\). By \cref{prop:phil-early-seeded-parent}, the
discarded conditional relative entropy at a deterministic activation time
\(\theta\) vanishes when \(\theta-\log(1/\xi)\to\infty\) and
\(\xi^4\theta\to0\), and diverges when
\(\log(1/\xi)-\theta\to\infty\); \cref{prop:phil-early-parent-dynamics} adds a
vanishing retained-path error on following windows short compared with
\(\mathrm T^{-1}\). By \cref{prop:phil-early-first-completion}, activation at
the first qualifying completion instead has diverging conditional error, also
when conditioned on the whole contrast event. These results concern the seeded
sub-event and leave the comparison of learned and frozen first blocks at
\(\epsilon_{\mathrm L}=\xi\) open.

The limit \(\epsilon_{\mathrm L}\to0\) at fixed \(\mathrm T\) places the
mobility of the rows below any positive mixing rate of the kinetics at fixed
rows, and \(t_{\mathrm L}\) is, up to its constant, the time in which an imprint
of order \(\epsilon_{\mathrm L}\) on the rows that grows at a rate of order
\(\epsilon_{\mathrm L}\) reaches order one. The asymmetry that learning would
amplify is supplied by the initial law, not by the generative model or the
kinetics: by part (c) of \cref{prop:phil-seeded-contrast-formation}, learning
from a Gaussian initial law that is invariant under the cycle automorphisms
leaves the rows uniform for all time. The theorem amplifies the asymmetry of
\(\mathcal N(\mathsf m_0,I_{12})\) and excludes a null set of means. The frozen
arm is the null against which selection is measured. This is a result about
the one-time assignment law at fixed temperature. The conditional block
events, a power-law vanishing-temperature regime, and conditional and dynamical parent
errors of \cref{open:rg-formation-adaptive-family} remain separate obligations.
The participatory reading requires collectives constituted by the coupled
inference, with a restoring interaction and a dependence on the preceding
state. Assignment selection alone does not establish that recursive claim. The noumenal reading buys structure behind the
agents' representations with posits that \cref{def:phil-idle-wheel} removes
unless an observable such as a loop holonomy carries them.

\openproblemheading{Physical-law identification}{claim:physical-law-identification}
Identifying an object produced by the coarse-graining of
\cref{ch:agent-network-rg}, or by its limits, with a physical law requires a
named target system, an observable and an estimator, a validation and
uncertainty protocol set before the data, with the error probabilities of its
decision, a baseline margin, probability-level closure, and held-out evidence that is robust across admissible blocking
schemes. None is supplied here, and the mathematics neither proves a physical
law nor supports ontological closure. For a named target system whose data the
model does not generate, observed through a declared record kernel that, under
the model's own law, separates two objects of the coarse-graining that differ
in a gauge-invariant datum, with an observable, an estimator, and a validation
and uncertainty protocol set before the data whose decision has two error
probabilities, stated in advance and summing to less than one, the problem is
settled by
held-out evidence, robust across admissible blocking schemes, that an object of
the coarse-graining meets a baseline margin with probability-level closure, or
by held-out evidence under the same protocol that the baseline margin is not
met.

Graph flatness (\cref{prop:geo-trivializing-criterion}) and the fixed-sector
condition of \cref{prop:cg-kernel-holonomy} check declared data but detect no
bundle topology or base curvature, so an interpretation must keep graph-link
holonomy, base-connection holonomy, bundle topology, and gauge redundancy apart.

\openproblemheading{Single-point testability}{open:phil-testability}
State an observation at \(c_*\) that could count against a claim of the program.
What holds is two model theorems, which no observation could count against:
\cref{thm:cg-loop-operational-holonomy} constructs a
population record and a frame-invariant statistic whose law separates two
independent-link configurations that differ only in the conjugacy class of a
loop holonomy, and \cref{conj:phil-contrast-formation} proves that the
continuous-time learning rule for the mixture proportions of the source prior
selects a pairing at the level of the law against a symmetric frozen-prior arm
in one declared family, a selection that does not separate the program from a
hierarchical model with the same learned prior. What is missing is a
target system: sampling units, a noise model, an estimator of gauge-typed
parameters with its uncertainty, a held-out statistic, an acceptance tolerance
and error probabilities set in advance, and baselines without link holonomy or
without formation. A
cross-scale test additionally needs two or more admissible partitions,
coarse channels, gauge alignment, and held-out coarse statistics predicted from
the fine fit without coarse refitting, and its failure would falsify the added
cross-scale hypothesis, not the finite-dimensional theorems. For a named
target system at \(c_*\) whose data the model does not generate, observed
through a declared record kernel, with sampling units and a noise model, that
reads the links of a closed walk or the coordinates of a candidate block, the
problem is settled by a test fixed before the data, with an estimator of
gauge-typed parameters, a held-out statistic, an acceptance tolerance, and
baselines without link holonomy or without formation, whose failure would
falsify the presence of nontrivial holonomy on that walk or of formation of
that block by the inference in the target, together with a proof that, under
the declared kernel and sample size, the test rejects under the baseline
without link holonomy (respectively without formation), and accepts under the
hypothesis, each with probability at least one minus an error probability
stated before the data, the two error probabilities summing to less than one;
or, negatively, by a proof that under the declared kernel every
statistic of the records has the same law for two link configurations that
differ only in the conjugacy class of the walk's holonomy, or the same law
whether or not the block is coherent at the declared resolution.

\paragraph{Physicist's summary.}
Beyond scaffolding, frame-change invariants select no realist position and
environment messages create no agents; coordinate ascent, learning, and delayed
feedback do not distinguish the program from hierarchical modeling, and the
pairing selection is an instance of learning; formation by the inference with a
counterfactual-handoff dependence would, if clustering dynamics without
inference were shown to lack it.

\chapter{Related Work}
\label{ch:related-work}

Where does this construction sit among existing work, and what does it owe to
it? The central comparison, with Bayesian renormalization in
\cref{sec:rg-bayesian-comparison}, rests on
\cref{prop:gen-no-distinguished-target} and
\cref{thm:cg-evidence-preserving-channel}, and the open bridge problem of
Supplement G builds on it.

No construction below is identified with this one or adopted as authority
beyond its cited statement, and qualifications of other works are this
manuscript's reading. \status{NOT-CLAIMED}

\section{Gauge theory, free energy, and multi-agent active inference}

\citet{Sengupta2016NeuronalGauge} propose, in an essay, a gauge theory of
neuronal dynamics in which variational free energy plays the role of a
Lagrangian, gauge fields arise from the information geometry of the brain's
probability distributions, and attention, the cognitive counterpart of
precision weighting, is a force produced by the curvature of that geometry.
Here the group acts on the sample fibers of one principal bundle
(\cref{ch:geometry}), links are independent group elements on edge copies at one
base point (\cref{def:geo-graph-links}), and attention is a latent source label
with a derived posterior (\cref{prop:obs-attention-elbo}) whose precision is
that of a record (\cref{prop:obs-selective-receiver-update}).

In federated inference \citep{friston2024federated}, agents share a generative
model of a shared world and communicate beliefs through a likelihood mapping to
overt expressions, and language emerges when agents share an econiche.
\citet{heins2024collective} obtain cohesion, milling, and directed motion from
agents minimizing surprise under generative models of their neighbors,
recovering social forces as suppressed prediction errors without a group-level
model. \citet{waade2025asone} treat a collective whose group-level Markov
blanket consists of its sensory agents' observations and its active agents'
actions as a group-level agent, fit a model of the same structure to simulated
collective behavior, and find the individual and group models nontrivially
related. Here agents may carry inequivalent representations and independent
links, their encoder, transmission, and decoder kernels enter one joint with one
evidence bound (\cref{def:obs-selective-interface,ch:elbo}), parent laws follow
from the fine joint by exact coarse channels
(\cref{thm:rg-exact-coarse-vfe,thm:rg-pointwise-parent-datum}), formation is
proved for a stated family (\cref{ch:finite-depth-formation}), and a low-rank
interface does not by itself give a Markov blanket
(\cref{prop:geo-selective-descent}).

\section{Connection Laplacians, sheaves, and equivariant networks}

\citet[Section 3 and Theorem 5.1]{SingerWu2012} build a graph connection
Laplacian from edge weights and orthogonal transformations aligning the tangent
spaces of neighboring sample points and prove its convergence to the connection
Laplacian of the underlying manifold; \citet[Section 3.6]{HansenGhrist2019}
identify it as the sheaf Laplacian of an orthogonal vector bundle over a graph,
in a spectral theory of sheaves whose discrete Hodge theorem identifies harmonic
cochains with sheaf cohomology \citep[Theorem 3.1]{HansenGhrist2019}. Our
quadratic energies with represented links are of this type, with
kernel the holonomy-fixed sector (\cref{prop:cg-kernel-holonomy}), but with
links in \(G\) acting through representations, semidefinite weights, and
normalized laws; the quadratic disagreement of
\cref{def:obs-selective-comparison} is the edge term of a weighted sheaf energy
with rectangular restriction maps, as in the discourse sheaves of selective
expression \citep{HansenGhrist2021Discourse}.
\citet[Section 4.2]{HansenGhrist2019} show that Kron reduction is in general
impossible for sheaves, since an eliminated vertex can impose constraints
between boundary vertices that no pairwise restriction maps express, and
possible for vertex stalks of dimension at most one, where the Laplacians have
factor width at most two \citep[Theorem 4.3]{HansenGhrist2019}; this is prior
art for the Kron obstructions of matrix-weighted Gaussian energies that
Supplement G classifies, whereas Galerkin closure
(\cref{thm:cg-rectangular-endpoint-closure}) restricts a form to a subspace.

Group-equivariant convolutional networks share weights across the orbits of a
global symmetry group \citep{Cohen2016}, gauge-equivariant ones make a network
on a manifold independent of the choice of local frames \citep{cohen2019gauge},
and a gauge-equivariant transformer applies multi-head self-attention to feature
fields on meshes with a parallel transport that makes the layer agnostic to the
orientation of local coordinate systems \citep{He2021GaugeEquivariantTransformer}.
These constrain learned maps over a base manifold with a connection. Here links
are independent dynamic data at one base point, a nonidentity per-edge twist has
no counterpart in an architecture with one shared key map
(\cref{tab:obs-transformer-correspondence}), and
\cref{cor:obs-selective-kernel-gauge,prop:geo-selective-covariance} describe how
declared data transform, not constraints on a learned layer.

\section{Averaging and path-space coarse-graining}

\citet[Section~2.1.4, Theorem~2.9]{HartmannNeureitherSharma2020}
bound path-space relative entropy for an effective diffusion under regularity,
conditional functional inequalities, and Novikov's condition, with the
resolved diffusion coefficient independent of unresolved coordinates. Their
fixed-time estimates in Theorems~2.4 and~2.7 have a different target. Neither
estimate supplies the growing-history, switching-source comparison requested
by \cref{open:rg-pointwise-parent-dynamics}. The one-time selection result
of \cref{conj:phil-contrast-formation} uses the separate entropy, transient
imprint, and saddle-escape estimates in \cref{sec:phil-pairing-selection-proof}.
A change in diffusion covariance can instead make continuous-path laws
singular. \citet[Remark~1 and Theorem~1]{BackhoffUnterberger2023}
study a partition-refinement relative entropy for martingale diffusions;
that quantity is distinct from ordinary path-space relative entropy. The
filter and memory constructions of \cref{sec:rg-memory-parent-mechanism}
retain the dynamical information that a stationary marginal alone omits.

\section{Information-geometric and Bayesian renormalization}
\label{sec:rg-bayesian-comparison}
\label{sec:obs-tau}

\citet{beny2015information} formulate the renormalization group as a family of
quantum channels connecting the microscopic world to the observable world at a
given scale, endow the space of Hamiltonians with a metric induced by an
information geometry on states, measure the information lost along a flow by the
distinguishability of thermal states, and introduce functions of two-point
correlations that do not increase along the flow; \cref{thm:cg-fisher-contraction}
is only the classical Fisher-information statement for one parameter-independent
kernel at one parameter value.

Bayesian renormalization, as formulated by Berman, Klinger, and Stapleton, builds
on the dynamical Bayesian inference of \citet{BermanHeckmanKlinger2022} and uses
Fisher distinguishability on model-parameter space to define an information
scale and to separate stiff from sloppy directions. In the late-data regime the
posterior is a Gaussian whose covariance is the inverse Fisher information times
the scale parameter \citep[Eq.~(38)]{BermanKlingerStapleton2023}. For a
differentiable forward map \(\mathfrak O\) from parameters to data they write the
pushforward of that posterior as a Gaussian
\citep[Eq.~(43)]{BermanKlingerStapleton2023} whose covariance is the scale
parameter times the pushed inverse Fisher metric
\begin{equation}
D^{\mathfrak O}_\theta
=\partial\mathfrak O(\theta)g^F(\theta)^{-1}\partial\mathfrak O(\theta)^\top
\label{eq:rg-bks-diffusion}
\end{equation}
at the posterior center \citep[Eq.~(44)]{BermanKlingerStapleton2023}, the
diffusion tensor of a Fokker--Planck equation and a stochastic flow on data
space \citep[Eqs.~(45) and~(46)]{BermanKlingerStapleton2023}. In our reading the
delta method makes their Gaussian form the leading local approximation of the
pushforward, exact when \(\mathfrak O\) is affine, and its continuation beyond
the late-data regime is a renormalization construction, not an exact nonlinear
pushforward or a Bayesian asymptotic theorem. It has been inverted into a
statistical-inference reading of exact renormalization flows
\citep{BermanKlinger2024Inverse} and applied to neural-network field theories
\citep{HowardKlingerMaitiStapleton2025}.

The information scale ranks which directions of a parent model presentation stay
distinguishable at a declared precision, but by itself it selects no graph
blocks, coarse-grains no nonflat links, and leaves the two-channel interaction
law open. Letting
a posterior flow does not make a generative kernel read its own posterior, so
\cref{req:gen-typing-prohibition} does not refute Bayesian renormalization, but
inserting a posterior from one scale as the generative top factor at another
declares a new model or an explicit transition between models, whose bounds
refer to different joints unless a common target and a comparison map are
supplied. By \cref{prop:gen-no-distinguished-target}, pointwise identities
designate no one member or evidence, and comparing two bounds shows an
improvement in one quantity only on an evidence level set; for two independently
declared joints with no evidence-preserving channel between them, the
difference of optimized bounds mixes evidence and tightness differences and
supplies no order, whereas an evidence-preserving recognition-independent Markov
kernel applied to the posterior and the recognition law orders the bounds by
data processing (\cref{thm:cg-evidence-preserving-channel,thm:cg-kl-dpi-extended}).
Relating the widening of the posterior along its scale to the addition of
precision under Gaussian aggregation would need a common state space, a map
between the scale parameters, and a transported monotone; Supplement G records
this bridge as open.

\citet{MehtaSchwab2014} map Kadanoff's variational renormalization group
exactly onto deep restricted Boltzmann machines; a deterministic
identification of coarse and fine coordinates does not supply its kernel, so the
Galerkin coarse form of \cref{sec:cg-galerkin} is not thereby a deep network.

\section{Network renormalization, graph coarsening, and multigrid}
\label{sec:grg-literature-boundary}
\label{sec:full-graph-vfe-literature}

Network renormalization addresses heterogeneous graphs, which have no lattice,
momentum, or geometric length \citep{gabrielli2025network}. The Laplacian
renormalization group of \citet{villegas2023laplacian} uses the diffusion
density operator \(e^{-tL}/\Tr e^{-tL}\) of a graph Laplacian \(L\), its
spectral entropy \(\mathcal S(t)\) normalized by the logarithm of the number of
nodes, and the specific heat
\begin{equation}
\chi(t)=-\frac{d\mathcal S}{d\log t},
\label{eq:grg-heat-capacity}
\end{equation}
whose peaks they use to identify characteristic resolution scales
\citep[Eqs.~(1)--(2)]{villegas2023laplacian}.

\remarkheading{Gauge behavior of heat-kernel diagnostics}{rem:grg-heat-kernel-gauge}
Let \(R\succ0\) and \(L\) be block matrices on \(\bigoplus_{i\in V}\R^{n_i}\), and
let a vertex reframing \(P=\operatorname{diag}(P_i)\) with invertible blocks
act by congruence, \(R\mapsto P^\top RP\) and \(L\mapsto P^\top LP\). Then the
generator \(R^{-1}L\) becomes \(P^{-1}R^{-1}LP\), so its spectrum, its heat trace
\(\Tr e^{-tR^{-1}L}\), and every spectral entropy built from them are invariant,
whereas the off-diagonal blocks of \(e^{-tR^{-1}L}\) become
\(P_i^{-1}(e^{-tR^{-1}L})_{ij}P_j\) and are not invariant scalar node affinities.
\status{ESTABLISHED}

\noindent\emph{Proof.}
\((P^\top RP)^{-1}P^\top LP=P^{-1}R^{-1}LP\), and the exponential of a similar
matrix is similar by the same matrix. \(\square\)

A gauge-invariant block selector built from heat-kernel data therefore needs
invariant statistics of the blocks, and a scale diagnostic alone selects no
unique partition (\cref{open:cg-natural-selector}).
\citet{Garuccio2023Multiscale} give a renormalization scheme for any hierarchy
of heterogeneous coarse-grainings and identify a class of scale-invariant random
graphs with a necessary dependence on additive node hidden variables, a
characterized family rather than arbitrary hidden-variable networks; for random
graphs whose link probabilities are conditioned on other links,
\citet{CatanzaroGarlaschelliPatil2026} derive a closed-form renormalization
transformation for maximum coordination number two and show that higher
coordination numbers admit none. These candidate partition and closure
mechanisms do not by themselves handle directed row-normalized attention, two
gauge representations, or nonflat link holonomy, and the aggregation of
\cref{ch:network-state} does not integrate out edge variables. A combined scheme
would infer a partition from gauge-invariant multiplex diagnostics of edge-event
laws through a structural channel with a named law and push the complete law
exactly.

Graph coarsening and pooling in machine learning add hypotheses of their own:
restricted spectral approximation controls spectral and cut properties
\citep{Loukas2019}, coarsening can carry message-passing guarantees
\citep{JolyKeriven2024}, differentiable pooling learns soft cluster assignments
end to end \citep{Ying2018DiffPool}, and deep graph convolutional networks whose
weights satisfy spectral conditions lose all information except connected
components and node degrees exponentially in depth \citep{OonoSuzuki2020}. The
Galerkin identity preserves the restricted quadratic energy, while spectral
approximation, message-passing commutation, and a graph-neural-network
realization need the further controls of \cref{sec:cg-galerkin}
(\cref{prop:cg-message-passing-residual}). That coarse form is the one of
algebraic multigrid: aggregation-based multigrid forms \(P^\top AP\) with a
piecewise-constant prolongation defined by disjoint aggregates
\citep[Section 2, Eqs.~(2.1)--(2.2)]{Notay2010Aggregation}, smoothed
aggregation smooths that tentative prolongation \citep{VanekMandelBrezina1996},
and lean algebraic multigrid for graph Laplacians combines piecewise-constant
interpolation, node aggregation by an affinity measure, and an energy correction
of the coarse systems \citep{LivneBrandt2012LAMG}. What is added here is the
placement of that operator inside normalized laws with gauge covariance and
exact evidence accounting.

\paragraph{Physicist's summary.}
Here the group acts on sample fibers with independent links at one point, and
parent laws are computed from the fine joint rather than shared, omitted, or
fitted as in multi-agent active inference. Sheaf Laplacians supply the linear
algebra and the Kron obstruction, and renormalization flows meet the exact
channels only through a common comparison target.


\appendix
\crefalias{chapter}{appendix}
\crefalias{section}{subappendix}
\chapter{Measure-Theoretic Lemmas and Integrability Witnesses}
\label{app:measure-lemmas}

This appendix proves the measure-theoretic statements of \Cref{ch:probability,ch:elbo} and separates the
hypotheses of \cref{hyp:elbo-evidence-domain} from factorwise log-integrability
(\cref{hyp:elbo-factorwise-integrability}) by witnesses.

\section{Kernel measurability and proofs for the probability chapter}

\theoremheading{Jointly measurable Radon--Nikodym version for a dominated kernel}{thm:prob-kernel-rn-measurable-version}
Let $(\mathsf E,\mathscr E)$ be a measurable space, and let $K$ and $L$ be probability kernels from $\mathsf E$
to $(\mathsf F,\mathscr F)$ with $K(x,\cdot)\ll L(x,\cdot)$ for every $x$. Then there is a finite
$\mathscr E\otimes\mathscr F$-measurable $r:\mathsf E\times\mathsf F\to[0,\infty)$ with
$K(x,B)=\int_Br(x,y)L(x,dy)$ for every $x\in\mathsf E$ and $B\in\mathscr F$. Selecting a derivative separately
for each $x$ does not give joint measurability; the theorem supplies one simultaneous version.
\status{ESTABLISHED}

\noindent\emph{Proof.}
Standard Borelness gives finite refining measurable partitions $(\mathcal P_n)_{n\geq1}$ whose generated
sigma-algebras increase to $\mathscr F$. With $0/0:=0$ put
\begin{equation}
r_n(x,y)=\sum_{A\in\mathcal P_n}\frac{K(x,A)}{L(x,A)}\mathbf 1_A(y),
\label{eq:prob-kernel-rn-partitions}
\end{equation}
and let $r=\limsup_nr_n$ where this is finite and $r=0$ elsewhere; like every $r_n$, $r$ is jointly
measurable. For fixed $x$, $r_n(x,\cdot)$ is the conditional expectation of $h_x=dK(x,\cdot)/dL(x,\cdot)$ given
$\sigma(\mathcal P_n)$ under $L(x,\cdot)$, so martingale convergence gives $r_n(x,\cdot)\to h_x<\infty$
$L(x,\cdot)$-almost everywhere, and $r(x,\cdot)$ is a derivative for every $x$. \hfill$\square$

\corollaryheading{Joint densities against one reference}{cor:prob-common-reference-joint-density}
Let $\nu$ be a sigma-finite measure on $(\mathsf F,\mathscr F)$ and $K$ a probability kernel from
$(\mathsf E,\mathscr E)$ to $\mathsf F$ with $K(x,\cdot)\ll\nu$ for every $x$. Then $K$ has a finite jointly
measurable density against $\nu$, and every finite or countable collection of kernels dominated by the same
$\nu$ has jointly measurable densities against that common reference; in particular this applies to the
generative and recognition kernels \eqref{eq:prob-generative-signature}--\eqref{eq:prob-recognition-signature}
when they are dominated by the product references. \status{ESTABLISHED}

\noindent\emph{Proof.}
Sigma-finiteness supplies measurable $w$ with $0<w<\infty$ everywhere and $\int w(y)\nu(dy)=1$, so the constant
kernel $L(x,dy)=w(y)\nu(dy)$ is equivalent to $\nu$. \Cref{thm:prob-kernel-rn-measurable-version} gives a finite
jointly measurable $r=dK/dL$, and $k=rw$ is the density. Apply this to each member of the collection.
\hfill$\square$

\propositionheading{Kernel integration preserves measurability}{prop:prob-kernel-integration-measurability}
Let $\kappa$ be an s-finite kernel from $(\mathsf E,\mathscr E)$ to $(\mathsf F,\mathscr F)$, that is, a
countable sum of finite kernels; probability kernels and fixed sigma-finite measures are instances. For
every $\mathscr E\otimes\mathscr F$-measurable $f:\mathsf E\times\mathsf F\to[0,\infty]$,
\label{prop:prob-sigmafinite-integration-measurability}
\begin{equation}
x\longmapsto\int_{\mathsf F}f(x,y)\kappa(x,dy)
\label{eq:prob-kernel-integration}
\end{equation}
is $\mathscr E$-measurable. For an extended signed integrand the same holds on the measurable set where the
positive-part and negative-part integrals are not both infinite, and no expression $+\infty-\infty$ is formed.
Consequently the evidence \eqref{eq:prob-evidence} and the $\mathbb Q_{V,o,X}$-expectations of \Cref{ch:elbo}
are measurable in $(o,X)$ whenever their integrands are jointly measurable. \status{ESTABLISHED}

\noindent\emph{Proof.}
For finite $\kappa$, the class of $C\in\mathscr E\otimes\mathscr F$ for which $x\mapsto\kappa(x,C_x)$ is
measurable, with $C_x=\{y:(x,y)\in C\}$, contains the rectangles, since
$\kappa(x,(A\times B)_x)=\mathbf 1_A(x)\kappa(x,B)$, and is closed under proper differences and increasing
unions because each $\kappa(x,\cdot)$ is finite. Rectangles form a $\pi$-system generating
$\mathscr E\otimes\mathscr F$, so the $\pi$--$\lambda$ theorem gives every $C$. Linearity and monotone
convergence pass to nonnegative measurable $f$, and countable sums to s-finite $\kappa$; a sigma-finite $\nu$
with a measurable partition into sets $F_m$ of finite measure is the constant s-finite kernel
$\sum_m\nu(\cdot\cap F_m)$. A signed integrand is split into $f^+$ and $f^-$, whose difference is measurable
where they are not both infinite. \hfill$\square$

\corollaryheading{Measurability of kernel relative entropy}{cor:prob-kernel-kl-measurability}
Let $\kappa$ and $\eta$ be probability kernels from $(\mathsf E,\mathscr E)$ to $(\mathsf F,\mathscr F)$. Then
$x\mapsto\KL(\kappa(x,\cdot)\Vert\eta(x,\cdot))$ is $\mathscr E$-measurable as a $[0,+\infty]$-valued map, and
for finite refining partitions $(\mathcal P_n)$ generating $\mathscr F$,
\begin{equation}
\KL(\kappa(x,\cdot)\Vert\eta(x,\cdot))
=\sup_{n\geq1}\sum_{A\in\mathcal P_n}\kappa(x,A)\log\frac{\kappa(x,A)}{\eta(x,A)},
\label{eq:prob-kernel-kl-partitions}
\end{equation}
with $0\log(0/b)=0$ and $a\log(a/0)=+\infty$ for $a>0$. \status{ESTABLISHED}

\noindent\emph{Proof.}
Put $L=(\kappa+\eta)/2$. \Cref{thm:prob-kernel-rn-measurable-version} applied to $(\kappa,L)$ and $(\eta,L)$
gives finite jointly measurable densities $a$ and $b$. Fix $x$, write $L_x=L(x,\cdot)$, and let $a_n,b_n$ be the
conditional expectations of $a(x,\cdot),b(x,\cdot)$ under $L_x$ given $\sigma(\mathcal P_n)$, whose cell values
are the cellwise ratios in \eqref{eq:prob-kernel-kl-partitions}. The perspective
$\phi_{\mathrm{KL}}(u,v)=u\log(u/v)-u+v$, with $\phi_{\mathrm{KL}}(0,v)=v$, $\phi_{\mathrm{KL}}(u,0)=+\infty$ for
$u>0$, and $\phi_{\mathrm{KL}}(0,0)=0$, is nonnegative, convex, and lower semicontinuous. Since both kernels have
unit mass, the $n$-th partition sum is $D_n(x)=\int\phi_{\mathrm{KL}}(a_n,b_n)dL_x$ and the divergence is
$\int\phi_{\mathrm{KL}}(a,b)dL_x$. Conditional Jensen makes $D_n(x)$ nondecreasing in $n$ and at most the
divergence, and martingale convergence $(a_n,b_n)\to(a,b)$ $L_x$-almost everywhere, with Fatou's lemma and
lower semicontinuity, gives $\int\phi_{\mathrm{KL}}(a,b)dL_x\leq\liminf_nD_n(x)$. Each $D_n$ is measurable in
$x$, so the countable supremum is measurable. \hfill$\square$

\noindent\emph{Proof of \Cref{thm:prob-mixed-family-common-domination}.}
If $A_*$ is countable, $\nu_{\mathrm{mix}}$ is sigma-finite, and a $\nu_{\mathrm{mix}}$-null set is Lebesgue null
and contains no point of $A_*$, so every term of every $\mu_t$ vanishes on it. Conversely, let a sigma-finite
$\xi$ dominate every $\mu_t$. Each $a\in A_*$ has $\mu_t(\{a\})>0$ for some $t$, whence $\xi(\{a\})>0$, and a
sigma-finite measure has at most countably many points of positive mass, as in the proof of
\Cref{sec:exp-dominated-nondominated}. The moving-atom family has $A_*=\R$. \hfill$\square$

\remarkheading{A regular-conditional version is not determined at a point}{rem:prob-version-witness}
On $\R^2$ with coordinates $(o,y)$, let $\phi$ be the standard normal density and
$\psi\neq\phi$ another probability density. Then $p(o,y)=\phi(o)\phi(y)$ and $\widetilde p$, equal to $p$ for
$o\neq0$ and to $\phi(0)\psi(y)$ for $o=0$, are versions of the same joint density; both evidence
representatives equal $\phi(0)$ at $o=0$, but their density-ratio conditionals there are $\phi(y)dy$ and
$\psi(y)dy$. The joint measure therefore determines neither pointwise posterior at $o=0$.
\status{ESTABLISHED}

\noindent\emph{Proof.}
The two functions differ only on the Lebesgue-null line $\{0\}\times\R$, and integrating out $y$ at $o=0$
gives $\phi(0)$ in both cases. \hfill$\square$

\noindent\emph{Proof of \Cref{lem:prob-kl-chain-rule}.}
(a) Measurability of the integrand is \Cref{cor:prob-kernel-kl-measurability} with $\mathsf E=\mathsf Y_1$.
Suppose first $Q\ll P$. Then $Q_1\ll P_1$, since $P_1(A)=0$ gives $P(A\times\mathsf Y_2)=0$. Let $r=dQ/dP$ and
$r_1=dQ_1/dP_1$. For measurable $A\subseteq\mathsf Y_1$ and $B\subseteq\mathsf Y_2$,
\[
\int_AQ_{2|1}(y_1,B)r_1(y_1)P_1(dy_1)=Q(A\times B)=\int_A\int_Br(y_1,y_2)P_{2|1}(y_1,dy_2)P_1(dy_1),
\]
so for each $B$ the integrands agree $P_1$-almost everywhere; over a countable generating algebra of
$\mathsf Y_2$ the exceptional sets form one $P_1$-null set $N$, off which both sides are finite measures in $B$
that agree on the algebra and hence everywhere. Off the $Q_1$-null set $N\cup\{r_1=0\}$, $Q_{2|1}(y_1,\cdot)$
therefore has the density $r(y_1,\cdot)/r_1(y_1)$ against $P_{2|1}(y_1,\cdot)$, and $\log r$ splits $Q$-almost
surely into $\log r_1$ and the logarithm of that density. Both summands have negative parts of $Q$-integral at
most $1/e$, as in \Cref{def:prob-log-density-ratio-convention}, so expectations add in $(-\infty,+\infty]$; the
first gives $\KL(Q_1\Vert P_1)$, and Tonelli's theorem for the positive and negative parts along
$Q=Q_1\otimes Q_{2|1}$ gives the integral of the conditional divergences. Suppose next $Q\not\ll P$. If
$Q_1\not\ll P_1$, the first term on the right is $+\infty$. If $Q_1\ll P_1$ and the conditional term were
finite, then $Q_{2|1}(y_1,\cdot)\ll P_{2|1}(y_1,\cdot)$ for $Q_1$-almost every $y_1$, so every $P$-null $C$, whose
sections are $P_{2|1}(y_1,\cdot)$-null for $P_1$-, hence $Q_1$-almost every $y_1$, would be $Q$-null. Conditional
versions are unique $Q_1$-almost everywhere when $Q_1\ll P_1$, and otherwise the right side is $+\infty$ for
every version.

(b) Quantize every block by finite refining partitions generating its sigma-algebra. On the finite product
alphabet a cell of positive $Q$-mass has marginal cells of positive $Q_b$-mass, so the first term on the right
is finite; if every marginal term is finite, every such cell has positive reference mass and the logarithm
splits cellwise, and otherwise a marginal cell of positive $Q_b$-mass and zero $\rho_b$-mass makes both sides
infinite. By \eqref{eq:prob-kernel-kl-partitions} for constant kernels, each of the three divergences is the
nondecreasing limit of its quantized values, which gives \eqref{eq:prob-kl-product-reference}.

(c) By (a), every $Q$ with $Q_1=\mu$ has $\KL(Q\Vert P)\geq\KL(\mu\Vert P_1)$. When
$\KL(\mu\Vert P_1)<\infty$, equality holds exactly when $Q_{2|1}=P_{2|1}$ $\mu$-almost everywhere,
which determines $Q$. When the marginal divergence is infinite, every member of the class has divergence $+\infty$.

(e) If $Q\ll P$ with $r=dQ/dP$, then $A_\#Q(B)=\int_{A^{-1}B}r(y)P(dy)=\int_Br(A^{-1}y')(A_\#P)(dy')$ for measurable
$B$, which is \eqref{eq:cg-rn-bimeasurable}, and integrating $\log(r\circ A^{-1})$ against $A_\#Q$ gives
$\int\log r(y)Q(dy)$. If $Q(S)>0=P(S)$, then $A(S)$ is measurable because $A^{-1}$ is, and
$A_\#Q(A(S))>0=A_\#P(A(S))$, so both divergences are $+\infty$.

(d) Form the laws $Q(dy)K(y,dz)$ and $P(dy)K(y,dz)$ on $\mathsf Y\times\mathsf Z$. By (a) with $y$ as the first
coordinate their divergence is $\KL(Q\Vert P)+\int\KL(K(y,\cdot)\Vert K(y,\cdot))Q(dy)=\KL(Q\Vert P)$. By (e) for the
coordinate swap and (a) with $z$ as the first coordinate, the same divergence is $\KL(QK\Vert PK)$ plus a term
in $[0,+\infty]$. \hfill$\square$

\section{The relative-log representation and integrability witnesses}
\label{sec:app-integrability-witnesses}

\propositionheading{Relative-log representation}{prop:elbo-relative-log-representation}
Under (H1)--(H2) of \Cref{hyp:elbo-evidence-domain}, let $M_{\theta,V,o,X}$ be the slice measure
\eqref{eq:elbo-slice-measure}. If $Q\ll M_{\theta,V,o,X}$ and $r=dQ/dM_{\theta,V,o,X}$, then
\begin{equation}
\begin{gathered}
\Lelbo_V^{\mathrm{ext}}(Q;X,o)=-\int_{\mathsf Y}\log r(y)Q(dy),\\
\int_{\{r<1\}}-\log r(y)Q(dy)=\int_{\{r<1\}}-r(y)\log r(y)M_{\theta,V,o,X}(dy)\leq\frac{p_{\theta,V}(o\given X)}{e},
\end{gathered}
\label{eq:elbo-relative-log}
\end{equation}
so the integral is well defined in $(-\infty,+\infty]$ and its negative lies in $[-\infty,+\infty)$. If
$Q\not\ll M_{\theta,V,o,X}$, then $\Lelbo_V^{\mathrm{ext}}(Q;X,o)=-\infty$. \status{ESTABLISHED}

\noindent\emph{Proof.}
Since $\boldsymbol\Pi_{\theta,V,o,X}=M_{\theta,V,o,X}/p_{\theta,V}(o\given X)$, the Radon--Nikodym chain rule gives
$dQ/d\boldsymbol\Pi_{\theta,V,o,X}=p_{\theta,V}(o\given X)r$. The negative part of $\log r$ is $Q$-integrable by the
displayed bound, because $-u\log u\leq1/e$ on $[0,1]$ and $M_{\theta,V,o,X}(\mathsf Y)=p_{\theta,V}(o\given X)$.
Hence $\KL(Q\Vert\boldsymbol\Pi_{\theta,V,o,X})=\log p_{\theta,V}(o\given X)+\int\log r(y)Q(dy)$ is a legitimate extended
identity, and substitution in \eqref{eq:elbo-extended-definition} gives the first branch. The slice measure and
the posterior have the same null sets, which gives the second. \hfill$\square$

\propositionheading{Integrability witnesses}{prop:elbo-integrability-witnesses}
The integrability hypotheses of \Cref{ch:elbo} are separated by the following witnesses.
\begin{enumerate}
\item[(a)] Clauses (H1)--(H3) of \Cref{hyp:elbo-evidence-domain} do not imply (H4). On $\R$ with a standard
Gaussian posterior and a standard Cauchy recognition law, (H1)--(H3) hold,
$\E_Q\log q=-\log4\pi$ is finite, $\E_Q\lvert\log p_{\theta,V}(o,Y\given X)\rvert=+\infty$, the divergence is
$+\infty$, and the extended bound is $-\infty$.
\item[(b)] Finite relative entropy does not imply (H4). On $\{2,3,\ldots\}$ with a trivial record, let the
posterior and the recognition law both equal $P(\{n\})=\mathrm c/(n(\log n)^2)$, with $\mathrm c$ normalizing.
The divergence is $0$, while $\sum_nP(\{n\})\lvert\log P(\{n\})\rvert=+\infty$, so both expectations in the
classical split equal $-\infty$ and their difference is undefined.
\item[(c)] Factorwise log-integrability (\cref{hyp:elbo-factorwise-integrability}) does not imply (H4). With
one latent $k\in\R$, generative factor $\mathcal N(0,1)$, and a record whose log density is bounded for
$k\in[0,1]$, let the recognition law spread mass $w_n=2^{-n}$ uniformly over disjoint intervals in $[0,1]$ of
widths $w_ne^{-1/w_n}$, whose total is below $e^{-2}$. Every factor logarithm is bounded on the support, while
$\log q=1/w_n$ on the $n$-th interval and $\E_Q\log q=\sum_nw_n/w_n=+\infty$.
\item[(d)] Factorwise log-integrability is strictly stronger than the generative half of (H4).
Take a root model coordinate $m\in\{2,3,\ldots\}$ with $p(m)=\mathrm c/(m(\log m)^2)$, a bridge
$k\given m\sim\mathcal N(0,\sigma_m^2)$ with $\sigma_m^2=p(m)^2/(2\pi)$, a trivial record, and the recognition
law equal to this law. Then $\log p(m)+\log p(k\given m)=-k^2/(2\sigma_m^2)$ has expectation $-\tfrac12$ and is
also the recognition log density, so (H4) holds, while $\E\lvert\log p(m)\rvert=+\infty$.
\end{enumerate}
\status{ESTABLISHED}

\noindent\emph{Proof.}
(a) Both laws have positive Lebesgue densities, so (H1)--(H3) hold. The log joint differs from $-y^2/2$ by a
constant and the Cauchy law has no second moment; the Cauchy differential entropy is $\log4\pi$; hence the
divergence $\E_Q\log q-\E_Q\log p$ is $+\infty$. (b) For large $n$, $P(\{n\})\leq1/n$, so
$\lvert\log P(\{n\})\rvert\geq\log n$ and the sum dominates $\mathrm c\sum_n1/(n\log n)=+\infty$. (c) The widths sum
to $\sum_n2^{-n}e^{-2^n}<e^{-2}$, so the intervals fit in $[0,1]$; the density on the $n$-th interval is
$w_n/(w_ne^{-1/w_n})=e^{1/w_n}$. (d) The bridge density at $k$ is
$(2\pi\sigma_m^2)^{-1/2}e^{-k^2/(2\sigma_m^2)}=p(m)^{-1}e^{-k^2/(2\sigma_m^2)}$, which gives the identity and
$\E\lvert\log p(m)+\log p(k\given m)\rvert=\tfrac12$; and $\E\lvert\log p(m)\rvert=+\infty$ as in (b). \hfill$\square$

\chapter{Deferred Proofs}
\label{app:deferred-proofs}

This appendix collects the longer proofs of \cref{ch:selective-communication,ch:relational-inference,ch:coarsegraining,ch:agent-network-rg}, in the order of the statements they prove.

\section{Proofs for Attention and Selective Communication}
\label{sec:proofs-attention}

\noindent\emph{Proof of \Cref{prop:obs-attention-elbo}.}
(a) The product label prior and exclusivity make the posterior of \(J\) factor
by rows, and Bayes' formula gives \eqref{eq:obs-attention-posterior}; a weight
\(\pi_{ij}e^{-D_{ij}/\tau_i}\) is positive exactly when \(\pi_{ij}>0\), since
the energies are finite, and with equal prior weights the largest weight has
the smallest energy, \(D_{ii}=0\). (b) With
\(Z_i^{\rm att}(y)=\sum_k\pi_{ik}e^{-D_{ik}(y)/\tau_i}\),
\(\KL(\beta^Q_i\Vert\beta^P_i(y))=\KL(\beta^Q_i\Vert\pi_i)
+\tau_i^{-1}\sum_j\beta^Q_{ij}D_{ij}(y)+\log Z_i^{\rm att}(y)\). Subtracting
\(\log Z_i^{\rm att}(y)\) and the source-independent \(\log\mathrm c_i(o_i,y)\)
and taking the recognition expectation proves
\eqref{eq:obs-attention-full-contribution} and \eqref{eq:obs-attention-vfe};
Lagrange multiplication under \(\sum_j\beta^Q_{ij}=1\) gives
\eqref{eq:obs-attention-recognition-optimum}, and strict convexity of
\(\KL(\cdot\Vert\pi_i)\) gives uniqueness, which the factor \(\tau_i>0\) does
not change. (c) The label part of the VFE is
\(\E_{\mathbb Q}[\log\mathbb Q_{J\mid Y}-\sum_i\log\pi_{iJ_i}
-\sum_i\log L_i(o_i\given Y,J_i)]\). The product-reference case of
\cref{lem:prob-kl-chain-rule} splits
\(\KL(\mathbb Q_{J\mid Y=y}\Vert\prod_i\pi_i)\) into the total correlation and
the row relative entropies, which with \eqref{eq:obs-attention-likelihood}
gives \eqref{eq:obs-attention-correlated-ledger}. At fixed \(y\) the expansion
of (b), applied to the product reference, equals the label terms other than
\(-\sum_i\log\mathrm c_i\) to
\(\KL(\mathbb Q_{J\mid Y=y}\Vert\prod_i\beta^P_i(y))-\sum_i\log Z_i^{\rm att}(y)\),
which the product of the posterior rows minimizes. \(\square\)

\noindent\emph{Proof of \Cref{prop:obs-attention-constant-row-gap}.}
Substituting \eqref{eq:obs-attention-recognition-optimum} into
\eqref{eq:obs-attention-vfe} gives \(R_i^Q\), the minimum by the expansion in
the proof of \cref{prop:obs-attention-elbo}. At fixed \(y\) the same expansion
gives
\(\KL(\mathbb Q_{J_i\mid Y=y}\Vert\pi_i)+\tau_i^{-1}\E_{\mathbb Q_{J_i\mid Y=y}}D_{iJ_i}(y)
=\KL(\mathbb Q_{J_i\mid Y=y}\Vert\beta_i^P(y))+R_i^P(y)\), minimized by
\(\beta_i^P(y)\), and the conditional total correlation in
\eqref{eq:obs-attention-correlated-ledger} is nonnegative and vanishes on the
product of these rows. Normalization of \(\pi_i\) and concavity of the
logarithm give \(\min_jD_{ij}/\tau_i\leq R_i^P\leq\sum_j\pi_{ij}D_{ij}/\tau_i\),
so \(R_i^P(Y)\) is integrable. The softmax formulas give
\(\log\beta^{Q\star}_{ij}-\log\beta^P_{ij}(y)
=(D_{ij}(y)-d_{ij})/\tau_i+R_i^Q-R_i^P(y)\), and averaging over
\(\beta_i^{Q\star}\) and then over \(\mathbb Q_Y\) removes the first term,
since \(\E_{\mathbb Q_Y}D_{ij}=d_{ij}\), which proves
\eqref{eq:obs-attention-jensen-gap}. The relative entropy vanishes at \(y\)
exactly when \(D_{ij}(y)-d_{ij}\) does not depend on \(j\), and almost-sure
constancy of the differences forces those constants to equal
\(d_{ij}-d_{ik}\). Under the affine hypothesis
\(d_{ij}-d_{ik}=(\widetilde D_{ij}-\widetilde D_{ik})(\bar\psi_i)\), and a
softmax row depends only on logit differences, which gives
\eqref{eq:obs-attention-affine-row}; template differences are affine in
\(y_i\). For the random source, uncorrelatedness gives
\(\E r_{ij}=\bar y_i-\rho_b(\Theta^b_{ij})\bar y_j\) and
\(\operatorname{Cov}r_{ij}=\operatorname{Cov}Y_i
+\rho_b(\Theta^b_{ij})\operatorname{Cov}(Y_j)\rho_b(\Theta^b_{ij})^\top\), so
\(d_{ij}\) exceeds \(D_{ij}\) at the recognition mean by
\(\frac12\Tr(W_i\operatorname{Cov}Y_i)+\frac12\Tr(W_i\rho_b(\Theta^b_{ij})
\operatorname{Cov}(Y_j)\rho_b(\Theta^b_{ij})^\top)\), whose first trace is
common to the row and cancels from the softmax. \(\square\)

\noindent\emph{Proof of \Cref{cor:obs-attention-replicator}.}
On the open simplex the Fisher metric is the Shahshahani metric, whose natural
gradient of a function with Euclidean gradient \(g\) is
\(-\beta_j(g_j-\sum_k\beta_kg_k)\). The Euclidean gradient of
\eqref{eq:obs-attention-vfe} is \(\mathsf c_{ij}+1\), and the constant cancels
on the simplex, which gives \eqref{eq:obs-attention-replicator}. Pairing the
gradient with the flow gives \(-\mathsf r_i\sum_j\beta_{ij}\mathsf c_{ij}
(\mathsf c_{ij}-\bar{\mathsf c}_i)=-\mathsf r_i\operatorname{Var}\), since
\(\sum_j\dot\beta_{ij}=0\); with moving \(\lambda\) the chain rule adds the
explicit \(\lambda\)-derivative, the second term of
\eqref{eq:obs-attention-joint-rate}. The quotient rule for \(\beta=\eta/\alpha\)
gives \eqref{eq:rg-attention-beta-event}, and differentiating
\(\log\beta_j=\log\pi_j-E_j/\tau-\log Z\) gives
\(\dot\beta_j/\beta_j=\chi_j-\dot Z/Z\) with \(\dot Z/Z=\sum_k\beta_k\chi_k\),
which is \eqref{eq:rg-attention-beta-softmax}. \(\square\)

\noindent\emph{Proof of \Cref{prop:obs-attention-reduced-hessian}.}
Differentiating the finite sum gives
\(\nabla\beta_{ij}=-\tau_i^{-1}\beta_{ij}(\nabla E_{ij}-\bar g)\) and the
gradient; once more,
\(\nabla^2R_i=\tau_i^{-1}\sum_j(\beta_{ij}\nabla^2E_{ij}+\nabla E_{ij}\nabla\beta_{ij}^\top)\),
and \(\sum_j\beta_{ij}\bar g(\nabla E_{ij}-\bar g)^\top=0\) turns the second
sum into \(-\tau_i^{-2}\operatorname{Cov}_{\beta_i}(\nabla E_i)\). The joint
Hessian of \(\mathscr J_i\) has blocks
\(\tau_i^{-1}\sum_j\beta_{ij}\nabla^2E_{ij}\), the mixed form
\((v,w)\mapsto\tau_i^{-1}\sum_jw_j\nabla E_{ij}\cdot v\), and
\(w\mapsto\sum_jw_j^2/\beta_{ij}\) on the tangent space; minimizing the joint
quadratic form over tangent \(w\) at fixed \(v\) gives
\(w_j=-\tau_i^{-1}\beta_{ij}(\nabla E_{ij}-\bar g)\cdot v\), which satisfies
\(\sum_jw_j=0\), with minimum
\(v^\top(\tau_i^{-1}\sum_j\beta_{ij}\nabla^2E_{ij})v
-\tau_i^{-2}\operatorname{Var}_{\beta_i}(\nabla E_{ij}\cdot v)\). The
constant-row identification is \cref{prop:obs-attention-constant-row-gap}; for
exact label integration, the product label prior and exclusivity give
\(\sum_j\pi_{ij}\mathrm c_i(o_i,y)e^{-D_{ij}(y)/\tau_i}
=\mathrm c_i(o_i,y)e^{-R_i(y)}\). In the witness, \(O\) given \(J=j\) is
\(\mathcal N(-j\delta,\tau(1+\omega^{-1}))\), whose density at zero is the
displayed evidence for both labels. At \(o=0\),
\(\frac12\sum_{j=\pm1}e^{-(x-j\delta)^2/(2\tau)}=e^{-(x^2+\delta^2)/(2\tau)}
\cosh(\delta x/\tau)\), which with the baseline gives
\eqref{eq:obs-attention-reduced-witness}, and
\(\frac{d^2}{dx^2}\log\cosh(\delta x/\tau)=\delta^2/\tau^2\) at \(x=0\). The
constant-row value follows from \eqref{eq:obs-attention-reduced-hessian} with
the affine energies, whose weights are equal at \(\mathsf m=0\), together with
the recognition entropy and the baseline term, whose \(\mathsf m\)-curvature is
\(\omega/\tau\). \(\square\)

\noindent\emph{Proof of \Cref{cor:obs-selective-kernel-gauge}.}
(a) Integrating \(o\), then \(\xi^+\), then \(\xi^-\) gives total mass one, and
measurability follows from measurable kernel integration
(\cref{ch:probability}). (b) Substituting \(\xi^-=(\phi_e^-)^{-1}\xi'\) and
\(\xi^+=(\phi_e^+)^{-1}\xi''\) in the composition cancels the message
recoordinatizations. (c) The first identity is the definition of \(F_e^v\), and
the represented link law \(\rho_x(\Theta^x_e)'=\mathsf g_i\rho_x(\Theta^x_e)
\mathsf g_j^{-1}\) combined with (b) gives the last. \(\square\)

\noindent\emph{Proof of \Cref{prop:geo-selective-covariance}.}
(a) The two internal basis changes cancel in
\(\mathsf L^{\rm dec\prime}_e\mathsf L^{\rm tr\prime}_e\mathsf L^{\rm enc\prime}_e\).
(b) The inclusions and the rank bound follow from composition. For
\(k\in\ker\mathsf L^{\rm enc}_e\) and any \(\mathsf u\),
\((\mathsf G_j^{-1}(\mathsf L^{\rm enc}_e)^\top\mathsf u)^\top\mathsf G_jk
=\mathsf u^\top\mathsf L^{\rm enc}_ek=0\), and
\(\dim\operatorname{im}(\mathsf L^{\rm enc}_e)^\top=\dim E_j-\dim\ker
\mathsf L^{\rm enc}_e\), which identifies the complement; the metric law is the
congruence induced by \(\mathsf g_j\). Coordinates on
\(\operatorname{im}\mathsf L_e\), the induced quotient as encoder, the identity
as transmission, and the inclusion of the image as decoder factor
\(\mathsf L_e\) through its rank, which proves existence, not canonicity of an
interface. (c) Rank is subadditive and bounded by the target dimension, and
the example is direct. \(\square\)

\noindent\emph{Proof of \Cref{prop:geo-selective-loops}.}
(a) The factorization holds by definition of \(F_e\), and
\(\mathsf L_e'=\mathsf g_i\mathsf L_e\mathsf g_j^{-1}\) with
\(\Theta_e'=\mathsf g_i\Theta_e\mathsf g_j^{-1}\) gives the law of \(F_e\); for
every choice of \(\Theta_e\) some \(F_e\) exists, so the link is not
determined, and a common \(F_i\) exists exactly when all
\(\mathsf L_e\Theta_e^{-1}\) coincide. (b) Consecutive factors
\(\mathsf g_{i_{a+1}}^{-1}\mathsf g_{i_{a+1}}\) cancel, which uses no inverse of
any \(\mathsf L_e\). (c) and (d) are direct matrix computations; the holonomy
of the loop through \(b\) and \(\bar a\) is
\(\Theta_a^{-1}\Theta_b=\Theta_b\ne I_2\). \(\square\)

\noindent\emph{Proof of \Cref{cor:obs-selective-softmax}.}
(a) Expanding \(\KL(\beta_i\Vert\beta_i^P(y))\) with
\eqref{eq:obs-attention-posterior} and \(D_{ij}=d_{ij}\) gives the identity,
and the relative entropy vanishes only at \(\beta_i^P(y)\); the point mass is
singular with respect to the posterior slice when the baseline gives \(\{y\}\)
zero mass. (b) The terms of
\(\KL(\mathcal N(\mu^Q_i,\Sigma^Q)\Vert\mathcal N(\mu^K_j,\Sigma^K))\) depending
on \(j\) are those of
\(\frac12(\mu^K_j-\mu^Q_i)^\top(\Sigma^K)^{-1}(\mu^K_j-\mu^Q_i)\), and expanding
the quadratic form gives \eqref{eq:obs-selective-dot-product}. \(\square\)

\noindent\emph{Proof of \Cref{prop:obs-selective-receiver-update}.}
(a) The product label prior and exclusivity make \(J_i\) independent of \(k_i\)
and of every other record given \(y_{-i}\), and the records are conditionally
independent given the latent state, so the joint posterior density of
\((k_i,j)\) is proportional to \(\pi_{ij}\mathcal N(k_i;\mu_i,\Sigma_i)
\mathcal N(0;S^Q_ik_i-S^K_ek_j,\widetilde W_i^{-1})\). For fixed \(j\), adding
the two quadratic forms in \(k_i\) gives the precision \((\Sigma_i^+)^{-1}\) and
the component mean. Integrating \(k_i\) instead gives the density at zero of
\(o_i\) given \(J_i=j\), an affine image of independent Gaussian variables with
mean \(r_{ij}\) and covariance \(\Sigma_i^{\rm pred}\), whose normalizer does
not depend on \(j\); Bayes' formula gives \(w_{ij}\). The mean follows from the
tower property, and multiplying
\((\Sigma_i^+)^{-1}=\Sigma_i^{-1}+S_i^{Q\top}\widetilde W_iS^Q_i\) by
\(\Sigma_i^+\) gives the last identity. (b) Here \(\Sigma_i^{\rm pred}=2\),
\(r_{i2}=-2\), and \(\Sigma_i^+=1/2\); substitution gives the displayed values.
(c) The identity follows from
\(\Sigma_i^+S_i^{Q\top}\widetilde W_i\Sigma_i^{\rm pred}
=\Sigma_i^+(S_i^{Q\top}+S_i^{Q\top}\widetilde W_iS^Q_i\Sigma_iS_i^{Q\top})
=\Sigma_i^+(\Sigma_i^+)^{-1}\Sigma_iS_i^{Q\top}\). Since
\(S^Q_i\Sigma_iS_i^{Q\top}\) is invertible,
\((\Sigma_i^{\rm pred})^{-1}\to(S^Q_i\Sigma_iS_i^{Q\top})^{-1}\) as
\(\widetilde W_i^{-1}\to0\), which gives the limit. Directly \(S^Q_iR_e=I\), and
any \(\delta\) with \(S^Q_i\delta=\mathsf y\) is \(R_e\mathsf y+w\) with
\(S^Q_iw=0\), for which the cross term
\((R_e\mathsf y)^\top\Sigma_i^{-1}w=\mathsf y^\top(S^Q_i\Sigma_iS_i^{Q\top})^{-1}S^Q_iw\)
vanishes, so \(R_e\mathsf y\) uniquely minimizes \(\delta^\top\Sigma_i^{-1}\delta\).
(d) Repeat (a) with the copy-dependent data. \(\square\)

\noindent\emph{Proof of \Cref{prop:obs-selective-head-coupling}.}
At fixed \(k_i\) the proof of \cref{prop:obs-attention-elbo} applies verbatim
to the rows \((i,1)\) and \((i,2)\). Given the labels, the two records at zero
are a Gaussian vector with mean \(-(k_j,k_l)\) and covariance
\(I_2+\mathbf 1\mathbf 1^\top\), whose inverse is
\(I_2-\frac13\mathbf 1\mathbf 1^\top\); completing the square gives
\eqref{eq:obs-selective-head-coupling}. The masses are \(1\) for \((1,1)\) and
\(e^{-a^2/3}\) for each of \((1,2)\), \((2,1)\), and \((2,2)\), so the odds
ratio is \(e^{-a^2/3}/e^{-2a^2/3}=e^{a^2/3}\). \(\square\)

\noindent\emph{Proof of \Cref{prop:obs-selective-information-chain}.}
The chain rule of \cref{lem:prob-kl-chain-rule}, applied to mutual information
as a relative entropy, gives
\(I(S_j;\xi^+,O\given\Upsilon)=I(S_j;\xi^+\given\Upsilon)
=I(S_j;O\given\Upsilon)+I(S_j;\xi^+\given O,\Upsilon)\), because
\(I(S_j;O\given\xi^+,\Upsilon)=0\) by the conditional Markov property;
nonnegativity proves the first inequality, including infinite values, and the
second is identical. For finite alphabets without side information this is the
data-processing inequality \citep[Section 2.8, Theorem 2.8.1]{CoverThomas2006}.
\(\square\)

\noindent\emph{Proof of \Cref{prop:geo-selective-descent}.}
(a) Necessity follows by evaluating the commuting identity on \(\ker M_j\). For
sufficiency define \(\overline{\mathsf L}_e(M_j\mathsf u)=M_i\mathsf L_e\mathsf u\);
the condition makes this independent of the representative, and surjectivity
of \(M_j\) gives its domain and its uniqueness. (b) If
\(M\mathcal A=\overline{\mathcal A}M\), then
\(M\mathcal A^n=\overline{\mathcal A}^nM\) for every \(n\), and summing the
exponential series gives \(Me^{t\mathcal A}=e^{t\overline{\mathcal A}}M\);
conversely, differentiating at \(t=0\) returns
\eqref{eq:geo-selective-flow-descent}, and \(M\mathcal A\) factors through
\(M\) exactly when it vanishes on \(\ker M\), which is (a) with
\(\mathsf L_e=\mathcal A\) and \(M_j=M_i=M\). (c) Along any solution
\(\tfrac{d}{dt}(M_t\mathsf u)=(\dot M_t+M_t\mathcal A_t)\mathsf u\), so
\eqref{eq:geo-selective-moving-descent} is sufficient. For necessity fix \(t\)
and \(\mathsf v\in E\); the linear equation with continuous coefficients has a
solution through \(\mathsf u(t)=\mathsf v\), so
\((\dot M_t+M_t\mathcal A_t)\mathsf v=\overline{\mathcal A}_tM_t\mathsf v\) for
every \(\mathsf v\), and a linear map \(E\to\overline E\) factors through the
surjection \(M_t\) exactly when it vanishes on \(\ker M_t\), with a unique
factor. Multiplying \eqref{eq:geo-selective-moving-descent} on the right by
\(M_t^+\) and using \(M_tM_t^+=I\) gives the formula; in fixed inner products
\(M_t^+=M_t^\top(M_tM_t^\top)^{-1}\) is continuous. (d) For the swap,
\(\ker M_j\) is spanned by the second coordinate, which the swap sends outside
\(\ker M_i\). For the flow, choose \(\mathsf u\in\ker M\) with
\(M\mathcal A\mathsf u\ne0\); the solutions from \(\mathsf u\) and from \(0\)
have retained value \(0\) at \(t=0\) and retained velocities
\(M\mathcal A\mathsf u\) and \(0\). (e) The conditional law is the image of the
law of \(\mathsf z^{\rm in}\) under the diagonal map, which is a product only
for a degenerate law. \(\square\)

\section{Proofs for Fiber Information Geometry and Inference Histories}
\label{sec:proofs-fiber-geometry}

\noindent\emph{Proof of \Cref{prop:pb-statistical-tensor-descent}.}
Each $\widehat\rho_x(g)$ is pushforward by the bimeasurable sample map $\rho_x(g)$, which
does not depend on the law parameter.  Pushforward is a unitary map
$L^2(p)\to L^2(\rho_x(g)_\#p)$, $f\mapsto f\circ\rho_x(g)^{-1}$, differentiability in
quadratic mean makes $\mathbf s_Z\circ\rho_x(g)^{-1}$ the score of the pushed direction,
and the change of variables preserves \eqref{eq:pb-fiber-fisher} and
\eqref{eq:pb-fiber-amari}; since $\iota_{u\cdot g}=\iota_u\circ\widehat\rho_x(g)$,
$(\iota_u^{-1})^*\tau$ is independent of $u$ exactly when
$\widehat\rho_x(g)^*\tau=\tau$ for every $g$.  For the counterexample,
$\mathrm X=\mu/\sqrt2$ and $\mathrm S=\varsigma$ turn $g^F$ into
$2(d\mathrm X^2+d\mathrm S^2)/\mathrm S^2$, twice the hyperbolic metric of the upper
half-plane, of which the real M\"obius map $z\mapsto(1+z)/(1-z)$ of determinant $2>0$ is
an isometry.  It fixes $z=\mathrm i$ with derivative $2/(1-\mathrm i)^2=\mathrm i$, a
quarter turn, so it sends the $g^F$-unit vector $\partial_\mu$ at $(0,1)$ to
$\partial_\varsigma/\sqrt2$.  There the only nonzero components of $T^F$ are
$T^F(\partial_\mu,\partial_\mu,\partial_\varsigma)=2$ and
$T^F(\partial_\varsigma,\partial_\varsigma,\partial_\varsigma)=8$, so the cubic values of
the two unit vectors are $0$ and $8/2^{3/2}=2\sqrt2$.  $\square$

\noindent\emph{Proof of \Cref{prop:pb-kl-divergence-jets}.}
Invariance is \eqref{eq:cg-kl-bimeasurable-invariance} for the sample maps $\rho(g)$.
With local parameters $\lambda$ and $\lambda'$ of the two arguments, primed derivatives
acting on $\lambda'$, and $\mathbf s_i=\partial_i\log p_\lambda$, the identities
$\E_\lambda[\mathbf s_i]=0$, $\partial_ip_\lambda=p_\lambda\mathbf s_i$, and
$\partial_i\partial_jp_\lambda=p_\lambda(\mathbf s_i\mathbf s_j+\partial_i\mathbf s_j)$
and differentiation at $\lambda'=\lambda$ under \cref{hyp:pb-regular-models} give
\begin{align}
-\partial_i\partial_j'\KL(p_\lambda\Vert p_{\lambda'})
&=\E_\lambda[\mathbf s_i\mathbf s_j],\notag\\
-\partial_i\partial_j\partial_k'\KL(p_\lambda\Vert p_{\lambda'})
&=\E_\lambda[(\mathbf s_i\mathbf s_j+\partial_i\mathbf s_j)\mathbf s_k],
\notag\\
-\partial_i'\partial_j'\partial_k\KL(p_\lambda\Vert p_{\lambda'})
&=\E_\lambda[\mathbf s_k\partial_i\mathbf s_j].
\label{eq:pb-kl-jet-calculation}
\end{align}
The second derivatives of the score are symmetric, so the last two lines differ by
$\E_\lambda[\mathbf s_i\mathbf s_j\mathbf s_k]$, and positivity of the metric jet is
positivity of $g^F$.  The same dualistic structure is derived from contrast functions in
\citet{AmariNagaoka2000,AyJostLeSchwachhofer2017}.  $\square$

\noindent\emph{Proof of \Cref{cor:pb-orbit-fisher-form}.}
Positive definiteness of $g^F$ identifies the radical of $\mathfrak k_q$ with the kernel
of $\xi\mapsto\xi_{\mathcal B}(q)$, which is $\mathfrak g_q$.  The curve
$\gamma(t)=\widehat\rho(\exp t\xi)q$ is the integral curve of $\xi_{\mathcal B}$ through
$q$, so by uniqueness of integral curves $\xi_{\mathcal B}(q)=0$ exactly when
$\gamma\equiv q$, that is, when $\exp t\xi\in G_q$ for every $t$; this gives the
vanishing clause, and since $G_q$ is closed, its Lie algebra is the set of $\xi$ with
$\exp t\xi\in G_q$ for every $t$, which is $\mathfrak g_q$.  The
function $\phi(t)=\mathscr D(q,\gamma(t))$ is $C^3$ with $\phi(0)=\phi'(0)=0$ because
$\mathscr D(q,\cdot)$ is minimal at $q$, and differentiating
$\partial'_j\mathscr D(\lambda,\lambda)=0$ along the diagonal gives
$\partial'_i\partial'_j\mathscr D=-\partial_i\partial'_j\mathscr D=(g_{\mathscr D})_{ij}$
there, so $\phi''(0)=g_{\mathscr D}(\xi_{\mathcal B}(q),\xi_{\mathcal B}(q))$ and
Taylor's theorem gives the expansion; for relative entropy $g_{\mathscr D}=g^F$ by
\cref{prop:pb-kl-divergence-jets}.  $\square$

\noindent\emph{Proof of \Cref{thm:cg-fisher-contraction}.}
Pollard's Le~Cam decomposition at $\lambda_0$ is
$P_{\lambda_0+\delta}=\mathsf p_\delta P_{\lambda_0}+P_\delta^\perp$ with
$P_\delta^\perp\perp P_{\lambda_0}$, $P_\delta^\perp(\mathsf Y)=o(\lVert\delta\rVert^2)$,
and $\int(\sqrt{\mathsf p_\delta}-1-\tfrac12\delta^\top\mathbf s_{\lambda_0})^2
dP_{\lambda_0}=o(\lVert\delta\rVert^2)$.  Hence
$P_{\lambda_0+\delta}\otimes K=\mathsf p_\delta(y)(P_{\lambda_0}\otimes
K)+P_\delta^\perp\otimes K$, where the second term has mass $o(\lVert\delta\rVert^2)$ and
is singular, because a $P_{\lambda_0}$-null set carrying $P_\delta^\perp$, multiplied by
$\bar{\mathsf Y}$, carries it and is $P_{\lambda_0}\otimes K$-null since each
$K(y,\cdot)$ is a probability, and the square-root remainder of the first term is
unchanged because the kernel integrates to one.  The joint experiment is therefore
differentiable in quadratic mean with score $\mathbf s_{\lambda_0}(Y)$, without a common
dominating measure, and Pollard's preservation theorem for measurable statistics
\citep[Theorem~7]{Pollard2013}, applied to $(y,\bar y)\mapsto\bar y$, gives
differentiability of $\{P_\lambda K\}$ with the score \eqref{eq:cg-score-projection} in
$L^2(P_{\lambda_0}K;\R^d)$; the projected singular term may acquire an absolutely
continuous part, of mass still $o(\lVert\delta\rVert^2)$.  The law of total covariance
gives \eqref{eq:cg-fisher-loss}, and since conditional expectation is an orthogonal
projection, equality in the direction $z$ holds exactly when
$z^\top\mathbf s_{\lambda_0}(Y)$ is $\bar Y$-measurable.  $\square$

\noindent\emph{Proof of \Cref{thm:hist-retained-fisher-quotient}.}
Since $h_Q(Z,\cdot)=g(d\operatorname{ret}_QZ,d\operatorname{ret}_Q\cdot)$, the
Cauchy--Schwarz inequality for the semidefinite $g$ makes $h_Q(Z,Z)=0$,
$h_Q(Z,\cdot)=0$, and $d\operatorname{ret}_QZ\in\operatorname{rad}g$ equivalent, which is
\eqref{eq:hist-cds-fisher-radical} and the transversality criterion.  The radical is the
kernel of the constant-rank $C^1$ bundle map $h^\flat:T\mathcal Q\to T^*\mathcal Q$,
hence a $C^1$ subbundle, and $h$ descends to the quotient because it annihilates its
radical.  When $\operatorname{rad}h=\ker d\operatorname{ret}$, $\operatorname{ret}$ has
zero derivative along, hence is constant on, each connected leaf and factors as $\bar r\circ\pi$ through the leaf space,
with $\pi$ a submersion, $\ker d\pi=\ker d\operatorname{ret}$, and
$\bar r=\operatorname{ret}\circ\iota$ smooth for smooth local sections $\iota$ of $\pi$;
then $h=\pi^*(\bar r^*g)$, $d\bar r$ is injective, and
$\operatorname{im}d\bar r=\operatorname{im}d\operatorname{ret}$ meets
$\operatorname{rad}g$ only in zero, so $\bar r^*g$ is positive definite.  For a finite
family, $h=\sum_y\partial p_y\partial p_y/p_y$ is the Gram matrix of the Jacobian of
$\operatorname{ret}$ weighted by $1/p_y$, hence $C^2$ and positive definite exactly
where the Jacobian has full column rank; an immersion is locally injective because the
inverse function theorem \citep[Ch.~4]{Lee2012} makes a coordinate projection of
$\operatorname{ret}$ with invertible square Jacobian block a local diffeomorphism, and
the Levi-Civita connection of a $C^2$ metric is $C^1$.  For softmax,
$Z^\top h(z)Z=\operatorname{Var}_p(Z)$ vanishes exactly on constant vectors, so
$\operatorname{rad}h=\R\mathbf 1=\ker d\operatorname{ret}$ with constant rank $k-1$, and
an objective of $p$ has $d\Fenergy(\mathbf 1)=0$, so
$\nabla\Fenergy\in\mathbf 1^\perp=\operatorname{im}h$, which gives the solution class and
$h^+\nabla\Fenergy\in\operatorname{im}h=\mathbf 1^\perp$.  $\square$

\noindent\emph{Proof of \Cref{prop:hist-general-score-residual}.}
The tower property gives $\E[\mathbf s\mathbf s_i^\top]=\mathsf G_i$ and its transpose,
so expanding $\mathfrak K$ gives $\mathfrak K=\mathsf G_{\mathrm J}-\sum_i\mathsf
G_i+\mathfrak C$, which rearranges to \eqref{eq:hist-joint-weighted-marginal-residual};
the criterion is its restriction, and \eqref{eq:hist-joint-versus-product} is its
quadratic form for $n=2$ and unit weights.  Under independence the joint score is the
sum of the block scores, so $\mathbf r=0$, $\mathfrak C=0$, and the weight condition is
$(1-w_1)\mathsf G_1+(1-w_2)\mathsf G_2=0$ in the stated direction.  For the normal
location family with precision $\Lambda$ the joint Fisher metric is $\Lambda$ and each
marginal is a location family of variance $1/(1-\mathrm c^2)$, so
$\mathsf G_1+\mathsf G_2=(1-\mathrm c^2)\mathrm I_2$, which gives the displayed difference
and its two values.  The lift metrics are the pullbacks of $d\lambda_1^2+d\lambda_2^2$
along the three maps.  $\square$

\noindent\emph{Proof of \Cref{thm:hist-three-agent-parity-lift}.}
Each atom has $p_0/\mathrm D=\prod_{y_{i,x}=1}(1-\lambda_{i,x})^{-1}
\prod_{y_{i,x}=0}\lambda_{i,x}^{-1}>1$, so
$p_\varkappa\geq p_0-|\varkappa|\mathrm D>(1-|\varkappa|)\mathrm D>0$.  Summing out any
one coordinate cancels the parity term, because $\mathrm D$ does not depend on $y$ and
$\operatorname{par}$ changes sign with that coordinate; this gives normalization and the
product proper marginals, hence $\operatorname{mar}\circ\iota_\varkappa=\operatorname{id}$,
which with polynomial dependence gives smoothness and injectivity of
$d\iota_\varkappa$.  The marginal map is the restriction to the open $63$-dimensional
simplex of a linear map of rank six, so its fibers are open subsets of $57$-dimensional
affine spaces, and $p_\varkappa$ is affine in $\varkappa$.  A paired complement changes
two bits, so it preserves $\operatorname{par}$ and $\mathrm D$ while $p_0$ transforms by
Bernoulli relabeling, and a typed permutation preserves all three expressions, which
proves \eqref{eq:hist-parity-paired-equivariance}; a single complement changes the sign
of $\operatorname{par}$ and sends $\lambda_{i,x}$ to $1-\lambda_{i,x}$, which preserves
$\mathrm D$ and proves \eqref{eq:hist-parity-odd-pseudoscalar}.  $\square$

\noindent\emph{Proof of \Cref{cor:hist-parity-fisher-residual}.}
Two-bit marginals are products by \cref{thm:hist-three-agent-parity-lift}, so centered
scores of distinct bits are uncorrelated.  With $\mathbf s_{i,x}$ the one-bit Bernoulli
score, differentiation gives the residual components
\begin{equation}
\mathbf r_{i,x}
=\frac{\varkappa\operatorname{par}(y)\mathrm D(\lambda)}{p_\varkappa(y;\lambda)}
\left(\frac{1-y_{i,x}}{\lambda_{i,x}}-\frac{y_{i,x}}{1-\lambda_{i,x}}\right).
\label{eq:hist-parity-score-residual-coordinate}
\end{equation}
If $Z^\top\mathbf r$ vanished on every atom, positivity and $\varkappa\mathrm D\ne0$
would make the $Z$-weighted sum of the six bracketed one-bit functions vanish on the
cube, and comparing two atoms that differ in one coordinate forces each component of $Z$
to vanish; so $\mathfrak K\succ0$.  At the center $\mathrm D=4^{-6}$, its derivative
vanishes, and the law is \eqref{eq:hist-parity-center-law},
so, each one-bit score having square $4$, averaging
$(1+\bar\varkappa\operatorname{par})^{-1}=(1-\bar\varkappa\operatorname{par})/(1-\bar\varkappa^2)$
over the cube gives \eqref{eq:hist-parity-center-fisher}, and $\bar\varkappa=1/128$ gives
\eqref{eq:hist-parity-half-fisher-residual}.  The one-bit tensors at the center are
$4e_{i,x}e_{i,x}^\top$, which gives \eqref{eq:hist-parity-center-weight-equality}, and the
last sentence compares $\varkappa=0$ with $\varkappa\ne0$.  $\square$

\noindent\emph{Proof of \Cref{prop:hist-parity-hyperedge-vfe}.}
Half of the atoms have each parity, so averaging \eqref{eq:hist-parity-hyperedge-kernel}
under $\operatorname{Unif}$ gives evidence $\tfrac12$, Bayes' rule gives the posterior
$64^{-1}(1+\bar\varkappa\operatorname{par})$, which is \eqref{eq:hist-parity-center-law},
and the gap identity of the free energy gives the tight value $-\log\tfrac12$.  Every
posterior one-bit marginal is uniform, so its product law is $\operatorname{Unif}$, and
since the two parities have equal mass,
$-\E_{\operatorname{Unif}}\log K_{\mathrm{hyp}}(1\given Y)
=\log2-\tfrac12\log(1-\bar\varkappa^2)$, which is
\eqref{eq:hist-parity-product-lift-excess}.  The posterior and $\operatorname{Unif}$ have
the same one-bit marginals, the posterior input and the evidence are unchanged, and the
two free-energy values differ when $\varkappa\ne0$.  $\square$

\noindent\emph{Proof of \Cref{thm:hist-fisher-clock-invariance}.}
The chain rule gives $\dot{\widetilde Q}=(\dot Q\circ\phi)\phi'$ with $\phi'>0$, and the
substitution $t=\phi(\widetilde t)$ proves length invariance and
\eqref{eq:hist-clock-reparameterization-law}.  Positive speed makes the clock invertible,
and dividing by $d\mathfrak t/dt$ gives unit speed.  Substituting
\eqref{eq:hist-vfe-ray-representative} into \eqref{eq:hist-fisher-speed} gives
\eqref{eq:hist-vfe-clock-rate}, and the other two identities follow from
\eqref{eq:hist-fisher-gradient-definition}.  For a second metric,
$d\Fenergy(Z)=\mathsf G'(\operatorname{grad}^{\mathsf G'}\Fenergy,Z)$, and the
Cauchy--Schwarz inequality gives \eqref{eq:hist-clock-metric-dissipation-bound}.  For a
block-diagonal metric the squared speeds add, so
$\dot{\mathfrak t}_i\leq\dot{\mathfrak t}\leq\sum_i\dot{\mathfrak t}_i$, and Minkowski's
integral inequality for the vector $(\dot{\mathfrak t}_i)_i$ gives
$\bigl(\sum_i(\int\dot{\mathfrak t}_i)^2\bigr)^{1/2}\leq\int\dot{\mathfrak t}$.  $\square$

\noindent\emph{Proof of \Cref{prop:hist-analytic-finite-duration}.}
Along the history $\tfrac{d}{dt}\Fenergy(\gamma)=-\lVert\dot\gamma\rVert_{\mathsf
G}^2\leq0$, so $\Fenergy(\gamma(t))$ decreases to $\Fenergy_*=\Fenergy(Q_*)$; if it
reaches $\Fenergy_*$ at some $t_1$, then $\dot\gamma=0$ on $[t_1,\infty)$ and the
duration is that of a compact segment.  Otherwise choose a closed Euclidean ball $\bar B$
of radius $\delta_0$ around $Q_*$ in the chart on which
$\mathrm c_1\mathrm I\preceq\mathsf G\preceq\mathrm c_2\mathrm I$ with
$0<\mathrm c_1\leq\mathrm c_2$ and the {\L}ojasiewicz gradient inequality
$\lVert\nabla\Fenergy\rVert\geq\mathrm c_0|\Fenergy-\Fenergy_*|^{\mathrm q}$ holds with
$\mathrm c_0>0$ and $\mathrm q\in[0,1)$ \citep[Lemma~2.1]{absil2005convergence}.  While
$\gamma$ stays in $\bar B$,
$\lVert\dot\gamma\rVert_{\mathsf G}^2=\nabla\Fenergy^\top\mathsf G^{-1}\nabla\Fenergy
\geq\mathrm c_2^{-1}\lVert\nabla\Fenergy\rVert^2$, hence
\begin{equation}
\frac{d}{dt}\left(\Fenergy(\gamma)-\Fenergy_*\right)^{1-\mathrm q}
=-(1-\mathrm q)\left(\Fenergy(\gamma)-\Fenergy_*\right)^{-\mathrm q}
\lVert\dot\gamma\rVert_{\mathsf G}^2
\leq-(1-\mathrm q)\mathrm c_0\mathrm c_2^{-1/2}\lVert\dot\gamma\rVert_{\mathsf G},
\label{eq:hist-lojasiewicz-length-estimate}
\end{equation}
so, with $\mathrm C=((1-\mathrm q)\mathrm c_0\mathrm c_2^{-1/2})^{-1}$, a segment
$\gamma([t_1,t_2])\subseteq\bar B$ has Fisher length at most
$\mathrm C(\Fenergy(\gamma(t_1))-\Fenergy_*)^{1-\mathrm q}$ and Euclidean length at most
$\mathrm c_1^{-1/2}$ times that.  For $t_1$ with $\lVert\gamma(t_1)-Q_*\rVert<\delta_0/2$
and $\mathrm c_1^{-1/2}\mathrm C(\Fenergy(\gamma(t_1))-\Fenergy_*)^{1-\mathrm
q}<\delta_0/2$, a first exit from $\bar B$ after $t_1$ would give a segment in $\bar B$ of
Euclidean length both at least and less than $\delta_0/2$; so $\gamma$ stays in $\bar B$,
has finite length on $[t_1,\infty)$, and is Cauchy, hence converges to $Q_*$.  If $Q_*$
were noncritical, $\tfrac{d}{dt}\Fenergy(\gamma)$ would tend to
$-\lVert\operatorname{grad}^{\mathsf G}\Fenergy(Q_*)\rVert^2_{\mathsf G}<0$,
contradicting convergence of $\Fenergy(\gamma)$.  This is the argument of
\citet[Theorem~2.2]{absil2005convergence}, with its angle condition supplied by the
metric bounds.  In the smooth counterexample the Euclidean plane is the Fisher geometry
of the unit-variance normal location family, and a curve of finite length in the plane
converges, so a history that does not converge has infinite duration.  $\square$

\noindent\emph{Proof of \Cref{prop:hist-invariant-history-horizontal}.}
Differentiating $\Fenergy(\exp(t\xi)\cdot Q)=\Fenergy(Q)$ at $t=0$ gives
$\mathsf G(\operatorname{grad}^{\mathsf G}\Fenergy,\xi_{\mathcal Q}(Q))=0$, and for a
velocity $Z\perp V_Q$ the quantity $\lVert Z-\xi_{\mathcal Q}(Q)\rVert^2=\lVert
Z\rVert^2+\lVert\xi_{\mathcal Q}(Q)\rVert^2$ is minimal where $\xi_{\mathcal Q}(Q)=0$.
In the free and proper case the orbit space is a smooth manifold and $\pi$ a surjective
submersion \citep[Theorem~21.10]{Lee2012}.  With
$\operatorname{Hor}_Q=(\ker T_Q\pi)^\perp$, $T_Q\pi$ maps $\operatorname{Hor}_Q$
isomorphically onto $T_{\pi(Q)}(\mathcal Q/\mathbb G)$; define $\bar{\mathsf G}$ by
transporting $\mathsf G|_{\operatorname{Hor}_Q}$ through this isomorphism.  Each
$L_g:Q\mapsto g\cdot Q$ is an isometry with $\pi\circ L_g=\pi$, so $TL_g$ maps
$\operatorname{Hor}_Q$ onto $\operatorname{Hor}_{gQ}$ with $T_{gQ}\pi\circ TL_g=T_Q\pi$
and the transported inner products agree along orbits; $\bar{\mathsf G}$ is smooth
because, through a smooth local section $\iota_\pi$ of $\pi$ given by the rank theorem,
its coefficients are those of $\mathsf G$ on the horizontal projections of the pushed
coordinate fields, and it is the unique metric for which $\pi$ is a Riemannian
submersion.  The invariant $\Fenergy$ is constant on the fibers of $\pi$, so
$\bar\Fenergy=\Fenergy\circ\iota_\pi$ is well defined and $C^1$ with
$\Fenergy=\bar\Fenergy\circ\pi$; hence $d\Fenergy$ vanishes on $\ker T\pi$, the gradient
is horizontal, and for horizontal $Z$,
$\mathsf G(\operatorname{grad}^{\mathsf G}\Fenergy,Z)=d\bar\Fenergy(T\pi Z)
=\bar{\mathsf G}(\operatorname{grad}^{\bar{\mathsf G}}\bar\Fenergy,T\pi Z)$, which gives
the gradient identity.  Applying $T\pi$ to \eqref{eq:hist-vfe-ray-representative} shows
that $\pi\circ Q$ is a natural-gradient history of $\bar\Fenergy$, and since $T\pi$ is
isometric on horizontal vectors and the velocity is horizontal, the two Fisher speeds
agree.  $\square$

\noindent\emph{Proof of \Cref{thm:hist-global-clock-exactness}.}
Exact forms are closed and have zero periods; conversely, under
\eqref{eq:hist-global-clock-period-criterion}, integrating $\bm\alpha$ from a base point
along curves in $\mathcal U$ is path independent and defines a function with
$d\mathfrak t^\perp=\bm\alpha$.  Differentiating $\bm\alpha=-\mathsf n^{-1}d\Fenergy$
gives \eqref{eq:hist-normalized-form-curvature}.  Since $d\Fenergy\ne0$, the form
$d\mathsf n\wedge d\Fenergy$ vanishes exactly when $d\mathsf n$ is locally a multiple of
$d\Fenergy$, that is, when $\mathsf n$ does not depend on the other coordinates of a chart
whose first coordinate is $\Fenergy$.  Then $d\mathfrak t^\perp\wedge
d\Fenergy=-\mathsf n^{-1}d\Fenergy\wedge d\Fenergy=0$, so $\mathfrak t^\perp$ is constant
on connected pieces of level sets and locally $\mathfrak t^\perp=\tau\circ\Fenergy$ with
$\tau'=-1/\mathsf n<0$.  Since $\bm\alpha(\mathbf e_{\Fenergy})=\mathsf G(\mathbf
e_{\Fenergy},\mathbf e_{\Fenergy})=1$, the unit rate follows, and the kernel of
$\bm\alpha=\mathbf e_{\Fenergy}^\flat$ is the orthogonal complement of
$\mathbf e_{\Fenergy}$.  In the example $\mathsf n=(z_1^2+z_2^2)^{1/2}$ and
$\bm\alpha=-(z_2dz_1+z_1dz_2)/\mathsf n$, whose exterior derivative is
\eqref{eq:hist-nonexact-clock-example}, with coefficient vanishing only on the diagonal.
$\square$

\noindent\emph{Proof of \Cref{prop:hist-cross-section-clock}.}
The unit field $\mathbf e_{\Fenergy}$ is $C^1$ because $\mathsf n>0$.  Let
$\mathcal W\subseteq\R\times\mathcal S$ be the open set of pairs $(t,q)$ for which
$\operatorname{Fl}^{\mathbf e_{\Fenergy}}_{t}(q)$ is defined in $\mathcal U$, and put
$\Xi(t,q)=\operatorname{Fl}^{\mathbf e_{\Fenergy}}_{t}(q)$, a $C^1$ map.  Each orbit meets
$\mathcal S$, so $\Xi$ is onto; $\Fenergy$ decreases strictly along orbits, so no orbit
is closed, and each orbit meets $\mathcal S$ once, so $\Xi$ is injective.  At $(t,q)$,
$T\Xi$ sends $\partial_{t}$ to $\mathbf e_{\Fenergy}$ and $T_q\mathcal S$ onto the image
of a complement of $\mathbf e_{\Fenergy}(q)$ under the invertible map
$T\operatorname{Fl}^{\mathbf e_{\Fenergy}}_{t}$, which also carries
$\mathbf e_{\Fenergy}(q)$ to $\mathbf e_{\Fenergy}$ at the image point, so transversality
makes $T\Xi$ invertible.  Hence $\Xi$ is a $C^1$ diffeomorphism, and
$\mathfrak t^{\mathcal S}=\operatorname{pr}_1\circ\Xi^{-1}$ satisfies
$d\mathfrak t^{\mathcal S}(\mathbf e_{\Fenergy})=1$, with $t$ the Fisher arc length
because $\mathbf e_{\Fenergy}$ has unit speed.  A level set is transversal and met once
because $d\Fenergy(\mathbf e_{\Fenergy})=-\mathsf n<0$.  If the level sets of
$\mathfrak t^{\mathcal S}$ are orthogonal to $\mathbf e_{\Fenergy}$, then
$d\mathfrak t^{\mathcal S}$ vanishes on the orthogonal complement and equals one on
$\mathbf e_{\Fenergy}$, so it equals $\bm\alpha$.  In the example the Euclidean descent
orbits $\dot z_1=-z_2$, $\dot z_2=-z_1$ in the quadrant are
$z_1=\mathrm c_1e^t+\mathrm c_2e^{-t}$, $z_2=-\mathrm c_1e^t+\mathrm c_2e^{-t}$ with
$\mathrm c_2>|\mathrm c_1|$, along which $\Fenergy=\mathrm c_2^2e^{-2t}-\mathrm
c_1^2e^{2t}$ decreases from $+\infty$ toward $0$, so each meets $\{z_1z_2=1\}$ exactly
once.  $\square$

\noindent\emph{Proof of \Cref{thm:hist-record-clock-contraction}.}
At each $t$, \cref{thm:cg-fisher-contraction} applied to the one-parameter family
$\{P_{t'}\}$ gives \eqref{eq:hist-record-score-projection} and the defect identity
\eqref{eq:hist-record-conditional-variance-defect} with its equality clause; taking
nonnegative square roots and integrating gives \eqref{eq:hist-record-clock-contraction}.
For a dominated channel, differentiating the joint density
$p_t(y)\mathsf k_t(\bar y\given y)$ gives the joint score
$\mathbf s_t(y)+\partial_t\log\mathsf k_t(\bar y\given y)$, whose projection onto the
record is \eqref{eq:hist-parameter-dependent-channel-score}.  In the example the fine
law does not move, and a Bernoulli law of success probability
$\mathrm b(t)=(1+e^{-t})^{-1}$ has Fisher speed
$|\dot{\mathrm b}|/\sqrt{\mathrm b(1-\mathrm b)}=\sqrt{\mathrm b(1-\mathrm b)}=(2\cosh(t/2))^{-1}$.
$\square$

\noindent\emph{Proof of \Cref{thm:hist-oriented-semiconjugacy}.}
Put $\gamma(t)=\mathsf F_\ell(\operatorname{Fl}^{\mathscr V_\ell}_tQ)$, so
$\dot\gamma=\mathsf a_\ell(\operatorname{Fl}^{\mathscr V_\ell}_tQ)\mathscr
V_{\ell+1}(\gamma)$.  The integrand of $\psi_Q$ is continuous and positive, so $\psi_Q$ is
a $C^1$ diffeomorphism onto an open interval.  The curve $\gamma\circ\psi_Q^{-1}$
satisfies $(\gamma\circ\psi_Q^{-1})'=\mathscr V_{\ell+1}(\gamma\circ\psi_Q^{-1})$ on
$\mathcal T'_Q$ and starts at $\mathsf F_\ell Q$, so by local Lipschitz continuity it is
contained in the maximal integral curve through $\mathsf F_\ell Q$; this gives the
inclusion and the flow identity, with the curve constructed before the coarse flow is
evaluated at $\psi_Q(t)$.  The endpoint formula is monotonicity of $\psi_Q$, and in the
complete case the limits are the two integrals.  In the examples
$\mathsf a_\ell=\mathsf F_\ell'$ because both fields are $\partial_z$.  The converse is
the chain rule.  $\square$

\noindent\emph{Proof of \Cref{prop:hist-natural-gradient-sufficiency}.}
In the first example, \eqref{eq:hist-oriented-vector-field-semiconjugacy} reads
$(z_1,2z_2)=\mathsf a(z_1,2z_2/\mathrm c)$, which forces $\mathsf a=1$ and $\mathrm c=1$
where $z_1z_2\ne0$.  Differentiating the compatibility identity gives
$d\Fenergy_\ell(Z)=\chi_\ell'd\Fenergy_{\ell+1}(T\mathsf F_\ell Z)=0$ for
$Z\in\ker T\mathsf F_\ell$, so $\Fenergy_\ell$ is constant along curves tangent to the
fibers, $Z_0=\operatorname{grad}^{\mathsf G_\ell}\Fenergy_\ell$ lies in
$\operatorname{Hor}_Q$, and for $Z\in\operatorname{Hor}_Q$
\[
\mathsf G_{\ell+1}(T\mathsf F_\ell Z_0,T\mathsf F_\ell Z)
=\mathsf d_\ell^2d\Fenergy_\ell[Z]
=\mathsf d_\ell^2\chi_\ell'\mathsf G_{\ell+1}\left(\operatorname{grad}^{\mathsf
G_{\ell+1}}\Fenergy_{\ell+1}\circ\mathsf F_\ell,T\mathsf F_\ell Z\right).
\]
Since $T\mathsf F_\ell(\operatorname{Hor}_Q)$ is the whole tangent space,
$T\mathsf F_\ell Z_0=\mathsf d_\ell^2\chi_\ell'\operatorname{grad}^{\mathsf
G_{\ell+1}}\Fenergy_{\ell+1}\circ\mathsf F_\ell$, and negating gives the claim.  The
quotient case is $\mathsf d_\ell\equiv1$ and $\chi_\ell=\operatorname{id}$.  In the
normal example the kernel sends $\mathcal N(\mu,1)$ to $\mathcal N(\mu,2)$, the two
Fisher metrics are $d\mu^2$ and $\tfrac12d\mu^2$, and $\mathscr V_\ell=-\mu$,
$\mathscr V_{\ell+1}=-2\mu$.  For the exact-contraction objective,
$\Fenergy_{\ell+1}\circ\mathsf F_\ell\leq\Fenergy_\ell$ with equality exactly on the set
$\mathcal A$ of fiberwise minimizers, so compatibility with
$\chi_\ell=\operatorname{id}$ holds exactly on $\mathcal A$; an open subset of
$\mathcal A$ would contain, with each of its points, a neighborhood of that point in its
fiber, which has more than one point, so $\mathcal A$ has empty interior.  In the last
example $\mathcal A=\{f=s\}$ and $d\Fenergy_\ell(\partial_f)=f-s$, which vanishes on no
open set, so the fiber condition fails for every $\chi_\ell$.  $\square$

\section{Proofs for Coarse Channels}
\label{sec:proofs-coarse-channels}

\noindent\emph{Proof of \Cref{prop:cg-finest-independent-partition}.}
Suppose $P$ factors over partitions $\mathfrak p$ and $\mathfrak q$. Marginalizing
the second factorization to a block $I\in\mathfrak p$ gives
\begin{equation*}
P_I=\bigotimes_{J\in\mathfrak q:I\cap J\ne\varnothing}P_{I\cap J}.
\end{equation*}
Substitution in $P=\bigotimes_{I\in\mathfrak p}P_I$ proves factorization over the
common refinement. The finite set of independent partitions is nonempty, since it
contains $\{V\}$, and is closed under common refinement. Its common finest refinement
is independent and uniquely refines every member. Every independent partition has
negative objective in \eqref{eq:cg-independent-selector-objective}, while every
dependent partition has positive objective. Among the former the finest has the
unique largest block count. If it is intermediate, refinement from $\{V\}$ toward it
strictly decreases the objective; any further strict refinement is dependent and
raises the objective across zero. Relabeling permutes the marginal products, and
separate bimeasurable changes push each factor to its corresponding marginal, so
independence and its finest partition have the asserted naturality and invariance.

For the conditional assertion there are finitely many partitions and pairs of
partitions. For each one with $\mathsf T_{P\mid\Upsilon}=0$, its conditional product
identity holds on a set of full $P_\Upsilon$-mass. Intersect these finitely many sets.
On that intersection the preceding marginalization argument proves closure under
common refinement, hence the unique finest conditional partition and its objective
characterization. The conditional laws given $\Upsilon=Y_V$ are point masses and
therefore products over every partition. These constructions choose a factorization
objective; no identity with the evidence bound has been introduced. $\square$

\noindent\emph{Proof of \Cref{prop:cg-independent-selector-instability}.}
The exponent in \eqref{eq:cg-independent-pair-law} is a sum of terms on disjoint
pairs, so its normalized law is the product of its pair marginals. Within a pair
the odds ratio is $e^{4\mathrm j}\ne1$, so the pair is not independent and cannot
split in an independent partition. This gives its finest partition. For
\eqref{eq:cg-independent-pair-perturbation} all edges of the path have positive
coupling. Any nontrivial partition cuts an edge $ij$. Conditional on all other
coordinates, the two-spin odds ratio across this edge is $e^{4\mathrm j_{ij}}\ne1$,
where $\mathrm j_{ij}$ is $\mathrm j$ or $\mathrm r$. A product over the partition
would instead make the two spins conditionally independent, with ratio one. Positivity
of the finite law permits conditioning on every outside assignment. Thus no nontrivial
independent partition exists. Its finitely many normalized atom probabilities are
continuous in $\mathrm r$ down to zero, which gives total-variation convergence. $\square$

\noindent\emph{Proof of \Cref{prop:cg-exchangeable-hard-selector}.}
Because every agent permutation leaves the input law unchanged, naturality forces the
output unlabeled partition to be invariant under every permutation. If it contained
a block $I$ with $2\leq|I|\leq N-1$, choose distinct $i,j\in I$ and $k\notin I$.
The transposition of $j,k$ fixes $i$ and sends $I$ to a different block containing $i$.
Two distinct blocks of a partition cannot intersect. Therefore no such block exists,
leaving only the singleton and one-block partitions. The argument concerns a
deterministic partition from this law alone and imposes no restriction on a
permutation-invariant probability distribution over partitions. $\square$

\noindent\emph{Proof of \Cref{prop:cg-partition-information-stability}.}
For a partition $\mathfrak p$ of $I$, the entropy numerator satisfies
\begin{equation*}
\sum_{J\in\mathfrak p}\mathrm H(P_J)-\mathrm H(P_I)
=\KL\Bigl(P_I\Bigm\Vert\bigotimes_{J\in\mathfrak p}P_J\Bigr)
\end{equation*}
by \cref{thm:elbo-total-correlation-chain}. It is nonnegative and is zero exactly at a
product. The minimum ranges over a finite set, so its zero case is a nontrivial
independent partition.

For each $J\subseteq I$, marginalization gives
$\operatorname{TV}(P_J,Q_J)\leq\delta$. Couple these two finite-alphabet variables
with mismatch probability $\delta_J\leq\delta$. Writing the mismatch indicator as
an auxiliary binary coordinate and applying the entropy chain rule gives
\begin{equation*}
|\mathrm H(P_J)-\mathrm H(Q_J)|
\leq h_2(\delta_J)+\delta_J\log D_J
\leq h_2(\delta)+\delta\log D_J.
\end{equation*}
The last step uses $\delta\leq1/2$. A singleton alphabet contributes zero. For a
partition with $k\geq2$ blocks, $\sum_{J\in\mathfrak p}\log D_J=\log D_I$;
therefore its normalized entropy expression changes by at most
\begin{equation*}
\frac{(k+1)h_2(\delta)+2\delta\log D_I}{k-1}
\leq3h_2(\delta)+2\delta\log D_I.
\end{equation*}
The minima of finitely many paired expressions differ by at most the largest
paired difference. This proves \eqref{eq:cg-partition-information-stability}.
The bound with $D_V\geq D_I$ proves the two threshold inclusions. A score at distance
greater than $\eta$ from the threshold cannot cross it, so the entire qualified
subset family and its maximal members are unchanged. Finally, the relative-entropy
representation of each numerator is invariant under separate bimeasurable changes
by \cref{thm:cg-dpi-equality} applied in both directions. The finite-alphabet entropy
bound supplies no corresponding bound on unrestricted continuous laws. $\square$

The proofs for \cref{sec:cg-retained-label-asymptotics} use the notation fixed at the head of
that section, including its convention on the constant $\mathrm C_{\mathsf K}$.

\noindent\emph{Proof of \Cref{prop:cg-descent-common-recovery}.}
For a law $P\ll P_0$ with density $r$ and bounded measurable $g$ on $\mathsf Y$,
$\int g d(PK)=\E_{\mathbb J_{P_0}}[g(Y)r(X)]=\E_{\mathbb J_{P_0}}[g(Y)\bar r(Y)]$ with
$\bar r(Y)=\E_{\mathbb J_{P_0}}[r(X)\given Y]$, because $r$ is $P_0$-integrable. This is the
computation in the proof of \Cref{thm:cg-dpi-equality}, which uses no finiteness of relative entropy,
and it shows that a measurable function of the coarse coordinate that agrees with $r$ almost everywhere
is $\bar r=d(PK)/d(P_0K)$. (i) If $r$ descends, disintegrating $\mathbb J_{P_0}$ over $y$ gives
$(PKR_{P_0})(A)=\int\bar r(y)R_{P_0}(y,A)(P_0K)(dy)=\E_{\mathbb J_{P_0}}[\mathbf 1_A(X)\bar r(Y)]
=\E_{\mathbb J_{P_0}}[\mathbf 1_A(X)r(X)]=P(A)$, and $r=1$ gives $P_0KR_{P_0}=P_0$. Conversely, let one
Markov kernel $R$ recover $P_0$ and $P$. Conditional Jensen for the concave square root gives
$\E_{\mathbb J_{P_0}}[\sqrt{r(X)}\given Y]\leq\sqrt{\bar r(Y)}$ almost surely, so the affinity
$\E_{P_0}[\sqrt r]$, which is at most $1$ by the Cauchy-Schwarz inequality, does not decrease through
$K$: $\E_{P_0}[\sqrt r]\leq\E_{P_0K}[\sqrt{\bar r}]$. The kernel $R$ sends the pair $(PK,P_0K)$, of density
$\bar r$, to $(P,P_0)$, of density $r$, so the same inequality through $R$ gives
$\E_{P_0K}[\sqrt{\bar r}]\leq\E_{P_0}[\sqrt r]$. The two finite affinities are therefore equal, the
conditional inequality is an almost sure equality, and
\begin{equation*}
\E_{\mathbb J_{P_0}}\bigl[(\sqrt{r(X)}-\sqrt{\bar r(Y)})^2\bigr]
=2-2\E_{\mathbb J_{P_0}}\bigl[\sqrt{\bar r(Y)}\E_{\mathbb J_{P_0}}[\sqrt{r(X)}\given Y]\bigr]
=2-2\E_{\mathbb J_{P_0}}[\bar r(Y)]=0,
\end{equation*}
where $\sqrt{r(X)\bar r(Y)}$ is integrable because it is at most $(r(X)+\bar r(Y))/2$. Hence $r$
descends. Applied to each $P_\theta$, and to a single law $P\ll P_0$, this proves the first sentence of
(i) and the characterization of the laws that $R_{P_0}$ recovers. The functions in $L^1(P_0)$ that
agree $\mathbb J_{P_0}$-almost everywhere with a measurable function of the coarse coordinate form a
linear subspace containing $1$ and, when every $r_\theta$ descends, every $r_\theta$; it is closed, so
it then contains their closed linear span: if $r^{(n)}\to r$ in $L^1(P_0)$ and $r^{(n)}(X)=\bar r^{(n)}(Y)$ almost surely, then, since $X$
has law $P_0$ under
$\mathbb J_{P_0}$, a subsequence converges $\mathbb J_{P_0}$-almost surely, and the limit superior of
$\bar r^{(n)}$ along it, set to zero where it is infinite, is a measurable function of $y$ equal to
$r(x)$ almost everywhere. For records, $d\bar{\mathsf m}_o/d\varrho=L_o/M_o$ by
\Cref{def:rg-measure-pair}, and $d(\bar{\mathsf m}_oC)/d(\varrho C)=L^c_o/M_o$ by
\eqref{eq:rg-reference-evidence-pair}. (ii) The graph $\operatorname{gr}c=\{(x,y):c(x)=y\}$ is
measurable because the diagonal of the standard Borel space $\mathsf Y$ is, so
$y\mapsto R(y,c^{-1}(\{y\}))=\int\mathbf 1_{\operatorname{gr}c}(x,y)R(y,dx)$ is measurable. Under $K_c$
the joint $\mathbb J_{P_0}$ is the law of $(X,c(X))$ and is carried by $\operatorname{gr}c$, which gives
the stated form of descent, and $\int R_{P_0}(y,c^{-1}(\{y\}))(c_\#P_0)(dy)=\mathbb J_{P_0}(\operatorname{gr}c)=1$,
so every version of $R_{P_0}$ is supported on fibers and recovers $P_0$. If $R$ is supported on fibers
and recovers $P_0$, the measure on $\mathsf X\times\mathsf Y$ that gives $A\times B$ the mass
$\int_BR(y,A)(c_\#P_0)(dy)$ is carried by $\operatorname{gr}c$ and has first marginal $P_0K_cR=P_0$, so
its value on $A\times B$ equals its value on $(A\cap c^{-1}(B))\times\mathsf Y$, which is
$P_0(A\cap c^{-1}(B))=\mathbb J_{P_0}(A\times B)$. Hence $R$ disintegrates $\mathbb J_{P_0}$ over $y$ and
is a version of $R_{P_0}$, and two versions agree $c_\#P_0$-almost everywhere because $\mathsf X$ is
countably generated. Since $P_\theta K_c\ll P_0K_c$, all versions give the same $P_\theta K_cR$, and (i)
gives the rest. (iii) The fiber of $c\times\operatorname{id}_{\mathsf X_2}$ over $(y,x_2)$ is
$c^{-1}(\{y\})\times\{x_2\}$, and the first marginal of $(c\times\operatorname{id}_{\mathsf X_2})_\#P_0$
is the law of $c(X_1)$, so $R$ is supported on fibers and, by (ii), recovers $P_0$ and every $P_\theta$
exactly when it recovers $P_0$ and every $r_\theta$ descends. For measurable $A_1\subseteq\mathsf X_1$,
$A_2\subseteq\mathsf X_2$, and $B\subseteq\mathsf Y$, fiber support gives
$R'(c(X_1),A_1\cap c^{-1}(B))=R'(c(X_1),A_1)\mathbf 1_B(c(X_1))$ almost surely, and measurable
rectangles generate the product $\sigma$-algebra, so $R$ recovers $P_0$ exactly when
\begin{equation*}
P_0\bigl(X_1\in A_1,c(X_1)\in B,X_2\in A_2\bigr)
=\E_{P_0}\bigl[R'(c(X_1),A_1)\mathbf 1_B(c(X_1))\mathbf 1_{A_2}(X_2)\bigr]
\end{equation*}
for all such sets. With $A_2=\mathsf X_2$ this says that $R'$ is a regular conditional law of $X_1$ given
$c(X_1)$. Granted that, conditioning both sides on $c(X_1)$ shows that the identity is equivalent to
$P_0(X_1\in A_1,X_2\in A_2\given c(X_1))=P_0(X_1\in A_1\given c(X_1))P_0(X_2\in A_2\given c(X_1))$
almost surely for all $A_1,A_2$, which is the conditional independence. These conditions on $P_0$
and the $r_\theta$ are equivalent to adequacy in the sense named after the proposition, that is, to a
conditional law of $X_1$ given $c(X_1)$ common to $P_0$ and every $P_\theta$ together with the
conditional independence of $X_1$ and $X_2$ given $c(X_1)$ under each of these laws. Conditional
independence under $P_0$ with descent $r_\theta(x_1,x_2)=\bar r_\theta(c(x_1),x_2)$ makes the
conditional law of $X_1$ given $(c(X_1),X_2)$ under $P_\theta$, whose density with respect to $P_0$ is a
function of $(c(X_1),X_2)$, equal to that under $P_0$, which is the conditional law of $X_1$ given
$c(X_1)$ under $P_0$, so the conditional independence holds under every $P_\theta$ with this common
conditional law; conversely, adequacy writes each member of the family as the law of
$(c(X_1),X_2)$ followed by the common kernel to $X_1$, so every $r_\theta$ is the density between
these laws of $(c(X_1),X_2)$ evaluated at $(c(x_1),x_2)$ and descends. In the witness every $R'$ is
supported on the one fiber, the density $1$ of the family descends,
$R_{P_0}((y,x_2),\cdot)=\delta_{(x_2,x_2)}$ recovers $P_0$,
and $R'(y,\cdot)\otimes\delta_{x_2}$ sends $P_0$ to the product of $R'(y,\cdot)$ with the uniform law on
$\{0,1\}$, which gives $\{x_1\neq x_2\}$ the mass $\tfrac12R'(y,\{1\})+\tfrac12R'(y,\{0\})=\tfrac12$.
$\square$

\noindent\emph{Proof of \Cref{thm:rg-exact-coarse-vfe}.}
(i) Since $C(y,\mathsf Z)=1$, a bounded function of $o$ alone has the same integral against
$\mathbb P^c$ and against $\mathbb P$, which gives normalization and the observation marginal. For
bounded measurable $f$ on $\mathsf Z$ and $g$ on $\mathsf O$ put $h(y)=\int f(z)C(y,dz)$. Then
\begin{align*}
\int g\Bigl(\int f d(\boldsymbol\Pi_oC)\Bigr)\mathbb P^c(do\times\mathsf Z)
&=\int g\Bigl(\int h d\boldsymbol\Pi_o\Bigr)\mathbb P(do\times\mathsf Y)\\
&=\int g(o)h(y)\mathbb P(do,dy)
=\int g(o)f(z)\mathbb P^c(do,dz),
\end{align*}
using the preserved marginal, the defining property of the fine version, and
\eqref{eq:rg-common-pushforward}. Hence $o\mapsto\boldsymbol\Pi_oC$ is a posterior version of
$\mathbb P^c$ at every observation where it is declared, because $C$ acts on the conditioned
variable only. If $\boldsymbol\Pi^c_o(N)=0$, then $C(\cdot,N)=0$ $\boldsymbol\Pi_o$-almost surely,
hence $\mathbb Q_o$-almost surely under $\mathbb Q_o\ll\boldsymbol\Pi_o$, and $\mathbb Q^c_o(N)=0$.
(ii) Disintegrated over $y$, both lifts have the marginals $\mathbb Q_o$ and $\boldsymbol\Pi_o$ and
the common conditional $C(y,\cdot)$, so the chain rule of \Cref{lem:prob-kl-chain-rule}, in
$[0,+\infty]$, gives $\KL(\widehat{\mathbb Q}_o\Vert\widehat{\boldsymbol\Pi}_o)
=\KL(\mathbb Q_o\Vert\boldsymbol\Pi_o)$; disintegrated over $z$, it gives
$\KL(\widehat{\mathbb Q}_o\Vert\widehat{\boldsymbol\Pi}_o)
=\KL(\mathbb Q^c_o\Vert\boldsymbol\Pi^c_o)+\Delta(o)$. Adding the finite real $-\log p(o)$, the
same at both scales by (i), and using that $\boldsymbol\Pi^c_o$ is the coarse posterior version gives
\eqref{eq:rg-vfe-chain-rule}; every term lies in $[0,+\infty]$ before the finite shift, so no
infinite quantity is subtracted. (iii) $\Delta(o)$ is the integral of a nonnegative measurable
function, which vanishes exactly when the integrand vanishes $\mathbb Q^c_o$-almost surely, and a
relative entropy vanishes exactly when its arguments agree. (iv) Finite fine divergence makes both
terms on the right of \eqref{eq:rg-vfe-chain-rule} finite. If $\Delta(o)=0$, the two divergences
are equal and finite, and \Cref{thm:cg-dpi-equality,cor:cg-pairwise-bayes-recovery} give recovery
of both laws by the Bayes kernel of $\boldsymbol\Pi_o$ under $C$. Conversely, a kernel $R$
recovering both laws gives $\KL(\mathbb Q_o\Vert\boldsymbol\Pi_o)\geq\KL(\mathbb Q^c_o\Vert\boldsymbol\Pi^c_o)
\geq\KL(\mathbb Q^c_oR\Vert\boldsymbol\Pi^c_oR)=\KL(\mathbb Q_o\Vert\boldsymbol\Pi_o)$ by data
processing through $C$ and then $R$, so $\Delta(o)=0$. $\square$

\noindent\emph{Proof of \Cref{thm:rg-pointwise-parent-datum}.}
(i) is \Cref{thm:rg-exact-coarse-vfe} with $\mathbb P=\mathbb P_I$, $C=C_A$, and the admitted
recognition law. (ii) On standard Borel spaces a disintegration of $\mathbb P_A(\cdot\given X)$
over $(m_A,\xi_A)$ exists and is jointly measurable (\Cref{ch:probability}); it satisfies the
compatibility identity by construction, while a family declared
independently of $\mathbb P_A$ satisfies it only if assumed. (iii) The pushforward formulas hold by
testing each parent law with functions of one coordinate. For two non-Dirac marginals,
\Cref{prop:prob-marginals-do-not-determine-joint} gives distinct couplings with the same pair, and
tensoring each with one fixed law on $\boldsymbol\Xi_A\times\mathsf H_A$ gives distinct laws on
$\mathsf Z_A$. For two points $h\neq h'$ of $\boldsymbol\Xi_A\times\mathsf H_A$ and the
$(b_A,m_A)$ marginal $R$ of the given law, $R\otimes\delta_h$ and $R\otimes\delta_{h'}$ share the
pair and differ, so one of them differs from the given law. If that space is a point and one
marginal is Dirac, every coupling assigns mass one to the Dirac coordinate. $\square$

\noindent\emph{Proof of \Cref{prop:cg-equivariant-bayes-versions}.}
(i) By equivariance,
$\mathbb J_{g_\#Q}(A\times B)=\mathbb J_Q(g^{-1}A\times g^{-1}B)$. Bimeasurability makes $R_Q^g$ a
Markov kernel, the $\mathsf Y$-marginal of $\mathbb J_{g_\#Q}$ is $g_\#(QK)$ by
\Cref{prop:cg-equivariant-channels}, and the change of variables $y=gy'$ gives
$\int\mathbf 1_B(y)R_Q^g(y,A)(g_\#(QK))(dy)=\mathbb J_{g_\#Q}(A\times B)$. (ii) With
$(g_\#P)K=g_\#(PK)$, $((g_\#P)KR_Q^g)(A)=(PKR_Q)(g^{-1}A)=(g_\#P)(A)$; the relative-entropy
identities are invariance under bimeasurable bijections, \eqref{eq:cg-rn-bimeasurable} and
\eqref{eq:cg-kl-bimeasurable-invariance}. (iii) Where $(QK)(\{y\})>0$ the first value is the value
of every version; the remaining points are $QK$-null, both branches transform as stated because
$(g_\#Q)(\{gx\})=Q(\{x\})$ and $K(gx,\{gy\})=K(x,\{y\})$, and the first branch is a rational
function of $Q$ on a Borel set. (iv) Each $R_Q^h$ with $h\in\mathbb G_Q$ is a version of $R_Q$ by
(i). Joint measurability makes $(h,y)\mapsto R_Q^h(y,A)$ measurable, so $\bar R_Q$ is a Markov
kernel by Tonelli's theorem and monotone convergence, and the defining identity, linear in the
kernel, passes to the mixture. Left invariance of $dh$ gives $\bar R_Q(ky,kA)=\bar R_Q(y,A)$ for
$k\in\mathbb G_Q$. If $g_\#Q=g'_\#Q$ then $g^{-1}g'\in\mathbb G_Q$, so $\bar R_Q^{g'}=\bar R_Q^g$;
finally $\bar R_Q^{kg}(ky,kA)=\bar R_Q^g(y,A)$. $\square$

\noindent\emph{Proof of \Cref{prop:rg-scale-functoriality}.}
(i) Functoriality forces the stated composites, and associativity gives
$C_{kr}\circ C_{\ell k}=C_{\ell r}$. (ii) is the functor axiom on two routes, and for a rule
family it is the axiom on the arrows indexed by $b_1$, $b_2$, and $b_1b_2$. (iii) Each composite of
configuration-keeping or additive arrows is the same statistic of the final partition, and the
witness is direct evaluation: by pairs $(0,0,1,0,1,0)\mapsto(0,1,1)\mapsto1$, by triples
$(0,0,1,0,1,0)\mapsto(0,0)\mapsto0$. (iv) If stochastic matrices $P,Q\geq0$ are mutually inverse,
rank gives them equal sizes. From $PQ=I$, every positive entry $P_{ik}$ forces row $k$ of $Q$ to
vanish off column $i$, so each column of $P$ has at most one positive entry, and unit row sums make
$P$ a permutation matrix. $\square$

\noindent\emph{Proof of \Cref{prop:rg-thermo-finite-labeled-fisher}.}
Differentiating $\sum_xp_\gamma(x)=1$ on the fixed support gives the branch-score identity, and
differentiating the softmax gives $\nabla\log W_\gamma=\xi_\gamma-\bar\xi$, hence
\eqref{eq:rg-thermo-finite-complete-score}. In the outer product of
\eqref{eq:rg-thermo-finite-complete-score}, summation over $x$ within each branch removes both cross
terms by the branch-score identity, which gives \eqref{eq:rg-thermo-finite-labeled-fisher}; the
label marginal has score $\xi_\gamma-\bar\xi$, so its Fisher tensor is
$\operatorname{Cov}_W(\xi_\gamma)$. Positive denominators on a finite support make the $C^3$ masses
give a $C^2$ Fisher tensor, and the Leibniz rule gives
\eqref{eq:rg-thermo-finite-weighted-fisher-derivatives}. A common shift cancels from
the weights and adds one vector to every $\xi_\gamma$, which leaves the covariance unchanged as a
function of $\theta$. Finally $v^\top g^F(P)v=\sum_zP(z)(v^\top\mathbf s_P(z))^2\geq0$. $\square$

\noindent\emph{Proof of \Cref{prop:rg-thermo-finite-missing-information}.}
On its positive support $\mathbf s^c=\mathbf s_P+\nabla\log K^{\mathrm{obs}}(o\given\gamma,x)$, and
each record row has fixed support and sums to one, so
$\sum_oK^{\mathrm{obs}}(o\given\gamma,x)\nabla\log K^{\mathrm{obs}}(o\given\gamma,x)=0$ and the cross
terms of the second moment vanish; this is
\eqref{eq:rg-thermo-finite-complete-observation-fisher}. The record is the image of the complete
sample under the parameter-independent map $(\gamma,x,o)\mapsto o$, with conditional law
$\boldsymbol\Pi_o$ given $o\in\mathsf O^+$. A finite fixed-support family with $C^3$ masses is
differentiable in quadratic mean, so \Cref{thm:cg-fisher-contraction} gives the score projection
and the Fisher loss in \eqref{eq:rg-thermo-finite-missing-information}. For a one-point alphabet
every record row equals one, so $p_\theta(o)=\sum_{\gamma,x}P_\theta(\gamma,x)=1$. The pointwise
observed-information form of this identity is due to \citet{Louis1982}; its average over the record
is \eqref{eq:rg-thermo-finite-missing-information}. $\square$

\noindent\emph{Proof of \Cref{prop:rg-thermo-account-word}.}
The score telescope comes first. Differentiating $P_j=P_{j-1}C_j$ at positive atoms gives
$\mathbf s_j(X_j)=\E[\mathbf s_{j-1}(X_{j-1})\given X_j]$, so the single-step loss is $D_j\succeq0$
by \Cref{thm:cg-fisher-contraction}, and the sum telescopes. The reversed path is Markov, so
$\E[\varepsilon^{\mathrm{sc}}_j\given X_j,\ldots,X_m]=0$; for $j<k$ the increment
$\varepsilon^{\mathrm{sc}}_k$ is measurable with respect to $X_j,\ldots,X_m$, which gives
orthogonality, and $\E[\varepsilon^{\mathrm{sc}}_j(\varepsilon^{\mathrm{sc}}_j)^\top]
=g_{j-1}-g_j$ because $\E[\mathbf s_{j-1}\mathbf s_j^\top]=g_j$. Domination persists under each
kernel on finite spaces, so every divergence is finite, and
\Cref{thm:rg-exact-coarse-vfe,thm:cg-evidence-preserving-channel} at each step give
$\Fenergy_{j-1}-\Fenergy_j=K^{\mathrm{KL}}_j$ at unchanged evidence; summing proves the free-energy
telescope, and division by $N_m$ gives the per-agent form.
$\square$

\noindent\emph{Proof of \Cref{thm:rg-thermo-label-strict}.}
Uniform convergence of $\lambda/N$ gives
$\lambda_{\gamma,N}-\lambda_{\gamma_*,N}\leq-\delta_{\mathsf K}N/2$ for each $\gamma\neq\gamma_*$ and
large $N$; the softmax denominator contains the winning mass, so
$W_{\gamma,N}\leq e^{-\delta_{\mathsf K}N/2}$. Suppress $N$ and write
$\bar\xi_i=\sum_\eta W_\eta\xi_{\eta,i}$. The softmax derivatives
$\partial_iW_\gamma=W_\gamma(\xi_{\gamma,i}-\bar\xi_i)$ and
$\partial_j\bar\xi_i=\sum_\eta W_\eta\partial_j\xi_{\eta,i}+\operatorname{Cov}_W(\xi_i,\xi_j)$ show
that, for $\lvert\beta\rvert\leq2$, $\partial^\beta W_\gamma$ is $W_\gamma$ times a polynomial in
the weights and in the components of the $\xi_\eta$ and $\partial\xi_\eta$, of degree at most
$\lvert\beta\rvert$ in the latter, each of which is $O(N)$ by
\eqref{eq:rg-thermo-label-effective-growth}. Hence
$\|\partial^\beta W_{\gamma,N}\|_{\mathsf K}\leq\mathrm C_{\mathsf K}N^{\lvert\beta\rvert}
e^{-\delta_{\mathsf K}N/2}$ for $\gamma\neq\gamma_*$ and $\lvert\beta\rvert\leq2$, and the same
bound holds for $W_{\gamma_*,N}-1$, which is minus the sum of the losing weights. With
$\delta\xi_\gamma=\xi_\gamma-\xi_{\gamma_*}$ and
$\bar d=\sum_{\gamma\neq\gamma_*}W_\gamma\delta\xi_\gamma$, exact centering gives
$\operatorname{Cov}_W(\xi)=\sum_{\gamma\neq\gamma_*}W_\gamma\delta\xi_\gamma\delta\xi_\gamma^\top-\bar d\bar d^\top$.
Each derivative of $\delta\xi_\gamma$ through order two is $O(N)$, so a term with $j$ derivatives in
the sum is bounded by $\mathrm C_{\mathsf K}N^{j+2}e^{-\delta_{\mathsf K}N/2}$: at most $j$
derivatives fall on a weight and the two score factors contribute at most $N$ each. The product
$\bar d\bar d^\top$ carries two exponential factors and obeys the same bound. Dividing by $N$ proves
\eqref{eq:rg-thermo-label-strict-covariance-bound}. The branch error equals
$\sum_{\gamma\neq\gamma_*}W_\gamma(\mathcal I_{\gamma,N}-\mathcal I_{\gamma_*,N})$, whose branch
differences have bounded $C^2$ norms, so the Leibniz sum with the weight bound gives the stated
bound. Adding the $C^2$ error of the
winning branch and the covariance error, via \eqref{eq:rg-thermo-finite-labeled-fisher}, proves the
$C^2$ convergence, since exponentials dominate the powers of $N$. For the floor, uniform convergence
in operator norm gives $\|g_N-\mathcal I_{\gamma_*}\|\leq\tfrac12\mathrm c_{\min}$ on $\mathsf K$ for
large $N$, so $z^\top g_Nz\geq z^\top\mathcal I_{\gamma_*}z-\tfrac12\mathrm c_{\min}\lvert z\rvert^2
\geq\tfrac12\mathrm c_{\min}\lvert z\rvert^2$. $\square$

\noindent\emph{Proof of \Cref{prop:rg-thermo-label-tie-scale}.}
The label set is finite, so there is a positive bulk gap from $S$ to every other label at
$\theta_0$; comparison with one maximizing label and convergence of $\lambda/N$ make every
nonmaximizing weight exponentially small, so its limit is zero. $C^1$ convergence gives
$\xi_{\gamma,N}/N\to\mathsf v_\gamma$. Dividing \eqref{eq:rg-thermo-finite-labeled-fisher} by $N^2$
makes the branch term $O(1/N)$ and leaves $\operatorname{Cov}_{W_N}(\xi_{\gamma,N}/N)$, which
converges by finite summation. For every $z$,
$z^\top\operatorname{Cov}_{\bar{\mathsf p}}(\mathsf v_\gamma)z
=\sum_\gamma\bar{\mathsf p}_\gamma[z\cdot(\mathsf v_\gamma-\bar{\mathsf v})]^2$, which vanishes
exactly when $z\cdot\mathsf v_\gamma$ is constant on the positive-weight labels, so the kernel is the
orthogonal complement of the stated span. A positive limiting quadratic form gives the asymptotic
after multiplication by $N^2$, and a zero one gives only the little-o conclusion. $\square$

\noindent\emph{Proof of \Cref{thm:rg-thermo-label-crossover}.}
A compact set of $\mathsf w$ is bounded, so its image under $\theta_N$ and the segments to
$\theta_0$ lie in the chosen neighborhood for large $N$. Put
$\lambda^{\mathrm{inc}}_{\gamma,N}(\mathsf w)=\lambda_{\gamma,N}(\theta_0+\mathsf w/N)-\lambda_{\gamma,N}(\theta_0)$.
The fundamental theorem of calculus gives
\begin{equation}
\lambda^{\mathrm{inc}}_{\gamma,N}(\mathsf w)
=\int_0^1N^{-1}\nabla\lambda_{\gamma,N}(\theta_0+\tau\mathsf w/N)\cdot\mathsf w d\tau,
\label{eq:rg-thermo-label-window-increment}
\end{equation}
whose integrand tends uniformly to $\mathsf v_\gamma\cdot\mathsf w$ by the $C^1$ part of
\eqref{eq:rg-thermo-label-effective-growth} and continuity of $\nabla\mathsf s_\gamma$.
Differentiation in $\mathsf w$ gives
\begin{equation}
D_{\mathsf w}\lambda^{\mathrm{inc}}_{\gamma,N}\to\mathsf v_\gamma,
\qquad
\|D^2_{\mathsf w}\lambda^{\mathrm{inc}}_{\gamma,N}\|=O(N^{-1}),
\qquad
\|D^3_{\mathsf w}\lambda^{\mathrm{inc}}_{\gamma,N}\|=O(N^{-2}),
\label{eq:rg-thermo-label-window-derivatives}
\end{equation}
so $\lambda^{\mathrm{inc}}_{\gamma,N}\to\mathsf w\cdot\mathsf v_\gamma$ in $C^3$. For tied labels,
the relative log masses $\lambda_{\gamma,N}-\lambda_{\gamma_0,N}$ at $\theta_N(\mathsf w)$ are the
logarithms of the center ratios plus the difference of increments, and converge in $C^3$ to
$\log\mathrm c_\gamma+\mathsf w\cdot(\mathsf v_\gamma-\mathsf v_{\gamma_0})$; exponentiation and the
finite softmax are smooth on the resulting compact ranges, with denominator at least the reference
mass one. For a nonmaximizing label the relative center log mass is at most $-\mathrm cN$ for some
constant $\mathrm c>0$, while its increment and $\mathsf w$-derivatives through order three are
bounded, so its relative exponential and those derivatives are $O(e^{-\mathrm cN})$. Normalizing
proves \eqref{eq:rg-thermo-label-limiting-weights} in $C^3$, including the zero coordinates outside
$S$. With $\hat\xi_{\gamma,N}(\mathsf w)=N^{-1}\xi_{\gamma,N}(\theta_N(\mathsf w))$,
\eqref{eq:rg-thermo-label-window-derivatives} gives $\hat\xi_{\gamma,N}\to\mathsf v_\gamma$ in $C^2$,
and the affine change $\theta_N$ has derivative $N^{-1}\mathrm{Id}$, so
\eqref{eq:rg-thermo-finite-labeled-fisher} becomes
\begin{equation}
g^F_{\mathsf w}(\mathcal E_N)(\mathsf w)
=\operatorname{Cov}_{W_N(\theta_N(\mathsf w))}\bigl(\hat\xi_{\gamma,N}(\mathsf w)\bigr)
+\frac1N\sum_\gamma W_{\gamma,N}(\theta_N(\mathsf w))\mathcal I_{\gamma,N}(\theta_N(\mathsf w)).
\label{eq:rg-thermo-label-window-decomposition}
\end{equation}
The covariance converges in $C^2$ by finite sums and products, without dividing by a weight that
tends to zero. The branch densities and their first two $\mathsf w$-derivatives are bounded, since a
$\mathsf w$-derivative is $N^{-1}$ times a $\theta$-derivative, and the weights have bounded $C^2$
norms, so every term of the Leibniz sum for the last term is $O(N^{-1})$. Finally the
log-sum-exponential has gradient $\sum_\gamma\bar{\mathsf p}_\gamma\mathsf v_\gamma$ and Hessian
$\operatorname{Cov}_{\bar{\mathsf p}}(\mathsf v_\gamma)$. The window is the finite-mixture
counterpart of the finite-volume rounding of a first-order transition, which
\citet{BorgsKotecky1990} derive from a lattice model as a hyperbolic-tangent crossover of width of
order the inverse volume. $\square$

\noindent\emph{Proof of \Cref{prop:rg-thermo-label-tie-counterexamples}.}
The common branch law cancels from the relative log mass and has zero score, so the label is the
only source of information. Then $W_1'=W_1(1-W_1)\omega'_N$, the centered label scores are
$(1-W_1)\omega'_N$ and $-W_1\omega'_N$, and averaging their squares gives
\eqref{eq:rg-thermo-label-binary-formula}, with $W_1(1-W_1)=\operatorname{bv}(\omega_N)$. In (a) the
$C^3$ norm of $\omega_N$ on a compact interval is at most $\mathrm CN$ and $\omega_N/N=\theta$; the
formulas follow by substitution. In (b) the derivatives vanish and $\omega_N/N\to0$ in $C^3$. In (c)
the $C^3$ norms are $O(N^{\mathrm c})$ and $O(\sqrt N)$, both $o(N)$, and substitution gives the two
statements, with $\operatorname{bv}(0)=1/4$ and $\operatorname{bv}(\log2)=2/9$. In (d)
$\omega_N/N=\theta-\mathrm cN^{-1}\log N\to\theta$ in $C^3$ and the $C^3$ norm is $O(N)$; at zero,
$\operatorname{bv}(-\mathrm c\log N)N^2/N=N^{1+\mathrm c}/(1+N^{\mathrm c})^2$. In (e) derivatives
through order three are bounded by $N^{-1/6}N=N^{5/6}$, so $\omega_N/N\to0$ in $C^3$;
$\lvert\omega_N\rvert\leq N^{-1/6}$ gives uniform convergence of the weights, and
$0\leq F_N\leq N^{-1}\tfrac14N^{1/3}$. The Leibniz rule gives
\eqref{eq:rg-thermo-label-oscillatory-derivative}. At zero $\omega_N=0$, $\omega'_N=N^{1/6}$,
$\omega''_N=0$, $\omega'''_N=-N^{5/6}$, and $\operatorname{bv}(0)=1/4$, $\operatorname{bv}'(0)=0$,
$\operatorname{bv}''(0)=-1/8$, which gives the value of $F''_N(0)$; the uniform limit has second
derivative zero. $\square$

\section{Proofs for Exact Agent-Network Coarse-Graining}
\label{sec:proofs-network-coarsegraining}

\noindent\emph{Proof of \Cref{thm:rg-bounded-action-calculus}.}
Positivity and unitality of $\mathsf U$ give the bounds in (a), the bound on
$\mathcal A(\psi)$, and the boundedness and positivity of $\mathsf U^\psi$;
substitution in \eqref{eq:rg-local-action-map} gives
\eqref{eq:rg-action-recentering-identity}. If $\lVert h\rVert_\infty<\log2$, then
$\lVert\mathsf U^\psi(e^{-h})-1\rVert_\infty\le e^{\lVert h\rVert_\infty}-1<1$, so
the exponential series in $L^\infty(\bar{\mathsf m}_o)$ and the series of
$\log(1+u)$ in the commutative Banach algebra $L^\infty(\bar{\mathsf m}_o^c)$
converge, and $h\mapsto\mathcal A(\psi+h)$ is real analytic on the ball of radius
$\log2$ about every $\psi$, which proves (a). The Fr\'echet chain rule for the
exponential, $\mathsf U^\psi$, and the logarithm gives
\eqref{eq:rg-action-first-frechet-general}; differentiating
$\int h(y)C^\dagger_{\psi+tk}(z,dy)$ at $t=0$ gives
\eqref{eq:rg-action-second-frechet-general}; and $\psi=0$ gives the origin. For (c),
$\psi\le\psi'+c$ implies $e^{-\psi}\ge e^{-c}e^{-\psi'}$, and positivity of
$\mathsf U$ gives the order statement; applying it in both directions with
$c=\lVert\psi-\psi'\rVert_\infty$ gives
\eqref{eq:rg-nonlinear-action-sup-contraction}, and
$\mathsf U(e^{-\psi-c})=e^{-c}\mathsf U(e^{-\psi})$ gives homogeneity, which with
$\mathsf U1=1$ gives (d). For (e),
$\widehat{\mathsf m}^\psi C=(e^{-\psi}\mathsf m_o)C/\mathsf m_o(e^{-\psi})$ and
$\bar{\mathsf m}_o^c=\mathsf m_oC/M_o$ give the density
$e^{-\mathcal A(\psi)}/\bar{\mathsf m}_o(e^{-\psi})$; the first two derivatives of
$t\mapsto\log\bar{\mathsf m}_o(e^{-th})$ at $t=0$ are $-\bar{\mathsf m}_o(h)$ and
$\operatorname{Var}_{\bar{\mathsf m}_o}(h)$, and polarization gives the covariance.
$\square$

\noindent\emph{Proof of \Cref{thm:rg-action-lp-contraction}.}
Conditional Jensen gives $\lvert\mathsf Uh\rvert^p\le\mathsf U\lvert h\rvert^p$ for
finite $p$, and integrating the disintegration
\eqref{eq:rg-normalized-reverse-kernel} proves
\eqref{eq:rg-action-lp-contraction} and preservation of the mean; $p=\infty$
follows from the pointwise bound. The conditional variance identity
\begin{equation}
\int h^2d\bar{\mathsf m}_o
=\int(\mathsf Uh)^2d\bar{\mathsf m}_o^c
+\int\operatorname{Var}_{C^\dagger(z,\cdot)}(h)\bar{\mathsf m}_o^c(dz),
\label{eq:rg-action-total-conditional-variance}
\end{equation}
the decomposition behind \Cref{thm:cg-fisher-contraction} applied to $h$, proves
\eqref{eq:rg-action-l2-conditional-defect}, and a nonnegative conditional variance
integrates to zero exactly when it vanishes almost everywhere. For the oscillation
bound, fix a representative of $h$ and let $N$ be the $\bar{\mathsf m}_o$-null set on which it
lies outside $[\operatorname{ess~inf}h,\operatorname{ess~sup}h]$. Disintegration gives
$\int C^\dagger(z,N)\bar{\mathsf m}_o^c(dz)=\bar{\mathsf m}_o(N)=0$, so
$C^\dagger(z,N)=0$ for almost every $z$. Thus this fixed $h$ lies between its essential bounds
almost surely under those kernels. For such a pair
$(z,z')$ the signed measure $\mu=C^\dagger(z,\cdot)-C^\dagger(z',\cdot)$ has total
mass zero and positive and negative parts of equal mass $t\le\delta$, so
$(\mathsf Uh)(z)-(\mathsf Uh)(z')=\int h(y)\mu^+(dy)-\int h(y)\mu^-(dy)
\le t(\operatorname{ess~sup}h-\operatorname{ess~inf}h)$. The sets where
$\mathsf Uh$ lies within any positive distance of its essential supremum or
infimum have positive $\bar{\mathsf m}_o^c$ measure, so the bound passes to
$\operatorname{ess~sup}\mathsf Uh-\operatorname{ess~inf}\mathsf Uh$; applying it
factor by factor gives the composition statement. $\square$

\noindent\emph{Proof of \Cref{thm:rg-hoeffding-action-isomorphism}.}
Fubini's theorem for the product $\nu_\ell$ gives commutation and
$\E_{\ell,B}\E_{\ell,B'}=\E_{\ell,B\cap B'}$; expanding the
product in \eqref{eq:rg-hoeffding-mobius-projectors} with this rule gives the
alternating sum, and M\"obius inversion on the Boolean lattice gives the
orthogonality relations and the resolution of the identity. Every term of
$\mathsf P_{\ell,B}f$ is a function of $y_B$, and the factor
$I-\E_{\ell,V_\ell\setminus\{i\}}$ removes its mean in coordinate $i$; hence
the nonempty components recover $\mathsf g$ and sum to
$f-\E_{\ell,\varnothing}f$, the class $[f]$. The triangle inequality and
$\lVert[\psi]\rVert_{\mathrm{osc}}\le\lVert\psi\rVert_\infty$ give
$\lVert\operatorname{asm}_\ell\rVert\le1$. For nonempty $B$ the projector
$\mathsf P_{\ell,B}$ annihilates constants and is an alternating sum of
$2^{\lvert B\rvert}$ contractions, so
$\lVert\mathsf P_{\ell,B}f\rVert_\infty\le2^{\lvert B\rvert}\lVert f-c\rVert_\infty$
for every constant $c$, and
$\sum_{B\ne\varnothing}2^{\lvert B\rvert}=3^{\lvert V_\ell\rvert}-1$. For the
biased coordinates, expanding $\prod_i\bigl((Y_i-m)+m\bigr)$ with
$m=\E Y_i=2p-1$ gives the components $m^{n-\lvert B\rvert}\prod_{i\in B}(Y_i-m)$,
whose supremum norms $m^{n-\lvert B\rvert}(1+m)^{\lvert B\rvert}$ sum to
$(1+2m)^n-m^n=(4p-1)^n-(2p-1)^n$. $\square$

\noindent\emph{Proof of \Cref{prop:rg-interaction-intra-block-invariance}.}
Because $\bar{\mathsf m}_\ell$ is a product over agents and hence over blocks, the
joint law $\bar{\mathsf m}_\ell(dy)C_\ell(y,dz)$ is the product of the block laws
$\bar{\mathsf m}_\ell^{I}(dy_I)C_I(y_I,dz_I)$; its $z$-marginal is the product of
the block marginals, and the product of the block reverse kernels
$C_I^\dagger(z_I,dy_I)$ is a version of the reverse kernel. For $\mathsf
g\in\mathcal G_\ell^{\mathrm{intra}}$, grouping components by block writes
$\psi_{\mathsf g}=\sum_I\psi_I$ with $\psi_I$ a function of $y_I$, so Fubini's
theorem gives $\mathcal A_\ell(\psi_{\mathsf g})=\sum_If_I(z_I)$ with
$f_I=-\log\mathsf U_I(e^{-\psi_I})$. For $\lvert B\rvert\ge2$ choose
$j\in B\setminus\{I\}$; the function $\E_{\ell+1,B}f_I$ does not depend on
$z_j$, so the factor $I-\E_{\ell+1,V_{\ell+1}\setminus\{j\}}$ of the product
form of \eqref{eq:rg-hoeffding-mobius-projectors} annihilates it. Hence
$\mathsf P_{\ell+1,B}f_I=0$, and linearity finishes the proof. $\square$

\noindent\emph{Proof of \Cref{prop:rg-interaction-rn-gauge-covariance}.}
Channel covariance gives $\widehat{\mathsf m}_\ell\widehat
C_\ell=(\kappa_{\ell+1})_\#(\mathsf m_\ell C_\ell)$ and
$(e^{-\mathsf K_\ell\psi}\widehat{\mathsf m}_\ell)\widehat
C_\ell=(\kappa_{\ell+1})_\#((e^{-\psi}\mathsf m_\ell)C_\ell)$. The Radon--Nikodym
derivative of two measures pushed through one Borel isomorphism is the transported
derivative, and applying $-\log$ gives \eqref{eq:rg-action-rn-gauge-covariance}
without selecting a pointwise reverse kernel. Differentiating at an arbitrary
bounded $\psi$ with part (b) of \Cref{thm:rg-bounded-action-calculus} gives
\eqref{eq:rg-tilted-derivative-gauge-covariance}. Composing with the intertwinings
\eqref{eq:rg-hoeffding-gauge-covariance} gives
\eqref{eq:rg-nonlinear-interaction-gauge-covariance}, and inserting the intertwining
of the retained projections into $I-\mathsf P^{\mathrm{ret}}_{\ell+1}$ and into
\Cref{prop:rg-interaction-residual} gives the residual formulas. $\square$

\noindent\emph{Proof of \Cref{prop:rg-hyperedge-generation}.}
(a) is Tonelli's theorem applied to the product of nonnegative factors. (b) is
M\"obius inversion on the Boolean lattice of subsets of $V^c$, and the component
indexed by $B$ reads only $z_B$ because every term does. (c) follows by
substituting $gz$ and $gz^\circ$ in \eqref{eq:rg-mobius-potentials}. For (d), write
$f(u)=-\log2\cosh u$, so $f'''=2\operatorname{sech}^2\tanh$. With
$x=\sum_r\xi_ry_r$ and $y_r^2=1$, only odd powers of $x$ contain the monomial
$y_1y_2y_3$; the cubic term $f'''(\xi_0)x^3/6$ contributes
$f'''(\xi_0)\xi_1\xi_2\xi_3$, and higher odd powers are of order five in the
couplings. At $\xi_0=0$ the function $f(x)$ is even in $y$, so every odd Walsh
coefficient vanishes. $\square$

\noindent\emph{Proof of \Cref{prop:rg-bridge-kernels}.}
Adjointness is \eqref{eq:rg-agent-meta-bayes-adjoint} integrated against
$f(y_j)g(z_I)$. For the factorization witness take two blocks with one
$\{-1,1\}$-valued agent each, a uniform product baseline, a channel that flips each
coordinate independently with probability $1/5$, and the likelihood
$e^{Jy_1y_2}$. The baseline reverse probability of $y_1=1$ given $z_1=1$ is $4/5$
whatever $z_2$ is, while with $x=e^{2J}$ the posterior reverse probability is
$4(x+4)/(8x+17)$ at $z_2=-1$ and $4(4x+1)/(17x+8)$ at $z_2=1$, which differ exactly
when $J\ne0$. For a deterministic statistic, $\mathsf C\mathsf
Jf(z)=\E[f(\chi(Y))\given\chi(Y)=z]=f(z)$ almost everywhere. A channel whose output
does not depend on $y$ gives $\mathsf C\mathsf Jf=\boldsymbol\Pi^c_o(f)$, which
differs from $f$ for nonconstant $f$. $\square$

\noindent\emph{Proof of \Cref{thm:rg-strong-lumpability}.}
Assume \eqref{eq:rg-lumped-kernel} and take $\mu=\delta_y$. Then
$\mathsf T(y,\chi^{-1}A)=\mathsf T^\chi(\chi(y),A)$ for every $y$ and Borel $A$,
whose right side depends on $y$ only through $\chi(y)$, and surjectivity determines
$\mathsf T^\chi$ at every $z$. Conversely, for fixed $z$ the map
$A\mapsto\mathsf T(\iota(z),\chi^{-1}A)$ is a probability measure because $\chi^{-1}$
commutes with countable disjoint unions, and for fixed $A$ it is Borel in $z$ as the
composition of $\iota$ with $y\mapsto\mathsf T(y,\chi^{-1}A)$. Since
$\chi(\iota(\chi(y)))=\chi(y)$, hypothesis \eqref{eq:rg-strong-lumpability} gives the
same pointwise identity, and integrating it against $\mu$ gives
\eqref{eq:rg-lumped-kernel}. $\square$

\noindent\emph{Proof of \Cref{thm:rg-memory-projection}.}
At $t=0$ the empty context has mass one. If the context at $t$ has its
$\mathbb P^{\mathrm{path}}_N$ marginal, multiplying by its row gives the correct
context-plus-symbol marginal, and marginalizing the oldest symbol gives the next
context marginal; by induction no positive prefix enters a null context, and
successive row sums give normalization. Expanding the logarithm of the products and
regrouping by contexts gives \eqref{eq:rg-memory-excess-risk}; a competitor that
excludes a positive word has infinite divergence. Equality forces every occupied row
to agree, and recursive reachability forces the same path law. Subtracting the
source chain-rule logarithm from the projected product leaves, in expectation, the
difference between entropy given the context and entropy given the full past, which
is the conditional mutual information \eqref{eq:rg-memory-information-defect} by the
chain rules of \citet[Sec.~2.5]{CoverThomas2006}. Conditioning on more history does
not increase entropy, so the defect decreases in $k$, and $k\ge N-1$ keeps the whole
past. The projection keeps its own defining marginals, and a shorter prefix only
clips the available history. The total-variation bound is Pinsker's inequality
\citep[Lemma~11.6.1]{CoverThomas2006}. For the cycle,
$\mathbb P(Z_0=0)=2/3$ and both one-step rows from $0$ put mass $1/2$ on $0$, so the
projection gives $000$ the mass $1/6$ although the cycle never produces it.
Subtracting the two finite log ratios gives
\eqref{eq:rg-memory-recognition-vfe-difference}. $\square$

\noindent\emph{Proof of \Cref{thm:rg-memory-duration}.}
Stationarity and nonnegativity force $\mathsf T_{yy'}=0$ when $\bar p_y>0=\bar
p_{y'}$, so the positive support is closed. For $k\ge1$ the positive $k$-context
chain shifts a context to its last $k$ symbols after appending a symbol, with row
$\mathbb P^{\mathrm{path}}_{k+1}(\mathbf zz)/\mathbb P^{\mathrm{path}}_k(\mathbf z)$; its incoming mass at a context
is the corresponding marginal of $\mathbb P^{\mathrm{path}}_{k+1}$, which equals $\mathbb P^{\mathrm{path}}_k$ of that
context by stationarity, and this gives stationarity and exact $(k+1)$-windows; for
$k=0$ the projection is the stationary product of the one-symbol marginal.
Conditioning and stationarity give $0\le\mathrm h_{j+1}\le\mathrm h_j$. The chain
rule gives $H(Z_{0:N-1})=\sum_{j<N}\mathrm h_j$, and for $N>k$ the projected
cross-entropy is $H(Z_{0:k-1})+(N-k)\mathrm h_k$, which gives the sum and its rate.
For two adjacent observed blocks of lengths $n_1$ and $n_2$, past and future are
conditionally independent given the intervening hidden state, so their mutual
information is at most $\mathrm c_Y$; expanding it by stationary entropies gives
$H(Z_{0:n_2-1})-\sum_{j=0}^{n_2-1}\mathrm h_{n_1+j}\le\mathrm c_Y$, and letting $n_1$
and then $n_2$ grow gives the series bound, whose monotone terms give
$(k+1)\mathrm h^{\mathrm{ex}}_k\le\mathrm c_Y$. Finally, for $N>k$,
$N\mathrm h^{\mathrm{ex}}_k-D^{(k)}_N=k\mathrm h^{\mathrm{ex}}_k
+\sum_{t=k}^{N-1}\mathrm h^{\mathrm{ex}}_t\le\mathrm c_Y$, for $N\le k$ the bound
$N\mathrm h^{\mathrm{ex}}_k\le\mathrm c_Y$ holds, and bounding each nonzero summand of
$D^{(k)}_N$ by $\mathrm h^{\mathrm{ex}}_k$ gives the last inequality. $\square$

\noindent\emph{Proof of \Cref{prop:hist-law-fisher-agentization-nogo}.}
Relabelings fix the uniform law and its Fisher metric, so a natural rule would be
permuted by them and would induce a homomorphism from the symmetric group $S_7$ to
the symmetric group $S_3$ of the three blocks. Its restriction to the alternating
group $A_7$, which is perfect, is trivial, because $S_3$ is solvable, so every block
would be $A_7$-invariant. The tangent space is the sum-zero representation of
$A_7$. Since $A_7$ is doubly transitive on the seven outcomes, the commutant of its
permutation representation is spanned by the identity and the all-ones matrix, so
every equivariant endomorphism of the sum-zero space is scalar, the representation
is irreducible, and it has no invariant two-dimensional subspace. $\square$


\chapter{Proofs for Finite-Depth Formation}
\label{app:finite-depth-formation-proofs}

This appendix proves the propositions of \cref{ch:finite-depth-formation},
the lemmas behind \Cref{thm:finite-depth-adaptive-formation},
\Cref{prop:rg-formation-endogenous-cycle}, and its extension
\Cref{prop:phil-seeded-contrast-formation} of \cref{ch:philosophy}, in the symbols of
\cref{tab:rg-formation-notation}; \Cref{tab:rg-formation-lemma-dependencies}
orders the lemmas. The depth \(L\ge2\) is fixed. Constants
\(\mathrm c_L,\mathrm C_L>0\) may change between displays; they depend on
\(L\), on the fixed data, and, once chosen, on \(k_L\) and \(\eta_L\), but not
on \(\epsilon\), \(\mathrm A\), \(\mathrm T\), or the strengths unless the
dependence is displayed, and all limits have \(\epsilon,\mathrm T>0\).
Densities are relative to Lebesgue measure times counting measure, values of
the complete assignment are written \(j,k\in\mathcal J\), and the independent
equilibrium root factor of every law is omitted from non-root integrals and
appended unchanged. On \(M_\ell\), \(H_\ell\) is the Gaussian Schur precision
and \(\mathrm d_\ell=\dim M_0-\dim M_\ell\) the number of eliminated
coordinates.

\begin{sciencetable}{Lemma dependencies of the finite-depth formation theorem.}
\label{tab:rg-formation-lemma-dependencies}
\footnotesize
\begin{tabularx}{\textwidth}{>{\raggedright\arraybackslash}p{0.11\textwidth}>{\raggedright\arraybackslash}p{0.27\textwidth}>{\raggedright\arraybackslash}X>{\raggedright\arraybackslash}p{0.23\textwidth}}
\toprule
\textbf{Lemma} & \textbf{Inputs} & \textbf{Output} & \textbf{Exponent or constant} \\
\midrule
\Cref{lem:formation-parent-score} & \Cref{lem:rg-formation-carrier} & parent potentials, rows, covector; force, Hessian, metric bounds & \(B_*=\delta\sqrt{2N}\), \(\delta_*<1/2\) \\
\Cref{lem:formation-preparation-margin} & \Cref{lem:rg-formation-carrier} & prepared direction, strict margin & \(\mathrm m_L>0\) \\
\Cref{lem:formation-prepared-continuation} & \Cref{lem:formation-parent-score} & well-posedness, moments, parameter bounds & \(\mathrm A^2\mathrm J_{\ell+1}+\mathrm T\varsigma_\epsilon^2\) \\
\Cref{lem:formation-parameter-comparison} & \Cref{lem:formation-prepared-continuation} & calibration-dominated parameter & \(r_\epsilon=O(\epsilon^{2\mathrm Q_1+20})\) \\
\Cref{lem:formation-entropy-tools} & cited inequalities & log-Sobolev, Herbst, duality, transport, Pinsker & curvature, \(\mathrm T\) explicit \\
\Cref{lem:formation-activation-entropy} & \Cref{lem:formation-parent-score,lem:formation-prepared-continuation,lem:formation-entropy-tools} & activation and Fisher budgets; score moments & \(\mathrm A^2\mathrm J_{\ell+1}+\mathrm T\varsigma_\epsilon^2\) \\
\Cref{lem:formation-conditional-entropy} & \Cref{lem:formation-entropy-tools,lem:formation-activation-entropy} & early discarded loss; covector error & \(\epsilon^{14}\); \(\tfrac12\mathrm c_{\mathrm{early},L}k_L>\mathrm Q_1+20\) \\
\Cref{lem:formation-terminal-joint-entropy} & \Cref{lem:formation-parent-score,lem:formation-activation-entropy} & terminal loss by joint relaxation & \(\epsilon^2\); \(\mathrm c_{\mathrm{joint},L}k_L>\mathrm Q_1+8\) \\
\Cref{lem:formation-fast-marginal} & \Cref{lem:formation-parent-score} & averaged row; environment errors & \(\mathrm e_\epsilon(y)\) \\
\Cref{lem:formation-fast-mixing} & \Cref{lem:formation-fast-marginal,lem:formation-prepared-continuation} & mixing to the moving fast equilibrium & rate \(\mathrm c_{\mathrm{fast},L}\mathrm J_1\) \\
\Cref{lem:formation-history-coupling} & \Cref{lem:formation-fast-marginal,lem:formation-fast-mixing,lem:formation-parameter-comparison} & assignment histories & \(\epsilon^6\); \(\mathrm c_{\mathrm{fast},L}k_L>\mathrm Q_1/2+12\) \\
\Cref{lem:formation-coherence} & \Cref{lem:formation-preparation-margin,lem:formation-prepared-continuation} & coherent components & \(e^{-\mathrm c_L/\epsilon^6}\); \(k_{\mathrm{geom},L}\), \(\eta_L^{\max}\) \\
\Cref{lem:formation-source-envelope} & \Cref{lem:formation-parameter-comparison} & source movement; kernel change & \(57/100\); \(1/200\) \\
\Cref{lem:formation-continuous-path} & \Cref{lem:formation-coherence,lem:formation-parameter-comparison} & common-path concentration & \(\mathrm A\epsilon^2\); \(e^{-\mathrm c_L/\epsilon^2}\) \\
\bottomrule
\end{tabularx}
\end{sciencetable}

\section{Exact marginals, rates, and parameter derivatives}
\label{sec:rg-formation-exact-proof}

\noindent\emph{Proof of \Cref{prop:rg-formation-normalized-experiment}.}
Under \(\nu_i^*\), \(k_i-m_i=\mathfrak d_i+\xi'_i\), whereas the physical
evaluator gives \(k_i-m_i=\xi'_i\) with the same law of \(m_i\). The
log-density ratio \([\mathfrak d_i(k_i-m_i)-\mathfrak d_i^2/2]/(1-\sigma^2_{\mathrm{ev}})\)
therefore has expectation \(\mathfrak d_i^2/(2(1-\sigma^2_{\mathrm{ev}}))\)
under every coupling of the two marginals, since
\(\E[k_i-m_i]=\bar k_i-\bar m_i\); this proves
\cref{eq:rg-formation-evaluator-optimum}, and summation gives
\(\mathbf x^\top R_{\mathrm{ev}}\mathbf x/2\). The rows \(\mathfrak d_i\) belong
to \(B_1\), so \(\mathrm J_1\mathfrak D^\top\mathfrak D\preceq H\), and
\(\mathrm T\le\mathrm J_1(1-\sigma^2_{\mathrm{ev}})/2\) gives
\(\mathrm TR_{\mathrm{ev}}\preceq H/2\); with \(\vartheta\ge1\) and the
orthogonal split of \Cref{lem:rg-formation-carrier},
\(\Omega_0(\vartheta)\succeq(\vartheta/2)H+P_{\mathcal K}\succ0\). Every factor
of \cref{eq:rg-formation-full-joint} is a positive normalized kernel of inputs
none of which is a recognition law, so iterated integration gives total mass
one and the records a positive density. \(\square\)

\noindent\emph{Proof of \Cref{prop:rg-formation-laws}.}
Since \(\widehat K_\vartheta-\mathrm TR_{\mathrm{ev}}\succeq(\vartheta/2)H
+(1+\kappa)P_{\mathcal K}\succ0\) and \(0\le\phi_S\le B_0\), both integrals are
finite and positive. At fixed \(\mathbf x\) both laws are proportional to
\(\pi_S(\vartheta)e^{-\phi_S(\mathbf x)}\) in \(S\), which gives
\cref{eq:rg-formation-full-source-row}, a product of receiver rows because
\(\phi_S\) is a sum and \(\pi_S\) a product over receivers. Define
\begin{equation}
\begin{aligned}
\mathrm n_{\mathcal K}&=\bigl(\kappa/(2\pi\mathrm T)\bigr)^{N/2},&
\mathrm n_{\mathrm{cal}}&=(2\pi\mathrm T)^{-1/2},\\
\mathrm B_{\mathrm T}(\vartheta)
 &=\mathrm T\log Z_0(\vartheta)
   -\mathrm T\log\bigl(\mathrm n_{\mathcal K}\mathrm n_{\mathrm{cal}}q_0^N\bigr)
   +\tfrac12\kappa|z_\star|^2+\tfrac12(\vartheta-2)^2,\\
\mathrm C_{\mathrm T}(\vartheta)&=\mathrm B_{\mathrm T}(\vartheta)
   -\mathrm T\log Z_{\mathrm{opt}}(\vartheta),&
p_\vartheta(\mathsf o)&=e^{-\mathrm B_{\mathrm T}(\vartheta)/\mathrm T}
                         Z_{\mathrm{phys}}(\vartheta).
\end{aligned}
\label{eq:rg-formation-evidence-scalar}
\end{equation}
As \(z_\star\in\mathcal K\), the log density of the joint at the records is
\(-[\mathrm E^{\mathrm{phys}}_\vartheta+\mathrm T\phi_S+\tfrac12\kappa|z_\star|^2
+\tfrac12(\vartheta-2)^2]/\mathrm T-\log Z_0
+\log(\mathrm n_{\mathcal K}\mathrm n_{\mathrm{cal}}q_0^N)+\log\pi_S\) plus the
physical evaluator density. Hence \(\mathrm T\E\log(q\prod_i\nu_i^*/p)\), in
which \cref{eq:rg-formation-evaluator-optimum} adds
\(\tfrac{\mathrm T}2\mathbf x^\top R_{\mathrm{ev}}\mathbf x
=\mathrm E_\vartheta-\mathrm E^{\mathrm{phys}}_\vartheta\), is the first line
of \cref{eq:rg-formation-restricted-vfe}, and completing the relative entropy
to \(q^\star_\vartheta\) gives the second. Integrating the joint over
\((\mathbf x,S)\) at the records gives \(p_\vartheta(\mathsf o)\), so
\(\mathrm C_{\mathrm T}+\mathrm T\log p_\vartheta(\mathsf o)
=\mathrm T\log(Z_{\mathrm{phys}}/Z_{\mathrm{opt}})\ge0\), as
\(\mathrm E^{\mathrm{phys}}_\vartheta\le\mathrm E_\vartheta\).

For the gap, \(Z_{\mathrm{phys}}/Z_{\mathrm{opt}}=
\E_{q^\star_\vartheta}e^{\mathbf x^\top R_{\mathrm{ev}}\mathbf x/2}\). The
Gaussian law \(\gamma^{\mathrm E}_\vartheta\propto e^{-\mathrm E_\vartheta/\mathrm T}\)
has its mean in \(\mathcal K\), which \(R_{\mathrm{ev}}\) annihilates, and since
\(R_{\mathrm{ev}}\) has rank \(N\) the eigenvalues \(\lambda_k\) of
\(\mathrm T\widehat K_\vartheta^{-1}R_{\mathrm{ev}}\) satisfy
\(\lambda_k\le\mathrm T/(\mathrm J_1(1-\sigma^2_{\mathrm{ev}}))\le1/2\) and
\(\sum_k\lambda_k\le N\mathrm T/(\mathrm J_1(1-\sigma^2_{\mathrm{ev}}))\le1\), so
\(\E_{\gamma^{\mathrm E}_\vartheta}e^{\mathbf x^\top R_{\mathrm{ev}}\mathbf x/2}
=\prod_k(1-\lambda_k)^{-1/2}\le e^{\sum_k\lambda_k}\). The density of
\(q^\star_\vartheta\) relative to \(\gamma^{\mathrm E}_\vartheta\otimes\pi(\vartheta)\)
is at most \(e^{B_0}\) and \(e^{\mathbf x^\top R_{\mathrm{ev}}\mathbf x/2}\ge1\),
so \(\E_{q^\star_\vartheta}[e^{\mathbf x^\top R_{\mathrm{ev}}\mathbf x/2}-1]
\le e^{B_0}(e^{\sum_k\lambda_k}-1)\le2e^{B_0}\sum_k\lambda_k\), and
\(\log(1+y)\le y\) gives the first bound of
\cref{eq:rg-formation-physical-optimizer-gap}. The second follows from
\(\KL(q^\star_\vartheta\Vert\boldsymbol\Pi_{\vartheta,\mathsf o})
=\log(Z_{\mathrm{phys}}/Z_{\mathrm{opt}})
-\tfrac12\E_{q^\star_\vartheta}\mathbf x^\top R_{\mathrm{ev}}\mathbf x\).
\(\square\)

\lemmaheading{Weighted geometry and exact parent score}{lem:formation-parent-score}
For the triangle-chain experiment, all marked marginal laws and inherited
proposal rows are normalized and compose exactly. With orthonormal flag
columns \(\mathsf Q_\ell\) spanning \(\mathcal U_{\ell-1}\cap\mathcal U_\ell^\perp\),
\(\mathsf B=B[\mathsf Q_1\cdots\mathsf Q_L]\), and
\(\mathbf x_{\mathcal U}=\sum_\ell\mathsf Q_\ell w_\ell\),
\begin{equation}
\|\mathsf B^{-1}\|_F^{-2}\sum_\ell\mathrm J_\ell|w_\ell|^2
\le \mathbf x_{\mathcal U}^{\top}H\mathbf x_{\mathcal U}
\le\|\mathsf B\|_F^2\sum_\ell\mathrm J_\ell|w_\ell|^2,
\label{eq:rg-formation-weighted-flag}
\end{equation}
and the same bounds hold for every Schur complement obtained by minimizing over
an initial group of normal layers. On \(1\le\vartheta\le21/10\), the parent
potentials are \cref{eq:rg-formation-parent-potential} with the tilt and
eliminated Gaussian scalar of
\cref{eq:rg-formation-parent-tilt-proof,eq:rg-formation-hidden-scalar-proof}
below, the inherited rows are those of \cref{eq:rg-formation-parent-generator}
and satisfy the detailed balance of \Cref{prop:rg-formation-rate-composition},
and the complete covector, including the source-prior score, is
\cref{eq:rg-formation-complete-score}. With \(B_*=\delta\sqrt{2N}\) and
\(\delta_*=2\delta+2N\delta^2<1/2\), the parent forces, Hessians, and scaled
metric obey
\begin{align}
 &|H_\ell^{-1/2}\mathrm T\nabla\psi_{\ell,j}|\le B_*\sqrt{\mathrm T},
 \qquad
 (\vartheta-\delta_*)H_\ell
 \preceq\nabla^2V_{\ell,j}
 \preceq(\vartheta+2\delta)H_\ell,
 \label{eq:rg-formation-parent-relative-bounds}\\
 &1\le\mathsf m(\vartheta)\le1+\mathrm C_L\mathrm T,\qquad
 |\mathsf m'(\vartheta)|\le\mathrm C_L\mathrm T,
 \label{eq:rg-formation-metric-bounds-proof}
\end{align}
for every fixed depth and all sufficiently small \(\epsilon\).
\status{ESTABLISHED}

\noindent\emph{Proof.}
By \Cref{lem:rg-formation-carrier}, rows of level at most \(\ell\) annihilate
\(\mathcal U_\ell\), so \(\mathsf B\) and \(\mathsf B^{-1}\) are block lower
triangular with invertible diagonal blocks. With
\(\mathsf D_{\mathrm J}=\operatorname{diag}(\sqrt{\mathrm J_\ell}I)\), each
nonzero block \((\ell,r)\) of
\(\mathsf D_{\mathrm J}\mathsf B^{\pm1}\mathsf D_{\mathrm J}^{-1}\) is that of
\(\mathsf B^{\pm1}\) times \(\sqrt{\mathrm J_\ell/\mathrm J_r}\le1\), which
proves \cref{eq:rg-formation-weighted-flag}, and \(H|_{M_0}=H_{\mathcal U}\oplus
\mathrm J_1I_{\mathcal H}\). The Schur quadratic is the infimum over the
eliminated coordinates, so minimizing both bounds proves the retained
versions; in particular \(H_\ell\succeq\mathrm h_LI\) for a fixed
\(\mathrm h_L>0\), since \(\mathrm J_L=1\), the normal part of \(H_\ell\) is at
most \(\mathrm C_L\mathrm J_{\ell+1}I\) with \(\mathrm J_{L+1}=0\), and the loop
part keeps the weight \(\mathrm J_1\).

For a cumulative hidden/retained split \((z,y)\) of \(M_0\) with blocks
\(H^{(\ell)}_{zz}\), \(H^{(\ell)}_{zy}\) of \(H\), completing the square gives
the lift in the first line below, with \(z\) embedded in the hidden subspace.
The centered hidden Gaussian \(\gamma^{\mathrm{hid}}_{\ell,\vartheta}\), of
covariance \((\mathrm T/\vartheta)(H^{(\ell)}_{zz})^{-1}\), has a mean
independent of \(\vartheta\), and integrating the source factor gives
\cref{eq:rg-formation-parent-potential} with the tilt, the scalar, and the
normalized optimizer of the other lines, \(\mathrm C_0^{\mathrm G}=0\):
\begin{align}
 \Xi_\ell y&=\bigl(-(H^{(\ell)}_{zz})^{-1}H^{(\ell)}_{zy}y,\ y\bigr),\qquad
 \Xi_\ell^\top H\Xi_\ell=H_\ell,\qquad \Xi_\ell^\top Hz=0,
 \label{eq:rg-formation-gaussian-lift}\\
 \psi_{\ell,j}(y,\vartheta)
 &=-\log\int e^{-\phi_j(\Xi_\ell y+z)}
                  \gamma^{\mathrm{hid}}_{\ell,\vartheta}(dz),
 \qquad 0\le\psi_{\ell,j}\le B_0,
 \label{eq:rg-formation-parent-tilt-proof}\\
 \mathrm C_\ell^{\mathrm G}(\vartheta)
 &=-\frac{\mathrm T\mathrm d_\ell}{2}\log\frac{2\pi\mathrm T}{\vartheta}
       +\frac{\mathrm T}2\log\det H^{(\ell)}_{zz},
 \label{eq:rg-formation-hidden-scalar-proof}\\
 q^{\star,\ell}_\vartheta(y,j)
 &=\frac{e^{-\psi_{\ell,j}(y,\vartheta)}
       \gamma_{\ell,\vartheta}(y)\pi_j(\vartheta)}{Z_\ell(\vartheta)},\quad
 \gamma_{\ell,\vartheta}=\mathcal N\bigl(0,(\mathrm T/\vartheta)H_\ell^{-1}\bigr),\quad
 e^{-B_0}\le Z_\ell\le1.
 \label{eq:rg-formation-normalized-parent-proof}
\end{align}
Positivity and Gaussian domination justify iterated integration and
differentiation on compact parameter intervals.
\Cref{prop:rg-formation-rate-composition} gives detailed balance, composition,
and \(\beta^\ell_{jk}\ge e^{-B_0}\pi_k(\vartheta)\) for the rows of
\cref{eq:rg-formation-parent-generator}, and the physical, optimizing, and
recognition laws compose by separate integration through the same channel.

Let \(\mathfrak B_{\mathrm{int}}\) have rows \(\sqrt{\mathrm J_1}\mathbf r_e\)
on internal edges, \(\mathbf r_e\) the covector of \(r_e\), so
\(\mathfrak B_{\mathrm{int}}^\top\mathfrak B_{\mathrm{int}}\preceq H\). An edge is
selected \(\mathrm n_e\in\{0,1,2\}\) times, \(\sum_e\mathrm n_e=N\), so the
force of \(\phi_j=\delta\sum_i\chi(\sqrt{\mathrm J_1/\mathrm T} r_{ij_i})\) is
\(\delta\sqrt{\mathrm T}\mathfrak B_{\mathrm{int}}^\top\mathfrak c\) with
\(\mathfrak c_e=\mathrm n_e\chi'(\sqrt{\mathrm J_1/\mathrm T} r_e)\) and, since
\(|\chi'|,|\chi''|\le1\), \(|\mathfrak c|^2\le2N\). This gives the first line
below on \(M_0\), and with the centered Gaussian variable of
\cref{eq:rg-formation-parent-tilt-proof} held fixed, the normalized tilted law
\(\nu_{\ell,y,j}\) gives the other two:
\begin{align}
 &|H^{-1/2}\mathrm T\nabla\phi_j|\le B_*\sqrt{\mathrm T},
 \qquad
 -2\delta H\preceq \mathrm T\nabla^2\phi_j\preceq2\delta H,
 \qquad B_*=\delta\sqrt{2N},
 \label{eq:rg-formation-original-relative-bounds}\\
 &\mathrm T\nabla\psi_{\ell,j}
 =\E_{\nu_{\ell,y,j}}\Xi_\ell^\top(\mathrm T\nabla\phi_j),
 \label{eq:rg-formation-marginal-force-proof}\\
 &\mathrm T\nabla^2\psi_{\ell,j}
 =\E_{\nu_{\ell,y,j}}\Xi_\ell^\top(\mathrm T\nabla^2\phi_j)\Xi_\ell
       -\mathrm T^{-1}\operatorname{Cov}_{\nu_{\ell,y,j}}
                    (\Xi_\ell^\top\mathrm T\nabla\phi_j).
 \label{eq:rg-formation-marginal-hessian-proof}
\end{align}
Through the isometry \(H^{1/2}\Xi_\ell H_\ell^{-1/2}\), the first line bounds
the force by \(B_*\sqrt{\mathrm T}\) in the \(H_\ell\) norm, the covariance by
\(B_*^2H_\ell=2N\delta^2H_\ell\), and the first Hessian term between
\(\pm2\delta H_\ell\), which proves
\cref{eq:rg-formation-parent-relative-bounds}. The source factor is even, so
centered Gaussian integration makes \(\psi_{\ell,j}\) even in \(y\), and every
non-root sheet has its unique minimum at zero.

The hidden Gaussian scale score
\(\mathfrak s_\ell=(\mathrm d_\ell-|\mathbf Z|^2)/(2\vartheta)\), with
\(\mathbf Z\) standard Gaussian, has mean zero and second moment
\(\mathrm d_\ell/(2\vartheta^2)\), and the tilted density ratio lies in
\([e^{-B_0},e^{B_0}]\); as the Gaussian conditional center is independent of
\(\vartheta\), uniformly in the retained input,
\begin{equation}
 \partial_\vartheta\psi_{\ell,j}=-\E_{\nu_{\ell,y,j}}\mathfrak s_\ell,\qquad
 |\partial_\vartheta\psi_{\ell,j}|
 \le(e^{B_0}-1)\frac{\sqrt{\mathrm d_\ell}}{\sqrt2 \vartheta}.
 \label{eq:rg-formation-tilt-score-proof}
\end{equation}
By \cref{eq:rg-formation-parent-potential} the parent free energy is
\(\mathrm T\E_q\log q+\E_q[V_{\ell,S}-\mathrm T\log\pi_S]+\mathrm B_{\mathrm T}\),
with \(\mathrm B_{\mathrm T}'=\vartheta-2-\tfrac{\mathrm T}2\Tr(\Omega_0^{-1}H)\),
and differentiation at a fixed recognition law, which does not differentiate
assignment masses, gives \cref{eq:rg-formation-complete-score}, whose source-prior term, present
although it has no continuous gradient, has absolute unscaled score at most
four for the distinguished binary prior. In the normalized joint the prior, calibration, and source-prior
scores have zero mean and zero cross expectations, the other kernels have no
explicit parameter dependence at fixed inputs, and the three scores have
variances \(\tfrac12\Tr[(\Omega_0^{-1}H)^2]\), \(1/\mathrm T\), and, for
\(4(\mathbf 1_{\{S_{i_0}=+\}}-\pi_{i_0,+})\), \(16\pi_{i_0,+}(1-\pi_{i_0,+})\),
which gives \cref{eq:rg-formation-model-fisher}. As \(R_{\mathrm{ev}}\) is
fixed, annihilates \(\mathcal K\), and is supported on the first normal layer,
\(\Omega_0\) is positive for small \(\mathrm T\), and conjugation by the
quadratic core bounds the trace terms and their derivatives uniformly;
calibration contributes exactly one to \(\mathsf m\) and the other two terms
are nonnegative and \(O(\mathrm T)\), which proves
\cref{eq:rg-formation-metric-bounds-proof}. \hfill\(\square\)

\noindent\emph{Proof of \Cref{prop:rg-formation-score-dissipation}.}
The metric and the covector are derived in \Cref{lem:formation-parent-score}.
The diffusion part of the generator is reversible with respect to
\(e^{-V_{\ell,S}/\mathrm T}\) on each sheet, which gives the first term of the
dissipation after integration by parts. For the jump part, detailed balance
(\Cref{prop:rg-formation-rate-composition}) turns
\(\mathrm T\int\sum_S\partial_tq_S\log\mathfrak r_S\) into
\(\mathrm T\int\sum_{S,S'}q^{\star,\ell}(y,S)\beta^\ell_{SS'}
(\mathfrak r_{S'}-\mathfrak r_S)\log\mathfrak r_S\), and symmetrizing in
\(S,S'\) gives the second term, whose integrand is nonnegative because the
logarithm is increasing. The parameter term is
\(g_\ell\dot\vartheta_\ell=-\epsilon\mathsf m^{-1}g_\ell^2\). At \(\ell=0\) the
row \(\beta_{S'}\) does not depend on \(S\), and with
\(\mathfrak r_S=(q_{\mathbf X}/q^\star_{\mathbf X}) q(S\mid\mathbf x)/\beta_S\)
the jump integrand sums to
\(2[\KL(q(\cdot\mid\mathbf x)\Vert\beta)+\KL(\beta\Vert q(\cdot\mid\mathbf x))]\).
These identities hold for positive smooth densities with finite entropy and
moments; the moment bounds of \Cref{lem:formation-prepared-continuation} and
the cutoff passage of \Cref{lem:formation-activation-entropy} justify them for
the admitted solutions. \(\square\)

\section{Prepared laws and coupled continuation}
\label{sec:rg-formation-continuation-proof}

\lemmaheading{Prepared direction with a strict margin}{lem:formation-preparation-margin}
Consider the list of linear functionals on \(\mathcal U\) consisting of each
internal edge covector and, for each bridge \(e\) and each \(\ell<\ell(e)\),
the functional \(r_e\circ P_\ell|_{\mathcal U}\). Every functional in the list
is nonzero, and there is a unit vector \(\mathbf v_\star\in\mathcal U\) with
\begin{equation}
\mathrm m_L=\min_{e,\ \ell<\ell(e)}
       |r_e(P_\ell\mathbf v_\star)|>0,
\label{eq:rg-formation-preparation-margin}
\end{equation}
where internal edges enter only at \(\ell=0\). A fixed neighborhood of
\(\mathbf v_\star\) keeps a margin of at least \(\mathrm m_L/2\).
\status{ESTABLISHED}

\noindent\emph{Proof.}
For a bridge \(e\) and \(\ell<\ell(e)\), \(r_e\) is independent of the rows of
\(B_1,\ldots,B_\ell\) (\Cref{lem:rg-formation-carrier}), so it does not vanish
on \(W_\ell\); it annihilates \(\mathcal H\oplus\mathcal K\), hence does not
vanish on \(\mathcal U_\ell\), which \(P_\ell\) fixes. An internal residual of
triangle \(q\) is a nonzero combination of the contrast rows
\(O_\triangle r_q\), which lie in \(\mathcal U\), plus a multiple of the loop
covector, so its restriction to \(\mathcal U\) is nonzero. Evaluated in an
orthonormal basis of \(\mathcal U\) on \((1,t,\ldots,t^{\dim\mathcal U-1})\),
each of the \(\mathrm n_0\) functionals is a nonzero polynomial of degree less
than \(\dim\mathcal U\), so some integer \(t\in[0,\mathrm n_0(\dim\mathcal U-1)]\)
is a simultaneous nonroot; normalizing gives \(\mathbf v_\star\), and
continuity gives the neighborhood. \(\square\)

\lemmaheading{Prepared continuation}{lem:formation-prepared-continuation}
For the preparation of \Cref{def:rg-formation-regime} and every finite
sequence of activations, the coupled diffusion, source proposals, and own-law
parameter equation have a unique nonexplosive solution through
\(t_*=2/\epsilon\), with positive smooth joint densities. On every active level
\(\ell\) and for all \(t\le t_*\),
\begin{align}
 \E_{q_t^\ell}y^\top H_\ell y
 &\le \mathrm C_L[\mathrm A^2\mathrm J_{\ell+1}+\mathrm T\varsigma_\epsilon^2],
 \label{eq:rg-formation-actual-energy}\\
 \E_{q_t^\ell}|y|^2
 &\le \mathrm C_L[\mathrm A^2+\mathrm T\varsigma_\epsilon^2],
 \qquad \mathrm J_1\E_{q_t^\ell}|h|^2\le \mathrm C_L\mathrm T,
 \label{eq:rg-formation-actual-moments}
\end{align}
and the parameter satisfies \(1\le\vartheta_\ell(t)\le21/10\) and
\(|\dot\vartheta_\ell(t)|\le\mathrm C_L\epsilon\); these bounds use no
conditional-entropy or history estimate. \status{ESTABLISHED}

\noindent\emph{Proof.}
For a fixed positive instance and a prescribed continuous parameter path near
\([1,21/10]\), each sheet drift has a bounded Hessian and linear growth by
\cref{eq:rg-formation-parent-relative-bounds}, and the proposal rows are
smooth with bounded derivatives and a positive lower bound, their
Gaussian-integral denominator being at least \(e^{-B_0}\); interlacing the
sheet diffusions on an independent rate-one Poisson grid gives a unique
nonexplosive process with all polynomial moments. For the parameter feedback
let \(\mathfrak u^\vartheta_j(t',\mathbf x)=\E_\vartheta[\mathsf f(\mathbf X_t,S_t)\mid
\mathbf X_{t'}=\mathbf x,S_{t'}=j]\), for smooth \(\mathsf f\) with polynomially
growing derivatives. Given the proposal times and labels, it composes
Lipschitz sheet flows with products of positive row weights; flow derivatives
are bounded on compact intervals, a differentiated weight is at most a
constant times the weight, and the proposal count has all exponential
moments, so \(\mathfrak u\) and \(\nabla\mathfrak u\) have polynomial bounds.
It\^o's formula for \(\mathfrak u^{\widetilde\vartheta}\) along the
\(\vartheta\) process, stopped in a ball and then using the moments, gives
\begin{equation}
 |\E_{\vartheta}\mathsf f(\mathbf X_t,S_t)
       -\E_{\widetilde\vartheta}\mathsf f(\mathbf X_t,S_t)|
 \le \mathrm C\int_0^t|\vartheta(t')-\widetilde\vartheta(t')| dt',
 \label{eq:rg-formation-local-law-dependence}
\end{equation}
since the second-order coefficients agree and the drift and row differences
are at most \(\mathrm C|\vartheta-\widetilde\vartheta|(1+|\mathbf x|)\) and
\(\mathrm C|\vartheta-\widetilde\vartheta|\). The observable terms of
\cref{eq:rg-formation-complete-score} meet these hypotheses, and their explicit
terminal parameter dependence adds a local Lipschitz term, so the integral map
for \(\dot\vartheta=-\epsilon\mathsf m(\vartheta)^{-1}g_\ell(q,\vartheta)\) is
a contraction on a short interval, which gives local coupled existence and
uniqueness. The step length is uniform: for parameter paths in
\([3/4,5/2]\), \(\vartheta-\delta_*\ge1/4\) and \(H_\ell\succeq\mathrm h_LI\)
make every non-root sheet drift uniformly confining toward zero, the root is a
fixed Ornstein--Uhlenbeck factor, proposals fix and projections contract the
configuration, and so the generator maps \(|\mathbf x|^{2k}\) below
\(\mathrm C_k-\mathrm c_k|\mathbf x|^{2k}\); all moments on \([0,t_*]\), hence
the constants of \cref{eq:rg-formation-local-law-dependence} and of the
covector, are uniform over such paths, and finitely many steps reach \(t_*\)
unless \(\vartheta\) leaves \([1,21/10]\), which the stopped argument below
excludes.

Stop at the first possible exit of \(\vartheta\) from \([1,21/10]\). The
Gaussian comparison
\(d\mathbf X^{\mathrm G}=-\vartheta(t)H_\ell\mathbf X^{\mathrm G}dt
+\sqrt{2\mathrm T} dW\), with the same parameter and Brownian increments on the
active level, starts from the reference of \Cref{def:rg-formation-regime}
coupled by \cref{eq:rg-formation-preparation-transport}; at an activation its
projected reference continues while the parent starts at the full retained
sample. For \(\mathbf e=\mathbf X-\mathbf X^{\mathrm G}\) the noise cancels and
proposals leave \(\mathbf e\) continuous, so
\cref{eq:rg-formation-parent-relative-bounds} gives, for every adapted label
history,
\begin{equation}
 \frac d{dt}|\mathbf e|^2\le-\mathbf e^\top H_\ell\mathbf e+\mathrm C_L\mathrm T,
 \qquad
 \sup_{t\le t_*}|\mathbf e_t|
 \le |\mathbf e_0|+\mathrm C_L\sqrt{\mathrm Tt_*}.
 \label{eq:rg-formation-path-force-bound}
\end{equation}
For the energy, variation of constants with the commuting propagator
\(\exp[-H_\ell\int_{t'}^t\vartheta]\) uses the first bound below, obtained by
splitting the integral at the inverse extreme eigenvalues, where the integrand
is at most \(\lambda_{\max}(H_\ell)\), \(1/(e\sigma)\), and
\(\lambda_{\min}(H_\ell)e^{-\lambda_{\min}(H_\ell)\sigma}\); Euclidean
projection and the Schur variational identity contract the two bounds at each
activation, and finitely many stages give the second:
\begin{equation}
 \int_0^\infty\|H_\ell e^{-H_\ell\sigma}\| d\sigma
 \le \mathrm C_L\left[1+\log
       \frac{\lambda_{\max}(H_\ell)}{\lambda_{\min}(H_\ell)}\right],
 \qquad
 \sup_t|H_\ell^{1/2}\mathbf e_t|\le \mathrm C_L\sqrt{\mathrm T}\varsigma_\epsilon.
 \label{eq:rg-formation-energy-force-bound}
\end{equation}
The Gaussian covariance barrier \(\Sigma^{\mathrm G}\preceq\mathrm C_L\mathrm TH_\ell^{-1}\)
needs only \(\vartheta\ge1\), by substitution into
\(\dot\Sigma^{\mathrm G}=-\vartheta(H_\ell\Sigma^{\mathrm G}+\Sigma^{\mathrm G}H_\ell)
+2\mathrm TI\), and survives retention because the retained block of
\(H_\ell^{-1}\) is the inverse Schur complement; the deterministic mean
contracts, has norm \(O_L(\mathrm A)\), and has zero loop part. This gives
\cref{eq:rg-formation-actual-energy,eq:rg-formation-actual-moments}, the loop
bound following from the core rate \(\vartheta\mathrm J_1\) and a force of at
most \(\mathrm C_L\sqrt{\mathrm T\mathrm J_1}\), whose forced error is at most
\(\mathrm C_L\sqrt{\mathrm T/\mathrm J_1}\).

Substituting these moments in the covector and metric bounds gives
\(\dot\vartheta=\epsilon(2-\vartheta)+O_L(\epsilon r_\epsilon)\), with
\(r_\epsilon\) of \cref{eq:rg-formation-parameter-remainder}, before any bound
on \(\dot\vartheta\) is used. The field points inward at both endpoints of
\([1,21/10]\) for small \(\epsilon\), which excludes the stopped exit; the
moments exclude explosion, so the solution reaches \(t_*\) with
\(|\dot\vartheta|\le\mathrm C_L\epsilon\). Positive-temperature elliptic
smoothing and positive finite rates keep densities positive and smooth, and
the construction restarts after an activation from the full retained law.
\hfill\(\square\)

\lemmaheading{Comparison of independently learned parameters}{lem:formation-parameter-comparison}
Every process or finite replacement lineage of
\Cref{lem:formation-prepared-continuation} satisfies
\cref{eq:rg-formation-parameter-bound} through \(t_*\), where every
independently evolving descendant uses its own law and copies only the
starting parameter. \status{ESTABLISHED}

\noindent\emph{Proof.}
The first two bounds are in \Cref{lem:formation-prepared-continuation}.
Subtracting \(\dot\vartheta_\circ=\epsilon(2-\vartheta_\circ)\),
\(\vartheta_\circ(0)=1\), from the equation of the process and integrating
with the factor \(e^{\epsilon t}\) bounds the forcing contribution by
\(\mathrm C_Lr_\epsilon\) through the horizon, where
\(r_\epsilon=\epsilon^{2\mathrm Q_1+20}(1+\epsilon^{\mathrm Q_1+6}\varsigma_\epsilon^2)\to0\)
by \cref{eq:rg-formation-parameter-remainder}, and copying the parameter at
finitely many activations adds a finite multiple of that error. The
comparison trajectory solves the calibration part of the own-law equation,
since the calibration record contributes \(\vartheta-2\) to the covector and
\(1/\mathrm T\) to \(\mathsf M\), and every law-dependent term is of order
\(\epsilon^{2\mathrm Q_1+20}\). As \(2-e^{-2}>7/4\), the coefficient exceeds
\(7/4\) on the last ten units of the horizon for small \(\epsilon\).
\hfill\(\square\)

\section{Entropy estimates}
\label{sec:rg-formation-entropy-proof}

\lemmaheading{Continuous entropy and score tools}{lem:formation-entropy-tools}
Let \(\mu(d\mathbf x)\propto e^{-\mathcal V(\mathbf x)/\mathrm T}d\mathbf x\) on
a finite-dimensional Euclidean space, where \(\mathcal V\) is smooth with a
globally bounded Hessian and \(\nabla^2\mathcal V\succeq\lambda I\),
\(\lambda>0\). For a positive density \(\mathsf f\) relative to \(\mu\), with
\(\operatorname{Ent}_\mu(\mathsf f)=\int\mathsf f\log(\mathsf f/\int\mathsf f d\mu) d\mu\),
for \(\mathfrak F\) with \(\operatorname{Lip}(\mathfrak F)\le\mathrm K\), for
every probability law \(\nu\), and whenever the right sides are finite,
\begin{align}
 \operatorname{Ent}_\mu(\mathsf f)
 &\le \frac{\mathrm T}{2\lambda}
           \int\frac{|\nabla\mathsf f|^2}{\mathsf f} d\mu,
 \label{eq:rg-formation-continuous-lsi}\\
 \log\E_\mu e^{\upsilon(\mathfrak F-\E_\mu\mathfrak F)}
 &\le\frac{\mathrm T\mathrm K^2\upsilon^2}{2\lambda},
 \qquad \upsilon\in\R,
 \label{eq:rg-formation-herbst-proof}\\
 \upsilon\E_\nu\mathfrak F&\le\KL(\nu\Vert\mu)+\log\E_\mu e^{\upsilon\mathfrak F},
 \label{eq:rg-formation-entropy-duality}
\end{align}
so \(|\E_\nu\mathfrak F-\E_\mu\mathfrak F|\le\mathrm K\sqrt{2\mathrm T\KL(\nu\Vert\mu)/\lambda}\).
A centered observable \(\xi\) that obeys the first line below obeys the second:
\begin{align}
 \log\E_\mu e^{\upsilon\xi}&\le
       \frac{\mathrm v\upsilon^2}{2(1-\mathrm b|\upsilon|)}\quad\text{for } |\upsilon|<\mathrm b^{-1},
 \label{eq:rg-formation-subgamma-proof}\\
 |\E_\nu\xi|
 &\le\sqrt{2\mathrm v\KL(\nu\Vert\mu)}+\mathrm b\KL(\nu\Vert\mu).
 \label{eq:rg-formation-score-duality}
\end{align}
Moreover \(\|\nu-\mu\|_{\mathrm{TV}}\le\sqrt{\KL(\nu\Vert\mu)/2}\). These
statements apply to every conditional sheet potential of the construction.
\status{ESTABLISHED}

\noindent\emph{Proof.}
For the semigroup \(P_t\) generated by
\(\mathrm T\Delta-\nabla\mathcal V\cdot\nabla\), the derivative of the
synchronous stochastic flow has operator norm at most \(e^{-\lambda t}\), so
weighted Cauchy--Schwarz gives
\(|\nabla P_t\mathsf f|^2/P_t\mathsf f\le e^{-2\lambda t}P_t(|\nabla\mathsf f|^2/\mathsf f)\)
and invariance of \(\mu\) gives
\(\mathcal I_\mu(P_t\mathsf f)\le e^{-2\lambda t}\mathcal I_\mu(\mathsf f)\) for
the Fisher information \(\mathcal I_\mu\). Integration by parts gives
\(d\operatorname{Ent}_\mu(P_t\mathsf f)/dt=-\mathrm T\mathcal I_\mu(P_t\mathsf f)\);
for bounded positive Lipschitz \(\mathsf f\), contraction against a stationary
initial sample gives \(P_t\mathsf f\to\int\mathsf f d\mu\) and, by bounded
convergence, entropy tending to zero, so integrating the Fisher bound proves
\cref{eq:rg-formation-continuous-lsi}. Cutoff, positive truncation, and short
semigroup smoothing extend it by lower semicontinuity, and Gaussian tails from
strong convexity justify the integrations. This is the curvature argument of
\citet[pp.~187, 194--200]{BakryEmery1985} with the diffusion coefficient made
explicit; its Gaussian case is \citet[Thm.~5.4]{BoucheronLugosiMassart2013}.
Applied to \(e^{\upsilon\mathfrak F}\), \cref{eq:rg-formation-continuous-lsi}
gives \(\operatorname{Ent}_\mu(e^{\upsilon\mathfrak F})\le
(\mathrm T\mathrm K^2\upsilon^2/(2\lambda))\E_\mu e^{\upsilon\mathfrak F}\), the
hypothesis of Herbst's argument \citep[Prop.~6.1]{BoucheronLugosiMassart2013}
with variance factor \(\mathrm T\mathrm K^2/\lambda\); applied to
\(\pm\mathfrak F\) it gives \cref{eq:rg-formation-herbst-proof}. The inequality
\cref{eq:rg-formation-entropy-duality} is the variational formula for relative
entropy \citep[Cor.~4.15]{BoucheronLugosiMassart2013}, and optimizing it over
\(\upsilon\) with \cref{eq:rg-formation-herbst-proof} gives the Lipschitz
consequence. Under \cref{eq:rg-formation-subgamma-proof}, the transportation
lemma \citep[Lemma~4.18]{BoucheronLugosiMassart2013} applied to \(\pm\xi\),
with the inverse convex conjugate \(\sqrt{2\mathrm vt}+\mathrm bt\) of the
sub-gamma function \citep[Sec.~2.4]{BoucheronLugosiMassart2013}, gives
\cref{eq:rg-formation-score-duality}; the case \(\KL=\infty\) is trivial. The
last inequality is Pinsker's inequality
\citep[Thm.~4.19]{BoucheronLugosiMassart2013}, with the constant \(1/2\) due to
\citet{Csiszar1967}. \hfill\(\square\)

\lemmaheading{Entropy at activations}{lem:formation-activation-entropy}
For the law of the active level \(\ell\), let
\(\mathfrak E_\ell(t)=\mathrm T\KL(q_t^\ell\Vert q^{\star,\ell}_{\vartheta_\ell(t)})\),
with the root factor canceled. At an activation time and on each ensuing
active interval,
\begin{align}
 \mathfrak E_\ell(t^{\mathrm{act}}_\ell)&\le \mathrm C_L[\mathrm A^2\mathrm J_{\ell+1}+\mathrm T\varsigma_\epsilon^2],
 \label{eq:rg-formation-activation-budget}\\
 \mathrm T^2\int_{t^{\mathrm{act}}_\ell}^{t}\mathcal I_\ell(\sigma) d\sigma
 &\le \mathrm C_L[\mathfrak E_\ell(t^{\mathrm{act}}_\ell)+\mathrm T],
 \label{eq:rg-formation-full-fisher-budget}
\end{align}
where \(\mathcal I_\ell\) is the full continuous relative Fisher information of
the joint law. The normalized parameter scores of the hidden Gaussian
conditionals and of the joint satisfy \cref{eq:rg-formation-subgamma-proof}
with constants \(\mathrm v_L,\mathrm b_L\) independent of the retained input,
of \(\mathrm T\), of the strengths, and of the source masses.
\status{ESTABLISHED}

\noindent\emph{Proof.}
The normalized hidden Gaussian conditionals have centers independent of
\(\vartheta\), and their logarithmic parameter scores, bounded marginal tilt
included, are centered versions of
\(-|\mathbf Z|^2/(2\vartheta)-\partial_\vartheta\psi_{\ell,j}\); the joint score
adds \(\partial_\vartheta\log\pi_j(\vartheta)\), of absolute value at most four.
A bounded tilt of a standard Gaussian satisfies
\begin{equation}
 \E|\mathbf Z|^{2m}
 \le e^{B_0}2^{\dim M_0/2}4^m m!,\qquad m\ge0,
 \label{eq:rg-formation-gaussian-score-moments}
\end{equation}
because \(r^m\le m!4^me^{r/4}\) and
\(\E e^{|\mathbf Z|^2/4}\le2^{\dim M_0/2}\), and the other score terms are
bounded by \cref{eq:rg-formation-tilt-score-proof}. Centering at most doubles
a geometric moment constant, and summing the exponential series without its
zero linear term, with \(\log(1+x)\le x\), gives
\cref{eq:rg-formation-subgamma-proof}, which conditional expectations keep by
Jensen's inequality.

For the differential and counting entropy
\(\mathcal S(q)=-\sum_j\int q_j\log q_j\) across changes of dimension, write
\(q_j=\lambda_jq_{\mid j}\), omitting terms of zero mass. The entropy--Fisher
bound \(\mathcal I(\mathsf f)\ge2\pi e \mathrm d e^{-2\mathcal S(\mathsf f)/\mathrm d}\)
of \citet{Stam1959}, equivalent to the Gaussian logarithmic Sobolev inequality
\citep{Carlen1991}, follows from the Gaussian case of
\cref{eq:rg-formation-continuous-lsi} by expanding both sides relative to
\(\mathcal N(0,\sigma^2I)\), integrating \(\int\mathbf x\cdot\nabla\mathsf f=-\mathrm d\),
canceling the second moments, and minimizing at
\(\sigma^2=\mathrm d/\mathcal I(\mathsf f)\); Jensen's inequality and
\(\mathcal S(q)=\sum_j\lambda_j\mathcal S(q_{\mid j})+\mathcal S(\lambda)\), with
\(\mathcal S(\lambda)\ge0\), extend it to
\(\sum_j\lambda_j\mathcal I(q_{\mid j})\ge2\pi e \mathrm d e^{-2\mathcal S(q)/\mathrm d}\).
At fixed \(\mathbf x\), a proposal moves the conditional label row
\(\mathfrak a\) in the direction \(\mathfrak a\beta^\ell-\mathfrak a\) and leaves
the continuous marginal unchanged, and concavity of categorical entropy bounds
the directional derivative below by
\(\mathcal S(\mathfrak a\beta^\ell)-\mathcal S(\mathfrak a)\ge-\log|\mathcal J|\).
On a \(\mathrm d\)-dimensional sheet, diffusion therefore gives the first line
below, and the two bounds together give the second:
\begin{align}
 \mathcal S'(q)_{\mathrm{diff}}
 &=\mathrm T\sum_j\lambda_j\mathcal I(q_{\mid j})-\E_q\Delta V_{\ell,S}
 \ge \mathrm T\sum_j\lambda_j\mathcal I(q_{\mid j})-\mathrm d\Lambda_\ell,
 \qquad \Lambda_\ell\le \mathrm C_L\mathrm J_1,
 \label{eq:rg-formation-actual-entropy-derivative}\\
 \mathcal S'(q)&\ge 2\pi e \mathrm d \mathrm T e^{-2\mathcal S(q)/\mathrm d}
                -\mathrm d\Lambda_\ell-\log|\mathcal J|.
 \label{eq:rg-formation-counting-entropy-floor}
\end{align}
Scalar comparison bounds \(\mathcal S(q)\) below by the lesser of its initial
value and
\(\frac{\mathrm d}2\log[2\pi e \mathrm d\mathrm T/(\mathrm d\Lambda_\ell+\log|\mathcal J|)]\).
At a projection dropping \(k\) coordinates \(z\), the chain rule and the
maximum-entropy Gaussian bound, applied conditionally on \((Y,S)\) and then
with Jensen's inequality, give
\begin{equation}
 \mathcal S(Y,S)\ge \mathcal S(\mathbf X,S)
       -\frac k2\log\frac{2\pi e \E|z-\bar z|^2}{k},
 \label{eq:rg-formation-entropy-projection}
\end{equation}
where \(\bar z\) can be the projected comparison mean, whose centered
covariance the moment proof of \Cref{lem:formation-prepared-continuation}
bounds by \(\mathrm C_L\mathrm T\varsigma_\epsilon^2H_\ell^{-1}\). Initial
entropies and all logarithms in these estimates are
\(O_L(\varsigma_\epsilon^2)\), because the covariance eigenvalues, \(\mathrm A\),
and \(\mathrm T\) are fixed powers of \(\epsilon\), the admitted initial density
ratios are bounded, and there are finitely many projections. Expanding
\(\mathrm T\KL(q\Vert\gamma_{\ell,\vartheta}\otimes\pi(\vartheta))\) and using
\cref{eq:rg-formation-actual-energy} proves
\cref{eq:rg-formation-activation-budget}; the bounded optimizer tilt adds at
most \(\mathrm TB_0\), and the source prior has a positive minimum on
\([1,21/10]\), so its logarithmic contribution does not depend on the masses
\(\lambda_j\).

For the time-integrated estimate put \(\mathcal V=y^\top H_\ell y/2\) and
\(\mathrm R_L=e^{B_0}2^{\dim M_0/2}\). The reference satisfies the first bound
below. Entropy duality gives
\(\vartheta\E_q\mathcal V\le2\mathfrak E_\ell+2\mathrm T\log\mathrm R_L\),
Jensen's inequality gives
\(\vartheta\E_{q^{\star,\ell}_\vartheta}\mathcal V\le2\mathrm T\log\mathrm R_L\),
and the other centered terms of the normalized full score are at most
\(2\mathrm C_{\mathrm{sc},L}\), which gives the second. Diffusion dissipates
\(\mathrm T^2\mathcal I_\ell\) and detailed balance makes the source
contribution nonpositive, which gives the third for a fixed
\(\mathrm C_{0,L}\):
\begin{align}
 \E_{q^{\star,\ell}_\vartheta}\exp\frac{\vartheta\mathcal V}{2\mathrm T}
 &\le \mathrm R_L,
 \label{eq:rg-formation-energy-exponential-moment}\\
 |\mathrm T\E_q\partial_\vartheta\log q^{\star,\ell}_\vartheta|
 &\le\frac{2\mathfrak E_\ell+4\mathrm T\log\mathrm R_L}{\vartheta}+2\mathrm T\mathrm C_{\mathrm{sc},L},
 \label{eq:rg-formation-entropy-parameter-budget}\\
 \mathfrak E_\ell'&\le -\mathrm T^2\mathcal I_\ell+
    \frac{2|\dot\vartheta|}{\vartheta}\bigl(\mathfrak E_\ell+\mathrm C_{0,L}\mathrm T\bigr).
 \label{eq:rg-formation-entropy-energy-inequality}
\end{align}
The integrating factor \(\exp[-2\int|\dot\vartheta|/\vartheta]\) has exponent
bounded over \(t_*\) by the speed bound, and integration proves
\cref{eq:rg-formation-full-fisher-budget}. At an activation the law and the
reference are pushed through the same channel at the same parameter, so the
chain rule makes \(\mathfrak E\) decrease along the finite lineage. All
entropy derivatives are first taken for positive truncated densities on
bounded domains, and the moment and entropy bounds remove the cutoffs.
\hfill\(\square\)

\lemmaheading{Conditional entropy and complete score error}{lem:formation-conditional-entropy}
There is a constant \(\mathrm c_{\mathrm{early},L}>0\) such that, for every
retention \(\ell-1\to\ell\), the conditional relative entropy \(\Delta_\ell\)
with respect to the optimizing hidden conditional satisfies
\begin{equation}
 \sup_{t\in I_\ell\text{ and later active times}}
 \Delta_\ell(t)
 \le
 \begin{cases}
 \mathrm C_L\epsilon^{14},&\ell<L,\\
 \mathrm C_L\epsilon^2,&\ell=L,
 \end{cases}
 \label{eq:rg-formation-conditional-final-bounds}
\end{equation}
where the early bound also applies, until the next activation, to the
corresponding unreplaced comparison process, and the terminal bound is
\Cref{lem:formation-terminal-joint-entropy}. At every such time the complete
fine/parent covector difference is at most
\(\mathrm C_L\mathrm T(\sqrt{\Delta_\ell}+\Delta_\ell)\). \status{ESTABLISHED}

\noindent\emph{Proof of the early bound and the covector estimate.}
For an early retention use the auxiliary continuous split of the first line
below. The true parent retains the loops, so the chain rule bounds its hidden
conditional defect \(\Delta_\ell\) by the auxiliary defect
\(\widehat\Delta_\ell\), which also integrates \(h\); the independent root
cancels, and conditional strong convexity in \(Z\) is at least
\(\lambda_\ell=\mathrm c_{\mathrm{early},L}\mathrm J_\ell\) after reducing the
constant if necessary. The mixed Hessian obeys the second line:
\begin{align}
 &Z=(\mathcal U_{\ell-1}\cap\mathcal U_\ell^\perp)\oplus\mathcal H,
 \qquad Y=\mathcal U_\ell,
 \qquad S\ \text{retained},
 \label{eq:rg-formation-auxiliary-split}\\
 &\|\nabla_Z\nabla_YV_{\ell-1,j}\|\le\Gamma_\ell=\mathrm C_L\mathrm J_{\ell+1},
 \qquad \Gamma_L=0.
 \label{eq:rg-formation-mixed-hessian-scale}
\end{align}
The nonlinear residuals lie in the first normal layer and the loop space, and
at \(\ell-1=0\) the nonlinear potential has no future normal dependence. For
\(\ell-1\ge1\), in coordinates \(r_E=B_Ew_E\) on the eliminated normal rows,
the Gaussian conditional precision
\(\Omega_E=\mathrm J_E+\mathcal C^\top\mathrm J_{\mathrm{fut}}\mathcal C\), with
the diagonal strengths \(\mathrm J_E\) and \(\mathrm J_{\mathrm{fut}}\) of the
eliminated and future rows, has \(\|\Omega_E-\mathrm J_E\|\le\mathrm C_L\mathrm J_\ell\)
and \(\|\Omega_E^{-1}\|\le\mathrm C_L/\mathrm J_{\ell-1}\), so the identity in the
first line below bounds the first block row of \(\Omega_E^{-1}\) by
\(\mathrm C_L/\mathrm J_1\). A future force column of level \(r\) has size
\(\mathrm C_L\mathrm J_r\), so the first conditional mean responds by
\(\mathrm C_L\mathrm J_r/\mathrm J_1\), with conditional covariance comparable to
\(\mathrm T/\mathrm J_1\). In the derivatives of the normalized Gaussian
integral defining \(\psi_{\ell-1,j}\), a first-layer or loop derivative of
\(\phi_j\) has size \(\mathrm C_L\sqrt{\mathrm J_1/\mathrm T}\) and a second
derivative size \(\mathrm C_L\mathrm J_1/\mathrm T\), each future derivative
brings that mean response, and the covariance term is bounded by products of
first-derivative bounds, which gives the second line:
\begin{align}
 \Omega_E^{-1}&=\mathrm J_E^{-1}
           -\mathrm J_E^{-1}(\Omega_E-\mathrm J_E)\Omega_E^{-1},
 \label{eq:rg-formation-first-row-resolvent}\\
 \mathrm T|\partial_h\partial_{w_r}\psi_{\ell-1,j}|&\le \mathrm C_L\mathrm J_r,\qquad
 \mathrm T|\partial_{w_i}\partial_{w_r}\psi_{\ell-1,j}|
       \le \mathrm C_L\frac{\mathrm J_i\mathrm J_r}{\mathrm J_1}.
 \label{eq:rg-formation-future-derivative-scales}
\end{align}
The Gaussian Schur cross block is at most \(\mathrm C_L\mathrm J_{\ell+1}\) by
the lower triangular residual form and its variational bound, and summing
finitely many blocks, generated conditional covariance included, proves
\cref{eq:rg-formation-mixed-hessian-scale}.

Write \(q=q_Y(y,j)q_{Z\mid Y}(z\mid y,j)\) and
\(q^\star=q^\star_Y(y,j)q^\star_{Z\mid Y}(z\mid y,j)\), and let
\(\mathfrak w_Y=\nabla_Y\log(q/q^\star)\),
\(\mathfrak v_Y=\nabla_Y\log(q_Y/q^\star_Y)\), and
\(\mathfrak e_Y=\E_{q_{Z\mid Y}}\nabla_Y\log q^\star_{Z\mid Y}\). Normalization
gives \(\E_{q^\star_{Z\mid Y}}\nabla_Y\log q^\star_{Z\mid Y}=0\) and
\(\E_{q_{Z\mid Y}}\mathfrak w_Y=\mathfrak v_Y-\mathfrak e_Y\), the marginal
flux is \(\mathrm Tq_Y(\mathfrak v_Y-\mathfrak e_Y)\), and subtracting the
marginal entropy derivative from the full one yields
\begin{align}
 \widehat\Delta_\ell'|_{\mathrm{cont}}
 ={}&-\mathrm T\mathcal I_Z
      -\mathrm T\E_{q_Y}\operatorname{Var}_{q_{Z\mid Y}}(\mathfrak w_Y)
      +\mathrm T\E_{q_Y}\bigl[\mathfrak e_Y\cdot\E_{q_{Z\mid Y}}\mathfrak w_Y\bigr]
 \notag\\
 &-\dot\vartheta \E_q\partial_\vartheta\log q^\star_{Z\mid Y}.
 \label{eq:rg-formation-conditional-balance}
\end{align}
The conditional logarithmic Sobolev inequality gives
\(\mathrm T\mathcal I_Z\ge2\lambda_\ell\widehat\Delta_\ell\). Applying
\cref{eq:rg-formation-herbst-proof} to each unit-vector component of
\(\nabla_Y\log q^\star_{Z\mid Y}\), whose Lipschitz constant in \(Z\) is at most
\(\Gamma_\ell/\mathrm T\), gives
\(|\mathfrak e_Y|\le\Gamma_\ell\sqrt{2\KL(q_{Z\mid Y}\Vert q^\star_{Z\mid Y})/(\mathrm T\lambda_\ell)}\),
so Cauchy--Schwarz and Young's inequality bound the mixed term by
\(\lambda_\ell\widehat\Delta_\ell/2+\mathrm T\Gamma_\ell^2\mathcal I_Y/\lambda_\ell^2\),
with \(\mathcal I_Y\) the \(Y\)-component of the full relative Fisher
information. The normalized conditional parameter score, sub-gamma by
\Cref{lem:formation-activation-entropy}, contributes at most
\(\lambda_\ell\widehat\Delta_\ell/2+\mathrm v_L|\dot\vartheta|^2/\lambda_\ell
+\mathrm b_L|\dot\vartheta|\widehat\Delta_\ell\).

For the source updates, at fixed \(Y=y\) each
\(\log(\gamma^{\mathrm{hid}}/q^\star_{Z\mid Y,j})\), with \(\gamma^{\mathrm{hid}}\)
the common label-independent hidden Gaussian, lies in \([-B_0,B_0]\), and the
conditional entropy relative to \(\gamma^{\mathrm{hid}}\) splits as
\(\KL(q_Z\Vert\gamma^{\mathrm{hid}})+\mathbb I(Z;S)\). A proposal leaves \(q_Z\)
unchanged, and the mutual information is convex in the conditional label
channel at fixed \(q_Z\) by joint convexity of relative entropy, so its
directional derivative toward the new channel is at most \(\log|\mathcal J|\);
the difference from the optimizing target, a linear expectation of
\(\log(\gamma^{\mathrm{hid}}/q^\star_{Z\mid Y,j})\), changes at rate at most
\(2B_0\). This divides by no source mass and extends to vanishing masses by
convex approximation. With \(\mathrm b_L|\dot\vartheta|\le\lambda_\ell/2\) for
small \(\epsilon\) and the negative variance term dropped,
\begin{align}
 \widehat\Delta_\ell'|_{\mathrm{jump}}&\le
 \mathrm C_{\mathcal J}:=\log|\mathcal J|+2B_0,
 \label{eq:rg-formation-conditional-jump-remainder}\\
 \widehat\Delta_\ell'
 &\le-\frac{\lambda_\ell}{2}\widehat\Delta_\ell
     +\frac{\mathrm T\Gamma_\ell^2}{\lambda_\ell^2}\mathcal I_Y
     +\frac{\mathrm v_L|\dot\vartheta|^2}{\lambda_\ell}
     +\mathrm C_{\mathcal J}.
 \label{eq:rg-formation-conditional-differential-bound}
\end{align}
The initial conditional entropy is at most
\(\mathfrak E_{\ell-1}(t^{\mathrm{act}}_{\ell-1})/\mathrm T\). Integrating with
\cref{eq:rg-formation-full-fisher-budget,eq:rg-formation-activation-budget}
and \(\mathrm A^2/\mathrm T=\epsilon^{-6}\) gives
\begin{align}
 \widehat\Delta_\ell(t)\le \mathrm C_L\bigg[
 &e^{-\mathrm c_{\mathrm{early},L}\mathrm J_\ell(t-t^{\mathrm{act}}_{\ell-1})}
                 (\epsilon^{-6}\mathrm J_\ell+\varsigma_\epsilon^2)
 \notag\\
 &+\epsilon^{-6}\frac{\mathrm J_{\ell+1}^2}{\mathrm J_\ell}
 +\varsigma_\epsilon^2\frac{\mathrm J_{\ell+1}^2}{\mathrm J_\ell^2}
 +\frac{\epsilon^2}{\mathrm J_\ell^2}+\frac{\mathrm C_{\mathcal J}}{\mathrm J_\ell}\bigg],
 \label{eq:rg-formation-early-entropy-bound}
\end{align}
with \(\mathrm c_{\mathrm{early},L}\) absorbing the factor \(1/2\). For
\(\ell<L\), \(\mathrm J_{\ell+1}^2/\mathrm J_\ell=\epsilon^{20}\) and
\(\mathrm J_\ell\ge\epsilon^{-20}\); on \(I_\ell\) the entrance time and
\(\mathrm c_{\mathrm{early},L}k_L/2>\mathrm Q_1+20\) bound the first term by
\(\mathrm C_L\epsilon^{14}\), and the other terms meet the same bound; at
\(\ell=L\) the last term would be \(\mathrm C_{\mathcal J}\).

For the true retention, \cref{eq:rg-formation-conditional-cost} with
\(\ell'=\ell-1\), differentiated at fixed recognition law and channel, gives
\(g_{\ell-1}-g_\ell=-\mathrm T\E_q\partial_\vartheta\log q^\star_{Z\mid Y,\mathrm{true}}\),
the common scalar and the source-prior derivative canceling. The normalized
conditional score has the same sub-gamma bound, so
\cref{eq:rg-formation-score-duality}, applied conditionally, with Jensen's
inequality gives the covector bound
\(|g_{\ell-1}-g_\ell|\le\mathrm C_L\mathrm T(\sqrt{\Delta_\ell}+\Delta_\ell)\).
\(\square\)

\lemmaheading{Joint relaxation at the weakest strength}{lem:formation-terminal-joint-entropy}
There is a constant \(\mathrm c_{\mathrm{joint},L}>0\), independent of
\(\epsilon\), such that every lineage satisfies
\begin{equation}
 \KL\bigl(q_t^\ell\Vert q^{\star,\ell}_{\vartheta_\ell(t)}\bigr)
 \le \mathrm C_L e^{-\mathrm c_{\mathrm{joint},L}t}
           [\epsilon^{-(\mathrm Q_1+6)}+\varsigma_\epsilon^2]
       +\mathrm C_L\epsilon^2.
 \label{eq:rg-formation-terminal-joint-bound}
\end{equation}
Consequently the terminal part of
\cref{eq:rg-formation-conditional-final-bounds} holds once
\(t\ge t_L=k_L\log(1/\epsilon)\). \status{ESTABLISHED}

\noindent\emph{Proof.}
Fix a level and a parameter, write \(q^\star=q^{\star,\ell}_\vartheta\) and
\(\mathfrak r_j=q_j/q^\star_j\). The bounded tilt gives
\(e^{-B_0}\le q^\star/(\gamma_{\ell,\vartheta}\otimes\pi)\le e^{B_0}\), and
\cref{eq:rg-formation-parent-generator} gives the conductance comparison in the
first line below. For a probability measure \(\mu\),
\(\operatorname{Ent}_\mu(\mathfrak r)=\inf_{r>0}\int[\mathfrak r\log(\mathfrak r/r)-\mathfrak r+r] d\mu\)
with a nonnegative integrand, so
\(\operatorname{Ent}_{q^\star}(\mathfrak r)\le e^{B_0}\operatorname{Ent}_{\gamma\otimes\pi}(\mathfrak r)\),
and the product chain rule and convexity give the second line, whose first
term the Gaussian entropy inequality bounds by
\(\mathrm T/(2\vartheta\mathrm h_L)\) times its continuous Fisher information:
\begin{align}
 q^\star_j(y)\beta_{jk}^\ell(y)
 &\ge e^{-2B_0}\gamma_{\ell,\vartheta}(y)\pi_j(\vartheta)\pi_k(\vartheta),
 \label{eq:rg-formation-conductance-comparison}\\
 \operatorname{Ent}_{\gamma\otimes\pi}(\mathfrak r)
 &\le\sum_j\pi_j\operatorname{Ent}_\gamma(\mathfrak r_j)
   +\E_\gamma\operatorname{Ent}_{\pi}(\mathfrak r(y)).
 \label{eq:rg-formation-product-entropy-bound}
\end{align}
For any positive finite vector \(\mathbf a\), Jensen's inequality for the
logarithm gives
\begin{equation}
 \frac12\sum_{j,k}\pi_j\pi_k
       (\mathbf a_j-\mathbf a_k)(\log\mathbf a_j-\log\mathbf a_k)
 =\sum_j\pi_j\mathbf a_j\log\mathbf a_j
       -\bar{\mathbf a}\sum_j\pi_j\log\mathbf a_j
 \ge \operatorname{Ent}_\pi(\mathbf a),
 \label{eq:rg-formation-categorical-entropy-bound}
\end{equation}
with \(\bar{\mathbf a}=\sum_j\pi_j\mathbf a_j\). Comparing the Fisher measures
and the conductances through their pointwise bounds, with the jump
dissipation of the first line below, gives the second:
\begin{align}
 \mathcal D_{\mathrm{jump}}
 &=\frac12\int\sum_{j,k}q^\star_j\beta_{jk}^\ell
           (\mathfrak r_j-\mathfrak r_k)(\log\mathfrak r_j-\log\mathfrak r_k) dy,
 \label{eq:rg-formation-jump-dissipation}\\
 \KL(q\Vert q^\star)&\le
 e^{2B_0}\frac{\mathrm T}{2\vartheta\mathrm h_L}\mathcal I_\ell
       +e^{3B_0}\mathcal D_{\mathrm{jump}}.
 \label{eq:rg-formation-joint-coercivity}
\end{align}
With \(\mathrm c_{0,L}=\min\{2\mathrm h_Le^{-2B_0},e^{-3B_0}\}>0\),
\(\mathrm T\mathcal I_\ell+\mathcal D_{\mathrm{jump}}\ge\mathrm c_{0,L}\KL\). The time-dependent normalized
full score satisfies \cref{eq:rg-formation-subgamma-proof}, hence, for small
\(\epsilon\),
\begin{equation}
 \KL'\le-\mathrm c_{0,L}\KL+
       \mathrm C_L\epsilon(\sqrt{\KL}+\KL)
 \le-\frac{\mathrm c_{0,L}}{2}\KL+\mathrm C_L\epsilon^2,
 \label{eq:rg-formation-joint-moving-reference}
\end{equation}
and \(\mathrm c_{\mathrm{joint},L}=\mathrm c_{0,L}/2\). Activations decrease the
relative entropy at the unchanged parameter, and each descendant then obeys
the same inequality with its own parameter speed. The initial relative entropy
is at most \(\mathrm C_L[\epsilon^{-(\mathrm Q_1+6)}+\varsigma_\epsilon^2]\), and
integration through finitely many decreases proves
\cref{eq:rg-formation-terminal-joint-bound}; the condition
\(\mathrm c_{\mathrm{joint},L}k_L>\mathrm Q_1+8\) makes its decaying term
\(O_L(\epsilon^2)\) at \(t_L\). Both the auxiliary and the true hidden
conditional entropies are bounded by the full joint relative entropy, which
proves the terminal assertion. \hfill\(\square\)

\section{Exact fast marginals and complete source histories}
\label{sec:rg-formation-history-proof}

\lemmaheading{Common fast marginal and environmental error}{lem:formation-fast-marginal}
In the flag coordinates \(w=(w_1,\ldots,w_L)\) of
\cref{eq:rg-formation-weighted-flag} and the loop coordinate \(h\), put
\(\Omega_*=\mathrm J_1\mathsf B_{11}^\top\mathsf B_{11}\) and
\begin{equation}
 \mu_{\vartheta,j}^*(dw_1,dh)
 =\frac{1}{Z_{\vartheta,j}^*}
   \exp\left\{-\frac{\vartheta}{2\mathrm T}
         (w_1^\top\Omega_*w_1+\mathrm J_1|h|^2)-\phi_j(w_1,h)\right\}dw_1 dh,
 \label{eq:rg-formation-fast-law}
\end{equation}
with loop marginal \(\mu_{\vartheta,j}^{*,h}\), first-normal conditional
\(\mu_{\vartheta,j}^{*,1}(dw_1\mid h)\), and, with environment scales
\(\varrho_\epsilon\), \(\varrho'_\epsilon\),
\begin{align}
 \beta_{jk}^{*,h}(h,\vartheta)
 &=\int\beta_k(w_1,h,\vartheta)
                   \mu_{\vartheta,j}^{*,1}(dw_1\mid h),
 \label{eq:rg-formation-fast-loop-row}\\
 \bar\beta_{jk}(\vartheta)
 &=\int\beta_k(w_1,h,\vartheta) \mu_{\vartheta,j}^*(dw_1,dh)
   =\int\beta_{jk}^{*,h}(h,\vartheta)
                    \mu_{\vartheta,j}^{*,h}(dh),
 \label{eq:rg-formation-common-averaged-row}\\
 \varrho_\epsilon&=\frac{\mathrm J_2}{\mathrm J_1},\qquad
 \varrho'_\epsilon=\frac{\mathrm J_2}{\sqrt{\mathrm J_1\mathrm T}},\qquad
 \mathrm e_\epsilon(y)=\varrho_\epsilon+\varrho'_\epsilon|y|.
 \label{eq:rg-formation-environment-scales}
\end{align}
Then \cref{eq:rg-formation-fast-law} is the exact \((w_1,h)\)-marginal of the
optimizer conditional on \(S=j\); averaging the inherited row
\(\beta^\ell_{jk}\) against the optimizer conditional on \(j\) gives the same
\(\bar\beta_{jk}(\vartheta)\) for every \(\ell\); and the exact parent rows and
loop forces satisfy
\cref{eq:rg-formation-parent-row-environment,eq:rg-formation-parent-force-environment,eq:rg-formation-row-fast-lipschitz}
below. \status{ESTABLISHED}

\noindent\emph{Proof.}
In the flag coordinates
\begin{equation}
 w^\top H_{\mathcal U}w
 =\mathrm J_1|\mathsf B_{11}w_1|^2
       +\sum_{\ell=2}^L\mathrm J_\ell
             \left|\sum_{r\le\ell}\mathsf B_{\ell r}w_r\right|^2,
 \label{eq:rg-formation-triangular-fast-quadratic}
\end{equation}
and \(\phi_j\) depends only on \((w_1,h)\). For fixed \(w_1,h\), change the
later \(w_\ell\) to the residuals \(r_\ell=\sum_{r\le\ell}\mathsf B_{\ell r}w_r\);
the domain stays the whole Euclidean product and the Jacobian is the constant
\(\prod_{\ell\ge2}|\det\mathsf B_{\ell\ell}|\), so integrating those Gaussian
residuals contributes a scalar independent of \(w_1,h,j\), which proves the
first claim, the prior factor \(\pi_j(\vartheta)\) canceling on conditioning.
The change of variables and the tower property give the second.

For an intermediate \(\ell\ge1\), keep \(y=(w_{\ell+1},\ldots,w_L)\) fixed and
integrate \(w_2,\ldots,w_\ell\). The unperturbed conditional first-normal law
\(\gamma^1_{\ell,\vartheta,y}\) and its divergence from
\(\gamma^{1,*}_\vartheta=\mathcal N(0,(\mathrm T/\vartheta)\Omega_*^{-1})\) are
\begin{align}
 &\gamma^1_{\ell,\vartheta,y}
 =\mathcal N\left(\mu_\ell(y),\frac{\mathrm T}{\vartheta}(\Omega^{(1)}_\ell)^{-1}\right),\qquad
 \Omega_*\preceq\Omega^{(1)}_\ell\preceq\Omega_*+\mathrm C_L\mathrm J_2I,\qquad
 |\mu_\ell(y)|\le \mathrm C_L\frac{\mathrm J_2}{\mathrm J_1}|y|,
 \label{eq:rg-formation-first-normal-conditional}\\
 &\frac12\left[
  \frac{\vartheta}{\mathrm T}\mu_\ell^\top\Omega_*\mu_\ell
  +\Tr\bigl(\Omega_*(\Omega^{(1)}_\ell)^{-1}\bigr)-\dim w_1
            -\log\det\bigl(\Omega_*(\Omega^{(1)}_\ell)^{-1}\bigr)\right]
 \le \mathrm C_L[\varrho_\epsilon^2+\varrho_\epsilon'^2|y|^2].
 \label{eq:rg-formation-first-normal-gaussian-kl}
\end{align}
For the lower precision bound, minimize the nonnegative later residual squares
at \(y=0\); for the upper, set the eliminated later coordinates to zero and use
\(\mathrm J_\ell\le\mathrm J_2\) for \(\ell\ge2\); at \(\ell=L\),
\(\Omega^{(1)}_L=\Omega_*\) and \(\mu_L=0\). The mean bound is the
first-block-row estimate \cref{eq:rg-formation-first-row-resolvent}, each
future coordinate of level \(\ell\) contributing at most
\(\mathrm C_L\mathrm J_\ell/\mathrm J_1\). The divergence bound holds because
the whitened covariance eigenvalues lie between
\((1+\mathrm C_L\varrho_\epsilon)^{-1}\) and one, where
\(t-1-\log t=O((t-1)^2)\), and Pinsker's inequality turns it into a total
variation at most \(\mathrm C_L\mathrm e_\epsilon(y)\).

Tilting two probability measures by one weight \(w\in[e^{-B_0},1]\) multiplies
their total-variation distance by at most \(2e^{B_0}\), since each term of
\(w\mu/\mu(w)-w\nu/\nu(w)=w(\mu-\nu)/\mu(w)+w\nu[1/\mu(w)-1/\nu(w)]\) has
total variation at most \(e^{B_0}\|\mu-\nu\|_{\mathrm{TV}}\). Applied to
\(w=e^{-\phi_j(w_1,h)}\) uniformly in \(h\), this gives the first bound below;
the second follows through the common Gaussian force
\(\vartheta\mathrm J_1h\) and \(\mathrm T\E\nabla_h\phi_j\), and the third
because \(\gamma^{1,*}_\vartheta\) does not depend on \(h\), so differentiation
under its normalized tilt gives
\[
\nabla_h\beta^{*,h}_{jk}
=\E_{\mu^{*,1}_{\vartheta,j}(\cdot\mid h)}\nabla_h\beta_k
-\operatorname{Cov}_{\mu^{*,1}_{\vartheta,j}(\cdot\mid h)}
 (\beta_k,\nabla_h\phi_j),
\]
and both terms have the asserted Lipschitz bound: the standardized source
potential has bounded first derivatives and \(\sum_k\beta_k=1\). Thus
\begin{align}
 \|\beta_{j\cdot}^\ell(y,h,\vartheta)
          -\beta_{j\cdot}^{*,h}(h,\vartheta)\|_{\mathrm{TV}}
 &\le \mathrm C_L\mathrm e_\epsilon(y),
 \label{eq:rg-formation-parent-row-environment}\\
 |\nabla_hV_{\ell,j}(y,h,\vartheta)
           -\nabla_hV_j^{*,h}(h,\vartheta)|
 &\le \mathrm C_L\sqrt{\mathrm T\mathrm J_1} \mathrm e_\epsilon(y),
 \label{eq:rg-formation-parent-force-environment}\\
 \operatorname{Lip}_{(w_1,h)}\beta_{j\cdot}
       +\operatorname{Lip}_h\beta_{j\cdot}^{*,h}
 &\le \mathrm C_L\sqrt{\mathrm J_1/\mathrm T}.
 \label{eq:rg-formation-row-fast-lipschitz}
\end{align}
\hfill\(\square\)

\lemmaheading{Fast mixing with a learned coefficient}{lem:formation-fast-mixing}
There is a constant \(\mathrm c_{\mathrm{fast},L}>0\) with the following property.
Let the auxiliary process have, on a proposal-free interval, fixed assignment
\(j\) and either the full fast potential of \cref{eq:rg-formation-fast-law} or
its exact loop marginal, and a deterministic coefficient with
\(1\le\vartheta\le21/10\) and \(|\dot\vartheta|\le\mathrm C_L\epsilon\). For a
point start \(\xi\) at time \(t'\) and \(t-t'\ge(\mathrm c_{\mathrm{fast},L}\mathrm J_1)^{-1}\),
\begin{equation}
 \|P^*_{t',t}(\xi,\cdot)-\mu^*_{\vartheta(t),j}\|_{\mathrm{TV}}
 \le \mathrm C_L\bigl(1+\sqrt{\mathrm J_1/\mathrm T}|\xi|\bigr)
              e^{-\mathrm c_{\mathrm{fast},L}\mathrm J_1(t-t')}
       +\mathrm C_L\epsilon/\mathrm J_1,
 \label{eq:rg-formation-fast-mixing-bound}
\end{equation}
where \(\mu^*\) means its loop marginal in the loop case. Also
\(\bar\beta_{jk}(\vartheta)\) is uniformly Lipschitz in \(\vartheta\).
\status{ESTABLISHED}

\noindent\emph{Proof.}
The relative Hessian bounds and the normalized marginal covariance calculation
give \(\mathrm c_L\mathrm J_1I\preceq\nabla^2V^*\preceq\mathrm C_L\mathrm J_1I\)
for both potentials, which are even, minimize at zero, and have equilibrium
second moments at most \(\mathrm C_L\mathrm T/\mathrm J_1\). For two point
starts \(\xi,\xi'\), reflection coupling before meeting gives their distance a
one-dimensional noise of size \(2\sqrt{2\mathrm T}\) and a drift at most
\(-\lambda\) times the distance, \(\lambda=\mathrm c_L\mathrm J_1\); multiplying
by \(e^{\lambda(t-t')}\) gives a time-changed Brownian motion of variance
\(4\mathrm T(e^{2\lambda(t-t')}-1)/\lambda\), and the reflection principle
bounds the probability of not meeting by
\begin{equation}
 \mathrm C|\xi-\xi'|\sqrt{\lambda/\mathrm T}
                 \bigl(e^{2\lambda(t-t')}-1\bigr)^{-1/2}.
 \label{eq:rg-formation-reflection-bound}
\end{equation}
The coupling reflects one Brownian motion orthogonally and predictably, so
both marginals have the required law; the drifts are globally Lipschitz for
each instance, so the copies are constructed away from the diagonal and joined
at meeting, the comparison and hitting bound following by stopping before the
diagonal and letting the stopping radius tend to zero. Averaging \(\xi'\)
against the equilibrium at time \(t'\) gives the first term of
\cref{eq:rg-formation-fast-mixing-bound}.

For the second, compare that evolving equilibrium start with the
instantaneous equilibrium at time \(t\). After standardizing \((w_1,h)\) by
\(\sqrt{\mathrm J_1/\mathrm T}\), the normalized parameter score is a centered
quadratic with uniform exponential moments, and the loop-marginal score, its
conditional expectation, has the same bound by Jensen's inequality. The
moving-reference entropy calculation with the continuous entropy inequality at
curvature \(\mathrm c_L\mathrm J_1\) gives, for the relative entropy
\(\mathrm D_{\mathrm{lag}}\) of the evolved law to the current equilibrium,
\begin{equation}
 \mathrm D_{\mathrm{lag}}'\le-\mathrm c_L\mathrm J_1\mathrm D_{\mathrm{lag}}
 +\mathrm C_L\epsilon^2/\mathrm J_1,\qquad \mathrm D_{\mathrm{lag}}(t')=0,
 \label{eq:rg-formation-fast-equilibrium-lag}
\end{equation}
the parameter term linear in \(\mathrm D_{\mathrm{lag}}\) being absorbed for
small \(\epsilon\); hence \(\mathrm D_{\mathrm{lag}}\le\mathrm C_L\epsilon^2/\mathrm J_1^2\),
and Pinsker's inequality gives the second term. Finally, differentiating
\(\bar\beta_{jk}=\E_{\mu_{\vartheta,j}^*}\beta_k\) gives the explicit row
derivative, whose absolute values sum to at most eight, and a covariance with
the centered scale score, whose standardized first moment is bounded, so
\(\sum_k|\partial_\vartheta\bar\beta_{jk}|\le\mathrm C_L\). \hfill\(\square\)

\lemmaheading{Coupling of complete source histories}{lem:formation-history-coupling}
The source-assignment histories of a process and of an activated exact parent
can be coupled, together with their common initial marked sample, so that the
probability of any disagreement through \(t_*\) is at most
\(\mathrm C_L\epsilon^6\). This holds through the finite activation lineage,
for the full finite-state path with arbitrary label correlations and the
rate-one proposal clock. \status{ESTABLISHED}

\noindent\emph{Proof.}
First bound a proposal row on one proposal-free interval. For a microscopic
process let \(\xi=(w_1,h)\), and couple its auxiliary fast diffusion to the
same initial fast point, Brownian increments, assignment, and parameter path.
The omitted Gaussian environmental force
\(\vartheta[(H_{11}-\Omega_*)w_1+H_{1,>1}w_{>1}]\) has norm at most
\(\mathrm C_L\mathrm J_2|w|\), and strong contraction from the most recent
proposal or activation time \(t_{\mathrm{prev}}\) gives the first bound below.
For a parent at level \(\ell\ge1\), coupling its retained loop to the
auxiliary loop process, \cref{eq:rg-formation-parent-force-environment} gives
the second:
\begin{align}
 |\xi_t-\xi_t^*|
 &\le \mathrm C_L\mathrm J_2\int_{t_{\mathrm{prev}}}^t
             e^{-\mathrm c_L\mathrm J_1(t-\sigma)}|w_\sigma| d\sigma,
 \label{eq:rg-formation-fast-environment-coupling}\\
 |h_t-h_t^*|
 &\le \mathrm C_L\sqrt{\mathrm T\mathrm J_1}
       \int_{t_{\mathrm{prev}}}^t e^{-\mathrm c_L\mathrm J_1(t-\sigma)}\mathrm e_\epsilon(y_\sigma) d\sigma,
 \label{eq:rg-formation-loop-environment-coupling}
\end{align}
with the adapted environment of the process. The row Lipschitz constant and
the direct parent row error \cref{eq:rg-formation-parent-row-environment}
convert these into proposal-row errors.

At random proposal times, condition on the whole independent future Poisson
grid, but not on its marks or on any
future continuous path, so that the Brownian law is unchanged. Evenness and
the uniform lower curvature on \(M_\ell\) give the first bound below on every
sheet and level; proposals preserve \(y\), projections contract its norm, and
the preparation has second moment at most \(\mathrm C_L(\mathrm A^2+\mathrm T)\),
which gives the second:
\begin{align}
 (\mathrm T\Delta-\nabla V_{\ell,j}\cdot\nabla)|y|^2
 &\le2\mathrm T\dim M_0-2\mathrm c_L|y|^2,
 \label{eq:rg-formation-grid-moment-generator}\\
 \sup_{t\le t_*}\E\bigl[|\mathbf X_{M_\ell,t}|^2\bigm|\text{Poisson grid}\bigr]
 &\le \mathrm C_L(\mathrm A^2+\mathrm T).
 \label{eq:rg-formation-grid-moment}
\end{align}
If the initial sample is fixed as well, the right side first contains its
squared norm, and integrating that sample returns the bound; the future grid
after a deterministic activation is independent of the past, so the argument
applies at each activation and at grid-selected times. Averaging the row
errors of
\cref{eq:rg-formation-fast-environment-coupling,eq:rg-formation-loop-environment-coupling}
with \cref{eq:rg-formation-grid-moment}, where the exponential convolution has
integral at most \(\mathrm C_L/\mathrm J_1\), over the distribution of every
past prefix, and dividing by the probability of no rare prefix, gives the
common expected per-proposal bound
\begin{equation}
 \eta_\epsilon
 =\mathrm C_L\left[\frac{\mathrm J_2}{\mathrm J_1}
       +\frac{\mathrm J_2}{\sqrt{\mathrm J_1}}
                       \left(1+\frac{\mathrm A}{\sqrt{\mathrm T}}\right)\right]
 =O_L(\epsilon^7),
 \label{eq:rg-formation-averaged-environment-error}
\end{equation}
since \(\mathrm Q_1=2\mathrm Q_2+20\) and \(\mathrm A/\sqrt{\mathrm T}=\epsilon^{-3}\).

Define one comparison source process on the same proposal grid, started at
the initial assignment and using the row \(\bar\beta_{j\cdot}(\vartheta_\circ(t))\),
\(\vartheta_\circ(t)=2-e^{-\epsilon t}\), a single comparison law common to both
processes while each process keeps its own law and coefficient. At a proposal
time \(t\), condition on the common starting sample, the full future grid, and
all previous marks, self marks included; the assignment and level are fixed
on \((t_{\mathrm{prev}},t)\), and no future continuous trajectory is conditioned
on. Starting the auxiliary diffusion from the fast point at
\(t_{\mathrm{prev}}\), the environmental coupling, then
\cref{eq:rg-formation-fast-mixing-bound}, then the Lipschitz comparison of
\(\vartheta(t)\) with \(\vartheta_\circ(t)\) bound the averaged conditional row
error, for gaps at least \(t_{\mathrm{gap}}\), by the second line below:
\begin{align}
 t_{\mathrm{gap}}&=\frac{k_L}{\mathrm J_1}\log(1/\epsilon),
 \label{eq:rg-formation-history-gap}\\
 &\mathrm C_L\left[
 \bigl(1+\mathrm A\sqrt{\mathrm J_1/\mathrm T}\bigr)e^{-\mathrm c_{\mathrm{fast},L}\mathrm J_1t_{\mathrm{gap}}}
       +\epsilon/\mathrm J_1+\eta_\epsilon+r_\epsilon\right];
 \label{eq:rg-formation-long-gap-row-error}
\end{align}
the moment of the point-start mixing factor, which includes
\(\sqrt{\mathrm J_1/\mathrm T}(\mathrm A+\sqrt{\mathrm T})\), is absorbed because
\(\mathrm A/\sqrt{\mathrm T}\ge1\), and averaging over the past marks returns
the grid-conditioned moments exactly. Short gaps use the trivial bound one.
For a rate-one Poisson grid the expected number of ordered pairs of proposal
times less than \(t_{\mathrm{gap}}\) apart is at most \(t_*t_{\mathrm{gap}}\),
by integrating the two-time intensity; proposals within \(t_{\mathrm{gap}}\)
after the initial time or one of the finitely many activations add at most
\((L+1)t_{\mathrm{gap}}\), and the expected total count is \(t_*\).

To pass from conditional row errors to the whole history, condition on a grid
with \(\mathrm n\) proposals, write the two joint mark laws as products of
their conditional kernels \(\mathsf K_m\) and \(\overline{\mathsf K}_m\), and
replace kernels in order, keeping the earlier prefix distributed according to
the process and integrating the comparison tail, whose mass is one. With
\(\mathsf j_m\in\mathcal J\) the \(m\)th proposal mark,
\begin{equation}
 \bigl\|\operatorname{Law}(\mathsf j_1,\ldots,\mathsf j_{\mathrm n})
          -\operatorname{Law}(\overline{\mathsf j}_1,\ldots,\overline{\mathsf j}_{\mathrm n})\bigr\|_{\mathrm{TV}}
 \le \sum_{m=1}^{\mathrm n}
       \E_{\text{prefix}}
       \|\mathsf K_m(\cdot\mid\mathsf j_{<m})
                 -\overline{\mathsf K}_m(\cdot\mid\mathsf j_{<m})\|_{\mathrm{TV}},
 \label{eq:rg-formation-history-telescoping}
\end{equation}
and integrating over the grid with the preceding bounds gives
\begin{align}
 \mathrm{Err}_{\mathrm{hist}}\le \mathrm C_L\big[
 &(t_*+L+1)t_{\mathrm{gap}}
 +t_*\bigl(1+\mathrm A\sqrt{\mathrm J_1/\mathrm T}\bigr)
                      e^{-\mathrm c_{\mathrm{fast},L}\mathrm J_1t_{\mathrm{gap}}}
 \notag\\
 &+t_*\epsilon/\mathrm J_1
       +t_*\eta_\epsilon+t_*r_\epsilon\big].
 \label{eq:rg-formation-full-history-error}
\end{align}
With \(\mathrm c_{\mathrm{fast},L}k_L>\mathrm Q_1/2+12\), the five contributions
are bounded by constants times \(\epsilon^{19}\varsigma_\epsilon\),
\(\epsilon^8\), \(\epsilon^{\mathrm Q_1}\), \(\epsilon^6\), and
\(\epsilon^{2\mathrm Q_1+19}+\epsilon^{3\mathrm Q_1+25}\varsigma_\epsilon^2\), so
\(\mathrm{Err}_{\mathrm{hist}}\le\mathrm C_L\epsilon^6\).

Fix the common initial sample and the grid, couple each mark sequence
maximally with the comparison sequence, and glue the two couplings along it;
the glued sequences disagree with probability at most
\(\mathrm C_L\epsilon^6\). The path spaces, continuous paths in a Euclidean
space and right-continuous paths in \(\mathcal J\), are standard Borel, so each
continuous trajectory can be drawn from its regular conditional law given the
initial sample, grid, and marks, which preserves both full-process marginals
and the parent projection of the initial sample. Once self-proposals are
discarded, each source path is a function of grid and marks, which proves the
first bound of \cref{eq:rg-formation-two-path-bounds}; finitely many
activations change only \(\mathrm C_L\). \hfill\(\square\)

\section{Coherence, source movement, and continuous paths}
\label{sec:rg-formation-dynamics-transfer}

\lemmaheading{Successive coherence with the original conditional marks}{lem:formation-coherence}
Under the preparation and regime of
\cref{eq:rg-formation-preparation-margin,eq:rg-formation-preparation-transport,eq:rg-formation-weakest-one-regime},
there are a finite \(k_{\mathrm{geom},L}\) and a tolerance bound
\(\eta_L^{\max}>0\), a fixed fraction of \(\mathrm m_L\), both depending only on
\(L\) and the fixed data, with the following property. For every
\(k_L>k_{\mathrm{geom},L}\) and every \(0<\eta_L\le\eta_L^{\max}\) there are
\(\epsilon_L,\mathrm c_L,\mathrm C_L>0\) such that, for \(0<\epsilon<\epsilon_L\),
conclusion (b) of \Cref{thm:finite-depth-adaptive-formation} holds for the fine
process and its successive parent descriptions on the windows
\(I_\ell=[t_\ell,2t_\ell]\), \(t_\ell=k_L\log(1/\epsilon)/\mathrm J_\ell\), with
the probability bound \cref{eq:rg-formation-coherence-probability}. On
\(I_\ell\) the discrepancies \(\mathfrak d_i\) and the residuals of level at
most \(\ell\) are below \(\eta_L\mathrm A\) and future bridges stay separated
by a fixed positive multiple of \(\mathrm A\), in the conditional
root-mean-square original residuals of
\cref{eq:rg-formation-inherited-edge-statistic} for eliminated coordinates,
separately for the physical posterior and the restricted optimizer.
\status{ESTABLISHED}

\noindent\emph{Proof.}
On an active level \(\ell'<\ell\), write the non-root Gaussian comparison of
\Cref{lem:formation-prepared-continuation} through a current normal component
\(z\in\mathcal U_{\ell'}\cap\mathcal U_\ell^\perp\), a future component
\(y\in\mathcal U_\ell\), and the separate loop component. The weighted flag
gives, for the current and future blocks of \(H_{\ell'}\),
\begin{equation}
 H_{zz}\succeq \mathrm c_L\mathrm J_\ell I,
 \qquad \|H_{zy}\|\le \mathrm C_L\mathrm J_{\ell+1},
 \qquad \|H_{yy}\|\le \mathrm C_L\mathrm J_{\ell+1},
 \qquad \mathrm J_{L+1}=0.
 \label{eq:rg-formation-current-future-blocks}
\end{equation}
Indeed rows of levels at most \(\ell\) annihilate \(\mathcal U_\ell\), so the
original cross and future blocks contain only strengths at most
\(\mathrm J_{\ell+1}\); in an earlier elimination through \(\ell''\), the
eliminated-current block is \(O_L(\mathrm J_{\ell''+1})\), the inverse hidden
block \(O_L(\mathrm J_{\ell''}^{-1})\), and the eliminated-future block
\(O_L(\mathrm J_{\ell+1})\), so the Schur correction is still
\(O_L(\mathrm J_{\ell+1})\); at level zero there is no correction, and
loop/normal cross blocks vanish in the Gaussian core.

The deterministic reference mean \(\mu^{\mathrm G}\) has norm at most
\(\mathrm C_L\mathrm A\) and zero loop component. Integrating
\(\dot\mu_z^{\mathrm G}=-\vartheta(t)(H_{zz}\mu_z^{\mathrm G}+H_{zy}\mu_y^{\mathrm G})\)
with \(1\le\vartheta\le21/10\) gives the first bound below; integrating the
full future row, of norm \(O_L(\mathrm J_{\ell+1})\), from the common center,
earlier projections preserving \(P_\ell\mu^{\mathrm G}\), gives the second for
a sequential process begun at \(t^{\mathrm{act}}_{\ell-1}\):
\begin{align}
 |\mu_z^{\mathrm G}(t)|
 &\le \mathrm C_L\mathrm A\left[e^{-\mathrm c_L\mathrm J_\ell(t-t_0)}+\frac{\mathrm J_{\ell+1}}{\mathrm J_\ell}\right],
 \label{eq:rg-formation-current-mean}\\
 |\mu^{\mathrm G}(t)-\mathrm AP_\ell\mathbf v_0|
 &\le \mathrm C_L\mathrm A\left[
 e^{-\mathrm c_L\mathrm J_\ell(t-t^{\mathrm{act}}_{\ell-1})}
 +\frac{\mathrm J_{\ell+1}}{\mathrm J_\ell}+\mathrm J_{\ell+1}t\right].
 \label{eq:rg-formation-window-mean}
\end{align}
The future terms are absent when \(\ell=L\), and a fine process uses \(t_0=0\).

For the fluctuations, on each active non-root space
\(\mathbf Z=\mathbf X^{\mathrm G}-\mu^{\mathrm G}\) obeys
\(d\mathbf Z=-\vartheta(t)H_{\ell'}\mathbf Zdt+\sqrt{2\mathrm T} dW\), with
\(\vartheta H_{\ell'}\succeq\mathrm h_LI\) and initial covariance at most
\(\mathrm C_L\mathrm TH_{\ell'}^{-1}\). Choose \(\mathrm a_*>0\) below
\(\mathrm h_L\) and small enough for the initial Gaussian exponential moment.
For \(\mathcal Y(\mathbf z)=e^{\upsilon}\) with
\(\upsilon=\mathrm a_*|\mathbf z|^2/(4\mathrm T)\), direct differentiation gives
the first line below; the initial expectation of \(\mathcal Y\) is bounded,
orthogonal retention decreases \(\mathcal Y\), and It\^o's formula up to the
first exit from a ball, localization, and the nonnegative stopped value of
\(\mathcal Y\) yield the second for every \(\varrho>0\), higher stiffnesses only
strengthening the negative term:
\begin{align}
 \mathscr L_t\mathcal Y
 &=\left[\frac{\mathrm a_*\mathrm d}{2}
 +\frac{\mathrm a_*^2|\mathbf z|^2}{4\mathrm T}
 -\frac{\mathrm a_*\vartheta\mathbf z^\top H_{\ell'}\mathbf z}{2\mathrm T}\right]\mathcal Y
 \le \mathrm a_*(\mathrm d/2-\upsilon)e^\upsilon\le \mathrm C_L,
 \label{eq:rg-formation-gaussian-exponential-generator}\\
 \Pr\Bigl\{\sup_{t\le t_*}|\mathbf X^{\mathrm G}_t-\mu^{\mathrm G}(t)|&>\varrho\mathrm A\Bigr\}
 \le \mathrm C_L(1+t_*)
       \exp\left[-\mathrm c_L\varrho^2\mathrm A^2/\mathrm T\right].
 \label{eq:rg-formation-gaussian-maximum}
\end{align}

Let \(\bar r_L\) bound the norms of the finitely many residual and discrepancy
covectors, let \(\eta_L^{\max}\) be a fixed fraction of \(\mathrm m_L\), fix
\(\eta_L\in(0,\eta_L^{\max}]\), and choose a fluctuation radius
\(\varrho\mathrm A\) with \(\bar r_L\varrho\mathrm A\) smaller than the reserved
tolerance margin, \(\mathrm c_L\) in the probability depending on \(\varrho\).
Take for example \(k_{\mathrm{geom},L}=2/\mathrm c_L\) for the restoring constant
of \cref{eq:rg-formation-current-mean}; increasing this finite constant is
allowed. For small \(\epsilon\), \(t^{\mathrm{act}}_{\ell-1}<t_\ell/2\), so the
exponential term of \cref{eq:rg-formation-window-mean} tends to zero uniformly
on \(I_\ell\), and
\begin{equation}
 \frac{\mathrm J_{\ell+1}}{\mathrm J_\ell}\longrightarrow0,
 \qquad \mathrm J_{\ell+1}t_\ell\longrightarrow0\quad(\ell<L),
 \qquad \frac{t^{\mathrm{act}}_L}{t_*}\longrightarrow0.
 \label{eq:rg-formation-geometric-entrance}
\end{equation}
The nonlinear correction of \cref{eq:rg-formation-path-force-bound} is at most
\(\mathrm C_L\mathrm A\epsilon^{5/2}\), including the admitted initial transport
and finitely many inherited errors; it enters through the source force only via
the uniform bound \(B_*\sqrt{\mathrm T}\), so it holds for every adapted
assignment path. Since \(\mathrm A^2/\mathrm T=\epsilon^{-6}\), the strict
prepared residual margins and \cref{eq:rg-formation-gaussian-maximum} give the
asserted probability and separation on every window. The stable current-block
equation keeps bounding formed residuals as their future forcing decreases,
and after the last normal block there is no future forcing, so the final
regime persists through \(t_*\) on the same event, with every loop retained.

For an eliminated edge and a cumulative elimination through \(\ell'\ge1\),
write a fine non-root state as \(\mathbf z+y\) with
\(\mathbf z\in W_{\ell'}^\perp\) and \(y\in M_{\ell'}\). The optimizing and
physical Gaussian precisions are \(\vartheta H\) and
\(\vartheta H-\mathrm TR_{\mathrm{ev}}\). Every row of \(\mathfrak D\) lies in the
first normal layer, so \(R_{\mathrm{ev}}\) annihilates \(\mathcal U_1\),
\(\mathcal H\), and \(\mathcal K\), and with \(\mathrm TR_{\mathrm{ev}}\preceq H/2\)
both laws have hidden precision at least \(\mathrm c_L\mathrm J_{\ell'}I\),
hidden/future block \(O_L(\mathrm J_{\ell'+1})\), and zero hidden/loop block.
The Gaussian conditional mean therefore has norm at most
\(\mathrm C_L(\mathrm J_{\ell'+1}/\mathrm J_{\ell'})|y_{\mathcal U_{\ell'}}|\) and
covariance trace at most \(\mathrm C_L\mathrm T/\mathrm J_{\ell'}\), and, given
the full assignment, each exact hidden density is a normalized tilt of its
Gaussian conditional with log oscillation at most \(B_0\), hence density ratio
at most \(e^{B_0}\). For
\(\mu_\vartheta\in\{\boldsymbol\Pi_{\vartheta,\mathsf o},q^\star_\vartheta\}\) and
every retained state with \(|y_{\mathcal U_{\ell'}}|\le\mathrm C_L\mathrm A\), this
gives the first bound below, and for every original residual covector the
triangle inequality and its reverse in conditional \(L^2\) give the second:
\begin{align}
 \E_{\mu_\vartheta}
   [|\mathbf z|^2\mid P_{\ell'}\mathbf X=y,S=j]
 &\le \mathrm C_Le^{B_0}\left[
       \mathrm A^2\left(\frac{\mathrm J_{\ell'+1}}{\mathrm J_{\ell'}}\right)^2
       +\frac{\mathrm T}{\mathrm J_{\ell'}}\right],
 \label{eq:rg-formation-hidden-original-mark-moment}\\
 \left|
 \sqrt{\E_{\mu_\vartheta}[r_e(\mathbf X)^2\mid P_{\ell'}\mathbf X=y,S=j]}
 -|r_e(y)|\right|
 &\le\|\mathbf r_e\|
      \sqrt{\E_{\mu_\vartheta}[|\mathbf z|^2\mid P_{\ell'}\mathbf X=y,S=j]}.
 \label{eq:rg-formation-original-mark-rms}
\end{align}
On the event already constructed the right side of
\cref{eq:rg-formation-original-mark-rms} is \(o(\mathrm A)\), uniformly in the
assignment and separately for the two laws, and the same holds for
\(\mathfrak d_i\); this transfers the small formed and the separated future
residuals to \cref{eq:rg-formation-inherited-edge-statistic}, with no erased
mark replaced by its value at \(P_{\ell'}\mathbf x\). The bridge levels join
successive pairs of already formed components, giving the stated counts.
\hfill\(\square\)

\lemmaheading{Source envelope and changing hidden kernels}{lem:formation-source-envelope}
Every process or finite replacement lineage of
\Cref{lem:formation-prepared-continuation} satisfies
\cref{eq:rg-formation-source-movement} for the distinguished binary source,
from any of its activation times. Between \(\vartheta=1\) and its final
parameter, the first hidden physical conditional and the first hidden
optimizing conditional each change by relative entropy greater than
\(1/200\), uniformly at the same retained input and assignment.
\status{ESTABLISHED}

\noindent\emph{Proof.}
The distinguished receiver's posterior row \(\beta_{i_0,\pm}\), a factor of
the product row \cref{eq:rg-formation-full-source-row}, has the log odds of the
first line below. Since \(0\le\chi\le1\), it lies between
\(\operatorname{sig}(4(\vartheta-3/2)\mp\delta)\), and summing the inherited
row over the assignments with distinguished label \(+\) and taking the
expectation preserves these bounds, so the generator applied to the indicator
of that label gives the second line, with \(\lambda(t)=\Pr\{S_{i_0}(t)=+\}\):
\begin{align}
 \log\frac{\beta_{i_0,+}(\mathbf x,\vartheta)}
               {\beta_{i_0,-}(\mathbf x,\vartheta)}
 &=4(\vartheta-3/2)
  -\delta\left[
 \chi\bigl(\sqrt{\mathrm J_1/\mathrm T} r_{i_0,+}(\mathbf x)\bigr)
 -\chi\bigl(\sqrt{\mathrm J_1/\mathrm T} r_{i_0,-}(\mathbf x)\bigr)\right],
 \label{eq:rg-formation-source-log-odds}\\
 \lambda'&=\bar\beta_\lambda-\lambda,
 \qquad
 \operatorname{sig}(4(\vartheta-3/2)-\delta)
 \le\bar\beta_\lambda(t)\le
 \operatorname{sig}(4(\vartheta-3/2)+\delta).
 \label{eq:rg-formation-source-forcing-envelope}
\end{align}
Neither diffusion nor a marked projection changes this indicator, so variation
of constants remains valid across every activation, even when the inherited row
depends on the current assignment and continuous state.

By \Cref{lem:formation-parameter-comparison}, uniformly before
\(t^{\mathrm{act}}_L=2k_L\log(1/\epsilon)\),
\(\vartheta\le1+\mathrm C_L\epsilon\log(1/\epsilon)+\mathrm C_Lr_\epsilon\). A
solution of \cref{eq:rg-formation-source-forcing-envelope} cannot exceed the
larger of its starting value and the supremum forcing, so the preparation
tolerance and \(\|\operatorname{sig}'\|_\infty\le1/4\) give the first bound
below. For small \(\epsilon\),
every activation precedes \(t_*-10\) and \(\vartheta\ge7/4\) throughout
\([t_*-10,t_*]\), so nonnegativity of the earlier contribution in variation of
constants yields the second; for \(\delta<10^{-4}\), the elementary bounds
\(e^{1-\delta}>249/100\), \(e^2>7\), and \(e^{10}>2^{10}\) give the third:
\begin{align}
 &\lambda(t^{\mathrm{act}}_\ell)
 \le\operatorname{sig}(-2)+10^{-2}+\delta/4
       +\mathrm C_L[\epsilon\log(1/\epsilon)+r_\epsilon],
 \qquad 0\le \ell\le L,
 \label{eq:rg-formation-source-at-activation}\\
 &\lambda(t_*)\ge(1-e^{-10})\operatorname{sig}(1-\delta),
 \label{eq:rg-formation-source-final-mass}\\
 &(1-e^{-10})\operatorname{sig}(1-\delta)-\operatorname{sig}(-2)
       -10^{-2}-\delta/4
 >\frac{1023}{1024} \frac{249}{349}
       -\frac18-\frac1{100}-\frac1{40000}
 >\frac{57}{100}.
 \label{eq:rg-formation-source-strict-margin}
\end{align}
The remaining error in \cref{eq:rg-formation-source-at-activation} tends to
zero, hence \(\lambda(t_*)-\lambda(t^{\mathrm{act}}_\ell)>1/2\) for every
activation, which for a binary row is its total-variation change.

For the hidden kernels, fix the retained input and the assignment, and first
omit the bounded source factor. On the first hidden normal space let
\(\Omega_1,\Omega_{\mathrm f}\) be the initial and final Gaussian precisions, with
covariances \(\mathrm T\Omega_1^{-1},\mathrm T\Omega_{\mathrm f}^{-1}\). For the
optimizer the precision scales with \(\vartheta\); for the physical law, on the
non-root space,
\begin{equation}
 \Omega_1=H-\mathrm TR_{\mathrm{ev}}\succ0,
 \qquad
 \Omega_{\mathrm f}=\Omega_1+(\vartheta_{\mathrm f}-1)H
       \succeq\vartheta_{\mathrm f}\Omega_1,
 \qquad \vartheta_{\mathrm f}\ge7/4.
 \label{eq:rg-formation-physical-precision-growth}
\end{equation}
Restriction to a conditional block preserves the order, and so does an earlier
Gaussian marginal, whose Schur quadratic is the minimum over the discarded
coordinates, so in either case \(\Omega_{\mathrm f}\succeq(5/4)\Omega_1\). For
the Gaussian laws
\(\gamma_1,\gamma_{\mathrm f}\) on a hidden space of dimension
\(\mathrm d_{\mathrm h}\ge1\), whitening \(\Omega_1\) and diagonalizing the
conjugated final precision, with eigenvalues \(\varrho_i\ge5/4\), and completing
the square in the Hellinger affinity give the first line below; each factor
decreases for \(\varrho_i\ge1\), so the affinity is at most
\([2\sqrt{5/4}/(1+5/4)]^{\mathrm d_{\mathrm h}/2}\). Each exact conditional
\(\nu_1,\nu_{\mathrm f}\) is a normalized source tilt of its Gaussian
conditional bounded above by \(e^{B_0}\), also after exact hidden integration,
so \(\mathcal A(\nu_1,\nu_{\mathrm f})\le e^{B_0}\mathcal A(\gamma_1,\gamma_{\mathrm f})\),
and Jensen's inequality under \(\nu_1\) gives
\(-\KL(\nu_1\Vert\nu_{\mathrm f})/2\le\log\mathcal A(\nu_1,\nu_{\mathrm f})\).
With \(\log(1+1/80)\ge1/81\) this gives the second line, separately for the
physical and optimizing conditionals and uniformly in the retained input and
assignment:
\begin{align}
 \mathcal A(\gamma_1,\gamma_{\mathrm f})
 &=\prod_{i=1}^{\mathrm d_{\mathrm h}}
       \left(\frac{2\sqrt{\varrho_i}}{1+\varrho_i}\right)^{1/2}
       e^{-\mathrm M},
 \qquad \mathrm M\ge0,
 \label{eq:rg-formation-gaussian-affinity}\\
 \KL(\nu_1\Vert\nu_{\mathrm f})
 &\ge\frac{\mathrm d_{\mathrm h}}{2}\log\frac{81}{80}-2B_0
 \ge\frac1{162}-\frac1{5000}>\frac1{200}.
 \label{eq:rg-formation-hidden-kernel-change}
\end{align}
\hfill\(\square\)

\lemmaheading{Concentration of continuous paths on a common deterministic path}{lem:formation-continuous-path}
Let a source process and its exact parent start from a common retained sample
and parameter at \(t^{\mathrm{act}}_\ell\). There is a deterministic path
\(\bar{\mathbf x}_\ell:[t^{\mathrm{act}}_\ell,t_*]\to W_\ell\), namely the
parent's Gaussian reference mean on the non-root coordinates and the constant
equilibrium root mean, such that
\begin{equation}
 \Pr\Bigl\{\sup_{t^{\mathrm{act}}_\ell\le t\le t_*}|P_\ell\mathbf X_t-\bar{\mathbf x}_\ell(t)|
   >\mathrm C_L\mathrm A\epsilon^2\Bigr\}
 +\Pr\Bigl\{\sup_{t^{\mathrm{act}}_\ell\le t\le t_*}|\overline{\mathbf X}_t-\bar{\mathbf x}_\ell(t)|
   >\mathrm C_L\mathrm A\epsilon^2\Bigr\}
 \le \mathrm C_L(1+t_*)e^{-\mathrm c_L/\epsilon^2}.
 \label{eq:rg-formation-common-path-proof}
\end{equation}
Consequently the displacement bound of \cref{eq:rg-formation-two-path-bounds}
holds under every coupling of the two processes, in particular under the
coupling of \Cref{lem:formation-history-coupling}. A finite number of
successive replacements changes only the constants. \status{ESTABLISHED}

\noindent\emph{Proof.}
First compare Gaussian means at one deterministic parameter path. In the split
of \cref{eq:rg-formation-current-future-blocks} the inherited normal precision
is \(H_\ell=H_{yy}-H_{zy}^\top H_{zz}^{-1}H_{zy}\). With
\(\mathfrak w=\mu_z^{\mathrm G}+H_{zz}^{-1}H_{zy}\mu_y^{\mathrm G}\), the future
mean of the source process satisfies the first line below, the bound following
from \cref{eq:rg-formation-current-mean} and
\(\|H_{zz}^{-1}H_{zy}\|\le\mathrm C_L\mathrm J_{\ell+1}/\mathrm J_\ell\), and
integrating the contractive future propagator from a matching retained mean
gives the second:
\begin{align}
 \dot\mu_y^{\mathrm G}&=-\vartheta H_\ell\mu_y^{\mathrm G}-\vartheta H_{zy}^\top\mathfrak w,
 \qquad
 |\mathfrak w(t)|\le \mathrm C_L\mathrm A\left[
 e^{-\mathrm c_L\mathrm J_\ell(t-t^{\mathrm{act}}_\ell)}+\frac{\mathrm J_{\ell+1}}{\mathrm J_\ell}\right],
 \label{eq:rg-formation-gaussian-mean-defect}\\
 \sup_{t^{\mathrm{act}}_\ell\le t\le t_*}
 |\mu_y^{\mathrm G}(t)-\bar\mu^{\mathrm G}(t)|
 &\le \mathrm C_L\mathrm A\left[
 \frac{\mathrm J_{\ell+1}}{\mathrm J_\ell}
 +t_*\frac{\mathrm J_{\ell+1}^2}{\mathrm J_\ell}\right].
 \label{eq:rg-formation-gaussian-mean-reduction}
\end{align}
There is no future normal contribution when \(\ell=L\), and the loop Gaussian
core is unchanged by the reduction and has zero mean. If the independently
learned parameters differ, with
\(\mathrm e_\vartheta=\sup|\vartheta-\bar\vartheta|\le\mathrm C_Lr_\epsilon\) by
\Cref{lem:formation-parameter-comparison}, Young's inequality in the \(H_\ell\)
quadratic form gives, for homogeneous parent means \(\mu,\bar\mu\) with
difference \(\mathfrak e\),
\begin{equation}
 \frac d{dt}|\mathfrak e|^2
 \le-\vartheta \mathfrak e^\top H_\ell\mathfrak e
       +\mathrm C_L\mathrm e_\vartheta^2 \bar\vartheta
                         \bar\mu^\top H_\ell\bar\mu,
 \label{eq:rg-formation-mean-parameter-energy}
\end{equation}
and since \(\frac d{dt}|\bar\mu|^2=-2\bar\vartheta\bar\mu^\top H_\ell\bar\mu\),
the time integral of the final energy is at most \(\mathrm C_L\mathrm A^2\), so
the extra displacement is at most \(\mathrm C_L\mathrm A \mathrm e_\vartheta\).

Attach to each process its own Gaussian comparison from
\Cref{lem:formation-prepared-continuation}; the discrepancy between the
parent's projected starting reference and its full retained starting sample is
included in \cref{eq:rg-formation-path-force-bound}. On the non-root
coordinates let \(\bar{\mathbf x}_\ell=\bar\mu^{\mathrm G}\) be the parent's
reference mean. For the parent,
\cref{eq:rg-formation-path-force-bound,eq:rg-formation-gaussian-maximum} with
\(\varrho=\epsilon^2\) give
\(\sup_t|\overline{\mathbf X}_t-\bar{\mathbf x}_\ell(t)|\le\mathrm C_L\mathrm A(\epsilon^2+\epsilon^{5/2})\)
outside a set of probability \(\mathrm C_L(1+t_*)e^{-\mathrm c_L/\epsilon^2}\).
For the source process the same two estimates for its own reference, then
\cref{eq:rg-formation-gaussian-mean-reduction} and the parameter comparison,
give, outside a set of the same form,
\begin{equation}
 \sup_{t^{\mathrm{act}}_\ell\le t\le t_*}|P_\ell\mathbf X_t-\bar{\mathbf x}_\ell(t)|
 \le \mathrm C_L\mathrm A\left[
 \epsilon^2+\epsilon^{5/2}+r_\epsilon
 +\frac{\mathrm J_{\ell+1}}{\mathrm J_\ell}
 +t_*\frac{\mathrm J_{\ell+1}^2}{\mathrm J_\ell}\right].
 \label{eq:rg-formation-continuous-error-terms}
\end{equation}
For \(\ell<L\) the last term equals \(2\epsilon^{19}\),
\(\mathrm J_{\ell+1}/\mathrm J_\ell=\epsilon^{\mathrm Q_{\ell+1}+20}\), and
\(r_\epsilon=o(\epsilon^2)\); both future terms vanish at \(\ell=L\). The root
coordinates of each process form the same autonomous Ornstein--Uhlenbeck
process in law, with stiffness \(1+\kappa\ge1\), started at its equilibrium
and independent of the source history, and the argument of
\cref{eq:rg-formation-gaussian-maximum} keeps each within \(\epsilon^2\mathrm A\)
of the constant mean \(\kappa z_\star/(1+\kappa)\) outside a set of the same
form. Taking that constant as the root component of \(\bar{\mathbf x}_\ell\)
proves \cref{eq:rg-formation-common-path-proof}. Each of its events concerns
the law of one process alone, so for any coupling the displacement event of
\cref{eq:rg-formation-two-path-bounds}, after doubling \(\mathrm C_L\), lies in
their union; a finite sequence of replacements uses a finite triangle
inequality and union bound. \hfill\(\square\)

\section{The Gaussian core across an early interface}
\label{sec:rg-formation-core-proof}

\noindent\emph{Proof of parts (b) and (c) of \Cref{prop:rg-formation-gaussian-core}.}
Fix \(1\le\ell\le L-1\). In this section \(\mathrm C\) and \(\mathrm c\)
denote positive constants depending only on \(L\) and the carrier, never on
the strengths, \(\mathrm A\), \(\mathrm T\), \(t\), or \(\epsilon\); put
\(\mathrm J_{L+1}=0\) and
\(\bar\varrho_\ell=\varrho_{\ell-1}+\varrho_\ell+\varrho_{\ell+1}\).

\emph{Flag layers and block scales.} In the flag coordinates
\(w=(w_1,\ldots,w_L)\) of \cref{eq:rg-formation-weighted-flag},
\(\mathbf x_{\mathcal U}=\sum_s\mathsf Q_sw_s\), the quadratic form
\(\mathbf x_{\mathcal U}^\top H\mathbf x_{\mathcal U}=\sum_s\mathrm J_s|\sum_{r\le s}\mathsf B_{sr}w_r|^2\)
has the blocks
\((H_{\mathcal U})_{rr'}=\sum_{s\ge\max(r,r')}\mathrm J_s\mathsf B_{sr}^\top\mathsf B_{sr'}\),
each of norm at most \(\mathrm C\mathrm J_{\max(r,r')}\); the loop block
\(\mathrm J_1I_{\mathcal H}\) is orthogonal to every layer and plays no part.
On the active level \(\ell-1\) the normal part of \(H_{\ell-1}\) is the Schur
complement of \(H_{\mathcal U}\) onto the layers \(r\ge\ell\), obtained by
eliminating the layers \(\mathcal E=\{1,\ldots,\ell-1\}\), with the discarded
block \(z=w_\ell\) and the retained block \(y=(w_{\ell+1},\ldots,w_L)\). The
lower bound of \cref{eq:rg-formation-weighted-flag} on \(\mathcal E\) gives
\(H_{\mathcal E\mathcal E}\succeq\mathrm c\mathrm J_{\ell-1}I\), so the Schur
correction to a block \((r,r')\) with \(r,r'\ge\ell\) has norm at most
\(\mathrm C\mathrm J_r\mathrm J_{r'}/\mathrm J_{\ell-1}\le\mathrm C\mathrm J_{\max(r,r')}\varrho_{\ell-1}\).
Since in each block the summand with the largest admissible strength leads
and the remainder carries the next strength,
\begin{equation}
\begin{aligned}
H_{zz}&=\mathrm J_\ell\bigl[\mathsf B_{\ell\ell}^\top\mathsf B_{\ell\ell}+O(\varrho_\ell+\varrho_{\ell-1})\bigr],\\
H_{z,\ell+1}&=\mathrm J_{\ell+1}\bigl[\mathsf B_{\ell+1,\ell}^\top\mathsf B_{\ell+1,\ell+1}+O(\varrho_{\ell+1}+\varrho_{\ell-1})\bigr],\\
H_{\ell+1,\ell+1}&=\mathrm J_{\ell+1}\bigl[\mathsf B_{\ell+1,\ell+1}^\top\mathsf B_{\ell+1,\ell+1}+O(\varrho_{\ell+1}+\varrho_{\ell-1})\bigr],
\end{aligned}
\label{eq:rg-formation-core-block-scales}
\end{equation}
with the constants of the remainders depending only on \(L\) and the carrier,
together with \(\|H_{z,r}\|\le\mathrm C\mathrm J_r\) and
\(\|H_{\ell+1,r}\|\le\mathrm C\mathrm J_{\ell+1}(\varrho_{\ell+1}+\varrho_{\ell-1})\)
for \(r\ge\ell+2\), and \(\|H_{r,\cdot}\|\le\mathrm C\mathrm J_r\) for every
row block \(r\ge\ell+1\), in particular
\(\|H_{zy}\|\le\mathrm C\mathrm J_{\ell+1}\). Hence, for
\(\bar\varrho_\ell\le\mathrm c_0\) with \(\mathrm c_0\) small,
\(\lambda_{\min}(H_{zz})\ge\tfrac34\mathfrak j_\ell\mathrm J_\ell\),
\(\|H_{zz}^{-1}\|\le\mathrm C/\mathrm J_\ell\), and
\(\|H_{zz}^{-1}H_{zy}\|\le\mathrm C\varrho_\ell\). The same computation at
every earlier active level \(p\le\ell-1\) bounds the rows of layer
\(\ell+1\) of \(H_p\) by \(\mathrm C\mathrm J_{\ell+1}\).

\emph{Full row rank.} Write \(\mathbf r_e\) for the covector of \(r_e\).
Since \(\mathcal U_\ell=\mathcal U_{\ell-1}\cap\bigcap_e\ker\mathbf r_e\) over
the rows \(e\) of \(B_\ell\), the layer spanned by \(\mathsf Q_\ell\) is the
span of the projections \(P_{\mathcal U_{\ell-1}}\mathbf r_e\), and
\(\mathsf B_{\ell+1,\ell}\) has full row rank exactly when the matrix
\(\mathrm M_{fe}=\mathbf r_f\cdot P_{\mathcal U_{\ell-1}}\mathbf r_e\), with
\(f\) over the rows of \(B_{\ell+1}\) and \(e\) over those of \(B_\ell\), has
full row rank, because
\(\mathrm M_{fe}=(P_{\mathcal U_{\ell-1}}\mathbf r_f)\cdot(P_{\mathcal U_{\ell-1}}\mathbf r_e)\),
the component of \(P_{\mathcal U_{\ell-1}}\mathbf r_f\) in \(\mathcal U_\ell\)
pairs to zero with the layer, and the second factors span the layer. For
\(\ell=1\), \(\mathcal U_0=\mathcal U\) contains every row of \(B\), so
\(\mathrm M_{fe}=\mathbf r_f\cdot\mathbf r_e\); in the column of the
discrepancy row of agent \(i\) the entry of a bridge \(f\) is
\([i=\operatorname{head}(f)]-[i=\operatorname{tail}(f)]\), and the level-two
bridges, those of odd index, join disjoint pairs of triangles and have
pairwise disjoint endpoint sets, so the columns of the discrepancy rows of
their heads form an identity submatrix. For \(\ell\ge2\) let \(S\) be the
set of bridges of level at least \(\ell\), let \(\boldsymbol\partial_f\in\R^V\)
be the vector with entry \(1\) at the head of the bridge \(f\), \(-1\) at its
tail, and \(0\) elsewhere, let \(\partial_S\) have the rows
\(\boldsymbol\partial_f^\top\), \(f\in S\), and let \(\partial\) be the
oriented incidence operator of the whole graph, so that
\(\partial^\top\partial\) is its Laplacian. For a vector \(\mathsf b\) on the
bridges put \(\mathbf x_{\mathrm{br}}(\mathsf b)=(0,0,-\mathsf b)\), the
configuration with only the bridge links, and
\(\mathbf x(\mathsf b)=P_{\mathcal K^\perp}\mathbf x_{\mathrm{br}}(\mathsf b)\).
A configuration of \(\mathcal U_1\) has \(\bar k=\bar m=:t\) and, in each
triangle, three equal residuals which sum to \(-\mathfrak h_q=0\) by
orthogonality to \(\mathcal H\), so its internal links are \(\partial t\) and
its bridge links are \((\partial t)_f-\mathsf b_f\) with \(\mathsf b_f=r_f\);
it is therefore \((t,t,\partial t)+\mathbf x_{\mathrm{br}}(\mathsf b)\) and,
being orthogonal to \(\mathcal K\), equals \(\mathbf x(\mathsf b)\).
Conversely \(\mathbf x_{\mathrm{br}}(\mathsf b)\in W_1\cap\mathcal H^\perp\)
and \(\mathcal K\subseteq W_1\cap\mathcal H^\perp\), so
\(\mathbf x(\mathsf b)\in\mathcal U_1\), and
\(r_f(\mathbf x(\mathsf b))=r_f(\mathbf x_{\mathrm{br}}(\mathsf b))=\mathsf b_f\)
because residuals annihilate \(\mathcal K\). Hence
\(\mathcal U_{\ell-1}=\{\mathbf x(\mathsf b):\operatorname{supp}\mathsf b\subseteq S\}\)
with \(\mathsf b\mapsto\mathbf x(\mathsf b)\) injective, and from
\(\langle(t,t,\partial t),(t',t',\partial t')\rangle=t^\top(2I+\partial^\top\partial)t'\)
and
\(\langle(t,t,\partial t),\mathbf x_{\mathrm{br}}(\mathsf b)\rangle=-t^\top\partial_S^\top\mathsf b\)
the Gram matrix is
\(\langle\mathbf x(\mathsf b),\mathbf x(\mathsf b')\rangle=\mathsf b^\top\mathrm G\mathsf b'\)
with \(\mathrm G=I-\partial_S(2I+\partial^\top\partial)^{-1}\partial_S^\top\)
on \(\R^S\). The functional \(\mathsf b\mapsto\mathsf b_f\), \(f\in S\), has
the Riesz representative \(\mathbf x(\mathsf b^{(f)})\) with
\(\mathsf b^{(f)}\) the \(f\)th column of \(\mathrm G^{-1}\), so
\(P_{\mathcal U_{\ell-1}}\mathbf r_f=\mathbf x(\mathsf b^{(f)})\) and
\(\mathrm M_{fe}=(\mathrm G^{-1})_{fe}\). Since
\(2I+\partial^\top\partial-\partial_S^\top\partial_S=2I+\partial_{\setminus S}^\top\partial_{\setminus S}\),
with \(\partial_{\setminus S}\) the incidence operator of the graph with the
bridges of \(S\) removed, is positive definite, the Woodbury identity gives
\(\mathrm M_{fe}=\delta_{fe}+\boldsymbol\partial_f^\top(2I+\partial_{\setminus S}^\top\partial_{\setminus S})^{-1}\boldsymbol\partial_e\).
The graph with the bridges of \(S\) removed is the disjoint union of \(2^{L-\ell+1}\)
connected blocks of \(2^{\ell-2}\) consecutive triangles, and the inverse is
block diagonal with entrywise strictly positive blocks: on a block \(B\) with
Laplacian \(\partial_B^\top\partial_B\) and maximal degree \(\mathrm D_B\),
the matrix \((\mathrm D_B+2)I-(2I+\partial_B^\top\partial_B)\) is entrywise
nonnegative, irreducible, and of spectral radius at most
\(\mathrm D_B<\mathrm D_B+2\), since the Laplacian's eigenvalues lie in
\([0,2\mathrm D_B]\), so the Neumann series of
\((2I+\partial_B^\top\partial_B)^{-1}\) in it converges and is entrywise
positive. The bridges of \(S\) are those of index \(j=\mathrm w2^{\ell-2}\),
\(1\le\mathrm w\le2^{L-\ell+1}-1\), of level \(\ell+v_2(\mathrm w)\), the
bridge at position \(\mathrm w\) joining block \(\mathrm w\) to block
\(\mathrm w+1\); level \(\ell\) means \(\mathrm w\) odd and level
\(\ell+1\) means \(\mathrm w\equiv2\pmod4\). For \(f\ne e\) in \(S\), with
positions \(\mathrm w_f,\mathrm w_e\), the entry
\(\boldsymbol\partial_f^\top(2I+\partial_{\setminus S}^\top\partial_{\setminus S})^{-1}\boldsymbol\partial_e\)
is a sum over the blocks both bridges touch, so it vanishes unless
\(|\mathrm w_f-\mathrm w_e|=1\), and then it is plus or minus one entry of
the positive inverse on the shared block, hence nonzero. Pair each
level-\((\ell+1)\) bridge \(f\) with the level-\(\ell\) bridge \(e(f)\) at
position \(\mathrm w_f+1\), which exists because \(4\) divides
\(2^{L-\ell+1}\); then \(\mathrm M_{f,e(f')}\ne0\) exactly when
\(|\mathrm w_f-\mathrm w_{f'}-1|=1\), that is
\(\mathrm w_f-\mathrm w_{f'}\in\{0,2\}\), and as
\(\mathrm w_f-\mathrm w_{f'}\in4\mathbb Z\) exactly when \(f=f'\). The
square submatrix \((\mathrm M_{f,e(f')})_{f,f'}\) is diagonal with nonzero
entries, so the rows of \(\mathrm M\) of level \(\ell+1\) are independent and
\(\mathsf B_{\ell+1,\ell}\) has full row rank; as
\(\mathsf B_{\ell+1,\ell+1}\mathsf B_{\ell+1,\ell+1}^\top\) is invertible and
\(\mathsf B_{\ell+1,\ell}^\top\) is injective, \(\mathfrak Y_\ell\) is
injective. Nothing used the choice of bridge endpoints.

\emph{Slaving on the active interval.} Let level \(\ell-1\) be active on
\([t_0,t^{\mathrm{act}}_\ell]\) with \(t_0=t^{\mathrm{act}}_{\ell-1}\), with
mean \(\bar{\mathbf x}(t)\) in the layers \(r\ge\ell\),
\(\dot{\bar{\mathbf x}}=-H_{\ell-1}\bar{\mathbf x}\), and
\(|\bar{\mathbf x}(t_0)|\le\mathrm C\mathrm A\), which holds because
\(|\bar{\mathbf x}(0)|=\mathrm A|\mathbf v_0|\), the flow contracts
\(|\bar{\mathbf x}|\), and the projections do not increase it. Write
\(\mathbf g=H_{\ell-1}\bar{\mathbf x}\), so that \(\mathfrak v=\mathbf g_z\)
and \(\mathrm T\Delta_{\ell-1:\ell}=\tfrac12\mathbf g_z^\top H_{zz}^{-1}\mathbf g_z\)
by \cref{eq:rg-formation-gaussian-core-loss}, and
\(\mathsf b_{\ell+1}(t)=B_{\ell+1}P_\ell\bar{\mathbf x}(t)=\mathsf B_{\ell+1,\ell+1}\bar{\mathbf x}_{\ell+1}(t)\),
since \(P_\ell\) removes the layer \(\ell\) and \(\mathsf B_{\ell+1,r}=0\) for
\(r>\ell+1\). The row bound gives
\(|\dot{\bar{\mathbf x}}_{\ell+1}|=|\mathbf g_{\ell+1}|\le\mathrm C\mathrm J_{\ell+1}\mathrm A\),
and the same bound at the earlier levels, whose projections \(P_p\),
\(p\le\ell\), leave the layer \(\ell+1\) unchanged, gives the second line of
\cref{eq:rg-formation-core-two-sided}. Since
\(\dot{\mathbf g}=-H_{\ell-1}\mathbf g\), the variable
\(\mathbf u=\mathbf g_z+H_{zz}^{-1}H_{zy}\mathbf g_y\) satisfies
\(\dot{\mathbf u}=-H_{zz}\mathbf u-H_{zz}^{-1}H_{zy}(H_{\ell-1}\mathbf g)_y\),
where \(|(H_{\ell-1}\mathbf g)_y|\le\mathrm C\mathrm J_{\ell+1}(|\mathbf g_z|+|\mathbf g_y|)\),
\(|\mathbf g_y|\le\mathrm C\mathrm J_{\ell+1}\mathrm A\), and
\(|\mathbf g_z|\le|\mathbf u|+\mathrm C\varrho_\ell\mathrm J_{\ell+1}\mathrm A\).
Variation of constants with \(\|e^{-H_{zz}s}\|\le e^{-\lambda s}\),
\(\lambda=\lambda_{\min}(H_{zz})\), gives
\begin{equation}
|\mathbf u(t)|\le e^{-\lambda(t-t_0)}|\mathbf u(t_0)|
+\int_{t_0}^te^{-\lambda(t-s)}\bigl(\mathrm C\varrho_\ell^2\mathrm J_\ell|\mathbf u(s)|
+\mathrm C\varrho_\ell\mathrm J_{\ell+1}^2\mathrm A\bigr)ds,
\label{eq:rg-formation-core-volterra}
\end{equation}
and iterating the positive Volterra operator of the right side compares
\(|\mathbf u|\) with the solution of the scalar equation with decay rate
\(\lambda-\mathrm C\varrho_\ell^2\mathrm J_\ell\), constant forcing
\(\mathrm C\varrho_\ell\mathrm J_{\ell+1}^2\mathrm A\), and initial value
\(|\mathbf u(t_0)|\), so that
\(|\mathbf u(t)|\le e^{-(\lambda-\mathrm C\varrho_\ell^2\mathrm J_\ell)(t-t_0)}|\mathbf u(t_0)|
+\mathrm C\varrho_\ell\mathrm J_{\ell+1}^2\mathrm A/(\lambda-\mathrm C\varrho_\ell^2\mathrm J_\ell)\).
For \(\bar\varrho_\ell\le\mathrm c_0\),
\(\lambda-\mathrm C\varrho_\ell^2\mathrm J_\ell\ge\tfrac12\mathfrak j_\ell\mathrm J_\ell\),
and \(|\mathbf u(t_0)|\le\mathrm C\mathrm J_\ell\mathrm A\), so
\(|\mathbf u(t)|\le\mathrm C(\mathrm J_{\ell+1}^2\mathrm A/\mathrm J_\ell)[\varrho_\ell+\mathfrak x_\ell(t)]\).
For the forcing,
\(H_{zy}\mathbf g_y=H_{z,\ell+1}\mathbf g_{\ell+1}+\sum_{r\ge\ell+2}H_{z,r}\mathbf g_r\),
the sum being at most
\(\mathrm C\mathrm J_{\ell+2}^2\mathrm A\le\mathrm C\mathrm J_{\ell+1}^2\varrho_{\ell+1}^2\mathrm A\);
from \(\mathbf g_z=H_{zz}\bar{\mathbf x}_z+H_{zy}\bar{\mathbf x}_y\),
\(|\bar{\mathbf x}_z|\le\mathrm C\mathrm A(\varrho_\ell+\mathfrak x_\ell)\),
and \cref{eq:rg-formation-core-block-scales} then gives
\(\mathbf g_{\ell+1}=\mathrm J_{\ell+1}[\mathsf B_{\ell+1,\ell+1}^\top\mathsf b_{\ell+1}+O(\mathrm A(\bar\varrho_\ell+\mathfrak x_\ell))]\)
and
\(H_{zy}\mathbf g_y=\mathrm J_{\ell+1}^2[\mathfrak Y_\ell\mathsf b_{\ell+1}+O(\mathrm A(\bar\varrho_\ell+\mathfrak x_\ell))]\).
Writing \(\mathbf g_z=-H_{zz}^{-1}H_{zy}\mathbf g_y+\mathbf u\) with
\(H_{zz}^{-1}=\mathrm J_\ell^{-1}[(\mathsf B_{\ell\ell}^\top\mathsf B_{\ell\ell})^{-1}+O(\bar\varrho_\ell)]\)
and inserting in \(\tfrac12\mathbf g_z^\top H_{zz}^{-1}\mathbf g_z\), with
\(|\mathsf b_{\ell+1}|\le\mathrm C\mathrm A\), gives
\(\mathrm T\Delta_{\ell-1:\ell}=(\mathrm A^2\mathrm J_{\ell+1}^4/\mathrm J_\ell^3)
[\tfrac12|(\mathsf B_{\ell\ell}^\top\mathsf B_{\ell\ell})^{-3/2}\mathfrak Y_\ell\mathsf b_{\ell+1}/\mathrm A|^2
+O(\bar\varrho_\ell+\mathfrak x_\ell+\mathfrak x_\ell^2)]\), which is the
first line of \cref{eq:rg-formation-core-two-sided} by the definition of
\(\mathfrak y_\ell\).

\emph{Windows.} On \(I_\ell\),
\(\mathrm J_\ell(t-t_0)\ge\mathrm J_\ell(t_\ell-2t_{\ell-1})=k_L\log(1/\epsilon)(1-2\varrho_{\ell-1})\),
so \(\mathfrak x_\ell\to0\) uniformly on \(I_\ell\) under the stated
condition, and \(\mathrm J_{\ell+1}t\le2k_L\log(1/\epsilon)\varrho_\ell\to0\)
gives \(\mathsf b_{\ell+1}(t)/\mathrm A=B_{\ell+1}P_\ell\mathbf v_0+o(1)\), so
the quadratic \(\mathfrak y_\ell\) satisfies
\(\mathfrak y_\ell(\bar{\mathbf x}(t)/\mathrm A)=\mathfrak y_\ell(\mathbf v_0)+o(1)\).
For the lower bound,
\(|(\mathsf B_{\ell\ell}^\top\mathsf B_{\ell\ell})^{-3/2}\mathfrak Y_\ell\mathsf b_{\ell+1}|
\ge\|\mathsf B_{\ell\ell}\|^{-3}\sigma_{\min}(\mathfrak Y_\ell)|\mathsf b_{\ell+1}|\),
and \(|B_{\ell+1}P_\ell\mathbf v_0|^2=\sum_er_e(P_\ell\mathbf v_0)^2\) over
the \(N_\triangle/2^\ell\) bridges \(e\) of level \(\ell+1\) is at least
\(N_\triangle\mathrm m_L^2/2^{\ell+2}\) by the margin of
\Cref{lem:formation-preparation-margin}, which gives
\cref{eq:rg-formation-core-window}.

\emph{Power laws.} With \(\mathrm J_s=\epsilon^{-\mathrm Q_s}\),
\(\varrho_s=\epsilon^{\mathrm Q_s-\mathrm Q_{s+1}}\to0\),
\(k_L\log(1/\epsilon)\varrho_\ell\to0\), and on \(I_\ell\)
\[
\mathfrak x_\ell\le\epsilon^{\mathfrak j_\ell k_L(1-2\varrho_{\ell-1})/2-2(\mathrm Q_\ell-\mathrm Q_{\ell+1})},
\]
which tends to zero exactly under the stated threshold on \(k_L\); the
window statement then places \(\Delta_{\ell-1:\ell}\) between
\(\tfrac12\mathfrak y_\ell(\mathbf v_0)\) and \(2\mathfrak y_\ell(\mathbf v_0)\)
times \((\mathrm A^2/\mathrm T)\mathrm J_{\ell+1}^4/\mathrm J_\ell^3
=\epsilon^{3\mathrm Q_\ell-4\mathrm Q_{\ell+1}-\mathrm Q_{\mathrm T}}\) on
\(I_\ell\) for small \(\epsilon\). It is \(O(\epsilon^{\mathrm Q_\Delta})\)
on every early window exactly when
\(3\mathrm Q_\ell-4\mathrm Q_{\ell+1}-\mathrm Q_{\mathrm T}\ge\mathrm Q_\Delta\)
for every \(\ell\le L-1\), that is
\(\mathrm Q_\ell+\mathrm Q_{\mathrm T}+\mathrm Q_\Delta\ge\tfrac43(\mathrm Q_{\ell+1}+\mathrm Q_{\mathrm T}+\mathrm Q_\Delta)\),
and iterating from \(\mathrm Q_L=0\) gives
\(\mathrm Q_1+\mathrm Q_{\mathrm T}+\mathrm Q_\Delta\ge(4/3)^{L-1}(\mathrm Q_{\mathrm T}+\mathrm Q_\Delta)\),
with equality along the recursion; vanishing is the case
\(\mathrm Q_\Delta=0\) with strict inequality. \(\square\)

\section{The endogenous four-cycle}
\label{sec:rg-formation-cycle-proof}

\noindent\emph{Proof of \Cref{prop:rg-formation-endogenous-cycle}.}
\emph{Normalization, free energy, and descent.} Every factor of the joint is a
normalized kernel, so the joint has mass one, \(E=\mathbf 1\) has positive
probability, and every candidate source has positive prior and posterior mass
at every finite configuration. The reference density is
\((2\pi)^{-6}e^{-|\mathbf x|^2/2}\), the four likelihood factors contribute
\(4\mathrm T\log2\), and \cref{eq:rg-formation-evaluator-optimum} with
\(\sigma^2_{\mathrm{ev}}=1/2\) adds
\(\tfrac{\mathrm T}2\mathbf x^\top R_{\mathrm{ev}}\mathbf x\), so the scaled
negative ELBO is
\(\mathrm T\E_q\log(q/\pi_S)+\E_qU_S+\mathrm B_{\mathrm T}
=\mathrm T\KL(q\Vert q^\star)+\mathrm B_{\mathrm T}-\mathrm T\log Z_{\mathrm{opt}}\),
and the evidence is \(e^{-\mathrm B_{\mathrm T}/\mathrm T}Z_{\mathrm{phys}}\),
the physical energy omitting \(R_{\mathrm{ev}}\). Every \(U_S\) is a positive
definite quadratic form with Hessian at least \(\mathrm TI\), whose gradient
contains both the sending and the receiving appearances of each agent.
Between proposals the process is a linear diffusion with one of sixteen
matrices, and proposals are finitely many on bounded intervals, so it exists,
is pathwise unique with common random inputs, and does not explode. Since
\(q^\star(\mathbf x,S)\beta_{S'}(\mathbf x)
=q^\star_{\mathbf X}(\mathbf x)\beta_S\beta_{S'}\) is symmetric in \(S,S'\), the
dissipation identity \cref{eq:rg-formation-dissipation} at level zero, without
a parameter term, shows that \(\Fenergy_{\mathrm T}\) does not increase. A cycle
automorphism, with link signs transformed under reflections, preserves the
reference, the kernels, the rows, the kinetics, and the preparation class, so
the law at every time is invariant.

\emph{Matching sector.} Fix \(S^{(0)}\), with selected edges \(e=e_0\) and
\(f=e_2\); \(S^{(1)}\) follows by rotation. For each edge the covectors
\(\mathbf a^b_g,\mathbf a^m_g\) of \(r^b_g,r^m_g\) have Gram matrix
\(\begin{pmatrix}3&1\\1&3\end{pmatrix}\), and the two selected edges have
disjoint supports. The interaction Hessian of \(V_{S^{(0)}}\) is
\(H_M=2\sum_{g\in\{e,f\}}(\mathbf a^b_g\mathbf a^{b\top}_g+\mathbf a^m_g\mathbf a^{m\top}_g)\),
and the unit covectors
\begin{equation}
 n^+_g=\frac{\mathbf a^b_g+\mathbf a^m_g}{2\sqrt2},\qquad
 n^-_g=\frac{\mathbf a^b_g-\mathbf a^m_g}{2},\qquad g\in\{e,f\},
 \label{eq:rg-formation-cycle-normals}
\end{equation}
are orthonormal eigenvectors with eigenvalues \(8\) and \(4\); the kernel of
\(H_M\) has dimension eight. With \(\sigma^2_{\mathrm{ev}}=1/2\),
\(R_{\mathrm{ev}}n^+_g=0\), \(R_{\mathrm{ev}}n^-_g=4n^-_g\), and
\(R_{\mathrm{ev}}\) leaves the kernel invariant, so the held-matching drift
matrix \(K_M=H_M+\mathrm T(I+R_{\mathrm{ev}})\) has normal rates
\(8+\mathrm T\) (twice) and \(4+5\mathrm T\) (twice), all rates are at least
\(\mathrm T\), and the physical matrix has normal rates \(8+\mathrm T\) and
\(4+\mathrm T\). With \(\mathbf p=\sum_{g\in\{e,f\}}(\mathbf a^b_g+\mathbf a^m_g)\),
the vector \(\mathbf x^\circ+\mathbf p/4\) lies in the kernel of \(H_M\),
\(H_M\mathbf p=8\mathbf p\), and \(R_{\mathrm{ev}}\) annihilates \(\mathbf x^\circ\)
and \(\mathbf p\), so the deterministic held path from \(\mathbf x^\circ\) is
\begin{equation}
 \bar{\mathbf x}(t)=e^{-\mathrm Tt}\Bigl[\mathbf x^\circ+\frac{1-e^{-8t}}4\mathbf p\Bigr],
 \qquad
 r_{\mathrm{int}}(\bar{\mathbf x}(t))=-e^{-(8+\mathrm T)t},
 \qquad
 r_{\mathrm{ext}}(\bar{\mathbf x}(t))=e^{-\mathrm Tt}\Bigl[-\frac32+\frac{e^{-8t}}2\Bigr],
 \label{eq:rg-formation-cycle-center}
\end{equation}
in both channels, where \(r_{\mathrm{int}}\) is a selected and
\(r_{\mathrm{ext}}\) an exterior residual; along it \(\bar k=\bar m\), so the
evaluator force vanishes.

\emph{Episode.} Let \(\mathscr B\) be the event that
\(|\mathbf X_0-\mathbf x^\circ|<\varrho_\circ\), \(S_0=S^{(0)}\), and no proposal
occurs on \([0,1/4]\). The centered Gaussian component has expected squared
norm \(12\sigma_\circ^2=\varrho_\circ^2/4\), so it puts mass at least \(3/4\) on
the ball, and with the mixture weight \(1/2\), the bound
\(1+\mathfrak z\ge1/2\), the assignment mass \(1/16\), and the independent
clock, \(\Pr(\mathscr B)\ge\tfrac3{256}e^{-1/4}\). On \([0,1/4]\) the process is
the held Ornstein--Uhlenbeck process
\(\widetilde{\mathbf X}_t=e^{-K_Mt}\mathbf X_0+\mathbf Z_t\),
\(\mathbf Z_t=\sqrt{2\mathrm T}\int_0^te^{-K_M(t-\sigma)}dW_\sigma\), and afterwards
the two coincide until the first proposal returning another assignment. In an
eigenbasis of \(K_M\), integration by parts gives
\(\sup_{t\le t^\circ_*}|Z^k_t|\le2\sqrt{2\mathrm T}\sup_{t\le t^\circ_*}|W^k_t|\),
and \(\E\sup_{t\le t^\circ_*}|W^k_t|^2\le4t^\circ_*\) by Doob's \(L^2\)
maximal inequality: with \(W^*=\sup|W^k|\), optional stopping at the hitting
time of a level \(\mathrm y>0\) gives
\(\mathrm y\Pr(W^*\ge\mathrm y)\le\E[|W^k_{t^\circ_*}|\mathbf 1_{W^*\ge\mathrm y}]\),
and integrating against \(2 d\mathrm y\) with truncation gives the bound.
Summing twelve coordinates, \(\E\sup|\mathbf Z|^2\le384\mathrm Tt^\circ_*\), so
with \(\delta_\epsilon=\epsilon^2/(8\sqrt3)\), Chebyshev's inequality gives
\(\Pr(\sup_{t\le t^\circ_*}|\mathbf Z_t|>\delta_\epsilon)\le73728 \epsilon^4t^\circ_*\).
On \(\mathscr B\) and the complementary noise event, contraction gives
\(|\widetilde{\mathbf X}_t-\bar{\mathbf x}(t)|\le\varrho_\circ+\delta_\epsilon\).
Every residual covector has norm \(\sqrt3\), so on \(\mathscr B\) each
residual lies within \(\sqrt3\varrho_\circ\) of \(-1\) at time zero and each
internal belief relative entropy exceeds \((1-\sqrt3/100)^2/2>1/4\), which is
the entrance clause of (a); the restoring rates of (a) are those of \(K_M\)
computed above. Since \(\mathrm Tt^\circ_*<1/100\),
for \(1/4\le t\le t^\circ_*\) the selected residuals have magnitude below
\(e^{-2}+\sqrt3(\varrho_\circ+\delta_\epsilon)<1/5\) and the exterior ones above
\(e^{-\mathrm Tt^\circ_*}(3/2-e^{-2}/2)-\sqrt3(\varrho_\circ+\delta_\epsilon)>13/10\).
Each receiver's exterior cost therefore exceeds its partner cost by more than
one, its posterior row gives its exterior source probability at most
\(e^{-1/\mathrm T}\), and a proposal leaves \(S^{(0)}\) with probability at most
\(4e^{-1/\mathrm T}\); for the rate-one clock the expected number of departing
proposals up to \(t^\circ_*\) is at most \(4t^\circ_*e^{-1/\mathrm T}\), with
self-proposals unrestricted. For \(t\ge t_\epsilon\), every residual
annihilates the kernel of \(H_M\), so a selected residual is at most
\(e^{-8t}+\sqrt3\varrho_\circ e^{-4t}+\sqrt3\delta_\epsilon\le
\epsilon^{16}+\sqrt3\varrho_\circ\epsilon^8+\epsilon^2/8\le\epsilon^2\), the belief
relative entropy \((r^b)^2/2\) and the model relative entropy
\((r^m)^2/(2\sigma^2_{\mathrm{ev}})\) are at most \(\epsilon^4\), and every
exterior belief relative entropy exceeds \((13/10)^2/2>1/2\), the model ones
being larger. Let \(\mathscr A_t\) be the prefix event that the noise stays
within \(\delta_\epsilon\) up to \(t\) and every proposal in \([1/4,t]\) returns
\(S^{(0)}\); then \(\Pr(\mathscr A_{t^\circ_*}^c\mid\mathscr B)\le p_\epsilon\).
Since \(\epsilon^4t^\circ_*\) increases on \((0,1/100]\), the first term of
\(p_\epsilon\) is at most \(73728\cdot10^{-8}(2\log100+1)<0.008\), and
\(e^{-1/\mathrm T}<\mathrm T\) makes the second below \(10^{-14}\), so
\(p_\epsilon<1/2\) and
\(\Pr(\mathscr B\cap\mathscr A_{t^\circ_*})\ge\tfrac3{256}e^{-1/4}(1-p_\epsilon)\ge p_0\).
The stopping rule is prefix-measurable, and the same argument holds for
\(S^{(1)}\) and at the reflected center, which proves the probability bound
of (a).

\emph{Stationary comparison.} On each sheet the diffusion with drift
\(-\nabla U_S\) and covariance \(2\mathrm TI\) is reversible for
\(e^{-U_S/\mathrm T}d\mathbf x\), and the symmetric measure
\(q^\star(\mathbf x,S)\beta_{S'}(\mathbf x)\) makes the proposal step
reversible for \(q^\star\), so \(q^\star\) is invariant. Under \(q^\star\),
given \(S_0=M\), the configuration is \(\mathcal N(0,\mathrm TK_M^{-1})\), and
for a matched edge \(e\) with belief covector \(\mathbf a^b_e\),
\(K_M\succeq H_M\succeq2\mathbf a^b_e\mathbf a^{b\top}_e\), so
\(r^b_e=\mathbf a^b_e\cdot\mathbf x\) has variance
\(\mathrm T\mathbf a^{b\top}_eK_M^{-1}\mathbf a^b_e\le\mathrm T/2\) and
\(|r^b_e|>1/\sqrt2\) has probability at most \(2e^{-1/(2\mathrm T)}\). The
event of (a) requires \(S_0\) to be a matching whose matched pairs carry
internal belief relative entropy above \(1/4\) at time zero, that is
\(|r^b_e|>1/\sqrt2\) on a matched edge, so under the process started from
\(q^\star\) it has probability at most \(2\cdot2e^{-1/(2\mathrm T)}\le4e^{-1/(4\mathrm T)}\),
and \(4e^{-1/(4\mathrm T)}<p_0/2\) because \(\mathrm T\le10^{-16}\). This
completes (a).

\emph{Parent channel.} The covector of \(\sum_i\omega_i\) pairs to zero with
\(n^-_g\) and to \(-1/\sqrt2\) with \(n^+_g\), so it annihilates \(Z_M\) and the
loop is a function of \(P_{W_M}\mathbf x\); \(W_M\), of dimension nine, contains
the kernel of \(H_M\) and the fast common-normal direction
\((n^+_e+n^+_f)/\sqrt2\). Both \(H_M\) and \(R_{\mathrm{ev}}\) preserve \(W_M\)
and \(Z_M\), so on a matching fiber the hidden optimizing conditional is
\(\nu_{\mathrm T}=\mathcal N(0,\mathrm TK_Z^{-1})\) with
\(K_Z=\operatorname{diag}(8+\mathrm T,4+5\mathrm T,4+5\mathrm T)\), independent of
the retained state, and the physical one has precision
\(\operatorname{diag}(8+\mathrm T,4+\mathrm T,4+\mathrm T)\). The channel \(C\)
is one deterministic measurable kernel onto the disjoint union
\(\coprod_S\{S\}\times W_S\), with Lebesgue measure on each \(W_S\) and counting
measure on \(S\); its pushforwards of the joint, the physical posterior, the
optimizer, and the recognition law are normalized, it keeps every assignment,
and the chain rule of \Cref{thm:rg-exact-coarse-vfe} splits
\(\KL(q\Vert q^\star)\) into the retained term and the mean discarded
conditional term. This proves (b).

\emph{Discarded conditional free energy.} Under the held process the hidden
coordinates are \(z_t=e^{-K_Zt}z_0+\xi_t\) with
\(\xi_t\sim\mathcal N(0,\mathrm TK_Z^{-1}(I-e^{-2K_Zt}))\) independent of the
retained coordinates and of the initial data, and \(|z_0|\le\varrho_\circ\) on
\(\mathscr B\) because \(\mathbf x^\circ\) has zero hidden component. Given the
retained state, the hidden law is a mixture over \(z_0\) of these Gaussians,
and for one coordinate with rate \(\lambda_k\) and \(a=e^{-2\lambda_kt}\) the
relative entropy to the stationary law is
\(\tfrac12[-a-\log(1-a)+\lambda_ka z_0^2/\mathrm T]\). Convexity of relative
entropy in its first argument, \(4\le\lambda_k\le9\), and
\(-a-\log(1-a)\le a^2/(2(1-a))\) bound the mean conditional relative entropy of
the held law conditioned on \(\mathscr B\) by
\(\frac{9\varrho_\circ^2}{2\mathrm T}e^{-8t}+\frac{3}{4(1-e^{-8})}e^{-16t}\le\eta_\epsilon\)
for \(t\ge t_\epsilon\). To pass to the episode, let \(\widehat Q\) be this held
law, \(\mathscr R=\widehat Q_Y\otimes\nu_{\mathrm T}\), and \(Q_{\mathscr A}\),
\(Q_{\mathscr A^c}\) its conditionals on \(\mathscr A_t\) and its complement, of
probabilities \(1-p\) and \(p\le p_\epsilon\). The relative entropy of the joint
law of \((\mathbf 1_{\mathscr A_t},Y,Z)\) to
\(\mathrm{Bernoulli}(1-p)\otimes\mathscr R\), computed in both orders, gives
\((1-p)\KL(Q_{\mathscr A}\Vert\mathscr R)+p\KL(Q_{\mathscr A^c}\Vert\mathscr R)
=\KL(\widehat Q\Vert\mathscr R)+\mathbb I(\mathbf 1_{\mathscr A_t};(Y,Z))\), and the
mutual information is at most \(\mathcal S_2(p)\). Dropping the second term on
the left and the nonnegative retained-marginal term of
\(\KL(Q_{\mathscr A}\Vert\mathscr R)\), and using that \(\mathcal S_2\) increases
on \([0,1/2]\), bounds the mean discarded optimizing conditional relative
entropy of \(Q_{\mathscr A}\) by \(\bar\eta_\epsilon\); on
\(\mathscr B\cap\mathscr A_t\) the held and native processes coincide, so
\(Q_{\mathscr A}\) is the episode law. The physical parent conditions the
discarded coordinates also on the evaluator samples, so compare with the
normalized law
\(\Pi(dz\mid y) P_{\mathrm{ev}}(dm,dk\mid\mathbf x)
\exp[(\mathfrak d_y\cdot(k-m)-|\mathfrak d_y|^2/2)/(1-\sigma^2_{\mathrm{ev}})]\),
where \(\mathfrak d_y=\mathfrak Dy\) and \(\mathfrak d_z=\mathfrak Dz\) are
orthogonal; the multiplier has expectation one and depends only on retained
variables, so this law has the physical hidden conditional given
\((y,m,k)\). The recognition evaluator differs from it by the mean shift
\(\mathfrak d_z\), at relative entropy
\(|\mathfrak d_z|^2/(2(1-\sigma^2_{\mathrm{ev}}))\), and the chain rule in both
orders bounds the mean discarded physical conditional relative entropy by
\(\KL(q(dz\mid y)\Vert\Pi(dz\mid y))+\E_q|\mathfrak d_z|^2/(2(1-\sigma^2_{\mathrm{ev}}))
=\KL(q(dz\mid y)\Vert\nu_{\mathrm T})+\log\frac{4+5\mathrm T}{4+\mathrm T}\), where
\(\log\frac{4+5\mathrm T}{4+\mathrm T}\le\frac{4\mathrm T}{4+\mathrm T}\le\mathrm T\).
This proves (c).

\emph{Parent dynamics.} The inherited proposal kernel
\[
\beta_C((S,y),\cdot)=\int q^\star(d\mathbf x\mid C=(S,y))\sum_{S'}\beta_{S'}(\mathbf x)
\delta_{(S',P_{W_{S'}}\mathbf x)}
\]
is a probability kernel with the rate-one clock, reversible with respect to
\(C_\#q^\star\) because the symmetric measure
\(q^\star(d\mathbf x,S)\beta_{S'}(\mathbf x)\) is pushed through \((C,C)\); on a
matching stratum the retained diffusion is the exact projection of the held
process, and elsewhere \(C\) is the identity. On a matching fiber, with
\(z_+=((n^+_e+n^+_f)/\sqrt2)\cdot y\) and the hidden coordinates
\(z_1,z_2,z_3\) of \(Z_M\), the selected costs are
\(\mathcal D_e=(z_++z_1)^2+z_2^2\) and \(\mathcal D_f=(z_+-z_1)^2+z_3^2\), and each
exterior common residual \((r^b_g+r^m_g)/2\) annihilates \(Z_M\) and satisfies
\(\mathcal D_g\ge((r^b_g+r^m_g)/2)^2\). If both exterior common residuals have
magnitude at least \(\mathrm r_\circ\), each receiver chooses its exterior
source with probability at most
\(\exp[(\mathcal D_{\mathrm{in}}-\mathcal D_{\mathrm{out}})/\mathrm T]\), and the
Gaussian identity
\(\E e^{(a+Z)^2/\mathrm T}=(1-2\sigma^2/\mathrm T)^{-1/2}\exp[a^2/(\mathrm T(1-2\sigma^2/\mathrm T))]\)
with \(\sigma^2=\mathrm T/(8+\mathrm T)\) and \(\mathrm T/(4+5\mathrm T)\), whose
factors are at most \(\sqrt{4/3}\) and \(\sqrt2\), gives
\(\beta_C((S^{(0)},y),\{S'\ne S^{(0)}\})\le4\sqrt{8/3}\exp[-(\mathrm r_\circ^2-4z_+^2/3)/\mathrm T]\).
On the episode \(\mathrm r_\circ=13/10\) and \(|z_+|\le\epsilon^2\) on
\([t_\epsilon,t^\circ_*]\), so this is at most \(4\sqrt{8/3} e^{-1/\mathrm T}\).
Condition on \(\mathscr B\cap\mathscr A_{t_\epsilon}\) and let
\(\mu_{\mathrm{act}}\) be the law of \(C(\mathbf X_{t_\epsilon},S_{t_\epsilon})\)
under this prefix. The Brownian increments after \(t_\epsilon\) are independent
of the prefix, and the native future equals the held process until a
departure or a tube exit, whose conditional probability is at most
\(\Pr(\mathscr A^c_{t^\circ_*}\mid\mathscr B)/\Pr(\mathscr A_{t_\epsilon}\mid\mathscr B)
\le p_\epsilon/(1-p_\epsilon)\). The inherited parent started from
\(\mu_{\mathrm{act}}\) and driven by the same increments has the retained drift
and diffusion of the projected held process until its first departure; the
retained tube fails with conditional probability at most
\(p_\epsilon/(1-p_\epsilon)\), and on the tube the parent departs during the
unit interval with probability at most \(4\sqrt{8/3} e^{-1/\mathrm T}\). The
union bound and the triangle inequality give (d). \(\square\)

\noindent\emph{Proof of \Cref{prop:phil-seeded-contrast-formation}.}
\emph{Rows.} For a continuous row path \(\pi\) whose entries lie in
\([\mathrm c_0,1-\mathrm c_0]\) for some \(\mathrm c_0\in(0,1/2]\), construct the
kinetics pathwise as in the proof above: the configuration drift
\(-\nabla U_S\) does not involve the prior, and at a clock time \(\tau\)
receiver \(i\) takes source \(i+1\) exactly when a uniform variable on
\([0,1]\), independent and fresh for each clock time and receiver, is at most
its posterior row
\[
\beta_{i,i+1}(\mathbf x;\pi(\tau))
=\frac{\pi_{i,i+1}(\tau)e^{-\mathcal D_{e(i,i+1)}(\mathbf x)/\mathrm T}}
{\sum_{j=i\pm1}\pi_{ij}(\tau)e^{-\mathcal D_{e(i,j)}(\mathbf x)/\mathrm T}}
\]
at \(\mathbf x=\mathbf X_{\tau-}\). Write \(q^\pi_t\) for the law at time
\(t\). Drive the constructions for two such row paths \(\pi,\pi'\) by the same
initial data, Brownian motion, clock, and uniform variables. They coincide
until the first clock time at which the uniform variable of some receiver falls
between its posterior rows under \(\pi\) and \(\pi'\), and, given agreement
before a clock time
\(\tau\), this happens at \(\tau\) with conditional probability at most
\(\sum_i|\beta_{i,i+1}(\mathbf x;\pi(\tau))-\beta_{i,i+1}(\mathbf x;\pi'(\tau))|\)
at the common configuration. With \(\mathrm p=\pi_{i,i+1}\) and
\(\pi_{i,i-1}=1-\mathrm p\), the derivative of a posterior row in \(\mathrm p\)
is \(\beta_{i,i+1}(1-\beta_{i,i+1})/(\mathrm p(1-\mathrm p))
\le1/(4\mathrm c_0(1-\mathrm c_0))\le1/(2\mathrm c_0)\), and the clock has rate
one, so
\[
\bigl|q^\pi_t(S_i=j)-q^{\pi'}_t(S_i=j)\bigr|
\le\frac2{\mathrm c_0}\int_0^t\|\pi(s)-\pi'(s)\|_\infty ds,
\]
where \(\|\cdot\|_\infty\) is the maximum over receivers and sources. A
solution of \cref{eq:phil-source-learning} that starts uniform satisfies
\[
\pi_{ij}(t)-\tfrac12=\epsilon\int_0^te^{-\epsilon(t-s)}
\bigl(q_s(S_i=j)-\tfrac12\bigr)ds,
\qquad
\bigl|\pi_{ij}(t)-\tfrac12\bigr|\le\tfrac12\bigl(1-e^{-\epsilon t}\bigr)
\le\tfrac{\epsilon t}2,
\]
because \(0\le q_s(S_i=j)\le1\); each row sums to one, and on \([0,t_1]\) its
entries lie in \([\mathrm c_0,1-\mathrm c_0]\) with
\(\mathrm c_0=e^{-\epsilon t_1}/2\). Fix \(t_1>0\) and let \(\Gamma\) send a row
path \(\pi\) to the right-hand side of the first identity with \(q_s=q^\pi_s\).
It maps the closed set of continuous row paths on \([0,t_1]\) with
\(|\pi_{ij}(t)-1/2|\le(1-e^{-\epsilon t})/2\) into itself, and
\[
\|\Gamma\pi(t)-\Gamma\pi'(t)\|_\infty
\le\frac{2\epsilon t_1}{\mathrm c_0}\int_0^t\|\pi(s)-\pi'(s)\|_\infty ds,
\qquad t\le t_1,
\]
so the \(n\)th iterate of \(\Gamma\) contracts the supremum distance by the
factor \((2\epsilon t_1^2/\mathrm c_0)^n/n!\), and \(\Gamma\) has exactly one
fixed point in that set. Every solution lies in the set, solutions on nested
intervals agree, and \(s\mapsto q^\pi_s(S_i=j)\) is continuous because the clock
has rate one, so the fixed point is differentiable and solves
\cref{eq:phil-source-learning}. The rows are therefore a deterministic
continuous function of time, and the row bound holds. This proves (a).

\emph{Seed events.} The function
\(\epsilon t^\circ_*=\epsilon(2\log(1/\epsilon)+1)\) increases on
\((0,1/100]\), with value \((1+4\log10)/100<0.1022\) at \(1/100\), so
\(t^\circ_*<2/\epsilon\) and, by (a), every row entry lies in
\([0.4489,0.5511]\) on \([0,t^\circ_*]\), whence
\(\pi_{ik}(t)/\pi_{ij}(t)\le0.5511/0.4489<1.23\) for the two entries of every
row. For a sign
\(\sigma\in\{1,-1\}\) and a perfect matching \(M\in\{S^{(0)},S^{(1)}\}\), let
\(\mathscr B_{\sigma,M}\) be the event that
\(|\mathbf X_0-\sigma\mathbf x^\circ|<\varrho_\circ\), \(S_0=M\), and no proposal
occurs on \([0,1/4]\). The initial configuration, the initial assignment, and
the clock are independent, so
\[
\Pr(\mathscr B_{\sigma,M})=\frac{e^{-1/4}}{16}\mathcal N(\mathsf m_0,I_{12})
\bigl\{|\mathbf x-\sigma\mathbf x^\circ|<\varrho_\circ\bigr\},
\]
and the four events are disjoint because \(|\mathbf x^\circ|=2>\varrho_\circ\).

\emph{Episode.} Condition on \(\mathscr B_{1,S^{(0)}}\). The configuration drift
does not involve the rows, and the rows are deterministic, so the paragraph
\emph{Episode} of the proof above applies line by line from every starting
point in the ball, whatever the law of \(\mathbf X_0\) on it: the held
Ornstein--Uhlenbeck process, whose noise bound involves only the Brownian
motion; the contraction to the held path \cref{eq:rg-formation-cycle-center};
the residual bounds below \(1/5\) and above \(13/10\) on \([1/4,t^\circ_*]\),
which make each receiver's exterior cost exceed its partner cost by more than
one; and the relative-entropy bounds for \(t\ge t_\epsilon\), all deterministic
on the noise event. The rows enter only the exterior source probability: for a
receiver \(i\) with partner \(j\) and exterior source \(k\),
\[
\beta_{ik}(\mathbf x;\pi(t))\le\frac{\pi_{ik}(t)}{\pi_{ij}(t)}
\exp\Bigl[-\frac{\mathcal D_{e(i,k)}(\mathbf x)-\mathcal D_{e(i,j)}(\mathbf x)}{\mathrm T}\Bigr]
<1.23 e^{-1/\mathrm T}
\]
in place of \(e^{-1/\mathrm T}\). A proposal therefore leaves \(S^{(0)}\) with
probability at most \(5e^{-1/\mathrm T}\), the expected number of departing
proposals up to \(t^\circ_*\) is at most \(5t^\circ_*e^{-1/\mathrm T}\), and the
prefix event \(\mathscr A_{t^\circ_*}\) of that paragraph fails with conditional
probability at most
\[
\tilde p_\epsilon=73728 \epsilon^4t^\circ_*+5t^\circ_*e^{-1/\mathrm T}
<0.008+10^{-14}<\tfrac12,
\]
by the same monotonicity of \(\epsilon^4t^\circ_*\) and
\(e^{-1/\mathrm T}<\mathrm T\). On
\(\mathscr B_{1,S^{(0)}}\cap\mathscr A_{t^\circ_*}\), the block \(\{0,1\}\), of
size two and joined to its complement by \(e_1\) and \(e_3\), carries on its
internal edge \(e_0\) a belief relative entropy above
\((1-\sqrt3\varrho_\circ)^2/2>1/4\) at time zero, because
\(|\mathbf X_0-\mathbf x^\circ|<\varrho_\circ\) and every residual covector
has norm \(\sqrt3\); it has internal belief and model relative entropies at
most \(\epsilon^4\) on \([t_\epsilon,t^\circ_*]\), an interval of length one
inside \([0,2/\epsilon]\), and both boundary belief relative entropies above
\(1/2\) on \([1/4,t^\circ_*]\), so contrast formation occurs.

\emph{The other seeds.} A cycle automorphism \(g\), with link signs transformed
under reflections, acts on configurations by an orthogonal signed permutation
\(\mathsf O_g\) with \(U_{gS}(\mathsf O_g\mathbf x)=U_S(\mathbf x)\) and
\(\mathcal D_{e(gi,gj)}(\mathsf O_g\mathbf x)=\mathcal D_{e(i,j)}(\mathbf x)\).
With \((g\pi)_{gi,gj}=\pi_{ij}\), it therefore maps the kinetics with row path
\(\pi\) from \((\mathbf x,S)\) to the kinetics with row path \(g\pi\) from
\((\mathsf O_g\mathbf x,gS)\), driven by the Brownian motion \(\mathsf O_gW\)
and by the receiver uniform variables relabeled by \(g\), the variable of
receiver \(i\) serving receiver \(gi\) and being replaced by one minus itself
when \(g\) is a reflection. Indeed a reflection sends the source \(i+1\) of
receiver \(i\) to the source \(gi-1\) of receiver \(gi\), and the image
receiver takes it exactly when the original variable is at most
\(\beta_{i,i+1}(\mathbf x;\pi)=\beta_{gi,gi-1}(\mathsf O_g\mathbf x;g\pi)
=1-\beta_{gi,gi+1}(\mathsf O_g\mathbf x;g\pi)\), which is the rule of the
image kinetics applied to one minus that variable, except on the null event
that the variable equals its row; under the rotation the sources are carried
along and the variable is kept. The relabeled family is again a family of
independent uniform variables independent of the other inputs, so the image
process has the law of the kinetics with rows \(g\pi\), and \(g\pi\) has the
same ratio bound. The rotation \(i\mapsto i+1\) fixes
\(\mathbf x^\circ\) and maps \(S^{(0)}\) to \(S^{(1)}\), the reflection
\(i\mapsto-i\) maps \((\mathbf x^\circ,S^{(0)})\) to
\((-\mathbf x^\circ,S^{(1)})\), and their composition maps
\((\mathbf x^\circ,S^{(0)})\) to \((-\mathbf x^\circ,S^{(0)})\), so the
conditional bound holds on the other three seed events, with the images of the
block. Summing over the four disjoint
events,
\[
\Pr(\text{contrast formation})\ge(1-\tilde p_\epsilon)\frac{2e^{-1/4}}{16}
\sum_\pm\mathcal N(\mathsf m_0,I_{12})
\bigl\{|\mathbf x\mp\mathbf x^\circ|<\varrho_\circ\bigr\},
\]
which gives the first bound of (b) because \(\tilde p_\epsilon<1/2\). The
Gaussian density on the ball about \(\mathbf x^\circ\) is at least
\((2\pi)^{-6}\exp[-(|\mathbf x^\circ-\mathsf m_0|+\varrho_\circ)^2/2]\), the ball
has volume \(\pi^6\varrho_\circ^{12}/6!\), and \((2\pi)^{-6}\pi^6/6!=1/46080\),
which gives the second bound after dropping the ball about
\(-\mathbf x^\circ\). Frozen rows have ratio one, so the argument and the bound
hold without learning. The sub-event named in (b) is the event of part (a) of
the proposition above with the block of the held matching. For fixed rows
\(\pi\) the proposal step is reversible for \(q^\star_\pi\) by the same
symmetry \(q^\star_\pi(\mathbf x,S)\beta_{S'}(\mathbf x;\pi)
=q^\star_{\pi,\mathbf X}(\mathbf x)\beta_S(\mathbf x;\pi)\beta_{S'}(\mathbf x;\pi)\),
the sheet diffusions do not involve the rows, and the conditional law
\(\mathcal N(0,\mathrm TK_M^{-1})\) of \(q^\star_\pi\) given \(S_0=M\) does not
depend on \(\pi\), so the paragraph \emph{Stationary comparison} above gives
the invariance and the bound \(4e^{-1/(4\mathrm T)}\) for every \(\pi\). This
proves (b).

\emph{Invariant initial laws.} Since \(\mathsf O_g\) is orthogonal, \(g\) maps
\(\mathcal N(\mathsf m_0,I_{12})\) to \(\mathcal N(\mathsf O_g\mathsf m_0,I_{12})\),
so the initial law is invariant exactly when \(\mathsf m_0\) is fixed by the
rotation and the reflection. The rotation fixes exactly the means whose belief,
model, and link blocks are constant, and the reflection reverses every edge and
hence the sign of a constant link block, which leaves
\((\mathrm c_1\mathbf 1_4,\mathrm c_2\mathbf 1_4,0)\). For such a mean, since
the uniform assignment law and the uniform rows are also invariant, \(g\) maps
the recognition law \(q_t\) of the rows \(\pi\) to the recognition law
\(g_\#q_t\) of the rows \(g\pi\) from the same initial law, by the pathwise
image of the paragraph \emph{The other seeds}, whose relabeled driving inputs
have the law of the original ones; and \(g\pi\) solves
\cref{eq:phil-source-learning} with \(g_\#q_t(S_{gi}=gj)=q_t(S_i=j)\), so
\((g\pi,g_\#q)\) is a second solution from the same data and uniqueness in (a)
gives \(g\pi=\pi\). The reflection \(j\mapsto2i-j\) fixes receiver \(i\) and
exchanges its sources \(i\pm1\), so \(\pi_{i,i+1}=\pi_{i,i-1}=1/2\). This
proves (c). \(\square\)

\section{Depth one under exclusive source gating}
\label{sec:rg-formation-depth-one-proof}

\noindent\emph{Proof of \Cref{prop:rg-gated-depth-one}.}
Throughout, \(F=\operatorname{att}(S)\) and \(W=W_F\) once a sheet \(S\) is
fixed, \(P_W\) is the orthogonal projection onto \(W\), and
\(\mathsf Q_W\), \(\mathsf Q_{W^\perp}\) are matrices whose columns form
orthonormal bases of \(W\) and of \(W^\perp\). Since \(H_S\) is the sum of the
positive semidefinite maps \(\mathbf a_{e(i,S_i)}^\top\mathbf a_{e(i,S_i)}\),
its kernel is \(\bigcap_{e\in F}\ker\mathbf a_e=W_F\) and its range is
\(W_F^\perp\), the span of the rows of \(\mathbf a_F\), so
\(\mathrm k_S=\operatorname{rank}\mathbf a_F\). For every edge the Gram
matrix of the two rows of \(\mathbf a_e\) is
\(\bigl(\begin{smallmatrix}3&1\\1&3\end{smallmatrix}\bigr)\), as in
\cref{eq:rg-formation-cycle-normals}, so \(\|\mathbf a_e\|=2\).

\emph{(a) Two bands.} Let \(M=H_S+\mathrm T\Omega\). By Weyl's inequality
\citep[Sec.~4.3]{HornJohnson2013}, with eigenvalues in increasing order,
\(\lambda_k(H_S)+\mathrm T\lambda_{\min}(\Omega)\le\lambda_k(M)
\le\lambda_k(H_S)+\mathrm T\|\Omega\|\) for every \(k\). The first
\(2N+|E|-\mathrm k_S\) eigenvalues of \(H_S\) vanish, the remaining
\(\mathrm k_S\) lie in \([\mathrm h_S,\|H_S\|]\), and
\(\lambda_{\min}(\Omega)=\|\Omega^{-1}\|^{-1}\). The range was identified
above, and \(\pi\) does not enter \(M\).

\emph{(b) Gating wells.} For each receiver, the row sums to one and every
weight is at least \(\min_j\pi_{ij}\), so
\(\min_j\pi_{ij}e^{-\min_j\mathcal D_{e(i,j)}/\mathrm T}
\le\sum_j\pi_{ij}e^{-\mathcal D_{e(i,j)}/\mathrm T}
\le e^{-\min_j\mathcal D_{e(i,j)}/\mathrm T}\); taking \(-\mathrm T\log\) and
summing over \(i\) gives \cref{eq:rg-gated-well-depth}. Since \(V_S\) is a sum
over receivers and \(\pi_S\) a product,
\[
\sum_S\pi_Se^{-V_S/\mathrm T}=\prod_i\sum_{j\in\mathcal N_i}\pi_{ij}
e^{-\mathcal D_{e(i,j)}/\mathrm T}=e^{-V_\pi/\mathrm T},
\]
which is the configuration marginal up to the reference factor. If every edge of a
covering set \(F\) has zero residual at \(\mathbf x\), every receiver has an
option of zero cost, so \(V_{\min}(\mathbf x)=0\); a common translation of the
belief and model means of the agents of one component of \((V,F)\), with the
links fixed, changes no residual of an edge of \(F\), because the two
endpoints of such an edge either both move or both stay, so \(V_{\min}\)
remains zero.

\emph{(c) Rank formula.} In the coordinates \((\bar k,\mathfrak d,\omega)\)
with \(\mathfrak d=\bar k-\bar m\), the rows of \(\mathbf a_e\) for
\(e:j\to i\) are \(r^b_e=\bar k_i-\bar k_j-\omega_e\) and
\(r^m_e=r^b_e-(\mathfrak d_i-\mathfrak d_j)\); replacing the second by
\(r^m_e-r^b_e=-(\partial\mathfrak d)_e\) is an invertible change within the
rows of \(e\). The rows \(r^b_e\), \(e\in F\), read only \((\bar k,\omega)\)
and each has coefficient \(-1\) on its own link coordinate and zero on every
other, so they are linearly independent. The rows \((\partial\mathfrak d)_e\),
\(e\in F\), read only \(\mathfrak d\) and form the oriented incidence map of
\((V,F)\) on \(\R^V\), whose kernel consists of the functions constant on each
component, so their rank is \(N-\operatorname{cc}(F)\). The two families act
on disjoint coordinate blocks, so
\(\operatorname{rank}\mathbf a_F=|F|+N-\operatorname{cc}(F)\).

\emph{Minimum edge covers.} Let \(F=\operatorname{att}(S)\). Every vertex
selects an incident edge, so \(F\) covers \(V\) and every component of
\((V,F)\) contains an edge; one edge from each component is a matching, so
\(\operatorname{cc}(F)\le\mathfrak n\), and a connected graph on
\(\mathrm c\) vertices has at least \(\mathrm c-1\) edges, so
\(|F|\ge N-\operatorname{cc}(F)\). Hence
\(\mathrm k_S\ge2(N-\operatorname{cc}(F))\ge2(N-\mathfrak n)\), with
equality exactly when \(F\) is a forest with \(\mathfrak n\) components and
therefore \(N-\mathfrak n\) edges. The same two inequalities show that every
edge cover has at least \(N-\mathfrak n\) edges, and a maximum matching
together with one incident edge at each unmatched vertex, which exists
because no vertex is isolated, is an edge cover with exactly
\(N-\mathfrak n\) edges, since the other endpoint of each added edge is
matched, or else the matching could be enlarged. An edge cover with
\(N-\mathfrak n\) edges has at least \(\mathfrak n\) components and hence
exactly \(\mathfrak n\), and is a forest; so the covering forests with
\(\mathfrak n\) components are exactly the minimum edge covers. In a
minimum edge cover every edge has an endpoint of degree one, since otherwise
removing it would leave a smaller cover, and a tree in which every edge has
a leaf endpoint is a star, because two vertices of degree at least two would
be joined by a path whose first edge has no leaf endpoint. Conversely, a
star with center \(z\) is attended when every leaf selects \(z\) and \(z\)
selects one leaf, which the source sets allow, so every minimum edge cover
is some \(\operatorname{att}(S)\). Adding to \(F\) an edge that joins two
components raises \(|F|\) by one and lowers \(\operatorname{cc}(F)\) by one,
hence raises the rank by two.

\emph{Sheet weights.} Since \(H_S/\mathrm T+\Omega\succ0\),
\(Z_S=(2\pi)^{(2N+|E|)/2}\det(H_S/\mathrm T+\Omega)^{-1/2}\). In the basis
\((\mathsf Q_{W^\perp},\mathsf Q_W)\) the matrix \(H_S/\mathrm T+\Omega\) has
the blocks \(H_1/\mathrm T+\Omega_{11}\), \(\Omega_{10}\), \(\Omega_{01}\),
\(\Omega_{00}\), with \(H_1=\mathsf Q_{W^\perp}^\top H_S\mathsf Q_{W^\perp}\succeq\mathrm h_SI\)
and \(\Omega_{00}=\Omega|_W\). Eliminating the invertible block
\(\Omega_{00}\) by block row and column operations gives
\(\det(H_S/\mathrm T+\Omega)=\det\Omega_{00}\det(H_1/\mathrm T+\Sigma)\) with
\(\Sigma=\Omega_{11}-\Omega_{10}\Omega_{00}^{-1}\Omega_{01}\), and
\(0\prec\Sigma\preceq\Omega_{11}\preceq\|\Omega\|I\). Then
\(\mathrm T\Sigma\preceq(\mathrm T\|\Omega\|/\mathrm h_S)H_1\), so
\(\det H_1\le\det(H_1+\mathrm T\Sigma)\le(1+\mathrm T\|\Omega\|/\mathrm h_S)^{\mathrm k_S}\det H_1\),
and \(\det H_1=\operatorname{pdet}H_S\). Since
\(\det(H_1/\mathrm T+\Sigma)=\mathrm T^{-\mathrm k_S}\det(H_1+\mathrm T\Sigma)\),
the determinant lies between \(\mathrm T^{-\mathrm k_S}\mathrm z_S^{-2}\) and
\((1+\mathrm T\|\Omega\|/\mathrm h_S)^{\mathrm k_S}\mathrm T^{-\mathrm k_S}\mathrm z_S^{-2}\),
which is \cref{eq:rg-gated-sheet-weight}. The same block form shows that
under \(q_\Omega(\cdot\mid S)\), which is
\(\mathcal N(0,(H_S/\mathrm T+\Omega)^{-1})\), the coordinates
\(y=\mathsf Q_W^\top\mathbf x\) given \(z=\mathsf Q_{W^\perp}^\top\mathbf x\)
are Gaussian with precision \(\Omega_{00}\), the \((W,W)\) block of the joint
precision, which does not depend on \(\mathrm T\).

\emph{Merge cost.} \(q_\Omega(S)=\pi_SZ_S/\sum_{S'}\pi_{S'}Z_{S'}
\le\pi_SZ_S/\sum_{S'\in\mathcal M}\pi_{S'}Z_{S'}\). The upper bound of
\cref{eq:rg-gated-sheet-weight} for \(S\) and its lower bound for each
\(S'\in\mathcal M\), where \(\mathrm k_{S'}=2(N-\mathfrak n)\) and
\(\mathrm h_{S'}\ge\mathrm h_*\), give
\(q_\Omega(S)\le\mathrm T^{(\mathrm k_S-2(N-\mathfrak n))/2}
(1+\mathrm T\|\Omega\|/\mathrm h_*)^{N-\mathfrak n}
\pi_S\mathrm z_S/\sum_{S'\in\mathcal M}\pi_{S'}\mathrm z_{S'}\). If
\(\operatorname{cc}(\operatorname{att}(S))\le\mathfrak n-m\), then
\(\mathrm k_S\ge2(N-\operatorname{cc}(\operatorname{att}(S)))\ge2(N-\mathfrak n)+2m\);
with \(\mathrm T\le1\), summing over such \(S\) gives
\cref{eq:rg-gated-merge-cost}.

\emph{(d) Block law.} The sheets with
\(\operatorname{cc}(\operatorname{att}(S))<\mathfrak n\) carry mass at most
\(\mathrm K_1\mathrm T\) by \cref{eq:rg-gated-merge-cost}. Fix a sheet with
\(\operatorname{cc}(F)=\mathfrak n\). No component of \((V,F)\) contains
two disjoint edges, since replacing that component's edge in the matching of
one edge per component would give \(\mathfrak n+1\) disjoint edges. A
connected graph in which no two edges are disjoint is a star or a triangle:
fix an edge \(uv\); every other edge meets \(\{u,v\}\); if there are edges
\(vw\) and \(uw'\) with \(w\ne u\) and \(w'\ne v\), then \(w=w'\), and a
further edge would have to meet \(\{u,v\}\) and would then be disjoint from
\(vw\) or from \(uw\), so the graph is the triangle \(uvw\); otherwise all
edges share one vertex. Replacing each triangle component by a two-edge path
on the same three vertices gives a covering forest of stars with
\(\mathfrak n\) components, a minimum edge cover by (c), so the partition
of \(V\) into the components of \(F\) is a minimum-edge-cover partition.
The coherent-block partition \(\mathfrak P_{\mathrm a}(\mathbf x)\) differs
from it only if some attended edge has \(|\mathbf a_e\mathbf x|>\mathrm a\) or
some edge joining two components of \(F\) has
\(|\mathbf a_e\mathbf x|\le\mathrm a\). Under \(q_\Omega(\cdot\mid S)\) the
configuration is \(\mathcal N(0,\mathrm T(H_S+\mathrm T\Omega)^{-1})\). For
\(e\in F\), \(M=H_S+\mathrm T\Omega\succeq\mathbf a_e^\top\mathbf a_e\), so
\(\|\mathbf a_eM^{-1/2}\|\le1\) and the covariance
\(\mathrm T\mathbf a_eM^{-1}\mathbf a_e^\top\) of \(\mathbf a_e\mathbf x\) is at
most \(\mathrm TI_2\); each of its two centered Gaussian coordinates has
variance at most \(\mathrm T\), and \(|\mathbf a_e\mathbf x|>\mathrm a\)
forces one of them beyond \(\mathrm a/\sqrt2\), an event of probability at
most \(2\cdot2e^{-\mathrm a^2/(4\mathrm T)}\). For an edge \(e\notin F\)
joining two components, \(\operatorname{rank}\mathbf a_{F\cup\{e\}}
=\operatorname{rank}\mathbf a_F+\operatorname{rank}(\mathbf a_e|_{W})\),
because \(\ker\mathbf a_{F\cup\{e\}}=W\cap\ker\mathbf a_e\); by (c) the
restriction has rank two, so \(\mathbf a^W_e=\mathbf a_e\mathsf Q_W\) has two nonzero
singular values, the least of which is at least \(\mathrm s_*\). Writing
\(\mathbf x=\mathsf Q_{W^\perp}z+\mathsf Q_Wy\), the conditional law of \(y\)
given \(z\) has precision \(\Omega_{00}\preceq\|\Omega\|I\), so \(\mathbf a^W_ey\) is
Gaussian in \(\R^2\) with covariance
\(\mathbf a^W_e\Omega_{00}^{-1}\mathbf a^{W\top}_e\succeq\|\Omega\|^{-1}\mathbf a^W_e\mathbf a^{W\top}_e\succeq(\mathrm s_*^2/\|\Omega\|)I_2\)
and density at most \((2\pi)^{-1}\|\Omega\|/\mathrm s_*^2\). The event
\(|\mathbf a_e\mathbf x|\le\mathrm a\) places \(\mathbf a^W_ey=\mathbf a_e\mathbf x-\mathbf a_e\mathsf Q_{W^\perp}z\)
in a disc of radius \(\mathrm a\) about a point determined by \(z\), so its
conditional probability is at most \(\pi\mathrm a^2\cdot\|\Omega\|/(2\pi\mathrm s_*^2)\).
A union bound over at most \(|E|\) edges of each kind, integrated over the
sheets, gives \cref{eq:rg-gated-block-law}. \(\square\)

\section{Proof of Pairing Selection}
\label{sec:phil-pairing-selection-proof}

\noindent\emph{Proof of \Cref{conj:phil-contrast-formation}.}
Write \(p_i=\pi_{i,i+1}\), encode the assignment by
\(\sigma_i=+1\) when receiver \(i\) selects \(i+1\), and put
\(\varepsilon=\epsilon_{\mathrm L}\). The state remains
\((\mathbf x,\sigma)\in\R^{12}\times\{\pm1\}^4\), with the shared edge
costs of \cref{eq:rg-formation-cycle-cost}. The multiplicity of edge
\(e\) is
\(m_e(\sigma)=1_{\{\sigma_e=+1\}}+1_{\{\sigma_{e+1}=-1\}}\).
In this proof define
\[
 \Omega=I+R_{\mathrm{ev}},\qquad
 \mathsf A_\sigma=H_\sigma+\mathrm T\Omega,\qquad
 H_\sigma=\nabla^2V_\sigma,\qquad
 \mathsf A_\sigma\succeq\mathrm TI.
\]
Let \(\mathsf w_\sigma=\det(\mathsf A_\sigma/\mathrm T)^{-1/2}\),
\(\mathsf Z(p)=\sum_\sigma\pi_\sigma(p)\mathsf w_\sigma\), and
\[
 \mu_p(\mathbf x,\sigma)
 =\frac{\pi_\sigma(p)}{(2\pi)^6\mathsf Z(p)}
   \exp\left[-\frac{\mathbf x^\top\mathsf A_\sigma\mathbf x}
                    {2\mathrm T}\right],\qquad
 \pi_\sigma(p)=\prod_i p_i^{1_{\{\sigma_i=+1\}}}
                         (1-p_i)^{1_{\{\sigma_i=-1\}}}.
\]
Each held-sheet diffusion preserves its Gaussian factor, and a clock ring
resamples \(\sigma\) from its conditional law under \(\mu_p\). Thus
\(\mu_p\) is invariant for the frozen rows, including boundary rows on their
supported sheets. Set
\begin{equation}
 \bar\beta_i(p)=\mu_p(\sigma_i=+1),\qquad
 \mathfrak g_i(p)=\partial_{p_i}\log\mathsf Z(p),\qquad
 \mathscr V_i(p)=\bar\beta_i(p)-p_i
               =p_i(1-p_i)\mathfrak g_i(p).
\label{eq:phil-selection-reduced-field}
\end{equation}
The last identity follows by differentiating the product prior. Since
\(\mathsf Z\) is a positive multilinear polynomial on the closed cube,
\(\mathfrak g\) and \(\mathscr V\) extend analytically to a neighborhood of
that cube and have bounded derivatives there. The well-posedness argument
for \cref{prop:phil-seeded-contrast-formation} applies at every fixed
\(\mathrm T>0\); it gives deterministic continuously differentiable rows and
\begin{equation}
 \min_i\{p_i(t),1-p_i(t)\}\ge\tfrac12e^{-\varepsilon t},\qquad
 \|p(t)-\tfrac12\|_\infty\le\tfrac12\varepsilon t.
\label{eq:phil-selection-row-bounds}
\end{equation}
All constants below may depend on the fixed temperature. Constants in
bounds on a fixed initial law may also depend on its mean, whereas the
window constant in the conclusion will depend only on the temperature.

\subsection*{Frozen convergence and the initial imprint}

Here \(\operatorname{Ent}_\nu(g)=\int g\log(g/\int g d\nu)d\nu\)
for nonnegative \(g\), and total variation is \(\sup_A|q(A)-\nu(A)|\).
For the time-one frozen kernel, the function
\(\mathsf V(\mathbf x)=|\mathbf x|^2\) satisfies
\begin{equation}
 \E_{\mathbf x}|\mathbf X_t|^2
 \le e^{-2\mathrm Tt}|\mathbf x|^2+12(1-e^{-2\mathrm Tt}),\qquad
 P_1^p\mathsf V\le e^{-2\mathrm T}\mathsf V+12.
\label{eq:phil-selection-moment}
\end{equation}
Indeed, jumps leave \(\mathbf X\) unchanged, and the diffusion generator
applied to its squared norm is at most \(24\mathrm T-2\mathrm T|\mathbf X|^2\).
Choose \(R>24/(1-e^{-2\mathrm T})\). For any frozen \(p\), some assignment
\(\sigma_*(p)\) has \(\pi_{\sigma_*}\ge1/16\). Restrict the rate-one
clock to exactly one ring in \([0,1/2]\) and none afterward. Its probability
is \(e^{-1}/2\). Starting with \(|\mathbf x|\le\sqrt R\), the first held
Gaussian motion is in the ball of radius
\(r'=\sqrt R+\sqrt{24\mathrm T}\) at that ring with probability at least
\(1/2\). On the ball, the probability to redraw \(\sigma_*\) is at least
\(16^{-1}e^{-v(r')/\mathrm T}\), where
\(v(r')=\max_{|\mathbf z|\le r',\sigma}V_\sigma(\mathbf z)<\infty\).
The subsequent held Gaussian transition lasts between \(1/2\) and \(1\).
Its means are bounded and its covariance matrices have common positive
lower and finite upper bounds. Its density consequently has a positive
common lower bound on the unit ball. Hence, for some \(\alpha_{\mathrm T}>0\)
independent of \(p\),
\[
 P_1^p((\mathbf x,\sigma),\cdot)
 \ge\alpha_{\mathrm T}\operatorname{Unif}(B_1)
                         \otimes\delta_{\sigma_*(p)}
 \quad\text{when }|\mathbf x|^2\le R.
\]
These are the drift and minorization hypotheses of
\citet[Assumptions 1--2 and Theorems 1.2--1.3]{HairerMattingly2011}.
The explicit dependence in their Theorem 1.3 makes the resulting constants
uniform in the frozen rows, even though the minorizing measure depends on
\(p\). Thus \(\mu_p\) is the unique invariant law and
\begin{equation}
 |P_t^p\phi(\mathbf x,\sigma)-\mu_p\phi|
 \le\mathrm K_{\mathrm T}e^{-\lambda_{\mathrm T}t}
       (1+|\mathbf x|^2)\|\phi\|_\infty.
\label{eq:phil-selection-harris}
\end{equation}
Intermediate times follow by applying the skeleton estimate to
\(P_r^p\phi\), \(0\le r<1\).

Put \(\mathbf v=(1,-1,1,-1)\),
\(f_{\mathbf v}(\sigma)=\frac12\sum_i\mathbf v_i\sigma_i\), and define
\begin{equation}
 \mathsf I_i(\mathsf m_0)
 =\int_0^\infty[q_t^0(\sigma_i=+1)-\tfrac12]dt,\qquad
 \mathsf j_{\mathrm T}(\mathsf m_0)
 =\mathbf v\cdot\mathsf I(\mathsf m_0).
\label{eq:phil-selection-imprint}
\end{equation}
At uniform rows, the stationary law is invariant under cycle rotations and
reflections, so each stationary row is \(1/2\) and
\(\mu_{1/2}f_{\mathbf v}=0\). Estimate
\cref{eq:phil-selection-harris} proves absolute convergence of these
integrals. It also gives a Poisson integral
\begin{equation}
 \mathfrak R_{\mathrm T}(\mathbf x,\sigma)
 =\int_0^\infty P_t^{1/2}f_{\mathbf v}(\mathbf x,\sigma)dt,\qquad
 |\mathfrak R_{\mathrm T}(\mathbf x,\sigma)|
 \le\mathrm K_{\mathrm T}(1+|\mathbf x|^2).
\label{eq:phil-selection-poisson}
\end{equation}
Convolution of this integral with the initial Gaussian expresses
\(\mathsf j_{\mathrm T}\) as an entire function of its complexified mean:
for every compact set of complex means, all differentiated integrands are
bounded by an integrable polynomial times a Gaussian. The same holds for
\(\mathsf I\). A rotation by one changes the sign of
\(\mathsf j_{\mathrm T}\).

For \(t\le1/(10\varepsilon)\), all row entries are in
\([\mathrm c_0,1-\mathrm c_0]\), \(\mathrm c_0=0.45\). A posterior receiver probability has derivative
at most \(1/[4\mathrm c_0(1-\mathrm c_0)]\le1/(2\mathrm c_0)\) with respect to its prior entry.
Coupling the common clocks and their four posterior rows therefore gives
\[
 \|q_t^\varepsilon-q_t^0\|_{\mathrm{TV}}
 \le\frac2{\mathrm c_0}\int_0^t\|p(s)-\tfrac12\|_\infty ds
 \le\frac{\varepsilon t^2}{2\mathrm c_0}.
\]
Variation of constants in the learning equation, followed by
\cref{eq:phil-selection-harris}, gives
\begin{equation}
 \left|p(t)-\tfrac12-\varepsilon\mathsf I(\mathsf m_0)\right|_\infty
 \le\mathrm K_{\mathrm T,\mathsf m_0}
       (\varepsilon e^{-\lambda_{\mathrm T}t}+\varepsilon^2t)
       +\frac{\varepsilon^2t^3}{6\mathrm c_0}.
\label{eq:phil-selection-imprint-expansion}
\end{equation}
In particular, at \(t_1=\varepsilon^{-1/4}\),
\(p(t_1)-1/2=\varepsilon\mathsf I(\mathsf m_0)+o(\varepsilon)\).

\subsection*{Nonvanishing of the matching imprint}

Take \(\mathbf u_\circ=(0_4,0_4,(0,1,0,-1))\), whose two residual vectors
are \((0,-1,0,1)\). In the noiseless switched motion the subspace
\(\bar k=\bar m\) is invariant. After removing the scalar factor
\(e^{-\mathrm Tt}\), its common residual vector satisfies
\begin{equation}
 \dot\rho=-\mathsf G M_\sigma\rho,\qquad
 \mathsf G=4I-A_{C_4},\qquad
 M_\sigma=\operatorname{diag}(m_e(\sigma)).
\label{eq:phil-selection-tied-residual}
\end{equation}
Here the cyclic difference matrix contributes
\(2I-A_{C_4}\) and the shared-link drift contributes \(2I\).
Of the sixteen assignments, six have \(m_1=m_3\). Their first held motion
is \(e^{-ct}(0,-1,0,1)\), with \(c=4m_1\in\{0,4,8\}\).
At their first clock ring every receiver prefers the even-edge matching
\(S^{(0)}\). Holding that matching then leaves the odd residuals constant
and the even residuals zero in the scalar-free coordinates.

Four further assignments have multiplicity patterns
\[
 (1,1,2,0),\quad(1,0,2,1),\quad(2,0,1,1),\quad(2,1,1,0).
\]
The signed swap of coordinates \(1,3\) and the rotation by two preserve
the even-edge matching and carry the initial residual to itself or its
negative. They reduce these four cases to \((1,1,2,0)\). Write its
scalar-free coordinates as \(\rho=(-a,-b,-d,h)\). Then
\[
 \begin{pmatrix}\dot a\\\dot b\\\dot d\end{pmatrix}
 =\begin{pmatrix}-4&1&0\\1&-4&2\\0&1&-8\end{pmatrix}
   \begin{pmatrix}a\\b\\d\end{pmatrix},\qquad
 (a,b,d)(0)=(0,1,0),\qquad \dot h=-(a+2d),\ h(0)=1.
\]
The nonnegative cone is invariant, and the negative row sums give
\(\max(a,b,d)\le1\), hence \(a\le1/4\), \(d\le1/8\).
The half-spaces
\(a-3b/4-11d/25\le0\) and \(d-b/4\le0\) are invariant: their
outward derivatives on the corresponding boundary are
\(-b/400-7d/100\) and \(-a/4-b/8\), respectively. Consequently
\(a\le43b/50\), \(d\le b/4\), and \(b\ge e^{-4t}\).
The attended block is similar to a negative definite symmetric matrix
because \(\mathsf G\succ0\). Solving
\(\rho(\infty)=\rho(0)-\mathsf G z\), with \(z_3=0\) and the first
three final coordinates zero, gives \(h(\infty)=6/7\). Thus
\(h\ge6/7\).

If the first ring occurs at \(t_{\mathrm r}\) and redraws \(S^{(0)}\),
the subsequent scalar-free held motion, in elapsed time \(r\), is
\[
 \begin{aligned}
 a_r&=a_0e^{-8r},&d_r&=d_0e^{-8r},\\
 b_r&=b_0+\tfrac14(a_0+d_0)(1-e^{-8r}),&
 h_r&=h_0-\tfrac14(a_0+d_0)(1-e^{-8r})\ge\tfrac{171}{224}.
 \end{aligned}
\]
Every adjacent odd-edge absolute residual exceeds its even-edge neighbor
by at least \((7/50)e^{-4t_{\mathrm r}}\). Combining these four cases
with the preceding six, the deterministic motion \(\mathbf Y\) from
\(s\mathbf u_\circ\) has, throughout the hypothesized matching hold,
\begin{equation}
 \sqrt{\mathcal D_{\mathrm{odd}}(\mathbf Y_t)}
 -\sqrt{\mathcal D_{\mathrm{even}}(\mathbf Y_t)}
 \ge\frac7{50}s e^{-\mathrm Tt-8t_{\mathrm r}},
 \qquad t\ge t_{\mathrm r}.
\label{eq:phil-selection-lock-margin}
\end{equation}
There are therefore ten good initial assignments. For the remaining six
only \(f_{\mathbf v}\ge-2\) will be used.

Set \(h_s=(\log s)/\mathrm T\), for \(s>e\), and write the initial
configuration as \(s\mathbf u_\circ+\xi\), \(\xi\sim\mathcal N(0,I)\).
Along the actual assignment path decompose
\(\mathbf X=\mathbf Y+\mathbf E+\mathbf Z\), where the first two terms
solve its switched deterministic equation from \(s\mathbf u_\circ\) and
\(\xi\), respectively, while
\[
 d\mathbf Z=-\mathsf A_{\sigma_t}\mathbf Zdt+\sqrt{2\mathrm T}dW_t,
 \qquad\mathbf Z_0=0.
\]
All three are adapted differential equations, and
\(|\mathbf E_t|\le e^{-\mathrm Tt}|\xi|\).
It\^o's formula yields
\[
 \mathscr L e^{|\mathbf Z|^2/4}
 \le\mathrm T(6-|\mathbf Z|^2/4)e^{|\mathbf Z|^2/4}
 \le\mathrm Te^5.
\]
Stopping at radius \(\log s\), and using
\(\E e^{|\xi|^2/4}=64\), bounds the probability that either
\(|\xi|>\log s\) or \(\sup_{t\le h_s}|\mathbf Z_t|>\log s\) by
\((65+\mathrm Te^5h_s)e^{-(\log s)^2/4}\).
On its complement each square-root edge cost changes from its
\(\mathbf Y\) value by at most \(2\mathrm C_N\log s\), where
\(\mathrm C_N\) is the largest norm of a residual-pair map divided by
\(\sqrt2\).

Define
\[
 L_s(t_{\mathrm r})=h_s-\frac8{\mathrm T}t_{\mathrm r}
                         -\frac2{\mathrm T}\log\log s,\qquad
 B_s=(L_s(t_{\mathrm r})-t_{\mathrm r})_+.
\]
For a good initial assignment, \cref{eq:phil-selection-lock-margin}
gives a square-root gap at least \((7/50)(\log s)^2\) on
\([t_{\mathrm r},L_s]\). For sufficiently large \(s\), while both noise bounds hold,
the actual energy gap is at least
\(\kappa_0(\log s)^4\), \(\kappa_0=(7/50)^2/4\).
Each receiver's wrong choice at a ring then has probability at most
\(e^{-\kappa_0(\log s)^4/\mathrm T}\). Apply this bound at the first
ring, then use the rate-one compensator up to the first wrong redraw or
exit from the noise bound, stopping at \(L_s\wedge h_s\).
The endpoint \(L_s\) is known after the first ring, and the gap estimate
holds until that first failure. Thus the matching holds throughout the
nonempty interval \([t_{\mathrm r},L_s]\) except on an event of
probability at most
\begin{equation}
 \delta_s=(65+\mathrm Te^5h_s)e^{-(\log s)^2/4}
       +4(1+h_s)e^{-\kappa_0(\log s)^4/\mathrm T},
 \qquad h_s\delta_s\longrightarrow0.
\label{eq:phil-selection-lock-failure}
\end{equation}
Since \(\E t_{\mathrm r}=1\), a good initial assignment contributes at
least
\[
 -2h_s+4\E B_s-4h_s\delta_s
 \ge2h_s-4(1+8/\mathrm T)-(8/\mathrm T)\log\log s-4h_s\delta_s
\]
to \(\E\int_0^{h_s}f_{\mathbf v}(\sigma_t)dt\). The other six contribute
at least \(-2h_s\). Averaging the sixteen uniform initial assignments gives
\[
 \E\int_0^{h_s}f_{\mathbf v}(\sigma_t)dt
 \ge\tfrac12h_s-\frac5{\mathrm T}\log\log s-O_{\mathrm T}(1).
\]
By \cref{eq:phil-selection-moment},
\(\E|\mathbf X_{h_s}|^2\le14\). The Markov property and
\cref{eq:phil-selection-poisson} bound the integral after \(h_s\) in
absolute value by a constant independent of \(s\). Therefore
\begin{equation}
 \mathsf j_{\mathrm T}(s\mathbf u_\circ)
 \ge\frac{\log s}{2\mathrm T}
       -\frac5{\mathrm T}\log\log s-O_{\mathrm T}(1).
\label{eq:phil-selection-nonvanishing}
\end{equation}
This proves that \(\mathsf j_{\mathrm T}\) is not identically zero at any
fixed \(\mathrm T>0\). Its zero set is Lebesgue-null. For completeness,
this last assertion follows by induction on dimension: a nonzero analytic
function in one variable has isolated zeros; in several variables at least
one coefficient \(\partial_n^kf(\mathbf x',0)\) is nonzero as an analytic
function, or analyticity would make \(f\) vanish everywhere. Outside the
null zero set of that coefficient, each one-dimensional slice has isolated
zeros, and Fubini's theorem concludes the induction.

\subsection*{Entropy balance for the evolving recognition law}

Fix a finite horizon \(h\). All rows are bounded away from zero on this
interval by \cref{eq:phil-selection-row-bounds}. Write \(q_\sigma(t,\mathbf x)\)
for the evolving density, \(Q=\sum_\sigma q_\sigma\), and
\(b_\sigma(t,\mathbf x)=\mu_{p(t)}(\mathbf x,\sigma)/
 \sum_\eta\mu_{p(t)}(\mathbf x,\eta)\).
For the forward Ornstein--Uhlenbeck semigroup \(K_\sigma\), the last-proposal
formula is
\begin{equation}
 q_\sigma(t)=e^{-t}K_\sigma(t)q_\sigma(0)
   +\int_0^te^{-(t-r)}K_\sigma(t-r)[b_\sigma(r)Q(r)]dr.
\label{eq:phil-selection-forward}
\end{equation}
Iterating this positive formula, and bounding each proposal factor by one,
reduces each term to a Gaussian motion along a deterministic sequence of
sheets. Put \(M=\max_\sigma\|\mathsf A_\sigma\|\).
Its means have norm at most \(|\mathsf m_0|\), and its covariances lie between
\(e^{-2Mh}I\) and \((1+2\mathrm Th)I\). Summing the at most \(16^n\)
sequences on each time simplex, of volume \(t^n/n!\), gives a finite upper
Gaussian envelope. The no-proposal term gives a lower Gaussian envelope.
Consequently, uniformly on \([0,h]\),
\begin{equation}
 0<\mathrm c_he^{-\mathrm C_h(1+|\mathbf x|^2)}
 \le q_\sigma(t,\mathbf x)\le\mathrm C_he^{-\mathrm c_h|\mathbf x|^2}.
\label{eq:phil-selection-density-envelopes}
\end{equation}
This bound does not assert Gaussianity conditional on a state-dependent
assignment history.

Differentiating one Ornstein--Uhlenbeck kernel against a function bounded
by a polynomial times a Gaussian gives a bound of the same type with
factor \((t-r)^{-1/2}\). To check it, use the Gaussian kernel covariance
bounds proportional to \(t-r\), absorb its linear derivative into a
Gaussian with twice the covariance, and complete the square in the
convolution. This singularity is integrable in
\cref{eq:phil-selection-forward}, proving a first-derivative Gaussian
bound. Moreover
\[
 \nabla b_\sigma=-\frac{b_\sigma}{\mathrm T}
       \left(H_\sigma\mathbf x-\sum_\eta b_\eta H_\eta\mathbf x\right),
 \qquad |\nabla b_\sigma|\le\mathrm K_{\mathrm T}|\mathbf x|.
\]
The source \(b_\sigma Q\) therefore has a first derivative with the same
type of bound. For the second derivative, transfer one derivative to this
source using
\(\nabla K_\sigma(r)g=e^{r\mathsf A_\sigma}K_\sigma(r)\nabla g\),
and differentiate the other kernel factor. The singularity is again only
\(r^{-1/2}\). The explicit initial Gaussian term supplies the bounds at
zero. The forward equation obtained by taking a time difference quotient
in \cref{eq:phil-selection-forward} then bounds the time derivative too:
\begin{equation}
 |\nabla q_\sigma|+|\nabla^2q_\sigma|+|\partial_tq_\sigma|
 \le\mathrm C_h(1+|\mathbf x|)^k e^{-\mathrm c_h|\mathbf x|^2}.
\label{eq:phil-selection-density-derivatives}
\end{equation}
The derivative transfer follows directly by integration by parts in the
Gaussian transition density; its inverse mean map is uniformly bounded
on this finite interval. These estimates justify the differentiations,
without assuming finite Fisher information.

Set \(f_\sigma=q_\sigma/\mu_{p(t),\sigma}\) and
\(\mathsf H(t)=\KL(q_t\Vert\mu_{p(t)})\).
The logarithms of both densities grow at most quadratically by
\cref{eq:phil-selection-density-envelopes}, and
\(\partial_t\log\mu_p\) is bounded in \(\mathbf x\), since only its prior
weight and normalizer vary. Together with
\cref{eq:phil-selection-density-derivatives}, this proves that
\(\mathsf H\) is absolutely continuous on \([0,h]\) and permits its
integral to be differentiated.
To integrate the diffusion by parts, first multiply by a smooth radial
cutoff \(\chi_R\), equal to one on \(B_R\), zero outside \(B_{2R}\),
with \(|\nabla\chi_R|\le\mathrm C/R\), increasing to one.
On each compact cylinder the densities have positive lower bounds, so
\(\log f_\sigma\) is a legitimate test function. The boundary term is
bounded by
\[
 \frac{\mathrm C_h}{R}\int_{R\le|\mathbf x|\le2R}
       (1+|\mathbf x|)^{k+3}e^{-\mathrm c_h|\mathbf x|^2}d\mathbf x
 \longrightarrow0.
\]
Detailed balance \(\mu_\sigma b_\eta=\mu_\eta b_\sigma\) symmetrizes the
finite jump sum. Both resulting production integrands are nonnegative,
so monotone convergence removes the cutoff and simultaneously proves
finiteness and local time integrability of the production. The result is
\begin{equation}
 \begin{aligned}
 \mathsf H'(t)&=-\mathscr D_{p(t)}(f)
               -\sum_\sigma\int q_\sigma\partial_t\log\mu_{p(t),\sigma},\\
 \mathscr D_p(f)&=\mathrm T\sum_\sigma\int\mu_{p,\sigma}
                         \frac{|\nabla f_\sigma|^2}{f_\sigma}\\
 &\quad+\frac12\sum_{\sigma,\eta}\int\mu_{p,\sigma}b_\eta
                  (f_\eta-f_\sigma)(\log f_\eta-\log f_\sigma).
 \end{aligned}
\label{eq:phil-selection-entropy-identity}
\end{equation}
The equality holds almost everywhere and in integrated form on every finite
interval. The jump term is integrable already from the Gaussian envelopes
and quadratic logarithmic bounds; thus the cutoff argument does not
presuppose a bound on the diffusion production. This also proves
\(\mathsf H(0)<\infty\).

\subsection*{Entropy contraction at frozen rows}

In the original twelve-dimensional coordinates put
\(\mathsf B_\sigma=\mathsf A_\sigma/\mathrm T\), the Gaussian sheet
precision, and set
\[
 \kappa=\min_\sigma\lambda_{\min}(\mathsf B_\sigma)>0,\qquad
 M=\max_\sigma\|\mathsf B_\sigma\|,\qquad
 q_\sigma^*=\mu_p(\sigma)
 =\frac{\pi_\sigma\mathsf w_\sigma}{\mathsf Z},
 \qquad q_{\min}^*=\min_\sigma q_\sigma^*.
\]
Let \(\gamma_\sigma=\mathcal N(0,\mathsf B_\sigma^{-1})\), and, for a
real function \(u\), define
\[
 \mathscr E_c(u)=\mathrm T\sum_\sigma\int\mu_{p,\sigma}|\nabla u_\sigma|^2,
 \qquad
 \mathscr E_j(u)=\frac12\sum_{\sigma,\eta}\int\mu_{p,\sigma}b_\eta
                                    (u_\sigma-u_\eta)^2.
\]
The Gaussian inequalities needed below follow from the explicit
Ornstein--Uhlenbeck semigroup with invariant law \(\gamma_\sigma\).
Its gradient commutes as
\(\nabla P_tg=e^{-t\mathsf B_\sigma}P_t\nabla g\). Differentiating
entropy and applying Cauchy--Schwarz inside this kernel gives
\[
 -\frac d{dt}\operatorname{Ent}_{\gamma_\sigma}(P_tg)
 =\int\frac{|\nabla P_tg|^2}{P_tg}d\gamma_\sigma
 \le e^{-2\kappa t}\int\frac{|\nabla g|^2}{g}d\gamma_\sigma.
\]
For bounded positive smooth \(g\), the explicit Gaussian formula and
dominated convergence make the entropy tend to zero. Integration, with
\(g=u^2\), gives
\begin{equation}
 \operatorname{Ent}_{\gamma_\sigma}(u^2)
 \le\frac2\kappa\int|\nabla u|^2d\gamma_\sigma,
 \qquad
 \operatorname{Var}_{\gamma_\sigma}(u)
 \le\frac1\kappa\int|\nabla u|^2d\gamma_\sigma.
\label{eq:phil-selection-gaussian-inequalities}
\end{equation}
The second estimate follows by differentiating variance in the same
semigroup. Positive regularization, truncation, and smooth cutoffs extend
both estimates to their form domains: the cutoff errors tend to zero by
integrability of the squared function and its gradient, and the entropy
passes to the limit by its variational formula.

Replace \(u\) by \(|u|\), which decreases both energies, and normalize
\(\int u^2d\mu_p=1\). Put
\(z_\sigma=(\int u_\sigma^2d\gamma_\sigma)^{1/2}\).
The entropy chain rule gives
\[
 \operatorname{Ent}_{\mu_p}(u^2)
 =\sum_\sigma q_\sigma^*
             \operatorname{Ent}_{\gamma_\sigma}(u_\sigma^2)
       +\operatorname{Ent}_{q^*}(z^2).
\]
Here \(0\le z_\sigma\le R=(q_{\min}^*)^{-1/2}\).
The function \(\psi(z)=z^2\log z^2-z^2+1\) satisfies
\(\psi(1)=\psi'(1)=0\) and \(\psi''(z)=4+4\log z\), whence
\(\psi(z)\le2(1+\log R)(z-1)^2\) on \([0,R]\).
Since \(\sum q_\sigma^*z_\sigma^2=1\),
\[
 \operatorname{Ent}_{q^*}(z^2)
 \le L_p\operatorname{Var}_{q^*}(z),\qquad
 L_p=4+2\log(1/q_{\min}^*).
\]
Indeed \(\sum q_\sigma^*(z_\sigma-1)^2
=2(1-\E_{q^*}z)\le2\operatorname{Var}_{q^*}(z)\).

Fix a ball \(B=B_\rho\) in these coordinates, and define
\[
 m_\sigma=|B|^{-1}\int_Bu_\sigma,
 \quad m_* =\sum_\sigma\pi_\sigma m_\sigma,
 \quad X_\sigma=
 \left[\frac{(2\pi)^6\mathsf w_\sigma
                   e^{\rho^2\|\mathsf B_\sigma\|/2}}{|B|}-1\right]_+,
 \quad B_{\mathrm T}=\max_\sigma(1+\sqrt{X_\sigma})^2.
\]
Variance about any center is an upper bound on the variance, so
\[
 \operatorname{Var}_{q^*}(z)
 \le2\sum_\sigma q_\sigma^*(z_\sigma-m_\sigma)^2
     +2\sum_\sigma q_\sigma^*(m_\sigma-m_*)^2.
\]
The first sum is at most \((B_{\mathrm T}/(\mathrm T\kappa))\mathscr E_c(u)\):
\(|z_\sigma-\gamma_\sigma u_\sigma|\le
\sqrt{\operatorname{Var}_{\gamma_\sigma}(u_\sigma)}\), and
Cauchy--Schwarz bounds \(|m_\sigma-\gamma_\sigma u_\sigma|\) by the same
standard deviation times
\(\sqrt{\chi^2(\operatorname{Unif}(B)\Vert\gamma_\sigma)}
\le\sqrt{X_\sigma}\).
For the second sum, Jensen's inequality gives
\[
 \sum_\sigma q_\sigma^*(m_\sigma-m_*)^2
 \le\frac{\mathsf w_{\max}}
              {\mathsf Z|B|}
       \sum_{\sigma,\eta}\pi_\sigma\pi_\eta
                          \int_B(u_\sigma-u_\eta)^2.
\]
On \(B\), the explicit Gaussian densities give
\[
 \mu_{p,\sigma}b_\eta
 \ge\frac{\pi_\sigma\pi_\eta e^{-\rho^2M}}
              {(2\pi)^6\mathsf Z}.
\]
The normalizer therefore cancels, and the preceding sum is at most
\(2(2\pi)^6\mathsf w_{\max}e^{\rho^2M}
 \mathscr E_j(u)/|B|\).
Combining these bounds proves
\[
 \operatorname{Ent}_{\mu_p}(u^2)
 \le\frac2{\mathrm T\kappa}(1+L_pB_{\mathrm T})\mathscr E_c(u)
    +\frac{4L_p(2\pi)^6\mathsf w_{\max}e^{\rho^2M}}{|B|}
                                                       \mathscr E_j(u).
\]
For \(u=\sqrt f\), the production is at least
\(4(\mathscr E_c(u)+\mathscr E_j(u))\). Its jump part uses
\((a-b)(\log a-\log b)\ge4(\sqrt a-\sqrt b)^2\), which follows by
Cauchy--Schwarz on the integral of \(r^{-1/2}\) between \(a\) and \(b\).
The inequality extends to the form domain by the preceding approximations.
Finally,
\[
 q_{\min}^*\ge p_{\min}^4
                 \frac{\mathsf w_{\min}}
                      {\mathsf w_{\max}},\qquad
 p_{\min}=\min_i\{p_i,1-p_i\}.
\]
Thus there are finite positive constants depending only on \(\mathrm T\)
with
\begin{equation}
 \mathscr D_p(f)\ge\lambda(\mathrm T,p)
                         \operatorname{Ent}_{\mu_p}(f),\qquad
 \lambda(\mathrm T,p)\ge
 \frac{\mathrm c_{\mathrm T}}
      {\mathrm K_{\mathrm T}+8\log(1/p_{\min})}.
\label{eq:phil-selection-entropy-contraction}
\end{equation}
The factor \(\mathrm T\) in \(\mathscr E_c\) is the same diffusion
coefficient as in \cref{eq:phil-selection-entropy-identity}.

\subsection*{The matching branch of the small-temperature landscape}

The needed sheet weights can be computed using two four-dimensional
determinants. Let \(D\) be the cyclic difference matrix, and
\(C_\sigma=\operatorname{diag}(m_e(\sigma))\). In the orthogonal
coordinates \((\bar k+\bar m)/\sqrt2\),
\((\bar k-\bar m)/\sqrt2\), and \(\omega\), the energy splits into
\(\frac12(Du-\sqrt2\omega)^\top C_\sigma(Du-\sqrt2\omega)
 +\frac12(Dv)^\top C_\sigma Dv\), while the discrepancy coordinate has
reference precision \(5I\). The determinant identity for the first block
then gives
\begin{equation}
 \det\mathsf A_\sigma
 =\mathrm T^4
  \det[\mathrm TI+C_\sigma(DD^\top+2I)]
  \det[5\mathrm TI+D^\top C_\sigma D].
\label{eq:phil-selection-sheet-determinants}
\end{equation}
The two matching assignments, eight end-path assignments, four middle-path
assignments, and two oriented cycles have the respective determinants
\[
 \begin{aligned}
 \mathsf d_M&=25\mathrm T^8(\mathrm T+8)^2(5\mathrm T+4)^2,\\
 \mathsf d_E&=5\mathrm T^6
   (\mathrm T^3+16\mathrm T^2+77\mathrm T+112)
   (125\mathrm T^3+200\mathrm T^2+85\mathrm T+8),\\
 \mathsf d_N&=5\mathrm T^6(\mathrm T+4)(5\mathrm T+2)
   (\mathrm T^2+12\mathrm T+28)(25\mathrm T^2+30\mathrm T+4),\\
 \mathsf d_C&=5\mathrm T^5(\mathrm T+2)(\mathrm T+4)^2(\mathrm T+6)
                                     (5\mathrm T+2)^2(5\mathrm T+4).
 \end{aligned}
\]
These follow in \cref{eq:phil-selection-sheet-determinants} from the
multiplicity representatives \((2,0,2,0)\), \((2,1,1,0)\),
\((1,2,1,0)\), and \((1,1,1,1)\), respectively.
Normalize the matching weight to one and write
\[
 \mathsf r_E=\sqrt{\mathsf d_M/\mathsf d_E},\qquad
 \mathsf r_N=\sqrt{\mathsf d_M/\mathsf d_N},\qquad
 \mathsf r_C=\sqrt{\mathsf d_M/\mathsf d_C}.
\]
Their leading terms are
\(\mathsf r_E,\mathsf r_N=\sqrt{40/7}\mathrm T+O(\mathrm T^2)\) and
\(\mathsf r_C=\sqrt{5/3}\mathrm T^{3/2}(1+O(\mathrm T))\).
Choose \(\mathrm T_0>0\) so that \(0<\mathsf r_E,\mathsf r_N,\mathsf r_C\le1/8\) for
\(0<\mathrm T\le\mathrm T_0\).

Use matching coordinates
\(x=(p_0,1-p_1,p_2,1-p_3)\). The normalized evidence is the
multilinear polynomial whose vertex weights are one at \(\mathbf0,\mathbf1\),
\(\mathsf r_E\) for subsets of cardinality one or three, \(\mathsf r_N\) for the four adjacent
pairs, and \(\mathsf r_C\) for the two opposite pairs. Complementation preserves
these weights, so the centered reduced field is odd. With
\(D_0=1+4\mathsf r_E+2\mathsf r_N+\mathsf r_C\), its Jacobian at \(x=\mathbf1/2\) is
\[
 \frac1{D_0}
 \begin{pmatrix}
 0&1-\mathsf r_C&1-2\mathsf r_N+\mathsf r_C&1-\mathsf r_C\\
 1-\mathsf r_C&0&1-\mathsf r_C&1-2\mathsf r_N+\mathsf r_C\\
 1-2\mathsf r_N+\mathsf r_C&1-\mathsf r_C&0&1-\mathsf r_C\\
 1-\mathsf r_C&1-2\mathsf r_N+\mathsf r_C&1-\mathsf r_C&0
 \end{pmatrix}.
\]
For example, the adjacent Hessian entry of the evidence is
\((1-\mathsf r_C)/2\), its uniform value is \(D_0/8\), and the row mobility is
\(1/4\), giving the displayed entry. The eigenvalues are
\begin{equation}
 \alpha=\frac{3-2\mathsf r_N-\mathsf r_C}{D_0}>0,\qquad
 \beta=\frac{-1-2\mathsf r_N+3\mathsf r_C}{D_0}<0,\qquad
 \eta=\frac{-1+2\mathsf r_N-\mathsf r_C}{D_0}<0\quad\text{twice}.
\label{eq:phil-selection-saddle-spectrum}
\end{equation}
The positive eigendirection is \(\mathbf1\) in matching coordinates,
or \(\mathbf v\) in the original rows. In fact the chosen bounds give
\(\alpha\ge7/5\), \(\beta\le-1/3\), and \(\eta\le-2/5\).

Cyclic invariance makes the matching diagonal \(x=z\mathbf1\) invariant.
Writing \(w=z(1-z)\), its normalized evidence is
\[
 Z_{\mathrm{diag}}(z)
 =1+4(\mathsf r_E-1)w+(2-8\mathsf r_E+4\mathsf r_N+2\mathsf r_C)w^2,
 \qquad
 \dot z=\frac{z(1-z)}{4Z_{\mathrm{diag}}(z)}Z_{\mathrm{diag}}'(z).
\]
The derivative is \((1-2z)[4(\mathsf r_E-1)+2(2-8\mathsf r_E+4\mathsf r_N+2\mathsf r_C)w]\).
The bracket is negative at both endpoints \(w=0,1/4\), where its values
are \(4(\mathsf r_E-1)\) and \(-3+2\mathsf r_N+\mathsf r_C\). It is affine in \(w\), hence negative
throughout. Thus the diagonal motion tends to \(\mathbf1\) from
\(z>1/2\), and to \(\mathbf0\) from \(z<1/2\), reaching every prescribed
endpoint neighborhood in finite slow time. At either vertex the Jacobian
is \((\mathsf r_E-1)I\): a single receiver flip changes weight one to \(\mathsf r_E\), the
mobility vanishes, and its diagonal derivative supplies the factor
\(-(1-\mathsf r_E)\). Both vertices are therefore sinks with positive inward rate
\(\nu=1-\mathsf r_E\).

\subsection*{Tracking, escape, and selection}

Let \(a_i(t)=q_t(\sigma_i=+1)\) and
\(\Delta_i(t)=a_i(t)-\bar\beta_i(p(t))\).
Differentiating the normalized product prior gives
\[
 \partial_{p_i}\log\mu_p
 =\frac{1_{\{\sigma_i=+1\}}-\bar\beta_i(p)}{p_i(1-p_i)}.
\]
Inserting the learning equation and
\cref{eq:phil-selection-reduced-field} in
\cref{eq:phil-selection-entropy-identity} yields
\begin{equation}
 \mathsf H'
 =-\mathscr D_p(f)-\varepsilon\sum_i\mathfrak g_i(p)\Delta_i
                  -\varepsilon\sum_i\frac{\Delta_i^2}{p_i(1-p_i)}
 \le-\lambda(\mathrm T,p)\mathsf H
      +\varepsilon\|\mathfrak g(p)\|_1\sqrt{\mathsf H/2}.
\label{eq:phil-selection-learning-entropy}
\end{equation}
The last bound uses \(|\Delta_i|\le\operatorname{TV}(q_t,\mu_p)
\le\sqrt{\mathsf H/2}\). The nonpositive square is kept until this step;
it removes any inverse-row factor from the forcing.
Writing \(\Lambda(t)=\int_0^t\lambda(\mathrm T,p(r))dr\), comparison
for \(\sqrt{\mathsf H+\delta^2}\), followed by \(\delta\downarrow0\), gives
\begin{equation}
 \sqrt{\mathsf H(t)}
 \le e^{-\Lambda(t)/2}\sqrt{\mathsf H(0)}
    +\frac{\varepsilon}{2\sqrt2}
       \int_0^te^{-(\Lambda(t)-\Lambda(r))/2}
                         \|\mathfrak g(p(r))\|_1dr.
\label{eq:phil-selection-tracking}
\end{equation}
The bounded fitness, \cref{eq:phil-selection-row-bounds}, and
\cref{eq:phil-selection-entropy-contraction} show that, on
\(t\le2\mathrm C_{\mathrm T}\varepsilon^{-1}\log(1/\varepsilon)\),
\begin{equation}
 \operatorname{TV}(q_t,\mu_{p(t)})
 \le e^{-\Lambda(t)/2}\sqrt{\mathsf H(0)/2}
       +O_{\mathrm T}(\varepsilon\log(1/\varepsilon)).
\label{eq:phil-selection-global-tracking}
\end{equation}
Indeed \(\lambda\) is bounded below on this horizon by a positive
constant divided by \(1+\log(1/\varepsilon)\).

For the saddle a finer bound is needed. Let
\(\|\mathfrak g(p)-\mathfrak g(p')\|_1
 \le L_{\mathrm T}\|p-p'\|_\infty\), and suppose
\(\lambda\ge\lambda_\flat>0\) on \([s_0,t]\).
The estimate \(|\dot p_i|\le\varepsilon\) in
\cref{eq:phil-selection-tracking} gives
\begin{equation}
 \operatorname{TV}(q_t,\mu_{p(t)})
 \le e^{-\lambda_\flat(t-s_0)/2}\sqrt{\mathsf H(s_0)/2}
   +\frac{\varepsilon}{2\lambda_\flat}\|\mathfrak g(p(t))\|_1
   +\frac{L_{\mathrm T}\varepsilon^2}{\lambda_\flat^2}.
\label{eq:phil-selection-refined-tracking}
\end{equation}
The constants follow by integrating \(1\) and \(t-r\) against
\(e^{-\lambda_\flat(t-r)/2}\).

Fix \(0<\mathrm T\le\mathrm T_0\) and a mean with
\(\mathsf j_{\mathrm T}(\mathsf m_0)\ne0\), and put
\(\mathbf e=\mathbf v/2\). In slow time
\(\mathfrak s=\varepsilon t\), write
\(y=p-1/2=a\mathbf e+z\), \(z\perp\mathbf e\).
By \cref{eq:phil-selection-imprint-expansion}, at
\(t_1=\varepsilon^{-1/4}\),
\[
 a(t_1)=\varepsilon(\mathsf j_{\mathrm T}(\mathsf m_0)/2+o(1)),
 \qquad |z(t_1)|=O(\varepsilon).
\]
On \([0,t_1]\), symmetry gives \(\mathfrak g(1/2)=0\), so
\(\|\mathfrak g(p(t))\|_1\le L_{\mathrm T}\varepsilon t/2\).
The bound in \cref{eq:phil-selection-tracking} gives
\(\sqrt{\mathsf H(t_1)}=O(\varepsilon^{7/4})\).
Consequently, until exit from any sufficiently small fixed ball,
\cref{eq:phil-selection-refined-tracking} implies
\begin{equation}
 |\Delta|_2\le\kappa_\varepsilon|y|+b_\varepsilon,
 \qquad \kappa_\varepsilon=O(\varepsilon),\quad
 b_\varepsilon=O(\varepsilon^{7/4})=o(|a(t_1)|).
\label{eq:phil-selection-relative-error}
\end{equation}
The constants here may depend on the fixed mean.

Let \(\beta_\perp=\max(\beta,\eta)<0\) and
\(\gamma=\alpha-\beta_\perp>0\).
Oddness and diagonal invariance of the analytic reduced field imply,
for \(|y|\le r\),
\[
 |\langle\mathscr V(1/2+y),\mathbf e\rangle-\alpha a|
       \le\mathrm C|y|^3,\qquad
 |P_\perp\mathscr V(1/2+y)-D\mathscr V(1/2)z|
       \le\mathrm C|y|^2|z|.
\]
For the second inequality, subtract the zero perpendicular field on
\(a\mathbf e\) and integrate its derivative along the transverse
segment. The derivative differs from its value at zero by
\(O(|y|^2)\) because the field is odd.
Writing \(w=|z|/|a|\), the upper right derivatives of norms give
\begin{equation}
 \begin{aligned}
 \frac{d|a|}{d\mathfrak s}
 &\ge |a|[\alpha-(\mathrm Cr^2+\kappa_\varepsilon)(1+w)]
                -b_\varepsilon,\\
 D^+w
 &\le-[\gamma-\mathrm Cr^2
             -(\mathrm Cr^2+\kappa_\varepsilon)(1+w)]w
        +(1+w)(\kappa_\varepsilon+b_\varepsilon/|a|).
 \end{aligned}
\label{eq:phil-selection-cone}
\end{equation}
Choose \(W\) larger than twice the initial limiting ratio, then choose
\(r>0\) so small that the terms \(\mathrm Cr^2(1+W)\) consume less than
one quarter of both positive rates. For small enough \(\varepsilon\),
\cref{eq:phil-selection-relative-error} controls the remaining terms.
A bootstrap then keeps \(w\le W\), makes \(|a|\) increase without changing
sign, and gives
\[
 w(\mathfrak s)
 \le w(\mathfrak s_1)e^{-\gamma(\mathfrak s-\mathfrak s_1)/2}
  +\frac{2(1+W)}\gamma
       (\kappa_\varepsilon+b_\varepsilon/|a(t_1)|).
\]
The lower growth bound for \(|a|\) yields an exit time no larger than
\[
 \frac{1+O(r^2)+o(1)}\alpha\log(1/\varepsilon)
\]
in slow time. The corresponding upper norm-growth bound forces the exit
time to tend to infinity. Hence \(w\to0\) at exit, and the exit point
tends to \(\operatorname{sign}(\mathsf j_{\mathrm T})r\mathbf e\).

The reduced diagonal orbit from that point enters any fixed neighborhood
of its matching vertex in finite slow time. The tracking error in
\cref{eq:phil-selection-global-tracking} tends to zero after \(t_1\);
finite-time Gr\"onwall comparison therefore carries the evolving rows into
the same neighborhood. At its selected vertex let
\(\zeta_i=|p_i-p_i^M|\). By the derivative computation above and continuity,
a sufficiently small neighborhood has
\((1-\zeta_i)|\mathfrak g_i(p)|\ge\nu/2\), where \(\nu=1-\mathsf r_E>0\), and
the fitness has the sign directed toward that vertex. Thus
\[
 D^+\zeta_i\le-\tfrac\nu2\zeta_i+
                   \operatorname{TV}(q_t,\mu_{p(t)})
\]
in slow time. The neighborhood remains invariant once the tracking floor
is small enough, and \(\zeta_i\) decays exponentially to a vanishing
floor.

Choose any \(\mathrm C_{\mathrm T}>1/\alpha\), independently of the mean,
and then choose the local radius so that its escape coefficient is less
than \(\mathrm C_{\mathrm T}\). The time remaining after the finite
transit tends to infinity uniformly on the slow-time window
\([\mathrm C_{\mathrm T}\log(1/\varepsilon),
  2\mathrm C_{\mathrm T}\log(1/\varepsilon)]\).
The row distances therefore vanish uniformly on that window. Every
normalized sheet weight is at most the matching weight one, so
\[
 \mu_p(M)\ge\pi_M\ge1-\sum_i\zeta_i,
 \qquad
 q_t(S\ne M)\le\sum_i\zeta_i+
                     \operatorname{TV}(q_t,\mu_{p(t)}).
\]
The last total variation tends to zero by
\cref{eq:phil-selection-global-tracking}; its integrated contraction
exponent at the beginning of the window grows at least as a positive
constant times \(\varepsilon^{-1}\).
This proves the first conclusion for every mean outside the null set in
\cref{eq:phil-selection-nonvanishing}. For the frozen arm,
\cref{eq:phil-selection-harris} and rotation symmetry prove the second
conclusion at every \(\mathrm T>0\) and every mean. The convergence is
pointwise in the fixed mean; no uniform bound near the zero set of the
imprint has been used. \(\square\)

\section{Proof of Joint Selection Along a Schedule}
\label{sec:phil-joint-selection-proof}

\noindent\emph{Proof of \Cref{cor:phil-joint-selection}.}
The saddle estimate \cref{eq:phil-selection-saddle-spectrum} gives
\(\alpha(\mathrm T)\ge7/5\) throughout the stated temperature range.
The proof of \Cref{conj:phil-contrast-formation} permits any window
constant greater than \(1/\alpha(\mathrm T)\), independently of the
initial mean. Thus the common choice \(2\) is valid. Write
\[
 \tau_\varepsilon=2\varepsilon^{-1}\log(1/\varepsilon),
 \qquad W_\varepsilon=[\tau_\varepsilon,2\tau_\varepsilon],
 \qquad 0<\varepsilon<e^{-1}.
\]
For a perfect matching \(M\), define
\[
 \nu_{\mathrm T}^M
 =\mathcal N(0,\mathrm T\mathsf A_M^{-1})\otimes\delta_M,
 \qquad
 \mu_{\mathrm T}=\mu_{\mathrm T,1/2}.
\]
Here \(\mu_{\mathrm T,p}\) is the frozen-row law used in the preceding
proof, with its temperature now displayed explicitly. The inverse
exists because \(\mathsf A_M\succeq\mathrm TI\). Let
\(M_{\mathrm T}(\mathsf m_0)\) be the matching selected by the sign of
\(\mathsf j_{\mathrm T}(\mathsf m_0)\), choosing either matching on its
null zero set.

First record the full one-time consequence of the fixed-temperature
proof. Conditional on assignment \(M\), the law
\(\mu_{\mathrm T,p}\) has the Gaussian factor of \(\nu_{\mathrm T}^M\),
independently of \(p\). The normalized sheet weights are at most the
matching weight one. For
\(\zeta_i=|p_i-p_i^M|\), therefore,
\[
 \mu_{\mathrm T,p}(M)\ge\pi_M\ge1-\sum_i\zeta_i,
 \qquad
 \operatorname{TV}(\mu_{\mathrm T,p},\nu_{\mathrm T}^M)
 =1-\mu_{\mathrm T,p}(M)\le\sum_i\zeta_i.
\]
The row convergence and
\cref{eq:phil-selection-global-tracking} imply
\[
 \sup_{t\in W_\varepsilon}
 \operatorname{TV}
 (q_t^{\mathrm T,\varepsilon,\mathsf m_0},
       \nu_{\mathrm T}^{M_{\mathrm T}(\mathsf m_0)})
 \longrightarrow0
\]
at fixed temperature outside the imprint's null zero set.
Moreover \cref{eq:phil-selection-tracking} itself gives convergence
of the learned relative entropy to zero, uniformly on this window.
For frozen rows, the entropy identity and contraction give
\[
 \KL(q_t^{\mathrm T,0,\mathsf m_0}\Vert\mu_{\mathrm T})
 \le e^{-\lambda(\mathrm T,1/2)t}
       \KL(q_0^{\mathrm T,0,\mathsf m_0}\Vert\mu_{\mathrm T}).
\]
The initial entropy is finite by the Gaussian density bounds, and
\(\lambda(\mathrm T,1/2)>0\). Frozen entropy and total variation
therefore tend to zero on the window for every fixed mean.

Use the explicit entropy majorants as follows. Let
\[
 G_{\mathrm T}=\sup_{p\in[0,1]^4}\|\mathfrak g(p)\|_1,\qquad
 \ell_{\mathrm T}(\varepsilon)
 =\frac{\mathrm c_{\mathrm T}}
 {\mathrm K_{\mathrm T}+8\log2+32\log(1/\varepsilon)},\qquad
 \lambda_u(\mathrm T)
 =\frac{\mathrm c_{\mathrm T}}{\mathrm K_{\mathrm T}+8\log2}.
\]
The first constant is finite. The row lower bound and
\cref{eq:phil-selection-entropy-contraction} make
\(\ell_{\mathrm T}(\varepsilon)\) a rate lower bound throughout
\(t\le2\tau_\varepsilon\), and \(\lambda_u(\mathrm T)\) a rate lower
bound at uniform rows. The common initial entropy equals
\[
 H_0(\mathrm T,\mathsf m_0)
 =\log\mathsf Z(1/2)-6+
 \frac{\operatorname{tr}\overline{\mathsf A}
              +\mathsf m_0^\top\overline{\mathsf A}\mathsf m_0}
      {2\mathrm T},
 \qquad
 \overline{\mathsf A}=\frac1{16}\sum_\sigma\mathsf A_\sigma.
\]
Indeed the log density ratio is
\(\log\mathsf Z(1/2)-|\mathbf x-\mathsf m_0|^2/2+
\mathbf x^\top\mathsf A_\sigma\mathbf x/(2\mathrm T)\);
taking its initial expectation proves the formula.
Thus \(H_0\) is continuous and nonnegative. Uniformly on
\(W_\varepsilon\), the entropy tracking estimate gives
\[
 \begin{aligned}
 \KL(q_t^{\mathrm T,\varepsilon,\mathsf m_0}
               \Vert\mu_{\mathrm T,p(t)})
 &\le B_L(\mathrm T,\varepsilon,\mathsf m_0)\\
 &:=\left[
 e^{-\ell_{\mathrm T}(\varepsilon)\tau_\varepsilon/2}\sqrt{H_0}
 +\frac{\varepsilon G_{\mathrm T}}
             {\sqrt2\ell_{\mathrm T}(\varepsilon)}
 \right]^2,\\
 \KL(q_t^{\mathrm T,0,\mathsf m_0}\Vert\mu_{\mathrm T})
 &\le B_F(\mathrm T,\varepsilon,\mathsf m_0)
 :=e^{-\lambda_u(\mathrm T)\tau_\varepsilon}H_0.
 \end{aligned}
\]
These majorants are continuous in positive mobility and mean and tend
to zero at fixed temperature and mean as \(\varepsilon\downarrow0\).

Let \(\mathcal E_{\mathrm T}(\varepsilon,\mathsf m_0)\) be the minimum
of one and the maximum of the following five nonnegative quantities:
\[
 \begin{aligned}
 &\sup_{t\in W_\varepsilon}
          \max_i|p_i(t)-p_i^{M_{\mathrm T}(\mathsf m_0)}|,\\
 &\sup_{t\in W_\varepsilon}
       \operatorname{TV}
       (q_t^{\mathrm T,\varepsilon,\mathsf m_0},
                   \nu_{\mathrm T}^{M_{\mathrm T}(\mathsf m_0)}),\\
 &B_L(\mathrm T,\varepsilon,\mathsf m_0),\\
 &\sup_{t\in W_\varepsilon}
       \operatorname{TV}
          (q_t^{\mathrm T,0,\mathsf m_0},\mu_{\mathrm T}),\\
 &B_F(\mathrm T,\varepsilon,\mathsf m_0).
 \end{aligned}
\]
The preceding estimates prove
\(\mathcal E_{\mathrm T}(\varepsilon,\mathsf m_0)\to0\)
for almost every fixed mean at each fixed temperature.

We will take a supremum over mobility, so measurability must be checked.
Fix a temperature and restrict the positive mobility and the mean to
compact sets, with the mobility bounded away from zero. The corresponding
horizons are bounded, and
\cref{eq:phil-selection-row-bounds} bounds every row away from zero and
one. On this row set the posterior redraw probabilities are uniformly
Lipschitz in the rows, independently of the configuration.
Maximally couple the initial laws and then use common Brownian motions
and proposal clocks until the first different redraw. The law difference
is bounded by the initial total variation plus a constant times the
integral of row differences. Subtracting the learning equations bounds
the row difference by integrals of row and law differences and by the
mobility difference. Gr\"onwall's inequality proves total-variation
continuity in the initial mean and positive mobility on each bounded
time interval. Time continuity follows from
\cref{eq:phil-selection-forward}.

Parametrizing the window by \(t=u\tau_\varepsilon\),
\(1\le u\le2\), proves continuity in positive mobility of each row or
total-variation window supremum. The explicit entropy majorants are
continuous by their formulas. The matching as a function of the mean is
Borel because its imprint is analytic. Thus the error is measurable
in the mean, and its supremum over an interval of positive mobilities
equals its supremum over rational mobilities. No uniform continuity at
zero mobility or separate entropy-continuity assertion is needed.

Fix the sequence \(\mathrm T_n\downarrow0\) in the stated temperature
interval and \(0<\eta_n<1\) with \(\eta_n\to0\). For the prescribed
exponent \(a>7/2\), put \(h_n=\mathrm T_n^a\). The following argument
in fact permits arbitrary positive caps \(h_n\).
On the mean space \(\R^{12}\), introduce the auxiliary probability
measure \(\rho=\mathcal N(0,I_{12})\). Its density is strictly positive;
it introduces no assumption that the physical mean is random.
Set
\[
 F_n(r,\mathsf m_0)
 =\sup_{0<\varepsilon\le r}
       \mathcal E_{\mathrm T_n}(\varepsilon,\mathsf m_0).
\]
The measurable functions \(F_n\) are bounded by one. At fixed \(n\),
they tend to zero for \(\rho\)-almost every mean as \(r\downarrow0\).
Dominated convergence yields
\[
 \int F_n(r,\mathsf m_0)d\rho(\mathsf m_0)\longrightarrow0.
\]
Choose \(d_n=2^{-k_n}\) recursively, where \(k_n\) is the smallest
positive integer satisfying
\[
 d_n\le\min(e^{-1},2^{-n},h_n,d_{n-1}/2),
 \qquad
 \int F_n(d_n,\mathsf m_0)d\rho(\mathsf m_0)
       \le\eta_n2^{-n},
\]
with the \(d_{n-1}/2\) condition omitted for \(n=1\). The preceding
limit proves existence. This defines deterministic thresholds
independent of the mean; it does not give numerical values or an explicit
power law.

Let \(B_n=\{\mathsf m_0:F_n(d_n,\mathsf m_0)>\eta_n\}\).
Markov's inequality and countable subadditivity give
\[
 \rho(B_n)\le2^{-n},
 \qquad
 \rho\left(\bigcup_{n\ge N}B_n\right)\le2^{1-N}.
\]
The decreasing intersection of these tail unions consequently has
measure zero. Add the countable union of the imprint zero sets. The
resulting exceptional set is Lebesgue-null because \(\rho\) and Lebesgue
measure have the same null sets. Outside it, for every sufficiently large
\(n\), all \(0<\varepsilon\le d_n\) satisfy
\(F_n(d_n,\mathsf m_0)\le\eta_n\). Since \(\eta_n<1\), clipping has
no effect on this inequality: the row and total-variation errors and both entropy majorants are at
most \(\eta_n\). The majorants bound the actual entropies, giving
all five stated conclusions.

The frozen matching probability is, exactly,
\[
 \mu_{\mathrm T}(M)
 =\frac1{2+8\mathsf r_E+4\mathsf r_N+2\mathsf r_C}
 =\frac1{2D_0(\mathrm T)}
 =\frac12+O(\mathrm T).
\]
Indeed \(\mathsf r_E,\mathsf r_N=O(\mathrm T)\) and
\(\mathsf r_C=O(\mathrm T^{3/2})\), as proved from the sheet
determinants. Pushforward cannot increase total variation. If
\(q^{S}\) denotes the assignment marginal, the two established
total-variation estimates therefore imply, uniformly on \(W_\varepsilon\),
\[
 \left|
 \operatorname{TV}
 (q_t^{\mathrm T_n,\varepsilon,\mathsf m_0,S},
                  q_t^{\mathrm T_n,0,\mathsf m_0,S})
 -\left(1-\frac1{2D_0(\mathrm T_n)}\right)
 \right|\le2\eta_n.
\]
Here the comparison distance is exactly
\(\operatorname{TV}(\delta_M,\mu_{\mathrm T_n}^S)
=1-\mu_{\mathrm T_n}(M)\), for either matching.
The assignment-law contrast thus tends to \(1/2\), even if the selected
matching changes with \(n\).

Finally, defining \(\mathrm T(\varepsilon)=\mathrm T_n\) on
\(d_{n+1}<\varepsilon\le d_n\) gives a single stepwise path with both
parameters tending to zero. The proved bounds hold for all sufficiently
small mobilities on this path. Arbitrary mobility caps are admissible;
for example \(h_n=\mathrm T_n^n\) forces
\(\varepsilon_n/\mathrm T_n^a\to0\) for each fixed \(a>0\) whenever
\(\varepsilon_n\le d_n\). This imposes an upper bound and does not certify
a prescribed equality between temperature and mobility.

Nothing in this argument makes the imprint sign eventually constant as
temperature vanishes. Its fixed-temperature nonvanishing and null zero
set give no quantitative lower bound at a fixed mean in that limit.
Accordingly the matching is allowed to depend on \(n\), and the result
concerns complete one-time laws, not complete histories or conditional
first-formation events. \(\square\)





\subsection*{An explicit temperature bound for tracking}

\propositionheading{Temperature-dependent entropy tracking}{prop:phil-temperature-tracking}
In the four-cycle of \cref{conj:phil-contrast-formation}, let
\(0<\mathrm T\le1/100\), \(p\in(0,1)^4\), and
\(p_{\min}=\min_i\{p_i,1-p_i\}\). Define
\[
 \mathrm K_0=210\sqrt{2\pi/5},\qquad
 \mathrm c_* =\min\left\{
 \frac2{1+4\mathrm K_0e^{13/8}},\frac1{\mathrm K_0e^{13/4}}\right\}.
\]
The entropy production of the full frozen law satisfies
\[
 \mathscr D_p(f)\ge
 \frac{\mathrm c_*\mathrm T^{5/2}}
 {4+3\log(1/\mathrm T)+8\log(1/p_{\min})}
 \operatorname{Ent}_{\mu_p}(f).
\]
Write \(\varepsilon=\epsilon_{\mathrm L}\),
\(t_{\mathrm L}=2\varepsilon^{-1}\log(1/\varepsilon)\), and
\[
 \mathrm D(\mathrm T,\varepsilon)
 =4+3\log(1/\mathrm T)+8\log2+32\log(1/\varepsilon),\qquad
 \ell=\frac{\mathrm c_*\mathrm T^{5/2}}
                {\mathrm D(\mathrm T,\varepsilon)}.
\]
For \(0<\varepsilon<e^{-1}\) and \(0\le t\le2t_{\mathrm L}\),
\begin{equation}
 \operatorname{TV}(q_t,\mu_{p(t)})
 \le e^{-\ell t/2}
       \sqrt{\frac{13(12+|\mathsf m_0|^2)}{4\mathrm T}}
    +\frac{2\varepsilon\mathrm D(\mathrm T,\varepsilon)}
                 {\mathrm c_*\mathrm T^{7/2}}.
\label{eq:phil-temperature-tracking}
\end{equation}
In particular, for \(\varepsilon=\mathrm T^a\) with any \(a>7/2\),
this total variation tends to zero uniformly on
\([t_{\mathrm L},2t_{\mathrm L}]\) and on bounded sets of initial means.
This conclusion compares one-time laws; it supplies no imprint lower bound.
\status{ESTABLISHED}

\paragraph{Proof.}
Use \(u=(\bar k+\bar m)/\sqrt2\),
\(v=(\bar k-\bar m)/\sqrt2\), and \(\omega\), and put
\(R=[D,-\sqrt2I]\). The common kernel
\(K=\ker R\oplus\ker D\) has dimension five. Every
\(\mathsf H_\sigma\) vanishes there, and the reference precision preserves
\(K\) and \(N=K^\perp\). Its restrictions have eigenvalues one and five.
The passive process on \(K\) is therefore an independent Gaussian
Ornstein--Uhlenbeck process with entropy contraction at least \(2\mathrm T\).
The active process on \(N\), together with the assignment, is autonomous:
all redraw costs depend only on its residuals. Rescaling its coordinate by
\(\mathrm T^{-1/2}\) gives unit diffusion and drift matrix
\(a_\sigma=(\mathsf H_\sigma+\mathrm T\mathsf\Omega)|_N\).
Since every edge multiplicity is at most two,
\[
 \mathrm TI\preceq a_\sigma\preceq13I.
\]
Indeed \(RR^\top=DD^\top+2I\) has eigenvalues \(2,4,4,6\),
while \(D^\top D\) has eigenvalues \(0,2,2,4\), so
\(\|\mathsf H_\sigma\|\le12\).

Divide the four determinants in
\cref{eq:phil-selection-sheet-determinants} by the passive determinant
\(5\mathrm T^5\). Their positive coefficients, evaluated at
\(\mathrm T=1/100\) for upper bounds, give the active bounds
\[
 \begin{aligned}
 5120\mathrm T^3&\le d_M\le5262\mathrm T^3,&
 896\mathrm T&\le d_E\le1017\mathrm T,\\
 896\mathrm T&\le d_N\le995\mathrm T,&
 3072&\le d_C<5120.
 \end{aligned}
\]
Thus \(w_M\) and \(w_C\) are the largest and smallest active sheet
weights, and \(w_M\le\mathrm T^{-3/2}/\sqrt{5120}\),
\(w_C/w_M\ge\mathrm T^{3/2}\). The finite-label constant in the
entropy proof consequently satisfies
\[
 L_p\le4+3\log(1/\mathrm T)+8\log(1/p_{\min}).
\]

Apply the Gaussian-sheet and prior-weighted-ball argument of
\cref{eq:phil-selection-gaussian-inequalities,eq:phil-selection-entropy-contraction}
in the seven active coordinates, using the ball of radius \(1/2\), whose
volume is \(\pi^3/840\). Its Gaussian comparison factor is at most
\(4\mathrm K_0e^{13/8}\mathrm T^{-3/2}\), because
\[
 \frac{(2\pi)^{7/2}}{(\pi^3/840)\sqrt{5120}}=\mathrm K_0.
\]
More explicitly, the two energy coefficients in that argument obey
\[
 C_c\le(2+8\mathrm K_0e^{13/8})L_p\mathrm T^{-5/2},\qquad
 C_j\le4\mathrm K_0e^{13/4}L_p\mathrm T^{-5/2}.
\]
The continuous production equals four times its quadratic form and the
jump production bounds four times its quadratic form. Hence the active
rate is at least \(\mathrm c_*\mathrm T^{5/2}/L_p\). Entropy decomposes
over the independent passive and active factors. Cauchy--Schwarz bounds
the energy of a marginal by the corresponding integrated energy, so the
full rate is at least the smaller of the two rates. The displayed active
lower bound is less than \(2\mathrm T\), proving the full-state inequality.
The form-domain extension uses the positive regularization and cutoffs
already used for \cref{eq:phil-selection-entropy-contraction}.

A one-receiver flip joins only the class pairs \(M\)--\(E\),
\(N\)--\(E\), and \(C\)--\(E\). The determinant bounds above make every
ratio of such adjacent weights at most \(1/\mathrm T\).
Writing the evidence as \(p_iZ_i^++(1-p_i)Z_i^-\), the ratio
\(Z_i^+/Z_i^-\) is between the extremal adjacent ratios. Therefore
\(|\partial_{p_i}\log Z|\le1/\mathrm T\), and
\(\|\mathfrak g\|_1\le4/\mathrm T\).
The row lower bound gives the rate floor \(\ell\) on the stated horizon.
The initial entropy is at most
\(13(12+|\mathsf m_0|^2)/(2\mathrm T)\): in the Gaussian density ratio,
\(\mathsf B_\sigma\preceq13I/\mathrm T\),
\(\mathsf Z\le1\), and the remaining quadratic term has nonpositive
expectation. Insert these bounds in
\cref{eq:phil-selection-tracking} and use Pinsker's inequality to obtain
\cref{eq:phil-temperature-tracking}.
For \(\varepsilon=\mathrm T^a\), its second term is
\(O(\mathrm T^{a-7/2}\log(1/\mathrm T))\), while
\(\ell t_{\mathrm L}\) grows as a positive constant times
\(\mathrm T^{5/2-a}\). The exponential therefore dominates the initial
\(\mathrm T^{-1/2}\) factor, proving the last assertion. \(\square\)
\section{Coherence and Persistence of the Selected Pairing}
\label{sec:phil-selected-coherence-proof}

\propositionheading{Coherence and persistence after pairing selection}{prop:phil-selected-pair-persistence}
In the four-cycle of \cref{conj:phil-contrast-formation}, let
\(0<\mathrm T\le1/100\), let \(M\) be one matching, and suppose
\([t_0,t_0+\mathsf h]\subseteq[t_{\mathrm L},2t_{\mathrm L}]\), with
\(\mathsf h>0\). Write
\(\nu_{\mathrm T}^M=\mathcal N(0,\mathrm T\mathsf A_M^{-1})\otimes\delta_M\).
Assume the learned and frozen recognition laws satisfy
\[
 \sup_{t_0\le t\le t_0+\mathsf h}
 \operatorname{TV}(q_t^{\mathrm L},\nu_{\mathrm T}^M)\le\delta,
 \qquad
 \operatorname{TV}(q_{t_0}^{0},\mu_{\mathrm T,1/2})\le\delta.
\]
Let \(\mathsf O_{\mathrm T,\mathsf h}^M\) be the stationary
Ornstein--Uhlenbeck path law with drift \(-\mathsf A_M\mathbf x\),
diffusion covariance \(2\mathrm TI\), and constant assignment \(M\).
For the configuration and assignment path laws on this interval,
\begin{align}
 \operatorname{TV}(\mathsf P^{\mathrm L}_{t_0,\mathsf h},
                    \mathsf O_{\mathrm T,\mathsf h}^M)
 &\le(\mathsf h+2)\delta,
 \label{eq:phil-selected-path-coupling}\\
 \operatorname{TV}\left(\mathsf P^{0}_{t_0,\mathsf h},
   \tfrac12\mathsf O_{\mathrm T,\mathsf h}^{M_0}
        +\tfrac12\mathsf O_{\mathrm T,\mathsf h}^{M_1}\right)
 &\le\delta+16\mathrm T+31\mathrm T\mathsf h.
 \label{eq:phil-frozen-path-mixture}
\end{align}
The frozen arm keeps its original uniform rows.

For resolution \(\mathsf a>0\), let
\(\mathscr E_M(\mathsf a,\mathsf h)\) require throughout the interval
that the assignment is \(M\), the stacked residual norm
\(\sqrt{(r^b)^2+(r^m)^2}\) on its two edges is at most
\(\mathsf a\), and that norm on each exterior edge exceeds
\(\mathsf a\).
Define
\[
 \mathsf b(\mathrm T,\mathsf a,\mathsf h)
 =8\sqrt{2/\pi}\mathsf a+
            (8+3400\mathsf h)\mathrm T/\mathsf a^2.
\]
Then
\begin{equation}
 \Pr_{\mathrm L}(\mathscr E_M^c)
 \le(\mathsf h+2)\delta+\mathsf b(\mathrm T,\mathsf a,\mathsf h).
 \label{eq:phil-selected-coherence-persistence}
\end{equation}
In particular, the choice
\(\mathsf a=\mathrm T^{1/4}\),
\(\mathsf h=\mathrm T^{-1/4}\), and \(\delta\le\mathrm T\)
makes this probability tend to zero.
The matching-window label, equal to \(M_j\) when
\(\mathscr E_{M_j}\) holds and to an additional symbol otherwise,
converges to \(\delta_M\) in the learned arm and to the equal matching
mixture in the frozen arm.

At one time, define \(\mathfrak P_{\mathsf a}(\mathbf x)\) as the
connected-component partition of edges whose stacked residual norm
\(\sqrt{(r^b)^2+(r^m)^2}\) is at most \(\mathsf a\).
For times when both
one-time hypotheses above hold, put
\[
 \mathsf e(\mathrm T,\mathsf a)
 =4\sqrt{2/\pi}\mathsf a+4\mathrm T/\mathsf a^2.
\]
The learned and frozen partition laws satisfy
\begin{equation}
 \left|\operatorname{TV}
   (q_t^{\mathrm L}\mathfrak P_{\mathsf a},
                 q_t^0\mathfrak P_{\mathsf a})-\tfrac12\right|
 \le2\delta+16\mathrm T+2\mathsf e(\mathrm T,\mathsf a).
 \label{eq:phil-coherent-partition-contrast}
\end{equation}
Thus the contrast persists for the coherent partition, not only the
source label.

These conclusions apply along \cref{cor:phil-joint-selection} with
tolerance \(\delta\le\mathrm T\), mobility cap
\(\epsilon_{\mathrm L}\le\mathrm T^2\), and \(t_0=t_{\mathrm L}\).
The matching can vary along that schedule. The internal restoring rates
are \(8+\mathrm T\) and \(4+5\mathrm T\), on an event whose probability
tends to one. The resolution therefore has restoring ratio at least
\(4\mathsf a^2/\mathrm T\to\infty\).
If the rows are also within \(\delta\) of their matching vertex,
the stationary law with those rows frozen realizes the same persistence
event with probability at least
\(1-4(\mathsf h+2)\delta-\mathsf b\).
This late coherence is restored equilibrium coherence; it supplies no
entrance or first-formation assertion.
\status{ESTABLISHED}

\noindent\emph{Proof.}
First retain the proposal mechanism in the comparison. Put
\(r(t)=q_t^{\mathrm L}(S\ne M)\), and let
\(b_{\ne M}(t,\mathbf x)\) be the probability of proposing a different
assignment at a clock ring. Since the clock has rate one,
\[
 r'(t)=\E b_{\ne M}(t,\mathbf X_t)-r(t).
\]
If \(N_{\ne M}\) counts proposals to an assignment other than \(M\)
on the interval, its compensator gives
\begin{equation}
 \E N_{\ne M}
 =r(t_0+\mathsf h)-r(t_0)+
                   \int_{t_0}^{t_0+\mathsf h}r(t)dt.
 \label{eq:phil-selected-proposal-occupation}
\end{equation}
Consequently
\(\Pr(\exists t\in[t_0,t_0+\mathsf h]:S_t\ne M)
 \le(\mathsf h+1)\delta\).
Couple the initial law maximally with \(\nu_{\mathrm T}^M\), and use
the same future Brownian motion for the native process and the
stationary matching process. Their paths agree until an initial
mismatch or a proposal to another assignment. The union bound and
\cref{eq:phil-selected-proposal-occupation} prove
\cref{eq:phil-selected-path-coupling}. This argument counts proposals
of the native process and does not condition its diffusion on an
assignment history.

For the frozen comparison, use the normalized weights
\(1,\mathsf r_E,\mathsf r_N,\mathsf r_C\) from
\cref{sec:phil-pairing-selection-proof}. The determinant bounds in the
proof of \cref{prop:phil-temperature-tracking} give
\[
 \mathsf r_E,\mathsf r_N\le\tfrac52\mathrm T,
 \qquad \mathsf r_C\le\tfrac32\mathrm T^{3/2}.
\]
The invariant mass \(\kappa_{\mathrm T}\) of nonmatching assignments
is therefore at most \(16\mathrm T\). Conditional on either matching,
its configuration law is the Gaussian factor of \(\nu_{\mathrm T}^M\).
The invariant law is consequently within \(\kappa_{\mathrm T}\) of
the equal mixture of these two laws.

Detailed balance bounds the stationary flux from \(M\) to all
nonmatching assignments by \(\kappa_{\mathrm T}\).
For the other matching \(M'\), write \(\mu_M(\mathbf x)\) for the
joint invariant density on the matching sheet. Pointwise,
\[
 \mu_M b_{M'}
 =\frac{\mu_M\mu_{M'}}{\sum_S\mu_S}
 \le\tfrac12\sqrt{\mu_M\mu_{M'}}.
\]
Since \((\mathsf A_M+\mathsf A_{M'})/2\) is the oriented-cycle
matrix, Gaussian integration gives
\(\int\sqrt{\mu_M\mu_{M'}}=\mu_{\mathrm T,1/2}(M)\mathsf r_C\).
Dividing the two flux bounds by the matching mass yields
\[
 \E_{\nu_{\mathrm T}^M}b_{\ne M}
 \le8\mathsf r_E+4\mathsf r_N+\tfrac52\mathsf r_C
 \le31\mathrm T.
\]
Thus the stationary flux out of the two matching sheets together is
at most \(31\mathrm T\). Couple the stationary native process at
its initial time with the equal matching mixture and use common
Brownian increments until the first departure from its coupled
matching. Its failure probability is at most
\(16\mathrm T+31\mathrm T\mathsf h\). The frozen native process
started at \(q_{t_0}^0\) differs from its stationary path law by at
most its initial total variation \(\delta\), because the frozen
transition kernel is the same. This proves
\cref{eq:phil-frozen-path-mixture}.

It remains to bound residuals on the comparison path.
The four normal coordinates
\(z_g^\pm=n_g^\pm\cdot\mathbf X\), \(g\in M\), from
\cref{eq:rg-formation-cycle-normals} have rates
\(8+\mathrm T\) and \(4+5\mathrm T\). Their stationary variances
are at most \(\mathrm T/4\), and
\[
 r_g^b=\sqrt2z_g^++z_g^-,\qquad
 r_g^m=\sqrt2z_g^+-z_g^-.
\]
For a scalar Ornstein--Uhlenbeck noise starting at zero, integration
by parts in its stochastic convolution gives
\[
 Z_t=\sqrt{2\mathrm T}
 \left[W_t-\lambda\int_0^te^{-\lambda(t-u)}W_u du\right],
 \qquad
 \E\sup_{t\le\mathsf h}|Z_t|^2\le32\mathrm T\mathsf h.
\]
The last inequality uses
\(\sup|Z_t|\le2\sqrt{2\mathrm T}\sup|W_t|\) and the Brownian
maximal estimate proved in the matching-episode argument of
\cref{sec:rg-formation-cycle-proof}.
For the four-dimensional stationary normal vector \(z\), it follows
that
\[
 \E\sup_{t\le\mathsf h}|z_t|^2
 \le2\mathrm T+256\mathrm T\mathsf h.
\]
Each selected stacked residual norm is at most \(2|z_t|\), since
\((r_g^b)^2+(r_g^m)^2=4(z_g^+)^2+2(z_g^-)^2\).
Thus the probability that any internal stacked norm exceeds
\(\mathsf a\) at any time is at most
\((8+1024\mathsf h)\mathrm T/\mathsf a^2\).

Every exterior link coordinate is absent from \(H_M\), and its
block in \(\mathsf A_M\) is the scalar \(\mathrm T\).
It is therefore an independent \(\mathcal N(0,1)\) coordinate in
the stationary Gaussian. Each exterior belief or model residual
has variance at least one. The probability that any of these four
residuals initially has absolute value at most \(2\mathsf a\)
is at most \(8\sqrt{2/\pi}\mathsf a\).

To control their increments, diagonalize the symmetric positive
matrix \(\mathsf A_M\). In a scalar eigencoordinate with rate
\(\lambda>0\), stationarity and the preceding noise estimate imply
\[
 \E\sup_{t\le\mathsf h}|X_t-X_0|^2
 \le\frac{2\mathrm T}{\lambda}(1-e^{-\lambda\mathsf h})^2
                       +64\mathrm T\mathsf h
 \le66\mathrm T\mathsf h,
\]
since \((1-e^{-v})^2\le v\) for \(v\ge0\).
Summing twelve coordinates gives
\begin{equation}
 \E\sup_{t\le\mathsf h}|\mathbf X_t-\mathbf X_0|^2
 \le792\mathrm T\mathsf h.
 \label{eq:phil-selected-ou-increments}
\end{equation}
Every residual covector has norm \(\sqrt3\). Except with
probability \(2376\mathrm T\mathsf h/\mathsf a^2\), all exterior
increments are therefore at most \(\mathsf a\) in absolute value.
The initial small-ball estimate and the two maximal estimates prove
\(\mathsf O_{\mathrm T,\mathsf h}^M(\mathscr E_M^c)\le\mathsf b\).
Together with \cref{eq:phil-selected-path-coupling}, this proves
\cref{eq:phil-selected-coherence-persistence}.
The frozen path comparison proves the stated matching-window label
limit whenever \(\mathsf b\to0\),
\(\mathsf h\delta\to0\), and \(\mathrm T\mathsf h\to0\), in
particular for the displayed choices.

For a single time no increment estimate is needed. The normal second
moment is at most \(\mathrm T\), so the internal failure probability
is at most \(4\mathrm T/\mathsf a^2\). The exterior small-ball
bound at radius \(\mathsf a\) is
\(4\sqrt{2/\pi}\mathsf a\). Thus the partition induced by
\(\nu_{\mathrm T}^{M_j}\) differs from the matching partition
\(\mathfrak P_j\) with probability at most \(\mathsf e\).
The learned partition law is within \(\delta+\mathsf e\) of
\(\delta_{\mathfrak P_M}\); the frozen one is within
\(\delta+16\mathrm T+\mathsf e\) of
\((\delta_{\mathfrak P_0}+\delta_{\mathfrak P_1})/2\).
These matching partitions are distinct, so the triangle inequality
for total variation proves \cref{eq:phil-coherent-partition-contrast}.

The schedule in \cref{cor:phil-joint-selection} permits tolerance
\(\mathrm T\) and mobility cap \(\mathrm T^2\).
Then \(\mathrm T^{-1/4}\le t_{\mathrm L}\) for sufficiently small
\(\mathrm T\), so the following persistence window lies within the
selection window. All errors just bounded tend to zero.
The normal rates establish the restoring ratio on the native event.
Finally, if \(\max_i|p_i-p_i^M|\le\delta\),
\(\operatorname{TV}(\mu_{\mathrm T,p},\nu_{\mathrm T}^M)
 =1-\mu_{\mathrm T,p}(M)\le4\delta\).
Applying the same occupation and coupling argument to this
stationary frozen-row process gives the stationary lower bound in
the statement. The persistence event is therefore already typical
at equilibrium for those rows and carries no first-entrance claim.
\hfill\(\square\)



\subsection{The first contrast on the original episode window}

\propositionheading{A first-block null under slower source learning}{prop:phil-early-block-null}
In the model of \cref{conj:phil-contrast-formation}, let
\(0<\mathrm T\le10^{-16}\), and let the initial law be
\(\mathcal N(\mathsf m_0,I_{12})\) times uniform assignments.
The learned and frozen arms have the same initial mean and differ only in
the source-learning mobility \(\epsilon_{\mathrm L}\).
Set
\[
 H_{\mathrm T}=2\mathrm T^{-1/8}.
\]
Keep the contrast event of \cref{prop:phil-seeded-contrast-formation}
with its regime parameter equal to \(\mathrm T^{1/8}\): an eligible
proper connected block has an internal belief relative entropy exceeding
\(1/4\) initially, and for a duration-one interval inside
\([0,H_{\mathrm T}]\) all its internal belief and model relative
entropies are at most \(\sqrt{\mathrm T}\), while its boundary belief
relative entropies exceed \(1/2\).
For each block take the infimum of its qualifying completion times,
with value infinity when there is no such interval. Let \(K_{\mathrm T}\)
be the block with least such infimum, using one fixed tie rule in both
arms, and give it a cemetery value when no block qualifies. Write
\(E_{\mathrm T}\) for the event that its value is not the cemetery value.
The infimum need not itself be a qualifying completion time.

Define the positive, temperature-independent constant
\[
 \mathsf c(\mathsf m_0)
 =\frac{e^{-1/4}\varrho_\circ^{12}}{737280}
   \exp\left[-\frac12
       (|\mathbf x^\circ-\mathsf m_0|+\varrho_\circ)^2\right],
 \qquad
 \mathsf d_{\mathrm T}
 =\min\{1,8\epsilon_{\mathrm L}\mathrm T^{-1/4}\}.
\]
Then
\[
 \Pr_0(E_{\mathrm T})\ge\mathsf c(\mathsf m_0),
 \qquad
 \Pr_{\epsilon_{\mathrm L}}(E_{\mathrm T})
       \ge\mathsf c(\mathsf m_0)-\mathsf d_{\mathrm T}.
\]
If \(\mathsf d_{\mathrm T}<\mathsf c(\mathsf m_0)\), the two
conditional block laws satisfy
\[
 \operatorname{TV}\left(
 \operatorname{Law}_{\epsilon_{\mathrm L}}(K_{\mathrm T}\mid E_{\mathrm T}),
 \operatorname{Law}_0(K_{\mathrm T}\mid E_{\mathrm T})\right)
 \le\frac{\mathsf d_{\mathrm T}}{\mathsf c(\mathsf m_0)}.
\]
In particular their difference vanishes for every fixed Gaussian mean
whenever
\(\mathrm T\downarrow0\) and
\(\epsilon_{\mathrm L}\mathrm T^{-1/4}\to0\).
This includes the caps
\(\epsilon_{\mathrm L}\le\mathrm T_n^a\), \(a>7/2\), in
\cref{cor:phil-joint-selection}.
\status{ESTABLISHED}

\noindent\emph{Proof.}
First prove a uniform comparison of the complete native path laws on
any finite horizon \(H\). The learned rows are deterministic and satisfy
\[
 |p_i(t)-1/2|\le\epsilon_{\mathrm L}t/2.
\]
At a common configuration, write \(A,B>0\) for a receiver's likelihood
factors for its two sources. Its posterior positive-source probabilities
in the two arms are
\[
 b(p)=\frac{pA}{pA+(1-p)B},
 \qquad b(1/2)=\frac A{A+B}.
\]
Direct subtraction gives
\[
 |b(p)-b(1/2)|
 =\frac{AB|2p-1|}{(pA+(1-p)B)(A+B)}
 \le|2p-1|.
\]
Indeed the denominator is
\(pA^2+AB+(1-p)B^2\ge AB\). This bound is independent of
temperature and the likelihood ratio, without an interior-row restriction.

Use identical initial samples, a common Brownian motion, and the same
rate-one proposal clock in the two arms. Until the first different redraw,
their configurations and assignments coincide: a held-sheet drift and its
noise do not depend on the prior rows. At each ring, couple each receiver's
two Bernoulli laws maximally. The union bound and the row estimate bound
the probability of any different receiver choice at time \(t\) by
\[
 \sum_i|b_i(p_i(t))-b_i(1/2)|
 \le4\epsilon_{\mathrm L}t.
\]
The compensator of the rate-one clock, stopped at the first mismatch,
bounds its probability before \(H\) by
\(\int_0^H4\epsilon_{\mathrm L}t dt=2\epsilon_{\mathrm L}H^2\).
Consequently,
\begin{equation}
 \operatorname{TV}\left(
 \operatorname{Law}_{\epsilon_{\mathrm L}}(\mathbf X_{[0,H]},S_{[0,H]}),
 \operatorname{Law}_0(\mathbf X_{[0,H]},S_{[0,H]})\right)
 \le\min\{1,2\epsilon_{\mathrm L}H^2\}.
\label{eq:phil-early-path-coupling}
\end{equation}
The bound follows because the complete paths are identical unless the
coupling has had a mismatch. Taking \(H=H_{\mathrm T}\) gives
\(\mathsf d_{\mathrm T}\).

The first-block mark is a measurable function of these paths. To see this
directly, for a fixed block let \(A_I(t)\) be the maximum of its internal
belief and model relative entropies on \([t-1,t]\), and let \(B_I(t)\)
be the minimum boundary belief relative entropy there. Configuration
paths are continuous, so these two window functions are continuous.
Replace the strict condition \(B_I(t)>1/2\) by the countable union of
conditions \(B_I(t)\ge1/2+1/k\), \(k\ge1\). For each \(k\) the
qualifying end times form a closed subset of the compact interval
\([1,H_{\mathrm T}]\). Its first time is measurable, as may be checked
by minimizing
\(\max(A_I(t)-\sqrt{\mathrm T},1/2+1/k-B_I(t))\) on each compact
initial subinterval. The infimum over \(k\), the initial eligibility
condition, and the finite comparison over blocks preserve measurability.
This construction implements the strict boundary convention without
asserting that its infimum is attained.

A measurable map of paths does not increase total variation. Let
\(a_I,b_I\) be the probabilities of the non-cemetery marks in the learned
and frozen arms, \(p=\sum_Ia_I\), \(q=\sum_Ib_I\), and \(d\) the
total-variation distance of the two cemetery-augmented mark laws. Then
\[
 2d=\sum_I|a_I-b_I|+|p-q|,
 \qquad d\le\mathsf d_{\mathrm T}.
\]
If \(p\ge q>0\), the triangle inequality gives
\[
 \frac12\sum_I\left|\frac{a_I}{p}-\frac{b_I}{q}\right|
 \le\frac{\sum_I|a_I-b_I|+(p-q)}{2p}
 =\frac d p.
\]
Interchanging the arms proves the bound \(d/\max(p,q)\).
The frozen matching-seed sub-event in
\cref{prop:phil-seeded-contrast-formation} belongs to \(E_{\mathrm T}\)
and has probability at least \(\mathsf c(\mathsf m_0)\).
Thus \(q\ge\mathsf c(\mathsf m_0)\) and
\(p\ge q-d\ge\mathsf c(\mathsf m_0)-\mathsf d_{\mathrm T}\).
The claimed conditional bound follows, and its limiting conclusion is
immediate. \(\square\)

The proposition keeps the original contrast event and uses its seeded
sub-event only as a probability lower bound. That sub-event already has
the restoring interaction and stationary comparison of
\cref{prop:rg-formation-endogenous-cycle}; the proof does not certify
every history in the broader contrast event under the coherence-episode
rule. Any common measurable admissibility restriction retaining the
seeded sub-event has the same comparison.

The result does not cover the original coupled choice
\(\epsilon_{\mathrm L}=\mathrm T^{1/8}\), for which its bound is
vacuous. It also does not compare different initial means or establish a
conditional parent error. On the slower schedules, late assignment
selection can coexist with a vanishing difference between the conditional
laws of the first contrast on \([0,H_{\mathrm T}]\). These are
different observables on different horizons.



\section{Activation from the Actual Conditioned Recognition Law}
\label{sec:phil-actual-conditioned-parent}

\propositionheading{Conditional information after a completed coherence prefix}{prop:phil-actual-conditioned-parent}
In the learned four-cycle, fix \(0<\mathrm T\le1/100\), a selected
matching \(M\), and \(s=t_{\mathrm L}\ge1\). Suppose, throughout
\([s,2s]\),
\[
 \KL(q_t\Vert\mu_{\mathrm T,p(t)})\le\eta,
 \qquad
 \operatorname{TV}(q_t,\nu_{\mathrm T}^M)\le\delta.
\]
Set \(\mathsf a=\mathrm T^{1/4}\), \(\theta=s+1\), and let
\(E=\mathscr E_M(\mathsf a,1)\) be the actual matching and coherence
event of \cref{prop:phil-selected-pair-persistence} on \([s,s+1]\).
Assume
\[
 \rho=3\delta+\mathsf b(\mathrm T,\mathsf a,1)<\tfrac12,
 \qquad
 h_b(r)=-r\log r-(1-r)\log(1-r),
 \qquad
 \kappa_E=\frac{\eta+h_b(\rho)}{1-\rho}.
\]
For all \(t\ge\theta\), write
\(q_t^E=\operatorname{Law}(\mathbf X_t,S_t\mid E)\), always
conditioning on this same completed prefix. Keep the single marked
channel \(C\) of \cref{prop:rg-formation-endogenous-cycle}(b), including
its identity action off the matching strata and its identity extension
on evaluator samples. The parent activation law is exactly
\[
 \mu_{\mathrm{act}}=C_\#q_\theta^E.
\]

Let \(\mathcal H_{\mathrm{opt}}(q)\) be the mean discarded conditional
relative entropy of \(q\) against the restricted optimizer under \(C\).
Let \(\mathcal H_{\mathrm{phys}}(q)\) be the corresponding mean
discarded conditional relative entropy against the full physical
posterior, after lifting \(q\) with its recognition evaluator and
retaining every evaluator sample. Then, for every \(t\in[\theta,2s]\),
\begin{align}
 \KL(q_t^E\Vert\mu_{\mathrm T,p(t)})&\le\kappa_E,
 \label{eq:phil-conditioned-full-entropy}\\
 \mathcal H_{\mathrm{opt}}(q_t^E)&\le\kappa_E,
 \label{eq:phil-conditioned-opt-entropy}\\
 \mathcal H_{\mathrm{phys}}(q_t^E)
 &\le\kappa_E+\log\frac{4+5\mathrm T}{4+\mathrm T}
 \le\kappa_E+\mathrm T.
 \label{eq:phil-conditioned-physical-entropy}
\end{align}
The two hidden-conditional bounds hold with any strictly positive
alternative source rows in their reference laws; the full-law bound
\cref{eq:phil-conditioned-full-entropy} refers to the native rows
\(p(t)\).

These are unscaled, dimensionless relative entropies. The corresponding
scaled optimizer free-energy loss is at most \(\mathrm T\kappa_E\);
the unscaled physical negative-ELBO loss is at most
\(\kappa_E+\mathrm T\), and its \(\mathrm T\)-scaled value is at most
\(\mathrm T\kappa_E+\mathrm T^2\).
Along the joint selection schedule with \(\eta,\delta\le\mathrm T\),
\[
 \rho=O(\mathrm T^{1/4}),
 \qquad
 \kappa_E=O(\mathrm T^{1/4}\log(1/\mathrm T)).
\]
The assumption \(\rho<1/2\) then holds for all sufficiently small
\(\mathrm T\), so these bounds are asymptotic.
\status{ESTABLISHED}

\noindent\emph{Proof.}
The learned-arm persistence bound
\cref{eq:phil-selected-coherence-persistence}, applied with
\(t_0=s\) and \(\mathsf h=1\), gives \(r:=\Pr(E^c)\le\rho\); its
proof uses only the learned-arm hypothesis. Thus \(E\) is a
\(\theta\)-measurable event of probability at least \(1-\rho\).
The law \(q_\theta^E\) is the actual conditional law and may have
retained--hidden correlations and a non-Gaussian hidden conditional.
Its marked pushforward is normalized by construction. At later times
no further event is conditioned upon. In particular, the deterministic
native row path is not recomputed from \(q_t^E\).

For a fixed \(t\in[\theta,2s]\), put
\(B=\mathbf1_E\) and \(R_t=(\mathbf X_t,S_t)\).
Compute the relative entropy of the joint law of \((B,R_t)\) against
\(\mathrm{Bernoulli}(1-r)\otimes\mu_{\mathrm T,p(t)}\) in two orders:
\begin{equation}
 \begin{aligned}
 &(1-r)\KL(q_t^E\Vert\mu_{\mathrm T,p(t)})
       +r\KL(q_t^{E^c}\Vert\mu_{\mathrm T,p(t)})\\
 &\qquad=\KL(q_t\Vert\mu_{\mathrm T,p(t)})
                       +\mathbb I(B;R_t)
 \le\eta+h_b(r).
 \end{aligned}
 \label{eq:phil-conditioned-event-information}
\end{equation}
When \(r=0\), use just the first term. The second term is nonnegative,
\(h_b\) is increasing on \([0,1/2]\), and \(1-r\ge1-\rho\).
This proves \cref{eq:phil-conditioned-full-entropy}.
It also establishes the needed absolute continuity and finiteness;
directly, \(q_t^E\le q_t/(1-r)\).

The channel's fibers can be checked explicitly. On a matching with
edges \(e,f\), let \(F_M\) have the orthonormal columns
\[
 (n_e^+-n_f^+)/\sqrt2,\qquad n_e^-,\qquad n_f^-.
\]
Write \(Z_M=\operatorname{im}F_M\), \(W_M=Z_M^\perp\), and
\(\mathbf x=y+F_Mz\).
The normal identities in \cref{eq:rg-formation-cycle-normals} give
\[
 \nu_{\mathrm T,Z}=\mathcal N(0,\mathrm TK_Z^{-1}),
 \qquad
 K_Z=\operatorname{diag}(8+\mathrm T,4+5\mathrm T,4+5\mathrm T)
\]
as the hidden optimizer conditional. Both the interaction Hessian
and the evaluator precision preserve \(W_M\) and \(Z_M\).
Consequently this conditional is independent of the retained
coordinate and of the prior rows. It is a reference conditional, not
an assertion that the actual conditioned recognition law has that
form. The same statement holds on the other matching by rotation;
off the two matchings the channel is the identity and has no
discarded conditional.

The target of \(C\) is a finite disjoint union of Euclidean strata,
so disintegration gives the exact entropy chain rule
\begin{equation}
 \KL(q_t^E\Vert\mu_{\mathrm T,p(t)})
 =\KL(C_\#q_t^E\Vert C_\#\mu_{\mathrm T,p(t)})
                         +\mathcal H_{\mathrm{opt}}(q_t^E).
 \label{eq:phil-conditioned-parent-chain}
\end{equation}
Dropping its retained term proves
\cref{eq:phil-conditioned-opt-entropy}. This argument uses the
complete actual law; it neither resets the hidden state nor assumes
Gaussianity conditional on a source history.

For the physical bound, work first on one matching fiber. Its physical
outer hidden Gaussian is
\[
 \Pi_{\mathrm T,Z}=\mathcal N(0,\mathrm T(K_Z^{\mathrm{phys}})^{-1}),
 \qquad
 K_Z^{\mathrm{phys}}
  =\operatorname{diag}(8+\mathrm T,4+\mathrm T,4+\mathrm T).
\]
All displayed matrices are positive definite. Write
\(v=1-\sigma_{\mathrm{ev}}^2=1/2\),
\(d_y=\mathfrak Dy\), and \(d_z=\mathfrak DF_Mz\).
Since \(R_{\mathrm{ev}}=\mathfrak D^\top\mathfrak D/v\)
preserves the split, \(d_y\cdot d_z=0\); its hidden eigenvalues
are \(0,4,4\).

Let \(U=(m,k)\) denote the full evaluator sample vectors.
At configuration \(\mathbf x\), the physical evaluator has
\(m\sim\mathcal N(\bar m,\sigma_{\mathrm{ev}}^2I)\) and independent
increment \(k-m\sim\mathcal N(0,vI)\). Its recognition counterpart,
\cref{eq:rg-formation-evaluator-lift}, has the same first law and
increment mean \(d_y+d_z\). At a fixed retained \(y\), form the
comparison probability law
\begin{equation}
 \overline P_y(dz,dU)
 =\Pi_{\mathrm T,Z}(dz)P_{\mathrm{ev}}(dU\mid y+F_Mz)
  \exp\left[
  \frac{d_y\cdot(k-m)-|d_y|^2/2}{v}\right].
 \label{eq:phil-conditioned-evaluator-tilt}
\end{equation}
Completing the Gaussian square in \(k-m\) proves that the multiplier
integrates to one for every \(z\). It changes the increment mean
to \(d_y\) and keeps \(\Pi_{\mathrm T,Z}\) as its \(z\) marginal.
The multiplier is retained-measurable in \((y,U)\), so
\(\overline P_y\) has the physical hidden conditional given these
retained variables. At a fixed \(z\), the recognition evaluator
differs by the additional increment mean \(d_z\), with relative
entropy \(|d_z|^2/(2v)\).

Compute the joint relative entropy first through \(z\), then through
\(U\), and drop the nonnegative evaluator-marginal term. For the
actual hidden law \(q(dz\mid M,y)\), this gives
\[
 \mathcal H_{\mathrm{phys}}(q\mid M,y)
 \le\KL(q(dz\mid M,y)\Vert\Pi_{\mathrm T,Z})
             +\E_q[|d_z|^2/(2v)\mid M,y].
\]
All terms are well-defined: the native process has finite second
moments, conditioning increases the second-moment bound by a factor
at most \((1-\rho)^{-1}\), and the Gaussian reference logarithms are quadratic.
The optimizer conditional is the exact tilt
\[
 \nu_{\mathrm T,Z}(dz)
 =\frac{\Pi_{\mathrm T,Z}(dz)e^{-|d_z|^2/(2v)}}{Z_{\mathrm{tilt}}},
 \qquad
 Z_{\mathrm{tilt}}=\frac{4+\mathrm T}{4+5\mathrm T}.
\]
For its normalizer, the common coordinate is unchanged and each of
the two discrepancy coordinates changes rate from \(4+\mathrm T\)
to \(4+5\mathrm T\); their two square-root Gaussian determinant
factors multiply to \(Z_{\mathrm{tilt}}\). Therefore
\[
 \KL(q(dz\mid M,y)\Vert\Pi_{\mathrm T,Z})
       +\E_q[|d_z|^2/(2v)\mid M,y]
 =\KL(q(dz\mid M,y)\Vert\nu_{\mathrm T,Z})
                         +\log\frac{4+5\mathrm T}{4+\mathrm T}.
\]
Integrate over the retained law. Matching strata have total
probability at most one, and all other strata have zero discarded
loss. Finally,
\[
 0\le\log\frac{4+5\mathrm T}{4+\mathrm T}
 =\log\left(1+\frac{4\mathrm T}{4+\mathrm T}\right)
 \le\frac{4\mathrm T}{4+\mathrm T}\le\mathrm T.
\]
This proves \cref{eq:phil-conditioned-physical-entropy}.
Only recognition has been conditioned on \(E\); the physical target
still uses the original observed records, with all evaluator samples
retained.

A source-prior factor \(\pi_S\) cancels in any conditional law on a
fiber because \(S\) is retained. The hidden optimizer and physical
reference conditionals are consequently the same for every strictly
positive choice of prior rows. The full-law entropy need not remain
small for different rows, because its retained term may change.
An imposed zero prior on an occupied stratum would remove reference
support and is excluded; ordinary finite-time learning from positive
rows preserves positivity.

The exact free-energy chain rules give the stated scaled and unscaled
losses. With \(\eta,\delta\le\mathrm T\) and
\(\mathsf a=\mathrm T^{1/4}\),
\[
 \rho\le3\mathrm T+8\sqrt{2/\pi}\mathrm T^{1/4}
                         +3408\mathrm T^{1/2}.
\]
Use \(h_b(\rho)\le\rho\log(e/\rho)\) and
\(1-\rho\ge1/2\) eventually to obtain the displayed order for
\(\kappa_E\). \(\square\)

\remarkheading{A high-probability prefix is needed}{rem:phil-conditioned-hidden-sign}
The high probability of the declared prefix has mathematical content.
Under the exact matching Gaussian \(\nu_{\mathrm T}^M\), a hidden
coordinate \(z_1\) is independent of the retained state and has zero
unconditioned hidden relative entropy. Conditioning on
\(\{z_1>0\}\), of probability \(1/2\), instead gives a density ratio
two on that half-space and hidden relative entropy \(\log2\) at every
temperature. Thus a probability merely bounded below does not suffice
for a vanishing conditional loss. Conditioning on the entire realized
fine history is different again: it fixes the hidden endpoint and
gives a Dirac conditional, singular relative to the nondegenerate
Gaussian fiber. The activation here uses the full joint law conditional
on \(E\), not a realized fine history.
\status{ESTABLISHED}

The proposition supplies conditional information at activation and for
the native future marginals under that same completed prefix. It does
not yet compare an independently evolving parent path. Its prefix is
the late restored coherence event, and carries no first-formation or
recursive-coalescence conclusion.



\subsection*{Inherited dynamics from the conditioned activation law}

\propositionheading{Dynamics of the conditioned parent}{prop:phil-actual-parent-dynamics}
Use the selected matching, channel, and conditioned activation of
\cref{prop:phil-actual-conditioned-parent}. More explicitly, let
\(0<\mathrm T\le1/100\), \(s=t_{\mathrm L}\), \(\theta=s+1\), and let \(\mathscr E\)
be the selected coherence event on \([s,\theta]\) at resolution
\(\mathsf a\) from \cref{prop:phil-selected-pair-persistence}.
Assume \(\Pr(\mathscr E^c)\le\rho<1\), and, throughout
\([s,\theta+\mathsf h]\subseteq[t_{\mathrm L},2t_{\mathrm L}]\),
\[
 \operatorname{TV}(q_t,\nu_{\mathrm T}^M)\le\delta,
 \qquad \max_i|p_i(t)-p_i^M|\le\delta.
\]
Write \(\varepsilon=\epsilon_{\mathrm L}\),
\(d=\delta+\varepsilon\mathsf h<1/4\), and use relative future time
\(u=t-\theta\). Initialize the parent with
\[
 q_0^C=C_\#\operatorname{Law}(\mathbf X_\theta,S_\theta\mid\mathscr E),
 \qquad p^C(0)=p(\theta).
\]
On the stratum \(\{S\}\times W_S\), its continuous generator is
\(-\mathsf A_S y\cdot\nabla_{W_S}+\mathrm T\Delta_{W_S}\).
Its rate-one proposal kernel is
\begin{equation}
 J_p((S,y),\cdot)=
 \int\kappa_S(d\mathbf x\mid y)
     \sum_{S'}b_p(S'\mid\mathbf x)
                   \delta_{(S',P_{W_{S'}}\mathbf x)},
 \label{eq:phil-actual-parent-kernel}
\end{equation}
where \(b_p(S'\mid\mathbf x)=\mu_{\mathrm T,p}(S'\mid\mathbf x)\) is the
product of receiver posterior rows at rows \(p\), and
\(\kappa_S\) is the optimizing hidden conditional: the three
independent Gaussian coordinates on a matching stratum, and a point
mass off the matchings. The parent rows evolve from its own law:
\begin{equation}
 \frac d{du}p_i^C(u)
 =\varepsilon\bigl(q_u^C(\sigma_i=+1)-p_i^C(u)\bigr).
 \label{eq:phil-actual-parent-learning}
\end{equation}
This coupled parent system has a unique global law with interior rows.

Let \(\mathsf P_{\theta,\mathsf h}^{\mathrm{fine},C}\) be the law of
\((C(\mathbf X_{\theta+u},S_{\theta+u}))_{0\le u\le\mathsf h}\)
conditioned only on \(\mathscr E\), with the fine process continuing
its original row path. Let \(\mathsf P_{\mathsf h}^C\) be the parent
path law from the displayed activation and its own learning equation.
Then
\begin{equation}
 \operatorname{TV}
 (\mathsf P_{\theta,\mathsf h}^{\mathrm{fine},C},
                                      \mathsf P_{\mathsf h}^C)
 \le2\rho+(\mathsf h+3)\delta+
                         \frac{4\mathsf h d}{1-4d}.
 \label{eq:phil-actual-parent-path-error}
\end{equation}
The comparison includes all retained coordinates and the assignment
history, hence also the loop history. Self-proposals remain part of
the rate-one kernel.

Along \cref{cor:phil-joint-selection}, choose
\(\delta\le\mathrm T\), \(\varepsilon\le\mathrm T^2\),
\(\mathsf a=\mathrm T^{1/4}\), and
\(\mathsf h=\mathrm T^{-1/4}\). One may take
\(\rho=3\delta+\mathsf b(\mathrm T,\mathsf a,1)\), where
\(\mathsf b\) is from \cref{prop:phil-selected-pair-persistence}.
The right side of \cref{eq:phil-actual-parent-path-error} is then
\(O(\mathrm T^{1/4})\). This is a conditional, finite-window
approximation after the selected coherence event; the activation law
is not replaced by a Gaussian law.
\status{ESTABLISHED}

\noindent\emph{Proof.}
For a matching \(M\), the decomposition
\(\R^{12}=W_M\oplus Z_M\) of
\cref{prop:rg-formation-endogenous-cycle} is preserved by
\(H_M\), \(\Omega\), and \(\mathsf A_M\).
The hidden drift matrix is
\(\operatorname{diag}(8+\mathrm T,4+5\mathrm T,4+5\mathrm T)\).
The prior weight \(\pi_S(p)\) cancels from the conditional law at
fixed \((S,y)\), so \(\kappa_S\) is independent of \(p\).
Off the matchings no coordinate is discarded. Thus the displayed
continuous generator is the projection of every held-sheet diffusion,
and \cref{eq:phil-actual-parent-kernel} is a probability kernel.
A self-proposal has \(S'=S\) and
\(P_{W_S}\mathbf x=y\), so it leaves the parent state unchanged.
At a different assignment it uses that assignment's projection,
including its twelve-dimensional stratum when appropriate.

For fixed rows, the pushed optimizing law \(C_\#\mu_{\mathrm T,p}\)
is invariant. Indeed each diffusion is reversible on its Gaussian
stratum, and
\(\mu_{\mathrm T,p}(d\mathbf x,S)b_p(S'\mid\mathbf x)\)
is symmetric in \(S,S'\). Pushing both endpoints through \(C\)
gives detailed balance for \(J_p\).
The loop covector annihilates each \(Z_M\), so this construction
retains the loop through the stratum changes as well.

We next check the existence assertion with evolving parent rows.
For any continuous interior row path, concatenate the linear
diffusions with the rate-one proposals, drawing \(\kappa_S\) at
each proposal. On every finite interval the clock has finitely many
rings and every inserted Gaussian is finite almost surely, so this
construction is nonexplosive. Its second moments remain finite:
projection cannot increase the norm, an inserted hidden draw adds
expected squared norm at most \(5\mathrm T/8\), and the diffusion
generator applied to \(|y|^2\) is at most
\(24\mathrm T-2\mathrm T|y|^2\). The full generator is consequently
bounded above by \(25\mathrm T-2\mathrm T|y|^2\).
The activation law has finite second moment because it conditions a
finite-moment fine law on an event of positive probability.

If two row paths have entries in \([\mathrm c_0,1-\mathrm c_0]\), couple their
parent processes by the same initial state, Brownian increments,
clocks, hidden Gaussian draws, and receiver uniforms until their
first different proposal. The derivative of a posterior receiver
probability with respect to its prior entry is at most
\(1/[4\mathrm c_0(1-\mathrm c_0)]\). Therefore
\[
 \operatorname{TV}(q_u^{C,p},q_u^{C,\widetilde p})
 \le\frac1{\mathrm c_0(1-\mathrm c_0)}
       \int_0^u\|p(v)-\widetilde p(v)\|_\infty dv.
\]
The variation-of-constants map for
\cref{eq:phil-actual-parent-learning} consequently has Lipschitz
constant at most \(\varepsilon u_*^2/[2\mathrm c_0(1-\mathrm c_0)]\) on an interval
of length \(u_*\). On a short enough interval it is a contraction
on continuous row paths in an interior cube. Its image stays in
that cube after further shortening the interval, since the row
velocity has absolute value at most \(\varepsilon\).
Moreover its solution satisfies
\[
 p_i^C(u)\ge p_i^C(0)e^{-\varepsilon u},\qquad
 1-p_i^C(u)\ge(1-p_i^C(0))e^{-\varepsilon u}.
\]
On any prescribed finite horizon these lower bounds give a common
interior cube and a positive contraction step length. Successive
steps give existence and uniqueness on that horizon, and then on
all finite horizons. The source marginal is continuous because its
jump rate is bounded, so the resulting rows are continuously
differentiable. This constructs learning from the parent's own law,
without imposing the fine row path after activation.

Let \(\overline\nu_{\mathrm T}^M\) be the retained Gaussian on
\(W_M\). Write \(\mathsf O^C_{\mathrm T,\mathsf h}\) for its
stationary held-\(M\) path law. The projection of
\(\mathsf O_{\mathrm T,\mathsf h}^M\) from
\cref{prop:phil-selected-pair-persistence} is exactly this law.
Only the learned-arm conclusions of that proposition are used, and
their proof does not use its frozen-arm hypothesis.
Conditioning any probability law on an event of probability
\(1-r\) changes its total variation by at most \(r\) after
any measurable projection. Thus the conditioned native retained
future obeys
\begin{equation}
 \operatorname{TV}
 (\mathsf P_{\theta,\mathsf h}^{\mathrm{fine},C},
                     \mathsf O^C_{\mathrm T,\mathsf h})
 \le\rho+(\mathsf h+2)\delta.
 \label{eq:phil-native-conditioned-held}
\end{equation}
This applies to a prefix event through \(\theta\). No condition is
placed on future coherence or future proposals. The same argument
at the activation time gives
\(\operatorname{TV}(q_0^C,
 \overline\nu_{\mathrm T}^M\otimes\delta_M)\le\rho+\delta\).

For the parent's own rows let \(\zeta_i^C\) be the probability in
row \(i\) assigned to the nonpartner of \(M\). Its equation and the
bound on any probability give
\[
 0\le\zeta_i^C(u)\le\zeta_i^C(0)+\varepsilon u
                  \le\delta+\varepsilon u\le d.
\]
The row-independent normalized sheet weights are at most the matching
weight one for \(0<\mathrm T\le1/100\) by
\cref{prop:phil-temperature-tracking}; hence, for any fixed row vector
with this bound,
\[
 \mu_{\mathrm T,p}(M)\ge\pi_M(p)
               \ge1-\sum_i\zeta_i\ge1-4d.
\]
Define the parent departure probability at a proposal by
\(k_p(y)=J_p((M,y),\{S'\ne M\})\).
Detailed balance, and the fact that the optimizing conditional
retained law on \(M\) is \(\overline\nu_{\mathrm T}^M\)
independently of the rows, give
\begin{equation}
 \int k_p(y)\overline\nu_{\mathrm T}^M(dy)
 \le\frac{1-\mu_{\mathrm T,p}(M)}{\mu_{\mathrm T,p}(M)}
 \le\frac{4d}{1-4d}.
 \label{eq:phil-parent-averaged-departure}
\end{equation}
The first inequality bounds stationary flux from \(M\) to its
complement by the complement's stationary mass.

Maximally couple the parent activation with the stationary retained
matching law and use common Brownian increments, clocks, and parent
proposal draws up to the first parent departure. Before departure,
the continuous paths agree because the matching drift splits, and
self-proposals do not move the retained coordinate. The row path
\(p^C(u)\) constructed above is deterministic, even though it was
computed from the parent's own law. Under the stationary comparison
path, its marginal at every time is
\(\overline\nu_{\mathrm T}^M\). Bound the first-departure
compensator with its stopped indicator by the compensator evaluated
on that whole stationary path. Applying
\cref{eq:phil-parent-averaged-departure} at each deterministic
\(p^C(u)\) gives
\[
 \Pr(\text{parent departure before }\mathsf h
                         \text{ while initially coupled})
 \le\frac{4\mathsf h d}{1-4d}.
\]
No Gaussian assertion is made about the parent's law conditioned on
survival or about the native hidden variables. We have proved
\[
 \operatorname{TV}(\mathsf P_{\mathsf h}^C,
                         \mathsf O^C_{\mathrm T,\mathsf h})
 \le\rho+\delta+\frac{4\mathsf h d}{1-4d}.
\]
Combining this with \cref{eq:phil-native-conditioned-held} proves
\cref{eq:phil-actual-parent-path-error}.

For the displayed joint choices,
\(d\le\mathrm T+\mathrm T^{7/4}\le2\mathrm T\),
\(\mathsf h\delta\le\mathrm T^{3/4}\), and
\(\rho=O(\mathrm T^{1/4})\).
The activation and following window fit inside the selection window
because \(1+\mathrm T^{-1/4}\le t_{\mathrm L}\) for sufficiently
small temperature under the mobility cap. The departure term is
\(O(\mathrm T^{3/4})\), giving the stated total error.
The Gaussian comparison laws served only to bound that error; the
parent's activation and subsequent learning remain the laws in the
statement. \hfill\(\square\)





\chapter{Notation}
\label{app:notation-contract}

\Cref{tab:notation} lists the symbols that recur across chapters, each with the one meaning it has in
this manuscript and the statement or section that introduces it. A superscript $x\in\{b,m\}$ names the
belief or the model channel. Symbols used inside a single chapter are declared at its opening, and those
of the formation family are collected in \Cref{tab:rg-formation-notation}. Metrics are sans-serif,
statement-local constants are upright, and $\mathcal N(\mu,\Sigma)$ is the normal law with mean $\mu$
and covariance $\Sigma$.

{\footnotesize
\renewcommand{\arraystretch}{1.12}
\begin{longtable}{@{}>{\raggedright\arraybackslash}p{0.29\textwidth}
>{\raggedright\arraybackslash}p{0.48\textwidth}
>{\raggedright\arraybackslash}p{0.17\textwidth}@{}}
\caption{Notation of the main manuscript.}
\label{tab:notation}\\
\toprule
\textbf{Symbol} & \textbf{Meaning} & \textbf{Introduced} \\
\midrule
\endfirsthead
\toprule
\textbf{Symbol} & \textbf{Meaning} & \textbf{Introduced} \\
\midrule
\endhead
\midrule
\multicolumn{3}{r@{}}{\emph{continued on the next page}} \\
\endfoot
\bottomrule
\endlastfoot
\multicolumn{3}{@{}l}{\textbf{Base point, bundles, and frames}} \\
$\mathcal C$, $c$, $c_*$ & contextual base, a context, the fixed base point & \cref{def:geo-context-base} \\
$\pi:\mathscr P_G\to\mathcal C$, $G$, $e_G$ & principal right $G$-bundle, its structure group, the identity of $G$ & \cref{def:geo-principal-systems} \\
$u_i$, $\sigma$, $U_i$ & common frame of agent $i$ at $c_*$ for both channels; a reference point of $(\mathscr P_G)_{c_*}$; the frame relative to it, $u_i=\sigma\cdot U_i^{-1}$ & \cref{def:geo-principal-systems} \\
$g_i$, $R_i^x=\rho_x(g_i)$ & passive re-expression $u_i'=u_i\cdot g_i^{-1}$ (vertex gauge); the represented sample-coordinate map & \cref{sec:geo-frame-gauge} \\
$\varphi_i$ & local function $F(u_i)=u_i\cdot\varphi_i$ of a fiber automorphism $F$ & \cref{prop:geo-active-image-point} \\
$T_{ij}$ & pointwise frame comparison, $u_j=u_i\cdot T_{ij}$, $T_{ij}=U_iU_j^{-1}$ & \cref{prop:geo-frame-comparisons} \\
$(\mathsf K,\mathscr K)$, $(\mathsf M,\mathscr M)$ & belief and model sample fibers & \cref{def:geo-principal-systems} \\
$\rho_b,\rho_m$; $\widehat\rho_x$ & representations of $G$ on the sample fibers; the pushforward action $\widehat\rho_x(g)q=(\rho_x(g))_\#q$ on laws & \cref{def:geo-principal-systems} \\
$\mathcal B_x$, $\mathcal E_x$, $\varpi_x$ & law fiber, associated bundle $\mathscr P_G\times_{\widehat\rho_x}\mathcal B_x$, and its projection & \cref{def:geo-associated-bundles} \\
$q_i^b$, $q_i^m$ & belief-law and model-law sections of agent $i$, and their coordinate laws at $c_*$ & \cref{def:geo-agent} \\
$\Phi$, $\widetilde\Phi$; $\Phi_i$, $\widetilde\Phi_i$ & cross morphisms $\mathcal E_b\to\mathcal E_m$ and $\mathcal E_m\to\mathcal E_b$; their representatives in the frame $u_i$ & \cref{def:geo-cross-morphisms} \\
\multicolumn{3}{@{}l}{\textbf{Interaction multigraph and links}} \\
$V$, $N$, $i,j$ & agent set, number of agents (also the extensive size of a finite labeled family), agent indices & \cref{def:geo-agent} \\
$\mathfrak I=(V,E,\bar{\ })$, $e:j\to i$, $\bar e$ & interaction multigraph at $c_*$, containing the reversal closure of the source-prior supports; oriented edge copy; reversal & \cref{def:geo-graph-links} \\
$\Theta^x_e$, $\mathsf t^x_e$ & independent link on $e$, with $\Theta^x_{\bar e}=(\Theta^x_e)^{-1}$; its link map $\mathsf t^x_e(u_j\cdot g)=u_i\cdot\Theta^x_eg$ & \cref{def:geo-graph-links} \\
$\mathring\Theta^x_e$ & per-edge twist $U_i^{-1}\Theta^x_eU_j$, whose conjugacy class is invariant & \cref{prop:geo-edge-twist-invariance} \\
$\Theta^x_p$, $\mathrm{Hol}^x(p)$ & path product along a walk $p$; holonomy of a closed walk & \cref{eq:geo-link-holonomy} \\
$\Theta^x_{I\leftarrow i}$, $\mathrm{Hol}^x_I$, $\widehat\Theta^x_e$ & tree transport to the root of a block $I$, based holonomy group of $I$, block-dressed boundary edge & \cref{def:rg-rooted-tree-data} \\
\multicolumn{3}{@{}l}{\textbf{Samples, records, and the joint}} \\
$k_i$, $m_i$; $\mathsf Y_i$ & belief sample and model presentation of agent $i$; its latent block $\mathsf K_i\times\mathsf M_i$ & \cref{def:prob-product-sample-spaces} \\
$y\in\mathsf Y$, $o\in\mathsf O$ & population latent sample and latent space; observation and observation space & \cref{def:prob-product-sample-spaces} \\
$X\in\mathsf X$, $\theta$ & structural data and their configuration space; generative parameters & \cref{def:prob-structural-kernel-signatures,def:gen-parameter-split} \\
$a\in\mathfrak A$, $\partial a$, $K_a$, $E_{a,o}$ & record, its scope, its normalized record kernel, and its negative log density at the record $o$ & \cref{def:prob-product-sample-spaces,hyp:local-interaction-kernels} \\
$\operatorname{ev}_i$, $K^X_{i,m_i}$ & evaluator of agent $i$; the evaluated mechanism of the presentation $m_i$ & \cref{def:prob-model-evaluation} \\
$\operatorname{Acc}_i$, $\mathbb Q_{i,o,X}$, $\operatorname{Cpl}(\mathbb Q_{\bullet,o,X})$ & access map, local recognition law, and the class of couplings of the local laws & \cref{def:prob-structural-kernel-signatures,def:prob-recognition-marginals} \\
$\mathfrak q_i^x$, $\mathsf Y^{\mathrm{law}}_i$; $\mathsf Y^{\mathrm{cfg}}$, $\operatorname{pres}$, $\kappa^{\mathrm{cfg}}$ & law points and the law-point block of agent $i$; configuration space, presentation map, presentation kernel & \cref{def:prob-description-levels} \\
$\mathbb P_{\theta,V}$, $\mathbb Q_{V,o,X}$, $\boldsymbol\Pi_{\theta,V,o,X}$ & constructed joint, selected recognition coupling, posterior slice & \cref{cons:gen-finite-directed-law,hyp:prob-measurable-coupling-selector} \\
$\widetilde p_X$ & unnormalized potential of an undirected model & \cref{def:gen-undirected-model} \\
$\Lelbo_V$, $\Fenergy_V=-\Lelbo_V$; $\Lelbo^{\mathrm{ext}}$, $\Fenergy^{\mathrm{ext}}$ & evidence bound and variational free energy; their extended versions, defined on every recognition law & \cref{def:elbo-extended} \\
$\KL$, $\operatorname{TC}$ & relative entropy with values in $[0,+\infty]$; total correlation & \cref{def:prob-log-density-ratio-convention,thm:elbo-total-correlation-chain} \\
\multicolumn{3}{@{}l}{\textbf{Attention and selective communication}} \\
$\pi_{ij}$, $\mathcal N_i$, $J_i$ & source prior of receiver $i$, its source set (the support of $\pi_i$), the source label & \cref{def:obs-attention-label-model} \\
$(i,h)$ & receiving row of head $h$ of agent $i$ & \cref{def:obs-attention-label-model} \\
$\beta^x_{ij}$, $\alpha^x_i$, $\eta^x_{ij}$ & attention row of receiver $i$, receiver occupancy, and event law $\eta=\alpha\beta$, with the role mark $P$, $\Pi$, or $Q$ of the law they are computed from ($\beta^m$ names the model channel); the event law carries no head index & \cref{def:obs-attention-event-law} \\
$\tau_i$, $\mathsf c_{ij}$, $\mathsf r_i$, $\Fenergy^{\mathrm{att},x}_i$ & attention temperature, replicator cost, row relaxation rate, row free energy & \cref{cor:obs-attention-replicator,def:full-graph-transported-divergences} \\
$S^Q_e$, $S^K_e$, $H^{\mathrm{cmp}}_e$; $\mu^Q_i$, $\mu^K_j$ & query and key comparison maps into the comparison space; projected query and key means & \cref{def:obs-selective-comparison} \\
$K^v_e$; $\mathsf L^{\mathrm{enc}}_e$, $\mathsf L^{\mathrm{tr}}_e$, $\mathsf L^{\mathrm{dec}}_e$, $\mathsf L_e$, $\mathsf L^{\mathrm{val}}_e$ & composed value kernel of an edge copy; encoder, transmission, decoder, composite, and value linear maps & \cref{def:obs-selective-interface,def:geo-selective-matrix-link} \\
$\Sigma^+_i$, $\Sigma^{\mathrm{pred}}_i$, $\widetilde W_i=W_i/\tau_i$ & receiver posterior and predictive covariances; tempered receiver precision & \cref{prop:obs-selective-receiver-update} \\
$S_j$, $\Upsilon$, $v$ & source sample; random receiver side information and its values & \cref{prop:obs-selective-information-chain} \\
\multicolumn{3}{@{}l}{\textbf{Fiber geometry, histories, and time scales}} \\
$g^F$, $T^F$, $\mathbf s$; $\lambda$ & Fisher tensor, Amari--Chentsov tensor, score; coordinates of a recognition family & \cref{def:pb-fiber-statistical-tensors} \\
$\mathfrak k_q$, $\mathfrak g_q$ & orbit Fisher form and isotropy algebra at a law $q$ & \cref{cor:pb-orbit-fisher-form} \\
$\mathcal Q_i$, $\mathcal Q_V$, $\mathsf G$ & configuration manifolds of an agent and of the population at $c_*$; a declared configuration metric & \cref{def:hist-configuration-manifold} \\
$\mathcal Q_\ell$, $\mathsf F_\ell$, $\mathscr V_\ell$, $\operatorname{Fl}^{\mathscr V_\ell}_t$ & configuration manifold at depth $\ell$, coarse configuration map, natural-gradient field, and its flow & \cref{def:hist-coarse-configuration-map} \\
$\mathfrak t$ & Fisher duration, the information clock of a selected history & \cref{def:hist-fisher-clock} \\
$\varkappa$; $\mathfrak K$, $\mathfrak C$; $\operatorname{ret}$ & parity amplitude; residual second-moment tensor and signed cross tensor of a joint score; retained-law map & \cref{def:hist-three-agent-parity-lift,prop:hist-general-score-residual,thm:hist-retained-fisher-quotient} \\
$f$, $s=(s_{\mathrm m},\theta)$ & fast and slow coordinate blocks & \cref{def:obs-fast-slow-blocks} \\
$\mathsf G_{\mathrm f}$, $\mathsf G_{\mathrm s}$, $\varepsilon$, $\mathsf G_\varepsilon$, $F_{\mathrm{red}}$ & block metrics, rate ratio, rate-separated metric $\varepsilon\mathsf G_{\mathrm f}\oplus\mathsf G_{\mathrm s}$, reduced functional & \cref{lem:obs-envelope-slow-flow,prop:obs-rate-separated-flow} \\
\multicolumn{3}{@{}l}{\textbf{Coarse channels and the network state}} \\
$\ell$ & depth of coarse-graining, and the level index of the formation family & \cref{def:rg-scale-diagram} \\
$I$, $J$; $A$ & child blocks; the parent label of a block & \cref{thm:rg-pointwise-parent-datum} \\
$C_A$, $\mathsf Z_A$, $\Delta_A$ & parent coarse kernel, parent state space, conditional-information defect of the coarse chain rule & \cref{thm:rg-pointwise-parent-datum,thm:rg-exact-coarse-vfe} \\
$\mathfrak X_\ell$, $C_{\ell k}$ & objects and arrows of a scale diagram & \cref{def:rg-scale-diagram} \\
$\gamma\in\mathfrak G$, $W_\gamma$ & label of a finite labeled family (a mixture or geometry label) and its weight & \cref{def:rg-thermo-label-family} \\
$\mathsf s_\gamma$, $\mathsf v_\gamma$, $\mathcal I_\gamma$, $\bar{\mathsf p}$, $\mathsf w$ & bulk log-mass density, tie slope $\nabla\mathsf s_\gamma$, bulk Fisher density, limiting label law, window coordinate $\mathsf w=N(\theta-\theta_0)$ & \cref{sec:cg-retained-label-asymptotics} \\
$L_I$, $\mathbb V_{I,\mathrm{fix}}$, $f_I$ & connection Laplacian of a block, its fixed fiber $\ker L_I$, and the retained rank $\dim\mathbb V_{I,\mathrm{fix}}$ & \cref{def:cg-connection-laplacian} \\
\multicolumn{3}{@{}l}{\textbf{Exact coarse-graining and the multiscale free energy}} \\
$V_\ell$, $\mathsf Y_\ell$, $C_\ell$, $C^\dagger$ & agents and state space at depth $\ell$, the depth arrow, the reverse kernel of a channel & \cref{def:rg-measure-pair} \\
$(\varrho_\ell,\mathsf m_\ell)$, $M_\ell$, $\bar{\mathsf m}_\ell$, $H_\ell$ & measure pair of a reference probability and a finite evidence measure; its mass, normalized law, and action & \cref{def:rg-measure-pair} \\
$\mathsf U$, $\mathcal A$ & reverse conditional expectation and the exact action increment of a channel & \cref{thm:rg-bounded-action-calculus} \\
$\mathcal G_\ell$, $\mathsf g$, $\mathcal T_\ell$ & space of Hoeffding interaction coordinates, an element, and the exact interaction step & \cref{def:rg-hoeffding-coordinates,def:rg-interaction-step} \\
$D_\ell$, $D_{n\leftarrow\ell}$ & finite-step multipliers of the interaction step & \cref{def:rg-interaction-modes} \\
$\mathsf P^{\mathrm{ret}}_\ell$, $r^{\mathcal G}_{\ell+1}$ & retained projection and closure residual & \cref{def:rg-retained-projection} \\
$Y_\ell=(Z_\ell,\mathfrak N_\ell)$, $\mathbf Y$ & state at depth $\ell$, made of coordinates and the graph datum; the whole tower & \cref{def:full-graph-tower-vfe} \\
$R_\ell\in\mathsf R_\ell$; $C^{\mathrm{struct}}_{\ell+1\leftarrow\ell}$, $\widehat{\mathsf Y}_{\ell+1\leftarrow\ell}$ & structural mark sampled by the structural channel, such as a random partition, and its space; the structural channel and its augmented transition state & \cref{def:full-graph-structural-channel} \\
$\mathfrak m_\ell$, $\mathfrak M_\ell$ & membership kernel and joint endpoint membership kernel & \cref{def:full-graph-membership-kernel} \\
\end{longtable}
}

\chapter{Open Problems}
\label{app:claim-ledger}

\Cref{tab:open-problems} collects the open problems of this manuscript in the
order of the chapters. Each row names the item, the numbered statement where it arises, which also
states what holds, and a condition that would settle it. Established results, refutations among
them, stay at their statements and are not listed. Open problems of the base and Gaussian settings
are collected in the two supplements.

{\footnotesize
\renewcommand{\arraystretch}{1.12}
\begin{longtable}{@{}>{\raggedright\arraybackslash}p{0.23\textwidth}
>{\raggedright\arraybackslash}p{0.16\textwidth}
>{\raggedright\arraybackslash}p{0.55\textwidth}@{}}
\caption{Open problems of the main manuscript.}
\label{tab:open-problems}\\
\toprule
\textbf{Item} & \textbf{Locator} & \textbf{Settling condition} \\
\midrule
\endfirsthead
\toprule
\textbf{Item} & \textbf{Locator} & \textbf{Settling condition} \\
\midrule
\endhead
\midrule
\multicolumn{3}{r@{}}{\emph{continued on the next page}} \\
\endfoot
\bottomrule
\endlastfoot
\multicolumn{3}{@{}l}{\textbf{Part I: single-point foundations}} \\
Nondominated likelihood families & \cref{open:prob-nondominated-likelihood} & A treatment of the evidence bound, its local disintegrations, and the coarse channels at the level of joint measures and regular conditional laws, with a rule for record-null versions, in which the exact identities of \cref{ch:elbo,ch:coarsegraining} are proved, or hypotheses that restore a jointly measurable likelihood; negatively, a record family for which those identities provably have no measure-level counterpart. \\
A normalized law-point prior & \cref{open:prob-law-point-prior} & For a class of law points, a sigma-finite reference on the law-point or presentation space, a finite normalizer for the exponentiated law-point energies, the transformation law of the prior under the active gauge action, record kernels that read law points, and a compatibility relation between the law-point posterior and the sample-level recognition marginals. \\
An intrinsic criterion for recognition structure & \cref{open:restrict-intrinsic-criterion} & A functional on admissible channels or partitions, expressible relative to one fixed joint and not monotone toward an endpoint, whose optimizer is nondegenerate; negatively, a proof of endpoint degeneracy for a precisely stated class of such functionals. \\
Convergence of the alternating scheme & \cref{open:elbo-alternating-convergence} & For models of \cref{hyp:local-interaction-kernels} with finitely many latent states, parameters in an open cone, and joint densities twice continuously differentiable in the parameters, under a full exact E phase alternating with the exact M-coordinate with a unique maximizer: conditions checkable on the model, including compactness of the upper level sets of the evidence, met by some model of the class with parameter dimension at least two whose stationary set at the limit value of an orbit is not discrete, under which every orbit converges to one stationary point of the extended bound; or a model of that class with compact upper level sets whose orbit has more than one limit point. \\
Discrete creation and deletion of interaction edges & \cref{open:obs-edge-creation} & For the threshold rule of the statement (source sets in multigraphs of bounded multiplicity; deletion at a rate that is a continuous nonincreasing function of the recognition row weight and creation at a rate that is a continuous nonincreasing function of the averaged comparison energy, each positive below a threshold; rows on the current support following the replicator flow): for some pair of rate functions, a normalized joint at $c_*$ with the source sets as a structural coordinate, declared with a named law and coupled to the link, interface, and label variables, in which the collective VFE decomposition, the event law, and its pushforward extend to the random support and the collective VFE is nonincreasing under the master equation of the rule from every initial support and recognition law; or a proof that for every pair of rate functions no joint of that class has this property. \\
Label-level multi-head attention & \cref{open:obs-multihead-attention} & For heads of one receiver with exclusive labels and comparison records whose projected queries are correlated under the receiver's conditional Gaussian law and whose integrated positive label table does not factor (the factorizing subclass, including untied channel-split heads under a channel-product receiver law, is excluded): an exact recursion with bounded head count, a family of comparison-record models closed under the meta-attention pushforward composed with a declared head-merging channel, in which the parent's per-head rows are the disintegrations of the merged event law and its posterior mean is the coarse coordinate of its children's posterior means; or a member and a partition for which no comparison-record model at the block with fewer heads than the block has rows reproduces the posterior mean of the block's coarse coordinate. \\
Formation with evolving selective interfaces & \cref{open:obs-selective-formation} & For a connected vector-tier block with Gaussian belief means as fast coordinates, factored linear communication through surjective interfaces of one fixed rank below the fiber dimension, comparison records read through the receiver's interface, the kinetics of \cref{def:rg-formation-cycle-experiment} at temperature $\mathrm T$, and interfaces learned as generative parameters by a natural-gradient flow at mobility $\epsilon_{\mathrm M}$: a parent that is a retained mechanism of fixed dimension below the block's, independent of $\mathrm T$, $\epsilon_{\mathrm M}$, and the horizon, whose path law of the communicated coordinates is within a relative entropy tending to zero as $\epsilon_{\mathrm M}\to0$ on horizons of order $1/\epsilon_{\mathrm M}$, uniformly over $\mathrm T\in(0,\mathrm T_0]$; or a member whose learned interface path has a rest point at which no such parent is projectable, or in which every such parent keeps that relative entropy bounded below as $\epsilon_{\mathrm M}$ and $\mathrm T$ tend to zero along a stated relation. \\
Global Fisher quotients and canonical block data & \cref{open:hist-retained-fisher-quotient} & A criterion, checkable on the configuration manifolds at $c_*$ and proved, under which the null foliation of a pulled-back Fisher tensor has a Hausdorff manifold leaf space, for instance through a free and proper group action, together with block data on the quotient constructed from a stated typing of agents and channels, or a counterexample showing that the typing does not determine them. \\
Scope beyond the finite parity witness & \cref{open:hist-parity-lift-scope} & For a stated regular family of agent laws and a stated class of morphisms containing kernels that do not preserve products, a smooth right inverse of the marginal map that is natural under the class, and correlated lifts equivariant under a connected structure group acting on continuous law fibers, each by a construction with proof or a no-go theorem. \\
Quantitative meta-agent clock ratio & \cref{open:hist-meta-clock-ratio} & A theorem computing or bounding the ratio of a meta-agent's Fisher duration to its constituents' collective duration, for a stated coarse map and a stated law of constituent velocities, including its dependence on $N$ for coherent and for incoherent histories. \\
Slow-manifold reduction of the collective free energy & \cref{open:obs-slow-manifold} & For the sampled four-cycle dynamics of \cref{def:rg-formation-cycle-experiment} with model means and links at the common slow rate $\varepsilon_{\mathrm m}$ of \cref{prop:obs-slow-model-link-dichotomy}, no parameter learning, and the split by mobility: a bound on the distance between the slow coordinates of the sampled process and the reduced slow flow of its assignment path over a bounded slow-time interval, explicit in $\varepsilon_{\mathrm m}$, $\mathrm T$, and the number of assignment changes, tending to zero with $\varepsilon_{\mathrm m}$ and $\mathrm T$, for every positive integer $\mathrm n$ with probability tending to one on paths with at most $\mathrm n$ changes per unit slow time; or a positive integer $\mathrm n$, mobility ratios and temperatures tending to zero, and initial states from which, with probability bounded below, the path has at most $\mathrm n$ changes per unit slow time and that distance stays bounded below. \\
\multicolumn{3}{@{}l}{\textbf{Part II: exact coarse-graining at one base point}} \\
Equivariant Bayes versions beyond the finite and compact cases & \cref{open:cg-equivariant-versions} & For a locally compact second countable group acting continuously on Polish spaces, an equivariant channel $K$, and laws $Q$ with noncompact stabilizer whose posterior has no weakly continuous version on the support of $QK$: an equivariant measurable selection theorem giving, for every member, an everywhere-defined version equivariant under the stabilizer, jointly measurable across every measurable family of such laws meeting several orbits; or a pair $(Q,K)$ of the class with no everywhere-defined version equivariant under the stabilizer. \\
A natural selector of partitions and membership laws & \cref{open:cg-natural-selector} & A relabeling-natural, gauge-compatible functional of the laws, not monotone under refinement, with a nondegenerate optimizer; or a theorem that no such selector exists in a stated class. \\
Invariant retained subspaces & \cref{open:rg-invariant-retained-subspace} & For finite agent spaces with fixed product references of full support, at least two blocks, and deterministic block-product channels whose block maps depend on every coordinate of their block, with base laws of full support in the class of product laws or of all laws: a proper invariant pair with fixed retained projections, invariant under the step of every base law in a nonempty open subset of the class, whose depth-$\ell$ range is not contained in the sum of the pullback sector and a grouping sector of at least two groups; or a proof that off a closed nowhere dense subset of the class every proper invariant pair has its depth-$\ell$ range in such a sum. \\
Dynamic closure of the pointwise parent & \cref{open:rg-pointwise-parent-dynamics} & For switching diffusions of the type of \cref{def:rg-formation-cycle-experiment} indexed by $\epsilon$, channels as in part (b) of \cref{prop:rg-formation-endogenous-cycle} that hide coordinates read by the receiver rows on an assignment of weight bounded below, and parents that are time-homogeneous Markov on the retained assignment and coordinates or carry in addition a filter in the explicitly declared regular class of \cref{def:rg-regular-finite-filter}, at two levels (windows of fixed length uniformly up to the switching time, and the complete history over horizons on which the expected number of assignment changes diverges): a proof that every member admits such a parent whose retained and assignment path law is within a total variation tending to zero with $\epsilon$ of the native one; or a member for which every such parent keeps that distance bounded below. \\
Accounting across episodes & \cref{open:full-graph-episode-accounting} & For label-sufficient latent routes of \cref{def:full-graph-episode-joint}, whose episode model is that of the parameter route with a finite label coordinate, the posterior slice as recognition coupling, and the aggregate update of constant step $\epsilon_{\mathrm e}$, whose read-out reads the label and an auxiliary preference coordinate only, and whose records are silent and informative or have likelihood vectors spanning a cone on which no nonzero polynomial vanishes: a route whose per-episode discrepancy from the aggregate update applied to its own filtered label law and label posterior is $o(\epsilon_{\mathrm e})$ uniformly over record histories and over episodes up to any fixed multiple of $1/\epsilon_{\mathrm e}$; or a lower bound of order $\epsilon_{\mathrm e}$ on that discrepancy for every such route. \\
Persistent multiscale hierarchy & \cref{open:full-graph-hierarchy} & For models at $c_*$ with structural marks from a finite partition catalog inside both laws, descending the multiscale free energy by a reversible diffusion at temperature $\mathrm T$ with rate-one redraws of the marks, with energies, precisions, prior ratios, hyperparameters, and mobilities independent of $\mathrm T$, as $\mathrm T\to0$: three results together, a lineage of at least two levels, formed with probability bounded below from every initial law with a density and a finite second moment, whose second level depends on the first (a declared handoff comparison with a margin bounded below over counterfactual handoffs, and no occurrence with probability tending to one uniformly over the handoff class) and is a coherence episode (a restoring clause at a declared resolution and a horizon clause), a nondegenerate intrinsic partition selector, and a separated slow manifold on which depth-$(\ell+1)$ variables feed back consistently into depth-$\ell$ kernels; or a proof that in this class no second level produced by descent, apart from a supplied-scale sequence whose event occurs with probability tending to one uniformly over the handoff class, meets the restoring and horizon conditions and exceeds its counterfactual statistic by such a margin. \\
\multicolumn{3}{@{}l}{\textbf{Part III: formation and interpretation}} \\
Sharpness of the doubly exponential regime & \cref{open:rg-formation-regime-sharpness} & In the family of \cref{thm:finite-depth-adaptive-formation} under (H1), (H3), and (H4), with power-law strengths satisfying $3\mathrm Q_\ell\ge4\mathrm Q_{\ell+1}+20$ and $\mathrm Q_\ell\ge14$ below the top, the regime formed from $\mathrm Q_1$, and an entrance constant above the thresholds of the theorem and of \cref{prop:rg-formation-gaussian-core}: conclusions (a) to (e) for every such schedule with constants and preparation neighborhoods independent of $L$, through an early estimate of the retained Fisher information by the future-layer energy, a history-environment bound without the top-interface recursion, and depth-uniform flag and jump constants; or an admitted preparation and a schedule of the class under which a conclusion fails, or a lower bound showing that a conclusion cannot hold with constants independent of $L$. \\
A family in which learning or holonomy selects the components & \cref{open:rg-formation-adaptive-family} & On the four-cycle of \cref{prop:phil-seeded-contrast-formation}, with its temperature, learning mobility, and horizon, Gaussian initial laws $\mathcal N(\mathsf m_0,I_{12})$ with $\mathsf m_0$ outside a Lebesgue-null set, and parents of the type of part (b) of \cref{prop:rg-formation-endogenous-cycle}: such parents with separately vanishing conditional and dynamical errors, together with a total-variation difference bounded below as $\epsilon\to0$ between the conditional laws, given contrast formation, of the first formed block in the learned and the frozen-row arm, or between those of two initial means that differ only in their loop coordinate; or a proof that for almost every $\mathsf m_0$ both conditional differences vanish as $\epsilon\to0$. \\
Repeated endogenous formation & \cref{open:rg-formation-general} & For models at $c_*$ with source-gated attention on a finite connected graph, outside the family of \cref{def:rg-gated-translation-family}, whose energies, precisions, prior ratios, hyperparameters, and mobilities do not depend on $\mathrm T$, as $\mathrm T\to0$, and a second coalescence declared by a handoff time, first- and second-stage events with their blocks, one statistic of the handoff state, and a handoff class fixed before the comparison: such a coalescence, with learned generative parameters and from initially nonmatching assignments, that does not occur with probability tending to one uniformly over the handoff class, whose statistic exceeds its counterfactual value by a margin bounded below, whose block is held at a declared resolution by a restoring interaction on horizons short compared with the stationary entrance time of the event, and whose parents have separately vanishing conditional and dynamical errors; or a proof that in this class every coalescence of first-stage blocks on such horizons fails the margin or the restoring clause or has a second-stage event occurring with probability tending to one uniformly over the handoff class. \\
Physical-law identification & \cref{claim:physical-law-identification} & For a named target system whose data the model does not generate, observed through a declared record kernel that under the model's own law separates two objects of the coarse-graining differing in a gauge-invariant datum, with an observable, an estimator, and a validation and uncertainty protocol set before the data: held-out evidence, robust across admissible blocking schemes, that an object of the coarse-graining meets a baseline margin with probability-level closure; or held-out evidence under the same protocol that the margin is not met. \\
Single-point testability & \cref{open:phil-testability} & For a named target system at $c_*$ whose data the model does not generate, observed with sampling units and a noise model through a declared record kernel that reads the links of a closed walk or the coordinates of a candidate block: a test fixed before the data, with an estimator of gauge-typed parameters, a held-out statistic, an acceptance tolerance, and baselines without link holonomy or without formation, whose failure would falsify nontrivial holonomy on that walk or formation of that block, with a proof that under the kernel the law of the tested statistic depends on that hypothesis; negatively, a proof that under the kernel every statistic of the records has the same law for two link configurations differing only in the conjugacy class of the walk's holonomy, or whether or not the block is coherent at the declared resolution. \\
\end{longtable}
}

\chapter{Numerical Index}
\label{app:numerical-provenance}

\Cref{tab:numerical-index} indexes the statements whose status is NUMERICAL in the sense of
\Cref{tab:status-taxonomy}, each with its check, what the check tests, its control, and its seed. No
statement of the main manuscript has that status: every lemma, proposition, theorem, corollary, and
claim-bearing remark is proved here or cited to a primary source, and computations that corroborate
its proved statements are not indexed here.

{\footnotesize
\renewcommand{\arraystretch}{1.12}
\begin{longtable}{@{}>{\raggedright\arraybackslash}p{0.16\textwidth}
>{\raggedright\arraybackslash}p{0.17\textwidth}
>{\raggedright\arraybackslash}p{0.28\textwidth}
>{\raggedright\arraybackslash}p{0.17\textwidth}
>{\raggedright\arraybackslash}p{0.11\textwidth}@{}}
\caption{Numerical index of the main manuscript.}
\label{tab:numerical-index}\\
\toprule
\textbf{Claim} & \textbf{Check} & \textbf{What is tested} & \textbf{Control} & \textbf{Seed} \\
\midrule
\endfirsthead
\bottomrule
\endlastfoot
\multicolumn{5}{@{}l}{No statement of the main manuscript has the status NUMERICAL.} \\
\end{longtable}
}


\bibliographystyle{plainnat}
\bibliography{references}
\end{document}
